% !TeX program = pdflatex
\documentclass[11pt,a4paper,oneside]{book}
\usepackage[T1]{fontenc}
\usepackage[utf8]{inputenc}
\usepackage[polish]{babel}
\usepackage{lmodern}
\usepackage{amsmath,amssymb,amsthm,mathtools}
\usepackage{microtype}
\usepackage{array,longtable}
\usepackage{etoolbox}
% Compatibility with the 2025-11 LaTeX output routine:
% longtable 4.24 appends \vss (infinite stretch AND shrink) before
% the kernel's box splitting. Retain stretch but remove infinite shrink
% at these two output locations. Do not suppress diagnostics.
\makeatletter
\IfFormatAtLeastTF{2025-11-01}{%
  \patchcmd{\LT@output}{\copy\LT@foot\vss}{\copy\LT@foot\vfil}{}{}%
  \patchcmd{\LT@output}{\copy\LT@foot\vss}{\copy\LT@foot\vfil}{}{}%
}{}
\makeatother
\usepackage{geometry}
\usepackage{enumitem}
\usepackage{xcolor}
\usepackage{tikz}
\usepackage[unicode]{hyperref}
\hypersetup{colorlinks=true,linkcolor=blue!45!black,urlcolor=blue!45!black}
\geometry{margin=2.6cm}
\setlength{\parindent}{0pt}
\setlength{\parskip}{4pt}
\setlength{\emergencystretch}{2em}
\setlist[enumerate]{topsep=4pt,itemsep=2pt}
\numberwithin{equation}{chapter}
\theoremstyle{plain}
\newtheorem{twierdzenie}{Twierdzenie}[chapter]
\newtheorem{lemat}[twierdzenie]{Lemat}
\newtheorem{wniosek}[twierdzenie]{Wniosek}
\theoremstyle{definition}
\newtheorem{definicja}[twierdzenie]{Definicja}
\newtheorem{przyklad}[twierdzenie]{Przykład}
\theoremstyle{remark}
\newtheorem{uwaga}[twierdzenie]{Uwaga}
\renewcommand{\proofname}{Dowód}
\newcommand{\dd}{\,\mathrm d}
% Funkcje charakterystyczne oznaczamy wyłącznie przez \chi_A.
\title{Teoria miary\\\large Rozdziały 1--32: od konstrukcji miary do całki Daniella--Stone'a}
\author{}
\date{2026}

% BEGIN GLOBAL CODE LISTINGS (generated from latex/pl/code-listings.tex)
\usepackage{xcolor}
\usepackage{listings}
\definecolor{codebg}{RGB}{248,249,251}
\definecolor{codeborder}{RGB}{180,191,203}
\definecolor{codecomment}{RGB}{61,113,79}
\definecolor{codekeyword}{RGB}{41,78,132}
\lstdefinelanguage{Julia}{
  morekeywords={abstract,baremodule,begin,break,catch,const,continue,do,else,elseif,
    end,export,false,finally,for,function,global,if,import,in,isa,let,local,macro,
    module,mutable,primitive,quote,return,struct,true,try,type,using,where,while,
    missing,nothing,println,print},
  sensitive=true,morecomment=[l]{\#},morecomment=[n]{\#=}{=\#},
  morestring=[b]",
}
\lstdefinestyle{skrypt}{
  backgroundcolor=\color{codebg},frame=single,rulecolor=\color{codeborder},
  framesep=5pt,frameround=tttt,framerule=.5pt,
  basicstyle=\ttfamily\footnotesize,keywordstyle=\color{codekeyword}\bfseries,
  commentstyle=\color{codecomment}\itshape,stringstyle=\color{purple!65!black},
  numbers=left,numberstyle=\tiny\color{gray!75!black},numbersep=8pt,
  stepnumber=1,breaklines=true,breakatwhitespace=false,
  columns=fullflexible,keepspaces=true,showstringspaces=false,tabsize=4,
  linewidth=.90\linewidth,captionpos=b,aboveskip=11pt,belowskip=12pt,
  xleftmargin=12pt,xrightmargin=2pt,
}
\lstdefinestyle{pythonpde}{style=skrypt,language=Python}
\lstdefinestyle{python}{style=skrypt,language=Python}
\lstdefinestyle{julia}{style=skrypt,language=Julia}
\lstset{style=skrypt}
\renewcommand{\lstlistingname}{Kod}
\renewcommand{\lstlistlistingname}{Spis kodów}
% END GLOBAL CODE LISTINGS
\begin{document}
\frontmatter
\hypersetup{pageanchor=false}
\maketitle
\hypersetup{pageanchor=true}
% Lokalnie zwarty spis treści, bez samotnej końcowej pozycji.
\begingroup
\setlength{\parskip}{2pt}
\tableofcontents
\endgroup
% Redakcyjny spis najczęstszych symboli (bez powtarzania definicji z rozdziałów).
\chapter*{Spis oznaczeń}
\markboth{SPIS OZNACZEŃ}{SPIS OZNACZEŃ}
\phantomsection\addcontentsline{toc}{chapter}{Spis oznaczeń}
\noindent Symbole mogą przyjmować dodatkowe założenia w poszczególnych twierdzeniach.
Pierwszeństwo mają zawsze definicje podane we wskazanym rozdziale.

\section*{Zbiory i mierzalność}
\begin{longtable}{@{}p{0.285\textwidth}p{0.555\textwidth}p{0.095\textwidth}@{}}
$X,\,\mathcal P(X)$ & Przestrzeń bazowa i zbiór jej wszystkich podzbiorów. & rozdz.~\ref{ch:przestrzenie-mierzalne} \\
$\mathcal A,\,\mathcal F$ & Rodziny zbiorów; w szczególności $\sigma$-algebry. & rozdz.~\ref{ch:przestrzenie-mierzalne} \\
$\sigma(\mathcal G)$ & Najmniejsza $\sigma$-algebra zawierająca rodzinę $\mathcal G$. & rozdz.~\ref{ch:przestrzenie-mierzalne} \\
$\mathcal B(X)$ & $\sigma$-algebra borelowska przestrzeni topologicznej $X$. & rozdz.~\ref{ch:przestrzenie-mierzalne} \\
$\chi_A$ & Funkcja charakterystyczna zbioru $A$; jedyna stosowana w książce notacja tego pojęcia. & rozdz.~\ref{ch:funkcje-mierzalne} \\
$\mu^*$ & Miara zewnętrzna (niekoniecznie zdefiniowana na $\sigma$-algebrze). & rozdz.~\ref{ch:caratheodory} \\
$\lambda_d$ & $d$-wymiarowa miara Lebesgue'a na $\mathbb R^d$. & rozdz.~\ref{ch:lebesgue} \\
$\mathrm d\mu(x)$ & Element notacji całkowania wskazujący miarę i zmienną całkowania. & rozdz.~\ref{ch:calkowanie-lebesgue} \\
\end{longtable}

\section*{Zbieżność i całkowanie}
\begin{longtable}{@{}p{0.285\textwidth}p{0.555\textwidth}p{0.095\textwidth}@{}}
$\mu\otimes\nu$ & Miara produktowa na mierzalnej przestrzeni produktowej. & rozdz.~\ref{ch:produkty-miary} \\
$L^p(\mu)$ & Przestrzeń klas funkcji $p$-całkowalnych; dla $p=\infty$ funkcji istotnie ograniczonych. & rozdz.~\ref{ch:przestrzenie-lp} \\
$\|f\|_{L^p}$ & Norma w przestrzeni $L^p(\mu)$. & rozdz.~\ref{ch:przestrzenie-lp} \\
$C_c(X)$ & Przestrzeń ciągłych funkcji o zwartym nośniku. & rozdz.~\ref{ch:radon-measures-20} \\
$C_0(X)$ & Przestrzeń ciągłych funkcji zanikających w nieskończoności. & rozdz.~\ref{ch:riesz-markov-21} \\
$F_*\mu$ & Miara obrazowa: $(F_*\mu)(A)=\mu(F^{-1}(A))$. & rozdz.~\ref{ch:rn-applications-19} \\
\end{longtable}

\section*{Miary i ich rozkłady}
\begin{longtable}{@{}p{0.285\textwidth}p{0.555\textwidth}p{0.095\textwidth}@{}}
$\nu\ll\mu$ & Bezwzględna ciągłość miary $\nu$ względem $\mu$. & rozdz.~\ref{ch:ac-singular-17} \\
$\nu\perp\mu$ & Wzajemna osobliwość miar. & rozdz.~\ref{ch:miary-ze-znakiem} \\
$\nu^+,\,\nu^-$ & Dodatnia i ujemna część rozkładu Jordana miary ze znakiem. & rozdz.~\ref{ch:miary-ze-znakiem} \\
$|\nu|$ & Wariacja całkowita miary ze znakiem, zespolonej lub wektorowej; zależnie od kontekstu. & rozdz.~\ref{ch:miary-ze-znakiem} \\
$\frac{\mathrm d\nu}{\mathrm d\mu}$ & Pochodna Radona--Nikodyma, a nie symbol całkowania. & rozdz.~\ref{ch:radon-nikodym-18} \\
$\operatorname{supp}\mu$ & Nośnik miary, z uwzględnieniem odpowiednich założeń topologicznych. & rozdz.~\ref{ch:radon-measures-20} \\
\end{longtable}

\section*{Miary geometryczne}
\begin{longtable}{@{}p{0.285\textwidth}p{0.555\textwidth}p{0.095\textwidth}@{}}
$\mathcal H^s_\delta,\,\mathcal H^s$ & Miary Hausdorffa na skali $\delta$ oraz po przejściu do granicy; z ustaloną normalizacją. & rozdz.~\ref{ch:hausdorff-26} \\
$\dim_{\mathrm H}(E)$ & Wymiar Hausdorffa zbioru $E$. & rozdz.~\ref{ch:hausdorff-26} \\
$\mu_g$ & Miara objętości związana z metryką riemannowską $g$. & rozdz.~\ref{ch:miary-rozmaitosci-27} \\
$J_k(A)$ & $k$-wymiarowy jakobian odwzorowania liniowego $A$. & rozdz.~\ref{ch:area-coarea-28} \\
$\operatorname{Jac}_m(B)$ & $m$-wymiarowy jakobian koobszaru odwzorowania liniowego $B$. & rozdz.~\ref{ch:area-coarea-28} \\
$BV(\Omega)$ & Przestrzeń funkcji o ograniczonej wariacji w $\Omega$. & rozdz.~\ref{ch:geometric-measure-29} \\
$BV_{\mathrm{loc}}(\Omega)$ & Przestrzeń funkcji o wariacji lokalnie skończonej. & rozdz.~\ref{ch:geometric-measure-29} \\
$Du,\,|Du|$ & Miara pochodnej dystrybucyjnej i jej wariacja. & rozdz.~\ref{ch:geometric-measure-29} \\
$P(E;\Omega)$ & Obwód zbioru $E$ względem otwartego zbioru $\Omega$. & rozdz.~\ref{ch:geometric-measure-29} \\
$\partial^*E,\,\nu_E$ & Brzeg zredukowany i zewnętrzna normalna miarowo-teoretyczna. & rozdz.~\ref{ch:geometric-measure-29} \\
\end{longtable}

\section*{Miary wektorowe}
\begin{longtable}{@{}p{0.285\textwidth}p{0.555\textwidth}p{0.095\textwidth}@{}}
$\nu:\mathcal A\to X$ & Miara o wartościach w przestrzeni Banacha $X$. & rozdz.~\ref{ch:vector-measures-30} \\
$\|\nu\|_{\mathrm{sv}}$ & Półwariacja miary wektorowej. & rozdz.~\ref{ch:vector-measures-30} \\
$\int_A g(x)\,\mathrm d\mu(x)$ & Całka funkcji wektorowej w sensie Bochnera, gdy $g$ spełnia odpowiednie założenia. & rozdz.~\ref{ch:vector-measures-30} \\
$\mathrm{RNP}$ & Własność Radona--Nikodyma przestrzeni Banacha. & rozdz.~\ref{ch:vector-measures-30} \\
\end{longtable}

\noindent\emph{Uwaga.} Notacja $\chi_A$ jest stosowana konsekwentnie
w całym skrypcie; miary geometryczne stosują wspólną kalibrację
ustaloną w rozdziale~\ref{ch:hausdorff-26}.


\section*{Dżety i przedłużanie danych}
\begin{longtable}{@{}p{0.285\textwidth}p{0.555\textwidth}p{0.095\textwidth}@{}}
$j_p^kf$ & Dżet $k$-tego rzędu funkcji w punkcie $p$. & rozdz.~\ref{ch:dzety} \\
$J^kE$ & Wiązka dżetów $k$-tego rzędu przekrojów $E\to M$. & rozdz.~\ref{ch:dzety} \\
$D_{\mathrm{ap}}f$ & Pochodna przybliżona, określona w odpowiednich punktach. & rozdz.~\ref{ch:dzety} \\
\end{longtable}

\mainmatter
% ROZDZIAŁ 1 — plik włączany przez main.tex
\chapter{Rodziny zbiorów i przestrzenie mierzalne}

\providecommand{\atlasref}[2]{#2}

% Konwencja obowiązująca w całym skrypcie: funkcję charakterystyczną zbioru A oznaczamy przez $\chi_A$.
\label{ch:przestrzenie-mierzalne}

\section{Dlaczego potrzebujemy \texorpdfstring{$\sigma$}{sigma}-algebr?}

W teorii miary chcemy przypisywać rozmiar zbiorom, a następnie całkować funkcje.
Aksjomaty teorii mają zapewniać zgodność tego rozmiaru z rozkładaniem zbioru
na przeliczalnie wiele rozłącznych części. Nie wystarczy zatem umieć mierzyć
pojedynczych zbiorów: potrzebujemy rodziny zbiorów zamkniętej na operacje,
które pojawiają się w granicach ciągów i szeregach. Właśnie tę rolę pełnią
$\sigma$-algebry. Istotne jest, że nie zakładamy mierzalności wszystkich
podzbiorów rozważanej przestrzeni.

W całym rozdziale $X$ i $Y$ oznaczają zbiory, $\mathcal P(X)$ --- zbiór
wszystkich podzbiorów $X$, natomiast $A^c=X\setminus A$ oznacza dopełnienie
zbioru $A\subseteq X$ względem ustalonej przestrzeni $X$. Symbol $\mathbb N$
oznacza zbiór liczb naturalnych dodatnich. Rodziny zbiorów oznaczamy
$\mathcal A,\mathcal F,\mathcal G,\mathcal M$.

\begin{definicja}[Pierścień, algebra, $\sigma$-algebra]
\label{def:algebry}
Niech $X$ będzie zbiorem, a $\mathcal R\subseteq\mathcal P(X)$.
\begin{enumerate}
\item $\mathcal R$ nazywamy \emph{pierścieniem zbiorów}, jeżeli
$\varnothing\in\mathcal R$ i dla dowolnych $A,B\in\mathcal R$ zachodzą
$A\cup B\in\mathcal R$ oraz $A\setminus B\in\mathcal R$.
\item $\mathcal R$ nazywamy \emph{algebrą zbiorów} na $X$, jeżeli jest
pierścieniem i $X\in\mathcal R$.
\item $\mathcal R$ nazywamy \emph{$\sigma$-algebrą} na $X$, jeżeli
$X\in\mathcal R$, $A^c\in\mathcal R$ dla każdego $A\in\mathcal R$ oraz
$\bigcup_{n=1}^{\infty}A_n\in\mathcal R$ dla każdego ciągu
$(A_n)_{n\geq1}$ elementów $\mathcal R$.
\end{enumerate}
Parę $(X,\mathcal F)$, gdzie $\mathcal F$ jest $\sigma$-algebrą na $X$,
nazywamy \emph{przestrzenią mierzalną}; elementy $\mathcal F$ to \emph{zbiory
mierzalne}.
\end{definicja}

\begin{uwaga}
Algebra jest zamknięta na skończone sumy, przekroje i różnice zbiorów.
$\sigma$-algebra jest ponadto zamknięta na przeliczalne przekroje, gdyż
z praw de Morgana
$\bigcap_{n\geq1}A_n=(\bigcup_{n\geq1}A_n^c)^c$.
Każda $\sigma$-algebra jest algebrą, ale nie odwrotnie.
\end{uwaga}

\begin{przyklad}[Najmniejsza i największa $\sigma$-algebra]
Dla dowolnego $X$ rodziny $\{\varnothing,X\}$ oraz $\mathcal P(X)$ są
$\sigma$-algebrami. Pierwszą nazywamy \emph{trywialną}, drugą ---
\emph{dyskretną}. Gdy $X=\varnothing$, są one równe.
\end{przyklad}

\begin{przyklad}[Zbiory skończone i ko-skończone]
\label{ex:finite-cofinite}
Niech $X$ będzie nieskończonym zbiorem. Rodzina
\[
\mathcal A=\{A\subseteq X:A\ \text{jest skończony lub }X\setminus A
\ \text{jest skończony}\}
\]
jest algebrą: zawiera $\varnothing$ i $X$. Suma dwóch jej elementów
jest skończona, gdy oba są skończone, a w przeciwnym przypadku
ma skończone dopełnienie. Dla $A,B\in\mathcal A$ różnica
$A\setminus B$ jest skończona, gdy $A$ jest skończony lub
$X\setminus B$ jest skończony. W pozostałym przypadku
$X\setminus A$ i $B$ są skończone, zatem
$X\setminus(A\setminus B)=(X\setminus A)\cup B$
jest skończony. Jeśli $X=\mathbb N$, algebra ta nie jest
$\sigma$-algebrą:
każdy singleton $\{2n\}$ należy do $\mathcal A$, lecz
$\bigcup_{n\geq1}\{2n\}$ ani nie jest skończony, ani nie ma skończonego
dopełnienia.
\end{przyklad}

\begin{przyklad}[Algebra przeliczalno-koprzeliczalna]
\label{ex:countable-cocountable}
Niech $X$ będzie dowolnym zbiorem. Rodzina
\[
\mathcal C_X=\{A\subseteq X:A\text{ jest przeliczalny lub }X\setminus A
\text{ jest przeliczalny}\}
\]
jest $\sigma$-algebrą. Istotnie, dopełnienia zachowują podaną własność.
Jeżeli w przeliczalnej rodzinie $(A_n)$ występuje zbiór o przeliczalnym
dopełnieniu, to dopełnienie $\bigcup_nA_n$ jest przeliczalne; jeśli zaś
wszystkie $A_n$ są przeliczalne, ich suma również jest przeliczalna.
Jeżeli $X$ jest nieprzeliczalny, to $\mathcal C_X\subsetneq\mathcal P(X)$.
Rzeczywiście, możemy wybrać dwa rozłączne nieprzeliczalne podzbiory
$A,B$ o sumie $X$ (wybór takiego podziału jest możliwy w zwykłych
założeniach teorii mnogości z aksjomatem wyboru). Wtedy ani $A$,
ani $X\setminus A=B$ nie jest przeliczalny, więc
$A\notin\mathcal C_X$.
\end{przyklad}

\begin{przyklad}[Zbiory określone przez podział]
Niech $(E_i)_{i\in I}$ będzie rozbiciem zbioru $X$ na niepuste, parami
rozłączne zbiory. Dla $J\subseteq I$ oznaczmy $E_J=\bigcup_{i\in J}E_i$.
Rodzina $\{E_J:J\subseteq I\}$ jest $\sigma$-algebrą:
$E_I=X$, a z rozłączności rozbicia otrzymujemy
$X\setminus E_J=E_{I\setminus J}$. Dla każdego ciągu
$(J_n)_{n\geq1}$ podzbiorów $I$ zachodzi
$\bigcup_n E_{J_n}=E_{\bigcup_n J_n}$.
W szczególności dla
$X=\{a,b,c\}$ podział $\{\{a,b\},\{c\}\}$ wyznacza $\sigma$-algebrę
$\{\varnothing,\{a,b\},\{c\},X\}$.
\end{przyklad}

\section{Generowanie \texorpdfstring{$\sigma$}{sigma}-algebr}

\begin{twierdzenie}[Przekroje $\sigma$-algebr]
\label{thm:intersection-sigma}
Niech $X$ będzie zbiorem i niech $(\mathcal F_i)_{i\in I}$ będzie
niepustą rodziną $\sigma$-algebr na $X$. Wówczas
$\bigcap_{i\in I}\mathcal F_i$ jest $\sigma$-algebrą na $X$.
\end{twierdzenie}
\begin{proof}
Wszystkie $\mathcal F_i$ zawierają $X$. Jeśli $A$ należy do każdego
$\mathcal F_i$, to jego dopełnienie również należy do każdego z nich.
Analogicznie, jeśli wszystkie $A_n$ należą do przekroju tych rodzin,
to $\bigcup_nA_n$ należy do każdej $\mathcal F_i$, a więc do ich przekroju.
\end{proof}

\begin{definicja}[$\sigma$-algebra generowana]
\label{def:sigma-generated}
Dla zbioru $X$ i rodziny $\mathcal E\subseteq\mathcal P(X)$ definiujemy
\[
\sigma_X(\mathcal E)=\bigcap\{\mathcal F:\mathcal F\text{ jest
$\sigma$-algebrą na }X\text{ i }\mathcal E\subseteq\mathcal F\}.
\]
Jest to \emph{$\sigma$-algebra generowana przez $\mathcal E$}.
Jeśli przestrzeń $X$ jest ustalona, piszemy po prostu $\sigma(\mathcal E)$.
\end{definicja}

Rodzina pod przekrojem jest niepusta, ponieważ zawiera $\mathcal P(X)$.
Z twierdzenia \ref{thm:intersection-sigma} wynika, że $\sigma(\mathcal E)$
jest $\sigma$-algebrą. Ponadto jest najmniejszą $\sigma$-algebrą
zawierającą $\mathcal E$, w sensie inkluzji.

\begin{przyklad}
Dla $E\subseteq X$ mamy
$\sigma(\{E\})=\{\varnothing,E,E^c,X\}$, z ewentualnymi powtórzeniami,
gdy $E=\varnothing$ lub $E=X$.
\end{przyklad}

\begin{uwaga}[Pułapka dotycząca sumy rodzin]
Nawet rosnąca suma $\sigma$-algebr nie musi być $\sigma$-algebrą.
Na $X=\mathbb N$ niech $\mathcal F_n$ oznacza $\sigma$-algebrę
generowaną przez singletony $\{1\},\ldots,\{n\}$.
Wówczas $\mathcal F_n\subseteq\mathcal F_{n+1}$. Dokładniej,
$\mathcal F_n$ składa się ze wszystkich sum elementów podziału
$\{\{1\},\ldots,\{n\},\mathbb N\setminus\{1,\ldots,n\}\}$:
każdy taki zbiór jest albo podzbiorem $\{1,\ldots,n\}$,
albo ma skończone dopełnienie zawarte w $\{1,\ldots,n\}$.
Odwrotnie, każdy skończony lub ko-skończony podzbiór $\mathbb N$
należy do pewnej $\mathcal F_n$, ponieważ jego skończona część
lub skończone dopełnienie mieści się w $\{1,\ldots,n\}$.
Stąd $\bigcup_n\mathcal F_n$ jest dokładnie algebrą zbiorów
skończonych i ko-skończonych z przykładu
\ref{ex:finite-cofinite}, a nie $\sigma$-algebrą.
\end{uwaga}

\begin{definicja}[$\sigma$-algebra śladowa]\label{idx:01:definicja:04}

Niech $(X,\mathcal F)$ będzie przestrzenią mierzalną i niech
$E\subseteq X$ będzie dowolnym podzbiorem. Definiujemy
$\mathcal F|_E=\{A\cap E:A\in\mathcal F\}$.
Rodzina ta jest $\sigma$-algebrą na $E$, nazywaną \emph{śladem}
$\mathcal F$ na $E$.
\end{definicja}
\begin{proof}[Uzasadnienie]
Do rodziny należy $E=X\cap E$. Dopełnienie w $E$ zbioru $A\cap E$
wynosi $A^c\cap E$, a przeliczalna suma zbiorów $A_n\cap E$ jest równa
$(\bigcup_n A_n)\cap E$.
\end{proof}

\section{Układy \texorpdfstring{$\pi$}{pi} i układy Dynkina}

\begin{definicja}[Układ $\pi$ i układ Dynkina]
\label{def:pi-lambda}
Niech $X$ będzie zbiorem. Rodzinę $\mathcal P\subseteq\mathcal P(X)$
nazywamy \emph{układem $\pi$}, jeżeli jest niepusta i
$A\cap B\in\mathcal P$ dla wszystkich $A,B\in\mathcal P$.
Rodzinę $\mathcal D\subseteq\mathcal P(X)$ nazywamy \emph{układem
Dynkina} (lub \emph{układem $\lambda$}), jeżeli:
\begin{enumerate}
\item $X\in\mathcal D$;
\item jeśli $A,B\in\mathcal D$ i $A\subseteq B$, to $B\setminus A\in\mathcal D$;
\item dla dowolnych parami rozłącznych $A_1,A_2,\ldots\in\mathcal D$
zachodzi $\bigcup_nA_n\in\mathcal D$.
\end{enumerate}
\end{definicja}

Każda $\sigma$-algebra jest układem Dynkina i układem $\pi$.
Samo bycie układem Dynkina nie gwarantuje zamknięcia na przekroje.
Przykładowo, jeśli $X=\{1,2,3,4\}$, to rodzina
\[
\mathcal D=\{\varnothing,X,\{1,2\},\{3,4\},\{1,3\},\{2,4\}\}
\]
jest układem Dynkina: każda z czterech niepustych właściwych
części ma w $\mathcal D$ swoje dopełnienie, a dwa rozłączne
niepuste elementy $\mathcal D$ inne niż $X$ są dopełnieniami
i mają sumę $X$. Zatem różnice zbiorów zagnieżdżonych
oraz przeliczalne sumy rozłączne pozostają w $\mathcal D$.
Nie jest jednak $\sigma$-algebrą, gdyż
$\{1,2\}\cap\{1,3\}=\{1\}\notin\mathcal D$.

\begin{lemat}
\label{lem:dynkin-intersection}
Układ Dynkina $\mathcal D$ na zbiorze $X$ jest $\sigma$-algebrą wtedy
i tylko wtedy, gdy jest zamknięty na skończone przekroje.
\end{lemat}
\begin{proof}
Konieczność jest oczywista. Załóżmy, że $\mathcal D$ jest układem
Dynkina zamkniętym na przekroje. Z warunku drugiego otrzymujemy
$A^c=X\setminus A\in\mathcal D$. Dla $A,B\in\mathcal D$ zachodzi
$B\setminus A=B\setminus(A\cap B)\in\mathcal D$, więc
$A\cup B=A\,\dot\cup\,(B\setminus A)\in\mathcal D$; symbol
$\dot\cup$ oznacza sumę rozłączną. Powtarzając ten krok, uzyskujemy
zamknięcie na sumy skończone. Dla dowolnego ciągu $(A_n)$ elementów
$\mathcal D$ zastępujemy $A_n$ przez
$B_n=A_n\setminus\bigcup_{k<n}A_k$. Dla $n=1$ mamy $B_1=A_1\in\mathcal D$. Dla $n>1$ suma
$C_n=\bigcup_{k<n}A_k$ należy do $\mathcal D$, więc
$B_n=A_n\setminus(A_n\cap C_n)\in\mathcal D$ na mocy drugiego
aksjomatu. Jeśli $m<n$, to $B_m\subseteq A_m\subseteq C_n$,
a $B_n\cap C_n=\varnothing$, zatem $B_m\cap B_n=\varnothing$.
Każdy punkt sumy $\bigcup_n A_n$ należy do $B_j$, gdzie $j$
jest najmniejszym indeksem takim, że punkt należy do $A_j$.
Ponieważ $B_n\subseteq A_n$, obie sumy są równe.
Trzeci aksjomat układu Dynkina kończy dowód.
\end{proof}

\begin{twierdzenie}[Twierdzenie $\pi$--$\lambda$ Dynkina]
\label{thm:pi-lambda}
Niech $X$ będzie zbiorem, $\mathcal P$ układem $\pi$ na $X$,
a $\mathcal D$ układem Dynkina na $X$. Jeśli
$\mathcal P\subseteq\mathcal D$, to
$\sigma(\mathcal P)\subseteq\mathcal D$.
\end{twierdzenie}
\begin{proof}
Przekrój dowolnej rodziny układów Dynkina na $X$ jest układem Dynkina.
Niech $\mathcal L$ będzie najmniejszym układem Dynkina zawierającym
$\mathcal P$. Wystarczy pokazać, że $\mathcal L$ jest $\sigma$-algebrą.

Dla ustalonego $A\in\mathcal P$ definiujemy
$\mathcal L_A=\{B\in\mathcal L:A\cap B\in\mathcal L\}$.
Rodzina $\mathcal L_A$ jest układem Dynkina: należy do niej $X$, gdyż
$A\cap X=A\in\mathcal L$; własność różnicy wynika z równości
$A\cap(B\setminus C)=(A\cap B)\setminus(A\cap C)$ dla $C\subseteq B$;
własność rozłącznych sum wynika z rozdzielności przekroju względem sumy.
Ponieważ $\mathcal P$ jest układem $\pi$, mamy
$\mathcal P\subseteq\mathcal L_A$. Z minimalności $\mathcal L$ wynika
$\mathcal L_A=\mathcal L$. Oznacza to, że
$A\cap B\in\mathcal L$ dla wszystkich $A\in\mathcal P$, $B\in\mathcal L$.

Teraz dla ustalonego $B\in\mathcal L$ rozważamy
$\mathcal K_B=\{A\in\mathcal L:A\cap B\in\mathcal L\}$.
Sprawdzamy trzy aksjomaty. Mamy $X\in\mathcal K_B$, ponieważ
$X\cap B=B\in\mathcal L$. Jeśli $A,C\in\mathcal K_B$ i $C\subseteq A$,
to $A\setminus C\in\mathcal L$ oraz
$(A\setminus C)\cap B=(A\cap B)\setminus(C\cap B)\in\mathcal L$,
gdyż $C\cap B\subseteq A\cap B$ i oba przekroje należą do $\mathcal L$.
Jeśli $(A_n)$ jest ciągiem parami rozłącznych elementów $\mathcal K_B$,
to $\bigcup_n A_n\in\mathcal L$, a zbiory $A_n\cap B$ są parami
rozłączne i należą do $\mathcal L$. Zatem
$(\bigcup_n A_n)\cap B=\bigcup_n(A_n\cap B)\in\mathcal L$.
Wobec tego $\mathcal K_B$ jest układem Dynkina. Wykazana własność daje
$\mathcal P\subseteq\mathcal K_B$, skąd
$\mathcal K_B=\mathcal L$. Zatem $\mathcal L$ jest zamknięty
na skończone przekroje. Na mocy lematu \ref{lem:dynkin-intersection}
jest $\sigma$-algebrą, a więc $\sigma(\mathcal P)\subseteq\mathcal L$.
Odwrotna inkluzja wynika z minimalności $\mathcal L$, ponieważ
$\sigma(\mathcal P)$ jest układem Dynkina zawierającym $\mathcal P$.
Ostatecznie $\mathcal L=\sigma(\mathcal P)\subseteq\mathcal D$.
\end{proof}

\begin{uwaga}
Twierdzenie $\pi$--$\lambda$ pozwala dowodzić własności wszystkich
zbiorów w $\sigma(\mathcal P)$, badając je najpierw na układzie
$\pi$. W rozdziale \ref{ch:miary-wlasnosci} zastosujemy je do
jednoznaczności miar; później również do jednoznaczności miar produktowych.
\end{uwaga}

\section{Klasy monotoniczne}

\begin{definicja}[Klasa monotoniczna]
\label{def:monotone-class}
Rodzinę $\mathcal M\subseteq\mathcal P(X)$ nazywamy \emph{klasą
monotoniczną}, jeżeli jest zamknięta na granice monotoniczne zbiorów:
\begin{align*}
A_n\in\mathcal M,\ A_n\uparrow A
&\Longrightarrow A\in\mathcal M,\\
B_n\in\mathcal M,\ B_n\downarrow B
&\Longrightarrow B\in\mathcal M.
\end{align*}
Zapis $A_n\uparrow A$ oznacza $A_n\subseteq A_{n+1}$ oraz
$A=\bigcup_nA_n$, a $B_n\downarrow B$ oznacza
$B_{n+1}\subseteq B_n$ oraz $B=\bigcap_nB_n$.
\end{definicja}

Przekrój klas monotonicznych jest klasą monotoniczną. Dlatego każda
rodzina $\mathcal E\subseteq\mathcal P(X)$ wyznacza najmniejszą klasę
monotoniczną $\mathcal M(\mathcal E)$, która ją zawiera.

\begin{twierdzenie}[Twierdzenie o klasie monotonicznej]
\label{thm:monotone-class}
Niech $X$ będzie zbiorem, a $\mathcal A$ algebrą podzbiorów $X$.
Wówczas $\mathcal M(\mathcal A)=\sigma(\mathcal A)$.
\end{twierdzenie}
\begin{proof}
Każda $\sigma$-algebra jest klasą monotoniczną, stąd
$\mathcal M(\mathcal A)\subseteq\sigma(\mathcal A)$.
Niech $\mathcal M=\mathcal M(\mathcal A)$. Pokażemy, że $\mathcal M$
jest $\sigma$-algebrą.

Najpierw rozważmy rodzinę
$\mathcal C=\{E\in\mathcal M:E^c\in\mathcal M\}$.
Jest ona klasą monotoniczną: dopełnienie zamienia ciągi rosnące
na malejące i odwrotnie. Ponieważ $\mathcal A$ jest algebrą,
$\mathcal A\subseteq\mathcal C$. Z minimalności $\mathcal M$
wynika $\mathcal C=\mathcal M$, czyli $\mathcal M$ jest zamknięta
na dopełnienia.

Ustalmy $A\in\mathcal A$ i połóżmy
$\mathcal M_A=\{E\in\mathcal M:E\cup A\in\mathcal M\}$.
Suma z ustalonym $A$ zachowuje granice monotoniczne, więc
$\mathcal M_A$ jest klasą monotoniczną. Algebraiczność $\mathcal A$
daje $\mathcal A\subseteq\mathcal M_A$, a zatem
$\mathcal M_A=\mathcal M$. Otrzymaliśmy
$E\cup A\in\mathcal M$ dla każdego $E\in\mathcal M$ i $A\in\mathcal A$.
Dla ustalonego $E\in\mathcal M$ rodzina
$\mathcal N_E=\{B\in\mathcal M:E\cup B\in\mathcal M\}$
jest więc klasą monotoniczną zawierającą $\mathcal A$, zatem
$\mathcal N_E=\mathcal M$. W konsekwencji $\mathcal M$ jest zamknięta
na sumy skończone.

Jeśli $E_1,E_2,\ldots\in\mathcal M$, to ich częściowe sumy
$F_n=\bigcup_{k=1}^nE_k$ należą do $\mathcal M$ i tworzą ciąg
rosnący. Stąd $\bigcup_nE_n\in\mathcal M$. Ponieważ
$X\in\mathcal A\subseteq\mathcal M$, otrzymaliśmy $\sigma$-algebrę.
Minimalność $\sigma(\mathcal A)$ daje drugą inkluzję.
\end{proof}

\begin{uwaga}[Dwa narzędzia rozszerzania własności]
Twierdzenia \ref{thm:pi-lambda} i \ref{thm:monotone-class} są podobne,
ale ich założenia nie są zamienne: twierdzenie Dynkina rozpoczyna się
od układu $\pi$, a twierdzenie o klasie monotonicznej --- od algebry.
Później użyjemy również wersji funkcyjnej twierdzenia o klasie
monotonicznej, po wprowadzeniu całki.
\end{uwaga}

\section{Zbiory borelowskie}

\noindent\emph{Powiązanie tematyczne w Bibliotece matematycznej:} \atlasref{GT-19::def:topog-polskie-borel}{GT-19: borelowska struktura mierzalna}.


\begin{definicja}[$\sigma$-algebra borelowska]
\label{def:borel}
Niech $(X,\tau)$ będzie przestrzenią topologiczną, gdzie $\tau$
stanowi rodzinę jej zbiorów otwartych. Definiujemy
$\mathcal B(X)=\sigma_X(\tau)$.
Jest to \emph{$\sigma$-algebra borelowska}; jej elementy nazywamy
\emph{zbiorami borelowskimi}. Gdy $X=\mathbb R^d$, używamy standardowej
topologii euklidesowej.
\end{definicja}

Każdy zbiór otwarty i domknięty jest borelowski, podobnie jak każdy
przeliczalny przekrój zbiorów otwartych i przeliczalna suma zbiorów
domkniętych. Zbiory tych ostatnich dwóch typów oznacza się odpowiednio
$G_\delta$ i $F_\sigma$. Nie każdy zbiór borelowski ma jedną z tych
dwóch postaci.

\begin{twierdzenie}[Przeliczalne generatory na $\mathbb R$]
\label{thm:borel-generators-R}
Dla standardowej topologii na $\mathbb R$ zachodzą równości
\[
\mathcal B(\mathbb R)
=\sigma\bigl(\{(a,b):a,b\in\mathbb Q,\ a<b\}\bigr)
=\sigma\bigl(\{(-\infty,q):q\in\mathbb Q\}\bigr)
=\sigma\bigl(\{(-\infty,q]:q\in\mathbb Q\}\bigr).
\]
\end{twierdzenie}
\begin{proof}
Wszystkie wymienione przedziały są borelowskie, więc każda z
wygenerowanych rodzin jest zawarta w $\mathcal B(\mathbb R)$.
Dla każdego $x\in U$, gdzie $U\subseteq\mathbb R$ jest otwarty,
istnieje $\varepsilon>0$ takie, że
$(x-\varepsilon,x+\varepsilon)\subseteq U$.
Z gęstości $\mathbb Q$ w $\mathbb R$ możemy wybrać
$a,b\in\mathbb Q$ spełniające
$x-\varepsilon<a<x<b<x+\varepsilon$.
Zatem $U$ jest sumą wszystkich przedziałów $(a,b)$
o końcach wymiernych zawartych w $U$.
Ponieważ takich przedziałów jest przeliczalnie wiele,
otrzymujemy pierwszą równość.

Dalej, dla $a<b$ mamy $(a,b)=(-\infty,b)\setminus(-\infty,a]$ oraz
$(-\infty,a]=\bigcap_{n\geq1}(-\infty,a+1/n)$.
Wystarczy brać wymierne $a,b$ w pierwszym wzorze; w drugim
$a+1/n$ jest wymierne dla wymiernego $a$. To dowodzi drugiej
równości. Wreszcie dla $q\in\mathbb Q$ zachodzą
$(-\infty,q]=\bigcap_{n\geq1}(-\infty,q+1/n)$ oraz
$(-\infty,q)=\bigcup_{n\geq1}(-\infty,q-1/n]$,
co daje trzecią równość.
\end{proof}

\begin{twierdzenie}[Borelowskie prostokąty]
\label{thm:borel-Rd}
Dla $d\in\mathbb N$ $\sigma$-algebra $\mathcal B(\mathbb R^d)$
jest generowana przez otwarte prostokąty
$\prod_{j=1}^d(a_j,b_j)$ o wymiernych końcach. Ponadto
\[
\mathcal B(\mathbb R^d)=\mathcal B(\mathbb R)^{\otimes d},
\]
gdzie prawa strona oznacza $\sigma$-algebrę generowaną przez
prostokąty $B_1\times\cdots\times B_d$ z $B_j\in\mathcal B(\mathbb R)$.
\end{twierdzenie}
\begin{proof}
Otwarte prostokąty o wymiernych końcach tworzą przeliczalną bazę
topologii $\mathbb R^d$. Zatem dowolny zbiór otwarty jest sumą
przeliczalnej podrodziny tej bazy, co daje pierwsze twierdzenie.
Oznaczmy przez $\mathcal F$ prawą stronę deklarowanej równości.
Każdy prostokąt bazowy należy do $\mathcal F$, skąd
$\mathcal B(\mathbb R^d)\subseteq\mathcal F$.
Z drugiej strony rzuty współrzędnych $p_j:\mathbb R^d\to\mathbb R$
są ciągłe, więc $p_j^{-1}(B_j)$ jest borelowski dla
$B_j\in\mathcal B(\mathbb R)$: rodzina zbiorów $B$ o tej
własności jest $\sigma$-algebrą zawierającą otwarte podzbiory
$\mathbb R$. Ponieważ
$\prod_{j=1}^dB_j=\bigcap_{j=1}^dp_j^{-1}(B_j)$,
każdy prostokąt borelowski należy do $\mathcal B(\mathbb R^d)$.
Stąd $\mathcal F\subseteq\mathcal B(\mathbb R^d)$.
\end{proof}

Dla dowolnego zbioru $A\subseteq X$ przez $\chi_A$ oznaczamy
funkcję charakterystyczną: $\chi_A(x)=1$ dla $x\in A$
oraz $\chi_A(x)=0$ dla $x\notin A$.

\begin{przyklad}[Racjonalne i irracjonalne]
$\mathbb Q$ jest podzbiorem borelowskim $\mathbb R$, bo jest
przeliczalną sumą singletonów, a singletony są domknięte.
Zatem $\mathbb R\setminus\mathbb Q$ również jest borelowski.
Jednak funkcja charakterystyczna $\chi_{\mathbb Q}$ jest nieciągła
w każdym punkcie $\mathbb R$, ponieważ każde otoczenie
każdego punktu zawiera zarówno liczbę wymierną, jak i niewymierną,
a funkcja przyjmuje na tych zbiorach różne wartości.
Mierzalność nie jest więc bliska
ciągłości w sensie punktowym.
\end{przyklad}

\section{Odwzorowania mierzalne i ich charakteryzacje}

\noindent\emph{Powiązanie tematyczne w Bibliotece matematycznej:} \atlasref{LST-08::def:image-preimage}{LST-08: obraz i przeciwobraz zbioru}.


\begin{definicja}[Odwzorowanie mierzalne]
\label{def:measurable-map}
Niech $(X,\mathcal F)$ i $(Y,\mathcal G)$ będą przestrzeniami
mierzalnymi. Odwzorowanie $f:X\to Y$ nazywamy
\emph{$(\mathcal F,\mathcal G)$-mierzalnym}, jeżeli
$f^{-1}(B)\in\mathcal F$ dla każdego $B\in\mathcal G$.
Jeśli $Y$ jest przestrzenią topologiczną z
$\mathcal G=\mathcal B(Y)$, mówimy krótko, że $f$ jest
\emph{borelowsko mierzalne} względem $\mathcal F$.
\end{definicja}

\begin{twierdzenie}[Kryterium przez generatory]
\label{thm:measurable-generator}
Niech $(X,\mathcal F)$ będzie przestrzenią mierzalną, $Y$ zbiorem,
$\mathcal E\subseteq\mathcal P(Y)$ rodziną zbiorów i
$\mathcal G=\sigma_Y(\mathcal E)$. Wówczas odwzorowanie
$f:X\to Y$ jest $(\mathcal F,\mathcal G)$-mierzalne wtedy i tylko
wtedy, gdy $f^{-1}(E)\in\mathcal F$ dla wszystkich $E\in\mathcal E$.
\end{twierdzenie}
\begin{proof}
Konieczność wynika z inkluzji $\mathcal E\subseteq\mathcal G$.
Dla dostateczności niech
$\mathcal H=\{B\subseteq Y:f^{-1}(B)\in\mathcal F\}$.
To $\sigma$-algebra na $Y$, ponieważ obrazy odwrotne zachowują
dopełnienia i przeliczalne sumy. Z założenia
$\mathcal E\subseteq\mathcal H$, zatem
$\mathcal G=\sigma_Y(\mathcal E)\subseteq\mathcal H$.
\end{proof}

\begin{twierdzenie}[Złożenia i funkcje ciągłe]
\label{thm:composition-measurable}
Niech $(X,\mathcal F)$, $(Y,\mathcal G)$ i $(Z,\mathcal H)$
będą przestrzeniami mierzalnymi. Jeżeli $f:X\to Y$ jest
$(\mathcal F,\mathcal G)$-mierzalne, a $g:Y\to Z$ jest
$(\mathcal G,\mathcal H)$-mierzalne, to $g\circ f$ jest
$(\mathcal F,\mathcal H)$-mierzalne.
Ponadto każde odwzorowanie ciągłe $g:Y\to Z$ pomiędzy
przestrzeniami topologicznymi jest
$(\mathcal B(Y),\mathcal B(Z))$-mierzalne.
\end{twierdzenie}
\begin{proof}
Dla $C\in\mathcal H$ mamy
$(g\circ f)^{-1}(C)=f^{-1}(g^{-1}(C))\in\mathcal F$.
Jeśli $g$ jest ciągłe, obrazy odwrotne zbiorów otwartych są otwarte;
z kryterium \ref{thm:measurable-generator} wynika jego mierzalność
borelowska.
\end{proof}

\begin{twierdzenie}[Funkcje rzeczywiste i działania]
\label{thm:measurable-real}
Niech $(X,\mathcal F)$ będzie przestrzenią mierzalną.
Dla funkcji $f:X\to\mathbb R$ następujące warunki są równoważne:
\begin{enumerate}
\item $f$ jest $(\mathcal F,\mathcal B(\mathbb R))$-mierzalna;
\item $\{x\in X:f(x)<q\}\in\mathcal F$ dla każdego $q\in\mathbb Q$;
\item $\{x\in X:f(x)\leq q\}\in\mathcal F$ dla każdego $q\in\mathbb Q$.
\end{enumerate}
Jeżeli $f,g:X\to\mathbb R$ są mierzalne, to funkcje
$f+g$, $fg$, $\max(f,g)$ i $\min(f,g)$ są mierzalne.
\end{twierdzenie}
\begin{proof}
Równoważność trzech warunków wynika z twierdzeń
\ref{thm:borel-generators-R} i \ref{thm:measurable-generator}.
Odwzorowanie $x\mapsto(f(x),g(x))$ do $\mathbb R^2$ jest mierzalne:
obrazy odwrotne prostokątów $B_1\times B_2$, gdzie
$B_1,B_2\in\mathcal B(\mathbb R)$, są przekrojami
$f^{-1}(B_1)\cap g^{-1}(B_2)$. Prostokąty te generują
$\mathcal B(\mathbb R^2)$ na mocy twierdzenia
\ref{thm:borel-Rd}. Funkcje na $\mathbb R^2$ określone przez
$(u,v)\mapsto u+v$, $(u,v)\mapsto uv$,
$(u,v)\mapsto\max(u,v)$ i $(u,v)\mapsto\min(u,v)$ są ciągłe.
Teza wynika więc z twierdzenia \ref{thm:composition-measurable}.
\end{proof}

\begin{przyklad}[Funkcje charakterystyczne zbiorów]
Dla $A\subseteq X$ definiujemy funkcję charakterystyczną
$\chi_A:X\to\mathbb R$ wzorem
$\chi_A(x)=1$ dla $x\in A$ i $\chi_A(x)=0$ dla $x\notin A$.
Wtedy $\chi_A$ jest mierzalna względem $\mathcal F$
dokładnie wtedy, gdy $A\in\mathcal F$.
Jeśli $A\in\mathcal F$, to dla każdego zbioru borelowskiego $B$
przeciwobraz $\chi_A^{-1}(B)$ jest równy $\varnothing$, $A$,
$X\setminus A$ lub $X$, zależnie od tego, czy $B$ nie zawiera
żadnej z liczb $0,1$, zawiera tylko $1$, tylko $0$, czy obie.
Wszystkie cztery zbiory należą do $\mathcal F$.
Konieczność wynika z równości
$A=\chi_A^{-1}((1/2,\infty))$.
\end{przyklad}

\begin{przyklad}[Mierzalność nie implikuje regularności geometrycznej]
Na $X=\mathbb R$ z $\mathcal F=\mathcal B(\mathbb R)$ funkcja
$\chi_{\mathbb Q}$ jest mierzalna, choć nie jest ciągła
w żadnym punkcie. Z drugiej strony, jeśli na nieprzeliczalnym
zbiorze $X$ przyjmiemy $\mathcal F=\{\varnothing,X\}$, to
każda mierzalna funkcja $X\to\mathbb R$ jest stała. Rzeczywiście,
gdyby przyjmowała wartości $a<b$, to dla pewnej liczby wymiernej
$q\in(a,b)$ zbiór $\{f<q\}$ byłby niepusty i różny od $X$.
\end{przyklad}

\begin{uwaga}[Co zostało, a co nie zostało wykazane]
Wprowadziliśmy pojęcie mierzalności bez definiowania jakiejkolwiek
miary. To rozróżnienie jest zasadnicze: przestrzeń mierzalna określa,
\emph{którym} zbiorom wolno przypisać rozmiar; miara określa
\emph{jaki} jest ten rozmiar. Miary pojawią się w rozdziale
\ref{ch:miary-wlasnosci}, a mierzalność funkcji o wartościach
w $[-\infty,\infty]$ omówimy szczegółowo przed zdefiniowaniem całki.
\end{uwaga}

% ROZDZIAŁ 2 — plik włączany przez main.tex
\chapter{Miary i ich podstawowe własności}
\label{ch:miary-wlasnosci}

\section{Definicja i pierwsze przykłady}

W rozdziale \ref{ch:przestrzenie-mierzalne} ustaliliśmy, które podzbiory
przestrzeni mają być mierzalne. Teraz przypiszemy im liczby z przedziału
$[0,\infty]$. Dopuszczenie wartości $\infty$ jest konieczne już dla
naturalnego pojęcia liczby elementów zbioru, a następnie dla miary
objętości na przestrzeniach nieograniczonych.

\begin{definicja}[Miara]
\label{def:measure}
Niech $(X,\mathcal F)$ będzie przestrzenią mierzalną.
\emph{Miarą} na $(X,\mathcal F)$ nazywamy funkcję
$\mu:\mathcal F\to[0,\infty]$ spełniającą:
\begin{enumerate}
\item $\mu(\varnothing)=0$;
\item dla każdego ciągu $(A_n)_{n\geq1}$ parami rozłącznych zbiorów
z $\mathcal F$ zachodzi \emph{przeliczalna addytywność}
\begin{equation}
\label{eq:countable-additivity}
\mu\left(\bigcup_{n=1}^{\infty}A_n\right)
=\sum_{n=1}^{\infty}\mu(A_n).
\end{equation}
\end{enumerate}
Trójkę $(X,\mathcal F,\mu)$ nazywamy \emph{przestrzenią z miarą}.
Wszystkie sumy liczb nieujemnych rozumiemy jako kres górny sum
częściowych, dopuszczając wartość $\infty$.
\end{definicja}

\begin{uwaga}
Nie jest wymagane, aby $\mu(X)$ było skończone. Aksjomat
$\mu(\varnothing)=0$ zapisujemy oddzielnie, ponieważ bez niego
funkcja stale równa $\infty$ spełniałaby równość
\eqref{eq:countable-additivity} dla każdej rodziny rozłącznej,
ale nie miałaby zamierzonych własności miary.
\end{uwaga}

\begin{przyklad}[Miara zerowa i miara Diraca]
\label{ex:dirac}
Na dowolnej przestrzeni mierzalnej $(X,\mathcal F)$ funkcja
$\mu(A)=0$ dla $A\in\mathcal F$ jest miarą zerową.
Dla ustalonego $x\in X$ definiujemy \emph{miarę Diraca} wzorem
\[
\delta_x(A)=\begin{cases}1,&x\in A,\\0,&x\notin A,\end{cases}
\qquad A\in\mathcal F.
\]
Jeśli $(A_n)$ są parami rozłączne, punkt $x$ należy co najwyżej
do jednego z tych zbiorów. Daje to przeliczalną addytywność.
Nie zakładamy przy tym, że singleton $\{x\}$ należy do $\mathcal F$.
\end{przyklad}

\begin{przyklad}[Miara licząca]
\label{ex:counting}
Na $(X,\mathcal P(X))$ definiujemy
\[
\#(A)=\begin{cases}
|A|,&A\text{ jest skończony},\\
\infty,&A\text{ jest nieskończony}.
\end{cases}
\]
Dla parami rozłącznych zbiorów $A_n$ liczność ich sumy jest sumą
liczności w podanym sensie. Jeśli $\bigcup_n A_n$ jest skończony,
tylko skończenie wiele zbiorów $A_n$ jest niepustych i każdy
jest skończony, więc równość addytywności wynika ze
zwykłego zliczania elementów. Jeśli $\bigcup_n A_n$
jest nieskończony, to albo pewien $A_n$ jest nieskończony,
albo nieskończenie wiele zbiorów $A_n$ jest niepustych.
W obu przypadkach $\sum_n\#(A_n)=\infty=\#(\bigcup_nA_n)$.
Jest to \emph{miara licząca}.
\end{przyklad}

\begin{przyklad}[Miara na $\sigma$-algebrze przeliczalno-koprzeliczalnej]
\label{ex:cocountable-measure}
Niech $X$ będzie nieprzeliczalnym zbiorem i niech
$\mathcal C_X$ będzie $\sigma$-algebrą z przykładu
\ref{ex:countable-cocountable}. Określmy
\[
\mu(A)=\begin{cases}
0,&A\text{ jest przeliczalny},\\
1,&X\setminus A\text{ jest przeliczalny}.
\end{cases}
\qquad A\in\mathcal C_X.
\]
Definicja jest jednoznaczna, ponieważ nieprzeliczalny zbiór
$X$ nie może być jednocześnie przeliczalny i mieć przeliczalne
dopełnienie. Dwie rozłączne części $X$ nie mogą obie mieć
przeliczalnego dopełnienia. Dla przeliczalnej rodziny
rozłącznych zbiorów z $\mathcal C_X$ albo wszystkie mają miarę
zero, a ich suma jest przeliczalna, albo dokładnie jeden ma
miarę jeden, a suma ma miarę jeden. Jest to zatem miara.
Ten przykład pokaże w sekcji \ref{sec:atoms}, że atomy
nie muszą być singletonami.
\end{przyklad}

\begin{przyklad}[Ważone miary dyskretne]
\label{ex:weighted-discrete}
Niech $X$ będzie zbiorem, $\mathcal F=\mathcal P(X)$,
$(x_n)_{n\geq1}$ ciągiem punktów $X$, a $(w_n)_{n\geq1}$
ciągiem liczb z $[0,\infty)$. Wtedy
\[
\mu(A)=\sum_{n=1}^{\infty}w_n\chi_A(x_n),
\qquad A\subseteq X,
\]
jest miarą. Dla parami rozłącznych $A_k$ mamy
$\chi_{\bigcup_k A_k}(x_n)=\sum_k\chi_{A_k}(x_n)$.
Możemy zamienić kolejność sumowania, ponieważ dla
$a_{n,k}=w_n\chi_{A_k}(x_n)\geq0$ obie sumy iterowane
$\sum_n\sum_k a_{n,k}$ i $\sum_k\sum_n a_{n,k}$
są równe kresowi górnemu sum
$\sum_{n\in F}\sum_{k\in G}a_{n,k}$
po wszystkich skończonych zbiorach indeksów $F,G\subseteq\mathbb N$.
Otrzymujemy więc przeliczalną addytywność, także gdy
któraś suma jest nieskończona. Nie wymagamy, aby punkty
$x_n$ były różne ani by $\sum_nw_n$ było skończone.
\end{przyklad}

\section{Monotoniczność i przeliczalna podaddytywność}

\begin{twierdzenie}[Podstawowe własności miary]
\label{thm:basic-measure}
Niech $(X,\mathcal F,\mu)$ będzie przestrzenią z miarą.
Wówczas:
\begin{enumerate}
\item dla dowolnych $A,B\in\mathcal F$ z $A\subseteq B$ zachodzi
$\mu(A)\leq\mu(B)$;
\item dla dowolnych $A,B\in\mathcal F$ zachodzi
$\mu(A\cup B)\leq\mu(A)+\mu(B)$;
\item dla dowolnych $A_n\in\mathcal F$, $n\geq1$, zachodzi
\begin{equation}
\label{eq:subadditivity}
\mu\left(\bigcup_{n=1}^{\infty}A_n\right)
\leq\sum_{n=1}^{\infty}\mu(A_n);
\end{equation}
\item jeśli $A,B\in\mathcal F$, $A\subseteq B$ i $\mu(A)<\infty$,
to $\mu(B\setminus A)=\mu(B)-\mu(A)$.
\end{enumerate}
\end{twierdzenie}
\begin{proof}
Zbiory $A$ i $B\setminus A$ są rozłączne, a ich suma jest równa $B$.
Przeliczalna addytywność, zastosowana do rodziny uzupełnionej
zbiorami pustymi, daje
$\mu(B)=\mu(A)+\mu(B\setminus A)$. Wynika stąd monotoniczność,
a jeżeli $\mu(A)<\infty$, również wzór na różnicę.

Ponadto $A\cup B=A\,\dot\cup\,(B\setminus A)$, więc
$\mu(A\cup B)\leq\mu(A)+\mu(B)$.
Dla dowolnego ciągu $(A_n)$ definiujemy
$B_1=A_1$ oraz
$B_n=A_n\setminus\bigcup_{k<n}A_k$ dla $n\geq2$.
Zbiory $B_n$ są parami rozłączne, $B_n\subseteq A_n$
i $\bigcup_nB_n=\bigcup_nA_n$. Stąd
$\mu(\bigcup_nA_n)=\sum_n\mu(B_n)\leq\sum_n\mu(A_n)$.
\end{proof}

\begin{uwaga}[Arytmetyka z nieskończonością]
Warunku $\mu(A)<\infty$ w ostatnim punkcie twierdzenia
\ref{thm:basic-measure} nie wolno usuwać. Jeśli
$\mu(A)=\mu(B)=\infty$, wyrażenie $\infty-\infty$
nie ma określonej wartości. Prawidłowym wzorem pozostaje
$\mu(B)=\mu(A)+\mu(B\setminus A)$, który nie wymaga odejmowania.
\end{uwaga}

\begin{wniosek}[Przeliczalne sumy zbiorów miary zero]
\label{cor:null-union}
Jeżeli $(X,\mathcal F,\mu)$ jest przestrzenią z miarą oraz
$N_n\in\mathcal F$ i $\mu(N_n)=0$ dla każdego $n\in\mathbb N$,
to $\mu(\bigcup_nN_n)=0$.
\end{wniosek}
\begin{proof}
Z podaddytywności mamy
$0\leq\mu(\bigcup_nN_n)\leq\sum_n\mu(N_n)=0$.
\end{proof}

\section{Ciągłość miary względem granic zbiorów}

\begin{twierdzenie}[Ciągłość z dołu i z góry]
\label{thm:measure-continuity}
Niech $(X,\mathcal F,\mu)$ będzie przestrzenią z miarą.
\begin{enumerate}
\item Jeśli $A_n\in\mathcal F$ oraz $A_n\uparrow A$, to
$A\in\mathcal F$ i
\begin{equation}
\label{eq:continuity-below}
\mu(A)=\lim_{n\to\infty}\mu(A_n).
\end{equation}
\item Jeśli $B_n\in\mathcal F$, $B_n\downarrow B$ oraz
$\mu(B_1)<\infty$, to $B\in\mathcal F$ i
\begin{equation}
\label{eq:continuity-above}
\mu(B)=\lim_{n\to\infty}\mu(B_n).
\end{equation}
\end{enumerate}
\end{twierdzenie}
\begin{proof}
Dla ciągu rosnącego przyjmujemy $E_1=A_1$ i
$E_n=A_n\setminus A_{n-1}$ dla $n\geq2$.
Zbiory $E_n$ są parami rozłączne oraz
$A=\bigcup_nE_n$ i $A_k=\bigcup_{n=1}^kE_n$.
Stąd
$\mu(A)=\sum_n\mu(E_n)=\lim_{k\to\infty}\sum_{n=1}^k\mu(E_n)
=\lim_k\mu(A_k)$.

Dla ciągu malejącego zbiory $C_n=B_1\setminus B_n$ rosną
do $B_1\setminus B$. Z pierwszej części mamy
$\mu(C_n)\to\mu(B_1\setminus B)$. Ponieważ
$\mu(B_1)<\infty$, wolno zastosować wzór na różnicę:
$\mu(C_n)=\mu(B_1)-\mu(B_n)$ oraz
$\mu(B_1\setminus B)=\mu(B_1)-\mu(B)$.
Po odjęciu od $\mu(B_1)$ otrzymujemy żądaną równość.
\end{proof}

\begin{uwaga}[Wystarczy skończoność miary jednego wyrazu]
W drugim punkcie twierdzenia~\ref{thm:measure-continuity} można
zastąpić założenie $\mu(B_1)<\infty$ słabszym
$\mu(B_m)<\infty$ dla pewnego $m\ge1$: wystarczy
zastosować twierdzenie do ogona ciągu
$(B_{m+k-1})_{k\ge1}$, którego przekrój jest nadal równy $B$.
Nie wystarcza natomiast sama skończoność $\mu(B)$.
\end{uwaga}

\begin{przyklad}[Dlaczego założenie skończoności jest niezbędne]
\label{ex:continuity-counterexample}
Na $X=\mathbb N$ z miarą liczącą zbiory
$B_n=\{n,n+1,n+2,\ldots\}$ tworzą ciąg malejący o pustym
przekroju. Każdy $B_n$ ma miarę $\infty$, ale
$\mu(\bigcap_nB_n)=\mu(\varnothing)=0$.
Tak więc ciągłość z góry nie zachodzi bez odpowiedniej
hipotezy skończoności.
\end{przyklad}

\begin{definicja}[Granica górna i dolna zbiorów]
\label{def:limsup-sets}
Dla ciągu $(A_n)_{n\geq1}$ podzbiorów $X$ definiujemy
\[
\liminf_{n\to\infty}A_n
=\bigcup_{m=1}^{\infty}\bigcap_{n\geq m}A_n,
\qquad
\limsup_{n\to\infty}A_n
=\bigcap_{m=1}^{\infty}\bigcup_{n\geq m}A_n.
\]
Pierwszy zbiór składa się z punktów należących do wszystkich
$A_n$ od pewnego miejsca, drugi --- z punktów należących do
nieskończenie wielu $A_n$. Jeśli oba zbiory są równe,
ich wspólną wartość oznaczamy przez $\lim_nA_n$.
\end{definicja}

\begin{twierdzenie}[Nierówności dla granic zbiorów]
\label{thm:fatou-for-sets}
Niech $(X,\mathcal F,\mu)$ będzie przestrzenią z miarą
oraz $A_n\in\mathcal F$ dla wszystkich $n\ge1$.
Wówczas
\[
 \mu(\liminf_{n\to\infty}A_n)
 \le \liminf_{n\to\infty}\mu(A_n).
\]
Jeżeli dodatkowo $\mu(X)<\infty$, to
\[
 \limsup_{n\to\infty}\mu(A_n)
 \le \mu(\limsup_{n\to\infty}A_n).
\]
\end{twierdzenie}
\begin{proof}
Niech $C_m=\bigcap_{n\ge m}A_n$.
Zbiory $C_m$ tworzą ciąg rosnący,
$\bigcup_m C_m=\liminf_n A_n$ oraz
$\mu(C_m)\le\inf_{n\ge m}\mu(A_n)$.
Ciągłość miary z dołu daje
$\mu(\liminf_n A_n)=\lim_m\mu(C_m)
\le\lim_m\inf_{n\ge m}\mu(A_n)$.
Dla drugiej nierówności kładziemy
$D_m=\bigcup_{n\ge m}A_n$.
Zbiory $D_m$ maleją do $\limsup_n A_n$,
a $\mu(D_m)\ge\sup_{n\ge m}\mu(A_n)$.
Ponieważ $\mu(D_1)\le\mu(X)<\infty$,
ciągłość miary z góry pozwala przejść do granicy:
$\mu(\limsup_n A_n)=\lim_m\mu(D_m)
\ge\lim_m\sup_{n\ge m}\mu(A_n)$.
\end{proof}
\begin{uwaga}
Założenie skończoności jest istotne w drugiej nierówności.
Dla miary liczącej na $\mathbb N$ i
$A_n=\{n,n+1,\ldots\}$ mamy
$\limsup_n\mu(A_n)=\infty$, lecz
$\limsup_n A_n=\varnothing$.
\end{uwaga}

\begin{twierdzenie}[Lemat Borela--Cantellego, część pierwsza]
\label{thm:borel-cantelli-1}
Niech $(X,\mathcal F,\mu)$ będzie przestrzenią z miarą
oraz $(A_n)_{n\geq1}$ ciągiem zbiorów z $\mathcal F$.
Jeżeli $\sum_{n=1}^{\infty}\mu(A_n)<\infty$, to
$\mu(\limsup_{n\to\infty}A_n)=0$.
\end{twierdzenie}
\begin{proof}
Zbiór $\limsup_nA_n$ jest mierzalny i dla każdego
$m\in\mathbb N$ zawiera się w $\bigcup_{n\geq m}A_n$.
Z monotoniczności i podaddytywności wynika
\[
0\leq\mu(\limsup_nA_n)
\leq\sum_{n\geq m}\mu(A_n).
\]
Ponieważ ogólny szereg jest zbieżny, jego ogon dąży do zera,
gdy $m\to\infty$. To kończy dowód. Nie użyliśmy ani
skończoności $\mu(X)$, ani niezależności zbiorów.
\end{proof}

\begin{przyklad}[Interpretacja]
Jeśli $\mu(A_n)\leq2^{-n}$ dla wszystkich $n\geq1$, to
$\mu(\limsup_nA_n)=0$. Prawie każdy punkt (w sensie, który
zaraz zdefiniujemy) może należeć tylko do skończenie wielu $A_n$.
W rachunku prawdopodobieństwa po wprowadzeniu niezależności
uzyskuje się drugą część lematu Borela--Cantellego, z innymi
założeniami i przeciwnym wnioskiem.
\end{przyklad}

\section{Skończoność, \texorpdfstring{$\sigma$}{sigma}-skończoność i półskończoność}

\begin{definicja}[Klasy miar]
\label{def:measure-finiteness}
Niech $(X,\mathcal F,\mu)$ będzie przestrzenią z miarą.
\begin{enumerate}
\item Miara $\mu$ jest \emph{skończona}, gdy $\mu(X)<\infty$.
\item Miara $\mu$ jest \emph{probabilistyczna}, gdy $\mu(X)=1$.
\item Miara $\mu$ jest \emph{$\sigma$-skończona}, gdy istnieją
$E_n\in\mathcal F$, $n\geq1$, takie że
$X=\bigcup_nE_n$ i $\mu(E_n)<\infty$ dla wszystkich $n$.
\item Miara $\mu$ jest \emph{półskończona}, gdy dla każdego
$A\in\mathcal F$ z $\mu(A)>0$ istnieje $B\in\mathcal F$
takie, że $B\subseteq A$ oraz $0<\mu(B)<\infty$.
\end{enumerate}
Jeśli $(X,\tau)$ jest przestrzenią topologiczną, a
$\mathcal F\supseteq\mathcal B(X)$, nazywamy $\mu$
\emph{lokalnie skończoną}, gdy każdy $x\in X$ posiada
otwarte otoczenie $U$ o $\mu(U)<\infty$.
\end{definicja}

\begin{twierdzenie}[Implikacje]
\label{thm:finiteness-implications}
Każda miara skończona na przestrzeni mierzalnej jest
$\sigma$-skończona, a każda miara $\sigma$-skończona
jest półskończona. Implikacje odwrotne są na ogół fałszywe.
\end{twierdzenie}
\begin{proof}
Jeśli $\mu(X)<\infty$, można w definicji $\sigma$-skończoności
przyjąć $E_n=X$. Niech teraz $X=\bigcup_nE_n$ z
$\mu(E_n)<\infty$. Dla $A\in\mathcal F$ o $\mu(A)>0$
nie wszystkie przekroje $A\cap E_n$ mogą mieć miarę zero:
w przeciwnym razie podaddytywność dałaby
$\mu(A)\leq\sum_n\mu(A\cap E_n)=0$.
Istnieje więc $n$ z $0<\mu(A\cap E_n)<\infty$.

Miara licząca na $\mathbb N$ jest $\sigma$-skończona,
ale nieskończona. Miara licząca na nieprzeliczalnym
zbiorze $X$ jest półskończona (każdy niepusty zbiór
zawiera singleton miary jeden), ale nie jest
$\sigma$-skończona, gdyż zbiory o skończonej mierze
są skończone, a przeliczalna suma zbiorów skończonych
jest przeliczalna.
\end{proof}

\begin{przyklad}[Miara, która nie jest półskończona]
Niech $X$ będzie niepustym zbiorem z $\mathcal F=\mathcal P(X)$
i niech $\mu(\varnothing)=0$ oraz $\mu(A)=\infty$
dla $A\ne\varnothing$. Jest to miara: dla rodziny
rozłącznej z niepustą sumą co najmniej jeden składnik
ma miarę $\infty$. Nie ma jednak żadnego podzbioru
o mierze dodatniej i skończonej.
\end{przyklad}

\begin{uwaga}[Lokalna skończoność a $\sigma$-skończoność]
Na przestrzeni topologicznej mającej przeliczalną bazę
lokalnie skończona miara borelowska jest $\sigma$-skończona.
Istotnie, otwarte otoczenia o skończonej mierze tworzą
pokrycie; własność Lindelöfa (wynikająca z istnienia
przeliczalnej bazy) pozwala wybrać pokrycie przeliczalne.
Bez założenia topologicznego takiej implikacji nie ma:
miara licząca na nieprzeliczalnej przestrzeni dyskretnej
jest lokalnie skończona, lecz nie jest $\sigma$-skończona.
\end{uwaga}

\section{Zbiory miary zero, własności prawie wszędzie i zupełność}

\begin{definicja}[Własność prawie wszędzie]
\label{def:almost-everywhere}
Niech $(X,\mathcal F,\mu)$ będzie przestrzenią z miarą.
Mówimy, że własność $P(x)$ zachodzi \emph{$\mu$-prawie wszędzie}
(skrót: $\mu$-p.w.), jeśli istnieje $N\in\mathcal F$
o $\mu(N)=0$ takie, że $P(x)$ zachodzi dla każdego
$x\in X\setminus N$. Nie wymagamy mierzalności zbioru
$\{x\in X:P(x)\text{ nie zachodzi}\}$.
Jeśli jest on mierzalny, jest to równoważne temu,
że ma miarę zero.
\end{definicja}

\begin{uwaga}
Dzięki wnioskowi \ref{cor:null-union} przeliczalnie wiele
własności zachodzących $\mu$-p.w. zachodzi jednocześnie
$\mu$-p.w. Bez przeliczalności nie można wyciągać tego
wniosku: przy miarze z przykładu \ref{ex:cocountable-measure}
każdy singleton ma miarę zero, lecz suma wszystkich singletonów
jest równa $X$ i ma miarę jeden.
\end{uwaga}

\begin{definicja}[Miara zupełna]
\label{def:complete-measure}
Miara $\mu$ na $(X,\mathcal F)$ jest \emph{zupełna}, jeżeli
dla każdego $N\in\mathcal F$ spełniającego $\mu(N)=0$
każdy podzbiór $A\subseteq N$ należy do $\mathcal F$.
Z monotoniczności wynika wówczas $\mu(A)=0$.
\end{definicja}

\begin{przyklad}[Miara niezupełna na trzech punktach]
\label{ex:incomplete-measure}
Niech $X=\{a,b,c\}$ oraz
$\mathcal F=\{\varnothing,\{a,b\},\{c\},X\}$.
Określmy $\mu(\{a,b\})=0$ i $\mu(\{c\})=1$,
a pozostałe wartości wyznaczmy przez addytywność:
$\mu(\varnothing)=0$, $\mu(X)=1$.
To dobrze określona miara na $\mathcal F$.
Zbiór $N=\{a,b\}$ ma miarę zero, ale jego podzbiór
$\{a\}$ nie jest mierzalny. Zatem $\mu$ nie jest zupełna.
\end{przyklad}

\begin{twierdzenie}[Uzupełnienie przestrzeni z miarą]
\label{thm:completion}
Niech $(X,\mathcal F,\mu)$ będzie przestrzenią z miarą.
Zdefiniujmy
\[
\overline{\mathcal F}
=\{E\subseteq X:\text{istnieją }A,B\in\mathcal F
\text{ takie, że }A\subseteq E\subseteq B
\text{ oraz }\mu(B\setminus A)=0\}.
\]
Wówczas $\overline{\mathcal F}$ jest $\sigma$-algebrą zawierającą
$\mathcal F$, a wzór
\[
\overline\mu(E)=\mu(A),\qquad
A\subseteq E\subseteq B,\quad A,B\in\mathcal F,
\quad\mu(B\setminus A)=0,
\]
definiuje na niej miarę zupełną. Miara $\overline\mu$
przedłuża $\mu$, a $\overline{\mathcal F}$ jest najmniejszą
$\sigma$-algebrą zawierającą $\mathcal F$ i wszystkie
podzbiory zbiorów z $\mathcal F$ o mierze zero.
\end{twierdzenie}
\begin{proof}
Najpierw sprawdzamy jednoznaczność definicji $\overline\mu$.
Jeśli $A_i\subseteq E\subseteq B_i$ i
$\mu(B_i\setminus A_i)=0$ dla $i=1,2$, to
$A_1\setminus A_2\subseteq B_2\setminus A_2$.
Z monotoniczności mamy $\mu(A_1\setminus A_2)=0$,
a stąd $\mu(A_1)\leq\mu(A_2)$, nawet jeśli
jedna z miar jest nieskończona. Zamiana indeksów
daje $\mu(A_1)=\mu(A_2)$. Ponadto
$\mu(A_i)=\mu(B_i)$, gdyż
$B_i=A_i\,\dot\cup\,(B_i\setminus A_i)$.

Jeśli $E\in\overline{\mathcal F}$ i
$A\subseteq E\subseteq B$ są jak w tezie, to
$X\setminus B\subseteq X\setminus E\subseteq X\setminus A$,
a różnica zbiorów skrajnych wynosi $B\setminus A$
i ma miarę zero. To dowodzi zamknięcia na dopełnienia.
Jeśli $E_n\in\overline{\mathcal F}$ i
$A_n\subseteq E_n\subseteq B_n$ z
$\mu(B_n\setminus A_n)=0$, to
\[
\bigcup_nA_n\subseteq\bigcup_nE_n\subseteq\bigcup_nB_n,
\qquad
\left(\bigcup_nB_n\right)\setminus\left(\bigcup_nA_n\right)
\subseteq\bigcup_n(B_n\setminus A_n).
\]
Ostatni zbiór ma miarę zero na mocy wniosku
\ref{cor:null-union}, więc przeliczalna suma należy do
$\overline{\mathcal F}$. Zawieranie $\mathcal F$ wynika
z wyboru $A=B=E\in\mathcal F$.

Dla parami rozłącznych $E_n\in\overline{\mathcal F}$
wybieramy $A_n\subseteq E_n$ jak wyżej. Zbiory $A_n$
są parami rozłączne i
$\overline\mu(\bigcup_nE_n)=\mu(\bigcup_nA_n)
=\sum_n\mu(A_n)=\sum_n\overline\mu(E_n)$.
Zatem $\overline\mu$ jest miarą i przedłuża $\mu$.

Aby wykazać zupełność, niech $E\in\overline{\mathcal F}$
i $\overline\mu(E)=0$. Wybierzmy $A,B\in\mathcal F$
takie, że $A\subseteq E\subseteq B$ i
$\mu(B\setminus A)=0$. Mamy $\mu(B)=\mu(A)=0$.
Dla dowolnego $T\subseteq E$ zachodzi
$\varnothing\subseteq T\subseteq B$, przy $\mu(B)=0$,
więc $T\in\overline{\mathcal F}$.

Wreszcie, jeśli $\mathcal G$ jest $\sigma$-algebrą zawierającą
$\mathcal F$ i wszystkie podzbiory jej zbiorów miary zero,
to dla $E\in\overline{\mathcal F}$ jak wyżej zbiór
$E\setminus A$ jest podzbiorem mierzalnego zbioru
$B\setminus A$ miary zero. Zatem $E\setminus A\in\mathcal G$
i $E=A\cup(E\setminus A)\in\mathcal G$. Dowodzi to minimalności.
\end{proof}

\begin{uwaga}[Jednoznaczność uzupełnienia]
Konstrukcja z twierdzenia \ref{thm:completion} nie polega
na arbitralnym wyborze wartości dla nowych podzbiorów.
Zbiory nowo dołączone różnią się od starych zbiorów
mierzalnych jedynie o części zbiorów miary zero.
Rozróżnienie między miarą borelowską a jej uzupełnieniem
będzie ważne przy konstrukcji miary Lebesgue'a.
\end{uwaga}

\section{Operacje na miarach}

\begin{twierdzenie}[Sumy i mnożenie przez stałą]
\label{thm:measure-operations}
Niech $(X,\mathcal F)$ będzie przestrzenią mierzalną,
$(\mu_n)_{n\geq1}$ ciągiem miar na $\mathcal F$,
a $(c_n)_{n\geq1}$ ciągiem liczb nieujemnych i skończonych.
Wówczas
$\mu(A)=\sum_{n=1}^{\infty}c_n\mu_n(A)$ dla $A\in\mathcal F$
jest miarą, przy konwencji, że $0\mu_n$ oznacza miarę
zerową również wtedy, gdy $\mu_n(A)=\infty$.
\end{twierdzenie}
\begin{proof}
Wszystkie składniki są nieujemne i wartość na zbiorze pustym
wynosi zero. Dla parami rozłącznych $A_k$ otrzymujemy
\[
\mu\left(\bigcup_kA_k\right)
=\sum_n c_n\sum_k\mu_n(A_k)
=\sum_k\sum_n c_n\mu_n(A_k)
=\sum_k\mu(A_k).
\]
Wolno przestawić kolejność sumowania, ponieważ obie strony
środkowej równości są kresami górnymi tych samych sum
skończonych wyrazów nieujemnych.
\end{proof}

\begin{definicja}[Miara śladowa i ograniczenie]
\label{def:restriction}
Niech $(X,\mathcal F,\mu)$ będzie przestrzenią z miarą
i $E\in\mathcal F$. Na przestrzeni mierzalnej
$(E,\mathcal F|_E)$ definiujemy \emph{miarę śladową}
$\mu|_E$ wzorem
$(\mu|_E)(B)=\mu(B)$ dla $B\in\mathcal F|_E$.
Ponieważ $E\in\mathcal F$, każdy zbiór śladowy
$B=A\cap E$ jest elementem $\mathcal F$.
Można też określić na $(X,\mathcal F)$ miarę
$\mu_E(A)=\mu(A\cap E)$; obie konstrukcje mają
różne dziedziny i nie należy ich utożsamiać.
\end{definicja}

\begin{proof}[Sprawdzenie aksjomatów]
Miary obu typów przypisują zbiorowi pustemu wartość zero.
Jeśli $B_n\subseteq E$ są parami rozłączne, to ich
suma jest rozłączna również w $X$, a przeliczalna
addytywność $\mu$ daje addytywność $\mu|_E$.
Dla $\mu_E$ korzystamy z tożsamości
$(\bigcup_nA_n)\cap E=\bigcup_n(A_n\cap E)$.
\end{proof}

\begin{definicja}[Miara obrazowa]
\label{def:pushforward}
Niech $(X,\mathcal F,\mu)$ będzie przestrzenią z miarą,
$(Y,\mathcal G)$ przestrzenią mierzalną, a
$T:X\to Y$ odwzorowaniem $(\mathcal F,\mathcal G)$-mierzalnym.
\emph{Miarą obrazową} $T_\#\mu$ na $(Y,\mathcal G)$ nazywamy
miarę określoną wzorem
\[
(T_\#\mu)(B)=\mu(T^{-1}(B)),\qquad B\in\mathcal G.
\]
\end{definicja}

\begin{proof}[Sprawdzenie aksjomatów]
Z mierzalności $T$ wszystkie obrazy odwrotne występujące
we wzorze są mierzalne. Obraz odwrotny zbioru pustego
jest pusty. Dla parami rozłącznych $B_n$ zbiory
$T^{-1}(B_n)$ są parami rozłączne i
$T^{-1}(\bigcup_nB_n)=\bigcup_nT^{-1}(B_n)$.
Równość przeliczalnej addytywności jest zatem konsekwencją
tego samego aksjomatu dla $\mu$.
\end{proof}

\begin{przyklad}[Obraz miary Diraca]
Jeżeli $x\in X$ i $T:X\to Y$ jest mierzalne, to
$T_\#\delta_x=\delta_{T(x)}$.
W szczególności obraz miary dyskretnej
$\sum_nw_n\delta_{x_n}$ jest równy
$\sum_nw_n\delta_{T(x_n)}$; punkty o tym samym obrazie
mogą powodować sumowanie wag.
\end{przyklad}

\begin{przyklad}[Normalizacja]
Jeżeli miara $\mu$ na $(X,\mathcal F)$ spełnia
$0<\mu(X)<\infty$, to $\mathbb P(A)=\mu(A)/\mu(X)$
jest miarą probabilistyczną. Nie można tego zrobić
w ten sposób dla miary o całkowitej masie nieskończonej.
\end{przyklad}

\section{Jednoznaczność miar na układach generujących}

\begin{twierdzenie}[Jednoznaczność miar skończonych]
\label{thm:uniqueness-finite}
Niech $X$ będzie zbiorem, $\mathcal P$ układem $\pi$
podzbiorów $X$, a $\mathcal F=\sigma(\mathcal P)$.
Jeżeli $\mu,\nu$ są miarami skończonymi na $(X,\mathcal F)$
spełniającymi $\mu(X)=\nu(X)$ i
$\mu(P)=\nu(P)$ dla każdego $P\in\mathcal P$,
to $\mu=\nu$ na całej $\mathcal F$.
\end{twierdzenie}
\begin{proof}
Niech $\mathcal D=\{A\in\mathcal F:\mu(A)=\nu(A)\}$.
Z założenia $X\in\mathcal D$ i $\mathcal P\subseteq\mathcal D$.
Jeżeli $A,B\in\mathcal D$ i $A\subseteq B$, to
$\mu(B\setminus A)=\mu(B)-\mu(A)
=\nu(B)-\nu(A)=\nu(B\setminus A)$.
Odejmowanie jest dopuszczalne, gdyż obie miary są skończone.
Dla parami rozłącznych $A_n\in\mathcal D$ mamy
$\mu(\bigcup_nA_n)=\sum_n\mu(A_n)
=\sum_n\nu(A_n)=\nu(\bigcup_nA_n)$.
Zatem $\mathcal D$ jest układem Dynkina zawierającym
$\mathcal P$. Twierdzenie \ref{thm:pi-lambda} daje
$\mathcal F\subseteq\mathcal D$.
\end{proof}

\begin{twierdzenie}[Jednoznaczność przy wspólnym wyczerpaniu]
\label{thm:uniqueness-sigma-finite}
Niech $X$ będzie zbiorem, $\mathcal P$ układem $\pi$ na $X$,
$\mathcal F=\sigma(\mathcal P)$, a $\mu,\nu$ miarami
na $(X,\mathcal F)$ zgodnymi na $\mathcal P$.
Załóżmy, że istnieją zbiory $E_n\in\mathcal P$ takie, że
$E_n\subseteq E_{n+1}$, $\bigcup_nE_n=X$ oraz
$\mu(E_n)=\nu(E_n)<\infty$ dla każdego $n$.
Wtedy $\mu=\nu$ na $\mathcal F$.
\end{twierdzenie}
\begin{proof}
Dla każdego $n$ definiujemy miary skończone na $\mathcal F$:
$\mu_n(A)=\mu(A\cap E_n)$ i
$\nu_n(A)=\nu(A\cap E_n)$.
Ponieważ $E_n\in\mathcal P$ i $\mathcal P$ jest układem $\pi$,
dla każdego $P\in\mathcal P$ zachodzi
$P\cap E_n\in\mathcal P$, a więc
$\mu_n(P)=\nu_n(P)$. Ponadto
$\mu_n(X)=\nu_n(X)<\infty$. Z twierdzenia
\ref{thm:uniqueness-finite} wynika $\mu_n=\nu_n$.
Dla dowolnego $A\in\mathcal F$ zbiory
$A\cap E_n$ rosną do $A$. Ciągłość miary z dołu
daje
$\mu(A)=\lim_n\mu(A\cap E_n)
=\lim_n\nu(A\cap E_n)=\nu(A)$.
\end{proof}

\begin{przyklad}[Miara określona wartościami dystrybuanty]
\label{ex:unique-cdf}
Niech $\mu$ i $\nu$ będą miarami probabilistycznymi
na $(\mathbb R,\mathcal B(\mathbb R))$.
Jeśli dla każdego $t\in\mathbb R$ zachodzi
$\mu((-\infty,t])=\nu((-\infty,t])$, to $\mu=\nu$.
Wystarczy bowiem zastosować twierdzenie
\ref{thm:uniqueness-finite} do układu $\pi$
$\{(-\infty,t]:t\in\mathbb R\}$.
Układ ten generuje $\mathcal B(\mathbb R)$,
a równość całkowitych mas wynika z probabilistyczności.
To wyjaśnia, dlaczego dystrybuanta jednoznacznie
wyznacza rozkład zmiennej losowej.
\end{przyklad}

\begin{przyklad}[Dlaczego wartości $\infty$ nie wystarczają]
\label{ex:uniqueness-fails}
Na $X=\mathbb R$ wyposażonym w $\mathcal B(\mathbb R)$
niech $\mu$ będzie miarą liczącą, a $\nu=2\mu$.
Obie miary przypisują wartość $\infty$ każdemu niepustemu
otwartemu przedziałowi o końcach wymiernych oraz wartość
zero zbiorowi pustemu. Zgadzają się zatem na układzie
$\pi$ tych przedziałów (razem ze zbiorem pustym),
który generuje $\mathcal B(\mathbb R)$.
Jednak $\mu(\{0\})=1$, a $\nu(\{0\})=2$.
Nie ma tu wspólnego wyczerpania zbiorami z układu $\pi$
o skończonych miarach. Przykład dowodzi, że zgodność
na generatorach bez odpowiednich warunków skończoności
nie wyznacza miary jednoznacznie.
\end{przyklad}

\section{Atomy miar}
\label{sec:atoms}

\begin{definicja}[Atom i miara bezatomowa]
\label{def:atom}
Niech $(X,\mathcal F,\mu)$ będzie przestrzenią z miarą.
Zbiór $A\in\mathcal F$ nazywamy \emph{atomem miary $\mu$},
jeżeli $\mu(A)>0$ oraz dla każdego $B\in\mathcal F$
zawartego w $A$ zachodzi $\mu(B)=0$ lub
$\mu(A\setminus B)=0$.
Miarę $\mu$ nazywamy \emph{bezatomową}, jeśli nie ma atomów.
\end{definicja}

\begin{przyklad}[Atomy miary dyskretnej]
Dla miary Diraca $\delta_x$ każdy mierzalny zbiór
zawierający $x$ jest atomem; wszystkie takie zbiory
są równoważne z dokładnością do zbiorów miary zero.
Dla miary liczącej na $\mathcal P(X)$ atomami są
dokładnie singletony: zbiór liczący co najmniej dwa punkty
można rozłożyć na dwie niepuste mierzalne części
o dodatniej mierze.
\end{przyklad}

\begin{przyklad}[Atom, który nie jest singletonem]
W przestrzeni z przykładu \ref{ex:cocountable-measure}
cały zbiór $X$ jest atomem: jego mierzalne podzbiory
mają miarę zero albo ich dopełnienia mają miarę zero.
Jednocześnie każdy singleton ma miarę zero.
Zatem nawet miara probabilistyczna może być atomowa,
chociaż żaden punkt nie ma dodatniej miary. Wniosek
``miara bez atomów wtedy i tylko wtedy, gdy wszystkie
singletony mają miarę zero'' jest bez dodatkowych
założeń fałszywy.
\end{przyklad}

\begin{twierdzenie}[Atomy miar $\sigma$-skończonych mają skończoną masę]
\label{thm:atom-sigma-finite}
Jeżeli $\mu$ jest miarą $\sigma$-skończoną na
$(X,\mathcal F)$ i $A\in\mathcal F$ jest atomem $\mu$,
to $0<\mu(A)<\infty$.
\end{twierdzenie}
\begin{proof}
Niech $X=\bigcup_nE_n$ z $\mu(E_n)<\infty$.
Gdyby wszystkie $A\cap E_n$ miały miarę zero,
to $\mu(A)=0$ na mocy przeliczalnej podaddytywności.
Istnieje więc $n$ z $\mu(A\cap E_n)>0$.
Atomowość $A$ implikuje $\mu(A\setminus E_n)=0$.
Otrzymujemy
$\mu(A)=\mu(A\cap E_n)+\mu(A\setminus E_n)
=\mu(A\cap E_n)\leq\mu(E_n)<\infty$.
\end{proof}

\begin{wniosek}[Przeliczalność rozłącznych atomów]
\label{cor:countably-many-disjoint-atoms}
Niech $\mu$ będzie miarą $\sigma$-skończoną na
$(X,\mathcal F)$, a $(A_i)_{i\in I}$ rodziną parami rozłącznych
atomów miary $\mu$. Wtedy zbiór indeksów $I$ jest co najwyżej
przeliczalny.
\end{wniosek}
\begin{proof}
Wybierzmy $X=\bigcup_{n\ge1}E_n$ z $\mu(E_n)<\infty$.
Z twierdzenia~\ref{thm:atom-sigma-finite} każdy $A_i$ ma
miarę dodatnią i skończoną. Istnieje więc $n$ takie, że
$\mu(A_i\cap E_n)>0$; z atomowości wynika wówczas
$\mu(A_i\cap E_n)=\mu(A_i)$. Dla ustalonych $n,k\ge1$
rodzina tych indeksów $i$, dla których
$\mu(A_i\cap E_n)\ge1/k$, jest skończona:
w przeciwnym przypadku suma miar dowolnie wielu
rozłącznych fragmentów przekroczyłaby $\mu(E_n)$.
Każdy indeks należy do przynajmniej jednej takiej
rodziny, więc $I$ jest przeliczalną sumą zbiorów
skończonych.
\end{proof}

\begin{uwaga}[Co daje $\sigma$-skończoność?]
$\sigma$-skończoność nie jest jedynie wygodnym
wzmocnieniem skończoności. Stanowi istotne założenie
jednoznaczności rozszerzenia miary, wielu wersji
twierdzenia Fubiniego oraz twierdzenia Radona--Nikodyma.
Nie należy jednak automatycznie zastępować
$\sigma$-skończoności półskończonością: są to różne hipotezy.
\end{uwaga}

\section{Synteza rozdziału i perspektywa dalszych konstrukcji}

Z rozdziałów \ref{ch:przestrzenie-mierzalne} i
\ref{ch:miary-wlasnosci} otrzymaliśmy następujący obraz.
$\sigma$-algebra kontroluje stabilność zbiorów mierzalnych
przy granicach przeliczalnych, zaś miara przekształca
rozłączne sumy w szeregi liczb nieujemnych. Te dwa
aksjomaty wystarczają do otrzymania monotoniczności,
podaddytywności, ciągłości z dołu i z góry oraz pierwszych
argumentów typu ``prawie wszędzie''.

Pozostaje jednak zasadnicze pytanie konstrukcyjne:
jak otrzymać miarę o zadanych wartościach na prostych
zbiorach, na przykład na prostokątach w $\mathbb R^d$?
Twierdzenia o jednoznaczności mówią, że pewne dane
wyznaczają miarę co najwyżej jedną; nie zapewniają
jej istnienia. W następnym rozdziale przejdziemy
do miar zewnętrznych i konstrukcji Carathéodory'ego,
która pozwoli przejść od prostych danych do prawdziwej
przeliczalnie addytywnej miary.

% ROZDZIAŁ 3 — plik włączany przez main.tex
\chapter{Miary zewnętrzne i konstrukcja Carathéodory'ego}
\label{ch:caratheodory}

\providecommand{\atlasref}[2]{#2}


\section{Problem konstrukcji miary}

W rozdziale \ref{ch:miary-wlasnosci} badaliśmy miarę już określoną na
$\sigma$-algebrze. Teraz odwrócimy pytanie: jeśli znamy ``rozmiary''
prostych zbiorów, to jak zbudować miarę na znacznie większej rodzinie?
Przykładowo na prostej chcemy przypisać przedziałowi $[a,b)$ długość
$b-a$, a następnie mierzyć zbiory otrzymane przez przeliczalne operacje.
Sam warunek skończonej addytywności jest za słaby. Kluczową rolę odgrywa
przeliczalna addytywność oraz możliwość pokrywania dowolnego zbioru
zbiorami prostszymi.

W całym rozdziale $X$ oznacza ustalony zbiór, a wszystkie funkcje
zbiorów przyjmują wartości w $[0,\infty]$, jeśli wyraźnie nie
zaznaczono inaczej. Przyjmujemy $\inf\varnothing=\infty$ oraz
$\sum_{n=1}^\infty a_n=\sup_{N\ge1}\sum_{n=1}^N a_n$ dla $a_n\ge0$.

\section{Miara zewnętrzna: definicja i własności}

\begin{definicja}[Miara zewnętrzna]
\label{def:outer-measure}
Funkcję $\mu^*:\mathcal P(X)\to[0,\infty]$ nazywamy
\emph{miarą zewnętrzną}, jeśli:
\begin{enumerate}[label=\textup{(\roman*)}]
\item $\mu^*(\varnothing)=0$;
\item $A\subseteq B$ implikuje $\mu^*(A)\le\mu^*(B)$;
\item dla dowolnego ciągu $(A_n)_{n\ge1}$ podzbiorów $X$ zachodzi
\begin{equation}\label{eq:outer-subadd}
\mu^*\Big(\bigcup_{n=1}^\infty A_n\Big)
\le\sum_{n=1}^\infty\mu^*(A_n).
\end{equation}
\end{enumerate}
Nie zakładamy, że $\mu^*$ jest miarą na $\mathcal P(X)$:
na ogół nie jest addytywna nawet dla rozłącznych zbiorów.
\end{definicja}

\begin{przyklad}[Najprostsze miary zewnętrzne]
Dla każdego zbioru $X$ funkcja $\mu^*(A)=0$ jest miarą zewnętrzną.
Miara licząca $\#(A)$ z przykładu \ref{ex:counting}, traktowana
jako funkcja na $\mathcal P(X)$, także jest miarą zewnętrzną.
Bardziej pouczający przykład otrzymujemy na $X=\{a,b\}$, gdy
$\mu^*(\varnothing)=0$ i $\mu^*(A)=1$ dla każdego $A\ne\varnothing$.
Wtedy $\mu^*(\{a\}\cup\{b\})=1\ne2=\mu^*(\{a\})+\mu^*(\{b\})$.
Wkrótce zobaczymy, że jej zbiory mierzalne w sensie Carathéodory'ego
to jedynie $\varnothing$ i $X$.
\end{przyklad}

\subsection{Miara zewnętrzna z kosztu pokryć}

\begin{twierdzenie}[Konstrukcja przez pokrycia]
\label{thm:outer-from-cover}
Niech $\mathcal C\subseteq\mathcal P(X)$ będzie rodziną zawierającą
$\varnothing$ oraz taką, że $X$ jest sumą przeliczalnie wielu zbiorów
z $\mathcal C$. Niech $c:\mathcal C\to[0,\infty]$ spełnia
$c(\varnothing)=0$. Dla każdego $E\subseteq X$ określmy
\begin{equation}\label{eq:outer-cover}
\mu_c^*(E):=\inf\Big\{\sum_{n=1}^\infty c(C_n):
C_n\in\mathcal C,\quad E\subseteq\bigcup_{n=1}^\infty C_n\Big\}.
\end{equation}
Wówczas $\mu_c^*$ jest miarą zewnętrzną.
\end{twierdzenie}

\begin{proof}
Pusty zbiór można pokryć ciągiem zbiorów pustych, więc
$0\le\mu_c^*(\varnothing)\le0$. Jeżeli $E\subseteq F$, to każde
pokrycie $F$ dopuszczalne w \eqref{eq:outer-cover} jest pokryciem $E$;
przy infimum po większej rodzinie nie można uzyskać większej wartości.

Pozostaje podaddytywność. Niech $E=\bigcup_{j\ge1}E_j$.
Jeżeli $\sum_j\mu_c^*(E_j)=\infty$, teza jest automatyczna.
W przeciwnym przypadku, przy ustalonym $\varepsilon>0$, dla każdego
$j$ wybieramy pokrycie $E_j\subseteq\bigcup_{k\ge1}C_{j,k}$, dla
którego
$\sum_{k\ge1}c(C_{j,k})<\mu_c^*(E_j)+\varepsilon2^{-j}$.
Taki wybór jest możliwy z definicji infimum; sumy skończone też
można uzupełnić pustymi zbiorami. Rodzina
$\{C_{j,k}:j,k\ge1\}$ jest przeliczalna i pokrywa $E$. Zatem
\[
\mu_c^*(E)\le \sum_{j,k\ge1}c(C_{j,k})
\le\sum_{j\ge1}\mu_c^*(E_j)+\varepsilon.
\]
Ponieważ $\varepsilon>0$ było dowolne, otrzymujemy
\eqref{eq:outer-subadd}. Przestawianie sum jest tu dozwolone,
gdyż wszystkie składniki są nieujemne: obie sumy to kresy górne
sum po skończonych podzbiorach $\mathbb N^2$.
\end{proof}

\begin{uwaga}[Rola założeń]
W konstrukcji wystarczy, aby każde $E\subseteq X$ miało jakieś
przeliczalne pokrycie zbiorami z $\mathcal C$; przyjęte założenie
zapewnia to jednocześnie wszystkim $E$. Nie wymagamy addytywności $c$.
Z drugiej strony samo pokrywanie nie gwarantuje, że
$\mu_c^*(C)=c(C)$ dla $C\in\mathcal C$.
Jeżeli natomiast $\mathcal C$ jest algebrą, a $c$ jest premiarą,
to równość $\mu_c^*(C)=c(C)$ zachodzi dla każdego
$C\in\mathcal C$, co udowodnimy dalej. Warunek ten jest
wystarczający, lecz nie jest konieczny.
\end{uwaga}

\section{Kryterium mierzalności Carathéodory'ego}

\begin{definicja}[Zbiór mierzalny w sensie Carathéodory'ego]
\label{def:caratheodory-measurable}
Niech $\mu^*$ będzie miarą zewnętrzną na $X$. Zbiór $E\subseteq X$
nazywamy $\mu^*$-\emph{mierzalnym}, jeśli dla każdego $T\subseteq X$
zachodzi
\begin{equation}\label{eq:caratheodory-criterion}
\mu^*(T)=\mu^*(T\cap E)+\mu^*(T\setminus E).
\end{equation}
Rodzinę takich zbiorów oznaczamy przez $\mathcal M_{\mu^*}$.
\end{definicja}

Z nierówności \eqref{eq:outer-subadd} zawsze wynika
$\mu^*(T)\le\mu^*(T\cap E)+\mu^*(T\setminus E)$.
Weryfikacja \eqref{eq:caratheodory-criterion} polega więc na
udowodnieniu \emph{odwrotnej} nierówności, także wtedy, gdy
$\mu^*(T)=\infty$.

\begin{twierdzenie}[Twierdzenie Carathéodory'ego]
\label{thm:caratheodory}
Jeśli $\mu^*$ jest miarą zewnętrzną na $X$, to
$\mathcal M_{\mu^*}$ jest $\sigma$-algebrą na $X$, zaś
$\mu:=\mu^*|_{\mathcal M_{\mu^*}}$ jest miarą. Ponadto miara
$\mu$ jest zupełna: każdy podzbiór zbioru $\mu$-miary zero jest
$\mu^*$-mierzalny i ma miarę zero.
\end{twierdzenie}

\begin{proof}
\emph{Krok 1: dopełnienia i zbiory trywialne.}
Dla $E=\varnothing$ równość
\eqref{eq:caratheodory-criterion} ma postać
$\mu^*(T)=0+\mu^*(T)$. Warunek jest symetryczny względem
$E$ i $X\setminus E$, toteż mierzalność $E$ pociąga mierzalność
jego dopełnienia. W szczególności $X\in\mathcal M_{\mu^*}$.

\emph{Krok 2: sumy dwóch zbiorów.}
Dla $E,F\in\mathcal M_{\mu^*}$ i dowolnego $T\subseteq X$
rozkładamy najpierw $T$ według $E$, a potem $T\setminus E$
według $F$:
\begin{align*}
\mu^*(T)
&=\mu^*(T\cap E)+\mu^*(T\setminus E)\\
&=\mu^*(T\cap E)+\mu^*((T\setminus E)\cap F)
  +\mu^*(T\setminus(E\cup F))\\
&\ge\mu^*(T\cap(E\cup F))
  +\mu^*(T\setminus(E\cup F)).
\end{align*}
Ostatnia nierówność wynika z podaddytywności, gdyż
$T\cap(E\cup F)=(T\cap E)\cup((T\setminus E)\cap F)$.
Nierówność w drugą stronę obowiązuje dla każdego zbioru, a zatem
$E\cup F$ jest mierzalny. Mamy też zamkniętość na skończone
przekroje, dzięki prawom de Morgana.

\emph{Krok 3: przeliczalne sumy rozłączne.}
Niech $(E_n)_{n\ge1}$ będzie parami rozłącznym ciągiem
zbiorów mierzalnych, a $E=\bigcup_nE_n$.
Wielokrotnie stosując kryterium do $E_1,\ldots,E_N$,
otrzymujemy dla dowolnego $T\subseteq X$
\[
\mu^*(T)=\sum_{n=1}^N\mu^*(T\cap E_n)
+\mu^*\Big(T\setminus\bigcup_{n=1}^N E_n\Big).
\]
Ponieważ $T\setminus E\subseteq T\setminus\bigcup_{n=1}^N E_n$,
monotoniczność daje
$\mu^*(T)\ge\sum_{n=1}^N\mu^*(T\cap E_n)
+\mu^*(T\setminus E)$.
Przechodząc do kresu górnego po $N$, otrzymujemy
\[
\mu^*(T)\ge\sum_{n=1}^\infty\mu^*(T\cap E_n)
+\mu^*(T\setminus E)
\ge\mu^*(T\cap E)+\mu^*(T\setminus E).
\]
Ostatnie przejście korzysta z podaddytywności dla
$T\cap E=\bigcup_n(T\cap E_n)$. To dowodzi mierzalności $E$.

\emph{Krok 4: dowolne sumy przeliczalne.}
Dla mierzalnych $A_n$ stosujemy rozłącznienie
$E_1=A_1$ oraz $E_n=A_n\setminus\bigcup_{k<n}A_k$ dla $n\ge2$.
Na mocy poprzednich kroków wszystkie $E_n$ są mierzalne i
$\bigcup_n A_n=\bigcup_n E_n$; zatem
$\mathcal M_{\mu^*}$ jest $\sigma$-algebrą.

\emph{Krok 5: przeliczalna addytywność.}
Dla rozłącznych mierzalnych $E_n$ podstawiamy $T=E=\bigcup_nE_n$
do nierówności z kroku 3. Otrzymujemy
$\mu^*(E)\ge\sum_n\mu^*(E_n)$, a odwrotną nierówność zapewnia
\eqref{eq:outer-subadd}. Ponadto $\mu^*(\varnothing)=0$,
więc ograniczenie $\mu$ jest miarą.

\emph{Krok 6: zupełność.}
Niech $N\in\mathcal M_{\mu^*}$, $\mu^*(N)=0$ i $S\subseteq N$.
Dla dowolnego $T$ z monotoniczności mamy
$\mu^*(T\cap S)=0$ i
$\mu^*(T\setminus S)\le\mu^*(T)$.
Stąd
$\mu^*(T\cap S)+\mu^*(T\setminus S)\le\mu^*(T)$.
Podaddytywność daje nierówność przeciwną. Zatem $S$ jest
mierzalny oraz $\mu^*(S)=0$.
\end{proof}

\begin{przyklad}[Miara zewnętrzna, która nie mierzy wszystkich zbiorów]
Niech $X=\{a,b\}$ oraz $\mu^*(A)=1$ dla $A\ne\varnothing$.
Gdyby $E=\{a\}$ było mierzalne, podstawienie $T=X$ dałoby
$1=\mu^*(X)=\mu^*(\{a\})+\mu^*(\{b\})=2$,
sprzeczność. Stąd $\mathcal M_{\mu^*}=\{\varnothing,X\}$.
To pokazuje, że kryterium Carathéodory'ego ma rzeczywistą treść.
\end{przyklad}

\section{Premiary na algebrach zbiorów}

\begin{definicja}[Premiara]
\label{def:premeasure}
Niech $\mathcal A$ będzie algebrą podzbiorów $X$.
Funkcję $\mu_0:\mathcal A\to[0,\infty]$ nazywamy
\emph{premiarą}, jeżeli $\mu_0(\varnothing)=0$ i dla każdego
ciągu parami rozłącznych $A_n\in\mathcal A$ takiego, że
$A=\bigcup_{n\ge1}A_n\in\mathcal A$, zachodzi
\begin{equation}\label{eq:premeasure}
\mu_0(A)=\sum_{n=1}^\infty\mu_0(A_n).
\end{equation}
Warunek nie zakłada, że $\bigcup_nA_n$ zawsze należy do
$\mathcal A$; algebra nie musi być $\sigma$-algebrą.
\end{definicja}

\begin{lemat}[Podaddytywność premiary na pokryciach]
\label{lem:premeasure-cover}
Jeżeli $\mu_0$ jest premiarą na $\mathcal A$, zaś
$A,A_n\in\mathcal A$ i $A\subseteq\bigcup_{n\ge1}A_n$,
to $\mu_0(A)\le\sum_{n\ge1}\mu_0(A_n)$.
\end{lemat}

\begin{proof}
Tworzymy rozłączne zbiory
$B_1=A\cap A_1$ i
$B_n=(A\cap A_n)\setminus\bigcup_{k<n}A_k$ dla $n\ge2$.
Wszystkie należą do algebry, a $A=\bigsqcup_nB_n$.
Z \eqref{eq:premeasure} oraz monotoniczności premiary
(wynikającej ze skończonej addytywności) dostajemy
$\mu_0(A)=\sum_n\mu_0(B_n)\le\sum_n\mu_0(A_n)$.
\end{proof}

\begin{przyklad}[Skończona addytywność nie wystarcza]
Na $X=\mathbb N$ niech $\mathcal A$ będzie algebrą
zbiorów skończonych i kofinitywnych. Określmy
$\nu_0(A)=0$, gdy $A$ jest skończony, oraz
$\nu_0(A)=1$, gdy $\mathbb N\setminus A$ jest skończony.
Jest to funkcja skończenie addytywna: dwa rozłączne
zbiory kofinitywne nie istnieją, a dołączenie zbioru
skończonego nie zmienia wartości. Nie jest premiarą, gdyż
$\mathbb N=\bigsqcup_{n\ge1}\{n\}$, natomiast
$\nu_0(\mathbb N)=1\ne\sum_n\nu_0(\{n\})=0$.
Każda miara rozszerzająca $\nu_0$ musiałaby spełniać
sprzeczną równość, toteż takie rozszerzenie nie istnieje.
\end{przyklad}

\begin{twierdzenie}[Kryterium ciągłości premiary skończonej]
\label{thm:premeasure-continuity}
Niech $\mathcal A$ będzie algebrą podzbiorów $X$, a
$\mu_0:\mathcal A\to[0,\infty)$ funkcją skończenie addytywną
spełniającą $\mu_0(X)<\infty$ i $\mu_0(\varnothing)=0$.
Wówczas $\mu_0$ jest premiarą wtedy i tylko wtedy, gdy
$A_n\in\mathcal A$, $A_n\downarrow\varnothing$ implikuje
$\mu_0(A_n)\to0$.
\end{twierdzenie}

\begin{proof}
Jeśli $\mu_0$ jest premiarą, to dla
$A_n\downarrow\varnothing$ zbiory $D_n=A_n\setminus A_{n+1}$
są rozłączne i $A_1=\bigsqcup_nD_n$.
W konsekwencji
$\mu_0(A_n)=\sum_{k\ge n}\mu_0(D_k)\to0$,
gdyż szereg ma skończoną sumę $\mu_0(A_1)\le\mu_0(X)$.

Odwrotnie, niech $A=\bigsqcup_n A_n$ i wszystkie te zbiory
należą do $\mathcal A$. Dla
$R_N=A\setminus\bigcup_{n=1}^N A_n$ mamy
$R_N\downarrow\varnothing$, a skończona addytywność daje
$\mu_0(A)=\sum_{n=1}^N\mu_0(A_n)+\mu_0(R_N)$.
Po przejściu do granicy otrzymujemy \eqref{eq:premeasure}.
\end{proof}

\begin{uwaga}
Nie wolno bez zmian stosować twierdzenia
\ref{thm:premeasure-continuity} przy $\mu_0(X)=\infty$.
Na przykład miara licząca na $\mathbb N$ jest premiarą,
lecz dla $A_n=\{n,n+1,\ldots\}$ mamy $A_n\downarrow\varnothing$
i $\mu_0(A_n)=\infty$ dla wszystkich $n$.
\end{uwaga}

\section{Twierdzenie o rozszerzaniu miar}

\noindent\emph{Powiązanie tematyczne w Bibliotece matematycznej:} \atlasref{P1-02::ch:probability}{P1-02: konstrukcje przestrzeni probabilistycznych}.


\begin{twierdzenie}[Rozszerzenie Carathéodory'ego]
\label{thm:extension-caratheodory}
Niech $\mathcal A$ będzie algebrą podzbiorów $X$, a
$\mu_0$ premiarą na $\mathcal A$. Zdefiniujmy $\mu^*$
przez \eqref{eq:outer-cover}, biorąc $\mathcal C=\mathcal A$
i $c=\mu_0$. Wówczas:
\begin{enumerate}[label=\textup{(\roman*)}]
\item $\mu^*$ jest miarą zewnętrzną;
\item każdy $A\in\mathcal A$ jest $\mu^*$-mierzalny;
\item $\mu^*(A)=\mu_0(A)$ dla wszystkich $A\in\mathcal A$;
\item $\mu=\mu^*|_{\mathcal M_{\mu^*}}$ jest zupełną miarą
rozszerzającą $\mu_0$, a jej ograniczenie do
$\sigma(\mathcal A)$ jest miarą na $\sigma(\mathcal A)$.
Miara na $\sigma(\mathcal A)$ nie musi jednak być zupełna.
\end{enumerate}
\end{twierdzenie}

\begin{proof}
Punkt (i) wynika z twierdzenia \ref{thm:outer-from-cover},
gdyż $X\in\mathcal A$.
Ustalmy $A\in\mathcal A$ i dowolny $T\subseteq X$.
Dla każdego pokrycia $T\subseteq\bigcup_nC_n$ przez
$C_n\in\mathcal A$ zbiory $C_n\cap A$ pokrywają $T\cap A$,
a $C_n\setminus A$ pokrywają $T\setminus A$. Dlatego
\begin{align*}
\mu^*(T\cap A)+\mu^*(T\setminus A)
&\le\sum_n\mu_0(C_n\cap A)+\sum_n\mu_0(C_n\setminus A)\\
&=\sum_n\mu_0(C_n).
\end{align*}
Równość korzysta ze skończonej addytywności $\mu_0$.
Biorąc infimum po pokryciach, otrzymujemy
$\mu^*(T\cap A)+\mu^*(T\setminus A)\le\mu^*(T)$.
W połączeniu z podaddytywnością jest to warunek
\eqref{eq:caratheodory-criterion}; dowodzi (ii).

Dla każdego $A\in\mathcal A$ pokrycie $A$ samym $A$
(i pozostałymi zbiorami pustymi) daje
$\mu^*(A)\le\mu_0(A)$. Z drugiej strony
lemat \ref{lem:premeasure-cover} zapewnia
$\mu_0(A)\le\sum_n\mu_0(C_n)$ dla dowolnego pokrycia
$A\subseteq\bigcup_nC_n$. Infimum daje
$\mu_0(A)\le\mu^*(A)$, a więc (iii).

Punkt (iv) jest już konsekwencją twierdzenia
\ref{thm:caratheodory}: $\mathcal M_{\mu^*}$ jest
$\sigma$-algebrą zawierającą $\mathcal A$, zatem
zawiera $\sigma(\mathcal A)$; ograniczenie $\mu^*$
jest na niej miarą, zaś na $\mathcal M_{\mu^*}$ miarą zupełną.
\end{proof}

\begin{definicja}[\texorpdfstring{$\sigma$}{sigma}-skończona premiara]\label{idx:03:definicja:09}

Premiarę $\mu_0$ na algebrze $\mathcal A$ nazywamy
$\sigma$-\emph{skończoną}, jeżeli istnieją
$A_n\in\mathcal A$ takie, że $X=\bigcup_nA_n$ oraz
$\mu_0(A_n)<\infty$ dla każdego $n$.
Możemy zawsze zastąpić $A_n$ sumami skończonymi
$A_1\cup\cdots\cup A_n$, otrzymując ciąg rosnący.
\end{definicja}

\begin{twierdzenie}[Jednoznaczność rozszerzenia]
\label{thm:extension-uniqueness}
Jeśli premiara $\mu_0$ na algebrze $\mathcal A$ jest
$\sigma$-skończona, to istnieje dokładnie jedna miara
na $\sigma(\mathcal A)$, której ograniczenie do
$\mathcal A$ jest równe $\mu_0$.
\end{twierdzenie}

\begin{proof}
Istnienie daje twierdzenie \ref{thm:extension-caratheodory}.
Niech $\mu$ i $\nu$ będą dwoma rozszerzeniami.
Wybierzmy rosnący ciąg $A_n\in\mathcal A$ o
$\bigcup_n A_n=X$ i $\mu_0(A_n)<\infty$.
Dla ustalonego $n$ funkcje
$\mu_n(E)=\mu(E\cap A_n)$ oraz $\nu_n(E)=\nu(E\cap A_n)$,
$E\in\sigma(\mathcal A)$, są miarami skończonymi.
Zgadzają się na algebrze $\mathcal A$, ponieważ
$E\cap A_n\in\mathcal A$ dla $E\in\mathcal A$.
Algebra jest układem $\pi$ i generuje
$\sigma(\mathcal A)$. Z twierdzenia
\ref{thm:uniqueness-finite} otrzymujemy
$\mu_n(E)=\nu_n(E)$ dla każdego $E\in\sigma(\mathcal A)$.
Ponieważ $E\cap A_n\uparrow E$, ciągłość miary z dołu
(twierdzenie \ref{thm:measure-continuity}) daje
$\mu(E)=\lim_n\mu_n(E)=\lim_n\nu_n(E)=\nu(E)$.
\end{proof}

\begin{przyklad}[Dlaczego potrzebna jest $\sigma$-skończoność]
\label{ex:nonunique-extension}
Na $X=\mathbb R$ niech $\mathcal A$ będzie algebrą
generowaną przez przedziały $[a,b)$, gdzie $a<b$ są
rzeczywiste. Każdy niepusty element $\mathcal A$ zawiera
niezdegenerowany przedział; ponadto
$\sigma(\mathcal A)=\mathcal B(\mathbb R)$.
Na zbiorach borelowskich rozważmy dwie miary:
\[
\mu(E)=\#E,\qquad
\nu(E)=\begin{cases}0,&E\text{ jest przeliczalny},\\
\infty,&E\text{ jest nieprzeliczalny}.
\end{cases}
\]
Funkcja $\nu$ jest miarą, gdyż przeliczalna suma zbiorów
przeliczalnych jest przeliczalna, a suma zawierająca
zbiór nieprzeliczalny jest nieprzeliczalna. Na każdym
niepustym elemencie $\mathcal A$ obie miary mają wartość
$\infty$, zaś na pustym obie są równe zero. Są to więc
różne rozszerzenia tej samej premiary:
$\mu(\{0\})=1$, natomiast $\nu(\{0\})=0$.
Premiara nie jest $\sigma$-skończona: jej jedynym
zbiorem o skończonej mierze jest $\varnothing$.
\end{przyklad}

\begin{uwaga}[Różne role założeń]
Premiara daje istnienie rozszerzenia na $\sigma(\mathcal A)$
bez założenia $\sigma$-skończoności. To założenie pojawia się
dopiero przy jednoznaczności rozszerzenia oraz przy
identyfikacji całej rodziny zbiorów Carathéodory'ego
z uzupełnieniem $\sigma(\mathcal A)$. W szczególności
stwierdzenie o istnieniu i stwierdzenie o jednoznaczności
nie są jednym i tym samym twierdzeniem.
\end{uwaga}

\section{Kiedy mierzalność Carathéodory'ego jest uzupełnieniem rozszerzenia?}

Rozszerzenie otrzymane z twierdzenia
\ref{thm:extension-caratheodory} jest określone na
$\mathcal M_{\mu^*}$, zwykle większej niż $\sigma(\mathcal A)$.
Nie należy utożsamiać tych rodzin bez dowodu.
Poniższy rezultat określa ich związek przy
$\sigma$-skończoności.

\begin{lemat}[Aproksymacja od góry]
\label{lem:outer-approx-algebra}
Niech $\mu_0$ będzie premiarą na algebrze $\mathcal A$ i
$\mu^*$ jej miarą zewnętrzną. Jeśli $E\subseteq X$ oraz
$\mu^*(E)<\infty$, to dla każdego $\varepsilon>0$
istnieje $B\in\sigma(\mathcal A)$, $E\subseteq B$, takie, że
$\mu^*(B)\le\mu^*(E)+\varepsilon$.
Jeżeli dodatkowo $E\in\mathcal M_{\mu^*}$, to
$\mu(B\setminus E)<\varepsilon$.
\end{lemat}

\begin{proof}
Wybieramy pokrycie $E\subseteq\bigcup_nA_n$ przez
$A_n\in\mathcal A$ z $\sum_n\mu_0(A_n)<\mu^*(E)+\varepsilon$.
Wystarczy wziąć $B=\bigcup_nA_n$. Z podaddytywności i
równości $\mu^*(A_n)=\mu_0(A_n)$ dostajemy
$\mu^*(B)\le\sum_n\mu_0(A_n)<\mu^*(E)+\varepsilon$.
Jeśli $E$ jest mierzalny, to $B\setminus E$ jest mierzalny
i $\mu(B)=\mu(E)+\mu(B\setminus E)$.
Obie pierwsze wartości są skończone, więc
$\mu(B\setminus E)<\varepsilon$.
\end{proof}

\begin{twierdzenie}[Uzupełnienie rozszerzenia]
\label{thm:caratheodory-completion}
Załóżmy, że premiara $\mu_0$ na algebrze $\mathcal A$
jest $\sigma$-skończona. Niech $\mu^*$ będzie wyznaczoną
przez nią miarą zewnętrzną, a $\mu_B=\mu^*|_{\sigma(\mathcal A)}$.
Wówczas przestrzeń
$(X,\mathcal M_{\mu^*},\mu^*|_{\mathcal M_{\mu^*}})$
jest uzupełnieniem przestrzeni
$(X,\sigma(\mathcal A),\mu_B)$ w sensie twierdzenia
\ref{thm:completion}. Dokładniej: dla każdego
$E\in\mathcal M_{\mu^*}$ istnieją
$C,D\in\sigma(\mathcal A)$ takie, że
$C\subseteq E\subseteq D$ i $\mu_B(D\setminus C)=0$.
\end{twierdzenie}

\begin{proof}
Wybierzmy rosnący ciąg $A_n\in\mathcal A$ pokrywający $X$,
przy czym $\mu_0(A_n)<\infty$. Dla ustalonego mierzalnego
$E$ zbiór $E_n=E\cap A_n$ jest mierzalny i ma skończoną miarę.
Dla każdego $k,n\ge1$ lemat \ref{lem:outer-approx-algebra}
daje $B_{n,k}\in\sigma(\mathcal A)$ z
$E_n\subseteq B_{n,k}$ oraz
$\mu(B_{n,k}\setminus E_n)<2^{-n}/k$.
Połóżmy $D_k=\bigcup_{n\ge1}B_{n,k}$.
Wtedy $D_k\in\sigma(\mathcal A)$,
$E\subseteq D_k$, gdyż $E=\bigcup_nE_n$, i
\[
\mu(D_k\setminus E)
\le\sum_n\mu(B_{n,k}\setminus E)
\le\sum_n\mu(B_{n,k}\setminus E_n)<\frac1k.
\]
Zbiór $D=\bigcap_{k\ge1}D_k$ należy do
$\sigma(\mathcal A)$, zawiera $E$ i spełnia
$\mu(D\setminus E)=0$, ponieważ
$D\setminus E\subseteq D_k\setminus E$ dla każdego $k$.
Powtarzamy konstrukcję dla mierzalnego zbioru
$X\setminus E$, otrzymując zbiór należący do
$\sigma(\mathcal A)$, mianowicie $H\supseteq X\setminus E$,
taki że $\mu(H\setminus(X\setminus E))=0$.
Wtedy $C=X\setminus H\in\sigma(\mathcal A)$,
$C\subseteq E$ i $\mu(E\setminus C)=0$.
Wobec $C\subseteq E\subseteq D$ mamy
$D\setminus C=(D\setminus E)\sqcup(E\setminus C)$,
zatem $\mu_B(D\setminus C)=0$.

Każdy zbiór z uzupełnienia $\sigma(\mathcal A)$ różni
się od zbioru z $\sigma(\mathcal A)$ o podzbiór
zbioru $\mu_B$-zerowego, a więc jest
$\mu^*$-mierzalny z zupełności twierdzenia
\ref{thm:caratheodory}. Przed chwilą udowodniliśmy
zawieranie przeciwne, kończąc dowód.
\end{proof}

\begin{uwaga}[Dlaczego to rozróżnienie jest ważne?]
W rozdziale \ref{ch:lebesgue} rozróżnimy zbiory borelowskie
i zbiory mierzalne w sensie Lebesgue'a. Konstrukcja
Carathéodory'ego od razu daje miarę zupełną; ograniczenie
do borelowskiej $\sigma$-algebry nie musi być zupełne.
\end{uwaga}

\section{Jeszcze jeden sposób interpretacji konstrukcji}

Dla zbioru $E\subseteq X$ liczba $\mu^*(E)$ jest
najmniejszym kosztem \emph{przeliczalnego pokrycia} $E$
przez zbiory z wyjściowej algebry. Mierzalność w sensie
Carathéodory'ego oznacza, że przecięcie z $E$ oraz z jego
dopełnieniem nie powoduje sztucznej utraty kosztu,
niezależnie od badanego zbioru $T$.

Następujące trzy kroki to logiczny szkielet rozdziału:
\begin{enumerate}
\item z premiary $\mu_0$ tworzymy $\mu^*$ na całej $\mathcal P(X)$;
\item warunek \eqref{eq:caratheodory-criterion} wyodrębnia
$\sigma$-algebrę, na której $\mu^*$ jest miarą;
\item przy $\sigma$-skończoności dostajemy jedyne
rozszerzenie na $\sigma(\mathcal A)$, a po uzupełnieniu
pełną dziedzinę mierzalności Carathéodory'ego.
\end{enumerate}
W następnym rozdziale zrealizujemy ten program dla
objętości prostopadłościanów w $\mathbb R^d$.

% ROZDZIAŁ 4 — plik włączany przez main.tex
\chapter{Miara Lebesgue'a na przestrzeni euklidesowej}
\label{ch:lebesgue}

\providecommand{\atlasref}[2]{#2}


\section{Co właściwie chcemy skonstruować?}

Długość przedziału $[a,b)$ to $b-a$, a objętość
$d$-wymiarowego prostopadłościanu
$Q=\prod_{j=1}^d[a_j,b_j)$ wynosi
$|Q|=\prod_{j=1}^d(b_j-a_j)$.
Chcemy rozciągnąć tę funkcję na bogatą rodzinę
podzbiorów $\mathbb R^d$, zachowując przeliczalną
addytywność i naturalną zgodność z geometrią.

Przy każdym $d\in\mathbb N$, $d\ge1$, przez
$\mathcal B(\mathbb R^d)$ oznaczamy $\sigma$-algebrę
borelowską z definicji \ref{def:borel}. Odróżniamy
\emph{miarę borelowską Lebesgue'a} $\lambda_d^B$
od jej uzupełnienia $\lambda_d$. Ich dziedziny nie są
jednakowe. Dla prostopadłościanów osiowych z dowolnym wyborem
otwartości i domknięcia ścian przez $|R|$ oznaczamy
\emph{elementarną objętość geometryczną}, czyli iloczyn
długości boków; nie zakładamy jeszcze, że jest ona miarą.
W szczególności dla $Q$ jest to wyjściowa objętość,
a dla zwartego $K$ lub otwartego $O$ ta sama liczba.
Symbol $|R|$ nie jest na tym etapie miarą dowolnego zbioru.

\section{Objętość elementarna i lemat o pokryciach}

\begin{definicja}[Prostopadłościan półotwarty]
\label{def:halfopen-box}
Przez prostopadłościan półotwarty rozumiemy zbiór
$Q=\prod_{j=1}^d[a_j,b_j)$, gdzie $a_j,b_j\in\mathbb R$,
$a_j<b_j$. Rodzinę takich \emph{ograniczonych} zbiorów
oraz zbioru pustego oznaczamy $\mathcal Q_d$.
Kładziemy $|\varnothing|=0$ i
$|Q|=\prod_{j=1}^d(b_j-a_j)$ dla niepustego $Q$.
\end{definicja}

Użycie przedziałów półotwartych nie jest kosmetyczne.
Dla $a<c<b$ mamy rozkład rozłączny
$[a,b)=[a,c)\sqcup[c,b)$. Analogicznie przecięcie
dwóch prostopadłościanów półotwartych jest
prostopadłościanem tego samego typu lub zbiorem pustym,
a różnica takich zbiorów jest skończoną sumą rozłączną
prostopadłościanów półotwartych.

\begin{lemat}[Skończony podział na siatkę]
\label{lem:rect-grid}
Niech $Q_1,\ldots,Q_N\in\mathcal Q_d$.
Istnieje skończona rodzina parami rozłącznych
prostopadłościanów półotwartych $P_1,\ldots,P_M$, taka że
każdy z $Q_i$ jest sumą pewnych $P_j$. W konsekwencji
można jednoznacznie określić na rodzinie $\mathcal R_d$
wszystkich skończonych sum zbiorów z $\mathcal Q_d$
funkcję $v$, przyjmując
$v(\bigsqcup_{i=1}^NQ_i)=\sum_{i=1}^N|Q_i|$
dla dowolnego rozłącznego przedstawienia.
Funkcja $v$ jest nieujemna i skończenie addytywna.
\end{lemat}

\begin{proof}
W każdej współrzędnej wypisujemy w porządku rosnącym
wszystkie końce przedziałów występujących w $Q_i$.
Sąsiednie końce wyznaczają przedziały półotwarte;
ich iloczyny kartezjańskie dają skończoną siatkę.
Każdy $Q_i$ składa się z całych komórek siatki.
Dla każdej komórki $P=\prod_j[c_j,e_j)$ jej objętość
to $\prod_j(e_j-c_j)$. Dodawanie objętości przy
podziale w jednej współrzędnej jest rozdzielnością
mnożenia względem dodawania; iterując po współrzędnych,
otrzymujemy $|Q_i|$ jako sumę objętości jej komórek.
Jeżeli ten sam zbiór ma dwa rozłączne przedstawienia,
wspólna siatka dla obu przedstawień pokazuje, że oba
sumowania dają tę samą liczbę. Tak samo dowodzi się
skończonej addytywności $v$.
\end{proof}

\begin{uwaga}
$\mathcal R_d$ jest pierścieniem zbiorów, ale nie
algebrą podzbiorów $\mathbb R^d$, ponieważ
$\mathbb R^d\notin\mathcal R_d$. Na ograniczonym
prostopadłościanie $H$ rodzina
$\mathcal R_d|_H=\{A\cap H:A\in\mathcal R_d\}$
jest natomiast algebrą podzbiorów $H$.
To lokalne rozróżnienie pozwoli nam zastosować
rozdział \ref{ch:caratheodory} bez udawania, że
nieograniczona przestrzeń jest ograniczonym prostokątem.
\end{uwaga}

\begin{lemat}[Nierówność objętościowa dla pokryć]
\label{lem:rect-cover-ineq}
Jeżeli $Q,Q_1,Q_2,\ldots\in\mathcal Q_d$ i
$Q\subseteq\bigcup_{n\ge1}Q_n$, to
\begin{equation}\label{eq:rect-cover-volume}
|Q|\le\sum_{n=1}^\infty|Q_n|.
\end{equation}
\end{lemat}

\begin{proof}
Dla $Q=\varnothing$ twierdzenie jest oczywiste.
Niech $Q=\prod_j[a_j,b_j)$ będzie niepusty.
Wystarczy rozważyć przypadek, w którym suma po prawej
stronie jest skończona. Ustalmy $\varepsilon>0$ oraz
$0<\delta<\frac12\min_j(b_j-a_j)$ i połóżmy
$K_\delta=\prod_j[a_j+\delta,b_j-\delta]$.
Jest to zwarty podzbiór $Q$.

Każdy $Q_n$ można zawrzeć w otwartym prostopadłościanie
$O_n$ o objętości elementarnej mniejszej niż
$|Q_n|+\varepsilon2^{-n}$: wystarczy nieznacznie
przesunąć wszystkie jego ściany na zewnątrz.
Zbiory $O_n$ tworzą otwarte pokrycie $K_\delta$.
Ze zwartości istnieją $n_1,\ldots,n_N$ takie, że
$K_\delta\subseteq\bigcup_{i=1}^N O_{n_i}$.

Wykażmy osobno elementarną nierówność dla skończonego
pokrycia zwartego prostopadłościanu
$K=\prod_{j=1}^d[\alpha_j,\beta_j]$, gdzie
$\alpha_j\le\beta_j$, przez otwarte prostopadłościany
$O_1,\ldots,O_m$:
$|K|\le\sum_{i=1}^m|O_i|$.
Jeśli któryś bok $K$ ma długość zero, nierówność jest
natychmiastowa. Dalej zakładamy $\alpha_j<\beta_j$
dla wszystkich $j$. Dla dostatecznie małego $\eta>0$ wprowadzamy
$P_\eta=\prod_j[\alpha_j+\eta,\beta_j-\eta)$.
W każdej współrzędnej wybieramy wszystkie końce
przedziałów tworzących $O_i$ i $P_\eta$.
Skończona siatka złożona z tych punktów rozkłada
$P_\eta$ na parami rozłączne półotwarte komórki $S_l$.
Wnętrze komórki nie przecina żadnej hiperpłaszczyzny,
która jest ścianą $O_i$. Wybierzmy
$x_l\in\operatorname{int}S_l\subseteq K$.
Punkt $x_l$ należy do pewnego $O_{i(l)}$.
Ponieważ granice współrzędnych $O_{i(l)}$ należą do
siatki, wynika stąd
$\operatorname{int}S_l\subseteq O_{i(l)}$, a więc
$S_l\subseteq\overline{O_{i(l)}}$.

Dla ustalonego $i$ oznaczmy przez $\mathcal S_i$ rodzinę
komórek przypisanych do $O_i$. Niech $H_i^{(\gamma)}$
będzie półotwartym prostopadłościanem zawierającym
$\overline O_i$, którego długości boków przewyższają
odpowiednie długości $O_i$ o $2\gamma$.
Każda komórka $S_l\in\mathcal S_i$ jest zawarta w
$H_i^{(\gamma)}$. Na mocy skończonej addytywności
i monotoniczności elementarnej objętości na pierścieniu
$\mathcal R_d$ mamy
$\sum_{S_l\in\mathcal S_i}|S_l|
\le |H_i^{(\gamma)}|$.
Gdy $\gamma\downarrow0$, prawa strona dąży do $|O_i|$.
Komórki $S_l$ rozkładają $P_\eta$ i każdą przypisaliśmy
jednemu z $O_i$, zatem
\[
 |P_\eta|=\sum_l|S_l|
 =\sum_{i=1}^m\sum_{S_l\in\mathcal S_i}|S_l|
 \le\sum_{i=1}^m|O_i|.
\]
Wreszcie $|P_\eta|\to|K|$ przy $\eta\downarrow0$.
Argument ten wykorzystuje wyłącznie skończone podziały
prostopadłościanów i nie zakłada wcześniej istnienia miary
Lebesgue'a.
Stąd
\[
|K_\delta|\le\sum_{i=1}^N|O_{n_i}|
\le\sum_{n=1}^\infty|Q_n|+\varepsilon.
\]
Przy $\delta\downarrow0$ mamy
$|K_\delta|=\prod_j(b_j-a_j-2\delta)\to|Q|$.
Następnie $\varepsilon\downarrow0$, co dowodzi
\eqref{eq:rect-cover-volume}.
\end{proof}

\begin{uwaga}[W którym miejscu używamy zwartości?]
Bez zwartości nie mielibyśmy prawa zastąpić
przeliczalnego pokrycia skończonym i skorzystać
z rozkładu na siatkę. To zasadniczy krok geometryczny
w dowodzie, że elementarna objętość przechodzi
z addytywności skończonej do przeliczalnej.
\end{uwaga}

\begin{twierdzenie}[Premiara na ograniczonym prostopadłościanie]
\label{thm:volume-premeasure}
Dla ustalonego $H=\prod_{j=1}^d[a_j,b_j)\in\mathcal Q_d$
funkcja $v$ z lematu \ref{lem:rect-grid}, ograniczona do
algebry $\mathcal R_d|_H$, jest skończoną premiarą.
Jej przedłużenie na $\sigma(\mathcal R_d|_H)$ istnieje
i jest jednoznaczne.
\end{twierdzenie}

\begin{proof}
Wszystkie zbiory z $\mathcal R_d|_H$ są skończonymi
rozłącznymi sumami prostopadłościanów półotwartych
zawartych w $H$; wspólna siatka zapewnia skończoną
addytywność. W szczególności $v(H)=|H|<\infty$.

Niech $A=\bigsqcup_{n\ge1}A_n$, gdzie
$A,A_n\in\mathcal R_d|_H$. Z addytywności skończonej
i monotoniczności dostajemy
$\sum_{n=1}^Nv(A_n)\le v(A)$ dla każdego $N$, więc
$\sum_nv(A_n)\le v(A)$.
Rozbijmy $A$ i każdy $A_n$ na skończenie wiele
rozłącznych prostopadłościanów. Każdy składnik rozkładu
$A$ jest pokryty składnikami wszystkich $A_n$.
Lemat \ref{lem:rect-cover-ineq} daje po zsumowaniu
$v(A)\le\sum_nv(A_n)$. Zachodzi więc addytywność
przeliczalna w sensie definicji \ref{def:premeasure}.
Istnienie przedłużenia wynika z twierdzenia
\ref{thm:extension-caratheodory}, a jednoznaczność
z twierdzenia \ref{thm:extension-uniqueness}.
\end{proof}

\section{Konstrukcja miary Lebesgue'a}

\begin{definicja}[Miara zewnętrzna Lebesgue'a]
\label{def:lebesgue-outer}
Dla $E\subseteq\mathbb R^d$ określamy
\begin{equation}\label{eq:lebesgue-outer}
\lambda_d^*(E)=\inf\Big\{\sum_{n=1}^\infty|Q_n|:
Q_n\in\mathcal Q_d,\quad E\subseteq\bigcup_{n=1}^\infty Q_n\Big\}.
\end{equation}
Rodzinę zbiorów $\lambda_d^*$-mierzalnych w sensie
Carathéodory'ego oznaczamy $\mathcal L_d$.
Miarę $\lambda_d=\lambda_d^*|_{\mathcal L_d}$
nazywamy \emph{$d$-wymiarową miarą Lebesgue'a}.
\end{definicja}

\begin{twierdzenie}[Istnienie i podstawowe własności]
\label{thm:lebesgue-exists}
Dla każdego $d\ge1$ zachodzą następujące fakty:
\begin{enumerate}[label=\textup{(\roman*)}]
\item $\lambda_d^*$ jest miarą zewnętrzną;
\item każdy $Q\in\mathcal Q_d$ należy do $\mathcal L_d$ i
$\lambda_d(Q)=|Q|$;
\item $\mathcal B(\mathbb R^d)\subseteq\mathcal L_d$;
\item $\lambda_d$ jest zupełną i $\sigma$-skończoną miarą;
\item jej ograniczenie
$\lambda_d^B=\lambda_d|_{\mathcal B(\mathbb R^d)}$
jest jedyną miarą borelowską, która przypisuje każdemu
$Q\in\mathcal Q_d$ liczbę $|Q|$.
\end{enumerate}
\end{twierdzenie}

\begin{proof}
\emph{(i)} Rodzina $\mathcal Q_d$ zawiera $\varnothing$,
a $\mathbb R^d$ ma przeliczalne pokrycie prostopadłościanami
$[-n,n)^d$. Twierdzenie \ref{thm:outer-from-cover}
gwarantuje zatem, że \eqref{eq:lebesgue-outer}
określa miarę zewnętrzną.

\emph{(ii), równość objętości.}
Pokrycie $Q$ jednym $Q$ daje $\lambda_d^*(Q)\le|Q|$.
Lemat \ref{lem:rect-cover-ineq} daje nierówność
przeciwną dla każdego pokrycia dopuszczalnego w
\eqref{eq:lebesgue-outer}; zatem $\lambda_d^*(Q)=|Q|$.

\emph{(ii), mierzalność.}
Ustalmy $Q\in\mathcal Q_d$ i $T\subseteq\mathbb R^d$.
Dla każdego pokrycia $T\subseteq\bigcup_nR_n$, gdzie
$R_n\in\mathcal Q_d$, dzielimy $R_n$ wzdłuż ścian $Q$.
Przecięcie $R_n\cap Q$ jest jednym prostopadłościanem
półotwartym lub zbiorem pustym, natomiast
$R_n\setminus Q$ jest skończoną rozłączną sumą
takich prostopadłościanów. Na mocy lematu
\ref{lem:rect-grid} łączna objętość obu części wynosi
$|R_n|$. Części pierwszego rodzaju pokrywają
$T\cap Q$, a części drugiego $T\setminus Q$, więc
\[
\lambda_d^*(T\cap Q)+\lambda_d^*(T\setminus Q)
\le\sum_{n=1}^\infty |R_n|.
\]
Biorąc infimum po pokryciach, uzyskujemy nierówność
$\le\lambda_d^*(T)$; przeciwna jest zawsze prawdziwa.
Wynika stąd $Q\in\mathcal L_d$.

\emph{(iii)} Twierdzenie \ref{thm:caratheodory} daje,
że $\mathcal L_d$ jest $\sigma$-algebrą. Ponieważ zawiera
wszystkie prostopadłościany półotwarte, zawiera również
ich $\sigma$-algebrę, czyli $\mathcal B(\mathbb R^d)$
(twierdzenie \ref{thm:borel-Rd}).

\emph{(iv)} Zupełność daje twierdzenie
\ref{thm:caratheodory}. Dla $H_n=[-n,n)^d$ mamy
$\mathbb R^d=\bigcup_nH_n$ i
$\lambda_d(H_n)=(2n)^d<\infty$, więc miara jest
$\sigma$-skończona.

\emph{(v)} Niech $\nu$ będzie inną miarą borelowską
o zadanych wartościach na $\mathcal Q_d$.
Rodzina $\mathcal Q_d$ wraz z $\varnothing$ jest
układem $\pi$, generuje $\mathcal B(\mathbb R^d)$,
a rosnące prostopadłościany $H_n$ pokrywają całą
przestrzeń i mają tę samą skończoną miarę dla
$\nu$ oraz $\lambda_d^B$. Twierdzenie
\ref{thm:uniqueness-sigma-finite} daje
$\nu=\lambda_d^B$.
\end{proof}

\begin{uwaga}[Zgodność obu konstrukcji]
Dla każdego ograniczonego prostopadłościanu $H$
miara z twierdzenia \ref{thm:volume-premeasure} jest
równa ograniczeniu $\lambda_d^B$ do
$\mathcal B(H)$. Obie są skończonymi miarami na
tej samej $\sigma$-algebrze i zgadzają się na
półotwartych prostopadłościanach zawartych w $H$.
W ten sposób globalna konstrukcja przez pokrycia
zgadza się z lokalnym rozszerzaniem premiary.
\end{uwaga}

\section{Niezmienniczość względem przesunięć i skalowania}

\begin{lemat}[Zbiory zawarte w hiperpłaszczyźnie]
\label{lem:hyperplane-null}
Dla każdego $c\in\mathbb R$ i $1\le j\le d$
hiperpłaszczyzna $\{x\in\mathbb R^d:x_j=c\}$ ma
miarę Lebesgue'a zero. W szczególności każda skończona
lub przeliczalna rodzina punktów ma miarę zero.
\end{lemat}

\begin{proof}
Przecięcie hiperpłaszczyzny z $[-N,N)^d$ zawiera się
w półotwartym prostopadłościanie, którego $j$-ty bok
ma długość $2\delta$, a pozostałe długość $2N$
(po ewentualnym nieznacznym powiększeniu boków).
Z definicji miary zewnętrznej
$\lambda_d^*(\{x:x_j=c\}\cap[-N,N)^d)
\le (2N)^{d-1}2\delta$ dla dowolnego $\delta>0$.
Przechodząc z $\delta\downarrow0$, dostajemy wartość zero.
Hiperpłaszczyzna jest przeliczalną sumą takich przecięć,
a więc ma miarę zero. Singleton jest podzbiorem
hiperpłaszczyzny, a przeliczalna suma zbiorów zerowych
ma miarę zero na mocy \eqref{eq:subadditivity}.
\end{proof}

\begin{twierdzenie}[Translacje i jednolite skalowanie]
\label{thm:invariance-lebesgue}
Dla każdego $E\in\mathcal L_d$, każdego
$a\in\mathbb R^d$ i każdego $t\in\mathbb R\setminus\{0\}$
zbiory $a+E$ i $tE$ należą do $\mathcal L_d$ oraz
\begin{equation}\label{eq:lebesgue-invariance}
\lambda_d(a+E)=\lambda_d(E),\qquad
\lambda_d(tE)=|t|^d\lambda_d(E).
\end{equation}
Przyjmujemy tu $|t|^d\cdot\infty=\infty$.
\end{twierdzenie}

\begin{proof}
Udowodnimy najpierw mocniejszy fakt dotyczący miary
\emph{zewnętrznej}: dla każdego $E\subseteq\mathbb R^d$
zachodzi
$\lambda_d^*(a+E)=\lambda_d^*(E)$ oraz
$\lambda_d^*(tE)=|t|^d\lambda_d^*(E)$.

Dla translacji argument jest bezpośredni: jeśli
$E\subseteq\bigcup_nQ_n$, to
$a+E\subseteq\bigcup_n(a+Q_n)$, a
$|a+Q_n|=|Q_n|$. Z definicji infimum otrzymujemy
$\lambda_d^*(a+E)\le\lambda_d^*(E)$.
Stosując tę samą nierówność do przesunięcia o $-a$,
dostajemy równość.

Jeżeli $t>0$, to dla każdego $Q\in\mathcal Q_d$
zbiór $tQ$ jest prostopadłościanem półotwartym
objętości $t^d|Q|$. Przeskalowanie dowolnego
pokrycia $E$ daje
$\lambda_d^*(tE)\le t^d\lambda_d^*(E)$;
powtarzając argument ze skalą $1/t$, otrzymujemy
równość. Jeśli $t<0$, obraz $tQ$ jest iloczynem
przedziałów typu $(tb_j,ta_j]$. Dla każdego
$\eta>0$ można go pokryć prostopadłościanem
półotwartym $R\in\mathcal Q_d$ o
$|R|<|t|^d|Q|+\eta$. Przeskalowanie pokrycia
$E\subseteq\bigcup_nQ_n$ i wybór błędów
$\eta_n=\varepsilon2^{-n}$ daje
$\lambda_d^*(tE)\le |t|^d\sum_n|Q_n|+\varepsilon$.
Po przejściu do infimum i $\varepsilon\downarrow0$
uzyskujemy nierówność $\le$; zastosowanie tego
argumentu do $1/t$ daje nierówność przeciwną.

Wreszcie niech $E\in\mathcal L_d$ i niech
$S(x)=a+tx$, gdzie $t\ne0$. Dla dowolnego
$T\subseteq\mathbb R^d$ możemy podstawić
$S^{-1}(T)$ w kryterium Carathéodory'ego
\eqref{eq:caratheodory-criterion}. Ponieważ $S$ jest
bijekcją oraz
$S(A\cap B)=S(A)\cap S(B)$,
$S(A\setminus B)=S(A)\setminus S(B)$,
a miara zewnętrzna skaluje się przez dodatnią
stałą $|t|^d$, otrzymujemy
\[
\lambda_d^*(T)=\lambda_d^*(T\cap S(E))
+\lambda_d^*(T\setminus S(E)).
\]
Zatem $S(E)\in\mathcal L_d$ i
$\lambda_d(S(E))=|t|^d\lambda_d(E)$.
Dla samej translacji postępujemy identycznie,
używając zachowania miary zewnętrznej bez czynnika.
\end{proof}

\begin{uwaga}[Ogólne przekształcenia liniowe]
Dla odwracalnego przekształcenia liniowego
$T:\mathbb R^d\to\mathbb R^d$ prawdziwy jest
silniejszy wzór $\lambda_d(TE)=|\det T|\lambda_d(E)$.
Jego pełny dowód, obejmujący rozkład przekształcenia
na operacje elementarne lub twierdzenie o zamianie
zmiennych, umieścimy w rozdziale o przekształcaniu całek.
Nie wykorzystujemy tego ogólnego wzoru w żadnym
dowodzie bieżącego rozdziału.
\end{uwaga}

\section{Regularność miary Lebesgue'a}

Przypomnijmy oznaczenia opisujące złożoność topologiczną:
zbiór typu $G_\delta$ jest przeliczalnym przekrojem zbiorów
otwartych, zaś zbiór typu $F_\sigma$ jest przeliczalną
sumą zbiorów domkniętych. Oba rodzaje zbiorów są borelowskie.
Wyniki poniżej wskazują, jak blisko nich leżą dowolne
zbiory mierzalne Lebesgue'a.

\begin{twierdzenie}[Aproksymacja otwartymi zbiorami]
\label{thm:lebesgue-outer-regular}
Dla każdego $E\subseteq\mathbb R^d$ zachodzi
\begin{equation}\label{eq:outer-open-regular}
\lambda_d^*(E)=\inf\{\lambda_d(U):
U\subseteq\mathbb R^d\ \text{otwarty},\ E\subseteq U\}.
\end{equation}
Jeśli ponadto $E\in\mathcal L_d$ i
$\lambda_d(E)<\infty$, to dla każdego
$\varepsilon>0$ można wybrać otwarty $U\supseteq E$
taki, że $\lambda_d(U\setminus E)<\varepsilon$.
\end{twierdzenie}

\begin{proof}
Każdy zbiór otwarty jest borelowski, a więc
mierzalny. Z monotoniczności
$\lambda_d^*(E)\le\lambda_d(U)$, skąd nierówność
$\le$ w \eqref{eq:outer-open-regular}.

Przypuśćmy najpierw, że $\lambda_d^*(E)<\infty$.
Dla $\varepsilon>0$ wybieramy pokrycie
$E\subseteq\bigcup_nQ_n$ przez półotwarte
prostopadłościany tak, że
$\sum_n|Q_n|<\lambda_d^*(E)+\varepsilon/2$.
Dla każdego $n$ wybieramy otwarty prostopadłościan
$O_n\supseteq Q_n$ taki, że
$|O_n|<|Q_n|+\varepsilon2^{-n-1}$.
Zbiór $U=\bigcup_nO_n$ jest otwarty, zawiera $E$ i
\[
\lambda_d(U)\le\sum_n\lambda_d(O_n)
=\sum_n|O_n|<\lambda_d^*(E)+\varepsilon.
\]
Objętość otwartego prostopadłościanu jest równa
iloczynowi długości boków: różni się on od
półotwartego tylko skończoną sumą fragmentów ścian,
które są zerowe na mocy lematu \ref{lem:hyperplane-null}.
Dla $\lambda_d^*(E)=\infty$ wszystkie otwarte
nadzbiory mają miarę nieskończoną, więc równość
\eqref{eq:outer-open-regular} również zachodzi.

Jeśli $E$ jest mierzalny i ma miarę skończoną,
$U\setminus E$ jest mierzalny, zaś
$\lambda_d(U)=\lambda_d(E)+\lambda_d(U\setminus E)$.
Wybór $U$ z poprzedniej części daje tezę.
\end{proof}

\begin{twierdzenie}[Aproksymacja zbiorami zwartymi]
\label{thm:lebesgue-inner-regular}
Dla każdego $E\in\mathcal L_d$ zachodzi
\begin{equation}\label{eq:inner-compact-regular}
\lambda_d(E)=\sup\{\lambda_d(K):K\subseteq E,
\ K\subseteq\mathbb R^d\ \text{zwarty}\}.
\end{equation}
Jeżeli $\lambda_d(E)<\infty$, to dla dowolnego
$\varepsilon>0$ istnieje zwarty $K\subseteq E$
spełniający $\lambda_d(E\setminus K)<\varepsilon$.
\end{twierdzenie}

\begin{proof}
Najpierw załóżmy $\lambda_d(E)<\infty$.
Z ciągłości miary z dołu dla
$E\cap[-N,N]^d\uparrow E$ wybieramy $N$ tak duże, że
$\lambda_d(E\setminus H)<\varepsilon/2$, gdzie
$H=[-N,N]^d$ jest zwarty.
Niech $F=H\setminus E$; jest to zbiór mierzalny
i $\lambda_d(F)\le\lambda_d(H)<\infty$.
Z twierdzenia \ref{thm:lebesgue-outer-regular}
wybieramy zbiór otwarty $U\supseteq F$ o
$\lambda_d(U\setminus F)<\varepsilon/2$.
Zbiór $K=H\setminus U$ jest zwarty i zawiera się
w $E$. Ponadto
$E\setminus K\subseteq(E\setminus H)
\cup((E\cap H)\cap U)$, a drugi składnik zawiera się
w $U\setminus F$. Z podaddytywności
$\lambda_d(E\setminus K)<\varepsilon$.

Jeśli $\lambda_d(E)=\infty$, to z
$E\cap[-N,N]^d\uparrow E$ wynika
$\sup_N\lambda_d(E\cap[-N,N]^d)=\infty$.
Każdy z tych przecięć ma miarę skończoną, więc
poprzedni argument daje zawarty w nim zwarty podzbiór
o mierze dowolnie bliskiej mierze przecięcia. Miary
zwartych podzbiorów $E$ nie są zatem ograniczone.
Z drugiej strony $\lambda_d(K)\le\lambda_d(E)$ dla
każdego $K\subseteq E$, co kończy dowód.
\end{proof}

\begin{twierdzenie}[Borelowskie aproksymacje i uzupełnienie]
\label{thm:lebesgue-completion}
Dla każdego $E\in\mathcal L_d$ istnieją
zbiór $G$ typu $G_\delta$ i zbiór $F$ typu $F_\sigma$
takie, że
\begin{equation}\label{eq:lebesgue-Borel-approx}
F\subseteq E\subseteq G,\qquad
\lambda_d(G\setminus F)=0.
\end{equation}
W konsekwencji $\mathcal L_d$ jest dokładnie
uzupełnieniem $\mathcal B(\mathbb R^d)$ względem
$\lambda_d^B$.
\end{twierdzenie}

\begin{proof}
Niech $H_n=[-n,n)^d$ i $E_n=E\cap H_n$.
Każdy $E_n$ ma miarę skończoną.
Dla $n,k\ge1$ z twierdzenia
\ref{thm:lebesgue-outer-regular} wybieramy
otwarty $U_{n,k}\supseteq E_n$ taki, że
$\lambda_d(U_{n,k}\setminus E_n)<2^{-n}/k$.
Kładziemy $U_k=\bigcup_{n\ge1}U_{n,k}$.
Zbiory $U_k$ są otwarte, zawierają $E$ i
\[
\lambda_d(U_k\setminus E)
\le\sum_{n\ge1}\lambda_d(U_{n,k}\setminus E)
\le\sum_{n\ge1}\lambda_d(U_{n,k}\setminus E_n)
<\frac1k.
\]
Zbiór $G=\bigcap_{k\ge1}U_k$ jest typu $G_\delta$,
zawiera $E$ i spełnia $\lambda_d(G\setminus E)=0$,
ponieważ $G\setminus E\subseteq U_k\setminus E$
dla każdego $k$.
Powtarzamy konstrukcję dla $\mathbb R^d\setminus E$,
otrzymując $G'\supseteq\mathbb R^d\setminus E$ typu
$G_\delta$ oraz $\lambda_d(G'\cap E)=0$.
Wówczas $F=\mathbb R^d\setminus G'$ jest typu
$F_\sigma$, $F\subseteq E$ i
$\lambda_d(E\setminus F)=0$. Zatem
$\lambda_d(G\setminus F)=0$.

Równość dziedziny $\mathcal L_d$ z uzupełnieniem
wynika z \eqref{eq:lebesgue-Borel-approx} i
zupełności $\lambda_d$: dowolny $E\in\mathcal L_d$
różni się od borelowskiego $F$ o podzbiór borelowskiego
zbioru zerowego $G\setminus F$. Odwrotnie, każdy
podzbiór borelowskiego zbioru zerowego jest
mierzalny w sensie Carathéodory'ego.
\end{proof}

\begin{wniosek}[Geometryczna charakterystyka mierzalności]
\label{cor:lebesgue-regularity-characterization}
Dla $E\subseteq\mathbb R^d$ następujące warunki
są równoważne:
\begin{enumerate}[label=\textup{(\roman*)}]
\item $E\in\mathcal L_d$;
\item dla każdego $\varepsilon>0$ istnieje zbiór
otwarty $U\supseteq E$, dla którego
$\lambda_d^*(U\setminus E)<\varepsilon$;
\item istnieje zbiór borelowski $B\supseteq E$
taki, że $\lambda_d^*(B\setminus E)=0$.
\end{enumerate}
\end{wniosek}

\begin{proof}
Implikacja (i)$\Rightarrow$(ii) dla zbiorów o miarze
nieskończonej wynika z zastosowania konstrukcji
$U_k$ z dowodu twierdzenia
\ref{thm:lebesgue-completion}; dla wszystkich $E$
zachodzi $\lambda_d(U_k\setminus E)<1/k$.
Implikacja (ii)$\Rightarrow$(iii) wymaga ostrożności:
wybieramy otwarte $U_k\supseteq E$ z
$\lambda_d^*(U_k\setminus E)<1/k$ i kładziemy
$B=\bigcap_{k\ge1}U_k$. Wtedy $B$ jest borelowski,
$E\subseteq B$ oraz
$\lambda_d^*(B\setminus E)\le1/k$ dla każdego $k$,
czyli wartość wynosi zero.
Dla (iii)$\Rightarrow$(i) zbiór $N=B\setminus E$
ma miarę zewnętrzną zero, zatem jest mierzalny
z twierdzenia \ref{thm:caratheodory}; wobec
$E=B\setminus N$ i mierzalności $B$ dostajemy tezę.
\end{proof}

\section{Pierwsze konsekwencje i przykłady}

\begin{przyklad}[Zbiory gęste miary zero]
Zbiór $\mathbb Q^d$ jest przeliczalny, a więc
$\lambda_d(\mathbb Q^d)=0$ mimo jego gęstości
w $\mathbb R^d$. Jednocześnie zbiór
$[0,1]^d\setminus\mathbb Q^d$ ma miarę $1$.
Gęstość topologiczna nie pozwala zatem wnioskować
o dodatniości miary ani o jej wielkości.
\end{przyklad}

\begin{przyklad}[Wykres funkcji Lipschitzowskiej]
\label{ex:lipschitz-graph-null}
Niech $a,b\in\mathbb R$, $a\le b$, $L\ge0$ oraz
niech $f:[a,b]\to\mathbb R$ spełnia warunek
$|f(x)-f(y)|\le L|x-y|$ dla wszystkich
$x,y\in[a,b]$.
Pokażemy, że wykres
$\Gamma_f=\{(x,f(x)):x\in[a,b]\}\subseteq\mathbb R^2$
ma miarę $\lambda_2$ równą zero.
Jeśli $a=b$, to $\Gamma_f$ jest singletonem, a więc
ma miarę zero na mocy lematu \ref{lem:hyperplane-null}.
Dalej zakładamy $a<b$ i dzielimy $[a,b]$ na $N$ przedziałów
$[x_{j-1},x_j]$ długości $\delta=(b-a)/N$.
Na $j$-tym przedziale zachodzi
$|f(x)-f(x_{j-1})|\le L\delta$, więc część wykresu
leży w prostokącie o poziomej szerokości $\delta$
i pionowej wysokości $2L\delta$.
Dla każdego $\eta>0$ pokrywamy tę część wykresu
półotwartym prostopadłościanem o bokach długości
$\delta+2\eta$ i $2L\delta+2\eta$.
Biorąc $\eta=\delta^2$, otrzymujemy pokrycie o
łącznej powierzchni co najwyżej
\[
N(\delta+2\delta^2)(2L\delta+2\delta^2)
=(b-a)(1+2\delta)(2L\delta+2\delta^2)
\xrightarrow[N\to\infty]{}0.
\]
Zatem $\lambda_2^*(\Gamma_f)=0$, a z zupełności
$\Gamma_f$ jest mierzalny i ma miarę zero.
Ten argument nie potrzebuje jeszcze twierdzenia Fubiniego.
\end{przyklad}

\begin{przyklad}[Miara kul]
Niech $B(x,r)=\{y\in\mathbb R^d:\|y-x\|_2<r\}$
dla $x\in\mathbb R^d$ i $r>0$.
Z twierdzenia \ref{thm:invariance-lebesgue}
wynika $\lambda_d(B(x,r))=\omega_dr^d$,
gdzie $\omega_d=\lambda_d(B(0,1))$.
Liczba $\omega_d$ jest dodatnia i skończona, ponieważ
$(-1/\sqrt d,1/\sqrt d)^d\subseteq B(0,1)
\subseteq[-1,1]^d$. Wzór
$\omega_d=\pi^{d/2}/\Gamma(d/2+1)$ wymaga dodatkowego
rachunku całkowego i zostanie wyprowadzony później.
\end{przyklad}

\section{Zbiór Cantora: miara zero bez przeliczalności}

\noindent\emph{Powiązanie tematyczne w Bibliotece matematycznej:} \atlasref{AM1-19::ma1:thm:cantor-properties}{AM1-19: własności zbioru Cantora}.


\begin{definicja}[Środkowotrzeci zbiór Cantora]
\label{def:cantor-set}
Niech $C_0=[0,1]$. Z $C_{n-1}$ otrzymujemy $C_n$,
usuwając z każdego z $2^{n-1}$ jego domkniętych
przedziałów składowych otwartą środkową trzecią część.
Definiujemy zbiór Cantora przez
$C=\bigcap_{n=0}^\infty C_n$.
\end{definicja}

\begin{center}
\begin{tikzpicture}[x=10.5cm,y=0.73cm]
\foreach \y/\lab in {0/{$C_0$},1/{$C_1$},2/{$C_2$},3/{$C_3$}} {
  \node[anchor=east] at (-0.045,-\y) {\lab};
}
\fill[blue!45] (0,-0.15) rectangle (1,0.15);
\foreach \a/\b in {0/0.333333,0.666667/1} {
  \fill[blue!45] (\a,-1.15) rectangle (\b,-0.85);
}
\foreach \a/\b in {0/0.111111,0.222222/0.333333,0.666667/0.777778,0.888889/1} {
  \fill[blue!45] (\a,-2.15) rectangle (\b,-1.85);
}
\foreach \a/\b in {0/0.037037,0.074074/0.111111,0.222222/0.259259,0.296296/0.333333,0.666667/0.703704,0.740741/0.777778,0.888889/0.925926,0.962963/1} {
  \fill[blue!45] (\a,-3.15) rectangle (\b,-2.85);
}
\node[anchor=north] at (0,-3.42) {$0$};
\node[anchor=north] at (1,-3.42) {$1$};
\end{tikzpicture}
\par\smallskip
\emph{Pierwsze cztery etapy konstrukcji środkowotrzeciego zbioru Cantora.
W każdym etapie pozostaje $2^n$ przedziałów długości $3^{-n}$.}
\end{center}

\begin{twierdzenie}[Podstawowe własności zbioru Cantora]
\label{thm:cantor-properties}
Zbiór Cantora $C\subseteq[0,1]$ jest niepusty, zwarty,
doskonały (zamknięty i bez punktów izolowanych),
nigdziegęsty (jego domknięcie ma puste wnętrze),
nieprzeliczalny i ma miarę Lebesgue'a zero.
Co więcej, jest w naturalny sposób homeomorficzny
z przestrzenią $\{0,1\}^{\mathbb N}$ wyposażoną
w topologię produktową. W tym przypadku topologia ta
jest generowana przez cylindry, czyli zbiory wszystkich
ciągów zero-jedynkowych o ustalonych pierwszych
$N$ wyrazach, dla dowolnego $N\ge1$.
\end{twierdzenie}

\begin{proof}
Każdy $C_n$ jest niepusty i zwarty, a zbiory są
malejące. Ich przekrój jest więc niepusty i zwarty.
Zbiór $C_n$ jest sumą $2^n$ domkniętych przedziałów
długości $3^{-n}$. Końce tych przedziałów należą do
$C$, ponieważ w dalszych etapach usuwamy wyłącznie
otwarte przedziały środkowe. Dla dowolnego $x\in C$
możemy znaleźć w dowolnie małym otoczeniu $x$
inny koniec przedziału składowego $C_n$.
Stąd $C$ nie ma punktów izolowanych.

Każdy przedział składowy $C_n$ ma długość $3^{-n}$.
Gdyby $C$ zawierał przedział otwarty dodatniej długości,
to musiałby on być zawarty w jednym z przedziałów
składowych każdego $C_n$, co jest niemożliwe dla
dostatecznie dużego $n$. Zatem wnętrze $C$ jest puste;
ponieważ $C$ jest zamknięty, jest nigdziegęsty.

Z monotoniczności miary i skończonej addytywności
otrzymujemy
$0\le\lambda_1(C)\le\lambda_1(C_n)=2^n3^{-n}=(2/3)^n$.
Przechodząc do granicy, dostajemy $\lambda_1(C)=0$.

Na każdym etapie wybrana przez punkt $x\in C$ lewa lub prawa
trzecia część przedziału wyznacza odpowiednio cyfrę $0$ lub $2$.
Długości przedziałów maleją do zera, a ich zagnieżdżenie
ma dokładnie jeden punkt wspólny. Zatem punkt $x\in C$
ma przedstawienie trójkowe
$x=\sum_{n\ge1}2\varepsilon_n3^{-n}$, gdzie
$\varepsilon_n\in\{0,1\}$. W obrębie cyfr $0$ i $2$
to przedstawienie jest jednoznaczne: jeżeli ciągi
różnią się po raz pierwszy na miejscu $k$, to różnica
składników z tego miejsca to $2\cdot3^{-k}$,
a możliwa suma różnic z późniejszych miejsc ma wartość
bezwzględną co najwyżej
$\sum_{n>k}2\cdot3^{-n}=3^{-k}$.
Odwzorowanie
$\Phi:\{0,1\}^{\mathbb N}\to C$ dane przez
$\Phi((\varepsilon_n))=\sum_n2\varepsilon_n3^{-n}$
jest więc bijekcją. Jeśli dwa ciągi mają pierwszych
$N$ cyfr jednakowych, ich obrazy różnią się o co
najwyżej $\sum_{n>N}2\cdot3^{-n}=3^{-N}$;
to dowodzi ciągłości $\Phi$ dla topologii produktowej.
Dziedzina jest zwarta. Można to stwierdzić bez użycia
twierdzenia Tychonoffa: metryka
$\rho(u,v)=\sum_{n\ge1}2^{-n}|u_n-v_n|$
wyznacza tę samą topologię co cylindry; z każdego ciągu
punktów można wybrać podciąg o stałej pierwszej
współrzędnej, następnie jego podciąg o stałej drugiej
współrzędnej itd. Przekątniowy podciąg jest zbieżny
w metryce $\rho$, co dowodzi zwartości metrycznej.
Ponieważ $C$ jest przestrzenią Hausdorffa,
ciągła bijekcja ze zwartej dziedziny jest homeomorfizmem.
Ponieważ $\{0,1\}^{\mathbb N}$ jest nieprzeliczalna,
tak samo jest z $C$.
\end{proof}

\begin{wniosek}[Lebesgue'owskie zbiory nieborelowskie]
\label{cor:nonborel-lebesgue}
Istnieją podzbiory $\mathbb R$ mierzalne w sensie
Lebesgue'a, które nie są borelowskie. W szczególności
$\mathcal B(\mathbb R)\subsetneq\mathcal L_1$.
\end{wniosek}

\begin{proof}
Zbiór Cantora $C$ jest borelowski, ma miarę zero
i ma moc continuum $\mathfrak c$.
Każdy jego podzbiór jest mierzalny w sensie Lebesgue'a,
ponieważ miara $\lambda_1$ jest zupełna.
Rodzina wszystkich podzbiorów $C$ ma moc
$2^{\mathfrak c}>\mathfrak c$ (twierdzenie Cantora),
podczas gdy rodzina zbiorów borelowskich prostej ma
moc co najwyżej $\mathfrak c$. Uzasadnijmy ostatnie
stwierdzenie. Rodzina przedziałów o końcach wymiernych
jest przeliczalną bazą topologii $\mathbb R$. Każdy
zbiór borelowski można opisać za pomocą
\emph{przeliczalnego kodu borelowskiego}: dobrze
ufundowanego drzewa, którego węzły są oznaczone
operacją dopełnienia, sumy przeliczalnej albo
jednym z przedziałów bazowych. Takie opisy otrzymuje
się indukcyjnie po etapach tworzenia najmniejszej
rodziny zamkniętej na dopełnienia i sumy przeliczalne.
Każde drzewo o przeliczalnie wielu węzłach jest
kodowane przez podzbiór zbioru wszystkich skończonych
ciągów liczb naturalnych wraz z przeliczalnym
etykietowaniem. Liczba takich kodów wynosi co najwyżej
$2^{\aleph_0}=\mathfrak c$. Wynika stąd, że zbiorów
borelowskich jest najwyżej $\mathfrak c$, zaś
podzbiorów $C$ jest $2^{\mathfrak c}$; na mocy
twierdzenia Cantora $2^{\mathfrak c}>\mathfrak c$.
Zatem nie każdy podzbiór $C$ jest borelowski.
\end{proof}

\section{Zbiór dodatniej miary bez wnętrza}

\begin{przyklad}[Zbiór Smitha--Volterry--Cantora]
\label{ex:fat-cantor}
Rozpoczynamy od $F_0=[0,1]$. Na etapie $n\ge1$
z każdego z $2^{n-1}$ przedziałów pozostawionych
w poprzednim etapie usuwamy otwarty przedział środkowy
długości $4^{-n}$; pozostały zbiór oznaczamy $F_n$.
Można to wykonać na każdym etapie: długość każdego
z dwóch potomnych przedziałów jest dodatnia.
W rzeczy samej, jeśli $\ell_n$ oznacza długość jednego
przedziału składowego $F_n$, to
$\ell_0=1$ i $\ell_n=(\ell_{n-1}-4^{-n})/2$.
Indukcja daje dokładny wzór
$\ell_n=2^{-n-1}+2^{-2n-1}>0$; co więcej,
$\ell_{n-1}>4^{-n}$, a zatem każde kolejne
usunięcie jest geometrycznie możliwe.

Łączna długość usuniętych przedziałów na etapie $n$
wynosi $2^{n-1}4^{-n}=2^{-n-1}$. Ponieważ otwarte
przedziały usuwane w różnych etapach są rozłączne,
dla $F=\bigcap_nF_n$ mamy
\[
\lambda_1(F)=1-\sum_{n=1}^\infty2^{-n-1}=\frac12.
\]
Zbiór $F$ jest zwarty. Ponieważ długości przedziałów
składowych $F_n$ dążą do zera, $F$ nie zawiera
niepustego przedziału otwartego, jest więc nigdziegęsty.
Końce każdego przedziału składowego pozostają w $F$,
a zatem, podobnie jak dla zbioru Cantora, żaden punkt
$F$ nie jest izolowany. Zbiór $F$ jest więc doskonały
i ma dodatnią miarę. A zatem nawet jednoczesna
zwartość, brak wnętrza i nigdziegęstość nie
wymuszają miary zero.
\end{przyklad}

\section{Funkcja Cantora i miara osobliwa: zapowiedź}

\begin{przyklad}[Funkcja Cantora]
\label{ex:cantor-function}
Dla $x\in C$ o przedstawieniu
$x=\sum_{n\ge1}2\varepsilon_n3^{-n}$ określamy
$F(x)=\sum_{n\ge1}\varepsilon_n2^{-n}$.
Na każdej składowej otwartej $[0,1]\setminus C$
przedłużamy $F$ jako funkcję stałą, przyjmując
wspólną wartość na końcach tej składowej.
Wartości końcowe są równe: końce przedziału
usuniętego na etapie $n$ mają przedstawienia
trójkowe, których pierwsze $n-1$ cyfr są jednakowe,
a pozostałe cyfry są odpowiednio $0,2,2,\ldots$
oraz $2,0,0,\ldots$, co w dwójkowym zapisie
prowadzi do tej samej liczby.

Tak uzyskana funkcja $F:[0,1]\to[0,1]$
jest ciągła i niemalejąca. Można to sprawdzić za
pomocą kolejnych ciągłych przybliżeń $F_n$:
na każdym przedziale składowym $C_n$ funkcja $F_n$
jest liniowa, a na przedziałach już usuniętych
stała, przy czym $F_n(0)=0$ oraz $F_n(1)=1$.
Dokładniej: na każdym z $2^n$ domkniętych przedziałów
składowych $C_n$ wartości na końcach różnią się o
$2^{-n}$, a na przedziałach już usuniętych $F_n$
pozostaje stała. Z konstrukcji wynika, że
$\|F_{n+1}-F_n\|_\infty\le 2^{-n}$.
Dla $m>n$ dostajemy
$\|F_m-F_n\|_\infty\le\sum_{k=n}^{m-1}2^{-k}$,
więc $(F_n)$ jest jednostajnym ciągiem Cauchy'ego.
Jego granica jest ciągła i niemalejąca. Na końcach
każdego usuniętego przedziału granica przyjmuje
wspólną, raz ustaloną wartość, a na $C$ zgadza się
z szeregiem definiującym $F$, ponieważ na przedziale
składowym $C_n$ zarówno $F_n$, jak i wartość tego
szeregu leżą między tymi samymi wartościami końcowymi,
oddalonymi o $2^{-n}$. To identyfikuje granicę
jednostajną z funkcją $F$ z definicji.

Ponadto $F$ jest stała na każdym przedziale
składowym $[0,1]\setminus C$, więc $F'(x)=0$
dla wszystkich $x\in(0,1)\setminus C$.
Ponieważ $\lambda_1(C)=0$, mamy $F'=0$
prawie wszędzie, ale $F(1)-F(0)=1$.
Nie jest to sprzeczność: funkcja Cantora nie jest
absolutnie ciągła. Przypomnijmy, że funkcja
$h:[a,b]\to\mathbb R$ jest absolutnie ciągła, jeśli
dla każdego $\varepsilon>0$ istnieje $\delta>0$
takie, że dla każdej skończonej rodziny parami
rozłącznych przedziałów $(a_k,b_k)\subseteq[a,b]$
warunek $\sum_k(b_k-a_k)<\delta$ implikuje
$\sum_k|h(b_k)-h(a_k)|<\varepsilon$.
Dla funkcji Cantora bierzemy $2^n$ przedziałów
składowych $C_n$. Suma ich długości wynosi
$(2/3)^n\to0$, podczas gdy suma przyrostów $F$
na tych przedziałach wynosi
$2^n\cdot2^{-n}=1$. Sprzeczność z powyższą definicją
dla $\varepsilon=1/2$ dowodzi braku absolutnej ciągłości.
W rozdziałach o całce Lebesgue'a i rozkładzie miar
zobaczymy, że $F$ wyznacza bezatomową miarę
osobliwą względem miary Lebesgue'a, to znaczy
skoncentrowaną na zbiorze o mierze Lebesgue'a zero.
\end{przyklad}

\section{Zbiór Vitaliego i ograniczenia mierzalności}

\noindent\emph{Powiązanie tematyczne w Bibliotece matematycznej:} \atlasref{LST-26::def:choice}{LST-26: aksjomat wyboru}.


\begin{twierdzenie}[Istnienie zbioru niemierzalnego]
\label{thm:vitali}
Przy aksjomacie wyboru istnieje podzbiór
$V\subseteq[0,1]$, który nie jest mierzalny
w sensie Lebesgue'a.
\end{twierdzenie}

\begin{proof}
Na $[0,1]$ wprowadzamy relację równoważności
$x\sim y$ wtedy i tylko wtedy, gdy $x-y\in\mathbb Q$.
Z każdej klasy wybieramy dokładnie jeden element,
tworząc zbiór reprezentantów $V\subseteq[0,1]$.
Ten krok wykorzystuje aksjomat wyboru.

Niech $Q=\mathbb Q\cap[-1,1]$. Zbiory $V+q$ dla
$q\in Q$ są parami rozłączne. Rzeczywiście, gdyby
$v_1+q_1=v_2+q_2$, to
$v_1-v_2=q_2-q_1\in\mathbb Q$, więc $v_1,v_2$
należałyby do tej samej klasy; wybór reprezentantów
daje $v_1=v_2$ i następnie $q_1=q_2$.

Mamy też
\begin{equation}\label{eq:vitali-sandwich}
[0,1]\subseteq\bigcup_{q\in Q}(V+q)\subseteq[-1,2].
\end{equation}
Pierwsze zawieranie wynika stąd, że dla dowolnego
$x\in[0,1]$ istnieje $v\in V$ z $x-v\in\mathbb Q$,
a jednocześnie $|x-v|\le1$.

Przypuśćmy, że $V$ jest mierzalny. Z niezmienniczości
translacyjnej $\lambda_1(V+q)=\lambda_1(V)$.
Zbiory $V+q$ są rozłączne, $Q$ jest przeliczalny i
nieskończony, a więc
\[
\lambda_1\Big(\bigcup_{q\in Q}(V+q)\Big)
=\sum_{q\in Q}\lambda_1(V)
=\begin{cases}0,&\lambda_1(V)=0,\\
\infty,&\lambda_1(V)>0.
\end{cases}
\]
Z \eqref{eq:vitali-sandwich} wynika jednak, że lewa
strona znajduje się między $1$ a $3$. To sprzeczność.
Zatem $V\notin\mathcal L_1$.
\end{proof}

\begin{wniosek}[Nie istnieje miara długości wszystkich podzbiorów]
\label{cor:no-all-subsets-invariant}
Nie istnieje przeliczalnie addytywna miara
$\nu:\mathcal P(\mathbb R)\to[0,\infty]$, która
jest niezmiennicza względem translacji i spełnia
$\nu([0,1])=1$.
\end{wniosek}

\begin{proof}
Najpierw musimy ustalić $\nu(\{0\})=0$.
Gdyby $\nu(\{0\})>0$, przeliczalnie wiele
rozłącznych translacji singletonu zawartych w $[0,1]$
miałoby miarę nieskończoną, sprzecznie z
$\nu([0,1])=1$. Zatem $\nu(\{0\})=0$,
a translacyjna niezmienniczość daje
$\nu(\{x\})=0$ dla każdego $x\in\mathbb R$.
Wobec tego $\nu([0,1))=\nu([0,1])=1$.
Trzy rozłączne przedziały $[-1,0)$, $[0,1)$
i $[1,2)$ mają po jednej jednostce miary;
po dodaniu punktu $2$ otrzymujemy
$\nu([-1,2])=3$. Teraz dowód Vitaliego stosuje się
bez zmian do $\nu$, ponieważ wszystkie podzbiory
w nim występujące należałyby do dziedziny $\nu$.
Otrzymujemy sprzeczność.
\end{proof}

\begin{uwaga}[Precyzyjny sens niemożliwości]
Powyższe twierdzenie nie zabrania istnienia miar na
$\mathcal P(\mathbb R)$, takich jak miara licząca
czy miara zerowa. Stwierdza niemożliwość jednoczesnego
zachowania \emph{przeliczalnej addytywności},
\emph{niezmienniczości translacyjnej} i
\emph{normalizacji długości przedziału}.
\end{uwaga}

\section{Podsumowanie i przygotowanie do całki}

Wykazaliśmy, że objętość prostopadłościanów
półotwartych wyznacza jedną miarę borelowską,
której uzupełnieniem jest miara Lebesgue'a. Zbiory
mierzalne mają aproksymacje otwarte i zwarte.
Gęstość, nigdziegęstość czy nieprzeliczalność
nie przesądzają o wielkości miary zbioru.

W następnych rozdziałach zdefiniujemy funkcje
mierzalne i całkę oraz wyjaśnimy rolę absolutnej
ciągłości w związku między całką a pochodną.

% ROZDZIAŁ 5 — plik włączany przez main.tex
\chapter{Struktura przestrzeni mierzalnych}
\label{ch:struktura-przestrzeni}

\providecommand{\atlasref}[2]{#2}


\section{Dlaczego sama definicja \texorpdfstring{$\sigma$}{sigma}-algebry nie wystarcza?}

Znamy już konstrukcję $\sigma$-algebr, miar i miary Lebesgue'a.
W dalszej teorii będziemy jednak przenosić miary przez odwzorowania,
budować przestrzenie produktowe i określać zbiory na podstawie
współrzędnych. Chcemy więc wiedzieć, jak z danych o prostych
zbiorach odzyskać mierzalność dowolnego odwzorowania.

W tym rozdziale przez $(X,\mathcal F)$ i $(Y,\mathcal G)$ rozumiemy
przestrzenie mierzalne. Nie zakładamy istnienia na nich miary,
jeśli nie podajemy jej wyraźnie. Zbiór $\mathbb N$ składa się z
dodatnich liczb całkowitych. Używamy oznaczeń
$\mathcal B(T)$ dla $\sigma$-algebry borelowskiej przestrzeni
topologicznej $T$ oraz $\mathcal F|_A=\{E\cap A:E\in\mathcal F\}$
dla śladu $\sigma$-algebry na $A\subseteq X$.

\section{Ślad \texorpdfstring{$\sigma$}{sigma}-algebry i podprzestrzenie}

Ślad $\mathcal F|_A$ zdefiniowaliśmy już w rozdziale~\ref{ch:przestrzenie-mierzalne}.
Teraz badamy jego zgodność z generowaniem $\sigma$-algebr
i topologią podprzestrzeni.


\begin{lemat}[Ślad i generowanie]
\label{lem:trace-generation}
Niech $X$ będzie zbiorem, $\mathcal E\subseteq\mathcal P(X)$,
$A\subseteq X$ i $\mathcal F=\sigma_X(\mathcal E)$.
Wtedy
\[
\mathcal F|_A=\sigma_A\bigl(\{A\cap E:E\in\mathcal E\}\bigr).
\]
\end{lemat}
\begin{proof}
Oznaczmy prawą stronę przez $\mathcal H$. Zbiory
$A\cap E$ dla $E\in\mathcal E$ należą do $\mathcal F|_A$;
ponieważ $\mathcal F|_A$ jest $\sigma$-algebrą na $A$,
mamy $\mathcal H\subseteq\mathcal F|_A$.
W drugą stronę rozważmy rodzinę
$\mathcal D=\{E\subseteq X:A\cap E\in\mathcal H\}$.
Zawiera ona $X$, gdyż $A\in\mathcal H$. Jeżeli
$E\in\mathcal D$, to
$A\cap(X\setminus E)=A\setminus(A\cap E)\in\mathcal H$.
Dla $E_n\in\mathcal D$ mamy
$A\cap\bigcup_nE_n=\bigcup_n(A\cap E_n)\in\mathcal H$.
Zatem $\mathcal D$ jest $\sigma$-algebrą na $X$, zawierającą
$\mathcal E$. Wynika stąd
$\mathcal F=\sigma_X(\mathcal E)\subseteq\mathcal D$, czyli
$\mathcal F|_A\subseteq\mathcal H$.
\end{proof}

\begin{twierdzenie}[Zbiory borelowskie podprzestrzeni]
\label{thm:borel-trace}
Niech $T$ będzie przestrzenią topologiczną i $A\subseteq T$
wyposażonym w topologię podprzestrzeni. Wtedy
$\mathcal B(A)=\mathcal B(T)|_A$.
\end{twierdzenie}
\begin{proof}
Z definicji topologii podprzestrzeni zbiory otwarte w $A$
mają postać $A\cap U$, gdzie $U$ jest otwarty w $T$.
Zastosowanie lematu \ref{lem:trace-generation} do rodziny
wszystkich zbiorów otwartych przestrzeni $T$ daje tezę.
\end{proof}

\begin{przyklad}[Ślad na zbiorze nieborelowskim]
Istnieje zbiór $A\subseteq\mathbb R$ nieborelowski
(wniosek \ref{cor:nonborel-lebesgue}). Mimo to możemy
rozpatrywać $(A,\mathcal B(\mathbb R)|_A)$.
To jest przestrzeń mierzalna; brak borelowskości $A$
w przestrzeni otaczającej nie uniemożliwia
zdefiniowania śladu. Nie wolno jednak utożsamiać
$\mathcal B(\mathbb R)|_A$ z rodziną
$\{E\in\mathcal B(\mathbb R):E\subseteq A\}$:
ta druga rodzina może nawet nie zawierać całego $A$.
\end{przyklad}

\section{Mierzalność odwzorowania i rodziny generujące}

Korzystamy z definicji mierzalności~\ref{def:measurable-map} i kryterium generatorów~\ref{thm:measurable-generator}, udowodnionych w rozdziale~\ref{ch:przestrzenie-mierzalne}. Ich podstawą są tożsamości dla przeciwobrazów: $f^{-1}(Y\setminus B)=X\setminus f^{-1}(B)$ oraz $f^{-1}(\bigcup_n B_n)=\bigcup_n f^{-1}(B_n)$. Nie wymagają one różnowartościowości ani surjektywności.

\begin{wniosek}[Prawostronny test mierzalności]
\label{cor:real-borel-tests}
Niech $(X,\mathcal F)$ będzie przestrzenią mierzalną,
a $f:X\to\mathbb R$ odwzorowaniem. Jest ono
$(\mathcal F,\mathcal B(\mathbb R))$-mierzalne wtedy i tylko
wtedy, gdy $\{x\in X:f(x)>q\}\in\mathcal F$
dla każdego $q\in\mathbb Q$.
\end{wniosek}
\begin{proof}
Rodzina prawych półprostych $(q,\infty)$, $q\in\mathbb Q$,
generuje $\mathcal B(\mathbb R)$: przez dopełnienia
otrzymujemy lewe półproste domknięte, a następnie
wszystkie otwarte przedziały. Stosujemy więc
kryterium generatorów~\ref{thm:measurable-generator}.
Warunki z lewymi półprostymi zostały już udowodnione
w twierdzeniu~\ref{thm:measurable-real}.
\end{proof}

\begin{twierdzenie}[Mierzalność ograniczenia]
\label{thm:restriction-meas-5}
Niech $(X,\mathcal F)$ i $(Y,\mathcal G)$ będą przestrzeniami
mierzalnymi, $f:X\to Y$ odwzorowaniem
$(\mathcal F,\mathcal G)$-mierzalnym oraz $A\subseteq X$.
Wówczas $f|_A:(A,\mathcal F|_A)\to(Y,\mathcal G)$ jest mierzalne.
\end{twierdzenie}
\begin{proof}
Dla każdego $B\in\mathcal G$ mamy
$(f|_A)^{-1}(B)=A\cap f^{-1}(B)\in\mathcal F|_A$.
Złożenia odwzorowań mierzalnych oraz mierzalność funkcji
ciągłych zostały już wykazane w twierdzeniu~\ref{thm:composition-measurable}.
\end{proof}

\section{Najmniejsza \texorpdfstring{$\sigma$}{sigma}-algebra, względem której funkcja jest mierzalna}

\begin{definicja}[$\sigma$-algebra generowana przez odwzorowanie]
\label{def:sigma-of-function}
Dla odwzorowania $f:X\to(Y,\mathcal G)$ określamy
\[
\sigma(f):=\{f^{-1}(B):B\in\mathcal G\}.
\]
Jeżeli $(f_i)_{i\in I}$ jest rodziną odwzorowań
$f_i:X\to(Y_i,\mathcal G_i)$, to
$\sigma(f_i:i\in I):=\sigma_X\bigl(\bigcup_{i\in I}\sigma(f_i)\bigr)$.
\end{definicja}

\begin{twierdzenie}[Minimalna struktura mierzalna]
\label{thm:initial-sigma}
Dla każdego $f:X\to(Y,\mathcal G)$ rodzina $\sigma(f)$
jest $\sigma$-algebrą na $X$ i jest najmniejszą
$\sigma$-algebrą, względem której $f$ jest mierzalne.
Dla rodziny $(f_i)_{i\in I}$ analogiczną własność
ma $\sigma(f_i:i\in I)$.
\end{twierdzenie}
\begin{proof}
Przeciwobraz całego $Y$ to $X$, przeciwobrazy zachowują
dopełnienia i przeliczalne sumy, więc $\sigma(f)$
jest $\sigma$-algebrą. Jeśli $f$ jest mierzalne
względem pewnej $\mathcal F$, to z definicji
$f^{-1}(B)\in\mathcal F$ dla każdego $B\in\mathcal G$,
toteż $\sigma(f)\subseteq\mathcal F$.
Dla rodziny odwzorowań powtarzamy argument dla
$\sigma$-algebry generowanej sumą wszystkich
rodzin przeciwobrazów.
\end{proof}

\begin{przyklad}[Rozróżnialność punktów]
Dla $X=\{a,b,c,d\}$ i funkcji $f:X\to(\mathbb R,\mathcal B(\mathbb R))$ o
$f(a)=f(b)=0$ oraz $f(c)=f(d)=1$ otrzymujemy
$\sigma(f)=\{\varnothing,\{a,b\},\{c,d\},X\}$.
Funkcja nie potrafi odróżnić $a$ od $b$ ani $c$ od $d$.
To wyjaśnia, dlaczego $\sigma(f)$ może być ściśle mniejsza
od $\mathcal P(X)$.
\end{przyklad}

\section{Iloczyn przestrzeni mierzalnych}

\noindent\emph{Powiązanie tematyczne w Bibliotece matematycznej:} \atlasref{GT-04::def:topog-iloczyn}{GT-04: topologia iloczynu (dla porównania konstrukcji)}.


\begin{definicja}[$\sigma$-algebra produktowa]
\label{def:product-sigma}
Dla przestrzeni mierzalnych $(X,\mathcal F)$ i
$(Y,\mathcal G)$ definiujemy
\[
\mathcal F\otimes\mathcal G
:=\sigma_{X\times Y}\bigl(\{A\times B:A\in\mathcal F,\ B\in\mathcal G\}\bigr).
\]
Parę $(X\times Y,\mathcal F\otimes\mathcal G)$
nazywamy \emph{iloczynem przestrzeni mierzalnych}.
Nie wprowadzamy w tym miejscu miary produktowej;
jej konstrukcja należy do planowanego rozdziału 12.
\end{definicja}

\begin{twierdzenie}[Charakteryzacja przez współrzędne]
\label{thm:product-projections}
Niech $(X,\mathcal F)$ i $(Y,\mathcal G)$ będą
przestrzeniami mierzalnymi, a
$\pi_X:X\times Y\to X$, $\pi_Y:X\times Y\to Y$
rzutowaniami. Wówczas
$\mathcal F\otimes\mathcal G=\sigma(\pi_X,\pi_Y)$.
Ponadto odwzorowanie $h:Z\to X\times Y$ z przestrzeni
mierzalnej $(Z,\mathcal H)$ jest
$(\mathcal H,\mathcal F\otimes\mathcal G)$-mierzalne
wtedy i tylko wtedy, gdy oba odwzorowania
$\pi_X\circ h:Z\to X$ i $\pi_Y\circ h:Z\to Y$ są mierzalne.
\end{twierdzenie}
\begin{proof}
Ponieważ $\pi_X^{-1}(A)=A\times Y$ i
$\pi_Y^{-1}(B)=X\times B$, każda z rodzin
$\sigma(\pi_X)$, $\sigma(\pi_Y)$ zawiera się w
$\mathcal F\otimes\mathcal G$. Z drugiej strony
$A\times B=\pi_X^{-1}(A)\cap\pi_Y^{-1}(B)$,
zatem każdy generator $\mathcal F\otimes\mathcal G$
należy do $\sigma(\pi_X,\pi_Y)$. To dowodzi równości.
Drugą część otrzymujemy, sprawdzając przeciwobrazy
generatorów $A\times B$:
$h^{-1}(A\times B)=(\pi_X\circ h)^{-1}(A)
\cap(\pi_Y\circ h)^{-1}(B)$.
Konieczność wynika z mierzalności rzutowań i złożenia.
\end{proof}

\begin{twierdzenie}[Produkt borelowski przestrzeni spełniających drugi aksjomat przeliczalności]
\label{thm:borel-product}
Niech $S,T$ będą przestrzeniami topologicznymi, z których
każda ma przeliczalną bazę topologii. Wtedy
\[
\mathcal B(S\times T)=\mathcal B(S)\otimes\mathcal B(T),
\]
gdzie na $S\times T$ rozważamy topologię produktową.
\end{twierdzenie}
\begin{proof}
Rzutowania $\pi_S,\pi_T$ są ciągłe, więc
$A\times T=\pi_S^{-1}(A)$ i $S\times B=\pi_T^{-1}(B)$
są borelowskie dla borelowskich $A,B$.
Stąd $\mathcal B(S)\otimes\mathcal B(T)\subseteq
\mathcal B(S\times T)$.
Wybierzmy przeliczalne bazy $\mathcal U,\mathcal V$.
Wszystkie $U\times V$, gdzie $U\in\mathcal U$ i
$V\in\mathcal V$, należą do
$\mathcal B(S)\otimes\mathcal B(T)$ i tworzą
przeliczalną bazę produktu. Każdy zbiór otwarty
w $S\times T$ jest więc przeliczalną sumą takich
prostokątów, czyli należy do $\mathcal B(S)\otimes\mathcal B(T)$.
Generując $\sigma$-algebrę, otrzymujemy zawieranie odwrotne.
\end{proof}

\begin{wniosek}[Borelowskie odwzorowania wektorowe]
\label{cor:coordinate-measurable}
Niech $(X,\mathcal F)$ będzie przestrzenią mierzalną,
$f=(f_1,\ldots,f_d):X\to\mathbb R^d$ i $d\ge1$.
Odwzorowanie $f$ jest $(\mathcal F,\mathcal B(\mathbb R^d))$-mierzalne
wtedy i tylko wtedy, gdy każda $f_j:X\to\mathbb R$
jest $(\mathcal F,\mathcal B(\mathbb R))$-mierzalna.
\end{wniosek}
\begin{proof}
Z twierdzenia \ref{thm:borel-product}, zastosowanego
indukcyjnie, wynika
$\mathcal B(\mathbb R^d)=\bigotimes_{j=1}^d\mathcal B(\mathbb R)$.
Wykorzystujemy teraz twierdzenie \ref{thm:product-projections}.
\end{proof}

\begin{uwaga}[Dlaczego potrzebna jest przeliczalna baza?]
Dla dowolnych przestrzeni topologicznych $S,T$ zachodzi
$\mathcal B(S)\otimes\mathcal B(T)\subseteq\mathcal B(S\times T)$,
ponieważ rzutowania są ciągłe. Równość może jednak zawieść.
Dla przykładu niech $D$ będzie zbiorem o mocy większej niż
$\mathfrak c=|\mathbb R|$, wyposażonym w topologię dyskretną.
Wtedy przekątna $\Delta=\{(x,x):x\in D\}$ jest zamknięta
w $D\times D$, ale nie należy do
$\mathcal P(D)\otimes\mathcal P(D)$.
Istotnie, każdy element $\sigma$-algebry generowanej
przez dowolną rodzinę zbiorów należy do $\sigma$-algebry
generowanej przez pewną jej przeliczalną podrodzinę:
rodzina zbiorów o tej własności sama jest $\sigma$-algebrą.
Gdyby $\Delta$ była mierzalna produktowo, istniałyby
prostokąty $A_n\times B_n$, $n\ge1$, generujące
$\sigma$-algebrę zawierającą $\Delta$.
Dla różnych $x,y\in D$ punkty $(x,x)$ i $(x,y)$
muszą być rozróżnione przez któryś z tych prostokątów,
inaczej wszystkie zbiory z wygenerowanej $\sigma$-algebry
zawierałyby oba albo żaden z nich. Ponieważ pierwsza
współrzędna obu punktów jest taka sama, pewien $B_n$
musi zawierać dokładnie jeden z punktów $x,y$.
Odwzorowanie $x\mapsto(\chi_{B_n}(x))_{n\ge1}$
byłoby więc iniekcją $D\to\{0,1\}^{\mathbb N}$, co
przeczy $|D|>\mathfrak c$.
\end{uwaga}

\section{Przestrzenie standardowe borelowskie}

\noindent\emph{Powiązanie tematyczne w Bibliotece matematycznej:} \atlasref{GT-19::def:topog-polskie}{GT-19: przestrzenie polskie}.


\begin{definicja}[Przestrzeń polska i standardowa borelowska]
\label{def:standard-borel}
Przestrzeń topologiczną nazywamy \emph{polską}, jeśli
jej topologię wyznacza pewna metryka zupełna, względem
której istnieje przeliczalny zbiór gęsty. Przestrzeń mierzalną $(X,\mathcal F)$ nazywamy
\emph{standardową borelowską}, jeżeli istnieje przestrzeń
polska $P$ i bijekcja $b:X\to P$, dla której
$b$ oraz $b^{-1}$ są mierzalne względem
$\mathcal F$ i $\mathcal B(P)$. Taką bijekcję
nazywamy \emph{izomorfizmem przestrzeni mierzalnych}.
\end{definicja}

\begin{uwaga}[Uwaga o definicji]
Równoważnie można zażądać, aby $\mathcal F$ była
$\sigma$-algebrą borelowską pewnej polskiej topologii na $X$.
Istotny jest charakter \emph{mierzalny}, a nie zgodność
z jakąkolwiek wcześniej wybraną topologią. W definicji
żądaliśmy jawnie mierzalności odwrotności; dla ogólnych
przestrzeni mierzalnych sama mierzalna bijekcja nie wystarcza.
\end{uwaga}

\begin{przyklad}[Klasyczne przestrzenie standardowe]
Przestrzenie $(\mathbb R^d,\mathcal B(\mathbb R^d))$
($d\ge1$), $(\mathbb N,\mathcal P(\mathbb N))$ oraz
$(\{0,1\}^{\mathbb N},\mathcal B(\{0,1\}^{\mathbb N}))$
są standardowe borelowskie. W pierwszych dwóch przypadkach
korzystamy ze zwykłej metryki euklidesowej lub metryki
dyskretnej. W trzecim stosujemy metrykę
$\rho(u,v)=\sum_{n\ge1}2^{-n}|u_n-v_n|$ z rozdziału
\ref{ch:lebesgue}. Zupełność wynika stąd, że w ciągu
Cauchy'ego każda współrzędna musi od pewnego miejsca
być stała; przeliczalnym zbiorem gęstym są ciągi
ostatecznie równe zeru. Metryka $\rho$ wyznacza topologię
produktową, więc mamy przestrzeń polską.
\end{przyklad}

\begin{twierdzenie}[Przeliczalne generowanie standardowych przestrzeni]
\label{thm:standard-countably-generated}
Jeżeli $(X,\mathcal F)$ jest standardową przestrzenią
borelowską, to istnieje przeliczalna rodzina
$\mathcal E\subseteq\mathcal F$ taka, że
$\mathcal F=\sigma_X(\mathcal E)$. Ponadto dla
każdej pary różnych punktów $x,y\in X$ istnieje
$E\in\mathcal E$ zawierający dokładnie jeden z nich.
\end{twierdzenie}
\begin{proof}
Niech $b:X\to P$ będzie izomorfizmem na przestrzeń polską.
Każda przestrzeń metryczna o przeliczalnym zbiorze gęstym
ma przeliczalną bazę: bierzemy kule o środkach w tym zbiorze
i dodatnich promieniach wymiernych. Niech
$\mathcal U=\{U_n:n\ge1\}$ będzie taką bazą $P$.
Rodzina $\mathcal E=\{b^{-1}(U_n):n\ge1\}$
jest przeliczalna i generuje $\mathcal F$, gdyż
$\mathcal B(P)=\sigma_P(\mathcal U)$, a przeciwobrazy
zachowują dopełnienia i sumy przeliczalne.
Jeśli $x\ne y$, to $b(x)\ne b(y)$; w przestrzeni
metrycznej istnieje zbiór otwarty zawierający $b(x)$,
lecz nie $b(y)$. Z własności bazy można wybrać
$U_n$ o tej samej własności, a $b^{-1}(U_n)$
oddziela $x$ od $y$.
\end{proof}

\begin{przyklad}[Uzupełnienie miary zmienia typ przestrzeni]
\label{ex:completion-not-standard}
Przestrzeń $(\mathbb R,\mathcal B(\mathbb R))$ jest
standardowa borelowska, ale
$(\mathbb R,\mathcal L_1)$ nie jest standardowa borelowska.
Z twierdzenia \ref{thm:standard-countably-generated}
każda standardowa $\sigma$-algebra jest przeliczalnie
generowana. Każda $\sigma$-algebra generowana przez
przeliczalną rodzinę ma co najwyżej $\mathfrak c$
elementów (kodowanie zbiorów borelowskich drzewami,
jak w dowodzie wniosku \ref{cor:nonborel-lebesgue}).
Tymczasem $\mathcal L_1$ zawiera wszystkie podzbiory
zbioru Cantora, a więc co najmniej $2^{\mathfrak c}>
\mathfrak c$ elementów. Nie może być przeliczalnie
generowana. Przykład pokazuje, że \emph{zupełność miary}
i \emph{standardowość borelowska} są innymi własnościami.
\end{przyklad}

\section{Struktura mierzalna a miara obrazowa}

Przypomnijmy, że jeśli $\mu$ jest miarą na $(X,\mathcal F)$,
a $f:(X,\mathcal F)\to(Y,\mathcal G)$ jest mierzalne,
to \emph{miara obrazowa} $f_\#\mu$ określona jest wzorem
$(f_\#\mu)(B)=\mu(f^{-1}(B))$, $B\in\mathcal G$
(definicja z rozdziału \ref{ch:miary-wlasnosci}).
To poprawne określenie miary, ponieważ przeciwobrazy
zachowują rozłączność i przeliczalne sumy.

\begin{twierdzenie}[Funktorialność miar obrazowych]
\label{thm:pushforward-composition}
Niech $\mu$ będzie miarą na $(X,\mathcal F)$,
$f:(X,\mathcal F)\to(Y,\mathcal G)$ i
$g:(Y,\mathcal G)\to(Z,\mathcal H)$ niech będą mierzalne.
Wtedy
\[
(g\circ f)_\#\mu=g_\#(f_\#\mu).
\]
\end{twierdzenie}
\begin{proof}
Dla każdego $C\in\mathcal H$ z definicji i
rachunku przeciwobrazów
$(g\circ f)_\#\mu(C)=\mu(f^{-1}(g^{-1}(C)))
=(f_\#\mu)(g^{-1}(C))=g_\#(f_\#\mu)(C)$.
Równość miar oznacza tu równość ich wartości
na wszystkich zbiorach mierzalnych.
\end{proof}

\begin{przyklad}[Miara obrazowa przez funkcję stałą]
Niech $(X,\mathcal F)$ i $(Y,\mathcal G)$ będą przestrzeniami
mierzalnymi, $y_0\in Y$, $\mu$ miarą na $(X,\mathcal F)$,
a $f:X\to Y$ funkcją stałą o wartości $y_0$.
Dla każdego $B\in\mathcal G$ mamy $(f_\#\mu)(B)=\mu(X)$,
gdy $y_0\in B$, oraz $(f_\#\mu)(B)=0$, gdy $y_0\notin B$.
Jeśli $\mu(X)<\infty$, można zapisać
$f_\#\mu=\mu(X)\delta_{y_0}$, gdzie $\delta_{y_0}$
wyraża miarę Diraca w punkcie $y_0$.
Gdy $\mu(X)=\infty$, pozostajemy przy podanym wzorze
zbiorowym: wyrażenie $\infty\cdot0$ nie jest w zwykłej
arytmetyce rozszerzonej określone.
\end{przyklad}

\section{Podsumowanie i powiązania}

Rodziny generujące upraszczają testy mierzalności.
Ślad opisuje przestrzenie indukowane, a produkt
sprowadza mierzalność wektorową do współrzędnych.
Ważnym rozróżnieniem jest różnica między przestrzenią
standardową borelowską a uzupełnioną przestrzenią
z miarą Lebesgue'a.

W następnym rozdziale wykorzystamy te narzędzia do
mierzalności funkcji rzeczywistych, ich aproksymacji
funkcjami prostymi i badania granic ciągów funkcji.

% ROZDZIAŁ 6 — plik włączany przez main.tex
\chapter{Funkcje mierzalne i ich aproksymacja}
\label{ch:funkcje-mierzalne}

\section{Funkcje jako podstawowy obiekt teorii miary}

Dotychczas mierzyliśmy zbiory. W praktyce chcemy mierzyć
również zbiory postaci $\{x:f(x)>a\}$, przybliżać
funkcje funkcjami o skończenie wielu wartościach, a następnie
określić całkę. Właśnie dlatego potrzebujemy pojęcia
funkcji mierzalnej. Ważne, by odróżniać \emph{mierzalność}
od \emph{całkowalności}: pierwsza pozwala dopiero
postawić drugie pytanie.

W całym rozdziale $(X,\mathcal F)$ oznacza przestrzeń mierzalną.
Jeśli potrzebna jest miara, zapisujemy $(X,\mathcal F,\mu)$.
Przez $\overline{\mathbb R}=[-\infty,\infty]$ rozumiemy
rozszerzoną prostą rzeczywistą wyposażoną w topologię
porządkową; jej $\sigma$-algebrą jest
$\mathcal B(\overline{\mathbb R})$. Każda funkcja
$\mathbb R$-wartościowa może być traktowana również jako
$\overline{\mathbb R}$-wartościowa.

\section{Funkcje rzeczywiste i rozszerzone rzeczywiste}

\begin{definicja}[Mierzalność funkcji rozszerzonej]
\label{def:extended-measurable}
Funkcję $f:X\to\overline{\mathbb R}$ nazywamy
\emph{mierzalną}, jeśli jest
$(\mathcal F,\mathcal B(\overline{\mathbb R}))$-mierzalna,
czyli $f^{-1}(B)\in\mathcal F$ dla każdego
$B\in\mathcal B(\overline{\mathbb R})$.
\end{definicja}

\begin{twierdzenie}[Testy mierzalności na półprostych]
\label{thm:extended-tests}
Niech $(X,\mathcal F)$ będzie przestrzenią mierzalną,
a $f:X\to\overline{\mathbb R}$. Następujące warunki są równoważne:
\begin{enumerate}[label=\textup{(\roman*)}]
\item $f$ jest mierzalna;
\item $\{f<a\}\in\mathcal F$ dla każdego $a\in\mathbb R$;
\item $\{f\le a\}\in\mathcal F$ dla każdego $a\in\mathbb R$;
\item $\{f>a\}\in\mathcal F$ dla każdego $a\in\mathbb R$;
\item $\{f\ge a\}\in\mathcal F$ dla każdego $a\in\mathbb R$.
\end{enumerate}
Wystarczy w każdym z warunków (ii)--(v) badać liczby wymierne.
\end{twierdzenie}
\begin{proof}
Zbiory postaci $[-\infty,a)$, $[-\infty,a]$,
$(a,\infty]$ oraz $[a,\infty]$ dla $a\in\mathbb Q$
generują tę samą $\sigma$-algebrę
$\mathcal B(\overline{\mathbb R})$. Istotnie, każda
z tych rodzin jest rozdzielona od pozostałych za pomocą
dopełnień i przeliczalnych granic. Na przykład
$[-\infty,a]=\bigcap_{n\ge1}[-\infty,a+1/n)$ oraz
$(a,\infty]=\overline{\mathbb R}\setminus[-\infty,a]$.
Każdy otwarty przedział w rozszerzonej prostej
jest przecięciem dwóch odpowiednich półprostych,
a każdy zbiór otwarty jest przeliczalną sumą takich
przedziałów, ponieważ $\overline{\mathbb R}$ ma bazę
przeliczalną. Teraz stosujemy twierdzenie
\ref{thm:measurable-generator}.
\end{proof}

\begin{uwaga}[Wartości nieskończone]
Dla mierzalnej funkcji $f:X\to\overline{\mathbb R}$
zbiory $\{f=+\infty\}=\bigcap_{n\ge1}\{f>n\}$ i
$\{f=-\infty\}=\bigcap_{n\ge1}\{f<-n\}$ są mierzalne.
Przy działaniach arytmetycznych musimy unikać
niezdefiniowanych wyrażeń $+\infty-\infty$,
$\infty/\infty$ lub $0\cdot\infty$, chyba że wcześniej
jawnie określimy przyjętą konwencję.
\end{uwaga}

\begin{przyklad}[Funkcja charakterystyczna zbioru]
\label{ex:indicator-measurable}
Dla zbioru $A\subseteq X$ definiujemy
$\chi_A(x)=1$ dla $x\in A$ oraz
$\chi_A(x)=0$ dla $x\notin A$.
Funkcja $\chi_A:X\to\mathbb R$ jest mierzalna
wtedy i tylko wtedy, gdy $A\in\mathcal F$.
Kierunek konieczny otrzymujemy z
$A=\chi_A^{-1}((1/2,\infty))$, a dostateczność
z faktu, że każdy przeciwobraz zbioru borelowskiego
przez $\chi_A$ jest jednym z czterech zbiorów:
$\varnothing$, $A$, $X\setminus A$, $X$.
\end{przyklad}

\section{Działania na funkcjach mierzalnych}

\begin{twierdzenie}[Działania algebraiczne]
\label{thm:measurable-algebra}
Niech $(X,\mathcal F)$ będzie przestrzenią mierzalną,
a $f,g:X\to\mathbb R$ funkcjami mierzalnymi.
Wtedy funkcje $f+g$, $f-g$, $fg$, $|f|$,
$\max(f,g)$ i $\min(f,g)$ są mierzalne.
Zbiór $\{g\ne0\}$ należy do $\mathcal F$, a
funkcja $f/g:\{g\ne0\}\to\mathbb R$ jest mierzalna
względem śladu $\sigma$-algebry. Jeżeli $g(x)\ne0$
dla każdego $x\in X$, to $f/g$ jest mierzalna na całym $X$.
\end{twierdzenie}
\begin{proof}
Odwzorowanie $(f,g):X\to\mathbb R^2$ jest mierzalne
na mocy wniosku \ref{cor:coordinate-measurable}.
Dodawanie $S(u,v)=u+v$, odejmowanie $D(u,v)=u-v$,
mnożenie $M(u,v)=uv$, maksimum i minimum są ciągłymi
funkcjami na $\mathbb R^2$. Moduł jest funkcją ciągłą
na $\mathbb R$. Ich złożenia z funkcją $(f,g)$ są
więc mierzalne na mocy twierdzenia
\ref{thm:composition-measurable}.
Zbiór $\{g\ne0\}=g^{-1}(\mathbb R\setminus\{0\})$
jest mierzalny; ograniczenie $(f,g)$ do niego jest
mierzalne. Funkcja $(u,v)\mapsto u/v$ jest ciągła
na otwartym zbiorze $\{(u,v):v\ne0\}$, co kończy dowód.
\end{proof}

\begin{twierdzenie}[Operacje na funkcjach rozszerzonych]
\label{thm:extended-operations}
Niech $(X,\mathcal F)$ będzie przestrzenią mierzalną,
a $f,g:X\to\overline{\mathbb R}$ funkcjami mierzalnymi.
Wtedy $-f$, $|f|$, $\max(f,g)$ oraz $\min(f,g)$
są mierzalne. Jeśli $f,g$ są nieujemne, ich suma
$f+g$ z naturalną regułą $a+\infty=\infty$
jest mierzalna. Jeżeli nie zachodzi w żadnym punkcie
jednocześnie $f=+\infty$, $g=-\infty$ ani
$f=-\infty$, $g=+\infty$, to również $f+g$ jest mierzalna.
\end{twierdzenie}
\begin{proof}
Negacja i moduł są ciągłe w topologii rozszerzonej
prostej. Funkcje maksimum i minimum są ciągłe na
$\overline{\mathbb R}^2$ z topologią produktową:
na przykład $\max(u,v)<a$ jest równoważne
$u<a$ i $v<a$, zaś $\max(u,v)>a$ zachodzi,
gdy $u>a$ lub $v>a$. Te testy wystarczają do
mierzalności z twierdzenia \ref{thm:extended-tests}.
Dla nieujemnych $f,g$ i dowolnego $a\in\mathbb R$
zachodzi
\[
\{f+g>a\}=\bigcup_{q\in\mathbb Q}
\bigl(\{f>q\}\cap\{g>a-q\}\bigr).
\]
Jeżeli $f(x)+g(x)>a$, można wybrać wymierne $q$
ściśle pomiędzy $a-g(x)$ i $f(x)$; gdy którakolwiek
wartość wynosi $+\infty$, odpowiedni przedział
wymiernych nadal jest niepusty. Implikacja odwrotna
wynika przez dodanie nierówności. Prawa strona
jest mierzalna, a test na półprostych kończy dowód.

W ogólnym przypadku rozdzielamy dziedzinę na mierzalne
zbiory
\[
\begin{aligned}
D_0&=\{|f|<\infty\}\cap\{|g|<\infty\},\\
D_+&=\{f=+\infty,\ g>-\infty\}
      \cup\{g=+\infty,\ f>-\infty\},\\
D_-&=\{f=-\infty,\ g<+\infty\}
      \cup\{g=-\infty,\ f<+\infty\}.
\end{aligned}
\]
W zapisach $|f|<\infty$ i $|g|<\infty$ chodzi o
skończoność wartości; dla punktów z $D_+$ lub $D_-$
suma jest odpowiednio równa $+\infty$ lub $-\infty$.
Założenie wykluczające przeciwne nieskończoności
sprawia, że $D_0,D_+,D_-$ tworzą rozłączny podział
$X$. Na $D_0$ sumę otrzymujemy przez złożenie
mierzalnego odwzorowania $(f,g)|_{D_0}$ z ciągłym
dodawaniem na $\mathbb R^2$. Na dwóch pozostałych
zbiorach funkcja jest stała. Przeciwobrazy
zbiorów borelowskich są zatem mierzalne, bez
odwoływania się do twierdzenia o granicach,
które dopiero zostanie udowodnione.
\end{proof}

\begin{uwaga}[Niezdefiniowane działania]
Dla $f(x)=+\infty$ i $g(x)=-\infty$ formalny zapis
$f(x)+g(x)$ nie określa liczby w $\overline{\mathbb R}$.
Nie można więc pytać o mierzalność takiej \emph{funkcji},
zanim nie ustali się jej wartości na zbiorze punktów
problematycznych. Ustalenie wartości stałej na tym
mierzalnym zbiorze daje jedną z możliwych mierzalnych
wersji rozszerzenia działania.
\end{uwaga}

\section{Granice punktowe funkcji mierzalnych}

\begin{twierdzenie}[Kresy i granice]
\label{thm:pointwise-limits}
Niech $(X,\mathcal F)$ będzie przestrzenią mierzalną,
a $(f_n)_{n\ge1}$ ciągiem mierzalnych funkcji
$f_n:X\to\overline{\mathbb R}$. Wówczas funkcje
\[
\sup_{n\ge1}f_n,\qquad \inf_{n\ge1}f_n,\qquad
\limsup_{n\to\infty}f_n,\qquad
\liminf_{n\to\infty}f_n
\]
są mierzalne. Zbiór
$D=\{x:\lim_{n\to\infty}f_n(x)\text{ istnieje w }
\overline{\mathbb R}\}$ należy do $\mathcal F$;
granica $f(x)=\lim_nf_n(x)$ określona na $D$ jest
mierzalna względem $\mathcal F|_D$.
\end{twierdzenie}
\begin{proof}
Dla dowolnego $a\in\mathbb R$ mamy
$\{\sup_nf_n\le a\}=\bigcap_n\{f_n\le a\}\in\mathcal F$,
a także $\{\inf_nf_n<a\}=\bigcup_n\{f_n<a\}$.
To dowodzi mierzalności supremum i infimum.
Ponieważ
$\limsup_nf_n=\inf_{m\ge1}\sup_{n\ge m}f_n$ i
$\liminf_nf_n=\sup_{m\ge1}\inf_{n\ge m}f_n$,
obie funkcje są mierzalne.
W rozszerzonej prostej granica istnieje wtedy
i tylko wtedy, gdy $\liminf_nf_n=\limsup_nf_n$.
Zbiór równości dwóch mierzalnych funkcji rozszerzonych
jest mierzalny: można na przykład zastosować
ciągłe, ściśle rosnące przekształcenie
$\theta:\overline{\mathbb R}\to[-1,1]$,
$\theta(t)=2\arctan(t)/\pi$ dla $t\in\mathbb R$,
$\theta(\pm\infty)=\pm1$, i następnie badać zbiór,
na którym $|\theta(\limsup_nf_n)-
\theta(\liminf_nf_n)|=0$.
Na $D$ granica jest ograniczeniem mierzalnej funkcji
$\limsup_nf_n$, więc jest mierzalna względem śladu.
\end{proof}

\begin{wniosek}[Monotoniczne ciągi funkcji]\label{idx:06:wniosek:06}
Niech $(X,\mathcal F)$ będzie przestrzenią mierzalną.
Jeżeli $(f_n)_{n\ge1}$ jest niemalejącym ciągiem
mierzalnych funkcji $f_n:X\to[0,\infty]$,
to $f=\sup_nf_n$
jest mierzalna i $f_n(x)\uparrow f(x)$ dla każdego $x$.
Analogicznie, jeżeli $(g_n)_{n\ge1}$ jest malejącym ciągiem
mierzalnych funkcji $g_n:X\to\overline{\mathbb R}$,
to funkcja $\inf_ng_n$ jest mierzalna oraz
$g_n(x)\downarrow\inf_ng_n(x)$ dla każdego $x\in X$.
\end{wniosek}

\begin{przyklad}[Granica punktowa funkcji ciągłych]
Funkcje $f_n(x)=x^n$ na $[0,1]$ są ciągłe,
a więc borelowsko mierzalne. Zbiegają punktowo do
$\chi_{\{1\}}$, która jest mierzalna, choć
nie jest ciągła w punkcie $1$ (dla topologii podprzestrzeni).
Granice punktowe zachowują mierzalność, lecz nie muszą
zachowywać ciągłości.
\end{przyklad}

\section{Funkcje proste}

\begin{definicja}[Funkcja prosta]
\label{def:simple-function}
Funkcję $s:X\to\mathbb R$ nazywamy
\emph{funkcją prostą mierzalną}, jeśli przyjmuje
skończenie wiele wartości i jest mierzalna.
Równoważnie istnieją $m\in\mathbb N$,
liczby $a_1,\ldots,a_m\in\mathbb R$ i zbiory
$A_1,\ldots,A_m\in\mathcal F$ tworzące podział $X$
takie, że
\[
s=\sum_{j=1}^m a_j\chi_{A_j}.
\]
Jeśli wszystkie $a_j\ge0$, funkcję prostą nazywamy
\emph{nieujemną}.
\end{definicja}

\begin{proof}[Dowód równoważności w definicji]
Jeżeli funkcja $s$ przyjmuje skończenie wiele różnych
wartości $a_1,\ldots,a_m$, to jej włókna
$A_j=\{x:s(x)=a_j\}$ są mierzalne i tworzą
podział $X$. Na $A_j$ funkcja ma wartość $a_j$,
co daje podaną reprezentację. Odwrotnie, podana suma
jest mierzalna jako suma skończona mierzalnych funkcji
charakterystycznych i ma zbiór wartości zawarty w
$\{a_1,\ldots,a_m\}$.
\end{proof}

\begin{uwaga}[Po co wymagamy podziału?]
Reprezentacja jako skończona kombinacja funkcji
charakterystycznych nie musi być jednoznaczna, jeśli
zbiory się nakładają. Zawsze można ją zastąpić
reprezentacją kanoniczną, w której $A_j$ są rozłączne
i odpowiadają dokładnie różnym wartościom $s$.
Wówczas wzór na całkę funkcji prostej będzie
niezależny od sposobu zapisania tej funkcji.
\end{uwaga}

\begin{twierdzenie}[Monotoniczna aproksymacja funkcji nieujemnej]
\label{thm:simple-up}
Niech $(X,\mathcal F)$ będzie przestrzenią mierzalną,
a $f:X\to[0,\infty]$ funkcją mierzalną. Dla
$n\ge1$ określmy
\begin{equation}\label{eq:canonical-simple}
s_n(x)=2^{-n}\Big\lfloor 2^n\min\{f(x),n\}\Big\rfloor,
\qquad x\in X.
\end{equation}
Wtedy każda $s_n$ jest nieujemną funkcją prostą mierzalną,
$0\le s_n\le s_{n+1}\le f$ oraz
$s_n(x)\to f(x)$ dla każdego $x\in X$,
przy czym granica może wynosić $+\infty$.
Dla $f(x)<\infty$ i $n>f(x)$ zachodzi
$0\le f(x)-s_n(x)<2^{-n}$.
\end{twierdzenie}
\begin{proof}
Wartość $\min(f,n)$ należy do $[0,n]$.
Po pomnożeniu przez $2^n$ i zastosowaniu części
całkowitej otrzymujemy jedną z liczb całkowitych
$0,1,\ldots,n2^n$. Funkcja $s_n$ przyjmuje zatem
skończenie wiele wartości. Dla $k=0,\ldots,n2^n-1$
zbiór $\{s_n=k2^{-n}\}$ jest opisany nierównościami
$k2^{-n}\le\min(f,n)<(k+1)2^{-n}$;
zbiór odpowiadający ostatniej wartości jest
$\{f\ge n\}$. Wszystkie te zbiory są mierzalne,
więc $s_n$ jest funkcją prostą mierzalną.

Wprowadźmy $u_n(x)=\min(f(x),n)$. Mamy
$u_n\le u_{n+1}$; funkcja
$t\mapsto 2^{-n}\lfloor2^nt\rfloor$ jest niemalejąca,
a ponadto dla każdego $t\ge0$
$2^{-(n+1)}\lfloor2^{n+1}t\rfloor\ge
2^{-n}\lfloor2^nt\rfloor$.
Ostatnia nierówność wynika z tego, że
$\lfloor2z\rfloor\ge2\lfloor z\rfloor$.
Wynika stąd $s_n\le s_{n+1}$. Nierówności
$0\le s_n\le u_n\le f$ wynikają z własności
części całkowitej. Gdy $f(x)$ jest skończona,
przy $n>f(x)$ mamy $u_n(x)=f(x)$ i
$0\le f(x)-s_n(x)<2^{-n}$.
Jeżeli $f(x)=\infty$, to $u_n(x)=n$ i
$s_n(x)=n\to\infty$. To kończy dowód.
\end{proof}

\begin{przyklad}[Aproksymacja z podziałem na poziomy]
Niech $X=[0,1]$, $f(x)=x$ i $n=2$.
Wtedy $s_2(x)=\frac14\lfloor4x\rfloor$ ma wartości
$0,1/4,1/2,3/4,1$ na kolejnych przedziałach
$[0,1/4)$, $[1/4,1/2)$, $[1/2,3/4)$,
$[3/4,1)$ i singletonie $\{1\}$.
Dalsze przybliżenia są coraz dokładniejszymi
schodkami leżącymi poniżej wykresu $f$.
\end{przyklad}

\begin{center}
\begin{tikzpicture}[x=9cm,y=4.0cm]
  \draw[->] (-0.015,0) -- (1.08,0) node[right] {$x$};
  \draw[->] (0,-0.025) -- (0,1.1) node[above] {$y$};
  \draw[blue!70!black,thick] (0,0)--(1,1) node[above left] {$f(x)=x$};
  \draw[red!75!black,thick] (0,0)--(.25,0) (.25,.25)--(.5,.25) (.5,.5)--(.75,.5) (.75,.75)--(1,.75);
  \foreach \x/\y in {.25/0,.5/.25,.75/.5,1/.75} {
    \draw[red!75!black,fill=white] (\x,\y) circle[radius=1.2pt];
  }
  \foreach \x/\y in {.25/.25,.5/.5,.75/.75,1/1} {
    \fill[red!75!black] (\x,\y) circle[radius=1.2pt];
  }
  \draw[dashed,gray!60] (1,0)--(1,1);
  \node[anchor=north] at (0,0) {$0$};
  \node[anchor=north] at (1,0) {$1$};
  \node[red!75!black,anchor=west] at (.42,.13) {$s_2$};
\end{tikzpicture}

\smallskip
\emph{Aproksymacja funkcji $f(x)=x$ na $[0,1]$ przez $s_2$.
Wartości w punktach skoku definiuje wzór \eqref{eq:canonical-simple}.}
\end{center}

\begin{twierdzenie}[Aproksymacja funkcji o wartościach rzeczywistych]
\label{thm:simple-real-approx}
Niech $(X,\mathcal F)$ będzie przestrzenią mierzalną,
a $f:X\to\mathbb R$ mierzalną funkcją skończoną.
Wówczas istnieje ciąg funkcji prostych mierzalnych
$(r_n)$, dla którego $r_n(x)\to f(x)$ dla każdego $x\in X$.
Jeśli $|f(x)|\le M$ dla wszystkich $x\in X$ i pewnego
$M<\infty$, można wybrać $r_n$ tak, by
$\|r_n-f\|_\infty\le2^{-n}$, gdzie
$\|h\|_\infty:=\sup_{x\in X}|h(x)|$.
\end{twierdzenie}
\begin{proof}
Niech $f^+=\max(f,0)$ i $f^-=\max(-f,0)$.
Obie funkcje są mierzalne i nieujemne.
Na mocy twierdzenia \ref{thm:simple-up} istnieją
niemalejące ciągi nieujemnych funkcji prostych
$s_n^+\uparrow f^+$ i $s_n^-\uparrow f^-$.
Funkcje $r_n=s_n^+-s_n^-$ są proste i mierzalne:
przyjmują skończenie wiele wartości i są różnicami
funkcji mierzalnych skończonych. Ponadto
$r_n(x)\to f^+(x)-f^-(x)=f(x)$ dla każdego $x$.
Jeśli $f$ jest ograniczona przez $M$, dla
$n>M$ oba błędy aproksymacji są mniejsze niż
$2^{-n}$, a ponieważ w każdym punkcie co najwyżej
jedna z liczb $f^+(x),f^-(x)$ jest dodatnia,
$|r_n(x)-f(x)|<2^{-n}$. Dla początkowych indeksów
można wybrać dowolne funkcje proste, gdyż nie
wpływa to na istnienie późniejszych przybliżeń;
aby uzyskać podaną nierówność dla wszystkich $n$,
bezpośrednio stosujemy obcięcie i kwantyzację
$r_n(x)=2^{-n}\lfloor2^nf(x)\rfloor$.
Przy ograniczonym $f$ funkcja ta jest prosta i
mierzalna, a dla każdego $x$ zachodzi
$0\le f(x)-r_n(x)<2^{-n}$. Dowodzi to ostatniej tezy.
\end{proof}

\begin{wniosek}[Charakteryzacja mierzalności przez funkcje proste]
\label{cor:meas-simple-characterization}
Niech $(X,\mathcal F)$ będzie przestrzenią mierzalną,
a $f:X\to\mathbb R$ funkcją. Wówczas $f$ jest
mierzalna wtedy i tylko wtedy, gdy istnieje
ciąg mierzalnych funkcji prostych $r_n:X\to\mathbb R$
zbieżny punktowo do $f$.
\end{wniosek}
\begin{proof}
Konieczność daje twierdzenie \ref{thm:simple-real-approx}.
Dostateczność wynika z twierdzenia
\ref{thm:pointwise-limits}, ponieważ punktowa granica
mierzalnych funkcji rzeczywistych jest mierzalna.
\end{proof}

\section{Mierzalność funkcji zespolonych i wektorowych}

\begin{definicja}[Funkcja zespolona mierzalna]\label{idx:06:definicja:11}

Funkcję $f:X\to\mathbb C$ nazywamy mierzalną,
jeśli jest $(\mathcal F,\mathcal B(\mathbb C))$-mierzalna,
gdzie $\mathbb C$ utożsamiamy z $\mathbb R^2$
przez $z\mapsto(\operatorname{Re}z,\operatorname{Im}z)$.
\end{definicja}

\begin{twierdzenie}[Test na częściach rzeczywistej i urojonej]
\label{thm:complex-test}
Niech $(X,\mathcal F)$ będzie przestrzenią mierzalną,
a $f:X\to\mathbb C$. Funkcja $f$ jest mierzalna
wtedy i tylko wtedy, gdy funkcje
$\operatorname{Re}f$ i $\operatorname{Im}f$ są
mierzalne jako funkcje rzeczywiste. W takim przypadku
$\overline f$ i $|f|$ także są mierzalne.
\end{twierdzenie}
\begin{proof}
Pierwsza część jest szczególnym przypadkiem
wniosku \ref{cor:coordinate-measurable} dla $d=2$.
Sprzężenie zespolone i moduł są ciągłe na $\mathbb C$,
a złożenie z funkcją mierzalną pozostaje mierzalne.
\end{proof}

\begin{przyklad}[Funkcje parametryczne]
Jeśli $u,v:X\to\mathbb R$ są mierzalne, to
$f(x)=u(x)+iv(x)$ jest mierzalna.
Nie potrzebujemy do tego żadnej dodatkowej miary
na $X$. Gdy zaś $u,v$ są funkcjami mierzalnymi
względem $\mathcal L_d$, ale niekoniecznie
borelowskimi, również $f$ jest mierzalna względem
$\mathcal L_d$.
\end{przyklad}

\section{Własności prawie wszędzie i uzupełnianie miar}

\begin{definicja}[Równość i zbieżność prawie wszędzie]
\label{def:ae-functions}
Niech $(X,\mathcal F,\mu)$ będzie przestrzenią z miarą.
Funkcje mierzalne $f,g:X\to\mathbb R$ są równe
\emph{$\mu$-prawie wszędzie}, jeśli istnieje
$N\in\mathcal F$ o $\mu(N)=0$ takie, że
$f(x)=g(x)$ dla każdego $x\in X\setminus N$.
Ciąg funkcji mierzalnych $f_n:X\to\mathbb R$
\emph{zbiega $\mu$-prawie wszędzie} do $f:X\to\mathbb R$,
co zapisujemy $f_n\to f$ $\mu$-p.w., jeśli istnieje
taki zbiór $N$, że $f_n(x)\to f(x)$ dla każdego
$x\in X\setminus N$.
\end{definicja}

\begin{uwaga}[Ostrożnie z niezupełną miarą]
Jeśli $\mu$ nie jest zupełna, podzbiór zbioru
mierzalnego miary zero nie musi być mierzalny.
Dlatego definicję wyraziliśmy przez istnienie
\emph{mierzalnego} zbioru zerowego zawierającego
zbiór punktów wyjątkowych. Dla funkcji mierzalnych
sam zbiór $\{f\ne g\}$ jest mierzalny, więc można
równoważnie wymagać $\mu(\{f\ne g\})=0$.
Podobnie dla mierzalnych $f_n,f$ zbiór braku
zbieżności jest mierzalny na mocy twierdzenia
\ref{thm:pointwise-limits}.
\end{uwaga}

\begin{twierdzenie}[Modyfikacje na zbiorach zerowych]
\label{thm:null-modifications}
Niech $(X,\mathcal F,\mu)$ będzie \emph{zupełną}
przestrzenią z miarą. Jeżeli $f:X\to\mathbb R$
jest mierzalna, $g:X\to\mathbb R$ dowolna,
a $\{x:f(x)\ne g(x)\}$ jest zawarty w pewnym
$N\in\mathcal F$ o $\mu(N)=0$, to $g$ jest mierzalna.
Bez założenia zupełności wniosek może być fałszywy.
\end{twierdzenie}
\begin{proof}
W przestrzeni zupełnej każdy podzbiór $N$ należy
do $\mathcal F$. Dla zbioru borelowskiego
$B\subseteq\mathbb R$ zbiory
$g^{-1}(B)\cap(X\setminus N)=f^{-1}(B)\cap(X\setminus N)$
i $g^{-1}(B)\cap N\subseteq N$ są mierzalne.
Ich suma jest równa $g^{-1}(B)$, zatem $g$ jest mierzalna.
Dla kontrprzykładu bierzemy nieuzupełnioną borelowską
miarę Lebesgue'a na $\mathbb R$ oraz podzbiór
$A$ zbioru Cantora, który nie jest borelowski
(wniosek \ref{cor:nonborel-lebesgue}).
Funkcja $f\equiv0$ jest mierzalna, ale
$g=\chi_A$ nie jest borelowsko mierzalna.
Jednocześnie $f=g$ poza borelowskim zbiorem
Cantora, którego miara wynosi zero.
\end{proof}

\begin{wniosek}[Wersje mierzalne funkcji określonych p.w.]\label{idx:06:wniosek:15}

Niech $(X,\mathcal F,\mu)$ będzie przestrzenią z miarą,
a $f_n:X\to\mathbb R$ mierzalnymi funkcjami.
Jeżeli ciąg $(f_n(x))$ jest zbieżny dla
$\mu$-prawie wszystkich $x$, to istnieje
mierzalna funkcja $f:X\to\mathbb R$, taka że
$f_n\to f$ $\mu$-p.w.
\end{wniosek}
\begin{proof}
Z twierdzenia \ref{thm:pointwise-limits} zbiór $D$
punktów, w których granica jest rzeczywista i
skończona, jest mierzalny, a granica na nim jest mierzalna.
Z założenia $X\setminus D$ ma miarę zero.
Definiujemy $f(x)=\lim_nf_n(x)$ na $D$ oraz
$f(x)=0$ na $X\setminus D$. Funkcja jest mierzalna,
gdyż dla borelowskiego $B$ jej przeciwobraz jest
sumą mierzalnego przeciwobrazu na $D$ i,
ewentualnie, całego $X\setminus D$.
\end{proof}

\begin{uwaga}[Zbieżność prawie wszędzie nie zapewnia mierzalności dowolnej wersji granicy]
\label{rem:noncomplete-ae-limit-6}
Jeśli nie wymagamy mierzalności funkcji granicznej,
z założenia $f_n\to f$ poza zbiorem zerowym nie można
wnioskować, że $f$ jest mierzalna. Na przestrzeni
$(\mathbb R,\mathcal B(\mathbb R),\lambda_1|_{\mathcal B})$
wybierzmy nieborelowski podzbiór $A$ zbioru Cantora.
Funkcje $f_n\equiv0$ są borelowskie i zbiegają do
$f=\chi_A$ poza borelowskim zbiorem Cantora miary zero,
lecz $f$ nie jest mierzalna względem $\mathcal B(\mathbb R)$.
Wniosek o mierzalnej wersji granicy gwarantuje istnienie
\emph{pewnego} mierzalnego przedstawiciela; nie oznacza,
że każdy wybór wartości na zbiorze wyjątkowym jest
mierzalny. W przestrzeni zupełnej takie modyfikacje są
mierzalne na mocy twierdzenia~\ref{thm:null-modifications}.
\end{uwaga}

\section{Zbieżność według miary: narzędzia potrzebne do całkowania}
Minimalna definicja i twierdzenie o podciągu pojawiają się tutaj,
ponieważ wykorzystamy je już w rozdziale 8; systematyczne porównanie
wszystkich rodzajów zbieżności pozostawiamy rozdziałowi 9.


\begin{definicja}[Zbieżność według miary]
\label{def:convergence-measure}
Niech $(X,\mathcal F,\mu)$ będzie przestrzenią z miarą,
a $f_n,f:X\to\mathbb R$ mierzalnymi funkcjami.
Mówimy, że $f_n$ \emph{zbiega do $f$ według miary $\mu$},
i piszemy $f_n\xrightarrow{\mu}f$, jeśli
\[
\forall\varepsilon>0\qquad
\mu\bigl(\{x\in X:|f_n(x)-f(x)|>\varepsilon\}\bigr)
\xrightarrow[n\to\infty]{}0.
\]
Definicja nie wymaga $\mu(X)<\infty$ ani
$\sigma$-skończoności.
\end{definicja}

\begin{twierdzenie}[Prawie wszędzie implikuje zbieżność według miary przy mierze skończonej]
\label{thm:ae-to-measure-finite}
Jeśli $(X,\mathcal F,\mu)$ jest przestrzenią
z $\mu(X)<\infty$, a $f_n,f:X\to\mathbb R$ są
mierzalne i $f_n\to f$ $\mu$-prawie wszędzie,
to $f_n\xrightarrow{\mu}f$.
\end{twierdzenie}
\begin{proof}
Ustalmy $\varepsilon>0$ i niech
$E_n=\bigcup_{k\ge n}\{|f_k-f|>\varepsilon\}$.
Wtedy $E_n\in\mathcal F$, $E_{n+1}\subseteq E_n$,
a $\bigcap_nE_n$ składa się z punktów, gdzie
nierówność $|f_k-f|>\varepsilon$ zachodzi
nieskończenie wiele razy. Zbieżność prawie wszędzie
gwarantuje $\mu(\bigcap_nE_n)=0$.
Ponieważ $\mu(E_1)\le\mu(X)<\infty$,
ciągłość miary z góry (twierdzenie
\ref{thm:measure-continuity}) daje $\mu(E_n)\to0$.
Wreszcie $\{|f_n-f|>\varepsilon\}\subseteq E_n$,
więc z monotoniczności $\mu(\{|f_n-f|>\varepsilon\})\to0$.
\end{proof}

\begin{przyklad}[Dlaczego skończoność miary jest ważna]
\label{ex:aenotmeasure}
Na $(\mathbb R,\mathcal L_1,\lambda_1)$ niech
$f_n=\chi_{[n,n+1]}$. Dla każdego ustalonego
$x\in\mathbb R$ mamy $f_n(x)=0$ dla wszystkich
dostatecznie dużych $n$, a więc $f_n\to0$ wszędzie.
Jednak dla $\varepsilon=1/2$ zachodzi
$\lambda_1(\{|f_n|>1/2\})=1$ dla wszystkich $n$.
Nie ma więc zbieżności według miary na całej $\mathbb R$.
\end{przyklad}

\begin{twierdzenie}[Podciąg zbieżny prawie wszędzie]
\label{thm:measure-subsequence}
Niech $(X,\mathcal F,\mu)$ będzie dowolną przestrzenią
z miarą, a $f_n,f:X\to\mathbb R$ mierzalnymi funkcjami.
Jeśli $f_n\xrightarrow{\mu}f$, to istnieje ściśle
rosnący ciąg indeksów $(n_k)_{k\ge1}$ taki, że
$f_{n_k}\to f$ $\mu$-prawie wszędzie.
\end{twierdzenie}
\begin{proof}
Z definicji zbieżności według miary wybieramy
indukcyjnie $n_k>n_{k-1}$ tak duże, aby
\[
\mu(A_k)<2^{-k},\qquad
A_k=\{|f_{n_k}-f|>2^{-k}\}.
\]
Wtedy $\sum_{k\ge1}\mu(A_k)\le1<\infty$.
Pierwsza część lematu Borela--Cantellego
(twierdzenie \ref{thm:borel-cantelli-1})
daje $\mu(\limsup_k A_k)=0$.
Poza tym zbiorem punkt $x$ należy do zaledwie
skończenie wielu $A_k$; dla dostatecznie dużych $k$
mamy więc $|f_{n_k}(x)-f(x)|\le2^{-k}$.
Prawa strona zbiega do zera, a zatem
$f_{n_k}(x)\to f(x)$.
\end{proof}

\begin{uwaga}[Dlaczego nie potrzebujemy skończoności miary?]
Wybór podciągu opiera się wyłącznie na sumowalności
$\sum_k\mu(A_k)$, a lemat Borela--Cantellego nie
zakłada $\mu(X)<\infty$. Jest to zasadniczo inny
argument niż ciągłość miary z góry zastosowana w
twierdzeniu~\ref{thm:ae-to-measure-finite}.
\end{uwaga}

\begin{wniosek}[Jednoznaczność granicy według miary]
\label{cor:measure-limit-unique}
Niech $(X,\mathcal F,\mu)$ będzie przestrzenią z miarą,
a $(f_n)_{n\ge1}$ ciągiem mierzalnych funkcji
$f_n:X\to\mathbb R$. Jeżeli mierzalne funkcje
$f,g:X\to\mathbb R$ spełniają $f_n\xrightarrow{\mu}f$
oraz $f_n\xrightarrow{\mu}g$, to $f=g$ $\mu$-prawie wszędzie.
\end{wniosek}
\begin{proof}
Dla $\varepsilon>0$ z nierówności trójkąta
$\{|f-g|>\varepsilon\}\subseteq
\{|f-f_n|>\varepsilon/2\}\cup
\{|f_n-g|>\varepsilon/2\}$.
Z podaddytywności i przejścia do granicy
otrzymujemy $\mu(\{|f-g|>\varepsilon\})=0$.
Ponieważ $\{f\ne g\}=\bigcup_{m\ge1}\{|f-g|>1/m\}$,
zbiór ten ma miarę zero.
\end{proof}

\begin{przyklad}[Zbieżność według miary bez zbieżności p.w.]
\label{ex:typewriter}
Na $X=[0,1)$ z miarą Lebesgue'a, dla każdego $k\ge0$
i $j=0,\ldots,2^k-1$, kładziemy
$f_{2^k+j}=\chi_{[j2^{-k},(j+1)2^{-k})}$.
Każdy indeks $n\ge1$ występuje w tym wzorze dokładnie
raz (blok indeksów $2^k,\ldots,2^{k+1}-1$).
Dla $n=2^k+j$ mamy
$\lambda_1(\{|f_n|>1/2\})=2^{-k}\to0$,
więc $f_n\to0$ według miary.
Z drugiej strony dla każdego $x\in[0,1)$ i każdego
$k$ dokładnie jeden z przedziałów poziomu $k$
zawiera $x$. Stąd w każdym bloku istnieje jeden indeks
$n$ z $f_n(x)=1$ oraz (dla $k\ge1$) inny indeks
z $f_n(x)=0$. Ciąg $f_n(x)$ nie ma granicy w żadnym
punkcie $[0,1)$.
\end{przyklad}

\begin{twierdzenie}[Lokalna zbieżność według miary]
\label{thm:local-convergence-measure}
Niech $(X,\mathcal F,\mu)$ będzie przestrzenią z miarą,
a $f_n,f:X\to\mathbb R$ funkcjami mierzalnymi,
przy czym $f_n\to f$ $\mu$-prawie wszędzie.
Wówczas dla każdego $A\in\mathcal F$ z $\mu(A)<\infty$
i każdego $\varepsilon>0$ zachodzi
\[
\mu\bigl(A\cap\{|f_n-f|>\varepsilon\}\bigr)\longrightarrow0.
\]
\end{twierdzenie}
\begin{proof}
Rozpatrujemy przestrzeń
$(A,\mathcal F|_A,\mu|_A)$ o skończonej mierze,
a następnie stosujemy twierdzenie
\ref{thm:ae-to-measure-finite} do funkcji ograniczonych
do $A$.
\end{proof}

\begin{uwaga}[Co jest, a co nie jest implikacją?]
Zbieżność prawie wszędzie implikuje zbieżność
według miary przy $\mu(X)<\infty$, lecz nie musi
implikować jej dla $\mu(X)=\infty$.
Zbieżność według miary zawsze pozwala wybrać podciąg
zbieżny prawie wszędzie, ale nie wymusza zbieżności
prawie wszędzie całego ciągu. Kontrprzykłady powyżej
pokazują niezależność obu problemów. Twierdzenia
Jegorowa i Riesza o podciągach oraz pełne zestawienie
rodzajów zbieżności rozwinie rozdział 9.
\end{uwaga}

\section{Przygotowanie do definicji całki Lebesgue'a}

Skończone kombinacje funkcji charakterystycznych są
pomostem pomiędzy mierzeniem zbiorów a całkowaniem.
Dla nieujemnej funkcji prostej $s=\sum_{j=1}^m
a_j\chi_{A_j}$, gdzie $A_j$ tworzą mierzalny
podział $X$, naturalny kandydat na całkę to
\[
\sum_{j=1}^m a_j\mu(A_j).
\]
Jeśli $a_j=0$ i $\mu(A_j)=\infty$, przyjmujemy
w tym nieujemnym wyrażeniu konwencję $0\cdot\infty=0$.
W następnym rozdziale dowiedziemy, że liczba ta
nie zależy od reprezentacji $s$, po czym określimy
całkę dowolnej funkcji nieujemnej jako supremum
po wszystkich prostych funkcjach mierzalnych
mniejszych od niej.

W ten sposób twierdzenie \ref{thm:simple-up}
ma fundamentalną rolę: zapewnia, że zbiór takich
przybliżeń jest dostatecznie bogaty, aby odtworzyć
wartości dowolnej funkcji mierzalnej.

% ROZDZIAŁ 7 — do włączenia w main.tex
\chapter{Konstrukcja całki Lebesgue'a}
\label{ch:calkowanie-lebesgue}

\section{Od mierzenia zbiorów do mierzenia wartości funkcji}

W klasycznej sumie Riemanna dzielimy \emph{dziedzinę} na małe
przedziały. Konstrukcja Lebesgue'a rozpoczyna się od podziału
\emph{zbioru wartości}: dla funkcji prostej $s$ mierzymy zbiory,
na których $s$ przyjmuje kolejne wartości, i sumujemy odpowiednie
iloczyny. Pozwala to całkować na dowolnych przestrzeniach z miarą,
bez porządku na dziedzinie i bez pojęcia przedziału.

W całym rozdziale $(X,\mathcal F,\mu)$ jest dowolną przestrzenią
z miarą. Nie zakładamy, że $\mu(X)<\infty$, ani nawet
$\sigma$-skończoności. Jeżeli nie zaznaczamy inaczej,
funkcje są mierzalne względem $\mathcal F$.
Oznaczenia $\overline{\mathbb R}=[-\infty,\infty]$,
$f^+=\max(f,0)$ i $f^-=\max(-f,0)$ mają zwykły sens;
$f=f^+-f^-$ w punktach, gdzie $f$ jest skończona.
Dla nieujemnych wielkości z $[0,\infty]$ umawiamy się,
że $0\cdot\infty=0$, również przy wyznaczaniu całki
nieujemnej funkcji prostej. Nie dopuszczamy odejmowania
$\infty-\infty$.

\section{Całka nieujemnej funkcji prostej}

\begin{definicja}[Całka funkcji prostej]
\label{def:integral-simple}
Niech $s:X\to[0,\infty)$ będzie mierzalną funkcją prostą.
Jeżeli $a_1,\ldots,a_m$ są wszystkimi różnymi wartościami $s$,
a $A_j=\{x\in X:s(x)=a_j\}$, określamy
\begin{equation}\label{eq:simple-integral}
\int_X s(x)\,\mathrm d\mu(x)
:=\sum_{j=1}^{m}a_j\mu(A_j)\in[0,\infty].
\end{equation}
Zbiory $A_j$ są mierzalne i tworzą rozłączny podział $X$.
Przy $a_j=0$ składnik jest równy zeru, niezależnie od
$\mu(A_j)$. Wartość całki może być nieskończona.
\end{definicja}

\begin{lemat}[Niezależność od reprezentacji]
\label{lem:simple-refinement}
Jeżeli $s=\sum_{i=1}^p a_i\chi_{A_i}$, gdzie
$a_i\ge0$, a zbiory $A_i\in\mathcal F$ tworzą
rozłączny podział $X$ (dopuszczamy powtarzające się
współczynniki), to
$\int_Xs\,\mathrm d\mu(x)=\sum_{i=1}^pa_i\mu(A_i)$.
Wartość ta nie zależy od wybranego podziału ani od
powtarzania współczynników.
\end{lemat}
\begin{proof}
Niech $a_1',\ldots,a_r'$ będą różnymi wartościami
funkcji $s$, a $C_j=\{s=a_j'\}$. Dla każdego $j$
zbiór $C_j$ jest rozłączną sumą tych $A_i$, dla których
$a_i=a_j'$. Zatem, na mocy skończonej addytywności miary,
$a_j'\mu(C_j)=\sum_{i:a_i=a_j'}a_i\mu(A_i)$.
Gdy $a_j'=0$, obie strony są zerem według przyjętej
konwencji; dla $a_j'>0$ równość wynika z rozdzielności
mnożenia w $[0,\infty]$. Sumując po $j$, otrzymujemy
tezę. Dla dwóch podziałów można również zastosować
wspólne rozdrobnienie $A_i\cap B_j$; dowodzi to
niezależności bez zakładania skończoności miary.
\end{proof}

\begin{twierdzenie}[Rachunek dla funkcji prostych]
\label{thm:simple-properties}
Niech $s,t:X\to[0,\infty)$ będą mierzalnymi funkcjami
prostymi, a $c\in[0,\infty)$. Wtedy:
\begin{enumerate}[label=\textup{(\roman*)}]
\item $\int_X0\,\mathrm d\mu(x)=0$ oraz
$\int_X\chi_A\,\mathrm d\mu(x)=\mu(A)$ dla każdego
$A\in\mathcal F$;
\item jeżeli $s\le t$, to
$\int_Xs\,\mathrm d\mu(x)\le\int_Xt\,\mathrm d\mu(x)$;
\item $\int_X(cs)\,\mathrm d\mu(x)=c\int_Xs\,\mathrm d\mu(x)$;
\item $\int_X(s+t)\,\mathrm d\mu(x)=
\int_Xs\,\mathrm d\mu(x)+\int_Xt\,\mathrm d\mu(x)$.
\end{enumerate}
W punktach (iii) i (iv) wszystkie działania odbywają się
w $[0,\infty]$ z konwencją $0\cdot\infty=0$.
\end{twierdzenie}
\begin{proof}
Pierwsza część wynika wprost z definicji. Niech
$s=\sum_i a_i\chi_{A_i}$ i
$t=\sum_j b_j\chi_{B_j}$ będą reprezentacjami
na rozłącznych podziałach $X$. Zbiory
$D_{ij}=A_i\cap B_j$ tworzą wspólne rozdrobnienie.
Na każdym $D_{ij}$ obie funkcje są stałe, dlatego
\[
\int_Xs\,\mathrm d\mu(x)=\sum_{i,j}a_i\mu(D_{ij}),\qquad
\int_Xt\,\mathrm d\mu(x)=\sum_{i,j}b_j\mu(D_{ij}).
\]
Jeśli $s\le t$, to dla każdego niepustego $D_{ij}$
mamy $a_i\le b_j$, więc każdy składnik pierwszej sumy
jest nie większy od odpowiedniego składnika drugiej.
Skalowanie i addytywność wynikają z równości
$(ca_i)\mu(D_{ij})=c(a_i\mu(D_{ij}))$ oraz
$(a_i+b_j)\mu(D_{ij})=a_i\mu(D_{ij})+b_j\mu(D_{ij})$.
Żadna operacja nie wymaga odejmowania nieskończoności.
\end{proof}

\begin{przyklad}[Mierzalna funkcja prosta na przestrzeni nieskończonej]
Na $(\mathbb R,\mathcal L_1,\lambda_1)$ funkcja
$s=3\chi_{[0,2]}+5\chi_{[4,5]}$ ma całkę
$3\cdot2+5\cdot1=11$. Natomiast
$t=\chi_{[0,\infty)}$ ma całkę $+\infty$.
Obie funkcje są proste i mierzalne, chociaż druga
nie jest całkowalna w sensie skończonej całki.
\end{przyklad}

\section{Definicja całki dowolnej funkcji nieujemnej}

\begin{definicja}[Całka Lebesgue'a funkcji nieujemnej]
\label{def:integral-positive}
Niech $f:X\to[0,\infty]$ będzie mierzalna. Definiujemy
\begin{equation}\label{eq:positive-integral}
\int_X f(x)\,\mathrm d\mu(x)
:=\sup\left\{\int_X s(x)\,\mathrm d\mu(x):
\begin{array}{l}
 s:X\to[0,\infty)\text{ jest mierzalna i prosta},\\[-2pt]
 0\le s(x)\le f(x)\text{ dla wszystkich }x\in X
\end{array}\right\}.
\end{equation}
Supremum należy do $[0,\infty]$. Rodzina funkcji
prostych w nawiasie jest niepusta, ponieważ zawiera zero.
\end{definicja}

\begin{uwaga}[Co dokładnie znaczy ta konstrukcja?]
Mierzalność funkcji $f$ nie wymusza skończoności całki.
Nie wybieramy też żadnego konkretnego przybliżenia:
rozpatrujemy \emph{wszystkie} nieujemne funkcje proste
nieprzekraczające $f$. Twierdzenie
\ref{thm:simple-up} zapewnia, że istnieją również
jawne, punktowo niemalejące przybliżenia $s_n\uparrow f$.
Dopiero twierdzenie Beppo Leviego w rozdziale
\ref{ch:twierdzenia-graniczne} wykaże, że całkę można
obliczać jako granicę całek takiej sekwencji.
\end{uwaga}

\begin{twierdzenie}[Pierwsze własności całki nieujemnej]
\label{thm:nonnegative-first-properties}
Niech $f,g:X\to[0,\infty]$ będą mierzalne,
$c\in[0,\infty)$ i $A\in\mathcal F$. Wtedy:
\begin{enumerate}[label=\textup{(\roman*)}]
\item gdy $f\le g$, zachodzi
$\int_X f\,\mathrm d\mu(x)\le\int_Xg\,\mathrm d\mu(x)$;
\item $\int_X cf\,\mathrm d\mu(x)=c\int_Xf\,\mathrm d\mu(x)$;
\item dla funkcji prostej $f$ definicje
\ref{def:integral-simple} i \ref{def:integral-positive}
są zgodne;
\item $\int_X f\chi_A\,\mathrm d\mu(x)$ zależy
wyłącznie od wartości $f$ na $A$ i dla
$A\subseteq B$, $A,B\in\mathcal F$, zachodzi
$\int_X f\chi_A\,\mathrm d\mu(x)\le
\int_Xf\chi_B\,\mathrm d\mu(x)$;
\item jeśli $\mu(A)=0$, to
$\int_X f\chi_A\,\mathrm d\mu(x)=0$,
nawet gdy $f(x)=+\infty$ w punktach $A$.
\end{enumerate}
\end{twierdzenie}
\begin{proof}
W (i) każda funkcja prosta $s\le f$ spełnia
również $s\le g$, więc supremum dla $f$
nie przekracza supremum dla $g$.
Dla $c>0$ podstawienie $t=cs$ daje bijekcję
między nieujemnymi prostymi minorantami $f$ i $cf$;
na mocy twierdzenia \ref{thm:simple-properties}
całki skalują się przez $c$. Przy $c=0$ obie
strony wynoszą zero według umowy $0\cdot\infty=0$.

Jeżeli $f$ jest prosta, sama należy do rodziny
minorant w \eqref{eq:positive-integral}. Dla każdego
$s\le f$ monotoniczność całki prostej daje
$\int_Xs\,\mathrm d\mu(x)\le\int_Xf\,\mathrm d\mu(x)$;
stąd zgodność obu definicji.
Część (iv) wynika z punktowej równości funkcji
$f\chi_A$ przy ustalonych wartościach na $A$
oraz z (i). Aby wykazać (v), ustalmy dowolną
nieujemną funkcję prostą $s\le f\chi_A$.
Poza $A$ mamy $s=0$; skoro $\mu(A)=0$,
każdy jej dodatni poziom ma miarę zero, więc
$\int_Xs\,\mathrm d\mu(x)=0$.
Biorąc supremum, otrzymujemy zero.
\end{proof}

\begin{definicja}[Całka po podzbiorze]
\label{def:integral-over-set}
Dla $A\in\mathcal F$ i mierzalnej $f:X\to[0,\infty]$
określamy
$\int_A f(x)\,\mathrm d\mu(x):=
\int_X f(x)\chi_A(x)\,\mathrm d\mu(x)$,
gdzie iloczyn $f\chi_A$ rozumiemy jako funkcję
równą $f$ na $A$ i zero na $X\setminus A$.
To ostatnie określenie usuwa dwuznaczność zapisu
$0\cdot\infty$. Równoważnie całkujemy funkcję
$f|_A$ względem ograniczenia miary $\mu|_A$
na $(A,\mathcal F|_A)$.
\end{definicja}

\begin{lemat}[Zgodność z ograniczeniem miary]
\label{lem:restriction-integral}
Dla dowolnych $A\in\mathcal F$ i mierzalnej
$f:X\to[0,\infty]$ zachodzi
\[
\int_A f(x)\,\mathrm d\mu(x)
=\int_A f|_A(x)\,\mathrm d(\mu|_A)(x).
\]
\end{lemat}
\begin{proof}
Dla funkcji prostej zachodzi to z definicji i
$\mu|_A(A\cap B)=\mu(A\cap B)$.
Jeżeli $t:A\to[0,\infty)$ jest prostą mierzalną
minorantą $f|_A$, to przedłużenie jej zerem poza $A$
jest prostą mierzalną funkcją na $X$, mniejszą od
$f\chi_A$. Odwrotnie, ograniczenie każdej
prostej minoranty $f\chi_A$ do $A$ jest
prostą minorantą $f|_A$. Równość całek prostych
i supremów daje tezę.
\end{proof}

\begin{twierdzenie}[Nierówność Markowa dla całki]
\label{thm:markov-nonnegative}
Niech $f:X\to[0,\infty]$ będzie mierzalna i $a>0$.
Wówczas
\begin{equation}\label{eq:markov}
\mu(\{x\in X:f(x)\ge a\})
\le \frac1a\int_X f(x)\,\mathrm d\mu(x).
\end{equation}
Jeżeli $\int_Xf\,\mathrm d\mu(x)<\infty$, to $f$ jest
skończona $\mu$-prawie wszędzie. Ponadto
$\int_Xf\,\mathrm d\mu(x)=0$ wtedy i tylko wtedy,
gdy $f=0$ $\mu$-prawie wszędzie.
\end{twierdzenie}
\begin{proof}
Dla $E_a=\{f\ge a\}$ prosta funkcja
$a\chi_{E_a}$ nie przekracza $f$. Zatem
$a\mu(E_a)=\int_X a\chi_{E_a}\,\mathrm d\mu(x)
\le\int_Xf\,\mathrm d\mu(x)$, co daje
\eqref{eq:markov}. Jeśli ostatnia całka jest skończona,
to $\mu(\{f=\infty\})\le\mu(\{f\ge n\})
\le n^{-1}\int_Xf\,\mathrm d\mu(x)\to0$.
Jeżeli całka wynosi zero, to
$\mu(\{f\ge 1/n\})=0$ dla wszystkich $n$, a zbiór
$\{f>0\}=\bigcup_{n\ge1}\{f\ge1/n\}$ jest zerowy.
Odwrotnie, jeśli $f=0$ poza mierzalnym zbiorem zerowym,
to część (v) twierdzenia
\ref{thm:nonnegative-first-properties} daje $\int_X f=0$.
\end{proof}

\begin{przyklad}[Zero prawie wszędzie, mimo dużych wartości]
Na $\mathbb R$ z miarą Lebesgue'a funkcja
$f=+\infty\cdot\chi_{\mathbb Q}$, rozumiana jako
$f(x)=+\infty$ dla $x\in\mathbb Q$ i $0$ dla
$x\notin\mathbb Q$, jest mierzalna i ma całkę zero.
To pokazuje, dlaczego nie wystarczy traktować
całki jako zwykłego sumowania wartości punktowych.
\end{przyklad}

\section{Funkcje ze znakiem i całkowalność}

\begin{twierdzenie}[Niezmienniczość całki przy zmianie na zbiorze zerowym]
\label{thm:ae-integral-invariance}
Niech $f,g:X\to[0,\infty]$ będą mierzalne i
$f=g$ $\mu$-prawie wszędzie. Wówczas
$\int_Xf\,\mathrm d\mu(x)=\int_Xg\,\mathrm d\mu(x)$.
Jeśli $f,g:X\to\overline{\mathbb R}$ mają określone
całki ze znakiem i są równe prawie wszędzie, ich
całki również są równe.
\end{twierdzenie}
\begin{proof}
Wybierzmy $N\in\mathcal F$ o $\mu(N)=0$ takie,
że $f=g$ poza $N$. Dla dowolnej prostej minoranty
$0\le s\le f$ funkcja $s\chi_{X\setminus N}$
jest prostą minorantą $g$, a
\[
\int_Xs\,\mathrm d\mu(x)
=\int_Xs\chi_{X\setminus N}\,\mathrm d\mu(x),
\]
ponieważ składniki odpowiadające częściom poziomic
zawartym w $N$ mają miarę zero. Stąd
$\int_X f\le\int_X g$ po przejściu do supremum.
Zamiana ról $f,g$ daje nierówność odwrotną.
W przypadku funkcji ze znakiem stosujemy część
pierwszą osobno do ich części dodatnich i ujemnych.
Równość wartości wyklucza możliwość, by tylko jedna
z odpowiadających całek była nieokreślona.
\end{proof}

\begin{definicja}[Części dodatnia i ujemna]
\label{def:positive-negative-parts}
Dla mierzalnej $f:X\to\overline{\mathbb R}$ definiujemy
$f^+(x)=\max\{f(x),0\}$ i
$f^-(x)=\max\{-f(x),0\}$. Są to mierzalne funkcje
$[0,\infty]$-wartościowe, a $|f|=f^++f^-$.
Jeśli co najmniej jedna z liczb
$\int_Xf^+\,\mathrm d\mu(x)$,
$\int_Xf^-\,\mathrm d\mu(x)$ jest skończona, definiujemy
\begin{equation}\label{eq:signed-integral}
\int_X f(x)\,\mathrm d\mu(x)
:=\int_Xf^+(x)\,\mathrm d\mu(x)
-\int_Xf^-(x)\,\mathrm d\mu(x).
\end{equation}
Jeśli obie liczby są nieskończone, całka funkcji $f$
\emph{nie jest określona}. Funkcję $f$ nazywamy
\emph{(absolutnie) całkowalną}, jeżeli
$\int_X|f|\,\mathrm d\mu(x)<\infty$.
\end{definicja}

\begin{uwaga}[Dlaczego nie używamy odejmowania nieskończoności?]
Niech $X=\mathbb R$ z miarą Lebesgue'a i
$f(x)=x$. Wtedy zarówno całka $f^+$, jak i $f^-$
wynosi $\infty$, więc całka $f$ na całej $\mathbb R$
nie jest określona. Symetria i formalne równanie
"$\infty-\infty=0$" nie zastępują definicji.
Wartość główna Cauchy'ego, jeśli jest rozważana,
stanowi odrębne pojęcie wymagające osobnej definicji.
\end{uwaga}

\begin{twierdzenie}[Kryterium absolutnej całkowalności]
\label{thm:integrability-criterion}
Dla mierzalnej $f:X\to\overline{\mathbb R}$ następujące
warunki są równoważne:
\begin{enumerate}[label=\textup{(\roman*)}]
\item $\int_X|f|\,\mathrm d\mu(x)<\infty$;
\item $\int_Xf^+\,\mathrm d\mu(x)<\infty$ oraz
$\int_Xf^-\,\mathrm d\mu(x)<\infty$.
\end{enumerate}
W takim przypadku $f$ jest skończona prawie wszędzie,
całka $\int_Xf\,\mathrm d\mu(x)$ jest rzeczywista i
$|\int_Xf\,\mathrm d\mu(x)|\le\int_X|f|\,\mathrm d\mu(x)$.
\end{twierdzenie}
\begin{proof}
Ponieważ $0\le f^\pm\le|f|$, monotoniczność całki
daje (i)$\Rightarrow$(ii). Dla implikacji odwrotnej
nie możemy jeszcze korzystać bez dowodu z addytywności
całki dla dowolnych funkcji nieujemnych. Do pełnej addytywności wrócimy w rozdziale
\ref{ch:twierdzenia-graniczne}. Tu wykorzystamy
wyłącznie rozdzielność nośników części dodatniej i ujemnej.
Bezpośredni argument jest następujący: jeśli
$s\le f^++f^-$ jest nieujemna i prosta, to
dzielimy dziedzinę na mierzalne zbiory
$P=\{f\ge0\}$ i $M=\{f<0\}$. Funkcje
$s\chi_P\le f^+$ i $s\chi_M\le f^-$
są proste, a
$\int_Xs\,\mathrm d\mu(x)=
\int_Xs\chi_P\,\mathrm d\mu(x)+
\int_Xs\chi_M\,\mathrm d\mu(x)$
na mocy twierdzenia \ref{thm:simple-properties}.
Przechodząc do supremum po $s\le|f|=f^++f^-$,
otrzymujemy
$\int_X|f|\,\mathrm d\mu(x)\le
\int_Xf^+\,\mathrm d\mu(x)+
\int_Xf^-\,\mathrm d\mu(x)<\infty$.
Z twierdzenia \ref{thm:markov-nonnegative} otrzymujemy
skończoność $|f|$ prawie wszędzie. Wreszcie, dla dowolnych prostych minorant
$0\le u\le f^+$ i $0\le v\le f^-$ zachodzi
$u+v\le|f|$. Stąd
$\int_Xu\,\mathrm d\mu(x)+\int_Xv\,\mathrm d\mu(x)
=\int_X(u+v)\,\mathrm d\mu(x)
\le\int_X|f|\,\mathrm d\mu(x)$.
Przechodząc niezależnie do supremum po $u$ i $v$,
uzyskujemy także nierówność
$\int_Xf^+\,\mathrm d\mu(x)+\int_Xf^-\,\mathrm d\mu(x)
\le\int_X|f|\,\mathrm d\mu(x)$.
Z obydwu nierówności wynika równość tych liczb.
W szczególności $|\int_Xf\,\mathrm d\mu(x)|
\le\int_Xf^+\,\mathrm d\mu(x)+
\int_Xf^-\,\mathrm d\mu(x)=\int_X|f|\,\mathrm d\mu(x)$.
\end{proof}

\begin{definicja}[Całka funkcji całkowalnej po zbiorze]
\label{def:signed-integral-over-set-7}
Niech $A\in\mathcal F$, a $f:X\to\mathbb R$ będzie
mierzalna oraz całkowalna, tj.
$\int_X|f(x)|\,\mathrm d\mu(x)<\infty$.
Definiujemy
\[
\int_A f(x)\,\mathrm d\mu(x)
:=\int_X\widetilde f_A(x)\,\mathrm d\mu(x),
\qquad
\widetilde f_A(x)=
\begin{cases}f(x),&x\in A,\\0,&x\notin A.\end{cases}
\]
Funkcja $\widetilde f_A$ jest mierzalna oraz
$|\widetilde f_A|\le|f|$, więc całka po prawej
jest skończona. Zgodnie z definicją całki ze znakiem
również
$\int_A f=\int_A f^+-\int_A f^-$;
obie całki są skończone. Definicja nie wymaga
żadnej konwencji dla iloczynu $0\cdot\infty$.
\end{definicja}

\begin{uwaga}[Szczegół przy wartości $+\infty$]
W dowodzie wprowadziliśmy rozbicie
$P=\{f\ge0\}$ i $M=\{f<0\}$.
Dla $f(x)=+\infty$ obowiązuje $x\in P$,
a dla $f(x)=-\infty$ obowiązuje $x\in M$.
Nie wykonujemy w tych punktach odejmowania
$\infty-\infty$: obie części dodatnia i ujemna
nie są jednocześnie dodatnie w żadnym punkcie.
\end{uwaga}

\begin{przyklad}[Miary dyskretne i sumy szeregów]
Niech $X=\mathbb N$, $\mathcal F=\mathcal P(\mathbb N)$
i $\mu(A)=\sum_{n\in A}w_n$, gdzie $w_n\ge0$.
Dla funkcji $f:\mathbb N\to[0,\infty]$ całka
jest równa zwykłej sumie nieujemnego szeregu
$\int_{\mathbb N} f(n)\,\mathrm d\mu(n)=
\sum_{n=1}^{\infty}f(n)w_n$.
Istotnie, suma po pierwszych $N$ punktach jest
całką prostej funkcji nieprzekraczającej $f$,
więc prawa strona nie przekracza całki.
Z drugiej strony, dla dowolnej prostej $s\le f$
przeliczalna addytywność miary pozwala zapisać
$\int_X s\,\mathrm d\mu(x)=
\sum_{n=1}^{\infty}s(n)w_n\le\sum_{n=1}^{\infty}f(n)w_n$.
Dla rzeczywistej $f$ warunek absolutnej całkowalności
jest równoważny
$\sum_n|f(n)|w_n<\infty$.
W szczególności miara licząca odpowiada $w_n=1$.
\end{przyklad}

\begin{przyklad}[Całka może istnieć jako $+\infty$]
Na $(0,1]$ z miarą Lebesgue'a połóżmy $f(x)=1/x$.
Dla $j=0,\ldots,n-1$ na przedziale
$I_j=(2^{-(j+1)},2^{-j}]$ mamy $f(x)\ge2^j$.
Funkcja prosta
$t_n=\sum_{j=0}^{n-1}2^j\chi_{I_j}$
spełnia $0\le t_n\le f$, a jej całka wynosi
\[
\int_{(0,1]}t_n(x)\,\mathrm d\lambda_1(x)
=\sum_{j=0}^{n-1}2^j\lambda_1(I_j)
=\sum_{j=0}^{n-1}2^j2^{-(j+1)}=n/2.
\]
Z definicji całki nieujemnej
$\int_{(0,1]}f(x)\,\mathrm d\lambda_1(x)\ge n/2$
dla każdego $n$, a więc całka jest równa $+\infty$.
Nie trzeba do tego rachunku twierdzenia Beppo Leviego.
\end{przyklad}

\begin{uwaga}[Istnienie całki ze znakiem nie jest tym samym co całkowalność]
Funkcja nieujemna $f(x)=1/x$ na $(0,1]$ ma
określoną całkę równą $+\infty$, ale nie jest
całkowalna w sensie $L^1$. Z kolei dla $f(x)=x$
na $\mathbb R$ obie całki części dodatniej i ujemnej
są nieskończone, a więc całka ze znakiem nie jest
nawet określona. Nie należy utożsamiać trzech sytuacji:
całki skończonej, całki nieskończonej i wyrażenia
nieokreślonego.
\end{uwaga}

\section{Wzajemny stosunek całki i mierzalności}

\begin{przyklad}[Funkcja Dirichleta]
Na przedziale $[0,1]$ funkcja
$f=\chi_{\mathbb Q\cap[0,1]}$ jest borelowsko
mierzalna i nieujemna. Ponieważ zbiór liczb
wymiernych jest przeliczalny, ma miarę Lebesgue'a
zero, a więc $\int_{[0,1]} f\,\mathrm d\lambda_1(x)=0$.
Jednocześnie funkcja jest nieciągła w każdym punkcie,
więc nie jest całkowalna w sensie Riemanna.
\end{przyklad}

\begin{przyklad}[Całkowalność na przestrzeni o nieskończonej mierze]
Na $\mathbb R$ z miarą Lebesgue'a funkcja
$f(x)=1/(1+x^2)$ jest nieujemna i całkowalna, choć
$\lambda_1(\mathbb R)=\infty$. Istotnie, na $[-1,1]$
mamy $f\le1$. Dla $n\ge1$ i $x\in[n,n+1)$ mamy
$f(x)\le n^{-2}$; analogicznie na ujemnej półosi.
Dla dowolnej funkcji prostej $0\le s\le f$,
przeliczalna addytywność miary na poziomicach $s$
pozwala rozłożyć jej całkę na rozłączne przedziały
$[-1,1]$, $[n,n+1)$ i $(-(n+1),-n]$, $n\ge1$.
Na każdym przedziale całka $s$ jest nie większa
od iloczynu jego długości i odpowiednio $1$ albo
$n^{-2}$, więc
$\int_{\mathbb R}s\,\mathrm d\lambda_1(x)
\le2+2\sum_{n\ge1}n^{-2}<\infty$.
Ta nierówność jest niezależna od wyboru $s$;
biorąc supremum w definicji całki, otrzymujemy
$\int_{\mathbb R} f(x)\,\mathrm d\lambda_1(x)<\infty$. Dokładną wartość $\pi$ wyprowadzimy
po omówieniu twierdzenia o zamianie zmiennych.
\end{przyklad}

\begin{uwaga}[Jeszcze nie wszystko zostało dowiedzione]
Na tym etapie znamy monotoniczność, jednorodność,
całki funkcji prostych i kryteria całkowalności.
Nie używaliśmy bez dowodu równości
$\int_X(f+g)=\int_X f+\int_X g$ dla \emph{dowolnych}
nieujemnych funkcji mierzalnych. Tę fundamentalną
własność oraz równość $\int_X\lim f_n=\lim\int_X f_n$
w odpowiednich warunkach wyprowadzimy wprost
z definicji w kolejnym rozdziale.
\end{uwaga}

% ROZDZIAŁ 8 — do włączenia w main.tex
\chapter{Twierdzenia graniczne dla całki Lebesgue'a}
\label{ch:twierdzenia-graniczne}

\providecommand{\atlasref}[2]{#2}


\section{Dlaczego nie wolno bezwarunkowo przestawiać granicy i całki?}

Punktem wyjścia jest pytanie, kiedy dla $f_n\to f$ wolno
zamienić kolejność granicy i całkowania:
\[
\int_X f(x)\,\mathrm d\mu(x)
=\lim_{n\to\infty}\int_X f_n(x)\,\mathrm d\mu(x).
\]
Zbieżność punktowa sama nie wystarcza. Na przykład na
$(0,1)$ z miarą Lebesgue'a funkcje
$f_n=n\chi_{(0,1/n)}$ zbiegają do zera w każdym
punkcie, ale ich całki wynoszą $1$.
Twierdzenia tego rozdziału podają precyzyjne warunki,
w których przejście do granicy jest uprawnione.

W całym rozdziale $(X,\mathcal F,\mu)$ jest przestrzenią
z miarą. Każde występujące całkowanie oznacza całkowanie
w sensie rozdziału \ref{ch:calkowanie-lebesgue}.
Nie zakładamy, że $\mu$ jest skończona ani zupełna,
jeśli nie zaznaczamy tego wyraźnie. Symbol $f_n\uparrow f$
oznacza niemalejącą zbieżność punktową do $f$.

\section{Fundamentalny lemat aproksymacyjny}

\begin{lemat}[Zbieżność całek prostej funkcji na zbiorach rosnących]
\label{lem:simple-exhaustion}
Niech $s:X\to[0,\infty)$ będzie mierzalną funkcją
prostą, a $(E_n)_{n\ge1}\subseteq\mathcal F$
niemalejącym ciągiem zbiorów. Jeśli
$E=\bigcup_{n\ge1}E_n$, to
\[
\lim_{n\to\infty}\int_{E_n}s(x)\,\mathrm d\mu(x)
=\int_Es(x)\,\mathrm d\mu(x).
\]
Równość zachodzi również wtedy, gdy prawa strona
wynosi $+\infty$.
\end{lemat}
\begin{proof}
Niech $s=\sum_{j=1}^m a_j\chi_{A_j}$
będzie reprezentacją na mierzalnym podziale, gdzie
$a_j\ge0$. Z definicji całki po zbiorze
$\int_{E_n}s\,\mathrm d\mu(x)=
\sum_{j=1}^m a_j\mu(A_j\cap E_n)$.
Zbiory $A_j\cap E_n$ rosną do $A_j\cap E$, więc
ciągłość miary z dołu (twierdzenie
\ref{thm:measure-continuity}) daje
$\mu(A_j\cap E_n)\uparrow\mu(A_j\cap E)$.
Ponieważ suma zawiera skończenie wiele nieujemnych
składników, możemy przejść do granicy w każdym z nich;
przy współczynniku $a_j=0$ obie strony są zerem.
Otrzymujemy wzór dla $E$.
\end{proof}

\section{Twierdzenie Beppo Leviego}

\begin{twierdzenie}[O zbieżności monotonicznej; Beppo Levi]
\label{thm:mct}
Niech $(f_n)_{n\ge1}$ będzie ciągiem mierzalnych funkcji
$f_n:X\to[0,\infty]$, takim że dla każdego $x\in X$
$f_n(x)\le f_{n+1}(x)$. Niech
$f(x)=\lim_{n\to\infty}f_n(x)=\sup_nf_n(x)$.
Wtedy $f$ jest mierzalna i
\begin{equation}\label{eq:mct}
\int_Xf(x)\,\mathrm d\mu(x)
=\lim_{n\to\infty}\int_Xf_n(x)\,\mathrm d\mu(x)
=\sup_{n\ge1}\int_Xf_n(x)\,\mathrm d\mu(x).
\end{equation}
Dopuszczamy wartość $+\infty$ i nie wymagamy
żadnego założenia o $\mu(X)$.
\end{twierdzenie}
\begin{proof}
Mierzalność $f$ wykazano w twierdzeniu
\ref{thm:pointwise-limits}. Z monotoniczności całki
(rozdział 7) mamy
$L:=\sup_n\int_Xf_n\,\mathrm d\mu(x)\le
\int_Xf\,\mathrm d\mu(x)$.
Trzeba wykazać trudniejszą nierówność odwrotną.

Ustalmy dowolną nieujemną mierzalną funkcję prostą
$s\le f$ oraz liczbę $\theta\in(0,1)$.
Definiujemy
$E_n=\{x\in X:f_n(x)\ge\theta s(x)\}$.
Zbiory $E_n$ są mierzalne i rosną.
Ponadto $\bigcup_nE_n=X$: jeżeli $s(x)=0$,
to $x\in E_n$ dla każdego $n$; jeżeli $s(x)>0$,
to $f(x)\ge s(x)>\theta s(x)$, więc z
$f_n(x)\uparrow f(x)$ wynika $f_n(x)\ge\theta s(x)$
dla dostatecznie dużego $n$. Argument działa także
gdy $f(x)=+\infty$, ponieważ $s(x)$ jest skończona.

Z nierówności $f_n\ge\theta s\chi_{E_n}$,
monotoniczności i jednorodności całki wynika
\[
\int_Xf_n\,\mathrm d\mu(x)\ge
\theta\int_{E_n}s(x)\,\mathrm d\mu(x).
\]
Przechodząc do granicy, stosując lemat
\ref{lem:simple-exhaustion} i $E_n\uparrow X$,
dostajemy $L\ge\theta\int_Xs\,\mathrm d\mu(x)$.
Niech $\theta\uparrow1$; dla całki skończonej
wynika $L\ge\int_Xs\,\mathrm d\mu(x)$, a jeśli
$\int_Xs=\infty$, to z powyższej nierówności
już dla jednego $\theta>0$ otrzymujemy $L=\infty$.
Na koniec bierzemy supremum po wszystkich
nieujemnych prostych $s\le f$. Na mocy definicji
\ref{def:integral-positive} daje to
$L\ge\int_Xf\,\mathrm d\mu(x)$. Obie nierówności
kończą dowód.
\end{proof}

\begin{wniosek}[Konstruktywne obliczanie całki]
\label{cor:simple-approx-integral}
Niech $(X,\mathcal F,\mu)$ będzie przestrzenią z miarą.
Jeżeli $f:X\to[0,\infty]$ jest mierzalna, a
$(s_n)$ jest \emph{dowolnym} ciągiem nieujemnych
funkcji prostych mierzalnych takim, że $s_n\uparrow f$
wszędzie, to
$\int_Xf\,\mathrm d\mu(x)=
\lim_{n\to\infty}\int_Xs_n\,\mathrm d\mu(x)$.
W szczególności można użyć funkcji z
\eqref{eq:canonical-simple}.
\end{wniosek}
\begin{proof}
To bezpośrednie zastosowanie twierdzenia
\ref{thm:mct} do ciągu $s_n$; całki funkcji prostych
znamy już z definicji \ref{def:integral-simple}.
\end{proof}

\begin{przyklad}[Ciąg monotoniczny o nieskończonej granicy całek]
Na $(0,1)$ z miarą Lebesgue'a funkcje
$f_n(x)=\min\{n,1/x\}$ są niemalejące i zbiegają
do $1/x$. Twierdzenie \ref{thm:mct} daje
$\int_{(0,1)}x^{-1}\,\mathrm d\lambda_1(x)=
\lim_n\int_{(0,1)}f_n\,\mathrm d\lambda_1(x)=+\infty$,
co jest zgodne z bezpośrednim rachunkiem z
rozdziału \ref{ch:calkowanie-lebesgue}.
Nie wolno wywnioskować, że same $f_n$ mają
nieograniczone wartości punktowe: każda z nich
jest ograniczona przez $n$.
\end{przyklad}

\begin{twierdzenie}[Zbieżność monotoniczna prawie wszędzie]
\label{thm:mct-ae}
Niech $f_n,f:X\to[0,\infty]$ będą mierzalne.
Jeśli dla każdego $n\ge1$ zachodzi $f_n\le f_{n+1}$ prawie wszędzie oraz $f_n\to f$
$\mu$-prawie wszędzie, to zachodzi równość
\eqref{eq:mct}.
\end{twierdzenie}
\begin{proof}
Dla każdego $n$ zbiór $\{f_n>f_{n+1}\}$ jest mierzalny i ma miarę zero.
Dołączamy do ich przeliczalnej sumy mierzalny zbiór miary zero,
poza którym zachodzi zbieżność do $f$, otrzymując wspólny
zbiór $N\in\mathcal F$ o $\mu(N)=0$.
Definiujemy $g_n=f_n$ oraz $g=f$ poza $N$,
a na $N$ nadajemy wszystkim tym funkcjom wartość zero.
W ten sposób nie trzeba interpretować iloczynu $0\cdot\infty$.
Funkcje $g_n$ są mierzalne, nieujemne i niemalejące wszędzie,
a ich granicą punktową jest $g$.
Stosujemy twierdzenie \ref{thm:mct} do $g_n$.
Równość całek $g_n$ i $f_n$ oraz całek $g$ i $f$
wynika z twierdzenia \ref{thm:ae-integral-invariance}.
Dzięki temu zastosowanie wersji punktowej
Beppo Leviego daje dokładnie żądaną równość.
\end{proof}

\section{Addytywność i liniowość całki}

\begin{twierdzenie}[Addytywność całki nieujemnej]
\label{thm:positive-additivity}
Dla dowolnych mierzalnych $f,g:X\to[0,\infty]$
zachodzi
\begin{equation}\label{eq:positive-additivity}
\int_X(f+g)(x)\,\mathrm d\mu(x)
=\int_Xf(x)\,\mathrm d\mu(x)+
\int_Xg(x)\,\mathrm d\mu(x).
\end{equation}
Obie strony należą do $[0,\infty]$; dopuszczamy
$+\infty$ po którejkolwiek stronie.
\end{twierdzenie}
\begin{proof}
Z twierdzenia \ref{thm:simple-up} wybierzmy
nieujemne funkcje proste $s_n\uparrow f$ i
$t_n\uparrow g$. Wtedy $s_n+t_n\uparrow f+g$.
Twierdzenie \ref{thm:mct}, a następnie addytywność
całki prostej, dają
\[
\int_X(f+g)\,\mathrm d\mu(x)
=\lim_n\int_X(s_n+t_n)\,\mathrm d\mu(x)
=\lim_n\left(\int_Xs_n\,\mathrm d\mu(x)+
\int_Xt_n\,\mathrm d\mu(x)\right).
\]
Ponowne zastosowanie twierdzenia \ref{thm:mct}
do $s_n$ i $t_n$ daje tezę; suma dwóch ciągów
niemalejących liczb nieujemnych ma granicę równą
sumie ich granic także w $[0,\infty]$.
\end{proof}

\begin{wniosek}[Rozkład całki na rozłączne zbiory]
\label{cor:integral-countable-partition}
Niech $f:X\to[0,\infty]$ będzie mierzalna, a
$(A_n)_{n\ge1}$ przeliczalną rodziną parami rozłącznych
zbiorów z $\mathcal F$. Wtedy
\begin{equation}\label{eq:integral-partition}
\int_{\bigcup_{n\ge1}A_n}f(x)\,\mathrm d\mu(x)
=\sum_{n=1}^{\infty}\int_{A_n}f(x)\,\mathrm d\mu(x).
\end{equation}
\end{wniosek}
\begin{proof}
Dla każdego $N$ zachodzi punktowo
$f\chi_{\cup_{n=1}^NA_n}=
\sum_{n=1}^Nf\chi_{A_n}$.
Addytywność całki daje wzór dla sumy skończonej.
Gdy $N\to\infty$, lewe funkcje rosną do
$f\chi_{\cup_{n\ge1}A_n}$. Zastosowanie
Beppo Leviego kończy dowód.
\end{proof}

\begin{twierdzenie}[Liniowość całki funkcji całkowalnych]
\label{thm:integral-linear-real}
Jeżeli $f,g:X\to\mathbb R$ są mierzalne i
$\int_X|f|\,\mathrm d\mu(x)<\infty$,
$\int_X|g|\,\mathrm d\mu(x)<\infty$, to dla
$a,b\in\mathbb R$ funkcja $af+bg$ jest całkowalna i
\[
\int_X(af+bg)(x)\,\mathrm d\mu(x)
=a\int_Xf(x)\,\mathrm d\mu(x)+
b\int_Xg(x)\,\mathrm d\mu(x).
\]
Ponadto $\bigl|\int_Xf\,\mathrm d\mu(x)\bigr|
\le\int_X|f|\,\mathrm d\mu(x)$.
\end{twierdzenie}
\begin{proof}
Dla $a\ge0$ rozkłady $(af)^+=af^+$ i
$(af)^-=af^-$ oraz jednorodność całki nieujemnej
dają $\int_X af=a\int_X f$. Dla $a<0$ korzystamy z
$(-f)^+=f^-$ i $(-f)^-=f^+$.

Z nierówności $|f+g|\le|f|+|g|$ i addytywności
nieujemnej wynika całkowalność sumy.
W każdym punkcie zachodzi tożsamość
\[
(f+g)^++f^-+g^-=(f+g)^-+f^++g^+.
\]
Wszystkie wyrazy mają skończone całki, bo są
ograniczone przez sumy $|f|+|g|$.
Całkujemy tę równość, stosujemy addytywność
nieujemną i przenosimy skończone składniki na
drugą stronę. Otrzymujemy
$\int_X(f+g)=\int_X f+\int_X g$.
Połączenie obu własności daje pełną liniowość.
Na koniec z $0\le f^\pm\le|f|$ i
$|\int_X f^+-\int_X f^-|\le\int_X f^++\int_X f^-
=\int_X|f|$ wynika oszacowanie.
\end{proof}

\begin{wniosek}[Nierówność trójkąta i struktura $L^1$]
\label{cor:integral-triangle}
Jeśli $f,g:X\to\mathbb R$ są całkowalne, to
$\int_X|f+g|\,\mathrm d\mu(x)\le
\int_X|f|\,\mathrm d\mu(x)+
\int_X|g|\,\mathrm d\mu(x)$.
Funkcjonał $f\mapsto\int_Xf\,\mathrm d\mu(x)$
jest liniowy na przestrzeni rzeczywistych
funkcji całkowalnych, a nierówność
$|\int_X f|\le\int_X|f|$ wyznacza jego ograniczenie.
Po utożsamieniu funkcji równych prawie wszędzie
otrzymamy przestrzeń normowaną $L^1$ w rozdziale 14.
\end{wniosek}
\begin{proof}
Całkujemy punktową nierówność $|f+g|\le|f|+|g|$
i stosujemy addytywność całki nieujemnej. Liniowość
i oszacowanie funkcjonału zostały udowodnione
w twierdzeniu \ref{thm:integral-linear-real}.
\end{proof}

\begin{definicja}[Całka funkcji zespolonej]
\label{def:complex-integral}
Funkcję mierzalną $f:X\to\mathbb C$ nazywamy
całkowalną, jeżeli $\int_X|f|\,\mathrm d\mu(x)<\infty$.
Wówczas określamy
$\int_Xf\,\mathrm d\mu(x)=
\int_X\operatorname{Re}f\,\mathrm d\mu(x)+
 i\int_X\operatorname{Im}f\,\mathrm d\mu(x)$.
\end{definicja}

\begin{twierdzenie}[Liniowość zespolona i oszacowanie modułu]
\label{thm:integral-complex-linear}
Jeśli $f,g:X\to\mathbb C$ są całkowalne i
$\alpha,\beta\in\mathbb C$, to $\alpha f+\beta g$
jest całkowalna,
$\int_X(\alpha f+\beta g)\,\mathrm d\mu(x)=
\alpha\int_Xf\,\mathrm d\mu(x)+
\beta\int_Xg\,\mathrm d\mu(x)$, a także
\[
\left|\int_Xf(x)\,\mathrm d\mu(x)\right|
\le\int_X|f(x)|\,\mathrm d\mu(x).
\]
\end{twierdzenie}
\begin{proof}
Z $|\alpha f+\beta g|\le|\alpha||f|+|\beta||g|$
wynika całkowalność. Liniowość otrzymujemy,
rozwijając części rzeczywiste i urojone i stosując
twierdzenie \ref{thm:integral-linear-real}.
Jeśli $I=\int_Xf\,\mathrm d\mu(x)\ne0$, wybierzmy
$\zeta=\overline I/|I|$, więc $|\zeta|=1$ i
$\zeta I=|I|$. Wtedy
$|I|=\int_X\operatorname{Re}(\zeta f)\,\mathrm d\mu(x)
\le\int_X|\operatorname{Re}(\zeta f)|\,\mathrm d\mu(x)
\le\int_X|f|\,\mathrm d\mu(x)$.
Dla $I=0$ nierówność jest oczywista.
\end{proof}

\section{Lemat Fatou}

\begin{twierdzenie}[Lemat Fatou]
\label{thm:fatou}
Niech $f_n:X\to[0,\infty]$ będą funkcjami mierzalnymi.
Wtedy
\begin{equation}\label{eq:fatou}
\int_X\liminf_{n\to\infty}f_n(x)\,\mathrm d\mu(x)
\le\liminf_{n\to\infty}\int_Xf_n(x)\,\mathrm d\mu(x).
\end{equation}
Nie wymagamy monotoniczności ciągu ani skończoności miary.
\end{twierdzenie}
\begin{proof}
Dla $m\ge1$ połóżmy
$h_m(x)=\inf_{n\ge m}f_n(x)$.
Na mocy twierdzenia \ref{thm:pointwise-limits}
każda $h_m$ jest mierzalna. Ponieważ z rosnącym $m$
infimum obejmuje coraz mniej wyrazów, funkcje $h_m$
są niemalejące. Z definicji granicy dolnej
$h_m(x)\uparrow\liminf_nf_n(x)$.
Twierdzenie \ref{thm:mct} daje
\[
\int_X\liminf_nf_n\,\mathrm d\mu(x)
=\lim_{m\to\infty}\int_Xh_m\,\mathrm d\mu(x).
\]
Dla wszystkich $n\ge m$ mamy $h_m\le f_n$, więc
$\int_Xh_m\,\mathrm d\mu(x)\le
\inf_{n\ge m}\int_Xf_n\,\mathrm d\mu(x)$.
Przechodzimy do granicy po $m$ i otrzymujemy
prawą stronę \eqref{eq:fatou}.
\end{proof}

\begin{przyklad}[Nierówność Fatou bywa ostra]
Na $\mathbb R$ z miarą Lebesgue'a niech
$f_n=\chi_{[n,n+1]}$. Dla każdego $x$ zachodzi
$f_n(x)\to0$, ale $\int_{\mathbb R} f_n(x)\,\mathrm d\lambda_1(x)=1$.
Zatem $\int_{\mathbb R}\liminf_nf_n(x)\,\mathrm d\lambda_1(x)=0$
i $\liminf_n\int_{\mathbb R} f_n(x)\,\mathrm d\lambda_1(x)=1$.
Masa może \emph{uciekać do nieskończoności}.
\end{przyklad}

\begin{wniosek}[Odwrotny lemat Fatou przy dominacji]
\label{cor:reverse-fatou}
Niech $f_n:X\to\mathbb R$ będą mierzalne i
$0\le f_n\le g$ $\mu$-prawie wszędzie, gdzie
$g:X\to[0,\infty]$ jest całkowalna. Wtedy
\[
\limsup_n\int_Xf_n\,\mathrm d\mu(x)
\le\int_X\limsup_nf_n\,\mathrm d\mu(x).
\]
\end{wniosek}
\begin{proof}
Ponieważ $g$ jest całkowalna, $g<\infty$ prawie
wszędzie. Usuwamy więc jeden wspólny mierzalny
zbiór zerowy, na którym może wystąpić $g=\infty$
lub naruszenie nierówności $0\le f_n\le g$.
Na tym zbiorze zastępujemy $f_n$ i $g$ zerem;
całki nie zmieniają się. Możemy teraz założyć,
że $g$ jest skończona wszędzie i $0\le f_n\le g$.
Funkcje $g-f_n$ są nieujemne i mierzalne, a
$\liminf_n(g-f_n)=g-\limsup_nf_n$ w punktach,
gdzie $g$ jest skończona. Ponieważ $g<\infty$
prawie wszędzie, lemat Fatou i liniowość całki
funkcji całkowalnych dają
$\int_X g-\int_X\limsup f_n\le
\int_X g-\limsup\int_X f_n$.
Odejmując skończoną liczbę $\int_X g$,
otrzymujemy tezę.
\end{proof}

\begin{twierdzenie}[Lemat Fatou z całkowalnym ograniczeniem od dołu]
\label{thm:fatou-integrable-lower-8}
Niech $f_n:X\to\overline{\mathbb R}$ będą mierzalne,
a $g:X\to[0,\infty]$ całkowalna. Załóżmy,
że $f_n\ge-g$ prawie wszędzie dla każdego $n$.
Wówczas każda całka $\int_Xf_n\,\mathrm d\mu(x)$
jest określona w $(-\infty,\infty]$, podobnie całka
funkcji $\liminf_nf_n$, oraz
\[
\int_X\liminf_{n\to\infty}f_n(x)\,\mathrm d\mu(x)
\le \liminf_{n\to\infty}\int_Xf_n(x)\,\mathrm d\mu(x).
\]
Nie jest wymagana skończoność miary $\mu(X)$.
\end{twierdzenie}
\begin{proof}
Bierzemy jeden wspólny mierzalny zbiór wyjątkowy miary zero,
na którym mogą zawodzić nierówności $f_n\ge-g$
lub skończoność $g$. Zmieniamy tam wszystkie
funkcje na zero; ich całki pozostają niezmienione.
Możemy więc zakładać, że $g<\infty$ wszędzie oraz
$f_n+g\ge0$ dla każdego $n$. Lemat Fatou daje
\[
\int_X(\liminf_nf_n+g)\,\mathrm d\mu(x)
=\int_X\liminf_n(f_n+g)\,\mathrm d\mu(x)
\le\liminf_n\int_X(f_n+g)\,\mathrm d\mu(x).
\]
Aby uzasadnić rachunek całek bez odejmowania
nieskończoności, zauważmy, że $f_n^-\le g$,
a więc $\int_X f_n^-<\infty$. Punktowa tożsamość
$(f_n+g)+f_n^-=f_n^++g$ i addytywność całki
nieujemnej dają
$\int_X(f_n+g)+\int_Xf_n^-=
\int_Xf_n^++\int_Xg$.
Wolno odjąć skończoną całkę $f_n^-$ i otrzymać
$\int_X(f_n+g)=\int_Xf_n+\int_Xg$.
Analogicznie postępujemy z $\liminf_nf_n$.
Odejmujemy tę samą skończoną liczbę $\int_Xg$
i otrzymujemy tezę. Żadne wyrażenie typu
$\infty-\infty$ nie występuje.
\end{proof}

\begin{przyklad}[Brak całkowalnego ograniczenia od dołu]
Na $(0,1)$ połóżmy $f_n=-n\chi_{(0,1/n)}$.
Mamy $f_n(x)\to0$ dla każdego $x\in(0,1)$,
ale $\int_{(0,1)}f_n(x)\,\mathrm d\lambda_1(x)=-1$.
Zatem nierówność z twierdzenia
\ref{thm:fatou-integrable-lower-8} byłaby fałszywa,
gdyby usunąć warunek dominacji od dołu.
\end{przyklad}

\section{Twierdzenie Lebesgue'a o zbieżności zmajoryzowanej}

\begin{twierdzenie}[O zbieżności zmajoryzowanej; Lebesgue]
\label{thm:dct}
Niech $f_n:X\to\mathbb R$ będą mierzalne i
$f_n(x)\to f(x)$ dla $\mu$-prawie wszystkich $x$, gdzie
$f:X\to\mathbb R$ jest mierzalna. Załóżmy, że
istnieje nieujemna całkowalna funkcja mierzalna
$g:X\to[0,\infty]$ taka, że
$|f_n(x)|\le g(x)$ dla wszystkich $n$ i
$\mu$-prawie wszystkich $x$, przy czym zbiór
wyjątkowy można wybrać wspólny dla całego ciągu.
Wówczas $f_n$ i $f$ są całkowalne oraz
\begin{equation}\label{eq:dct}
\lim_{n\to\infty}\int_Xf_n(x)\,\mathrm d\mu(x)
=\int_Xf(x)\,\mathrm d\mu(x).
\end{equation}
\end{twierdzenie}
\begin{proof}
Ponieważ $g$ jest całkowalna, zbiór $\{g=\infty\}$
jest zerowy. Dołączamy go do mierzalnego zbioru zerowego obejmującego
wszystkie wyjątki od zbieżności oraz nierówności dominacji.
Na tak otrzymanym zbiorze możemy zastąpić wszystkie
$f_n$, funkcję $f$ i $g$ zerem. Całki pozostają
niezmienione, a teraz dla każdego $x$ zachodzi
$f_n(x)\to f(x)$, $|f_n(x)|\le g(x)<\infty$.
W szczególności operacje $g\pm f_n$ są wszędzie
określone i nieujemne.
Z ciągłości wartości bezwzględnej wynika $|f|\le g$.
Monotoniczność daje
$\int_X|f_n|\le\int_X g<\infty$ oraz
$\int_X|f|\le\int_X g<\infty$.

Funkcje $g+f_n$ i $g-f_n$ są nieujemne.
Stosując do nich lemat Fatou, otrzymujemy
\[
\int_X(g+f)\,\mathrm d\mu(x)
\le\liminf_n\int_X(g+f_n)\,\mathrm d\mu(x)
=\int_Xg\,\mathrm d\mu(x)+
\liminf_n\int_Xf_n\,\mathrm d\mu(x),
\]
\[
\int_X(g-f)\,\mathrm d\mu(x)
\le\liminf_n\int_X(g-f_n)\,\mathrm d\mu(x)
=\int_Xg\,\mathrm d\mu(x)-
\limsup_n\int_Xf_n\,\mathrm d\mu(x).
\]
W drugim przejściu wszystkie całki $f_n$ są
skończone i jednostajnie ograniczone przez $\int_X g$.
Po skróceniu skończonej liczby $\int_X g$ pierwsza
nierówność daje
$\int_X f\le\liminf_n\int_X f_n$, a druga
$\limsup_n\int_X f_n\le\int_X f$.
Ponieważ $\liminf\le\limsup$, obie granice
są równe $\int_X f$.
\end{proof}

\begin{wniosek}[Zbieżność zmajoryzowana w $L^1$]
\label{cor:dct-l1}
Przy założeniach twierdzenia \ref{thm:dct} zachodzi
\[
\lim_{n\to\infty}\int_X|f_n(x)-f(x)|\,\mathrm d\mu(x)=0.
\]
Dla funkcji zespolonych wniosek i twierdzenie
\ref{thm:dct} także zachodzą, jeśli zastąpić
wartość bezwzględną modułem.
\end{wniosek}
\begin{proof}
Mamy $|f_n-f|\to0$ prawie wszędzie oraz
$|f_n-f|\le2g$ prawie wszędzie.
Stosujemy twierdzenie \ref{thm:dct} do nieujemnych
funkcji $|f_n-f|$. Dla funkcji zespolonych najpierw otrzymujemy $|f|\le g$
prawie wszędzie, więc wszystkie $f_n$ i $f$ są całkowalne.
Stosujemy rzeczywistą wersję twierdzenia \ref{thm:dct}
do nieujemnych funkcji $|f_n-f|$, a z
$\bigl|\int_X(f_n-f)\bigr|\le\int_X|f_n-f|$
wynika zbieżność całek.
\end{proof}

\begin{wniosek}[Twierdzenie o zbieżności ograniczonej]
\label{cor:bct}
Niech $\mu(X)<\infty$ oraz $f_n,f:X\to\mathbb R$
będą mierzalne. Jeśli $f_n\to f$ prawie wszędzie
i istnieje $M\in[0,\infty)$ takie, że $|f_n|\le M$
prawie wszędzie dla wszystkich $n$, to wszystkie $f_n$ oraz $f$ są całkowalne i
$\int_X f_n\,\mathrm d\mu\to\int_X f\,\mathrm d\mu$ oraz $\int_X|f_n-f|\,\mathrm d\mu\to0$.
\end{wniosek}
\begin{proof}
Funkcja stała $g\equiv M$ jest całkowalna,
gdyż $\int_X g\,\mathrm d\mu(x)=M\mu(X)<\infty$.
Stosujemy twierdzenie \ref{thm:dct} i wniosek
\ref{cor:dct-l1}.
\end{proof}

\begin{przyklad}[Dlaczego potrzebna jest dominacja]
Na $(0,1)$ weźmy
$f_n=n\chi_{(0,1/n)}$.
Dla każdego ustalonego $x>0$ zachodzi $f_n(x)\to0$,
lecz $\int_{(0,1)}f_n\,\mathrm d\lambda_1(x)=1$.
Gdyby istniała całkowalna $g\ge f_n$ dla wszystkich
$n$, wówczas twierdzenie o zbieżności zmajoryzowanej
dawałoby sprzeczność. Problem nie leży zatem
w zbieżności punktowej, lecz w braku wspólnej
całkowalnej funkcji dominującej.
\end{przyklad}

\begin{przyklad}[Dlaczego skończoność miary ma znaczenie]
Na $\mathbb R$ funkcje $f_n=\chi_{[n,n+1]}$
są ograniczone wspólną stałą $1$, a mimo to
$\int_{\mathbb R}f_n(x)\,\mathrm d\lambda_1(x)=1\not\to0$, chociaż $f_n\to0$ wszędzie.
Stała $1$ nie jest całkowalna na $\mathbb R$,
więc twierdzenie \ref{cor:bct} nie ma zastosowania.
\end{przyklad}

\section{Szeregi funkcji i całek}

\begin{twierdzenie}[Szeregi nieujemnych funkcji mierzalnych]
\label{thm:sum-nonnegative-integrals}
Niech $f_n:X\to[0,\infty]$ będą mierzalne. Wtedy
\begin{equation}\label{eq:sum-nonnegative-integrals}
\int_X\sum_{n=1}^{\infty}f_n(x)\,\mathrm d\mu(x)
=\sum_{n=1}^{\infty}\int_Xf_n(x)\,\mathrm d\mu(x),
\end{equation}
gdzie obie strony mogą być nieskończone.
\end{twierdzenie}
\begin{proof}
Sumy częściowe $s_N=\sum_{n=1}^N f_n$ są
mierzalne, nieujemne i rosną do
$s=\sum_{n\ge1}f_n$. Z twierdzenia \ref{thm:mct}
$\int_Xs\,\mathrm d\mu(x)=\lim_N\int_Xs_N\,\mathrm d\mu(x)$.
Dla każdego $N$ stosujemy skończoną addytywność
całki nieujemnej z twierdzenia
\ref{thm:positive-additivity}; granica jest
z definicji sumą szeregu nieujemnego.
\end{proof}

\begin{twierdzenie}[Całkowanie szeregów bezwzględnie sumowalnych]
\label{thm:sum-absolutely-integrable}
Niech $f_n:X\to\mathbb C$ będą mierzalne i
całkowalne oraz niech
$\sum_{n\ge1}\int_X|f_n(x)|\,\mathrm d\mu(x)<\infty$.
Wówczas $\sum_nf_n(x)$ jest bezwzględnie zbieżny
$\mu$-prawie wszędzie. Jeśli $F$ oznacza jego sumę
na zbiorze bezwzględnej zbieżności i zero na jego dopełnieniu,
to $F$ jest mierzalna i całkowalna oraz
\begin{equation}\label{eq:interchange-series}
\int_XF(x)\,\mathrm d\mu(x)
=\sum_{n=1}^{\infty}\int_Xf_n(x)\,\mathrm d\mu(x).
\end{equation}
Szereg całek jest bezwzględnie zbieżny, a dla sum
częściowych $S_N=\sum_{n=1}^N f_n$ zachodzi
\[
 \int_X|F-S_N|\,\mathrm d\mu(x)
 \le\sum_{n>N}\int_X|f_n|\,\mathrm d\mu(x)
 \longrightarrow0.
\]
\end{twierdzenie}
\begin{proof}
Niech $G=\sum_{n\ge1}|f_n|$, mierzalna suma
nieujemna. Z twierdzenia
\ref{thm:sum-nonnegative-integrals}
$\int_XG\,\mathrm d\mu(x)=
\sum_n\int_X|f_n|\,\mathrm d\mu(x)<\infty$.
Na mocy twierdzenia \ref{thm:markov-nonnegative}
$G(x)<\infty$ prawie wszędzie. Poza zbiorem
zerowym szereg $\sum_nf_n(x)$ zbiega więc
bezwzględnie. Zbiór bezwzględnej zbieżności $A=\{G<\infty\}$
jest mierzalny i $\mu(X\setminus A)=0$.
Funkcje równe $S_N=\sum_{n=1}^N f_n$ na $A$ i zero na $X\setminus A$
są mierzalne i zbiegają punktowo na całym $X$ do $F$.
Stąd $F$ jest mierzalna. Ponadto $|F|\le G$ prawie wszędzie,
więc $F$ jest całkowalna (także gdy $\mu$ nie jest zupełna,
ponieważ $F$ zdefiniowano na całym $X$ jako funkcję mierzalną).
Dla sum częściowych $S_N=\sum_{n=1}^Nf_n$
mamy $|S_N|\le G$ oraz $S_N\to F$ prawie wszędzie.
Zespolona wersja twierdzenia o zbieżności zmajoryzowanej
(wniosek \ref{cor:dct-l1}) daje $\int_XS_N\to\int_XF$. Liniowość całki
zespolonej pozwala zapisać
$\int_XS_N=\sum_{n=1}^N\int_Xf_n$, co daje równość całek.
Poza mierzalnym zbiorem zerowym
$|F-S_N|\le\sum_{n>N}|f_n|$.
Całkujemy i stosujemy twierdzenie
\ref{thm:sum-nonnegative-integrals}, otrzymując
oszacowanie ogona. Bezwzględna zbieżność szeregu całek
wynika z $|\int_Xf_n\,\mathrm d\mu|\le\int_X|f_n|\,\mathrm d\mu$.
\end{proof}

\begin{twierdzenie}[Wzór warstwowy na całkę nieujemną]
\label{thm:layer-cake-8}
Niech $f:X\to[0,\infty]$ będzie mierzalna.
Wtedy funkcja $t\mapsto\mu(\{x:f(x)>t\})$
jest borelowska na $[0,\infty)$ i zachodzi
\begin{equation}\label{eq:layer-cake-8}
\int_Xf(x)\,\mathrm d\mu(x)
=\int_{[0,\infty)}\mu(\{x:f(x)>t\})\,
\mathrm d\lambda_1(t).
\end{equation}
Obie strony mogą być nieskończone. Dowód nie korzysta
z twierdzenia Tonellego, które pojawi się później.
\end{twierdzenie}
\begin{proof}
Funkcja $t\mapsto\mu(\{f>t\})$ jest \emph{nierosnąca},
a każda funkcja monotoniczna
o wartościach w $[0,\infty]$ jest borelowsko mierzalna.
Najpierw niech $s=\sum_{j=1}^m a_j\chi_{A_j}$
będzie nieujemną funkcją prostą, gdzie zbiory $A_j$
tworzą mierzalny podział $X$. Dla $t\ge0$
\[
\mu(\{s>t\})=\sum_{j=1}^m
\mu(A_j)\chi_{[0,a_j)}(t).
\]
Całkowanie skończonej sumy po $t$ daje
$\sum_ja_j\mu(A_j)=\int_Xs\,\mathrm d\mu(x)$;
przy $a_j=0$ stosujemy umowę $0\cdot\infty=0$.
Dla ogólnej $f$ weźmy $s_n\uparrow f$ z
twierdzenia~\ref{thm:simple-up}. Dla każdego $t\ge0$
mamy $\{s_n>t\}\uparrow\{f>t\}$. Ciągłość miary
z dołu daje
$\mu(\{s_n>t\})\uparrow\mu(\{f>t\})$.
Stosujemy Beppo Leviego na $[0,\infty)$ i na $X$:
\[
\int_{[0,\infty)}\mu(\{f>t\})\,
\mathrm d\lambda_1(t)
=\lim_n\int_Xs_n\,\mathrm d\mu(x)
=\int_Xf\,\mathrm d\mu(x).
\]
\end{proof}

\begin{wniosek}[Uogólniona formuła rozkładu na poziomice]\label{idx:08:wniosek:20}

Dla każdego $p>0$ i mierzalnej funkcji
$f:X\to\overline{\mathbb R}$ (lub $f:X\to\mathbb C$)
zachodzi
\[
\int_X|f(x)|^p\,\mathrm d\mu(x)
=p\int_{[0,\infty)}t^{p-1}
\mu(\{x:|f(x)|>t\})\,\mathrm d\lambda_1(t).
\]
\end{wniosek}
\begin{proof}
Niech najpierw $s=\sum_{j=1}^m a_j\chi_{A_j}$,
gdzie $a_j\ge0$, a $(A_j)$ jest mierzalnym podziałem.
Dla $t\ge0$ zachodzi
$\mu(\{s>t\})=\sum_j\mu(A_j)\chi_{[0,a_j)}(t)$.
Z addytywności całki nieujemnej i elementarnego wzoru
$\int_0^{a}pt^{p-1}\,\mathrm dt=a^p$ dla $a\ge0$
(w sensie całki Riemanna na przedziałach
$[\delta,a]$, a następnie przy $\delta\downarrow0$)
otrzymujemy
\[
p\int_{[0,\infty)}t^{p-1}\mu(\{s>t\})\,
\mathrm d\lambda_1(t)=\sum_j a_j^p\mu(A_j)
=\int_Xs(x)^p\,\mathrm d\mu(x).
\]
Wartość całego integrandu w $t=0$ przyjmujemy równą
zero, niezależnie od $p$ i masy zbioru $\{|f|>0\}$.
Nie zmienia to całki. Dla $t>0$ nie ma nieoznaczonego
iloczynu, ponieważ $t^{p-1}$ jest dodatnie i skończone.
Dla ogólnej funkcji $|f|$ wybieramy $s_n\uparrow|f|$
ze zbioru nieujemnych funkcji prostych.
Wtedy $s_n^p\uparrow|f|^p$ i
$\mu(\{s_n>t\})\uparrow\mu(\{|f|>t\})$.
Dwukrotne zastosowanie twierdzenia Beppo Leviego
daje żądaną równość bez używania twierdzenia Tonellego.
\end{proof}

\section{Przykłady obliczeń i kryteria całkowalności}

\begin{przyklad}[Całka potęgi blisko zera]
\label{ex:power-integral}
Niech $\alpha\in\mathbb R$ i
$f(x)=x^{-\alpha}$ na $(0,1]$.
Jeżeli $\alpha<1$, to funkcja jest całkowalna oraz
\begin{equation}\label{eq:power-integral}
\int_{(0,1]}x^{-\alpha}\,\mathrm d\lambda_1(x)
=\frac1{1-\alpha}.
\end{equation}
Jeśli $\alpha\ge1$, całka wynosi $+\infty$.
\end{przyklad}
\begin{proof}
Dla $\alpha\ge1$ mamy $x^{-\alpha}\ge x^{-1}$
na $(0,1]$, a rozbieżność drugiej całki
wykazaliśmy w rozdziale 7. Dla $\alpha<1$
rozważmy $f_n=f\chi_{[1/n,1]}$.
Wtedy $f_n\uparrow f$ na $(0,1]$.
Funkcja $f_n$ jest ograniczona i ciągła na
$[1/n,1]$, więc jej całka Lebesgue'a zgadza się
z całką Riemanna na tym przedziale (zob.
twierdzenie \ref{thm:riemann-to-lebesgue} poniżej).
Z rachunku elementarnego
$\int_{[1/n,1]}x^{-\alpha}\,\mathrm d\lambda_1(x)
=(1-n^{\alpha-1})/(1-\alpha)$.
Twierdzenie \ref{thm:mct} daje po przejściu do
granicy \eqref{eq:power-integral}.
\end{proof}

\begin{przyklad}[Kryterium logarytmiczne]
Dla $p\in\mathbb R$ rozważmy na $(0,e^{-1}]$
funkcję $f_p(x)=\bigl(x(\log(1/x))^p\bigr)^{-1}$.
Całka z $f_p$ jest skończona dokładnie wtedy, gdy
$p>1$. Rzeczywiście, podstawienie klasycznej
całki Riemanna $u=\log(1/x)$ na
$[e^{-R},e^{-1}]$ (a następnie twierdzenie
\ref{thm:mct}) daje
$\int_{[e^{-R},e^{-1}]}f_p\,\mathrm d\lambda_1(x)
=\int_1^R u^{-p}\,\mathrm d\lambda_1(u)$.
Całka po prawej ma skończoną granicę przy
$R\to\infty$ dokładnie wtedy, gdy $p>1$.
Samo użycie podstawienia w tym przykładzie
można też zastąpić oszacowaniami dyadycznymi;
pełne twierdzenie o zamianie zmiennych
udowodnimy później.
\end{przyklad}

\section{Całki zależne od parametru}

\begin{twierdzenie}[Ciągłość pod znakiem całki]
\label{thm:parameter-continuity}
Niech $T$ będzie przestrzenią metryczną, $t_0\in T$
i $F:T\times X\to\mathbb C$ funkcją taką, że
$x\mapsto F(t,x)$ jest mierzalna dla każdego $t\in T$.
Załóżmy, że istnieje otoczenie $V$ punktu $t_0$
oraz całkowalna $g:X\to[0,\infty]$ takie, że
$|F(t,x)|\le g(x)$ dla wszystkich $t\in V$ i
prawie wszystkich $x$, poza jednym wspólnym zbiorem
zerowym, oraz że $F(t,x)\to F(t_0,x)$, gdy
$t\to t_0$, dla prawie wszystkich $x$.
Wówczas funkcja
$\Phi(t)=\int_XF(t,x)\,\mathrm d\mu(x)$
jest określona dla $t\in V$ i ciągła w $t_0$.
\end{twierdzenie}
\begin{proof}
Dla każdego $t\in V$ mierzalność przekroju i dominacja
przez $g$ zapewniają jego całkowalność. Ustalmy dowolny
ciąg $t_n\to t_0$ w $V$. Funkcje
$F(t_n,\cdot)$ są zdominowane przez $g$ poza
wspólnym zbiorem zerowym i zbiegają prawie wszędzie
do $F(t_0,\cdot)$. Zespolona wersja twierdzenia
\ref{thm:dct} daje $\Phi(t_n)\to\Phi(t_0)$.
Ponieważ na przestrzeni metrycznej ciągowa
ciągłość jest równoważna ciągłości, otrzymujemy tezę.
\end{proof}

\begin{twierdzenie}[Różniczkowanie pod znakiem całki]
\label{thm:differentiating-integrals}
Niech $I\subseteq\mathbb R$ będzie przedziałem otwartym,
$t_0\in I$, a $F:I\times X\to\mathbb R$ funkcją,
dla której $x\mapsto F(t,x)$ jest mierzalna
dla każdego $t\in I$. Załóżmy, że:
\begin{enumerate}[label=\textup{(\roman*)}]
\item $F(t_0,\cdot)$ jest całkowalna;
\item istnieją $\delta>0$ o
$[t_0-\delta,t_0+\delta]\subset I$ i całkowalna
$g:X\to[0,\infty]$ takie, że dla prawie wszystkich
$x$ funkcja $t\mapsto F(t,x)$ jest różniczkowalna
na $(t_0-\delta,t_0+\delta)$ i
$|\partial_tF(t,x)|\le g(x)$ dla wszystkich $t$ z
tego przedziału (zbiór wyjątkowy jest wspólny).
\end{enumerate}
Wówczas $F(t,\cdot)$ jest całkowalna dla
$t$ dostatecznie bliskich $t_0$, funkcja
$\Phi(t)=\int_XF(t,x)\,\mathrm d\mu(x)$ jest
różniczkowalna w $t_0$ i
\begin{equation}\label{eq:differentiation-under-integral}
\Phi'(t_0)=\int_X\partial_tF(t_0,x)\,\mathrm d\mu(x).
\end{equation}
W całce pochodną $\partial_tF(t_0,\cdot)$ przyjmujemy
równą zeru na wspólnym mierzalnym zbiorze wyjątkowym,
na którym nie założono różniczkowalności.
\end{twierdzenie}
\begin{proof}
Poza wspólnym zbiorem zerowym zastosujmy twierdzenie
o wartości średniej do funkcji $t\mapsto F(t,x)$.
Dla $|h|<\delta$ mamy
$|F(t_0+h,x)-F(t_0,x)|\le |h|g(x)$.
Zatem $|F(t_0+h,x)|\le|F(t_0,x)|+|h|g(x)$,
a prawa strona jest całkowalna. To dowodzi
istnienia całek w otoczeniu punktu $t_0$.
Funkcje mierzalne
$Q_h(x)=(F(t_0+h,x)-F(t_0,x))/h$ spełniają
$|Q_h(x)|\le g(x)$ dla wszystkich małych $h\ne0$
i zbiegają dla prawie wszystkich $x$ do
$\partial_tF(t_0,x)$. Wybieramy $n_0$ takie, że $1/n_0<\delta$.
Niech $N$ będzie wspólnym mierzalnym zbiorem wyjątkowym
miary zero. Dla $n\ge n_0$ definiujemy mierzalną funkcję
$\widetilde Q_n$ równą $Q_{1/n}$ poza $N$ i zero na $N$.
Funkcje $\widetilde Q_n$ zbiegają punktowo na całym $X$
do pochodnej zdefiniowanej jako zero na $N$.
Z twierdzenia \ref{thm:pointwise-limits} wynika mierzalność
tej pochodnej, bez założenia zupełności miary. Zbieżność zmajoryzowana,
zastosowana do dowolnego ciągu $h_n\to0$,
$h_n\ne0$, daje
$\lim_n\int_XQ_{h_n}\,\mathrm d\mu(x)
=\int_X\partial_tF(t_0,x)\,\mathrm d\mu(x)$.
Z liniowości całki lewa strona jest granicą
ilorazów różnicowych funkcji $\Phi$; niezależność
od wybranego ciągu dowodzi istnienia pochodnej.
\end{proof}

\begin{przyklad}[Transformata Laplace'a funkcji ograniczonej]
Niech $u:[0,\infty)\to\mathbb R$ będzie mierzalna,
$|u(x)|\le M<\infty$, i dla $t>0$ określmy
$L(t)=\int_{[0,\infty)}e^{-tx}u(x)\,\mathrm d\lambda_1(x)$.
Całka jest skończona, gdyż $|e^{-tx}u(x)|\le Me^{-tx}$,
a całka z $e^{-tx}$ wynosi $1/t$ (z granicy
całek Riemanna na $[0,R]$, używając Beppo Leviego).
Dla ustalonego $t_0>0$ i $t\in(t_0/2,3t_0/2)$
mamy
$|\partial_t(e^{-tx}u(x))|=xe^{-tx}|u(x)|
\le Mxe^{-t_0x/2}$.
Funkcja dominująca jest całkowalna, np. z oszacowania
$x e^{-t_0x/2}\le C e^{-t_0x/4}$, gdzie $C<\infty$.
Twierdzenie \ref{thm:differentiating-integrals} daje
$L'(t)=-\int_{[0,\infty)}xe^{-tx}u(x)\,\mathrm d\lambda_1(x)$.
Analogicznie, przez indukcję, otrzymamy
$L^{(k)}(t)=(-1)^k\int_{[0,\infty)}x^ke^{-tx}u(x)
\,\mathrm d\lambda_1(x)$ dla wszystkich $k\ge1$.
Istotnie, dla ustalonych $t_0>0$ i $k\ge0$
wszystkie potrzebne pochodne są lokalnie
zdominowane przez $Mx^{k+1}e^{-t_0x/2}$;
funkcja ta jest całkowalna, ponieważ dla
$d=t_0/4$ iloczyn $x^{k+1}e^{-dx}$ jest
ograniczony na $[0,\infty)$.
\end{przyklad}

\section{Jednostajna całkowalność i zbieżność według miary}

\begin{definicja}[Jednostajna całkowalność]
\label{def:uniform-integrability}
Niech $\mu(X)<\infty$. Niepustą rodzinę $\mathcal U$
mierzalnych, całkowalnych funkcji rzeczywistych
na $X$ nazywamy \emph{jednostajnie całkowalną},
jeśli
\[
\lim_{K\to\infty}\sup_{u\in\mathcal U}
\int_{\{|u|>K\}}|u(x)|\,\mathrm d\mu(x)=0.
\]
Warunek kontroluje masę części funkcji o dużych
wartościach, a nie jedynie wielkość miary poziomic.
\end{definicja}

\begin{lemat}[Kontrola całek na małych zbiorach]
\label{lem:uniform-absolute-continuity}
Jeśli $\mu(X)<\infty$ i $\mathcal U$ jest jednostajnie
całkowalną rodziną funkcji rzeczywistych, to:
\begin{enumerate}[label=\textup{(\roman*)}]
\item $\sup_{u\in\mathcal U}\int_X|u|\,\mathrm d\mu(x)<\infty$;
\item dla każdego $\varepsilon>0$ istnieje $\delta>0$
takie, że dla każdego $A\in\mathcal F$ z
$\mu(A)<\delta$ zachodzi
$\sup_{u\in\mathcal U}\int_A|u|\,\mathrm d\mu(x)<\varepsilon$.
\end{enumerate}
\end{lemat}
\begin{proof}
Ustalmy $\varepsilon>0$ i wybierzmy $K>0$ takie, że
$\sup_{u\in\mathcal U}\int_{\{|u|>K\}}|u|<\varepsilon/2$.
Dla każdego $u\in\mathcal U$ i $A\in\mathcal F$
z podziału na $\{|u|\le K\}$ i $\{|u|>K\}$
wynika
$\int_A|u|\,\mathrm d\mu(x)\le K\mu(A)+\varepsilon/2$.
Biorąc $A=X$, dowodzimy (i) przy skończonej
mierze $\mu(X)$. Wybór
$\delta=\varepsilon/(2K)$ dla $K>0$ daje (ii);
przy $K=0$ oszacowanie jest nawet prostsze.
\end{proof}

\begin{twierdzenie}[Dwa równoważne opisy jednostajnej całkowalności]
\label{thm:ui-two-criteria-8}
Niech $\mu(X)<\infty$ i niech $\mathcal U$ będzie
\emph{niepustą} rodziną rzeczywistych funkcji mierzalnych
$u:X\to\mathbb R$, spełniających
$\int_X|u|\,\mathrm d\mu(x)<\infty$.
Równość funkcji prawie wszędzie nie ma wpływu
na żadną z rozważanych własności.
Wówczas $\mathcal U$ jest jednostajnie całkowalna
w sensie definicji~\ref{def:uniform-integrability}
wtedy i tylko wtedy, gdy jednocześnie:
\begin{enumerate}[label=\textup{(\roman*)}]
\item $M:=\sup_{u\in\mathcal U}\int_X|u|\,\mathrm d\mu(x)<\infty$;
\item dla każdego $\varepsilon>0$ istnieje $\delta>0$
takie, że $\mu(A)<\delta$ implikuje
$\sup_{u\in\mathcal U}\int_A|u|\,\mathrm d\mu(x)<\varepsilon$.
\end{enumerate}
\end{twierdzenie}
\begin{proof}
Konieczność to lemat~\ref{lem:uniform-absolute-continuity}.
Odwrotnie, przy $M=0$ wszystkie funkcje są zerowe
prawie wszędzie, więc teza jest natychmiastowa.
Dla $M>0$ i zadanego $\varepsilon>0$ wybieramy
$\delta$ z (ii). Nierówność Markowa daje
$\mu(\{|u|>K\})\le M/K$ dla każdego $u\in\mathcal U$.
Gdy $K>M/\delta$, zbiór $\{|u|>K\}$ ma miarę
mniejszą niż $\delta$, zatem (ii) daje
$\int_{\{|u|>K\}}|u|<\varepsilon$ jednostajnie
względem $u$. Jest to dokładnie warunek z
definicji~\ref{def:uniform-integrability}.
\end{proof}

\begin{twierdzenie}[Zbieżność Vitaliego na przestrzeni o mierze skończonej]
\label{thm:vitali-l1}
Niech $\mu(X)<\infty$, a $f_n:X\to\mathbb R$ będą
mierzalnymi funkcjami tworzącymi rodzinę
jednostajnie całkowalną. Jeśli
$f_n\xrightarrow{\mu}f$ dla pewnej mierzalnej
$f:X\to\mathbb R$, to $f$ jest całkowalna oraz
\[
\lim_{n\to\infty}\int_X|f_n(x)-f(x)|\,\mathrm d\mu(x)=0.
\]
W szczególności $\int_Xf_n\,\mathrm d\mu(x)\to
\int_Xf\,\mathrm d\mu(x)$.
\end{twierdzenie}
\begin{proof}
Z twierdzenia \ref{thm:measure-subsequence}
można wybrać podciąg $f_{n_k}\to f$ prawie wszędzie.
Na mocy lematu Fatou i części (i) lematu
\ref{lem:uniform-absolute-continuity}
$\int_X|f|\,\mathrm d\mu(x)\le
\liminf_k\int_X|f_{n_k}|\,\mathrm d\mu(x)<\infty$.

Ustalmy $\varepsilon>0$. Z lematu
\ref{lem:uniform-absolute-continuity} wybieramy
$\delta_1>0$ takie, że
$\mu(A)<\delta_1\Rightarrow
\sup_n\int_A|f_n|\,\mathrm d\mu(x)<\varepsilon/3$.
Ponieważ $f$ jest całkowalna, również jej całka
jest bezwzględnie ciągła względem miary w następującym
sensie: istnieje $\delta_2>0$ takie, że
$\mu(A)<\delta_2\Rightarrow\int_A|f|<\varepsilon/3$.
Istotnie, wystarczy obciąć $|f|$ na poziomie $K$,
tak by $\int_{\{|f|>K\}}|f|<\varepsilon/6$
(co wynika z twierdzenia \ref{thm:dct}, zastosowanego do
$|f|\chi_{\{|f|>K\}}$ dla całkowitych $K\to\infty$),
i oszacować pozostałą część przez $K\mu(A)$.
Dla $K>0$ wystarczy $\delta_2=\varepsilon/(6K)$.

Niech $\eta>0$ będzie tak małe, aby
$\eta\mu(X)<\varepsilon/3$. Ze zbieżności według
miary wynika, że dla dużych $n$ zbiór
$A_n=\{|f_n-f|>\eta\}$ ma miarę mniejszą od
$\min(\delta_1,\delta_2)$. Dzieląc całkę na
$A_n$ i dopełnienie, dostajemy
\[
\int_X|f_n-f|\,\mathrm d\mu(x)
\le\eta\mu(X)+\int_{A_n}|f_n|\,\mathrm d\mu(x)
+\int_{A_n}|f|\,\mathrm d\mu(x)<\varepsilon.
\]
To dowodzi zbieżności w całce z modułu różnicy.
Równość granic całek wynika z oszacowania
$|\int_X f_n-\int_X f|\le\int_X|f_n-f|$.
\end{proof}

\begin{przyklad}[Masę można skoncentrować na małych zbiorach]
Na $(0,1)$ niech $f_n=n\chi_{(0,1/n)}$.
Ciąg zbiega do zera według miary, ponieważ
$\lambda_1(\{|f_n|>\varepsilon\})\le1/n$
dla dużych $n$. Rodzina nie jest jednak jednostajnie
całkowalna: dla każdego $K>0$ i każdego $n>K$
zachodzi
$\int_{\{|f_n|>K\}}|f_n|\,\mathrm d\lambda_1(x)=1$.
Twierdzenie Vitaliego zatem nie ma zastosowania.
\end{przyklad}

\section{Związek z całką Riemanna}

\noindent\emph{Powiązanie tematyczne w Bibliotece matematycznej:} \atlasref{AM1-29::ma1:thm:riemann-darboux-equivalence}{AM1-29: równoważność konstrukcji Riemanna i Darboux}.


\begin{twierdzenie}[Zgodność całki Riemanna i Lebesgue'a]
\label{thm:riemann-to-lebesgue}
Niech $a<b$ i $f:[a,b]\to\mathbb R$ będzie funkcją
ograniczoną, całkowalną w sensie Riemanna.
Wówczas $f$ jest mierzalna względem uzupełnionej
miary Lebesgue'a na $[a,b]$, jest całkowalna
w sensie Lebesgue'a, a obie całki są równe:
\[
\int_a^b f(x)\,\mathrm dx
=\int_{[a,b]}f(x)\,\mathrm d\lambda_1(x).
\]
Po lewej stronie występuje tradycyjna całka Riemanna,
a po prawej całka względem miary Lebesgue'a.
\end{twierdzenie}
\begin{proof}
Niech $M$ będzie ograniczeniem $|f|$.
Z kryterium Darboux dla całkowalności Riemanna
(\atlasref{AM1-29::ma1:thm:darboux-criterion}{AM I, rozdział 29})
możemy wybrać podziały $P_n$ takie, że
$U(f,P_n)-L(f,P_n)<2^{-2n}$,
gdzie $U,L$ są odpowiednio górnymi i dolnymi
sumami Darboux. Z każdego podziału budujemy
proste funkcje borelowskie $u_n,\ell_n$:
na rozłącznych półotwartych częściach podziału
przypisujemy odpowiednio supremum i infimum $f$
na odpowiadających im \emph{domkniętych} przedziałach
podziału (ostatni przedział zawiera prawy koniec).
Dzięki temu $\ell_n\le f\le u_n$ w każdym punkcie,
a $\int_{[a,b]}u_n(x)\,\mathrm d\lambda_1(x)=U(f,P_n)$,
$\int_{[a,b]}\ell_n(x)\,\mathrm d\lambda_1(x)=L(f,P_n)$.

Ponieważ $u_n-\ell_n\ge0$ i
$\int_{[a,b]}(u_n-\ell_n)\,\mathrm d\lambda_1(x)
<2^{-2n}$, nierówność Markowa daje
$\lambda_1(\{u_n-\ell_n\ge2^{-n}\})\le2^{-n}$.
Pierwsza część lematu Borela--Cantellego (twierdzenie \ref{thm:borel-cantelli-1}) implikuje,
że dla prawie wszystkich $x$ istnieje $N(x)$ takie,
że $u_n(x)-\ell_n(x)<2^{-n}$ dla $n\ge N(x)$.
Z nierówności $\ell_n\le f\le u_n$ otrzymujemy
$u_n(x)\to f(x)$ poza mierzalnym zbiorem zerowym.
Granica mierzalnych funkcji jest mierzalna poza
tym zbiorem; dzięki zupełności miary Lebesgue'a
sama $f$ jest mierzalna. Ponadto $|f|\le M$
i $\lambda_1([a,b])<\infty$, więc $f$ jest całkowalna.

Z monotoniczności całki i równości dla prostych
funkcji otrzymujemy
$L(f,P_n)\le\int_{[a,b]}f\,\mathrm d\lambda_1(x)
\le U(f,P_n)$. Obie skrajne liczby dążą do
całki Riemanna, bo ich różnica dąży do zera,
a całka Riemanna leży pomiędzy nimi.
Twierdzenie o trzech ciągach kończy dowód.
\end{proof}

\begin{twierdzenie}[Kryterium Lebesgue'a całkowalności Riemanna]
\label{thm:riemann-criterion}
Niech $f:[a,b]\to\mathbb R$ będzie ograniczona, gdzie
$a<b$. Wówczas $f$ jest całkowalna w sensie Riemanna
tedy i tylko wtedy, gdy zbiór punktów nieciągłości $f$
ma miarę Lebesgue'a zero.
\end{twierdzenie}
\begin{proof}
Dla $x\in[a,b]$ oznaczmy przez $\omega_f(x)$
oscylację $f$ w punkcie, czyli
\[
\omega_f(x):=\inf_{r>0}
\left(\sup_{y\in[a,b]\cap(x-r,x+r)}f(y)
-\inf_{y\in[a,b]\cap(x-r,x+r)}f(y)\right).
\]
Jeżeli $f$ jest ciągła w $x$, to dla dowolnego $\eta>0$ na dostatecznie małym otoczeniu względnym zachodzi $|f(y)-f(x)|<\eta/3$, więc jego oscylacja nie przekracza $2\eta/3$. Stąd $\omega_f(x)=0$. Odwrotnie, jeśli $\omega_f(x)=0$, to dla każdego $\eta>0$ istnieje otoczenie względne o oscylacji mniejszej od $\eta$. Ponieważ zawiera ono $x$, dla każdego należącego do niego $y$ zachodzi $|f(y)-f(x)|<\eta$. To daje ciągłość. Dla $\varepsilon>0$ zbiór
$D_\varepsilon=\{x:\omega_f(x)\ge\varepsilon\}$
jest domknięty: jeśli $\omega_f(x)<\varepsilon$,
to dla pewnego $r>0$ oscylacja $f$ na $[a,b]\cap B(x,r)$ jest mniejsza od $\varepsilon$. Dla każdego $y\in[a,b]\cap B(x,r/2)$ otoczenie $[a,b]\cap B(y,r/2)$ zawiera się w poprzednim, więc $\omega_f(y)<\varepsilon$. Zatem $D_\varepsilon$ jest zwarty.

Załóżmy najpierw całkowalność Riemanna i
ustalmy $\varepsilon>0$. Weźmy dowolny podział
$P=(a=x_0<\cdots<x_N=b)$. Jeśli
$x\in D_\varepsilon$ nie jest punktem podziału,
to oscylacja $f$ na domkniętym przedziale
$[x_{j-1},x_j]$ zawierającym $x$ jest co najmniej
$\varepsilon$. Dlatego $D_\varepsilon$ mieści się
w sumie skończenie wielu punktów podziału i tych
przedziałów, których oscylacja jest nie mniejsza
niż $\varepsilon$. Suma długości tych przedziałów
jest nie większa niż
$\varepsilon^{-1}(U(f,P)-L(f,P))$.
Ponieważ dla funkcji całkowalnej w sensie Riemanna
różnicę tę można uczynić dowolnie małą na mocy kryterium Darboux
(\atlasref{AM1-29::ma1:thm:darboux-criterion}{AM I, rozdział 29}),
$\lambda_1(D_\varepsilon)=0$.
Zbiór nieciągłości jest sumą
$\bigcup_{m\ge1}D_{1/m}$, więc jest zerowy.

Odwrotnie, niech zbiór nieciągłości $D$ ma miarę
zero, a $|f|\le M$. Ustalmy $\varepsilon>0$.
Zbiór $D_\varepsilon$ jest zwarty i zerowy, więc
istnieje zbiór otwarty $O$ zawierający $D_\varepsilon$
o dowolnie małej mierze Lebesgue'a (twierdzenie \ref{thm:lebesgue-outer-regular}). Wybieramy przy tym
$\lambda_1(O)<b-a$, dzięki czemu zwarty zbiór
$K=[a,b]\setminus O$ jest niepusty. Dla każdego
$x\in K$ mamy $\omega_f(x)<\varepsilon$,
zatem istnieje promień $r_x>0$, na którego kuli
$B(x,r_x)\cap[a,b]$ oscylacja $f$ jest
mniejsza niż $\varepsilon$. Zwartość $K$ pozwala
wybrać skończenie wiele kul $B(x_j,r_{x_j}/2)$
pokrywających $K$. Połóżmy
$\delta=\frac12\min_j r_{x_j}>0$.
Jeżeli domknięty przedział $J\subseteq[a,b]$
o długości mniejszej niż $\delta$ przecina $K$,
wybieramy $y\in J\cap K$ i indeks $j$ taki, że
$y\in B(x_j,r_{x_j}/2)$. Dla każdego $z\in J$
mamy
$|z-x_j|\le|z-y|+|y-x_j|<\delta+r_{x_j}/2
\le r_{x_j}$.
Zatem $J\subseteq B(x_j,r_{x_j})$, gdzie
oscylacja $f$ jest mniejsza niż $\varepsilon$.

Weźmy podział $P$ o średnicy mniejszej niż $\delta$.
Na przedziałach przecinających $K$ oscylacja $f$
jest mniejsza niż $\varepsilon$. Pozostałe przedziały, łącznie z końcami, zawierają się w $O$, ponieważ nie przecinają $K$. Mają parami rozłączne wnętrza, a wspólne końce mają miarę zero, więc suma ich długości nie przekracza $\lambda_1(O)$. Otrzymujemy
\[
U(f,P)-L(f,P)
\le\varepsilon(b-a)+2M\lambda_1(O).
\]
Dla zadanej dokładności $\eta>0$ bierzemy $\varepsilon=\eta/(2(b-a))$ i wybieramy $O$ o mierze mniejszej niż $\min\{b-a,\eta/(4(M+1))\}$. Wówczas prawa strona jest mniejsza od $\eta$; wybór działa także dla $M=0$ i $D_\varepsilon=\varnothing$. Kryterium Darboux dowodzi
całkowalności Riemanna.
\end{proof}

\begin{uwaga}[Dlaczego oba warunki mają znaczenie?]
Funkcja Dirichleta $\chi_{\mathbb Q\cap[0,1]}$
jest całkowalna w sensie Lebesgue'a, ale
nie jest całkowalna w sensie Riemanna, gdyż
zbiór jej punktów nieciągłości to cały $[0,1]$.
Z kolei każda ograniczona funkcja o przeliczalnym
zbiorze punktów nieciągłości jest całkowalna
w sensie Riemanna. Warunek ograniczoności jest
w kryterium konieczny w jego klasycznej postaci:
na przykład funkcja na $[0,1]$ określona przez $f(0)=0$
i $f(x)=1/x$ dla $x>0$ jest nieograniczona, choć
jej jedynym punktem nieciągłości jest $0$.
\end{uwaga}

\section{Podsumowanie logiczne}

Twierdzenie Beppo Leviego opiera się na definicji
całki jako supremum po funkcjach prostych i na
ciągłości miary z dołu. Z niego wynikają
addytywność całki nieujemnej i możliwość całkowania
szeregów nieujemnych. Lemat Fatou to zastosowanie
Beppo Leviego do ciągu infimów ogonowych, a twierdzenie
o zbieżności zmajoryzowanej uzyskujemy z lematu Fatou
przez rozważenie nieujemnych funkcji $g\pm f_n$.
Wszystkie te rezultaty pozostają ważne dla miar
nieskończonych; wyjątkiem jest twierdzenie o zbieżności
ograniczonej, w którym skończoność miary
zapewnia całkowalność stałej funkcji dominującej.

% ROZDZIAŁ 9 — samodzielny fragment do włączenia w main.tex
\chapter{Rodzaje zbieżności funkcji mierzalnych}
\label{ch:rodzaje-zbieznosci}

\section{Dlaczego potrzebujemy kilku pojęć zbieżności?}

Zbieżność punktowa opisuje zachowanie funkcji w każdym ustalonym
punkcie, lecz nie kontroluje zbiorów, na których funkcje przyjmują
duże wartości. Zbieżność według miary kontroluje rozmiar takich
zbiorów, ale dopuszcza ich przemieszczanie. Zbieżność w całce
kontroluje dodatkowo całkowitą wielkość błędu. To właśnie różnice
między tymi warunkami wyjaśniają, dlaczego nie można dowolnie
zamieniać kolejności granicy i całkowania.

Przez cały rozdział $(X,\mathcal F,\mu)$ jest przestrzenią
z miarą; nie wymagamy zupełności ani $\sigma$-skończoności, chyba
że zaznaczono inaczej. Wszystkie funkcje $f_n,f$ są mierzalne i
rzeczywiste, o ile nie wprowadzono innych założeń. Mówimy, że
własność zachodzi \emph{$\mu$-prawie wszędzie}, gdy zbiór punktów,
w których może zawodzić, zawiera się w pewnym zbiorze należącym
do $\mathcal F$ o mierze zero. Jest to ważne przy miarach
niezupełnych.

\section{Definicje i pierwsze relacje}

Zbieżność według miary zdefiniowaliśmy już w
Definicji~\ref{def:convergence-measure}, a twierdzenie o podciągu
udowodniliśmy w Twierdzeniu~\ref{thm:measure-subsequence}.
Teraz porównujemy ją z pozostałymi pojęciami zbieżności.


\begin{definicja}[Pięć rodzajów zbieżności]
\label{def:five-convergences}
Niech $(X,\mathcal F,\mu)$ będzie przestrzenią z miarą,
a $f_n,f:X\to\mathbb R$ funkcjami mierzalnymi.
\begin{enumerate}[label=\textup{(\roman*)}]
\item $f_n\to f$ \emph{punktowo}, jeśli
$f_n(x)\to f(x)$ dla każdego $x\in X$.
\item $f_n\to f$ \emph{$\mu$-prawie wszędzie}, w skrócie
$f_n\xrightarrow{\mathrm{p.w.}}f$, jeśli zbieżność
punktowa zachodzi poza pewnym mierzalnym zbiorem zerowym.
\item $f_n\to f$ \emph{jednostajnie}, jeśli dla każdego
$\varepsilon>0$ istnieje $N$, takie że
$|f_n(x)-f(x)|<\varepsilon$ dla wszystkich $n\ge N$
i $x\in X$.
\item $f_n\to f$ \emph{prawie jednostajnie}, jeśli dla każdego
$\eta>0$ istnieje $E\in\mathcal F$ z $\mu(E)<\eta$
takie, że $f_n\to f$ jednostajnie na $X\setminus E$.
\item $f_n\to f$ \emph{według miary}, co oznaczamy
$f_n\xrightarrow{\mu}f$, jeśli dla każdego $\varepsilon>0$
zachodzi
\begin{equation}\label{eq:conv-in-measure-9}
\lim_{n\to\infty}\mu\bigl(\{x\in X:
|f_n(x)-f(x)|>\varepsilon\}\bigr)=0.
\end{equation}
\end{enumerate}
Nie wymagamy, aby zbiory poziomicowe w
\eqref{eq:conv-in-measure-9} miały skończoną miarę dla
poszczególnych $n$; wymagana jest ich granica równa zeru.
\end{definicja}

\begin{uwaga}[Dwie kwantyfikacje, których nie wolno mylić]
W zbieżności jednostajnej kolejność kwantyfikatorów jest
następująca: dla każdego $\varepsilon>0$ istnieje jedno $N$
działające \emph{jednocześnie} dla wszystkich $x$. W zbieżności
punktowej indeks $N$ może zależeć od $x$. Zbieżność prawie
jednostajna pozwala usunąć zbiór o dowolnie małej mierze,
ale zbiór ten zależy od wymaganego dopuszczalnego błędu w mierze.
Zbiór wyjątkowy nie może natomiast zależeć od indeksu $n$
w definicji jednostajności na jego dopełnieniu.
\end{uwaga}

\begin{twierdzenie}[Podstawowe implikacje]
\label{thm:conv-implications-9}
Dla dowolnej przestrzeni z miarą $(X,\mathcal F,\mu)$
i mierzalnych funkcji $f_n,f:X\to\mathbb R$ zachodzą:
\begin{enumerate}[label=\textup{(\roman*)}]
\item zbieżność jednostajna pociąga zbieżność punktową
i według miary;
\item zbieżność prawie jednostajna pociąga zbieżność
prawie wszędzie i według miary;
\item zbieżność punktowa pociąga zbieżność prawie wszędzie.
\end{enumerate}
Jeżeli dodatkowo $\mu(X)<\infty$, to zbieżność prawie
wszędzie pociąga zbieżność według miary.
\end{twierdzenie}
\begin{proof}
W (i) dla ustalonego $\varepsilon>0$ i wszystkich dużych
$n$ zbiór $\{|f_n-f|>\varepsilon\}$ jest pusty.
Implikacja punktowa jest oczywista, podobnie (iii).

Dla (ii) ustalmy $\varepsilon,\eta>0$ i wybierzmy
$E\in\mathcal F$ o $\mu(E)<\eta$, poza którym zbieżność
jest jednostajna. Dla dużych $n$ zbiór
$\{|f_n-f|>\varepsilon\}$ zawiera się w $E$,
zatem jego miara jest nie większa od $\eta$.
Ponieważ $\eta$ jest dowolne, otrzymujemy
\eqref{eq:conv-in-measure-9}.

Niech teraz $N$ będzie zbiorem, na którym zbieżność
punktowa nie zachodzi. Ponieważ
\[
N=\bigcup_{m\ge1}\bigcap_{N_0\ge1}
\bigcup_{n\ge N_0}\{|f_n-f|\ge 1/m\},
\]
zbiór $N$ jest mierzalny. Jeśli zbieżność jest prawie
jednostajna, to dla każdego $\eta>0$ istnieje zbiór
$E$ o mierze mniejszej niż $\eta$, poza którym zbieżność
zachodzi; stąd $N\subseteq E$ i $\mu(N)\le\eta$.
Wynika $\mu(N)=0$, co dowodzi drugiej części (ii).

Ostatnia implikacja jest dokładnie twierdzeniem~\ref{thm:ae-to-measure-finite}, już udowodnionym w rozdziale~\ref{ch:funkcje-mierzalne}. To kończy dowód.
\end{proof}

\begin{przyklad}[Dlaczego w ostatniej implikacji potrzebna jest skończoność miary]
\label{ex:ae-not-measure-infinite-9}
Na $\mathbb R$ z miarą Lebesgue'a niech
$f_n=\chi_{[n,\infty)}$. Dla każdego ustalonego
$x\in\mathbb R$ mamy $f_n(x)\to0$, natomiast
dla $0<\varepsilon<1$ zachodzi
$\lambda_1(\{|f_n|>\varepsilon\})=\infty$ dla
każdego $n$. Nie wystarcza nawet zbieżność punktowa
na całej przestrzeni; problemem jest brak kontroli
masy poza zbiorami skończonej miary.
\end{przyklad}

\section{Twierdzenie Jegorowa}

\begin{twierdzenie}[Jegorow]
\label{thm:egorov-9}
Niech $(X,\mathcal F,\mu)$ będzie przestrzenią z miarą
skończoną, tj. $\mu(X)<\infty$. Jeśli
$f_n,f:X\to\mathbb R$ są mierzalne oraz
$f_n\to f$ $\mu$-prawie wszędzie, to
$f_n\to f$ prawie jednostajnie. Dokładniej: dla każdego
$\eta>0$ istnieje $E\in\mathcal F$, $\mu(E)<\eta$,
takie, że
$\sup_{x\in X\setminus E}|f_n(x)-f(x)|\to0$.
Jeśli $X\setminus E=\varnothing$, przyjmujemy supremum
nieujemnej funkcji na zbiorze pustym równe $0$;
zbieżność jednostajna na zbiorze pustym zachodzi.
\end{twierdzenie}
\begin{proof}
Niech $N\in\mathcal F$ będzie zbiorem miary zero,
poza którym $f_n(x)\to f(x)$. Dla $m,N_0\in\mathbb N$
określmy
\[
A_{m,N_0}=\bigcup_{n\ge N_0}
\{x\in X:|f_n(x)-f(x)|\ge1/m\}.
\]
Zbiory $A_{m,N_0}$ są mierzalne i dla ustalonego $m$
maleją wraz z $N_0$. Zbieżność poza $N$ oznacza,
że $\bigcap_{N_0\ge1}A_{m,N_0}\subseteq N$.
Zatem miara tego przecięcia wynosi zero. Skończoność
miary $X$ pozwala zastosować ciągłość z góry:
$\mu(A_{m,N_0})\to0$ dla $N_0\to\infty$.

Ustalmy $\eta>0$. Dla każdego $m$ możemy wybrać
$N_m$ tak duże, by
$\mu(A_{m,N_m})<\eta 2^{-(m+1)}$.
Połóżmy $E=\bigcup_{m\ge1}A_{m,N_m}$.
Przeliczalna podaddytywność daje
$\mu(E)<\eta$.
Dla $x\notin E$ i każdego $m$ oraz $n\ge N_m$
mamy $|f_n(x)-f(x)|<1/m$. Jeśli teraz
$\varepsilon>0$, wybieramy $m$ z $1/m<\varepsilon$;
wówczas nierówność z indeksem $N_m$ zachodzi
jednocześnie dla wszystkich $x\in X\setminus E$.
To jest dokładnie zbieżność jednostajna na
$X\setminus E$.
\end{proof}

\begin{wniosek}[Wersja lokalna dla przestrzeni $\sigma$-skończonej]
\label{cor:local-egorov-9}
Niech $(X,\mathcal F,\mu)$ będzie przestrzenią z miarą,
a $f_n,f:X\to\mathbb R$ funkcjami mierzalnymi.
Załóżmy, że $X=\bigcup_{j\ge1}X_j$, gdzie $X_j\in\mathcal F$,
$X_j\subseteq X_{j+1}$ i $\mu(X_j)<\infty$.
Jeżeli $f_n\to f$ prawie wszędzie, to dla każdego
$\eta>0$ istnieje $E\in\mathcal F$ o $\mu(E)<\eta$
takie, że $f_n\to f$ jednostajnie na
$X_j\setminus E$ dla każdego ustalonego $j$.
Nie twierdzimy, że zbieżność jest jednostajna
na całym $X\setminus E$.
\end{wniosek}
\begin{proof}
Dla każdego $j$ stosujemy twierdzenie
\ref{thm:egorov-9} do ograniczenia miary i funkcji
na $X_j$, z tolerancją $\eta 2^{-(j+1)}$.
Otrzymujemy $E_j\subseteq X_j$ takie, że
$\mu(E_j)<\eta 2^{-(j+1)}$ i zbieżność jest
jednostajna na $X_j\setminus E_j$.
Niech $E=\bigcup_jE_j$. Wtedy $\mu(E)<\eta$, a
$X_j\setminus E\subseteq X_j\setminus E_j$,
więc zbieżność jest jednostajna na pierwszym z tych
zbiorów.
\end{proof}

\begin{przyklad}[Nie istnieje globalna wersja Jegorowa na $\mathbb R$]
Niech $f_n=\chi_{[n,n+1)}$ na $\mathbb R$.
Funkcje zbiegają do zera wszędzie, ale jeśli
$\lambda_1(E)<\infty$, to nie może zachodzić
$f_n\to0$ jednostajnie na $\mathbb R\setminus E$.
W przeciwnym razie istniałoby $N$ takie, że
$[n,n+1)\subseteq E$ dla wszystkich $n\ge N$
(gdyż $f_n$ przyjmuje wyłącznie wartości $0,1$).
Wynikałoby $\lambda_1(E)=\infty$, sprzeczność.
\end{przyklad}

\section{Zbieżność według miary i twierdzenie Riesza o podciągu}

Jednoznaczność granicy według miary, z dokładnością do zbioru zerowego,
została wykazana w Wniosku~\ref{cor:measure-limit-unique}.
Ponadto twierdzenie Riesza o podciągu jest już w pełni
udowodnione w Twierdzeniu~\ref{thm:measure-subsequence};
w dalszych wynikach korzystamy właśnie z tej wcześniejszej wersji.

\begin{wniosek}[Kryterium za pomocą wszystkich podciągów]
\label{cor:every-subseq-9}
Niech $(X,\mathcal F,\mu)$ będzie przestrzenią z miarą,
$\mu(X)<\infty$. Dla mierzalnych
$f_n,f:X\to\mathbb R$ następujące
warunki są równoważne:
\begin{enumerate}[label=\textup{(\roman*)}]
\item $f_n\xrightarrow{\mu}f$;
\item każdy podciąg $(f_{n_k})$ ma dalszy podciąg,
który zbiega do $f$ $\mu$-prawie wszędzie.
\end{enumerate}
\end{wniosek}
\begin{proof}
Jeżeli zachodzi (i), każdy podciąg także zbiega
według miary, więc stosujemy twierdzenie
\ref{thm:measure-subsequence}.
Załóżmy (ii) i przypuśćmy, że (i) nie zachodzi.
Wówczas istnieją $\varepsilon,\delta>0$ oraz
podciąg $(f_{n_k})$ taki, że
$\mu(\{|f_{n_k}-f|>\varepsilon\})\ge\delta$
dla każdego $k$. Na mocy (ii) ma on dalszy
podciąg zbieżny do $f$ prawie wszędzie.
Ponieważ $\mu(X)<\infty$, twierdzenie
\ref{thm:conv-implications-9} daje jego zbieżność
według miary, wbrew ustalonemu oszacowaniu.
\end{proof}

\begin{uwaga}[Nie można skreślić założenia $\mu(X)<\infty$]
Na przestrzeni o mierze nieskończonej warunek (ii)
nie wystarcza: ciąg $\chi_{[n,\infty)}$ na
$\mathbb R$ zbiega wszędzie do zera, lecz nie
według miary. Właśnie dlatego równoważność
\ref{cor:every-subseq-9} została sformułowana
dla miar skończonych.
\end{uwaga}

\section{Zbieżność w $L^p$ i oszacowania Czebyszewa}

\begin{definicja}[Zbieżność w średniej rzędu $p$]
\label{def:lp-convergence-9}
Niech $(X,\mathcal F,\mu)$ będzie przestrzenią z miarą
i niech $1\le p<\infty$. Mierzalną funkcję
$h:X\to\mathbb R$ nazywamy \emph{$p$-całkowalną},
gdy $\int_X|h|^p\,\mathrm d\mu(x)<\infty$.
Dla $p$-całkowalnych funkcji $f_n,f:X\to\mathbb R$
mówimy, że $f_n\to f$ \emph{w $L^p(\mu)$}, jeśli
$\int_X|f_n(x)-f(x)|^p\,\mathrm d\mu(x)\to0$.
Pełną konstrukcję przestrzeni normowanej $L^p$
na klasach funkcji równych prawie wszędzie,
wraz z nierównościami H\"oldera i Minkowskiego,
przedstawimy w rozdziale 14.
\end{definicja}

\begin{twierdzenie}[Nierówność Czebyszewa i zbieżność według miary]
\label{thm:lp-implies-measure-9}
Niech $(X,\mathcal F,\mu)$ będzie przestrzenią z miarą,
$1\le p<\infty$, a $f_n,f:X\to\mathbb R$
funkcjami mierzalnymi. Dla każdego $\varepsilon>0$ zachodzi
\begin{equation}\label{eq:chebyshev-lp-9}
\mu\bigl(\{|f_n-f|>\varepsilon\}\bigr)
\le\varepsilon^{-p}
\int_X|f_n(x)-f(x)|^p\,\mathrm d\mu(x).
\end{equation}
W szczególności zbieżność w $L^p(\mu)$ pociąga
zbieżność według miary.
\end{twierdzenie}
\begin{proof}
Stosujemy nierówność Markowa z rozdziału 7
do nieujemnej mierzalnej funkcji $|f_n-f|^p$
i poziomu $\varepsilon^p$. Jeśli jej całka dąży
do zera, prawa strona \eqref{eq:chebyshev-lp-9}
dąży do zera dla każdego ustalonego $\varepsilon>0$.
\end{proof}

\begin{przyklad}[Zbieżność w $L^p$ nie wymusza zbieżności prawie wszędzie]
\label{ex:typewriter-lp-9}
Dla $k\ge0$ i $0\le j<2^k$ niech
$g_{2^k+j}=\chi_{[j2^{-k},(j+1)2^{-k})}$
na $[0,1)$ z miarą Lebesgue'a. Dla każdego
$1\le p<\infty$ mamy
$\int_{[0,1)}|g_{2^k+j}|^p\,\mathrm d\lambda_1(x)=2^{-k}$,
więc $g_n\to0$ w $L^p$, a także według miary.
Jednak dla każdego $x\in[0,1)$ w każdym bloku o długości
$2^k$ dokładnie jedna funkcja przyjmuje w $x$
wartość $1$, a pozostałe wartość $0$ (dla $k\ge1$).
Stąd $g_n(x)$ nie ma granicy w żadnym punkcie $[0,1)$.
Jest to klasyczna \emph{maszyna do pisania}.
\end{przyklad}

\begin{przyklad}[Zbieżność prawie wszędzie i według miary bez zbieżności w $L^1$]
Niech $f_n=n\chi_{(0,1/n)}$ na $(0,1)$.
Wtedy $f_n\to0$ wszędzie i według miary, lecz
$\int_{(0,1)}|f_n|\,\mathrm d\lambda_1(x)=1$
dla każdego $n$. Brakuje jednostajnej całkowalności;
por. twierdzenie Vitaliego z rozdziału 8.
\end{przyklad}

\section{Ciągi Cauchy'ego według miary i zupełność}

\begin{definicja}[Ciąg Cauchy'ego według miary]
\label{def:cauchy-measure-9}
Niech $(X,\mathcal F,\mu)$ będzie przestrzenią z miarą.
Ciąg mierzalnych funkcji rzeczywistych
$f_n:X\to\mathbb R$ nazywamy \emph{ciągiem Cauchy'ego według miary}, jeśli
dla wszystkich $\varepsilon,\delta>0$ istnieje
$N\in\mathbb N$ takie, że dla $m,n\ge N$
zachodzi $\mu(\{|f_n-f_m|>\varepsilon\})<\delta$.
\end{definicja}

\begin{twierdzenie}[Zupełność zbieżności według miary]
\label{thm:cauchy-measure-complete-9}
Niech $(X,\mathcal F,\mu)$ będzie dowolną przestrzenią
z miarą. Jeżeli $(f_n)$ jest ciągiem Cauchy'ego według
miary złożonym z mierzalnych funkcji rzeczywistych,
to istnieje mierzalna funkcja $f:X\to\mathbb R$ taka,
że $f_n\xrightarrow{\mu}f$. Funkcja $f$ jest wyznaczona
jednoznacznie z dokładnością do równości prawie wszędzie.
\end{twierdzenie}
\begin{proof}
Wybierzmy ściśle rosnący ciąg indeksów $n_k$
rekurencyjnie tak, aby dla wszystkich $m,n\ge n_k$
zachodziło
$\mu(\{|f_m-f_n|>2^{-k}\})<2^{-k}$.
Możemy jednocześnie wymagać $n_{k+1}>n_k$.
Niech
$A_k=\{|f_{n_{k+1}}-f_{n_k}|>2^{-k}\}$.
Wówczas $\sum_k\mu(A_k)<\infty$ i lemat
Borela--Cantellego zapewnia, że poza mierzalnym
zbiorem zerowym $N$ każdy punkt należy tylko do
skończenie wielu $A_k$.

Dla $x\notin N$ i dużych $k$ mamy
$|f_{n_{k+1}}(x)-f_{n_k}(x)|\le2^{-k}$.
Dla $j>k$ nierówność trójkąta daje
\[
|f_{n_j}(x)-f_{n_k}(x)|
\le\sum_{\ell=k}^{j-1}2^{-\ell}\le2^{1-k}.
\]
Zatem $(f_{n_k}(x))$ jest ciągiem Cauchy'ego w
$\mathbb R$ i zbiega do pewnej skończonej liczby
$f(x)$. Na $N$ definiujemy $f(x)=0$.
Mierzalność $f$ wynika z twierdzenia o mierzalności
granic punktowych po przedłużeniu funkcji granicznej
zerem na $N$. Zbieżność $f_{n_k}\to f$ prawie
wszędzie nie gwarantuje sama przez się zbieżności
według miary na przestrzeni nieskończonej.

Tutaj jednak mamy silniejsze oszacowanie: dla
ustalonego $k$ i dowolnego $j>k$, poza zbiorem
$\bigcup_{\ell\ge k}A_\ell$, zachodzi
$|f_{n_j}-f_{n_k}|\le2^{1-k}$.
Przechodząc do granicy $j\to\infty$ poza $N$,
dostajemy
\[
\mu\bigl(\{|f-f_{n_k}|>2^{1-k}\}\bigr)
\le\sum_{\ell\ge k}\mu(A_\ell)
\le\sum_{\ell\ge k}2^{-\ell}\longrightarrow0.
\]
Zatem $f_{n_k}\to f$ według miary.
Dla ustalonych $\varepsilon,\delta>0$ warunek
Cauchy'ego daje $N_C$ takie, że dla $m,n\ge N_C$
zachodzi
$\mu(\{|f_n-f_m|>\varepsilon/2\})<\delta/2$.
Wybieramy teraz $k$ \emph{jednocześnie} tak duże,
aby $n_k\ge N_C$ i
$\mu(\{|f_{n_k}-f|>\varepsilon/2\})<\delta/2$.
Wówczas dla każdego $n\ge N_C$ mamy
$\mu(\{|f_n-f_{n_k}|>\varepsilon/2\})<\delta/2$.
Wówczas $\mu(\{|f_n-f|>\varepsilon\})<\delta$.
Jednoznaczność wynika z twierdzenia
\ref{cor:measure-limit-unique}.
\end{proof}

\section{Metryzacja zbieżności według miary na przestrzeni skończonej miary}

\begin{twierdzenie}[Metryka zbieżności według miary]
\label{thm:measure-metric-9}
Niech $(X,\mathcal F,\mu)$ będzie przestrzenią z miarą,
$\mu(X)<\infty$. Oznaczmy przez
$\mathcal M(X,\mathcal F)$ zbiór mierzalnych funkcji
rzeczywistych na $X$, a przez $\mathcal M/\!\sim$
jego iloraz przez relację
$f\sim g\Longleftrightarrow f=g$ $\mu$-prawie wszędzie.
Dla klas $[f],[g]\in\mathcal M/\!\sim$ połóżmy
\begin{equation}\label{eq:measure-metric-9}
\rho([f],[g])=
\int_X\min\{1,|f(x)-g(x)|\}\,\mathrm d\mu(x).
\end{equation}
Wówczas $\rho$ jest metryką na $\mathcal M/\!\sim$.
Ponadto $f_n\xrightarrow{\mu}f$ wtedy i tylko wtedy,
gdy $\rho([f_n],[f])\to0$.
\end{twierdzenie}
\begin{proof}
Całka w \eqref{eq:measure-metric-9} jest skończona,
bo całkujemy funkcję ograniczoną przez $1$,
a $\mu(X)<\infty$. Równość funkcji prawie
wszędzie nie zmienia wartości całki, więc wzór
zależy tylko od klas. Symetria i nieujemność
są oczywiste. Jeśli $\rho([f],[g])=0$, to
nieujemna funkcja $\min\{1,|f-g|\}$ ma całkę
zero, a więc znika prawie wszędzie; wynika
$[f]=[g]$. Nierówność trójkąta jest konsekwencją
punktowej nierówności
\[
\min\{1,|f-h|\}\le
\min\{1,|f-g|\}+\min\{1,|g-h|\},
\]
którą otrzymujemy z nierówności trójkąta dla
wartości bezwzględnej i oszacowania
$\min\{1,a+b\}\le\min\{1,a\}+\min\{1,b\}$
dla $a,b\ge0$. Całkujemy obie strony.

Załóżmy $\rho([f_n],[f])\to0$. Dla $\varepsilon>0$
zachodzi punktowo
$\min\{1,|f_n-f|\}\ge
\min\{1,\varepsilon\}\chi_{\{|f_n-f|>\varepsilon\}}$.
Zatem
\[
\mu(\{|f_n-f|>\varepsilon\})\le
\frac{\rho([f_n],[f])}{\min\{1,\varepsilon\}}
\longrightarrow0.
\]

Odwrotnie, jeśli $f_n\to f$ według miary,
to dla każdego $0<\varepsilon<1$ mamy
\[
\rho([f_n],[f])\le\varepsilon\mu(X)+
\mu(\{|f_n-f|>\varepsilon\}).
\]
Dla ustalonego $\varepsilon$ drugi składnik
dąży do zera, zatem
$\limsup_n\rho([f_n],[f])\le\varepsilon\mu(X)$.
Pozwalając $\varepsilon\downarrow0$, otrzymujemy
$\rho([f_n],[f])\to0$.
\end{proof}

\begin{wniosek}[Zupełność metryki zbieżności według miary]
\label{cor:measure-metric-complete-9}
Niech $(X,\mathcal F,\mu)$ będzie przestrzenią z miarą,
$\mu(X)<\infty$, a $(\mathcal M/\!\sim,\rho)$ przestrzenią
metryczną klas mierzalnych funkcji rzeczywistych z metryką
zdefiniowaną w twierdzeniu \ref{thm:measure-metric-9}.
Wówczas $(\mathcal M/\!\sim,\rho)$ jest zupełna.
\end{wniosek}
\begin{proof}
Niech $([f_n])$ będzie ciągiem Cauchy'ego w metryce
$\rho$. Dla $\varepsilon>0$ mamy
$\mu(\{|f_n-f_m|>\varepsilon\})\le
\rho([f_n],[f_m])/\min\{1,\varepsilon\}$.
Zatem $(f_n)$ jest ciągiem Cauchy'ego według miary.
Z twierdzenia~\ref{thm:cauchy-measure-complete-9}
istnieje mierzalna $f$ taka, że $f_n\to f$
według miary. Twierdzenie~\ref{thm:measure-metric-9}
daje $\rho([f_n],[f])\to0$.
\end{proof}

\begin{uwaga}[Skończoność miary jest istotna dla tego wzoru]
Jeśli $\mu(X)=\infty$, wartość
$\int_X\min\{1,|f-g|\}\,\mathrm d\mu(x)$
może być nieskończona, więc nie musi określać
metryki o wartościach skończonych.
Na przestrzeniach $\sigma$-skończonych można budować
metryki z wagami po wyczerpaniu przestrzeni zbiorami
o skończonej mierze, ale taka metryka na ogół
opisuje zbieżność \emph{lokalną} według miary,
nie globalną zbieżność w znaczeniu
\eqref{eq:conv-in-measure-9}.
\end{uwaga}


\begin{przyklad}[Zbieżność lokalna według miary nie jest globalną]
\label{ex:local-vs-global-9}
Na $\mathbb R$ z miarą Lebesgue'a rozważmy
$f_n=\chi_{[n,n+1)}$ i $f=0$. Dla każdego
$R>0$ mamy $f_n|_{[-R,R]}=0$ dla dostatecznie
dużych $n$; funkcje zbiegają zatem według miary na
każdym zbiorze $[-R,R]$. Jednak dla $0<\varepsilon<1$
zachodzi $\lambda_1(\{|f_n-f|>\varepsilon\})=1$
dla wszystkich $n$. Nie ma globalnej zbieżności
według miary. Ten przykład objaśnia różnicę między
metryzacją zbieżności lokalnej na przestrzeni
$\sigma$-skończonej i metryką \eqref{eq:measure-metric-9}
na przestrzeni o mierze skończonej.
\end{przyklad}

\begin{twierdzenie}[Kontrola ucieczki masy przy zbieżności lokalnej]
\label{thm:local-to-global-measure-9}
Niech $(X,\mathcal F,\mu)$ będzie przestrzenią z miarą,
a $(X_j)_{j\ge1}$ ciągiem zbiorów mierzalnych.
Załóżmy, że $X_j\in\mathcal F$, $X_j\subseteq X_{j+1}$,
$\bigcup_jX_j=X$ i $\mu(X_j)<\infty$.
Niech $f_n,f:X\to\mathbb R$ będą mierzalne.
Wówczas $f_n\to f$ według miary na całym $X$
wtedy i tylko wtedy, gdy dla każdego $\varepsilon>0$:
\begin{enumerate}[label=\textup{(\roman*)}]
\item $\mu(X_j\cap\{|f_n-f|>\varepsilon\})\to0$
dla każdego ustalonego $j$;
\item $\displaystyle \lim_{j\to\infty}\limsup_{n\to\infty}
\mu((X\setminus X_j)\cap\{|f_n-f|>\varepsilon\})=0$.
\end{enumerate}
\end{twierdzenie}
\begin{proof}
Oznaczmy $B_n=\{|f_n-f|>\varepsilon\}$.
Jeśli $\mu(B_n)\to0$, oba warunki są natychmiastowe.
Odwrotnie, dla każdego $j,n$ z podaddytywności
\[
\mu(B_n)\le\mu(X_j\cap B_n)+
\mu((X\setminus X_j)\cap B_n).
\]
Przechodząc do $\limsup$ po $n$ i korzystając
z (i), dostajemy
$\limsup_n\mu(B_n)\le
\limsup_n\mu((X\setminus X_j)\cap B_n)$.
Na mocy (ii) prawa strona dąży do zera wraz
z $j\to\infty$, co dowodzi tezy.
\end{proof}

\section{Użyteczne kryteria graniczne}

\begin{twierdzenie}[Lemat Scheffégo]
\label{thm:scheffe-9}
Niech $(X,\mathcal F,\mu)$ będzie przestrzenią z miarą,
a $f_n,f:X\to[0,\infty)$ funkcjami mierzalnymi,
całkowalne oraz $f_n\to f$ prawie wszędzie.
Jeżeli
$\int_Xf_n(x)\,\mathrm d\mu(x)\to
\int_Xf(x)\,\mathrm d\mu(x)<\infty$,
to
$\int_X|f_n(x)-f(x)|\,\mathrm d\mu(x)\to0$.
\end{twierdzenie}
\begin{proof}
Mamy punktową tożsamość
$|f_n-f|=f_n+f-2\min(f_n,f)$.
Ponieważ $0\le\min(f_n,f)\le f$ oraz
$\min(f_n,f)\to f$ prawie wszędzie,
twierdzenie o zbieżności zmajoryzowanej daje
$\int_X\min(f_n,f)\,\mathrm d\mu(x)\to\int_Xf\,\mathrm d\mu(x)$.
Całkując tożsamość i korzystając z założenia
o granicach całek, otrzymujemy tezę.
\end{proof}

\begin{wniosek}[Wersja dla gęstości prawdopodobieństwa]\label{idx:09:wniosek:14}
Niech $(X,\mathcal F,\mu)$ będzie przestrzenią z miarą,
a $f_n,f:X\to[0,\infty)$ funkcjami mierzalnymi.
Jeżeli
$\int_X f_n\,\mathrm d\mu(x)=
\int_X f\,\mathrm d\mu(x)=1$ oraz $f_n\to f$ prawie
wszędzie, to $\int_X|f_n-f|\,\mathrm d\mu(x)\to0$.
W zastosowaniach do prawdopodobieństwa jest to
podstawowy sposób przejścia od zbieżności gęstości
prawie wszędzie do zbieżności w wariacji całkowitej.
\end{wniosek}

\begin{uwaga}[Co naprawdę wynika z poszczególnych pojęć?]
\label{rem:conv-summary-9}
Na przestrzeni skończonej miary zbieżność prawie
wszędzie pociąga zbieżność według miary; na dowolnej
przestrzeni zbieżność według miary zapewnia podciąg
zbieżny prawie wszędzie. Ani zbieżność według miary,
ani w $L^p$ nie muszą zapewniać zbieżności prawie
wszędzie całego ciągu. Sama zbieżność prawie wszędzie
nie zapewnia zbieżności całek. Dodatkową kontrolę
zapewniają dominacja, monotoniczność lub jednostajna
całkowalność.
\end{uwaga}

\begin{center}
\begin{tikzpicture}[>=stealth, every node/.style={font=\small},
box/.style={draw,rounded corners,align=center,inner sep=5pt}]
\node[box] (u) at (0,2.5) {jednostajnie};
\node[box] (au) at (4.4,2.5) {prawie jednostajnie};
\node[box] (ae) at (4.4,0.7) {prawie wszędzie};
\node[box] (lp) at (-0.1,-1) {w $L^p$, $1\le p<\infty$};
\node[box] (mu) at (4.4,-1) {według miary};
\draw[->] (u) -- (au);
\draw[->] (au) -- (ae);
\draw[->] (au.east) to[out=0,in=0] (mu.east);
\draw[->] (lp) -- (mu);
\draw[->,dashed] (ae) -- node[right,font=\scriptsize]{gdy $\mu(X)<\infty$} (mu);
\draw[->,dashed,bend left=24] (ae.west) to node[above,font=\scriptsize]{Jegorow: $\mu(X)<\infty$} (au.south west);
\end{tikzpicture}
\par\smallskip\textit{Schemat implikacji. Strzałki przerywane wymagają skończoności miary.}
\end{center}

% ROZDZIAŁ 10 — do włączenia w main.tex
\chapter[Regularność funkcji mierzalnych]{Regularność funkcji mierzalnych i aproksymacja funkcjami ciągłymi}
\label{ch:regularnosc-funkcji}

\section{Co oznacza regularność funkcji mierzalnej?}

Funkcja mierzalna może być nieciągła we wszystkich punktach,
jak funkcja Dirichleta. Mimo to mierzalność jest blisko
związana z ciągłością \emph{w sensie miary}: na zbiorze o
skończonej mierze można usunąć zbiór o dowolnie małej mierze
i otrzymać ciągłość na pozostałym odpowiednio dobranym
zbiorze zwartym. Taką treść ma twierdzenie Łuzina.

W odróżnieniu od twierdzenia Jegorowa, które dotyczy
zbieżności \emph{ciągu} funkcji, twierdzenie Łuzina
przybliża pojedynczą funkcję mierzalną przez funkcje
ciągłe. Oba twierdzenia współpracują: najpierw
aproksymujemy funkcję mierzalną funkcjami prostymi,
a następnie wymuszamy jednostajność zbieżności
poza niewielkim zbiorem wyjątkowym.

W całym rozdziale $d\ge1$ i $\lambda_d$ oznacza miarę
Lebesgue'a na $\mathbb R^d$. Przez $\mathcal L_d$
oznaczamy jej $\sigma$-algebrę zbiorów mierzalnych.
Topologię na podzbiorach $\mathbb R^d$ rozumiemy jako
topologię podprzestrzeni. W szczególności sformułowanie
``$f|_K$ jest ciągła'' \emph{nie} oznacza,
że $f$ jest ciągła w punktach $K$ jako funkcja na
całej dziedzinie.

\begin{uwaga}[Dwa różne symbole]
Funkcję charakterystyczną zbioru $A$ oznaczamy
wyłącznie przez $\chi_A$. Gładkie funkcje odcinające
oznaczamy literą $\vartheta$ lub $\psi$; nie są one
funkcjami charakterystycznymi i przyjmują na ogół
również wartości pośrednie między zerem a jedynką.
\end{uwaga}

\section{Dwa narzędzia: regularność miary i funkcje odcinające}

\begin{lemat}[Przybliżenie mierzalnego zbioru przez zbiór zwarty]
\label{lem:compact-inside-10}
Niech $A\in\mathcal L_d$ oraz $\lambda_d(A)<\infty$.
Dla każdego $\varepsilon>0$ istnieją zbiór
zwarty $K$ i zbiór otwarty $O$ w $\mathbb R^d$ takie,
że $K\subseteq A\subseteq O$ i
\begin{equation}\label{eq:compact-open-10}
\lambda_d(O\setminus K)<\varepsilon.
\end{equation}
\end{lemat}
\begin{proof}
Regularność wewnętrzna miary Lebesgue'a daje
zwarty $K\subseteq A$ z
$\lambda_d(A\setminus K)<\varepsilon/2$.
Regularność zewnętrzna daje otwarty $O\supseteq A$
z $\lambda_d(O\setminus A)<\varepsilon/2$.
Ponieważ $O\setminus K$ jest rozłączną sumą
$(O\setminus A)\cup(A\setminus K)$,
addytywność miary daje \eqref{eq:compact-open-10}.
Jeśli chcemy jedynie otrzymać otwarte otoczenie
zbioru zwartego $K$ o małej różnicy miar,
możemy je dodatkowo wybrać ograniczone:
wystarczy przeciąć $O$ z kulą zawierającą $K$.
\end{proof}

\begin{uwaga}[Niedopuszczalna zamiana kwantyfikatorów]
Zbiór zwarty $K$ zależy od $\varepsilon$.
Na ogół nie istnieje jeden zbiór zwarty zawierający
wszystkie poza mierzalnym zbiorem zerowym punkty
$A$ (w szczególności gdy $A$ nie jest ograniczony).
Nie wolno również zastępować założenia
$\lambda_d(A)<\infty$ samym założeniem
$A\in\mathcal L_d$ w żądaniu
$\lambda_d(A\setminus K)<\varepsilon$.
\end{uwaga}

\begin{lemat}[Gładka funkcja odcinająca]
\label{lem:smooth-cutoff-10}
Niech $K\subseteq\mathbb R^d$ będzie zwarty,
a $O\subseteq\mathbb R^d$ otwarty i zawierający $K$.
Istnieje funkcja $\vartheta\in C_c^\infty(\mathbb R^d)$
taka, że $0\le\vartheta\le1$, $\vartheta=1$ na $K$
oraz $\operatorname{supp}\vartheta\subseteq O$.
Przez $\operatorname{supp}\vartheta$ rozumiemy domknięcie
zbioru $\{\vartheta\ne0\}$, a $C_c^\infty$ oznacza
funkcje gładkie o zwartym nośniku.
\end{lemat}
\begin{proof}
Jeśli $K=\varnothing$, bierzemy $\vartheta=0$.
W przeciwnym razie dla każdego $a\in K$ wybieramy
$r_a>0$ tak, żeby $\overline{B(a,r_a)}\subseteq O$.
Zwartość $K$ pozwala wybrać skończenie wiele
kul $B(a_j,r_j/2)$ pokrywających $K$.

Niech $\eta:\mathbb R\to[0,\infty)$ będzie określona
wzorem $\eta(t)=e^{-1/t}$ dla $t>0$ i $\eta(t)=0$
dla $t\le0$. Wzór ten określa funkcję klasy
$C^\infty$; każda jej pochodna dla $t>0$ jest
sumą wyrazów $P(1/t)e^{-1/t}$ i ma granicę zero
przy $t\downarrow0$, co dowodzi gładkości także
w zerze. Zdefiniujmy
\[
\psi_j(x)=
\frac{\eta(r_j^2-|x-a_j|^2)}
{\eta(r_j^2-|x-a_j|^2)+
 \eta(|x-a_j|^2-r_j^2/4)}.
\]
Mianownik jest dodatni dla każdego $x$:
dla $|x-a_j|<r_j$ dodatni jest pierwszy składnik,
a dla $|x-a_j|\ge r_j$ drugi. Widzimy, że
$\psi_j=1$ na $\overline{B(a_j,r_j/2)}$,
$\psi_j=0$ poza $B(a_j,r_j)$ i
$0\le\psi_j\le1$.
Kładziemy $\vartheta=1-\prod_j(1-\psi_j)$.
Jest gładka, równa $1$ na $K$ i ma nośnik zawarty
w skończonej sumie zwartych kul
$\overline{B(a_j,r_j)}\subseteq O$.
\end{proof}

\begin{wniosek}[Przybliżanie funkcji charakterystycznej]
\label{cor:indicator-smooth-10}
Niech $A\in\mathcal L_d$ i $\lambda_d(A)<\infty$.
Dla każdego $1\le p<\infty$ i $\varepsilon>0$
istnieje $\vartheta\in C_c^\infty(\mathbb R^d)$ taka,
że $0\le\vartheta\le1$ oraz
\begin{equation}\label{eq:indicator-smooth-10}
\int_{\mathbb R^d}|\chi_A(x)-\vartheta(x)|^p
\,\mathrm d\lambda_d(x)<\varepsilon.
\end{equation}
\end{wniosek}
\begin{proof}
Z lematu \ref{lem:compact-inside-10} wybieramy
$K\subseteq A\subseteq O$ z
$\lambda_d(O\setminus K)<\varepsilon$.
Stosujemy lemat \ref{lem:smooth-cutoff-10}.
Na $K$ funkcje $\chi_A$ i $\vartheta$ są równe $1$,
a poza $O$ są równe $0$. Na pozostałym zbiorze
$O\setminus K$ ich różnica ma moduł co najwyżej $1$.
Stąd całka po $\mathbb R^d$ jest ograniczona przez
$\lambda_d(O\setminus K)<\varepsilon$.
\end{proof}

\begin{center}
\begin{tikzpicture}[scale=.9,every node/.style={font=\small}]
\draw[rounded corners=10pt,thick,fill=blue!7,draw=blue!65!black]
(0,0) .. controls (-.7,1) and (-.2,2.3) .. (1.7,2.4)
.. controls (3.6,2.9) and (6,2.15) .. (6.4,1.25)
.. controls (6.9,-.3) and (3.6,-.5) .. (0,0);
\draw[rounded corners=8pt,thick,fill=green!10,draw=green!55!black]
(.6,.35) .. controls (.4,1.2) and (1,2.1) .. (2.05,2.05)
.. controls (4.2,2.5) and (5.55,1.85) .. (5.75,1.2)
.. controls (6,.35) and (2.5,.05) .. (.6,.35);
\draw[rounded corners=8pt,thick,fill=orange!18,draw=orange!75!black]
(1.25,.7) .. controls (1.1,1.25) and (1.55,1.8) .. (2.4,1.72)
.. controls (3.6,1.98) and (4.9,1.58) .. (5.02,1.04)
.. controls (4.5,.45) and (2.2,.4) .. (1.25,.7);
\node at (3.2,1.2) {$K$};
\node at (5.45,1.75) {$A$};
\node at (6.06,2.10) {$O$};
\end{tikzpicture}
\par\smallskip
\textit{Schemat regularności miary: $K\subseteq A\subseteq O$,
przy czym $\lambda_d(O\setminus K)$ jest dowolnie mała.
Rysunek ma charakter poglądowy.}
\end{center}

\section{Funkcje proste ciągłe na dużym zbiorze zwartym}

\begin{lemat}[Łuzin dla funkcji prostych]
\label{lem:lusin-simple-10}
Niech $E\in\mathcal L_d$, $\lambda_d(E)<\infty$,
a $s:E\to\mathbb R$ będzie mierzalną funkcją prostą.
Dla każdego $\varepsilon>0$ istnieje zwarty
$K\subseteq E$ taki, że
$\lambda_d(E\setminus K)<\varepsilon$ oraz
$s|_K$ jest ciągła.
\end{lemat}
\begin{proof}
Jeżeli $E=\varnothing$, bierzemy $K=\varnothing$;
ograniczenie $s|_K$ jest ciągłe, a
$\lambda_d(E\setminus K)=0<\varepsilon$.
Dalej zakładamy $E\ne\varnothing$, zatem liczba
różnych wartości $m$ jest dodatnia.
Niech $a_1,\dots,a_m$ będą różnymi wartościami $s$
i $A_j=\{x\in E:s(x)=a_j\}$.
Są to parami rozłączne zbiory mierzalne, których
suma jest równa $E$. Wszystkie mają skończoną miarę.
Dla każdego $j$ regularność wewnętrzna daje
zwarty $K_j\subseteq A_j$ taki, że
$\lambda_d(A_j\setminus K_j)<\varepsilon/m$.
Zbiór $K=\bigcup_{j=1}^mK_j$ jest zwarty,
$K\subseteq E$ i
$\lambda_d(E\setminus K)<\varepsilon$.

Ponieważ zbiory $K_j$ są parami rozłączne i zwarte,
każdy z nich jest otwarty i domknięty w podprzestrzeni
$K$: jego dopełnienie w $K$ jest skończoną sumą
zbiorów zwartych, a więc domkniętych.
Funkcja $s|_K$, stała na każdym z tych zbiorów,
jest zatem ciągła. Nie wyciągamy błędnego wniosku,
że $s$ jest ciągła na całym $E$.
\end{proof}

\section{Twierdzenie Łuzina}

\begin{twierdzenie}[Łuzin]
\label{thm:lusin-10}
Niech $E\in\mathcal L_d$ będzie zbiorem o
$\lambda_d(E)<\infty$, a $f:E\to\mathbb R$
funkcją mierzalną względem śladu
$\mathcal L_d|_E$. Wtedy dla każdego
$\varepsilon>0$ istnieje zwarty zbiór
$K\subseteq E$ taki, że
\begin{equation}\label{eq:lusin-measure-10}
\lambda_d(E\setminus K)<\varepsilon
\end{equation}
oraz $f|_K:K\to\mathbb R$ jest ciągła.
\end{twierdzenie}
\begin{proof}
\textit{Krok 1: aproksymacja funkcją prostą.}
Na mocy twierdzenia o aproksymacji funkcji
mierzalnych z rozdziału 6 istnieją mierzalne
funkcje proste $s_n:E\to\mathbb R$ takie, że
$s_n(x)\to f(x)$ dla każdego $x\in E$.
Możemy na przykład stosować obcięcia wartości
$f$ i ich kwantyzację na przedziałach długości
$2^{-n}$, najpierw do części dodatniej i ujemnej.
Nie korzystamy w tym kroku z twierdzenia Łuzina.

\textit{Krok 2: prawie jednostajność.}
Ponieważ $\lambda_d(E)<\infty$, twierdzenie
Jegorowa \ref{thm:egorov-9}, zastosowane do
przestrzeni z miarą $E$, daje mierzalny zbiór
$Z\subseteq E$ taki, że
$\lambda_d(Z)<\varepsilon/4$ i $s_n\to f$
jednostajnie na $E\setminus Z$.
Z regularności wewnętrznej wybieramy zwarty
$K_0\subseteq E\setminus Z$, dla którego
$\lambda_d((E\setminus Z)\setminus K_0)<\varepsilon/4$.
Na $K_0$ zbieżność jest jednostajna.

\textit{Krok 3: ciągłość wszystkich przybliżeń.}
Dla każdego $n\ge1$ stosujemy lemat
\ref{lem:lusin-simple-10} do $s_n$ na $E$
z tolerancją $\varepsilon 2^{-(n+2)}$.
Otrzymujemy zwarty $K_n\subseteq E$ taki, że
$\lambda_d(E\setminus K_n)<
\varepsilon2^{-(n+2)}$ i $s_n|_{K_n}$ jest ciągła.
Kładziemy
$K=K_0\cap\bigcap_{n\ge1}K_n$.
Jest to zbiór zwarty jako domknięty podzbiór
zbioru zwartego $K_0$. Ponadto
\[
\lambda_d(E\setminus K)
\le\lambda_d(Z)+
\lambda_d((E\setminus Z)\setminus K_0)+
\sum_{n\ge1}\lambda_d(E\setminus K_n)
<\frac\varepsilon4+\frac\varepsilon4+
\frac\varepsilon4<\varepsilon.
\]
Dla każdego $n$ funkcja $s_n|_K$ jest ciągła,
a ciąg $(s_n|_K)$ zbiega jednostajnie do $f|_K$.
Jednostajna granica funkcji ciągłych na przestrzeni
metrycznej jest ciągła; stąd $f|_K$ jest ciągła.
\end{proof}

\begin{wniosek}[Funkcje skończone prawie wszędzie]
\label{cor:lusin-finite-ae-10}
Niech $E\in\mathcal L_d$ ma skończoną miarę,
a $f:E\to\overline{\mathbb R}$ jest mierzalna względem
śladu $\mathcal L_d|_E$ i skończona $\lambda_d$-prawie
wszędzie. Dla każdego $\varepsilon>0$ istnieje
zwarty $K\subseteq E$ taki, że
$\lambda_d(E\setminus K)<\varepsilon$ oraz
$f|_K:K\to\mathbb R$ jest ciągła.
\end{wniosek}
\begin{proof}
Zbiór $N=\{x\in E:|f(x)|=\infty\}$ jest
mierzalny i zerowy. Na $E\setminus N$ funkcja
jest rzeczywista i mierzalna. Stosujemy twierdzenie
Łuzina do $E\setminus N$, którego miara jest
równa $\lambda_d(E)$. Otrzymany zwarty
$K\subseteq E\setminus N$ ma żądane własności.
\end{proof}

\begin{wniosek}[Jednoczesne twierdzenie Łuzina dla przeliczalnej rodziny]
\label{cor:countable-lusin-10}
Niech $E\in\mathcal L_d$ ma skończoną miarę,
a $f_j:E\to\mathbb R$, $j\ge1$, będą
funkcjami mierzalnymi. Dla każdego $\varepsilon>0$
istnieje zwarty $K\subseteq E$ taki, że
$\lambda_d(E\setminus K)<\varepsilon$ i
\emph{wszystkie} ograniczenia $f_j|_K$ są ciągłe.
\end{wniosek}
\begin{proof}
Dla każdego $j$ stosujemy twierdzenie Łuzina
z tolerancją $\varepsilon2^{-(j+1)}$.
Otrzymujemy zwarte $K_j\subseteq E$ takie,
że $\lambda_d(E\setminus K_j)<\varepsilon2^{-(j+1)}$
i $f_j|_{K_j}$ jest ciągła.
Zbiór $K=\bigcap_{j\ge1}K_j$ jest zwarty
(jako domknięty podzbiór $K_1$), a
\[
\lambda_d(E\setminus K)
\le\sum_{j\ge1}\lambda_d(E\setminus K_j)
<\varepsilon.
\]
Każda $f_j|_K$ jest ciągła jako ograniczenie
funkcji ciągłej na $K_j$.
\end{proof}

\begin{uwaga}[Granice twierdzenia]
\label{rem:lusin-limitations-10}
Twierdzenie Łuzina nie orzeka, że $f$ jest
ciągła \emph{w punktach} $K$ względem całej
dziedziny $E$. Na przykład $f=\chi_{\mathbb Q}$
na $[0,1]$ jest nieciągła w każdym punkcie, ale
jej ograniczenie do dowolnego zbioru zwartego
niezawierającego liczb wymiernych jest funkcją
stałą, a więc ciągłą. Możemy znaleźć takie
zbiory zwarte o mierze dowolnie bliskiej $1$.

Skończoność miary $E$ jest istotna w podanej formie.
Dla $E=\mathbb R^d$ żaden zwarty podzbiór $K$
nie spełnia $\lambda_d(E\setminus K)<1$,
niezależnie od wyboru funkcji $f$.
\end{uwaga}

\section{Przedłużanie funkcji ciągłej ze zbioru zwartego}

\begin{lemat}[Lipschitzowska aproksymacja na zbiorze zwartym]
\label{lem:lipschitz-approx-compact-10}
Niech $K\subseteq\mathbb R^d$ będzie niepustym
zbiorem zwartym, a $h:K\to\mathbb R$ funkcją ciągłą.
Dla każdego $\varepsilon>0$ istnieje ograniczona
funkcja Lipschitzowska $H:\mathbb R^d\to\mathbb R$
taka, że
$\sup_{x\in K}|H(x)-h(x)|<\varepsilon$
i $\sup_{\mathbb R^d}|H|\le\sup_K|h|$.
\end{lemat}
\begin{proof}
Niech $M=\sup_K|h|<\infty$. Jeśli $M=0$,
wystarczy $H=0$. W przeciwnym razie z jednostajnej
ciągłości $h$ na $K$ wybieramy $\delta>0$ takie, że
$|h(x)-h(y)|<\varepsilon/2$ przy
$x,y\in K$, $|x-y|<\delta$.
Wybieramy $L>2M/\delta$ i definiujemy na
całym $\mathbb R^d$
\[
U_L(x)=\inf_{y\in K}\bigl(h(y)+L|x-y|\bigr).
\]
Dla dowolnych $x,z\in\mathbb R^d$ z nierówności
$|x-y|\le|x-z|+|z-y|$ wynika
$U_L(x)\le U_L(z)+L|x-z|$. Zamieniając $x,z$,
dostajemy $|U_L(x)-U_L(z)|\le L|x-z|$.

Dla $x\in K$ mamy $U_L(x)\le h(x)$ (wybierając
$y=x$). Jeśli $|y-x|<\delta$, to
$h(y)+L|x-y|>h(x)-\varepsilon/2$.
Jeżeli $|y-x|\ge\delta$, to
$h(y)+L|x-y|\ge-M+L\delta>M\ge h(x)$.
Wynika $h(x)-\varepsilon/2\le U_L(x)\le h(x)$.
Funkcja
$H(x)=\max(-M,\min(M,U_L(x)))$
jest ograniczona i Lipschitzowska, a obcięcie do
$[-M,M]$ nie zwiększa błędu na $K$.
\end{proof}

\begin{twierdzenie}[Przedłużenie z kompaktu w $\mathbb R^d$]
\label{thm:compact-extension-10}
Jeśli $K\subseteq\mathbb R^d$ jest zwarty,
a $h:K\to\mathbb R$ ciągła, to istnieje
$H\in C(\mathbb R^d)$ taka, że $H|_K=h$.
Można dodatkowo wymagać
$\sup_{\mathbb R^d}|H|\le\sup_K|h|$
dla $K\ne\varnothing$. Jeśli $K=\varnothing$,
wybieramy $H=0$; nie używamy wtedy supremum
rzeczywistych wartości funkcji na zbiorze pustym.
\end{twierdzenie}
\begin{proof}
Przy $K=\varnothing$ bierzemy $H=0$.
Niech $K\ne\varnothing$ i $M_0=\sup_K|h|$.
Konstruujemy funkcje Lipschitzowskie $u_n$ i
reszty $r_n:K\to\mathbb R$ rekurencyjnie.
Kładziemy $r_0=h$. Jeśli
$M_n=\sup_K|r_n|>0$, stosujemy lemat
\ref{lem:lipschitz-approx-compact-10} z błędem
$M_n/2$, uzyskując $u_n$ o
$\sup_{\mathbb R^d}|u_n|\le M_n$ i
$\sup_K|u_n-r_n|<M_n/2$.
Definiujemy $r_{n+1}=r_n-u_n|_K$.
Wtedy $M_{n+1}\le M_n/2$. Jeśli $M_n=0$,
bierzemy $u_n=0$ i wszystkie dalsze reszty
równe zero.

Z oszacowania
$\sup_{\mathbb R^d}|u_n|\le M_0 2^{-n}$
wynika jednostajna zbieżność szeregu
$G=\sum_{n\ge0}u_n$ na całym $\mathbb R^d$.
Jest to funkcja ciągła. Na $K$ mamy dla każdego
$N$ tożsamość
$h=\sum_{n=0}^Nu_n|_K+r_{N+1}$;
po przejściu do granicy otrzymujemy $G|_K=h$.
Na koniec funkcja $H=\max(-M_0,\min(M_0,G))$
pozostaje ciągła, zgadza się z $h$ na $K$
i ma żądane ograniczenie.
\end{proof}

\begin{wniosek}[Wersja Łuzina przez funkcję ciągłą na całej przestrzeni]
\label{cor:lusin-extension-10}
Niech $E\in\mathcal L_d$, $\lambda_d(E)<\infty$,
i $f:E\to\mathbb R$ będzie mierzalna. Dla każdego
$\varepsilon>0$ istnieją $K\subseteq E$ zwarty
i $H\in C(\mathbb R^d)$ takie, że
$\lambda_d(E\setminus K)<\varepsilon$ oraz
$H(x)=f(x)$ dla wszystkich $x\in K$.
\end{wniosek}
\begin{proof}
Stosujemy twierdzenie \ref{thm:lusin-10},
a następnie twierdzenie \ref{thm:compact-extension-10}
do funkcji ciągłej $f|_K$.
\end{proof}

\section{Aproksymacja w mierze przez funkcje ciągłe i gładkie}

\begin{twierdzenie}[Aproksymacja funkcji o nośniku mierzalnym skończonej miary]
\label{thm:cc-approx-in-measure-10}
Niech $E\in\mathcal L_d$, $\lambda_d(E)<\infty$,
i niech $f:E\to\mathbb R$ będzie mierzalna.
Przedłużmy $f$ zerem poza $E$, otrzymując
$\widetilde f:\mathbb R^d\to\mathbb R$.
Dla każdych $\varepsilon,\delta>0$ istnieje
$\varphi\in C_c(\mathbb R^d)$ takie, że
\begin{equation}\label{eq:cc-approx-measure-10}
\lambda_d\bigl(\{x\in\mathbb R^d:
|\widetilde f(x)-\varphi(x)|>\delta\}\bigr)
<\varepsilon.
\end{equation}
Można dodatkowo wybrać $\varphi$ tak, aby
$\varphi=\widetilde f$ \emph{dokładnie} na pewnym
zbiorze zwartym $K\subseteq E$ z
$\lambda_d(E\setminus K)<\varepsilon/2$.
\end{twierdzenie}
\begin{proof}
Z twierdzenia Łuzina wybieramy zwarty
$K\subseteq E$ taki, że
$\lambda_d(E\setminus K)<\varepsilon/2$
oraz $f|_K$ jest ciągła.
Z regularności zewnętrznej \emph{zbioru zwartego}
$K$ wybieramy otwarty $O$ zawierający $K$ tak,
że $\lambda_d(O\setminus K)<\varepsilon/2$.
Nie musimy wymagać, aby $O$ zawierał całe $E$.

Przedłużamy $f|_K$ do funkcji ciągłej
$H:\mathbb R^d\to\mathbb R$ na mocy
\ref{thm:compact-extension-10}. Na mocy lematu
\ref{lem:smooth-cutoff-10} istnieje
$\vartheta\in C_c^\infty(\mathbb R^d)$, $\vartheta=1$ na $K$,
z $\operatorname{supp}\vartheta\subseteq O$.
Kładziemy $\varphi=\vartheta H$. Wtedy
$\varphi\in C_c(\mathbb R^d)$ i na $K$ zachodzi
$\varphi=f$. Poza $E\cup O$ obie funkcje
$\varphi$ i $\widetilde f$ są zerowe. Stąd
zbiór, na którym funkcje mogą się różnić,
jest zawarty w $(E\setminus K)\cup(O\setminus K)$,
a jego miara jest mniejsza niż $\varepsilon$.
Implikuje to \eqref{eq:cc-approx-measure-10}
dla każdego $\delta>0$.
\end{proof}

\begin{uwaga}[Dlaczego nie istnieje takie twierdzenie dla dowolnego $f$ na $\mathbb R^d$?]
Dla funkcji $f\equiv1$ na $\mathbb R^d$ i dowolnej
$\varphi\in C_c(\mathbb R^d)$ zbiór
$\{|f-\varphi|>1/2\}$ zawiera dopełnienie
zwartego nośnika $\varphi$, a więc ma nieskończoną
miarę Lebesgue'a. Założenie o skończonej mierze
zbioru podtrzymującego funkcję nie jest czysto
technicznym dodatkiem.
\end{uwaga}

\begin{lemat}[Wygładzanie funkcji ciągłej o zwartym nośniku]
\label{lem:mollify-cc-10}
Niech $h\in C_c(\mathbb R^d)$. Istnieje ciąg
$h_n\in C_c^\infty(\mathbb R^d)$ taki, że
$\sup_{x\in\mathbb R^d}|h_n(x)-h(x)|\to0$.
\end{lemat}
\begin{proof}
Z lematu \ref{lem:smooth-cutoff-10} otrzymujemy
nieujemną, niezerową funkcję
$\rho\in C_c^\infty(\mathbb R^d)$ o nośniku
zawartym w kuli $B(0,1)$. Po podzieleniu przez
$\int_{\mathbb R^d}\rho(x)\,\mathrm d\lambda_d(x)>0$
możemy przyjąć, że ta całka jest równa $1$.
Dla $t>0$ określamy
$\rho_t(x)=t^{-d}\rho(x/t)$ oraz
\[
h_t(x)=\int_{\mathbb R^d}
 h(y)\rho_t(x-y)\,\mathrm d\lambda_d(y).
\]
Z niezmienniczości i skalowania miary Lebesgue'a
(rozdział 4) wynika
$\int_{\mathbb R^d}\rho_t(x-y)\,\mathrm d\lambda_d(y)=1$.

Funkcja $h_t$ ma nośnik zawarty w
$\operatorname{supp}h+\overline{B(0,t)}$,
który jest zwarty. Jest gładka: pochodne względem
$x$ dowolnego rzędu można przenosić pod znak
całki dzięki ograniczoności pochodnych $\rho_t$
i temu, że integrujemy po ustalonym zwartym
nośniku $h$. Formalnie stosujemy iteracyjnie
twierdzenie o różniczkowaniu pod znakiem całki
z rozdziału 8, osobno do każdej współrzędnej.

Ponadto dla każdego $x$ mamy
\[
|h_t(x)-h(x)|
\le\int_{\mathbb R^d}|h(y)-h(x)|\rho_t(x-y)
\,\mathrm d\lambda_d(y)
\le\sup_{|x-y|\le t}|h(y)-h(x)|.
\]
Funkcja $h$ jest jednostajnie ciągła na
$\mathbb R^d$, bo jest ciągła i ma zwarty nośnik.
Ostatnie supremum dąży jednostajnie do zera dla
$t\downarrow0$. Wybieramy $h_n=h_{1/n}$.
\end{proof}

\begin{wniosek}[Aproksymacja w mierze funkcjami gładkimi]
\label{cor:smooth-approx-in-measure-10}
Niech $E\in\mathcal L_d$ ma skończoną miarę
i niech $f:E\to\mathbb R$ będzie mierzalna
względem śladu. Oznaczmy przez
$\widetilde f:\mathbb R^d\to\mathbb R$
przedłużenie funkcji $f$ zerem poza $E$.
Dla dowolnych $\varepsilon,\delta>0$ istnieje
$\psi\in C_c^\infty(\mathbb R^d)$ taka, że
\[
\lambda_d\bigl(\{|\widetilde f-\psi|>\delta\}\bigr)
<\varepsilon.
\]
\end{wniosek}
\begin{proof}
Z twierdzenia \ref{thm:cc-approx-in-measure-10}
wybieramy $\varphi\in C_c(\mathbb R^d)$ taką,
że $\lambda_d(\{|\widetilde f-\varphi|>
\delta/2\})<\varepsilon$.
Z lematu \ref{lem:mollify-cc-10} wybieramy
$\psi\in C_c^\infty$ taką, że
$\sup|\psi-\varphi|<\delta/2$.
Wtedy
$\{|\widetilde f-\psi|>\delta\}\subseteq
\{|\widetilde f-\varphi|>\delta/2\}$.
\end{proof}

\section{Gęstość $C_c^\infty$ w przestrzeniach $L^p$}

\begin{twierdzenie}[Gęstość funkcji gładkich]
\label{thm:cc-in-lp-10}
Niech $d\ge1$ i $1\le p<\infty$.
Dla każdej funkcji mierzalnej
$f:\mathbb R^d\to\mathbb R$ spełniającej
$\int_{\mathbb R^d}|f(x)|^p\,\mathrm d\lambda_d(x)
<\infty$ oraz dla każdego $\varepsilon>0$
istnieje $\varphi\in C_c^\infty(\mathbb R^d)$ taka, że
\begin{equation}\label{eq:lp-smooth-dense-10}
\int_{\mathbb R^d}|f(x)-\varphi(x)|^p
\,\mathrm d\lambda_d(x)<\varepsilon.
\end{equation}
W szczególności $C_c^\infty(\mathbb R^d)$ jest
gęsta w $L^p(\mathbb R^d,\lambda_d)$.
\end{twierdzenie}
\begin{proof}
\textit{Krok 1: obcięcie nośnika i wartości.}
Niech $B_R=B(0,R)$. Kładziemy
\[
f_{R,M}(x)=\max\{-M,\min\{f(x),M\}\}\chi_{B_R}(x).
\]
Przy $R,M\to\infty$ (na przykład po przekątnej
$R=M=n$) mamy $f_{R,M}(x)\to f(x)$ dla prawie
wszystkich $x$ oraz
$|f(x)-f_{R,M}(x)|^p\le 2^p|f(x)|^p$.
Twierdzenie o zbieżności zmajoryzowanej daje
$\int_{\mathbb R^d}|f-f_{R,M}|^p\,\mathrm d\lambda_d(x)\to0$. Możemy zatem przybliżać
$f$ funkcjami ograniczonymi, których nośniki są
zawarte w kulach o skończonej mierze.

\textit{Krok 2: przybliżenie funkcją prostą.}
Ustalmy ograniczoną mierzalną funkcję $g$
zerową poza kulą $B_R$. Dla $h>0$ przyjmijmy
$s_h(x)=h\lfloor g(x)/h\rfloor$ na $B_R$ i
$s_h(x)=0$ poza $B_R$. Funkcja $s_h$ jest prosta,
ponieważ $g$ jest ograniczona, a jej poziomice
są mierzalne. Na $B_R$ mamy $|s_h-g|\le h$,
więc
$\int_{\mathbb R^d}|s_h-g|^p\,\mathrm d\lambda_d(x)
\le h^p\lambda_d(B_R)\to0$.
Po ewentualnym usunięciu zerowych współczynników
zapisujemy $s_h=\sum_{j=1}^m a_j\chi_{A_j}$,
gdzie wszystkie $A_j$ są mierzalne i mają skończoną
miarę (jako podzbiory $B_R$).

\textit{Krok 3: przybliżenie funkcji prostej.}
Dla każdego $j$ stosujemy wniosek
\ref{cor:indicator-smooth-10}, aby znaleźć
$\vartheta_j\in C_c^\infty$ z dowolnie małą całką
$\int_{\mathbb R^d}|\chi_{A_j}-\vartheta_j|^p\,\mathrm d\lambda_d(x)$.
Kładziemy $\varphi=\sum_{j=1}^m a_j\vartheta_j$.
Jest to funkcja gładka o zwartym nośniku.
Z nierówności elementarnej
$|z_1+\dots+z_m|^p\le m^{p-1}
\sum_{j=1}^m|z_j|^p$ dla $p\ge1$ wynika
\[
\int_{\mathbb R^d}|s_h-\varphi|^p
\,\mathrm d\lambda_d(x)
\le m^{p-1}\sum_{j=1}^m|a_j|^p
\int_{\mathbb R^d}|\chi_{A_j}-\vartheta_j|^p
\,\mathrm d\lambda_d(x).
\]
Prawą stronę możemy uczynić dowolnie małą.

\textit{Krok 4: połączenie oszacowań.}
Używamy trzykrotnie nierówności
$|u+v|^p\le 2^{p-1}(|u|^p+|v|^p)$,
aby oszacować $|f-\varphi|^p$ przez stałą
razy sumę błędów z poprzednich kroków.
Wybierając kolejno $R,M$, następnie $h$,
a na końcu $\vartheta_j$, możemy uczynić
$\int_{\mathbb R^d}|f-\varphi|^p\,\mathrm d\lambda_d(x)$ mniejszą od zadanego
$\varepsilon$. Nie korzystamy z nierówności
Minkowskiego, której pełny dowód będzie
przedmiotem rozdziału 14.
\end{proof}

\begin{przyklad}[Dlaczego nie można zastąpić $p<\infty$ przez $p=\infty$]
Niech $f=\chi_{(0,\infty)}$ na $\mathbb R$.
Nie istnieje funkcja ciągła $g:\mathbb R\to\mathbb R$
taka, że $\operatorname*{ess\,sup}_{x\in\mathbb R}
|f(x)-g(x)|<1/2$. W przeciwnym razie dla prawie
wszystkich $x<0$ zachodziłoby $g(x)<1/2$,
a dla prawie wszystkich $x>0$ zachodziłoby
$g(x)>1/2$. Z ciągłości otrzymalibyśmy
$g(0)\le1/2$ i $g(0)\ge1/2$; przy ścisłym
oszacowaniu błędu $q<1/2$ dostalibyśmy nawet
$g(0)\le q<1-q\le g(0)$, sprzeczność.
Zatem funkcje ciągłe nie są gęste w $L^\infty$
w normie supremum istotnego.
\end{przyklad}

\section{Granice prawie wszędzie funkcji gładkich}

\begin{twierdzenie}[Prawie wszędzie gładka aproksymacja funkcji mierzalnej]
\label{thm:ae-smooth-approx-10}
Dla dowolnej funkcji mierzalnej
$f:\mathbb R^d\to\mathbb R$ istnieje ciąg
$(\varphi_n)\subseteq C_c^\infty(\mathbb R^d)$
taki, że $\varphi_n(x)\to f(x)$ dla
$\lambda_d$-prawie wszystkich $x\in\mathbb R^d$.
Nie zakładamy, że $f$ jest całkowalna.
\end{twierdzenie}
\begin{proof}
Niech $B_n=B(0,n)$. Stosujemy wniosek
\ref{cor:smooth-approx-in-measure-10} do ograniczenia
$f|_{B_n}$, tolerancji błędu $1/n$ i miary
wyjątkowego zbioru $2^{-n}$. Otrzymujemy
$\varphi_n\in C_c^\infty$ takie, że
\[
\lambda_d\bigl(B_n\cap
\{|f-\varphi_n|>1/n\}\bigr)<2^{-n}.
\]
Niech $A_n=B_n\cap\{|f-\varphi_n|>1/n\}$.
Zbiór $\limsup_n A_n$ ma miarę zero na mocy
lematu Borela--Cantellego. Dla $x$ poza nim
i dostatecznie dużego $n$ zachodzi $x\in B_n$
oraz $|f(x)-\varphi_n(x)|\le1/n$.
Stąd $\varphi_n(x)\to f(x)$.
\end{proof}

\begin{wniosek}[Borelowska reprezentacja funkcji Lebesgue'a]
\label{cor:borel-representative-10}
Jeżeli $f:\mathbb R^d\to\mathbb R$ jest
$\mathcal L_d$-mierzalna, to istnieje funkcja
$g:\mathbb R^d\to\mathbb R$ borelowsko mierzalna
taka, że $f=g$ $\lambda_d$-prawie wszędzie.
\end{wniosek}
\begin{proof}
Niech $(\varphi_n)$ będzie ciągiem z twierdzenia
\ref{thm:ae-smooth-approx-10}.
Zbiór $D$, na którym $\varphi_n(x)$ zbiega
do skończonej liczby, jest borelowski, ponieważ
istnienie granicy skończonej można wyrazić przez
warunek Cauchy'ego z przeliczalnymi kwantyfikatorami.
Definiujemy $g(x)=\lim_n\varphi_n(x)$ dla $x\in D$
i $g(x)=0$ dla $x\notin D$.
Jest to funkcja borelowska i zgadza się z $f$
poza zbiorem miary zero.
\end{proof}

\begin{uwaga}[Granica punktowa wszędzie a granica prawie wszędzie]
Nie każda funkcja mierzalna Lebesgue'a jest granicą
punktową \emph{wszędzie} funkcji ciągłych.
Istnieją funkcje Lebesgue'a, które nie są nawet
borelowsko mierzalne, natomiast każda granica
punktowa wszędzie funkcji borelowskich jest
borelowska. Wnioski
\ref{thm:ae-smooth-approx-10} i
\ref{cor:borel-representative-10} celowo mówią
o równości prawie wszędzie.
\end{uwaga}


\begin{uwaga}[Przedstawiciele borelowscy a zbiory wyjątkowe]
\label{rem:borel-null-representatives-10}
Wniosek~\ref{cor:borel-representative-10} zapewnia
istnienie funkcji borelowskiej $g$ równej $f$
$\lambda_d$-prawie wszędzie. Nie zapewnia natomiast,
że zbiór $\{f\ne g\}$ jest borelowski: jeśli $f$
jest jedynie mierzalna w sensie Lebesgue'a,
zbiór ten może być nieborelowski. Istnieje za to
\emph{borelowski zbiór zerowy} $N$ zawierający
$\{f\ne g\}$, ponieważ miara Lebesgue'a jest
uzupełnieniem miary borelowskiej i każdy mierzalny
zbiór zerowy ma borelowską otoczkę miary zero.
To rozróżnienie będzie potrzebne przy porównywaniu
miary Lebesgue'a z produktem miar w rozdziale
\ref{ch:produkty-miary}.
\end{uwaga}

\section{Podsumowanie i zastosowania}

Z regularności miary Lebesgue'a wynika, że zbiory
mierzalne o skończonej mierze można przybliżać
zbiorami zwartymi i otwartymi. W połączeniu z
twierdzeniem Jegorowa i aproksymacją funkcjami
prostymi prowadzi to do twierdzenia Łuzina.
Pozwala ono przedłużyć funkcję z dużego zbioru
zwartego do funkcji ciągłej na całej przestrzeni.
Z kolei przybliżanie funkcji charakterystycznych
przez gładkie funkcje odcinające daje gęstość
$C_c^\infty$ w $L^p$ dla $1\le p<\infty$.
Te konstrukcje przygotowują późniejsze rozdziały
o całkach produktowych, funkcjonałach na $L^p$,
miarach Radona oraz słabych pochodnych.

% ROZDZIAŁ 11 — do włączenia w main.tex
\chapter{Całki zależne od parametru}
\label{ch:parametry-calki}

\section{Problem i jego dokładne sformułowanie}

Niech $(X,\mathcal F,\mu)$ będzie przestrzenią z miarą, $T$ przestrzenią
parametrów, a $f:T\times X\to\mathbb C$ rodziną funkcji. Zajmujemy się
funkcją $F:T\to\mathbb C$ określoną wzorem
$F(t)=\int_X f(t,x)\,\mathrm d\mu(x)$, o ile $f(t,\cdot)$ jest całkowalna.
Pytamy, kiedy własności funkcji $t\mapsto f(t,x)$ przechodzą na $F$:
ciągłość, różniczkowalność, wyższa regularność i analityczność. Problem
polega na tym, że własność zachodząca osobno dla każdego $x$ nie
kontroluje automatycznie masy rozłożonej na całym $X$.

\begin{przyklad}[Koncentracja masy przy zerze]
\label{ex:params-escape-11}
Na $X=(0,\infty)$ z miarą Lebesgue'a połóżmy dla $t>0$
$f(t,x)=t\chi_{(0,1/t)}(x)$, a $f(0,x)=0$.
Wówczas $f(t,x)\to0$ dla każdego $x>0$ przy $t\downarrow0$,
ale $\int_X f(t,x)\,\mathrm d\lambda_1(x)=1$ dla każdego $t>0$.
Zatem $F$ nie jest ciągła w zerze, mimo ciągłości $t\mapsto f(t,x)$
w punkcie $0$ dla każdego $x$. Nie istnieje też jedna funkcja
$g\in L^1(X)$ dominująca wszystkie $f(t,\cdot)$ dla dostatecznie
małych $t>0$: w przeciwnym razie twierdzenie o zbieżności
zmajoryzowanej dawałoby sprzeczność.
\end{przyklad}

\begin{uwaga}[Zakres rozdziału]
W rozdziale 8 udowodniono podstawowe twierdzenia o całkach z parametrem.
Tutaj formułujemy ich wersje \emph{lokalne} i \emph{jednostajne},
wyprowadzamy twierdzenia wyższego rzędu i badamy kilka użytecznych
transformacji. Nie używamy twierdzenia Fubiniego, które zostanie
udowodnione dopiero w rozdziale~\ref{ch:produkty-miary}.
\end{uwaga}

\section{Dominacja lokalna i ciągłość w parametrze}

\begin{definicja}[Lokalna dominacja całkowalna]
\label{def:local-majorant-11}
Niech $(X,\mathcal F,\mu)$ będzie przestrzenią z miarą,
$(T,d)$ przestrzenią metryczną, $t_0\in T$ oraz
$f:T\times X\to\mathbb C$. Mówimy, że rodzina $f$ jest
\emph{lokalnie zdominowana całkowalnie w $t_0$}, jeżeli istnieją
otoczenie $U$ punktu $t_0$ w $T$ oraz mierzalna funkcja
$g:X\to[0,\infty)$ z $\int_Xg\,\mathrm d\mu(x)<\infty$ takie, że
istnieje $N\in\mathcal F$, $\mu(N)=0$, taki, że
$|f(t,x)|\le g(x)$ dla wszystkich $t\in U$ i $x\in X\setminus N$.
Warunek ten odróżniamy od słabszego wymagania, by dla każdego $t$
istniał osobny zbiór zerowy $N_t$. W twierdzeniu o ciągłości
słabszy wariant również wystarcza: dla ustalonego ciągu
$t_n\to t_0$ zbiór $N_{t_0}\cup\bigcup_n N_{t_n}$
pozostaje zerowy, toteż nadal można użyć twierdzenia
o zbieżności zmajoryzowanej. W pozostałych twierdzeniach
sprawdzamy potrzebny wariant oddzielnie.
\end{definicja}

\begin{twierdzenie}[Lokalne kryterium ciągłości]
\label{thm:local-continuity-11}
Niech $(X,\mathcal F,\mu)$ będzie przestrzenią z miarą,
$(T,d)$ przestrzenią metryczną, $t_0\in T$,
a $f:T\times X\to\mathbb C$ spełnia następujące warunki:
\begin{enumerate}[label=\textup{(\roman*)}]
\item dla każdego $t\in T$ funkcja $x\mapsto f(t,x)$ jest mierzalna;
\item dla $\mu$-prawie każdego $x$ odwzorowanie $t\mapsto f(t,x)$
      jest ciągłe w $t_0$;
\item istnieje całkowalna dominacja w pewnym otoczeniu $t_0$
      w sensie definicji~\ref{def:local-majorant-11}.
\end{enumerate}
Wtedy $F(t)=\int_X f(t,x)\,\mathrm d\mu(x)$ jest określona
w pewnym otoczeniu $t_0$ i ciągła w $t_0$.
Jeśli warunki te zachodzą w każdym $t_0\in T$, to $F$ jest
ciągła na $T$.
\end{twierdzenie}
\begin{proof}
Ustalmy otoczenie $U$ i funkcję $g$ z (iii). Dla $t\in U$
mamy $\int_X|f(t,x)|\,\mathrm d\mu(x)\le\int_Xg\,\mathrm d\mu(x)$,
więc $F(t)$ istnieje jako całka bezwzględnie zbieżna. Niech
$t_n\to t_0$. Dla dużych $n$ mamy $t_n\in U$, a funkcje
$f(t_n,\cdot)$ zbiegają do $f(t_0,\cdot)$ prawie wszędzie i
są zdominowane przez $g$. Twierdzenie o zbieżności zmajoryzowanej,
zastosowane do części rzeczywistej i urojonej lub bezpośrednio
do modułu różnicy, daje
$\int_X |f(t_n,x)-f(t_0,x)|\,\mathrm d\mu(x)\to0$.
W konsekwencji $|F(t_n)-F(t_0)|\to0$. Ciągowa charakteryzacja
ciągłości w przestrzeni metrycznej kończy dowód.
\end{proof}

\begin{wniosek}[Ciągłość w normie $L^1$]
\label{cor:l1-param-11}
Niech $(X,\mathcal F,\mu)$ będzie przestrzenią z miarą,
$(T,d)$ przestrzenią metryczną, $t_0\in T$,
a $f:T\times X\to\mathbb C$ rodziną o mierzalnych
przekrojach. Załóżmy, że $t\mapsto f(t,x)$ jest ciągła
w $t_0$ dla $\mu$-prawie każdego $x$ oraz że istnieją
otoczenie $U$ punktu $t_0$, mierzalny zbiór zerowy
$N\in\mathcal F$ i nieujemna całkowalna funkcja
$g:X\to[0,\infty)$ takie, że
$|f(t,x)|\le g(x)$ dla wszystkich $t\in U$
i $x\in X\setminus N$. Wtedy
\[
\int_X|f(t,x)-f(t_0,x)|\,\mathrm d\mu(x)\longrightarrow0
\qquad(t\to t_0).
\]
W szczególności odwzorowanie
$U\ni t\mapsto[f(t,\cdot)]\in L^1(\mu)$ jest
ciągłe w $t_0$.
\end{wniosek>

\begin{przyklad}[Brak potrzeby globalnej dominacji]
Niech $f(t,x)=e^{-tx}$ dla $x>0$ i $t>0$. Jeśli $t_0>0$,
to dla $t\in(t_0/2,3t_0/2)$ zachodzi
$0\le f(t,x)\le e^{-t_0x/2}$, a funkcja po prawej stronie
jest całkowalna na $(0,\infty)$. Zatem
$F(t)=\int_{(0,\infty)}e^{-tx}\,\mathrm d\lambda_1(x)=1/t$
jest ciągła na $(0,\infty)$. Nie istnieje jednak całkowalna
dominanta dla całej rodziny $\{f(t,\cdot):t>0\}$, gdyż
$\sup_{t>0} e^{-tx}=1$ dla każdego $x>0$.
\end{przyklad}

\section{Jednostajność zbieżności całek w parametrze}

\begin{twierdzenie}[Zbieżność jednostajna po scałkowaniu]
\label{thm:uniform-parameter-11}
Niech $(X,\mathcal F,\mu)$ będzie przestrzenią z miarą,
$T$ niepustym zbiorem, a $f_n,f:T\times X\to\mathbb C$
rodzinami funkcji o mierzalnych przekrojach.
Załóżmy, że istnieją mierzalne funkcje
$h_n,g:X\to[0,\infty)$, przy czym $g$ jest całkowalna, takie, że
$|f_n(t,x)-f(t,x)|\le h_n(x)\le g(x)$
dla wszystkich $t\in T$, wszystkich $n$ oraz prawie wszystkich
$x$ poza wspólnym zbiorem zerowym, a także $h_n(x)\to0$
prawie wszędzie. Załóżmy dodatkowo, że dla pewnego
$t_*\in T$ funkcja $f(t_*,\cdot)$ jest całkowalna oraz że
istnieje mierzalna funkcja całkowalna
$g_0:X\to[0,\infty)$ taka, że
$|f(t,x)-f(t_*,x)|\le g_0(x)$ dla każdego $t\in T$
i dla $\mu$-prawie każdego $x$ (zbiór wyjątkowy
może zależeć od $t$). Wtedy
\[
\sup_{t\in T}\left|\int_X f_n(t,x)\,\mathrm d\mu(x)
-\int_Xf(t,x)\,\mathrm d\mu(x)\right|\longrightarrow0.
\]
\end{twierdzenie}
\begin{proof}
Dla każdego $t\in T$ z nierówności dla całki wynika
$|\int_X(f_n-f)(t,x)\,\mathrm d\mu(x)|\le\int_Xh_n\,\mathrm d\mu(x)$.
Oszacowanie po prawej stronie nie zależy od $t$.
Twierdzenie Lebesgue'a o zbieżności zmajoryzowanej daje
$\int_Xh_n\,\mathrm d\mu(x)\to0$. Przechodząc do supremum
po $t$, otrzymujemy tezę.
\end{proof}

\begin{uwaga}[Mierzalność supremum nie jest automatyczna]
Jeśli chciałoby się wybrać $h_n(x)=\sup_{t\in T}|f_n(t,x)-f(t,x)|$,
trzeba najpierw sprawdzić mierzalność tej funkcji. Wystarcza
na przykład, aby $T$ było separowalną przestrzenią metryczną,
a odwzorowania $t\mapsto f_n(t,x)-f(t,x)$ były ciągłe dla
wszystkich $x$: wtedy supremum po $T$ równa się supremum
po ustalonym przeliczalnym podzbiorze gęstym. Jeśli nie
mamy tego warunku, lepiej założyć istnienie mierzalnej
majoranty $h_n$ bez odwoływania się do supremum.
\end{uwaga}

\section{Różniczkowanie pod znakiem całki}

\begin{twierdzenie}[Pierwsza pochodna — wersja lokalna]
\label{thm:local-differentiation-11}
Niech $(X,\mathcal F,\mu)$ będzie przestrzenią z miarą,
$I\subseteq\mathbb R$ przedziałem otwartym,
$t_0\in I$, a $f:I\times X\to\mathbb C$. Załóżmy, że:
\begin{enumerate}[label=\textup{(\roman*)}]
\item dla każdego $t\in I$ przekrój $f(t,\cdot)$ jest mierzalny;
\item dla prawie wszystkich $x$ funkcja $t\mapsto f(t,x)$ jest
      różniczkowalna na $I$;
\item istnieje $r>0$ z $[t_0-r,t_0+r]\subset I$ i funkcja
      $g\in L^1(\mu)$, $g\ge0$, taka, że
      $|\partial_t f(t,x)|\le g(x)$ dla wszystkich
      $t\in[t_0-r,t_0+r]$ poza wspólnym zbiorem zerowym;
\item funkcja $f(t_0,\cdot)$ jest całkowalna.
\end{enumerate}
Wówczas w otoczeniu $t_0$ funkcja
$F(t)=\int_Xf(t,x)\,\mathrm d\mu(x)$ jest dobrze określona,
różniczkowalna w $t_0$ i spełnia
\begin{equation}\label{eq:diff-param-11}
F'(t_0)=\int_X\partial_t f(t_0,x)\,\mathrm d\mu(x).
\end{equation}
Na wspólnym mierzalnym zbiorze wyjątkowym, na którym
nie zakładamy różniczkowalności, przyjmujemy
$\partial_t f(t_0,x)=0$.
\end{twierdzenie}
\begin{proof}
Poza wspólnym zbiorem zerowym ustalmy $x$ oraz $t$
i zastosujmy twierdzenie o wartości średniej do
rzeczywistej funkcji
$s\mapsto\operatorname{Re}(\overline{u}f(s,x))$,
gdzie $u\in\mathbb C$, $|u|=1$. Jej pochodna
ma moduł co najwyżej $g(x)$, zatem
$|\operatorname{Re}(\overline{u}(f(t,x)-f(t_0,x)))|
\le |t-t_0|g(x)$. Biorąc supremum po
$|u|=1$, otrzymujemy dokładnie
$|f(t,x)-f(t_0,x)|\le |t-t_0|g(x)$.
Argument nie wymaga ciągłości pochodnej ani
nieuzasadnionego zastosowania wzoru Newtona--Leibniza.
Stąd $f(t,\cdot)$ jest całkowalna dla $|t-t_0|<r$.
Dla $0<|h|<r$ iloraz
$q_h(x)=(f(t_0+h,x)-f(t_0,x))/h$ jest mierzalny
jako różnica funkcji mierzalnych; poza zbiorem zerowym
$q_h(x)\to\partial_tf(t_0,x)$ przy $h\to0$,
a $|q_h(x)|\le g(x)$. Pochodna $x\mapsto\partial_tf(t_0,x)$
jest mierzalna jako granica ilorazów różnicowych, na przykład
dla ciągu $h=1/n$ dostatecznie dużych $n$.
Twierdzenie o zbieżności zmajoryzowanej daje
$\int_Xq_h(x)\,\mathrm d\mu(x)\to
\int_X\partial_tf(t_0,x)\,\mathrm d\mu(x)$,
gdy $h\to0$. Z liniowości całki
$\int_Xq_h\,\mathrm d\mu(x)=(F(t_0+h)-F(t_0))/h$,
co jest dokładnie definicją pochodnej $F$.
\end{proof}

\begin{uwaga}[Dlaczego samo istnienie pochodnej nie wystarcza]
Dla $f(t,x)=t\chi_{(0,1/t)}(x)$ z pierwszego przykładu
nie można nawet otrzymać ciągłości całki w $t=0$.
Dominacja \emph{ilorazów różnicowych} lub lokalna dominacja
pochodnej są niezależnymi, istotnymi warunkami.
\end{uwaga}

\begin{twierdzenie}[Pochodne wyższego rzędu]
\label{thm:higher-differentiation-11}
Niech $(X,\mathcal F,\mu)$ będzie przestrzenią z miarą,
$I\subseteq\mathbb R$ otwartym przedziałem,
$m\ge1$ liczbą naturalną, a $f:I\times X\to\mathbb C$.
Załóżmy, że dla każdego $t\in I$ funkcja $f(t,\cdot)$ jest mierzalna,
a dla prawie wszystkich $x$ odwzorowanie $t\mapsto f(t,x)$
ma pochodne do rzędu $m$ włącznie, wszystkie ciągłe na $I$.
Załóżmy ponadto, że dla każdego zwartego przedziału
$J\subset I$ oraz każdego $1\le k\le m$
istnieje nieujemna mierzalna funkcja całkowalna
$g_{J,k}:X\to[0,\infty)$ taka, że
$|\partial_t^k f(t,x)|\le g_{J,k}(x)$ dla $t\in J$
poza wspólnym zbiorem zerowym, a $f(t_*,\cdot)$ jest
całkowalna dla pewnego $t_*\in I$.
Wtedy $F(t)=\int_Xf(t,x)\,\mathrm d\mu(x)$ należy
do $C^m(I)$ oraz dla wszystkich $t\in I$ i $1\le k\le m$
\begin{equation}\label{eq:higher-param-11}
F^{(k)}(t)=\int_X\partial_t^k f(t,x)\,\mathrm d\mu(x).
\end{equation}
Wszystkie pochodne podcałkowe określamy jako zero na
wspólnym mierzalnym zbiorze wyjątkowym, na którym
nie założono istnienia pochodnych.
\end{twierdzenie}
\begin{proof}
Dla $t\in I$ odcinek o końcach $t_*$ i $t$ jest zwarty
i zawiera się w $I$. Oszacowanie pierwszej pochodnej
zapewnia $|f(t,x)|\le|f(t_*,x)|+|t-t_*|g_{J,1}(x)$,
więc $f(t,\cdot)$ jest całkowalna. Dla ustalonego
$t\in I$ zastosujmy twierdzenie~\ref{thm:local-differentiation-11}
w małym otoczeniu $t$: otrzymujemy istnienie $F'$ i wzór
z $k=1$. Gdy $1\le k<m$, funkcja $\partial_t^kf(t,\cdot)$
jest mierzalna jako punktowa granica ilorazów różnicowych.
Dzięki majorantom $g_{J,k}$ jest całkowalna. Stosując
twierdzenie~\ref{thm:local-differentiation-11} do
$\partial_t^kf$, otrzymujemy pochodną $(k+1)$-szego rzędu.
Ciągłość każdej funkcji $t\mapsto\int_X\partial_t^kf(t,x)\,\mathrm d\mu(x)$
wynika z twierdzenia~\ref{thm:local-continuity-11},
gdyż $\partial_t^kf(t,x)$ jest punktowo ciągła w $t$
i lokalnie zdominowana. Zatem $F\in C^m(I)$.
\end{proof}

\begin{przyklad}[Całka Laplace'a i pochodne]
\label{ex:laplace-11}
Dla $t>0$ kładziemy $F(t)=\int_{(0,\infty)}e^{-tx}\,\mathrm d\lambda_1(x)$.
Mamy $F(t)=1/t$. Ponieważ
$\partial_t^k e^{-tx}=(-x)^k e^{-tx}$, dla
$t$ w zwartym podprzedziale $(0,\infty)$ funkcja
$x^k e^{-tx}$ jest zdominowana przez
$x^ke^{-ax}$, gdzie $a>0$ zależy od podprzedziału.
Ta ostatnia funkcja jest całkowalna, więc
\[
\int_{(0,\infty)}x^ke^{-tx}\,\mathrm d\lambda_1(x)
=\frac{k!}{t^{k+1}}\qquad(k\in\mathbb N_0).
\]
W szczególności wzór na momenty rozkładu wykładniczego
wynika z różniczkowania pod znakiem całki, a nie z
wielokrotnego całkowania przez części.
\end{przyklad}

\section{Zbieżność według miary zamiast dominacji punktowej}

\begin{twierdzenie}[Kryterium Vitaliego dla parametru]
\label{thm:param-vitali-11}
Niech $(X,\mathcal F,\mu)$ będzie przestrzenią z miarą
skończoną, $(T,d)$ przestrzenią metryczną, $t_0\in T$,
a $f:T\times X\to\mathbb C$ rodziną funkcji o mierzalnych
przekrojach. Załóżmy, że istnieje otoczenie $U$ punktu
$t_0$ takie, że rodzina modułów
$\{|f(t,\cdot)|:t\in U\}$ jest jednostajnie całkowalna
w sensie Definicji~\ref{def:uniform-integrability}
oraz dla każdego $\varepsilon>0$ zachodzi
\[
\mu(\{x\in X:|f(t,x)-f(t_0,x)|>\varepsilon\})
\longrightarrow0\qquad(t\to t_0).
\]
Jest to naturalne rozszerzenie zbieżności według miary
z Definicji~\ref{def:convergence-measure} na funkcje
zespolone. Wówczas
$\|f(t,\cdot)-f(t_0,\cdot)\|_{L^1(\mu)}\to0$,
a więc funkcja $t\mapsto\int_X f(t,x)\,\mathrm d\mu(x)$
jest ciągła w $t_0$.
\end{twierdzenie}
\begin{proof}
Załóżmy przeciwnie, że granica norm nie jest zerem.
Istnieją wówczas $\varepsilon>0$ i ciąg $t_n\to t_0$,
$t_n\in U$, dla którego
$\|f(t_n,\cdot)-f(t_0,\cdot)\|_1\ge\varepsilon$.
Funkcje $f(t_n,\cdot)$ zbiegają według miary do
$f(t_0,\cdot)$, a rodziny ich części rzeczywistych i urojonych
są jednostajnie całkowalne. Twierdzenie Vitaliego
z rozdziału~\ref{ch:twierdzenia-graniczne}
(dowodzone osobno dla części rzeczywistych i urojonych)
implikuje zbieżność w $L^1$, sprzeczną z wybraną
nierównością. W szczególności wszystkie całki
są określone i ich ciągłość wynika z nierówności
$|\int_X(f(t,x)-f(t_0,x))\,\mathrm d\mu(x)|
\le\|f(t,\cdot)-f(t_0,\cdot)\|_1$.
\end{proof}

\begin{uwaga}
Jednostajna całkowalność nie wymaga wspólnej
całkowalnej dominanty. W przeciwieństwie do twierdzenia
o dominacji lokalnej założenie o zbieżności jest tutaj
formułowane według miary, a nie punktowo. Dla $\mu(X)=\infty$
potrzebna jest dodatkowa kontrola ogonów; sam warunek
jednostajnej całkowalności na małych zbiorach nie wystarcza.
\end{uwaga}

\section{Zmienna dziedzina całkowania}

\begin{twierdzenie}[Całkowanie po zależnej od parametru dziedzinie]
\label{thm:moving-domains-11}
Niech $(X,\mathcal F,\mu)$ będzie przestrzenią z miarą,
$(T,d)$ przestrzenią metryczną, $t_0\in T$,
a $A_t\in\mathcal F$ dla każdego $t\in T$.
Niech $u:T\times X\to\mathbb C$
ma mierzalne przekroje $u(t,\cdot)$ i załóżmy, że dla prawie
wszystkich $x$ zachodzi
$\chi_{A_t}(x)u(t,x)\to\chi_{A_{t_0}}(x)u(t_0,x)$
dla $t\to t_0$. Jeżeli w pewnym otoczeniu $U$ punktu $t_0$
istnieje nieujemna funkcja $g\in L^1(\mu)$ taka, że
\emph{dla każdego} $t\in U$ nierówność
$|\chi_{A_t}(x)u(t,x)|\le g(x)$
zachodzi dla $\mu$-prawie każdego $x$ (zbiór wyjątkowy
może zależeć od $t$), to
$t\mapsto\int_{A_t}u(t,x)\,\mathrm d\mu(x)$
jest ciągła w $t_0$.
\end{twierdzenie}
\begin{proof}
Definiujemy na całym $X$ funkcję
$f(t,x)=u(t,x)$ dla $x\in A_t$ i $f(t,x)=0$
dla $x\notin A_t$. Jest ona mierzalna w $x$ dla każdego $t$;
wyrażenia $\chi_{A_t}u$ należy tak interpretować,
aby uniknąć nieoznaczonych iloczynów $0\cdot\infty$.
Ustalmy ciąg $t_n\to t_0$. Po pominięciu skończenie
wielu wyrazów możemy założyć $t_n\in U$. Dla każdego
$n\ge0$ wybierzmy zbiór zerowy $N_n$, poza którym
$|f(t_n,x)|\le g(x)$; w tym zapisie $t_0$ odpowiada
$n=0$. Zbiór $N=\bigcup_{n\ge0}N_n$ ma miarę zero.
Poza $N$ oraz zbiorem wyjątkowym z założenia o zbieżności
punktowej zachodzi jednocześnie $f(t_n,x)\to f(t_0,x)$
i $|f(t_n,x)|\le g(x)$ dla wszystkich $n\ge0$.
Twierdzenie Lebesgue'a o zbieżności zmajoryzowanej
zapewnia więc zbieżność całek wzdłuż każdego ciągu
$t_n\to t_0$. Ciągowa charakteryzacja ciągłości
w przestrzeni metrycznej kończy dowód. W ten sposób
nie użyliśmy silniejszego, niewymaganego założenia
o wspólnym zbiorze wyjątkowym dla całej rodziny parametrów.
\end{proof}

\begin{przyklad}[Przedziały ruchome i brak dominacji]
Niech $u(t,x)=1$ na $(0,\infty)$ i $A_t=(0,t)$ dla $t>0$.
Dla $t$ w otoczeniu $t_0>0$ wszystkie $A_t$ mieszczą się
w $(0,2t_0)$; funkcja $g=\chi_{(0,2t_0)}$ jest
całkowalną dominantą, a dla każdego $x\ne t_0$
zachodzi $\chi_{A_t}(x)\to\chi_{A_{t_0}}(x)$.
Stąd $\int_{A_t}1\,\mathrm d\lambda_1(x)=t$ zależy ciągle od $t$.
Ten sam argument przestaje działać, jeśli dziedziny rozrastają
się tak, że masa ucieka poza każdy zbiór skończonej miary.
\end{przyklad}

\section{Całki niewłaściwe i jednostajna kontrola ogonów}

\begin{twierdzenie}[Kryterium jednostajnej kontroli ogonów]
\label{thm:tails-11}
Niech $t_0\in T$, gdzie $(T,d)$ jest przestrzenią metryczną,
a $(X,\mathcal F,\mu)$ przestrzenią z miarą. Załóżmy, że
istnieje wyczerpanie $X_1\subseteq X_2\subseteq\dots$,
$\bigcup_{n\ge1}X_n=X$, zbiory $X_n\in\mathcal F$
o skończonej mierze, oraz $f:T\times X\to\mathbb C$,
dla której $f(t,\cdot)$ jest całkowalna na każdym $X_n$.
Załóżmy ponadto, że $t\mapsto\int_{X_n}f(t,x)\,\mathrm d\mu(x)$
jest ciągła w $t_0$ dla każdego $n$, oraz że istnieje
otoczenie $U$ punktu $t_0$ takie, że
\begin{equation}\label{eq:uniform-tails-11}
\lim_{n\to\infty}\sup_{t\in U}
\int_{X\setminus X_n}|f(t,x)|\,\mathrm d\mu(x)=0.
\end{equation}
Wtedy wszystkie $f(t,\cdot)$ dla $t\in U$ są całkowalne,
a $t\mapsto\int_Xf(t,x)\,\mathrm d\mu(x)$ jest ciągła w $t_0$.
\end{twierdzenie}
\begin{proof}
Warunek \eqref{eq:uniform-tails-11} gwarantuje, że dla
pewnego $n$ i wszystkich $t\in U$ całki modułów po
$X\setminus X_n$ są skończone; ich suma z całką po $X_n$
jest więc skończona. Ustalmy $\varepsilon>0$ i wybierzmy
$n$ tak duże, aby supremum w
\eqref{eq:uniform-tails-11} było mniejsze niż $\varepsilon/3$.
Dla $t\in U$ mamy
\begin{align*}
\left|\int_Xf(t,x)\,\mathrm d\mu(x)-\int_Xf(t_0,x)\,\mathrm d\mu(x)\right|
&\le\left|\int_{X_n}(f(t,x)-f(t_0,x))\,\mathrm d\mu(x)\right|\\
&\quad+\int_{X\setminus X_n}|f(t,x)|\,\mathrm d\mu(x)
+\int_{X\setminus X_n}|f(t_0,x)|\,\mathrm d\mu(x).
\end{align*}
Dwa ostatnie składniki są mniejsze niż $2\varepsilon/3$.
Pierwszy jest mniejszy niż $\varepsilon/3$, gdy $t$
jest dostatecznie blisko $t_0$ dzięki ciągłości
całki na $X_n$. To kończy dowód.
\end{proof}

\begin{uwaga}[Ograniczenie kryterium]
Warunek \eqref{eq:uniform-tails-11} jest wystarczający, lecz
niekonieczny dla ciągłości całki: przy funkcjach zmieniających
znak mogą występować znoszące się wkłady z ogona.
Kryterium jest użyteczne zwłaszcza dla całek niewłaściwych,
gdzie trudno znaleźć jedną majorantę całkowalną na całym $X$.
\end{uwaga}

\section{Transformata Laplace'a i funkcje holomorficzne}

\begin{twierdzenie}[Holomorficzność transformaty Laplace'a]
\label{thm:laplace-holo-11}
Niech $a\in\mathbb R$ i niech $h:(0,\infty)\to\mathbb C$
będzie mierzalna oraz
$\int_{(0,\infty)}e^{-\sigma x}|h(x)|\,\mathrm d\lambda_1(x)<\infty$
dla każdego $\sigma>a$. Dla $z\in\mathbb C$,
$\operatorname{Re}z>a$, określmy
$L_h(z)=\int_{(0,\infty)}e^{-zx}h(x)\,\mathrm d\lambda_1(x)$.
Wtedy $L_h$ jest holomorficzna w półpłaszczyźnie
$\{\operatorname{Re}z>a\}$ i dla każdego $k\in\mathbb N_0$
\begin{equation}\label{eq:laplace-derivative-11}
L_h^{(k)}(z)=\int_{(0,\infty)}(-x)^k e^{-zx}h(x)
\,\mathrm d\lambda_1(x).
\end{equation}
\end{twierdzenie}
\begin{proof}
Ustalmy $z_0$ z $\operatorname{Re}z_0>a$ i liczbę
$\sigma$ spełniającą $a<\sigma<\operatorname{Re}z_0$.
W pewnym otoczeniu $z_0$ zachodzi $\operatorname{Re}z>\sigma$.
Dla $k\ge0$ istnieje $C_{k,\sigma,z_0}<\infty$ taka, że
$x^k e^{-\operatorname{Re}z\,x}\le
C_{k,\sigma,z_0}e^{-\sigma x}$ w nieco mniejszym
otoczeniu $z_0$: wybieramy pośrednią liczbę
$\sigma'>\sigma$ z $\operatorname{Re}z>\sigma'$
i korzystamy z ograniczoności $x^ke^{-(\sigma'-\sigma)x}$.
Zatem całki po prawej stronie
\eqref{eq:laplace-derivative-11} są zbieżne bezwzględnie.

Aby udowodnić różniczkowalność zespoloną, dla małego
$w\in\mathbb C\setminus\{0\}$ rozważamy
\[
\frac{e^{-(z_0+w)x}-e^{-z_0x}}{w}
=-x\int_{[0,1]}e^{-(z_0+sw)x}\,\mathrm d\lambda_1(s).
\]
Tożsamość z prawą stroną jest klasyczną
całką po rzeczywistym odcinku $s\in[0,1]$,
a nie zamianą kolejności dwóch całkowań.
Moduł ilorazu jest ograniczony przez
$x e^{-\sigma'x}$ dla odpowiednio małych $w$,
co po pomnożeniu przez $|h(x)|$ jest zdominowane
przez stałą razy $e^{-\sigma x}|h(x)|$.
Dla każdego $x$ iloraz dąży do $-xe^{-z_0x}$.
Twierdzenie o zbieżności zmajoryzowanej daje
$L_h'(z_0)=-\int_{(0,\infty)}xe^{-z_0x}h(x)
\,\mathrm d\lambda_1(x)$. Powtarzając argument dla
$(-x)^kh(x)$, otrzymujemy wszystkie pochodne.
Ponieważ pochodna zespolona istnieje w każdym punkcie
półpłaszczyzny, $L_h$ jest holomorficzna.
\end{proof}

\begin{przyklad}[Funkcja Gamma]
\label{ex:gamma-param-11}
Dla $s>0$ definiujemy
$\Gamma(s)=\int_{(0,\infty)}x^{s-1}e^{-x}\,\mathrm d\lambda_1(x)$.
Całka jest skończona: przy zerze $x^{s-1}$ jest
całkowalne dla $s>0$, a w nieskończoności czynnik
$e^{-x}$ dominuje każdą potęgę $x$.
Dla każdego $s_0>0$ wybieramy $0<a<s_0<b$;
na $s\in[a,b]$ pochodne rzędu $k$ funkcji podcałkowej
wynoszą $x^{s-1}(\log x)^k e^{-x}$ i mają całkowalną
majorantę
\[
\chi_{(0,1]}(x)x^{a-1}|\log x|^k
+\chi_{(1,\infty)}(x)x^{b-1}(1+\log x)^k e^{-x}.
\]
W konsekwencji $\Gamma\in C^\infty((0,\infty))$ i
$\Gamma^{(k)}(s)=\int_{(0,\infty)}
 x^{s-1}(\log x)^k e^{-x}\,\mathrm d\lambda_1(x)$.
Całkując przez części na $[\varepsilon,R]$ i przechodząc
do granic $\varepsilon\downarrow0$, $R\uparrow\infty$,
dostajemy $\Gamma(s+1)=s\Gamma(s)$.
Ponadto $\Gamma(1)=1$, więc $\Gamma(n+1)=n!$ dla $n\ge0$.
\end{przyklad}

\section{Podsumowanie}
Dominacja lokalna pozwala przenosić ciągłość i
różniczkowalność pod znak całki bez narzucania
jednej globalnej majoranty. Jednostajna kontrola
błędów oraz ogonów pozwala badać rodziny funkcji,
dla których pojedyncza dominanta bywa niedostępna.
Wersja zespolona prowadzi do holomorficzności
transformaty Laplace'a i umożliwia obliczanie jej
pochodnych. W kolejnym rozdziale skonstruujemy
miary produktowe, dzięki którym można zamieniać
kolejność całkowania względem dwóch niezależnych
miar.

% ROZDZIAŁ 12 — do włączenia w main.tex
\chapter{Miary produktowe. Twierdzenia Tonellego i Fubiniego}
\label{ch:produkty-miary}

\section{Po co konstruujemy miarę produktową?}

Niech $(X,\mathcal F,\mu)$ i $(Y,\mathcal G,\nu)$ będą
przestrzeniami z miarami. Chcemy przypisać mierzalnemu
zbiorowi $E\subseteq X\times Y$ wielkość, która dla
prostokąta mierzalnego $A\times B$ wynosi $\mu(A)\nu(B)$.
Potrzebujemy następnie twierdzenia, które pozwoli obliczyć
całkę po $X\times Y$ jako dwie kolejne całki.

Już na etapie formułowania pojawiają się trzy pytania:
\emph{(i)} jaka jest odpowiednia $\sigma$-algebra
na $X\times Y$; \emph{(ii)} czy przypisanie
$A\times B\mapsto\mu(A)\nu(B)$ wyznacza miarę;
\emph{(iii)} kiedy można zmieniać kolejność całkowania.
Będziemy rozróżniać trzy odpowiedzi: konstrukcję miary
produktowej, twierdzenie Tonellego dla funkcji nieujemnych
i twierdzenie Fubiniego dla funkcji całkowalnych.

\section{Przestrzeń produktowa i przekroje}

Iloczyn $\sigma$-algebr $\mathcal F\otimes\mathcal G$
został zdefiniowany w Definicji~\ref{def:product-sigma}.
Zbiory $A\times B$ dla $A\in\mathcal F$ i $B\in\mathcal G$
nazywamy \emph{prostokątami mierzalnymi}.
Dla $E\subseteq X\times Y$ \emph{przekrojem pionowym}
przez $x\in X$ nazywamy $E_x=\{y\in Y:(x,y)\in E\}$,
a \emph{przekrojem poziomym} przez $y\in Y$ zbiór
$E^y=\{x\in X:(x,y)\in E\}$.
Dla funkcji $f:X\times Y\to\overline{\mathbb R}$ piszemy
$f_x(y)=f(x,y)$ i $f^y(x)=f(x,y)$.

\begin{lemat}[Mierzalność wszystkich przekrojów]
\label{lem:sections-measurable-12}
Niech $E\in\mathcal F\otimes\mathcal G$.
Wtedy $E_x\in\mathcal G$ dla każdego $x\in X$
oraz $E^y\in\mathcal F$ dla każdego $y\in Y$.
Jeżeli $f:(X\times Y,\mathcal F\otimes\mathcal G)
\to(\overline{\mathbb R},\mathcal B(\overline{\mathbb R}))$
jest mierzalna, to wszystkie funkcje $f_x$ oraz $f^y$
są mierzalne.
\end{lemat}
\begin{proof}
Ustalmy $x\in X$. Niech $\mathcal C_x$ oznacza rodzinę
tych zbiorów $E\subseteq X\times Y$, których przekrój
$E_x$ należy do $\mathcal G$. Ponieważ
$(X\times Y)_x=Y$, $(E^c)_x=Y\setminus E_x$ oraz
$(\bigcup_nE_n)_x=\bigcup_n(E_n)_x$, rodzina
$\mathcal C_x$ jest $\sigma$-algebrą. Dla prostokąta
$A\times B$ przekrój ma postać $B$, gdy $x\in A$,
oraz $\varnothing$, gdy $x\notin A$. Zatem
$\mathcal C_x$ zawiera wszystkie prostokąty i
$\mathcal F\otimes\mathcal G\subseteq\mathcal C_x$.
Argument dla $E^y$ jest symetryczny. Wreszcie
$\{y:f_x(y)\in H\}=(f^{-1}(H))_x$ dla każdego
zbioru borelowskiego $H\subseteq\overline{\mathbb R}$;
poprzednia część dowodu zapewnia jego mierzalność.
Analogicznie postępujemy z $f^y$.
\end{proof}

\begin{uwaga}[Nie odwracamy twierdzenia o przekrojach]
Z mierzalności wszystkich przekrojów nie wynika
mierzalność zbioru w $\mathcal F\otimes\mathcal G$.
Kontrprzykład z podzbiorem przekątnej podamy w dalszej
części rozdziału. Mierzalność przekrojów jest
\emph{konsekwencją} mierzalności produktowej, a nie
jej ogólnym kryterium.
\end{uwaga}

\section{Lemat o podwójnym liczeniu miary zbioru}

W tym miejscu nie mamy jeszcze zdefiniowanej miary na
$X\times Y$. Najpierw udowodnimy, że wielkości
$\int_X\nu(E_x)\,\mathrm d\mu(x)$ oraz
$\int_Y\mu(E^y)\,\mathrm d\nu(y)$ mają sens
i są równe, gdy obie miary są $\sigma$-skończone.
Argument jest podstawą konstrukcji miary produktowej.

\begin{lemat}[Przekroje na prostokącie skończonej miary]
\label{lem:finite-rectangle-sections-12}
Niech $C\in\mathcal F$, $D\in\mathcal G$ oraz
$\mu(C)<\infty$, $\nu(D)<\infty$. Jeżeli
$E\in\mathcal F\otimes\mathcal G$ i $E\subseteq C\times D$,
to funkcje $x\mapsto\nu(E_x)$ oraz
$y\mapsto\mu(E^y)$ są mierzalne i
\begin{equation}\label{eq:set-fubini-finite-12}
\int_C\nu(E_x)\,\mathrm d\mu(x)
=\int_D\mu(E^y)\,\mathrm d\nu(y).
\end{equation}
\end{lemat}
\begin{proof}
Rozpatrujemy przestrzeń mierzalną
$R=C\times D$ z $\sigma$-algebrą
$\mathcal H=(\mathcal F\otimes\mathcal G)|_R$.
Jest ona generowana przez prostokąty
$A\times B$ z $A\in\mathcal F|_C$ i
$B\in\mathcal G|_D$. Niech $\mathcal M$ będzie
rodziną tych $E\in\mathcal H$, dla których
obie funkcje przekrojowe są mierzalne oraz zachodzi
\eqref{eq:set-fubini-finite-12}.

\emph{Krok 1: algebra prostokątów.}
Dla $E=A\times B\subseteq R$ mamy
$\nu(E_x)=\nu(B)\chi_A(x)$ i
$\mu(E^y)=\mu(A)\chi_B(y)$. Obie całki
są równe $\mu(A)\nu(B)$.
Ta sama własność obowiązuje dla skończonej
\emph{rozłącznej} sumy prostokątów, z addytywności
miar i liniowości całki. Każdą skończoną sumę
prostokątów można przedstawić jako skończoną
rozłączną sumę prostokątów: wystarczy rozbić $C$
i $D$ na skończone rozłączne podziały wyznaczone
przez występujące pierwsze i drugie czynniki.
Zatem $\mathcal M$ zawiera algebrę $\mathcal A$
skończonych sum prostokątów.

\emph{Krok 2: ciągi rosnące.}
Jeżeli $E_n\in\mathcal M$ i $E_n\uparrow E$,
to dla każdego $x$ zachodzi $(E_n)_x\uparrow E_x$.
Ciągłość miary $\nu$ z dołu daje
$\nu(E_x)=\lim_n\nu((E_n)_x)$; zatem funkcja
$x\mapsto\nu(E_x)$ jest mierzalna jako granica
funkcji mierzalnych. Podobnie dla drugiego przekroju.
Twierdzenie Beppo Leviego zastosowane po obu stronach
\eqref{eq:set-fubini-finite-12} daje tę samą
równość dla $E$.

\emph{Krok 3: ciągi malejące.}
Jeżeli $E_n\in\mathcal M$ i $E_n\downarrow E$,
to $\nu((E_n)_x)\downarrow\nu(E_x)$, ponieważ
$\nu((E_1)_x)\le\nu(D)<\infty$. Ponadto funkcje
$\nu((E_n)_x)$ są zdominowane przez
$\nu(D)\chi_C(x)$, która ma skończoną
całkę $\mu(C)\nu(D)$ (w przypadku wartości zero
stosujemy zwykłą konwencję dla całki funkcji zerowej).
Twierdzenie o zbieżności zmajoryzowanej daje zbieżność
całek; analogicznie z drugiej strony, z dominantą
$\mu(C)\chi_D$. Stąd równość dla $E$.

Z kroków 2--3 wynika, że $\mathcal M$ jest klasą
monotoniczną. Zawiera algebrę $\mathcal A$, więc
na mocy twierdzenia o klasie monotonicznej
$\mathcal H=\sigma(\mathcal A)\subseteq\mathcal M$.
To kończy dowód bez użycia Fubiniego lub Tonellego.
\end{proof}

\begin{twierdzenie}[Tożsamość przekrojowa przy $\sigma$-skończoności]
\label{thm:section-identity-12}
Niech $\mu$ i $\nu$ będą $\sigma$-skończone.
Dla każdego $E\in\mathcal F\otimes\mathcal G$
funkcje $x\mapsto\nu(E_x)$ i $y\mapsto\mu(E^y)$
są mierzalne oraz
\begin{equation}\label{eq:section-identity-12}
\int_X\nu(E_x)\,\mathrm d\mu(x)
=\int_Y\mu(E^y)\,\mathrm d\nu(y)
\quad\text{w }[0,\infty].
\end{equation}
\end{twierdzenie}
\begin{proof}
Wybieramy rosnące wyczerpania
$X_n\uparrow X$ i $Y_n\uparrow Y$ zbiorami mierzalnymi
o skończonych miarach. Można je otrzymać, zastępując
zbiory z dowolnych wyczerpań ich skończonymi sumami.
Niech $R_n=X_n\times Y_n$ oraz $E_n=E\cap R_n$.
Lemat~\ref{lem:finite-rectangle-sections-12}
zastosowany do $E_n$ zapewnia mierzalność obu
funkcji przekrojowych dla $E_n$ i równość odpowiednich
całek. Ponieważ $E_n\uparrow E$, dla każdego $x$
zachodzi $(E_n)_x\uparrow E_x$, a dla każdego $y$
zachodzi $(E_n)^y\uparrow E^y$. Z ciągłości miary
z dołu wynika
$\nu(E_x)=\lim_n\nu((E_n)_x)$ i
$\mu(E^y)=\lim_n\mu((E_n)^y)$. Granice są
mierzalne. Stosując dwukrotnie twierdzenie
Beppo Leviego do całek po $X$ i $Y$, uzyskujemy
\eqref{eq:section-identity-12}.
\end{proof}

\section{Istnienie i jednoznaczność miary produktowej}

\begin{twierdzenie}[Konstrukcja miary produktowej]
\label{thm:product-measure-12}
Niech $(X,\mathcal F,\mu)$ i $(Y,\mathcal G,\nu)$
będą przestrzeniami z miarami $\sigma$-skończonymi.
Istnieje dokładnie jedna miara $\pi$ na
$(X\times Y,\mathcal F\otimes\mathcal G)$ taka, że
\begin{equation}\label{eq:rect-product-12}
\pi(A\times B)=\mu(A)\nu(B),
\qquad A\in\mathcal F,\quad B\in\mathcal G,
\end{equation}
przy umowie $0\cdot\infty=\infty\cdot0=0$.
Oznaczamy ją przez $\mu\otimes\nu$.
Dla każdego $E\in\mathcal F\otimes\mathcal G$ zachodzi
\begin{equation}\label{eq:measure-section-12}
(\mu\otimes\nu)(E)=
\int_X\nu(E_x)\,\mathrm d\mu(x)
=\int_Y\mu(E^y)\,\mathrm d\nu(y).
\end{equation}
\end{twierdzenie}
\begin{proof}
\emph{Istnienie.}
Definiujemy $\pi(E)$ jako pierwszą całkę
w \eqref{eq:measure-section-12}; na mocy
\ref{thm:section-identity-12} całkowana funkcja jest
mierzalna i można też użyć drugiej całki.
Oczywiście $\pi(\varnothing)=0$.
Jeśli $(E_j)_{j\ge1}$ jest ciągiem parami
rozłącznych zbiorów w $\mathcal F\otimes\mathcal G$,
to dla każdego $x$ przekroje $(E_j)_x$ są parami
rozłączne i
\[
\nu\left(\left(\bigcup_{j\ge1}E_j\right)_x\right)
=\sum_{j\ge1}\nu((E_j)_x).
\]
Stosując twierdzenie Beppo Leviego do sum częściowych,
ostatecznie otrzymujemy
$\pi(\bigcup_jE_j)=\sum_j\pi(E_j)$.
Zatem $\pi$ jest miarą.
Dla $E=A\times B$ funkcja przekrojowa jest równa
$\nu(B)\chi_A$ (w rozumieniu funkcji przyjmującej
wartość zero poza $A$), więc
$\pi(A\times B)=\mu(A)\nu(B)$ także w przypadkach
$\mu(A)=0$ lub $\nu(B)=0$.

\emph{Jednoznaczność.}
Niech $\rho$ będzie inną miarą spełniającą
\eqref{eq:rect-product-12}. Z wyczerpań użytych
w poprzednim twierdzeniu wybieramy prostokąty
$R_n=X_n\times Y_n\uparrow X\times Y$, gdzie
$\mu(X_n)\nu(Y_n)<\infty$.
Na $\sigma$-algebrze śladu na $R_n$ miary
$\rho$ i $\pi$ są skończone i zgadzają się na
prostokątach $A\times B\subseteq R_n$.
Przecięcia dwóch prostokątów są prostokątami,
a prostokąty generują $\sigma$-algebrę śladu;
twierdzenie $\pi$--$\lambda$ z rozdziału 1 daje
równość miar na całym śladzie.
W szczególności dla dowolnego
$E\in\mathcal F\otimes\mathcal G$ mamy
$\rho(E\cap R_n)=\pi(E\cap R_n)$.
Przechodząc do granicy $n\to\infty$ przez
ciągłość miar z dołu, dostajemy $\rho(E)=\pi(E)$.
\end{proof}

\begin{uwaga}[Co daje $\sigma$-skończoność?]
W powyższym dowodzie jest potrzebna do
wyczerpania przestrzeni prostokątami skończonej
miary, a w konsekwencji do jednoznaczności.
Bez $\sigma$-skończoności równość iterowanych całek
i jednoznaczność rozszerzenia nie muszą zachodzić.
Nie należy traktować jej jako dekoracyjnego założenia.
\end{uwaga}

\section{Twierdzenie Tonellego}

\begin{twierdzenie}[Tonelli]
\label{thm:tonelli-12}
Niech $\mu$ i $\nu$ będą miarami $\sigma$-skończonymi,
a $f:X\times Y\to[0,\infty]$ funkcją
$\mathcal F\otimes\mathcal G$-mierzalną.
Wtedy funkcje
\[
x\longmapsto\int_Y f(x,y)\,\mathrm d\nu(y),
\qquad
y\longmapsto\int_X f(x,y)\,\mathrm d\mu(x)
\]
są mierzalne, a trzy całki są równe
(wartość $+\infty$ jest dozwolona):
\begin{equation}\label{eq:tonelli-12}
\int_{X\times Y}f(x,y)\,\mathrm d(\mu\otimes\nu)(x,y)
=\int_X\left(\int_Yf(x,y)\,\mathrm d\nu(y)\right)
\mathrm d\mu(x)
=\int_Y\left(\int_Xf(x,y)\,\mathrm d\mu(x)\right)
\mathrm d\nu(y).
\end{equation}
\end{twierdzenie}
\begin{proof}
\emph{Krok 1: funkcje charakterystyczne.}
Dla $f=\chi_E$, $E\in\mathcal F\otimes\mathcal G$,
mamy $\int_Yf(x,y)\,\mathrm d\nu(y)=\nu(E_x)$.
Obie równości w \eqref{eq:tonelli-12} są dokładnie
wzorem \eqref{eq:measure-section-12}.

\emph{Krok 2: funkcje proste nieujemne.}
Jeśli $s=\sum_{j=1}^m a_j\chi_{E_j}$,
gdzie $a_j\ge0$, $E_j$ są mierzalne i parami
rozłączne, to liniowość całki skończonej sumy
funkcji nieujemnych daje
$\int_Y s(x,y)\,\mathrm d\nu(y)
=\sum_j a_j\nu((E_j)_x)$.
Wyrażenie po prawej jest mierzalne w $x$.
Całkując po $X$ i używając kroku 1,
otrzymujemy pierwszą równość dla $s$;
analogicznie postępujemy z odwrotną kolejnością.

\emph{Krok 3: dowolne funkcje nieujemne.}
Na mocy konstrukcji funkcji prostych z rozdziału 6
istnieje ciąg $s_n\uparrow f$ nieujemnych,
prostych i mierzalnych na $X\times Y$.
Dla każdego $x$ mamy $(s_n)_x\uparrow f_x$,
więc twierdzenie Beppo Leviego względem $\nu$
daje $\int_Y(s_n)_x\,\mathrm d\nu(y)\uparrow
\int_Yf_x\,\mathrm d\nu(y)$. Granica funkcji
mierzalnych jest mierzalna. Ponownie stosujemy
Beppo Leviego, tym razem względem $\mu$:
\[
\int_X\left(\int_Y f_x\,\mathrm d\nu(y)\right)\mathrm d\mu(x)
=\lim_n\int_X\left(\int_Y (s_n)_x\,\mathrm d\nu(y)\right)\mathrm d\mu(x)
=\lim_n\int_{X\times Y}s_n\,\mathrm d(\mu\otimes\nu)(x,y).
\]
Ostatnia granica jest równa
$\int_{X\times Y}f\,\mathrm d(\mu\otimes\nu)(x,y)$,
jeszcze raz z twierdzenia Beppo Leviego.
Argument z całkowaniem najpierw względem $\mu$
jest identyczny.
\end{proof}

\begin{wniosek}[Kryterium nieujemnej całkowalności]
\label{cor:tonelli-integrability-12}
W sytuacji twierdzenia~\ref{thm:tonelli-12}
niech $f\ge0$. Funkcja $f$ ma skończoną całkę
względem $\mu\otimes\nu$ wtedy i tylko wtedy,
gdy dowolna (a zatem każda) z iterowanych całek
w \eqref{eq:tonelli-12} jest skończona.
\end{wniosek}
\begin{proof}
Wszystkie trzy całki są równe jako elementy
$[0,\infty]$. Skończoność jednej jest zatem
równoważna skończoności pozostałych.
\end{proof}

\section{Twierdzenie Fubiniego dla funkcji całkowalnych}

\begin{twierdzenie}[Fubini]
\label{thm:fubini-12}
Niech $\mu$ i $\nu$ będą miarami $\sigma$-skończonymi,
a $f:X\times Y\to\mathbb R$ lub $\mathbb C$ funkcją
$\mathcal F\otimes\mathcal G$-mierzalną taką, że
\begin{equation}\label{eq:abs-fubini-12}
\int_{X\times Y}|f(x,y)|\,\mathrm d(\mu\otimes\nu)(x,y)<\infty.
\end{equation}
Wtedy:
\begin{enumerate}[label=\textup{(\roman*)}]
\item dla $\mu$-prawie każdego $x$ przekrój $f_x$ jest
      całkowalny względem $\nu$; analogicznie dla
      $\nu$-prawie każdego $y$ przekrój $f^y$ jest
      całkowalny względem $\mu$;
\item przyjmując wartość zero na odpowiednich
      mierzalnych zbiorach wyjątkowych otrzymujemy
      całkowalne funkcje
      $F(x)=\int_Y f(x,y)\,\mathrm d\nu(y)$ i
      $G(y)=\int_X f(x,y)\,\mathrm d\mu(x)$;
\item zachodzi
\begin{equation}\label{eq:fubini-12}
\int_{X\times Y}f(x,y)\,\mathrm d(\mu\otimes\nu)(x,y)
=\int_XF(x)\,\mathrm d\mu(x)
=\int_YG(y)\,\mathrm d\nu(y).
\end{equation}
\end{enumerate}
\end{twierdzenie}
\begin{proof}
\emph{Krok 1: całkowalność przekrojów prawie wszędzie.}
Stosując twierdzenie Tonellego do $|f|$,
otrzymujemy
\[
\int_X\left(\int_Y|f(x,y)|\,\mathrm d\nu(y)\right)
\mathrm d\mu(x)
=\int_{X\times Y}|f|\,\mathrm d(\mu\otimes\nu)(x,y)<\infty.
\]
Funkcja nieujemna $H(x)=\int_Y|f(x,y)|\,\mathrm d\nu(y)$
jest mierzalna i ma skończoną całkę. Zatem
$N=\{x:H(x)=\infty\}$ jest mierzalnym zbiorem
miary zero. Dla $x\notin N$ przekrój $f_x$
jest całkowalny; analogicznie otrzymujemy zbiór
wyjątkowy $M\subseteq Y$ miary zero.

\emph{Krok 2: mierzalność całek przekrojowych.}
Dla funkcji rzeczywistej zapisujemy $f=f^+-f^-$.
Funkcje $x\mapsto\int_Y f^+(x,y)\,\mathrm d\nu(y)$
i $x\mapsto\int_Y f^-(x,y)\,\mathrm d\nu(y)$
są mierzalne z Tonellego. Dla $x\notin N$
obie wartości są skończone, więc ich różnica
jest dobrze określona i mierzalna. Przedłużając
różnicę zerem na $N$, otrzymujemy mierzalną
funkcję $F$. Dla funkcji zespolonej powtarzamy
tę konstrukcję osobno dla części rzeczywistej
i urojonej. Analogicznie postępujemy z $G$.

\emph{Krok 3: całkowalność i równość całek.}
Dla $x\notin N$ nierówność trójkąta dla całki
daje $|F(x)|\le H(x)$, a na $N$ funkcja $F$
jest równa zero. Zatem $F\in L^1(\mu)$;
podobnie $G\in L^1(\nu)$.
Do funkcji $f^+$ i $f^-$ stosujemy Tonellego,
a następnie odejmujemy skończone liczby;
uzyskujemy pierwszą i drugą równość
\eqref{eq:fubini-12}. Dla funkcji zespolonej
stosujemy ten argument do obu części składowych.
\end{proof}

\begin{wniosek}[Kryterium praktyczne]
\label{cor:practical-fubini-12}
Jeśli dla mierzalnej funkcji $f:X\times Y\to\mathbb C$
możemy wykazać
$\int_X\bigl(\int_Y |f(x,y)|\,\mathrm d\nu(y)\bigr)
\mathrm d\mu(x)<\infty$, to $f$ jest całkowalna
na przestrzeni produktowej i wolno zamienić kolejność
całkowania w \eqref{eq:fubini-12}.
\end{wniosek}
\begin{proof}
Twierdzenie Tonellego zastosowane do $|f|$ zapewnia
równość podanej iterowanej całki i całki
$\int_{X\times Y}|f|\,\mathrm d(\mu\otimes\nu)(x,y)$.
Skończoność tej ostatniej oznacza spełnienie
warunków Fubiniego.
\end{proof}

\section{Przykłady rachunkowe}

\begin{przyklad}[Pole trójkąta bez całkowania po powierzchni]
Niech $T=\{(x,y)\in[0,1]^2:0\le y\le x\}$.
Przekrój pionowy $T_x$ jest przedziałem $[0,x]$
dla $x\in[0,1]$, a poza tym zbiorem jest pusty.
Dlatego Tonelli daje
\[
(\lambda_1\otimes\lambda_1)(T)=\int_{[0,1]}\left(
\int_{[0,1]}\chi_{\{y\le x\}}\,\mathrm d\lambda_1(y)
\right)\mathrm d\lambda_1(x)
=\int_{[0,1]}x\,\mathrm d\lambda_1(x)=\frac12.
\]
Z twierdzenia~\ref{thm:lebesgue-products-12} poniżej wynika także
$\lambda_2(T)=(\lambda_1\otimes\lambda_1)(T)=1/2$.
W odwrotnej kolejności przekrój poziomy ma długość
$1-y$, a całka wynosi
$\int_{[0,1]}(1-y)\,\mathrm d\lambda_1(y)=1/2$.
Nawet w tym elementarnym przykładzie widać dwie
różne parametryzacje tej samej miary.
\end{przyklad}

\begin{przyklad}[Całka z nieograniczoną funkcją nieujemną]
Na $(0,1)^2$ określmy $f(x,y)=\frac{1}{x+y}$.
Funkcja jest mierzalna i nieujemna; dlatego wolno
najpierw całkować po $y$, mimo że jeszcze nie
wiemy, czy podwójna całka jest skończona:
\[
\int_{(0,1)^2}\frac{1}{x+y}\,\mathrm d\lambda_2(x,y)
=\int_{(0,1)}\bigl(\log(1+x)-\log x\bigr)
\,\mathrm d\lambda_1(x)=2\log2.
\]
Ostatnia równość wynika z elementarnego
całkowania funkcji $\log(1+x)$ i $\log x$
na $(0,1)$ oraz przejścia do granicy przy zerze.
\end{przyklad}

\begin{przyklad}[Funkcja zmieniająca znak]
Niech $f(x,y)=e^{-x-y}(x-y)$ na $(0,\infty)^2$.
Ponieważ $|x-y|\le x+y$, Tonelli oraz wzory
$\int_{(0,\infty)}e^{-x}\,\mathrm d\lambda_1(x)=1$ i
$\int_{(0,\infty)}xe^{-x}\,\mathrm d\lambda_1(x)=1$
z rozdziału~\ref{ch:parametry-calki} dają
\begin{align*}
\int_{(0,\infty)^2}|f|\,
\mathrm d(\lambda_1\otimes\lambda_1)
&\le\int_{(0,\infty)^2}(x+y)e^{-x-y}
\,\mathrm d(\lambda_1\otimes\lambda_1)\\
&=1\cdot1+1\cdot1=2.
\end{align*}
Zatem wolno użyć Fubiniego. Zamiana $x$ z $y$
zmienia znak całkowanej funkcji i zachowuje
miarę produktową, a więc
$\int_{(0,\infty)^2}f\,\mathrm d(\lambda_1\otimes\lambda_1)=0$.
\end{przyklad}

\section{Dlaczego założenia twierdzeń są ważne?}

\begin{przyklad}[Dwie iterowane całki, dwa różne wyniki]
\label{ex:nonabs-fubini-12}
Niech $X=Y=\mathbb N$ z $\sigma$-algebrą potęgową
i miarą liczącą $\kappa$. Definiujemy
\[
a_{m,n}=\begin{cases}
1,&n=m,\\
-1,&n=m+1,\\
0,&\text{w pozostałych przypadkach}.
\end{cases}
\]
Dla każdego $m$ szereg $\sum_{n=1}^\infty a_{m,n}$
ma tylko dwa niezerowe wyrazy i wynosi zero. Natomiast
dla $n=1$ szereg $\sum_{m=1}^\infty a_{m,1}$
ma sumę $1$, a dla $n\ge2$ sumę zero.
Wobec tego
\[
\sum_{m=1}^\infty\left(\sum_{n=1}^\infty a_{m,n}\right)=0,
\qquad
\sum_{n=1}^\infty\left(\sum_{m=1}^\infty a_{m,n}\right)=1.
\]
Nie ma sprzeczności z Fubinim, ponieważ
$\sum_{m,n}|a_{m,n}|=\infty$.
W rzeczywistości $a^+$ i $a^-$ mają nieskończoną
całkę względem $\kappa\otimes\kappa$, więc
całka podpisana po całej przestrzeni nie istnieje.
\end{przyklad}

\begin{przyklad}[Bez $\sigma$-skończoności Tonelli może zawieść]
\label{ex:non-sigma-tonelli-12}
Niech $X=Y=[0,1]$, przy czym na $X$ bierzemy
$\mathcal F=\mathcal P([0,1])$ i miarę liczącą $\kappa$,
a na $Y$ miarę Lebesgue'a $\lambda_1$ ze zwykłą
$\sigma$-algebrą. Miara $\kappa$ nie jest
$\sigma$-skończona, gdyż jej zbiory skończonej
miary są skończone, a przeliczalna suma zbiorów
skończonych jest przeliczalna.
Zbiór przekątnej $D=\{(x,y):x=y\}$ jest
mierzalny w $\mathcal F\otimes\mathcal L_1$,
gdyż jest domknięty w zwykłej topologii produktu,
a $\mathcal B([0,1]^2)\subseteq\mathcal P([0,1])\otimes\mathcal L_1$:
każdy prostokąt otwarty należy do tej $\sigma$-algebry,
a otwarte podzbiory $[0,1]^2$ są przeliczalnymi sumami
takich prostokątów.
Dla każdego $x$ przekrój $D_x=\{x\}$ ma
$\lambda_1$-miarę zero, natomiast dla każdego $y$
przekrój $D^y=\{y\}$ ma miarę liczącą $1$.
Dlatego
\[
\int_X\left(\int_Y\chi_D(x,y)\,
\mathrm d\lambda_1(y)\right)\mathrm d\kappa(x)=0,
\qquad
\int_Y\left(\int_X\chi_D(x,y)\,
\mathrm d\kappa(x)\right)\mathrm d\lambda_1(y)=1.
\]
Nie można bez dodatkowych warunków konstruować
jednej miary produktowej przez dowolną z dwóch
kolejności całkowania. Ten przykład pokazuje
dokładnie miejsce, w którym używaliśmy
$\sigma$-skończoności w dowodzie.
\end{przyklad}

\section{Produkt miar Lebesgue'a i problem zupełności}

\begin{twierdzenie}[Produkt miar Lebesgue'a]
\label{thm:lebesgue-products-12}
Niech $m,n\ge1$. Miara $\lambda_m\otimes\lambda_n$
na $\mathcal L_m\otimes\mathcal L_n$ jest ograniczeniem
miary Lebesgue'a $\lambda_{m+n}$ do tej
$\sigma$-algebry. Ponadto uzupełnienie
$(\mathcal L_m\otimes\mathcal L_n,
\lambda_m\otimes\lambda_n)$ daje dokładnie
$(\mathcal L_{m+n},\lambda_{m+n})$.
W ogólności $\mathcal L_m\otimes\mathcal L_n$
nie jest zupełna.
\end{twierdzenie}
\begin{proof}
\emph{Krok 1: relacja między $\sigma$-algebrami.}
Dla $A\in\mathcal L_m$ istnieje zbiór borelowski
$B\subseteq\mathbb R^m$ taki, że
$A\mathbin\triangle B\subseteq N$ dla pewnego
borelowskiego zbioru zerowego $N$.
Analogicznie postępujemy z $C\in\mathcal L_n$.
Prostokąt $A\times C$ różni się od
borelowskiego prostokąta $B\times D$ co najwyżej
na podzbiorze
$(N\times\mathbb R^n)\cup(\mathbb R^m\times M)$,
gdzie $N$ i $M$ są borelowskie i mają miarę zero.
Zbiory te są $\lambda_{m+n}$-zerowe:
po przecięciu z ograniczonymi prostopadłościanami
wynika to z pokrywania małymi prostokątami
lub ze znanej już konstrukcji miary Lebesgue'a,
a następnie z przeliczalnej sumy po rosnących
prostopadłościanach. Stąd każdy prostokąt
$A\times C$ jest $\lambda_{m+n}$-mierzalny,
i $\mathcal L_m\otimes\mathcal L_n\subseteq\mathcal L_{m+n}$.

\emph{Krok 2: zgodność miar.}
Niech najpierw $Q\subseteq\mathbb R^m$ będzie ograniczonym
półotwartym prostopadłościanem. Dla borelowskich
$D\subseteq\mathbb R^n$ obie funkcje
$D\mapsto\lambda_{m+n}(Q\times D)$ oraz
$D\mapsto\lambda_m(Q)\lambda_n(D)$ są miarami
$\sigma$-skończonymi na $\mathcal B(\mathbb R^n)$;
zgadzają się na półotwartych prostopadłościanach
w $\mathbb R^n$, ponieważ ich produkty z $Q$ są
prostopadłościanami w $\mathbb R^{m+n}$.
Twierdzenie o jednoznaczności miar z rozdziału 2
daje równość dla każdego borelowskiego $D$.

Teraz ustalmy borelowski $D\subseteq\mathbb R^n$
o $\lambda_n(D)<\infty$. Obie funkcje
$B\mapsto\lambda_{m+n}(B\times D)$ oraz
$B\mapsto\lambda_n(D)\lambda_m(B)$ są
$\sigma$-skończonymi miarami borelowskimi na
$\mathbb R^m$ (gdy $\lambda_n(D)=0$, obie są
zerowe). Dzięki poprzedniemu akapitowi zgadzają
się na ograniczonych półotwartych prostopadłościanach
$B=Q$. Ponowne zastosowanie jednoznaczności
daje $\lambda_{m+n}(B\times D)=\lambda_m(B)\lambda_n(D)$
dla wszystkich borelowskich $B$ i $D$ ze
skończoną miarą drugiego czynnika. Dla dowolnego
borelowskiego $D$ wyczerpujemy go zbiorami
$D\cap[-k,k]^n$ i stosujemy ciągłość miary z dołu.
Tak otrzymujemy formułę także przy nieskończonych
miarach czynników, z umową $0\cdot\infty=0$.

Dla dowolnych $A,C$ jak wyżej wartość miary
prostokąta nie zmienia się po zastąpieniu go
odpowiednim prostokątem borelowskim, z wyjątkiem
przypadków nieskończonych miar, gdzie stosujemy
wyczerpanie obydwu czynników zbiorami ograniczonymi
i ciągłość miar z dołu. Zatem
$\lambda_{m+n}(A\times C)=\lambda_m(A)\lambda_n(C)$
(z konwencją dla iloczynu z zerem).
Jednoznaczność miary produktowej daje zgodność
miar na całej $\mathcal L_m\otimes\mathcal L_n$.

\emph{Krok 3: uzupełnienie.}
$\mathcal B(\mathbb R^{m+n})$ jest podzbiorem
$\mathcal L_m\otimes\mathcal L_n$, ponieważ
borelowska $\sigma$-algebra produktu euklidesowego
jest generowana przez prostokąty otwarte.
Każdy zbiór z $\mathcal L_{m+n}$ znajduje się
między dwoma zbiorami borelowskimi $B\subseteq E\subseteq C$
ze $\lambda_{m+n}(C\setminus B)=0$
(dzięki regularności miary Lebesgue'a i jej
zupełności). Zatem $E$ należy do uzupełnienia
$\mathcal L_m\otimes\mathcal L_n$.
W drugą stronę uzupełnienie pozostaje podzbiorem
$\mathcal L_{m+n}$, ponieważ ta ostatnia
$\sigma$-algebra jest zupełna. To kończy dowód.
\end{proof}

\begin{przyklad}[Niezmierzalny w produkcie podzbiór przekątnej]
\label{ex:diagonal-not-product-12}
Niech $V\subseteq[0,1]$ będzie zbiorem niemierzalnym
w sensie Lebesgue'a, np. zbiorem Vitaliego, i niech
$E=\{(x,x):x\in V\}\subseteq\mathbb R^2$.
Zbiór $E$ jest podzbiorem przekątnej
$D=\{(x,x):x\in[0,1]\}$, a $\lambda_2(D)=0$.
Zupełność $\lambda_2$ daje $E\in\mathcal L_2$.
Gdyby natomiast $E\in\mathcal L_1\otimes\mathcal L_1$,
to przeciwobraz $\Delta^{-1}(E)$ przez
$\Delta(x)=(x,x)$ byłby Lebesgue'a-mierzalny.
Rzeczywiście oba odwzorowania współrzędne
$\operatorname{pr}_j\circ\Delta$ są identycznością,
a więc $\Delta$ jest mierzalna jako odwzorowanie
do produktu przestrzeni mierzalnych. Jednak
$\Delta^{-1}(E)=V$, sprzeczność. Stąd
$\mathcal L_1\otimes\mathcal L_1\subsetneq\mathcal L_2$.
Co ciekawe, każdy pionowy i poziomy przekrój $E$
jest pusty lub jednoelementowy, a więc mierzalny:
to dodatkowo dowodzi, że nie można odwracać
lematu~\ref{lem:sections-measurable-12}.
\end{przyklad}

\begin{uwaga}[Pełne zbiory mierzalne i przekroje]
Z lematu~\ref{lem:sections-measurable-12} nie wynika,
że \emph{wszystkie} przekroje dowolnego zbioru
$E\in\mathcal L_{m+n}$ są mierzalne. Na przykład
dla niemierzalnego $V\subseteq[0,1]$ zbiór
$V\times\{0\}$ jest $\lambda_2$-mierzalny jako
podzbiór prostej poziomej miary zero, ale jego
przekrój przy $y=0$ jest równy $V$.
Wersja Fubiniego dla uzupełnionej miary mówi o
mierzalności i całkowalności przekrojów
\emph{prawie wszędzie}; do tego potrzebne jest
osobne uzasadnienie po zastąpieniu funkcji
przez mierzalnego przedstawiciela produktowego.
\end{uwaga}


\begin{twierdzenie}[Tonelli i Fubini dla uzupełnionej miary produktowej]
\label{thm:completed-product-fubini-12}
Niech $(X,\mathcal F,\mu)$ i $(Y,\mathcal G,\nu)$
będą przestrzeniami z miarami $\sigma$-skończonymi
i \emph{zupełnymi}. Niech $(\widehat{\mathcal H},
\widehat\pi)$ będzie uzupełnieniem przestrzeni
$(X\times Y,\mathcal F\otimes\mathcal G,\mu\otimes\nu)$.
Jeśli $f:X\times Y\to[0,\infty]$ jest
$\widehat{\mathcal H}$-mierzalna, to dla $\mu$-prawie
każdego $x$ przekrój $f_x$ jest $\mathcal G$-mierzalny,
a dla $\nu$-prawie każdego $y$ przekrój $f^y$ jest
$\mathcal F$-mierzalny. Po nadaniu wartości zero
na wyjątkowych mierzalnych zbiorach parametrów
funkcje powstałe przez całkowanie przekrojów są
mierzalne i spełniają wzór Tonellego z całką
względem $\widehat\pi$.

Jeśli $f$ jest rzeczywista lub zespolona,
$\widehat{\mathcal H}$-mierzalna i
$\int_{X\times Y}|f|\,\mathrm d\widehat\pi<\infty$,
to przekroje są całkowalne prawie wszędzie,
a wzór Fubiniego również zachodzi po pominięciu
zbiorów wyjątkowych.
\end{twierdzenie}
\begin{proof}
\emph{Krok 1: przedstawiciel mierzalny względem
$\mathcal F\otimes\mathcal G$.}
Z definicji uzupełnienia miary każda funkcja
$\widehat{\mathcal H}$-mierzalna o wartościach
w $\overline{\mathbb R}$ jest równa
$\widehat\pi$-prawie wszędzie pewnej funkcji
$g$ mierzalnej względem $\mathcal F\otimes\mathcal G$.
Aby to zobaczyć, dla $f\ge0$ wybieramy proste
mierzalne przybliżenia $s_n\uparrow f$, a każdy
zbiór poziomicowy każdej funkcji $s_n$ zastępujemy
zbiorem z $\mathcal F\otimes\mathcal G$ różniącym
się od niego tylko na podzbiorze mierzalnego
$\pi$-zbioru zerowego. Przeliczalna suma takich
zbiorów zerowych nadal ma $\pi$-miarę zero;
poza nią zastąpione funkcje proste mają tę samą
granicę co $f$. W przypadku funkcji ze znakiem
i zespolonych postępujemy osobno z częściami
rzeczywistymi i urojonymi. Daje to $g$ i zbiór
$N\in\mathcal F\otimes\mathcal G$ z $\pi(N)=0$
taki, że $f=g$ poza $N$.

\emph{Krok 2: przekroje zbioru zerowego.}
Tonelli dla $\chi_N$ daje
$0=\pi(N)=\int_X\nu(N_x)\,\mathrm d\mu(x)$.
Wobec tego $\nu(N_x)=0$ dla $\mu$-prawie każdego
$x$. Analogicznie $\mu(N^y)=0$ dla $\nu$-prawie
każdego $y$. Dla takich $x$ przekroje $f_x$ i $g_x$
różnią się jedynie na podzbiorze $N_x$.
Ponieważ $\nu$ jest zupełna, $f_x$ jest
$\mathcal G$-mierzalna i ma tę samą całkę co $g_x$;
podobnie z przekrojami poziomymi. Poza zbiorami
wyjątkowymi można zatem całkować $g$ zamiast $f$.

\emph{Krok 3: równość całek.}
Z równości $f=g$ $\widehat\pi$-prawie wszędzie
wynika równość ich całek względem $\widehat\pi$
w przypadku funkcji nieujemnej lub całkowalnej.
Na $g$ stosujemy odpowiednio twierdzenie
\ref{thm:tonelli-12} lub \ref{thm:fubini-12}.
Z kroku 2 całki przekrojowe $f$ i $g$ zgadzają
się prawie wszędzie, więc ich całki iterowane
są równe. To daje obie tezy.
\end{proof}

\section{Zbiory zerowe w produkcie i wyjątki przekrojowe}

\begin{wniosek}[Prawie wszystkie przekroje zbioru zerowego]
\label{cor:null-sections-12}
Niech $\mu,\nu$ będą miarami $\sigma$-skończonymi,
a $N\in\mathcal F\otimes\mathcal G$. Wówczas
$(\mu\otimes\nu)(N)=0$ wtedy i tylko wtedy, gdy
$\nu(N_x)=0$ dla $\mu$-prawie każdego $x$, a równoważnie
$\mu(N^y)=0$ dla $\nu$-prawie każdego $y$.
Nie wolno w tym stwierdzeniu zastępować
\emph{prawie każdego} przez \emph{każdego}.
\end{wniosek}
\begin{proof}
Na mocy twierdzenia~\ref{thm:section-identity-12}
obie funkcje przekrojowe są mierzalne i mają
nieujemne wartości. Z konstrukcji miary produktowej
\eqref{eq:measure-section-12} ich całki są równe
$(\mu\otimes\nu)(N)$. Całka funkcji nieujemnej
jest równa zeru dokładnie wtedy, gdy ta funkcja
znika prawie wszędzie. To daje obie równoważności.
Dla przykładu, gdy $X=Y=[0,1]$ z miarami Lebesgue'a,
zbiór $N=\{0\}\times[0,1]$ ma miarę produktową zero,
lecz $\lambda_1(N_0)=1$.
\end{proof}

\section{Więcej niż dwa czynniki i zastosowanie do prawdopodobieństwa}

\begin{wniosek}[Łączność skończonych produktów]
\label{cor:associative-product-12}
Niech $(X_j,\mathcal F_j,\mu_j)$, $j=1,2,3$,
będą przestrzeniami z miarami $\sigma$-skończonymi.
Przy naturalnej identyfikacji
$(X_1\times X_2)\times X_3=X_1\times(X_2\times X_3)$
zachodzi
$(\mu_1\otimes\mu_2)\otimes\mu_3
=\mu_1\otimes(\mu_2\otimes\mu_3)$.
W szczególności dla każdej nieujemnej funkcji
mierzalnej na iloczynie trzech przestrzeni
kolejność całkowania jest dowolna, a dla funkcji
bezwzględnie całkowalnej można stosować Fubiniego
w dowolnej kolejności.
\end{wniosek}
\begin{proof}
Oba konstrukty są miarami na tej samej
$\sigma$-algebrze produktowej, ponieważ
$\sigma$-algebra generowana przez trójwymiarowe
prostokąty nie zależy od umieszczenia nawiasów.
Dla każdego prostokąta $A_1\times A_2\times A_3$
obie miary przyjmują wartość
$\mu_1(A_1)\mu_2(A_2)\mu_3(A_3)$
(z konwencją dla zer i nieskończoności).
Miary są $\sigma$-skończone jako produkty miar
$\sigma$-skończonych. Jednoznaczność rozszerzenia
na generującym układzie prostokątów zapewnia ich
równość. Twierdzenia o całkach wynikają przez
kolejne zastosowanie Tonellego i Fubiniego.
\end{proof}

\begin{przyklad}[Niezależność dwóch zmiennych losowych]
Niech $P$ i $Q$ będą miarami probabilistycznymi
na $(X,\mathcal F)$ i $(Y,\mathcal G)$.
Na produkcie $(X\times Y,\mathcal F\otimes\mathcal G)$
rozważmy $P\otimes Q$ i odwzorowania współrzędne
$U:X\times Y\to X$, $U(x,y)=x$, oraz
$V:X\times Y\to Y$, $V(x,y)=y$.
Oba są mierzalne względem odpowiednich $\sigma$-algebr.
Dla $A\in\mathcal F$ i $B\in\mathcal G$
mamy
$(P\otimes Q)(\{U\in A,V\in B\})=P(A)Q(B)$.
To jest definicja niezależności zmiennych losowych
$U$ i $V$. Jeżeli $\varphi:X\to[0,\infty]$ oraz
$\psi:Y\to[0,\infty]$ są mierzalne, Tonelli daje
\[
\int_{X\times Y}\varphi(x)\psi(y)
\,\mathrm d(P\otimes Q)(x,y)
=\left(\int_X\varphi(x)\,\mathrm dP(x)\right)
\left(\int_Y\psi(y)\,\mathrm dQ(y)\right).
\]
Przy funkcjach zespolonych tożsamość obowiązuje,
jeżeli są całkowalne względem odpowiednich miar.
\end{przyklad}

\section{Podsumowanie i dalszy kierunek}
Najpierw udowodniliśmy mierzalność przekrojów,
a następnie \emph{bez użycia Fubiniego} wykazaliśmy
równość całek przekrojowych w prostokątach o
skończonej mierze. Twierdzenie o klasie monotonicznej
oraz wyczerpanie przez prostokąty $\sigma$-skończone
umożliwiły konstrukcję jednoznacznej miary
$\mu\otimes\nu$. Dopiero z tej konstrukcji
wyprowadziliśmy Tonellego i Fubiniego.

Warto zapamiętać rozdział ról: Tonelli pozwala
przenosić kolejność dla funkcji \emph{nieujemnej},
także o nieskończonej całce. Fubini zapewnia
pełną równość całek funkcji zmieniających znak,
lecz wymaga kontroli $\int|f|$. Przypadki
miar nie-$\sigma$-skończonych i niezupełność
$\sigma$-algebry produktowej ostrzegają przed
nadmiernym uogólnianiem tych wyników.

% ROZDZIAŁ 13 — samodzielny rozdział do projektu teoria_miary
\chapter{Zamiana zmiennych w całce Lebesgue'a}
\label{ch:zamiana-zmiennych}

\section{Miary obrazowe: co właściwie zmienia się w całce?}
Niech $(X,\mathcal F,\mu)$ będzie przestrzenią z miarą, $(Y,\mathcal G)$
przestrzenią mierzalną, a $T:X\to Y$ odwzorowaniem mierzalnym.
\emph{Miara obrazowa} $T_\#\mu$ jest określona na $\mathcal G$ wzorem
$(T_\#\mu)(B)=\mu(T^{-1}(B))$. Nazwa ta nie oznacza, że
$\mu$ jest funkcją i można ją w zwykły sposób złożyć z $T$.

\begin{twierdzenie}[Całkowanie względem miary obrazowej]
\label{thm:pushforward-integral-13}
Niech $(X,\mathcal F,\mu)$ będzie przestrzenią z miarą,
$(Y,\mathcal G)$ przestrzenią mierzalną, a
$T:X\to Y$ odwzorowaniem $\mathcal F/\mathcal G$-mierzalnym.
Jeśli $h:Y\to[0,\infty]$ jest $\mathcal G$-mierzalna, to
\begin{equation}\label{eq:pushforward-13}
\int_Y h(y)\,\mathrm d(T_\#\mu)(y)
=\int_X h(T(x))\,\mathrm d\mu(x).
\end{equation}
Ta sama równość zachodzi dla funkcji $h:Y\to\mathbb C$
całkowalnych względem $T_\#\mu$.
\end{twierdzenie}
\begin{proof}
Dla $h=\chi_B$ obie strony są równe $\mu(T^{-1}(B))$.
Z liniowości wynika teza dla nieujemnych funkcji prostych.
Dla dowolnej $h\ge0$ wybieramy funkcje proste $h_n\uparrow h$;
$h_n\circ T\uparrow h\circ T$, więc dwukrotne zastosowanie
twierdzenia Beppo Leviego daje \eqref{eq:pushforward-13}.
Jeśli $h$ jest zespolona i całkowalna, poprzedni przypadek dla
$|h|$ dowodzi całkowalności $h\circ T$. Następnie rozdzielamy
części rzeczywiste i urojone oraz ich części dodatnie i ujemne.
\end{proof}

\begin{uwaga}
Równość \eqref{eq:pushforward-13} nie wymaga, aby $T$ było
różnowartościowe, ciągłe, różniczkowalne, ani nawet aby przestrzenie
$X,Y$ były topologiczne. Natomiast wzór z wyznacznikiem jakobianu,
który udowodnimy poniżej, wykorzystuje dodatkowo lokalną strukturę
przestrzeni euklidesowej i różnowartościowość $T$.
\end{uwaga}

\begin{przyklad}[Miara obrazu a rozkład zmiennej losowej]
Na $X=(0,1)$ z miarą Lebesgue'a rozważmy $T(x)=x^2$.
Dla $0\le t\le1$ mamy
$(T_\#\lambda_1)((0,t])=\sqrt t$; miara $T_\#\lambda_1$
ma zatem gęstość $y\mapsto(2\sqrt y)^{-1}$ na $(0,1)$.
W istocie wzór
$\int_{(0,1)}h(x^2)\,\mathrm d\lambda_1(x)
=\int_{(0,1)}h(y)/(2\sqrt y)\,\mathrm d\lambda_1(y)$
poznamy jako szczególny przypadek zamiany zmiennych.
Na razie wyprowadziliśmy równość miar na przedziałach $(0,t]$;
jednoznaczność miar skończonych na generującym układzie przedziałów
uzasadnia stwierdzenie o gęstości.
\end{przyklad}

\section{Zamiana liniowa i geometryczne znaczenie wyznacznika}

\begin{lemat}[Translacje, permutacje i ścinanie]
\label{lem:shear-13}
Niech $d\ge1$, $b\in\mathbb R^d$ i $E\subset\mathbb R^d$
będzie zbiorem borelowskim. Translacja $E\mapsto E+b$ zachowuje
miarę $\lambda_d$, podobnie jak dowolna permutacja współrzędnych.
Dla $d\ge2$, dowolnych
$i,j\in\{1,\ldots,d\}$, $i\ne j$, oraz $a\in\mathbb R$
miarę zachowuje także ścinanie postaci
\[
S(x_1,\ldots,x_d)=(x_1,\ldots,x_{j-1},
 x_j+a x_i,x_{j+1},\ldots,x_d).
\]
\end{lemat}
\begin{proof}
Niezmienniczość translacyjna została udowodniona w rozdziale 4.
Permutacje współrzędnych zachowują miarę prostokątów
półotwartych, zatem — przez jednoznaczność miary
$\sigma$-skończonej — wszystkich zbiorów borelowskich.
Dla ścinania ustalmy $d-1$ współrzędnych innych niż $x_j$.
W takim przekroju $S$ jest translacją zmiennej $x_j$
o stałą $a x_i$. Przekrój $S(E)$ jest więc translacją
odpowiedniego przekroju $E$. Ponieważ $S$ jest homeomorfizmem,
$S(E)$ jest borelowski. Twierdzenie Tonellego
\ref{thm:tonelli-12} zastosowane do funkcji charakterystycznych $\chi_E$ i $\chi_{S(E)}$
i niezmienniczość $\lambda_1$ dają równość ich miar.
\end{proof}

\begin{lemat}[Odwzorowania Lipschitzowskie zachowują zbiory zerowe]
\label{lem:lipschitz-null-13}
Niech $U\subset\mathbb R^d$ będzie otwarty, a
$F:U\to\mathbb R^d$ lokalnie Lipschitzowskie.
Jeżeli $N\subset U$ ma miarę Lebesgue'a zero, to
$\lambda_d(F(N))=0$ (w szczególności $F(N)$ jest mierzalny).
\end{lemat}
\begin{proof}
Zbiór $U$ można pokryć przeliczalnie domkniętymi sześcianami
$Q_j$ takimi, że każdy $Q_j$ zawiera się we wnętrzu
większego sześcianu położonego w $U$, na którym $F$
jest Lipschitzowska ze stałą $L_j<\infty$.
Takie pokrycie istnieje dzięki lokalnej własności
Lipschitza i przeliczalności bazy przestrzeni euklidesowej. Wystarczy rozważyć $N\cap Q_j$.
Dla ustalonego $j$ i $\varepsilon>0$ pokrywamy ten zbiór
przeliczalnie sześcianami $I_k$ o tak małych średnicach,
że leżą w odpowiednim otoczeniu $Q_j$ oraz
$\sum_k\lambda_d(I_k)<\varepsilon$.
Obraz $F(I_k)$ jest zawarty w sześcianie o boku
co najwyżej $2L_j\sqrt d$ razy bok $I_k$ (wokół
obrazu wybranego punktu $I_k$). Stąd
$\lambda_d^*(F(N\cap Q_j))\le C_d L_j^d\varepsilon$.
Przechodząc do $\varepsilon\downarrow0$, otrzymujemy
miarę zewnętrzną zero; zupełność miary Lebesgue'a
zapewnia mierzalność. Przeliczalna suma obrazów
$F(N\cap Q_j)$ jest więc zerowa.
\end{proof}

\begin{twierdzenie}[Liniowy wzór na objętość]
\label{thm:linear-change-13}
Jeżeli $A\in\mathrm{GL}(d,\mathbb R)$, $b\in\mathbb R^d$
i $E\subset\mathbb R^d$ jest borelowski, to
\begin{equation}\label{eq:linear-volume-13}
\lambda_d(AE+b)=|\det A|\lambda_d(E).
\end{equation}
Równość obowiązuje również dla dowolnego zbioru
$\lambda_d$-mierzalnego $E$, z wartościami w $[0,\infty]$.
\end{twierdzenie}
\begin{proof}
Mnożenie jednej współrzędnej przez $c\ne0$ mnoży miarę
przez $|c|$: w jednej zmiennej jest to skalowanie miary
Lebesgue'a, a w pozostałych stosujemy Tonellego.
Lemat~\ref{lem:shear-13} pokazuje, że operacje elementarne
polegające na dodaniu wielokrotności jednego wiersza do drugiego,
permutacji współrzędnych i skalowaniu pojedynczej współrzędnej
zmieniają miarę dokładnie przez moduł swego wyznacznika.
Eliminacja Gaussa przedstawia $A$ jako iloczyn skończenie
wielu takich macierzy. Kolejne zastosowanie ustalonych wzorów
i multiplikatywność wyznacznika dają \eqref{eq:linear-volume-13}
dla zbiorów borelowskich.

Macierz $A$ i jej odwrotność są globalnie Lipschitzowskie.
Z lematu~\ref{lem:lipschitz-null-13} wynika, że odwzorowanie lokalnie
Lipschitzowskie przenosi zbiory zerowe na zbiory miary zero.
Każdy zbiór Lebesgue'a-mierzalny różni się od zbioru borelowskiego
co najwyżej na podzbiorze borelowskiego zbioru zerowego;
zatem $AE+b$ jest mierzalny, a jego miara jest wyznaczona
przez część borelowską. To daje rozszerzenie na miarę zupełną.
\end{proof}

\begin{uwaga}[Wymiar ma znaczenie]
Twierdzenie~\ref{thm:linear-change-13} wymaga
macierzy kwadratowej nieosobliwej. W przypadku odwzorowań
$\mathbb R^m\to\mathbb R^n$, $m\ne n$, pojawiają się
inne wielkości (jakobiany wymiarowe, miary Hausdorffa),
które wprowadzimy w geometrycznej części skryptu.
\end{uwaga}

\section{Lokalna deformacja sześcianu}

\begin{lemat}[Oszacowanie obrazu małego sześcianu]
\label{lem:small-cube-13}
Niech $U\subset\mathbb R^d$ będzie otwarty, a
$\Phi:U\to\mathbb R^d$ odwzorowaniem klasy $C^1$.
Niech $x_0\in U$ oraz załóżmy, że
$A=D\Phi(x_0)$ jest odwracalna.
Dla $r>0$ oznaczmy $Q_r=[-r,r]^d$ i przyjmijmy
$x_0+Q_r\subset U$. Istnieje funkcja $\delta(r)\to0$
przy $r\downarrow0$ taka, że dla dostatecznie małych $r$
\begin{equation}\label{eq:cube-est-13}
(1-\delta(r))^d|\det A|\lambda_d(Q_r)
\le \lambda_d(\Phi(x_0+Q_r))
\le(1+\delta(r))^d|\det A|\lambda_d(Q_r).
\end{equation}
W szczególności iloraz miary obrazu i objętości sześcianu
dąży do $|\det D\Phi(x_0)|$.
\end{lemat}
\begin{proof}
Wprowadzamy normę maksimum $\|z\|_\infty$ w $\mathbb R^d$
i funkcję $H(z)=A^{-1}(\Phi(x_0+z)-\Phi(x_0))$.
Zachodzi $H(0)=0$ i $DH(0)=I$.
Niech $R(z)=H(z)-z$. Z ciągłości $DH$ wynika, że
\[
\delta(r):=\sup_{z\in Q_r}\|DR(z)\|_{\infty\to\infty}
\longrightarrow0.
\]
Dla $z,w\in Q_r$ twierdzenie o przyrostach skończonych
wzdłuż odcinka $[z,w]$ daje
$\|R(z)-R(w)\|_\infty\le\delta(r)\|z-w\|_\infty$,
a po podstawieniu $w=0$ dostajemy
$\|R(z)\|_\infty\le\delta(r)r$.
Zatem $H(Q_r)\subseteq Q_{(1+\delta(r))r}$.

Dla $y\in Q_{(1-\delta(r))r}$ rozważmy
$T_y(z)=y-R(z)$ na pełnej, domkniętej przestrzeni
metrycznej $Q_r$. Funkcja ta odwzorowuje $Q_r$ w siebie,
gdyż $\|T_y(z)\|_\infty\le(1-\delta)r+\delta r=r$,
i jest zwężeniem o stałej $\delta(r)<1$.
Twierdzenie o punkcie stałym Banacha (iteracja
$z_{n+1}=T_y(z_n)$ i oszacowanie geometryczne różnic)
daje $z\in Q_r$ z $T_y(z)=z$, czyli $H(z)=y$.
Wynika stąd drugie zawieranie
$Q_{(1-\delta(r))r}\subseteq H(Q_r)$.
Po zastosowaniu $A$ i translacji otrzymujemy
\[
\Phi(x_0)+A Q_{(1-\delta)r}
\subseteq\Phi(x_0+Q_r)
\subseteq\Phi(x_0)+A Q_{(1+\delta)r}.
\]
Miary skrajnych zbiorów liczymy za pomocą
\eqref{eq:linear-volume-13}; monotoniczność miary
daje \eqref{eq:cube-est-13}.
\end{proof}

\begin{uwaga}
W dowodzie nie użyto nieudowodnionego wzoru o zamianie zmiennych,
lecz tylko przypadku liniowego i twierdzenia o punkcie stałym.
Co więcej, na zbiorze zwartym, na którym $D\Phi$ jest ciągła
i odwracalna, parametry $\delta(r)$ można dobrać
\emph{jednostajnie} względem środka sześcianu, o ile sześciany
leżą w ustalonym otwartym otoczeniu tego zbioru.
Istotne są tu jednostajna ciągłość $D\Phi$ i ograniczoność
normy $(D\Phi)^{-1}$ na rozważanym zbiorze zwartym.
\end{uwaga}

\section{Główne twierdzenie o zamianie zmiennych}

\begin{twierdzenie}[Zamiana zmiennych dla dyfeomorfizmu $C^1$]
\label{thm:change-variables-13}
Niech $U,V\subset\mathbb R^d$ będą otwarte i niech
$\Phi:U\to V$ będzie dyfeomorfizmem klasy $C^1$
(tzn. bijekcją, dla której $\Phi$ i $\Phi^{-1}$
są klasy $C^1$). Wtedy:
\begin{enumerate}[label=\textup{(\roman*)}]
\item dla każdego zbioru Lebesgue'a-mierzalnego $E\subset U$
zbiór $\Phi(E)$ jest Lebesgue'a-mierzalny i
\begin{equation}\label{eq:change-sets-13}
\lambda_d(\Phi(E))=\int_E |\det D\Phi(x)|\,
\mathrm d\lambda_d(x);
\end{equation}
\item dla każdej Lebesgue'a-mierzalnej
$f:V\to[0,\infty]$ zachodzi
\begin{equation}\label{eq:change-main-13}
\int_V f(y)\,\mathrm d\lambda_d(y)
=\int_U f(\Phi(x))|\det D\Phi(x)|\,
\mathrm d\lambda_d(x);
\end{equation}
\item ta sama równość zachodzi dla funkcji
$f:V\to\mathbb C$ całkowalnych względem $\lambda_d$.
\end{enumerate}
\end{twierdzenie}
\begin{proof}
\emph{Krok 1. Miara obrazu dla zbiorów borelowskich.}
Dla borelowskiego $E\subset U$ określamy
$\tau(E)=\lambda_d(\Phi(E))$. Jest to miara na
$\mathcal B(U)$, ponieważ $\Phi$ jest homeomorfizmem:
obrazy zbiorów borelowskich są borelowskie w $V$, a obrazy
rozłącznych zbiorów są rozłączne. Określamy drugą miarę
$\rho(E)=\int_E J(x)\,\mathrm d\lambda_d(x)$,
gdzie $J(x)=|\det D\Phi(x)|$. Funkcja $J$ jest dodatnia
i ciągła; obie miary są $\sigma$-skończone, ponieważ
$U$ jest przeliczalną sumą zwartych podzbiorów,
a obrazy tych zwartych zbiorów są zwarte.

\emph{Krok 2. Równość na sześcianach.}
Niech $Q\subset U$ będzie domkniętym sześcianem.
Dzielimy go na $n^d$ przystających, domkniętych
sześcianów $Q_{n,j}$, $1\le j\le n^d$, o rozłącznych wnętrzach.
Na mocy lematu~\ref{lem:lipschitz-null-13} obraz brzegu
każdego $Q_{n,j}$ ma miarę zero. Zatem
$\tau(Q)=\sum_j\tau(Q_{n,j})$.
Wybieramy środek $x_{n,j}$ sześcianu $Q_{n,j}$.
Ponieważ $Q\Subset U$, funkcje $D\Phi$ i
$(D\Phi)^{-1}$ są na zwartym zbiorze $Q$ odpowiednio
jednostajnie ciągłe i ograniczone; stąd w oszacowaniach
wszystkich $Q_{n,j}$ można przyjąć jeden parametr
$\delta_n\to0$, niezależny od $j$.
Oszacowanie z lematu~\ref{lem:small-cube-13},
jednostajne względem tych środków, prowadzi do
\[
(1-\delta_n)^d\sum_jJ(x_{n,j})\lambda_d(Q_{n,j})
\le\tau(Q)\le
(1+\delta_n)^d\sum_jJ(x_{n,j})\lambda_d(Q_{n,j}),
\qquad\delta_n\to0.
\]
Sumy $\sum_jJ(x_{n,j})\lambda_d(Q_{n,j})$ zbiegają
do $\int_QJ(x)\,\mathrm d\lambda_d(x)$ jako sumy Riemanna
funkcji ciągłej na zwartym sześcianie. Otrzymujemy
$\tau(Q)=\rho(Q)$. Ta sama równość dotyczy wnętrza
sześcianu i półotwartego sześcianu, bo ich brzegi
oraz obrazy tych brzegów mają miarę zero.

\emph{Krok 3. Równość miar na $\mathcal B(U)$.}
Niech $\mathcal P$ będzie rodziną wszystkich półotwartych
sześcianów diadycznych, których domknięcia zawierają się
w $U$, wraz ze zbiorem pustym. Dwa sześciany diadyczne
albo są rozłączne, albo jeden zawiera się w drugim.
Zatem $\mathcal P$ jest układem typu $\pi$:
dwa sześciany w tej samej siatce diadycznej albo są
rozłączne, albo identyczne, a sześciany w siatkach
o różnych skalach są rozłączne lub jeden zawiera się
w drugim. Dokładniej używamy półotwartych sześcianów
$\prod_{j=1}^d[k_j2^{-n},(k_j+1)2^{-n})$,
$n\in\mathbb N_0$, $k_j\in\mathbb Z$, których domknięcia
są zawarte w $U$. Rodzina ta jest przeliczalna,
pokrywa $U$ i generuje
$\mathcal B(U)$: dla każdego $x\in U$ oraz otwartego
otoczenia $W\subset U$ punktu $x$ można wybrać
sześcian diadyczny $Q$ z $x\in Q$ i $\overline Q\subset W$.
Na każdym $C\in\mathcal P$ mamy
$\tau(C)=\rho(C)<\infty$. Aby poprawnie zastosować
argument $\pi$--$\lambda$, dla ustalonego takiego $C$
rozważamy dwie \emph{skończone} miary
$\tau_C(E)=\tau(E\cap C)$ oraz
$\rho_C(E)=\rho(E\cap C)$ na $\mathcal B(U)$.
Na $\mathcal P$ są równe, ponieważ $E\cap C\in\mathcal P$,
a z twierdzenia $\pi$--$\lambda$ wynika
$\tau(E\cap C)=\rho(E\cap C)$ dla dowolnego
$E\in\mathcal B(U)$. Wreszcie numerujemy sześciany
$C_1,C_2,\dots$ pokrywające $U$ i rozłączamy je przez
$D_1=C_1$, $D_j=C_j\setminus\bigcup_{i<j}C_i$.
Dla każdego borelowskiego $E\subset U$ równość miar
na $E\cap D_j$ wynika ze skończonej addytywności
równości na $E\cap C_j$ i $E\cap C_j\cap C_i$
(lub z rozdzielenia przez wcześniejsze sześciany).
Po zsumowaniu po $j$ dostajemy
$\tau(E)=\rho(E)$. To rozpisanie jest potrzebne:
rodzina sześcianów $\mathcal P$ sama nie zawiera
na ogół rosnącego wyczerpania $U$.

\emph{Krok 4. Uzupełnienie i przejście do całek.}
Ponieważ $\Phi$ oraz $\Phi^{-1}$ są lokalnie
Lipschitzowskie, lemat~\ref{lem:lipschitz-null-13}
pokazuje, że obrazy i przeciwobrazy zbiorów
$\lambda_d$-zerowych są zerowe. Każdy zbiór Lebesgue'a-
mierzalny różni się od borelowskiego na zbiorze zerowym,
a więc tożsamość $\tau(E)=\rho(E)$ rozciąga się na
wszystkie Lebesgue'a-mierzalne $E\subset U$; daje to (i).

Wzór (ii) sprawdzamy najpierw dla funkcji $f=\chi_B$,
$B\subset V$ mierzalnego: jego lewa strona to
$\lambda_d(B)$, a prawa, na mocy (i), to
$\int_{\Phi^{-1}(B)}J(x)\,\mathrm d\lambda_d(x)$,
co jest równe $\lambda_d(B)$.
Następnie korzystamy z liniowości dla funkcji prostych
i z twierdzenia Beppo Leviego dla funkcji nieujemnych.
W przypadku funkcji zespolonej całkowalnej stosujemy
(ii) do $|f|$, aby zapewnić całkowalność prawej strony,
a następnie rozkładamy $f$ na cztery części nieujemne.
\end{proof}

\begin{wniosek}[Zamiana na mierzalnym podzbiorze]
\label{cor:change-subdomain-13}
Niech $U,V\subset\mathbb R^d$ będą otwarte,
a $\Phi:U\to V$ dyfeomorfizmem klasy $C^1$.
Niech $E\subset U$ będzie zbiorem Lebesgue'a-mierzalnym,
a $g:\Phi(E)\to[0,\infty]$ funkcją mierzalną.
Wtedy, przy oznaczeniu
$J(x)=|\det D\Phi(x)|$, zachodzi
\[
\int_{\Phi(E)}g(y)\,\mathrm d\lambda_d(y)
=\int_E g(\Phi(x))J(x)\,\mathrm d\lambda_d(x).
\]
Ta sama równość zachodzi dla zespolonych funkcji $g$
całkowalnych na $\Phi(E)$.
\end{wniosek}
\begin{proof}
Przedłużamy $g$ przez zero na $V\setminus\Phi(E)$
i stosujemy poprzednie twierdzenie.
\end{proof}

\begin{uwaga}[Skąd bierze się moduł?]
Gdy $\det D\Phi<0$, odwzorowanie zmienia orientację,
lecz miara Lebesgue'a jest nieujemna. Z tego powodu
w formule występuje $|\det D\Phi|$, nie sam wyznacznik.
W całkowaniu form różniczkowych bez modułu uwzględnia się
orientację i otrzymuje inną regułę.
\end{uwaga}

\section{Przykłady: współrzędne biegunowe, walcowe i sferyczne}

\begin{przyklad}[Współrzędne biegunowe]
\label{ex:polar-13}
Rozważmy $\Phi:(0,\infty)\times(0,2\pi)\to\mathbb R^2$,
$\Phi(r,\theta)=(r\cos\theta,r\sin\theta)$.
Jest to dyfeomorfizm na płaszczyznę bez dodatniej
półosi poziomej i punktu początkowego, czyli na
$\mathbb R^2\setminus\{(x,0):x\ge0\}$.
Jego macierz pochodnej ma kolumny
$(\cos\theta,\sin\theta)^T$ i
$(-r\sin\theta,r\cos\theta)^T$, stąd
$\det D\Phi(r,\theta)=r$. Usunięta półprosta ma
$\lambda_2$-miarę zero. Dlatego dla każdej
mierzalnej funkcji $f:\mathbb R^2\to[0,\infty]$
\begin{equation}\label{eq:polar-13}
\int_{\mathbb R^2}f(x,y)\,\mathrm d\lambda_2(x,y)
=\int_{(0,2\pi)}\left(\int_{(0,\infty)}
 f(r\cos\theta,r\sin\theta)r\,\mathrm d\lambda_1(r)\right)
\mathrm d\lambda_1(\theta).
\end{equation}
Kolejność całkowania uzasadnia twierdzenie Tonellego.
Dla $f=\chi_{\{x^2+y^2<R^2\}}$ otrzymujemy
$\lambda_2(B(0,R))=\int_{(0,2\pi)}\int_{(0,R)}r\,
\mathrm d\lambda_1(r)\,\mathrm d\lambda_1(\theta)=\pi R^2$.
\end{przyklad}

\begin{center}
\begin{tikzpicture}[scale=1.0,>=stealth]
\draw[->] (-0.25,0)--(3.45,0) node[right] {$x$};
\draw[->] (0,-0.15)--(0,2.55) node[above] {$y$};
\filldraw[fill=blue!13,draw=blue!55!black,thick]
(25:1.55) -- (25:2.15) arc[start angle=25,end angle=42,radius=2.15]
-- (42:1.55) arc[start angle=42,end angle=25,radius=1.55] -- cycle;
\draw[dashed] (0,0)--(25:2.6);
\draw[dashed] (0,0)--(42:2.6);
\draw[->] (0.7,0) arc[start angle=0,end angle=25,radius=0.7];
\node at (12:1.1) {$\theta$};
\node[anchor=west] at (2.2,1.85) {$\Delta r$};
\draw[->] (2.18,1.75)--(42:2.07);
\node at (60:2.6) {$\Delta\theta$};
\node at (15:2.7) {$r$};
\node[anchor=west] at (0.1,-0.35) {Mały wycinek pierścienia};
\end{tikzpicture}
\end{center}
\noindent
Dla małych przyrostów $\Delta r$ i $\Delta\theta$ pole
zaznaczonego wycinka jest w pierwszym przybliżeniu
$r\,\Delta r\,\Delta\theta$. Jest to geometryczna intuicja
za jakobianem $r$; ścisłym uzasadnieniem jest jednak
twierdzenie~\ref{thm:change-variables-13}.

\begin{przyklad}[Współrzędne walcowe]
Dla $r>0$, $\theta\in(0,2\pi)$, $z\in\mathbb R$
połóżmy $\Psi(r,\theta,z)=(r\cos\theta,r\sin\theta,z)$.
Wyznacznik jej macierzy pochodnej wynosi $r$.
Po odrzuceniu półpłaszczyzny miary Lebesgue'a zero
otrzymujemy dla każdej mierzalnej
$f:\mathbb R^3\to[0,\infty]$
\[
\begin{aligned}
\int_{\mathbb R^3}f(x,y,z)\,\mathrm d\lambda_3(x,y,z)
&=\int_{\mathbb R}\int_{(0,2\pi)}\int_{(0,\infty)}
 f(r\cos\theta,r\sin\theta,z)r\,
\mathrm d\lambda_1(r)\,\mathrm d\lambda_1(\theta)
\,\mathrm d\lambda_1(z).
\end{aligned}
\]
Równość obowiązuje również dla zespolonych funkcji
$f\in L^1(\mathbb R^3,\lambda_3)$.
\end{przyklad}

\begin{przyklad}[Współrzędne sferyczne]
Dla $r>0$, $\vartheta\in(0,\pi)$, $\varphi\in(0,2\pi)$
rozważmy
$\Psi(r,\vartheta,\varphi)=
(r\sin\vartheta\cos\varphi,
 r\sin\vartheta\sin\varphi,r\cos\vartheta)$.
Bezpośrednie obliczenie wyznacznika daje
$|\det D\Psi|=r^2\sin\vartheta$.
Po odrzuceniu odpowiedniej półpłaszczyzny i osi,
które mają miarę zero, otrzymujemy dla każdej mierzalnej
$f:\mathbb R^3\to[0,\infty]$
\begin{align*}
\int_{\mathbb R^3}f(x)\,\mathrm d\lambda_3(x)
&=\int_{(0,\infty)}\int_{(0,\pi)}\int_{(0,2\pi)}
 f(\Psi(r,\vartheta,\varphi))r^2\sin\vartheta\,
 \mathrm d\lambda_1(\varphi)\,\mathrm d\lambda_1(\vartheta)
 \,\mathrm d\lambda_1(r).
\end{align*}
Ta sama równość zachodzi dla zespolonych funkcji
$f\in L^1(\mathbb R^3,\lambda_3)$.
Dla kuli $B(0,R)$ daje to $\lambda_3(B(0,R))=4\pi R^3/3$.
\end{przyklad}

\section{Obrazy miar posiadających gęstości}
\begin{twierdzenie}[Transformacja gęstości]
\label{thm:density-transform-13}
Niech $U,V\subset\mathbb R^d$ będą otwarte,
a $\Phi:U\to V$ dyfeomorfizmem klasy $C^1$.
Niech $w:U\to[0,\infty]$ będzie mierzalna i
niech $\mu$ będzie miarą na zbiorach Lebesgue'a-mierzalnych
$E\subset U$, określoną wzorem
$\mu(E)=\int_Ew(x)\,\mathrm d\lambda_d(x)$.
Wtedy dla każdego mierzalnego $B\subset V$
\[
(\Phi_\#\mu)(B)=\int_B
w(\Phi^{-1}(y))|\det D\Phi^{-1}(y)|\,
\mathrm d\lambda_d(y).
\]
\end{twierdzenie}
\begin{proof}
Z definicji miary obrazowej mamy
$(\Phi_\#\mu)(B)=\int_{\Phi^{-1}(B)}w(x)\,
\mathrm d\lambda_d(x)$.
Stosujemy twierdzenie~\ref{thm:change-variables-13}
do dyfeomorfizmu $\Phi^{-1}:V\to U$ oraz funkcji
$h:U\to[0,\infty]$, określonej przez
$h(x)=\chi_{\Phi^{-1}(B)}(x)w(x)$.
Ponieważ
$h(\Phi^{-1}(y))=\chi_B(y)w(\Phi^{-1}(y))$,
otrzymujemy
\[
\int_{\Phi^{-1}(B)}w(x)\,\mathrm d\lambda_d(x)
=\int_B w(\Phi^{-1}(y))|\det D\Phi^{-1}(y)|
\,\mathrm d\lambda_d(y).
\]
W obu całkach iloczyn $0\cdot\infty$ interpretujemy
jako $0$, zgodnie z konwencją całkowania funkcji
nieujemnych o wartościach rozszerzonych.
\end{proof}

\begin{przyklad}[Rozkład jednostajny i kwadratowanie]
Jeśli $X$ ma rozkład jednostajny na $(0,1)$ i
$Y=X^2$, to $Y$ ma gęstość na $\mathbb R$ daną wzorem
\[
p_Y(y)=
\begin{cases}
\dfrac{1}{2\sqrt y},&0<y<1,\\
0,&y\notin(0,1).
\end{cases}
\]
To dokładnie miara obrazowa rozważana na początku rozdziału.
Nie ma potrzeby zakładać, że sam $X$ jest określony
na $(0,1)$: wystarczy, aby jego rozkład był jednostajny.
\end{przyklad}

\section{Założenia, które naprawdę są potrzebne}
\begin{przyklad}[Dlaczego nie wolno pomijać różnowartościowości]
Dla $\Phi:\mathbb R\to\mathbb R$, $\Phi(x)=x^2$,
nie istnieje globalna odwrotność. Gdyby bez dodatkowych
poprawek użyć wzoru z jakobianem na całej prostej, dla
$f=\chi_{(0,1)}$ otrzymalibyśmy po prawej stronie
$\int_{(-1,1)}2|x|\,\mathrm d\lambda_1(x)=2$,
zaś po lewej $\int_{(0,1)}1\,\mathrm d\lambda_1(y)=1$.
Rozbieżność wynika z dwukrotnego liczenia punktów obrazu.
Dla nieiniektywnych odwzorowań właściwym narzędziem
jest wzór z krotnością włókien, a w większej ogólności
wzór na pole powierzchni (późniejszy rozdział).
\end{przyklad}

\begin{przyklad}[Wymagana regularność]
Z samej różnowartościowości i ciągłości $\Phi$ nie wynika
wzór z pochodną. Na przykład istnieją ściśle rosnące
homeomorfizmy przedziału, których pochodna znika
prawie wszędzie (homeomorfizmy osobliwe).
Przy takiej funkcji całka pochodnej nie daje długości obrazu.
Nie wystarczy więc samo istnienie pochodnej prawie wszędzie:
w odpowiedniej wersji uogólnionej potrzebne są dodatkowe
własności, np. absolutna ciągłość w wymiarze jeden.
\end{przyklad}

\section{Podsumowanie}
Wzór na całkę względem miary obrazowej jest czysto
mierzalny. Natomiast wyznacznik jakobianu opisuje
lokalną zmianę objętości pod dyfeomorfizmem $C^1$.
Dowód przeszedł od macierzy elementarnych do lokalnych
sześcianów, a następnie do równości miar i funkcji
mierzalnych. Formuły biegunowe i sferyczne wynikają
z jednego twierdzenia, a nie z odrębnych reguł rachunkowych.

% ROZDZIAŁ 14 — samodzielny rozdział do projektu teoria_miary
\chapter{Przestrzenie $L^p$ i ich podstawowe własności}
\label{ch:przestrzenie-lp}

\section{Dlaczego utożsamiamy funkcje równe prawie wszędzie?}
Niech $(X,\mathcal F,\mu)$ będzie dowolną przestrzenią z miarą.
Nie zakładamy, że $\mu(X)<\infty$, ani że $\mu$ jest $\sigma$-skończona.
Dla $1\le p<\infty$ rozważmy mierzalne funkcje
$f:X\to\mathbb C$ spełniające
$\int_X|f(x)|^p\,\mathrm d\mu(x)<\infty$.
Naturalna wielkość
$(\int_X|f|^p\,\mathrm d\mu(x))^{1/p}$
może znikać dla funkcji nieidentycznie zerowej:
wystarczy, że $f$ różni się od zera jedynie na zbiorze
miary zero. Aby otrzymać normę, trzeba przejść
od funkcji do ich klas równoważności.

\begin{definicja}[Równość prawie wszędzie i przestrzeń $L^p$]
\label{def:lp-14}
Niech $(X,\mathcal F,\mu)$ będzie przestrzenią z miarą.
Dwie mierzalne funkcje zespolone $f,g$ na $X$
uznajemy za równoważne, jeśli $\mu(\{x:f(x)\ne g(x)\})=0$.
Dla $1\le p<\infty$ przez $L^p(X,\mathcal F,\mu)$, lub
krótko $L^p(\mu)$, rozumiemy przestrzeń klas
równoważności mierzalnych funkcji $f:X\to\mathbb C$,
dla których $\int_X|f(x)|^p\,\mathrm d\mu(x)<\infty$.
Wartość wyrażenia $\int_X|f|^p\,\mathrm d\mu(x)$
nie zależy od przedstawiciela klasy. Kładziemy
\begin{equation}\label{eq:p-norm-14}
\|f\|_p:=\left(\int_X|f(x)|^p\,\mathrm d\mu(x)\right)^{1/p}.
\end{equation}
Dla funkcji rzeczywistych używamy oznaczenia $L^p(\mu;\mathbb R)$.
\end{definicja}

\begin{definicja}[Kres górny istotny i przestrzeń $L^\infty$]
\label{def:linf-14}
Niech $(X,\mathcal F,\mu)$ będzie przestrzenią z miarą.
Dla mierzalnej funkcji $f:X\to\mathbb C$ określamy
\[
\|f\|_\infty:=\inf\{M\ge0:\mu(\{|f|>M\})=0\},
\]
z umową $\inf\varnothing=\infty$.
Przestrzeń $L^\infty(\mu)$ jest przestrzenią klas
funkcji o skończonej wartości $\|f\|_\infty$.
Liczbę tę nazywamy \emph{kresem górnym istotnym} funkcji $|f|$.
\end{definicja}

\begin{lemat}[Znaczenie kresu istotnego]
\label{lem:esssup-14}
Niech $(X,\mathcal F,\mu)$ będzie przestrzenią z miarą.
Jeśli $f:X\to\mathbb C$ jest mierzalna
i $\|f\|_\infty<\infty$,
to $|f(x)|\le\|f\|_\infty$ dla prawie wszystkich $x$.
Ponadto $\|f\|_\infty=0$ wtedy i tylko wtedy, gdy
$f=0$ prawie wszędzie.
\end{lemat}
\begin{proof}
Niech $M=\|f\|_\infty$. Dla każdego $n$ istnieje
$M_n<M+1/n$ spełniające $\mu(\{|f|>M_n\})=0$.
Poza przeliczalną sumą tych zbiorów zachodzi
$|f|\le M_n<M+1/n$ dla wszystkich $n$, więc $|f|\le M$.
W przypadku $M=0$ otrzymujemy $f=0$ poza zbiorem zerowym.
Implikacja odwrotna wynika wprost z definicji.
\end{proof}

\begin{uwaga}[Przedstawiciele a klasy]
Zapis $f(x)$ dla $f\in L^p(\mu)$ jest skrótem:
ściśle rzecz biorąc $f$ jest klasą równoważności, więc
wartość w pojedynczym punkcie nie musi być określona.
Ilekroć piszemy $f(x)$, wybieramy przedstawiciela;
tezy o wartościach punktowych formułujemy prawie wszędzie.
Operacje algebraiczne na klasach są poprawnie określone,
ponieważ suma skończonej liczby zbiorów zerowych jest zerowa.
\end{uwaga}

\section{Elementarna nierówność Younga}

\begin{lemat}[Young dla iloczynu dwóch liczb]
\label{lem:young-numbers-14}
Niech $1<p,q<\infty$ oraz $1/p+1/q=1$.
Dla $a,b\ge0$ zachodzi
\begin{equation}\label{eq:young-14}
ab\le\frac{a^p}{p}+\frac{b^q}{q}.
\end{equation}
Równość zachodzi dokładnie wtedy, gdy $a^p=b^q$.
\end{lemat}
\begin{proof}
Dla $b>0$ funkcja $\varphi(a)=a^p/p-ab+b^q/q$
ma pochodną $\varphi'(a)=a^{p-1}-b$.
Jej minimum na $[0,\infty)$ wypada w punkcie
$a=b^{1/(p-1)}=b^{q-1}$, gdzie
$\varphi(a)=0$, ponieważ $(q-1)p=q$.
Dla $b=0$ nierówność jest natychmiastowa;
równość zachodzi wtedy jedynie przy $a=0$.
\end{proof}

\begin{uwaga}
Wykładniki $p$ i $q$ w lemacie nazywamy
\emph{sprzężonymi wykładnikami Höldera}.
Dla $p=1$ przyjmujemy formalnie $q=\infty$ i odwrotnie.
Przypadki brzegowe rozpatrujemy osobno.
\end{uwaga}

\section{Nierówność Höldera}
\begin{twierdzenie}[Hölder]
\label{thm:holder-14}
Niech $(X,\mathcal F,\mu)$ będzie przestrzenią z miarą,
$1\le p,q\le\infty$ i $1/p+1/q=1$, przy czym $1/\infty=0$.
Jeśli $f\in L^p(\mu)$ i $g\in L^q(\mu)$,
to $fg\in L^1(\mu)$ oraz
\begin{equation}\label{eq:holder-14}
\int_X|f(x)g(x)|\,\mathrm d\mu(x)
\le\|f\|_p\|g\|_q.
\end{equation}
\end{twierdzenie}
\begin{proof}
Załóżmy najpierw $1<p,q<\infty$.
Jeśli $\|f\|_p=0$ lub $\|g\|_q=0$,
to odpowiedni czynnik znika prawie wszędzie,
a więc teza jest oczywista. W przeciwnym przypadku
kładziemy $F=|f|/\|f\|_p$ oraz $G=|g|/\|g\|_q$.
Zachodzi $\int_XF^p\,\mathrm d\mu(x)=
\int_XG^q\,\mathrm d\mu(x)=1$.
Z lematu~\ref{lem:young-numbers-14} mamy punktowo
$FG\le F^p/p+G^q/q$; całkowanie daje
$\int_XFG\,\mathrm d\mu(x)\le1/p+1/q=1$.
Po pomnożeniu przez $\|f\|_p\|g\|_q$ otrzymujemy
\eqref{eq:holder-14} i zarazem całkowalność iloczynu.

Gdy $p=1$ i $q=\infty$, poza zbiorem zerowym zachodzi
$|fg|\le\|g\|_\infty|f|$ na mocy
lematu~\ref{lem:esssup-14}. Całkowanie daje nierówność;
przypadek $p=\infty$, $q=1$ jest symetryczny.
\end{proof}

\begin{wniosek}[Uogólniony Hölder]
\label{cor:generalized-holder-14}
Niech $(X,\mathcal F,\mu)$ będzie przestrzenią z miarą,
$m\ge2$, $1\le p_j\le\infty$ oraz
$\sum_{j=1}^m1/p_j=1$. Jeżeli $f_j\in L^{p_j}(\mu)$,
to $\prod_{j=1}^m f_j\in L^1(\mu)$ i
$\int_X\prod_{j=1}^m|f_j(x)|\,\mathrm d\mu(x)
\le\prod_{j=1}^m\|f_j\|_{p_j}$.
\end{wniosek}
\begin{proof}
Czynniki z $p_j=\infty$ są ograniczone prawie wszędzie;
można je oszacować przez $\|f_j\|_\infty$ i usunąć
z rozważań. Jeśli pozostał tylko jeden czynnik,
jego wykładnik musi być równy $1$ i wynik jest oczywisty.
Dla większej liczby czynników dowodzimy twierdzenia indukcyjnie.
Niech $p_m<\infty$ i $1/r=\sum_{j<m}1/p_j=1-1/p_m$.
Jeżeli $r=\infty$, wszystkie pierwsze czynniki
są ograniczone i dostajemy tezę bezpośrednio.
Jeśli $1<r<\infty$, to $p_j\ge r$ dla $j<m$,
a liczby $s_j=p_j/r\ge1$ spełniają
$\sum_{j<m}1/s_j=1$. Stosujemy krok indukcyjny
do funkcji $|f_j|^r\in L^{s_j}$:
\[
\int_X\prod_{j<m}|f_j(x)|^r\,\mathrm d\mu(x)
\le\prod_{j<m}\left(\int_X|f_j(x)|^{p_j}\,
\mathrm d\mu(x)\right)^{r/p_j}.
\]
Po wyciągnięciu pierwiastka stopnia $r$ uzyskujemy
$\|\prod_{j<m}f_j\|_r\le\prod_{j<m}\|f_j\|_{p_j}$.
Ostatnie zastosowanie Höldera do tego iloczynu i $f_m$
ze sprzężonymi wykładnikami $r,p_m$ kończy dowód.
\end{proof}

\begin{przyklad}[Oszacowanie iloczynu]
Na $X=(0,1)$ z miarą Lebesgue'a rozważmy
$f(x)=x^{-1/4}$ i $g(x)=x^{-1/3}$.
Wtedy $f\in L^2(0,1)$, $g\in L^2(0,1)$,
ponieważ $\int_{(0,1)}x^{-1/2}\,\mathrm d\lambda_1(x)<\infty$
i $\int_{(0,1)}x^{-2/3}\,\mathrm d\lambda_1(x)<\infty$.
Z Höldera wynika, że $fg=x^{-7/12}\in L^1(0,1)$.
Można to również sprawdzić bezpośrednio.
\end{przyklad}

\section{Nierówność Minkowskiego i struktura normowana}
\begin{twierdzenie}[Minkowski]
\label{thm:minkowski-14}
Dla każdej przestrzeni z miarą $(X,\mathcal F,\mu)$,
każdego $1\le p\le\infty$ oraz $f,g\in L^p(\mu)$
zachodzi $f+g\in L^p(\mu)$ oraz
\begin{equation}\label{eq:minkowski-14}
\|f+g\|_p\le\|f\|_p+\|g\|_p.
\end{equation}
\end{twierdzenie}
\begin{proof}
Dla $p=1$ z nierówności trójkąta dla wartości
bezwzględnej i liniowości całki otrzymujemy
$\int_X|f+g|\,\mathrm d\mu(x)\le\int_X|f|\,\mathrm d\mu(x)+
\int_X|g|\,\mathrm d\mu(x)$.
Dla $p=\infty$ z
$|f+g|\le\|f\|_\infty+\|g\|_\infty$ prawie wszędzie
wynika teza.

Niech $1<p<\infty$, a $q=p/(p-1)$.
Na początek mamy
$|f+g|^p\le 2^{p-1}(|f|^p+|g|^p)$,
co wynika z wypukłości $t\mapsto t^p$ na $[0,\infty)$.
Ta nierówność zapewnia $f+g\in L^p$, zanim zastosujemy Höldera.
Następnie dla $s=f+g$ korzystamy z
$|s|^p\le |f||s|^{p-1}+|g||s|^{p-1}$.
Ponieważ $(p-1)q=p$, funkcja $|s|^{p-1}$ należy
do $L^q$ i $\||s|^{p-1}\|_q=\|s\|_p^{p-1}$.
Hölder daje
\[
\|s\|_p^p
\le(\|f\|_p+\|g\|_p)\|s\|_p^{p-1}.
\]
Gdy $\|s\|_p=0$, teza jest natychmiastowa;
w pozostałym przypadku dzielimy przez $\|s\|_p^{p-1}$.
\end{proof}

\begin{wniosek}[Przestrzeń normowana]
\label{cor:lp-normed-14}
Niech $(X,\mathcal F,\mu)$ będzie przestrzenią z miarą.
Dla $1\le p\le\infty$ wzory z
definicji~\ref{def:lp-14} i \ref{def:linf-14}
określają normę na zespolonej przestrzeni liniowej $L^p(\mu)$.
\end{wniosek}
\begin{proof}
Nieujemność jest oczywista; rozdzielność wynika
dla $p<\infty$ z faktu, że całka funkcji nieujemnej
równa zeru oznacza znikanie jej prawie wszędzie,
a dla $p=\infty$ z lematu~\ref{lem:esssup-14}.
Jednorodność $\|cf\|_p=|c|\|f\|_p$ wynika
z jednorodności całki lub kresu istotnego.
Nierówność trójkąta jest treścią twierdzenia
\ref{thm:minkowski-14}. Operacje na klasach
są niezależne od wyboru przedstawiciela.
\end{proof}

\section{Twierdzenie Riesza--Fischera: zupełność przestrzeni $L^p$}

\begin{twierdzenie}[Kryterium zbieżności szeregów w $L^p$]
\label{thm:series-lp-14}
Niech $(X,\mathcal F,\mu)$ będzie przestrzenią z miarą,
$1\le p\le\infty$ i $(u_n)_{n\ge1}\subset L^p(\mu)$.
Jeżeli $\sum_{n=1}^\infty\|u_n\|_p<\infty$,
to istnieje zbiór mierzalny $N$ o $\mu(N)=0$,
poza którym szereg $\sum_{n=1}^\infty u_n(x)$
jest bezwzględnie zbieżny. Jego suma określa element $u\in L^p(\mu)$
i $\sum_{n=1}^N u_n\to u$ w normie $L^p$.
\end{twierdzenie}
\begin{proof}
Dla $p<\infty$ kładziemy $v_N(x)=\sum_{n=1}^N|u_n(x)|$.
Nierówność Minkowskiego daje
$\|v_N\|_p\le\sum_{n=1}^N\|u_n\|_p\le M$,
gdzie $M=\sum_{n\ge1}\|u_n\|_p<\infty$.
Ponieważ $v_N\uparrow v:=\sum_{n\ge1}|u_n|$,
twierdzenie Beppo Leviego zastosowane do $v_N^p$ daje
$\int_Xv^p\,\mathrm d\mu(x)\le M^p$.
W konsekwencji $v<\infty$ prawie wszędzie.
Dla $x$ poza tym zbiorem wyjątkowym szereg
$\sum_nu_n(x)$ zbiega bezwzględnie do $u(x)$.
Wszystkie sumy częściowe są mierzalne, więc $u$
można określić jako funkcję mierzalną poza zbiorem
wyjątkowym i równą zero na nim. Mamy $|u|\le v$,
zatem $u\in L^p$.

Dla ogona $u-\sum_{n\le N}u_n$ mamy punktowo
$|u-\sum_{n\le N}u_n|\le\sum_{n>N}|u_n|$.
Stosując Minkowskiego do skończonych ogonów
i Fatou po przejściu z końcem ogona do nieskończoności,
otrzymujemy
$\|u-\sum_{n\le N}u_n\|_p\le\sum_{n>N}\|u_n\|_p\to0$.

Dla $p=\infty$ wybieramy przedstawicieli $u_n$
i jeden wspólny zbiór zerowy $N$, poza którym
$|u_n(x)|\le\|u_n\|_\infty$ dla każdego $n$.
Na $X\setminus N$ szereg zbiega bezwzględnie
oraz jednostajnie, gdyż suma $\sum_n\|u_n\|_\infty$
jest skończona. Otrzymujemy mierzalną sumę $u$,
$\|u\|_\infty\le M$ i
$\|u-\sum_{n\le N}u_n\|_\infty\le\sum_{n>N}\|u_n\|_\infty$.
\end{proof}

\begin{twierdzenie}[Riesz--Fischer: zupełność $L^p$]
\label{thm:riesz-fischer-14}
Dla każdej przestrzeni z miarą $(X,\mathcal F,\mu)$
i każdego $1\le p\le\infty$ przestrzeń normowana
$L^p(\mu)$ jest zupełna, a więc jest przestrzenią Banacha.
Dla $p=2$ jest ponadto przestrzenią Hilberta
z iloczynem skalarnym
$\langle f,g\rangle:=\int_X f(x)\overline{g(x)}\,\mathrm d\mu(x)$.
\end{twierdzenie}
\begin{proof}
Niech $(f_n)$ będzie ciągiem Cauchy'ego w $L^p$.
Indukcyjnie wybieramy rosnący ciąg indeksów $(n_k)$
taki, że $\|f_{n_{k+1}}-f_{n_k}\|_p\le2^{-k}$.
Wówczas szereg $\sum_{k\ge1}(f_{n_{k+1}}-f_{n_k})$
ma sumę norm mniejszą niż $\sum_k2^{-k}<\infty$.
Na mocy twierdzenia~\ref{thm:series-lp-14}
zbiega więc w $L^p$ do pewnego elementu $w$,
a zatem $f_{n_k}\to f:=f_{n_1}+w$ w $L^p$.
Aby wykazać zbieżność całego ciągu, ustalamy
$\varepsilon>0$ i wybieramy $N$ tak, aby
$\|f_m-f_n\|_p<\varepsilon/2$ dla $m,n\ge N$.
Wybieramy także $k$ z $n_k\ge N$ i
$\|f_{n_k}-f\|_p<\varepsilon/2$.
Wtedy dla $n\ge N$ nierówność trójkąta daje
$\|f_n-f\|_p\le\|f_n-f_{n_k}\|_p+
\|f_{n_k}-f\|_p<\varepsilon$.

Dla $p=2$ Hölder gwarantuje, że podana całka
jest skończona. Jest ona liniowa w pierwszym
argumencie, antyliniowa w drugim, hermitowska
i dodatnio określona na klasach. Wytworzona przez nią
norma to $\|f\|_2$, a udowodniona już zupełność
oznacza zupełność Hilberta.
\end{proof}

\begin{uwaga}[Zbieżność punktowa w dowodzie]
Istotnym elementem dowodu jest istnienie
\emph{podciągu} zbieżnego prawie wszędzie.
Nie wynika stąd, że dowolny ciąg zbieżny w $L^p$
zbiega punktowo prawie wszędzie; zobaczymy niżej
kontrprzykład związany z rodzajami zbieżności.
\end{uwaga}

\section{Nierówność Jensena i jej znaczenie}
\begin{twierdzenie}[Jensen dla miary probabilistycznej]
\label{thm:jensen-14}
Niech $(X,\mathcal F,P)$ będzie przestrzenią
probabilistyczną, $f\in L^1(P;\mathbb R)$,
a $\varphi:\mathbb R\to\mathbb R$ funkcją wypukłą.
Wtedy $\varphi\circ f$ jest mierzalna, jej część ujemna
jest całkowalna, a więc całka
$\int_X\varphi(f(x))\,\mathrm dP(x)\in(-\infty,\infty]$
jest dobrze określona i
\begin{equation}\label{eq:jensen-14}
\varphi\left(\int_Xf(x)\,\mathrm dP(x)\right)
\le\int_X\varphi(f(x))\,\mathrm dP(x).
\end{equation}
\end{twierdzenie}
\begin{proof}
Niech $m=\int_Xf\,\mathrm dP(x)$.
Funkcja wypukła skończona na całej prostej jest ciągła,
a więc $\varphi\circ f$ jest mierzalna.
Dla funkcji wypukłej na $\mathbb R$ w każdym punkcie $m$
istnieje prosta podpierająca: jest liczba $a\in\mathbb R$
taka, że $\varphi(t)\ge\varphi(m)+a(t-m)$
dla wszystkich $t\in\mathbb R$.
Istnienie takiego $a$ wynika z nierówności między
lewym i prawym nachyleniem siecznych funkcji wypukłej;
można wybrać dowolną liczbę między jej lewą
i prawą pochodną jednostronną w $m$.
Po podstawieniu $t=f(x)$ dostajemy
$\varphi(f(x))\ge\varphi(m)+a(f(x)-m)$.
Prawa strona jest całkowalna, więc część ujemna
$\varphi(f)$ ma skończoną całkę. Możemy całkować
nierówność, otrzymując
$\int_X\varphi(f)\,\mathrm dP(x)
\ge\varphi(m)+a(\int_Xf\,\mathrm dP(x)-m)=\varphi(m)$.
\end{proof}

\begin{przyklad}[Wartość średnia a momenty]
Dla $1\le p<\infty$ funkcja $\varphi(t)=|t|^p$
jest wypukła. Gdy $P(X)=1$ i $f\in L^p(P)$,
Jensen daje
$|\int_X f\,\mathrm dP(x)|^p\le
\int_X|f(x)|^p\,\mathrm dP(x)$ dla $f$ rzeczywistej.
Dla $f$ zespolonej analogiczna nierówność wynika
z nierówności trójkąta dla całki i Höldera
z funkcją stałą $1$ (należącą do $L^q(P)$).
\end{przyklad}

\section{Zależności między przestrzeniami $L^p$}

\begin{twierdzenie}[Włożenia przy skończonej mierze]
\label{thm:lp-embedding-finite-14}
Niech $(X,\mathcal F,\mu)$ będzie przestrzenią z miarą,
$\mu(X)<\infty$ i $1\le p<q\le\infty$.
Wtedy $L^q(\mu)\subseteq L^p(\mu)$ oraz
\begin{equation}\label{eq:finite-embedding-14}
\|f\|_p\le\mu(X)^{\,1/p-1/q}\|f\|_q
\quad(f\in L^q(\mu)),
\end{equation}
przy umowie $1/\infty=0$ i naturalnej interpretacji
w przypadku $\mu(X)=0$.
\end{twierdzenie}
\begin{proof}
Jeśli $\mu(X)=0$, wszystkie przestrzenie są
zerowe i teza jest natychmiastowa. Załóżmy więc
$0<\mu(X)<\infty$.
Dla $q<\infty$ stosujemy Höldera do
$|f|^p\cdot1$ z wykładnikami $q/p$ i $q/(q-p)$:
\[
\int_X|f|^p\,\mathrm d\mu(x)
\le\left(\int_X|f|^q\,\mathrm d\mu(x)\right)^{p/q}
\mu(X)^{1-p/q}.
\]
Podnosimy obie strony do potęgi $1/p$.
Dla $q=\infty$ nierówność
$|f|^p\le\|f\|_\infty^p$ prawie wszędzie
i całkowanie dają wynik.
\end{proof}

\begin{przyklad}[Na nieskończonej mierze włożenia zawodzą]
Na $(1,\infty)$ z miarą Lebesgue'a
funkcja $f(x)=x^{-\alpha}$ spełnia
$f\in L^p$ dokładnie wtedy, gdy $\alpha p>1$.
Na $(0,1)$ funkcja $g(x)=x^{-\beta}$
należy do $L^p$ dokładnie wtedy, gdy $\beta p<1$.
Dla $1\le p<q<\infty$ wybór
$1/q<\alpha\le1/p$ daje na $(1,\infty)$
funkcję z $L^q\setminus L^p$,
a wybór $1/q\le\beta<1/p$ daje na $(0,1)$
funkcję z $L^p\setminus L^q$.
Po przedłużeniu przez zero odpowiednie przykłady
działają na całej $\mathbb R$.
Zatem na przestrzeni nieskończonej miary
w ogólności nie zachodzi żadna z inkluzji między
różnymi $L^p$.
\end{przyklad}

\begin{twierdzenie}[Zbieżność w $L^p$ a zbieżność według miary]
\label{thm:lp-to-measure-14}
Niech $(X,\mathcal F,\mu)$ będzie przestrzenią z miarą,
$1\le p<\infty$ i niech $f_n,f\in L^p(\mu)$.
Jeżeli $\|f_n-f\|_p\to0$, to $f_n\to f$
według miary $\mu$, tj. dla każdego $\varepsilon>0$
$\mu(\{|f_n-f|>\varepsilon\})\to0$.
Dla $p=\infty$ zachodzi silniejsze stwierdzenie:
$\|f_n-f\|_\infty\to0$ implikuje
$\mu(\{|f_n-f|>\varepsilon\})=0$ dla dostatecznie
dużych $n$.
\end{twierdzenie}
\begin{proof}
Dla $p<\infty$ nierówność Czebyszewa daje
$\varepsilon^p\mu(\{|f_n-f|>\varepsilon\})
\le\int_X|f_n-f|^p\,\mathrm d\mu(x)=\|f_n-f\|_p^p$.
Prawa strona zbiega do zera.
Dla $p=\infty$ wybieramy $n$ tak duże, że
$\|f_n-f\|_\infty<\varepsilon$ i korzystamy
z definicji kresu istotnego.
\end{proof}

\begin{przyklad}[Zbieżność według miary nie wystarcza]
Na $(0,1)$ z miarą Lebesgue'a połóżmy
$f_n(x)=n^{1/p}\chi_{(0,1/n)}(x)$ dla $p<\infty$.
Wtedy $f_n\to0$ według miary i prawie wszędzie,
ale $\|f_n\|_p=1$ dla każdego $n$.
Zbieżność według miary nie pociąga więc zbieżności
w normie, nawet na przestrzeni probabilistycznej.
\end{przyklad}

\section{Interpolacja pomiędzy normami $L^p$}

\begin{twierdzenie}[Nierówność interpolacyjna]
\label{thm:lp-interpolation-14}
Niech $(X,\mathcal F,\mu)$ będzie przestrzenią z miarą,
$1\le p<r<q\le\infty$, i niech $\theta\in(0,1)$
będzie jedyną liczbą spełniającą
$1/r=\theta/p+(1-\theta)/q$, przy czym $1/\infty=0$.
Dla każdej funkcji $f\in L^p(\mu)\cap L^q(\mu)$
zachodzi $f\in L^r(\mu)$ oraz
\begin{equation}\label{eq:lp-interpolation-14}
\|f\|_r\le\|f\|_p^{\theta}\|f\|_q^{1-\theta}.
\end{equation}
W szczególności nie zakładamy $\mu(X)<\infty$.
\end{twierdzenie}
\begin{proof}
Najpierw załóżmy $q<\infty$ i obie normy po prawej
dodatnie. Wybieramy sprzężone wykładniki
$a=p/(\theta r)>1$ oraz $b=q/((1-\theta)r)>1$;
warunek na $\theta$ jest równoważny $1/a+1/b=1$.
Z twierdzenia Höldera dla czynników
$|f|^{\theta r}$ i $|f|^{(1-\theta)r}$ otrzymujemy
\[
\int_X|f(x)|^r\,\mathrm d\mu(x)
\le\left(\int_X|f(x)|^p\,\mathrm d\mu(x)\right)^{\theta r/p}
\left(\int_X|f(x)|^q\,\mathrm d\mu(x)\right)^{(1-\theta)r/q}.
\]
Podnosząc obie strony do potęgi $1/r$, dostajemy tezę.
Gdy $q=\infty$, warunek na $\theta$ daje
$\theta r=p$. Wówczas
$|f|^r\le\|f\|_\infty^{r-p}|f|^p$
prawie wszędzie, z czego po całkowaniu wynika
\eqref{eq:lp-interpolation-14}. Jeżeli któraś
z norm po prawej jest zerowa, funkcja znika
prawie wszędzie i nierówność jest oczywista.
\end{proof}

\begin{przyklad}[Interpolacja bez skończoności miary]
Na $X=\mathbb R^d$ założenie $f\in L^1\cap L^\infty$
gwarantuje $f\in L^r$ dla wszystkich $1<r<\infty$,
a dokładniej $\|f\|_r\le\|f\|_1^{1/r}
\|f\|_\infty^{1-1/r}$. Nie wolno jednak wnioskować
z samego $f\in L^q$ o przynależności do $L^p$
przy $p<q$ bez dodatkowych założeń o mierze.
\end{przyklad}

\section{Aproksymacja i gęstość wybranych klas funkcji}

\begin{twierdzenie}[Gęstość funkcji prostych]
\label{thm:simple-dense-lp-14}
Niech $(X,\mathcal F,\mu)$ będzie dowolną przestrzenią
z miarą i $1\le p<\infty$. Wtedy w $L^p(\mu)$ gęste są
klasy funkcji prostych, przyjmujących skończenie wiele
wartości i znikających poza mierzalnym zbiorem skończonej miary.
Dla $p=\infty$ funkcje proste są gęste w $L^\infty(\mu)$,
ale bez dodatkowych założeń nie można żądać, aby
znikały poza zbiorami skończonej miary.
\end{twierdzenie}
\begin{proof}
Najpierw niech $f\ge0$ i $f\in L^p$ dla $p<\infty$.
Dla $n\ge1$ kładziemy $E_n=\{f>1/n\}$.
Z nierówności Czebyszewa mamy
$\mu(E_n)\le n^p\int_Xf^p\,\mathrm d\mu(x)<\infty$.
Funkcje $f\chi_{E_n}\chi_{\{f\le n\}}$
zbiegają punktowo do $f$ poza zbiorem, na którym
$f=\infty$ (zerowym, gdy $f\in L^p$), a ich
różnice z $f$ są zdominowane przez $f$.
Twierdzenie Lebesgue'a o zbieżności zmajoryzowanej
zastosowane do $p$-tych potęg modułów różnic pokazuje,
że obcięcia zbiegają do $f$ w $L^p$.

Dla ustalonych $n$ obcięta funkcja jest ograniczona
przez $n$ i znika poza $E_n$ o skończonej mierze.
Aproksymujemy ją z dołu funkcjami prostymi
$o$ skoku wartości co najwyżej $1/m$; błąd
jest co najwyżej $1/m$ na $E_n$, a zatem jego
norma $L^p$ jest co najwyżej
$\mu(E_n)^{1/p}/m\to0$. Dla funkcji rzeczywistej
o dowolnym znaku stosujemy to rozumowanie
osobno do jej części dodatniej i ujemnej;
obie należą do $L^p(\mu)$, a różnica ich
aproksymacji ma skończony zbiór wartości
i znika poza sumą dwóch zbiorów
skończonej miary. Dla funkcji zespolonej
wykonujemy tę konstrukcję osobno dla części
rzeczywistej i urojonej, a następnie
korzystamy z nierówności Minkowskiego.

Dla $p=\infty$ każdą funkcję istotnie ograniczoną
można poza zbiorem zerowym aproksymować
jednostajnie funkcjami prostymi, przez podział
ograniczonego prostokąta wartości w $\mathbb C$
na kwadraty o dowolnie małych średnicach.
Jeżeli $\mu(X)=\infty$, funkcja $1\in L^\infty(\mu)$,
ale odległość w normie $L^\infty$ między nią
a każdą funkcją zerową poza zbiorem skończonej
miary wynosi co najmniej $1$.
\end{proof}

\begin{twierdzenie}[Gęstość funkcji gładkich w przestrzeni euklidesowej]
\label{thm:cc-dense-lp-14}
Dla $d\ge1$ oraz $1\le p<\infty$ zbiór
$C_c^\infty(\mathbb R^d)$ jest gęsty w
$L^p(\mathbb R^d,\mathcal L_d,\lambda_d)$.
Nie jest gęsty w $L^\infty(\mathbb R^d)$.
\end{twierdzenie}
\begin{proof}
Przypadek $1\le p<\infty$ został wykazany
w rozdziale~\ref{ch:regularnosc-funkcji},
w twierdzeniu~\ref{thm:cc-in-lp-10}:
korzysta on z regularności miary Lebesgue'a,
aproksymacji funkcji charakterystycznych zbiorów mierzalnych
funkcjami ciągłymi o zwartym nośniku,
a następnie z wygładzania przez mollifikatory.
Przytoczenie tego wyniku nie wykorzystuje
żadnego twierdzenia z późniejszych rozdziałów.

Dla $p=\infty$ rozważmy funkcję stałą $1$.
Dla dowolnej $\varphi\in C_c^\infty(\mathbb R^d)$
poza zwartym nośnikiem $\varphi$ zachodzi $\varphi=0$,
a zbiór ten ma dodatnią miarę. Stąd
$\|1-\varphi\|_\infty\ge1$, więc gęstość nie zachodzi.
\end{proof}

\begin{twierdzenie}[Ciągłość translacji w $L^p(\mathbb R^d)$]
\label{thm:translation-lp-14}
Niech $d\ge1$, $1\le p<\infty$ oraz
$f\in L^p(\mathbb R^d,\mathcal L_d,\lambda_d)$.
Dla $h\in\mathbb R^d$ kładziemy
$(\tau_hf)(x)=f(x-h)$.
Wówczas $\|\tau_hf-f\|_p\to0$ przy $h\to0$.
\end{twierdzenie}
\begin{proof}
Niezmienniczość miary Lebesgue'a względem translacji
daje $\|\tau_hg\|_p=\|g\|_p$.
Jeżeli $g\in C_c^\infty(\mathbb R^d)$, to jest
jednostajnie ciągła, a dla $|h|\le1$
nośniki $g(\cdot-h)$ oraz $g$ leżą we wspólnym
zbiorze zwartym $K$ o skończonej mierze.
Mamy zatem
$\|\tau_hg-g\|_p\le
\lambda_d(K)^{1/p}\sup_x|g(x-h)-g(x)|\to0$.

Dla dowolnego $f\in L^p$ i $\varepsilon>0$
wybieramy z gęstości $g\in C_c^\infty$ tak, że
$\|f-g\|_p<\varepsilon$.
Nierówność trójkąta prowadzi do
$\|\tau_hf-f\|_p\le
\|\tau_h(f-g)\|_p+\|\tau_hg-g\|_p+\|g-f\|_p
<2\varepsilon+\|\tau_hg-g\|_p$.
Biorąc granicę górną przy $h\to0$ i następnie
$\varepsilon\downarrow0$, otrzymujemy tezę.
\end{proof}

\begin{przyklad}[Dlaczego $p=\infty$ jest inny]
Dla $f=\chi_{(0,\infty)}$ na $\mathbb R$
$\|\tau_hf-f\|_\infty=1$ dla każdego $h\ne0$.
Zatem translacje nie są silnie ciągłe w
$L^\infty(\mathbb R)$.
\end{przyklad}

\section{Separowalność: kiedy istnieje przeliczalna rodzina gęsta?}

\begin{twierdzenie}[Separowalność $L^p$]
\label{thm:separable-lp-14}
Niech $(X,\mathcal F,\mu)$ będzie przestrzenią
z miarą $\sigma$-skończoną, przy czym $\mathcal F$
jest generowana przez przeliczalną rodzinę zbiorów.
Wówczas $L^p(\mu)$ jest separowalna dla
każdego $1\le p<\infty$.
\end{twierdzenie}
\begin{proof}
Wybieramy rosnące wyczerpanie $X_n\uparrow X$
zbiorami o skończonej mierze. Do przeliczalnej
rodziny generującej $\mathcal F$ dodajemy zbiory $X_n$
i oznaczamy przez $\mathcal A$ algebrę przez nią generowaną.
Ponieważ wszystkie operacje w algebrze są skończone,
$\mathcal A$ jest przeliczalna i $\sigma(\mathcal A)=\mathcal F$.

Na każdym $X_n$ rozważamy klasę zbiorów $E\in\mathcal F|_{X_n}$,
dla których dla każdego $\varepsilon>0$ istnieje
$A\in\mathcal A|_{X_n}$ z
$\mu(E\mathbin\triangle A)<\varepsilon$.
Klasa ta jest $\sigma$-algebrą: dla dwóch zbiorów
aproksymowanych zbiorami $A,B$ z algebry zachodzi
$\mu((E\cup F)\mathbin\triangle(A\cup B))\le
\mu(E\mathbin\triangle A)+\mu(F\mathbin\triangle B)$;
analogicznie traktujemy dopełnienia względem $X_n$.
Dla $E_k\uparrow E$ mamy
$\mu(E\setminus E_k)\to0$ dzięki
$\mu(X_n)<\infty$. Stąd przeliczalne sumy
aproksymujemy najpierw skończonymi, a następnie
zbiorami z algebry. Klasa zawiera $\mathcal A|_{X_n}$,
a więc jest całą $\mathcal F|_{X_n}$.

Z twierdzenia~\ref{thm:simple-dense-lp-14}
wystarczy aproksymować funkcje proste znikające
poza zbiorem skończonej miary. Ponieważ $X_n\uparrow X$, dla każdej
$f\in L^p(\mu)$ mamy
$\int_X |f|^p\chi_{X_n}\,\mathrm d\mu
\uparrow\int_X|f|^p\,\mathrm d\mu$
na mocy twierdzenia Beppo Leviego.
Stąd $\|f\chi_{X\setminus X_n}\|_p\to0$;
w szczególności zastąpienie funkcji prostej
przez jej iloczyn z $\chi_{X_n}$ ma dowolnie mały błąd. Następnie każdą
funkcję charakterystyczną na $X_n$ aproksymujemy
w $L^p$ przez funkcje charakterystyczne zbiorów z
$\mathcal A|_{X_n}$, bo
$\|\chi_E-\chi_A\|_p
=\mu(E\mathbin\triangle A)^{1/p}$.
Wreszcie współczynniki zespolone aproksymujemy
liczbami o wymiernych częściach rzeczywistych
i urojonych. Zbiór wszystkich takich skończonych
kombinacji funkcji charakterystycznych zbiorów z przeliczalnej
algebry (z obcięciem do wybranych $X_n$) jest
przeliczalny i gęsty w $L^p$.
\end{proof}

\begin{przyklad}[Dwie przeszkody dla separowalności]
Dla $X=[0,1]$ z miarą Lebesgue'a przestrzeń
$L^\infty(X)$ nie jest separowalna: funkcje
$\chi_{(0,t)}$, $t\in(0,1)$, są parami
oddalone w normie $L^\infty$ o $1$.
Zbiór nieprzeliczalny, którego różne elementy są
od siebie oddalone co najmniej o $1$, nie może
mieścić się w przestrzeni separowalnej.

Również dla $p<\infty$ istnieją przestrzenie
nieseparowalne: jeśli $I$ jest nieprzeliczalny,
a $\kappa$ jest miarą liczącą na $\mathcal P(I)$,
to funkcje $e_i=\chi_{\{i\}}\in L^p(\kappa)$
dla $i\in I$ mają parami odległość $2^{1/p}$.
Tutaj założenie o przeliczalnym generowaniu
$\sigma$-algebry zawodzi.
\end{przyklad}

\section{Co dalej: dualność i słaba zbieżność}
Dzięki zupełności możemy traktować $L^p$ jako
przestrzenie Banacha, a $L^2$ jako przestrzeń Hilberta.
Nie przesądza to jednak ich struktury dualnej.
W następnym rozdziale zbadamy, dla jakich
wykładników i przestrzeni z miarą każdy ciągły
funkcjonał liniowy na $L^p$ można zapisać jako
$f\mapsto\int_Xf(x)g(x)\,\mathrm d\mu(x)$.
Będziemy tam precyzyjnie odróżniać
przestrzeń dualną od samego zbioru funkcji reprezentujących,
a także omówimy refleksywność i zbieżność słabą.

% ROZDZIAŁ 15 — pełny, z dowodami; nie zakłada twierdzenia Radona–Nikodyma
\chapter[Dualność i geometria przestrzeni $L^p$]{Dualność, wypukłość i słaba zbieżność w przestrzeniach $L^p$}
\label{ch:dualnosc-lp}

\section{Przestrzeń dualna i podstawowy problem reprezentacji}
Niech $(X,\mathcal F,\mu)$ będzie przestrzenią z miarą. Przestrzeń
$L^p(\mu)$ zbudowaliśmy w rozdziale~\ref{ch:przestrzenie-lp} jako
przestrzeń Banacha. Celem jest znalezienie \emph{wszystkich} ciągłych
funkcjonałów liniowych na tej przestrzeni. Należy odróżnić
wykazanie, że pewien wzór daje funkcjonał, od twierdzenia, że
każdy funkcjonał ma tę postać.

\begin{definicja}[Dualność i norma funkcjonału]
\label{def:dual-lp-15}
Dla zespolonej przestrzeni Banacha $E$ jej \emph{przestrzenią dualną}
$E^*$ nazywamy przestrzeń wszystkich ciągłych \emph{zespolono-liniowych}
odwzorowań $\Lambda:E\to\mathbb C$, z normą
$\|\Lambda\|:=\sup\{|\Lambda(f)|:\|f\|_E\le1\}$.
W całym rozdziale parowanie $L^p$ z $L^q$ zapisujemy
$\langle f,g\rangle=\int_X f(x)g(x)\,\mathrm d\mu(x)$;
\emph{nie} wstawiamy tu sprzężenia zespolonego. Jest to parowanie
dwuliniowe, a nie iloczyn skalarny na $L^2$.
\end{definicja}

\begin{lemat}[Izometryczne zanurzenie do dualnej]
\label{lem:canonical-lpq-15}
Niech $1\le p<\infty$ i $q\in(1,\infty]$ będzie wykładnikiem
sprzężonym ($1/p+1/q=1$). Jeśli $p=1$, dodatkowo
zakładamy $\sigma$-skończoność $\mu$. Dla każdego $g\in L^q(\mu)$ wzór
$\Lambda_g(f)=\int_X f(x)g(x)\,\mathrm d\mu(x)$ definiuje element
$(L^p(\mu))^*$ i zachodzi $\|\Lambda_g\|=\|g\|_q$.
Dla $1<p<\infty$ nie zakładamy $\sigma$-skończoności.
\end{lemat}
\begin{proof}
Z nierówności Höldera (twierdzenie~\ref{thm:holder-14}) mamy
$|\Lambda_g(f)|\le\|f\|_p\|g\|_q$, a więc operator jest ciągły.
Dla $1<p<\infty$, jeśli $g\ne0$, definiujemy
$f=\overline{g}|g|^{q-2}$ w punktach, gdzie $g\ne0$, i $f=0$
poza nimi. Gdy $q<2$, zapis w punktach $g=0$ zastępujemy wartością zero;
moduł tak określonej funkcji wynosi $|g|^{q-1}$.
Ponieważ $(q-1)p=q$, mamy $f\in L^p$,
$\Lambda_g(f)=\|g\|_q^q$ i $\|f\|_p=\|g\|_q^{q-1}$.
Wynika stąd nierówność odwrotna $\|\Lambda_g\|\ge\|g\|_q$.

Jeśli $p=1$, kładziemy $M=\|g\|_\infty$. Dla dowolnego
$0<a<M$ zbiór $A=\{|g|>a\}$ ma dodatnią miarę.
Z założonej $\sigma$-skończoności wynika, że istnieje
mierzalny zbiór $B\subseteq A$ o $0<\mu(B)<\infty$:
przecinamy $A$ z jednym ze zbiorów wyczerpania
$X=\bigcup_nX_n$, $\mu(X_n)<\infty$.
Wybierając $f=\chi_B\overline g/(|g|\mu(B))$ na $B$
i $f=0$ poza $B$, otrzymujemy $\|f\|_1=1$ i
$|\Lambda_g(f)|\ge a$. Przechodzimy z $a\uparrow M$.
\end{proof}

\begin{uwaga}[Dlaczego wymagane jest doprecyzowanie dla $p=1$?]
W przestrzeni z miarą, która nie jest półskończona, funkcja $g$
może mieć niezerowy kres istotny, ale całkować się do zera
z każdą funkcją $f\in L^1$. Na przykład na jednoelementowej
przestrzeni $\{*\}$ z miarą $\mu(\{*\})=\infty$ mamy
$L^1=\{0\}$ i $L^\infty\cong\mathbb C$.
Zatem ogólne twierdzenie o izometrycznym zanurzeniu z poprzedniego
lematu należy dla $p=1$ rozumieć tylko przy założeniu
$\sigma$-skończoności. Jest to istotna granica twierdzenia;
przy $1<p<\infty$ problem nie występuje.
\end{uwaga}

\section{Nierówności wypukłości w $L^p$}
W szczególności skonstruujemy reprezentację funkcjonałów
\emph{bez użycia twierdzenia Radona--Nikodyma}. Narzędziem
będzie ścisła wypukłość potęgi normy.

\begin{lemat}[Dwie nierówności dla punktu środkowego]
\label{lem:midpoint-lp-15}
Niech $f,h\in L^p(\mu)$ i $1<p<\infty$. Piszemy
$m=(f+h)/2$ oraz $v=(f-h)/2$.
\begin{enumerate}[label=\textup{(\roman*)}]
\item Jeśli $p\ge2$, to
$\|m\|_p^p+\|v\|_p^p\le(\|f\|_p^p+\|h\|_p^p)/2$.
\item Jeśli $1<p<2$ i $\|f\|_p+\|h\|_p>0$, to
\begin{equation}\label{eq:midpoint-low-15}
\frac{\|f\|_p^p+\|h\|_p^p}{2}-\|m\|_p^p
\ge \frac{p(p-1)}2\,
\frac{\|v\|_p^2}{(\|f\|_p+\|h\|_p)^{2-p}}.
\end{equation}
\end{enumerate}
\end{lemat}
\begin{proof}
Dla $p\ge2$ wystarczy wykazać dla $a,b\in\mathbb C$ nierówność
$|a+b|^p+|a-b|^p\le2^{p-1}(|a|^p+|b|^p)$.
Ponieważ $r=p/2\ge1$ i $u^r+v^r\le(u+v)^r$ dla $u,v\ge0$,
\[
|a+b|^p+|a-b|^p
\le (|a+b|^2+|a-b|^2)^{p/2}
=2^{p/2}(|a|^2+|b|^2)^{p/2}.
\]
Wypukłość $t\mapsto t^{p/2}$ daje
$(|a|^2+|b|^2)^{p/2}\le2^{p/2-1}(|a|^p+|b|^p)$.
Po podzieleniu przez $2^p$ i scałkowaniu otrzymujemy (i).

Dla $1<p<2$ identyfikujemy $\mathbb C$ z $\mathbb R^2$.
Dla $z\ne0$ hesjan funkcji $\psi(z)=|z|^p$ ma postać
$D^2\psi(z)=p|z|^{p-2}I+p(p-2)|z|^{p-4}z\otimes z$.
Najmniejsza wartość własna wynosi
$p(p-1)|z|^{p-2}$. Niech $a,b\in\mathbb C$,
$u=(a+b)/2$, $v=(a-b)/2$ i $S=|a|+|b|$.
Wzdłuż odcinka $t\mapsto u+tv$, $-1\le t\le1$,
zachodzi $|u+tv|\le S$. Dla $v\ne0$ funkcja
$t\mapsto |u+tv|^p$ ma bezwzględnie ciągłą pochodną
na $[-1,1]$. Istotnie, poza ewentualnym pojedynczym
punktem $t_0$, w którym $u+t_0v=0$, jest ona klasy $C^2$,
a jeśli taki punkt istnieje, to druga pochodna jest
w jego otoczeniu oszacowana przez wielokrotność
$|t-t_0|^{p-2}$. Jest to funkcja całkowalna, gdyż
$p-2>-1$; ponadto pierwsza pochodna jest ciągła
również w $t_0$. Można więc zastosować dwukrotnie
wzór Newtona--Leibniza na odcinku, rozbijając go
w razie potrzeby w punkcie $t_0$. Z nierówności
$D^2\psi(z)[v,v]\ge p(p-1)|z|^{p-2}|v|^2$
oraz $p-2<0$ otrzymujemy
\[
\frac{|a|^p+|b|^p}{2}-|u|^p
\ge \frac{p(p-1)}2\,|v|^2S^{p-2}.
\]
W przypadku $S=0$ obie strony rozumiemy jako zero.
Przy $S>0$ wyrażenie z prawej jest nieujemne i mierzalne.
Aby oszacować jego całkę od dołu, używamy Höldera
z wykładnikami $2/p$ i $2/(2-p)$:
\[
\int_X|v|^p\,\mathrm d\mu(x)
\le\left(\int_X |v|^2S^{p-2}\,\mathrm d\mu(x)\right)^{p/2}
\left(\int_XS^p\,\mathrm d\mu(x)\right)^{(2-p)/2}.
\]
Jeżeli pierwsza całka jest nieskończona, teza jest oczywista.
W przeciwnym razie podnosimy obie strony do potęgi $2/p$.
Minkowski zapewnia $\|S\|_p\le\|f\|_p+\|h\|_p$.
Otrzymujemy (ii).
\end{proof}

\begin{twierdzenie}[Jednostajna wypukłość]
\label{thm:uniform-convexity-15}
Dla dowolnej przestrzeni z miarą $(X,\mathcal F,\mu)$
i każdego $1<p<\infty$ przestrzeń $L^p(\mu)$ jest
\emph{jednostajnie wypukła}: dla każdego $\varepsilon>0$
istnieje $\delta>0$ taka, że jeżeli $\|f\|_p=\|h\|_p=1$
i $\|f-h\|_p\ge\varepsilon$, to
$\|(f+h)/2\|_p\le1-\delta$.
\end{twierdzenie}
\begin{proof}
Dla $p\ge2$ lemat~\ref{lem:midpoint-lp-15} daje
$\|(f+h)/2\|_p^p\le1-(\varepsilon/2)^p$.
Dla $1<p<2$ wzór \eqref{eq:midpoint-low-15} daje
$\|(f+h)/2\|_p^p\le1-c_p\varepsilon^2$,
gdzie $c_p=p(p-1)/32>0$ jest dopuszczalną stałą.
Ponieważ $\|f-h\|_p\le2$, dla $0<\varepsilon\le2$
prawa strona jest dodatnia. W obu przypadkach
możemy przyjąć $\delta=1-(1-c\varepsilon^r)^{1/p}>0$
z odpowiednimi $c,r$, jak wyżej; dla $\varepsilon>2$
warunek $\|f-h\|_p\ge\varepsilon$ jest niemożliwy.
\end{proof}

\section{Metoda wariacyjna i dualność dla $1<p<\infty$}
\begin{lemat}[Minimum energii]
\label{lem:variational-lp-15}
Niech $1<p<\infty$ i $\Lambda\in(L^p(\mu))^*$.
Funkcjonał rzeczywisty
$J(f)=\|f\|_p^p/p-\operatorname{Re}\Lambda(f)$
osiąga minimum w pewnym $u\in L^p(\mu)$.
\end{lemat}
\begin{proof}
Z oszacowania $J(f)\ge\|f\|_p^p/p-\|\Lambda\|\|f\|_p$
wynika, że $J$ jest ograniczony z dołu, a każdy ciąg
minimalizujący $(f_n)$ jest ograniczony w $L^p$:
$\|f_n\|_p\le R$ dla pewnego $R<\infty$.
Ponieważ człon liniowy zachowuje średnie, z
lematu~\ref{lem:midpoint-lp-15} wynika przy $p\ge2$
\[
\frac{J(f_n)+J(f_m)}2-J\left(\frac{f_n+f_m}{2}\right)
\ge\frac1p\left\|\frac{f_n-f_m}{2}\right\|_p^p.
\]
Dla $1<p<2$ analogiczna różnica jest co najmniej
$c(p,R)\|f_n-f_m\|_p^2$, gdzie $c(p,R)>0$;
przy $R=0$ ciąg jest zerowy. Niech $a=\inf J$.
Skoro $J((f_n+f_m)/2)\ge a$, a $J(f_n),J(f_m)\to a$,
prawe strony powyższych nierówności zmierzają do zera.
Zatem $(f_n)$ jest ciągiem Cauchy'ego w $L^p$.
Na mocy twierdzenia Riesza--Fischera
\ref{thm:riesz-fischer-14} zbiega do $u\in L^p$.
Ciągłość normy i $\Lambda$ daje $J(u)=a$.
\end{proof}

\begin{twierdzenie}[Dualność $L^p$ dla $1<p<\infty$]
\label{thm:dual-lp-15}
Niech $(X,\mathcal F,\mu)$ będzie \emph{dowolną} przestrzenią
z miarą, $1<p<\infty$, a $q=p/(p-1)$.
Wówczas każdy $\Lambda\in (L^p(\mu))^*$ ma \emph{jedyną}
reprezentację
\begin{equation}\label{eq:dual-lp-15}
\Lambda(f)=\int_X f(x)g(x)\,\mathrm d\mu(x)
\qquad(f\in L^p(\mu))
\end{equation}
przez pewne $g\in L^q(\mu)$. Zachodzi
$\|\Lambda\|=\|g\|_q$. Nie jest potrzebna
$\sigma$-skończoność ani twierdzenie Radona--Nikodyma.
\end{twierdzenie}
\begin{proof}
Z lematu~\ref{lem:variational-lp-15} istnieje minimum $u$
funkcjonału $J(f)=\|f\|_p^p/p-\operatorname{Re}\Lambda(f)$.
Dla $v\in L^p$ i $t\in\mathbb R$ funkcja $t\mapsto J(u+tv)$
ma minimum przy $t=0$. Jej pochodna w zerze wynosi
\[
\int_X |u(x)|^{p-2}\operatorname{Re}(\overline{u(x)}v(x))
\,\mathrm d\mu(x)-\operatorname{Re}\Lambda(v).
\]
Istnienie tej pochodnej uzasadniamy następująco:
dla $|t|\le1$ iloraz różnicowy $|u+tv|^p/p$ jest
oszacowany przez $(|u|+|v|)^{p-1}|v|$,
a funkcja ta jest całkowalna na mocy Höldera,
ponieważ $u,v\in L^p$. Możemy więc zastosować
twierdzenie o zbieżności zmajoryzowanej. Z minimum
wynika $\operatorname{Re}\Lambda(v)=
\operatorname{Re}\int_Xv\overline u|u|^{p-2}\,\mathrm d\mu(x)$.
Zastępując $v$ przez $iv$, otrzymujemy także
równość części urojonych. Zatem \eqref{eq:dual-lp-15}
zachodzi dla $g=\overline u|u|^{p-2}$ (przyjmujemy $g=0$
tam, gdzie $u=0$). Ponieważ $(p-1)q=p$,
$\int_X|g|^q\,\mathrm d\mu(x)=\int_X|u|^p\,\mathrm d\mu(x)<\infty$.

Lemat~\ref{lem:canonical-lpq-15} daje równość norm.
Jeśli $g_1,g_2$ reprezentują ten sam funkcjonał,
to funkcjonał $\Lambda_{g_1-g_2}$ jest zerowy,
a z izometrii wynika $\|g_1-g_2\|_q=0$.
Jest to jednoznaczność w sensie równości prawie wszędzie.
\end{proof}

\begin{przyklad}[Norma funkcjonału całkowego]
Niech $1<p<\infty$, $X=(0,1)$ z miarą Lebesgue'a
oraz $g(x)=x^{-\alpha}$ dla $0<\alpha<1/q$.
Funkcjonał $\Lambda(f)=\int_{(0,1)}f(x)x^{-\alpha}
\,\mathrm d\lambda_1(x)$ należy do $(L^p)^*$, a
$\|\Lambda\|=(1-\alpha q)^{-1/q}$.
Jeśli $\alpha q\ge1$, funkcja $g$ nie należy do $L^q$;
nie istnieje wtedy \emph{ciągłe na całym $L^p$} rozszerzenie
tego funkcjonału całkowego zachowujące tę samą postać.
\end{przyklad}

\section{Przypadek $p=1$ i szczególna rola $L^\infty$}
\begin{twierdzenie}[Dualność $L^1$ przy $\sigma$-skończoności]
\label{thm:dual-l1-15}
Jeżeli $(X,\mathcal F,\mu)$ jest $\sigma$-skończona, to
każdy $\Lambda\in(L^1(\mu))^*$ ma jednoznaczną reprezentację
$\Lambda(f)=\int_X f(x)g(x)\,\mathrm d\mu(x)$
przez $g\in L^\infty(\mu)$, a $\|\Lambda\|=\|g\|_\infty$.
\end{twierdzenie}
\begin{proof}
Wybieramy mierzalne rozłączne $E_n$, których suma jest $X$
i $0\le\mu(E_n)<\infty$. Pomijając zbiory zerowe, kładziemy
$w(x)=\sum_{n\ge1}2^{-n}(1+\mu(E_n))^{-1}\chi_{E_n}(x)$.
Mamy $w>0$ prawie wszędzie i $\int_Xw\,\mathrm d\mu(x)\le1$.
Określamy skończoną miarę $P(A)=\int_Aw\,\mathrm d\mu(x)$.
Dla $h\in L^2(P)$ funkcja $wh$ należy do $L^1(\mu)$, gdyż
$\int_X|wh|\,\mathrm d\mu(x)=\int_X|h|\,\mathrm dP(x)
\le P(X)^{1/2}\|h\|_{L^2(P)}$. Zatem
$T(h):=\Lambda(wh)$ jest ciągłym funkcjonałem na
przestrzeni Hilberta $L^2(P)$. Z twierdzenia o reprezentacji
Riesza dla przestrzeni Hilberta istnieje $g\in L^2(P)$
takie, że $T(h)=\int_X h(x)g(x)\,\mathrm dP(x)$.
Uzasadnienie wykorzystanej reprezentacji jest elementarne:
dla $T\ne0$ wybieramy $u$ prostopadły do $\ker T$ o
$T(u)=1$, po rzutowaniu ortogonalnym na domkniętą
podprzestrzeń $\ker T$; wtedy
$T(h)=\langle h,u\rangle_2/\|u\|_2^2$ przy odpowiedniej
konwencji liniowości. Dla samodzielności przypomnijmy dowód potrzebnego
przypadku twierdzenia o rzucie. Jeśli $K\subset H$ jest
domkniętą podprzestrzenią przestrzeni Hilberta,
a $h\in H$, wybieramy $z_n\in K$ tak, że
$\|h-z_n\|^2\to d^2:=\inf_{z\in K}\|h-z\|^2$.
Tożsamość równoległoboku daje
\[
\|z_n-z_m\|^2
=2\|h-z_n\|^2+2\|h-z_m\|^2
-4\left\|h-\frac{z_n+z_m}{2}\right\|^2
\le2\|h-z_n\|^2+2\|h-z_m\|^2-4d^2.
\]
Stąd $(z_n)$ jest ciągiem Cauchy'ego i ma granicę
$z\in K$. Dla każdego $k\in K$ funkcja
$t\mapsto\|h-z-tk\|^2$ ma minimum przy $t=0$;
biorąc również $tk$ z czysto urojonym współczynnikiem,
otrzymujemy $h-z\perp K$. Tak otrzymujemy potrzebny rzut
bez odwołania do zewnętrznego twierdzenia.

Dla każdego $h\in L^2(P)\cap L^1(P)$ mamy ponadto
$|T(h)|\le\|\Lambda\|\int_X|h|\,\mathrm dP(x)$.
W szczególności, podstawiając
$h=\chi_A\overline g/|g|$ (równe zero przy $g=0$),
otrzymujemy
$\int_A|g|\,\mathrm dP(x)\le\|\Lambda\|P(A)$
dla dowolnego $A\in\mathcal F$. Stąd
$|g|\le\|\Lambda\|$ $P$-prawie wszędzie; ponieważ
$w>0$ $\mu$-prawie wszędzie, także $\mu$-prawie wszędzie.

Jeśli $f\in L^1(\mu)$, to $h=f/w\in L^1(P)$
(zera $w$ pomijamy). Funkcje
$h_n=h\chi_{\{|h|\le n\}}$ należą do
$L^2(P)$, a $wh_n\to f$ w $L^1(\mu)$, co wynika
z dominacji przez $|f|$. Zatem
$\Lambda(f)=\lim_nT(h_n)=\lim_n\int_Xh_ng\,\mathrm dP(x)
=\int_Xfg\,\mathrm d\mu(x)$.
Izometrię i jednoznaczność otrzymujemy z
lematu~\ref{lem:canonical-lpq-15} w wariancie $p=1$
i z $\sigma$-skończonością.
\end{proof}

\begin{uwaga}[Dlaczego dualna $L^\infty$ jest większa?]
Nie jest prawdą, że $(L^\infty)^*=L^1$ nawet dla
miary liczącej na $\mathbb N$. Wtedy $L^\infty=\ell^\infty$.
Na podprzestrzeni $c$ ciągów zbieżnych funkcjonał
$T((x_n))=\lim_nx_n$ jest ciągły i ma normę $1$.
Twierdzenie Hahna--Banacha przedłuża go do
$\widetilde T\in(\ell^\infty)^*$ o tej samej normie.
Jeżeli byłby postaci $\widetilde T(x)=\sum_n a_nx_n$
z $(a_n)\in\ell^1$, to dla wektorów jednostkowych $e_n$
mielibyśmy $a_n=\widetilde T(e_n)=0$ (każdy $e_n\in c_0$),
lecz $\widetilde T(\chi_X)=1$. Sprzeczność pokazuje,
że taki funkcjonał nie pochodzi z $\ell^1$.
W rozdziałach o miarach ze znakiem i Radonie--Nikodymie
poznamy dodatkowo reprezentacje przez miary, lecz nie
należy utożsamiać dualnej $L^\infty$ z miarami
\emph{przeliczalnie} addytywnymi.
\end{uwaga}

\begin{wniosek}[Dokładne maksimum w nierówności Höldera]
\label{cor:holder-sharp-15}
Niech $1<p<\infty$, $q=p/(p-1)$, a
$g\in L^q(\mu)$, $g\neq0$. Istnieje
$f\in L^p(\mu)$ o $\|f\|_p=1$ takie, że
\[
\int_X f(x)g(x)\,\mathrm d\mu(x)=\|g\|_q.
\]
Można wybrać
$f(x)=\overline{g(x)}|g(x)|^{q-2}/\|g\|_q^{q-1}$
przy $g(x)\neq0$ oraz $f(x)=0$ przy $g(x)=0$.
\end{wniosek}
\begin{proof}
Mamy $|f|^p=|g|^q/\|g\|_q^q$,
ponieważ $(q-1)p=q$, zatem $\|f\|_p=1$.
Ponadto $fg=|g|^q/\|g\|_q^{q-1}$, więc
jego całka wynosi $\|g\|_q$. Te rachunki
pokazują, że stała $1$ w nierówności Höldera
jest optymalna niezależnie od przestrzeni z miarą.
\end{proof}

\section{Refleksywność, zbieżność słaba i słaba-*}
\begin{definicja}[Refleksywność i zbieżność słaba]
\label{def:weak-lp-15}
Dla przestrzeni Banacha $E$ zanurzenie kanoniczne
$J:E\to E^{**}$ jest dane przez $(Jx)(\Lambda)=\Lambda(x)$.
Przestrzeń $E$ jest \emph{refleksywna}, jeżeli $J$ jest
surjekcją. Ciąg $x_n$ \emph{zbiega słabo} do $x$,
co zapisujemy $x_n\rightharpoonup x$, jeżeli
$\Lambda(x_n)\to\Lambda(x)$ dla każdego $\Lambda\in E^*$.
Na dualnej $E^*$ mówimy o \emph{zbieżności słabej-*}
$\Lambda_n\overset{*}{\rightharpoonup}\Lambda$,
gdy $\Lambda_n(x)\to\Lambda(x)$ dla każdego $x\in E$.
Są to dwie na ogół różne topologie na $E^*$.
\end{definicja}

\begin{twierdzenie}[Refleksywność $L^p$]
\label{thm:reflexive-lp-15}
Dla dowolnej przestrzeni z miarą $(X,\mathcal F,\mu)$
i $1<p<\infty$ przestrzeń $L^p(\mu)$ jest refleksywna.
\end{twierdzenie}
\begin{proof}
Niech $q=p/(p-1)$. Z twierdzenia~\ref{thm:dual-lp-15}
izometrycznie identyfikujemy $(L^p)^*$ z $L^q$ przez
$g\mapsto\Lambda_g$, gdzie $\Lambda_g(f)=\int_Xfg\,\mathrm d\mu(x)$.
Ponieważ $1<q<\infty$, to samo twierdzenie zastosowane
z zamienionymi rolami $p,q$ utożsamia $(L^q)^*$ z $L^p$.
Pod tymi identyfikacjami kanoniczny obraz $J(f)$
działa na $g\in L^q$ jako $J(f)(g)=\int_Xfg\,\mathrm d\mu(x)$.
Każdy element $(L^q)^*$ ma tę postać, co dowodzi
surjektywności $J$.
\end{proof}

\begin{twierdzenie}[Kryterium słabej zbieżności w $L^p$]
\label{thm:weak-lp-15}
Niech $1<p<\infty$, $q=p/(p-1)$ i $f_n,f\in L^p(\mu)$.
Wtedy $f_n\rightharpoonup f$ w $L^p$ dokładnie wtedy, gdy
\[
\int_Xf_n(x)g(x)\,\mathrm d\mu(x)
\longrightarrow\int_Xf(x)g(x)\,\mathrm d\mu(x)
\quad\text{dla każdego }g\in L^q(\mu).
\]
Jeśli $f_n\to f$ w normie $L^p$, to $f_n\rightharpoonup f$.
\end{twierdzenie}
\begin{proof}
Pierwsza równoważność jest dokładnie połączeniem
definicji~\ref{def:weak-lp-15} z
reprezentacją~\ref{thm:dual-lp-15}.
Druga wynika z $|\Lambda(f_n-f)|\le\|\Lambda\|\|f_n-f\|_p$.
\end{proof}

\begin{twierdzenie}[Dolna półciągłość normy]
\label{thm:weak-lsc-15}
W dowolnej przestrzeni Banacha $E$, jeśli $x_n\rightharpoonup x$,
to $\|x\|\le\liminf_n\|x_n\|$.
\end{twierdzenie}
\begin{proof}
Dla każdego $\Lambda\in E^*$ z $\|\Lambda\|\le1$
mamy $|\Lambda(x)|=\lim_n|\Lambda(x_n)|
\le\liminf_n\|x_n\|$.
Na mocy twierdzenia Hahna--Banacha istnieje
funkcjonał normujący wektor $x$, co daje
$\|x\|\le\liminf_n\|x_n\|$.
W przypadku $E=L^p$, $1<p<\infty$, można uniknąć
Hahna--Banacha: z lematu~\ref{lem:canonical-lpq-15}
$\|x\|_p=\sup_{\|g\|_q\le1}|\int_Xxg\,\mathrm d\mu(x)|$.
\end{proof}

\begin{przyklad}[Słaba zbieżność bez zbieżności normowej]
Na $X=(0,1)$ z miarą Lebesgue'a niech
$r_n(x)=\operatorname{sign}(\sin(2^{n+1}\pi x))$,
przy dowolnym ustaleniu wartości na zbiorze miejsc zerowych.
Funkcja $r_n$ przyjmuje naprzemiennie wartości $1$ i $-1$
na kolejnych przedziałach diadycznych długości $2^{-n-1}$.
Dla każdego $1<p<\infty$ mamy $\|r_n\|_p=1$, lecz
$r_n\rightharpoonup0$ w $L^p(0,1)$. Istotnie, całka
$r_n$ po dowolnym przedziale diadycznym ustalonego
rzędu jest zerowa dla wszystkich dostatecznie dużych $n$.
W konsekwencji $\int r_ng\to0$ dla funkcji schodkowych
na przedziałach diadycznych. Ich gęstość w $L^q(0,1)$
i Hölder pozwalają przenieść zbieżność na każde $g\in L^q$.
\end{przyklad}

\begin{twierdzenie}[Podciąg słabo zbieżny: wariant separowalny]
\label{thm:weak-subseq-15}
Niech $1<p<\infty$, a $L^q(\mu)$ będzie separowalna.
Wtedy z każdego ciągu ograniczonego $(f_n)$ w $L^p(\mu)$
można wybrać podciąg słabo zbieżny do pewnego $f\in L^p(\mu)$.
\end{twierdzenie}
\begin{proof}
Niech $(g_j)_{j\ge1}$ będzie przeliczalnym gęstym
zbiorem w $L^q$. Dla każdego $j$ ciąg liczb
$\int_Xf_n g_j\,\mathrm d\mu(x)$ jest ograniczony przez
$M\|g_j\|_q$, gdzie $M=\sup_n\|f_n\|_p$.
Kolejno wybieramy podciągi, na których te liczby
są zbieżne, i przyjmujemy podciąg diagonalny
$(f_{n_k})$. Granica definiuje funkcjonał liniowy na
liniowej powłoce $(g_j)$, ograniczony przez $M\|g\|_q$.
Przedłużamy go przez ciągłość na całe $L^q$.
Z twierdzenia~\ref{thm:dual-lp-15} zastosowanego do $q$
funkcjonał ten ma postać $g\mapsto\int_Xfg\,\mathrm d\mu(x)$
dla pewnego $f\in L^p$. Gęstość i jednolite
oszacowanie $M$ zapewniają zbieżność całek
$\int_Xf_{n_k}g\to\int_Xfg$ dla wszystkich $g\in L^q$.
Kryterium~\ref{thm:weak-lp-15} daje tezę.
\end{proof}

\begin{uwaga}[Zastosowanie i ograniczenie]
W ogólnym przypadku refleksywności przestrzeni Banacha
odpowiednikiem ostatniego twierdzenia jest twierdzenie
Eberleina--\v Smuliana (zwartość słaba i zwartość ciągowa).
Nie dowodzimy go tutaj: nie jest twierdzeniem czysto
miarowym. Dowód diagonalny powyżej jest samodzielny
i wystarcza między innymi dla $L^p(\mathbb R^d)$,
$1<p<\infty$.
\end{uwaga}

\section{Zastosowanie: minimalizacja funkcjonałów całkowych}
\begin{twierdzenie}[Minimalizacja wypukłej energii]
\label{thm:minimize-lp-15}
Niech $1<p<\infty$, $g\in L^q(\mu)$ i $\alpha>0$.
Funkcjonał $\mathcal E(f)=\alpha\|f\|_p^p/p
-\operatorname{Re}\int_X f g\,\mathrm d\mu(x)$
ma dokładnie jedno minimum w $L^p(\mu)$. Jego
punkt minimalny jest dany prawie wszędzie przez
\[
f_*(x)=\alpha^{-1/(p-1)}\,\overline{g(x)}\,|g(x)|^{q-2},
\]
przy umowie $f_*(x)=0$ dla $g(x)=0$.
\end{twierdzenie}
\begin{proof}
Istnienie minimum dla $\alpha>0$ wynika z dowodu
lematu~\ref{lem:variational-lp-15} (stały dodatni
współczynnik nie zmienia argumentu). Funkcja
$z\mapsto|z|^p$ jest ściśle wypukła na $\mathbb C$,
a więc całka jest ściśle wypukła na klasach $L^p$;
stąd jednoznaczność. Warunek pierwszego rzędu daje
$\alpha|f_*|^{p-2}\overline{f_*}=g$.
Rozwiązanie tego równania w punktach $g\ne0$
ma wskazaną postać, ponieważ $q-1=1/(p-1)$;
$|f_*|^p=\alpha^{-q}|g|^q$ jest całkowalne.
To dowodzi zarówno istnienia wskazanej funkcji,
jak i spełnienia warunku minimalności (wypukłość
$\mathcal E$ sprawia, że warunek jest wystarczający).
\end{proof}

\section{Mapa zależności}
Dowiedliśmy, że dla $1<p<\infty$ kolejno zachodzą:
nierówności punktu środkowego, jednostajna wypukłość,
istnienie minimizatora, reprezentacja dualna i
refleksywność. Przypadek $p=1$ wykorzystał zamiast tego
reprezentację Hilbertowską po zmianie miary na skończoną.
Nie zakładaliśmy rozkładu Radona--Nikodyma; w rozdziale 18
zobaczymy jego szerszą wersję dla samych miar.

% ROZDZIAŁ 16 — skończone miary ze znakiem i zespolone
\chapter[Miary ze znakiem i zespolone]{Miary ze znakiem, miary zespolone i wariacja całkowita}
\label{ch:miary-ze-znakiem}

\section{Od miar dodatnich do miar przyjmujących wartości ujemne}
Dotychczas miara $\mu$ przyjmowała wartości w $[0,\infty]$.
W zastosowaniach pojawiają się różnice miar, np.
$\nu(E)=\int_E f(x)\,\mathrm d\mu(x)$ dla rzeczywistej
funkcji $f\in L^1(\mu)$, oraz miary zespolone
powstające przez całkowanie funkcji zespolonych.
W tym rozdziale zaczynamy od \emph{skończonych wartości},
aby nie tworzyć niezdefiniowanego wyrażenia $+\infty-\infty$.
O miarach ze znakiem przyjmujących jedną z wartości
nieskończonych wspomnimy osobno na końcu rozdziału;
główne twierdzenia udowadniamy dla miar skończenie wartościowych.

\begin{definicja}[Skończona miara ze znakiem]
\label{def:signed-measure-16}
Niech $(X,\mathcal F)$ będzie przestrzenią mierzalną.
\emph{Skończoną (skończenie wartościową) miarą ze znakiem}
nazywamy funkcję $\nu:\mathcal F\to\mathbb R$, dla której
$\nu(\varnothing)=0$ oraz dla dowolnych parami rozłącznych
$(E_n)_{n\ge1}\subset\mathcal F$ zachodzi
$\nu(\bigcup_{n\ge1}E_n)=\sum_{n\ge1}\nu(E_n)$.
Szereg po prawej jest wówczas zbieżny i — jako że równość
obowiązuje przy dowolnej permutacji zbiorów — bezwzględnie
zbieżny. Nie żądamy monotoniczności $\nu$.
\end{definicja}

\begin{przyklad}[Różnica skończonych miar]
Jeżeli $\mu_1,\mu_2$ są dodatnimi miarami na $\mathcal F$
i $\mu_1(X),\mu_2(X)<\infty$, to $\nu=\mu_1-\mu_2$
jest skończoną miarą ze znakiem.
Jeżeli $f\in L^1(\mu;\mathbb R)$, to
$\nu_f(E)=\int_E f\,\mathrm d\mu(x)$ jest skończoną miarą ze znakiem.
Jej przeliczalna addytywność wynika z całkowania sumy
funkcji charakterystycznych rozłącznych zbiorów, osobno dla $f^+$ i $f^-$.
\end{przyklad}

\begin{uwaga}[Ciągłość z góry i z dołu]
Dla skończonej miary ze znakiem $\nu$, jeśli
$E_n\uparrow E$, to $\nu(E_n)\to\nu(E)$, natomiast
jeśli $E_n\downarrow E$, to $\nu(E_n)\to\nu(E)$.
Pierwszy fakt otrzymujemy, rozkładając $E$ na rozłączne
przyrosty $E_1,E_2\setminus E_1,\dots$; drugi wynika
z pierwszego dla dopełnień oraz skończonej addytywności.
W odróżnieniu od miar nieujemnych nie jest potrzebne
założenie $\nu(E_1)<\infty$, bo wszystkie wartości
$\nu$ są rzeczywiste.
\end{uwaga}

\section{Dlaczego skończona miara ze znakiem jest ograniczona?}
\begin{lemat}[Ograniczoność zakresu]
\label{lem:bounded-signed-16}
Jeśli $\nu:\mathcal F\to\mathbb R$ jest skończoną
miarą ze znakiem, to $M:=\sup_{E\in\mathcal F}|\nu(E)|<\infty$.
\end{lemat}
\begin{proof}
Przypuśćmy przeciwnie, że wartości $\nu$ są nieograniczone
na podzbiorach $X$. Powiemy, że zbiór $B\in\mathcal F$
jest \emph{nieograniczony}, jeżeli
$\sup\{|\nu(A)|:A\in\mathcal F,A\subseteq B\}=\infty$.
Mamy więc nieograniczony zbiór $B_1=X$.
Wybieramy $A\subseteq B_1$ z
$|\nu(A)|>|\nu(B_1)|+2$. Ponieważ
$\nu(B_1\setminus A)=\nu(B_1)-\nu(A)$,
oba rozłączne zbiory $A$ i $B_1\setminus A$ mają
wartość $\nu$ o module większym niż $2$.
Co najmniej jeden z nich jest nieograniczony: gdyby oba
były ograniczone, dla każdego $C\subseteq B_1$ mielibyśmy
$\nu(C)=\nu(C\cap A)+\nu(C\cap(B_1\setminus A))$,
co dawałoby ograniczoność na $B_1$.
Jako $B_2$ wybieramy część nieograniczoną, a drugą
nazywamy $D_1$. Wtedy $|\nu(D_1)|>2$.
Powtarzając konstrukcję w $B_2$, otrzymujemy kolejne
nieograniczone $B_n$ i parami rozłączne $D_n$
o $|\nu(D_n)|>2$. Jednak przeliczalna addytywność
wymaga zbieżności szeregu $\sum_n\nu(D_n)$,
a więc w szczególności $\nu(D_n)\to0$.
Otrzymaliśmy sprzeczność.
\end{proof}

\section{Rozkład Hahna}
\begin{definicja}[Zbiór dodatni i ujemny]
\label{def:positive-negative-16}
Niech $\nu$ będzie skończoną miarą ze znakiem.
Zbiór $P\in\mathcal F$ nazywamy $\nu$-\emph{dodatnim},
jeżeli $\nu(A)\ge0$ dla każdego mierzalnego $A\subseteq P$.
Zbiór $N\in\mathcal F$ nazywamy $\nu$-\emph{ujemnym},
jeżeli $\nu(A)\le0$ dla każdego mierzalnego $A\subseteq N$.
Nierówność dla samego $\nu(P)$, bez kwantyfikacji po
wszystkich podzbiorach, nie wystarcza.
\end{definicja}

\begin{twierdzenie}[Rozkład Hahna]
\label{thm:hahn-16}
Dla każdej skończonej miary ze znakiem $\nu$ na
$(X,\mathcal F)$ istnieją rozłączne mierzalne zbiory
$P,N\in\mathcal F$ takie, że $X=P\cup N$,
$P$ jest $\nu$-dodatni, a $N$ jest $\nu$-ujemny.
Ponadto jeśli $(P',N')$ jest drugim takim rozkładem,
to $P\mathbin\triangle P'$ jest $\nu$-zerowy
w silnym sensie: $\nu(A)=0$ dla każdego mierzalnego
$A\subseteq P\mathbin\triangle P'$.
\end{twierdzenie}
\begin{proof}
\emph{Istnienie.} Na mocy lematu~\ref{lem:bounded-signed-16}
liczba $\alpha:=\sup\{\nu(A):A\in\mathcal F\}$ jest
skończona i $\alpha\ge0$. Dla $n\ge1$ wybieramy
$A_n\in\mathcal F$ z $\nu(A_n)>\alpha-2^{-n}$.
Dla każdego mierzalnego $B\subseteq A_n$ zachodzi
$\nu(B)\ge-2^{-n}$: w przeciwnym przypadku
$\nu(A_n\setminus B)=\nu(A_n)-\nu(B)>\alpha$,
sprzecznie z definicją $\alpha$.

Niech $F_m=\bigcup_{n\ge m}A_n$. Dla $k\ge m$,
przy dokładaniu $A_{k+1}$ do $\bigcup_{n=m}^kA_n$
dokładamy zbiór mierzalny zawarty w $A_{k+1}$,
a więc jego wartość $\nu$ jest nie mniejsza
niż $-2^{-(k+1)}$. Zatem
\[
\nu\left(\bigcup_{n=m}^kA_n\right)
\ge\nu(A_m)-\sum_{n=m+1}^k2^{-n}
>\alpha-2^{1-m}.
\]
Dla $k\to\infty$ ciągłość $\nu$ z dołu daje
$\nu(F_m)\ge\alpha-2^{1-m}$; z definicji
$\alpha$ mamy również $\nu(F_m)\le\alpha$.
Zbiory $F_m$ maleją, więc dla
$P:=\bigcap_{m\ge1}F_m$ ciągłość $\nu$ z góry daje
$\nu(P)=\lim_m\nu(F_m)=\alpha$.

Jeśli $B\subseteq P$ jest mierzalny, to
$\nu(P\setminus B)\le\alpha$ i
$\nu(B)=\nu(P)-\nu(P\setminus B)\ge0$.
Zatem $P$ jest dodatni. Jeżeli $B\subseteq N:=X\setminus P$,
to $\nu(P\cup B)=\alpha+\nu(B)\le\alpha$,
więc $\nu(B)\le0$. To dowodzi, że $N$ jest ujemny.

\emph{Jednoznaczność.} Zbiór $P\cap N'$ jest jednocześnie
podzbiorem zbioru dodatniego i ujemnego. Każdy jego
mierzalny podzbiór ma zatem miarę ze znakiem równą zero.
To samo dotyczy $N\cap P'$. Ich suma
$(P\cap N')\cup(N\cap P')=P\mathbin\triangle P'$ ma
więc wskazaną własność.
\end{proof}

\begin{przyklad}[Rozkład Hahna nie zależy tylko od znaku całkowitej masy]
Na $X=[0,1]$ z miarą Lebesgue'a połóżmy
$\nu(E)=\int_E(2x-1)\,\mathrm d\lambda_1(x)$.
Wtedy $P=[1/2,1]$ oraz $N=[0,1/2)$ tworzą rozkład
Hahna, chociaż $\nu(X)=0$. Dowolny fragment
zbioru $\{1/2\}$ można przenieść z $P$ do $N$;
nie zmienia to wartości $\nu$ na żadnym podzbiorze.
\end{przyklad}

\begin{center}
\begin{tikzpicture}[scale=0.92]
\draw[rounded corners=8pt,thick] (0,0) rectangle (9,2.55);
\fill[blue!12,rounded corners=8pt] (0.04,0.04) rectangle (4.5,2.51);
\fill[red!11,rounded corners=8pt] (4.5,0.04) rectangle (8.96,2.51);
\draw[thick] (4.5,0.05)--(4.5,2.5);
\node at (2.25,1.85) {$P$ -- zbiór dodatni};
\node at (6.75,1.85) {$N$ -- zbiór ujemny};
\node at (2.25,0.85) {$\nu(A)\ge0$ dla $A\subseteq P$};
\node at (6.75,0.85) {$\nu(B)\le0$ dla $B\subseteq N$};
\node[above] at (4.5,2.56) {$X=P\sqcup N$};
\end{tikzpicture}
\end{center}
\noindent Rysunek ilustruje kwantyfikację po \emph{wszystkich}
zbiorach mierzalnych zawartych w częściach $P$ i $N$;
nie wystarcza sam znak $\nu(P)$ oraz $\nu(N)$.

\section{Rozkład Jordana i wariacja całkowita}
\begin{definicja}[Wzajemna osobliwość]
\label{def:singular-16}
Dwie nieujemne miary $\mu_1,\mu_2$ na
$(X,\mathcal F)$ są \emph{wzajemnie osobliwe}, co zapisujemy
$\mu_1\perp\mu_2$, jeżeli istnieje $P\in\mathcal F$
taki, że $\mu_1(X\setminus P)=0$ oraz $\mu_2(P)=0$.
\end{definicja}

\begin{twierdzenie}[Rozkład Jordana]
\label{thm:jordan-16}
Niech $\nu$ będzie skończoną miarą ze znakiem.
Istnieje dokładnie jedna para skończonych miar nieujemnych
$(\nu^+,\nu^-)$ taka, że
$\nu=\nu^+-\nu^-$ i $\nu^+\perp\nu^-$.
Jeżeli $(P,N)$ jest rozkładem Hahna, to
\[
\nu^+(E)=\nu(E\cap P),\qquad
\nu^-(E)=-\nu(E\cap N)\quad(E\in\mathcal F).
\]
Ponadto $\nu^+(X)=\sup_E\nu(E)$ i
$\nu^-(X)=-\inf_E\nu(E)$.
\end{twierdzenie}
\begin{proof}
Dla rozkładu Hahna $P,N$ podane funkcje
są nieujemne i przeliczalnie addytywne, bo ich
argumenty $E\cap P$ i $E\cap N$ zachowują rozłączność.
Mamy $\nu(E)=\nu(E\cap P)+\nu(E\cap N)
=\nu^+(E)-\nu^-(E)$ oraz
$\nu^+(N)=\nu^-(P)=0$. Obie miary są skończone,
bo $\nu$ przyjmuje wartości rzeczywiste.

Jeżeli $\nu=\alpha-\beta$, gdzie $\alpha,\beta$ są
nieujemnymi miarami skończonymi i $\alpha\perp\beta$,
to istnieje $P'$ z $\alpha(X\setminus P')=0$
i $\beta(P')=0$. Każdy podzbiór $P'$ ma nieujemną
wartość $\nu$, a każdy podzbiór $X\setminus P'$ ujemną.
Jest to więc inny rozkład Hahna. Z jednoznaczności
w twierdzeniu~\ref{thm:hahn-16} zbiór
$Z=P\mathbin\triangle P'$ ma własność $\nu(C)=0$
dla każdego mierzalnego $C\subseteq Z$.
Na $P\cap N'$ mamy $\alpha=0$ i $\nu=-\beta$;
zerowość $\nu$ na podzbiorach $P\cap N'$ wymusza
$\beta(P\cap N')=0$. Na $N\cap P'$ mamy analogicznie
$\beta=0$, $\nu=\alpha$, a więc
$\alpha(N\cap P')=0$. Dla dowolnego $E\in\mathcal F$
możemy teraz policzyć
$\alpha(E)=\alpha(E\cap P')=\alpha(E\cap P\cap P')
=\nu(E\cap P\cap P')=\nu(E\cap P)=\nu^+(E)$.
Tak samo $\beta(E)=-\nu(E\cap N)=\nu^-(E)$.
Stąd jednoznaczność pary Jordana.
Wreszcie dla każdego $E$ mamy
$\nu(E)=\nu^+(E)-\nu^-(E)\le\nu^+(X)$,
a równość jest osiągana na $P$. Analogicznie
$\nu(E)\ge-\nu^-(X)$, z równością na $N$.
\end{proof}

\begin{definicja}[Wariacja całkowita]
\label{def:variation-signed-16}
Dla skończonej miary ze znakiem $\nu$ kładziemy
$|\nu|:=\nu^++\nu^-$. Jest to nieujemna miara
skończona, zwana \emph{wariacją całkowitą} $\nu$.
Liczba $\|\nu\|_{\mathrm{TV}}:=|\nu|(X)$ jest
\emph{normą w wariacji całkowitej}.
\end{definicja}

\begin{twierdzenie}[Opis wariacji przez podziały]
\label{thm:partition-variation-16}
Dla skończonej miary ze znakiem $\nu$ i $E\in\mathcal F$
zachodzi
\begin{equation}\label{eq:variation-partitions-16}
|\nu|(E)=\sup\left\{\sum_{j=1}^m|\nu(E_j)|:
E=\bigsqcup_{j=1}^mE_j,\ E_j\in\mathcal F,\ m\ge1\right\}.
\end{equation}
Ponadto $|\nu(E)|\le|\nu|(E)$ dla każdego $E$.
\end{twierdzenie}
\begin{proof}
Dla dowolnego podziału $E=\bigsqcup_jE_j$ mamy
$|\nu(E_j)|=|\nu^+(E_j)-\nu^-(E_j)|
\le\nu^+(E_j)+\nu^-(E_j)$.
Po zsumowaniu otrzymujemy oszacowanie supremum
przez $|\nu|(E)$. Podział $E=(E\cap P)\sqcup(E\cap N)$
daje sumę $\nu(E\cap P)-\nu(E\cap N)=|\nu|(E)$,
więc supremum jest osiągnięte. Ostatnia nierówność
wynika z dopuszczenia podziału jednoelementowego.
\end{proof}

\begin{wniosek}[Najmniejsza miara dominująca]
\label{cor:minimal-variation-16}
Jeśli $\rho$ jest miarą nieujemną na $(X,\mathcal F)$
i $|\nu(E)|\le\rho(E)$ dla wszystkich $E$, to
$|\nu|(E)\le\rho(E)$ dla każdego $E\in\mathcal F$.
\end{wniosek}
\begin{proof}
Dla każdego skończonego podziału
$E=\bigsqcup_jE_j$ suma $\sum_j|\nu(E_j)|$
nie przekracza $\sum_j\rho(E_j)=\rho(E)$.
Bierzemy supremum w \eqref{eq:variation-partitions-16}.
\end{proof}

\section{Miary zespolone}
\begin{definicja}[Miara zespolona]
\label{def:complex-measure-16}
Na przestrzeni mierzalnej $(X,\mathcal F)$
\emph{miarą zespoloną} nazywamy przeliczalnie addytywną
funkcję $\eta:\mathcal F\to\mathbb C$ z $\eta(\varnothing)=0$.
Nie dopuszczamy wartości $\infty$; szeregi na rozłącznych
rodzinach zbiorów są zatem bezwarunkowo zbieżne.
Jej część rzeczywista i urojona są skończonymi
miarami ze znakiem.
\end{definicja}

\begin{twierdzenie}[Wariacja miary zespolonej]
\label{thm:complex-variation-16}
Dla miary zespolonej $\eta$ definiujemy
\[
|\eta|(E)=\sup\left\{\sum_{j=1}^m|\eta(E_j)|:
E=\bigsqcup_{j=1}^mE_j,\ E_j\in\mathcal F\right\}.
\]
Jest to skończona miara nieujemna. Ponadto
$|\eta(E)|\le|\eta|(E)$ i
$|\eta|(E)\le|\operatorname{Re}\eta|(E)+
|\operatorname{Im}\eta|(E)$, gdzie po prawej występują
wariacje skończonych miar ze znakiem.
\end{twierdzenie}
\begin{proof}
Nierówność $|\eta(A)|\le|\operatorname{Re}\eta(A)|+
|\operatorname{Im}\eta(A)|$ i twierdzenie
\ref{thm:partition-variation-16} dają skończoność
$|\eta|(X)$ oraz podane oszacowanie dla każdego $E$.
Jeżeli $E,F$ są rozłączne, każdy ich skończony podział
można uszczegółowić, przecinając części z $E$ i $F$;
z nierówności trójkąta dostajemy
$|\eta|(E\cup F)\le|\eta|(E)+|\eta|(F)$.
W drugą stronę sklejamy dowolne podziały $E$ i $F$;
stąd wynika równość.

Niech teraz $E_n$ będą parami rozłączne, a
$E=\bigcup_nE_n$. Dla dowolnego skończonego
podziału $E=\bigsqcup_{j=1}^m A_j$ mamy
\[
\sum_{j=1}^m|\eta(A_j)|
=\sum_j\left|\sum_n\eta(A_j\cap E_n)\right|
\le\sum_n\sum_j|\eta(A_j\cap E_n)|
\le\sum_n|\eta|(E_n).
\]
Biorąc supremum po podziałach $E$, otrzymujemy
$|\eta|(E)\le\sum_n|\eta|(E_n)$. Z drugiej strony
skończona addytywność i monotoniczność $|\eta|$
dają $|\eta|(E)\ge\sum_{n=1}^N|\eta|(E_n)$
dla każdego $N$. Przechodząc do granicy, uzyskujemy
przeliczalną addytywność.
\end{proof}

\begin{przyklad}[Miara atomowa na $\mathbb N$]
Dla $(a_n)\in\ell^1(\mathbb N;\mathbb C)$ wzór
$\eta(A)=\sum_{n\in A}a_n$ definiuje miarę zespoloną
na $\mathcal P(\mathbb N)$. Jej wariacja jest równa
$|\eta|(A)=\sum_{n\in A}|a_n|$, a
$\|\eta\|_{\mathrm{TV}}=\sum_n|a_n|$.
Odwrotnie, każda miara zespolona na zbiorze
przeliczalnym ma tę postać: wystarczy przyjąć
$a_n=\eta(\{n\})$ i zastosować przeliczalną
addytywność oraz skończoność wariacji.
\end{przyklad}

\section{Miary o gęstości i warunek bezwzględnej całkowalności}
\begin{twierdzenie}[Wariacja miary określonej całką]
\label{thm:density-variation-16}
Niech $(X,\mathcal F,\mu)$ będzie dowolną przestrzenią
z miarą, a $f\in L^1(\mu;\mathbb C)$. Wzór
$\eta_f(E)=\int_E f(x)\,\mathrm d\mu(x)$
definiuje miarę zespoloną. Jej wariacja jest dana przez
\begin{equation}\label{eq:density-var-16}
|\eta_f|(E)=\int_E|f(x)|\,\mathrm d\mu(x)
\quad(E\in\mathcal F).
\end{equation}
Gdy $f$ jest rzeczywista, rozkład Jordana ma postać
$\eta_f^+(E)=\int_E f^+\,\mathrm d\mu(x)$ oraz
$\eta_f^-(E)=\int_E f^-\,\mathrm d\mu(x)$.
\end{twierdzenie}
\begin{proof}
Przeliczalna addytywność $\eta_f$ wynika z liniowości
całki i twierdzenia o zbieżności zmajoryzowanej dla
$\sum_{n=1}^N\chi_{E_n}f$ przy rozłącznych $E_n$.
Dla każdego skończonego podziału $E=\bigsqcup_jE_j$
\[
\sum_j\left|\int_{E_j}f\,\mathrm d\mu(x)\right|
\le\sum_j\int_{E_j}|f|\,\mathrm d\mu(x)
=\int_E|f|\,\mathrm d\mu(x).
\]
Daje to nierówność $|\eta_f|(E)\le\int_E|f|\,\mathrm d\mu(x)$.
Dla odwrotnej niech $u=\overline f/|f|$ przy $f\ne0$
i $u=0$ przy $f=0$. Wtedy $|u|\le1$ i $uf=|f|$.
Aproksymujemy $u\chi_E$ punktowo funkcjami prostymi
$u_k=\sum_jc_{k,j}\chi_{A_{k,j}}$, tak aby
$|c_{k,j}|\le1$, wszystkie $A_{k,j}\subseteq E$ i zbiory
$A_{k,j}$ były parami rozłączne.
Istnienie takich aproksymacji wynika z podziału dysku
jednostkowego na skończenie wiele zbiorów o malejących
średnicach. Mamy
\[
\left|\int_Xu_kf\,\mathrm d\mu(x)\right|
=\left|\sum_jc_{k,j}\eta_f(A_{k,j})\right|
\le\sum_j|\eta_f(A_{k,j})|\le|\eta_f|(E).
\]
Z twierdzenia o zbieżności zmajoryzowanej
$\int_Xu_kf\,\mathrm d\mu(x)\to\int_E|f|\,\mathrm d\mu(x)$.
To daje równość. Dla rzeczywistego $f$ zbiory
$\{f\ge0\}$ i $\{f<0\}$ tworzą rozkład Hahna,
a więc wzory na części Jordana wynikają z
twierdzenia~\ref{thm:jordan-16}.
\end{proof}

\begin{uwaga}[Odróżnienie warunku $\eta\ll\mu$ od istnienia gęstości]
Z wzoru \eqref{eq:density-var-16} wynika, że
$\mu(E)=0$ implikuje $\eta_f(E)=0$. W rozdziale 17
nazwiemy tę relację \emph{bezwzględną ciągłością miar}.
Nie każda miara, nawet skończona, ma gęstość względem
zadanej miary $\mu$: na przykład miara Diraca $\delta_0$
nie ma gęstości względem miary Lebesgue'a na $\mathbb R$.
Dokładne kryterium istnienia gęstości poda twierdzenie
Radona--Nikodyma w rozdziale 18.
\end{uwaga}

\section{Przestrzeń miar skończonych jako przestrzeń Banacha}
\begin{definicja}[Przestrzeń $\mathcal M(X,\mathcal F)$]
\label{def:complex-measures-space-16}
Niech $\mathcal M(X,\mathcal F)$ oznacza zespoloną
przestrzeń wektorową wszystkich miar zespolonych
na $(X,\mathcal F)$, z normą
$\|\eta\|_{\mathrm{TV}}=|\eta|(X)$.
\end{definicja}

\begin{twierdzenie}[Zupełność przestrzeni miar]
\label{thm:complete-measures-16}
Dla każdej przestrzeni mierzalnej $(X,\mathcal F)$
przestrzeń $(\mathcal M(X,\mathcal F),\|\cdot\|_{\mathrm{TV}})$
jest przestrzenią Banacha.
\end{twierdzenie}
\begin{proof}
Wariacja jest nieujemna i jednorodna. Nierówność
$\|\eta+\xi\|_{\mathrm{TV}}\le\|\eta\|_{\mathrm{TV}}+
\|\xi\|_{\mathrm{TV}}$ wynika z nierówności trójkąta
zastosowanej do sum po dowolnym podziale $X$.
Jeśli $\|\eta\|_{\mathrm{TV}}=0$, to
$|\eta(E)|\le|\eta|(E)\le|\eta|(X)=0$ dla wszystkich $E$.
Jest to więc norma.

Niech $(\eta_n)$ będzie ciągiem Cauchy'ego w tej normie.
Dla każdego $E\in\mathcal F$ z nierówności
$|\eta_n(E)-\eta_m(E)|\le\|\eta_n-\eta_m\|_{\mathrm{TV}}$
wynika istnienie granicy $\eta(E)=\lim_n\eta_n(E)$.
Zachowana jest skończona addytywność oraz
$\eta(\varnothing)=0$. Wykażemy przeliczalną addytywność.
Niech $E_j$ będą parami rozłączne i
$R_K=\bigcup_{j>K}E_j$. Dla dowolnego ustalonego $n$
\[
|\eta(R_K)|\le|\eta(R_K)-\eta_n(R_K)|+|\eta_n|(R_K).
\]
Pierwszy składnik jest nie większy niż
$\limsup_{m\to\infty}\|\eta_m-\eta_n\|_{\mathrm{TV}}$,
a drugi dąży do zera przy $K\to\infty$ z ciągłości
z góry skończonej miary $|\eta_n|$. Dobieramy najpierw
$n$ tak duże, aby pierwszy składnik był dowolnie mały,
a następnie $K$. Wtedy $\eta(R_K)\to0$,
co wraz ze skończoną addytywnością dowodzi
$\eta(\bigcup_jE_j)=\sum_j\eta(E_j)$.

Niech $\varepsilon>0$ i $n$ będzie tak duże, że
$\|\eta_m-\eta_n\|_{\mathrm{TV}}<\varepsilon$ dla $m\ge n$.
Dla każdego skończonego podziału $X=\bigsqcup_{j=1}^rA_j$
\[
\sum_j|\eta(A_j)-\eta_n(A_j)|
=\lim_{m\to\infty}\sum_j|\eta_m(A_j)-\eta_n(A_j)|
\le\varepsilon.
\]
Biorąc supremum po podziałach i korzystając z definicji
wariacji, dostajemy $\|\eta-\eta_n\|_{\mathrm{TV}}
\le\varepsilon$. Oznacza to zbieżność w normie.
\end{proof}

\begin{twierdzenie}[Oszacowanie całki względem miary zespolonej]
\label{thm:bounded-integral-complex-16}
Niech $\eta$ będzie miarą zespoloną na $(X,\mathcal F)$,
a $f:X\to\mathbb C$ funkcją ograniczoną i mierzalną.
Całka $\int_Xf\,\mathrm d\eta$, zdefiniowana przez
cztery części Jordana miar
$\operatorname{Re}\eta$ i $\operatorname{Im}\eta$, jest dobrze
określona i
$|\int_Xf\,\mathrm d\eta|\le\|f\|_{L^\infty(|\eta|)}\,|\eta|(X)$,
gdzie kres istotny liczony jest względem miary $|\eta|$.
Dokładniej, dla każdego $E\in\mathcal F$ zachodzi
$|\int_Ef\,\mathrm d\eta|\le\int_E|f|\,\mathrm d|\eta|$.
\end{twierdzenie}
\begin{proof}
Dla funkcji prostej $f=\sum_jc_j\chi_{A_j}$
o parami rozłącznych $A_j$ mamy
$|\int_Ef\,\mathrm d\eta|
=|\sum_jc_j\eta(E\cap A_j)|
\le\sum_j|c_j||\eta|(E\cap A_j)
=\int_E|f|\,\mathrm d|\eta|$.
Każdą funkcję ograniczoną mierzalną aproksymujemy
jednostajnie funkcjami prostymi. Całka po stronie
lewej jest ciągła przy tej aproksymacji, gdyż
$\eta$ jest kombinacją czterech skończonych miar
nieujemnych; również całka po prawej jest ciągła.
Przechodzimy do granicy. Ostatnie oszacowanie wynika z nierówności
$|f|\le\|f\|_{L^\infty(|\eta|)}$ prawie wszędzie
względem $|\eta|$.
\end{proof}

\begin{twierdzenie}[Dualny opis normy w wariacji całkowitej]
\label{thm:tv-dual-tests-16}
Niech $(X,\mathcal F)$ będzie przestrzenią mierzalną,
a $\eta$ skończoną miarą zespoloną na $\mathcal F$.
Wówczas
\begin{equation}\label{eq:tv-dual-tests-16}
 |\eta|(X)=\sup\left\{\left|\int_X \varphi(x)\,\mathrm d\eta(x)\right|:
 \varphi:X\to\mathbb C\ \text{mierzalna},\quad |\varphi|\le1\right\}.
\end{equation}
W supremum wystarczy dopuścić funkcje proste o module
nieprzekraczającym $1$.
\end{twierdzenie}
\begin{proof}
Z twierdzenia~\ref{thm:bounded-integral-complex-16}
każda dopuszczalna funkcja spełnia
$|\int_X\varphi(x)\,\mathrm d\eta(x)|\le|\eta|(X)$.
Dla dowolnego $\varepsilon>0$ definicja wariacji daje
skończony podział $X=\bigsqcup_{j=1}^mE_j$ z
$\sum_{j=1}^m|\eta(E_j)|>|\eta|(X)-\varepsilon$.
Ustalmy
$c_j=\overline{\eta(E_j)}/|\eta(E_j)|$, jeżeli
$\eta(E_j)\neq0$, i $c_j=0$ w przeciwnym przypadku.
Wtedy $\varphi=\sum_{j=1}^mc_j\chi_{E_j}$ jest
mierzalną funkcją prostą o module co najwyżej $1$, oraz
\[
 \int_X\varphi(x)\,\mathrm d\eta(x)
 =\sum_{j=1}^mc_j\eta(E_j)
 =\sum_{j=1}^m|\eta(E_j)|>|\eta|(X)-\varepsilon.
\]
Przechodząc z $\varepsilon\downarrow0$, otrzymujemy
odwrotną nierówność i równość supremów.
\end{proof}

\begin{uwaga}[Dlaczego jest to norma silniejsza od zbieżności na zbiorach?]
Równość~\eqref{eq:tv-dual-tests-16} pokazuje, że zbieżność
w wariacji całkowitej kontroluje całki \emph{jednostajnie}
po wszystkich mierzalnych funkcjach testowych o module
co najwyżej $1$. Sama zbieżność $\eta_n(E)\to0$ dla
każdego ustalonego zbioru $E$ nie daje takiej kontroli.
\end{uwaga}

\section{Przykłady i kontrprzykłady}
\begin{przyklad}[Miara całkowita równa zero, wariacja dodatnia]
Na $X=\{a,b\}$ z $\mathcal F=\mathcal P(X)$ połóżmy
$\nu(\{a\})=1$, $\nu(\{b\})=-1$. Wtedy
$\nu(X)=0$, ale $|\nu|(X)=2$.
Masa $\nu(X)$ nie jest więc normą miary ze znakiem.
\end{przyklad}

\begin{przyklad}[Oscylujące gęstości]
Na $X=(0,2\pi)$ z miarą Lebesgue'a definiujemy
$\nu_n(E)=\int_E\sin(nx)\,\mathrm d\lambda_1(x)$.
Wtedy $\nu_n(X)=0$ dla każdego $n$, natomiast
$\|\nu_n\|_{\mathrm{TV}}=
\int_{(0,2\pi)}|\sin(nx)|\,\mathrm d\lambda_1(x)=4$.
Zatem nawet zanikanie całkowitej masy wszystkich miar
nie kontroluje wariacji całkowitej. Poza tym dla każdej
funkcji ciągłej $f$ na $[0,2\pi]$ zachodzi
$\int f\,\mathrm d\nu(x)_n\to0$ (można to sprawdzić
najpierw dla funkcji schodkowych, a potem przez
aproksymację jednostajną). Nie zachodzi jednak
zbieżność $\nu_n\to0$ w normie wariacji całkowitej.
\end{przyklad}

\begin{uwaga}[Miary o wartościach nieskończonych]
W części literatury ``miarą ze znakiem'' nazywa się
również przeliczalnie addytywną funkcję
$\nu:\mathcal F\to[-\infty,\infty]$, która nie przyjmuje
jednocześnie wartości $+\infty$ i $-\infty$.
Dla takiej miary prawdziwe są odpowiednie wersje
rozkładów Hahna i Jordana, ale \emph{co najwyżej jedna}
z części $\nu^+$, $\nu^-$ może mieć nieskończoną masę.
Wtedy $\nu=\nu^+-\nu^-$ jest dobrze określone,
natomiast $|\nu|(X)$ nie musi być skończone.
Pełne dowody przedstawione w tym rozdziale dotyczą
skończenie wartościowych miar ze znakiem; nie wolno
bez sprawdzenia przenosić kroków opartych na
ciągłości z góry i finitości $\alpha$ na przypadek
nieskończony.
\end{uwaga}

% ROZDZIAŁ 17 — rozkład Lebesgue'a dowodzony niezależnie od RN
\chapter[Bezwzględna ciągłość i rozkład Lebesgue'a]{Bezwzględna ciągłość, osobliwość i rozkład Lebesgue'a miar}
\label{ch:ac-singular-17}

\providecommand{\atlasref}[2]{#2}


\section{Porównywanie miar: co znaczy, że jedna miara nie widzi więcej niż druga?}

\noindent\emph{Powiązania tematyczne w Bibliotece matematycznej:} \atlasref{P3-08::ch:girsanov}{P3-08: zastosowanie bezwzględnej ciągłości i równoważności miar w analizie stochastycznej}.

Niech $(X,\mathcal F)$ będzie przestrzenią mierzalną. Dwie miary mogą
przypisywać tym samym zbiorom różne wartości, ale istotne jest również,
czy mają te same zbiory miary zero. Zagadnienie to prowadzi do dwóch
przeciwstawnych pojęć: bezwzględnej ciągłości i wzajemnej osobliwości.
Na tym etapie nie będziemy zakładać istnienia gęstości jednej miary
względem drugiej: właśnie możliwość zbudowania takich gęstości
stanowi treść następnego rozdziału.

\begin{definicja}[Bezwzględna ciągłość miary]
\label{def:ac-measure-17}
Niech $\mu$ i $\nu$ będą miarami nieujemnymi na $(X,\mathcal F)$.
Mówimy, że $\nu$ jest \emph{bezwzględnie ciągła względem $\mu$},
i piszemy $\nu\ll\mu$, jeżeli dla każdego $E\in\mathcal F$
z warunku $\mu(E)=0$ wynika $\nu(E)=0$.
Dla skończonej miary zespolonej $\eta$ przyjmujemy tę samą definicję,
równoważnie $|\eta|\ll\mu$, gdzie $|\eta|$ jest wariacją całkowitą. Istotnie, gdy $\mu(E)=0$, wszystkie mierzalne podzbiory $E$ są zerowe względem $\mu$, więc wszystkie składniki sum wyznaczających $|\eta|(E)$ są zerowe. Kierunek odwrotny wynika z $|\eta(E)|\le|\eta|(E)$.
\end{definicja}

\begin{uwaga}[Dlaczego warunki dla miary zespolonej są równoważne?]
Jeżeli $\mu(E)=0$, to $\mu(A)=0$ dla każdego mierzalnego $A\subseteq E$.
Stąd $\eta(A)=0$, a z definicji wariacji jako supremum sum po podziałach
wynika $|\eta|(E)=0$. Odwrotna implikacja wynika z
$|\eta(E)|\le|\eta|(E)$. Sam warunek $\eta(E)=0$ dla jednego,
ustalonego $E$ nie oznacza jednak $|\eta|(E)=0$.
\end{uwaga}

\begin{przyklad}[Miara określona gęstością]
\label{ex:ac-density-17}
Dla dowolnej miary $\mu$ i mierzalnej funkcji $f:X\to[0,\infty]$
miara $\nu_f(E)=\int_E f(x)\,\mathrm d\mu(x)$ jest bezwzględnie
ciągła względem $\mu$, nawet jeśli $\nu_f(X)=\infty$.
W przypadku zespolonego $f\in L^1(\mu)$ analogiczna własność
wynika z twierdzenia~\ref{thm:density-variation-16}.
Nie wyciągamy jeszcze wniosku odwrotnego: $\nu\ll\mu$
samo w sobie nie jest \emph{definicją} istnienia gęstości.
\end{przyklad}

\begin{przyklad}[Miary Diraca a miara Lebesgue'a]
Na $(\mathbb R,\mathcal L_1)$ oznaczmy przez $\delta_a$ miarę Diraca
w punkcie $a\in\mathbb R$. Wówczas $\delta_a\not\ll\lambda_1$,
ponieważ $\lambda_1(\{a\})=0$, lecz $\delta_a(\{a\})=1$.
Z drugiej strony $\lambda_1\not\ll\delta_a$, gdyż
$\delta_a(\mathbb R\setminus\{a\})=0$ i
$\lambda_1(\mathbb R\setminus\{a\})=\infty$.
\end{przyklad}

\begin{definicja}[Wzajemna równoważność miar]
\label{def:equivalent-measures-17}
Dwie nieujemne miary $\mu,\nu$ na $(X,\mathcal F)$ nazywamy
\emph{równoważnymi}, co zapisujemy $\mu\sim\nu$,
jeśli $\mu\ll\nu$ oraz $\nu\ll\mu$.
Znaczy to dokładnie, że mają te same mierzalne zbiory miary zero.
\end{definicja}

\begin{przyklad}[Równoważne miary nie muszą być równe]
Na $X=(0,1)$ miary $\mu=\lambda_1|_X$ i
$\nu(E)=\int_E x\,\mathrm d\lambda_1(x)$ są równoważne,
ponieważ $x>0$ na $X$. Istotnie, $\mu(E)=0$
implikuje $\nu(E)=0$. Odwrotnie, jeśli $\nu(E)=0$,
to dla $E_n=E\cap\{x\ge1/n\}$ mamy
$\mu(E_n)\le n\int_{E_n}x\,\mathrm d\lambda_1(x)=0$.
Ponieważ $E=\bigcup_{n\ge1}E_n$, również $\mu(E)=0$.
Miary te nie są równe: $\mu(X)=1$, zaś $\nu(X)=1/2$.
Zmiana wagi może zachować wszystkie zbiory zerowe,
lecz zmienić wartości miar.
\end{przyklad}

\section{Wzajemna osobliwość i przykłady}
Pojęcie wzajemnej osobliwości wprowadziliśmy już w
definicji~\ref{def:singular-16}. Powtórzymy je w postaci dogodnej
do badania rozkładów.

Przypominamy definicję~\ref{def:singular-16}:
dwie nieujemne miary $\nu,\mu$ na $(X,\mathcal F)$ są
\emph{wzajemnie osobliwe}, $\nu\perp\mu$, gdy istnieje
$S\in\mathcal F$ takie, że $\mu(S)=0$ i
$\nu(X\setminus S)=0$. Dla miary zespolonej $\eta$
piszemy $\eta\perp\mu$, gdy $|\eta|\perp\mu$.
Jest to rozszerzenie tej samej relacji na miary zespolone,
a nie nowa definicja dla miar nieujemnych.

Zbiór $S$ w definicji jest \emph{zbiorem skupiającym}
miarę $\nu$ (można też mówić o mierzalnym nośniku). Nie należy mylić go z
\emph{nośnikiem topologicznym}: ten drugi wymaga topologii
i jest definiowany za pomocą otoczeń punktów.

\begin{lemat}[Elementarne własności osobliwości]
\label{lem:singularity-17}
Niech $\mu,\nu,\rho$ będą miarami nieujemnymi na $(X,\mathcal F)$.
Wtedy:
\begin{enumerate}[label=\textup{(\roman*)}]
\item $\nu\perp\mu$ wtedy i tylko wtedy, gdy $\mu\perp\nu$;
\item jeśli $\nu\ll\mu$ i $\nu\perp\mu$, to $\nu=0$;
\item jeśli $\nu\perp\mu$ oraz $\rho\perp\mu$, to
$\nu+\rho\perp\mu$;
\item dla dowolnego $S\in\mathcal F$ miary
$\mu|_S(E):=\mu(E\cap S)$ i
$\mu|_{X\setminus S}(E):=\mu(E\setminus S)$ są wzajemnie osobliwe.
\end{enumerate}
\end{lemat}
\begin{proof}
(i) Jeśli $\mu(S)=0$ i $\nu(X\setminus S)=0$, to
$X\setminus S$ świadczy o relacji odwrotnej.
(ii) Z $\nu\ll\mu$ wynika $\nu(S)=0$, a z osobliwości
$\nu(X\setminus S)=0$, więc $\nu(X)=0$ i $\nu=0$.
(iii) Jeśli $\nu$ jest skupiona na $S$, a $\rho$ na $T$,
to $\mu(S\cup T)=0$ i $(\nu+\rho)(X\setminus(S\cup T))=0$.
(iv) Pierwsza miara jest skupiona na $S$, a druga na $X\setminus S$.
\end{proof}

\begin{przyklad}[Miara dyskretna i miara ciągła]
Niech $(a_n)_{n\ge1}$ będzie ciągiem punktów w $\mathbb R$, a
$(c_n)_{n\ge1}$ ciągiem liczb nieujemnych o
$\sum_nc_n<\infty$. Miara $\nu=\sum_nc_n\delta_{a_n}$
jest osobliwa względem $\lambda_1$: wystarczy
$S=\{a_n:n\ge1\}$, zbiór przeliczalny, więc zerowy
względem miary Lebesgue'a. Nie ma znaczenia, czy $S$ jest
gęsty w $\mathbb R$.
\end{przyklad}

\section{Kryterium epsilon--delta i rola skończoności}
Warunek $\nu\ll\mu$ mówi jedynie o zbiorach dokładnie zerowej
miary. Gdy $\nu$ jest skończona, można go wzmocnić do
\emph{jednostajnej} kontroli małych zbiorów.

\begin{twierdzenie}[Kryterium epsilon--delta]
\label{thm:eps-delta-ac-17}
Niech $\mu$ będzie dowolną miarą nieujemną na $(X,\mathcal F)$,
a $\nu$ miarą nieujemną o $\nu(X)<\infty$.
Następujące warunki są równoważne:
\begin{enumerate}[label=\textup{(\roman*)}]
\item $\nu\ll\mu$;
\item dla każdego $\varepsilon>0$ istnieje $\delta>0$ taka,
że dla każdego $E\in\mathcal F$ zachodzi implikacja
$\mu(E)<\delta\implies\nu(E)<\varepsilon$.
\end{enumerate}
Dla skończonej miary zespolonej $\eta$ równoważność zachodzi po
zastąpieniu $\nu(E)$ przez $|\eta|(E)$.
\end{twierdzenie}
\begin{proof}
(ii)$\Rightarrow$(i). Jeśli $\mu(E)=0$, to
$\mu(E)<\delta$ dla każdej liczby $\delta>0$.
Stąd $\nu(E)<\varepsilon$ dla każdego $\varepsilon>0$,
a więc $\nu(E)=0$.

(i)$\Rightarrow$(ii). Gdyby (ii) było fałszywe, istniałoby
$\varepsilon_0>0$ takie, że dla każdego $n\ge1$
można wybrać $E_n\in\mathcal F$ z
$\mu(E_n)<2^{-n}$ oraz $\nu(E_n)\ge\varepsilon_0$.
Niech $F_m=\bigcup_{n\ge m}E_n$ i
$F=\bigcap_{m\ge1}F_m=\limsup_nE_n$.
Z przeliczalnej podaddytywności otrzymujemy
$\mu(F_m)\le\sum_{n\ge m}2^{-n}$,
zatem $\mu(F)=0$. Każdy $F_m$ zawiera $E_m$, więc
$\nu(F_m)\ge\varepsilon_0$. Ponieważ $F_m\downarrow F$
i $\nu(F_1)\le\nu(X)<\infty$, ciągłość miary z góry daje
$\nu(F)=\lim_m\nu(F_m)\ge\varepsilon_0$.
Przeczy to $\nu\ll\mu$. W przypadku zespolonym
stosujemy wynik do skończonej miary $|\eta|$.
\end{proof}

\begin{uwaga}[Dlaczego skończoności $\nu$ nie wolno pominąć?]
Na $X=(0,1)$ z miarą $\mu=\lambda_1$ określmy
$\nu(E)=0$ dla $\mu(E)=0$ oraz $\nu(E)=\infty$
dla $\mu(E)>0$. Jest to miara: dla rozłącznych
$(E_n)$ suma $\sum_n\nu(E_n)$ jest równa $0$ wtedy
i tylko wtedy, gdy wszystkie $E_n$ są zerowe; w przeciwnym
razie jest równa $\infty$. Mamy $\nu\ll\mu$,
ale nie zachodzi warunek epsilon--delta, bo przedziały
o dowolnie małej dodatniej długości mają $\nu(E)=\infty$.
\end{uwaga}

\begin{wniosek}[Ciągłość całki względem zbioru całkowania]
\label{cor:uniform-abs-cont-int-17}
Niech $f\in L^1(\mu)$ na dowolnej przestrzeni z miarą.
Wtedy dla każdego $\varepsilon>0$ istnieje $\delta>0$ taka,
że $\mu(E)<\delta$ implikuje
$\int_E|f(x)|\,\mathrm d\mu(x)<\varepsilon$.
\end{wniosek}
\begin{proof}
Miara $E\mapsto\int_E|f|\,\mathrm d\mu(x)$ jest skończona
i bezwzględnie ciągła względem $\mu$, zatem stosujemy
poprzednie twierdzenie.
\end{proof}

\section{Rozkład Lebesgue'a: przypadek miar skończonych}
\begin{twierdzenie}[Rozkład Lebesgue'a dla miar skończonych]
\label{thm:lebesgue-decomp-finite-17}
Niech $\mu$ i $\nu$ będą skończonymi miarami nieujemnymi
na tej samej przestrzeni mierzalnej $(X,\mathcal F)$.
Istnieje dokładnie jedna para miar nieujemnych
$(\nu_{\mathrm{ac}},\nu_{\mathrm{s}})$ taka, że
\begin{equation}\label{eq:decomp-finite-17}
\nu=\nu_{\mathrm{ac}}+\nu_{\mathrm{s}},\qquad
\nu_{\mathrm{ac}}\ll\mu,\qquad \nu_{\mathrm{s}}\perp\mu.
\end{equation}
Obie składowe są skończone. W dowodzie nie korzystamy
z twierdzenia Radona--Nikodyma.
\end{twierdzenie}
\begin{proof}
\emph{Konstrukcja.} Dla każdego $k\ge1$ stosujemy rozkład
Hahna (twierdzenie~\ref{thm:hahn-16}) do skończonej miary
ze znakiem $\nu-2^k\mu$. Otrzymujemy rozłączne zbiory
$P_k,N_k\in\mathcal F$ takie, że
$P_k\cup N_k=X$ oraz
\begin{equation}\label{eq:ineq-hahn-decomp-17}
\nu(A)\ge2^k\mu(A)\quad(A\subseteq P_k),\qquad
\nu(B)\le2^k\mu(B)\quad(B\subseteq N_k)
\end{equation}
dla wszystkich mierzalnych $A,B$. W szczególności
$\mu(P_k)\le2^{-k}\nu(P_k)\le2^{-k}\nu(X)$.
Definiujemy
\[
S=\limsup_{k\to\infty}P_k
  =\bigcap_{m\ge1}\bigcup_{k\ge m}P_k.
\]
Ponieważ $\sum_k\mu(P_k)<\infty$, lemat
Borela--Cantellego (lub bezpośrednia podaddytywność)
daje $\mu(S)=0$.

Dla każdego $x\in X\setminus S$ istnieje $m$ takie,
że $x\notin P_k$ dla wszystkich $k\ge m$.
W szczególności
$X\setminus S\subseteq\bigcup_{k\ge1}N_k$.
Jeśli $E\in\mathcal F$ i $\mu(E)=0$, to z drugiej
nierówności w \eqref{eq:ineq-hahn-decomp-17} wynika
$\nu(E\cap N_k)=0$ dla każdego $k$.
Zatem $\nu(E\setminus S)=0$.
Kładziemy
$\nu_{\mathrm{s}}(A)=\nu(A\cap S)$ oraz
$\nu_{\mathrm{ac}}(A)=\nu(A\setminus S)$
dla wszystkich $A\in\mathcal F$.
Obie funkcje są skończonymi miarami nieujemnymi,
a ich suma jest równa $\nu$.
Miara $\nu_{\mathrm{s}}$ jest skupiona na $S$,
przy $\mu(S)=0$, a z poprzedniego akapitu
wynika $\nu_{\mathrm{ac}}\ll\mu$.

\emph{Jednoznaczność.} Załóżmy, że
$\nu=\alpha+\sigma=\alpha'+\sigma'$,
gdzie $\alpha,\alpha'\ll\mu$ oraz
$\sigma,\sigma'\perp\mu$. Niech $T,T'$ będą
mierzalnymi zbiorami $\mu$-zerowymi skupiającymi
odpowiednio $\sigma,\sigma'$, i niech $Z=T\cup T'$.
Wówczas $\mu(Z)=0$, zatem $\alpha(Z)=\alpha'(Z)=0$.
Na $X\setminus Z$ mamy $\sigma=\sigma'=0$,
a więc dla każdego $A\in\mathcal F$
$\alpha(A\setminus Z)=\nu(A\setminus Z)=\alpha'(A\setminus Z)$.
Ponieważ obie miary $\alpha,\alpha'$ zerują się na $Z$,
wynika $\alpha=\alpha'$ i następnie $\sigma=\sigma'$.
\end{proof}

\begin{uwaga}[Dlaczego ten dowód nie jest kołowy?]
Użyliśmy wyłącznie rozkładu Hahna dla skończonych miar ze
znakiem, oszacowania szeregu $\sum_k2^{-k}$ oraz ciągłości
miary. Nie korzystaliśmy z reprezentacji przez gęstość.
Dopiero twierdzenie Radona--Nikodyma pokaże, że składową
$\nu_{\mathrm{ac}}$ można zapisać jako całkę funkcji.
\end{uwaga}

\section{Rozkład Lebesgue'a dla miar sigma-skończonych}
\begin{twierdzenie}[Rozkład Lebesgue'a w przypadku $\sigma$-skończonym]
\label{thm:lebesgue-decomp-sigfin-17}
Niech $\mu$ i $\nu$ będą miarami nieujemnymi,
$\sigma$-skończonymi na $(X,\mathcal F)$.
Wtedy istnieje dokładnie jedna para miar nieujemnych
$\nu_{\mathrm{ac}},\nu_{\mathrm{s}}$, spełniająca
$\nu=\nu_{\mathrm{ac}}+\nu_{\mathrm{s}}$,
$\nu_{\mathrm{ac}}\ll\mu$ i $\nu_{\mathrm{s}}\perp\mu$.
Obie składowe są $\sigma$-skończone.
\end{twierdzenie}
\begin{proof}
Wybieramy mierzalne pokrycia
$X=\bigcup_iA_i=\bigcup_jB_j$ takie, że
$\mu(A_i)<\infty$ i $\nu(B_j)<\infty$.
Rodzina $(A_i\cap B_j)_{i,j\ge1}$ jest przeliczalnym
pokryciem $X$ przez zbiory mające skończoną miarę
względem obu miar. Po uporządkowaniu tej rodziny
i odejmowaniu wcześniejszych zbiorów uzyskujemy
mierzalny, parami rozłączny podział $X=\bigsqcup_{n\ge1}C_n$
z $\mu(C_n),\nu(C_n)<\infty$.

Na każdym $C_n$, wyposażonym w $\sigma$-algebrę śladową,
stosujemy twierdzenie~\ref{thm:lebesgue-decomp-finite-17}.
Niech $S_n\subseteq C_n$ będzie zbiorem o
$\mu(S_n)=0$, na którym skupia się część osobliwa
ograniczenia $\nu|_{C_n}$, a na $C_n\setminus S_n$
miara $\nu|_{C_n}$ jest bezwzględnie ciągła względem
$\mu|_{C_n}$. Przyjmujemy $S=\bigcup_nS_n$.
Ponieważ $\mu(S)=0$, miara
$\nu_{\mathrm{s}}(E):=\nu(E\cap S)$ jest osobliwa.
Jeśli $\mu(E)=0$, to
$\nu((E\setminus S)\cap C_n)=0$ dla każdego $n$;
przeliczalna addytywność daje $\nu(E\setminus S)=0$.
Zatem $\nu_{\mathrm{ac}}(E):=\nu(E\setminus S)$
jest bezwzględnie ciągła względem $\mu$.
Składowe są $\sigma$-skończone, gdyż dominują je
ograniczenia $\nu$ na zbiorach $B_j$.
Jednoznaczność dowodzi się dokładnie tak jak
w przypadku skończonym, bez użycia skończoności masy.
\end{proof}

\begin{wniosek}[Rozkład miary zespolonej]
\label{cor:complex-lebesgue-decomp-17}
Niech $\mu$ będzie nieujemną miarą $\sigma$-skończoną,
a $\eta$ skończoną miarą zespoloną na $(X,\mathcal F)$.
Wtedy istnieje dokładnie jedna para miar zespolonych
$\eta_{\mathrm{ac}}$ i $\eta_{\mathrm{s}}$ taka, że
$\eta=\eta_{\mathrm{ac}}+\eta_{\mathrm{s}}$,
$\eta_{\mathrm{ac}}\ll\mu$ oraz $\eta_{\mathrm{s}}\perp\mu$.
\end{wniosek}
\begin{proof}
Stosujemy twierdzenie~\ref{thm:lebesgue-decomp-sigfin-17}
do nieujemnej skończonej miary $|\eta|$, otrzymując
zbiór $S$ o $\mu(S)=0$, taki że
$|\eta||_{X\setminus S}\ll\mu$.
Definiujemy $\eta_{\mathrm{s}}(E)=\eta(E\cap S)$
oraz $\eta_{\mathrm{ac}}(E)=\eta(E\setminus S)$.
Dla każdego $E$ wariacja ograniczenia $\eta|_S$
nie przekracza $|\eta|(E\cap S)$, co dowodzi osobliwości;
podobnie $|\eta_{\mathrm{ac}}|(E)\le|\eta|(E\setminus S)$,
co daje bezwzględną ciągłość.
Dla jednoznaczności bierzemy sumę zbiorów
$\mu$-zerowych skupiających dwie hipotetyczne części
osobliwe i powtarzamy argument z dowodu twierdzenia
\ref{thm:lebesgue-decomp-finite-17}, z uwzględnieniem
zerowości części bezwzględnie ciągłych na tych zbiorach.
\end{proof}

\section{Co może być składową osobliwą?}
\begin{przyklad}[Rozkład dyskretno-ciągły]
Na $X=(0,1)$ z miarą $\mu=\lambda_1|_X$ określmy
\[
\nu(E)=\int_E(1+x)\,\mathrm d\lambda_1(x)
       +3\delta_{1/3}(E)+2\delta_{2/3}(E).
\]
Wtedy $\nu_{\mathrm{ac}}(E)=\int_E(1+x)\,\mathrm d\lambda_1(x)$,
a $\nu_{\mathrm{s}}=3\delta_{1/3}+2\delta_{2/3}$.
Jest to rozkład Lebesgue'a: część osobliwa jest skupiona
na dwuelementowym zbiorze miary Lebesgue'a zero.
\end{przyklad}

\begin{przyklad}[Miara Cantora — osobliwość bez atomów]
\label{ex:cantor-measure-17}
Podamy konstrukcję niewymagającą wcześniejszego twierdzenia
o miarach Lebesgue'a--Stieltjesa. Dla $x\in[0,1)$ niech
\[
b_n(x)=\lfloor2^nx\rfloor-2\lfloor2^{n-1}x\rfloor\in\{0,1\},
\qquad n\ge1.
\]
Są to cyfry rozwinięcia dwójkowego, przy czym dla liczb
dwójkowo wymiernych przyjęliśmy rozwinięcie zakończone
samymi zerami. Określamy
$T(x)=\sum_{n\ge1}2b_n(x)3^{-n}$ dla $x\in[0,1)$
oraz $T(1)=1$. Każda funkcja $b_n$ jest borelowska,
a szereg definiujący $T$ jest jednostajnie zbieżny,
więc $T:[0,1]\to[0,1]$ jest borelowska. Jej wartości
należą do środkowego zbioru Cantora $C$, ponieważ ich
rozwinięcia trójkowe używają wyłącznie cyfr $0$ i $2$.

Niech $\kappa=T_*(\lambda_1|_{[0,1]})$. Jest to miara
probabilistyczna, $\kappa(C)=1$. Na poziomie $n$ konstrukcji zbioru Cantora
każdy z $2^n$ rozłącznych przedziałów podstawowych
$I_{\varepsilon_1,\ldots,\varepsilon_n}$ odpowiada
słowu $(\varepsilon_1,\ldots,\varepsilon_n)\in\{0,1\}^n$:
jego lewy koniec to $\ell=\sum_{j=1}^n2\varepsilon_j3^{-j}$,
a prawy to $\ell+3^{-n}$. Jeśli pierwsze $n$ cyfr
$b_j(x)$ tworzą to słowo, pozostały ogon szeregu dla
$T(x)$ należy do $[0,3^{-n}]$, więc $T(x)\in
I_{\varepsilon_1,\ldots,\varepsilon_n}$. Odwrotnie,
gdy słowo jest inne, $T(x)$ należy do innego przedziału
poziomu $n$, rozłącznego z rozpatrywanym.
Dla $x<1$ teleskopowanie definicji cyfr daje
$\lfloor2^nx\rfloor=\sum_{j=1}^n b_j(x)2^{n-j}$,
gdyż $\lfloor x\rfloor=0$.
Zatem przeciwobraz danego przedziału podstawowego
w $[0,1)$ jest dokładnie przedziałem dwójkowym
$[k2^{-n},(k+1)2^{-n})$, gdzie
$k=\sum_{j=1}^n\varepsilon_j2^{n-j}$.
Punkt $1$ nie wpływa na miarę przeciwobrazu. Stąd
$\kappa(I_{\varepsilon_1,\ldots,\varepsilon_n})=2^{-n}$.
Każdy $y\in C$ należy na każdym poziomie do jednego
takiego przedziału, zatem
$\kappa(\{y\})\le2^{-n}$ dla wszystkich $n$ i
$\kappa(\{y\})=0$. Dla $y\notin C$ mamy
$\kappa(\{y\})=0$, gdyż $\kappa$ jest skupiona na $C$.
Z twierdzenia \ref{thm:atomic-cont-singular-17} wynika więc, że $\kappa$ jest bezatomowa. Ponieważ
$\lambda_1(C)=0$ (zbiór $C$ zawiera się na każdym poziomie $n$ w $2^n$ przedziałach długości $3^{-n}$, o łącznej długości $(2/3)^n$), otrzymujemy $\kappa\perp\lambda_1$.
Pokazuje to, że \emph{osobliwość} nie jest tym samym co
\emph{dyskretność}. Dystrybuanta $x\mapsto\kappa((-\infty,x])$ jest ciągła dzięki brakowi atomów, niemalejąca, ma wartości $0$ i $1$ na końcach $[0,1]$, a na każdej luce zbioru Cantora jest stała. Jeśli $I=[\ell,r]$ jest $k$-tym przedziałem
podstawowym poziomu $n$, licząc od lewej od $k=0$,
to dystrybuanta przyjmuje wartości $k2^{-n}$ w $\ell$
i $(k+1)2^{-n}$ w $r$: przedziały wcześniejsze mają
łącznie masę $k2^{-n}$, a singletony mają masę zero.
Są to wartości klasycznej funkcji Cantora. Końce
przedziałów podstawowych są gęste w zbiorze Cantora,
gdyż każdy jego punkt znajduje się w przedziale
długości $3^{-n}$ na każdym poziomie. Zatem wartości
na końcach wyznaczają dystrybuantę jednoznacznie
przez ciągłość na zbiorze Cantora i stałość w lukach.
Pochodna dystrybuanty jest zerowa w każdej otwartej
luce, a suma luk ma pełną miarę Lebesgue'a w $[0,1]$.
Dystrybuanta nie jest jednak bezwzględnie ciągła:
gdyby była, wzór Newtona--Leibniza dla funkcji
bezwzględnie ciągłych
(twierdzenie \atlasref{AM1-37::it:thm:ac-lebesgue-primitive}{AM I-36: wzór Newtona--Leibniza})
dawałby różnicę wartości końcowych równą zero zamiast $1$.
\end{przyklad}

W dalszych przykładach oznaczamy funkcję Cantora przez $C$, środkowotrzeci zbiór Cantora przez $K$, a skonstruowaną miarę przez $\mu_C=\kappa$.

\begin{figure}[htbp]
\centering
\begin{tikzpicture}[x=8.4cm,y=0.85cm,font=\small]
 \foreach \n/\yy in {0/0,1/-0.65,2/-1.3,3/-1.95}{
   \node[anchor=east] at (-0.04,\yy) {$K_{\n}$};
   \draw[gray!35,line width=0.7pt] (0,\yy)--(1,\yy);
 }
 \foreach \left/\right/\yy in {0/1/0,0/{1/3}/-0.65,{2/3}/1/-0.65,
  0/{1/9}/-1.3,{2/9}/{1/3}/-1.3,{2/3}/{7/9}/-1.3,{8/9}/1/-1.3,
  0/{1/27}/-1.95,{2/27}/{3/27}/-1.95,{6/27}/{7/27}/-1.95,{8/27}/{9/27}/-1.95,
  {18/27}/{19/27}/-1.95,{20/27}/{21/27}/-1.95,{24/27}/{25/27}/-1.95,{26/27}/1/-1.95}{
    \draw[blue!65!black,line width=2.6pt] ({\left},\yy)--({\right},\yy);
 }
 \foreach \xx/\txt in {0/0,{1/3}/{1/3},{2/3}/{2/3},1/1}{
  \draw[gray!50,densely dotted] ({\xx},0.16)--({\xx},-2.18);
  \node[below] at ({\xx},-2.22) {$\txt$};
 }
 \node[anchor=west,align=left] at (1.04,-0.65) {masa $1/2$};
 \node[anchor=west,align=left] at (1.04,-1.3) {masa $1/4$};
 \node[anchor=west,align=left] at (1.04,-1.95) {masa $1/8$};
 \begin{scope}[shift={(0,-7.0)},yscale=2.7]
 \draw[->] (-0.025,0)--(1.05,0) node[right] {$t$};
 \draw[->] (0,-0.05)--(0,1.13) node[above] {$C_3(t)$};
 \foreach \yy/\txt in {0/0,{1/4}/{1/4},{1/2}/{1/2},{3/4}/{3/4},1/1}{
  \draw[gray!40,densely dotted] (0,{\yy})--(1,{\yy});
  \node[left] at (0,{\yy}) {$\txt$};
 }
 \foreach \xx/\txt in {{1/3}/{1/3},{2/3}/{2/3},1/1}{
  \node[below] at ({\xx},0) {$\txt$};
 }
 \draw[blue!65!black,line width=1.1pt]
 (0,0)--({1/27},{1/8})--({2/27},{1/8})--({3/27},{1/4})
 --({6/27},{1/4})--({7/27},{3/8})--({8/27},{3/8})--({9/27},{1/2})
 --({18/27},{1/2})--({19/27},{5/8})--({20/27},{5/8})--({21/27},{3/4})
 --({24/27},{3/4})--({25/27},{7/8})--({26/27},{7/8})--(1,1);
 \end{scope}
\end{tikzpicture}
\caption{Pierwsze etapy konstrukcji zbioru Cantora $K_n$ (u góry);
podana masa dotyczy pojedynczego niebieskiego przedziału.
U dołu: $C_3$, ciągłe odcinkami afiniczne przybliżenie funkcji Cantora $C$,
zgodne z nią na końcach przedziałów poziomu trzeciego i stałe na
usuniętych lukach. Nie jest to wykres samej funkcji granicznej.}
\label{it:fig:cantor-construction}
\end{figure}

Niech $K_n$ oznacza sumę przedziałów podstawowych
poziomu $n$ (w szczególności $K_0=[0,1]$).
Na rysunku \ref{it:fig:cantor-construction} przedstawiono
pierwsze trzy etapy. Funkcja $C_n$ zgodna z $C$ na
końcach przedziałów poziomu $n$, afiniczna na tych
przedziałach i stała na lukach, przybliża $C$
jednostajnie: na każdym przedziale podstawowym
obie funkcje przyjmują wartości pomiędzy tymi samymi
wartościami końcowymi, których różnica wynosi $2^{-n}$,
a w lukach są równe. Stąd $\|C_n-C\|_\infty\le 2^{-n}$.

\begin{twierdzenie}[Samopodobieństwo i całkowanie względem miary Cantora]
\label{it:thm:cantor-selfsimilar-integrals}
Niech $S_0(t)=t/3$ oraz $S_1(t)=(t+2)/3$ dla $t\in[0,1]$.
Wówczas
\[
 \mu_C=\tfrac12(S_0)_*\mu_C+\tfrac12(S_1)_*\mu_C.
\]
Dla każdej nieujemnej funkcji borelowskiej $f:[0,1]\to[0,+\infty]$
oraz dla każdej funkcji borelowskiej całkowalnej względem
$\mu_C$ zachodzi
\[
 \int f\,\mathrm d\mu_C=
 \tfrac12\int f(t/3)\,\mathrm d\mu_C(t)
 +\tfrac12\int f((t+2)/3)\,\mathrm d\mu_C(t).
\]
W przypadku funkcji całkowalnej oba składniki po prawej
stronie są skończonymi całkami.
\end{twierdzenie}
\begin{proof}
Dla $x\in[0,1)$ kładziemy $j=\lfloor2x\rfloor$ oraz
$u=2x-j\in[0,1)$. Z definicji cyfr dwójkowych
$b_1(x)=j$ i $b_{n+1}(x)=b_n(u)$ dla $n\ge1$:
w różnicy
$\lfloor2^{n+1}x\rfloor-2\lfloor2^nx\rfloor$
składniki zawierające $j$ znoszą się po podstawieniu
$x=(u+j)/2$. Stąd
$T(x)=(2j+T(u))/3=S_j(T(u))$.
Dla borelowskiego $A\subset[0,1]$ przeciwobraz
$T^{-1}(A)$ na połowie $[j/2,(j+1)/2)$ jest obrazem
zbioru $\{u\in[0,1): S_j(T(u))\in A\}$ przez
odwzorowanie $u\mapsto(u+j)/2$. Miara Lebesgue'a
tego obrazu jest połową miary zbioru wyjściowego.
Punkt $1$ ma miarę zero, więc sumując po $j=0,1$,
otrzymujemy
\[
 \mu_C(A)=\tfrac12\mu_C(S_0^{-1}(A))
          +\tfrac12\mu_C(S_1^{-1}(A)).
\]
To jest żądana równość miar.
Wzór całkowy dla funkcji charakterystycznych jest
dokładnie tą równością. Dla funkcji prostych nieujemnych
wynika z liniowości, a dla dowolnych funkcji nieujemnych
z aproksymacji rosnącym ciągiem funkcji prostych i
twierdzenia \atlasref{ME-08::thm:mct}{ME-08: twierdzenie Beppo Leviego} o zbieżności monotonicznej.
Dla funkcji całkowalnych stosujemy najpierw wzór
do $|f|$, aby uzyskać całkowalność obu złożeń.
Rozkład na części dodatnią i ujemną, a dla funkcji
zespolonych również na rzeczywistą i urojoną,
daje wzór w pozostałych przypadkach.
\end{proof}

\begin{przyklad}[Momenty miary osobliwej]
\label{it:ex:cantor-moments}
Oznaczmy $m_n=\int_{[0,1]}t^n\,\mathrm d\mu_C(t)$
dla $n\ge0$. Wszystkie momenty są skończone, gdyż
$0\le t^n\le1$ i $\mu_C$ jest probabilistyczna.
Mamy $m_0=1$. Z twierdzenia
\ref{it:thm:cantor-selfsimilar-integrals} i wzoru dwumianowego
dla $n\ge1$ otrzymujemy
\[
 m_n=\frac{1}{2\cdot3^n}
       \left(m_n+\sum_{k=0}^n\binom{n}{k}2^{n-k}m_k\right),
 \qquad
 m_n=\frac{\sum_{k=0}^{n-1}\binom{n}{k}2^{n-k}m_k}
               {2(3^n-1)}.
\]
W drugim wzorze przenieśliśmy oba wyrazy zawierające
$m_n$ na lewą stronę i podzieliliśmy przez dodatni
współczynnik. Rekurencja daje $m_1=2/4=1/2$,
$m_2=(4+4m_1)/16=3/8$ oraz
$m_3=(8+12m_1+6m_2)/52=5/16$.
Z kryterium całkowalności Riemanna--Stieltjesa (\atlasref{MA1-40::it:thm:rs-oscillation}{MA1-40}) oraz zgodności z całką względem miary (\atlasref{MA1-42::tw:ks-measure}{MA1-42}) dla funkcji
ciągłych całki Riemanna--Stieltjesa względem $C$
są równe całkom względem $\mu_C$ (brak masy w $0$
pozwala całkować po $[0,1]$ zamiast $(0,1]$).
W szczególności
\[
\int_0^1 1\,\mathrm dC=1,\qquad \int_0^1 x\,\mathrm dC=\frac12,\qquad \int_0^1 x^2\,\mathrm dC=\frac38.
\]
Ponadto z całkowania przez części (twierdzenie \atlasref{MA1-40::it:thm:rs-tools}{MA1-40}), ponieważ $C$ jest ciągła i ma wahanie ograniczone, wynika $2\int_0^1 C\,\mathrm dC=C(1)^2-C(0)^2=1$. Choć $C'=0$ prawie wszędzie, nie wolno zastępować $\mathrm dC$ przez $C'(x)\,\mathrm dx$: zgubilibyśmy część osobliwą miary. Wszystkie funkcje podcałkowe w tym przykładzie są
ciągłe, więc podane równości obowiązują dla całek
Riemanna--Stieltjesa, Lebesgue'a--Stieltjesa i
Kurzweila--Stieltjesa (twierdzenie \atlasref{MA1-42::tw:ks-measure}{MA1-42}).
\end{przyklad}


\begin{uwaga}[Trzy składniki klasycznego rozkładu na prostej]
Dla skończonej miary borelowskiej na $\mathbb R$ część
osobliwą można dalej rozdzielić na część atomową,
skupioną na przeliczalnym zbiorze, oraz część osobliwą
bezatomową. Zatem otrzymujemy składowe: bezwzględnie
ciągłą, osobliwą atomową i osobliwą bezatomową.
To \emph{dalszy} podział względem miary Lebesgue'a,
nie część definicji rozkładu Lebesgue'a.
\end{uwaga}

\begin{twierdzenie}[Dalszy rozkład miary osobliwej na prostej]
\label{thm:atomic-cont-singular-17}
Niech $\nu$ będzie skończoną miarą borelowską na
$\mathbb R$. Zbiór
$A=\{x\in\mathbb R:\nu(\{x\})>0\}$ jest co najwyżej
przeliczalny, a miary
\[
 \nu_{\mathrm{at}}(E):=\nu(E\cap A),\qquad
 \nu_{\mathrm{na}}(E):=\nu(E\setminus A)
 \quad(E\in\mathcal B(\mathbb R))
\]
są jednoznacznie określone. Pierwsza ma postać
$\nu_{\mathrm{at}}=\sum_{a\in A}\nu(\{a\})\delta_a$,
a druga jest bezatomowa. Jeśli $\nu\perp\lambda_1$,
to również $\nu_{\mathrm{at}}\perp\lambda_1$ i
$\nu_{\mathrm{na}}\perp\lambda_1$.
\end{twierdzenie}
\begin{proof}
Dla każdego $n\ge1$ zbiór
$A_n=\{x:\nu(\{x\})\ge1/n\}$ jest skończony:
gdyby zawierał $m$ różnych punktów, mielibyśmy
$\nu(\mathbb R)\ge m/n$. Ponieważ
$A=\bigcup_{n\ge1}A_n$, jest przeliczalny.
Przeliczalna addytywność daje podany wzór na
$\nu_{\mathrm{at}}$ oraz rozkład $\nu$.
Z definicji $\nu_{\mathrm{na}}(\{x\})=0$ dla
każdego $x$. Wykażemy, że miara ta nie ma także
żadnych atomów w sensie teorii miary.

Przypuśćmy, że mierzalny zbiór $B$ jest atomem
$\nu_{\mathrm{na}}$ o dodatniej masie $b$:
każdy mierzalny podzbiór $B$ ma masę $0$ lub $b$.
Ponieważ $\mathbb R=\bigcup_{m\ge1}[-m,m]$,
pewien domknięty ograniczony przedział $I_0$
spełnia $\nu_{\mathrm{na}}(B\cap I_0)=b$.
Dzielimy $I_0$ na dwa domknięte przedziały równej
długości. Co najmniej jeden z nich, nazwany $I_1$,
ma $\nu_{\mathrm{na}}(B\cap I_1)>0$, a więc masę $b$.
Powtarzając, otrzymujemy zstępujący ciąg przedziałów
zamkniętych $I_n$ o długościach malejących do zera,
dla których $\nu_{\mathrm{na}}(B\cap I_n)=b$.
Ich przecięcie jest singletonem $\{x\}$.
Z ciągłości skończonej miary z góry mamy
$\nu_{\mathrm{na}}(B\cap\{x\})=b>0$,
sprzecznie z zerowością wszystkich singletonów.
Zatem $\nu_{\mathrm{na}}$ jest bezatomowa.

Rozkład jest jednoznaczny również wśród wszystkich
rozkładów na dodatnią miarę skupioną na zbiorze
co najwyżej przeliczalnym i dodatnią miarę bezatomową.
Istotnie, w każdym takim rozkładzie część bezatomowa
ma zerową masę każdego singletonu. Część atomowa
musi zatem mieć w punkcie $x$ masę $\nu(\{x\})$,
a ponieważ jest skupiona na zbiorze przeliczalnym,
masy te wyznaczają ją na każdym zbiorze borelowskim.
Pozostały składnik otrzymujemy przez odjęcie.

Jeśli $\nu$ skupia się na zbiorze borelowskim $S$
o $\lambda_1(S)=0$, to obie części jako ograniczenia
$\nu$ również skupiają się na $S$, więc są osobliwe.
\end{proof}

\section{Główna idea na przyszłość}
W rozdziale~\ref{ch:radon-nikodym-18} udowodnimy, że przy
$\sigma$-skończoności obie mierzone przestrzenie można porównać
jeszcze precyzyjniej: istnieje mierzalna funkcja
$f\ge0$ taka, że
$\nu_{\mathrm{ac}}(E)=\int_E f\,\mathrm d\mu(x)$.
Twierdzenie o rozkładzie Lebesgue'a mówi zatem,
\emph{którą część miary warto próbować opisywać gęstością};
twierdzenie Radona--Nikodyma daje tę gęstość i
zapewnia jej jednoznaczność prawie wszędzie.

% ROZDZIAŁ 18 — Radon-Nikodym, pełny dowód przez przestrzeń Hilberta
\chapter[Twierdzenie Radona--Nikodyma]{Twierdzenie Radona--Nikodyma i pochodne miar}
\label{ch:radon-nikodym-18}

\providecommand{\atlasref}[2]{#2}


\section{Problem reprezentacji miary przez funkcję}
Niech $(X,\mathcal F,\mu)$ będzie przestrzenią z miarą,
a $\nu$ drugą miarą nieujemną na tej samej $\sigma$-algebrze.
Chcielibyśmy znaleźć mierzalną funkcję $f\ge0$, dla której
\begin{equation}\label{eq:RN-question-18}
\nu(E)=\int_E f(x)\,\mathrm d\mu(x)
\qquad(E\in\mathcal F).
\end{equation}
Wynikałoby stąd natychmiast $\nu\ll\mu$.
Czy warunek ten jest wystarczający? Odpowiedź jest twierdząca
przy odpowiednich założeniach o skończoności miar.

\begin{definicja}[Pochodna Radona--Nikodyma]
\label{def:RN-derivative-18}
Jeśli $\nu$ i $\mu$ są miarami nieujemnymi na $(X,\mathcal F)$
i istnieje mierzalna funkcja $f:X\to[0,\infty]$
spełniająca \eqref{eq:RN-question-18}, nazywamy ją
\emph{gęstością} lub \emph{pochodną Radona--Nikodyma}
miary $\nu$ względem $\mu$ i piszemy
$f=\mathrm d\nu/\mathrm d\mu$. Jeżeli taka pochodna
jest jednoznaczna z dokładnością do równości $\mu$-prawie
wszędzie, symbol ten oznacza jej klasę równoważności.
W tej definicji symbol $\mathrm d\nu/\mathrm d\mu$
nie jest zwykłym ilorazem dwóch funkcji.
\end{definicja}

\begin{lemat}[Całkowanie względem miary o gęstości]
\label{lem:integration-density-18}
Niech $\lambda$ będzie miarą nieujemną na $(X,\mathcal F)$,
$w:X\to[0,\infty]$ funkcją mierzalną. Wówczas wzór
$\mu(E)=\int_Ew\,\mathrm d\lambda$ dla $E\in\mathcal F$
określa miarę nieujemną na $(X,\mathcal F)$.
Ponadto dla każdej mierzalnej $\varphi:X\to[0,\infty]$
zachodzi
\begin{equation}\label{eq:integration-density-18}
\int_X\varphi(x)\,\mathrm d\mu(x)
=\int_X\varphi(x)w(x)\,\mathrm d\lambda(x),
\end{equation}
przy konwencji $0\cdot\infty=0$ dla iloczynu
nieujemnych funkcji mierzalnych. Jeśli $\varphi:X\to\mathbb C$ jest mierzalna, to
$\varphi\inL^1(\mu)$ wtedy i tylko wtedy, gdy
\[
 \int_X|\varphi|w\,\mathrm d\lambda<\infty.
\]
W tym przypadku analogiczna równość zachodzi dla całek
zespolonych. W iloczynie $\varphi w$ przyjmujemy wartość
zero tam, gdzie $\varphi=0$, także gdy $w=+\infty$;
na ewentualnym $\lambda$-zerowym zbiorze
$\{\varphi\ne0,\ w=+\infty\}$ wybieramy wartość zero
jako skończonego reprezentanta iloczynu.
\end{lemat}
\begin{proof}
Najpierw sprawdzamy, że $\mu$ jest miarą. Całka
po zbiorze pustym wynosi zero. Jeżeli $(A_j)_{j\ge1}$
jest rodziną parami rozłącznych zbiorów mierzalnych, to
przy konwencji $0\cdot\infty=0$ mamy
\[
 \sum_{j=1}^{N}w\chi_{A_j}
 \uparrow w\chi_{\bigcup_{j\ge1}A_j}.
\]
Z addytywności całki nieujemnej dla skończonych sum
oraz z twierdzenia \atlasref{ME-08::thm:mct}{ME-08: twierdzenie Beppo Leviego} o zbieżności
monotonicznej otrzymujemy
\[
 \mu\Bigl(\bigcup_{j\ge1}A_j\Bigr)
 =\lim_{N\to\infty}\sum_{j=1}^{N}
          \int_X w\chi_{A_j}\,\mathrm d\lambda
 =\sum_{j\ge1}\mu(A_j).
\]
Dowód nie wymaga $\sigma$-skończoności $\lambda$
ani skończoności $w$. Dla $\varphi=\chi_A$
tożsamość całkowa jest definicją $\mu$.
Z liniowości dla funkcji prostych nieujemnych otrzymujemy
równość dla wszystkich takich funkcji. Każdą funkcję
$\varphi\ge0$ aproksymujemy z dołu ciągiem mierzalnych
funkcji prostych $\varphi_n\uparrow\varphi$.
Dwukrotne zastosowanie twierdzenia \atlasref{ME-08::thm:mct}{ME-08: twierdzenie Beppo Leviego} o zbieżności monotonicznej
daje \eqref{eq:integration-density-18}. W przypadku
funkcji zespolonej stosujemy tę tożsamość najpierw do $|\varphi|$.
Otrzymujemy równość całek modułu, która daje obie
implikacje kryterium całkowalności. Jeżeli są skończone,
to mierzalny zbiór $\{\varphi\ne0,\ w=+\infty\}$ jest
$\lambda$-zerowy, ponieważ $|\varphi|w=+\infty$ na nim.
Zatem opisany w tezie skończony reprezentant iloczynu
jest mierzalny i całkowalny. Następnie stosujemy równość do części
dodatnich i ujemnych części rzeczywistej i urojonej $\varphi$,
które są całkowalne względem $\mu$ dzięki założeniu, oraz łączymy cztery równości.
\end{proof}

\section{Krok pierwszy: miary skończone}
\begin{twierdzenie}[Radon--Nikodym dla miar skończonych]
\label{thm:RN-finite-18}
Niech $\mu,\nu$ będą skończonymi miarami nieujemnymi
na $(X,\mathcal F)$ i niech $\nu\ll\mu$.
Istnieje mierzalna funkcja $f:X\to[0,\infty)$ taka, że
\begin{equation}\label{eq:RN-finite-18}
\nu(E)=\int_E f(x)\,\mathrm d\mu(x)
\quad\text{dla każdego }E\in\mathcal F.
\end{equation}
Funkcja $f$ należy do $L^1(\mu)$, jest jednoznaczna
$\mu$-prawie wszędzie i
$\int_Xf\,\mathrm d\mu(x)=\nu(X)$.
\end{twierdzenie}
\begin{proof}
\emph{Etap 1: wspólna miara dominująca.}
Definiujemy $\lambda=\mu+\nu$. Jest to miara skończona
na $\mathcal F$. Dla $u\in L^2(\lambda;\mathbb C)$
kładziemy
\[
T(u)=\int_X u(x)\,\mathrm d\nu(x).
\]
Funkcjonał ten jest dobrze określony: ponieważ
$\nu\le\lambda$ w sensie wartości na zbiorach,
$\int_X|u|^2\,\mathrm d\nu(x)\le
\int_X|u|^2\,\mathrm d\lambda$.
Nierówność Cauchy'ego--Schwarza daje
\[
|T(u)|\le\nu(X)^{1/2}
\left(\int_X|u|^2\,\mathrm d\nu(x)\right)^{1/2}
\le\nu(X)^{1/2}\|u\|_{L^2(\lambda)}.
\]
Jeśli dwa przedstawiciele $u$ są równe $\lambda$-prawie
wszędzie, to są równe $\nu$-prawie wszędzie, więc
$T$ jest funkcjonałem na klasach $L^2(\lambda)$.

\emph{Etap 2: reprezentacja Hilbertowska.}
Wykorzystujemy twierdzenie Riesza o reprezentacji
funkcjonałów na przestrzeni Hilberta $L^2(\lambda)$:
istnieje $h\in L^2(\lambda)$ takie, że
\begin{equation}\label{eq:RN-lambda-h-18}
T(u)=\int_Xu(x)h(x)\,\mathrm d\lambda(x)
\qquad(u\in L^2(\lambda)).
\end{equation}
Dla kompletności przypomnijmy uzasadnienie użytej
reprezentacji. Jeśli $T=0$, bierzemy $h=0$.
W przeciwnym razie $K=\ker T$ jest domkniętą
podprzestrzenią Hilberta; każdy $v$ ma jednoznaczny
rzut ortogonalny $v_K$ na $K$, którego istnienie wynika
z zupełności i tożsamości równoległoboku: ciąg
minimalizujący odległość od $K$ jest ciągiem Cauchy'ego.
Wybieramy $v\in L^2(\lambda)$ z $T(v)=1$ i rzutujemy
$v$ na $K^\perp$. Otrzymany wektor $u$ spełnia
$u\in K^\perp$ i $T(u)=1$, ponieważ $T$ zeruje się na $K$.
Dla dowolnego $v$ różnica $v-T(v)u$ należy do $K$,
więc przy iloczynie skalarnym liniowym w pierwszym
argumencie mamy
$\langle v,u\rangle=T(v)\|u\|_2^2$.
Zatem $T(v)=\langle v,u\rangle/\|u\|_2^2$,
a reprezentantem w zapisie \eqref{eq:RN-lambda-h-18}
jest $h=\overline u/\|u\|_2^2$.
Nie użyliśmy w tym miejscu twierdzenia o gęstościach miar.

\emph{Etap 3: oszacowanie $0\le h\le1$.}
Ponieważ $\lambda(X)<\infty$, funkcje charakterystyczne $\chi_E$
należą do $L^2(\lambda)$. Z \eqref{eq:RN-lambda-h-18}
otrzymujemy
\begin{equation}\label{eq:RN-mu-nu-h-18}
\nu(E)=\int_Eh\,\mathrm d\lambda,\qquad
\mu(E)=\int_E(1-h)\,\mathrm d\lambda.
\end{equation}
Pierwsza całka jest rzeczywista i nieujemna dla
każdego $E$. Wybierając kolejno zbiory, na których
$\operatorname{Im}h$ jest dodatnia lub ujemna,
dostajemy $\operatorname{Im}h=0$ $\lambda$-prawie wszędzie.
Podobnie z nieujemności obu całek w
\eqref{eq:RN-mu-nu-h-18} otrzymujemy
$h\ge0$ i $1-h\ge0$ prawie wszędzie.
Po modyfikacji na zbiorze $\lambda$-zerowym możemy
więc założyć, że $0\le h\le1$ wszędzie.

\emph{Etap 4: wykluczenie obszaru $h=1$.}
Niech $S=\{x:h(x)=1\}$. Druga równość
w \eqref{eq:RN-mu-nu-h-18} daje $\mu(S)=0$.
Z założenia $\nu\ll\mu$ wynika $\nu(S)=0$.
Zatem $\lambda(S)=\mu(S)+\nu(S)=0$.
Poza $S$ mianownik $1-h$ jest dodatni.
Definiujemy
\[
f(x)=\begin{cases}
\dfrac{h(x)}{1-h(x)},&h(x)<1,\\[4pt]
0,&h(x)=1.
\end{cases}
\]
Funkcja $f$ jest mierzalna i nieujemna. Na mocy
lematu~\ref{lem:integration-density-18}, zastosowanego
do $\mu(E)=\int_E(1-h)\,\mathrm d\lambda$,
\[
\int_E f\,\mathrm d\mu(x)
 =\int_E f(1-h)\,\mathrm d\lambda
 =\int_Eh\,\mathrm d\lambda
 =\nu(E).
\]
W środkowej równości ewentualne punkty $S$ nie
mają znaczenia, gdyż $\lambda(S)=0$.
Dla $E=X$ mamy $\int_Xf\,\mathrm d\mu(x)=\nu(X)<\infty$,
więc $f$ jest skończona $\mu$-prawie wszędzie
i należy do $L^1(\mu)$.

\emph{Etap 5: jednoznaczność.}
Jeżeli $f,g\in L^1(\mu)$, $f,g\ge0$, spełniają
\eqref{eq:RN-finite-18}, to dla każdego $E\in\mathcal F$
$\int_E(f-g)\,\mathrm d\mu(x)=0$.
Biorąc $E=\{f>g\}$, otrzymujemy
$\int_{\{f>g\}}(f-g)\,\mathrm d\mu(x)=0$.
Całkowana funkcja jest nieujemna, więc
$\mu(\{f>g\})=0$. Zamiana $f$ z $g$ daje
$\mu(\{g>f\})=0$. Stąd $f=g$ prawie wszędzie.
\end{proof}

\begin{uwaga}[Rola pomocniczej miary $\lambda=\mu+\nu$]
Dowód nie zakłada z góry, że $\nu$ ma gęstość względem
$\mu$. Najpierw reprezentujemy funkcjonał całkowy w
$L^2(\mu+\nu)$, a następnie uzyskujemy dwie gęstości
$h$ i $1-h$ względem tej samej miary $\lambda$.
Dopiero warunek $\nu\ll\mu$ pozwala wykluczyć
zbiór $\{h=1\}$ i podzielić te gęstości.
W ten sposób znika potencjalne błędne koło.
\end{uwaga}

\section{Przypadek sigma-skończony: twierdzenie główne}
\begin{twierdzenie}[Twierdzenie Radona--Nikodyma]
\label{thm:RN-sigfin-18}
Niech $\mu$ i $\nu$ będą miarami nieujemnymi,
$\sigma$-skończonymi na przestrzeni mierzalnej $(X,\mathcal F)$.
Jeśli $\nu\ll\mu$, to istnieje mierzalna funkcja
$f:X\to[0,\infty)$, skończona $\mu$-prawie wszędzie,
taka, że
\begin{equation}\label{eq:RN-general-18}
\nu(E)=\int_E f(x)\,\mathrm d\mu(x)
\qquad(E\in\mathcal F).
\end{equation}
Funkcja $f$ jest jednoznaczna $\mu$-prawie wszędzie.
Dla każdej mierzalnej funkcji $\varphi\ge0$ zachodzi
\begin{equation}\label{eq:RN-integral-change-18}
\int_X\varphi(x)\,\mathrm d\nu(x)
=\int_X\varphi(x)f(x)\,\mathrm d\mu(x).
\end{equation}
W szczególności, jeśli $\nu(X)<\infty$, to
$f\in L^1(\mu)$ i $\int_Xf\,\mathrm d\mu(x)=\nu(X)$.
\end{twierdzenie}
\begin{proof}
Ponieważ obie miary są $\sigma$-skończone,
można wybrać mierzalny, parami rozłączny podział
$X=\bigsqcup_{n\ge1}C_n$ taki, że
$\mu(C_n)<\infty$ i $\nu(C_n)<\infty$.
Jego konstrukcję podaliśmy w dowodzie twierdzenia
\ref{thm:lebesgue-decomp-sigfin-17}.
Dla ograniczeń obu miar do $C_n$ zachodzi
$\nu|_{C_n}\ll\mu|_{C_n}$. Z twierdzenia
\ref{thm:RN-finite-18} otrzymujemy mierzalną
$\mathcal F|_{C_n}$-funkcję $f_n\ge0$ taką, że
$\nu(E\cap C_n)=\int_{E\cap C_n}f_n\,\mathrm d\mu(x)$
dla wszystkich $E\in\mathcal F$.
Funkcja $f$, określona przez $f|_{C_n}=f_n$,
jest mierzalna, gdyż $C_n$ tworzą przeliczalny
podział mierzalny. Dla każdego $E$ twierdzenie
Beppo Leviego i przeliczalna addytywność $\nu$ dają
\[
\int_E f\,\mathrm d\mu(x)
=\sum_{n\ge1}\int_{E\cap C_n}f_n\,\mathrm d\mu(x)
=\sum_{n\ge1}\nu(E\cap C_n)=\nu(E).
\]
Z sumowalności $\int_{C_n}f_n\,\mathrm d\mu(x)=\nu(C_n)<\infty$
wynika, że $f$ jest skończona prawie wszędzie na każdym
$C_n$, a więc także na $X$.

Jeśli $f,g$ reprezentują $\nu$, to na każdym $C_n$
obie funkcje są całkowalne i z jednoznaczności w
przypadku skończonym są równe $\mu$-prawie wszędzie.
Po zsumowaniu przeliczalnej rodziny zbiorów wyjątkowych
otrzymujemy jednoznaczność na $X$.
Wzór \eqref{eq:RN-integral-change-18} jest zastosowaniem
lematu~\ref{lem:integration-density-18}.
\end{proof}

\begin{uwaga}[Co dokładnie zapewnia $\sigma$-skończoność?]
Założenie dotyczy \emph{obu} miar i wystarcza, aby rozbić
$X$ na przeliczalnie wiele części, na których
$\mu$ i $\nu$ są skończone. Nie jest ono warunkiem
koniecznym w każdej konkretnej przestrzeni;
istnieją szersze klasy miar, dla których reprezentacja
Radona--Nikodyma również działa. Bez dalszych warunków
nie można go jednak po prostu wykreślić.
\end{uwaga}

\begin{przyklad}[Brak $\sigma$-skończoności może zniszczyć wniosek]
\label{ex:RN-failure-nonsigma-18}
Niech $X=[0,1]$ z $\sigma$-algebrą borelowską,
$\mu$ miarą liczącą, a $\nu$ ograniczeniem miary
Lebesgue'a do zbiorów borelowskich. Miara $\mu$
nie jest $\sigma$-skończona, gdyż każdy jej zbiór
skończonej miary jest skończony, a przeliczalna suma
takich zbiorów nie pokrywa nieprzeliczalnego $X$.
Ponieważ jedynym zbiorem $\mu$-zerowym jest $\varnothing$,
mamy $\nu\ll\mu$. Przypuśćmy, że
$\nu(E)=\int_Ef\,\mathrm d\mu(x)$ dla pewnej
mierzalnej $f\ge0$. Dla każdego $x\in X$:
\[
0=\nu(\{x\})=\int_{\{x\}}f\,\mathrm d\mu(x)=f(x).
\]
Zatem $f=0$ wszędzie, wbrew
$\nu(X)=1$. Jest to nawet przykład z miarą $\mu$
\emph{półskończoną}: każdy zbiór nieskończonej miary
zawiera niepusty zbiór skończonej dodatniej miary.
\end{przyklad}

\begin{przyklad}[Dlaczego $\sigma$-skończoność samej $\mu$ nie wystarcza?]
Na $X=(0,1)$ z $\mu=\lambda_1|_X$ definiujemy
$\nu(E)=0$ dla $\mu(E)=0$ oraz
$\nu(E)=\infty$ dla $\mu(E)>0$.
Jest to miara nieujemna, bezwzględnie ciągła względem
$\mu$, lecz nie jest $\sigma$-skończona: każdy zbiór
skończonej miary $\nu$ jest $\mu$-zerowy, a przeliczalna
suma takich zbiorów nie może pokryć $X$.

Nie istnieje gęstość $f:X\to[0,\infty)$, skończona
$\mu$-prawie wszędzie, reprezentująca $\nu$.
Istotnie, gdyby istniała, to dla pewnego $m\ge1$
zbiór $E_m=\{x:f(x)\le m\}$ miałby dodatnią miarę
$\mu$, ponieważ $\bigcup_{m\ge1}E_m$ ma pełną miarę.
Wtedy
$\int_{E_m}f(x)\,\mathrm d\mu(x)\le m\mu(E_m)<\infty$,
podczas gdy $\nu(E_m)=\infty$. Jest to sprzeczność.
Dopuszczenie gęstości przyjmującej wartość $\infty$
prowadzi do innego sformułowania i nie daje wniosku
o skończoności pochodnej prawie wszędzie.
\end{przyklad}

\section{Pochodna miary ze znakiem i miary zespolonej}
\begin{twierdzenie}[Radon--Nikodym dla miar zespolonych]
\label{thm:RN-complex-18}
Niech $\mu$ będzie nieujemną miarą $\sigma$-skończoną,
a $\eta$ skończoną miarą zespoloną na $(X,\mathcal F)$.
Jeżeli $\eta\ll\mu$, to istnieje dokładnie jedna
(z dokładnością do równości $\mu$-prawie wszędzie)
funkcja $f\in L^1(\mu;\mathbb C)$, dla której
\begin{equation}\label{eq:RN-complex-18}
\eta(E)=\int_E f(x)\,\mathrm d\mu(x)
\qquad(E\in\mathcal F).
\end{equation}
Ponadto $|\eta|(E)=\int_E|f|\,\mathrm d\mu(x)$,
a w szczególności $\|\eta\|_{\mathrm{TV}}=\|f\|_1$.
Dla rzeczywistej miary ze znakiem $\eta$ funkcję $f$
można wybrać rzeczywistą.
\end{twierdzenie}
\begin{proof}
Jeśli $\eta$ jest rzeczywistą skończoną miarą ze znakiem,
rozważamy jej rozkład Jordana
$\eta=\eta^+-\eta^-$ z rozdziału~\ref{ch:miary-ze-znakiem}.
Ponieważ $\eta\ll\mu$, dla każdego mierzalnego
zbioru $E$ o $\mu(E)=0$ również wszystkie podzbiory
mierzalne $A\subseteq E$ spełniają $\eta(A)=0$.
Z wzorów na części Jordana wynika
$\eta^+(E)=\eta^-(E)=0$. Obie miary dodatnie są
skończone, a więc $\sigma$-skończone, i są
bezwzględnie ciągłe względem $\mu$.
Stosujemy twierdzenie~\ref{thm:RN-sigfin-18} do każdej
z nich, otrzymując $f_+,f_-\in L^1(\mu)$.
Wtedy $f=f_+-f_-$ jest żądaną reprezentacją.
Dla miary zespolonej stosujemy ten argument do
$\operatorname{Re}\eta$ i $\operatorname{Im}\eta$
i otrzymujemy $f=f_1+if_2\in L^1(\mu;\mathbb C)$.
Jednoznaczność wynika z jednoznaczności części
rzeczywistych i urojonych, a wzór na wariację
jest twierdzeniem~\ref{thm:density-variation-16}.
\end{proof}

\begin{wniosek}[Pełny rozkład miary na część z gęstością i osobliwą]
\label{cor:RN-lebesgue-sum-18}
Niech $\mu,\nu$ będą nieujemnymi miarami
$\sigma$-skończonymi na $(X,\mathcal F)$.
Istnieją dokładnie jedna (z dokładnością do
$\mu$-prawie wszędzie) funkcja mierzalna
$f:X\to[0,\infty)$ i dokładnie jedna miara
$\nu_{\mathrm{s}}\perp\mu$ takie, że
\begin{equation}\label{eq:RN-Lebesgue-decomp-18}
\nu(E)=\int_E f(x)\,\mathrm d\mu(x)+\nu_{\mathrm{s}}(E)
\qquad(E\in\mathcal F).
\end{equation}
Dla skończonej miary zespolonej $\eta$ i
$\sigma$-skończonej $\mu$ analogiczna reprezentacja ma
$f\in L^1(\mu;\mathbb C)$ i osobliwą miarę zespoloną.
\end{wniosek}
\begin{proof}
Twierdzenie~\ref{thm:lebesgue-decomp-sigfin-17}
daje jednoznaczny rozkład $\nu=\nu_{\mathrm{ac}}+
\nu_{\mathrm{s}}$. Obie części są $\sigma$-skończone,
a więc \ref{thm:RN-sigfin-18} daje
$\nu_{\mathrm{ac}}(E)=\int_Ef\,\mathrm d\mu(x)$.
Jednoznaczność wynika z jednoznaczności obu składników
i pochodnej. Wersję zespoloną otrzymujemy z
wniosku~\ref{cor:complex-lebesgue-decomp-17}
i twierdzenia~\ref{thm:RN-complex-18}.
\end{proof}

\begin{twierdzenie}[Rozkład biegunowy miary zespolonej]
\label{thm:polar-complex-measure-18}
Niech $\eta$ będzie skończoną miarą zespoloną
na $(X,\mathcal F)$, a $|\eta|$ jej wariacją
całkowitą. Istnieje funkcja mierzalna
$u:X\to\mathbb C$, określona jednoznacznie
$|\eta|$-prawie wszędzie, taka, że
\begin{equation}\label{eq:polar-complex-measure-18}
 |u(x)|=1\quad |\eta|\text{-prawie wszędzie},\qquad
 \eta(E)=\int_Eu(x)\,\mathrm d|\eta|(x)
 \quad(E\in\mathcal F).
\end{equation}
Jeśli $\eta$ jest rzeczywistą miarą ze znakiem, to
$u$ można wybrać o wartościach w $\{-1,1\}$.
\end{twierdzenie}
\begin{proof}
Z definicji wariacji wynika $\eta\ll|\eta|$.
Miara $|\eta|$ jest skończona, toteż twierdzenie
\ref{thm:RN-complex-18} daje $u\in L^1(|\eta|;\mathbb C)$
spełniającą drugą równość w~\eqref{eq:polar-complex-measure-18}.
Twierdzenie~\ref{thm:density-variation-16} zastosowane
do tej reprezentacji daje
$|\eta|(E)=\int_E|u(x)|\,\mathrm d|\eta|(x)$.
Z drugiej strony $|\eta|(E)=\int_E1\,\mathrm d|\eta|(x)$.
Jednoznaczność pochodnej Radona--Nikodyma implikuje
$|u|=1$ $|\eta|$-prawie wszędzie.
W przypadku rzeczywistym gęstość jest rzeczywista,
więc warunek $|u|=1$ daje $u\in\{-1,1\}$
prawie wszędzie. Ustalając wartość $1$ na pozostałym
zbiorze zerowym, otrzymujemy żądaną funkcję mierzalną.
\end{proof}

\begin{przyklad}[Rozkład biegunowy jako znak gęstości]
Niech $\eta(E)=\int_E f(x)\,\mathrm d\mu(x)$,
gdzie $f\in L^1(\mu;\mathbb C)$. Wówczas
$|\eta|(E)=\int_E|f(x)|\,\mathrm d\mu(x)$,
a funkcję z poprzedniego twierdzenia można przyjąć
równą $u=f/|f|$ na $\{f\neq0\}$ oraz $u=1$
na $\{f=0\}$. W szczególności dla rzeczywistego $f$
jest to odpowiednik znaku funkcji.
\end{przyklad}

\section{Reguły rachunku dla pochodnych miar}
\begin{twierdzenie}[Reguła łańcuchowa]
\label{thm:RN-chain-18}
Niech $\mu,\nu,\rho$ będą nieujemnymi miarami
$\sigma$-skończonymi na $(X,\mathcal F)$
i niech $\rho\ll\nu\ll\mu$.
Oznaczmy $f=\mathrm d\nu/\mathrm d\mu$
i $g=\mathrm d\rho/\mathrm d\nu$.
Wtedy $\rho\ll\mu$ i
\begin{equation}\label{eq:RN-chain-18}
\frac{\mathrm d\rho}{\mathrm d\mu}
=g f\qquad\mu\text{-prawie wszędzie},
\end{equation}
przy czym na zbiorze $\{f=0\}$ przyjmujemy $gf=0$
niezależnie od wartości reprezentanta $g$.
\end{twierdzenie}
\begin{proof}
Bezwzględna ciągłość $\rho\ll\mu$ wynika
z przechodniości implikacji dla zbiorów zerowych.
Z lematu~\ref{lem:integration-density-18}, dla każdego
$E\in\mathcal F$,
\[
\rho(E)=\int_Eg\,\mathrm d\nu(x)
=\int_Eg f\,\mathrm d\mu(x).
\]
Ponieważ zmiana $g$ na zbiorze $\nu$-zerowym nie zmienia
całki względem $\nu$, a $\nu(\{f=0\})=0$,
przyjęta konwencja jest spójna. Jednoznaczność
pochodnej daje \eqref{eq:RN-chain-18}.
\end{proof}

\begin{twierdzenie}[Pochodna odwrotna]
\label{thm:RN-inverse-18}
Niech $\mu,\nu$ będą nieujemnymi miarami
$\sigma$-skończonymi i $\mu\sim\nu$.
Jeżeli $f=\mathrm d\nu/\mathrm d\mu$, to
$f>0$ $\mu$-prawie wszędzie i
\begin{equation}\label{eq:RN-inverse-18}
\frac{\mathrm d\mu}{\mathrm d\nu}=\frac1f
\qquad\nu\text{-prawie wszędzie}.
\end{equation}
\end{twierdzenie}
\begin{proof}
Z równoważności $\mu\ll\nu$ i wzoru
$\nu(\{f=0\})=\int_{\{f=0\}}f\,\mathrm d\mu(x)=0$
wynika $\mu(\{f=0\})=0$. Zatem $1/f$ jest
zdefiniowana i skończona prawie wszędzie.
Dla każdego $E\in\mathcal F$ lemat
\ref{lem:integration-density-18} daje
\[
\int_E\frac1f\,\mathrm d\nu(x)
=\int_E\frac1f f\,\mathrm d\mu(x)=\mu(E).
\]
Jednoznaczność pochodnej kończy dowód.
\end{proof}

\begin{wniosek}[Różnica i suma miar z gęstościami]
\label{cor:RN-linearity-18}
Niech $\mu$ będzie miarą $\sigma$-skończoną,
a $\eta_1,\eta_2$ skończonymi miarami zespolonymi,
bezwzględnie ciągłymi względem $\mu$.
Dla $a,b\in\mathbb C$ zachodzi
\[
\frac{\mathrm d(a\eta_1+b\eta_2)}{\mathrm d\mu}
=a\frac{\mathrm d\eta_1}{\mathrm d\mu}
+b\frac{\mathrm d\eta_2}{\mathrm d\mu}
\quad\mu\text{-prawie wszędzie}.
\]
\end{wniosek}
\begin{proof}
Wstawiamy reprezentacje całkowe z twierdzenia
\ref{thm:RN-complex-18}, korzystamy z liniowości całki,
a następnie z jednoznaczności gęstości.
\end{proof}

\section{Zastosowanie: warunkowa wartość oczekiwana}

\noindent\emph{Powiązanie tematyczne w Bibliotece matematycznej:} \atlasref{P1-10::ch:condexp}{P1-10: warunkowa wartość oczekiwana i rzut ortogonalny}.

Przedstawimy zastosowanie pokazujące, że twierdzenie
Radona--Nikodyma nie jest jedynie wynikiem o porównywaniu
miar. Pozwala ono skonstruować jeden z podstawowych
operatorów teorii prawdopodobieństwa.

\begin{definicja}[Warunkowa wartość oczekiwana]
\label{def:conditional-exp-18}
Niech $(\Omega,\mathcal F,\mathbb P)$ będzie przestrzenią
probabilistyczną, $\mathcal G\subseteq\mathcal F$
pod-$\sigma$-algebrą, a $Y\in L^1(\mathbb P;\mathbb R)$.
\emph{Warunkową wartością oczekiwaną} $Y$ względem
$\mathcal G$ nazywamy klasę funkcji
$Z\in L^1(\mathbb P;\mathbb R)$, która jest
$\mathcal G$-mierzalna i spełnia
\begin{equation}\label{eq:CE-def-18}
\int_A Z(\omega)\,\mathrm d\mathbb P(\omega)
=\int_A Y(\omega)\,\mathrm d\mathbb P(\omega)
\qquad(A\in\mathcal G).
\end{equation}
Piszemy $Z=\mathbb E[Y\mid\mathcal G]$.
Równość jest rozumiana z dokładnością do
$\mathbb P$-prawie wszystkich punktów. Ściślej, pracujemy
z klasami funkcji mierzalnych względem $\mathcal G$;
nie każda modyfikacja na zbiorze $\mathbb P$-zerowym musi
pozostać $\mathcal G$-mierzalna, gdy $\mathcal G$ nie jest zupełna.
\end{definicja}

\begin{twierdzenie}[Istnienie, jednoznaczność i własności podstawowe]
\label{thm:CE-existence-18}
Dla każdej przestrzeni probabilistycznej
$(\Omega,\mathcal F,\mathbb P)$, każdej
pod-$\sigma$-algebry $\mathcal G$ oraz $Y\in L^1(\mathbb P;\mathbb R)$
istnieje dokładnie jedna warunkowa wartość oczekiwana
$Z=\mathbb E[Y\mid\mathcal G]$ w sensie równości
$\mathbb P$-prawie wszędzie. Ponadto:
\begin{enumerate}[label=\textup{(\roman*)}]
\item odwzorowanie $Y\mapsto\mathbb E[Y\mid\mathcal G]$
jest liniowe i $\|Z\|_{L^1}\le\|Y\|_{L^1}$;
\item $Y\ge0$ prawie wszędzie implikuje $Z\ge0$
prawie wszędzie;
\item jeśli $Y$ jest $\mathcal G$-mierzalna,
to $Z=Y$ prawie wszędzie;
\item jeśli $\mathcal H\subseteq\mathcal G\subseteq\mathcal F$,
to $\mathbb E[\mathbb E[Y\mid\mathcal G]\mid\mathcal H]
=\mathbb E[Y\mid\mathcal H]$ prawie wszędzie.
\end{enumerate}
\end{twierdzenie}
\begin{proof}
Na przestrzeni mierzalnej $(\Omega,\mathcal G)$
określamy skończoną miarę ze znakiem
$\eta(A)=\int_A Y\,\mathrm d\mathbb P$.
Jest ona bezwzględnie ciągła względem
$\mu=\mathbb P|_{\mathcal G}$, a $\mu$ jest miarą
skończoną. Twierdzenie~\ref{thm:RN-complex-18}
(dla rzeczywistych miar) daje $\mathcal G$-mierzalną
$Z\in L^1(\mu)$, która spełnia
\eqref{eq:CE-def-18}. Jednoznaczność wynika z
jednoznaczności pochodnej Radona--Nikodyma.

Liniowość wynika z liniowości całki po $A$ i
jednoznaczności. Jeśli $Y\ge0$, to $\eta$ jest
miarą nieujemną, a więc jej pochodna $Z$ jest
nieujemna. Jeśli $Y$ jest $\mathcal G$-mierzalna,
sama spełnia \eqref{eq:CE-def-18}, co daje (iii).
Aby dowieść oszacowania $L^1$, korzystamy z
$|\eta|(\Omega)=\int_\Omega|Z|\,\mathrm d\mathbb P$
oraz z faktu, że dla każdego skończonego podziału
$\Omega=\bigsqcup_jA_j$, $A_j\in\mathcal G$,
\[
\sum_j|\eta(A_j)|
\le\sum_j\int_{A_j}|Y|\,\mathrm d\mathbb P
=\int_\Omega|Y|\,\mathrm d\mathbb P.
\]
Biorąc supremum po podziałach otrzymujemy (i).
Dla (iv) niech $A\in\mathcal H$; wówczas także
$A\in\mathcal G$, więc
\[
\int_A\mathbb E[Y\mid\mathcal G] \,\mathrm d\mathbb P
=\int_A Y\,\mathrm d\mathbb P.
\]
Zatem obie strony tożsamości (iv) spełniają
warunek definiujący wartość oczekiwaną względem
$\mathcal H$; jednoznaczność kończy dowód.
\end{proof}

\begin{twierdzenie}[Przypadek $L^2$: rzut ortogonalny]
\label{thm:CE-orthogonal-18}
W założeniach definicji~\ref{def:conditional-exp-18},
jeśli $Y\in L^2(\mathbb P;\mathbb R)$, to
$Z=\mathbb E[Y\mid\mathcal G]\in L^2(\mathbb P)$,
$\|Z\|_2\le\|Y\|_2$ i
\[
\int_\Omega(Y-Z)W\,\mathrm d\mathbb P=0
\qquad\text{dla każdego }W\in L^2(\Omega,\mathcal G,\mathbb P).
\]
Oznacza to, że $Z$ jest rzutem ortogonalnym $Y$
na domkniętą podprzestrzeń funkcji $\mathcal G$-mierzalnych.
\end{twierdzenie}
\begin{proof}
Nie wiemy jeszcze, czy $Z\in L^2$.
Niech $Z_n=\max(-n,\min(Z,n))$; jest to ograniczona
funkcja $\mathcal G$-mierzalna, więc z definicji
warunkowej wartości oczekiwanej (najpierw dla funkcji
prostych, potem przez zbieżność zmajoryzowaną) mamy
$\int_\Omega ZZ_n\,\mathrm d\mathbb P
=\int_\Omega YZ_n\,\mathrm d\mathbb P$.
Ponieważ $ZZ_n\ge Z_n^2$ punktowo,
\[
\|Z_n\|_2^2\le\int_\Omega ZZ_n\,\mathrm d\mathbb P
\le\|Y\|_2\|Z_n\|_2.
\]
Zatem $\|Z_n\|_2\le\|Y\|_2$.
Ponieważ $Z_n^2\uparrow Z^2$, Beppo Levi daje
$\|Z\|_2\le\|Y\|_2$. Definicyjna równość
$\int_A(Y-Z)\,\mathrm d\mathbb P=0$ dla $A\in\mathcal G$
rozciąga się liniowo na funkcje proste i przez
zbieżność zmajoryzowaną na ograniczone funkcje
$\mathcal G$-mierzalne. Trunkacje dowolnego
$W\in L^2(\mathcal G)$ zbliżają się do niego w $L^2$;
Cauchy--Schwarz pozwala przejść do granicy.
Podprzestrzeń $L^2(\mathcal G)$ jest domknięta:
jeśli $W_n$ są $\mathcal G$-mierzalne i zbieżne w $L^2$
do $W$, to wybieramy podciąg $(W_{n_k})$ taki, że
$\|W_{n_k}-W\|_2<2^{-2k}$. Z nierówności
Czebyszewa i lematu Borela--Cantellego wynika
$W_{n_k}\to W$ prawie wszędzie względem $\mathbb P$.
Zbiór $D$ punktów, w których ciąg $(W_{n_k})$
ma granicę skończoną, jest $\mathcal G$-mierzalny,
gdyż warunek Cauchy'ego i skończoność wartości
opisuje się przeliczalnymi działaniami mierzalnymi.
Funkcja $\widetilde W=\lim_kW_{n_k}$ na $D$
i $\widetilde W=0$ poza $D$ jest
$\mathcal G$-mierzalna i równa $W$ prawie wszędzie.
To konstruuje przedstawiciela w $L^2(\mathcal G)$,
bez zakładania zupełności $\mathcal G$.
Równanie ortogonalności identyfikuje $Z$ z rzutem.
\end{proof}

\section{Przykłady obliczania pochodnych}
\begin{przyklad}[Nieskończona gęstość przy punkcie brzegowym]
Na $X=(0,1)$ z miarą Lebesgue'a określmy
$\nu(E)=\int_E x^{-1/2}\,\mathrm d\lambda_1(x)$.
Wtedy $\nu\ll\lambda_1$, a
$\mathrm d\nu/\mathrm d\lambda_1=x^{-1/2}$.
Choć gęstość rośnie bez ograniczenia przy $0$,
jest całkowalna i $\nu(X)=2$.
\end{przyklad}

\begin{przyklad}[Miara osobliwa nie ma pochodnej reprezentującej masę]
Niech $X=[0,1]$ i $\mu=\lambda_1|_X$ na
$\mathcal B(X)$. Dla miary Cantora $\kappa$ z przykładu
\ref{ex:cantor-measure-17} zachodzi $\kappa\perp\mu$.
Nie istnieje $f\in L^1(\mu)$ spełniająca
$\kappa(E)=\int_Ef(x)\,\mathrm d\mu(x)$ dla wszystkich
$E\in\mathcal B(X)$, bo przy $E=C$ prawa strona
jest równa zero, a $\kappa(C)=1$.
Fakt, że pochodna klasycznej funkcji Cantora wynosi
zero prawie wszędzie, nie oznacza, że jej miara
Stieltjesa jest zerowa. Jest ona osobliwa względem
miary Lebesgue'a.
\end{przyklad}

\section{Rozkład miary Lebesgue'a--Stieltjesa}

\begin{twierdzenie}[Kryterium oparte na rozkładzie miary]\label{it:thm:stieltjes-decomposition}
Niech $\mu_\alpha$ będzie skończoną miarą Lebesgue'a--Stieltjesa ograniczoną do $(a,b]$. Istnieje rozkład
\[
\mu_\alpha=w\lambda_1+\mu_{\mathrm{sc}}+\sum_{k\in J}c_k\delta_{t_k},
\]
gdzie wszystkie miary rozpatrujemy na $(a,b]$, $w\ge0$ jest mierzalna, $\mu_{\mathrm{sc}}\perp\lambda_1$ jest bezatomowa, $J$ jest zbiorem co najwyżej przeliczalnym, $t_k$ są różne, $c_k>0$ i $\sum_{k\in J}c_k<\infty$. Dla każdej funkcji borelowskiej $f:(a,b]\to\mathbb{R}$ zachodzi
\[
\int_{(a,b]}|f|\,\mathrm d\mu_\alpha
=\int_{(a,b]}|f|w\,\mathrm d\lambda_1+
\int_{(a,b]}|f|\,\mathrm d\mu_{\mathrm{sc}}+
\sum_{k\in J}c_k|f(t_k)|.
\]
Funkcja jest bezwzględnie całkowalna względem $\mu_\alpha$ wtedy i tylko wtedy, gdy prawa strona jest skończona. W tym przypadku
\[
 \int_{(a,b]}f\,\mathrm d\mu_\alpha
 =\int_{(a,b]}fw\,\mathrm d\lambda_1+
 \int_{(a,b]}f\,\mathrm d\mu_{\mathrm{sc}}+
 \sum_{k\in J}c_k f(t_k),
\]
a ostatni szereg jest bezwzględnie zbieżny.
\end{twierdzenie}
\begin{proof}
Obie miary $\mu_\alpha$ i $\lambda_1|_{(a,b]}$ są skończone, więc twierdzenie o rozkładzie Lebesgue'a daje część bezwzględnie ciągłą i osobliwą. Twierdzenie \ref{thm:RN-sigfin-18} Radona--Nikodyma daje nieujemną gęstość $w$ pierwszej części względem $\lambda_1$. Drugą część przedłużamy przez zero na $\mathbb{R}$ i stosujemy twierdzenie \atlasref{ME-17::thm:atomic-cont-singular-17}{ME-17: rozkład części osobliwej}. Masa atomów sumuje się do masy części atomowej, nie większej od $\mu_\alpha((a,b])$.
Ponieważ część bezwzględnie ciągła nie ma mas punktowych,
każdy współczynnik $c_k$ jest równy
$\mu_\alpha(\{t_k\})$. Dla generatora prawostronnie
ciągłego oznacza to $c_k=\alpha(t_k)-\alpha(t_k-)$.
Jednoznaczność rozkładu Lebesgue'a i gęstości
Radona--Nikodyma daje jednoznaczność $w$ prawie wszędzie
względem $\lambda_1$ oraz części osobliwej.
Masy singletonów wyznaczają jej część atomową,
a odjęcie tej części wyznacza $\mu_{\mathrm{sc}}$.
Dla nieujemnej borelowskiej funkcji prostej równość całek względem sumy miar wynika z definicji całki i przeliczalnej addytywności sum nieujemnych. Dla dowolnej nieujemnej funkcji mierzalnej przybliżamy ją od dołu rosnącym ciągiem funkcji prostych i stosujemy twierdzenie \atlasref{ME-08::thm:mct}{ME-08: twierdzenie Beppo Leviego} o zbieżności monotonicznej do każdej miary oraz do sumy po $J$ (miara licząca). Lemat \ref{lem:integration-density-18} identyfikuje całkę względem $w\lambda_1$, a całka względem $\delta_{t_k}$ jest wartością w $t_k$. Stosujemy otrzymaną równość do $|f|$. Jeśli prawa strona jest skończona, stosujemy ją również osobno do $f^+$ i $f^-$ i odejmujemy skończone wartości. Z oszacowania dla $|f|$ wynika bezwzględna zbieżność szeregu z wartościami $f(t_k)$.
\end{proof}

\begin{twierdzenie}[Rozkład funkcji generującej miarę Stieltjesa]
\label{it:thm:stieltjes-generator-decomposition}
Niech $a<b$, a $g\in\operatorname{BV}[a,b]$ będzie prawostronnie ciągła
na $[a,b)$. Istnieją jednoznacznie określone funkcje
$g_{\mathrm{ac}},g_{\mathrm{sc}},g_{\mathrm{at}}\in\operatorname{BV}[a,b]$,
zerujące się w $a$, takie że
\[
 g(t)=g(a)+g_{\mathrm{ac}}(t)+g_{\mathrm{sc}}(t)+g_{\mathrm{at}}(t),
\]
gdzie
\[
 g_{\mathrm{ac}}(t)=\int_a^t w\,\mathrm d\lambda_1,\qquad
 g_{\mathrm{at}}(t)=\sum_{a<c\le t}d_c,
 \qquad w\inL^1[a,b],\quad \sum_c|d_c|<\infty.
\]
Funkcja $g_{\mathrm{sc}}$ jest ciągła, o ograniczonym
wahaniu, a jej pochodna znika prawie wszędzie.
Ponadto $w=g'$ prawie wszędzie i $d_c=g(c)-g(c-)$.
Sumy po $c$ obejmują punkty z $(a,b]$ o niezerowym
współczynniku; jest ich co najwyżej przeliczalnie wiele.
W zapisie miary $g'\lambda_1$ oraz we wzorach całkowych
wybieramy skończonego borelowskiego reprezentanta
pochodnej, określonej prawie wszędzie.
Zachodzi dokładna równość
\[
 V_a^b(g)=\int_a^b|g'|\,\mathrm d\lambda_1+
 V_a^b(g_{\mathrm{sc}})+\sum_{a<c\le b}|d_c|.
\]
Jeżeli $g$ jest niemalejąca, wszystkie trzy składniki
są niemalejące.
\end{twierdzenie}
\begin{proof}
Niech $\nu=\nu_g$ będzie miarą z twierdzenia
\atlasref{ME-16::it:thm:stieltjes-signed-bv}{ME-16}. Do dodatnich części
$\nu^+,\nu^-$ jej rozkładu Jordana stosujemy rozkład
z twierdzenia \ref{it:thm:stieltjes-decomposition}:
każda z tych miar jest miarą Stieltjesa dzięki twierdzeniu
\atlasref{ME-04::it:thm:stieltjes-construction}{ME-04} (przedłużamy ją przez
zero poza $(a,b]$). Odejmując odpowiadające sobie części,
otrzymujemy
\[
 \nu=w\lambda_1+\sigma+\sum_{c\in J}d_c\delta_c,
\]
gdzie $\sigma$ jest skończoną miarą ze znakiem osobliwą
względem $\lambda_1$ i jej wahanie jest bezatomowe.
Istotnie, wahanie różnicy dwóch dodatnich miar jest
ograniczone przez ich sumę, która w tym przypadku
jest osobliwa i bezatomowa. Na prostej bezatomowość
wynika z zerowych mas wszystkich singletonów
(twierdzenie \atlasref{ME-17::thm:atomic-cont-singular-17}{ME-17}).
Zbiór $J$ jest co najwyżej przeliczalny, a
$d_c=\nu(\{c\})=g(c)-g(c-)$, ponieważ pozostałe
części nie mają atomów. Stąd
$\sum_{c\in J}|d_c|\le|\nu|((a,b])$.

Wybieramy borelowski zbiór zerowy $S$, na którym skupia się
$|\sigma|$, i dołączamy $J$ do $S$. Trzy miary ze znakiem
są skupione odpowiednio na rozłącznych zbiorach
$(a,b]\setminus S$, $S\setminus J$, $J$.
Ich wahania zatem sumują się do $|\nu|$:
na każdym z tych trzech zbiorów ograniczenie $\nu$
zgadza się z odpowiednim składnikiem, a wahanie
ograniczenia jest ograniczeniem wahania.
Jednoznaczność części atomowej wynika z wartości na
singletonach. Po jej odjęciu, gdyby istniały dwa rozkłady
$\nu-\sum_c d_c\delta_c=w_i\lambda_1+\sigma_i$,
$i=1,2$, to $(w_1-w_2)\lambda_1=\sigma_2-\sigma_1$.
Miara po lewej jest bezwzględnie ciągła, a po prawej
skupiona na sumie dwóch borelowskich zbiorów zerowych.
Miara ze znakiem o obu tych własnościach jest zerowa:
każdy zbiór borelowski rozcinamy na część w tym zbiorze
zerowym i część poza nim. Stąd $\sigma_1=\sigma_2$ oraz $w_1=w_2$ prawie
wszędzie. Ostatnia implikacja wynika z jednoznaczności
gęstości Radona--Nikodyma lub bezpośrednio przez
całkowanie różnicy na zbiorach, gdzie jest ona dodatnia
i gdzie jest ujemna. Gdy gęstości są jedynie mierzalne
w sensie Lebesgue'a, wybieramy borelowskie wersje
zgodne z nimi prawie wszędzie.

Definiujemy składniki funkcji przez masy $(a,t]$
odpowiednich miar; w $a$ przyjmujemy zero.
Twierdzenie \atlasref{ME-16::it:thm:stieltjes-signed-bv}{ME-16} zapewnia
prawostronną ciągłość, ograniczone wahanie i jego dokładną
wartość dla każdego składnika. Dla $\sigma$ dodatkowo
$g_{\mathrm{sc}}(t)-g_{\mathrm{sc}}(t-)=\sigma(\{t\})=0$,
więc $g_{\mathrm{sc}}$ jest ciągła także lewostronnie.
Twierdzenie \atlasref{ME-25::thm:cdf-derivative-25}{ME-25: pochodna funkcji rozkładu}, zastosowane
do dodatnich części miar przedłużonych przez zero na
$\mathbb R$, daje $g'=w$ prawie wszędzie oraz zerową pochodną
części osobliwej i atomowej prawie wszędzie.
W szczególności $g_{\mathrm{sc}}'=0$ prawie wszędzie.
W powyższej konstrukcji $w$ jest różnicą borelowskich
gęstości Radona--Nikodyma dla miar na borelowskich
podzbiorach $(a,b]$, zatem można użyć $w$ jako
reprezentanta $g'$ (zmieniając wartości nieskończone
na zbiorze zerowym oraz wybierając wartość w $a$).
Zapisy z $g'$ w dalszych wzorach nie zależą od
wyboru takiego reprezentanta, ponieważ dotyczą
wyłącznie całkowania względem $\lambda_1$.

W celu sprawdzenia jednoznaczności zwróćmy uwagę,
że każda funkcja postaci
$g_{\mathrm{at}}(t)=\sum_{a<c\le t}d_c$ jest prawostronnie
ciągła. Jest bowiem jednostajną granicą skończonych sum
funkcji $d_c\chi_{[c,b]}$: norma supremum ogona jest
ograniczona przez ogon zbieżnego szeregu $\sum_c|d_c|$.
Każda skończona suma jest prawostronnie ciągła.
Pozostałe dwa składniki są ciągłe z założeń.
Wszystkie spełniają zatem hipotezy twierdzenia
\atlasref{ME-16::it:thm:stieltjes-signed-bv}{ME-16}.
Jednoznaczność funkcji wynika z jednoznaczności rozkładu
miary: bezwzględnie ciągły składnik generuje miarę
$w\lambda_1$, ciągły składnik o zerowej pochodnej
generuje miarę osobliwą bezatomową: ciągłość wyklucza
atomy na mocy twierdzenia \atlasref{ME-16::it:thm:stieltjes-signed-bv}{ME-16},
a twierdzenie \atlasref{ME-25::thm:cdf-derivative-25}{ME-25: pochodna funkcji rozkładu} zastosowane
do obu części Jordana daje gęstość części bezwzględnie
ciągłej równą różnicy ich gęstości, czyli pochodnej
funkcji generującej, tutaj zerowej prawie wszędzie.
Składnik skokowy generuje miarę atomową, gdyż jego
przyrost na $(s,t]$ jest sumą współczynników
$d_c$ dla $s<c\le t$ i szereg jest bezwzględnie zbieżny.
Jednoznaczność miary na takich przedziałach kończy
identyfikację.
Normalizacja w $a$ usuwa dowolność stałych.
Doprecyzujmy przejście od rozkładu miar do równości
wahań. Dla borelowskiego $B\subset(a,b]$ zachodzi
\[
 |w\lambda_1|(B)=\int_B|w|\,\mathrm d\lambda_1.
\]
Istotnie, dla dowolnego skończonego mierzalnego podziału
$B$ suma modułów całek $w$ nie przekracza całki $|w|$.
Podział $B$ na części, na których $w\ge0$ i $w<0$,
daje równość. Podobnie dla miary atomowej
$\tau=\sum_c d_c\delta_c$ mamy
$|\tau|(B)=\sum_{c\in B}|d_c|$: nierówność górna wynika
z nierówności trójkąta dla bezwzględnie zbieżnych
szeregów, a dolna z podziału na zbiory punktów
o dodatnich i ujemnych współczynnikach oraz ich
dopełnienie. Z rozłączności zbiorów skupienia
trzech składników wynika zatem
\[
 |\nu|(B)=\int_B|w|\,\mathrm d\lambda_1
           +|\sigma|(B)+\sum_{c\in B}|d_c|.
\]
Dla $B=(a,b]$ stosujemy twierdzenie
\atlasref{ME-16::it:thm:stieltjes-signed-bv}{ME-16} do $g$ i $g_{\mathrm{sc}}$,
a następnie $w=g'$ prawie wszędzie. Otrzymujemy
dokładnie żądaną równość wahań funkcji. Dla niemalejącego $g$ miara $\nu$ jest
nieujemna i jej trzy składniki również są nieujemne.
\end{proof}

\begin{wniosek}[Kryterium całkowalności względem miary Stieltjesa ze znakiem]
\label{it:cor:stieltjes-signed-integrability}
W hipotezach twierdzenia \ref{it:thm:stieltjes-generator-decomposition}
niech $\nu_g=g'\lambda_1+\sigma+\sum_c d_c\delta_c$
będzie otrzymanym rozkładem miary.
Dla każdej rzeczywistej lub zespolonej funkcji borelowskiej
$f$ na $(a,b]$ zachodzi
\[
 \int_{(a,b]}|f|\,\mathrm d|\nu_g|
 =\int_{(a,b]}|f|\,|g'|\,\mathrm d\lambda_1
  +\int_{(a,b]}|f|\,\mathrm d|\sigma|
  +\sum_c|d_c|\,|f(c)|.
\]
Funkcja $f$ jest całkowalna względem $\nu_g$ wtedy i tylko
wtedy, gdy prawa strona jest skończona. Wówczas
\[
 \int_{(a,b]}f\,\mathrm d\nu_g
 =\int_{(a,b]}fg'\,\mathrm d\lambda_1
  +\int_{(a,b]}f\,\mathrm d\sigma+\sum_c d_cf(c),
\]
przy czym szereg jest bezwzględnie zbieżny.
\end{wniosek}
\begin{proof}
Równość miar dodatnich
$|\nu_g|=|g'|\lambda_1+|\sigma|+\sum_c|d_c|\delta_c$
wykazano w dowodzie poprzedniego twierdzenia.
Argument dla całek nieujemnych z twierdzenia
\ref{it:thm:stieltjes-decomposition} daje pierwszą równość,
w szczególności kryterium całkowalności. Jeżeli prawa
strona jest skończona, to wszystkie trzy całki względem
składników miary ze znakiem istnieją, a szereg wartości
jest bezwzględnie zbieżny. Dla rzeczywistych funkcji
prostych równość wynika z addytywności rozkładu miary.
Dla funkcji nieujemnych całkowalnych przechodzimy
do granicy rosnących funkcji prostych, stosując
zbieżność monotoniczną osobno do dodatnich i ujemnych
części każdej miary. Wszystkie otrzymane granice są
skończone dzięki pierwszej równości. Rozkład funkcji
rzeczywistej na część dodatnią i ujemną oraz rozkład
funkcji zespolonej na część rzeczywistą i urojoną
kończą dowód.
\end{proof}

\section{Mapa implikacji}
Przy $\sigma$-skończoności $\mu,\nu$ uzyskaliśmy
następujący ciąg rezultatów: rozkład Lebesgue'a
$\nu=\nu_{\mathrm{ac}}+\nu_{\mathrm{s}}$,
reprezentację $\nu_{\mathrm{ac}}=f\mu$,
jednoznaczność gęstości, reguły rachunkowe i
konstrukcję warunkowej wartości oczekiwanej.
Przykład~\ref{ex:RN-failure-nonsigma-18} ostrzega,
że sama implikacja między zbiorami zerowymi
nie zastępuje hipotez twierdzenia.

% ROZDZIAŁ 19 — zastosowania pochodnej Radona--Nikodyma
\chapter{Zastosowania twierdzenia Radona--Nikodyma}
\label{ch:rn-applications-19}

\section{Gęstość jako opis zmiany miary}
Twierdzenie Radona--Nikodyma z rozdziału~\ref{ch:radon-nikodym-18}
zamienia porównywanie miar w rachunek funkcyjny. Nie wystarczy jednak
bezrefleksyjnie dzielić gęstości: przy każdej zmianie miary należy
wskazać, względem której miary zachodzi równość prawie wszędzie,
a także sprawdzić, czy wyrażenia są całkowalne.

\begin{twierdzenie}[Całkowanie po zmianie miary]
\label{thm:change-measure-19}
Niech $(X,\mathcal F)$ będzie przestrzenią mierzalną,
$\mu,\nu$ nieujemnymi miarami $\sigma$-skończonymi na niej,
$\nu\ll\mu$, a $f=\mathrm d\nu/\mathrm d\mu$.
Dla każdej mierzalnej funkcji $h:X\to[0,\infty]$ zachodzi
\begin{equation}\label{eq:change-measure-19}
 \int_X h(x)\,\mathrm d\nu(x)
 =\int_X h(x)f(x)\,\mathrm d\mu(x).
\end{equation}
Jeśli $h:X\to\mathbb C$ jest mierzalna, to $h\in L^1(\nu)$
wtedy i tylko wtedy, gdy $hf\in L^1(\mu)$; wówczas
\eqref{eq:change-measure-19} pozostaje prawdziwy.
\end{twierdzenie}
\begin{proof}
Dla $h=\chi_A$ równość jest definicją pochodnej.
Przez liniowość otrzymujemy ją dla nieujemnych funkcji prostych,
a następnie przez aproksymację monotoniczną i twierdzenie Beppo
Leviego dla każdej $h\ge0$. Zastosowanie wyniku do $|h|$
daje równoważność całkowalności. Rozłożenie funkcji zespolonej
na dodatnie i ujemne części jej części rzeczywistej i urojonej
kończy dowód.
\end{proof}

\begin{przyklad}[Ważenie miary]
Niech $\mu=\lambda_1|_{(0,1)}$ i $f(x)=2x$.
Wówczas $\nu(E)=\int_E 2x\,\mathrm d\lambda_1(x)$ jest
miarą probabilistyczną. Dla $h(x)=x^2$ otrzymujemy
$\int_{(0,1)}h\,\mathrm d\nu(x)
=\int_{(0,1)}2x^3\,\mathrm d\lambda_1(x)=1/2$.
Z kolei $\mathrm d\mu/\mathrm d\nu=(2x)^{-1}$
$\nu$-prawie wszędzie. Mimo nieograniczoności tej pochodnej
przy zerze zachodzi
$\int_{(0,1)}(2x)^{-1}\,\mathrm d\nu(x)=1$.
\end{przyklad}

\section{Miary obrazowe oraz gęstości obserwacji}
Dla mierzalnego $T:X\to Y$ i miary $\mu$ na $(X,\mathcal F)$
przypominamy miarę obrazową
$T_*\mu(B)=\mu(T^{-1}(B))$, $B\in\mathcal G$,
gdzie $(Y,\mathcal G)$ jest przestrzenią mierzalną.
Nie wymagamy, aby $T$ była odwracalna.

\begin{twierdzenie}[Bezwzględna ciągłość zachowana przez obraz]
\label{thm:image-ac-19}
Niech $T:(X,\mathcal F)\to(Y,\mathcal G)$ będzie mierzalne,
a $\mu,\nu$ miarami nieujemnymi na $\mathcal F$.
Jeżeli $\nu\ll\mu$, to $T_*\nu\ll T_*\mu$.
Jeżeli ponadto $\mu,\nu$ są miarami probabilistycznymi,
$f=\mathrm d\nu/\mathrm d\mu$ i
$g=\mathrm d(T_*\nu)/\mathrm d(T_*\mu)$,
to $g\circ T$ jest wersją warunkowej wartości oczekiwanej
$\mathbb E_\mu[f\mid\sigma(T)]$, gdzie
$\sigma(T)=\{T^{-1}(B):B\in\mathcal G\}$.
\end{twierdzenie}
\begin{proof}
Jeżeli $(T_*\mu)(B)=0$, to $\mu(T^{-1}(B))=0$,
zatem $\nu(T^{-1}(B))=0$, co jest równoważne
$(T_*\nu)(B)=0$. Obie miary obrazowe są skończone,
więc pochodna $g$ istnieje. Dla każdego $B\in\mathcal G$
z twierdzenia o całkowaniu względem miary obrazowej
(definicję miary obrazowej z rozdziału~\ref{ch:struktura-przestrzeni})
uzyskujemy
\[
\int_{T^{-1}(B)}g(T(x))\,\mathrm d\mu(x)
=\int_Bg(y)\,\mathrm d(T_*\mu)(y)
=(T_*\nu)(B)
=\int_{T^{-1}(B)}f(x)\,\mathrm d\mu(x).
\]
Funkcja $g\circ T$ jest $\sigma(T)$-mierzalna. Jest również
całkowalna, ponieważ
\[
\int_X|g(T(x))|\,\mathrm d\mu(x)
=\int_Y|g(y)|\,\mathrm d(T_*\mu)(y)<\infty.
\]
Spełnia więc warunek charakteryzujący warunkową wartość oczekiwaną
z definicji~\ref{def:conditional-exp-18}.
\end{proof}

\begin{uwaga}[Kiedy informacja znika?]
W ogólności $g(T(x))\neq f(x)$. Po zastosowaniu
odwzorowania $T$ pozostaje tylko średnia warunkowa $f$
względem informacji dostępnej poprzez $T$. Jeśli $T$
jest stałe, a $\mu,\nu$ są miarami probabilistycznymi,
to $T_*\mu=T_*\nu$, choć $\mu$ i $\nu$ mogą być różne.
Nie wolno też z samego $\nu\ll\mu$ wnioskować, że dowolne
odwzorowanie $T$ zachowuje osobliwość miar.
\end{uwaga}

\section{Wariacja całkowita a testowanie różnic między rozkładami}
Dla miar probabilistycznych $P,Q$ na $(X,\mathcal F)$
wprowadzamy odległość
\begin{equation}\label{eq:tv-prob-19}
d_{\mathrm{TV}}(P,Q):=\sup_{A\in\mathcal F}|P(A)-Q(A)|.
\end{equation}
Należy odróżnić ją od normy wariacji całkowitej miary ze znakiem
$P-Q$, oznaczonej w rozdziale~\ref{ch:miary-ze-znakiem} przez
$\|P-Q\|_{\mathrm{TV}}=|P-Q|(X)$.

\begin{twierdzenie}[Formuła wariacyjna dla dwóch prawdopodobieństw]
\label{thm:total-variation-formula-19}
Niech $P,Q$ będą miarami probabilistycznymi na $(X,\mathcal F)$.
Niech $\lambda=P+Q$, $p=\mathrm dP/\mathrm d\lambda$
i $q=\mathrm dQ/\mathrm d\lambda$. Wówczas
\begin{equation}\label{eq:tv-formula-19}
 d_{\mathrm{TV}}(P,Q)
 =\frac12\int_X|p(x)-q(x)|\,\mathrm d\lambda(x)
 =\frac12\|P-Q\|_{\mathrm{TV}}
 =1-\int_X\min(p,q)\,\mathrm d\lambda.
\end{equation}
Supremum z \eqref{eq:tv-prob-19} jest osiągane na zbiorze
$A_+=\{p>q\}$. W szczególności $0\le d_{\mathrm{TV}}(P,Q)\le1$,
a równość $d_{\mathrm{TV}}(P,Q)=1$ zachodzi dokładnie wtedy,
gdy $P\perp Q$.
\end{twierdzenie}
\begin{proof}
Miara ze znakiem $\eta=P-Q$ ma gęstość $p-q$
względem $\lambda$. Zatem
$\eta(A)=\int_A(p-q)\,\mathrm d\lambda$ i
$\eta(X)=0$. Rozkład na części dodatnią i ujemną
pokazuje, że $\eta^+(X)=\eta^-(X)=|\eta|(X)/2$.
Dla dowolnego $A$ mamy
$-\eta^-(X)\le\eta(A)\le\eta^+(X)$,
a $\eta(A_+)=\eta^+(X)$; stąd pierwsze dwie równości.
Ponieważ $\int_X(p+q)\,\mathrm d\lambda=2$ i punktowo
$|p-q|=p+q-2\min(p,q)$, dostajemy ostatnią równość.
Jeżeli $d_{\mathrm{TV}}=1$, to $\min(p,q)=0$
$\lambda$-prawie wszędzie, czyli $p q=0$ prawie wszędzie.
Zbiory $\{p>0\}$ i $\{q>0\}$ są wówczas rozłączne
z dokładnością do zbioru $\lambda$-zerowego i skupiają
odpowiednio $P$ i $Q$; stąd $P\perp Q$.
Implikacja odwrotna wynika bezpośrednio z definicji osobliwości.
\end{proof}

\begin{wniosek}[Przetwarzanie danych nie zwiększa wariacji całkowitej]
\label{cor:tv-data-processing-19}
W założeniach twierdzenia~\ref{thm:image-ac-19},
bez wymogu $\nu\ll\mu$, jeśli $P,Q$ są prawdopodobieństwami,
to
$ d_{\mathrm{TV}}(T_*P,T_*Q)\le d_{\mathrm{TV}}(P,Q)$.
\end{wniosek}
\begin{proof}
Dla dowolnego $B\in\mathcal G$ zachodzi
$|(T_*P)(B)-(T_*Q)(B)|
=|P(T^{-1}B)-Q(T^{-1}B)|\le d_{\mathrm{TV}}(P,Q)$.
Bierzemy supremum po $B$.
\end{proof}

\begin{przyklad}[Zbieżność słaba nie wymusza zbieżności w wariacji]
Na $\mathbb R$ rozważmy $P_n=\delta_{1/n}$ i $P=\delta_0$.
Dla każdej funkcji ciągłej i ograniczonej $h$
zachodzi $\int_{\mathbb R}h\,\mathrm dP_n\to
\int_{\mathbb R}h\,\mathrm dP$, co jest zbieżnością wąską.
Jednak $d_{\mathrm{TV}}(P_n,P)=1$ dla każdego $n$,
ponieważ $P_n(\{1/n\})=1$ i $P(\{1/n\})=0$.
Precyzyjną teorię zbieżności miar rozwiniemy w rozdziale~22.
\end{przyklad}

\section{Lemat Neymana--Pearsona: optymalność testu progowego}
Przedstawimy rezultat stanowiący użyteczne, a zarazem nieintuicyjne
zastosowanie pochodnej Radona--Nikodyma. Z dwóch prostych hipotez
można skonstruować test optymalny bez żadnych założeń o istnieniu
gęstości względem miary Lebesgue'a.

\begin{definicja}[Test losowy i jego moc]
\label{def:test-19}
Niech $P,Q$ będą prawdopodobieństwami na $(X,\mathcal F)$.
Testem losowym jest mierzalna funkcja $\varphi:X\to[0,1]$.
Liczbę $\int_X\varphi\,\mathrm dP$ nazywamy
prawdopodobieństwem odrzucenia hipotezy $P$ pod $P$,
a $\int_X\varphi\,\mathrm dQ$ jego mocą pod $Q$.
Test ma poziom nieprzekraczający $\alpha\in[0,1]$,
gdy pierwsza z tych liczb nie przekracza $\alpha$.
\end{definicja}

\begin{twierdzenie}[Lemat Neymana--Pearsona]
\label{thm:neyman-pearson-19}
Niech $P,Q$ będą prawdopodobieństwami na $(X,\mathcal F)$,
$\lambda=P+Q$, $p=\mathrm dP/\mathrm d\lambda$
i $q=\mathrm dQ/\mathrm d\lambda$.
Dla $c\ge0$ i $\gamma\in[0,1]$ ustalmy test
\[
\varphi_*=\chi_{\{q>cp\}}+
 \gamma\chi_{\{q=cp\}}.
\]
Niech $\alpha=\int_X\varphi_*\,\mathrm dP$.
Wtedy dla każdego testu $\varphi:X\to[0,1]$ o poziomie
nieprzekraczającym $\alpha$ zachodzi
\[
 \int_X\varphi\,\mathrm dQ\le
 \int_X\varphi_*\,\mathrm dQ.
\]
Dla $c>0$ dowód nie wymaga założenia, że
$P$ i $Q$ są wzajemnie bezwzględnie ciągłe.
W razie $c=0$ twierdzenie pozostaje prawdziwe;
nie wymaga istnienia gęstości ani $P$, ani $Q$ względem
żadnej wcześniej zadanej miary odniesienia.
\end{twierdzenie}
\begin{proof}
Na zbiorze $\{q>cp\}$ mamy $\varphi_*=1$, więc
$\varphi-\varphi_*\le0$; na $\{q<cp\}$ mamy
$\varphi_*=0$, więc $\varphi-\varphi_*\ge0$.
Na zbiorze $\{q=cp\}$ iloczyn jest zerem. W każdym
przypadku zatem
$(\varphi-\varphi_*)(q-cp)\le0$ punktowo.
Całkujemy względem $\lambda$ i dostajemy
\[
\int_X(\varphi-\varphi_*)\,\mathrm dQ
\le c\int_X(\varphi-\varphi_*)\,\mathrm dP\le0,
\]
gdzie ostatnia nierówność wynika z poziomu testu
i nieujemności $c$. To dokładnie teza.
\end{proof}

\begin{uwaga}[Rola losowania na granicy]
Wprowadzony parametr $\gamma$ pozwala kontrolować poziom
testu, gdy zbiór $\{q=cp\}$ ma dodatnią miarę $P$.
W ogólności nie każdy poziom $\alpha$ da się otrzymać
bez losowania. Na zbiorze $\{p=0,q>0\}$ test progowy
zawsze ma wartość $1$: odrzucenie $P$ nic tam nie kosztuje.
\end{uwaga}

\begin{wniosek}[Konstrukcja testu optymalnego o zadanym poziomie]
\label{cor:NP-exact-alpha-19}
Niech $P,Q,\lambda,p,q$ mają znaczenie z twierdzenia
\ref{thm:neyman-pearson-19}. Dla każdego
$\alpha\in(0,1)$ istnieją $c\ge0$ i
$\gamma\in[0,1]$, dla których test
$\varphi_*=\chi_{\{q>cp\}}+\gamma\chi_{\{q=cp\}}$
ma poziom dokładnie $\alpha$, a spośród wszystkich
testów poziomu co najwyżej $\alpha$ ma największą moc.
Dla $\alpha=0$ testem optymalnym jest
$\chi_{\{p=0,\,q>0\}}$, a dla $\alpha=1$
można przyjąć test stały równy $1$.
\end{wniosek}
\begin{proof}
Niech $L=q/p$ na $\{p>0\}$ i $L=0$ na $\{p=0\}$.
Ponieważ $P(\{p=0\})=0$, mamy $L<\infty$
$P$-prawie wszędzie. Dla $0<\alpha<1$ kładziemy
\[
 c=\inf\{t\ge0:P(L>t)\le\alpha\}.
\]
Z nierówności $P(L>t)\downarrow0$ przy $t\uparrow\infty$
wynika, że zbiór pod infimum jest niepusty, a $c$ jest
skończone. Dla $t_n\downarrow c$ i $t_n>c$ z ciągłości
miary z dołu otrzymujemy $P(L>c)\le\alpha$.
Jeśli $c>0$, to z definicji infimum
$P(L>t)>\alpha$ dla $0\le t<c$;
przechodząc z $t\uparrow c$ i używając ciągłości miary
z góry, dostajemy $P(L\ge c)\ge\alpha$.
Gdy $c=0$, ostatnia nierówność wynika z
$P(L\ge0)=1$.
W konsekwencji
\[
 P(L>c)\le\alpha\le P(L\ge c).
\]
Jeśli $P(L=c)>0$, przyjmujemy
$\gamma=(\alpha-P(L>c))/P(L=c)\in[0,1]$;
w przeciwnym przypadku nierówności dają
$P(L>c)=\alpha$ i wybieramy $\gamma=0$.
Zbiory $\{L>c\}$, $\{L=c\}$ zgadzają się
$P$-prawie wszędzie z $\{q>cp\}$ i $\{q=cp\}$,
więc $\int_X\varphi_*(x)\,\mathrm dP(x)=\alpha$.
Optymalność wynika z twierdzenia
\ref{thm:neyman-pearson-19}.

Dla $\alpha=0$ dopuszczalny test $\varphi$ jest zerowy
$P$-prawie wszędzie, a więc $\varphi=0$
$\lambda$-prawie wszędzie na $\{p>0\}$.
Stąd jego moc nie przekracza
$Q(\{p=0\})$, osiąganej przez wskazany test.
Przypadek $\alpha=1$ jest natychmiastowy.
\end{proof}

\begin{przyklad}[Dlaczego losowanie na granicy testu bywa konieczne?]
\label{ex:np-randomization-19}
Niech $X=\{a,b\}$, $\mathcal F=\mathcal P(X)$ oraz
$P(\{a\})=7/10$, $P(\{b\})=3/10$,
$Q(\{a\})=1/5$ i $Q(\{b\})=4/5$.
Stosunek wiarygodności $q/p$ dla $\lambda=P+Q$
przyjmuje wartości $2/7$ w punkcie $a$ i $8/3$
w punkcie $b$. Dla poziomu $\alpha=3/20$
przyjmujemy $c=8/3$ oraz $\gamma=1/2$.
Wówczas $\{q>cp\}=\varnothing$, a $\{q=cp\}=\{b\}$,
zatem test progowy ma postać
$\varphi_*=\frac12\chi_{\{b\}}$. Jego poziom jest
równy $\int_X\varphi_*\,\mathrm dP=3/20$, a moc
$\int_X\varphi_*\,\mathrm dQ=2/5$.

Każdy test \emph{nielosowy} przyjmujący tylko wartości
$0$ i $1$, którego poziom nie przekracza $3/20$,
musi być identycznie zerowy: odrzucenie hipotezy
w którymkolwiek punkcie kosztuje co najmniej
$P(\{b\})=3/10$. Ma więc moc równą $0$.
Losowanie na $\{q=cp\}$ rzeczywiście zwiększa
moc dopuszczalnego testu.
\end{przyklad}

\section{Wzór Bayesa jako twierdzenie o gęstościach}
Wynik sformułujemy bez założenia, że zmienna obserwowana
posiada klasyczną gęstość na $\mathbb R^d$.

\begin{twierdzenie}[Bayes dla skończonej przestrzeni parametrów]
\label{thm:bayes-finite-19}
Niech $m\ge1$, $\pi_i>0$, $\sum_{i=1}^m\pi_i=1$.
Niech $(Y,\mathcal G)$ będzie przestrzenią mierzalną,
a $P_1,\ldots,P_m$ miarami probabilistycznymi na $Y$.
Na przestrzeni $\{1,\ldots,m\}\times Y$ definiujemy
prawdopodobieństwo wspólne przez
$\mathbb P(\{i\}\times B)=\pi_iP_i(B)$.
Niech $\lambda=\sum_iP_i$, $p_i=\mathrm dP_i/\mathrm d\lambda$,
a $r=\sum_i\pi_i p_i$. Dla $r(y)>0$ określmy
\[
 \beta_i(y)=\frac{\pi_i p_i(y)}{r(y)}.
\]
Na zbiorze $\{r=0\}$ możemy przyjąć dowolne mierzalne
wartości $\beta_i\in[0,1]$. Wtedy dla obserwacji
$Y(i,y)=y$ i etykiety $\Theta(i,y)=i$ zachodzi
\begin{equation}\label{eq:bayes-finite-19}
 \beta_i(Y)=\mathbb P(\Theta=i\mid\sigma(Y))
 \quad\mathbb P\text{-prawie wszędzie}.
\end{equation}
Ponadto $\sum_i\beta_i=1$ $P_Y$-prawie wszędzie,
gdzie $P_Y=\sum_i\pi_iP_i$ jest prawem obserwacji.
\end{twierdzenie}
\begin{proof}
Dla dowolnego $B\in\mathcal G$ definicja rozkładu
wspólnego daje
$\mathbb P(\Theta=i,\,Y\in B)=\pi_iP_i(B)
=\int_B\pi_ip_i(y)\,\mathrm d\lambda(y)$.
Tymczasem $P_Y(B)=\int_Br(y)\,\mathrm d\lambda(y)$,
więc stosując twierdzenie~\ref{thm:change-measure-19}
uzyskujemy
\[
\int_{Y^{-1}(B)}\beta_i(Y)\,\mathrm d\mathbb P
=\int_B\beta_i(y)r(y)\,\mathrm d\lambda(y)
=\pi_iP_i(B).
\]
Zbiór $\{r=0\}$ ma miarę $P_Y$ równą zero, więc nie
wpływa na wzory. Funkcja $\beta_i(Y)$ jest
$\sigma(Y)$-mierzalna i ograniczona. Jest zatem
warunkową wartością oczekiwaną funkcji charakterystycznej
$\chi_{\{\Theta=i\}}$ względem $\sigma(Y)$.
Równość sumy wynika z $\sum_i\pi_ip_i=r$ na $\{r>0\}$.
\end{proof}

\begin{przyklad}[Dwie obserwacje Bernoulliego]
Niech $\Theta\in\{0,1\}$ z prawdopodobieństwami
$\mathbb P(\Theta=1)=\pi$ i $\mathbb P(\Theta=0)=1-\pi$.
Warunkowo względem $\Theta=i$ zaobserwujmy
niezależną parę $(Y_1,Y_2)$ zmiennych Bernoulliego
o parametrze $\theta_i\in(0,1)$. Dla obserwacji
$(Y_1,Y_2)=(1,0)$ wzór Bayesa daje
\[
\mathbb P(\Theta=1\mid Y_1=1,Y_2=0)
=\frac{\pi\theta_1(1-\theta_1)}
{\pi\theta_1(1-\theta_1)+(1-\pi)\theta_0(1-\theta_0)}.
\]
Nie potrzeba tu miary Lebesgue'a: wspólną miarą
dominującą jest miara licząca na czteroelementowym zbiorze.
\end{przyklad}

\section{Warunkowa wartość oczekiwana i nierówności}
Warunkową wartość oczekiwaną zdefiniowaliśmy w
rozdziale~\ref{ch:radon-nikodym-18}. Teraz dowiedziemy
jej własności przydatnych w analizie oraz prawdopodobieństwie.

\begin{lemat}[Wyciąganie czynnika mierzalnego]
\label{lem:conditional-pullout-19}
Niech $(\Omega,\mathcal F,\mathbb P)$ będzie przestrzenią
probabilistyczną, $\mathcal G\subseteq\mathcal F$,
$Y\in L^1(\mathbb P)$ oraz $H$ ograniczoną rzeczywistą
funkcją $\mathcal G$-mierzalną. Wówczas
\[
 \mathbb E[HY\mid\mathcal G]
 =H\mathbb E[Y\mid\mathcal G]
 \quad\mathbb P\text{-prawie wszędzie}.
\]
\end{lemat}
\begin{proof}
Dla $H=\chi_B$, $B\in\mathcal G$, tożsamość wynika
z warunku definiującego warunkową wartość oczekiwaną,
ponieważ dla dowolnego $A\in\mathcal G$ całkujemy po
$A\cap B$. Z liniowości wynika przypadek prostych $H$.
Dla ograniczonego $H$ istnieją proste $H_n$, mierzalne
względem $\mathcal G$, takie, że $H_n\to H$ punktowo
i $|H_n|\le\|H\|_\infty+1$. Przez zbieżność zmajoryzowaną
$H_nY\to HY$ i $H_n\mathbb E[Y\mid\mathcal G]
\to H\mathbb E[Y\mid\mathcal G]$ w $L^1$.
Całkując po każdym $A\in\mathcal G$, przechodzimy
do granicy; jednoznaczność daje tezę.
\end{proof}

\begin{twierdzenie}[Warunkowa nierówność Jensena]
\label{thm:conditional-jensen-19}
Niech $(\Omega,\mathcal F,\mathbb P)$ będzie przestrzenią
probabilistyczną, $\mathcal G\subseteq\mathcal F$,
$Y\in L^1(\mathbb P;\mathbb R)$, a
$\Phi:\mathbb R\to\mathbb R$ funkcją wypukłą, wszędzie
skończoną. Załóżmy ponadto, że $\Phi(Y)\in L^1(\mathbb P)$.
Wówczas
\[
 \Phi\bigl(\mathbb E[Y\mid\mathcal G]\bigr)
 \le\mathbb E[\Phi(Y)\mid\mathcal G]
 \quad\mathbb P\text{-prawie wszędzie}.
\]
\end{twierdzenie}
\begin{proof}
Dla każdego $t\in\mathbb Q$ wybieramy liczbę
$a_t\in\partial\Phi(t)$, gdzie
$\partial\Phi(t)$ oznacza zbiór nachyleń prostych
podpierających wykres $\Phi$ w punkcie $t$.
Istnienie takiego nachylenia wynika z monotoniczności
ilorazów różnicowych funkcji wypukłej; możemy np. przyjąć
prawostronną pochodną $\Phi'_+(t)$, która jest skończona.
Wówczas
$\Phi(x)\ge a_tx+\Phi(t)-a_tt$ dla każdego $x\in\mathbb R$.
Różnica
$\Phi(Y)-a_tY-\Phi(t)+a_tt$ jest nieujemna,
więc dodatniość i liniowość warunkowej wartości oczekiwanej
dają
\[
\mathbb E[\Phi(Y)\mid\mathcal G]
\ge a_t\mathbb E[Y\mid\mathcal G]+\Phi(t)-a_tt.
\]
Ponieważ zbiór liczb wymiernych jest przeliczalny,
można wybrać jeden wspólny zbiór pełnej miary, na którym
wszystkie nierówności obowiązują jednocześnie.
Ciągłość $\Phi$ i jej lokalna lipschitzowskość
wynikają z wypukłości i skończoności na całej prostej.
Pozwala to przejść z $t\in\mathbb Q$ do dowolnego $t$:
$\Phi(x)=\sup_{t\in\mathbb Q}
(a_tx+\Phi(t)-a_tt)$.
Podstawiając $x=\mathbb E[Y\mid\mathcal G]$
i biorąc supremum otrzymujemy tezę.
\end{proof}

\begin{wniosek}[Kontrakcja w $L^p$]
\label{cor:conditional-lp-19}
Jeśli $1\le p<\infty$, $Y\in L^p(\mathbb P;\mathbb R)$,
to $\mathbb E[Y\mid\mathcal G]\in L^p(\mathbb P)$ oraz
$\|\mathbb E[Y\mid\mathcal G]\|_p\le\|Y\|_p$.
Dla $p=\infty$ zachodzi analogiczne oszacowanie
normy istotnego supremum.
\end{wniosek}
\begin{proof}
Dla $1\le p<\infty$ stosujemy twierdzenie
\ref{thm:conditional-jensen-19} do $\Phi(t)=|t|^p$.
Całkując otrzymaną nierówność względem $\mathbb P$
i korzystając z własności zachowania całki warunkowej,
dostajemy oszacowanie normy. Jeżeli $|Y|\le M$ prawie
wszędzie, to z monotoniczności
$-M\le\mathbb E[Y\mid\mathcal G]\le M$,
co daje przypadek $p=\infty$.
\end{proof}

\section{Informacja ilościowa: entropia względna}
Dla prawdopodobieństw $Q\ll P$ określamy
\emph{entropię względną} (dywergencję Kullbacka--Leiblera)
\[
 D(Q\Vert P):=\int_X f(x)\log f(x)\,\mathrm dP(x),
 \qquad f=\frac{\mathrm dQ}{\mathrm dP},
\]
przy konwencji $0\log0=0$. Całka ma wartości
w $[0,\infty]$: z $t\log t\ge -e^{-1}$
wynika, że dla prawdopodobieństwa $P$ część ujemna
całkowanej funkcji jest całkowalna. Gdy $Q\not\ll P$,
przyjmujemy $D(Q\Vert P)=+\infty$.

\begin{twierdzenie}[Nieujemność entropii względnej]
\label{thm:gibbs-19}
Dla dowolnych prawdopodobieństw $P,Q$ na $(X,\mathcal F)$
zachodzi $D(Q\Vert P)\ge0$. Jeżeli $Q\ll P$, to
$D(Q\Vert P)=0$ wtedy i tylko wtedy, gdy $Q=P$.
\end{twierdzenie}
\begin{proof}
Dla $t\ge0$ funkcja $\psi(t)=t\log t-t+1$
(z $\psi(0)=1$) jest nieujemna i zeruje się wyłącznie
w $t=1$: na $(0,\infty)$ pochodna jest równa $\log t$,
więc minimum jest w $1$. Jeśli $Q\ll P$ i $f=\mathrm dQ/\mathrm dP$,
to $\int_Xf\,\mathrm dP=1$, a zatem
\[
 D(Q\Vert P)=\int_X\psi(f(x))\,\mathrm dP(x)\ge0.
\]
Jeśli całka jest zerowa, to $\psi(f)=0$ $P$-prawie wszędzie,
a więc $f=1$ $P$-prawie wszędzie, czyli $Q=P$.
W pozostałym przypadku definicja daje nieskończoność.
\end{proof}

\begin{twierdzenie}[Zasada przetwarzania danych dla entropii]
\label{thm:kl-data-processing-19}
Niech $P,Q$ będą prawdopodobieństwami na $(X,\mathcal F)$,
$Q\ll P$, a $T:(X,\mathcal F)\to(Y,\mathcal G)$
odwzorowaniem mierzalnym. Wówczas
\[
 D(T_*Q\Vert T_*P)\le D(Q\Vert P).
\]
\end{twierdzenie}
\begin{proof}
Niech $f=\mathrm dQ/\mathrm dP$, a
$g=\mathrm d(T_*Q)/\mathrm d(T_*P)$.
Z twierdzenia~\ref{thm:image-ac-19} wynika
$g(T)=\mathbb E_P[f\mid\sigma(T)]$ prawie wszędzie.
Dla funkcji wypukłej
$\psi(t)=t\log t-t+1$ na $[0,\infty)$
warunkowa nierówność Jensena pozostaje prawdziwa
również wtedy, gdy $\psi(f)$ nie jest całkowalna:
stosujemy ją do funkcji nieujemnej $\psi$ po odjęciu
wspierającej funkcji afinicznej, lub aproksymujemy
$\psi$ supremum przeliczalnej rodziny prostych
podpierających. Wówczas
\[
 \psi(g(T))\le\mathbb E_P[\psi(f)\mid\sigma(T)]
\]
w rozumieniu rozszerzonej warunkowej wartości oczekiwanej
funkcji nieujemnej. Jeżeli $D(Q\Vert P)=\infty$,
nierówność jest natychmiastowa. W przypadku skończonym
$\psi(f)\in L^1(P)$, a powyższą nierówność można
uzyskać wprost z dowodu twierdzenia
\ref{thm:conditional-jensen-19}, z rozszerzeniem na funkcję
$\psi$ na $[0,\infty)$; nieujemność $f$ i $g$
pozwala używać wyłącznie podpór w punktach dodatnich.
Całkując względem $P$ i stosując definicję miary obrazowej,
otrzymujemy
\[
 D(T_*Q\Vert T_*P)
=\int_Y\psi(g(y))\,\mathrm d(T_*P)(y)
\le\int_X\psi(f(x))\,\mathrm dP(x)=D(Q\Vert P).
\]
\end{proof}

\begin{uwaga}[Interpretacja]
Wzory dla wariacji całkowitej i entropii względnej
pokazują dwa różne sensy zasady, że samo przekształcenie
obserwacji nie może zwiększyć informacji o tym,
z którego z dwóch rozkładów pochodzi wynik. Są to
nierówności matematyczne; nie wymagają aparatu teorii
informacji ani wcześniejszej znajomości statystyki.
\end{uwaga}

\section{Ograniczenia i dalsze kierunki}
Pochodna Radona--Nikodyma jest określona prawie wszędzie,
a nie punktowo kanonicznie. Wartość $f(x)$ w danym punkcie
nie musi być znacząca; w zastosowaniach liczą się całki,
zbiory poziomicowe z dokładnością do zbiorów zerowych
oraz odpowiednie klasy funkcji. Warunkowe rozkłady
jako rodziny miar $P(\,\cdot\mid Y=y)$ wymagają dodatkowych
założeń, np. standardowej borelowskości przestrzeni,
czym zajmiemy się w module o dezintegracji.

\begin{przyklad}[Brak punktowej jednoznaczności rozkładu warunkowego]
Niech $\Omega=(0,1)$, $\mathcal F=\mathcal B((0,1))$,
$\mathbb P=\lambda_1|_{(0,1)}$, a $Y(\omega)=\Theta(\omega)=\omega$.
Dla $y\in(0,1)$ i $A\in\mathcal B((0,1))$ zdefiniujmy
$K(y,A)=\chi_A(y)$. Dla każdego $y$ odwzorowanie
$A\mapsto K(y,A)$ jest miarą Diraca $\delta_y$,
zaś dla każdego $A$ funkcja $y\mapsto K(y,A)$ jest
mierzalna. Ponadto dla dowolnych zbiorów borelowskich
$A,B\subseteq(0,1)$ zachodzi
\[
 \int_B K(y,A)\,\mathrm dP_Y(y)
 =\lambda_1(A\cap B)
 =\mathbb P(\Theta\in A,\,Y\in B).
\]
Zatem $K$ jest rozkładem warunkowym zmiennej
$\Theta$ względem obserwacji $Y$.

Ustalmy teraz $y_0\in(0,1)$ i miarę
probabilistyczną $\rho\ne\delta_{y_0}$ na $(0,1)$. Jeżeli zastąpimy
$K(y_0,\cdot)$ przez $\rho$, pozostawiając
$K(y,\cdot)=\delta_y$ dla $y\ne y_0$, otrzymamy
inne jądro mierzalne spełniające tę samą tożsamość
całkową: zbiór $\{y_0\}$ ma miarę $P_Y$ równą zero.
Warunkowy rozkład jest więc określony jedynie
$P_Y$-prawie wszędzie, nawet w tak prostej przestrzeni
standardowej borelowskiej.
\end{przyklad}

% ROZDZIAŁ 20 — miary Radona i ich regularność
\chapter{Miary borelowskie, regularność i miary Radona}
\label{ch:radon-measures-20}

\providecommand{\atlasref}[2]{#2}


\section{Dlaczego miara potrzebuje topologii?}

\noindent\emph{Powiązania tematyczne w Bibliotece matematycznej:} \atlasref{GT-19::def:topog-polskie-borel}{GT-19: zbiory borelowskie}.

Do tej pory wszystkie miary były określone na przestrzeniach
mierzalnych. Jeśli jednak przestrzeń ma topologię, można pytać,
czy wartości miary zbioru dają się odtworzyć z miar zbiorów
otwartych, domkniętych lub zwartych. Są to trzy różne własności:
regularność zewnętrzna, regularność wewnętrzna i ciasność.
Rozróżnianie ich będzie niezbędne w twierdzeniu reprezentacyjnym
Riesza--Markowa--Kakutaniego.

\begin{definicja}[Miara borelowska i lokalna skończoność]
\label{def:borel-locfinite-20}
Niech $(X,\tau)$ będzie przestrzenią topologiczną,
a $\mathcal B(X)$ jej $\sigma$-algebrą borelowską.
\emph{Miarą borelowską} na $X$ nazywamy nieujemną miarę
$\mu:\mathcal B(X)\to[0,\infty]$.
Miara $\mu$ jest \emph{lokalnie skończona}, jeśli dla
każdego $x\in X$ istnieje otwarte otoczenie $U_x$ punktu $x$
takie, że $\mu(U_x)<\infty$.
\end{definicja}

\begin{uwaga}[Skończoność na zbiorach zwartych]
\label{rem:compact-finiteness-20}
Na dowolnej przestrzeni topologicznej lokalna skończoność miary borelowskiej
implikuje, że $\mu(K)<\infty$ dla każdego zwartego
\emph{borelowskiego} zbioru $K\subseteq X$ (w przestrzeni Hausdorffa każdy zbiór zwarty jest domknięty, a więc borelowski): ze skończonego podpokrycia
$K\subseteq U_{x_1}\cup\cdots\cup U_{x_m}$ dostajemy
$\mu(K)\le\sum_i\mu(U_{x_i})<\infty$.
Na przestrzeni lokalnie zwartej Hausdorffa zachodzi
również implikacja odwrotna: każdy punkt posiada otwarte
otoczenie o zwartym domknięciu, więc jego miara nie
przekracza skończonej miary tego domknięcia.
W przestrzeni, która nie jest przestrzenią Hausdorffa, zbiór zwarty nie musi być borelowski, dlatego samo wyrażenie $\mu(K)$ może nie być określone. W dalszej części rozdziału przestrzenie są Hausdorffa. Bez lokalnej zwartości powyższy dowód implikacji odwrotnej nie ma zastosowania.
\end{uwaga}

\begin{definicja}[Regularność względem topologii]
\label{def:regularity-20}
Niech $X$ będzie przestrzenią Hausdorffa, a $\mu$ miarą
borelowską na $X$. Mówimy, że jest
\emph{zewnętrznie regularna}, gdy dla każdego
$E\in\mathcal B(X)$ zachodzi
\begin{equation}\label{eq:outer-regular-20}
 \mu(E)=\inf\{\mu(U):E\subseteq U,\ U\subseteq X
 \text{ otwarty}\}.
\end{equation}
Mówimy, że jest \emph{wewnętrznie regularna względem
zbiorów domkniętych}, gdy dla każdego $E\in\mathcal B(X)$
\[
 \mu(E)=\sup\{\mu(F):F\subseteq E,\ F\subseteq X
 \text{ domknięty}\}.
\]
Jest \emph{wewnętrznie regularna względem zbiorów zwartych},
gdy dla każdego $E\in\mathcal B(X)$
\begin{equation}\label{eq:inner-compact-20}
 \mu(E)=\sup\{\mu(K):K\subseteq E,\ K\subseteq X
 \text{ zwarty}\}.
\end{equation}
Przyjmujemy $\inf\varnothing=+\infty$ i
$\sup\varnothing=0$ dla podanych rodzin miar nieujemnych.
\end{definicja}

\begin{definicja}[Miara Radona --- przyjęta konwencja]
\label{def:radon-measure-20}
Niech $X$ będzie lokalnie zwartą przestrzenią Hausdorffa.
\emph{Miarą Radona} na $X$ nazywamy miarę borelowską $\mu$,
która jest skończona na zbiorach zwartych, regularna
zewnętrznie na wszystkich zbiorach borelowskich oraz
\emph{regularna wewnętrznie na zbiorach otwartych}, to znaczy
\[
 \mu(U)=\sup\{\mu(K): K\subseteq U,\ K\text{ zwarty}\}
 \qquad(U\subseteq X\text{ otwarty}).
\]
Skończoność na zbiorach zwartych jest na $X$ równoważna
lokalnej skończoności (uwaga~\ref{rem:compact-finiteness-20}).
\end{definicja}

\begin{uwaga}[Dwie konwencje regularności]
W definicji~\ref{def:radon-measure-20} regularność
wewnętrzną wymagamy na zbiorach otwartych, a nie na wszystkich
borelowskich. To rozróżnienie jest konieczne na ogólnych
przestrzeniach lokalnie zwartych Hausdorffa. Silniejsza
regularność na wszystkich zbiorach borelowskich zachodzi,
jeśli przestrzeń jest $\sigma$-zwarta lub jeśli rozważany
zbiór ma miarę $\sigma$-skończoną.
W rozdziale~\ref{ch:haar-23} podamy jawny przykład
miary Haara, dla której silniejszy warunek zawodzi.
Na lokalnie zwartych przestrzeniach metrycznych ośrodkowych
obie konwencje są równoważne.
\end{uwaga}

\begin{lemat}[Regularność na zbiorach $\sigma$-skończonych]
\label{lem:radon-sigmafinite-inner-20}
Niech $X$ będzie lokalnie zwartą przestrzenią Hausdorffa,
a $\mu$ miarą Radona według definicji
\ref{def:radon-measure-20}. Jeżeli $E\in\mathcal B(X)$
i $\mu(E)<\infty$, to
\[
 \mu(E)=\sup\{\mu(K):K\subseteq E,\ K\text{ zwarty}\}.
\]
Równość ta zachodzi także dla każdego borelowskiego
$E$ o mierze $\sigma$-skończonej, czyli dla którego
istnieją borelowskie $E_j\subseteq E$ ze
$\mu(E_j)<\infty$ i $E=\bigcup_jE_j$.
\end{lemat}
\begin{proof}
Niech $\mu(E)<\infty$ i $\varepsilon>0$.
Z regularności zewnętrznej wybieramy otwarty
$U\supseteq E$ z $\mu(U)<\mu(E)+\varepsilon$.
Mamy $\mu(U\setminus E)<\varepsilon$.
Następnie wybieramy otwarty $W\supseteq U\setminus E$
z $\mu(W)<\mu(U\setminus E)+\varepsilon<2\varepsilon$,
a z regularności wewnętrznej na otwartym $U$ zbiór zwarty
$F\subseteq U$ z $\mu(F)>\mu(U)-\varepsilon$.
Wówczas $K=F\setminus W$ jest zwarty jako domknięty
podzbiór zwartego $F$. Ponieważ $F\subseteq U$
oraz $U\setminus E\subseteq W$, mamy $K\subseteq E$.
Ze skończoności $\mu(F)$ i $\mu(W)$ oraz z rozkładu
$F=K\sqcup(F\cap W)$ dostajemy
\[
 \mu(K)=\mu(F)-\mu(F\cap W)
 \ge\mu(F)-\mu(W)>\mu(E)-3\varepsilon.
\]
Dowolność $\varepsilon$ daje nierówność od dołu
dla supremum miar zbiorów zwartych zawartych w $E$,
a nierówność odwrotna wynika z monotoniczności miary.
Dla zbioru o mierze $\sigma$-skończonej zastępujemy dane pokrycie zbiorami skończonej miary jego rosnącymi skończonymi sumami $E_j\subset E$. Jeśli $\mu(E)=\infty$, dla dowolnego $M>0$ wybieramy $j$ z $\mu(E_j)>M+1$ na mocy ciągłości miary z dołu, a następnie zwarty $K\subseteq E_j$ o mierze większej od $\mu(E_j)-1$. Dostajemy $\mu(K)>M$. Jeśli $\mu(E)<\infty$, stosujemy pierwszy przypadek bezpośrednio. W szczególności gdy $X$ jest $\sigma$-zwarta, każdy borelowski $E$ ma miarę $\sigma$-skończoną: przecinamy go z przeliczalnym pokryciem zwartymi zbiorami skończonej miary.
\end{proof}

\section{Regularność skończonych miar borelowskich na przestrzeniach metrycznych}
\begin{twierdzenie}[Aproksymacja zbiorów borelowskich]
\label{thm:finite-metric-regular-20}
Niech $(X,d)$ będzie dowolną przestrzenią metryczną,
a $\mu$ skończoną miarą borelowską na $X$.
Dla każdego $E\in\mathcal B(X)$ i $\varepsilon>0$
istnieją zbiór domknięty $F$ i zbiór otwarty $U$ takie,
że $F\subseteq E\subseteq U$ oraz
\begin{equation}\label{eq:finite-regular-error-20}
 \mu(U\setminus F)<\varepsilon.
\end{equation}
W szczególności $\mu$ jest zewnętrznie regularna
i wewnętrznie regularna względem zbiorów domkniętych.
Nie twierdzimy w tym miejscu, że $F$ można wybrać zwarty.
\end{twierdzenie}
\begin{proof}
\emph{Krok 1: zbiory otwarte.}
Dla zbioru otwartego $U\neq X$ definiujemy
$F_n=\{x\in X:d(x,X\setminus U)\ge1/n\}$.
Funkcja $x\mapsto d(x,X\setminus U)$ jest
$1$-lipschitzowska, więc $F_n$ jest domknięty.
Z otwartości $U$ wynika, że każdy $x\in U$
ma dodatnią odległość od $X\setminus U$;
stąd $F_n\uparrow U$. Ponieważ $\mu(U)<\infty$,
ciągłość miary z dołu daje
$\mu(U\setminus F_n)\to0$. Jeżeli $U=X$,
wystarczy przyjąć $F_n=X$. Zatem każdy zbiór otwarty
spełnia warunek z tezy dla dowolnego $\varepsilon>0$.

\emph{Krok 2: rodzina zbiorów aproksymowalnych.}
Oznaczmy przez $\mathcal D$ rodzinę tych zbiorów
$E\in\mathcal B(X)$, dla których warunek tezy zachodzi
dla każdego $\varepsilon>0$. Rodzina ta zawiera wszystkie
zbiory otwarte. Wykażemy, że jest $\sigma$-algebrą.
Jeżeli $F\subseteq E\subseteq U$ i
$\mu(U\setminus F)<\varepsilon$, to
$X\setminus U\subseteq X\setminus E\subseteq X\setminus F$,
a zbiór $X\setminus U$ jest domknięty,
$X\setminus F$ zaś otwarty. Różnica między nimi
jest równa $U\setminus F$, więc $X\setminus E\in\mathcal D$.

Niech teraz $E=\bigcup_{n\ge1}E_n$, gdzie $E_n\in\mathcal D$.
Dla danego $\varepsilon>0$ dobieramy
$F_n\subseteq E_n\subseteq U_n$, z $F_n$ domkniętym,
$U_n$ otwartym oraz
$\mu(U_n\setminus F_n)<\varepsilon2^{-n-3}$.
Ponieważ $\mu(E)<\infty$ i
$\bigcup_{n\le N}E_n\uparrow E$, istnieje $N$ takie,
że $\mu(E\setminus\bigcup_{n\le N}E_n)<\varepsilon/4$.
Stawiamy $F=\bigcup_{n\le N}F_n$ i $U=\bigcup_{n\ge1}U_n$.
Zbiór $F$ jest domknięty jako suma skończona,
$U$ jest otwarty, a $F\subseteq E\subseteq U$.
Ponieważ
\[
U\setminus F\subseteq
(U\setminus E)\cup(E\setminus F),
\quad
U\setminus E\subseteq\bigcup_{n\ge1}(U_n\setminus E_n),
\]
oraz
$E\setminus F\subseteq
(E\setminus\bigcup_{n\le N}E_n)
\cup\bigcup_{n\le N}(E_n\setminus F_n)$,
podaddytywność daje jawne oszacowanie
\[
 \mu(U\setminus F)\le
 \sum_{n\ge1}\varepsilon2^{-n-3}+
 \varepsilon/4+\sum_{n\le N}\varepsilon2^{-n-3}
 \le\varepsilon/2<\varepsilon.
\]
W konsekwencji $\mathcal D$ jest $\sigma$-algebrą.

\emph{Krok 3: zakończenie.}
Ponieważ $\mathcal D$ zawiera wszystkie zbiory otwarte,
zawiera także $\sigma$-algebrę przez nie generowaną,
czyli $\mathcal B(X)$. Równości w definicjach regularności
wynikają z monotoniczności miary oraz
$\mu(E\setminus F),\mu(U\setminus E)\le\mu(U\setminus F)$.
\end{proof}

\begin{uwaga}[Dlaczego nie można automatycznie zastąpić $F$ zbiorem zwartym?]
W ogólnej przestrzeni metrycznej zbiór domknięty
i ograniczony nie musi być zwarty. Na przykład zamknięta
kula jednostkowa w nieskończenie wymiarowej przestrzeni
Hilberta nie jest zwarta. Istotnie, wybieramy indukcyjnie
ciąg ortonormalny $(e_n)$: na każdym kroku dopełnienie
ortogonalne przestrzeni rozpiętej przez poprzednie wektory
zawiera niezerowy wektor, który normalizujemy. Dla $n\ne m$
mamy $\|e_n-e_m\|^2=2$, więc żaden podciąg tego ciągu
nie jest ciągiem Cauchy'ego, a tym bardziej nie jest zbieżny.
Nawet lokalna zwartość nie wystarcza do zwartości
zbiorów domkniętych i ograniczonych: nieskończony zbiór
z metryką dyskretną $d(x,y)=1$ dla $x\ne y$
jest lokalnie zwarty (singletony są otwarte i zwarte),
domknięty i ograniczony, lecz pokrycie singletonami
nie ma skończonego podpokrycia.
Przykłady wyjaśniają jedynie, dlaczego nie wolno
zastąpić domkniętości i ograniczoności zwartością;
nie wykluczają regularności względem zbiorów zwartych
miar na tych przestrzeniach. Aproksymacja przez
zbiory zwarte wymaga osobnego argumentu.
\end{uwaga}

\section{Ciasność i miary na przestrzeniach polskich}

\noindent\emph{Powiązania tematyczne w Bibliotece matematycznej:} \atlasref{GT-19::def:topog-polskie}{GT-19: przestrzenie polskie}.

\begin{definicja}[Ciasność pojedynczej miary]
\label{def:tight-single-20}
Niech $X$ będzie przestrzenią Hausdorffa, a $\mu$
skończoną miarą borelowską. Mówimy, że $\mu$ jest
\emph{ciasna}, jeżeli dla każdego $\varepsilon>0$
istnieje zwarty zbiór $K\subseteq X$ taki,
że $\mu(X\setminus K)<\varepsilon$.
\end{definicja}

\begin{twierdzenie}[Ciasność skończonych miar na przestrzeniach polskich]
\label{thm:polish-tight-20}
Niech $X$ będzie przestrzenią polską, tj. przestrzenią
topologiczną z metryką zupełną i ośrodkiem przeliczalnym,
zgodną z jej topologią. Każda skończona miara borelowska
$\mu$ na $X$ jest ciasna. Ponadto dla każdego
$E\in\mathcal B(X)$ zachodzi
\[
 \mu(E)=\sup\{\mu(K):K\subseteq E,\ K\text{ zwarty}\}.
\]
\end{twierdzenie}
\begin{proof}
Jeżeli $X=\varnothing$, wybieramy $K=\varnothing$ i teza jest
spełniona. W przeciwnym razie wybieramy zgodną metrykę
zupełną $d$ oraz niepusty co najwyżej przeliczalny
zbiór gęsty $D=\{x_j:j\ge1\}$ (w przypadku skończonym
możemy powtarzać jego elementy). Ustalmy $\varepsilon>0$.
Dla każdego $n\ge1$ kule otwarte
$B(x_j,2^{-n})$, $j\ge1$, pokrywają $X$.
Z ciągłości miary z dołu (twierdzenie \ref{thm:measure-continuity}),
stosowanej do ich skończonych sum, oraz ze skończoności
$\mu(X)$ wynika, że istnieje $m_n$ takie, że
\[
\mu\left(X\setminus\bigcup_{j=1}^{m_n}B(x_j,2^{-n})\right)
<\varepsilon 2^{-n-1}.
\]
Niech $C_n=\bigcup_{j=1}^{m_n}\overline{B}(x_j,2^{-n})$.
Zbiór $C_n$ jest domknięty i
$\mu(X\setminus C_n)<\varepsilon2^{-n-1}$.
Kładziemy $K=\bigcap_{n\ge1}C_n$.
Zbiór $K$ jest domknięty, a dla każdego $n$
zawiera się w sumie skończenie wielu kul o promieniu
$2^{-n}$. Jest zatem całkowicie ograniczony:
dla każdego $\delta>0$ wybieramy $n$ tak duże,
by $2^{1-n}<\delta$, i otrzymujemy skończone
$\delta$-pokrycie $K$ przez kule o środkach w $K$:
z każdej kuli domkniętej w $C_n$, która przecina $K$,
wybieramy punkt $y\in K$; jej przecięcie z $K$ zawiera się
w $B(y,\delta)$ z nierówności trójkąta i $2^{1-n}<\delta$.
Ponieważ $K$ jest domknięty w zupełnym $X$, jest zupełny.
Zupełność i całkowita ograniczoność implikują zwartość:
dla dowolnego ciągu w $K$ wybieramy kolejno nieskończone
podciągi zawarte w jednej z kul skończonego pokrycia
promieniami $2^{-m}$, $m=1,2,\ldots$. Wybór przekątniowy
daje podciąg Cauchy'ego, ponieważ odległość dwóch wyrazów
od $m$-tego miejsca nie przekracza $2^{1-m}$.
Jest on zbieżny w $K$ dzięki zupełności. W przestrzeni
metrycznej zwartość ciągowa jest równoważna zwartości. Wreszcie
\[
 \mu(X\setminus K)
 \le\sum_{n\ge1}\mu(X\setminus C_n)<\varepsilon.
\]
Udowodniliśmy ciasność.

Teraz ustalmy $E\in\mathcal B(X)$. Z twierdzenia
\ref{thm:finite-metric-regular-20} wybieramy domknięty
$F\subseteq E$ taki, że $\mu(E\setminus F)<\varepsilon/2$.
Z ciasności wybieramy zwarty $K_0\subseteq X$
z $\mu(X\setminus K_0)<\varepsilon/2$.
Zbiór $K=F\cap K_0$ jest zwarty, zawiera się w $E$, a
$\mu(E\setminus K)\le\mu(E\setminus F)+\mu(X\setminus K_0)
<\varepsilon$. Monotoniczność daje żądaną równość.
\end{proof}

\begin{przyklad}[Rozkłady gaussowskie w nieskończonym wymiarze]
Twierdzenie~\ref{thm:polish-tight-20} ma zastosowanie
do każdej skończonej miary borelowskiej na ośrodkowej
przestrzeni Banacha z metryką normową, również w nieskończonym
wymiarze. Nie oznacza to, że domknięte kule są zwarte;
zwarty zbiór zapewniający ciasność może mieć znacznie
bardziej skomplikowaną geometrię. W szczególności
przestrzeń $\ell^2$ jest polska, więc jej skończone
miary borelowskie są ciasne.
\end{przyklad}

\section{Miary Radona na przestrzeniach lokalnie zwartych}

\noindent\emph{Powiązania tematyczne w Bibliotece matematycznej:} \atlasref{GT-14::def:topog-baire-lokalna-zwartosc}{GT-14: lokalna zwartość}.

\begin{twierdzenie}[Lokalnie skończone miary borelowskie]
\label{thm:lch-secondcountable-radon-20}
Niech $X$ będzie lokalnie zwartą, ośrodkową
przestrzenią metryczną, a $\mu$ lokalnie skończoną
miarą borelowską na $X$. Wówczas $\mu$ jest miarą
Radona w sensie definicji~\ref{def:radon-measure-20}.
Ponadto $\mu$ jest $\sigma$-skończona.
\end{twierdzenie}
\begin{proof}
\emph{Krok 1: otwarte pokrycie skończonej miary.}
Jeśli $X=\varnothing$, teza zachodzi dla miary zerowej.
W przypadku niepustym skończone podpokrycie można
uzupełnić zbiorami pustymi do ciągu.
Dla każdego $x\in X$ istnieje otwarte otoczenie $W_x$
z $\mu(W_x)<\infty$. Wybieramy otwarte otoczenie $O_x$ o zwartym domknięciu, a następnie promień $r>0$ taki, że $B(x,2r)\subset O_x\cap W_x$. Dla $V_x=B(x,r)$ domknięcie $\overline{V_x}$ zawiera się w $B(x,2r)\subset O_x\cap W_x$ i jako domknięty podzbiór zwartego $\overline{O_x}$ jest zwarte.
W szczególności $\mu(\overline{V_x})<\infty$.
Ponieważ przestrzeń metryczna ośrodkowa spełnia drugi
aksjomat przeliczalności i ma własność Lindelöfa
(por. \atlasref{GT-09::thm:topog-metryczne-trzy-warunki}{GT-09: równoważność ośrodkowości, drugiego aksjomatu przeliczalności i własności Lindelöfa}),
z pokrycia $(V_x)_{x\in X}$ wybieramy przeliczalne
podpokrycie $(V_n)_{n\ge1}$. Dla każdego $n$
$\overline{V_n}$ jest zwarty i ma skończoną miarę.
Już stąd wynika $\sigma$-skończoność $\mu$.
Ponadto każdy zwarty $K\subseteq X$ ma skończoną
miarę: z pokrycia $(V_n)$ wybieramy skończone
podpokrycie $K\subseteq V_{n_1}\cup\dots\cup V_{n_m}$,
a z podaddytywności
\[
 \mu(K)\le\sum_{j=1}^m\mu(V_{n_j})<\infty.
\]
Zwartość w przestrzeni Hausdorffa zapewnia przy tym,
że $K$ jest domknięty, a więc borelowski.

\emph{Krok 2: regularność zewnętrzna.}
Rozbijamy $X$ na parami rozłączne zbiory borelowskie
$B_1=V_1$, $B_n=V_n\setminus\bigcup_{j<n}V_j$.
Dla danego $E\in\mathcal B(X)$ i $\varepsilon>0$
rozważamy skończoną miarę
$\mu_n(A)=\mu(A\cap V_n)$ na przestrzeni metrycznej $X$.
Z twierdzenia~\ref{thm:finite-metric-regular-20}
wybieramy otwarty $U_n\supseteq E\cap B_n$ taki, że
$\mu_n(U_n\setminus(E\cap B_n))<\varepsilon2^{-n-1}$.
Niech $U=\bigcup_{n\ge1}(U_n\cap V_n)$.
Wówczas $U$ jest otwarty, $E\subseteq U$, ponieważ
$E\cap B_n\subseteq U_n\cap V_n$. Zbiory $B_n$ przypisują każdemu punktowi dokładnie jeden indeks pokrycia. Rozłączność nie jest jednak konieczna dla oszacowania regularności: analogicznie można aproksymować $E\cap V_n$, gdyż błędy sumuje się za pomocą podaddytywności.
Ponadto $(U_n\cap V_n)\setminus E\subseteq (U_n\cap V_n)\setminus(E\cap B_n)$, a miara tego drugiego zbioru wynosi $\mu_n(U_n\setminus(E\cap B_n))$. Stąd
\[
\mu(U\setminus E)
\le\sum_{n\ge1}\mu((U_n\cap V_n)\setminus(E\cap B_n))
<\varepsilon.
\]
Jeżeli $\mu(E)<\infty$, daje to
$\mu(U)<\mu(E)+\varepsilon$; jeżeli $\mu(E)=\infty$,
regularność zewnętrzna jest natychmiastowa,
ponieważ każdy otwarty nadzbiór ma miarę nieskończoną.

\emph{Krok 3: regularność wewnętrzna względem zwartych.}
Niech $E\in\mathcal B(X)$ i $\varepsilon>0$.
Dla każdego $n$ rozważamy przestrzeń zwartą
$\overline{V_n}$ z miarą skończoną
$\mu|_{\overline{V_n}}$ i metryką indukowaną.
Zbiór $E\cap V_n$ jest borelowski w podprzestrzeni $\overline{V_n}$, ponieważ $\mathcal B(\overline{V_n})$ jest śladem $\mathcal B(X)$ (otwarte zbiory podprzestrzeni są przecięciami ze zbiorami otwartymi w $X$, a ślad zachowuje dopełnienia względne i przeliczalne sumy). Z twierdzenia~\ref{thm:finite-metric-regular-20}
dla zbioru $E\cap V_n$ wybieramy zbiór
$K_n$ domknięty w $\overline{V_n}$, zawarty w
$E\cap V_n$, o własności
$\mu((E\cap V_n)\setminus K_n)<\varepsilon2^{-n-1}$.
Każdy $K_n$ jest zwarty w $X$.
Dla $E$ o skończonej mierze wybieramy $N$ takie, że
$\mu(E\setminus\bigcup_{n\le N}V_n)<\varepsilon/2$.
Zbiór $K=\bigcup_{n\le N}K_n$ jest zwarty i zawiera się
w $E$, a
\[
 \mu(E\setminus K)
 \le\mu(E\setminus\bigcup_{n\le N}V_n)
 +\sum_{n\le N}\mu((E\cap V_n)\setminus K_n)
 <\varepsilon.
\]
Gdy $\mu(E)=\infty$, z ciągłości miary z dołu
$\mu(E\cap\bigcup_{n\le N}V_n)\uparrow\infty$.
Dla dowolnego $M>0$ wybieramy $N$ z miarą tej części
większą niż $M+1$ i ponawiamy aproksymację przez
$K_n$ dla $1\le n\le N$, z błędami sumującymi się do mniej niż $1$. Zbiór $E\cap\bigcup_{n\le N}V_n$ ma skończoną miarę, a różnica między nim i $K=\bigcup_{n\le N}K_n$ ma miarę mniejszą niż $1$. Dostajemy zwarty $K\subseteq E$ z $\mu(K)>M$.
To dowodzi \eqref{eq:inner-compact-20} także w przypadku
nieskończonej masy. Razem z krokami 1--2 daje tezę.
\end{proof}

\begin{uwaga}[Na jakich przestrzeniach twierdzenie obowiązuje?]
Założenie ośrodkowości i metryczności nie jest częścią
definicji miary Radona; zostało wprowadzone w powyższym
twierdzeniu, aby można było podać samodzielny dowód.
Na ogólnych lokalnie zwartych przestrzeniach Hausdorffa
regularność nie każdej miary borelowskiej jest automatyczna.
Nie wystarczy również samo założenie lokalnej skończoności.
\end{uwaga}

\begin{przyklad}[Lokalna skończoność nie wystarcza]
\label{ex:nonradon-discrete-20}
Niech $X$ będzie nieprzeliczalnym zbiorem z topologią
dyskretną. Jest to lokalnie zwarta przestrzeń metryczna,
ale nieośrodkowa. Wszystkie podzbiory są borelowskie.
Definiujemy
\[
\mu(A)=\begin{cases}0,&A\text{ jest co najwyżej przeliczalny},\\
+\infty,&A\text{ jest nieprzeliczalny}.
\end{cases}
\]
To rzeczywiście miara: $\mu(\varnothing)=0$, a dla
przeliczalnej rodziny zbiorów rozłącznych ich suma jest
nieprzeliczalna wtedy i tylko wtedy, gdy któryś ze zbiorów
jest nieprzeliczalny. W takim przypadku obie strony
równości addytywności są nieskończone, a w przeciwnym
przypadku obie są zerowe.
Miara jest lokalnie skończona, ponieważ każde otwarte
jednoelementowe otoczenie ma miarę zero. Tymczasem
zbiory zwarte w przestrzeni dyskretnej są skończone,
więc wszystkie mają miarę zero, choć $\mu(X)=\infty$.
Regularność wewnętrzna względem zwartych zawodzi.
\end{przyklad}

\section{Przykłady miar Radona}

\begin{przyklad}[Dwie konwencje regularności na nieprzeliczalnej sumie przestrzeni]
\label{it:ex:radon-two-regularities}
Niech $J$ będzie zbiorem nieprzeliczalnym, a
$X=J\times[0,1]$, gdzie $J$ wyposażamy w topologię
dyskretną, a $[0,1]$ w zwykłą topologię.
Przestrzeń $X$ jest lokalnie zwarta i metryzowalna
(na różnych składowych przyjmujemy odległość $2$,
a na tej samej zwykłą odległość w $[0,1]$).
Każda składowa $X_j=\{j\}\times[0,1]$ jest otwarta
i domknięta oraz zwarta, co daje lokalną zwartość.
Metryka opisana powyżej spełnia nierówność trójkąta:
jeśli końce leżą w różnych składowych, co najmniej
jeden odcinek pośredni ma długość $2$; jeśli leżą
w jednej składowej, stosujemy zwykłą nierówność
trójkąta na $[0,1]$ albo zauważamy, że przejście
przez inną składową ma długość $4$.
Kule o promieniu mniejszym od $1$ generują topologię
na każdej składowej, więc metryka daje zadaną topologię.
Przestrzeń nie jest $\sigma$-zwarta: każdy zwarty podzbiór
przecina tylko skończenie wiele składowych, ponieważ
składowe tworzą otwarte pokrycie $X$. Przeliczalna
suma zwartych zbiorów przecina więc co najwyżej
przeliczalnie wiele składowych.

Dla $E\in\mathcal B(X)$ oznaczmy
$E_j=\{t\in[0,1]:(j,t)\in E\}$ oraz
$J(E)=\{j\in J:E_j\ne\varnothing\}$.
Każdy $E_j$ jest borelowski jako przeciwobraz
$E$ przez ciągłe odwzorowanie $t\mapsto(j,t)$.
Określmy dwie funkcje zbioru:
\[
 \nu(E)=\sum_{j\in J}\lambda_1(E_j),\qquad
 \mu(E)=
 \begin{cases}
 \nu(E),& J(E)\text{ jest przeliczalny},\\
 +\infty,& J(E)\text{ jest nieprzeliczalny}.
 \end{cases}
\]
Sumę nieprzeliczalnej rodziny liczb nieujemnych
rozumiemy jako supremum sum po skończonych
podrodzinach.

Funkcja $\nu$ jest miarą. Istotnie, dla rozłącznych
$E_n$ i skończonego $F\subseteq J$ mamy
\begin{align*}
 \sum_{j\in F}\lambda_1((\bigcup_n E_n)_j)
 &=\sum_n\sum_{j\in F}\lambda_1((E_n)_j)\\
 &\le\sum_n\nu(E_n).
\end{align*}
Supremum po $F$ daje jedną nierówność.
Dla każdego $N$ i skończonych $F_1,\ldots,F_N\subseteq J$
mamy
$\sum_{n=1}^N\sum_{j\in F_n}\lambda_1((E_n)_j)
\le\nu(\bigcup_n E_n)$, stosując skończoną sumę
po $F_1\cup\cdots\cup F_N$ i addytywność na składowych.
Niezależne supremum po $F_n$ i następnie po $N$
daje nierówność odwrotną.
Funkcja $\mu$ również jest miarą: jeżeli choć jeden
$J(E_n)$ jest nieprzeliczalny, obie strony warunku
addytywności są nieskończone. W przeciwnym razie
$J(\bigcup_n E_n)=\bigcup_nJ(E_n)$ jest przeliczalny,
więc addytywność wynika z addytywności $\nu$.
Obie miary są lokalnie skończone, ponieważ każda
składowa $X_j$ jest otwarta i ma masę $1$.

Dla otwartego $U\subseteq X$ zachodzi $\mu(U)=\nu(U)$.
Gdy $J(U)$ jest nieprzeliczalny, wszystkie niepuste
$U_j$ są względnie otwarte i mają dodatnią miarę
Lebesgue'a. Ponieważ rodzina dodatnich liczb o skończonej
sumie ma co najwyżej przeliczalnie wiele wyrazów
(dla każdego $n\ge1$ tylko skończenie wiele może być
większych od $1/n$), w tym przypadku $\nu(U)=\infty$.
Dla przeliczalnego $J(U)$ równość wynika z definicji.

Miara $\mu$ jest zewnętrznie regularna. Jeśli
$\mu(E)=\infty$, wynika to z monotoniczności.
Jeżeli $\mu(E)<\infty$, to $J(E)$ jest przeliczalny;
numerujemy jego elementy i dla każdego z nich
(z numeracją rozpoczynającą się od $1$) aproksymujemy $E_j$ od góry zbiorem względnie otwartym
$U_j$ w $[0,1]$ o błędzie mniejszym niż
$\varepsilon2^{-n}$ dla $j$ o numerze $n$,
korzystając z twierdzenia
\ref{thm:finite-metric-regular-20}.
Suma otwartych zbiorów $\{j\}\times U_j$ zawiera $E$
i ma miarę $\mu$ mniejszą od $\mu(E)+\varepsilon$
(dla $E=\varnothing$ wybieramy zbiór pusty).
Miara $\mu$ jest też wewnętrznie regularna na otwartych:
dla skończonego $F\subseteq J$ przybliżamy każde $U_j$,
$j\in F$, zwartym $K_j\subseteq U_j$ na mocy
twierdzenia \ref{thm:finite-metric-regular-20}
na zwartym $[0,1]$, z błędem mniejszym niż
$\varepsilon/(1+|F|)$ na każdej składowej. Skończona suma
$K=\bigcup_{j\in F}(\{j\}\times K_j)$ jest zwarta
i $\mu(K)=\sum_{j\in F}\lambda_1(K_j)$.
Dla każdego skończonego $F$ mamy więc
$\mu(K)>\sum_{j\in F}\lambda_1(U_j)-\varepsilon$.
Dowolność błędu oraz supremum po $F$ dają
$\sup_{K\subseteq U,\ K\text{ zwarty}}\mu(K)
=\nu(U)=\mu(U)$, także dla nieskończonej masy.
Skończoność na zwartych wynika z lokalnej skończoności.
Zatem $\mu$ jest miarą Radona według definicji
\ref{def:radon-measure-20}.

Rozważmy jednak domknięty zbiór $E=J\times\{0\}$.
Mamy $\mu(E)=+\infty$, natomiast każdy zwarty
$K\subseteq E$ jest skończony i $\mu(K)=0$.
Regularność wewnętrzna względem zwartych nie zachodzi
więc na wszystkich zbiorach borelowskich.
Z kolei $\nu(E)=0$, lecz każdy otwarty nadzbiór $U$
zbioru $E$ ma $\nu(U)=+\infty$, zatem $\nu$ nie jest
zewnętrznie regularna, mimo lokalnej skończoności.
Obie miary zgadzają się na całkach funkcji z $C_c(X)$,
ponieważ ich nośniki przecinają tylko skończenie
wiele składowych, na których miary są identyczne.
To pokazuje znaczenie regularności w twierdzeniu
o jednoznaczności miary Radona.
\end{przyklad}

\begin{przyklad}[Miara Lebesgue'a jako miara Radona]
Na $\mathbb R^d$, $d\ge1$, ograniczenie miary Lebesgue'a
$\lambda_d$ do $\mathcal B(\mathbb R^d)$ jest miarą
lokalnie skończoną: każdy punkt ma otoczenie będące
kulą o skończonej objętości. Przestrzeń $\mathbb R^d$
jest lokalnie zwarta, metryczna i ośrodkowa.
Na mocy twierdzenia~\ref{thm:lch-secondcountable-radon-20}
miara ta jest Radona. Jej uzupełnienie na
$\mathcal L_d$ zawiera także zbiory nieborelowskie;
trzeba rozróżniać tę miarę zupełną od jej ograniczenia
borelowskiego.
\end{przyklad}

\begin{przyklad}[Miary z lokalnie całkowalną gęstością]
Niech $w:\mathbb R^d\to[0,\infty]$ będzie borelowska
i lokalnie całkowalna względem $\lambda_d$,
tj. $\int_Kw\,\mathrm d\lambda_d<\infty$ dla każdego
zwartego $K\subseteq\mathbb R^d$. Wówczas
$\mu_w(E)=\int_Ew(x)\,\mathrm d\lambda_d(x)$
dla $E\in\mathcal B(\mathbb R^d)$ określa miarę na mocy
lematu \ref{lem:integration-density-18}. Warunek lokalnej
całkowalności gwarantuje lokalną skończoność, ponieważ
kula otwarta o skończonym promieniu zawiera się w zwartej
kuli domkniętej. Regularność wynika z twierdzenia
\ref{thm:lch-secondcountable-radon-20}. Na przykład $w(x)=\exp(|x|^2)$ daje miarę
Radona o nieskończonej masie całkowitej.
\end{przyklad}

\begin{przyklad}[Atomy gęste i nośnik całej prostej]
Niech $(q_n)_{n\ge1}$ będzie numeracją zbioru
$\mathbb Q$ bez powtórzeń. Na $\mathbb R$ określmy
$\mu=\sum_{n\ge1}2^{-n}\delta_{q_n}$.
Jest to skończona miara borelowska, a więc miara Radona.
Każdy niepusty przedział otwarty zawiera pewien $q_n$
i ma dodatnią miarę $\mu$, choć $\mu(\mathbb R\setminus
\mathbb Q)=0$. W konsekwencji nośnikiem topologicznym
$\mu$ jest cała prosta. Pełny nośnik topologiczny
nie oznacza bezwzględnej ciągłości względem miary Lebesgue'a:
w tym przykładzie $\mu\perp\lambda_1$.
\end{przyklad}

\section{Nośnik topologiczny}
\begin{definicja}[Nośnik miary]
\label{def:support-radon-20}
Dla miary borelowskiej $\mu$ na przestrzeni topologicznej
$X$ określamy jej \emph{nośnik topologiczny} przez
\[
\operatorname{supp}\mu
:=\{x\in X:\mu(U)>0\text{ dla każdego otwartego }
U\ni x\}.
\]
Jest to zbiór domknięty. Nie jest to z definicji
wybrany mierzalny zbiór skupiający miarę w sensie
rozdziału~\ref{ch:ac-singular-17}.
\end{definicja}

\begin{twierdzenie}[Nośnik na przestrzeni spełniającej drugi aksjomat przeliczalności]
\label{thm:support-fullmass-20}
Niech $X$ będzie przestrzenią spełniającą drugi aksjomat przeliczalności,
a $\mu$ dowolną miarą borelowską na $X$. Wtedy
$\operatorname{supp}\mu$ jest domknięty,
$\mu(X\setminus\operatorname{supp}\mu)=0$ i jest
najmniejszym domkniętym zbiorem o pełnej mierze
w sensie zawierania.
\end{twierdzenie}
\begin{proof}
Z definicji
$X\setminus\operatorname{supp}\mu$
jest sumą wszystkich otwartych zbiorów o mierze zero,
a więc jest otwarty. Niech $(B_n)$ będzie przeliczalną
bazą topologii. Jeżeli $x$ należy do tej sumy, to
istnieje otwarty $U\ni x$ o $\mu(U)=0$.
Można wybrać $B_n$ takie, że $x\in B_n\subseteq U$.
Wówczas $\mu(B_n)=0$, a więc
\[
X\setminus\operatorname{supp}\mu
=\bigcup_{\{n:\mu(B_n)=0\}}B_n.
\]
Przeliczalna podaddytywność daje miarę zero tej sumy.
Niech $F\subseteq X$ będzie domknięty i
$\mu(X\setminus F)=0$. Każdy $x\notin F$ ma
otwarte otoczenie $X\setminus F$ miary zero,
zatem $x\notin\operatorname{supp}\mu$.
W konsekwencji $\operatorname{supp}\mu\subseteq F$.
\end{proof}

\begin{uwaga}[Dlaczego założenie przeliczalności ma znaczenie?]
W przykładzie~\ref{ex:nonradon-discrete-20} każda
jednopunktowa okolica ma miarę zero, zatem
$\operatorname{supp}\mu=\varnothing$, choć
$\mu(X)=+\infty$. Problem wynika z faktu, że
nieprzeliczalna suma zbiorów zerowych nie musi być
zerowa. Pokazuje to, dlaczego nie należy dowodzić
pełności miary nośnika przez nieuzasadnione przejście
od dowolnego pokrycia do przeliczalnej sumy.
Dla miary Radona w przyjętej konwencji drugi aksjomat
przeliczalności nie jest jednak potrzebny: otwarty zbiór
$U=X\setminus\operatorname{supp}\mu$ jest sumą otwartych
zbiorów zerowych. Każdy zwarty $K\subseteq U$ ma skończone
podpokrycie takimi zbiorami, zatem $\mu(K)=0$.
Regularność wewnętrzna na otwartych daje $\mu(U)=0$.
Argument minimalności z poprzedniego dowodu pozostaje ważny.
\end{uwaga}

\section{Funkcje ciągłe o zwartym nośniku}
Dla lokalnie zwartej przestrzeni Hausdorffa $X$ oznaczamy
przez $C_c(X)$ przestrzeń ciągłych funkcji
$f:X\to\mathbb R$, których \emph{nośnik funkcji}
$\operatorname{supp}f=\overline{\{x:f(x)\neq0\}}$
jest zwarty. W analogiczny sposób definiuje się
$C_c(X;\mathbb C)$. Nie należy mylić nośnika funkcji
z nośnikiem topologicznym miary.

\begin{lemat}[Funkcja odcinająca]
\label{lem:cutoff-radon-20}
Niech $X$ będzie lokalnie zwartą przestrzenią Hausdorffa,
$K\subseteq X$ zwartym zbiorem, a $U\subseteq X$
otwartym zbiorem zawierającym $K$. Wtedy istnieje
$f\in C_c(X)$ taka, że $0\le f\le1$, $f|_K=1$
oraz $\operatorname{supp}f\subseteq U$.
\end{lemat}
\begin{proof}
Jeśli $K=\varnothing$, wystarczy przyjąć $f=0$.
W dalszej części zakładamy $K\neq\varnothing$.
Z lematu \atlasref{GT-14::lem:topog-baire-zwarte-domkniecie}{GT-14: otoczenie o zwartym domknięciu} dla każdego $x\in K$
wybieramy otwarte $W_x\ni x$ o zwartym domknięciu zawartym
w $U$. Ze zwartości $K$ otrzymujemy skończone podpokrycie
$W_{x_1},\ldots,W_{x_m}$. Dla
$V=\bigcup_{i=1}^m W_{x_i}$ mamy
$K\subseteq V\subseteq\overline V=\bigcup_{i=1}^m\overline{W_{x_i}}\subseteq U$,
a $\overline V$ jest zwarty jako skończona suma zwartych zbiorów.

Przestrzeń $\overline V$ jest normalna na mocy twierdzenia
\atlasref{GT-07::thm:topog-zwartosc-hausdorff-normalna}{GT-07: normalność zwartych przestrzeni Hausdorffa}.
Zbiory $K$ i $\overline V\setminus V$ są w niej domknięte i rozłączne.
Z lematu \atlasref{GT-05::thm:topog-urysohn}{GT-05: lemat Urysohna} dla przestrzeni normalnych otrzymujemy ciągłą funkcję $g:\overline V\to[0,1]$, równą $1$ na $K$ oraz $0$ na $\overline V\setminus V$ (gdy drugi zbiór jest pusty, można przyjąć $g=1$). Przedłużamy $g$ zerem na $X\setminus V$. Definicje zgadzają się na przecięciu dwóch zbiorów domkniętych $\overline V$ i $X\setminus V$, które pokrywają $X$, więc funkcja $f$ jest ciągła. Jej nośnik zawiera się w $\overline V\subseteq U$, co daje tezę.

W przypadku metrycznym można zamiast lematu Urysohna użyć jawnego wzoru: jeżeli $V=X$, przyjmujemy $f=1$, a w przeciwnym razie określamy
\[
 f(x)=\frac{d(x,X\setminus V)}
 {d(x,X\setminus V)+d(x,K)}.
\]
Mianownik nie zeruje się nigdzie: dla $x\in K$
pierwszy składnik jest dodatni, a poza $K$ drugi
składnik jest dodatni (ponieważ $K$ jest domknięty).
Obie funkcje odległości są ciągłe, stąd $f$ jest
ciągła; na $K$ ma wartość $1$, poza $V$ jest zerowa,
a $\operatorname{supp}f\subseteq\overline V\subseteq U$.
\end{proof}

\begin{wniosek}[Charakteryzacja nośnika przez funkcje testowe]
\label{cor:support-test-20}
Niech $X$ będzie lokalnie zwartą przestrzenią Hausdorffa,
a $\mu$ miarą Radona na $X$. Dla każdego $x\in X$
następujące warunki są równoważne:
\begin{enumerate}[label=\textup{(\roman*)}]
\item $x\in\operatorname{supp}\mu$;
\item dla każdej funkcji $f\in C_c(X)$ takiej, że
$f\ge0$ i $f(x)>0$, zachodzi
$\int_X f(y)\,\mathrm d\mu(y)>0$.
\end{enumerate}
\end{wniosek}
\begin{proof}
Załóżmy (i) i wybierzmy $f$ jak w (ii).
Z ciągłości $f$ wynika, że istnieją $\delta>0$
oraz otwarte otoczenie $V$ punktu $x$ takie,
że $f(y)\ge\delta$ dla $y\in V$. Zatem
$\int_Xf\,\mathrm d\mu\ge\delta\mu(V)>0$.
Całka jest skończona, ponieważ $f$ jest
ograniczona i ma zwarty nośnik, a $\mu$
jest skończona na zbiorach zwartych.

Jeśli natomiast $x\notin\operatorname{supp}\mu$,
to istnieje otwarty zbiór $U\ni x$ z $\mu(U)=0$.
Lemat~\ref{lem:cutoff-radon-20}, zastosowany do
$K=\{x\}\subseteq U$, daje $f\in C_c(X)$
spełniającą $f\ge0$, $f(x)=1$ i
$\operatorname{supp}f\subseteq U$.
Stąd $\int_X f\,\mathrm d\mu=0$, co przeczy (ii).
\end{proof}

\begin{twierdzenie}[Jednoznaczność miary Radona przez całki]
\label{thm:uniqueness-cc-20}\label{tw:radon-unikalnosc}
Niech $X$ będzie lokalnie zwartą przestrzenią Hausdorffa,
a $\mu,\nu$ miarami Radona na $X$. Jeżeli
\[
\int_X f(x)\,\mathrm d\mu(x)
=\int_Xf(x)\,\mathrm d\nu(x)
\qquad\text{dla każdego }f\in C_c(X),
\]
to $\mu=\nu$ na $\mathcal B(X)$.
\end{twierdzenie}
\begin{proof}
Dla każdego otwartego $U\subseteq X$ i dowolnej
miary Radona $\rho$ na $X$ zachodzi
\begin{equation}\label{eq:cc-open-sup-20}
\rho(U)=\sup\left\{\int_Xf\,\mathrm d\rho:
 f\in C_c(X),\ 0\le f\le1,\
 \operatorname{supp}f\subseteq U\right\}.
\end{equation}
Istotnie, każda z powyższych całek jest co najwyżej
$\rho(U)$. Z drugiej strony dla każdego zwartego
$K\subseteq U$ lemat~\ref{lem:cutoff-radon-20}
daje $f$ równe $1$ na $K$, więc
$\int_Xf\,\mathrm d\rho\ge\rho(K)$.
Z regularności wewnętrznej miary $\rho$
uzyskujemy \eqref{eq:cc-open-sup-20} także dla
$\rho(U)=+\infty$. Przy założonej równości całek
zbiory funkcji w supremum są te same dla obu miar,
a więc $\mu(U)=\nu(U)$ dla wszystkich otwartych $U$.
Wreszcie dla dowolnego zbioru borelowskiego
$E\subseteq X$ z regularności zewnętrznej
$\mu(E)=\inf_{U\supseteq E}\mu(U)
=\inf_{U\supseteq E}\nu(U)=\nu(E)$.
\end{proof}

\begin{uwaga}[Związek z twierdzeniem reprezentacyjnym]
Twierdzenie~\ref{thm:uniqueness-cc-20} daje
\emph{jednoznaczność} miary reprezentującej funkcjonał
$L(f)=\int_Xf\,\mathrm d\mu$ na $C_c(X)$.
Istnienie takiej miary dla dodatniego funkcjonału
jest przedmiotem następnego rozdziału, poświęconego
twierdzeniu Riesza--Markowa--Kakutaniego. W przestrzeni
niezwartej funkcjonał całkowania względem miary
Radona o nieskończonej masie nie musi być ograniczony
w normie supremum na całym $C_c(X)$; dlatego rozróżnimy
wtedy $C_c(X)$ i $C_0(X)$.
\end{uwaga}

\begin{twierdzenie}[Gęstość $C_c(X)$ w przestrzeniach $L^p$]
\label{thm:cc-dense-lp-radon-20}
Niech $X$ będzie lokalnie zwartą przestrzenią Hausdorffa,
a $\mu$ miarą Radona w sensie definicji
\ref{def:radon-measure-20}. Wówczas dla każdego
$1\le p<\infty$ przestrzeń $C_c(X;\mathbb C)$ jest
gęsta w $L^p(X,\mathcal B(X),\mu;\mathbb C)$.
Analogiczny wniosek zachodzi dla funkcji rzeczywistych.
\end{twierdzenie}
\begin{proof}
Najpierw ustalmy $E\in\mathcal B(X)$ takie, że
$\mu(E)<\infty$. Dla $\varepsilon>0$ z lematu \ref{lem:radon-sigmafinite-inner-20}
(stosowanego do zbioru skończonej miary $E$) i regularności
zewnętrznej wybieramy zwarty $K$
i otwarty $U$ takie, że
$K\subseteq E\subseteq U$ oraz
$\mu(E\setminus K)<\varepsilon/2$,
$\mu(U\setminus E)<\varepsilon/2$.
Lemat~\ref{lem:cutoff-radon-20} daje $g\in C_c(X)$
taką, że $0\le g\le1$, $g|_K=1$ oraz
$\operatorname{supp}g\subseteq U$.
Na $K$ funkcje $\chi_E$ i $g$ są równe, poza $U$
obie znikają, a w pozostałej części ich różnica
ma moduł co najwyżej $1$. Zatem
\[
 \int_X|\chi_E(x)-g(x)|^p\,\mathrm d\mu(x)
 \le\mu(E\setminus K)+\mu(U\setminus E)
 <\varepsilon.
\]
Otrzymujemy aproksymację funkcji charakterystycznej
każdego zbioru o skończonej mierze.

Z twierdzenia \ref{thm:simple-dense-lp-14} (które nie
wymaga $\sigma$-skończoności miary) dla każdego $f\in L^p(\mu)$
i $\eta>0$ wybieramy funkcję prostą
$s=\sum_{j=1}^m a_j\chi_{E_j}$, gdzie $\mu(E_j)<\infty$,
z $\|f-s\|_p<\eta/2$. Jeżeli wszystkie $a_j$ są zerowe,
wystarczy funkcja zerowa. W przeciwnym razie z poprzedniego
akapitu wybieramy $g_j\in C_c(X)$ tak, aby
$\|\chi_{E_j}-g_j\|_p<\eta/(2(1+\sum_{k=1}^m|a_k|))$.
Nierówność Minkowskiego (twierdzenie \ref{thm:minkowski-14}) daje
\[
 \left\|f-\sum_{j=1}^m a_jg_j\right\|_p
 \le \|f-s\|_p+\sum_{j=1}^m|a_j|\|\chi_{E_j}-g_j\|_p
 <\eta.
\]
Suma $\sum_{j=1}^m a_jg_j$ jest więc żądaną aproksymacją. Jest ona ciągła, a jej nośnik zawiera się w skończonej sumie
zwartych nośników funkcji $g_j$, więc jest zwarty.
To dowodzi żądanej gęstości.
\end{proof}

\begin{uwaga}[Znaczenie założenia $p<\infty$]
Gęstość w normie $L^p$ nie rozciąga się bez dodatkowych
warunków na $p=\infty$. Na przykład względem miary
Lebesgue'a na $[0,1]$ funkcji $\chi_{[0,1/2]}$
nie da się aproksymować w normie $L^\infty$
funkcjami ciągłymi. Istotnie, gdyby dla ciągłej $g$
zachodziło $\|g-\chi_{[0,1/2]}\|_\infty<1/2$, wybralibyśmy
$r<1/2$ większe od tej normy. Wówczas $|g(x)-1|\le r$
na całym $(0,1/2)$ oraz $|g(x)|\le r$ na całym $(1/2,1)$:
naruszenie którejkolwiek nierówności w punkcie, przez ciągłość,
zachodziłoby na otwartym przedziale dodatniej miary.
Przechodząc do granicy w $1/2$ z obu stron, otrzymujemy
$1\le|1-g(1/2)|+|g(1/2)|\le2r<1$, sprzeczność.
\end{uwaga}

\begin{lemat}[Regularność miar zdominowanych]
\label{it:lem:radon-dominated}
Niech $X$ będzie lokalnie zwartą przestrzenią Hausdorffa.
Jeżeli $\sigma$ jest skończoną miarą Radona na $X$,
a $\tau$ nieujemną miarą borelowską z $\tau\le\sigma$,
to $\tau$ jest skończoną miarą Radona.
Skończona suma skończonych miar Radona jest również
skończoną miarą Radona.
\end{lemat}
\begin{proof}
Dla borelowskiego $E$ i $\varepsilon>0$ regularność
$\sigma$ i lemat \ref{lem:radon-sigmafinite-inner-20}
dają zwarty $K\subseteq E$ oraz otwarty $U\supseteq E$
z $\sigma(U\setminus K)<\varepsilon$ (wybieramy
błędy aproksymacji od wewnątrz i z zewnątrz
mniejsze niż $\varepsilon/2$). Wtedy
$\tau(U\setminus K)<\varepsilon$. Ponieważ
$\tau(X)\le\sigma(X)<\infty$, możemy odejmować
skończone wartości i otrzymujemy
$\tau(K)>\tau(E)-\varepsilon$ oraz
$\tau(U)<\tau(E)+\varepsilon$. Dowolność
$\varepsilon$ daje obie regularności dla $\tau$. Dla skończonej sumy $m$ miar wybieramy
dla każdej z nich analogiczną parę $K_i,U_i$ z
błędem mniejszym niż $\varepsilon/m$. Kładąc
$K=\bigcup_iK_i$, $U=\bigcap_iU_i$, dostajemy
zwarty podzbiór $E$, otwarty nadzbiór $E$ oraz
$U\setminus K\subseteq U_i\setminus K_i$ dla każdego $i$.
Zsumowanie błędów daje tezę.
\end{proof}

\section{Regularność a uzupełnienie miary}
Jeżeli $\mu$ jest miarą borelowską, jej uzupełnienie
(wprowadzone w rozdziale~2) może zawierać zbiory
nieborelowskie, ale nie zmienia wartości miary
na zbiorach borelowskich. Na przykład miara Lebesgue'a
na $\mathcal L_d$ jest uzupełnieniem miary borelowskiej
$\lambda_d|_{\mathcal B(\mathbb R^d)}$. Z drugiej strony
nie każda miara Radona jest zupełna na swojej
$\sigma$-algebrze borelowskiej: jeśli borelowski zbiór
zerowy ma podzbiór nieborelowski, to uzupełnienie
ma większą dziedzinę. Jest to rozróżnienie
istotne przy formułowaniu twierdzeń o obrazach miar,
całkowaniu i odsyłaczach do $L^p$.

\section{Co przygotowaliśmy na następny rozdział?}
Przeprowadziliśmy dwa kierunki regularności:
miary skończone na przestrzeniach metrycznych są
aproksymowane od wewnątrz przez zbiory domknięte,
a na przestrzeniach polskich przez zbiory zwarte;
miary lokalnie skończone na lokalnie zwartych
przestrzeniach metrycznych ośrodkowych są miarami Radona.
Ustaliliśmy ponadto, że wszystkie całki z funkcji
$C_c(X)$ jednoznacznie wyznaczają miarę Radona.
Następny rozdział odwróci tę konstrukcję:
zacznie od dodatniego funkcjonału na funkcjach ciągłych
i skonstruuje reprezentującą go miarę.

% ROZDZIAŁ 21 — RIESZ–MARKOW–KAKUTANI
\chapter{Twierdzenie Riesza--Markowa--Kakutaniego}
\label{ch:riesz-markov-21}

\providecommand{\atlasref}[2]{#2}


\section{Sformułowanie i rola twierdzenia}
Niech $X$ będzie lokalnie zwartą przestrzenią Hausdorffa,
a $\mathbb K\in\{\mathbb R,\mathbb C\}$.
Przestrzeń $C_c(X;\mathbb K)$ i pojęcie zwartego nośnika
wprowadzono w rozdziale~\ref{ch:radon-measures-20}.
Przez $C_0(X;\mathbb K)$ oznaczamy przestrzeń funkcji ciągłych
\emph{znikających w nieskończoności}, to znaczy takich, że
dla każdego $\varepsilon>0$ zbiór
$\{x:|f(x)|\ge\varepsilon\}$ jest zwarty.
Przestrzeń $C_0(X;\mathbb K)$ z normą supremum
jest przestrzenią Banacha. Najpierw zauważmy, że
$C_c(X;\mathbb K)\subseteq C_0(X;\mathbb K)$,
gdyż dodatnie zbiory poziomicowe modułu funkcji
o zwartym nośniku są domkniętymi podzbiorami tego
nośnika. Z nierówności trójkąta
$\{|f+g|\ge\varepsilon\}\subseteq
\{|f|\ge\varepsilon/2\}\cup\{|g|\ge\varepsilon/2\}$
wynika zwartość zbioru po lewej, gdyż jest domknięty.
Dla skalara $a\ne0$ mamy
$\{|af|\ge\varepsilon\}=\{|f|\ge\varepsilon/|a|\}$,
a dla $a=0$ zbiór jest pusty. Otrzymujemy przestrzeń
wektorową. Każda jej funkcja jest ograniczona:
jeśli zwarty zbiór $\{|f|\ge1\}$ jest niepusty,
moduł osiąga na nim skończone maksimum, a poza nim
jest mniejszy od $1$; jeśli jest pusty, moduł jest
wszędzie mniejszy od $1$. Dla pustego $X$
przyjmujemy normę jedynej funkcji równą zero. Dla ciągu Cauchy'ego $(f_n)$ w normie supremum
każdy ciąg $(f_n(x))$ jest ciągiem Cauchy'ego w
zupełnym ciele $\mathbb K$, ma więc granicę $f(x)$.
Zbieżność do $f$ jest jednostajna, ponieważ po ustaleniu $\varepsilon>0$ w nierówności $\|f_n-f_m\|_\infty\le\varepsilon$ dla $n,m$ dostatecznie dużych przechodzimy z $m\to\infty$. Granica $f$ jest ciągła. Jeżeli $\|f-f_n\|_\infty<\varepsilon/2$, to domknięty zbiór $\{|f|\ge\varepsilon\}$ zawiera się w zwartym zbiorze $\{|f_n|\ge\varepsilon/2\}$, jest więc zwarty. To daje zupełność.
Zupełność $C(K;\mathbb K)$ dla zwartej przestrzeni Hausdorffa $K$ omówiono w
\atlasref{FA-12::prop:ck-banach12}{FA-12: zupełność przestrzeni funkcji ciągłych}. Dla dowolnej lokalnie zwartej przestrzeni Hausdorffa udowodnimy gęstość
$C_c(X;\mathbb K)$ w $C_0(X;\mathbb K)$.

Konstrukcję miary z dodatniego funkcjonału całkowego
na kracie funkcji przedstawiono z pełnym dowodem w
\atlasref{ME-32::it:thm:daniell-stone}{ME-32: twierdzenie Daniella--Stone'a}
oraz w wersji bez jedynki w
\atlasref{ME-32::it:thm:daniell-stone-general}{ME-32: warunek Stone'a i ogólna reprezentacja}.
Ciągłość porządkowa oznacza tam malejącą zbieżność
punktową funkcji do zera, a nie tylko zerowe infimum
w kracie. W przypadku $C_c(X;\mathbb R)$ ciągłość ta
wynika z twierdzenia Diniego na zwartym nośniku
pierwszego wyrazu i oszacowania przez funkcję odcinającą.
Dla lokalnie zwartej ośrodkowej przestrzeni metrycznej
przeliczalne pokrycie przez zbiory niezerowości funkcji
z $C_c(X)$ wynika z $\sigma$-zwartości.
W \atlasref{ME-32::it:ex:daniell-cc}{ME-32: zastosowanie do funkcji o zwartym nośniku}
wyprowadzono w ten sposób reprezentację w przypadku metrycznym.
Poniżej podajemy konstrukcję ogólną, a następnie alternatywny
dowód metryczny. Porównanie dziedzin i wartości miar przedstawia
\atlasref{ME-32::it:thm:daniell-radon-comparison}{ME-32: porównanie miar Daniella i Radona}.

\begin{twierdzenie}[Reprezentacja dodatnich funkcjonałów]
\label{thm:riesz-positive-21}
Niech $X$ będzie \emph{lokalnie zwartą przestrzenią Hausdorffa}, a $L:C_c(X;\mathbb R)\to\mathbb R$ funkcjonałem
liniowym takim, że $L(f)\ge0$ dla $f\ge0$.
Istnieje dokładnie jedna miara Radona $\mu$ na $X$ spełniająca
\begin{equation}\label{eq:riesz-representation-21}
 L(f)=\int_X f(x)\,\mathrm d\mu(x),\qquad f\in C_c(X;\mathbb R).
\end{equation}
Przy czym $L$ nie musi być ograniczony w normie supremum na
całej przestrzeni $C_c(X)$.
\end{twierdzenie}

\begin{uwaga}[Zakres twierdzenia]
Wersja ogólna wykorzystuje konwencję Radona z definicji \ref{def:radon-measure-20}, czyli regularność wewnętrzną na otwartych. Nie wymagamy $\sigma$-skończoności miary ani metryczności przestrzeni. Dodatkowa regularność na wszystkich zbiorach borelowskich wynika z lematu \ref{lem:radon-sigmafinite-inner-20}, gdy $X$ jest $\sigma$-zwarta. W szczególności dotyczy to lokalnie zwartych ośrodkowych przestrzeni metrycznych. Poniżej podajemy także wcześniejszą alternatywną konstrukcję metryczną.
\end{uwaga}

\section{Dowód ogólny: lokalnie zwarte przestrzenie Hausdorffa}
\begin{lemat}[Lokalne ograniczenie funkcjonału]\label{it:lem:radon-local-bound}
Niech $X$ będzie lokalnie zwartą przestrzenią Hausdorffa
i niech $L:C_c(X;\mathbb R)\to\mathbb R$ będzie
dodatnim funkcjonałem liniowym.
Dla każdego zwartego $H\subseteq X$ istnieje $\eta\in C_c(X)$, $0\le\eta\le1$, $\eta|_H=1$, taka, że
\[
 |L(f)|\le L(\eta)\|f\|_\infty
 \quad\text{dla }f\in C_c(X),\quad\operatorname{supp}f\subset H.
\]
\end{lemat}
\begin{proof}
Lemat \ref{lem:cutoff-radon-20} zastosowany do $H\subset X$ daje $\eta$. Punktowo $-\|f\|_\infty\eta\le f\le\|f\|_\infty\eta$. Dodatniość i liniowość $L$ dają oszacowanie. Nie używamy tutaj istnienia miary.
\end{proof}

\begin{proof}[Dowód twierdzenia \ref{thm:riesz-positive-21} w ogólności]
Jeśli $X=\varnothing$, jedyną miarą jest miara zerowa i teza zachodzi. W dalszej części nie zakładamy, że $X$ ma przeliczalne pokrycie zbiorami zwartymi.

\emph{Krok 1: zawartość zbiorów otwartych.}
Dla otwartego $U\subset X$ definiujemy
\begin{equation}\label{it:eq:radon-open-content-general}
 a(U)=\sup\{L(f): f\in C_c(X),\ 0\le f\le1,\
 \operatorname{supp}f\subset U\}.
\end{equation}
Funkcja zerowa jest zawsze dopuszczalna, więc $a$ przyjmuje wartości w $[0,\infty]$, jest monotoniczna oraz $a(\varnothing)=0$. Dodatniość i liniowość $L$ implikują monotoniczność na $C_c(X)$, ponieważ dla $f\le g$ mamy $L(g)-L(f)=L(g-f)\ge0$. Jeśli $\overline U$ jest zwarty, z lematu \ref{it:lem:radon-local-bound} mamy $a(U)\le L(\eta)<\infty$. Dla rozłącznych otwartych $U,V$ i funkcji dopuszczalnych $f,g$ odpowiednio dla $U,V$ mamy $0\le f+g\le1$ i $\operatorname{supp}(f+g)\subset U\cup V$. Zatem $a(U\cup V)\ge L(f)+L(g)$. Biorąc niezależne suprema, otrzymujemy $a(U\cup V)\ge a(U)+a(V)$, również jeśli jeden ze składników jest nieskończony.

\emph{Krok 2: podaddytywność na otwartych.}
Jeśli $U\subset\bigcup_{j\ge1}U_j$ dla otwartych $U_j$, ustalmy funkcję dopuszczalną $f$ dla $U$. Jeśli $f=0$, oszacowanie jest natychmiastowe. W przeciwnym przypadku dla każdego $x$ ze zwartego $H=\operatorname{supp}f$ wybieramy $j$ i otwarte otoczenie $V_x$ o zwartym domknięciu zawartym w $U_j$. Istnienie wynika z regularności i lokalnej zwartości, jak w dowodzie lematu \ref{lem:cutoff-radon-20}. Ze zwartości $H$ otrzymujemy skończone pokrycie $V_1,\ldots,V_m$ z $\overline V_i\subset U_{j_i}$. Lemat odcinający daje $\psi_i\in C_c(X)$, $0\le\psi_i\le1$, $\psi_i=1$ na $\overline V_i$ i $\operatorname{supp}\psi_i\subset U_{j_i}$. Wówczas $\psi=\sum_i\psi_i$ jest dodatnia na $H$. Definiujemy
\[
 f_i(x)=
 \begin{cases}
 f(x)\psi_i(x)/\psi(x),&\psi(x)>0,\\
 0,&\psi(x)=0.
 \end{cases}
\]
Na otwartym zbiorze $\{\psi>0\}$ funkcja $f_i$ jest ciągła. W każdym punkcie $x$ z $\psi(x)=0$ mamy $x\notin H$, a zatem $f=0$ w pewnym otoczeniu $x$, co daje ciągłość także tam. Ponadto $0\le f_i\le f$ i zbiór punktów, w których $f_i\ne0$, zawiera się w przecięciu zbiorów, w których $f\ne0$ i $\psi_i\ne0$. Po wzięciu domknięć otrzymujemy $\operatorname{supp}f_i\subseteq H\cap\operatorname{supp}\psi_i\subseteq U_{j_i}$. Ponieważ $H$ jest zwarty, każdy $f_i$ ma zwarty nośnik. Na $\{\psi>0\}$ suma ilorazów wynosi $f$, a na $\{\psi=0\}$ mamy $f=0$; stąd $\sum_i f_i=f$ na całym $X$. Grupujemy powtarzające się indeksy: dla każdego występującego $j$ kładziemy $g_j=\sum_{i:j_i=j}f_i$. Wówczas $0\le g_j\le f\le1$ i $\operatorname{supp}g_j\subset U_j$. Z liniowości $L(f)=\sum_jL(g_j)\le\sum_{j\ge1}a(U_j)$, gdzie lewa suma jest skończona. Biorąc supremum po $f$, otrzymujemy $a(U)\le\sum_ja(U_j)$. W połączeniu z krokiem 1 daje to addytywność na dwóch rozłącznych zbiorach otwartych.

\emph{Krok 3: miara zewnętrzna i zwarta aproksymacja przed mierzalnością.}
Określamy dla dowolnego $E\subset X$
\begin{equation}\label{it:eq:radon-outer-general}
 m^*(E)=\inf\{a(U):E\subset U,\ U\text{ otwarty}\}.
\end{equation}
Mamy $m^*(\varnothing)=0$ oraz monotoniczność. Jeśli $E\subset\bigcup_jE_j$ i $\sum_jm^*(E_j)<\infty$, dla $\varepsilon>0$ wybieramy otwarte $U_j\supset E_j$ z $a(U_j)<m^*(E_j)+\varepsilon2^{-j}$. Wtedy
$m^*(E)\le a(\bigcup_jU_j)\le\sum_jm^*(E_j)+\varepsilon$.
Przechodząc z $\varepsilon\downarrow0$, otrzymujemy przeliczalną podaddytywność. Dla nieskończonej sumy podaddytywność jest automatyczna. Stąd $m^*$ jest miarą zewnętrzną. Dla otwartego $U$ zachodzi $m^*(U)=a(U)$: sam zbiór $U$ jest dopuszczalny w infimum, a każdy jego otwarty nadzbiór $V$ spełnia $a(V)\ge a(U)$. Każdy zwarty $H$ ma otwarte otoczenie o zwartym domknięciu (skończona suma takich otoczeń punktów $H$), więc $m^*(H)<\infty$.

Jeżeli $f$ jest dopuszczalna dla $U$, to dla $H=\operatorname{supp}f\subset U$ i każdego otwartego $V\supset H$ zachodzi $a(V)\ge L(f)$. Stąd $m^*(H)\ge L(f)$. Biorąc supremum po $f$, dostajemy
\begin{equation}\label{it:eq:radon-outer-inner-general}
 m^*(U)=a(U)=\sup\{m^*(H):H\subset U,\ H\text{ zwarty}\}.
\end{equation}
Odwrotna nierówność wynika z monotoniczności $m^*$; argument działa także dla nieskończonego $a(U)$. Na tym etapie nie zakładamy mierzalności $H$ ani $U$.

\emph{Krok 4: kryterium Carathéodory'ego dla otwartych.}
Ustalmy otwarty $U\subset X$ oraz dowolny $E\subset X$. Jeśli $m^*(E)=+\infty$, podaddytywność wymusza, że suma $m^*(E\cap U)+m^*(E\setminus U)$ jest nieskończona, więc kryterium zachodzi. Niech $m^*(E)<\infty$ i $\varepsilon>0$. Wybieramy otwarty $W\supset E$ z $a(W)<m^*(E)+\varepsilon$. Mamy $a(W\cap U)\le a(W)<\infty$. Z \eqref{it:eq:radon-outer-inner-general} wybieramy zwarty $H\subset W\cap U$ z
$m^*(H)>a(W\cap U)-\varepsilon$ (jeśli $a(W\cap U)=0$, wystarczy $H=\varnothing$). Z lematu \atlasref{GT-14::lem:topog-baire-zwarte-domkniecie}{GT-14: otoczenie o zwartym domknięciu} i skończonego podpokrycia zwartego $H$ wybieramy otwarty $V$ taki, że $H\subseteq V\subseteq\overline V\subseteq W\cap U$ i $\overline V$ jest zwarty; dla pustego $H$ można przyjąć $V=\varnothing$. Wówczas $V$ i $W\setminus\overline V$ są rozłącznymi otwartymi podzbiorami $W$, a $E\setminus U\subset W\setminus\overline V$. Z kroku 1 i monotoniczności mamy
\begin{align*}
 m^*(E)+\varepsilon
 &>a(W)\ge a(V)+a(W\setminus\overline V)\\
 &\ge m^*(H)+m^*(E\setminus U)\\
 &>a(W\cap U)-\varepsilon+m^*(E\setminus U)\\
 &\ge m^*(E\cap U)+m^*(E\setminus U)-\varepsilon .
\end{align*}
Zatem $m^*(E\cap U)+m^*(E\setminus U)<m^*(E)+2\varepsilon$. Z dowolności $\varepsilon$ otrzymujemy nierówność potrzebną w kryterium Carathéodory'ego; przeciwna jest podaddytywnością. Wszystkie zbiory otwarte są więc mierzalne względem $m^*$. Twierdzenie \ref{thm:caratheodory} daje mierzalność wszystkich zbiorów borelowskich i przeliczalną addytywność $\mu=m^*|_{\mathcal B(X)}$. Z kroków 3--4 wynika skończoność na zwartych, regularność zewnętrzna i regularność wewnętrzna na otwartych. Zatem $\mu$ jest miarą Radona według definicji \ref{def:radon-measure-20}.

\emph{Krok 5: identyfikacja funkcjonału z całką.}
Niech $f\in C_c(X)$, $f\ge0$, $f\ne0$. Oznaczmy $H=\operatorname{supp}f$, $M=\|f\|_\infty>0$, wybierzmy $\eta$ z lematu \ref{it:lem:radon-local-bound}, a dla $N\ge1$ połóżmy $\delta=M/N$ i $U_i=\{f>i\delta\}$, $0\le i\le N$. Wszystkie $U_i$ są otwarte, zawarte w $H$, zatem mają skończoną miarę. Dla $0\le i<N$ zbiór $H_i=\{f\ge(i+1)\delta\}$ jest zwartym podzbiorem $U_i$. Z lematu \ref{lem:cutoff-radon-20} wybieramy $h_i\in C_c(X)$, $0\le h_i\le1$, $h_i|_{H_i}=1$ i $\operatorname{supp}h_i\subset U_i$. Punktowo $f\le\delta\eta+\delta\sum_{i=0}^{N-1}h_i$: dla $x\in H$ i $q=\lfloor f(x)/\delta\rfloor$ mamy $0\le q\le N$, $\eta(x)=1$ oraz $h_i(x)=1$ dla $0\le i<q$. Zatem prawa strona jest nie mniejsza od $(q+1)\delta\ge f(x)$. Poza $H$ lewa strona jest zerowa, a prawa nieujemna. Dodatniość daje
\[
 L(f)\le\delta L(\eta)+\delta\sum_{i=0}^{N-1}\mu(U_i).
\]
Dla dowolnych dopuszczalnych funkcji $g_i$ dla $U_i$, $1\le i\le N$, mamy $\delta\sum_i g_i\le f$, ponieważ liczba progów $i\delta<f(x)$ nie przekracza $f(x)/\delta$. Z dodatniości i wzięcia niezależnych supremów (skończona liczba skończonych supremów) wynika $\delta\sum_{i=1}^N\mu(U_i)\le L(f)$. Dla $0\le y\le M=N\delta$ mamy
$\delta\sum_{i=1}^{N}\chi_{(i\delta,\infty)}(y)\le y
\le\delta\sum_{i=0}^{N-1}\chi_{(i\delta,\infty)}(y)$:
przy $y=0$ obie sumy są zerowe; dla
$j\delta<y\le(j+1)\delta$, $0\le j<N$,
dolna suma wynosi $j\delta$, a górna $(j+1)\delta$.
Stosując te nierówności do $y=f(x)$ i całkując
mierzalne funkcje proste, otrzymujemy
\[
 \delta\sum_{i=1}^{N}\mu(U_i)
 \le\int_X f\,\mathrm d\mu
 \le\delta\sum_{i=0}^{N-1}\mu(U_i).
\]
Różnica sum wynosi $\delta\mu(U_0)\le\delta\mu(H)$, ponieważ $U_N=\varnothing$. Stąd
\[
 \left|L(f)-\int_X f\,\mathrm d\mu\right|
 \le\delta\bigl(L(\eta)+\mu(H)\bigr)\longrightarrow0.
\]
Zatem zachodzi reprezentacja dla $f\ge0$ (dla $f=0$ obie strony są zerowe). Dla dowolnej rzeczywistej $f\in C_c(X)$ stosujemy wynik do $f^+$ i $f^-$, które również mają zwarty nośnik, i odejmujemy skończone całki. Jednoznaczność wynika z twierdzenia \ref{thm:uniqueness-cc-20}. 
\end{proof}



\begin{wniosek}[Odzyskiwanie miary zbioru zwartego z funkcjonału]
\label{it:cor:radon-compact-inf}
Niech $L$ i $\mu$ spełniają tezę twierdzenia
\ref{thm:riesz-positive-21}. Dla każdego zwartego
$K\subseteq X$ zachodzi
\[
 \mu(K)=\inf\{L(f):f\in C_c(X),\
             0\le f\le1,\ f|_K=1\}.
\]
Wraz ze wzorem \eqref{eq:cc-open-sup-20} daje to
dwie formuły odzyskiwania miary bez znajomości
jej wartości na dowolnych zbiorach borelowskich.
\end{wniosek}
\begin{proof}
Każda dopuszczalna funkcja $f$ spełnia
$\chi_K\le f$. Z reprezentacji i monotoniczności
całki otrzymujemy $\mu(K)\le L(f)$.
Dla $\varepsilon>0$ regularność zewnętrzna
i skończoność $\mu(K)$ dają otwarty $U\supseteq K$
z $\mu(U)<\mu(K)+\varepsilon$.
Lemat \ref{lem:cutoff-radon-20} daje dopuszczalną
funkcję $f$ o nośniku zawartym w $U$.
Mamy $f\le\chi_U$, zatem
$L(f)=\int_X f\,\mathrm d\mu\le\mu(U)
<\mu(K)+\varepsilon$.
Dowolność $\varepsilon$ dowodzi równości.
Dla $K=\varnothing$ można też bezpośrednio
przyjąć $f=0$.
\end{proof}

\section{Lemat o metrycznych miarach zewnętrznych}
\begin{definicja}[Miara zewnętrzna metryczna]
\label{def:metric-outer-21}
Niech $(Z,d)$ będzie przestrzenią metryczną.
Miara zewnętrzna $m^*:\mathcal P(Z)\to[0,\infty]$
jest \emph{metryczna}, gdy dla dowolnych $A,B\subseteq Z$
o dodatniej odległości $d(A,B)>0$ zachodzi
$m^*(A\cup B)=m^*(A)+m^*(B)$.
Dla niepustych zbiorów przyjmujemy
$d(A,B)=\inf\{d(x,y):x\in A,\ y\in B\}$,
a jeśli co najmniej jeden ze zbiorów jest pusty,
przyjmujemy $d(A,B)=+\infty$.
\end{definicja}

\begin{lemat}[Borelowska mierzalność]
\label{lem:metric-outer-borel-21}
Jeżeli $m^*$ jest metryczną miarą zewnętrzną na $(Z,d)$,
to każdy zbiór borelowski w $Z$ jest mierzalny w sensie
Carathéodory'ego.
\end{lemat}
\begin{proof}
Wystarczy wykazać mierzalność zbiorów domkniętych, gdyż
rodzina zbiorów mierzalnych Carathéodory'ego jest
$\sigma$-algebrą (twierdzenie \atlasref{ME-03::thm:caratheodory}{ME-03: twierdzenie Carathéodory'ego}).
Niech $F\subseteq Z$ będzie domknięty i $E\subseteq Z$.
Musimy dowieść
$m^*(E)\ge m^*(E\cap F)+m^*(E\setminus F)$;
nierówność odwrotna wynika z podaddytywności.
Jeśli $F=\varnothing$ lub $F=Z$, mierzalność jest oczywista.
Jeśli $m^*(E)=+\infty$, podaddytywność wymusza nieskończoną sumę po prawej, więc kryterium zachodzi. Załóżmy
$m^*(E)<\infty$ i oznaczmy $r(x)=d(x,F)$.
Dla każdego $n\ge1$ zbiory $E\cap F$ oraz
$E\cap\{r>1/n\}$ mają dodatnią odległość, gdyż
$r$ jest $1$-lipschitzowska. Zatem
\begin{equation}\label{eq:metric-outer-sep-21}
 m^*(E)\ge m^*(E\cap F)+m^*(E\cap\{r>1/n\}).
\end{equation}
Wykażemy, że drugi składnik po prawej zbiega do
$m^*(E\setminus F)$; nie wolno tutaj korzystać z ciągłości
z dołu, ponieważ $m^*$ jest dopiero miarą zewnętrzną.
Niech $A_0=E\cap\{r>1\}$ oraz
$A_j=E\cap\{1/(j+1)<r\le1/j\}$, $j\ge1$.
Dla $j\ge1$ i $k\ge j+2$, jeśli $x\in A_j$,
$y\in A_k$, to z $1$-lipschitzowskości $r$ wynika
\[
 d(x,y)\ge r(x)-r(y)>
 \frac1{j+1}-\frac1{j+2}>0.
\]
Zatem $d(A_j,A_k)\ge1/(j+1)-1/(j+2)>0$.
Także $d(A_0,A_k)\ge1/2$ dla $k\ge2$.
Dla skończonej rodziny warstw tej samej parzystości
odległość każdej warstwy od sumy pozostałych
jest nie mniejsza od minimum skończenie wielu
dodatnich odległości (puste warstwy można pominąć).
Metryczność i indukcja dają więc
$\sum m^*(A_j)=m^*(\bigcup A_j)\le m^*(E)$.
W konsekwencji oba szeregi, po indeksach parzystych
i nieparzystych, są zbieżne; zatem
$\sum_{j\ge1}m^*(A_j)<\infty$. Ponieważ
\[
 E\cap\{0<r\le1/n\}\subseteq\bigcup_{j\ge n}A_j,
\]
miara zewnętrzna części po lewej dąży do zera.
Z podaddytywności
\[
 m^*(E\setminus F)\le
 m^*(E\cap\{r>1/n\})+
 m^*(E\cap\{0<r\le1/n\}),
\]
a z monotoniczności wynika przeciwna nierówność
$m^*(E\cap\{r>1/n\})\le m^*(E\setminus F)$.
Przechodząc do granicy w \eqref{eq:metric-outer-sep-21},
otrzymujemy żądane kryterium Carathéodory'ego.
\end{proof}

\section{Przypadek zwarty: konstrukcja od podstaw}

\noindent\emph{Powiązania tematyczne w Bibliotece matematycznej:} \atlasref{FA-12::thm:riesz-markov12}{FA-12: reprezentacja funkcjonałów na przestrzeni funkcji ciągłych}.

\begin{twierdzenie}[Reprezentacja na zwartej przestrzeni metrycznej]
\label{thm:compact-riesz-21}
Niech $K$ będzie niepustą zwartą przestrzenią metryczną,
a $L:C(K;\mathbb R)\to\mathbb R$ dodatnim funkcjonałem
liniowym. Istnieje dokładnie jedna skończona regularna
miara borelowska $\mu$ na $K$ taka, że
$L(f)=\int_Kf\,\mathrm d\mu$ dla wszystkich $f\in C(K)$.
Ponadto $\mu(K)=L(\chi_K)$ i
$\|L\|=L(\chi_K)$.
\end{twierdzenie}
\begin{proof}
\emph{Krok 1: lokalne ograniczenie funkcjonału.}
Z $-\|f\|_\infty\chi_K\le f\le\|f\|_\infty\chi_K$
wynika $|L(f)|\le L(\chi_K)\|f\|_\infty$.
Dla $f=\chi_K$ mamy równość w normie.

\emph{Krok 2: przypisanie liczb zbiorom otwartym.}
Dla zbioru otwartego $U\subseteq K$ definiujemy
\begin{equation}\label{eq:open-content-21}
 a(U)=\sup\{L(f): f\in C(K),\ 0\le f\le1,\
          \operatorname{supp}f\subseteq U\}.
\end{equation}
W szczególności $a(\varnothing)=0$,
$0\le a(U)\le L(\chi_K)$, $a(K)=L(\chi_K)$,
a $a$ jest monotoniczna. Jeżeli $U,V$ są rozłączne
i otwarte, to $a(U\cup V)\ge a(U)+a(V)$:
wybieramy funkcje $f,g$ dopuszczalne odpowiednio dla
$U,V$; ponieważ ich nośniki są rozłączne,
$0\le f+g\le1$ i $L(f+g)=L(f)+L(g)$.

\emph{Krok 3: przeliczalna podaddytywność na otwartych.}
Jeżeli $U\subseteq\bigcup_{j\ge1}U_j$ i $f$ jest
dopuszczalna w \eqref{eq:open-content-21}, przypadek
$f=0$ daje od razu $L(f)=0\le\sum_j a(U_j)$.
W przeciwnym razie $\operatorname{supp}f$ ma skończone pokrycie
$U_{j_1},\ldots,U_{j_N}$. Zbudujmy ciągłe funkcje
$\theta_i\ge0$ o nośnikach zawartych w $U_{j_i}$,
takie, że $\sum_i\theta_i=1$ na
$H:=\operatorname{supp}f$. Dla każdego $x\in H$
wybieramy otwarty zbiór $V_x$ zawierający $x$,
którego domknięcie w $K$ mieści się w pewnym $U_{j_i}$.
Ze zwartości $H$ wybieramy skończone pokrycie
takimi $V_x$; po ewentualnym powtórzeniu indeksów
oznaczmy je przez $V_i$, tak aby
$\overline V_i\subseteq U_{j_i}$.
Lemat~\ref{lem:cutoff-radon-20} daje funkcje
$\psi_i\in C(K)$ spełniające
$0\le\psi_i\le1$, $\psi_i=1$ na $\overline V_i$ i
$\operatorname{supp}\psi_i\subseteq U_{j_i}$.
Dla $\psi=\sum_i\psi_i$ mamy $\psi>0$ na $H$.
Ponownie stosujemy lemat o funkcji odcinającej,
uzyskując $\chi\in C(K)$ taką, że $0\le\chi\le1$,
$\chi=1$ na $H$ oraz
$\operatorname{supp}\chi\subseteq\{\psi>0\}$.
Wówczas $\theta_i=\chi\psi_i/\psi$ na $\{\psi>0\}$
i $\theta_i=0$ poza tym zbiorem. Funkcje
$\theta_i$ są ciągłe również na brzegu $\{\psi>0\}$,
ponieważ nośnik $\chi$ jest w nim zawarty,
a ich suma jest równa $\chi=1$ na $H$. Każde $f\theta_i$ jest nieujemne,
nie przekracza $1$ i ma nośnik w $U_{j_i}$.
Jeżeli indeks $j_i$ występuje więcej niż raz, grupujemy odpowiadające mu funkcje: $h_j=\sum_{i:j_i=j}f\theta_i$. Mamy $0\le h_j\le f\le1$ (ponieważ suma wszystkich składników jest równa $f$), a jego nośnik zawiera się w $U_j$, jako skończona suma odpowiednich nośników. Wobec tego
$L(f)=\sum_j L(h_j)\le\sum_{j\ge1}a(U_j)$.
Po wzięciu supremum otrzymujemy
$a(U)\le\sum_j a(U_j)$.

\emph{Krok 4: miara zewnętrzna.}
Dla dowolnego $A\subseteq K$ definiujemy
\begin{equation}\label{eq:riesz-outer-21}
 m^*(A)=\inf\{a(U):A\subseteq U,\ U\text{ otwarty}\}.
\end{equation}
Jeżeli $A\subseteq\bigcup_j A_j$, dobieramy otwarte
$U_j\supseteq A_j$ z
$a(U_j)\le m^*(A_j)+\varepsilon2^{-j}$ (przyjmujemy
osobno trywialny przypadek nieskończonej sumy).
Krok 3 daje
$m^*(A)\le a(\bigcup_jU_j)\le\sum_jm^*(A_j)+\varepsilon$.
Przechodząc z $\varepsilon\downarrow0$ otrzymujemy
przeliczalną podaddytywność. Ponadto
$m^*(\varnothing)=0$ z $a(\varnothing)=0$,
a monotoniczność wynika z zawierania się rodzin
dopuszczalnych otwartych nadzbiorów. Zatem $m^*$
jest miarą zewnętrzną.
Dla zbioru otwartego $U$ mamy $m^*(U)=a(U)$, ponieważ
$U\subseteq V$ implikuje $a(U)\le a(V)$.
Jeśli $A$ lub $B$ jest pusty, metryczna addytywność zachodzi wprost. W przeciwnym razie, jeśli $d(A,B)=\delta>0$, to dla każdego otwartego
$W\supseteq A\cup B$ zbiory
$W_A=W\cap\{x:d(x,A)<\delta/3\}$ i
$W_B=W\cap\{x:d(x,B)<\delta/3\}$ są otwarte,
rozłączne i zawierają odpowiednio $A,B$. Stąd
$a(W)\ge a(W_A)+a(W_B)\ge m^*(A)+m^*(B)$.
Rozłączność $W_A,W_B$ wynika z nierówności
$d(A,B)\le d(x,A)+d(x,B)$, która zachodzi dla
każdego $x$ na mocy nierówności trójkąta i wzięcia
infimum po punktach obu zbiorów. Gdyby $x$ należał
do obu otoczeń, otrzymalibyśmy $\delta<2\delta/3$.
Biorąc infimum po $W$ otrzymujemy
$m^*(A\cup B)\ge m^*(A)+m^*(B)$; nierówność
przeciwna wynika z podaddytywności. To dowodzi
metryczności $m^*$.
Lemat~\ref{lem:metric-outer-borel-21} zapewnia, że
wszystkie zbiory borelowskie są mierzalne. Ograniczenie
$\mu=m^*|_{\mathcal B(K)}$ jest więc miarą, a
$\mu(K)=L(\chi_K)<\infty$.

\emph{Krok 5: regularność.}
Regularność zewnętrzna wynika z \eqref{eq:riesz-outer-21}.
Dla otwartego $U$ i $\varepsilon>0$ dobieramy
$0\le f\le1$ z $\operatorname{supp}f\subseteq U$ i
$L(f)>\mu(U)-\varepsilon$. Przyjmijmy
$H=\operatorname{supp}f$, który jest zwarty.
Każdy otwarty $V\supseteq H$ spełnia
$L(f)\le a(V)$, zatem
$\mu(H)=\inf_{V\supseteq H}a(V)\ge L(f)$.
To dowodzi regularności wewnętrznej dla otwartych.
Dla dowolnego borelowskiego $E$ i $\varepsilon>0$
wybieramy otwarty $U\supseteq E$ z
$\mu(U\setminus E)<\varepsilon$.
Następnie wybieramy zwarty $H\subseteq U$ z
$\mu(U\setminus H)<\varepsilon$ oraz otwarty
$W\supseteq U\setminus E$ z
$\mu(W)<\mu(U\setminus E)+\varepsilon$.
Zbiór $H\setminus W$ jest zwarty, zawiera się w $E$,
a $\mu(E\setminus(H\setminus W))<3\varepsilon$,
ponieważ $E\setminus(H\setminus W)\subseteq(U\setminus H)\cup W$
i $\mu(W)<2\varepsilon$.
Otrzymujemy regularność wewnętrzną względem zwartych.

\emph{Krok 6: identyfikacja całki z $L$.}
Niech $f\in C(K)$, $f\ge0$. Dla $f=0$ równość jest natychmiastowa; w pozostałym przypadku $M=\|f\|_\infty>0$.
Dla $N\ge1$ połóżmy $\delta=M/N$ i
$U_i=\{f>i\delta\}$, $i=0,\ldots,N$.
Dla $i=0,\ldots,N-1$ zbiór
$H_i=\{f\ge(i+1)\delta\}$ jest zwarty i zawiera się
w $U_i$. Istnieje $h_i\in C(K)$, $0\le h_i\le1$,
$h_i=1$ na $H_i$, $\operatorname{supp}h_i\subseteq U_i$;
wynika to z lematu o funkcji odcinającej
\ref{lem:cutoff-radon-20}. Punktowo
$f\le\delta\chi_K+\delta\sum_{i=0}^{N-1}h_i$.
Z dodatniości $L$ dostajemy
\begin{equation}\label{eq:riesz-upper-sum-21}
 L(f)\le\delta L(\chi_K)+
           \delta\sum_{i=0}^{N-1}\mu(U_i).
\end{equation}
Z drugiej strony dla dowolnych funkcji
$g_i\in C(K)$, $0\le g_i\le1$,
$\operatorname{supp}g_i\subseteq U_i$ ($1\le i\le N$),
zachodzi $\delta\sum_{i=1}^{N}g_i\le f$ punktowo.
Zatem, biorąc niezależnie supremum po $g_i$,
\begin{equation}\label{eq:riesz-lower-sum-21}
 \delta\sum_{i=1}^{N}\mu(U_i)\le L(f).
\end{equation}
Natomiast dla $y=f(x)$ równego zero obie sumy
progowe są zerowe. Jeżeli $j\delta<y\le(j+1)\delta$
dla $0\le j<N$, to dokładnie $j$ progów spośród
$\delta,2\delta,\ldots,N\delta$ i dokładnie $j+1$
progów spośród $0,\delta,\ldots,(N-1)\delta$
jest mniejszych od $y$. Otrzymujemy punktowe
oszacowanie przez dolną i górną sumę progową;
po całkowaniu względem $\mu$ daje ono
\[
 \delta\sum_{i=1}^{N}\mu(U_i)
 \le\int_Kf(x)\,\mathrm d\mu(x)
 \le\delta\sum_{i=0}^{N-1}\mu(U_i).
\]
Różnica między lewą i prawą sumą wynosi
$\delta\mu(U_0)\le\delta L(\chi_K)$.
Z \eqref{eq:riesz-upper-sum-21} i
\eqref{eq:riesz-lower-sum-21} wynika
$|L(f)-\int_K f\,\mathrm d\mu|\le2\delta L(\chi_K)$.
Przechodząc do granicy przy $N\to\infty$, otrzymujemy reprezentację
$f\ge0$. Ogólny $f$ rozkładamy na części dodatnią
i ujemną.

\emph{Krok 7: jednoznaczność.}
Dla dowolnej miary regularnej $\rho$ reprezentującej $L$
i otwartego $U$ regularność oraz funkcje odcinające
dają $\rho(U)=a(U)$ (por.
\eqref{eq:cc-open-sup-20}). Stąd z regularności
zewnętrznej $\rho=\mu$ na wszystkich zbiorach
borelowskich.
\end{proof}

\section{Przestrzenie lokalnie zwarte: lokalizacja i sklejanie}
\begin{lemat}[Wyczerpanie przez zbiory relatywnie zwarte]
\label{lem:exhaustion-21}
Niech $X$ będzie lokalnie zwartą, $\sigma$-zwartą przestrzenią Hausdorffa. Istnieją otwarte zbiory $U_n$ takie, że
$U_n\subseteq\overline U_n\subseteq U_{n+1}$,
$\overline U_n$ jest zwarty i $\bigcup_nU_n=X$.
Ponadto istnieją $\chi_n\in C_c(X)$ spełniające
$0\le\chi_n\le1$, $\chi_n=1$ na $\overline U_n$
oraz $\operatorname{supp}\chi_n\subseteq U_{n+1}$.
\end{lemat}
\begin{proof}
Z $\sigma$-zwartości wybieramy zwarte zbiory $D_n$ pokrywające $X$. Kładziemy $U_0=\varnothing$. Jeśli $U_{n-1}$ ma zwarte domknięcie, to $H_n=\overline U_{n-1}\cup D_n$ jest zwarty. Wybieramy skończone pokrycie $H_n$ otwartymi otoczeniami o zwartych domknięciach i za $U_n$ przyjmujemy ich sumę (dla pustego $H_n$ przyjmujemy $U_n=\varnothing$). Domknięcie skończonej sumy jest sumą domknięć, więc $\overline U_n$ jest zwarte. Ponadto $\overline U_{n-1}\subset U_n$ i $D_n\subset U_n$, skąd $\bigcup_n U_n=X$. W szczególności każda lokalnie zwarta ośrodkowa przestrzeń metryczna jest $\sigma$-zwarta: własność Lindelöfa (por. \atlasref{GT-09::thm:topog-metryczne-trzy-warunki}{GT-09: własność Lindelöfa przestrzeni metrycznej ośrodkowej}) daje przeliczalne pokrycie relatywnie zwartymi otoczeniami, a ich domknięcia tworzą pokrycie zwarte.
Funkcje $\chi_n$ istnieją na mocy lematu
\ref{lem:cutoff-radon-20}.
\end{proof}

\begin{proof}[Alternatywny dowód twierdzenia~\ref{thm:riesz-positive-21} w przypadku metrycznym]
Zakładamy w tym dowodzie, że $X$ jest lokalnie zwartą ośrodkową przestrzenią metryczną.
Jeśli $X=\varnothing$, teza zachodzi dla miary zerowej.
\emph{Krok 1: dodatniość lokalna i zwartość.}
Dla dowolnego zwartego $H\subseteq X$ wybieramy
$\eta\in C_c(X)$, $\eta\ge0$, $\eta=1$ na $H$.
Jeśli $\operatorname{supp}f\subseteq H$, to
$-\|f\|_\infty\eta\le f\le\|f\|_\infty\eta$,
więc $|L(f)|\le\|f\|_\infty L(\eta)$.
Wybierzmy wyczerpanie i funkcje $\chi_n$ z lematu
\ref{lem:exhaustion-21}; oznaczmy $K_n=\overline U_n$
i $J_n=\overline U_{n+1}$.
Dla $h\in C(J_n)$ funkcję $\chi_nh$ na $J_n$
przedłużamy przez zero na $X\setminus J_n$.
Jest ona ciągła, ponieważ
$\operatorname{supp}\chi_n\subseteq U_{n+1}$.
Definiujemy $L_n(h)=L(\chi_nh)$; jest to dodatni
funkcjonał na $C(J_n)$.
Jeżeli $J_n=\varnothing$, przyjmujemy miarę zerową.
W przeciwnym razie z twierdzenia~\ref{thm:compact-riesz-21}
istnieje skończona regularna miara $\rho_n$
na $J_n$ reprezentująca $L_n$. Przedłużamy ją
przez zero na dopełnienie $J_n$, otrzymując miarę
borelowską na $X$ (także oznaczaną $\rho_n$).

\emph{Krok 2: zgodność miar na wyczerpaniu.}
Jeżeli $m\ge n$, to $\chi_m=\chi_n=1$ na $U_n$.
Dla dowolnego $f\in C_c(U_n)$ zachodzi zatem
\[
 \int_Xf\,\mathrm d\rho_m=L(f)
  =\int_Xf\,\mathrm d\rho_n.
\]
Przedłużone miary $\rho_j$ są skończonymi miarami
borelowskimi na $X$, więc z twierdzenia
\ref{thm:lch-secondcountable-radon-20} są miarami Radona.
Ich ograniczenia do $U_n$ są skończonymi miarami
borelowskimi na lokalnie zwartej ośrodkowej przestrzeni
metrycznej $U_n$, więc z tego samego twierdzenia są
miarami Radona.
Jednoznaczność z twierdzenia~\ref{thm:uniqueness-cc-20}
daje $\rho_m|_{U_n}=\rho_n|_{U_n}$.


\emph{Krok 3: sklejenie.}
Stawiamy $U_0=\varnothing$, $B_n=U_n\setminus U_{n-1}$.
Dla borelowskiego $E\subseteq X$ określamy
\begin{equation}\label{eq:gluing-measures-21}
 \mu(E)=\sum_{n\ge1}\rho_n(E\cap B_n).
\end{equation}
Ponieważ $B_n$ są parami rozłączne, a suma miar jest
miarą, $\mu$ jest miarą borelowską. Dla $E\subseteq U_m$
przeliczalna addytywność i krok 2 dają
\[
 \mu(E)=\sum_{n=1}^m\rho_n(E\cap B_n)
 =\sum_{n=1}^m\rho_m(E\cap B_n)=\rho_m(E),
\]
ponieważ $B_n\subseteq U_n$ i $E\cap B_n=\varnothing$
dla $n>m$. Dotyczy to wszystkich $E\in\mathcal B(U_m)$,
które są borelowskie również w $X$, gdyż $U_m$ jest otwarty. W szczególności
$\mu(U_m)=\rho_m(U_m)<\infty$.
Zatem $\mu$ jest lokalnie skończona. Rozdział
\ref{ch:radon-measures-20} pokazuje, że na lokalnie
zwartej ośrodkowej przestrzeni metrycznej jest ona
miarą Radona.

\emph{Krok 4: reprezentacja.}
Jeśli $f\in C_c(X)$, jej nośnik mieści się w pewnym
$U_m$ (ponieważ $(U_m)$ jest rosnącym pokryciem
zwartego nośnika). Wtedy
\[
 \int_X f(x)\,\mathrm d\mu(x)
 =\int_X f(x)\,\mathrm d\rho_m(x)
 =L(\chi_m f)=L(f).
\]
Jednoznaczność wynika z twierdzenia
\ref{thm:uniqueness-cc-20}.
\end{proof}

\section{Przestrzeń \texorpdfstring{$C_0(X)$}{C\string_0(X)} oraz miary o skończonej wariacji}
\begin{lemat}[Gęstość $C_c(X)$ w $C_0(X)$]
\label{lem:density-cc-c0-21}
Dla dowolnej lokalnie zwartej przestrzeni Hausdorffa $X$
przestrzeń $C_c(X;\mathbb K)$ jest gęsta w
$C_0(X;\mathbb K)$ w normie supremum, gdzie
$\mathbb K=\mathbb R$ lub $\mathbb C$.
\end{lemat}
\begin{proof}
Dla $f\in C_0(X;\mathbb K)$ i $\varepsilon>0$
zbiór $K=\{|f|\ge\varepsilon\}$ jest zwarty.
Z lematu \ref{lem:cutoff-radon-20} wybieramy
$\eta\in C_c(X)$, $0\le\eta\le1$, $\eta|_K=1$.
Funkcja $\eta f$ ma nośnik zawarty w zwartym
nośniku $\eta$ i jest ciągła. Na $K$ błąd jest
zerowy, poza $K$ mamy $|f-\eta f|\le |f|<\varepsilon$.
Stąd $\|f-\eta f\|_\infty\le\varepsilon$.
\end{proof}

\begin{twierdzenie}[Ograniczone dodatnie funkcjonały na $C_0$]
\label{thm:riesz-c0-positive-21}
Niech $X$ będzie lokalnie zwartą przestrzenią Hausdorffa.
Dla dodatniego liniowego funkcjonału $L$ na $C_c(X;\mathbb R)$
i jego miary reprezentującej $\mu$ następujące warunki
są równoważne: $L$ jest ograniczony w normie supremum;
$\mu(X)<\infty$; $L$ ma ciągłe liniowe rozszerzenie
na $C_0(X;\mathbb R)$. Rozszerzenie jest wtedy
jednoznaczne, dodatnie, dane przez całkę względem $\mu$
i ma normę $\mu(X)$.
Każdy ograniczony dodatni funkcjonał na $C_0(X;\mathbb R)$
powstaje w ten sposób.
\end{twierdzenie}
\begin{proof}
Jeżeli $L$ jest ograniczony, ze wzoru
\eqref{eq:cc-open-sup-20} dla $U=X$ otrzymujemy
$\mu(X)\le\|L\|<\infty$. Nie wymagamy tu
przeliczalnego wyczerpania przestrzeni.
Jeśli $\mu(X)<\infty$, to
$|\int_X f\,\mathrm d\mu|\le\mu(X)\|f\|_\infty$
dla wszystkich $f\in C_0(X;\mathbb R)$.
Całka jest zatem ciągłym dodatnim rozszerzeniem
$L$. Jednoznaczność wynika z gęstości
(lemat \ref{lem:density-cc-c0-21}).
Ograniczenie ciągłego funkcjonału do $C_c(X)$
jest ograniczone, co zamyka cykl implikacji.
Norma rozszerzenia jest nie większa od $\mu(X)$,
a oszacowanie ze wzoru \eqref{eq:cc-open-sup-20}
daje przeciwną nierówność. Dla dowolnego
ograniczonego dodatniego funkcjonału na $C_0(X)$
stosujemy twierdzenie \ref{thm:riesz-positive-21}
do jego ograniczenia do $C_c(X)$, a następnie
korzystamy z jednoznaczności rozszerzenia.
\end{proof}

\begin{lemat}[Aproksymacja względem skończonej miary Radona]
\label{it:lem:radon-cc-l1}
Niech $X$ będzie lokalnie zwartą przestrzenią Hausdorffa.
Jeśli $\sigma$ jest skończoną miarą Radona na $X$, to
$C_c(X;\mathbb K)$ jest gęsta w $L^1(\sigma;\mathbb K)$,
także po uzupełnieniu miary.
Dla mierzalnej funkcji $u$ z $|u|\le1$ można wybrać
przybliżenia $f_n\in C_c(X;\mathbb K)$ spełniające
$\|f_n\|_\infty\le1$.
\end{lemat}
\begin{proof}
Dla miary na $\mathcal B(X)$ gęstość wynika bezpośrednio
z twierdzenia \ref{thm:cc-dense-lp-radon-20} dla $p=1$,
które obowiązuje na dowolnej lokalnie zwartej przestrzeni
Hausdorffa. Niech $\overline\sigma$ oznacza uzupełnienie $\sigma$.
Dla każdego zbioru $E$ mierzalnego względem $\overline\sigma$
istnieje borelowski $B$ taki, że $\overline\sigma(E\triangle B)=0$.
Każdą funkcję prostą względem uzupełnienia można zatem
zastąpić borelowską funkcją prostą równą jej prawie wszędzie,
zastępując w ten sposób wszystkie zbiory występujące
w skończonej reprezentacji. Funkcje proste są gęste w
$L^1(\overline\sigma)$ na mocy twierdzenia \ref{thm:simple-dense-lp-14}.
Aproksymujemy więc najpierw przez taką funkcję prostą,
a następnie przez funkcję z $C_c(X;\mathbb K)$,
stosując nierówność trójkąta dla błędów.
Dla $|u|\le1$ zastępujemy dowolne przybliżenia $g_n$
funkcjami $f_n=P\circ g_n$, gdzie
$P(z)=z/\max(1,|z|)$.
Jest to odwzorowanie ciągłe, $P(0)=0$, więc
$f_n\in C_c(X;\mathbb K)$ i $|f_n|\le1$.
Jeżeli $z=rv$, $r>1$, $|v|=1$, $|w|\le1$, to
\[
 |z-w|^2-|P(z)-w|^2
 =(r-1)(r+1-2\operatorname{Re}(v\overline w))\ge0.
\]
Dla $|z|\le1$ mamy $P(z)=z$.
Stąd $|f_n-u|\le|g_n-u|$, co dowodzi tezy.
\end{proof}

\begin{definicja}[Skończona miara Radona ze znakiem lub zespolona]
\label{idx:21:definicja:08}
Niech $X$ będzie lokalnie zwartą przestrzenią Hausdorffa.
Przeliczalnie addytywną miarę borelowską $\nu$ o wartościach
w $\mathbb K=\mathbb R$ lub $\mathbb C$ nazywamy skończoną
miarą Radona, gdy jej całkowite wahanie $|\nu|$ jest
skończoną miarą Radona. Oznaczamy przestrzeń tych miar
przez $M(X;\mathbb K)$ i kładziemy $\|\nu\|=|\nu|(X)$.
Całkowite wahanie jest określone wzorem
\[
 |\nu|(E)=\sup\left\{\sum_{j=1}^m|\nu(E_j)|:
 E=\bigsqcup_{j=1}^m E_j,\ E_j\in\mathcal B(X)\right\}.
\]
\end{definicja}

\begin{uwaga}[Własności miar ze znakiem i zespolonych]
Korzystamy z wyników teorii miary:
\atlasref{ME-16::thm:jordan-16}{rozkładu Jordana},
\atlasref{ME-16::thm:partition-variation-16}{wzoru na wahanie miary ze znakiem}
oraz \atlasref{ME-16::thm:complex-variation-16}{twierdzenia o wahaniu miary zespolonej}.
Jeżeli $\nu$ jest rzeczywista, to
$|\nu|=\nu^++\nu^-$, gdzie $\nu^+,\nu^-$ są
wzajemnie osobliwymi nieujemnymi częściami Jordana.
Dla miary zespolonej całkę definiujemy liniowo przez
cztery nieujemne części Jordana jej części rzeczywistej
i urojonej. Dla ograniczonej mierzalnej $f$ mamy
\[
 \left|\int_X f\,\mathrm d\nu\right|
 \le\int_X |f|\,\mathrm d|\nu|.
\]
Istotnie, dla funkcji prostej $f=\sum_j a_j\chi_{E_j}$
o rozłącznych zbiorach $E_j$ wynika to z nierówności
$|\nu(E_j)|\le|\nu|(E_j)$ i nierówności trójkąta.
Dla ograniczonych funkcji mierzalnych przechodzimy do
granicy w jednostajnej aproksymacji funkcjami prostymi;
wszystkie cztery miary dodatnie są skończone.
Części Jordana miar $\operatorname{Re}\nu$ i
$\operatorname{Im}\nu$ są zdominowane przez $|\nu|$,
gdyż ich całkowite wahania są nie większe od $|\nu|$
na mocy definicji przez podziały. Są więc Radona
z lematu \ref{it:lem:radon-dominated}. Odwrotnie,
każda kombinacja liniowa skończonych dodatnich miar
Radona jest miarą Radona: jej całkowite wahanie
jest zdominowane przez sumę tych miar pomnożonych
przez moduły współczynników, co wynika z definicji
przez podziały. Lemat \ref{it:lem:radon-dominated}
daje regularność. W szczególności $M(X;\mathbb K)$
jest przestrzenią wektorową. Przyporządkowanie
$\nu\mapsto|\nu|(X)$ jest normą: z definicji przez
podziały wynika $|a\nu|(X)=|a||\nu|(X)$ i
$|\nu+\xi|(X)\le|\nu|(X)+|\xi|(X)$, ponieważ dla
każdego podziału stosujemy nierówność trójkąta
do każdego składnika i bierzemy supremum.
Jeżeli $|\nu|(X)=0$, to dla każdego borelowskiego $E$
mamy $|\nu(E)|\le|\nu|(E)\le|\nu|(X)=0$,
a więc $\nu=0$.
\end{uwaga}

\begin{twierdzenie}[Reprezentacja funkcjonałów rzeczywistych i zespolonych]
\label{thm:riesz-complex-21}
Niech $X$ będzie dowolną lokalnie zwartą przestrzenią
Hausdorffa. Dla $\mathbb K=\mathbb R$ lub $\mathbb C$
każdy ograniczony liniowy funkcjonał $T$ na
$C_0(X;\mathbb K)$ ma dokładnie jedną reprezentację
\[
 T(f)=\int_X f\,\mathrm d\nu,\qquad
 \nu\in M(X;\mathbb K).
\]
Ponadto $\|T\|=|\nu|(X)$, a przyporządkowanie
$\nu\mapsto T_\nu$ jest liniowym izometrycznym
izomorfizmem $M(X;\mathbb K)\cong C_0(X;\mathbb K)'$.
\end{twierdzenie}
\begin{proof}
\emph{Krok 1: rozkład funkcjonału rzeczywistego.}
Dla $f\in C_0(X;\mathbb R)$, $f\ge0$, kładziemy
\[
 T^+(f)=\sup\{T(g):g\in C_0(X;\mathbb R),\ 0\le g\le f\}.
\]
Mamy $0\le T^+(f)\le\|T\|\|f\|_\infty$.
Dodatnia jednorodność wynika z przeskalowania zbioru
dopuszczalnych $g$ (dla skalaru zero wprost).
Dla $f_1,f_2\ge0$ sumy dopuszczalnych $g_1,g_2$
są dopuszczalne dla $f_1+f_2$, więc
$T^+(f_1+f_2)\ge T^+(f_1)+T^+(f_2)$.
Odwrotnie, jeśli $0\le g\le f_1+f_2$, kładziemy
$g_1=g\wedge f_1$, $g_2=g-g_1$. Są to nieujemne
funkcje z $C_0(X)$, $g_1\le f_1$, $g_2\le f_2$.
Własność znikania w nieskończoności wynika z
domkniętości zbiorów poziomicowych i dominacji.
Zatem $T(g)\le T^+(f_1)+T^+(f_2)$.
Biorąc supremum po $g$, otrzymujemy nierówność
przeciwną, a zatem addytywność $T^+$ na stożku.

Definiujemy $T^+(u-v)=T^+(u)-T^+(v)$ dla $u,v\ge0$.
Jeżeli $u-v=u'-v'$, to $u+v'=u'+v$, więc
addytywność dowodzi niezależności od rozkładu.
Dodawanie rozkładów i skalowanie (dla ujemnego
skalaru zamiana składników) dają liniowość.
Otrzymujemy funkcjonał dodatni, a
$|T^+(f)|\le T^+(|f|)\le\|T\|\|f\|_\infty$.
Funkcjonał $T^-=T^+-T$ jest dodatni, bo dla $f\ge0$
można wybrać $g=f$ w supremum. Jest też ograniczony
jako różnica ograniczonych funkcjonałów.
Twierdzenie \ref{thm:riesz-c0-positive-21} daje
skończone miary dodatnie $\alpha,\beta$ reprezentujące
$T^+,T^-$. Stąd $\nu=\alpha-\beta$ reprezentuje $T$
i jest miarą Radona z lematu \ref{it:lem:radon-dominated}.

Sprawdzimy jeszcze, że $\alpha,\beta$ są dokładnie
częściami Jordana $\nu^+,\nu^-$, a nie tylko pewnymi
miarami dodatnimi, których różnicą jest $\nu$.
Niech $P$ będzie dodatnią częścią borelowskiego
rozkładu Hahna miary $\nu$.
Dla $f\in C_0(X;\mathbb R)$, $f\ge0$, oraz
$0\le g\le f$ mamy
\[
 T(g)=\int_X g\,\mathrm d\nu
 \le\int_X f\,\mathrm d\nu^+.
\]
Stąd $T^+(f)\le\int_X f\,\mathrm d\nu^+$.
Z lematu \ref{it:lem:radon-cc-l1} zastosowanego
do skończonej miary Radona $|\nu|$ otrzymujemy
$h_n\in C_c(X;\mathbb R)$ z
$\|h_n-\chi_P\|_{L^1(|\nu|)}\to0$.
Obcinając $h_n$ do przedziału $[0,1]$ przez
$h_n\mapsto0\vee(h_n\wedge1)$, zachowujemy
ciągłość, zwarty nośnik i nie zwiększamy błędu
względem $\chi_P$. Możemy więc założyć
$0\le h_n\le1$. Funkcje $g_n=fh_n$ są
dopuszczalne w definicji $T^+(f)$ i
\[
 \left|\int_X g_n\,\mathrm d\nu
             -\int_X f\chi_P\,\mathrm d\nu\right|
 \le\|f\|_\infty\|h_n-\chi_P\|_{L^1(|\nu|)}
 \longrightarrow0.
\]
Ponieważ $\int_X f\chi_P\,\mathrm d\nu
=\int_X f\,\mathrm d\nu^+$, otrzymujemy równość
$T^+(f)=\int_X f\,\mathrm d\nu^+$.
Części Jordana są miarami Radona przez dominację
przez $|\nu|$. Liniowość i rozkład funkcji na części
dodatnią i ujemną dają równość całek dla wszystkich
funkcji rzeczywistych z $C_0(X)$. Z jednoznaczności
reprezentacji dodatniej
(twierdzenie \ref{thm:riesz-positive-21}) mamy
$\alpha=\nu^+$, a z $\nu=\alpha-\beta$ wynika
$\beta=\nu^-$.

\emph{Krok 2: przypadek zespolony.}
Na rzeczywistych funkcjach z $C_0(X)$ definiujemy
$T_1(f)=\operatorname{Re}T(f)$ i
$T_2(f)=\operatorname{Im}T(f)$. Są rzeczywistymi
ograniczonymi funkcjonałami liniowymi, więc krok 1
daje reprezentujące je miary rzeczywiste $\nu_1,\nu_2$.
Miara $\nu=\nu_1+i\nu_2$ jest Radona dzięki dominacji
$|\nu|\le|\nu_1|+|\nu_2|$. Dla $f=u+iv$ z $u,v$
rzeczywistymi mamy
\[
 T(f)=T_1(u)-T_2(v)+i\bigl(T_2(u)+T_1(v)\bigr)
 =\int_X f\,\mathrm d\nu.
\]

\emph{Krok 3: norma.}
Dla każdej miary $\nu\in M(X;\mathbb K)$ oszacowanie
całki daje $\|T_\nu\|\le|\nu|(X)$.
Dla $\varepsilon>0$ wybieramy skończony borelowski
podział $X=\bigsqcup_{j=1}^m E_j$ z
$\sum_j|\nu(E_j)|>|\nu|(X)-\varepsilon$.
Niech $c_j=\overline{\nu(E_j)}/|\nu(E_j)|$ gdy
$\nu(E_j)\ne0$, a $c_j=0$ w przeciwnym razie.
Funkcja $u=\sum_jc_j\chi_{E_j}$ ma moduł co najwyżej
$1$ i $\int_Xu\,\mathrm d\nu=\sum_j|\nu(E_j)|$.
Dla miary rzeczywistej współczynniki są rzeczywiste.
Lemat \ref{it:lem:radon-cc-l1} dla $\sigma=|\nu|$
daje funkcje $f_n\in C_c(X;\mathbb K)$ o normie
supremum co najwyżej $1$ i
$\int_X|f_n-u|\,\mathrm d|\nu|\to0$. Zatem
\[
 |\nu|(X)-\varepsilon
 <\left|\int_Xu\,\mathrm d\nu\right|
 =\lim_n|T_\nu(f_n)|\le\|T_\nu\|.
\]
Dowolność $\varepsilon$ daje równość norm.
Nie korzystaliśmy z twierdzenia Radona--Nikodyma.

\emph{Krok 4: jednoznaczność i izomorfizm.}
Każda skończona miara Radona definiuje ograniczony
funkcjonał przez całkę. Jeżeli dwie miary definiują
ten sam funkcjonał, ich różnica $\xi$ jest Radona
i $T_\xi=0$. Z kroku 3 mamy $|\xi|(X)=0$,
więc $|\xi(E)|\le|\xi|(X)=0$ dla każdego $E$.
To dowodzi jednoznaczności. Liniowość względem miary wynika najpierw dla
funkcji prostych z równości
$\int_X s\,\mathrm d\nu=\sum_j a_j\nu(E_j)$
dla $s=\sum_j a_j\chi_{E_j}$ na rozłącznych zbiorach,
a następnie dla funkcji ograniczonych przez jednostajną
aproksymację funkcjami prostymi i oszacowanie całki
przez całkowite wahanie każdej z rozważanych miar.
Surjektywność wynika z kroków 1--2, a izometryczność
z kroku 3.
\end{proof}

\begin{proof}[Alternatywny dowód w przypadku zespolonym]
Dla nieujemnego $g\in C_c(X;\mathbb R)$ zdefiniujmy
\[
 M(g)=\sup\{|T(h)|:h\in C_c(X;\mathbb C),\ |h|\le g\}.
\]
Jest to skończone, bo $M(g)\le\|T\|\|g\|_\infty$.
Dodatnia jednorodność $M(ag)=aM(g)$ dla $a>0$ wynika
z bijekcji $h\mapsto ah$ pomiędzy zbiorami dopuszczalnymi;
dla $a=0$ jedyną funkcją dopuszczalną jest zero.
Wykażemy addytywność $M(g_1+g_2)=M(g_1)+M(g_2)$.
Dla nierówności „$\le$” rozkładamy dowolną $h$,
$|h|\le g_1+g_2$, na $h_i=h g_i/(g_1+g_2)$ na
$\{g_1+g_2>0\}$, a w pozostałych punktach przyjmujemy
$h_i=0$. Ciągłość przy zerach mianownika wynika
z $|h_i|\le g_i$; nośnik $h_i$ zawiera się w zwartym
nośniku $g_i$, więc $h_i\in C_c(X;\mathbb C)$.
Mamy $h_1+h_2=h$, stąd
$|T(h)|\le |T(h_1)|+|T(h_2)|\le M(g_1)+M(g_2)$.
Supremum po $h$ daje pierwszą nierówność.
Dla nierówności „$\ge$” ustalmy $\varepsilon>0$ i wybierzmy
$h_i$ z $|h_i|\le g_i$ oraz $|T(h_i)|>M(g_i)-\varepsilon/2$.
Mnożymy $h_i$ przez $\overline{T(h_i)}/|T(h_i)|$,
gdy $T(h_i)\ne0$, a przez $1$ w przeciwnym przypadku.
Nie zmieniamy modułów funkcji, a dzięki liniowości zespolonej
otrzymujemy $T(h_i)=|T(h_i)|$. Wtedy $|h_1+h_2|\le g_1+g_2$
i $M(g_1+g_2)\ge T(h_1+h_2)>M(g_1)+M(g_2)-\varepsilon$.
Dowolność $\varepsilon$ kończy dowód addytywności.
Z addytywności i dodatniej jednorodności $M$ przedłuża
się jednoznacznie do dodatniego rzeczywistego funkcjonału
na $C_c(X;\mathbb R)$: $M(f)=M(f^+)-M(f^-)$.
Dla dwóch rozkładów $f=u-v=u'-v'$ na funkcje nieujemne
mamy $u+v'=u'+v$; addytywność zapewnia niezależność
różnicy $M(u)-M(v)$ od rozkładu. Dodawanie i skalowanie
rozkładów dają liniowość przedłużenia. Z dodatniości
mamy $|M(f)|\le M(|f|)\le\|T\|\|f\|_\infty$,
więc przedłużony funkcjonał jest ograniczony i rozciąga się
na $C_0(X;\mathbb R)$ na mocy lematu
\ref{lem:density-cc-c0-21}.
Twierdzenie~\ref{thm:riesz-c0-positive-21} daje
skończoną miarę Radona $\sigma$ taką, że
$M(g)=\int_Xg\,\mathrm d\sigma$ dla $g\ge0$.
W szczególności dla $h\in C_c(X;\mathbb C)$ zachodzi
\begin{equation}\label{eq:riesz-domination-21}
 |T(h)|\le \int_X|h(x)|\,\mathrm d\sigma(x).
\end{equation}
Z \eqref{eq:riesz-domination-21} wynika, że jeśli $h_1=h_2$
$\sigma$-prawie wszędzie, to $T(h_1-h_2)=0$.
Funkcjonał jest zatem dobrze określony na klasach
równoważności funkcji z $C_c(X;\mathbb C)$ i ograniczony
w normie $L^1(\sigma)$. Klasy te są gęste w
$L^1(\sigma;\mathbb C)$ na mocy lematu
\ref{it:lem:radon-cc-l1}. Zatem $T$ jednoznacznie
przedłuża się do funkcjonału na $L^1(\sigma)$
o normie co najwyżej $1$.
Twierdzenie \ref{thm:dual-l1-15} o dualności
$L^1$--$L^\infty$ (miara $\sigma$ jest skończona) daje $w\in L^\infty(\sigma)$,
$|w|\le1$ $\sigma$-prawie wszędzie, takie że
$T(h)=\int_X h w\,\mathrm d\sigma$.
Kładziemy $\nu(E)=\int_E w\,\mathrm d\sigma$.
Ponieważ $\sigma$ jest skończona i $w\in L^\infty(\sigma)$,
mamy $w\in L^1(\sigma)$. Z twierdzenia \ref{thm:density-variation-16}
wynika, że $\nu$ jest miarą zespoloną oraz
$|\nu|(E)=\int_E|w|\,\mathrm d\sigma\le\sigma(E)$.
Jest to więc skończona miara Radona, gdyż
regularność przechodzi na miary zdominowane (lemat \ref{it:lem:radon-dominated}).
Reprezentacja obowiązuje dla $C_c$, a przez gęstość
także dla $C_0$.

Pozostaje identyczność norm. Nierówność
$\|T\|\le|\nu|(X)$ jest bezpośrednia. Dla odwrotnej
korzystamy z twierdzenia \ref{thm:polar-complex-measure-18}:
jeśli $v$ jest gęstością $\nu$ względem $|\nu|$, to
$|v|=1$ $|\nu|$-prawie wszędzie. Wybieramy mierzalną
funkcję $u=\overline v$ na $\{|v|=1\}$ i $u=0$ poza tym
zbiorem. Mamy $|u|\le1$ oraz $u\,\mathrm d\nu=\mathrm d|\nu|$.
Z regularności $|\nu|$ i aproksymacji funkcji
mierzalnych przez $C_c$ w $L^1(|\nu|)$ istnieją
$f_n\in C_c(X;\mathbb C)$, $\|f_n\|_\infty\le1$,
$f_n\to u$ w $L^1(|\nu|)$. Ograniczenie normy
zapewnia obcięcie aproksymantów do domkniętego
dysku jednostkowego, które nie zwiększa błędu
punktowego względem $u$. Stąd
$T(f_n)\to\int_Xu\,\mathrm d\nu=|\nu|(X)$,
więc $\|T\|\ge|\nu|(X)$.
Dla jednoznaczności niech $\nu,\widetilde\nu$ reprezentują
$T$ i niech $\xi=\nu-\widetilde\nu$. Miara $\xi$ jest
skończoną miarą Radona, a jej funkcjonał całkowy jest zerowy.
Udowodniona równość norm dla funkcjonału całkowego
(także dla $\xi$) daje $|\xi|(X)=0$. Zatem dla każdego
borelowskiego $E$ mamy $|\xi(E)|\le|\xi|(X)=0$,
czyli $\nu=\widetilde\nu$.
\end{proof}



\begin{przyklad}[Funkcjonał ewaluacji]
Dla ustalonego $x_0\in X$ funkcjonał $T(f)=f(x_0)$
na $C_0(X)$ jest ciągły i ma normę $1$.
Jego miarą reprezentującą jest $\delta_{x_0}$.
\end{przyklad}

\begin{przyklad}[Miara nieskończona i funkcjonał lokalny]
Dla $X=\mathbb R$ i $L(f)=\int_{\mathbb R} f(x)
\,\mathrm d\lambda_1(x)$, $f\in C_c(\mathbb R)$,
otrzymujemy dodatni funkcjonał nieograniczony w normie
supremum na całym $C_c$: funkcje odcinające równe $1$
na $[-n,n]$ mają normę supremum $1$, ale całki
co najmniej $2n$. Dlatego $L$ nie przedłuża się
ciągle do $C_0(\mathbb R)$.
\end{przyklad}

\begin{przyklad}[Gęstość i atom w jednym funkcjonale]
Na $C_c(\mathbb R)$ definiujemy $\Lambda(\varphi)=\int_{[0,1]}\varphi(x)\,\mathrm d\lambda_1(x)+2\varphi(0)$. Jest to dodatni funkcjonał liniowy. Miara $\mu(E)=\lambda_1(E\cap[0,1])+2\delta_0(E)$ jest skończoną miarą borelowską, więc z twierdzenia \ref{thm:lch-secondcountable-radon-20} jest Radona. Całka względem sumy miar jest sumą całek, najpierw dla funkcji prostych, potem przez aproksymację monotoniczną i rozkład na części dodatnią i ujemną. Stąd $\mu$ reprezentuje $\Lambda$ i
\[
 \int_{\mathbb R}1\,\mathrm d\mu=3,\quad
 \int_\mathbb R x^2\,\mathrm d\mu(x)=\int_0^1 x^2\,\mathrm dx+2\cdot0^2=1/3,\quad
 \int_\mathbb R\chi_{\{0\}}\,\mathrm d\mu=2.
\]
W rachunku z $x^2$ użyliśmy zgodności całek Riemanna
i Lebesgue'a (twierdzenie \atlasref{ME-08::thm:riemann-to-lebesgue}{ME-08: zgodność całek Riemanna i Lebesgue'a}).
Wszystkie trzy funkcje są całkowalne względem $\mu$,
choć nie należą do $C_c(\mathbb R)$: miara jest skupiona na
$[0,1]$, na którym są ograniczone. Rachunki dotyczą rozszerzonej całki.
\end{przyklad}
\begin{przyklad}[Całka skupiona na okręgu]
Dla $\varphi\in C_c(\mathbb R^2)$ definiujemy $\Lambda(\varphi)=\int_0^{2\pi}\varphi(\cos t,\sin t)\,\mathrm dt$. Funkcjonał jest dodatni i liniowy. Niech $\gamma(t)=(\cos t,\sin t)$ i $\mu=\gamma_*(\lambda_1|_{[0,2\pi]})$. Jest to skończona miara borelowska na $\mathbb R^2$, więc Radona z twierdzenia \ref{thm:lch-secondcountable-radon-20}. Twierdzenie \atlasref{ME-13::thm:pushforward-integral-13}{ME-13: całka względem miary obrazowej}
o całkowaniu względem miary obrazowej oraz twierdzenie
\atlasref{ME-08::thm:riemann-to-lebesgue}{ME-08: zgodność całek Riemanna i Lebesgue'a} o zgodności całki Riemanna
z Lebesgue'a dają reprezentację $\Lambda$. Miara opisuje długość łuku, gdyż $|\gamma'(t)|=1$, a podwójnie przyjmowany punkt końcowy odpowiada zbiorowi parametrów miary zero. Mamy $\mu(S^1)=2\pi$ oraz
\[
 \int_{S^1}x_1^2\,\mathrm d\mu(x)=\int_0^{2\pi}\cos^2t\,\mathrm dt
 =\tfrac12\int_0^{2\pi}(1+\cos(2t))\,\mathrm dt=\pi.
\]
Natomiast $\lambda_2(S^1)=0$. Okrąg jest domknięty
w $\mathbb R^2$, a więc należy do produktowej
$\sigma$-algebry borelowskiej, gdyż $\mathbb R$
ma przeliczalną bazę i otwarte podzbiory produktu
są przeliczalnymi sumami otwartych prostokątów.
Każdy pionowy przekrój okręgu ma co najwyżej dwa punkty,
a więc jednowymiarową miarę Lebesgue'a zero.
Z twierdzeń \atlasref{ME-12::thm:tonelli-12}{ME-12: twierdzenie Tonellego} i
\atlasref{ME-12::thm:lebesgue-products-12}{ME-12: produkty miar Lebesgue'a} otrzymujemy
\[
 \lambda_2(S^1)
 =\int_{\mathbb R}\lambda_1(\{t:(x_1,t)\in S^1\})
      \,\mathrm d\lambda_1(x_1)=0.
\]
\end{przyklad}
\begin{przyklad}[Skończona suma ewaluacji punktowych]
Niech $X$ będzie lokalnie zwartą przestrzenią Hausdorffa. Dla punktów $x_1,\dots,x_m\in X$ oraz liczb $c_1,\dots,c_m>0$ funkcjonał $\Lambda(f)=\sum_{j=1}^m c_jf(x_j)$ na $C_c(X)$ jest dodatni i liniowy. Reprezentuje go miara Radona $\mu=\sum_{j=1}^m c_j\delta_{x_j}$. Istotnie, dla dowolnego borelowskiego $E$ zwarty zbiór $\{x_j:x_j\in E\}$ jest podzbiorem $E$ o tej samej mierze, a otwarty zbiór $X\setminus\{x_j:x_j\notin E\}$ jest nadzbiorem $E$ o tej samej mierze. To daje obie regularności; skończoność wynika z sumy wag. Reprezentacja wynika z definicji całki względem mas Diraca i addytywności względem miary. Twierdzenie \ref{tw:radon-unikalnosc} pokazuje, że żadna inna regularna miara borelowska nie daje tych samych wartości $\Lambda(f)$ dla wszystkich $f\in C_c(X)$.
\end{przyklad}

\begin{wniosek}[Porządek dodatnich funkcjonałów i miar]
\label{cor:riesz-order-21}
Niech $X$ będzie lokalnie zwartą przestrzenią
Hausdorffa, a $L_1,L_2:C_c(X;\mathbb R)\to\mathbb R$
dodatnimi funkcjonałami liniowymi, reprezentowanymi
odpowiednio przez miary Radona $\mu_1,\mu_2$.
Wówczas następujące warunki są równoważne:
\begin{enumerate}[label=(\roman*)]
\item $L_1(f)\le L_2(f)$ dla każdego nieujemnego $f\in C_c(X)$;
\item $\mu_1(E)\le\mu_2(E)$ dla każdego zbioru borelowskiego
$E\subseteq X$.
\end{enumerate}
\end{wniosek}
\begin{proof}
Implikacja (ii)$\Rightarrow$(i) wynika z monotoniczności całki.
Załóżmy (i). Dla otwartego $U\subseteq X$ zachodzi
\[
 \mu_1(U)=\sup\{L_1(f):f\in C_c(X),\ 0\le f\le1,
 \ \operatorname{supp}f\subseteq U\}.
\]
Wzór ten udowodniono w rozdziale~\ref{ch:radon-measures-20};
obejmuje także przypadek $\mu_1(U)=+\infty$.
Nierówność dla funkcjonałów daje
$\mu_1(U)\le\mu_2(U)$. Dla dowolnego borelowskiego
$E$ z regularności zewnętrznej otrzymujemy
$\mu_1(E)=\inf_{U\supseteq E}\mu_1(U)
\le\inf_{U\supseteq E}\mu_2(U)=\mu_2(E)$.
\end{proof}

\begin{uwaga}[Gdzie użyliśmy twierdzenia Tichonowa?]
W konstrukcji na zwartej przestrzeni metrycznej nie
potrzebowaliśmy twierdzenia Tichonowa: miarę zewnętrzną
zbudowaliśmy bezpośrednio z dodatniego funkcjonału.
W rozdziale~\ref{ch:haar-23}, poświęconym mierze Haara,
pojawi się natomiast zwartość produktu nieprzeliczalnie
wielu przedziałów.
Nie należy jej mylić ze zwartą metryczną przestrzenią,
dla której wystarczy argument diagonalny.
\end{uwaga}

\section{Konsekwencje: oddzielenie miar i problem momentów}
\begin{wniosek}[Miary odróżniane przez funkcje ciągłe]
\label{cor:measures-separate-21}
Niech $\mu,\nu$ będą miarami Radona na lokalnie zwartej
przestrzeni Hausdorffa $X$. Wtedy
$\mu=\nu$ wtedy i tylko wtedy, gdy
$\int_Xf\,\mathrm d\mu=\int_Xf\,\mathrm d\nu$
dla wszystkich $f\in C_c(X;\mathbb R)$.
\end{wniosek}
\begin{proof}
Implikacja w prawo jest oczywista, a odwrotna to
\ref{thm:uniqueness-cc-20}.
\end{proof}

\begin{uwaga}[Problem momentów jako ostrzeżenie]
Jeżeli na $\mathbb R$ znamy całki
$\int_{\mathbb R}x^k\,\mathrm d\mu(x)$ dla wszystkich
$k\ge0$, nie zawsze wyznacza to jednoznacznie miarę.
Funkcje wielomianowe nie należą do $C_0(\mathbb R)$
i nie ma automatycznej gęstości wielomianów w normie
supremum na całej prostej. Na zwartym przedziale
sytuacja jest inna: twierdzenie Stone'a--Weierstrassa
przybliża dowolną funkcję ciągłą wielomianami,
więc momenty jednoznacznie wyznaczają skończoną miarę.
\end{uwaga}

\begin{przyklad}[Norma funkcjonału atomowego]
\label{ex:riesz-complex-atoms-21}
Niech $X$ będzie lokalnie zwartą przestrzenią Hausdorffa,
$x_1,\ldots,x_m\in X$ parami różnymi punktami,
a $a_1,\ldots,a_m\in\mathbb C$. Na $C_0(X;\mathbb C)$
określmy funkcjonał
\[
 T(f)=\sum_{j=1}^ma_jf(x_j).
\]
Reprezentuje go miara zespolona
$\nu=\sum_{j=1}^ma_j\delta_{x_j}$, a jej wariacja
całkowita wynosi
$|\nu|=\sum_{j=1}^m|a_j|\delta_{x_j}$.
W szczególności
$\|T\|=\sum_{j=1}^m|a_j|$.

Sprawdźmy bezpośrednio ostatnią równość. Z nierówności
trójkąta otrzymujemy
$|T(f)|\le\|f\|_\infty\sum_j|a_j|$.
Ponieważ punktów jest skończenie wiele i są różne,
dla każdej pary punktów wybieramy rozłączne otwarte
otoczenia z aksjomatu Hausdorffa. Dla każdego punktu
przecinamy skończenie wiele przypisanych mu otoczeń.
Otrzymujemy parami rozłączne otwarte otoczenia, które
możemy zmniejszyć do otwartych $U_j$ o zwartych domknięciach
na mocy lematu \atlasref{GT-14::lem:topog-baire-zwarte-domkniecie}{GT-14: otoczenie o zwartym domknięciu}. Z lematu
\ref{lem:cutoff-radon-20} istnieją
$\varphi_j\in C_c(X)$, $0\le\varphi_j\le1$,
$\varphi_j(x_j)=1$ i
$\operatorname{supp}\varphi_j\subseteq U_j$.
Dla $a_j\ne0$ kładziemy
$b_j=\overline{a_j}/|a_j|$, a dla $a_j=0$
przyjmujemy $b_j=0$. Wtedy
$f=\sum_{j=1}^m b_j\varphi_j$ należy do
$C_c(X;\mathbb C)$ i, dzięki rozłączności nośników,
$\|f\|_\infty\le1$. Ponadto
$T(f)=\sum_j a_jb_j=\sum_j|a_j|$,
co dowodzi wzoru na normę.
\end{przyklad}

\section{Podsumowanie i przygotowanie do zbieżności miar}
Reprezentacja dodatnich funkcjonałów pozwala przejść
od analizy na $C_c(X)$ do teorii miar Radona. Dla
funkcjonałów na $C_0(X)$ skończoność normy odpowiada
skończoności masy, zaś dla zespolonych funkcjonałów
norma jest wariacją całkowitą. W rozdziale~\ref{ch:weak-measures-22}
zastosujemy zwarty przypadek twierdzenia do wykazania
ciągowej zwartości przestrzeni prawdopodobieństw na
zwartej przestrzeni metrycznej.

% ROZDZIAŁ 22 — ZBIEŻNOŚĆ MIAR I PROCHOROW
\chapter[Zbieżność miar i twierdzenie Prochorowa]{Zbieżność miar, twierdzenie Portmanteau i zwartość Prochorowa}
\label{ch:weak-measures-22}

\section{Trzy topologie, których nie należy utożsamiać}
Niech $(X,d)$ będzie przestrzenią metryczną,
$\mathcal B(X)$ jej $\sigma$-algebrą borelowską,
a $\mathcal P(X)$ zbiorem miar probabilistycznych na $X$.
Przestrzeń $C_b(X)$ składa się z ograniczonych rzeczywistych
funkcji ciągłych. Jeśli $X$ jest dodatkowo lokalnie zwarta,
$C_c(X)$ oznacza przestrzeń zdefiniowaną w
rozdziale~\ref{ch:radon-measures-20}, a $C_0(X)$
przestrzeń wprowadzoną w rozdziale~\ref{ch:riesz-markov-21}.

\begin{definicja}[Zbieżność wąska]
\label{def:narrow-measures-22}
Ciąg $(\mu_n)\subseteq\mathcal P(X)$ jest
\emph{wąsko zbieżny} do $\mu\in\mathcal P(X)$,
co zapisujemy $\mu_n\Rightarrow\mu$, jeśli
\begin{equation}\label{eq:narrow-22}
 \int_X f(x)\,\mathrm d\mu_n(x)
 \longrightarrow\int_X f(x)\,\mathrm d\mu(x)
 \quad\text{dla każdej }f\in C_b(X).
\end{equation}
\end{definicja}

\begin{definicja}[Zbieżność słaba-* i mglista]
\label{def:vague-weakstar-22}
Niech $X$ będzie lokalnie zwartą ośrodkową przestrzenią
metryczną, a $\mu_n,\mu$ miarami Radona.
Nazywamy zbieżność $\mu_n\to\mu$ \emph{mglistą},
gdy $\int_Xf\,\mathrm d\mu_n\to\int_Xf\,\mathrm d\mu$
dla każdego $f\in C_c(X)$.
Jeżeli miary są skończone, zbieżność
$\sigma(C_0(X)^*,C_0(X))$ (\emph{słaba-*})
oznacza zbieżność całek dla wszystkich $f\in C_0(X)$.
W przypadku miar probabilistycznych zbieżność wąska
implikuje obie te zbieżności. Na przestrzeni niezwartej
implikacje odwrotne nie muszą zachodzić.
\end{definicja}

\begin{przyklad}[Ucieczka masy]
\label{ex:mass-escape-22}
Niech $X=\mathbb R$, $\mu_n=\delta_n$.
Dla każdego $f\in C_0(\mathbb R)$ zachodzi $f(n)\to0$,
zatem $\mu_n\to0$ słabo-* w dualnej przestrzeni
$C_0(\mathbb R)^*$. Także $\mu_n\to0$ mgliście.
Jednak $\mu_n$ nie zbiega wąsko do żadnej miary
probabilistycznej: stała funkcja $1\in C_b(\mathbb R)$
ma całkę równą $1$ względem każdej $\mu_n$.
W szczególności określenie „zbieżność słaba miar” bez
wskazania klasy funkcji testowych bywa dwuznaczne.
\end{przyklad}

\section{Twierdzenie Portmanteau}
\begin{definicja}[Zbiór ciągłości miary]
\label{def:continuity-set-22}
Niech $\mu\in\mathcal P(X)$. Zbiór borelowski
$A\subseteq X$ nazywamy \emph{zbiorem ciągłości miary}
$\mu$, jeżeli $\mu(\partial A)=0$, gdzie
$\partial A=\overline A\setminus A^\circ$ jest jego brzegiem.
\end{definicja}

\begin{twierdzenie}[Portmanteau]
\label{thm:portmanteau-22}
Niech $(X,d)$ będzie dowolną przestrzenią metryczną,
a $\mu_n,\mu\in\mathcal P(X)$. Następujące warunki
są równoważne:
\begin{enumerate}[label=(\roman*)]
\item $\mu_n\Rightarrow\mu$;
\item dla każdej funkcji ograniczonej, dolnie półciągłej
$f:X\to\mathbb R$ zachodzi
$\int_Xf\,\mathrm d\mu\le\liminf_n\int_Xf\,\mathrm d\mu_n$;
\item dla każdego otwartego $G\subseteq X$ zachodzi
$\mu(G)\le\liminf_n\mu_n(G)$;
\item dla każdego domkniętego $F\subseteq X$ zachodzi
$\limsup_n\mu_n(F)\le\mu(F)$;
\item dla każdego zbioru ciągłości $A$ miary $\mu$
zachodzi $\mu_n(A)\to\mu(A)$.
\end{enumerate}
W warunku (ii) można także dopuścić dolnie półciągłe
funkcje ograniczone z dołu, przy całkach w $(-\infty,+\infty]$.
\end{twierdzenie}
\begin{proof}
\emph{(i)$\Rightarrow$(iii).}
Dla otwartego $G$ określamy
$f_k(x)=\min\{1,kd(x,X\setminus G)\}$, gdy
$G\neq X$, oraz $f_k\equiv1$, gdy $G=X$.
Funkcje $f_k$ są ciągłe, $0\le f_k\le\chi_G$,
a $f_k\uparrow\chi_G$. Zatem dla każdego $k$
\[
\liminf_n\mu_n(G)\ge\lim_n\int_X f_k\,\mathrm d\mu_n
=\int_X f_k\,\mathrm d\mu.
\]
Z twierdzenia Beppo Leviego, po przejściu $k\to\infty$,
otrzymujemy (iii).

\emph{(iii)$\Leftrightarrow$(iv).}
Dla $F$ domkniętego zbiór $X\setminus F$ jest otwarty,
a $\mu_n(X\setminus F)=1-\mu_n(F)$ i podobnie
dla $\mu$. Pozwala to zamienić dolną granicę na górną.

\emph{(iii)$\Rightarrow$(ii).}
Najpierw niech $f$ będzie dolnie półciągła i
$0\le f\le M$. Dla każdego $t\in[0,M]$ zbiór
$G_t=\{x:f(x)>t\}$ jest otwarty. Z warunku (iii)
$\mu(G_t)\le\liminf_n\mu_n(G_t)$.
Z reprezentacji warstwowej całki
$\int_X f\,\mathrm d\mu=\int_{[0,M]}\mu(G_t)
\,\mathrm d\lambda_1(t)$, twierdzenia Tonellego
i lematu Fatou otrzymujemy
\[
\int_X f\,\mathrm d\mu
\le\int_{[0,M]}\liminf_n\mu_n(G_t)\,\mathrm d\lambda_1(t)
\le\liminf_n\int_Xf\,\mathrm d\mu_n.
\]
Jeśli $f\ge a$ i nie jest ograniczona z góry,
stosujemy poprzedni argument do
$f_k=\min\{f-a,k\}$, a następnie Beppo Leviego;
odjęta stała $a$ ma ten sam całkowy wkład dla
wszystkich miar probabilistycznych.

\emph{(ii)$\Rightarrow$(i).}
Dla dowolnej $h\in C_b(X)$ zarówno $h$, jak i $-h$
są dolnie półciągłe i ograniczone z dołu. Stąd
$\int h\,\mathrm d\mu\le\liminf_n\int h\,\mathrm d\mu_n$
oraz $-\int h\,\mathrm d\mu\le
-\limsup_n\int h\,\mathrm d\mu_n$.
Obie nierówności wymuszają zbieżność całek.

\emph{(iii), (iv)$\Rightarrow$(v).}
Dla $A$ spełniającego $\mu(\partial A)=0$ mamy
$A^\circ\subseteq A\subseteq\overline A$ oraz
$\mu(A^\circ)=\mu(A)=\mu(\overline A)$.
Wobec tego
\[
\mu(A)\le\liminf_n\mu_n(A^\circ)
\le\liminf_n\mu_n(A)
\le\limsup_n\mu_n(A)
\le\limsup_n\mu_n(\overline A)\le\mu(A).
\]

\emph{(v)$\Rightarrow$(i).}
Niech $h\in C_b(X)$, $\|h\|_\infty\le M$.
Miara obrazowa $h_*\mu$ na $[-M,M]$ ma co najwyżej
przeliczalnie wiele atomów: dla każdego $k$ liczba
atomów o masie co najmniej $1/k$ jest skończona.
Zatem można wybrać podział
$-M=t_0<\cdots<t_N=M$ o dostatecznie małych oczkach,
którego punkty wewnętrzne są nieatomowe dla $h_*\mu$.
Aby uniknąć atomów w końcach, wolno wybrać
$t_0<-M$ i $t_N>M$ oraz wymagać, by również one
były nieatomowe. Dla przedziałów $(t_{j-1},t_j]$
zbiory $A_j=\{t_{j-1}<h\le t_j\}$ mają
$\partial A_j\subseteq\{h=t_{j-1}\}\cup\{h=t_j\}$,
więc są zbiorami ciągłości miary. Z (v) wynika
$\mu_n(A_j)\to\mu(A_j)$. Aproksymując $h$
na każdym $A_j$ przez stałą wybraną z przedziału,
z błędem co najwyżej długość przedziału, otrzymujemy
\[
 \limsup_n\left|\int_Xh\,\mathrm d\mu_n-
 \int_Xh\,\mathrm d\mu\right|
 \le2\max_j(t_j-t_{j-1}).
\]
Wybierając dowolnie drobne podziały kończymy dowód.
\end{proof}

\begin{przyklad}[Dlaczego przedziały nie zawsze są zbiorami ciągłości?]
Dla $\mu_n=\delta_{1/n}$ i $\mu=\delta_0$ na
$\mathbb R$ mamy $\mu_n\Rightarrow\mu$.
Tymczasem $\mu_n(\{0\})=0$ i $\mu(\{0\})=1$.
Brzeg zbioru $\{0\}$ ma dodatnią miarę $\mu$,
więc Portmanteau nie gwarantuje zbieżności
wartości miar tego zbioru.
\end{przyklad}

\begin{wniosek}[Kiedy zbieżność mglista jest słabą-*?]
\label{cor:vague-c0-bounded-22}
Niech $X$ będzie lokalnie zwartą ośrodkową przestrzenią
metryczną i niech $\mu_n,\mu$ będą skończonymi dodatnimi
miarami Radona. Jeśli $\sup_n\mu_n(X)<\infty$, to
zbieżność mglista $\mu_n\to\mu$ jest równoważna zbieżności
$\sigma(C_0(X)^*,C_0(X))$.
Bez ograniczenia mas implikacja od zbieżności mglistej
do słabej-* może zawieść.
\end{wniosek}
\begin{proof}
Zbieżność słaba-* daje mglistą, gdyż $C_c(X)\subset C_0(X)$.
Załóżmy zbieżność mglistą oraz $C:=\sup_n\mu_n(X)<\infty$.
Niech $f\in C_0(X)$ i $\varepsilon>0$. Z gęstości
$C_c(X)$ w $C_0(X)$ (lemat~\ref{lem:density-cc-c0-21})
wybieramy $g\in C_c(X)$ takie, że
$\|f-g\|_\infty<\varepsilon$. Wówczas
\[
 \left|\int_X f\,\mathrm d\mu_n-\int_X f\,\mathrm d\mu\right|
 \le\left|\int_Xg\,\mathrm d\mu_n-\int_Xg\,\mathrm d\mu\right|
 +(C+\mu(X))\varepsilon.
\]
Pierwszy składnik dąży do zera, a $\varepsilon$ jest dowolne.

Aby wykazać konieczność kontroli mas, weźmy
$X=\mathbb R$, $\mu_n=n\delta_n$ i $\mu=0$.
Dla każdego $g\in C_c(\mathbb R)$ zachodzi
$\int g\,\mathrm d\mu_n=0$ dla dostatecznie dużych $n$,
więc zachodzi zbieżność mglista. Funkcja
$f(x)=(1+|x|)^{-1/2}$ należy do $C_0(\mathbb R)$,
a $\int f\,\mathrm d\mu_n=n/\sqrt{1+n}\to+\infty$.
Zbieżności słabej-* nie ma.
\end{proof}

\section{Zbieżność względem odległości w wariacji całkowitej}
\begin{twierdzenie}[Zbieżność w wariacji implikuje wąską]
\label{thm:tv-implies-narrow-22}
Jeśli $\mu_n,\mu\in\mathcal P(X)$ i
$d_{\mathrm{TV}}(\mu_n,\mu)\to0$, to
$\mu_n\Rightarrow\mu$.
\end{twierdzenie}
\begin{proof}
Dla dowolnej $h\in C_b(X)$, korzystając z normy
wariacji miary ze znakiem, mamy
\[
\left|\int_Xh\,\mathrm d\mu_n-\int_Xh\,\mathrm d\mu\right|
\le \|h\|_\infty\|\mu_n-\mu\|_{\mathrm{TV}}
=2\|h\|_\infty d_{\mathrm{TV}}(\mu_n,\mu)\to0.
\]
\end{proof}

\begin{uwaga}[Brak implikacji odwrotnej]
Przykład $\delta_{1/n}\Rightarrow\delta_0$ pokazuje,
że wąska zbieżność nie implikuje zbieżności w wariacji:
$d_{\mathrm{TV}}(\delta_{1/n},\delta_0)=1$.
\end{uwaga}

\section{Ciasność rodzin miar}
\begin{definicja}[Jednostajna ciasność]
\label{def:uniform-tight-22}
Niech $\mathcal F\subseteq\mathcal P(X)$.
Rodzina $\mathcal F$ jest \emph{jednostajnie ciasna},
jeżeli dla każdego $\varepsilon>0$ istnieje
zwarty zbiór $K\subseteq X$ spełniający
$\inf_{\mu\in\mathcal F}\mu(K)>1-\varepsilon$.
\end{definicja}

\begin{lemat}[Przybliżanie lipschitzowskie i przedłużanie]
\label{lem:lipschitz-extend-22}
Niech $(Y,D)$ będzie przestrzenią metryczną, $K\subseteq Y$
niepustym zbiorem zwartym, a $h:K\to\mathbb R$ funkcją
ciągłą z $\|h\|_\infty\le M$. Dla każdego $\eta>0$
istnieje ograniczona funkcja Lipschitzowska
$g:Y\to\mathbb R$ spełniająca
$\|g\|_\infty\le M$ oraz
$\sup_{x\in K}|g(x)-h(x)|<\eta$.
\end{lemat}
\begin{proof}
Funkcja $h$ jest jednostajnie ciągła na $K$.
Dla $L>0$ definiujemy
$h_L(x)=\inf_{z\in K}(h(z)+L D(x,z))$, $x\in K$.
Z nierówności trójkąta wynika
$|h_L(x)-h_L(y)|\le L D(x,y)$; ponadto
$h_L\le h$ i $-M\le h_L\le M$.
Wybierzmy $\delta>0$ takie, że
$|h(z)-h(x)|<\eta/2$ dla $D(x,z)<\delta$.
Jeżeli $L>2M/\delta$, to dla $D(x,z)\ge\delta$
mamy $h(z)+LD(x,z)\ge M\ge h(x)$.
Dla $D(x,z)<\delta$ natomiast
$h(z)+LD(x,z)\ge h(x)-\eta/2$.
Zatem $0\le h-h_L\le\eta/2$ na $K$.

Funkcję $h_L$ przedłużamy na $Y$ wzorem
$G(y)=\inf_{z\in K}(h_L(z)+L D(y,z))$.
Ponownie $G$ jest $L$-lipschitzowska. Na $K$
mamy $G=h_L$: nierówność $G\le h_L$ wynika
z podstawienia $z=y$, a odwrotna z własności
Lipschitza funkcji $h_L$. Funkcja
$g(y)=\max\{-M,\min\{M,G(y)\}\}$ jest Lipschitzowska,
ma normę co najwyżej $M$ i na $K$ jest równa $h_L$.
Wobec tego $\sup_K|g-h|\le\eta/2<\eta$.
\end{proof}

\begin{lemat}[Zwartość prawdopodobieństw na przestrzeni zwartej]
\label{lem:compact-prob-22}
Jeśli $Y$ jest zwartą przestrzenią metryczną,
to każdy ciąg w $\mathcal P(Y)$ ma podciąg wąsko
zbieżny do pewnego elementu $\mathcal P(Y)$.
\end{lemat}
\begin{proof}
Przestrzeń $C(Y)$ jest separowalna w normie supremum.
Aby to udowodnić bez odwołania do twierdzenia
Stone'a--Weierstrassa, wybierzmy przeliczalny zbiór
gęsty $(y_j)$ i rozważmy przeliczalną rodzinę
funkcji postaci
$F(x)=\min_{1\le i\le N}\{q_i+L d(x,y_{j_i})\}$,
gdzie $N,j_i\in\mathbb N$, $q_i\in\mathbb Q$,
$L\in\mathbb Q_{>0}$.
Pokażemy, że jest ona gęsta. Dla dowolnej
$h\in C(Y)$ i $\varepsilon>0$ lemat
\ref{lem:lipschitz-extend-22} (z $K=Y$) daje
funkcję $L$-lipschitzowską $v$ z
$\|h-v\|_\infty<\varepsilon/3$.
Zwiększając $L$, możemy przyjąć go wymiernym.
Wybieramy skończoną $\delta$-siatkę
$y_{j_1},\ldots,y_{j_N}$ oraz wymierne $q_i$
bliskie $v(y_{j_i})$, tak by
$|q_i-v(y_{j_i})|<\varepsilon/6$ i
$2L\delta<\varepsilon/6$.
Z lipschitzowskości wynika
$F(x)\ge v(x)-\varepsilon/6$ oraz
$F(x)\le v(x)+\varepsilon/6+2L\delta$.
W konsekwencji $\|F-h\|_\infty<\varepsilon$.
Niech $(f_j)$ będzie taką gęstą
rodziną. Dla danego ciągu $(\mu_n)$ każdy ciąg liczb
$\bigl(\int_Y f_j\,\mathrm d\mu_n\bigr)_{n\ge1}$
jest ograniczony.
Metoda diagonalna dostarcza podciągu $(\mu_{n_k})$,
dla którego wszystkie te całki mają granice.
Ponieważ
$|\int h\,\mathrm d\mu_{n_k}-\int f_j\,\mathrm d\mu_{n_k}|
\le\|h-f_j\|_\infty$, te granice rozszerzają się
jednoznacznie na ciągły dodatni funkcjonał $L$ na
$C(Y)$ o normie $1$ i $L(\chi_Y)=1$.
Na mocy twierdzenia~\ref{thm:compact-riesz-21}
istnieje prawdopodobieństwo $\mu$ reprezentujące $L$.
Definicja $L$ daje $\mu_{n_k}\Rightarrow\mu$.
\end{proof}

\begin{lemat}[Włożenie separowalnej przestrzeni metrycznej
w zwartą przestrzeń metryczną]
\label{lem:hilbert-cube-22}
Dla każdej separowalnej przestrzeni metrycznej $(X,d)$
istnieje zwarta przestrzeń metryczna $Y$ i homeomorfizm
$\iota:X\to\iota(X)\subseteq Y$.
\end{lemat}
\begin{proof}
Niech $(x_j)$ będzie przeliczalnym zbiorem gęstym.
Definiujemy
\[
\iota(x)=\left(\frac{d(x,x_j)}{1+d(x,x_j)}\right)_{j\ge1}
\in[0,1]^{\mathbb N}.
\]
Przestrzeń kostki Hilberta $[0,1]^{\mathbb N}$ z metryką
$D(a,b)=\sum_{j\ge1}2^{-j}|a_j-b_j|$ jest zwarta:
metoda diagonalna zapewnia zbieżny podciąg współrzędnych,
a oszacowanie ogona sumy daje zbieżność w $D$.
Każda współrzędna $\iota$ jest ciągła, więc
$\iota$ jest ciągła. Jeśli $\iota(x)=\iota(y)$,
to $d(x,x_j)=d(y,x_j)$ dla wszystkich $j$;
przez gęstość wynika $x=y$.
Aby wykazać ciągłość odwrotności na obrazie, załóżmy,
że $\iota(x_n)\to\iota(x)$, i ustalmy $\varepsilon>0$.
Wybierzmy $x_j$ z $d(x,x_j)<\varepsilon/3$.
Dla dużych $n$ współrzędna $j$ daje
$|d(x_n,x_j)-d(x,x_j)|<\varepsilon/3$,
zatem $d(x_n,x)\le d(x_n,x_j)+d(x_j,x)<\varepsilon$.
Wystarczy przyjąć $Y=\overline{\iota(X)}$ w kostce.
\end{proof}

\begin{twierdzenie}[Prochorow: kierunek ciasność $\Rightarrow$ zwartość]
\label{thm:prokhorov-forward-22}
Niech $X$ będzie separowalną przestrzenią metryczną,
a $\mathcal F\subseteq\mathcal P(X)$ rodziną jednostajnie
ciasną. Wtedy każdy ciąg w $\mathcal F$ ma podciąg
wąsko zbieżny do prawdopodobieństwa na $X$.
\end{twierdzenie}
\begin{proof}
Niech $(\mu_n)\subseteq\mathcal F$.
Z jednostajnej ciasności wybieramy rosnący ciąg
zwartych zbiorów $K_m\subseteq X$ taki, że
\begin{equation}\label{eq:tight-compacts-22}
 \inf_n\mu_n(K_m)\ge1-2^{-m},\qquad m\ge1.
\end{equation}
Stosujemy włożenie $\iota:X\to Y$ z lematu
\ref{lem:hilbert-cube-22}. Miary
$\nu_n=\iota_*\mu_n$ są probabilistyczne na
zwartej metrycznej $Y$. Z lematu
\ref{lem:compact-prob-22} wybieramy podciąg,
ponownie oznaczany $(\nu_n)$, wąsko zbieżny
do $\nu\in\mathcal P(Y)$.
Dla każdego $m$ zbiór $\iota(K_m)$ jest zwarty,
a zatem domknięty w $Y$. Z Portmanteau
\[
\nu(\iota(K_m))\ge\limsup_n\nu_n(\iota(K_m))
\ge1-2^{-m}.
\]
Wobec rosnącości $K_m$ zbiór
$C=\bigcup_m\iota(K_m)$ jest borelowski w $Y$,
a $\nu(C)=1$.

Definiujemy teraz prawdopodobieństwo $\mu$ na $X$.
Dla $A\in\mathcal B(X)$ kładziemy
\begin{equation}\label{eq:pullback-limit-22}
\mu(A)=\nu\bigl(\iota(A\cap\bigcup_mK_m)\bigr).
\end{equation}
Ta definicja jest poprawna: $A\cap K_m$ jest
borelowski w zwartej przestrzeni $K_m$,
a $\iota|_{K_m}$ jest homeomorfizmem na zwarty
podzbiór $Y$. Zatem $\iota(A\cap K_m)$ jest
borelowski w $Y$, podobnie jak ich przeliczalna suma.
Iniektywność $\iota$ i addytywność $\nu$ dowodzą,
że $\mu$ jest miarą; ponadto $\mu(X)=\nu(C)=1$
i $\mu(K_m)=\nu(\iota(K_m))$.

Niech $f\in C_b(X)$, $M=\|f\|_\infty$.
Dla ustalonego $m$ funkcja $f\circ\iota^{-1}$
jest ciągła na zwartym $\iota(K_m)$.
Lemat~\ref{lem:lipschitz-extend-22}, zastosowany
z dowolnym $\eta>0$, daje funkcję
$g_{m,\eta}\in C_b(Y)$, $\|g_{m,\eta}\|_\infty\le M$,
która na $\iota(K_m)$ różni się od
$f\circ\iota^{-1}$ o mniej niż $\eta$.
Stąd, z \eqref{eq:tight-compacts-22},
\[
\left|\int_X f\,\mathrm d\mu_n-
\int_Y g_{m,\eta}\,\mathrm d\nu_n\right|\le2M2^{-m}+\eta,
\]
a przez analogiczny rachunek
$|\int_Xf\,\mathrm d\mu-\int_Yg_{m,\eta}\,\mathrm d\nu|
\le2M2^{-m}+\eta$.
Dla ustalonych $m,\eta$ granica całek z $g_{m,\eta}$ wynika
ze zbieżności wąskiej na $Y$; dlatego
\[
\limsup_n\left|\int_X f\,\mathrm d\mu_n-
\int_X f\,\mathrm d\mu\right|\le4M2^{-m}+2\eta.
\]
Przechodząc najpierw $\eta\to0$, a następnie
$m\to\infty$, kończymy dowód.
\end{proof}

\begin{lemat}[Jednostajna ciasność ciągu wąsko zbieżnego]
\label{lem:weak-sequence-tight-22}
Jeśli $X$ jest polską przestrzenią metryczną,
$\mu_n\Rightarrow\mu$ i $\mu_n,\mu\in\mathcal P(X)$,
to rodzina $\{\mu,\mu_1,\mu_2,\ldots\}$ jest
jednostajnie ciasna.
\end{lemat}
\begin{proof}
Niech $d$ będzie zupełną, zgodną metryką, a $(x_j)$
przeliczalnym zbiorem gęstym. Ustalmy $\varepsilon>0$.
Dla każdego $m$ przeliczalna rodzina kul otwartych
$B(x_j,2^{-m})$ pokrywa $X$. Istnieje skończona
suma tych kul $O_m$, dla której
$\mu(O_m)>1-\varepsilon2^{-m-2}$.
Z otwartej wersji Portmanteau wynika, że
$\mu_n(O_m)>1-\varepsilon2^{-m-1}$ dla wszystkich
$n$ od pewnego $N_m$.
Powiększając skończoną rodzinę kul w $O_m$,
możemy zapewnić tę samą nierówność dla
$\mu_1,\ldots,\mu_{N_m-1}$; pozostaje prawdziwa
także dla wszystkich późniejszych miar.
Niech $F_m$ będzie sumą domkniętych kul
$\overline B(x_j,2^{-m})$ o środkach użytych w $O_m$.
Kładziemy $K=\bigcap_mF_m$. Zbiór $K$ jest domknięty,
a dla każdego $m$ mieści się w skończonej sumie kul
o promieniu $2^{-m}$, jest zatem całkowicie ograniczony.
Zupełność $X$ daje zwartość $K$. Ponadto
$\mu_n(X\setminus K)\le\sum_{m\ge1}\mu_n(X\setminus F_m)
<\sum_{m\ge1}\varepsilon2^{-m-1}<\varepsilon$
dla wszystkich $n$. Podobne oszacowanie zachodzi dla $\mu$.
\end{proof}

\begin{twierdzenie}[Prochorow na przestrzeni polskiej]
\label{thm:prokhorov-22}
Niech $X$ będzie przestrzenią polską, a
$\mathcal F\subseteq\mathcal P(X)$.
Następujące warunki są równoważne:
\begin{enumerate}[label=(\roman*)]
\item $\mathcal F$ jest jednostajnie ciasna;
\item każdy ciąg w $\mathcal F$ ma podciąg wąsko
zbieżny do pewnego prawdopodobieństwa na $X$.
\end{enumerate}
Warunek (ii) oznacza względną zwartość
ciągową dla topologii zbieżności wąskiej.
\end{twierdzenie}
\begin{proof}
Implikację (i)$\Rightarrow$(ii) udowodniliśmy w twierdzeniu
\ref{thm:prokhorov-forward-22}. Dla odwrotnej
ustalmy $\varepsilon>0$, zgodną metrykę zupełną $d$
i przeliczalny zbiór gęsty $\{x_j:j\ge1\}$.
Dla każdego $m\ge1$ oznaczmy
$V_{m,k}=\bigcup_{j=1}^k B(x_j,2^{-m})$.
Dla ustalonego $m$ zbiory $V_{m,k}$ rosną do $X$.
Wykażemy, że istnieje $k_m$ takie, że
\begin{equation}\label{eq:prokhorov-uniform-balls-22}
\inf_{\rho\in\mathcal F}\rho(V_{m,k_m})
>1-\varepsilon2^{-m-1}.
\end{equation}
Gdyby tak nie było, dla każdego $k$ wybralibyśmy
$\rho_k\in\mathcal F$ z
$\rho_k(V_{m,k})\le1-\varepsilon2^{-m-1}$.
Z założenia (ii) ciąg ten ma podciąg
$\rho_{k_l}\Rightarrow\rho$ z granicą probabilistyczną.
Ponieważ $V_{m,k}\uparrow X$, istnieje $k_0$ takie,
że $\rho(V_{m,k_0})>1-\varepsilon2^{-m-1}$.
Portmanteau dla otwartego $V_{m,k_0}$ daje
$\rho_{k_l}(V_{m,k_0})>1-\varepsilon2^{-m-1}$
dla dostatecznie dużych $l$. Ponieważ $k_l\ge k_0$,
monotoniczność wymusza wówczas
$\rho_{k_l}(V_{m,k_l})>1-\varepsilon2^{-m-1}$,
sprzecznie z konstrukcją $\rho_{k_l}$.
Dowiedliśmy \eqref{eq:prokhorov-uniform-balls-22}
bez odwoływania się do metryzowalności
przestrzeni prawdopodobieństw.

Niech $F_m=\bigcup_{j\le k_m}\overline B(x_j,2^{-m})$
i $K=\bigcap_{m\ge1}F_m$. Zbiór $K$ jest domknięty,
a dla każdego $m$ mieści się w skończonej sumie kul
o promieniu $2^{-m}$, więc jest całkowicie ograniczony.
Zupełność $X$ zapewnia zwartość $K$.
Z \eqref{eq:prokhorov-uniform-balls-22} dla wszystkich
$\rho\in\mathcal F$ uzyskujemy
\[
 \rho(X\setminus K)\le\sum_{m\ge1}\rho(X\setminus F_m)
 <\sum_{m\ge1}\varepsilon2^{-m-1}<\varepsilon.
\]
Jest to jednostajna ciasność $\mathcal F$.
\end{proof}

\begin{uwaga}[Metryzowalność a dowód Prochorowa]
Na przestrzeni polskiej topologia zbieżności wąskiej
na $\mathcal P(X)$ jest metryzowalna, na przykład przez
odległość opartą na ograniczonych funkcjach Lipschitzowskich.
Nie było jednak potrzeby przyjmowania tego wyniku
w dowodzie twierdzenia~\ref{thm:prokhorov-22}:
własność wyrażona przez podciągi i argument ze
wzrastającymi skończonymi pokryciami kulowymi
dały bezpośrednią konstrukcję wspólnego zwartego zbioru.
\end{uwaga}

\begin{wniosek}[Ciasność z oszacowania momentów]
\label{cor:coercive-tight-22}
Niech $X$ będzie przestrzenią polską i
$V:X\to[0,+\infty]$ funkcją dolnie półciągłą,
której zbiory podpoziomicowe
$K_R=\{x\in X:V(x)\le R\}$ są zwarte dla każdego
$R<\infty$. Jeśli rodzina $\mathcal F\subseteq\mathcal P(X)$
spełnia
$\sup_{\rho\in\mathcal F}\int_X V(x)\,\mathrm d\rho(x)
\le C<\infty$, to jest jednostajnie ciasna.
W szczególności na $\mathbb R^d$ wystarcza jednostajne
ograniczenie dodatniego momentu rzędu $p>0$.
\end{wniosek}
\begin{proof}
Dla $R>0$ nierówność Czebyszewa daje
\[
 \rho(X\setminus K_R)=\rho(\{V>R\})
 \le R^{-1}\int_X V(x)\,\mathrm d\rho(x)\le C/R.
\]
Wybierając $R>C/\varepsilon$ otrzymujemy
$\inf_{\rho\in\mathcal F}\rho(K_R)>1-\varepsilon$.
Na $\mathbb R^d$ funkcja $V(x)=|x|^p$ jest ciągła,
a jej zbiory podpoziomicowe są zwarte.
\end{proof}

\begin{przyklad}[Ograniczony moment nie musi być zbieżny]
\label{ex:moment-not-continuous-22}
Niech $X=\mathbb R$, a dla $n\ge1$ połóżmy
\[
 \mu_n=\left(1-\frac1n\right)\delta_0+
 \frac1n\delta_n,\qquad \mu=\delta_0.
\]
Ponieważ $\mu_n-\mu=n^{-1}(\delta_n-\delta_0)$,
mamy $d_{\mathrm{TV}}(\mu_n,\mu)=1/n\to0$.
W szczególności $\mu_n\Rightarrow\mu$, a rodzina
$\{\mu_n:n\ge1\}$ jest jednostajnie ciasna.
Mimo to
\[
 \int_{\mathbb R}|x|\,\mathrm d\mu_n(x)=1
 \quad\text{dla każdego }n,\qquad
 \int_{\mathbb R}|x|\,\mathrm d\mu(x)=0.
\]
Wobec tego jednostajne ograniczenie pierwszego momentu
wystarcza do ciasności, ale nie do zbieżności wartości
tego momentu, nawet gdy zachodzi zbieżność w wariacji
całkowitej. Funkcja $x\mapsto|x|$ jest nieograniczona,
więc nie jest dopuszczalną funkcją testową w definicji
zbieżności wąskiej.
\end{przyklad}

\section{Następstwa i kryteria praktyczne}
\begin{wniosek}[Zbieżność mglista bez ucieczki masy]
\label{cor:vague-narrow-22}
Niech $X$ będzie lokalnie zwartą ośrodkową
przestrzenią metryczną. Jeśli $\mu_n,\mu\in\mathcal P(X)$,
$\mu_n$ zbiega mgliście do $\mu$, to
$\mu_n\Rightarrow\mu$.
\end{wniosek}
\begin{proof}
Ustalmy $\varepsilon>0$. Z rozdziału
\ref{ch:radon-measures-20} miara $\mu$ jest ciasna.
Wybieramy zwarty $K$ o $\mu(K)>1-\varepsilon$
oraz $\chi\in C_c(X)$ z $0\le\chi\le1$,
$\chi|_K=1$. Wówczas
$\int_X\chi\,\mathrm d\mu>1-\varepsilon$,
a ze zbieżności mglistej
$\int_X\chi\,\mathrm d\mu_n>1-2\varepsilon$
dla dużych $n$. Dla $h\in C_b(X)$
funkcja $h\chi\in C_c(X)$, a więc
\[
\left|\int_Xh\,\mathrm d\mu_n-
\int_Xh\chi\,\mathrm d\mu_n\right|
\le\|h\|_\infty\int_X(1-\chi)\,\mathrm d\mu_n
<2\varepsilon\|h\|_\infty.
\]
Analogiczne oszacowanie zachodzi dla $\mu$,
z błędem $\varepsilon\|h\|_\infty$.
Zbieżność całek z $h\chi$ i dowolność $\varepsilon$
kończą dowód.
\end{proof}

\begin{przyklad}[Rozkłady jednostajnie skupione na zwartym przedziale]
Niech $X=\mathbb R$, a $\mu_n$ dowolnymi
prawdopodobieństwami spełniającymi
$\mu_n([-R,R])=1$ dla stałego $R<\infty$.
Rodzina $\{\mu_n\}$ jest jednostajnie ciasna.
Z twierdzenia Prochorowa każdy jej ciąg ma podciąg
wąsko zbieżny. Nie wynika stąd zbieżność całego ciągu:
na przykład miary $\delta_{-1},\delta_1,\delta_{-1},\ldots$
mają dwa różne punkty skupienia.
\end{przyklad}

\begin{przyklad}[Brak ciasności]
Rodzina $\{\delta_n:n\ge1\}\subseteq\mathcal P(\mathbb R)$
nie jest jednostajnie ciasna: dla każdego zwartego
$K\subseteq\mathbb R$ pewne $n$ nie należy do $K$.
Nie istnieje podciąg zbieżny wąsko do prawdopodobieństwa.
\end{przyklad}

\begin{uwaga}[Gdzie pojawi się ta teoria?]
Zbieżność wąska, ciasność i twierdzenie Prochorowa
stanowią narzędzia w teorii procesów stochastycznych,
rachunku prawdopodobieństwa, analizie wariacyjnej
oraz w zagadnieniach transportu optymalnego.
Dalszy rozdział skryptu poświęcimy mierze Haara
na grupach lokalnie zwartych, gdzie niezmienniczość
miary dostarczy kolejnego sposobu jej scharakteryzowania.
\end{uwaga}

% ROZDZIAŁ 23 — MIARA HAARA
\chapter{Miara Haara na grupach lokalnie zwartych}
\label{ch:haar-23}

\providecommand{\atlasref}[2]{#2}


\section{Grupy topologiczne, translacje i miary niezmiennicze}

\noindent\emph{Powiązania tematyczne w Bibliotece matematycznej:} \atlasref{GT-24::def:topog24-group}{GT-24: grupy topologiczne}.

\begin{definicja}[Grupa topologiczna]
\label{def:topological-group-23}
Grupę $G$, wyposażoną w topologię Hausdorffa, nazywamy
\emph{grupą topologiczną}, gdy odwzorowania $(x,y)\mapsto xy$
oraz $x\mapsto x^{-1}$ są ciągłe. Element neutralny oznaczamy $e$.
Grupa jest \emph{lokalnie zwarta}, gdy każdy punkt ma otoczenie
o zwartej domkniętej powłoce. Piszemy $xA=\{xa:a\in A\}$
i $Ax=\{ax:a\in A\}$.
\end{definicja}

\begin{uwaga}[Dlaczego lokalna zwartość?]
W przestrzeni $\mathbb R^d$ miara Lebesgue'a jest niezmiennicza
względem przesunięć. Dla grup nieprzemiennych należy rozróżniać
translacje \emph{lewe} i \emph{prawe}. Zwartość zbioru pozwala pokryć go skończenie wieloma
translacjami dowolnego otwartego otoczenia $e$.
Lokalna zwartość grupy zapewnia ponadto istnienie dowolnie
małych otoczeń o zwartym domknięciu; obie własności
wykorzystamy w konstrukcji.
\end{uwaga}

\begin{definicja}[Lewostronna i prawostronna miara Haara]
\label{def:haar-23}
Niech $G$ będzie lokalnie zwartą grupą Hausdorffa.
\emph{Lewostronną miarą Haara} nazywamy niezerową miarę Radona
$\mu$ na $\mathcal B(G)$, w sensie definicji~\ref{def:radon-measure-20},
dla której $\mu(gE)=\mu(E)$ dla wszystkich $g\in G$
i $E\in\mathcal B(G)$. Analogicznie definiujemy
\emph{prawostronną miarę Haara} przez warunek
$\mu(Eg)=\mu(E)$. Zakres wymaganej regularności
przypominamy poniżej ze względu na brak założeń przeliczalności.
\end{definicja}

\begin{twierdzenie}[Haara: istnienie i jednoznaczność]
\label{thm:haar-main-23}
Na każdej lokalnie zwartej grupie Hausdorffa istnieje
lewostronna miara Haara. Każde dwie takie miary różnią się
dodatnim czynnikiem stałym. To samo dotyczy prawostronnych miar
Haara.
\end{twierdzenie}

\begin{uwaga}[Zakres twierdzenia i regularność]
W twierdzeniu~\ref{thm:haar-main-23} \emph{nie} zakładamy ani
metryzowalności, ani $\sigma$-zwartości, ani drugiego aksjomatu przeliczalności.
W ogólnej przestrzeni lokalnie zwartej Hausdorffa rozróżniamy
trzy warunki dla miary borelowskiej $\mu$:
\begin{enumerate}[label=(\roman*)]
\item $\mu(K)<\infty$ dla każdego kompaktu $K$;
\item $\mu(E)=\inf\{\mu(U):E\subseteq U,\ U\text{ otwarty}\}$
      dla każdego zbioru borelowskiego $E$;
\item $\mu(U)=\sup\{\mu(K):K\subseteq U,\ K\text{ zwarty}\}$
      dla każdego zbioru otwartego $U$.
\end{enumerate}
To właśnie te trzy własności przyjmujemy w definicji miary Radona
na dowolnej grupie lokalnie zwartej. Nie wymagamy regularności
wewnętrznej na \emph{każdym} zbiorze borelowskim: bez
$\sigma$-zwartości grupy jest to założenie zbyt silne.
Dla borelowskich zbiorów o mierze $\sigma$-skończonej
regularność wewnętrzna wyniknie z (i)--(iii).
Wyjaśnimy różnicę na konkretnym przykładzie
\ref{ex:haar-non-sigma-compact-23}.
\end{uwaga}

\section{Lematy topologiczne poprzedzające konstrukcję}

\noindent\emph{Powiązania tematyczne w Bibliotece matematycznej:} \atlasref{GT-14::def:topog-baire-lokalna-zwartosc}{GT-14: lokalna zwartość}.

\begin{lemat}[Małe otoczenia, zwarte rozdrobnienia i pokrycia]
\label{lem:haar-neigh-23}
Niech $G$ będzie lokalnie zwartą grupą Hausdorffa.
\begin{enumerate}[label=(\roman*)]
\item Dla każdego otoczenia $O$ elementu $e$ istnieje symetryczne,
  relatywnie zwarte otoczenie otwarte $V$ takie, że
  $V^{-1}V\subseteq O$ oraz $\overline V\subseteq O$.
  Relatywnie zwarte otoczenia otwarte
  tworzą bazę otoczeń $e$.
\item Dla każdego zbioru zwartego $K\subseteq G$ i otwartego otoczenia
  $V$ punktu $e$ istnieje skończone pokrycie $K$
  przez lewe przesunięcia $V$.
\item Jeśli $K,L$ są rozłącznymi zbiorami zwartymi, to istnieje otoczenie
  $V$ punktu $e$ takie, że żaden zbiór $xV$, $x\in G$,
  nie przecina jednocześnie $K$ i $L$.
\item W dowolnej lokalnie zwartej przestrzeni Hausdorffa
  $X$ (niekoniecznie grupie), jeśli
  $K\subseteq\bigcup_{j=1}^r O_j$ jest zwarty,
  gdzie $O_j$ są otwarte w $X$, to istnieją zbiory zwarte
  $K_j\subseteq O_j$ spełniające $K=\bigcup_{j=1}^rK_j$.
\end{enumerate}
\end{lemat}
\begin{proof}
Dla (i) wybieramy otwarte otoczenie $O'$ elementu $e$
zawarte w $O$. Z ciągłości mnożenia w $(e,e)$
istnieją otwarte otoczenia $A,B$ elementu $e$
takie, że $AB\subseteq O'$. Kładziemy
$W=A\cap B\cap A^{-1}\cap B^{-1}$; jest ono
otwarte i symetryczne oraz $W^2\subseteq AB\subseteq O$.
Z lokalnej zwartości istnieje zwarty $C$ z
$e\in\operatorname{int}C$. Kładziemy
$V=W\cap\operatorname{int}C\cap(\operatorname{int}C)^{-1}$.
Wówczas $V$ jest otwarty i symetryczny,
$\overline V\subseteq C$ oraz
$V^{-1}V\subseteq W^2\subseteq O$.
Ponadto $\overline W\subseteq W^2$: jeśli $x\in\overline W$,
to otwarty zbiór $xW$ przecina $W$, a więc
$x=uw^{-1}$ dla pewnych $u,w\in W$.
Stąd także $\overline V\subseteq O$.
Dobór $O$ wewnątrz dowolnego zadanego otoczenia
daje bazę relatywnie zwartych otoczeń.

W (ii) rodzina $(xV)_{x\in K}$ jest pokryciem otwartym
zbioru zwartego $K$, skąd istnieje podpokrycie skończone.
Dla (iii) zbiór $K^{-1}L$ jest zwarty jako obraz
$K\times L$ przez ciągłe odwzorowanie
$(k,l)\mapsto k^{-1}l$, a więc jest domknięty.
Nie zawiera $e$, gdyż $K\cap L=\varnothing$;
stosujemy (i) do otwartego otoczenia
$G\setminus K^{-1}L$. Gdyby $xv\in K$ i $xw\in L$
przy $v,w\in V$, otrzymalibyśmy
$v^{-1}w\in K^{-1}L\cap V^{-1}V$, sprzeczność.

W (iv) wykorzystamy następujący fakt o każdej przestrzeni
lokalnie zwartej Hausdorffa $X$: jeśli $x\in O$ i $O$
jest otwarty, to istnieje otwarty $W\ni x$ ze
zwartym $\overline W\subseteq O$. Dla dowodu wybieramy
zwartą okolicę $C$ punktu $x$, z $x\in\operatorname{int}C$.
Zbiór $F=C\setminus O$ jest zwartym podzbiorem
przestrzeni Hausdorffa $C$. Dla każdego $y\in F$
rozdzielamy $x$ i $y$ rozłącznymi otoczeniami w $C$.
Zwartość $F$ pozwala wybrać skończenie wiele z nich;
przecięcie otoczeń $x$ daje otwarty $Z\subseteq C$
taki, że $x\in Z$ i $\overline Z^{\,C}\cap F=\varnothing$.
Bierzemy otwarty $\widetilde Z\subseteq X$ z
$\widetilde Z\cap C=Z$ i definiujemy
$W=\widetilde Z\cap\operatorname{int}C$.
Wtedy $\overline W\subseteq C\cap\overline Z^{\,C}
\subseteq O$ oraz $\overline W$ jest zwarty.

Przypadek pustego $K$ uzyskujemy, przyjmując wszystkie $K_j=\varnothing$.

Dla każdego $x\in K$ wybieramy zatem indeks $j(x)$
oraz takie otoczenie $W_x$, aby
$x\in W_x\subseteq\overline W_x\subseteq O_{j(x)}$.
Wybierzmy $x_1,\dots,x_s$ tak, aby zbiory $W_{x_i}$
pokrywały $K$, i połóżmy
$K_j=K\cap\bigcup_{i:\,j(x_i)=j}\overline W_{x_i}$.
Wówczas każde $K_j$ jest zwarte, gdyż jest domkniętym
podzbiorem $K$, i zawiera się w $O_j$.
Każdy punkt $K$ należy do pewnego $W_{x_i}$, a więc
do $K_{j(x_i)}$, co dowodzi równości
$K=\bigcup_j K_j$.
\end{proof}

\begin{lemat}[Funkcja odcinająca bez metryzowalności]
\label{lem:haar-cutoff-23}
Niech $X$ będzie lokalnie zwartą przestrzenią Hausdorffa,
$K\subseteq X$ zbiorem zwartym i $U\supseteq K$ zbiorem
otwartym. Istnieje $h\in C_c(X)$ taka, że
$0\le h\le1$, $h|_K=1$ i $\operatorname{supp}h\subseteq U$.
\end{lemat}
\begin{proof}
Dla $K=\varnothing$ bierzemy $h=0$. W przeciwnym razie
z argumentu topologicznego w lemacie~\ref{lem:haar-neigh-23}(iv)
i skończonego pokrycia $K$ wybieramy otwarty zbiór $W$
taki, że $K\subseteq W\subseteq C:=\overline W\subseteq U$,
gdzie $C$ jest zwarty. Zwarta przestrzeń Hausdorffa $C$
jest normalna. Z lematu Urysohna dla rozłącznych domkniętych
podzbiorów $K$ i $C\setminus W$ otrzymujemy ciągłą
$h_C:C\to[0,1]$ równą $1$ na $K$ i $0$ na $C\setminus W$.
Na domkniętym zbiorze $X\setminus W$ kładziemy $h=0$,
a na $C$ kładziemy $h=h_C$. Oba przepisy są zgodne
na przecięciu, zatem lemat o sklejaniu daje ciągłość $h$.
Jej nośnik zawiera się w $C\subseteq U$.
\end{proof}

\section{Zwarta przestrzeń danych bez założenia przeliczalności}
\begin{lemat}[Zwartość dowolnego produktu przedziałów]
\label{lem:haar-product-compact-23}
Niech $I$ będzie dowolnym zbiorem indeksów, a
$a_i\in[0,\infty)$ dla $i\in I$. W topologii produktowej
przestrzeń $P=\prod_{i\in I}[0,a_i]$ jest zwarta.
\end{lemat}
\begin{proof}
Dla zupełności podajemy argument ultrafiltrowy, nie wymagający
uprzedniego użycia twierdzenia Tichonowa. Korzystamy jedynie
z lematu Kuratowskiego--Zorna, czyli ze standardowej postaci
aksjomatu wyboru. Na niepustym zbiorze $S$, jeśli rodzina
zbiorów ma własność skończonych przecięć, to jej skończone
przecięcia (włącznie z $S$ jako pustym przecięciem) tworzą
bazę filtru właściwego. Rodzina filtrów właściwych ją
zawierających jest więc niepusta. Każdy niepusty łańcuch
tych filtrów ma ograniczenie górne, którym jest ich suma
zbiorowa: dowolne dwa jej elementy należą do wspólnego
filtru z łańcucha, więc ich przecięcie również należy do
sumy; domknięcie na nadzbiory jest zachowane i zbiór
pusty do sumy nie należy. Dla pustego łańcucha
wystarczy filtr wygenerowany przez rodzinę wyjściową. Lemat Zorna dostarcza filtru
maksymalnego $\mathcal U$. Jest to ultrafiltr: gdy $A\notin\mathcal U$, maksymalność oznacza, że
rodzina $\mathcal U\cup\{A\}$ nie ma własności
skończonych przecięć (w przeciwnym razie generowałaby
ściśle większy filtr właściwy). Istnieje więc skończone
przecięcie $B$ elementów $\mathcal U$ takie, że
$B\cap A=\varnothing$. Mamy $B\in\mathcal U$, a stąd
$S\setminus A\in\mathcal U$ przez domknięcie na nadzbiory.

Aby dowieść zwartości $P$, niech $(F_\alpha)_{\alpha\in J}$
będzie rodziną domkniętych podzbiorów $P$ o własności
skończonych przecięć. Rozszerzamy ją do ultrafiltru
$\mathcal U$ na $P$ (produkt jest niepusty, gdyż zawiera
rodzinę samych zer; dla pustego $I$ jest jednopunktowy).
Dla każdego $i\in I$ definiujemy ultrafiltr obrazowy
\[
 \mathcal U_i=\{A\subseteq[0,a_i]:
                  \pi_i^{-1}(A)\in\mathcal U\}.
\]
Jest to filtr właściwy, ponieważ przeciwobrazy zachowują
inkluzje, przecięcia, zbiór pusty i cały zbiór. Dla
każdego $A\subseteq[0,a_i]$ dokładnie jeden z przeciwobrazów
$A$ i jego dopełnienia należy do $\mathcal U$,
więc $\mathcal U_i$ jest ultrafiltrem.
Domknięcia elementów $\mathcal U_i$ mają własność
skończonych przecięć, gdyż już same te elementy mają
niepuste skończone przecięcia. Ze zwartości przedziału
$C_i:=\bigcap_{A\in\mathcal U_i}\overline A$
jest niepustym zbiorem zwartym.
Wybieramy $t_i=\min C_i$ (nie trzeba dokonywać dodatkowego
wyboru punktów). Każde otwarte otoczenie $W$ punktu $t_i$
należy do $\mathcal U_i$: w przeciwnym razie
$[0,a_i]\setminus W$ należałoby do ultrafiltru,
co wykluczałoby $t_i\in C_i$.

Punkt $t=(t_i)_{i\in I}$ należy do domknięcia każdego
zbioru z $\mathcal U$. Istotnie, podstawowe otoczenie
$t$ ma postać skończonego przecięcia przeciwobrazów
otwartych otoczeń odpowiednich $t_i$. Każdy z tych
przeciwobrazów należy do $\mathcal U$ z definicji
$\mathcal U_i$, więc ich przecięcie również należy
do $\mathcal U$ i musi przecinać każdy jego element. W szczególności $t\in F_\alpha$ dla
każdego $\alpha$, ponieważ $F_\alpha$ jest domknięty
i należy do $\mathcal U$. Wobec tego każda rodzina
zbiorów domkniętych o własności skończonych przecięć
ma niepuste przecięcie, co dowodzi zwartości $P$.
\end{proof}

\section{Konstrukcja miary przez liczby pokrycia}
Niech odtąd $G$ będzie \emph{dowolną} lokalnie zwartą
grupą Hausdorffa. Wybieramy kompakt $K_0$ będący otoczeniem $e$
oraz otwarte otoczenie $V_0$ elementu $e$,
spełniające $V_0\subseteq K_0$. Przez $\mathcal N$
oznaczamy rodzinę wszystkich relatywnie zwartych otwartych
otoczeń $e$; jest ona skierowana przez odwrotną inkluzję.
Dla zwartego $K$ i $V\in\mathcal N$ definiujemy
\begin{equation}\label{eq:haar-cover-number-23}
 N(K,V)=\min\left\{r\in\mathbb N_0:\
 K\subseteq\bigcup_{j=1}^{r}x_jV
 \text{ dla pewnych }x_1,\dots,x_r\in G\right\}.
\end{equation}
W szczególności $N(\varnothing,V)=0$, a dla
$K\ne\varnothing$ zachodzi $1\le N(K,V)<\infty$.

\begin{lemat}[Jednolite oszacowania ilorazów]
\label{lem:haar-ratio-23}
Niech $C_K=N(K,V_0)$. Dla każdego $V\in\mathcal N$ mamy
\[
 0\le q_V(K):=\frac{N(K,V)}{N(K_0,V)}\le C_K,
 \qquad q_V(K_0)=1.
\]
Ponadto $q_V(gK)=q_V(K)$, funkcje $q_V$ są monotoniczne
i podaddytywne dla sum zwartych, zaś dla rozłącznych
zbiorów zwartych $K,L$ istnieje $W\in\mathcal N$ takie, że
$q_V(K\cup L)=q_V(K)+q_V(L)$ dla każdego $V\subseteq W$.
\end{lemat}
\begin{proof}
Zbiory $K$ można pokryć przez $C_K$ przesunięć $V_0$,
a każde $xV_0\subseteq xK_0$. Pokrywając $K_0$
przez $N(K_0,V)$ przesunięć $V$ i następnie przesuwając
to pokrycie, otrzymujemy $N(K,V)\le C_KN(K_0,V)$.
Ponieważ $K_0$ jest niepusty, mianownik jest dodatni.
Przesunięcie minimalnego pokrycia $K$ przez $g$ daje
$N(gK,V)\le N(K,V)$; przesunięcie przez $g^{-1}$
daje nierówność odwrotną. Jeżeli $K\subseteq L$,
każde pokrycie $L$ pokrywa $K$, więc
$N(K,V)\le N(L,V)$. Połączenie pokryć $K$ i $L$
daje $N(K\cup L,V)\le N(K,V)+N(L,V)$.
Dzieląc przez dodatnie $N(K_0,V)$, otrzymujemy
odpowiednie własności $q_V$.

Dla rozłącznych $K,L$ wybierzmy $W$ według
lematu~\ref{lem:haar-neigh-23}(iii). Gdy $V\subseteq W$,
żadne przesunięcie $V$ nie przecina zarazem $K$ i $L$.
Jeśli jeden ze zbiorów jest pusty, równość wynika
z $N(\varnothing,V)=0$. W przeciwnym razie żaden
element minimalnego pokrycia $K\cup L$ nie jest
rozłączny z $K\cup L$, gdyż można byłoby go usunąć.
Każdy element przecina więc dokładnie jeden z dwóch
zbiorów. Podział na dwie rozłączne rodziny daje
odpowiednio pokrycia $K$ i $L$. Stąd
$N(K\cup L,V)\ge N(K,V)+N(L,V)$.
Nierówność odwrotna wynika z łączenia pokryć.
\end{proof}

\begin{lemat}[Zawartość niezmiennicza na zbiorach zwartych]
\label{lem:haar-content-23}
Istnieje skończona funkcja nieujemna $q$ na rodzinie
wszystkich zbiorów zwartych $K\subseteq G$, taka że
\begin{enumerate}[label=(\roman*)]
\item $q(\varnothing)=0$, $q(K_0)=1$ oraz $q(gK)=q(K)$;
\item $K\subseteq L$ implikuje $q(K)\le q(L)$;
\item $q(K\cup L)\le q(K)+q(L)$;
\item jeśli $K\cap L=\varnothing$, to
      $q(K\cup L)=q(K)+q(L)$.
\end{enumerate}
\end{lemat}
\begin{proof}
Niech $\mathcal K$ będzie zbiorem wszystkich zwartych
podzbiorów $G$. Na mocy
lematu~\ref{lem:haar-product-compact-23} produkt
$P=\prod_{K\in\mathcal K}[0,C_K]$ jest zwarty,
bez względu na moc $\mathcal K$.
Każde $q_V=(q_V(K))_{K\in\mathcal K}$ należy do $P$.
Dla $W\in\mathcal N$ połóżmy
$F_W=\overline{\{q_V:V\in\mathcal N,\ V\subseteq W\}}$.
Domknięcie bierzemy w $P$. Jest to niepusty zbiór
domknięty, ponieważ zawiera $q_W$.
Przecięcie skończenie wielu elementów $\mathcal N$
jest otwartym otoczeniem $e$ o zwartym domknięciu:
domknięcie przecięcia jest domkniętym podzbiorem
zwartego domknięcia dowolnego z nich.
Dla niepustej skończonej rodziny $W_1,\dots,W_r$ mamy
$F_{W_1\cap\dots\cap W_r}\subseteq\bigcap_{j=1}^rF_{W_j}$.
Zwartość $P$ pozwala wybrać
$q\in\bigcap_{W\in\mathcal N}F_W$.

Dla ustalonych $g,K,L$ odpowiednie warunki (i)--(iii)
są równościami albo nierównościami ciągłych
funkcji współrzędnych na $P$. Na przykład
$\{p\in P:p(K\cup L)-p(K)-p(L)\le0\}$ jest
domknięty jako przeciwobraz przedziału $(-\infty,0]$.
Każdy taki zbiór zawiera wszystkie $q_V$,
a więc także każde $F_W$ i punkt $q$.
Warunki obowiązują dla dowolnego wyboru
$g,K,L$, co daje (i)--(iii).
Jeżeli $K\cap L=\varnothing$, z poprzedniego lematu
mamy otoczenie $W$, poniżej którego wszystkie $q_V$
spełniają równość addytywności dla $K,L$.
Jej zbiór rozwiązań jest domknięty i zawiera $F_W$,
a więc również $q$. To dowodzi (iv).
\end{proof}

\begin{lemat}[Od zawartości na zbiorach zwartych do miary borelowskiej]
\label{lem:haar-content-to-measure-23}
Niech $X$ będzie lokalnie zwartą przestrzenią Hausdorffa,
a $q:\mathcal K(X)\to[0,\infty)$ funkcją, gdzie
$\mathcal K(X)$ jest rodziną wszystkich zwartych podzbiorów $X$.
Załóżmy, że $q(\varnothing)=0$, $q$ jest monotoniczna,
$q(K\cup L)\le q(K)+q(L)$ dla wszystkich $K,L\in\mathcal K(X)$
oraz $q(K\cup L)=q(K)+q(L)$, gdy $K\cap L=\varnothing$. Dla otwartego $U$
i dowolnego $A\subseteq X$ definiujemy
\begin{equation}\label{eq:haar-from-content-23}
 t(U)=\sup_{K\subseteq U,\ K\ \mathrm{zwarty}}q(K),
 \qquad
 m^*(A)=\inf_{U\supseteq A,\ U\ \mathrm{otwarty}}t(U).
\end{equation}
Wówczas $m^*$ jest miarą zewnętrzną,
wszystkie zbiory borelowskie są mierzalne w sensie
Carathéodory'ego, a $\mu=m^*|_{\mathcal B(X)}$
spełnia warunki (i)--(iii) regularności podane
na początku rozdziału. Ponadto $\mu(K)\ge q(K)$.
Jeśli homeomorfizm $T:X\to X$ zachowuje $q$ na
zbiorach zwartych, to $\mu(TA)=\mu(A)$ dla każdego
zbioru borelowskiego $A$.
\end{lemat}
\begin{proof}
Wprost z definicji $m^*(U)=t(U)$, gdy $U$ jest otwarty.
Najpierw dowodzimy podaddytywności $t$ na pokryciach
otwartych. Jeśli $K\subseteq\bigcup_{j\ge1}U_j$ jest
zwarty, możemy wybrać skończone podpokrycie.
Z lematu~\ref{lem:haar-neigh-23}(iv) (którego dowód
jest topologiczny i działa w każdej przestrzeni lokalnie
zwartej Hausdorffa) dostajemy zbiory zwarte
$K_j\subseteq U_j$ pokrywające $K$. Podaddytywność $q$
daje $q(K)\le\sum_jt(U_j)$. Biorąc supremum po $K$,
uzyskujemy
$t(\bigcup_{j\ge1}U_j)\le\sum_jt(U_j)$.
Aby uzyskać przeliczalną podaddytywność $m^*$,
dla zbiorów $A_j$ o skończonej sumie $\sum_jm^*(A_j)$
i $\varepsilon>0$ wybieramy otwarte $U_j\supseteq A_j$
z $t(U_j)<m^*(A_j)+\varepsilon2^{-j}$.
Wtedy $m^*(\bigcup_jA_j)\le\sum_jm^*(A_j)+\varepsilon$;
dowolność $\varepsilon$ daje przeliczalną podaddytywność.
Gdy suma po prawej stronie jest nieskończona, nierówność
zachodzi automatycznie. Ponadto $m^*(\varnothing)=t(\varnothing)=0$,
a monotoniczność $m^*$ wynika z tego, że dla $A\subseteq B$
każdy otwarty nadzbiór $B$ jest także otwartym nadzbiorem $A$.

Najważniejszym krokiem jest mierzalność \emph{otwartych}
zbiorów bez odwoływania się do jakiejkolwiek metryki.
Ustalmy otwarty $O\subseteq X$, dowolny zbiór $A$
o $m^*(A)<\infty$ i $\varepsilon>0$.
Wybieramy otwarty $V\supseteq A$ z
$t(V)<m^*(A)+\varepsilon$. Ponieważ
$t(V\cap O)\le t(V)<\infty$, istnieje zwarty
$K\subseteq V\cap O$ taki, że
$q(K)>t(V\cap O)-\varepsilon$; jeśli supremum
jest zerowe, bierzemy $K=\varnothing$.
Dla dowolnego zbioru zwartego $L\subseteq V\setminus K$
addytywność na rozłącznych zbiorach zwartych daje
$q(K)+q(L)=q(K\cup L)\le t(V)$.
Stąd $t(V\setminus K)\le t(V)-q(K)$.
Ponieważ $A\setminus O\subseteq V\setminus K$,
a $A\cap O\subseteq V\cap O$, dostajemy
\[
 m^*(A\setminus O)+m^*(A\cap O)
 \le t(V\setminus K)+t(V\cap O)
 \le t(V)+\varepsilon<m^*(A)+2\varepsilon.
\]
Zbiór $V\setminus K$ jest otwarty, gdyż $K$ jest
zwarty w przestrzeni Hausdorffa, a więc domknięty.
Wszystkie użyte w oszacowaniu wartości są skończone,
ponieważ nie przekraczają $t(V)$.
Przy $m^*(A)=\infty$ nierówność
$m^*(A)\ge m^*(A\setminus O)+m^*(A\cap O)$
jest automatyczna; nierówność odwrotna w obu
przypadkach wynika z podaddytywności.
Po przejściu z $\varepsilon$ do zera warunek
Carathéodory'ego jest spełniony. Z twierdzenia
\atlasref{ME-03::thm:caratheodory}{ME-03: twierdzenie Carathéodory'ego}
wynika, że rodzina zbiorów mierzalnych w sensie
Carathéodory'ego jest $\sigma$-algebrą i że
ograniczenie do niej $m^*$ jest miarą. Ponieważ
zawiera wszystkie otwarte zbiory, zawiera
$\mathcal B(X)$. Stąd $\mu$ jest miarą borelowską.

Regularność zewnętrzna jest zawarta w definicji $m^*$.
Dla zbioru zwartego $K$ wybieramy otwarty $W\supseteq K$
o zwartej domkniętej powłoce $C$. Wówczas
$\mu(K)\le t(W)\le q(C)<\infty$.
Dla każdego otwartego $U\supseteq K$ mamy
$t(U)\ge q(K)$, więc $\mu(K)\ge q(K)$.
Dla otwartego $O$ zachodzi
$\mu(O)=t(O)=\sup_{K\subseteq O}q(K)
\le\sup_{K\subseteq O}\mu(K)\le\mu(O)$,
co dowodzi regularności wewnętrznej na otwartych.

Na koniec, jeśli $T$ jest wskazanym homeomorfizmem,
to $T$ przekształca bijektywnie zbiory zwarte zawarte w $U$
na zbiory zwarte zawarte w $TU$, skąd $t(TU)=t(U)$.
Ponadto $U\mapsto TU$ jest bijekcją między rodzinami
otwartych nadzbiorów $A$ i $TA$, ponieważ
$T^{-1}$ jest homeomorfizmem. Biorąc infimum
po tych rodzinach, dostajemy $m^*(TA)=m^*(A)$.
Homeomorfizm zachowuje borelowskość, więc równość
przechodzi na miarę $\mu$. Nie twierdzimy, że
$\mu(K)=q(K)$ dla wszystkich zwartych $K$:
do konstrukcji wystarcza nierówność $\mu(K)\ge q(K)$.
\end{proof}

\begin{proof}[Dowód istnienia w twierdzeniu~\ref{thm:haar-main-23}
bez założeń przeliczalności]
Z lematu~\ref{lem:haar-content-23} otrzymujemy
zawartość $q$. Lemat~\ref{lem:haar-content-to-measure-23}
daje borelowską miarę Radona $\mu$ niezmienniczą
względem wszystkich lewych translacji.
Ponieważ $\mu(K_0)\ge q(K_0)=1$, jest niezerowa.
Jej skończoność na kompaktach wynika z konstrukcji.
Dla miary prawej stosujemy inwersję:
$\mu_R(E)=\mu(E^{-1})$.
Jest to miara Radona, bo inwersja jest homeomorfizmem,
a $\mu_R(Eg)=\mu(g^{-1}E^{-1})=\mu(E^{-1})$.
Nie było potrzebne żadne założenie o przeliczalnej bazie
otoczeń $e$ ani o $\sigma$-zwartości $G$.
\end{proof}

\begin{wniosek}[Dodatniość miary zbiorów otwartych]
\label{cor:haar-open-positive-23}
Jeśli $\mu$ jest niezerową lewą miarą Haara
na lokalnie zwartej grupie Hausdorffa $G$, to
$\mu(U)>0$ dla każdego niepustego otwartego $U\subseteq G$.
\end{wniosek}
\begin{proof}
Jeśli $\mu(U)=0$ dla pewnego niepustego $U$, wybierzmy
$a\in U$. Dla każdego zbioru zwartego $K\subseteq G$ skończenie
wiele zbiorów $xa^{-1}U$ pokrywa $K$, ponieważ każdy
z nich jest przesunięciem otwartego otoczenia $e$.
Z niezmienniczości każdy zbiór w skończonym
pokryciu ma miarę zero, zatem z monotoniczności
i skończonej podaddytywności $\mu(K)=0$.
Regularność wewnętrzna na otwartych daje teraz
$\mu(G)=0$, wbrew niezerowości miary.
\end{proof}

\section{Jednoznaczność: argument z funkcją przybliżającą jedynkę}
\begin{lemat}[Jednostajna ciągłość translacji funkcji]
\label{lem:haar-uniform-23}
Niech $G$ będzie lokalnie zwartą grupą Hausdorffa.
Jeśli $f\in C_c(G)$, to dla każdego $\varepsilon>0$
istnieje otoczenie $V$ elementu $e$ takie, że
$|f(yz)-f(y)|<\varepsilon$ dla wszystkich
$y\in G$ i $z\in V$.
\end{lemat}
\begin{proof}
Dla $f=0$ teza jest natychmiastowa. Niech
$K=\operatorname{supp}f$ i wybierzmy symetryczne
relatywnie zwarte otwarte otoczenie $C$ elementu $e$
z lematu \ref{lem:haar-neigh-23}(i).
Zbiór $M=K\overline C$ jest zwarty jako obraz
zwartego produktu $K\times\overline C$ przez ciągłe
mnożenie. Ponadto $K\subseteq M$, gdyż $e\in C$. Ciągłość
$(y,z)\mapsto f(yz)-f(y)$ w każdym $(y,e)$,
$y\in M$, pozwala wybrać otoczenia $W_y$ punktu $y$
i $V_y$ elementu $e$ takie, że
$|f(uz)-f(u)|<\varepsilon$ dla $u\in W_y$, $z\in V_y$.
Ze zwartego $M$ wybieramy skończone pokrycie
$W_{y_1},\dots,W_{y_r}$; połóżmy
$V=C\cap V_{y_1}\cap\cdots\cap V_{y_r}$,
zmniejszając je w razie potrzeby do symetrycznego
otoczenia. Dla $u\in M$, $z\in V$ teza wynika
z konstrukcji. Jeśli $u\notin M$, to $f(u)=0$,
a $uz\in K$ prowadziłoby do
$u\in K V^{-1}\subseteq K\overline C=M$.
Zatem również $f(uz)=0$.
\end{proof}

\begin{lemat}[Fubini dla ciągłych jąder o zwartym nośniku]
\label{lem:haar-fubini-compact-23}
Niech $X,Y$ będą lokalnie zwartymi przestrzeniami Hausdorffa,
a $\alpha,\beta$ nieujemnymi miarami borelowskimi
skończonymi na zbiorach zwartych. Jeśli
$H\in C_c(X\times Y;\mathbb K)$, gdzie
$\mathbb K=\mathbb R$ lub $\mathbb C$, to obie całki iterowane
\[
 \int_X\!\left(\int_Y H(x,y)\,\mathrm d\beta(y)\right)
 \mathrm d\alpha(x),\qquad
 \int_Y\!\left(\int_X H(x,y)\,\mathrm d\alpha(x)\right)
 \mathrm d\beta(y)
\]
są dobrze określone, skończone i równe, nawet gdy
$X,Y$ nie są metryzowalne ani $\alpha,\beta$
nie są $\sigma$-skończone.
\end{lemat}
\begin{proof}
Dla $H=0$ obie całki wynoszą zero. W przeciwnym
razie rzuty zwartego nośnika $H$ są niepustymi
zbiorami zwartymi $K\subseteq X$ i $L\subseteq Y$,
jako obrazy przez ciągłe rzuty. Są domknięte w
odpowiednich przestrzeniach Hausdorffa, a więc
borelowskie. Funkcja $H$ znika poza $K\times L$,
więc każdą całkę wewnętrzną można ograniczyć
do odpowiedniego zbioru zwartego. Przekroje są
ciągłe i ograniczone przez $\|H\|_\infty$;
ich całki są bezwzględnie zbieżne, bo
$\alpha(K),\beta(L)<\infty$.
Ograniczenia miar do $K$ i $L$ są miarami na
ich $\sigma$-algebrach borelowskich, ponieważ
borelowskie zbiory podprzestrzeni domkniętej
są borelowskie w przestrzeni wyjściowej.
Poniżej wykażemy mierzalność funkcji otrzymanych
po całkowaniu wewnętrznym, co uzasadni również
ograniczenie całek zewnętrznych.
Na iloczynie tych zwartych przestrzeni każdą funkcję
ciągłą przybliża się jednostajnie przez sumy
$\sum_{j=1}^Nc_j a_j(x)b_j(y)$, gdzie
$c_j\in\mathbb K$, $a_j\in C(K;\mathbb R)$,
$b_j\in C(L;\mathbb R)$. Oto bezpośrednie
uzasadnienie tej aproksymacji. Dla danego $\varepsilon>0$
z ciągłości $H$ i zwartości iloczynu można wybrać
skończone otwarte pokrycia $(U_i)$ zbioru $K$ oraz
$(V_j)$ zbioru $L$ takie, że oscylacja $H$ na każdym
$U_i\times V_j$ jest mniejsza niż $\varepsilon$.
Istnienie tych pokryć sprawdzamy bez metryzowalności:
dla ustalonego $x\in K$ i każdego $y\in L$ ciągłość
$H$ daje otoczenia $U_{x,y}\ni x$, $V_{x,y}\ni y$
takie, że $|H(x',y')-H(x,y)|<\varepsilon/4$
dla $(x',y')\in U_{x,y}\times V_{x,y}$.
Zwartość $L$ pozwala wybrać skończenie wiele
$V_{x,y}$ pokrywających $L$; przecięcie
odpowiednich $U_{x,y}$ daje otoczenie $U_x$,
na którym dla dowolnych $x',x''\in U_x$ i $y\in L$
mamy $|H(x',y)-H(x'',y)|<\varepsilon/2$.
Analogicznie dobieramy otoczenia $V_y$ z oszacowaniem
$|H(x,y')-H(x,y'')|<\varepsilon/2$ jednostajnym
względem $x\in K$. Skończone podpokrycia
$(U_i)$, $(V_j)$ dają żądane oszacowanie oscylacji
na prostokątach $U_i\times V_j$.
Na zwartych przestrzeniach Hausdorffa konstruujemy
ciągłe podziały jedności $(a_i)$, $(b_j)$ podporządkowane
tym pokryciom: stosujemy lemat~\ref{lem:haar-neigh-23}(iv)
do każdego z pokryć, a następnie lemat~\ref{lem:haar-cutoff-23}
do otrzymanych zwartych podzbiorów. Sumy funkcji odcinających
są wszędzie co najmniej $1$; dzieląc każdą funkcję przez
odpowiednią sumę, uzyskujemy żądane podziały jedności.
Puste elementy pokryć pomijamy. Wybierając
$c_{ij}$ jako wartość $H$ w dowolnym punkcie
$U_i\times V_j$, zauważamy, że jeśli
$a_i(x)b_j(y)>0$, to $(x,y)\in U_i\times V_j$,
więc $|H(x,y)-c_{ij}|<\varepsilon$.
Ponieważ $a_i,b_j$ są nieujemne i
$\sum_{i,j}a_i(x)b_j(y)=1$, nierówność trójkąta
daje
\[
 \left\|H-\sum_{i,j}c_{ij}a_i\otimes b_j\right\|_\infty
 \le\varepsilon.
\]
Dla sumy elementarnej $S=\sum_{i,j}c_{ij}a_i\otimes b_j$
obie całki iterowane wynoszą
\[
 \sum_{i,j}c_{ij}
       \left(\int_K a_i\,\mathrm d\alpha\right)
       \left(\int_L b_j\,\mathrm d\beta\right).
\]
Wszystkie składniki są skończone, więc stosujemy
wyłącznie liniowość skończonych całek i skończonych
sum, bez twierdzenia Fubiniego. Funkcja $x\mapsto\int_L H(x,y)\,\mathrm d\beta(y)$
na $K$ jest granicą jednostajną ciągłych funkcji otrzymanych
z tych sum, ponieważ błąd jest ograniczony przez
$\varepsilon\beta(L)$. Jest zatem ciągła na $K$,
a po przedłużeniu przez zero poza domknięty $K$ jest
borelowska na $X$: przeciwobraz zbioru otwartego
w $\mathbb K$ jest borelowskim podzbiorem $K$,
uzupełnionym o $X\setminus K$, gdy zbiór otwarty
zawiera zero. Przedłużenie to jest właśnie
$x\mapsto\int_Y H(x,y)\,\mathrm d\beta(y)$,
ponieważ $H$ znika poza $K\times L$. Analogicznie otrzymujemy mierzalność
drugiej całki wewnętrznej. Skończoność $\alpha(K)$
i $\beta(L)$ pozwala przejść do granicy w obu całkach
zewnętrznych, z błędem co najwyżej $\varepsilon\alpha(K)\beta(L)$.
Dowodzi to tezy również dla funkcji zespolonych.
\end{proof}

\begin{twierdzenie}[Jednoznaczność miary Haara]
\label{thm:haar-unique-23}
Jeśli $\mu,\nu$ są dwiema lewostronnymi miarami Haara
na lokalnie zwartej grupie Hausdorffa $G$, to istnieje
$c>0$ takie, że $\nu=c\mu$ na $\mathcal B(G)$.
\end{twierdzenie}
\begin{proof}
Dla dowolnej funkcji $h\in C_c(G)$ oznaczamy
$\mu(h)=\int_G h\,\mathrm d\mu$ i analogicznie $\nu(h)$.
Jeśli $h\ge0$ i $h\ne0$, to dla pewnego $a>0$
zbiór $\{h>a\}$ jest niepusty i otwarty. Z wniosku
\ref{cor:haar-open-positive-23} mamy
$\mu(h)\ge a\mu(\{h>a\})>0$ i analogicznie $\nu(h)>0$;
obie całki są skończone przez zwartość nośnika.
Takie funkcje istnieją z lematu \ref{lem:haar-cutoff-23}
dla zbioru $\{e\}$ i relatywnie zwartego otoczenia $e$.
Niech $f\in C_c(G)$, a $\varphi\in C_c(G)$ będzie
nieujemna, niezerowa i znormalizowana przez
$\mu(\varphi)=1$.
Funkcja $(x,y)\mapsto f(x)\varphi(y^{-1}x)$
ma nośnik zawarty w zwartym iloczynie
$\operatorname{supp}f\times
(\operatorname{supp}f)(\operatorname{supp}\varphi)^{-1}$.
Stosujemy zatem lemat~\ref{lem:haar-fubini-compact-23},
który nie wymaga $\sigma$-skończoności miary na $G$.
Otrzymujemy
\begin{align*}
 I&=\int_G\!\int_G f(x)\varphi(y^{-1}x)
       \,\mathrm d\nu(y)\,\mathrm d\mu(x)\\
  &=c_\varphi\int_Gf(x)\,\mathrm d\mu(x),\qquad
 c_\varphi:=\int_G\varphi(z^{-1})\,\mathrm d\nu(z),
\end{align*}
bo po podstawieniu $y=xz$ mamy
$y^{-1}x=z^{-1}$ i wewnętrzna całka nie zależy
od $x$. Podstawienie jest uzasadnione twierdzeniem
\atlasref{ME-13::thm:pushforward-integral-13}{ME-13: całka względem miary obrazowej}: obraz $\nu$ przez
lewą translację jest równy $\nu$ z niezmienniczości.
Funkcja $z\mapsto\varphi(z^{-1})$ jest ciągła,
nieujemna, niezerowa i ma zwarty nośnik, więc
$0<c_\varphi<\infty$.
Zamieniając kolejność całkowania i podstawiając
$x=yz$ (lewostronna niezmienniczość $\mu$), dostajemy
\begin{equation}\label{eq:haar-uniqueness-23}
 c_\varphi\int_G f(x)\,\mathrm d\mu(x)
 =\int_G\!\int_G f(yz)\varphi(z)
       \,\mathrm d\mu(z)\,\mathrm d\nu(y).
\end{equation}
Ustalmy symetryczne relatywnie zwarte otwarte otoczenie $C$
elementu $e$. Rodzina otwartych otoczeń $V\subseteq C$
elementu $e$ jest skierowana przez odwrotną inkluzję.
Dla każdego $V$ z lematu~\ref{lem:haar-cutoff-23}
wybieramy nieujemną $\varphi_V\in C_c(G)$ równą $1$
w $e$ i o nośniku zawartym w $V$, a następnie dzielimy
ją przez jej całkę względem $\mu$, dodatnią na mocy
wniosku~\ref{cor:haar-open-positive-23}.
Otrzymujemy sieć funkcji z $\mu(\varphi_V)=1$,
nie zakładając istnienia przeliczalnej bazy otoczeń.
Na mocy lematu~\ref{lem:haar-uniform-23}
dla każdego $\varepsilon>0$ i dostatecznie małego $V$ mamy
\[
 \left|\int_G f(yz)\varphi_V(z)\,\mathrm d\mu(z)-f(y)\right|
 \le\int_G |f(yz)-f(y)|\varphi_V(z)\,\mathrm d\mu(z)
 \le\varepsilon
\]
dla wszystkich $y\in G$. Obie funkcje znikają poza wspólnym
zwartym zbiorem $\operatorname{supp}f\,\overline C$,
więc błąd ich całek względem $\nu$ jest ograniczony przez
$\varepsilon\nu(\operatorname{supp}f\,\overline C)$.
Zatem prawa strona dąży do $\int_Gf(y)\,\mathrm d\nu(y)$.
Wybierając ustaloną niezerową funkcję
$f_0\in C_c(G)$, $f_0\ge0$, mamy
$0<\mu(f_0),\nu(f_0)<\infty$ na mocy początkowego
oszacowania. Stosujemy \eqref{eq:haar-uniqueness-23}
do $f_0$ i $\varphi_V$ oraz dzielimy przez
$\mu(f_0)$. Otrzymujemy zbieżność sieci
$c_{\varphi_V}\to c:=\nu(f_0)/\mu(f_0)>0$.
Wstawienie dowolnego $f$ do
\eqref{eq:haar-uniqueness-23} daje
$\nu(f)=c\mu(f)$. Połóżmy $\widetilde\mu=c\mu$.
Pozostaje wyjaśnić, dlaczego równość całek względem
$\widetilde\mu$ i $\nu$ dla wszystkich funkcji z $C_c(G)$
oznacza równość tych miar również bez metryzowalności $G$.
Dla każdego zwartego $K\subseteq U$, gdzie $U$ jest otwarty,
istnieje funkcja $h\in C_c(G)$ taka, że
$0\le h\le1$, $h|_K=1$ i $\operatorname{supp}h\subseteq U$
na mocy lematu~\ref{lem:haar-cutoff-23}.
Wówczas $\widetilde\mu(K)\le\int_Gh(x)\,\mathrm d\widetilde\mu(x)
=\int_Gh(x)\,\mathrm d\nu(x)\le\nu(U)$. Z regularności
wewnętrznej $\widetilde\mu(U)\le\nu(U)$. Po zamianie miar
otrzymujemy równość na otwartych, a regularność zewnętrzna
daje równość na wszystkich zbiorach borelowskich.
Stąd $\nu=c\mu$. W przypadku dwóch prawostronnych miar
Haara ich obrazy przez inwersję są lewostronnymi miarami
Haara, a więc różnią się dodatnim czynnikiem stałym.
Ponowne zastosowanie inwersji daje tę samą zależność
dla wyjściowych miar prawostronnych. Kończy to również
dowód jednoznaczności w twierdzeniu~\ref{thm:haar-main-23}.
\end{proof}

\begin{przyklad}[Dlaczego nie wymagamy regularności wewnętrznej wszędzie]
\label{ex:haar-non-sigma-compact-23}
Niech $D$ będzie nieprzeliczalną grupą dyskretną,
a $G=\mathbb R\times D$ grupą z topologią produktową.
Rozważmy domknięty zbiór $Y=\{0\}\times D$.
Każdy zbiór zwarty $K\subseteq Y$ jest skończony, ponieważ
rzutowanie $K$ na dyskretną przestrzeń $D$ daje zbiór
zwarty, a więc skończony. Punkty mają miarę Haara zero: wszystkie jednopunktowe
zbiory mają tę samą miarę $a=\mu(\{(0,e)\})$
na mocy niezmienniczości lewostronnej.
W zwartym zbiorze $[0,1]\times\{e\}$ można wybrać
dowolnie wiele różnych punktów, więc dla każdego
$n\ge1$ mamy $na\le\mu([0,1]\times\{e\})<\infty$.
Stąd $a=0$.
Wobec tego $\mu(K)=0$ dla każdego zbioru zwartego $K\subseteq Y$.

Z drugiej strony, jeśli $U\supseteq Y$ jest otwarty,
to $U\cap(\mathbb R\times\{d\})$ jest niepustym
zbiorem otwartym dla każdego $d\in D$, a więc ma
ściśle dodatnią miarę $\mu$ na mocy wniosku~\ref{cor:haar-open-positive-23}, gdyż składowa jest otwarta w $G$.
Składowe są parami rozłączne i jest ich nieprzeliczalnie
wiele, zatem $\mu(U)=\infty$: gdyby $\mu(U)$ było
skończone, najwyżej przeliczalnie wiele rozłącznych
zbiorów mierzalnych mogłoby mieć dodatnią miarę.
Istotnie, dla każdego $n\ge1$ tylko skończenie wiele
rozłącznych części $U$ mogłoby mieć miarę co najmniej
$1/n$, gdyż dowolna ich skończona suma ma miarę
nie większą niż $\mu(U)$. Każda część o dodatniej
mierze należałaby do jednej z tych skończonych rodzin,
co przeczy nieprzeliczalności $D$.
Regularność zewnętrzna daje $\mu(Y)=\infty$, chociaż
$\sup_{K\subseteq Y,\,K\text{ zwarty}}\mu(K)=0$.
Takiego zjawiska nie ma dla zbiorów o mierze
$\sigma$-skończonej, zgodnie z lematem
\atlasref{ME-20::lem:radon-sigmafinite-inner-20}{ME-20: regularność wewnętrzna na zbiorach o mierze $\sigma$-skończonej}.
\end{przyklad}

\section{Całka względem miary Haara i normalizacja}
\begin{definicja}[Całka względem miary Haara]\label{def:haar-integral-23}
Dla nieujemnej funkcji borelowskiej $f:G\to[0,\infty]$ definiujemy $\int_G f\,\mathrm d\mu$ jako supremum całek nieujemnych funkcji prostych $s\le f$. Dla rzeczywistej mierzalnej $f$ całka jest skończona wtedy i tylko wtedy, gdy $\int_G|f|\,\mathrm d\mu<\infty$; wówczas $\int_Gf\,\mathrm d\mu=\int_Gf^+\,\mathrm d\mu-\int_Gf^-\,\mathrm d\mu$. To dokładnie całka Lebesgue'a na przestrzeni z miarą $(G,\mathcal B(G),\mu)$.
\end{definicja}
Z lewostronnej niezmienniczości oraz twierdzenia
\atlasref{ME-13::thm:pushforward-integral-13}{ME-13: całka względem miary obrazowej} wynika
$\int_G f(ax)\,\mathrm d\mu(x)=\int_Gf(x)\,\mathrm d\mu(x)$
dla $a\in G$ i funkcji nieujemnych borelowskich lub
bezwzględnie całkowalnych. W przypadku zespolonym
całkujemy osobno część rzeczywistą i urojoną. Każda funkcja z $C_c(G)$ jest całkowalna; ciągłość bez zwartego nośnika nie zapewnia tego na grupie niezwartej.
Dla ustalonej niezerowej $f_0\in C_c(G)$, $f_0\ge0$, całka jest dodatnia i skończona na mocy wniosku \ref{cor:haar-open-positive-23} i zwartości nośnika. Jednoznaczność miary Haara pozwala zatem wybrać dokładnie jedną miarę spełniającą $\int_G f_0\,\mathrm d\mu=1$. Na grupie zwartej analogicznie normalizujemy $\mu(G)=1$.

\section{Prawe miary Haara, funkcja modularna i unimodularność}
\begin{twierdzenie}[Prawa miara Haara i modularność]
\label{thm:haar-modular-23}
Niech $\mu$ będzie lewostronną miarą Haara.
Wówczas $\mu_R(E):=\mu(E^{-1})$ jest prawostronną
miarą Haara. Istnieje dokładnie jeden ciągły
homomorfizm $c:G\to(0,\infty)$ taki, że
\begin{equation}\label{eq:haar-modular-23}
 \mu(Eg)=c(g)\mu(E)\qquad(E\in\mathcal B(G),\ g\in G).
\end{equation}
Grupę nazywamy \emph{unimodularną}, gdy $c\equiv1$.
\end{twierdzenie}
\begin{proof}
Inwersja jest homeomorfizmem, toteż $\mu_R$ jest
miarą Radona. Ponadto
$\mu_R(Eg)=\mu((Eg)^{-1})=\mu(g^{-1}E^{-1})
=\mu(E^{-1})=\mu_R(E)$.
Dla ustalonego $g$ miara $\mu_g(E)=\mu(Eg)$
jest obrazem $\mu$ przez homeomorfizm $R_{g^{-1}}$.
Jest niezerową miarą Radona: homeomorfizm zachowuje
zbiory zwarte i otwarte oraz ich inkluzje, więc
zachowuje oba warunki regularności. Ponadto
$\mu_g(aE)=\mu(aEg)=\mu(Eg)=\mu_g(E)$.
Na mocy twierdzenia~\ref{thm:haar-unique-23} istnieje
dokładnie jedna stała $c(g)>0$ taka, że
$\mu_g=c(g)\mu$.
Wybieramy zwarty zbiór $K$ o $0<\mu(K)<\infty$,
na przykład domknięcie relatywnie zwartego otwartego
otoczenia $e$ (wniosek \ref{cor:haar-open-positive-23}).
Wówczas
\[
 c(gh)\mu(K)=\mu(Kgh)=c(h)\mu(Kg)
 =c(h)c(g)\mu(K).
\]
Dzieląc przez $\mu(K)$, otrzymujemy
$c(gh)=c(g)c(h)$. Analogicznie $c(e)=1$ i
$c(g^{-1})=c(g)^{-1}$. Zastąpienie $\mu$ przez
dodatnią wielokrotność nie zmienia $c$.
Dla niezerowej $f\in C_c(G)$, $f\ge0$,
mamy $\int_G f\,\mathrm d\mu>0$ oraz, z twierdzenia
\atlasref{ME-13::thm:pushforward-integral-13}{ME-13: całka względem miary obrazowej} zastosowanego do
$R_{g^{-1}}$, $c(g)$ jest ilorazem całek
\[
 c(g)=\frac{\int_G f(xg^{-1})\,\mathrm d\mu(x)}
 {\int_G f(x)\,\mathrm d\mu(x)}.
\]
Dla ustalonego $g_0\in G$ wybieramy relatywnie zwarte
otwarte otoczenie $W$ punktu $g_0$. Dla $g\in W$ funkcja
$x\mapsto f(xg^{-1})$ znika poza wspólnym zwartym zbiorem
$M=(\operatorname{supp}f)\overline W$.
Z lematu~\ref{lem:haar-uniform-23}, zastosowanego do
$y=xg_0^{-1}$ i $z=g_0g^{-1}\to e$, wynika jednostajna
zbieżność $f(xg^{-1})\to f(xg_0^{-1})$ względem $x$.
Dla $g\in W$ moduł różnicy liczników nie przekracza
\[
 \mu(M)\sup_{x\in G}|f(xg^{-1})-f(xg_0^{-1})|,
\]
co dąży do zera dla $g\to g_0$ w sensie otoczeń.
Zatem licznik, a więc również $c$, jest ciągły.
Nie odwołujemy się do ciągów jako kryterium ciągłości.
\end{proof}

\begin{uwaga}[Dwie konwencje]
W literaturze spotyka się dwie odwrotne konwencje
funkcji modularnej. W tym skrypcie konsekwentnie
oznaczamy $\Delta=c$ i przyjmujemy
$\mu(Eg)=\Delta(g)\mu(E)$. W konwencji przeciwnej
wszystkie czynniki $\Delta$ należy zastąpić ich
odwrotnościami. W szczególności unimodularność
nie zależy od wyboru konwencji.
\end{uwaga}

\begin{wniosek}[Wzory całkowe i inwersja]\label{cor:haar-integral-inversion-23}
Oznaczmy $\Delta=c$ w konwencji \eqref{eq:haar-modular-23}. Dla funkcji nieujemnych borelowskich zachodzi
\begin{equation}\label{eq:modularna}
\int_G f(xa)\,\mathrm d\mu(x)=\Delta(a)^{-1}\int_G f(x)\,\mathrm d\mu(x).
\end{equation}
Ponadto
\begin{equation}\label{eq:inwersja}
\int_G f(x^{-1})\,\mathrm d\mu(x)=\int_G f(x)\Delta(x)^{-1}\,\mathrm d\mu(x).
\end{equation}
Wzory obowiązują również dla funkcji rzeczywistych
i zespolonych, jeśli odpowiednie całki są bezwzględnie
zbieżne. W szczególności w \eqref{eq:inwersja}
warunkiem jest $\int_G |f(x)|\Delta(x)^{-1}
\,\mathrm d\mu(x)<\infty$, a nie sama przynależność
$f$ do $L^1(G,\mu)$.
\end{wniosek}
\begin{proof}
Dla $f=\chi_E$ lewa strona \eqref{eq:modularna}
jest równa $\mu(Ea^{-1})=\Delta(a)^{-1}\mu(E)$.
Liniowość dla nieujemnych funkcji prostych oraz
twierdzenie \ref{thm:mct} dają wzór dla każdej
nieujemnej borelowskiej $f$.
W szczególności dla nieujemnej borelowskiej $h$
i borelowskiego $A$ zachodzi
\[
 \int_{Aa} h(y)\,\mathrm d\mu(y)
 =\Delta(a)\int_A h(xa)\,\mathrm d\mu(x),
\]
gdyż stosujemy \eqref{eq:modularna} do $h\chi_{Aa}$.

Niech $w=\Delta^{-1}$ oraz
$\rho(A)=\int_A w\,\mathrm d\mu$.
Jest to miara borelowska na mocy przeliczalnej
addytywności całki nieujemnej (twierdzenie
\ref{thm:mct} zastosowane do skończonych sum
funkcji charakterystycznych rozłącznych zbiorów).
Dla każdej nieujemnej borelowskiej $F$ zachodzi
$\int_G F\,\mathrm d\rho=\int_G Fw\,\mathrm d\mu$:
dla funkcji charakterystycznych jest to definicja
$\rho$, dla nieujemnych funkcji prostych wynika
z liniowości, a w ogólności z twierdzenia
\ref{thm:mct} zastosowanego po obu stronach do
rosnących przybliżeń prostych (także po pomnożeniu
przez $w$).
Ponieważ $w(xa)=w(x)w(a)$, dostajemy
\[
 \rho(Aa)=\Delta(a)\int_A w(xa)\,\mathrm d\mu(x)
 =\rho(A).
\]
Na zbiorze zwartym $K$ funkcja $w$ jest ograniczona,
więc $\rho(K)\le\sup_K w\,\mu(K)<\infty$ dla
niepustego $K$; dla pustego zbioru wartość wynosi zero.
Na niepustym zwartym otoczeniu $e$ o dodatniej
mierze $\mu$ funkcja $w$ ma dodatnie minimum,
stąd $\rho$ jest niezerowa.

\emph{Regularność wewnętrzna.}
Dla otwartego $U$ i $n\ge1$ kładziemy
\[
 s_n=2^{-n}\sum_{k=1}^{n2^n}
             \chi_{U\cap\{w>k2^{-n}\}}.
\]
Są to nieujemne funkcje proste, $s_n\le w\chi_U$,
oraz $s_n(x)\to w(x)\chi_U(x)$: dla $x\in U$
i $n>w(x)$ błąd jest nie większy niż $2^{-n}$.
Z lematu Fatou (twierdzenie \ref{thm:fatou})
i monotoniczności całki wynika
$\int_G s_n\,\mathrm d\mu\to\rho(U)$, także
gdy $\rho(U)=+\infty$.
Jeśli $0\le M<\rho(U)$, wybieramy $n$ z
$\int_G s_n\,\mathrm d\mu>M$. Każdy zbiór
$O_k=U\cap\{w>k2^{-n}\}$ jest otwarty.
Z regularności wewnętrznej $\mu$ na otwartych
możemy dobrać zwarte $K_k\subseteq O_k$ tak, aby
$2^{-n}\sum_k\mu(K_k)>M$: jeśli któraś
$\mu(O_k)=\infty$, wystarczy jedna odpowiednio
duża skończona wartość; w przeciwnym razie
przybliżamy skończoną liczbę skończonych wartości
z dostatecznie małym łącznym błędem.
Dla zwartego $K=\bigcup_kK_k\subseteq U$ mamy
$2^{-n}\sum_k\chi_{K_k}\le s_n\chi_K\le w\chi_K$.
Zatem $\rho(K)>M$. Przypadek $\rho(U)=0$
wynika z monotoniczności. To dowodzi regularności
wewnętrznej $\rho$ na otwartych.

\emph{Regularność zewnętrzna.}
Dla borelowskiego $A$ o $\rho(A)<\infty$ rozbijamy
$A$ na rozłączne zbiory
$A_j=A\cap\{2^j\le w<2^{j+1}\}$, $j\in\mathbb Z$,
które pokrywają $A$, gdyż $0<w<\infty$.
Mamy $\mu(A_j)\le2^{-j}\rho(A_j)<\infty$.
Dla $\varepsilon>0$ wybieramy otwarte $O_j\supseteq A_j$
z
\[
 \mu(O_j\setminus A_j)
 <\frac{\varepsilon 2^{-|j|}}{6\cdot2^{j+2}}.
\]
Jest to możliwe z regularności zewnętrznej $\mu$
i skończoności $\mu(A_j)$.
Niech $U_j=O_j\cap\{2^{j-1}<w<2^{j+2}\}$.
Jest to otwarty nadzbiór $A_j$, a
$\rho(U_j\setminus A_j)<\varepsilon2^{-|j|}/6$.
Dla $U=\bigcup_{j\in\mathbb Z}U_j$ mamy $A\subseteq U$
oraz $U\setminus A\subseteq\bigcup_j(U_j\setminus A_j)$.
Ponieważ $\sum_{j\in\mathbb Z}2^{-|j|}=3$,
otrzymujemy $\rho(U\setminus A)\le\varepsilon/2$.
Ponieważ $U=A\sqcup(U\setminus A)$, mamy
$\rho(A)\le\inf_{O\supseteq A,\ O\text{ otwarty}}\rho(O)
\le\rho(U)\le\rho(A)+\varepsilon/2$.
Dowolność $\varepsilon$ daje regularność zewnętrzną.
Dla $\rho(A)=\infty$ regularność zewnętrzna wynika
z monotoniczności. Dowodzi to, że $\rho$ jest
prawostronną miarą Haara, bez żadnego globalnego
założenia $\sigma$-skończoności.

\emph{Stała normalizacyjna.}
Jednoznaczność prawej miary Haara daje
$\mu_R=k\rho$ dla pewnego $k>0$. Zatem dla każdej
nieujemnej borelowskiej $F$ wzór na całkę względem
miary obrazowej (twierdzenie
\atlasref{ME-13::thm:pushforward-integral-13}{ME-13: całka względem miary obrazowej}) daje
\[
 \int_G F(x^{-1})\,\mathrm d\mu(x)
 = k\int_G F(x)w(x)\,\mathrm d\mu(x).
\]
Stosując go najpierw do $F(x)=h(x^{-1})$,
a następnie do $F(y)=h(y)w(y^{-1})$, otrzymujemy
\[
 \int_G h\,\mathrm d\mu
 =k\int_G h(x^{-1})w(x)\,\mathrm d\mu(x)
 =k^2\int_G h(y)w(y^{-1})w(y)\,\mathrm d\mu(y)
 =k^2\int_G h\,\mathrm d\mu .
\]
Wybieramy niezerową nieujemną $h\in C_c(G)$,
dla której całka jest dodatnia i skończona,
zatem $k^2=1$, czyli $k=1$.
To daje \eqref{eq:inwersja}. Dla funkcji rzeczywistych
lub zespolonych stosujemy oba wzory najpierw do
modułu, aby sprawdzić równoważność odpowiednich
warunków całkowalności, a potem do części
dodatnich i ujemnych części rzeczywistej i urojonej.
\end{proof}

\begin{wniosek}[Równoważne warunki unimodularności]
\label{it:cor:haar-unimodular}
Następujące warunki są równoważne:
\begin{enumerate}[label=(\roman*)]
\item $\Delta\equiv1$;
\item lewostronna miara Haara $\mu$ jest także prawostronna;
\item $\mu(E^{-1})=\mu(E)$ dla każdego $E\in\mathcal B(G)$.
\end{enumerate}
Każda przemienna grupa lokalnie zwarta jest unimodularna.
\end{wniosek}
\begin{proof}
Z \eqref{eq:haar-modular-23} wynika (i)$\Rightarrow$(ii).
Jeśli (ii) zachodzi, porównujemy $\mu(Kg)=\mu(K)$
z $\mu(Kg)=\Delta(g)\mu(K)$ dla zwartego $K$
o dodatniej i skończonej mierze i dzielimy przez
$\mu(K)$, uzyskując (i).
Ze wzoru \eqref{eq:inwersja} dla $\chi_E$
wynika (i)$\Rightarrow$(iii).
Jeśli (iii) zachodzi, to $\mu=\mu_R$, która
jest prawostronna, co daje (ii).
W grupie przemiennej lewe i prawe translacje są
identyczne, zatem (ii) jest spełniony.
\end{proof}

\begin{lemat}[Miara Haara z lewostronnie niezmienniczej formy]
\label{it:lem:haar-lie-volume}
Niech $G$ będzie $n$-wymiarową grupą Liego
(w standardowej konwencji rozmaitości Hausdorffa
o przeliczalnej bazie). Wybór niezerowej formy
$\omega_e\in\bigwedge^n(T_eG)^*$ wyznacza
gładką, nigdzie nieznikającą, lewostronnie
niezmienniczą formę $\omega$ na $G$.
Miara objętości wyznaczona przez gęstość $|\omega|$
jest lewostronną miarą Haara.
W szczególności każda grupa Liego jest orientowalna.
\end{lemat}
\begin{proof}
Definiujemy dla $v_1,\ldots,v_n\in T_xG$
\[
 \omega_x(v_1,\ldots,v_n)=
 \omega_e\bigl(\mathrm d(L_{x^{-1}})_xv_1,\ldots,
                \mathrm d(L_{x^{-1}})_xv_n\bigr).
\]
Różniczki lewych translacji są izomorfizmami,
więc forma nigdzie nie znika. Jej gładkość wynika
z gładkości odwzorowania $(x,y)\mapsto x^{-1}y$:
w lokalnych mapach współczynniki różniczki względem
$y$, obliczone w $y=x$, są gładkimi funkcjami $x$.
Dla $a,x\in G$ mamy
$L_{(ax)^{-1}}\circ L_a=L_{x^{-1}}$,
zatem reguła łańcuchowa daje $L_a^*\omega=\omega$.
Nigdzie nieznikająca forma najwyższego stopnia
wyznacza orientację, również na grupie niespójnej.

Dla $n=0$ grupa jest dyskretna, a forma stopnia
zero jest stałą funkcją o wartości $\omega_e\ne0$.
Odpowiadająca jej gęstość wyznacza miarę liczącą
pomnożoną przez $|\omega_e|$, niezerową i regularną
(zbiory zwarte są skończone, a wszystkie podzbiory
są otwarte). Translacje są bijekcjami, więc miara
jest lewostronnie niezmiennicza. To kończy ten przypadek.

Niech dalej $n\ge1$. W mapie $\psi:U\to V\subseteq\mathbb R^n$
forma ma postać $\omega|_U=\psi^*(b\,\mathrm dy_1
\wedge\cdots\wedge\mathrm dy_n)$, gdzie $b$ jest
gładka i nigdzie nie znika. Lokalną miarę na $U$
określamy przez
\[
 \mu_U(E)=\int_{\psi(E)}|b(y)|\,\mathrm d\lambda_n(y),
 \qquad E\in\mathcal B(U).
\]
Dla drugiej mapy $\widetilde\psi$ ze współczynnikiem
$\widetilde b$ niech
$F=\widetilde\psi\circ\psi^{-1}$ na obrazie przecięcia
dziedzin. Z definicji cofnięcia formy najwyższego
stopnia mamy $b(y)=\widetilde b(F(y))\det DF(y)$.
Dla borelowskiego $E$ w przecięciu dziedzin map
twierdzenie \atlasref{ME-13::thm:change-variables-13}{ME-13: zamiana zmiennych} daje
\[
 \int_{\psi(E)}|b(y)|\,\mathrm d\lambda_n(y)
 =\int_{\widetilde\psi(E)}|\widetilde b(z)|
     \,\mathrm d\lambda_n(z).
\]
Jest to zgodność lokalnych miar. Z przeliczalnej bazy wybieramy
przeliczalne pokrycie dziedzinami map $U_j$.
Dla $B_j=U_j\setminus\bigcup_{k<j}U_k$ kładziemy
$\mu_\omega(E)=\sum_j\mu_{U_j}(E\cap B_j)$.
Jest to miara borelowska jako suma miar; na każdej
dziedzinie mapy zgadza się z lokalną miarą, ponieważ
zgodność na przecięciach pozwala zastąpić każdy
składnik miarą z tej mapy i zastosować addytywność.
Każdy punkt ma otoczenie, którego obraz ma zwarte
domknięcie w dziedzinie mapy, więc gęstość jest
na nim ograniczona i miara skończona. Niepusty
otwarty podzbiór mapy ma dodatnią miarę, bo gęstość
jest dodatnia. Rozmaitość o przeliczalnej bazie
jest lokalnie zwarta i metryzowalna, a więc
ośrodkowa; twierdzenie
\ref{thm:lch-secondcountable-radon-20}
daje regularność miary $\mu_\omega$.
Z równości $L_a^*\omega=\omega$ i wzoru zamiany
zmiennych w mapach wynika
$\mu_\omega(aE)=\mu_\omega(E)$: najpierw dla
borelowskich $E$ zawartych w jednej dziedzinie mapy
i takich, że $aE$ także mieści się w jednej dziedzinie,
a następnie dla dowolnego $E$ przez rozłączny
borelowski podział podporządkowany przeliczalnemu
pokryciu przecięciami $U_i\cap a^{-1}U_j$.
Zatem $\mu_\omega$ jest lewostronną miarą Haara,
a twierdzenie \ref{thm:haar-unique-23} pokazuje,
że każda inna jest jej dodatnią wielokrotnością.
\end{proof}

\begin{wniosek}[Modularność grupy Liego]
\label{cor:haar-adjoint-23}
Niech $G$ będzie skończenie wymiarową grupą Liego,
$\mathfrak g=T_eG$ jej algebrą Liego i
$\operatorname{Ad}_g:\mathfrak g\to\mathfrak g$
różniczką w $e$ sprzężenia $x\mapsto gxg^{-1}$.
Dla konwencji modularnej z~\eqref{eq:haar-modular-23}
zachodzi
\[
 c(g)=|\det(\operatorname{Ad}_{g^{-1}})|
       =|\det(\operatorname{Ad}_{g})|^{-1}.
\]
W szczególności grupa Liego jest unimodularna wtedy
i tylko wtedy, gdy $|\det\operatorname{Ad}_g|=1$
dla każdego $g\in G$.
\end{wniosek}
\begin{proof}
Niech $R_g(x)=xg$, $L_x(y)=xy$.
Przez lewą translację identyfikujemy $T_xG$ i
$T_{xg}G$ z $\mathfrak g$. W tych identyfikacjach
różniczka $R_g$ ma postać
\[
 \mathrm d(L_{(xg)^{-1}})_{xg}\circ
 \mathrm d(R_g)_x\circ\mathrm d(L_x)_e
 =\operatorname{Ad}_{g^{-1}},
\]
co sprawdzamy, podstawiając krzywą $h(t)$ wychodzącą
z $e$: $(xg)^{-1}(xh(t)g)=g^{-1}h(t)g$.
Różniczkowanie w $t=0$ daje wskazane odwzorowanie
na dowolnym wektorze $h'(0)\in\mathfrak g$.
Jeśli $\dim G=0$, grupa jest dyskretna i jej miara
licząca jest również prawostronnie niezmiennicza.
Oba wyznaczniki na zerowej przestrzeni wektorowej
są równe $1$, zgodnie z tezą. Dalej zakładamy
$\dim G\ge1$.
Z lematu \ref{it:lem:haar-lie-volume} możemy
obliczyć czynnik modularny dla miary $\mu_\omega$
indukowanej przez lewą formę $\omega$.
Dla $v_i=\mathrm d(L_x)_e u_i$ otrzymujemy
\[
 (R_g^*\omega)_x(v_1,\ldots,v_n)
 =\omega_e(\operatorname{Ad}_{g^{-1}}u_1,\ldots,
            \operatorname{Ad}_{g^{-1}}u_n)
 =\det(\operatorname{Ad}_{g^{-1}})\,
     \omega_x(v_1,\ldots,v_n).
\]
Wzięcie modułu daje czynnik transformacji gęstości
$|\det\operatorname{Ad}_{g^{-1}}|$.
Stosujemy wzór zamiany zmiennych w mapach, a następnie
przeliczalny rozłączny podział jak w poprzednim lemacie:
$\mu_\omega(Eg)=|\det\operatorname{Ad}_{g^{-1}}|
\mu_\omega(E)$ dla każdego borelowskiego $E$.
Każda lewa miara Haara jest dodatnią wielokrotnością
$\mu_\omega$, zatem ma ten sam czynnik modularny.
Ponieważ sprzężenia przez $g$ i $g^{-1}$ są wzajemnie
odwrotne, ich różniczki w $e$ są wzajemnie odwrotne,
co daje drugą równość w tezie.
\end{proof}

\begin{przyklad}[Grupy zwarte i dyskretne]
\label{ex:haar-compact-discrete-23}
Na grupie zwartej $G$ mamy $0<\mu(G)<\infty$,
więc normalizujemy miarę Haara przez $\mu(G)=1$.
Dla każdego $g$ równość $\mu(Gg)=c(g)\mu(G)$
daje $c(g)=1$.
Odwrotnie, jeżeli $G$ nie jest zwarta, wybieramy niepuste
relatywnie zwarte otwarte otoczenie $V$ elementu $e$
i kładziemy $K=\overline V$. Mamy $0<\mu(K)<\infty$.
Istnieje ciąg parami rozłącznych lewych translacji $g_nK$:
po wybraniu $g_1,\ldots,g_n$ wybieramy $g_{n+1}$ poza
zwartym zbiorem $\bigcup_{j=1}^n g_jKK^{-1}$, który nie
może być równy niezwartej grupie $G$. Jeśli
$g_{n+1}K\cap g_jK\ne\varnothing$, to
$g_{n+1}\in g_jKK^{-1}$, wbrew wyborowi.
Z niezmienniczości i przeliczalnej addytywności wynika
$\mu(G)\ge\sum_{n\ge1}\mu(g_nK)=+\infty$.
Zatem lewostronną miarę Haara można unormować do
prawdopodobieństwa dokładnie wtedy, gdy grupa jest zwarta.
Na grupie dyskretnej miarą Haara jest miara licząca:
wszystkie zbiory są jednocześnie otwarte i domknięte,
a jej wartość na dowolnym zbiorze jest supremum
liczności skończonych podzbiorów, czyli zbiorów zwartych.
Translacje są bijekcjami, więc zachowują liczność.
Otrzymujemy także $c\equiv1$. Jeżeli grupa jest skończona,
można ją unormować do prawdopodobieństwa przez
podzielenie miary liczącej przez $|G|$.
\end{przyklad}

\begin{przyklad}[$\mathbb R^d$, torus i grupy macierzowe]
\label{ex:haar-standard-23}
Na $(\mathbb R^d,+)$ lewostronną i prawostronną miarą
Haara jest ograniczenie miary Lebesgue'a $\lambda_d$
do zbiorów borelowskich: jest lokalnie skończona,
więc Radona z twierdzenia
\ref{thm:lch-secondcountable-radon-20}, a translacje
zachowują ją na mocy twierdzenia
\atlasref{ME-13::thm:linear-change-13}{ME-13: liniowa zamiana zmiennych}.
Na torusie $\mathbb T^d=\mathbb R^d/\mathbb Z^d$
miarę Haara dostajemy przez przeniesienie
$\lambda_d|_{[0,1)^d}$ przez odwzorowanie ilorazowe
$\pi:\mathbb R^d\to\mathbb T^d$; jest to miara
probabilistyczna. Dla dowolnego $u\in\mathbb R^d$
i zbioru borelowskiego $E\subseteq\mathbb T^d$
funkcja $\chi_{\pi^{-1}(E)}$ jest $\mathbb Z^d$-okresowa.
Jej całka po $[0,1)^d+u$ jest równa całce po $[0,1)^d$.
Istotnie, dla nieujemnej mierzalnej funkcji
$1$-okresowej $h$ oraz $u=m+r$, $m\in\mathbb Z$,
$0\le r<1$, zachodzi
\[
 \int_{[u,u+1)}h\,\mathrm d\lambda_1
 =\int_{[r,1)}h\,\mathrm d\lambda_1
  +\int_{[0,r)}h\,\mathrm d\lambda_1
 =\int_{[0,1)}h\,\mathrm d\lambda_1.
\]
Przesuwamy najpierw o $-m$, potem rozcinamy przedział
w $1$ i przesuwamy drugi fragment o $-1$.
Dla $\chi_{\pi^{-1}(E)}$ stosujemy ten rachunek
kolejno do każdej współrzędnej, korzystając z twierdzeń
\atlasref{ME-12::thm:tonelli-12}{ME-12: twierdzenie Tonellego} i \atlasref{ME-12::thm:lebesgue-products-12}{ME-12: produkty miar Lebesgue'a}.
Z niezmienniczości miary Lebesgue'a i tej równości
wynika niezmienniczość miary obrazowej względem
translacji o $\pi(u)$. Jest ona miarą Radona na mocy
twierdzenia~\ref{thm:lch-secondcountable-radon-20}. Na grupie
$\operatorname{GL}_d(\mathbb R)$ lewa i prawa miara Haara
mają gęstość $|\det A|^{-d}$ względem miary
Lebesgue'a na przestrzeni $M_d(\mathbb R)\cong\mathbb R^{d^2}$.
Zbiór $\operatorname{GL}_d(\mathbb R)$ jest otwarty
w przestrzeni macierzy; dodatnia ciągła gęstość
jest ograniczona na każdym zbiorze zwartym,
więc miara jest lokalnie skończona i Radona z
twierdzenia \ref{thm:lch-secondcountable-radon-20}.
Dla $B\in\operatorname{GL}_d(\mathbb R)$
odwzorowanie $A\mapsto BA$ działa przez $B$
na każdej z $d$ kolumn, a zatem ma wyznacznik
$({\det B})^d$. Z twierdzenia
\atlasref{ME-13::thm:change-variables-13}{ME-13: zamiana zmiennych} otrzymujemy
\[
 \int_{BE}|\det A|^{-d}\,\mathrm d\lambda_{d^2}(A)
 =\int_E|\det(BA)|^{-d}|\det B|^d
       \,\mathrm d\lambda_{d^2}(A)
 =\int_E|\det A|^{-d}\,\mathrm d\lambda_{d^2}(A).
\]
Mnożenie $A\mapsto AB$ działa analogicznie na
wierszach przez macierz $B^{\mathsf T}$ w zapisie
kolumnowym; ma ten sam moduł wyznacznika.
Otrzymujemy również prawostronną niezmienniczość,
a więc $\operatorname{GL}_d(\mathbb R)$ jest
unimodularna.
\end{przyklad}

\begin{przyklad}[Nieunimodularna grupa afiniczna]
\label{ex:haar-affine-23}
Rozważmy grupę $G=\mathbb R\rtimes(0,\infty)$
z mnożeniem
$(b,a)(c,t)=(b+ac,at)$. Lewą miarą Haara jest
\[
 \mathrm d\mu(b,a)=\frac{\mathrm d\lambda_2(b,a)}{a^2},
\]
a prawą miarą Haara
$\mathrm d\mu_R(b,a)=a^{-1}\mathrm d\lambda_2(b,a)$.
Obie gęstości są dodatnie i ciągłe na otwartej
półpłaszczyźnie, więc odpowiadające miary są
niezerowe, lokalnie skończone i Radona.
Dla lewego i prawego przesunięcia mamy
\[
 L_{(c,t)}(b,a)=(c+tb,ta),\qquad
 R_{(c,t)}(b,a)=(b+ac,at).
\]
Ich macierze różniczek to odpowiednio
$\begin{pmatrix}t&0\\0&t\end{pmatrix}$
i $\begin{pmatrix}1&c\\0&t\end{pmatrix}$,
o wyznacznikach $t^2$ i $t$.
Z twierdzenia \atlasref{ME-13::thm:change-variables-13}{ME-13: zamiana zmiennych}
wynika
\[
 \mu((c,t)E)=\int_E\frac{t^2}{(ta)^2}
       \,\mathrm d\lambda_2(b,a)=\mu(E),
 \qquad
 \mu(E(c,t))=\int_E\frac{t}{(at)^2}
       \,\mathrm d\lambda_2(b,a)=t^{-1}\mu(E).
\]
Stąd $\Delta(c,t)=t^{-1}$, a wzór
\eqref{eq:inwersja} potwierdza
$\mathrm d\mu_R(b,a)=a\,\mathrm d\mu(b,a)
=a^{-1}\mathrm d\lambda_2(b,a)$.
Można też sprawdzić wprost prawostronną
niezmienniczość: $t/(at)=1/a$.
Grupa nie jest unimodularna, bo $\Delta(0,2)=1/2$.
\end{przyklad}



\section{Dodatkowe przykłady całkowania}
\begin{przyklad}[Grupa dyskretna: całka jest sumą]
Na dowolnej grupie dyskretnej zbiory zwarte są dokładnie zbiorami skończonymi, a wszystkie podzbiory są borelowskie. Miarą Haara jest miara licząca $\mu(A)=\#A$ dla skończonego $A$, a $+\infty$ dla nieskończonego. Dla $f\ge0$ mamy $\int_Gf\,\mathrm d\mu=\sum_{x\in G}f(x)$, rozumianą jako supremum skończonych sum; $f\in L^1(G,\mu)$ wtedy i tylko wtedy, gdy $\sum_{x\in G}|f(x)|<\infty$. Równość całki i sumy wynika z definicji całki:
dla każdego skończonego $F\subseteq G$ całka
$f\chi_F$ jest sumą po $F$ (również dla wartości
nieskończonych), co daje dolne oszacowanie.
Jeśli supremum tych sum jest skończone, dla każdego
$n\ge1$ tylko skończenie wiele punktów może spełniać
$f(x)\ge1/n$. Wtedy $f$ znika poza zbiorem
przeliczalnym i z twierdzenia \ref{thm:mct}
dla rosnących skończonych podzbiorów jego nośnika
otrzymujemy równość. Jeśli supremum jest nieskończone,
dolne oszacowanie już wystarcza.
Stosując to do $|f|$, otrzymujemy kryterium $L^1$.
Sama miara licząca na nieprzeliczalnej grupie nie jest
$\sigma$-skończona, ponieważ zbiory skończonej miary
są skończone, a ich przeliczalna suma jest przeliczalna. Dla grupy skończonej normalizacja $\mu(G)=1$ daje średnią $|G|^{-1}\sum_{x\in G}f(x)$.
\end{przyklad}
\begin{przyklad}[Mnożenie liczb dodatnich]
Grupa $G=((0,\infty),\cdot)$ jest przez logarytm izomorficzna topologicznie z $(\mathbb{R},+)$. Obraz miary Lebesgue'a przez $\exp$ jest miarą
Haara: izomorfizm topologiczny zachowuje regularność
i przenosi translacje na mnożenia. Z twierdzenia
\atlasref{ME-13::thm:change-variables-13}{ME-13: zamiana zmiennych} dla logarytmu wynika
$\mathrm d\mu(x)=x^{-1}\,\mathrm d\lambda_1(x)$,
więc $\mu([a,b])=\log(b/a)$ dla $0<a<b$. W szczególności $\int_0^\infty e^{-x}\,\mathrm dx=1$, lecz $\int_{(0,\infty)}e^{-x}\,\mathrm d\mu(x)=\int_0^\infty e^{-x}\,\mathrm dx/x=\infty$. Rozbieżność wynika z $e^{-x}/x\ge e^{-1}/x$
dla $0<x\le1$. Zmiana grupy i miary zmienia
całkowalność tej samej funkcji.
\end{przyklad}
\begin{przyklad}[Okrąg i torus]
Niech $\mathbb T=\mathbb{R}/(2\pi\mathbb Z)$ będzie grupą dodawania kątów modulo $2\pi$. Jest zwarta i jej znormalizowana miara Haara spełnia
$\int_{\mathbb T}f\,\mathrm d\mu=\frac1{2\pi}\int_0^{2\pi}f([t])\,\mathrm dt$. W szczególności $\int_{\mathbb T}\cos(kt)\,\mathrm d\mu=0$ dla $k\in\mathbb Z\setminus\{0\}$ oraz $\int_{\mathbb T}\cos^2t\,\mathrm d\mu=1/2$. Na torusie kątowym $\mathbb T^d=(\mathbb R/(2\pi\mathbb Z))^d$
bierzemy iloczyn znormalizowanych miar Haara.
Z przykładu \ref{ex:haar-standard-23} po przeskalowaniu
współrzędnych i z twierdzenia
\atlasref{ME-13::thm:pushforward-integral-13}{ME-13: całka względem miary obrazowej} mamy
\[
 \int_{\mathbb T^d}f\,\mathrm d\mu
 =(2\pi)^{-d}\int_{[0,2\pi)^d}f([t])\,\mathrm d\lambda_d(t)
\]
dla nieujemnych lub bezwzględnie całkowalnych $f$.
Nawias $[t]$ oznacza klasę w torusie; domknięcie
dziedziny całkowania nie zmienia wartości całki,
ponieważ jej brzeg ma miarę Lebesgue'a zero.
\end{przyklad}
\begin{przyklad}[Kanoniczny rozkład losowego obrotu]
Grupa $\mathrm{SO}(n)=\{R\in\mathbb{R}^{n\times n}:R^{\mathsf T}R=I,\ \det R=1\}$ jest domknięta jako przeciwobraz warunków
$R^{\mathsf T}R=I$ i $\det R=1$ przez odwzorowania
ciągłe oraz ograniczona, ponieważ każda kolumna
ma normę $1$. Jest więc zwarta. Mnożenie i inwersja
($R^{-1}=R^{\mathsf T}$) są ciągłe. Ma więc jedyną miarę Haara o masie $1$, niezmienniczą z obu stron. Dla $n\ge2$ i $v\in\mathbb{R}^n$ zachodzi $\int_{\mathrm{SO}(n)}Rv\,\mathrm d\mu(R)=0$, gdzie całkujemy współrzędnie. Jeśli $m=\int Rv\,\mathrm d\mu(R)$, to dla
każdego $Q\in\mathrm{SO}(n)$ liniowość całki
współrzędnej oraz lewostronna niezmienniczość dają
$Qm=\int QRv\,\mathrm d\mu(R)=m$.
Gdyby $m\ne0$, obrót o kąt $\pi$ w płaszczyźnie
rozpiętej przez $m$ i dowolny niezerowy wektor
prostopadły do $m$ (identyczność na jej dopełnieniu
ortogonalnym) należałby do $\mathrm{SO}(n)$
i przekształcał $m$ na $-m$. Otrzymalibyśmy
sprzeczność. Zatem $m=0$.
\end{przyklad}
\begin{przyklad}[Obliczenia na grupie afinicznej]
Współrzędne $(b,a)$ i miary jak w przykładzie \ref{ex:haar-affine-23}. Dla $A=[0,1]\times[1,2]$, z twierdzenia
\atlasref{ME-12::thm:tonelli-12}{ME-12: twierdzenie Tonellego}, otrzymujemy
\[
 \mu(A)=\int_1^2 a^{-2}\,\mathrm d\lambda_1(a)=1/2,
 \qquad
 \mu_R(A)=\int_1^2 a^{-1}\,\mathrm d\lambda_1(a)=\log2.
\]
Prawe przesunięcie przez $(0,2)$ ma postać
$(b,a)\mapsto(b,2a)$, więc $A(0,2)=[0,1]\times[2,4]$.
Analogicznie całkujemy po $[2,4]$, uzyskując
$\mu(A(0,2))=1/4$ i $\mu_R(A(0,2))=\log2$,
zgodnie z \eqref{eq:haar-modular-23} oraz prawą
niezmienniczością $\mu_R$. 
\end{przyklad}
\begin{przyklad}[Całka Gaussa]
Na grupie $(\mathbb{R},+)$ z miarą Lebesgue'a $\int_\mathbb{R} e^{-x^2}\,\mathrm d\lambda_1(x)=\sqrt\pi$, natomiast $\int_\R1\,\mathrm d\lambda_1=\infty$. Dla kompletności oznaczmy pierwszą całkę przez $I$.
Jest skończona, gdyż dla $|x|\ge1$ mamy
$e^{-x^2}\le e^{-|x|}$. Z twierdzeń
\atlasref{ME-12::thm:tonelli-12}{ME-12: twierdzenie Tonellego} i \atlasref{ME-12::thm:lebesgue-products-12}{ME-12: produkty miar Lebesgue'a}
mamy $I^2=\int_{\mathbb R^2}e^{-(x^2+y^2)}
\,\mathrm d\lambda_2(x,y)$. Odwzorowanie
$(r,\theta)\mapsto(r\cos\theta,r\sin\theta)$
jest dyfeomorfizmem z $(0,\infty)\times(0,2\pi)$
na płaszczyznę bez nieujemnej półosi poziomej;
jego jakobian ma moduł $r$. Pominięta półoś ma
miarę zero z twierdzenia Tonellego. Z twierdzenia
\atlasref{ME-13::thm:change-variables-13}{ME-13: zamiana zmiennych} otrzymujemy
\[
 I^2=\int_0^{2\pi}\int_0^\infty re^{-r^2}
    \,\mathrm d\lambda_1(r)\,\mathrm d\lambda_1(\theta)
 =2\pi\cdot\tfrac12=\pi,
\]
więc $I=\sqrt\pi$ z dodatniości $I$.
\end{przyklad}

\section{Splot na grupie lokalnie zwartej}
\begin{definicja}[Splot na grupie]
\label{it:def:haar-convolution}
Niech $G$ będzie lokalnie zwartą grupą Hausdorffa,
$\mu$ jej lewostronną miarą Haara, a
$\mathbb K=\mathbb R$ lub $\mathbb C$.
Dla $f,g\in C_c(G;\mathbb K)$ definiujemy
\[
 (f*g)(x)=\int_G f(y)g(y^{-1}x)\,\mathrm d\mu(y).
\]
Dla ustalonego $x$ funkcja podcałkowa jest ciągła
i znika poza zwartym nośnikiem $f$, więc całka
istnieje i jest skończona. Na grupie addytywnej
$\mathbb R^d$ otrzymujemy zwykły splot
$(f*g)(x)=\int_{\mathbb R^d}f(y)g(x-y)
\,\mathrm d\lambda_d(y)$.
\end{definicja}

\begin{twierdzenie}[Podstawowe własności splotu]
\label{it:prop:haar-convolution}
Dla $f,g,h\in C_c(G;\mathbb K)$ splot jest
dwuliniowy, $f*g\in C_c(G;\mathbb K)$ oraz
$\operatorname{supp}(f*g)\subseteq
(\operatorname{supp}f)(\operatorname{supp}g)$.
Ponadto
\[
 \int_G(f*g)\,\mathrm d\mu
 =\left(\int_G f\,\mathrm d\mu\right)
  \left(\int_G g\,\mathrm d\mu\right),
 \qquad
 \|f*g\|_1\le\|f\|_1\|g\|_1,
\]
\[
 \|f*g\|_\infty\le\|f\|_1\|g\|_\infty,\qquad
 (f*g)*h=f*(g*h).
\]
\end{twierdzenie}
\begin{proof}
Dwuliniowość wynika z liniowości całki.
Oznaczmy $K=\operatorname{supp}f$,
$L=\operatorname{supp}g$ i $M=\operatorname{supp}h$.
Jeżeli $f(y)g(y^{-1}x)\ne0$, to $y\in K$,
$y^{-1}x\in L$, a więc $x\in KL$. Zbiór $KL$
jest zwarty jako obraz $K\times L$ przez mnożenie,
a zatem domknięty. Stąd nośnik splotu zawiera się
w $KL$. Dla $a\to e$ mamy
\[
 |(f*g)(xa)-(f*g)(x)|
 \le\|f\|_1\sup_{z\in G}|g(za)-g(z)|\longrightarrow0
\]
na mocy lematu \ref{lem:haar-uniform-23},
który w przypadku zespolonym stosujemy do części
rzeczywistej i urojonej. Ponieważ każde $x'$
ma postać $xa$ z $a=x^{-1}x'$, otrzymujemy
ciągłość splotu w każdym $x$.
Nierówność z normą supremum wynika bezpośrednio
z oszacowania $|g(y^{-1}x)|\le\|g\|_\infty$.

Jądro $(y,x)\mapsto f(y)g(y^{-1}x)$ jest ciągłe,
o nośniku zawartym w zwartym $K\times KL$.
Z lematu \ref{lem:haar-fubini-compact-23}
i lewostronnej niezmienniczości mamy
\[
 \int_G(f*g)(x)\,\mathrm d\mu(x)
 =\int_G f(y)\left(\int_G g(y^{-1}x)
                 \,\mathrm d\mu(x)\right)\mathrm d\mu(y)
 =\int_G f\,\mathrm d\mu\int_G g\,\mathrm d\mu.
\]
Punktowa nierówność $|f*g|\le |f|*|g|$ oraz
powyższy wzór dla funkcji nieujemnych dają
nierówność norm $L^1$.

Dla ustalonego $x$ jądro
$(y,z)\mapsto f(y)g(y^{-1}z)h(z^{-1}x)$
jest ciągłe i znika poza $K\times KL$.
Możemy zatem zamienić kolejność całkowania
tym samym lematem Fubiniego. Po podstawieniu
$z=yu$ w całce względem $z$ dostajemy
\begin{align*}
 ((f*g)*h)(x)
 &=\int_G\int_G f(y)g(y^{-1}z)h(z^{-1}x)
       \,\mathrm d\mu(y)\,\mathrm d\mu(z)\\
 &=\int_G f(y)\int_G g(u)h(u^{-1}y^{-1}x)
       \,\mathrm d\mu(u)\,\mathrm d\mu(y)\\
 &= (f*(g*h))(x).
\end{align*}
Nie zakładamy $\sigma$-skończoności całej grupy.
\end{proof}

\begin{wniosek}[Zmiana postaci splotu]
\label{it:cor:haar-convolution-change}
W przyjętej konwencji \eqref{eq:haar-modular-23}
dla $f,g\in C_c(G;\mathbb K)$ zachodzi
\[
 (f*g)(x)=\int_G f(xz^{-1})g(z)\Delta(z)^{-1}
                  \,\mathrm d\mu(z).
\]
Jeżeli grupa jest przemienna, to $f*g=g*f$.
\end{wniosek}
\begin{proof}
W definicji najpierw podstawiamy $y=xu$
(lewa translacja), a następnie stosujemy
wzór inwersji \eqref{eq:inwersja} do funkcji
$z\mapsto f(xz^{-1})g(z)$. Ma ona zwarty nośnik
zawarty w nośniku $g$, więc obie całki są skończone.
W grupie przemiennej $\Delta=1$ oraz
$xz^{-1}=z^{-1}x$, co daje przemienność splotu.
\end{proof}

\begin{przyklad}[Splot nieprzemienny]
Na skończonej nieprzemiennej grupie $G$ z miarą
liczącą wybierzmy $a,b$ takie, że $ab\ne ba$.
Wówczas $\chi_{\{a\}}*\chi_{\{b\}}=\chi_{\{ab\}}$,
ponieważ jedyny niezerowy składnik sumy odpowiada
$y=a$ i warunkowi $a^{-1}x=b$.
W odwrotnej kolejności otrzymujemy
$\chi_{\{ba\}}$, a zatem inny splot.
\end{przyklad}

\begin{twierdzenie}[Algebra splotu $L^1$]
\label{it:thm:haar-l1-convolution}
Splot na $C_c(G;\mathbb K)$ ma dokładnie jedno
ograniczone dwuliniowe rozszerzenie do
$L^1(G,\mu;\mathbb K)\times L^1(G,\mu;\mathbb K)$
o wartościach w $L^1(G,\mu;\mathbb K)$.
Rozszerzenie jest łączne i spełnia
$\|f*g\|_1\le\|f\|_1\|g\|_1$.
Zatem $L^1(G,\mu;\mathbb K)$ ze splotem jest
algebrą Banacha, także gdy $\mu$ nie jest
$\sigma$-skończona.
\end{twierdzenie}
\begin{proof}
Miara Haara jest miarą Radona, więc z wniosku
\atlasref{ME-32::it:cor:radon-cc-l1-general}{ME-32: gęstość funkcji o zwartym nośniku} funkcje z $C_c(G)$
są gęste w $L^1(G,\mu)$. Ich klasy w $L^1$
jednoznacznie wyznaczają funkcje ciągłe:
jeżeli ciągła funkcja $f$ ma normę $L^1$ zero,
to dla każdego $\varepsilon>0$ otwarty zbiór
$\{|f|>\varepsilon\}$ ma miarę zero, a z wniosku
\ref{cor:haar-open-positive-23} musi być pusty.

Wybieramy ciągi $f_n,g_n\in C_c(G;\mathbb K)$
zbieżne w $L^1$ odpowiednio do $f,g$.
Z twierdzenia \ref{it:prop:haar-convolution}
mamy
\[
 \|f_n*g_n-f_m*g_m\|_1
 \le\|f_n-f_m\|_1\|g_n\|_1+
    \|f_m\|_1\|g_n-g_m\|_1.
\]
Ciągi norm są ograniczone, więc prawa strona
dąży do zera przy $m,n\to\infty$. Z zupełności
$L^1$ (twierdzenie \atlasref{ME-14::thm:riesz-fischer-14}{ME-14: zupełność przestrzeni $L^1$})
istnieje granica, którą oznaczamy $f*g$.
To samo oszacowanie dla dwóch różnych par
ciągów przybliżających dowodzi niezależności
od wyboru tych ciągów. Przechodząc do granicy
w nierówności norm, otrzymujemy
$\|f*g\|_1\le\|f\|_1\|g\|_1$.
Dwuliniowość wynika z dwuliniowości splotu
na $C_c(G)$ i liniowości granicy.
Jeśli także $h_n\to h$ w $L^1$, to z otrzymanej
nierówności i dwuliniowości
$(f_n*g_n)*h_n\to(f*g)*h$ oraz
$f_n*(g_n*h_n)\to f*(g*h)$ w $L^1$.
Równość wyrazów obu ciągów daje łączność.
Każde ograniczone dwuliniowe rozszerzenie
musi zachowywać te granice, co dowodzi
jednoznaczności. Zupełność algebry wynika
z zupełności $L^1$.

W tej ogólności definiujemy splot klas $L^1$
przez rozszerzenie ciągłe, nie przez nieuzasadnione
użycie globalnego twierdzenia Fubiniego.
\end{proof}

\section{Dlaczego sama grupa topologiczna nie wystarcza?}
\begin{przyklad}[Brak lokalnie skończonej miary niezmienniczej w $\ell^2$]
\label{it:ex:haar-ell2}
Przestrzeń Hilberta $\ell^2$ nad $\mathbb R$, z dodawaniem
i topologią normy, jest grupą topologiczną Hausdorffa,
ale nie jest lokalnie zwarta. Nie istnieje na niej
niezerowa nieujemna miara borelowska niezmiennicza
względem wszystkich przesunięć i lokalnie skończona.

Istotnie, dodawanie i branie elementu przeciwnego są
ciągłe, ponieważ
$\|(x+y)-(x'+y')\|\le\|x-x'\|+\|y-y'\|$
oraz $\|(-x)-(-x')\|=\|x-x'\|$.
Przestrzeń normowana jest Hausdorffa.
Przypuśćmy, że istnieje miara $\mu$ o wymienionych
własnościach. Z lokalnej skończoności w zerze
wybieramy $r>0$ takie, że $\mu(B(0,r))<\infty$.
Niech $(e_n)$ będzie standardowym ciągiem
ortonormalnym w $\ell^2$, $z_n=(r/2)e_n$ oraz
$\rho=r/8$. Dla $n\ne m$ zachodzi
\[
 \|z_n-z_m\|=\frac{r}{\sqrt2}>2\rho.
\]
Kule $B(z_n,\rho)$ są więc parami rozłączne
(z nierówności trójkąta). Są zawarte w $B(0,r)$,
gdyż dla $x\in B(z_n,\rho)$ mamy
$\|x\|<r/2+r/8<r$.
Z niezmienniczości translacyjnej oraz
przeliczalnej addytywności otrzymujemy
\[
 \infty>\mu(B(0,r))
 \ge\sum_{n=1}^\infty\mu(B(z_n,\rho))
 =\sum_{n=1}^\infty\mu(B(0,\rho)).
\]
Stąd $\mu(B(0,\rho))=0$.

Zbiór $D$ ciągów o skończonym nośniku i wymiernych
współrzędnych jest przeliczalny jako przeliczalna
suma kopii $\mathbb Q^n$. Jest gęsty: dla
$x=(x_j)\in\ell^2$ i $\varepsilon>0$ wybieramy
$N$ z $\sum_{j>N}|x_j|^2<\varepsilon^2/4$,
a następnie wymierne $q_j$ dla $1\le j\le N$
takie, że $\sum_{j=1}^N|q_j-x_j|^2<\varepsilon^2/4$.
Wektor $(q_1,\ldots,q_N,0,\ldots)$ jest odległy
od $x$ o mniej niż $\varepsilon$.
Zatem $\ell^2=\bigcup_{d\in D}(d+B(0,\rho))$.
Wszystkie zbiory tej przeliczalnej rodziny
mają miarę zero, więc $\mu(\ell^2)=0$, wbrew
niezerowości miary. Dowód nie wymaga regularności
ani zupełności miary.

Aby wykazać brak lokalnej zwartości, niech $U$
będzie dowolnym otoczeniem zera i $r>0$ takie,
że $B(0,r)\subseteq U$. Ciąg $((r/2)e_n)$
leży w $U$, a odległość każdych dwóch różnych
jego wyrazów wynosi $r/\sqrt2$. Nie ma więc
podciągu Cauchy'ego, tym bardziej zbieżnego.
Gdyby $\overline U$ było zwarte, to jako
zwarta przestrzeń metryczna byłoby ciągowo zwarte,
co daje sprzeczność. Żadne otoczenie zera nie ma
zatem zwartego domknięcia.
\end{przyklad}

Lokalna zwartość zapewnia funkcje odcinające i kontrolę
na zbiorach zwartych; pokrycia przesunięciami oraz
zwartość iloczynu pozwoliły nam skonstruować
niezmienniczą zawartość na zbiorach zwartych,
a następnie bezpośrednio miarę Radona przez
kryterium Carathéodory'ego.
Całka Lebesgue'a względem tej miary daje całkę Haara.
Twierdzenie Riesza--Markowa--Kakutaniego
(\ref{thm:riesz-positive-21}) daje równoważne
funkcjonałowe ujęcie tej całki. Nie trzeba zakładać
na grupie struktury różniczkowej.

\section{Zmiana miary Haara przy automorfizmie grupy}
\begin{twierdzenie}[Moduł automorfizmu]
\label{tw:haar-automorfizm}
Niech $G$ będzie lokalnie zwartą grupą Hausdorffa
z lewostronną miarą Haara $\mu$, a $T:G\to G$
automorfizmem topologicznym (automorfizmem grupy
będącym homeomorfizmem).
Istnieje jednoznaczna liczba $c_T>0$ taka, że
\[
 \mu(T(A))=c_T\mu(A),\qquad A\in\mathcal B(G).
\]
Liczba $c_T$ nie zależy od normalizacji $\mu$.
Dla dowolnych automorfizmów topologicznych $S,T$
zachodzą równości
\[
 c_{T\circ S}=c_Tc_S,\qquad c_{\operatorname{id}_G}=1,
 \qquad c_{T^{-1}}=c_T^{-1}.
\]
Ponadto dla nieujemnej funkcji borelowskiej $f$
\begin{equation}\label{it:eq:haar-automorphism-integral}
 \int_G f(T(x))\,\mathrm d\mu(x)
 =c_T^{-1}\int_G f(y)\,\mathrm d\mu(y).
\end{equation}
Wzór obowiązuje także dla funkcji rzeczywistych
i zespolonych bezwzględnie całkowalnych.
\end{twierdzenie}
\begin{proof}
Miara $\nu(A)=\mu(T(A))$ jest obrazem $\mu$ przez
$T^{-1}$, a więc jest niezerową miarą borelowską.
Homeomorfizm przekształca bijektywnie zwarte
podzbiory zbioru otwartego $U$ na zwarte podzbiory
$T(U)$, a otwarte nadzbiory $A$ na otwarte
nadzbiory $T(A)$. Zachowuje zatem regularność
wewnętrzną na otwartych i regularność zewnętrzną
na borelowskich. Skończoność $\nu$ na zwartych
wynika ze zwartości ich obrazów przez $T$.
Dla $g\in G$ mamy
$\nu(gA)=\mu(T(g)T(A))=\mu(T(A))=\nu(A)$.
Zatem $\nu$ jest lewostronną miarą Haara.
Twierdzenie \ref{thm:haar-unique-23} daje
$\nu=c_T\mu$ dla dokładnie jednej stałej $c_T>0$.

Wybieramy zwarty zbiór $K$ o dodatniej skończonej
mierze (domknięcie niepustego relatywnie zwartego
otoczenia $e$, wniosek
\ref{cor:haar-open-positive-23}). Zatem
$c_T=\mu(T(K))/\mu(K)$, co pokazuje niezależność
od przemnożenia $\mu$ przez dodatnią stałą.
Dla złożenia automorfizmów mamy
\[
 c_{T\circ S}\mu(K)=\mu(T(S(K)))
 =c_T\mu(S(K))=c_Tc_S\mu(K).
\]
Dzieląc przez $\mu(K)$, otrzymujemy multiplikatywność.
Dla identyczności mamy $c_{\operatorname{id}_G}=1$,
a stosując multiplikatywność do $T\circ T^{-1}$,
otrzymujemy $c_{T^{-1}}=c_T^{-1}$.
W ten sposób moduł jest homomorfizmem grupy
automorfizmów topologicznych do $(0,\infty)$;
nie zakładamy tutaj żadnej topologii na grupie
automorfizmów.

Obraz miary $\mu$ przez $T$ spełnia
$T_\#\mu(A)=\mu(T^{-1}(A))=c_T^{-1}\mu(A)$.
Twierdzenie \atlasref{ME-13::thm:pushforward-integral-13}{ME-13: całka względem miary obrazowej}
daje więc \eqref{it:eq:haar-automorphism-integral}.
Stosując wzór do $|f|$, otrzymujemy równoważność
całkowalności $f\circ T$ i $f$; następnie stosujemy
go do części dodatnich i ujemnych części
rzeczywistej i urojonej.
\end{proof}

\begin{wniosek}[Automorfizmy wewnętrzne i funkcja modularna]
\label{it:cor:haar-inner-module}
Niech $I_g(x)=gxg^{-1}$. W konwencji
\eqref{eq:haar-modular-23} zachodzi
\[
 c_{I_g}=\Delta(g)^{-1}.
\]
Ponadto dla każdego automorfizmu topologicznego $T$
mamy $\Delta(T(g))=\Delta(g)$.
\end{wniosek}
\begin{proof}
Z lewostronnej niezmienniczości i definicji
$\Delta$ otrzymujemy
$\mu(I_g(A))=\mu(Ag^{-1})=\Delta(g)^{-1}\mu(A)$,
co daje pierwszy wzór.
Dla zwartego $K$ o dodatniej skończonej mierze
zbiór $T(K)$ także ma dodatnią skończoną miarę,
gdyż $\mu(T(K))=c_T\mu(K)$. Porównujemy
\[
 c_T\Delta(g)\mu(K)=\mu(T(Kg))
 =\mu(T(K)T(g))
 =\Delta(T(g))c_T\mu(K).
\]
Dzieląc przez dodatnią skończoną liczbę
$c_T\mu(K)$, dostajemy żądaną równość.
\end{proof}

\begin{wniosek}[Moduł gładkiego automorfizmu grupy Liego]
\label{it:cor:haar-lie-automorphism}
Jeżeli $G$ jest skończenie wymiarową grupą Liego,
a $T$ jej gładkim automorfizmem z gładką odwrotnością,
to
\[
 c_T=|\det(\mathrm dT_e)|.
\]
\end{wniosek}
\begin{proof}
Jeżeli $\dim G=0$, automorfizm jest bijekcją grupy
dyskretnej, więc zachowuje miarę liczącą i $c_T=1$.
Wyznacznik na zerowej przestrzeni wektorowej jest
również równy $1$. W pozostałych przypadkach
korzystamy z lewostronnie niezmienniczej formy
najwyższego stopnia $\omega$ z lematu
\ref{it:lem:haar-lie-volume}.
Ponieważ $T\circ L_x=L_{T(x)}\circ T$,
reguła łańcuchowa daje
$\mathrm dT_x\circ\mathrm d(L_x)_e
=\mathrm d(L_{T(x)})_e\circ\mathrm dT_e$.
Dla dowolnych $u_i\in T_eG$ mamy więc
\[
 (T^*\omega)_x\bigl(\mathrm d(L_x)_e u_1,\ldots,
                       \mathrm d(L_x)_e u_n\bigr)
 =\omega_e(\mathrm dT_e u_1,\ldots,\mathrm dT_e u_n)
 =\det(\mathrm dT_e)\omega_e(u_1,\ldots,u_n).
\]
Stąd $T^*\omega=\det(\mathrm dT_e)\omega$.
Wzięcie modułu i zamiana zmiennych w mapach,
tak jak w dowodzie wniosku \ref{cor:haar-adjoint-23},
dają $\mu_\omega(T(A))=|\det(\mathrm dT_e)|\mu_\omega(A)$.
Niezależność modułu od normalizacji miary Haara
kończy dowód.
\end{proof}

\begin{przyklad}[Jakobian jako moduł automorfizmu]
Dla grupy $(\mathbb R^d,+)$ i odwracalnej macierzy
$A\in\mathbb R^{d\times d}$ automorfizm $T_A(x)=Ax$
spełnia $c_{T_A}=|\det A|$ na mocy twierdzenia
\atlasref{ME-13::thm:linear-change-13}{ME-13: liniowa zamiana zmiennych} (lub wniosku
\ref{it:cor:haar-lie-automorphism}). Zatem
\[
 \int_{\mathbb R^d} f(Ax)\,\mathrm d\lambda_d(x)
 =|\det A|^{-1}\int_{\mathbb R^d} f(y)\,\mathrm d\lambda_d(y).
\]
Nie należy mylić modułu automorfizmu z funkcją
modularną $\Delta$, która dotyczy prawych translacji:
na grupie przemiennej $(\mathbb R^d,+)$ mamy
$\Delta\equiv1$, mimo że moduły automorfizmów
mogą przyjmować dowolne dodatnie wartości
(dla $d\ge1$ wystarczy $A=tI$, $t>0$).
\end{przyklad}

% ROZDZIAŁ 24 — TWIERDZENIA POKRYCIOWE
\chapter{Twierdzenia pokryciowe Vitaliego i Besicovitcha}
\label{ch:covering-24}

\section{Motywacja i terminologia}
W poprzednich rozdziałach zbiory mierzalne przybliżaliśmy
zbiorami otwartymi i zwartymi. W zastosowaniach do
różniczkowania całek i miar potrzebujemy mocniejszego narzędzia:
z dowolnej rodziny małych kul należy wybrać takie kule, których
nakładanie można kontrolować. W tym rozdziale podajemy trzy
odmienne rezultaty: lemat $5r$, twierdzenie Vitaliego i twierdzenie
Besicovitcha. Nie wolno ich utożsamiać: pierwsze dotyczy
\emph{powiększonych} kul, drugie pokrycia \emph{z dokładnością do
zbioru miary zero}, a trzecie pokrycia kulami \emph{bez powiększania},
z ograniczoną krotnością.

\begin{definicja}[Kule i krotność pokrycia]
\label{def:balls-24}
Niech $(X,d)$ będzie przestrzenią metryczną. Piszemy
$B(x,r)=\{y\in X:d(x,y)<r\}$ dla $r>0$.
Kulę traktujemy jako zbiór wraz z zadanym środkiem i promieniem:
w ogólnej przestrzeni metrycznej różne pary $(x,r)$ mogą
wyznaczać ten sam zbiór. Jeśli $B=B(x,r)$ i $a>0$, to
$aB:=B(x,ar)$ dla tej konkretnej pary.
Rodzina $(B_i)_{i\in I}$ jest \emph{parami rozłączna}, jeżeli
$B_i\cap B_j=\varnothing$ dla $i\ne j$.
Jej \emph{krotność w punkcie} $x$ to liczba indeksów $i$ takich,
że $x\in B_i$ (dopuszczamy $+\infty$).
Funkcja charakterystyczna dowolnego zbioru $A$ jest w całym
skrypcie oznaczana przez $\chi_A$.
\end{definicja}

\begin{uwaga}
W całym rozdziale $d\in\mathbb N$, $d\ge1$,
gdy oznacza wymiar przestrzeni euklidesowej.
Twierdzenia sformułowane dla kul w $\mathbb R^d$ odnoszą się
zawsze do metryki euklidesowej i miary Lebesgue'a $\lambda_d$.
Przypadki metryczne bez miary rozpatrujemy oddzielnie.
Symbol $v_d:=\lambda_d(B(0,1))$ oznacza objętość kuli jednostkowej;
$\lambda_d(B(x,r))=v_dr^d$.
\end{uwaga}

\section{Lemat $5r$: wybór kul bez zakładania miary}
\begin{twierdzenie}[Lemat o pięciokrotnym powiększeniu]
\label{thm:5r-24}
Niech $(X,d)$ będzie przestrzenią metryczną, a $\mathcal B$
niepustą rodziną otwartych kul o promieniach ograniczonych
z góry przez pewną liczbę $R<\infty$.
Istnieje parami rozłączna podrodzina $\mathcal D\subseteq\mathcal B$
taka, że
\[
 \bigcup_{B\in\mathcal B}B\subseteq\bigcup_{D\in\mathcal D}5D.
\]
Jeżeli $X$ jest ośrodkowa, rodzina $\mathcal D$ jest co najwyżej
przeliczalna.
\end{twierdzenie}
\begin{proof}
Rozdzielmy kule na klasy według promienia:
$\mathcal B_k=\{B(x,r)\in\mathcal B:
2^{-k-1}R<r\le2^{-k}R\}$, $k=0,1,\ldots$.
W klasie $\mathcal B_0$ wybieramy maksymalną względem inkluzji
rodzinę parami rozłącznych kul. Indukcyjnie, po dokonaniu wyborów
z klas wcześniejszych, w $\mathcal B_k$ wybieramy maksymalną
rodzinę parami rozłącznych kul, rozłącznych również z kulami
wybranymi wcześniej. Maksymalność można uzyskać z lematu
Zorna: uporządkowujemy przez inkluzję rodziny spełniające
warunek rozłączności, a suma łańcucha pozostaje taką rodziną.
Niech $\mathcal D$ będzie sumą wszystkich wybranych rodzin.

Weźmy kulę $B=B(x,r)\in\mathcal B_k$, która nie została
wybrana. Na odpowiednim etapie była albo niedopuszczalna
wskutek przecięcia z wcześniej wybraną kulą, albo przecinała
kulę z maksymalnego wyboru w swojej klasie. W obu przypadkach
istnieje $D=B(y,s)\in\mathcal D$ taka, że $B\cap D\ne\varnothing$
i $s>2^{-k-1}R\ge r/2$. Dla każdego $z\in B$ mamy
\[
 d(z,y)\le d(z,x)+d(x,y)<r+(r+s)=2r+s<5s.
\]
Stąd $B\subseteq5D$. Teza jest oczywista także dla kul
wybranych, bo $D\subseteq5D$.

Jeżeli $X$ ma przeliczalny zbiór gęsty $Q$, to każda wybrana
kula zawiera punkt $Q$. Kule z $\mathcal D$ są rozłączne,
więc przyporządkowanie im różnych punktów $Q$ dowodzi
przeliczalności $\mathcal D$.
\end{proof}

\begin{wniosek}[Oszacowanie miary sumy kul]
\label{cor:5r-volume-24}
Niech $\mathcal B$ będzie rodziną kul euklidesowych o promieniach
ograniczonych z góry. Istnieje parami rozłączna podrodzina
$(B_i)$ taka, że
\[
 \lambda_d\Bigl(\bigcup_{B\in\mathcal B}B\Bigr)
 \le 5^d\sum_i\lambda_d(B_i).
\]
\end{wniosek}
\begin{proof}
Stosujemy twierdzenie~\ref{thm:5r-24}, przeliczalną
podaddytywność oraz równość $\lambda_d(5B_i)=5^d\lambda_d(B_i)$.
Ponieważ suma rozłącznych kul jest mierzalna, szereg po prawej
jest dobrze określony również wtedy, gdy ma wartość $+\infty$.
\end{proof}

\begin{uwaga}[Inne stałe]
W przypadku kul w $\mathbb R^d$ często formułuje się wariant
$3r$, jeżeli promienie porządkuje się dokładniej. Stała $5$
jest wygodna, gdy nie zakłada się osiągania supremum promieni;
jej konkretna wartość nie jest istotna w większości zastosowań.
Nie należy jednak z lematu $5r$ wyciągać wniosku, że same
wybrane kule pokrywają całą początkową rodzinę.
\end{uwaga}

\begin{center}
\begin{tikzpicture}[scale=0.90]
  \draw[blue!50!black,densely dashed] (0,0) circle (2);
  \draw[blue!70!black,thick] (0,0) circle (0.4);
  \draw[gray!80,thick] (0.8,0) circle (0.5);
  \fill (0,0) circle (1.6pt);
  \fill (0.8,0) circle (1.6pt);
  \node[font=\small,below] at (0,-0.4) {$D$};
  \node[font=\small,below] at (0.8,-0.5) {$B$};
  \node[font=\small,above] at (0,2) {$5D$};
  \draw[blue!70!black,thick] (3.9,0) circle (0.35);
  \draw[blue!50!black,densely dashed] (3.9,0) circle (1.75);
  \fill (3.9,0) circle (1.6pt);
  \node[font=\small,below] at (3.9,-0.35) {$D'$};
  \node[font=\small,above] at (3.9,1.75) {$5D'$};
\end{tikzpicture}
\end{center}
\noindent\textit{Idea lematu $5r$:} wybrane kule $D,D'$
są rozłączne. Odrzucona kula $B$ przecina $D$, a jej
promień jest mniejszy od dwukrotności promienia $D$.
Stąd $B\subseteq5D$. Kreskowane okręgi mają dokładnie
pięciokrotnie większe promienie niż odpowiadające im
okręgi niebieskie ciągłe.

\section{Pokrycia drobne i twierdzenie Vitaliego}
\begin{definicja}[Drobne pokrycie Vitaliego]
\label{def:vitali-fine-24}
Niech $E\subseteq\mathbb R^d$ i $\mathcal V$ będzie rodziną
otwartych kul euklidesowych. Mówimy, że $\mathcal V$ jest
\emph{drobna w każdym punkcie $E$}, jeżeli dla wszystkich
$x\in E$ i $\varepsilon>0$ istnieje kula
$B(x,r)\in\mathcal V$ ze środkiem w $x$ i $0<r<\varepsilon$.
Wymaganie środka w punkcie $x$ jest tutaj częścią definicji.
\end{definicja}

\begin{twierdzenie}[Pokryciowe twierdzenie Vitaliego]
\label{thm:vitali-cover-24}
Niech $E\subseteq\mathbb R^d$ będzie ograniczony,
$O\subseteq\mathbb R^d$ otwarty i ograniczony, $E\subseteq O$,
a $\mathcal V$ rodziną kul drobną w każdym punkcie $E$.
Wówczas można wybrać skończoną lub przeliczalną rodzinę
parami rozłącznych kul $(B_j)$ z $\mathcal V$, spełniającą
$B_j\subseteq O$ oraz
\[
 \lambda_d^*\Bigl(E\setminus\bigcup_jB_j\Bigr)=0,
 \qquad \sum_j\lambda_d(B_j)\le\lambda_d(O).
\]
Zapis $\lambda_d^*$ oznacza zewnętrzną miarę Lebesgue'a;
nie zakładamy z góry mierzalności zbioru $E$.
\end{twierdzenie}
\begin{proof}
Ograniczamy rodzinę do tych kul, których domknięcia mieszczą
się w $O$. Nadal jest ona drobna w każdym punkcie $E$:
przy ustalonym $x\in E$ wybieramy promień mniejszy od
odległości $x$ od $\mathbb R^d\setminus O$.
Wszystkie takie kule mieszczą się w $O$, więc ich promienie
są jednostajnie ograniczone.

Budujemy indukcyjnie rozłączne kule. Jeżeli po wybraniu
$B_1,\ldots,B_{n-1}$ pozostaje choć jedna kula naszej rodziny
rozłączna ze wszystkimi wcześniej wybranymi, oznaczmy przez
$s_n$ supremum promieni takich kul. Wybierzmy jedną z nich,
$B_n=B(x_n,r_n)$, o promieniu $r_n>s_n/2$. Jeżeli nie ma
dopuszczalnych kul, kończymy. Konstrukcja daje
$\sum_jv_dr_j^d\le\lambda_d(O)<\infty$, a zatem
$r_j\to0$ w przypadku nieskończonego ciągu.

Jeżeli konstrukcja kończy się po $m$ wyborach (także gdy
$m=0$), to $E\subseteq\bigcup_{j=1}^m\overline{B_j}$.
Istotnie, punkt $x$ spoza tej skończonej sumy domkniętych
kul ma otoczenie rozłączne ze wszystkimi wybranymi kulami.
Drobność daje w nim dopuszczalną kulę o środku $x$ i domknięciu
zawartym w $O$, wbrew zakończeniu konstrukcji.
W tym przypadku zbiór niepokryty kulami otwartymi zawiera
się w skończonej sumie ich sfer i ma miarę zewnętrzną zero.

Załóżmy dalej, że konstrukcja trwa nieskończenie długo.
Niech $N$ będzie dowolną dodatnią liczbą całkowitą.
Pokażemy inkluzję
\begin{equation}\label{eq:vitali-tail-24}
 E\setminus\bigcup_j\overline{B_j}
 \subseteq\bigcup_{j>N}5B_j.
\end{equation}
Rzeczywiście, ustalmy $x\in E\setminus\bigcup_j\overline{B_j}$.
Jest on w dodatniej odległości od skończonej sumy
$\bigcup_{j\le N}\overline{B_j}$. Z drobności pokrycia
wybieramy kulę $B=B(x,r)\in\mathcal V$ o domknięciu
w $O$, rozłączną z $B_1,\ldots,B_N$.
Kula $B$ musi przecinać którąś z późniejszych kul. Gdyby bowiem była rozłączna
ze wszystkimi $B_j$, pozostawałaby dopuszczalna na
każdym etapie. Otrzymalibyśmy $s_j\ge r$ i
$r_j>s_j/2\ge r/2$ dla wszystkich $j>N$,
sprzecznie z $\sum_jr_j^d<\infty$.
Niech zatem $j>N$ będzie pierwszym indeksem, dla
którego $B\cap B_j\ne\varnothing$.
Przed krokiem $j$ kula $B$ pozostawała dopuszczalna,
zatem $r\le s_j<2r_j$. Odległość centrów jest mniejsza
od $r+r_j$, a więc $d(x,x_j)<3r_j<5r_j$.
Otrzymaliśmy $x\in5B_j$, co daje \eqref{eq:vitali-tail-24}.

Niech $Z=\bigcup_{j\ge1}\partial B_j$.
Każda sfera $\partial B_j$ ma zerową miarę Lebesgue'a:
dla $0<\delta<r_j$ zawiera się w powłoce o mierze
$v_d((r_j+\delta)^d-(r_j-\delta)^d)$, dążącej do zera
przy $\delta\to0$. Z przeliczalnej podaddytywności
otrzymujemy $\lambda_d(Z)=0$.
Ze \eqref{eq:vitali-tail-24}, podaddytywności i skalowania
miary wynika
\[
 \lambda_d^*\Bigl(E\setminus\bigcup_jB_j\Bigr)
 \le\lambda_d(Z)+5^d\sum_{j>N}\lambda_d(B_j).
\]
Po przejściu $N\to\infty$ prawa strona dąży do zera.
Suma objętości nie przekracza $\lambda_d(O)$, bo kule są
parami rozłączne i zawarte w $O$. To kończy dowód.
\end{proof}

\begin{uwaga}[Zakres założeń twierdzenia Vitaliego]
Drobność w każdym punkcie $E$ można zastąpić drobnością
poza zbiorem $Z\subseteq E$ o $\lambda_d^*(Z)=0$:
stosujemy twierdzenie do $E\setminus Z$, a następnie
korzystamy z podaddytywności miary zewnętrznej.
Ograniczoność $O$ zapewnia skończoność $\lambda_d(O)$
oraz ograniczenie promieni. Jeśli $O$ jest nieograniczony,
ale ma skończoną miarę, te same dwa warunki nadal zachodzą:
$v_dr^d\le\lambda_d(O)$ dla każdej kuli $B(x,r)\subseteq O$.
Dla dowolnego otwartego $O$ można dodatkowo zastosować
twierdzenie osobno w rozłącznych otwartych warstwach
$O\cap\{x:n<|x|<n+1\}$, $n=0,1,\ldots$.
Pominięte sfery $\{|x|=n\}$ tworzą zbiór miary zero,
a kule z różnych warstw są rozłączne. Powyższa wersja
wystarcza do różniczkowania miary Lebesgue'a.
\end{uwaga}

\begin{wniosek}[Przybliżanie zbioru rozłącznymi kulami]
\label{cor:vitali-open-24}
Jeżeli $O\subseteq\mathbb R^d$ jest niepustym,
ograniczonym zbiorem otwartym, to istnieją parami rozłączne
kule $B_j\subseteq O$ takie, że
$\lambda_d(O\setminus\bigcup_jB_j)=0$.
\end{wniosek}
\begin{proof}
Stosujemy twierdzenie~\ref{thm:vitali-cover-24} do $E=O$
i rodziny wszystkich kul o domknięciu zawartym w $O$.
Ta rodzina jest drobna w każdym punkcie $O$.
\end{proof}

\section{Geometria ograniczonej krotności: lemat kątowy}
\begin{lemat}[Skończoność konfiguracji Besicovitcha]
\label{lem:besicovitch-angle-24}
Niech $B_i=B(x_i,r_i)$, $i=1,\ldots,m$, będą kulami
w $\mathbb R^d$, których część wspólna zawiera punkt $p$.
Załóżmy ponadto, że dla $i\ne j$ zachodzi
$x_i\notin B_j$ i $x_j\notin B_i$.
Wówczas $m\le N_d$, gdzie $N_d$ zależy wyłącznie od $d$;
można przyjąć $N_d=3^d$.
\end{lemat}
\begin{proof}
Jeżeli $x_i=p$ dla pewnego $i$, to dla każdego $j\ne i$
mielibyśmy $x_i=p\in B_j$, wbrew założeniu.
W tym przypadku $m=1$ i teza zachodzi.
Możemy więc założyć, że $x_i\ne p$ dla każdego $i$.
Dla pozostałych stawiamy $u_i=(x_i-p)/|x_i-p|\in S^{d-1}$.
Ponieważ $p\in B_i$, mamy $|x_i-p|<r_i$.
Z warunku braku centrów w innych kulach wynika
$|x_i-x_j|\ge\max\{r_i,r_j\}>
\max\{|x_i-p|,|x_j-p|\}$.
Z twierdzenia cosinusów
\[
 2\langle x_i-p,x_j-p\rangle
 =|x_i-p|^2+|x_j-p|^2-|x_i-x_j|^2
 <\min\{|x_i-p|^2,|x_j-p|^2\}.
\]
Dzieląc przez $2|x_i-p||x_j-p|>0$, otrzymujemy
\[
 \langle u_i,u_j\rangle
 <\frac{\min\{|x_i-p|,|x_j-p|\}}
 {2\max\{|x_i-p|,|x_j-p|\}}\le\frac12.
\]
Zatem
$|u_i-u_j|>1$. Kulki $B(u_i,1/2)$ w $\mathbb R^d$
są parami rozłączne i leżą w $B(0,3/2)$.
Porównanie objętości daje
$m v_d(1/2)^d\le v_d(3/2)^d$, skąd $m\le3^d$.
\end{proof}

\begin{twierdzenie}[Skończona postać lematu Besicovitcha]
\label{thm:besicovitch-finite-24}
Niech $x_1,\ldots,x_M$ będą punktami $\mathbb R^d$,
a $r_1,\ldots,r_M>0$. Istnieje podrodzina kul
$B(x_i,r_i)$, która pokrywa wszystkie punkty
$x_1,\ldots,x_M$ i której krotność w żadnym punkcie
$\mathbb R^d$ nie przekracza $N_d=3^d$.
\end{twierdzenie}
\begin{proof}
Porządkujemy kule nierosnąco według promienia (dla remisów
wybieramy dowolny porządek). Przechodzimy po kolei przez
ich listę i wybieramy kulę, jeśli jej środek nie należy
do sumy kul wybranych wcześniej. Każdy środek jest ostatecznie
pokryty: albo przez własną wybraną kulę, albo przez jedną
z poprzednich. Jeśli wybrano $B_i$ wcześniej niż $B_j$,
to $x_j\notin B_i$, a stąd $|x_i-x_j|\ge r_i\ge r_j$
i także $x_i\notin B_j$. Wybrane kule spełniają więc
warunek lematu~\ref{lem:besicovitch-angle-24}.
Dla dowolnego punktu $p$ kule wybrane zawierające $p$
tworzą konfigurację z tego lematu, zatem jest ich najwyżej
$N_d$.
\end{proof}

\begin{wniosek}[Oszacowanie sumy miar w rodzinie skończonej]
\label{cor:finite-besicovitch-measure-24}
W założeniach twierdzenia~\ref{thm:besicovitch-finite-24}
niech $B_{i_1},\ldots,B_{i_s}$ będzie wybraną podrodziną.
Dla dowolnej dodatniej miary borelowskiej $\rho$
na $\mathbb R^d$ zachodzi
\[
 \sum_{j=1}^s\rho(B_{i_j})
 \le N_d\,\rho\Bigl(\bigcup_{j=1}^s B_{i_j}\Bigr).
\]
Teza pozostaje prawdziwa, gdy prawa strona jest nieskończona.
\end{wniosek}
\begin{proof}
Definicja krotności i skończona addytywność całki dają
\[
 \sum_{j=1}^s\rho(B_{i_j})
 =\int_{\mathbb R^d}\sum_{j=1}^s\chi_{B_{i_j}}(x)
 \,\mathrm d\rho(x)
 \le N_d\int_{\mathbb R^d}
 \chi_{\cup_jB_{i_j}}(x)\,\mathrm d\rho(x).
\]
To ostatnie wyrażenie jest równe prawej stronie tezy.
\end{proof}

\begin{twierdzenie}[Twierdzenie pokryciowe Besicovitcha -- pełny dowód]
\label{thm:besicovitch-24}
Niech $d\ge1$, niech $E\subseteq\mathbb R^d$ będzie
zbiorem ograniczonym i niech każdemu $x\in E$
przyporządkowano promień $r_x\in(0,R]$, gdzie
$R<\infty$. Wówczas z rodziny kul $\{B(x,r_x):x\in E\}$
można wybrać co najwyżej przeliczalną podrodzinę
$(B_j)_{j\in J}$ taką, że
\begin{equation}\label{eq:besicovitch-infinite-24}
 E\subseteq\bigcup_{j\in J}B_j,\qquad
 \sum_{j\in J}\chi_{B_j}(z)\le15^d
 \quad(z\in\mathbb R^d).
\end{equation}
Jeżeli dla każdego środka $x$ zadano kilka dopuszczalnych
kul o promieniach nie większych niż $R$, wolno najpierw
wybrać jedną z nich i zastosować twierdzenie.
Nie zakładamy mierzalności zbioru $E$.
\end{twierdzenie}
\begin{proof}
\emph{Krok 1: podział na skale.}
Dla $k=0,1,\ldots$ oznaczmy $a_k=2^{-k-1}R$ i
$E_k^0=\{x\in E:a_k<r_x\le2a_k\}$. Zbiory
$E_k^0$ tworzą rozłączny podział $E$.
Będziemy kolejno wybierać skończone zbiory środków
$S_k\subseteq E_k^0$; wybranej rodzinie w skali $k$
odpowiadają kule $\{B(x,r_x):x\in S_k\}$.

Przypuśćmy, że dokonano już wyboru dla skal
$0,\ldots,k-1$, i niech $O_{k-1}$ oznacza sumę
wszystkich wcześniej wybranych kul (dla $k=0$
przyjmujemy $O_{-1}=\varnothing$).
Niech $E_k=E_k^0\setminus O_{k-1}$. Spośród środków
$E_k$ wybieramy maksymalny zbiór $S_k$ o własności
\begin{equation}\label{eq:besicovitch-separated-24}
 |x-y|\ge a_k\quad(x,y\in S_k,\ x\ne y).
\end{equation}
Jeżeli $E_k=\varnothing$, przyjmujemy $S_k=\varnothing$.
W przeciwnym razie taki zbiór można zbudować zachłannie,
bez lematu Zorna: zaczynamy od dowolnego środka, a następnie
dodajemy kolejny, jeśli jego odległość od każdego
wcześniej wybranego jest co najmniej $a_k$.
Ponieważ $E$ jest ograniczony, każde dwa punkty
wybranego zbioru są oddalone o co najmniej $a_k$,
a kule $B(x,a_k/2)$ o tych środkach są rozłączne
i jeśli $E\subseteq B(0,L)$ dla $L>0$, leżą w
$B(0,L+a_k/2)$. Po $n$ wyborach porównanie objętości
daje $n(a_k/2)^d\le (L+a_k/2)^d$, a więc
$n\le(1+2L/a_k)^d$. Ogranicza to z góry liczbę kroków;
procedura kończy się więc po skończenie wielu
wyborach, jeżeli $E_k\ne\varnothing$.

Z maksymalności wynika, że każdy $x\in E_k$
albo należy do $S_k$, albo istnieje $y\in S_k$
z $|x-y|<a_k<r_y$. Toteż wszystkie środki z
$E_k$ należą do kul wybranych w skali $k$.
W konsekwencji każdy $x\in E$ zostaje pokryty:
jeżeli $x\in E_k^0$, to albo znajdował się już
w $O_{k-1}$, albo zostaje pokryty w kroku $k$.
Suma $\bigcup_{k\ge0}S_k$ jest przeliczalna,
gdyż wszystkie $S_k$ są skończone.

\emph{Krok 2: liczba kul jednej skali.}
Ustalmy $z\in\mathbb R^d$ i skalę $k$.
Jeżeli $B(x,r_x)$ jest wybraną kulą skali $k$
zawierającą $z$, to $|x-z|<r_x\le2a_k$.
Z \eqref{eq:besicovitch-separated-24} wynika, że
kule $B(x,a_k/2)$ o takich środkach są rozłączne,
a każda z nich mieści się w $B(z,5a_k/2)$.
Porównując ich objętości, otrzymujemy oszacowanie
\begin{equation}\label{eq:besicovitch-single-scale-24}
 \#\{x\in S_k:z\in B(x,r_x)\}\le5^d.
\end{equation}

\emph{Krok 3: liczba skal występujących w punkcie.}
Spośród każdej skali $k$, dla której jakaś wybrana
kula zawiera $z$, wybierzmy jedną taką kulę
$B(x_k,r_{x_k})$. Gdy $k<\ell$, środek
$x_\ell$ nie należy do żadnej kuli wybranej w
skali $k$ (tak zdefiniowano $E_\ell$), więc
$|x_k-x_\ell|\ge r_{x_k}$. Ponadto
$r_{x_\ell}\le2a_\ell=2^{1-(\ell-k)}a_k\le a_k<r_{x_k}$,
a zatem
$x_k\notin B(x_\ell,r_{x_\ell})$.
Tak wybrane kule wszystkich występujących skal
mają wspólny punkt $z$ i spełniają oba warunki
wyłączenia środków z
Lematu~\ref{lem:besicovitch-angle-24}.
Dowolna \emph{skończona} liczba tych skal jest
więc ograniczona przez $3^d$. Gdyby skal było
więcej niż $3^d$, otrzymalibyśmy sprzeczność
już dla pewnej skończonej podrodziny.

Każda skala wnosi najwyżej $5^d$ kul
zawierających $z$, a skal jest najwyżej $3^d$.
Daje to drugą nierówność w
\eqref{eq:besicovitch-infinite-24} i kończy dowód.
\end{proof}

\begin{uwaga}[Dlaczego metoda działa dla rodzin nieskończonych?]
Kluczowe jest oddzielenie dwóch rodzajów nakładania.
W jednej skali wystarcza pakowanie oddalonych
środków, natomiast pomiędzy skalami obowiązuje
silniejsza reguła: środek późniejszej kuli nie leży
w żadnej kuli wcześniej wybranej. Dopiero
połączenie obu argumentów daje jednostajną
krotność $15^d$. Nie przechodzimy tu od pokryć
skończonych do granicy i nie powiększamy żadnej kuli.
Stała nie jest optymalna; istotna jest jej
zależność wyłącznie od wymiaru.
\end{uwaga}

\begin{wniosek}[Besicovitch dla dodatnich miar borelowskich]
\label{cor:besicovitch-measure-infinite-24}
Przy założeniach Twierdzenia~\ref{thm:besicovitch-24},
dla dowolnej dodatniej miary borelowskiej $\rho$
na $\mathbb R^d$ wybrana rodzina kul spełnia
\[
 \sum_{j\in J}\rho(B_j)
 \le15^d\,\rho\Bigl(\bigcup_{j\in J}B_j\Bigr).
\]
\end{wniosek}
\begin{proof}
Stosujemy twierdzenie Beppo Leviego do rosnących
skończonych sum funkcji $\chi_{B_j}$ (bez założenia
$\sigma$-skończoności $\rho$). Otrzymujemy
\[
 \sum_{j\in J}\rho(B_j)
 =\int_{\mathbb R^d}\sum_{j\in J}\chi_{B_j}(z)\,\mathrm d\rho(z)
 \le15^d\,\rho\Bigl(\bigcup_{j\in J}B_j\Bigr).
\]
\end{proof}

\section{Zastosowania do oszacowań typu maksymalnego}
\begin{definicja}[Funkcja maksymalna Hardy'ego--Littlewooda]
\label{def:HL-max-24}
Dla $f\in L^1_{\mathrm{loc}}(\mathbb R^d)$ definiujemy
\emph{centrowaną funkcję maksymalną}
\[
 Mf(x):=\sup_{r>0}\frac1{\lambda_d(B(x,r))}
 \int_{B(x,r)}|f(y)|\,\mathrm d\lambda_d(y)
 \in[0,+\infty].
\]
Definicja nie zależy od reprezentanta klasy $f$ w
$L^1_{\mathrm{loc}}$, ponieważ całki funkcji równych
prawie wszędzie są równe na każdej kuli.
Funkcja $Mf$ jest mierzalna jako supremum przeliczalnej
rodziny ciągłych średnich dla dodatnich promieni wymiernych. Dokładniej, gdy $(x_n,r_n)\to(x,r)$ z $r>0$,
wszystkie kule $B(x_n,r_n)$ od pewnego indeksu
mieszczą się w jednej ograniczonej kuli $D$.
Dla $y\notin\partial B(x,r)$ zachodzi
$\chi_{B(x_n,r_n)}(y)\to\chi_{B(x,r)}(y)$.
Sfera ma miarę zero, a funkcja $|f|\chi_D$
jest całkowalna. Twierdzenie Lebesgue'a o zbieżności
zdominowanej daje ciągłość licznika; mianownik
$v_dr_n^d\to v_dr^d>0$. W szczególności można
ograniczyć supremum do dodatnich promieni wymiernych.
\end{definicja}

\begin{twierdzenie}[Nierówność słabego typu $(1,1)$]
\label{thm:HL-weak-24}
Jeśli $f\in L^1(\mathbb R^d)$, to dla każdego $t>0$
\[
 \lambda_d\bigl(\{x\in\mathbb R^d:Mf(x)>t\}\bigr)
 \le\frac{5^d}{t}\int_{\mathbb R^d}|f(x)|\,
 \mathrm d\lambda_d(x).
\]
\end{twierdzenie}
\begin{proof}
Jeśli $\{Mf>t\}=\varnothing$, teza jest oczywista.
W przeciwnym razie ustalmy niepusty zwarty zbiór
$K\subseteq\{Mf>t\}$ (dla pustego $K$ żądane
oszacowanie jest natychmiastowe).
Dla każdego $x\in K$ wybieramy promień $r_x>0$ taki,
że $\int_{B(x,r_x)}|f|\,\mathrm d\lambda_d>
 t\lambda_d(B(x,r_x))$. Promienie takich kul są
ograniczone z góry: ponieważ ich objętość jest mniejsza
od $\|f\|_1/t$, zachodzi $r_x^d<\|f\|_1/(tv_d)$.
Z lematu $5r$ wybieramy rozłączne kule $(B_j)$, których
pięciokrotne powiększenia pokrywają $K$. Wówczas
\[
 \lambda_d(K)\le5^d\sum_j\lambda_d(B_j)
 <\frac{5^d}{t}\sum_j\int_{B_j}|f|\,\mathrm d\lambda_d
 \le\frac{5^d}{t}\|f\|_1.
\]
Zbiór $\{Mf>t\}$ jest otwarty: dla dowolnego $x$
warunek średniej większej od $t$ zachowuje się przy małej
zmianie środka kuli dzięki ciągłości całki z $|f|$.
Jego miara jest supremum miar zwartych podzbiorów,
na mocy regularności miary Lebesgue'a. Biorąc supremum
po $K$, kończymy dowód.
\end{proof}

\begin{wniosek}[Kontrola zbioru, na którym średnie są duże]
\label{cor:HL-level-24}
Dla $f\in L^1(\mathbb R^d)$ oraz $a>0$ funkcja $Mf$
jest skończona prawie wszędzie, a
$\lambda_d(\{Mf>a\})\le5^d\|f\|_1/a$.
\end{wniosek}
\begin{proof}
Oszacowanie wynika z twierdzenia~\ref{thm:HL-weak-24}.
Przechodząc do granicy $a\to\infty$, dostajemy
$\lambda_d(\{Mf=+\infty\})=0$.
\end{proof}

\begin{uwaga}[Co zostawiamy do następnego rozdziału?]
W rozdziale~\ref{ch:covering-24} skonstruowaliśmy narzędzie
pokryciowe i udowodniliśmy nierówność słabego typu.
W kolejnym rozdziale użyjemy ich do dowodu twierdzenia
Lebesgue'a o różniczkowaniu całek i miar, bez zakładania
go w argumentach niniejszego rozdziału.
\end{uwaga}

% ROZDZIAŁ 25 — RÓŻNICZKOWANIE CAŁEK I MIAR
\chapter{Różniczkowanie całek i miar}
\label{ch:differentiation-25}

\providecommand{\atlasref}[2]{#2}


\section{Dlaczego różniczkowanie jest twierdzeniem o średnich?}
Niech $d\in\mathbb N$, $d\ge1$, a $\lambda_d$ oznacza miarę Lebesgue'a na
$\mathbb R^d$. Dla $x\in\mathbb R^d$ i $r>0$ piszemy
$B(x,r)=\{y:|y-x|<r\}$ oraz $v_d=\lambda_d(B(0,1))$.
Jeżeli $f\in L^1_{\mathrm{loc}}(\mathbb R^d)$, to na pierwszy
rzut oka naturalne jest oczekiwanie, że średnia $f$ po kuli
malejącej do punktu $x$ dąży do $f(x)$. Nie zachodzi to
we wszystkich punktach ani dla dowolnego przedstawiciela
funkcji. Poniższa teoria precyzyjnie określa, w jakim
sensie odtwarzamy funkcję z jej całek.

\begin{definicja}[Wartość średnia i punkt Lebesgue'a]
\label{def:lebesgue-point-25}
Dla $f\in L^1_{\mathrm{loc}}(\mathbb R^d)$ definiujemy
$A_rf(x)=\lambda_d(B(x,r))^{-1}\int_{B(x,r)}f(y)\,\mathrm d\lambda_d(y)$.
Punkt $x$ nazywamy \emph{punktem Lebesgue'a} wybranego
przedstawiciela $f$, jeżeli
\begin{equation}\label{eq:lebesgue-point-25}
 \lim_{r\downarrow0}\frac1{v_dr^d}
 \int_{B(x,r)}|f(y)-f(x)|\,\mathrm d\lambda_d(y)=0.
\end{equation}
W szczególności $A_rf(x)\to f(x)$ w takim punkcie.
Przypominamy, że wartości $f(x)$ na zbiorze miary zero
można zmieniać bez zmiany klasy elementu $L^1_{\mathrm{loc}}$.
\end{definicja}

\begin{przyklad}[Zależność punktów Lebesgue'a od przedstawiciela]
Jeżeli $f=\chi_{\mathbb Q^d}$, to $A_rf(x)=0$ dla
wszystkich $x,r$, ponieważ $\lambda_d(\mathbb Q^d)=0$.
Żaden punkt należący do $\mathbb Q^d$ nie jest punktem Lebesgue'a
tego konkretnego przedstawiciela, natomiast każdy punkt spoza
$\mathbb Q^d$ nim jest.
Równość prawie wszędzie nie może więc zostać zastąpiona
równością w każdym punkcie.
\end{przyklad}

\section{Wzmocnienie nierówności maksymalnej}
\begin{definicja}[Maksymalna funkcja miary]
\label{def:max-measure-25}
Jeżeli $\nu$ jest skończoną dodatnią miarą borelowską
na $\mathbb R^d$, określamy
$M\nu(x)=\sup_{r>0}\nu(B(x,r))/(v_dr^d)$.
Dla miary ze znakiem bądź zespolonej podstawiamy jej
wariację całkowitą $|\nu|$.
\end{definicja}

\begin{twierdzenie}[Oszacowanie słabego typu dla miar]
\label{thm:weak-measure-25}
Niech $\nu$ będzie skończoną dodatnią miarą borelowską
na $\mathbb R^d$. Dla każdego $t>0$ zachodzi
\begin{equation}\label{eq:weak-measure-25}
 \lambda_d\bigl(\{x:M\nu(x)>t\}\bigr)
 \le \frac{5^d}{t}\nu(\mathbb R^d).
\end{equation}
To samo oszacowanie obowiązuje z $\nu$ zastąpioną przez
$|\nu|$, gdy $\nu$ jest miarą ze znakiem lub zespoloną.
\end{twierdzenie}
\begin{proof}
Dla ustalonego $r>0$ funkcja $x\mapsto\nu(B(x,r))$
jest półciągła z dołu: jeśli $x_n\to x$, to dla każdego
$y\in B(x,r)$ ostatecznie $y\in B(x_n,r)$, zatem lemat
Fatou zastosowany do miary $\nu$ daje
$\nu(B(x,r))\le\liminf_n\nu(B(x_n,r))$. Stąd
$\{M\nu>t\}$ jest otwarty. Jeżeli zbiór ten jest pusty, teza zachodzi. W przeciwnym
razie niech $K$ będzie jego niepustym zwartym podzbiorem
(dla pustego $K$ oszacowanie jest automatyczne). Dla każdego $x\in K$ istnieje
kula $B_x=B(x,r_x)$ taka, że
$\nu(B_x)>tv_dr_x^d$. Ponieważ $\nu$ jest skończona,
promienie są ograniczone z góry. Z lematu $5r$
(twierdzenie~\ref{thm:5r-24}) wybieramy rozłączne
kule $B_j$ o pokryciu $K\subseteq\bigcup_j5B_j$.
Wtedy
\[
\lambda_d(K)\le5^d\sum_jv_dr_j^d
 <\frac{5^d}{t}\sum_j\nu(B_j)
 \le\frac{5^d}{t}\nu(\mathbb R^d).
\]
Korzystając z regularności miary Lebesgue'a na zbiorze
otwartym $\{M\nu>t\}$, przechodzimy do supremum po
jego zwartych podzbiorach. W przypadku miary zespolonej
lub ze znakiem stosujemy otrzymany wynik do $|\nu|$.
\end{proof}

\begin{uwaga}[Nie mylić dwóch operatorów maksymalnych]
Dla $f\in L^1(\mathbb R^d)$ miara
$\nu_f(A)=\int_A|f(y)|\,\mathrm d\lambda_d(y)$ spełnia
$M\nu_f=Mf$, gdzie $Mf$ zdefiniowano w
\ref{def:HL-max-24}. Operator $M\nu$ pozwala jednak
również badać miary osobliwe, nieposiadające gęstości.
\end{uwaga}

\section{Twierdzenie Lebesgue'a o różniczkowaniu całek}

\noindent\emph{Powiązanie tematyczne w Bibliotece matematycznej:} \atlasref{FA-15B::thm:bochner-lebesgue-differentiation15b}{FA-15B: różniczkowanie całki Bochnera (uogólnienie wektorowe)}.

\begin{lemat}[Przybliżanie lokalne funkcją ciągłą]
\label{lem:continuous-dense-25}
Dla dowolnego $f\in L^1(\mathbb R^d)$ i
$\varepsilon>0$ istnieje $\varphi\in C_c(\mathbb R^d)$
taka, że $\|f-\varphi\|_1<\varepsilon$.
\end{lemat}
\begin{proof}
Jest to szczególny przypadek twierdzenia o gęstości
$C_c^\infty(\mathbb R^d)$ w $L^1$, dowiedzionego w
rozdziale~\ref{ch:regularnosc-funkcji}. W szczególności nie
zakładamy tutaj różniczkowania całek przy konstrukcji
aproksymacji.
\end{proof}

\begin{twierdzenie}[Lebesgue'a o różniczkowaniu]
\label{thm:lebesgue-diff-25}
Jeżeli $f\in L^1_{\mathrm{loc}}(\mathbb R^d)$, to
prawie każdy $x\in\mathbb R^d$ jest punktem Lebesgue'a
$f$, tzn. spełnia \eqref{eq:lebesgue-point-25}.
W szczególności
$\lim_{r\downarrow0}A_rf(x)=f(x)$ dla prawie każdego $x$.
\end{twierdzenie}
\begin{proof}
Najpierw załóżmy $f\in L^1(\mathbb R^d)$ i wybierzmy
wszędzie skończony mierzalny reprezentant $f$.
Zmiana wartości na zbiorze zerowym nie wpływa na tezę
prawie wszędzie. Ustalmy
$\varphi\in C_c(\mathbb R^d)$ i połóżmy $g=f-\varphi$.
Dla każdego $x$, dla którego $f(x)$ jest określona,
z nierówności trójkąta otrzymujemy
\begin{align*}
 \limsup_{r\downarrow0}\frac1{v_dr^d}
 \int_{B(x,r)}|f(y)-f(x)|\,\mathrm d\lambda_d(y)
 &\le Mg(x)+|g(x)|,
\end{align*}
gdzie $Mg=M(|g|\lambda_d)$. Składnik z funkcją
$\varphi$ znika wskutek jej ciągłości w $x$.
Niech $D_f(x)$ oznacza powyższy $\limsup$, a $t>0$.
Dla każdego $\varphi$ zachodzi inkluzja
$\{D_f>2t\}\subseteq\{Mg>t\}\cup\{|g|>t\}$.
Dla ustalonego $r>0$ funkcja
\[
 (x,c)\longmapsto\int_{B(x,r)}|f(y)-c|\,\mathrm d\lambda_d(y)
\]
jest ciągła: ciągłość w $x$ przy ustalonym $c$ wynika
ze zbieżności zdominowanej (brzeg kuli ma miarę zero),
a zmiana $c$ na $c'$ zmienia całkę najwyżej o
$v_dr^d|c-c'|$. Podstawiając $c=f(x)$, uzyskujemy
mierzalność średniej oscylacji. Dla każdego $x$
jest ona ciągła względem $r>0$ na mocy tego samego
argumentu ze zbieżnością zdominowaną.
Zatem $D_f$ jest mierzalna, gdyż granicę górną można
zapisać jako infimum po $n\ge1$ z supremów po
wymiernych $r\in(0,1/n)$.
Nierówność słabego typu i nierówność Czebyszewa dają
\[
 \lambda_d(\{D_f>2t\})
 \le\frac{5^d+1}{t}\|f-\varphi\|_1.
\]
Lemat~\ref{lem:continuous-dense-25} pozwala uczynić
prawą stronę dowolnie małą. Zatem
$\lambda_d(\{D_f>2t\})=0$ dla każdego $t>0$.
Ponieważ $D_f\ge0$, otrzymujemy $D_f=0$ prawie
wszędzie.

Jeżeli $f$ jest tylko lokalnie całkowalna, dla każdego
$n\ge1$ rozważamy $f_n=f\chi_{B(0,n+2)}\in L^1$.
Dla $x\in B(0,n)$ i $r<1$ średnie funkcji
$f_n$ i $f$ są jednakowe, a wartości obu funkcji
w $x$ są równe. Zastosowanie pierwszej części do
wszystkich $n$ i przeliczalna suma zbiorów zerowych
dają tezę.
\end{proof}

\begin{wniosek}[Wartość funkcji z otaczających ją średnich]
\label{cor:diff-basis-25}
Jeżeli $f\in L^1_{\mathrm{loc}}(\mathbb R^d)$,
to istnieje zbiór $N$ o $\lambda_d(N)=0$, taki że
$\lim_{r\downarrow0}A_rf(x)=f(x)$ dla wszystkich
$x\in\mathbb R^d\setminus N$. Zbioru wyjątkowego
nie można w ogólności usunąć.
\end{wniosek}
\begin{proof}
Pierwsza część wynika z twierdzenia
\ref{thm:lebesgue-diff-25}; druga z przykładu
$f=\chi_{\mathbb Q^d}$.
\end{proof}

\begin{wniosek}[Twierdzenie Lebesgue'a o gęstości]
\label{cor:lebesgue-density-25}
Niech $E\subseteq\mathbb R^d$ będzie zbiorem mierzalnym
w sensie Lebesgue'a. Dla $\lambda_d$-prawie każdego $x$ zachodzi
\[
 \lim_{r\downarrow0}\frac{\lambda_d(E\cap B(x,r))}{v_dr^d}
 =\chi_E(x).
\]
W szczególności prawie każdy punkt $E$ jest punktem gęstości
$1$, a prawie każdy punkt dopełnienia ma gęstość $0$.
\end{wniosek}
\begin{proof}
Funkcja $\chi_E$ jest lokalnie całkowalna, nawet gdy
$\lambda_d(E)=+\infty$. Stosujemy
twierdzenie~\ref{thm:lebesgue-diff-25} i zauważamy, że
$A_r\chi_E(x)=\lambda_d(E\cap B(x,r))/(v_dr^d)$.
\end{proof}

\begin{przyklad}[Przedstawiciel kanoniczny w punktach Lebesgue'a]
Jeżeli $f\in L^1_{\mathrm{loc}}(\mathbb R^d)$, to
na zbiorze punktów, w których istnieje skończona granica
$\lim_{r\downarrow0}A_rf(x)$, można ją przyjąć za
wartość przedstawiciela $\widetilde f(x)$. W pozostałych
punktach przyjmujemy $0$. Otrzymujemy funkcję mierzalną
$\widetilde f=f$ prawie wszędzie. Zmiana przedstawiciela
$f$ na zbiorze zerowym nie zmienia wartości żadnej średniej,
więc $\widetilde f$ nie zależy od tego wyboru.
Istnienie granicy samych średnich nie wystarcza jednak
do stwierdzenia, że dany punkt jest punktem Lebesgue'a.
Na przykład dla $f:\mathbb R\to\mathbb R$ określonej przez
$f(y)=1$ dla $y>0$, $f(y)=-1$ dla $y<0$ i $f(0)=0$
mamy $A_rf(0)=0$ dla każdego $r>0$, ale
$A_r|f-f(0)|(0)=1$. Nawet kanoniczna wartość w zerze
nie usuwa więc oscylacji.
\end{przyklad}

\section{Różniczkowanie miar: część absolutnie ciągła i osobliwa}
\begin{lemat}[Zanik gęstości miary osobliwej]
\label{lem:singular-density-25}
Niech $\nu$ będzie skończoną dodatnią miarą borelowską
na $\mathbb R^d$ i $\nu\perp\lambda_d$. Wówczas
\begin{equation}\label{eq:singular-density-25}
 \lim_{r\downarrow0}\frac{\nu(B(x,r))}{v_dr^d}=0
 \quad\text{dla }\lambda_d\text{-prawie każdego }x.
\end{equation}
\end{lemat}
\begin{proof}
Wybierzmy borelowski $N$ taki, że $\lambda_d(N)=0$
i $\nu(\mathbb R^d\setminus N)=0$. Regularność
skończonej miary borelowskiej daje dla każdego
$\varepsilon>0$ zwarty $K\subseteq N$ taki, że
$\nu(\mathbb R^d\setminus K)<\varepsilon$.
Oznaczmy $\nu_0=\nu|_K$, $\nu_1=\nu|_{\mathbb R^d\setminus K}$.
Dla każdego $x\notin K$ odległość $\operatorname{dist}(x,K)$
jest dodatnia, zatem $\nu_0(B(x,r))=0$ dla dostatecznie
małych $r$. W konsekwencji, poza $K$,
\[
 \limsup_{r\downarrow0}\frac{\nu(B(x,r))}{v_dr^d}
 \le M\nu_1(x).
\]
Dla ustalonego $r>0$ iloraz jest półciągły z dołu
względem $x$. Ponadto dla $r_n\uparrow r$, $r_n<r$,
z ciągłości miary od dołu ilorazy zbiegają do ilorazu
dla $r$. Możemy więc w granicy górnej ograniczyć się
do dodatnich promieni wymiernych, co dowodzi jej
borelowskiej mierzalności.
Niech $t>0$. Z twierdzenia~\ref{thm:weak-measure-25}
i z $\lambda_d(K)=0$ otrzymujemy
\[
 \lambda_d\left(\left\{x:\limsup_{r\downarrow0}
 \frac{\nu(B(x,r))}{v_dr^d}>t\right\}\right)
 \le\frac{5^d\varepsilon}{t}.
\]
Ponieważ $\varepsilon$ jest dowolne, zbiór ten ma
miarę zero. Dla $t=1/n$, $n\ge1$, dostajemy
zerowanie granicy górnej prawie wszędzie. Ilorazy są
nieujemne, więc istnieje granica równa zero.
\end{proof}

\begin{twierdzenie}[Różniczkowanie miary względem Lebesgue'a]
\label{thm:measure-differentiation-25}
Niech $\nu$ będzie dodatnią miarą Radona na
$\mathbb R^d$. Jej rozkład Lebesgue'a względem
$\lambda_d$ ma postać $\nu=f\lambda_d+\nu_s$,
gdzie $f\ge0$ jest lokalnie całkowalna, a
$\nu_s\perp\lambda_d$. Wtedy
\begin{equation}\label{eq:measure-density-25}
 \lim_{r\downarrow0}\frac{\nu(B(x,r))}{v_dr^d}
 =f(x)\qquad\text{dla }\lambda_d\text{-prawie każdego }x.
\end{equation}
Dla skończonej miary ze znakiem lub miary zespolonej
$\nu$ (w sensie rozdziału~\ref{ch:miary-ze-znakiem})
ta sama teza zachodzi z
$f=\mathrm d\nu_{\mathrm{ac}}/\mathrm d\lambda_d$,
gdzie $\nu_{\mathrm{ac}}$ oznacza część absolutnie ciągłą
miary $\nu$. Wynik stosuje się też lokalnie do różnic
dodatnich miar Radona (oraz zespolonych kombinacji takich
różnic), interpretując ich wartości tylko na ograniczonych
zbiorach borelowskich.
\end{twierdzenie}
\begin{proof}
Z rozkładu Lebesgue'a i twierdzenia Radona--Nikodyma
(rozdziały~\ref{ch:ac-singular-17} i
\ref{ch:radon-nikodym-18}) otrzymujemy
$\nu(B)=\int_Bf(y)\,\mathrm d\lambda_d(y)+\nu_s(B)$.
Pierwszy składnik po podzieleniu przez $v_dr^d$
zmierza do $f(x)$ prawie wszędzie z twierdzenia
\ref{thm:lebesgue-diff-25}. Dla składnika osobliwego
lokalizujemy miarę na kulach $B(0,n)$ i stosujemy
lemat~\ref{lem:singular-density-25}; przy
$x\in B(0,n-1)$ obie miary mają takie same wartości
na dostatecznie małych kulach. Przeliczalna suma
zbiorów wyjątkowych jest zerowa. W przypadku skończonej miary ze znakiem lub zespolonej
stosujemy ten sam argument do części absolutnie ciągłej.
Dla części osobliwej korzystamy z nierówności
$|\nu_s(B(x,r))|\le|\nu_s|(B(x,r))$ i z lematu
\ref{lem:singular-density-25}, ponieważ
$|\nu_s|\perp\lambda_d$. W wersji lokalnej ograniczamy
każdą z dodatnich miar składowych do $B(0,n+2)$.
Na $B(0,n)$ i dla $r<1$ ilorazy się nie zmieniają;
liniowość granicy oraz jednoznaczność lokalnych gęstości
na przecięciach dają tezę. Nie odejmujemy nieskończonych
mas na całej przestrzeni.
\end{proof}

\begin{uwaga}[Ważne rozróżnienie]
Istnienie granicy w \eqref{eq:measure-density-25} dla
$\lambda_d$-prawie wszystkich $x$ \emph{nie} oznacza,
że $\nu_s=0$. Dla miary Cantora pochodna gęstości
względem $\lambda_1$ jest zerowa prawie wszędzie,
chociaż sama miara jest niezerowa.
\end{uwaga}

\begin{przyklad}[Miara Diraca]
Niech $\nu=\delta_0$ na $\mathbb R^d$. Jeżeli $x\ne0$,
to dla wszystkich dostatecznie małych $r$ zachodzi
$\delta_0(B(x,r))=0$. Natomiast
$\delta_0(B(0,r))/(v_dr^d)=1/(v_dr^d)\to+\infty$.
Miara osobliwa może mieć więc nieskończoną gęstość
w swoich punktach koncentracji.
\end{przyklad}

\section{Pochodna całki w jednym wymiarze}
\begin{twierdzenie}[Podstawowe twierdzenie rachunku całkowego dla całki Lebesgue'a]
\label{thm:FTC-lebesgue-25}
Niech $a,b\in\mathbb R$, $a<b$, i niech
$f\in L^1((a,b),\lambda_1)$ będzie funkcją rzeczywistą. Funkcja
$F:[a,b]\to\mathbb R$ określona wzorem
$F(x)=\int_{(a,x)}f(t)\,\mathrm d\lambda_1(t)$
jest absolutnie ciągła i spełnia
$F'(x)=f(x)$ dla prawie każdego $x\in(a,b)$.
\end{twierdzenie}
\begin{proof}
Dla skończonej rodziny parami rozłącznych przedziałów
$(a_j,b_j)\subseteq[a,b]$ mamy
$\sum_j|F(b_j)-F(a_j)|
\le\int_{\cup_j(a_j,b_j)}|f|\,\mathrm d\lambda_1$.
Bezwzględna ciągłość całki zapewnia, że dla danego
$\varepsilon>0$ prawa strona jest mniejsza od
$\varepsilon$, gdy suma długości przedziałów jest
dostatecznie mała. Jest to definicja absolutnej
ciągłości funkcji $F$.

Niech $x\in(a,b)$ będzie punktem Lebesgue'a $f$.
Przedłużamy $f$ przez zero na $\mathbb R\setminus(a,b)$.
Dla $0<h<\min\{x-a,b-x\}$ zachodzi
\[
 \left|\frac{F(x+h)-F(x)}h-f(x)\right|
 \le\frac1h\int_{(x,x+h)}|f(t)-f(x)|\,
 \mathrm d\lambda_1(t)
 \le\frac2{2h}\int_{(x-h,x+h)}|f(t)-f(x)|\,
 \mathrm d\lambda_1(t).
\]
Ostatnia średnia zmierza do zera.
Dla $h<0$ kładziemy $k=-h$. Wówczas
$(F(x+h)-F(x))/h=k^{-1}\int_{(x-k,x)}f(t)\,\mathrm d\lambda_1(t)$.
Błąd jest ograniczony przez dwa razy średnią oscylację
na $(x-k,x+k)$, która również dąży do zera.
Prawie wszystkie punkty są
punktami Lebesgue'a na mocy twierdzenia
\ref{thm:lebesgue-diff-25}, zatem $F'=f$ prawie
wszędzie.
\end{proof}

\begin{przyklad}[Dlaczego sama pochodna prawie wszędzie nie wystarcza?]
Klasyczna funkcja Cantora $C:[0,1]\to[0,1]$ jest
ciągła i niemalejąca, a $C'=0$ prawie wszędzie, lecz
$C(1)-C(0)=1$. Nie jest absolutnie ciągła.
Dlatego z samej równości $F'=f$ prawie wszędzie nie
wynika $F(b)-F(a)=\int_{(a,b)}f\,\mathrm d\lambda_1$.
\end{przyklad}

\begin{twierdzenie}[Funkcja rozkładu a absolutna ciągłość miary]
\label{thm:cdf-derivative-25}
Niech $\nu$ będzie skończoną dodatnią miarą borelowską
na $\mathbb R$. Określmy $F:\mathbb R\to[0,\nu(\mathbb R)]$
wzorem $F(x)=\nu(( -\infty,x])$.
Jeśli $\nu=f\lambda_1+\nu_s$ jest rozkładem Lebesgue'a
miary $\nu$ względem $\lambda_1$, to
$F'(x)=f(x)$ dla $\lambda_1$-prawie każdego $x$.
Ponadto następujące warunki są równoważne:
\begin{enumerate}[label=(\roman*)]
\item $\nu\ll\lambda_1$;
\item $F$ jest absolutnie ciągła na każdym zwartym
przedziale;
\item dla wszystkich $a<b$ zachodzi
$F(b)-F(a)=\int_{(a,b)}F'(t)\,\mathrm d\lambda_1(t)$.
\end{enumerate}
W (iii) pochodna istnieje prawie wszędzie.
\end{twierdzenie}
\begin{proof}
Niech $x$ będzie punktem Lebesgue'a $f$, w którym
$\nu_s(B(x,r))/(2r)\to0$. Punkty takie tworzą zbiór
o dopełnieniu miary Lebesgue'a zero na mocy twierdzenia
\ref{thm:lebesgue-diff-25} i lematu~\ref{lem:singular-density-25}.
Nie trzeba osobno wykluczać atomów: w takim punkcie
$\nu_s(\{x\})=0$ wynika już z zaniku ilorazu.
Dla $h>0$ mamy
\[
 \frac{F(x+h)-F(x)}h
 =\frac1h\int_{(x,x+h]}f(t)\,\mathrm d\lambda_1(t)
 +\frac{\nu_s((x,x+h])}h.
\]
Pierwszy składnik dąży do $f(x)$, ponieważ średnia
wartości $|f(t)-f(x)|$ po $(x,x+h)$ jest ograniczona
przez dwa razy analogiczną średnią po $(x-h,x+h)$.
Drugi składnik, dla $h>0$, jest ograniczony przez
$\nu_s(B(x,2h))/h$ i zatem dąży do zera.
Dla $h<0$ rozpatrujemy przedział $(x+h,x]$;
jest on zawarty w $B(x,2|h|)$, więc składnik osobliwy
podzielony przez $|h|$ także dąży do zera.
Składnik absolutnie ciągły szacujemy przez średnią
po $B(x,|h|)$, jak dla $h>0$. Dowodzi to $F'=f$ prawie wszędzie.

Jeśli (i) zachodzi, to $\nu=f\lambda_1$ z
$f\in L^1(\mathbb R)$, a więc
$F(x)=\int_{(-\infty,x]}f(t)\,\mathrm d\lambda_1(t)$.
Bezwzględna ciągłość całki na zbiorach o małej mierze
daje (ii), zaś podstawowe twierdzenie rachunku
całkowego (twierdzenie~\ref{thm:FTC-lebesgue-25})
daje (iii).

Załóżmy (ii) i wykażmy najpierw (i), bez powoływania
się na twierdzenie o całkowaniu pochodnej dowolnej funkcji
absolutnie ciągłej. Ustalmy $n\ge1$ i zbiór borelowski
$E\subseteq[-n,n]$ z $\lambda_1(E)=0$. Funkcja $F$
jest absolutnie ciągła na $[-n-1,n+1]$, więc dla
każdego $\varepsilon>0$ istnieje $\delta>0$ takie,
że dla każdej skończonej rodziny rozłącznych przedziałów
zawartych w $[-n-1,n+1]$, o łącznej długości poniżej
$\delta$, suma przyrostów $F$ jest mniejsza od
$\varepsilon$. Pokryjmy $E$ otwartym zbiorem
$O\subset(-n-1,n+1)$ o $\lambda_1(O)<\delta$.
Składowymi $O$ są parami rozłączne przedziały otwarte
$(a_j,b_j)$. Ponieważ $F$ jest ciągła,
$\nu(\{x\})=F(x)-F(x^-)=0$ dla każdego $x$
wewnątrz $[-n-1,n+1]$, a więc
$\nu((a_j,b_j))=F(b_j)-F(a_j)$.
Stosując nierówność absolutnej ciągłości do skończonych
sum i przechodząc do granicy z przeliczalnej addytywności,
dostajemy $\nu(E)\le\nu(O)\le\varepsilon$.
Zatem $\nu(E)=0$. Ponieważ $n$ był dowolny,
$\nu\ll\lambda_1$ na całej prostej. Wykazane już
(i)$\Rightarrow$(iii) daje również (ii)$\Rightarrow$(iii).
Załóżmy wreszcie (iii). Ponieważ
$\nu((a,b])=F(b)-F(a)$, dla każdych $a<b$ mamy
$\nu((a,b])=\int_{(a,b)}F'(t)\,\mathrm d\lambda_1(t)$;
przy całce nie ma znaczenia wybór końców.
Z udowodnionej równości $F'=f$ prawie wszędzie
wynika, że pochodna ma nieujemną mierzalną wersję,
daną przez gęstość $f$ części bezwzględnie ciągłej.
Jest ona całkowalna, gdyż
$\int_{\mathbb R}f\,\mathrm d\lambda_1\le\nu(\mathbb R)<\infty$.
Obie miary są skończone na zbiorach
ograniczonych i zgadzają się na przedziałach $(a,b]$.
Jednoznaczność miary na generującym $\pi$-układzie
i przejście przez zbiory $(-n,n]$ pokazują
$\nu(E)=\int_EF'(t)\,\mathrm d\lambda_1(t)$
dla wszystkich $E\in\mathcal B(\mathbb R)$.
To daje (i).
\end{proof}

\begin{przyklad}[Pochodna nie wykrywa składnika osobliwego]
W przypadku miary Cantora funkcja rozkładu $F$ jest
ciągła, niemalejąca i spełnia $F'=0$ prawie wszędzie,
ale nie jest absolutnie ciągła. Twierdzenie
\ref{thm:cdf-derivative-25} wyjaśnia przyczynę:
pochodna odzyskuje wyłącznie część absolutnie ciągłą
miary, a nie składnik osobliwy.
\end{przyklad}

\section{Przenoszenie różniczkowania na inne rodziny zbiorów}
\begin{wniosek}[Średnie po kostkach]
\label{cor:cubes-differentiation-25}
Niech $f\in L^1_{\mathrm{loc}}(\mathbb R^d)$. Dla prawie
każdego $x$ średnia $f$ po dowolnych kostkach osiowych
$Q_r(x)=x+(-r,r)^d$ dąży do $f(x)$ przy $r\downarrow0$.
Ponadto $\lambda_d(Q_r(x))^{-1}
\int_{Q_r(x)}|f(y)-f(x)|\,\mathrm d\lambda_d(y)\to0$.
\end{wniosek}
\begin{proof}
Mamy $Q_r(x)\subseteq B(x,\sqrt d\,r)$ oraz
$\lambda_d(B(x,\sqrt d\,r))/\lambda_d(Q_r(x))
=v_dd^{d/2}/2^d$, niezależnie od $r$. Dodatnia średnia
po kostce jest więc ograniczona przez stałą razy
średnia po większej kuli, która zmierza do zera
w każdym punkcie Lebesgue'a.
\end{proof}

\section{Różniczkowanie względem dowolnej miary Radona}

\begin{lemat}[Nierówność maksymalna dla dwóch miar]
\label{lem:relative-maximal-25}
Niech $\mu$ będzie dodatnią miarą Radona na $\mathbb R^d$,
a $\eta$ skończoną dodatnią miarą borelowską na $\mathbb R^d$.
Dla $R>0$ określamy
\[
 M_{\mu,R}\eta(x)=
 \sup_{\substack{0<r\le R\\\mu(B(x,r))>0}}
 \frac{\eta(B(x,r))}{\mu(B(x,r))},
\]
przyjmując supremum pustej rodziny równym zero.
Funkcja ta jest borelowska i dla każdego $t>0$
\[
 \mu(\{M_{\mu,R}\eta>t\})\le\frac{15^d}{t}\eta(\mathbb R^d).
\]
\end{lemat}
\begin{proof}
Dla ustalonego $r$ funkcje $x\mapsto\eta(B(x,r))$
oraz $x\mapsto\mu(B(x,r))$ są półciągłe z dołu
na mocy argumentu z lematem Fatou z dowodu
twierdzenia~\ref{thm:weak-measure-25}. Mianownik jest
skończony z lokalnej skończoności $\mu$.
Jeżeli $\mu(B(x,r))>0$, dla wymiernych $r_n\uparrow r$,
$r_n<r$, ciągłość miar od dołu daje zbieżność obu mas,
a mianowniki są dodatnie od pewnego indeksu. Ilorazy
zbiegają więc do ilorazu dla $r$. W supremum wystarczą
dodatnie promienie wymierne nie większe niż $R$, a każdy
iloraz przedłużony przez zero na zbiór zerowego
mianownika jest borelowski.

Niech $E_n=B(0,n)\cap\{M_{\mu,R}\eta>t\}$.
Dla każdego $x\in E_n$ wybieramy $0<r_x\le R$ takie, że
$0<t\mu(B(x,r_x))<\eta(B(x,r_x))$.
Z twierdzenia~\ref{thm:besicovitch-24} istnieje
co najwyżej przeliczalne pokrycie $E_n$ wybranymi
kulami $(B_j)$ o krotności co najwyżej $15^d$.
Nie powiększamy kul. Z wniosku
\ref{cor:besicovitch-measure-infinite-24} wynika
\[
 \mu(E_n)\le\sum_j\mu(B_j)
 \le\frac1t\sum_j\eta(B_j)
 \le\frac{15^d}{t}\eta(\mathbb R^d).
\]
Ponieważ $E_n$ rosną do całego zbioru poziomicowego,
ciągłość $\mu$ od dołu daje tezę.
\end{proof}

\begin{twierdzenie}[Różniczkowanie całek względem miary Radona]
\label{thm:relative-lebesgue-25}
Niech $\mu$ będzie dodatnią miarą Radona na $\mathbb R^d$
i $f\in L^1_{\mathrm{loc}}(\mu)$. Dla $\mu$-prawie każdego
$x$ wszystkie kule $B(x,r)$, $r>0$, mają dodatnią miarę
i zachodzi
\[
 \lim_{r\downarrow0}
 \frac{1}{\mu(B(x,r))}
 \int_{B(x,r)}|f(y)-f(x)|\,\mathrm d\mu(y)=0.
\]
W szczególności średnie $f$ względem $\mu$ dążą do $f(x)$.
\end{twierdzenie}
\begin{proof}
Niech $S=\operatorname{supp}\mu$ oznacza zbiór punktów,
których każde otwarte otoczenie ma dodatnią miarę.
Dopełnienie $S$ jest sumą otwartych zbiorów miary zero.
Z przeliczalnej bazy $\mathbb R^d$ wynika, że wystarczy
przeliczalnie wiele takich zbiorów; zatem $\mu(S^c)=0$.
W dalszym ciągu rozpatrujemy tylko $x\in S$.

Wybieramy wszędzie skończony borelowski reprezentant $f$;
dla funkcji mierzalnej względem uzupełnienia miary taki
reprezentant otrzymujemy przez zmianę na zbiorze miary zero.
Najpierw niech $f\in L^1(\mu)$. Funkcje $C_c(\mathbb R^d)$
są gęste w $L^1(\mu)$: przybliżamy $f$ przez całkowalne
funkcje proste o zbiorach skończonej miary, a dla takiego
zbioru $A$ regularność daje zwarty $K\subseteq A$
i otwarty $U\supseteq A$ z $\mu(U\setminus K)$ dowolnie małą.
Funkcja odcinająca $h\in C_c$, $0\le h\le1$,
$h|_K=1$, $\operatorname{supp}h\subseteq U$ spełnia
$\|\chi_A-h\|_{L^1(\mu)}\le\mu(U\setminus K)$.
Dla zbiorów mierzalnych w uzupełnieniu miary wybieramy
najpierw borelowskie reprezentanty modulo zbiory zerowe.

Dla $\varphi\in C_c$ kładziemy $g=f-\varphi$.
Ciągłość $\varphi$ w $x$ i nierówność trójkąta dają
\[
 D_{\mu,f}(x):=\limsup_{r\downarrow0}
 \frac{\int_{B(x,r)}|f(y)-f(x)|\,\mathrm d\mu(y)}
 {\mu(B(x,r))}
 \le M_{\mu,1}(|g|\mu)(x)+|g(x)|.
\]
Uzasadnijmy również mierzalność $D_{\mu,f}$ na $S$.
Dla ustalonego $r>0$ i stałej $c$ funkcja
$x\mapsto\int_{B(x,r)}|f(y)-c|\,\mathrm d\mu(y)$
jest półciągła z dołu na mocy lematu Fatou.
Dla ustalonego $x$ zmiana $c$ na $c'$ zmienia tę
całkę najwyżej o $\mu(B(x,r))|c-c'|$.
Przybliżając borelowską funkcję $f(x)$ przez borelowskie
funkcje proste o co najwyżej przeliczalnym zbiorze
wartości, otrzymujemy mierzalność licznika z $c=f(x)$.
Mianownik na $S$ jest dodatni i skończony.
Dla ustalonego $x$ i $r_n\uparrow r$ oba składniki
ilorazu zbiegają z twierdzenia Beppo Leviego (dla masy
kuli jest to ciągłość miary od dołu). Granicę górną
możemy więc obliczać jako infimum po $n\ge1$ z supremów
po dodatnich wymiernych promieniach mniejszych niż $1/n$.
Z lematu~\ref{lem:relative-maximal-25}
i nierówności Czebyszewa otrzymujemy dla $t>0$
\[
 \mu(\{x\in S:D_{\mu,f}(x)>2t\})
 \le\frac{15^d+1}{t}\|f-\varphi\|_{L^1(\mu)}.
\]
Ponieważ $\varphi$ można wybrać z dowolnie małym
błędem w $L^1(\mu)$, zbiór po lewej stronie ma miarę zero.
Przeliczalna suma takich zbiorów dla $t=1/n$ zawiera
wszystkie punkty $S$ z $D_{\mu,f}(x)>0$, więc ma miarę zero.
Dla funkcji lokalnie całkowalnej stosujemy wynik do
$f\chi_{B(0,n+2)}$ i rozpatrujemy $x\in B(0,n)$,
$r<1$. Przeliczalna suma zbiorów wyjątkowych ma
miarę zewnętrzną zero.
\end{proof}

\begin{twierdzenie}[Różniczkowanie dwóch miar Radona]
\label{thm:relative-radon-25}
Niech $\mu,\nu$ będą dodatnimi miarami Radona na $\mathbb R^d$,
a $\nu=f\mu+\nu_s$ rozkładem Lebesgue'a względem $\mu$,
gdzie $\nu_s\perp\mu$. Wówczas
\[
 \lim_{r\downarrow0}\frac{\nu(B(x,r))}{\mu(B(x,r))}
 =f(x)\qquad\text{dla }\mu\text{-prawie każdego }x.
\]
Iloraz rozpatrujemy w punktach nośnika $\mu$.
\end{twierdzenie}
\begin{proof}
Miary są $\sigma$-skończone, gdyż przestrzeń jest
przeliczalną sumą kul o zwartym domknięciu.
Niech $S=\operatorname{supp}\mu$; jak w poprzednim dowodzie,
$\mu(S^c)=0$ i wszystkie kule o środku w $S$ mają dodatnią miarę.
Rozkład Lebesgue'a i twierdzenie Radona--Nikodyma
dają $f\in L^1_{\mathrm{loc}}(\mu)$, ponieważ
$\int_K f\,\mathrm d\mu\le\nu(K)<\infty$ dla zwartych $K$.
Twierdzenie~\ref{thm:relative-lebesgue-25} rozstrzyga
składnik absolutnie ciągły.

Najpierw załóżmy, że $\nu_s$ jest skończona.
Istnieje borelowski $N$ z $\mu(N)=0$ i $\nu_s(N^c)=0$.
Dla $\varepsilon>0$ wybieramy zwarty $K\subseteq N$
z $\nu_s(K^c)<\varepsilon$ i kładziemy
$\eta=\nu_s|_{K^c}$. Dla $x\in S\setminus K$ małe kule
nie przecinają $K$, więc
\[
 \limsup_{r\downarrow0}\frac{\nu_s(B(x,r))}{\mu(B(x,r))}
 \le M_{\mu,1}\eta(x).
\]
Dla $t>0$ miara zewnętrzna zbioru punktów $x\in S$,
w których granica górna przekracza $t$, nie przekracza
$\mu(K)+15^d\varepsilon/t=15^d\varepsilon/t$.
Dowolność $\varepsilon$ i przeliczalna suma dla $t=1/n$
dają zanik ilorazu $\mu$-prawie wszędzie.
W przypadku ogólnym stosujemy ten argument do
$\nu_s|_{B(0,n+2)}$, a następnie do punktów $x\in B(0,n)$
i promieni $r<1$. Otrzymujemy żądany wynik lokalnie
i w konsekwencji na całej przestrzeni.
\end{proof}

\begin{wniosek}[Różniczkowanie miar ze znakiem i zespolonych]
\label{cor:relative-signed-25}
Niech $\mu$ będzie dodatnią miarą Radona na $\mathbb R^d$,
a $\nu$ skończoną borelowską miarą ze znakiem lub miarą
zespoloną. Jeśli $\nu=f\mu+\nu_s$ jest rozkładem
Lebesgue'a względem $\mu$, gdzie $|\nu_s|\perp\mu$,
to dla $\mu$-prawie każdego $x\in\operatorname{supp}\mu$
zachodzi
\[
 \lim_{r\downarrow0}\frac{\nu(B(x,r))}{\mu(B(x,r))}
 =f(x).
\]
\end{wniosek}
\begin{proof}
Rozkład i pochodną rozumiemy w sensie rozdziałów
\ref{ch:ac-singular-17} i~\ref{ch:radon-nikodym-18}.
Część absolutnie ciągłą różniczkujemy za pomocą
twierdzenia~\ref{thm:relative-lebesgue-25}.
Dodatnia miara skończona $|\nu_s|$ jest miarą Radona
na $\mathbb R^d$ i osobliwą względem $\mu$. Z twierdzenia
\ref{thm:relative-radon-25} jej iloraz mas dąży do zera
dla $\mu$-prawie każdego $x$. Ponieważ
$|\nu_s(B(x,r))|\le|\nu_s|(B(x,r))$, to samo zachodzi
dla ilorazu części osobliwej $\nu_s$. Dodanie obu granic
kończy dowód.
\end{proof}

\begin{uwaga}[Znaczenie twierdzenia Besicovitcha]
W ogólności nie istnieje stała $C$ taka, że
$\mu(5B)\le C\mu(B)$ dla wszystkich kul, nawet o dodatniej
mierze. Na prostej weźmy skończoną miarę Radona
$\mu=\sum_{n\ge1}(2^{-2n}\delta_{10n}+2^{-n}\delta_{10n+2})$
i kule $B_n=B(10n,1)$. Wówczas
$\mu(B_n)=2^{-2n}$, natomiast
$\mu(5B_n)=2^{-2n}+2^{-n}$, więc iloraz tych mas wynosi
$1+2^n\to+\infty$.
Dlatego w dowodzie nierówności dla dwóch miar
stosujemy pokrycie bez powiększania kul, z ograniczoną
krotnością. Kule zerowej miary pomijamy; poza nośnikiem
$\mu$ nie ma potrzeby definiowania ilorazów dla twierdzenia
o różniczkowaniu.
\end{uwaga}

% ROZDZIAŁ 26 — MIARA I WYMIAR HAUSDORFFA, TWIERDZENIE FROSTMANA
\chapter{Miara Hausdorffa i wymiar Hausdorffa}
\label{ch:hausdorff-26}

\providecommand{\atlasref}[2]{#2}


\section{Dlaczego miara Lebesgue'a nie wystarcza?}
Obraz ograniczonego przedziału przez odwzorowanie
Lipschitzowskie w $\mathbb R^2$ ma zerową miarę
$\lambda_2$, choć może mieć dodatnią długość.
Dla dowolnych krzywych ciągłych nie można tego
twierdzić (istnieją krzywe wypełniające kwadrat).
Zbiór Cantora ma zerową długość, ale posiada nietrywialną
strukturę w skali ułamkowego wymiaru. Celem tego
rozdziału jest skonstruowanie miar $\mathcal H^s$,
które mierzą zbiory odpowiednio do ich wymiaru, oraz
wyprowadzenie skutecznych kryteriów tego wymiaru.

\section{Konstrukcja miary Hausdorffa od podstaw}
\label{sec:hausdorff-konstrukcja-26}

\noindent\emph{Powiązanie tematyczne w Bibliotece matematycznej:} \atlasref{DG-11::def:miara-hausdorffa}{DG-11: miara Hausdorffa w geometrii różniczkowej}.


W tej sekcji przeprowadzimy całą konstrukcję, łącznie z
mierzalnością zbiorów borelowskich. Nie będziemy zakładać
twierdzenia, że metryczna miara zewnętrzna automatycznie
mierzy zbiory borelowskie: udowodnimy potrzebny fakt.
Wykorzystamy jedynie elementarne własności kresów,
przeliczalnych pokryć i odległości w przestrzeni metrycznej.

\begin{definicja}[Koszt pokrycia i miara Hausdorffa]
\label{def:hausdorff-measure-26}
Niech $(X,d)$ będzie przestrzenią metryczną, $s\geq0$
i $\delta>0$. Dla $U\subseteq X$ definiujemy
$\operatorname{diam}U:=\sup\{d(x,y):x,y\in U\}$,
przyjmując $\operatorname{diam}\varnothing=0$.
Ustalamy $c_s>0$, przy czym $c_0=1$.
Dla $s>0$ koszt zbioru to $h_s(U)=c_s(\operatorname{diam}U)^s$.
Dla $s=0$ przyjmujemy jawnie
$h_0(U)=0$ dla $U=\varnothing$ oraz $h_0(U)=1$
dla $U\neq\varnothing$, niezależnie od średnicy.
Definiujemy
\begin{equation}\label{eq:hausdorff-delta-26}
 \mathcal H^s_\delta(E)
 :=\inf\left\{\sum_{j=1}^\infty h_s(U_j):
 E\subseteq\bigcup_{j=1}^\infty U_j,\quad
 \operatorname{diam}U_j\leq\delta\right\}.
\end{equation}
Infimum pustej rodziny dopuszczalnych pokryć wynosi
$+\infty$; zbiór pusty ma koszt $0$ (pokrycie pustymi
zbiorami). \emph{Miarę Hausdorffa wymiaru $s$} określamy
wzorem
\begin{equation}\label{eq:hausdorff-26}
 \mathcal H^s(E):=\lim_{\delta\downarrow0}
 \mathcal H^s_\delta(E)
 =\sup_{\delta>0}\mathcal H^s_\delta(E).
\end{equation}
Jest to dobrze określona wartość z $[0,+\infty]$,
gdyż $0<\delta'<\delta$ pociąga za sobą
$\mathcal H^s_{\delta'}(E)\geq\mathcal H^s_\delta(E)$.
Dla $s>0$ popularna normalizacja to
$c_s=\omega_s2^{-s}$, gdzie
$\omega_s=\pi^{s/2}/\Gamma(1+s/2)$;
$\Gamma$ oznacza funkcję gamma Eulera, tj.
$\Gamma(t)=\int_{(0,+\infty)}u^{t-1}e^{-u}\,\mathrm d\lambda_1(u)$
dla $t>0$. Można jednak przyjąć $c_s=1$, nie zmieniając
własności mierzalności ani wymiaru Hausdorffa.
W przypadku $s=1$ standardowy wybór daje $c_1=1$.
\end{definicja}

Ważne jest rozdzielenie dwóch operacji: najpierw
wybieramy \emph{dowolne przeliczalne} pokrycie zbiorami
o średnicy co najwyżej $\delta$ i minimalizujemy koszt;
dopiero potem wymuszamy $\delta\downarrow0$.
Nie ograniczamy pokryć do kul ani do zbiorów mierzalnych.
W przestrzeni nieośrodkowej może nie istnieć przeliczalne
pokrycie danego $E$ o małych średnicach; wtedy
odpowiednia wartość jest po prostu nieskończona.

\begin{uwaga}[Miara a zawartość Hausdorffa]
Gdy w~\eqref{eq:hausdorff-delta-26} nie narzucamy
ograniczenia średnic, otrzymujemy \emph{zawartość
Hausdorffa} $\mathcal H^s_\infty(E)$. Sama zawartość
nie musi być addytywna nawet dla rozłącznych
zbiorów odległych od siebie: jedno duże pokrycie
może być tańsze od dwóch oddzielnych.
Dlatego warunek $\delta\downarrow0$ jest istotny.
\end{uwaga}

\begin{lemat}[Zewnętrzne własności kosztu pokryć]
\label{lem:hausdorff-outer-26}
Dla każdego $s\geq0$ i $\delta>0$ funkcja
$E\mapsto\mathcal H^s_\delta(E)$ spełnia:
$\mathcal H^s_\delta(\varnothing)=0$,
monotoniczność i przeliczalną podaddytywność.
To samo dotyczy funkcji $E\mapsto\mathcal H^s(E)$.
\end{lemat}
\begin{proof}
Pusta rodzina punktów nie wymaga żadnego pokrycia,
a pokrycie złożone ze zbiorów pustych ma koszt zero.
Jeśli $A\subseteq B$, każde pokrycie $B$ pokrywa $A$,
więc infimum dla $A$ nie przekracza infimum dla $B$.

Niech $E\subseteq\bigcup_{k\ge1}E_k$. Jeśli suma
$\sum_k\mathcal H^s_\delta(E_k)$ jest nieskończona,
nierówność podaddytywności jest automatyczna.
W przeciwnym razie wszystkie składniki są skończone.
Dla $\varepsilon>0$ i każdego $k$ wybieramy dopuszczalne
pokrycie $(U_{kj})_{j\ge1}$ zbioru $E_k$ o koszcie
nieprzekraczającym
$\mathcal H^s_\delta(E_k)+\varepsilon2^{-k}$.
Rodzina wszystkich $U_{kj}$ jest przeliczalna,
ma średnice co najwyżej $\delta$ i pokrywa $E$.
W konsekwencji
$\mathcal H^s_\delta(E)\leq
\sum_k\mathcal H^s_\delta(E_k)+\varepsilon$.
Przechodząc $\varepsilon\downarrow0$, otrzymujemy
podaddytywność na ustalonej skali.

Przy przejściu do granicy \emph{nie zamieniamy}
supremum z nieskończoną sumą. Dla każdego $\delta$
z poprzedniej nierówności i z definicji supremum mamy
$\mathcal H^s_\delta(E)\leq
\sum_k\mathcal H^s_\delta(E_k)
\leq\sum_k\mathcal H^s(E_k)$.
Wzięcie supremum po $\delta>0$ kończy dowód.
\end{proof}

\begin{lemat}[Addytywność na zbiorach oddalonych]
\label{lem:hausdorff-metricity-26}
Jeżeli $A,B\subseteq X$ i
$\operatorname{dist}(A,B):=
\inf\{d(a,b):a\in A,\ b\in B\}>0$,
to $\mathcal H^s(A\cup B)=\mathcal H^s(A)+\mathcal H^s(B)$.
Wniosek pozostaje prawdziwy dla każdej skończonej
rodziny parami dodatnio oddalonych zbiorów.
\end{lemat}
\begin{proof}
Ustalmy $0<\delta<\operatorname{dist}(A,B)$.
Żaden zbiór o średnicy nieprzekraczającej $\delta$
nie może przecinać jednocześnie $A$ i $B$.
Każde pokrycie $A\cup B$ rozdzielamy zatem na
podrodzinę pokrywającą $A$ i podrodzinę pokrywającą
$B$; zbiory nieprzecinające żadnego z nich pomijamy.
Suma kosztów pierwotnego pokrycia jest co najmniej
$\mathcal H^s_\delta(A)+\mathcal H^s_\delta(B)$.
Biorąc infimum, otrzymujemy nierówność
$\mathcal H^s_\delta(A\cup B)\geq
\mathcal H^s_\delta(A)+\mathcal H^s_\delta(B)$.
Po przejściu do granicy $\delta\downarrow0$
i użyciu lematu~\ref{lem:hausdorff-outer-26}
uzyskujemy równość. Dla skończonej rodziny
stosujemy ten argument indukcyjnie: suma skończonej
liczby zbiorów odległych od pozostałych ma nadal
od nich dodatnią odległość.
\end{proof}

\begin{lemat}[Kryterium Carathéodory'ego, z dowodem]
\label{lem:hausdorff-caratheodory-selfcontained-26}
Niech $m^*:2^X\to[0,+\infty]$ będzie dowolną
miarą zewnętrzną, czyli $m^*(\varnothing)=0$,
$m^*$ jest monotoniczna i przeliczalnie
podaddytywna. Zbiór $A\subseteq X$ nazywamy
$m^*$-mierzalnym, jeśli dla każdego $E\subseteq X$
\begin{equation}\label{eq:hausdorff-caratheodory-test-26}
 m^*(E)=m^*(E\cap A)+m^*(E\setminus A).
\end{equation}
Rodzina $\mathcal M_{m^*}$ wszystkich takich zbiorów
jest $\sigma$-algebrą, a $m^*|_{\mathcal M_{m^*}}$
jest miarą zupełną.
\end{lemat}
\begin{proof}
Zbiór pusty i $X$ spełniają warunek, a zamiana $A$
na jego dopełnienie jedynie przestawia składniki.
Jeżeli $A,B$ spełniają~\eqref{eq:hausdorff-caratheodory-test-26},
to dwukrotnie rozcinając $E$ otrzymujemy
\[
 m^*(E)=m^*(E\cap A\cap B)
 +m^*(E\cap A\setminus B)+m^*(E\setminus A)
 \geq m^*(E\cap A\cap B)+m^*(E\setminus(A\cap B)).
\]
Odwrotną nierówność daje podaddytywność, więc
$A\cap B$ jest mierzalny. Dopełnienia dają
zamkniętość na skończone sumy i różnice.

Niech $(A_j)$ będzie ciągiem \emph{rozłącznych}
zbiorów mierzalnych, $B=\bigcup_{j\ge1}A_j$
i $B_N=\bigcup_{j=1}^NA_j$. Wielokrotne
zastosowanie kryterium daje dla każdego $E\subseteq X$
\[
 m^*(E)=\sum_{j=1}^Nm^*(E\cap A_j)
 +m^*(E\setminus B_N)
 \geq\sum_{j=1}^Nm^*(E\cap A_j)+m^*(E\setminus B).
\]
Przechodzimy $N\to\infty$, a następnie stosujemy
podaddytywność do $E\cap B=\bigcup_j(E\cap A_j)$.
Otrzymujemy nierówność $m^*(E)\geq
m^*(E\cap B)+m^*(E\setminus B)$;
odwrotna wynika z podaddytywności. Zatem $B$
jest mierzalny. Dla dowolnego ciągu zbiorów
mierzalnych rozłączność uzyskujemy, zastępując
$A_j$ przez $A_j\setminus\bigcup_{i<j}A_i$.
To dowodzi, że $\mathcal M_{m^*}$ jest $\sigma$-algebrą.

Dla rozłącznych mierzalnych $A_j$ stosujemy
poprzedni rachunek do $E=B$; po przejściu
do granicy otrzymujemy
$m^*(B)\geq\sum_jm^*(A_j)$.
Podaddytywność daje nierówność odwrotną,
więc mamy przeliczalną addytywność.
Wreszcie, gdy $m^*(N)=0$, dla dowolnego $A\subseteq N$
i $E\subseteq X$ zachodzi
$m^*(E\cap A)=0$ oraz
$m^*(E\setminus A)\leq m^*(E)$.
Podaddytywność wymusza~\eqref{eq:hausdorff-caratheodory-test-26}
dla $A$. Każdy podzbiór zbioru zerowego jest zatem
mierzalny i zerowy: miara jest zupełna.
\end{proof}

\begin{lemat}[Metryczna miara zewnętrzna a zbiory domknięte]
\label{lem:hausdorff-metric-open-26}
Załóżmy, że miara zewnętrzna $m^*$ spełnia
$m^*(A\cup B)=m^*(A)+m^*(B)$ zawsze, gdy
$\operatorname{dist}(A,B)>0$. Wtedy każdy zbiór
domknięty $F\subseteq X$ spełnia kryterium
Carathéodory'ego~\eqref{eq:hausdorff-caratheodory-test-26}.
W szczególności wszystkie zbiory otwarte są mierzalne.
\end{lemat}
\begin{proof}
Dla $F=\varnothing$ lub $F=X$ teza jest oczywista.
Załóżmy, że $F$ jest niepusty i właściwy, i ustalmy
$E\subseteq X$. Jeśli $m^*(E)=+\infty$,
podaddytywność już daje równość z kryterium,
więc możemy przyjąć $m^*(E)<\infty$.
Funkcja $\rho(x)=\operatorname{dist}(x,F)$
jest $1$-Lipschitzowska, $\rho^{-1}(0)=F$.
Dla $k\ge0$ wprowadzamy warstwy
$S_k=E\cap\{2^{-(k+1)}<\rho\leq2^{-k}\}$.
Warstwy $S_0,S_2,S_4,\ldots$ są parami dodatnio
oddalone: między zakresami wartości $\rho$ dwóch
kolejnych parzystych warstw leży niepusty przedział,
a $|\rho(x)-\rho(y)|\leq d(x,y)$.
To samo dotyczy warstw o nieparzystych indeksach.
Z addytywności na skończonych rodzinach oddalonych
i z monotoniczności wynika zatem
$\sum_{j\ge0}m^*(S_{2j})\leq m^*(E)$ oraz
$\sum_{j\ge0}m^*(S_{2j+1})\leq m^*(E)$.
Wobec tego szereg $\sum_km^*(S_k)$ jest zbieżny,
a jego ogony dążą do zera.

Połóżmy $E_n=E\cap\{\rho>2^{-n}\}$. Mamy
$E\setminus F=\bigcup_nE_n$ oraz
$(E\setminus F)\setminus E_n\subseteq\bigcup_{k\ge n}S_k$.
Z podaddytywności i monotoniczności dostajemy
\[
 m^*(E_n)\leq m^*(E\setminus F)
 \leq m^*(E_n)+\sum_{k\ge n}m^*(S_k).
\]
W szczególności $m^*(E_n)\to m^*(E\setminus F)$.
Ponieważ zbiory $E\cap F$ i $E_n$ mają odległość
co najmniej $2^{-n}$, metryczna addytywność daje
$m^*(E)\geq m^*(E\cap F)+m^*(E_n)$.
Przechodząc do granicy, otrzymujemy
$m^*(E)\geq m^*(E\cap F)+m^*(E\setminus F)$.
Podaddytywność daje nierówność odwrotną.
Zatem $F$ jest mierzalny, a jego dopełnienie
jest mierzalne na mocy poprzedniego lematu.
\end{proof}

\noindent\emph{Porównanie w Bibliotece matematycznej:} \atlasref{CAL-10::thm:hausdorff-borel-26}{CAL-10: miara Hausdorffa w teorii całki}.

\begin{twierdzenie}[Samodzielna konstrukcja miary Hausdorffa]
\label{thm:hausdorff-borel-26}
Dla dowolnej przestrzeni metrycznej $(X,d)$ i $s\geq0$
funkcja $\mathcal H^s$ z~\eqref{eq:hausdorff-26}
jest metryczną miarą zewnętrzną. Rodzina jej zbiorów
mierzalnych w sensie Carathéodory'ego jest
$\sigma$-algebrą zawierającą $\mathcal B(X)$,
a ograniczenie miary zewnętrznej do tej rodziny jest zupełne.
Ograniczenie $\mathcal H^s$ do $\mathcal B(X)$ jest
miarą borelowską (niekoniecznie zupełną).
\end{twierdzenie}
\begin{proof}
Lemat~\ref{lem:hausdorff-outer-26} daje
wszystkie aksjomaty miary zewnętrznej,
a lemat~\ref{lem:hausdorff-metricity-26}
jej metryczną addytywność.
Lemat~\ref{lem:hausdorff-metric-open-26} mówi, że
każdy zbiór domknięty, a więc również każdy otwarty,
jest mierzalny. Rodzina zbiorów mierzalnych jest
$\sigma$-algebrą na mocy
lematu~\ref{lem:hausdorff-caratheodory-selfcontained-26};
zawiera więc $\sigma$-algebrę generowaną przez
zbiory otwarte, czyli $\mathcal B(X)$.
Ten sam lemat daje przeliczalną addytywność
i zupełność na pełnej $\sigma$-algebrze
Carathéodory'ego. Nie użyliśmy żadnego twierdzenia
o mierzalności borelowskiej metrycznych miar zewnętrznych
jako przesłanki.
\end{proof}

\begin{twierdzenie}[Regularizacja przez zbiory $G_\delta$]
\label{thm:hausdorff-gdelta-26}\label{it:thm:hausdorff-borel-hull-zero}
Niech $(X,d)$ będzie przestrzenią metryczną i $s\ge0$.
Dla dowolnego $E\subseteq X$ istnieje zbiór
$G\in\mathcal B(X)$ będący przeliczalnym
przecięciem zbiorów otwartych, taki że
$E\subseteq G$ i $\mathcal H^s(G)=\mathcal H^s(E)$.
W szczególności miara zewnętrzna Hausdorffa ma
borelowskie nadzbiory o tej samej mierze.
\end{twierdzenie}
\begin{proof}
Gdy $\mathcal H^s(E)=+\infty$, wystarczy wziąć $G=X$.
Niech $M:=\mathcal H^s(E)<\infty$ i $\delta_n=2^{-n}$.
Z definicji infimum wybieramy dla każdego $n$
pokrycie $E$ zbiorami $U_{nj}$ o średnicach
nieprzekraczających $\delta_n/2$, o łącznym
koszcie co najwyżej $M+1/(2n)$.
Puste zbiory pozostawiamy puste. Każdy niepusty
$U_{nj}$ powiększamy do otwartego otoczenia
$V_{nj}=\{x:\operatorname{dist}(x,U_{nj})<\eta_{nj}\}$
o średnicy
mniejszej niż $\delta_n$, wybierając grubość
otoczenia tak małą, by jego koszt przekraczał
koszt $U_{nj}$ najwyżej o $2^{-j}/(2n)$.
Jest to możliwe, ponieważ
$\operatorname{diam}V_{nj}\le\operatorname{diam}U_{nj}+2\eta_{nj}$
(z nierówności trójkąta), a dla $s>0$ funkcja
$t\mapsto c_st^s$ jest ciągła, a dla $s=0$
koszt każdego niepustego zbioru pozostaje równy $1$.
Połóżmy $O_n=\bigcup_jV_{nj}$ i $G=\bigcap_nO_n$.
Wtedy $E\subseteq G$, a $O_n$ jest otwarty i
\[
 \mathcal H^s_{\delta_n}(O_n)
 \leq\sum_jh_s(V_{nj})\leq M+1/n.
\]
Dla dowolnego $\delta>0$ i każdego dostatecznie
dużego $n$ (gdy $\delta_n<\delta$) mamy
$\mathcal H^s_\delta(G)\leq
\mathcal H^s_{\delta_n}(O_n)\leq M+1/n$.
Po przejściu $n\to\infty$ otrzymujemy
$\mathcal H^s_\delta(G)\leq M$.
Biorąc supremum po $\delta>0$,
dostajemy $\mathcal H^s(G)\leq M$;
monotoniczność daje równość.
\end{proof}

\begin{uwaga}[Alternatywna otoczka borelowska]
Można również wybrać nadzbiór $B$ typu $F_{\sigma\delta}$ o tej samej mierze.
Gdy $\mathcal H^s(E)=\infty$, wybieramy $B=X$;
monotoniczność daje żądaną równość. W przeciwnym
razie dla każdego $n\ge1$ wybieramy pokrycie $(U_{nj})_{j\ge1}$
o średnicach co najwyżej $1/n$ i koszcie nie większym
niż $\mathcal H^s_{1/n}(E)+1/n$.
Domknięcie nie zmienia średnicy: z ciągłości metryki
każdą odległość dwóch punktów domknięcia można
przybliżyć odległościami punktów zbioru.
Domknięcie zachowuje też niepustość.
Zatem pokrycia przez $\overline{U_{nj}}$ mają ten sam koszt.
Kładziemy
\[
 B=\bigcap_{n\ge1}\bigcup_{j\ge1}\overline{U_{nj}}.
\]
Jest to zbiór borelowski zawierający $E$.
Dla każdego $n$ odpowiednia rodzina domkniętych zbiorów
pokrywa również $B$, stąd
$\mathcal H^s_{1/n}(B)\le\mathcal H^s(E)+1/n$.
Przejście do granicy oraz monotoniczność dają równość.
\end{uwaga}

\begin{uwaga}
Własność z twierdzenia~\ref{thm:hausdorff-gdelta-26}
nie oznacza regularności względem zbiorów otwartych.
Na przykład dla $\mathcal H^0$ na $\mathbb R$
singleton ma miarę $1$, ale każdy jego otwarty
nadzbiór ma miarę $+\infty$.
\end{uwaga}

\begin{przyklad}[Wymiar zero: miara licząca]
\label{prz:hausdorff-counting-26}
Dla każdej przestrzeni metrycznej i zbioru $E\subseteq X$
miara $\mathcal H^0(E)$ jest liczbą punktów $E$,
z wartością $+\infty$, gdy $E$ jest nieskończony.
Istotnie, dla $E=\{x_1,\ldots,x_N\}$
wybieramy $\delta$ mniejszą niż odległość
każdej pary różnych punktów.
Każdy element pokrycia o średnicy $\leq\delta$
może wtedy zawierać najwyżej jeden punkt $E$;
pokrycie singletonami osiąga koszt $N$.
Jeżeli $E$ jest nieskończony, dla każdego $N$
wybieramy $N$ różnych punktów i stosujemy
monotoniczność: $\mathcal H^0(E)\geq N$.
\end{przyklad}

\begin{twierdzenie}[W wymiarze jeden otrzymujemy miarę Lebesgue'a]
\label{thm:hausdorff-line-26}
Dla standardowej normalizacji $c_1=1$ oraz
każdego $E\subseteq\mathbb R$ zachodzi
$\mathcal H^1(E)=\lambda_1^*(E)$.
Tutaj $\lambda_1^*(E)$ oznacza infimum
$\sum_j(b_j-a_j)$ po wszystkich przeliczalnych
pokryciach $E$ otwartymi przedziałami $(a_j,b_j)$.
Jest to konstrukcja miary zewnętrznej Lebesgue'a
na prostej, podana również w rozdziale czwartym.
\end{twierdzenie}
\begin{proof}
Każdy niepusty zbiór $U\subseteq\mathbb R$
o skończonej średnicy $a$ zawiera się w przedziale
domkniętym długości $a$ (między swoim infimum
i supremum). Dla dowolnego dopuszczalnego pokrycia
$E\subseteq\bigcup_jU_j$ i $\varepsilon>0$ każdy niepusty $U_j$ pokrywamy otwartym przedziałem długości nie większej niż $\operatorname{diam}U_j+\varepsilon 2^{-j}$, również gdy średnica wynosi zero. Z definicji miary zewnętrznej
$\lambda_1^*(E)\leq\sum_j\operatorname{diam}U_j+\varepsilon$. Dowolność $\varepsilon$ daje nierówność bez składnika $\varepsilon$.
Po przejściu do infimum i granicy
$\lambda_1^*(E)\leq\mathcal H^1(E)$.

W drugą stronę, jeśli $\lambda_1^*(E)<\infty$,
z definicji miary zewnętrznej Lebesgue'a
wybieramy przeliczalne pokrycie $E$ otwartymi
przedziałami $I_j$ o łącznej długości
mniejszej od $\lambda_1^*(E)+\varepsilon$.
Domknięcie każdego $I_j$ dzielimy na skończenie wiele
przedziałów domkniętych o długościach
nieprzekraczających $\delta$; suma długości
tych odcinków jest równa długości $I_j$.
Otrzymujemy dopuszczalne pokrycie Hausdorffa,
zatem $\mathcal H^1_\delta(E)\leq
\lambda_1^*(E)+\varepsilon$.
Ponieważ $\delta,\varepsilon>0$ były dowolne,
$\mathcal H^1(E)\leq\lambda_1^*(E)$.
W przypadku $\lambda_1^*(E)=+\infty$
pierwsza nierówność już kończy dowód.
\end{proof}

\begin{twierdzenie}[Skalowanie i odwzorowania Lipschitzowskie]
\label{thm:hausdorff-lipschitz-26}\label{it:thm:hausdorff-lipschitz-dimension}
Niech $(X,d_X)$ i $(Y,d_Y)$ będą przestrzeniami
metrycznymi, $s\ge0$, a $T:X\to Y$ będzie
odwzorowaniem $L$-Lipschitzowskim, tj.
$d_Y(Tx,Ty)\le Ld_X(x,y)$ dla wszystkich $x,y$,
gdzie $L>0$. Wówczas dla każdego $E\subseteq X$
$\mathcal H_Y^s(T(E))\le L^s\mathcal H_X^s(E)$.
W szczególności dla podobieństwa $T$ przestrzeni
euklidesowej o skali $L$ zachodzi
$\mathcal H^s(T(E))=L^s\mathcal H^s(E)$.
\end{twierdzenie}
\begin{proof}
Jeżeli $E\subseteq\bigcup_jU_j$ i
$\operatorname{diam}U_j\le\delta$, to
$T(E)\subseteq\bigcup_jT(U_j)$ i
$\operatorname{diam}T(U_j)\le L\operatorname{diam}U_j
\le L\delta$. Wobec tego
$\mathcal H^s_{L\delta}(T(E))\le
L^s\mathcal H^s_\delta(E)$.
Przechodzimy $\delta\downarrow0$. W przypadku
podobieństwa stosujemy tę samą nierówność także do
$T^{-1}$, którego stała Lipschitza wynosi $L^{-1}$.
\end{proof}

\begin{wniosek}[Punkty i zbiory przeliczalne]
\label{cor:hausdorff-singletons-auditVI}
Niech $(X,d)$ będzie przestrzenią metryczną i $s>0$.
Wtedy $\mathcal H^s(\{x\})=0$ dla każdego $x\in X$.
W konsekwencji $\mathcal H^s(A)=0$ dla każdego zbioru
przeliczalnego $A\subseteq X$. Jeśli $T:X\to Y$ jest
odwzorowaniem lokalnie Lipschitzowskim w sensie pokrycia
$X$ przeliczalnie wieloma zbiorami, na których $T$
jest Lipschitzowskie, to
$\mathcal H_Y^s(T(A))=0$ zawsze, gdy $\mathcal H_X^s(A)=0$.
\end{wniosek}
\begin{proof}
Dla każdego $\delta>0$ zbiór jednoelementowy
można pokryć nim samym, którego koszt w wymiarze
$s>0$ wynosi zero. Przeliczalna podaddytywność
miary zewnętrznej dowodzi drugiego zdania.
W ostatnim zdaniu piszemy $X=\bigcup_jX_j$,
przy czym ograniczenie $T|_{X_j}$ ma stałą
Lipschitza $L_j$. Na mocy
Twierdzenia~\ref{thm:hausdorff-lipschitz-26}
$\mathcal H_Y^s(T(A\cap X_j))\le
L_j^s\mathcal H_X^s(A\cap X_j)=0$.
Ponownie stosujemy podaddytywność.
\end{proof}

\section{Wymiar Hausdorffa jako próg zanikania miar}
\begin{lemat}[Próg zerowania]
\label{lem:hausdorff-threshold-26}
Niech $E\subseteq X$ oraz $0\le s<t$. Jeżeli
$\mathcal H^s(E)<\infty$, to $\mathcal H^t(E)=0$.
Jeżeli $\mathcal H^t(E)>0$, to $\mathcal H^s(E)=+\infty$.
\end{lemat}
\begin{proof}
Rozważmy niepusty zbiór $U_j$ w pokryciu,
którego średnica nie przekracza $\delta$.
Z definicji kosztu wynika
\[
 h_t(U_j)\le\frac{c_t}{c_s}\,\delta^{t-s}h_s(U_j).
\]
Dla $s=0$ stosujemy przy tym $h_0(U_j)=1$.
Dla zbioru pustego obie strony są zerowe.
Sumując koszty i biorąc infimum, otrzymujemy
\[
\mathcal H^t_\delta(E)\le
\frac{c_t}{c_s}\,\delta^{t-s}\mathcal H^s_\delta(E).
\]
Jeżeli $\mathcal H^s(E)<\infty$, prawa strona dąży
do zera. Druga implikacja jest kontrapozycją pierwszej.
\end{proof}

\begin{definicja}[Wymiar Hausdorffa]
\label{def:hausdorff-dimension-26}
Dla $E\subseteq X$ definiujemy
\[
 \dim_H E:=\inf\{s\ge0:\mathcal H^s(E)=0\}
 =\sup\{s\ge0:\mathcal H^s(E)=+\infty\},
\]
z konwencją $\sup\varnothing=0$ (w szczególności
$\dim_H\varnothing=0$) i $\inf\varnothing=+\infty$.
Równość obu kresów wynika z lematu
\ref{lem:hausdorff-threshold-26}: jeśli $a$ jest
pierwszym kresem, to dla $0\le s<a$ wybieramy
$s<t<a$ (także gdy $a=+\infty$). Wtedy
$\mathcal H^t(E)>0$, więc $\mathcal H^s(E)=+\infty$.
Jeżeli $a<+\infty$ i $s>a$, istnieje $t<s$
z $\mathcal H^t(E)=0$, więc $\mathcal H^s(E)=0$.
Wraz z przyjętymi konwencjami dowodzi to równości.
Przy skończonej wartości granicznej $s=\dim_H E$ miara
$\mathcal H^s(E)$ może być zerowa, skończona dodatnia
lub nieskończona.
\end{definicja}

\begin{wniosek}[Wymiar sumy przeliczalnej]
\label{cor:hausdorff-union-26}
Dla dowolnych zbiorów $E_n\subseteq X$ zachodzi
$\dim_H(\bigcup_{n\ge1}E_n)=\sup_n\dim_HE_n$.
W szczególności zbiory przeliczalne mają wymiar zero.
\end{wniosek}
\begin{proof}
Monotoniczność daje nierówność $\ge$.
Jeżeli $s>\sup_n\dim_HE_n$, to
$\mathcal H^s(E_n)=0$ dla każdego $n$; przeliczalna
podaddytywność daje $\mathcal H^s(\bigcup_nE_n)=0$.
Jeżeli supremum wymiarów jest skończone, przechodzimy
z $s$ do tego supremum z góry, otrzymując nierówność
odwrotną; dla supremum nieskończonego pierwsza
nierówność już daje tezę. Każdy pojedynczy punkt ma miarę
$\mathcal H^s$ równą zero dla $s>0$, gdyż można
pokryć go zbiorem o dowolnie małej średnicy.
\end{proof}

\begin{uwaga}[Wymiar obrazu i metryka indukowana]
Z twierdzenia \ref{thm:hausdorff-lipschitz-26} wynika, że
$\dim_H T(E)\le\dim_H E$: jeżeli $\dim_HE<\infty$,
to dla każdego $s>\dim_HE$ zachodzi $\mathcal H^s(E)=0$,
a więc $\mathcal H^s(T(E))=0$.
Przypadek $\dim_HE=\infty$ nie wymaga dodatkowego argumentu.
Jeżeli ograniczenie $T|_E$ jest bi-lipschitzowskie, stosujemy
nierówność także do odwrotności, uzyskując równość wymiarów.
Miara Hausdorffa podzbioru podprzestrzeni metrycznej jest
taka sama jak w przestrzeni otaczającej: pokrycia w podprzestrzeni
są dopuszczalne w przestrzeni otaczającej, a w przeciwnym kierunku
przecinamy każdy zbiór pokrycia z podprzestrzenią,
nie zwiększając średnic ani kosztu.
Nie zakładamy mierzalności ani ośrodkowości rozważanych zbiorów.
\end{uwaga}

\begin{twierdzenie}[Wymiar przestrzeni euklidesowej]
\label{thm:euclidean-dimension-26}
Jeżeli $E\subseteq\mathbb R^d$ jest mierzalny
w sensie Lebesgue'a i $\lambda_d(E)>0$, to $\dim_HE=d$. Każdy podzbiór
$\mathbb R^d$ ma wymiar Hausdorffa co najwyżej $d$.
Co więcej, $\mathcal H^d$ i $\lambda_d$ są wzajemnie
porównywalne na zbiorach borelowskich, z dodatnimi
stałymi zależnymi tylko od $d$.
\end{twierdzenie}
\begin{proof}
Niech $d$ będzie dodatnią liczbą całkowitą i niech
$v_d=\lambda_d(B(0,1))$. Najpierw niech $A\subseteq\mathbb R^d$
będzie dowolnym zbiorem. Każdy niepusty zbiór $U$ o skończonej średnicy $\rho$ mieści się w domkniętej
kuli o promieniu $\rho$ (o środku w dowolnym punkcie $U$),
a zatem $\lambda_d^*(U)\le v_d\rho^d$.
Dla pokrycia $A\subseteq\bigcup_jU_j$ zbiorami o
średnicach $\le\delta$ mamy
$\lambda_d^*(A)\le v_d\sum_j(\operatorname{diam}U_j)^d$.
Stąd $\mathcal H^d(A)\ge(c_d/v_d)\lambda_d^*(A)$.

Dla ograniczonego $A$ wybierzmy otwarty ograniczony
zbiór $O\supseteq A$ o mierze dowolnie bliskiej
$\lambda_d^*(A)$. Z ustalonej siatki diadycznej wybieramy maksymalne
kostki półotwarte $Q$, których domknięcia mieszczą się
w $O$ i których boki nie przekraczają $\delta/\sqrt d$.
Dla każdego $x\in O$ istnieje dostatecznie mała kostka
siatki zawierająca $x$ i mająca domknięcie w $O$;
każdą taką kostkę można powiększać do kolejnych
przodków w siatce, aż dalsze powiększenie naruszy
ograniczenie długości boku lub warunek zawierania
domknięcia w $O$. Uzyskujemy kostkę maksymalną.
Dwie kostki tej samej siatki diadycznej są rozłączne
lub jedna zawiera się w drugiej, więc wybrane kostki
półotwarte są parami rozłączne. Ich domknięcia pokrywają $O$.
Wnętrza maksymalnych kostek są parami rozłączne
i zawarte w $O$, toteż suma ich objętości nie
przekracza $\lambda_d(O)$. Każda domknięta kostka
ma średnicę równą $\sqrt d$ razy długość boku,
więc jej koszt przy $s=d$ jest równy $c_dd^{d/2}$
razy jej objętości. Z tego wynika
$\mathcal H^d_\delta(A)\le c_dd^{d/2}\lambda_d(O)$.
Przechodząc do infimum po $O$ i do granicy po $\delta$,
dostajemy porównywalność na ograniczonych zbiorach.
Dla zbioru borelowskiego nieograniczonego stosujemy
ją do $A\cap B(0,n)$ i przechodzimy do granicy
z ciągłości obu miar od dołu. W szczególności $\mathcal H^d(E)>0$, jeśli
$\lambda_d(E)>0$.

Dla $s>d$ i $L>0$ dzielimy kostkę $[0,L]^d$
na $m^d$ domkniętych kostek, gdzie
$m=\lceil L\sqrt d/\delta\rceil$ (dla $\delta$ małego).
Każda ma średnicę $L\sqrt d/m\le\delta$.
Stąd
\[
 \mathcal H^s_\delta([0,L]^d)
 \le c_s m^d(L\sqrt d/m)^s
 =c_s(L\sqrt d)^s m^{d-s}\longrightarrow0
 \quad(\delta\downarrow0).
\]
Całe $\mathbb R^d$ jest przeliczalną sumą translacji takich
kostek, których miara Hausdorffa pozostaje niezmieniona
na mocy twierdzenia~\ref{thm:hausdorff-lipschitz-26}, więc $\dim_HE\le d$ dla każdego $E$.
Połączenie z dodatniością $\mathcal H^d(E)$ daje
$\dim_HE=d$ przy $\lambda_d(E)>0$.
\end{proof}

\begin{twierdzenie}[Kalibracja wymiaru całkowitego bez twierdzeń geometrycznych]
\label{thm:hausdorff-calibration-26}
Niech $d\ge1$ będzie liczbą całkowitą. Zapiszmy
$\widetilde{\mathcal H}^{d}$ dla miary Hausdorffa
na $\mathbb R^d$ skonstruowanej powyżej z $c_d=1$
i połóżmy $\kappa_d=\widetilde{\mathcal H}^{d}([0,1]^d)$.
Wtedy $0<\kappa_d<\infty$ oraz dla każdego
zbioru borelowskiego $E\subseteq\mathbb R^d$
zachodzi
\[
 \widetilde{\mathcal H}^{d}(E)=\kappa_d\lambda_d(E).
\]
W szczególności wybór $c_d=\kappa_d^{-1}$ daje
miarę Hausdorffa spełniającą
$\mathcal H^d|_{\mathcal B(\mathbb R^d)}=\lambda_d$.
Nie potrzebujemy do tego nierówności izodiametrycznej
ani twierdzenia o mierze Haara.
\end{twierdzenie}
\begin{proof}
Dodatniość i skończoność $\kappa_d$ wynikają z
udowodnionych wyżej dwustronnych oszacowań
$\widetilde{\mathcal H}^{d}$ przez miarę Lebesgue'a.
Dla podobieństwa $S(x)=ax+b$, gdzie $a>0$,
obrazy pokryć o średnicach co najwyżej $\delta$
są pokryciami o średnicach co najwyżej $a\delta$
i kosztach pomnożonych przez $a^d$.
Zastosowanie tego samego argumentu do $S^{-1}$
daje równość
$\widetilde{\mathcal H}^{d}_{a\delta}(S(E))
=a^d\widetilde{\mathcal H}^{d}_\delta(E)$.
Przy $\delta\downarrow0$ otrzymujemy
$\widetilde{\mathcal H}^{d}(S(E))
=a^d\widetilde{\mathcal H}^{d}(E)$,
w szczególności niezmienniczość względem translacji.

Wykażmy najpierw, że brzeg sześcianu ma zerową
$\widetilde{\mathcal H}^{d}$-miarę.
Dla $d=1$ brzeg składa się z dwóch singletonów,
które mają zerową miarę Hausdorffa wymiaru jeden.
Dla $d\ge2$ każdą z $2d$ ścian sześcianu jednostkowego
dzielimy na $m^{d-1}$ domkniętych kostek w jej
$(d-1)$-wymiarowej płaszczyźnie, gdzie
$m=\lceil\sqrt{d-1}/\delta\rceil$ dla $0<\delta<1$.
Każda ma średnicę $\sqrt{d-1}/m\le\delta$.
Koszt całego pokrycia brzegu spełnia oszacowanie
\[
 \widetilde{\mathcal H}^{d}_\delta(\partial[0,1]^d)
 \le 2d\,m^{d-1}
          \left(\frac{\sqrt{d-1}}m\right)^d
 =\frac{2d(d-1)^{d/2}}m\longrightarrow0.
\]
W obu przypadkach brzeg ma zerową miarę. Wobec tego
$\widetilde{\mathcal H}^{d}([0,1)^d)=\kappa_d$.
Z niezmienniczości i skalowania wynika, że każdy
półotwarty sześcian o boku $a$ ma miarę
$\kappa_da^d$.

Niech $R=\prod_{i=1}^d[a_i,b_i)$ będzie
prostokątem o wymiernych końcach.
Wybieramy dodatnią liczbę wymierną $h$, która
dzieli długość każdego boku $b_i-a_i$.
Podział na skończenie wiele półotwartych sześcianów
o boku $h$ jest rozłączny, zatem przeliczalna
addytywność daje
$\widetilde{\mathcal H}^{d}(R)
=\kappa_d\prod_{i=1}^d(b_i-a_i)
=\kappa_d\lambda_d(R)$.

Rodzina takich prostokątów wraz ze zbiorem pustym
jest układem $\pi$ generującym $\mathcal B(\mathbb R^d)$.
Obie miary są skończone na prostokątach ograniczonych,
a $\mathbb R^d$ wyczerpuje rosnący ciąg prostokątów
o wymiernych końcach. Stosując twierdzenie
$\pi$--$\lambda$ o jednoznaczności miar (którego dowód podaliśmy w rozdziale pierwszym) do
każdego z tych prostokątów, otrzymujemy równość
na wszystkich zbiorach borelowskich.
Równość rozciąga się nawet na miary zewnętrzne wszystkich
podzbiorów $A\subseteq\mathbb R^d$:
dla $\widetilde{\mathcal H}^d$ stosujemy twierdzenie
\ref{thm:hausdorff-gdelta-26}, a dla $\lambda_d^*$
bierzemy borelowski nadzbiór o tej samej mierze zewnętrznej
(otrzymany jako przecięcie otwartych nadzbiorów
o miarach dążących do $\lambda_d^*(A)$).
Jeśli odpowiednia miara zewnętrzna jest nieskończona,
można wziąć całą przestrzeń. W obu kierunkach
monotoniczność i wykazana równość na zbiorach
borelowskich dają
$\widetilde{\mathcal H}^d(A)=\kappa_d\lambda_d^*(A)$.
\end{proof}

\begin{uwaga}[Konwencja normalizacji w dalszych rozdziałach]
\label{rem:hausdorff-convention-26}
Od tego miejsca w całym skrypcie, dla każdego całkowitego
$k\ge1$, symbol $\mathcal H^k$ oznacza miarę Hausdorffa
zdefiniowaną w Definicji~\ref{def:hausdorff-measure-26}
z wyborem $c_k=\kappa_k^{-1}$ ustalonym w
Twierdzeniu~\ref{thm:hausdorff-calibration-26}.
W szczególności $\mathcal H^k$ na $\mathbb R^k$
zgadza się z $\lambda_k$. Jeśli $P$ jest izometrycznie
osadzoną $k$-płaszczyzną w $\mathbb R^N$, dla $E\subseteq P$
miara obliczona wewnątrz $P$ jest równa mierze
obliczonej w $\mathbb R^N$: każde pokrycie w $P$
jest pokryciem w przestrzeni otaczającej, a każde
pokrycie w przestrzeni otaczającej można przeciąć z $P$,
nie zwiększając średnic ani kosztów.
Izometria $\mathbb R^k\to P$ zachowuje miarę Hausdorffa
na mocy twierdzenia~\ref{thm:hausdorff-lipschitz-26}
zastosowanego w obu kierunkach.
Dla $k=0$ zachowujemy $\mathcal H^0$ jako miarę liczącą.
Dla niecałkowitych $s>0$ pozostajemy przy dowolnej,
raz ustalonej dodatniej stałej $c_s$; jej wybór nie
zmienia wymiaru Hausdorffa. W konsekwencji wzory
na pole powierzchni, koobszar i obwód w rozdziałach
\ref{ch:area-coarea-28}--\ref{ch:geometric-measure-29}
nie zawierają niejawnych czynników normalizacyjnych.
\end{uwaga}

\begin{uwaga}[Stała geometryczna a stała kalibracyjna]
Często stosuje się jawny współczynnik
$c_d=\omega_d2^{-d}$, gdzie $\omega_d$ jest
objętością euklidesowej kuli jednostkowej.
Wykazanie, że ten współczynnik jest równy
$\kappa_d^{-1}$, wymaga jeszcze nierówności
izodiametrycznej. Nie korzystamy z tej
identyczności: powyższa kalibracja jest
zdefiniowana w pełni wewnętrznie i wystarcza
do wzorów geometrycznych w rozdziale 28.
\end{uwaga}

\begin{wniosek}[Całkowanie na afinicznej płaszczyźnie]
\label{it:cor:hausdorff-affine-integration}
Niech $1\le k\le N$ i niech $P\subset\mathbb R^N$
będzie $k$-wymiarową płaszczyzną afiniczną.
Niech $T:\mathbb R^k\to P$ będzie izometrią afiniczną
na $P$. Przyjmujemy kalibrację
$c_k=\kappa_k^{-1}$ z twierdzenia
\ref{thm:hausdorff-calibration-26} (dla $k=1$ mamy $c_1=1$).
Dla borelowskiego $A\subset P$ i nieujemnej borelowskiej
funkcji $f$ na $A$ zachodzi
\[
 \int_A f\,\mathrm d\mathcal H^k
 =\int_{T^{-1}(A)} f(T(u))\,\mathrm d\lambda_k(u).
\]
Dla borelowskiej funkcji rzeczywistej lub zespolonej
równość zachodzi, jeśli jedna ze stron jest bezwzględnie
całkowalna; wówczas druga również jest.
Ponadto $\lambda_N(P)=0$, jeżeli $k<N$.
\end{wniosek}
\begin{proof}
Pokrycia podzbiorów $P$ można rozpatrywać w metryce
indukowanej, nie zmieniając kosztu: pokrycie w otaczającej
przestrzeni przecinamy z $P$, a pokrycie w $P$ jest
także pokryciem w $\mathbb R^N$.
Izometria $T$ i jej odwrotność zachowują średnice,
więc z definicji miary Hausdorffa i kalibracji
$\mathcal H^k(A)=\lambda_k(T^{-1}(A))$.
Inaczej mówiąc, $\mathcal H^k|_P=T_*\lambda_k$
na borelowskich podzbiorach $P$.
Wzór całkowy wynika najpierw dla funkcji charakterystycznych,
następnie prostych nieujemnych, a przez twierdzenie
\ref{thm:mct} dla dowolnych funkcji nieujemnych.
Stosując go do $|f|$, uzyskujemy równoważność
całkowalności; rozkład na części dodatnią i ujemną
oraz rzeczywistą i urojoną daje tezę dla całek
rzeczywistych i zespolonych.
Jeśli $k<N$, to z nierówności lipschitzowskiej
\ref{thm:hausdorff-lipschitz-26} i
twierdzenia \ref{thm:euclidean-dimension-26} wynika
$\dim_H P\le k<N$. Z definicji wymiaru
$\mathcal H^N(P)=0$. Kalibracja w wymiarze $N$
daje $\lambda_N(P)=0$.
\end{proof}

\begin{przyklad}[Całka na odcinku w przestrzeni euklidesowej]
Niech $p,q\in\mathbb R^N$, $p\ne q$, $L=|q-p|$,
a $A=[p,q]$ będzie odcinkiem o tych końcach.
Izometria $t\mapsto p+t(q-p)/L$ identyfikuje
$[0,L]$ z $A$, więc poprzedni wniosek i liniowe
podstawienie $t=Lu$ dają
\[
 \int_A f\,\mathrm d\mathcal H^1
 =L\int_0^1 f(p+u(q-p))\,\mathrm d\lambda_1(u).
\]
Dla $f=1$ otrzymujemy $\mathcal H^1(A)=L$,
a dla $f(x)=|x-p|^2$ całka wynosi
$L^3\int_0^1u^2\,\mathrm d\lambda_1(u)=L^3/3$.
Jeżeli $N>1$, to jednocześnie $\lambda_N(A)=0$.
\end{przyklad}

\section{Zasada rozkładu masy}
\begin{twierdzenie}[Zasada rozkładu masy]
\label{thm:mass-distribution-26}
Niech $E\subseteq\mathbb R^d$ będzie zbiorem borelowskim,
$s>0$, a $\mu$ skończoną dodatnią miarą borelowską
spełniającą $\mu(E)>0$ i $\mu(\mathbb R^d\setminus E)=0$.
Załóżmy, że istnieje $C>0$ takie, iż
$\mu(B(x,r))\le Cr^s$ dla wszystkich
$x\in\mathbb R^d$ i $r>0$. Wtedy
$\mathcal H^s(E)\ge(c_s/(2^sC))\mu(E)>0$,
a więc $\dim_HE\ge s$.
\end{twierdzenie}
\begin{proof}
Niech $E\subseteq\bigcup_jU_j$ będzie dowolnym
przeliczalnym pokryciem o średnicach nieprzekraczających
$\delta$; dopuszczamy także zbiory o średnicy zero.
Jeśli $U_j\cap E\ne\varnothing$, wybierzmy
$x_j\in U_j\cap E$. Dla dodatniej średnicy
$U_j\subseteq B(x_j,2\operatorname{diam}U_j)$.
Nie musimy zakładać mierzalności $U_j$: zbiory
$E\cap U_j$ są zawarte w wymienionych kulach,
a funkcja $\mu$ jest przeliczalnie podaddytywna
na pokryciach przez kule borelowskie. Zbiory pokrycia
o zerowej średnicy są puste lub jednoelementowe.
Dla każdego $x$ mamy $\mu(\{x\})\le Cr^s$ dla
każdego $r>0$, a więc $\mu(\{x\})=0$.
Przeliczalna suma tych singletonów jest borelowska
i ma miarę zero. Ostatecznie
$\mu(E)\le C2^s\sum_j(\operatorname{diam}U_j)^s$.
Biorąc infimum po wszystkich dopuszczalnych pokryciach,
otrzymujemy $\mu(E)\le(C2^s/c_s)\mathcal H^s_\delta(E)$. W granicy
$\delta\downarrow0$ otrzymujemy tezę.
\end{proof}

\section{Twierdzenie Frostmana — konstrukcja miary}
\begin{lemat}[Dodatnia zawartość zbioru zwartego]
\label{lem:positive-content-26}
Niech $K\subseteq\mathbb R^d$ będzie zwarty i $s>0$.
Jeżeli $\mathcal H^s(K)>0$, to
$\mathcal H^s_\infty(K)>0$.
\end{lemat}
\begin{proof}
Załóżmy przeciwnie, że $\mathcal H^s_\infty(K)=0$.
Dla każdego $\varepsilon>0$ można pokryć $K$
zbiorami $U_j$ takimi, że
$\sum_j(\operatorname{diam}U_j)^s<\varepsilon/c_s$.
Każda średnica w tym pokryciu jest mniejsza niż
$(\varepsilon/c_s)^{1/s}$. Wybierając
$\varepsilon$ dowolnie małe, otrzymujemy
$\mathcal H^s_\delta(K)=0$ dla każdego $\delta>0$,
a więc $\mathcal H^s(K)=0$, sprzeczność.
\end{proof}

\begin{twierdzenie}[Frostmana, wersja dla zbiorów zwartych]
\label{thm:frostman-26}
Niech $K\subseteq\mathbb R^d$ będzie niepustym
zbiorem zwartym i $s>0$. Następujące warunki
są równoważne:
\begin{enumerate}[label=(\roman*)]
 \item $\mathcal H^s(K)>0$;
 \item istnieje niezerowa skończona dodatnia miara
 borelowska $\mu$, skoncentrowana na $K$, oraz stała
 $C>0$ takie, że
 $\mu(B(x,r))\le Cr^s$ dla wszystkich
 $x\in\mathbb R^d$ i $r>0$.
\end{enumerate}
Ponadto w (ii) miarę można znormalizować do
prawdopodobieństwa na $K$, zastępując stałą $C$
przez $C/\mu(K)$.
\end{twierdzenie}
\begin{proof}
Implikacja (ii)$\Rightarrow$(i) wynika z zasady
rozkładu masy~\ref{thm:mass-distribution-26}.
Pokażemy konstruktywnie implikację odwrotną, nie
zakładając twierdzenia o przepływach w grafach.

\emph{Krok 1: drzewo kostek.} Wybierzmy półotwartą
kostkę osiową $Q_0$ o boku $L>0$, tak aby
$K\subset\operatorname{int}Q_0$. Dzielimy ją na
$2^d$ kostek potomnych o połowie boku i powtarzamy
tę operację do poziomu $n$. Każda kostka na poziomie
$j$ ma bok $L2^{-j}$. Zachowujemy także kostki
nieprzecinające $K$, ale nadajemy im zerową pojemność.
Każdej kostce $Q$ przypisujemy pojemność
\[
 w(Q)=\begin{cases}
 (\sqrt d\,\operatorname{side}(Q))^s,&Q\cap K\ne\varnothing,\\
 0,&Q\cap K=\varnothing.
 \end{cases}
\]
Zdefiniujmy liczby $a_n(Q)$ rekurencyjnie od poziomu
$n$ do korzenia: dla liści $a_n(Q)=w(Q)$, a dla
pozostałych węzłów
\begin{equation}\label{eq:frostman-recursion-26}
 a_n(Q)=\min\left\{w(Q),\sum_{Q'\text{ potomek bezpośredni }Q}
 a_n(Q')\right\}.
\end{equation}
Indukcja od liści pokazuje, że $a_n(Q)$ jest minimalną
sumą pojemności kostek z poziomów od poziomu $Q$
do $n$, których suma pokrywa $K\cap Q$ w strukturze
drzewa. Istotnie, wybieramy albo samą kostkę $Q$
(koszt $w(Q)$), albo pokrywamy oddzielnie każdy z jej
$2^d$ potomków (koszt sumy minimów). W przypadku pustego $K\cap Q$ koszt minimalny
wynosi zero (możemy użyć pustej rodziny).
W szczególności każde pokrycie uzyskane dla korzenia
wyznacza pokrycie zbioru $K$ przez domknięcia
wybranych kostek, których średnice nie przekraczają
wartości użytych w kosztach. Stąd, na mocy lematu
\ref{lem:positive-content-26},
\begin{equation}\label{eq:frostman-positive-26}
 a_n(Q_0)\ge\mathcal H^s_\infty(K)/c_s=:m_0>0.
\end{equation}
Równocześnie $a_n(Q_0)\le w(Q_0)$.

\emph{Krok 2: rozdzielanie masy po drzewie.} Nadajemy
korzeniowi masę $m_n(Q_0)=a_n(Q_0)$.
Jeżeli masa $m_n(Q)$ jest już określona i
$m_n(Q)\le a_n(Q)$, to ponieważ
$\sum_{Q'}a_n(Q')\ge a_n(Q)\ge m_n(Q)$,
można rozdzielić $m_n(Q)$ między potomków $Q'$
tak, by $0\le m_n(Q')\le a_n(Q')$ oraz
$\sum_{Q'}m_n(Q')=m_n(Q)$. Można to zrobić
kolejno, napełniając potomków do ich pojemności.
Na poziomie $n$ wybieramy dla każdej kostki
$Q$ o dodatniej masie punkt $x_Q\in K\cap Q$
i definiujemy miarę atomową
$\mu_n=\sum_{Q\text{ liść},\,m_n(Q)>0}m_n(Q)\delta_{x_Q}$.
Dzięki półotwartości kostek każdy punkt liścia należy
do dokładnie jednego węzła na każdym wcześniejszym
poziomie. Z zachowania masy w rozgałęzieniach wynika
$\mu_n(Q)=m_n(Q)$ dla kostek o poziomie co najwyżej $n$.
Z konstrukcji otrzymujemy
\begin{equation}\label{eq:frostman-dyadic-bound-26}
 m_0\le\mu_n(K)\le w(Q_0),\qquad
 \mu_n(Q)\le w(Q)
 \quad\text{dla każdej kostki }Q\text{ o poziomie }\le n.
\end{equation}

\emph{Krok 3: oszacowanie kul.} Niech $0<r<L/2$
i wybierzmy poziom $j$ taki, że
$r\le\ell_j=L2^{-j}<2r$ (przy $r$ wystarczająco
małym taki $j\ge0$ istnieje). Kula $B(x,r)$
przecina co najwyżej $4^d$ półotwartych kostek
diadycznych poziomu $j$: w każdym kierunku jej
rzut ma długość $2r\le2\ell_j$, więc przecina
co najwyżej cztery przedziały siatki.
Dla każdego $n\ge j$ z
\eqref{eq:frostman-dyadic-bound-26} wynika
\[
 \mu_n(B(x,r))\le4^d(\sqrt d\,\ell_j)^s
 \le4^d(2\sqrt d)^s r^s.
\]
Dla większych $r$ analogiczne oszacowanie wynika z
$\mu_n(K)\le w(Q_0)$ po powiększeniu stałej.
Otrzymaliśmy stałą $C_{d,s,L}$ niezależną od $x,r,n$,
taką że dla każdego ustalonego $r>0$ i każdego $x$
oszacowanie $\mu_n(B(x,r))\le C_{d,s,L}r^s$
zachodzi dla wszystkich dostatecznie dużych $n$
(dla małych $r$ wymagamy $n\ge j(r)$).
Nie zachodzi ono jednocześnie dla wszystkich $n$ i $r$:
miary $\mu_n$ są atomowe. Do przejścia do granicy
przy ustalonej kuli wystarczy właśnie oszacowanie ogona ciągu.

\emph{Krok 4: granica słaba.} Miary $\mu_n$ są
skoncentrowane na zwartym $K$ i mają jednostajnie
ograniczone masy. Twierdzenie o zwartości miar na
przestrzeni zwartej, dowiedzione w rozdziale
\ref{ch:weak-measures-22}, pozwala wybrać podciąg
zbieżny wąsko do skończonej miary $\mu$ na $K$.
Testując funkcję stałą $1$, otrzymujemy
$\mu(K)=\lim\mu_{n_k}(K)\ge m_0>0$.
Ustalmy dowolne $x\in\mathbb R^d$ i $r>0$.
Zbiór $B(x,r)\cap K$ jest otwarty w $K$. Względem
miar na $K$ twierdzenie Portmanteau daje (dla tej
ustalonej kuli oszacowanie z kroku 3 zachodzi
od pewnego indeksu)
$\mu(B(x,r))\le\liminf_k\mu_{n_k}(B(x,r))
\le C_{d,s,L}r^s$. Jest to (ii).
Normalizując przez $\mu(K)>0$, otrzymujemy miarę
prawdopodobieństwa. Cała konstrukcja została
przeprowadzona wewnątrz zwartego $K$.
\end{proof}

\begin{uwaga}[Zakres twierdzenia]
Dowiedliśmy pełnej wersji dla \emph{zbiorów zwartych}.
Rozszerzenie na dowolne zbiory analityczne lub borelowskie
korzysta z dodatkowego twierdzenia o pojemnościach
(regularności od wewnątrz w sensie pojemności).
Nie wolno automatycznie zastępować w powyższym dowodzie
zwartości przez samą mierzalność.
\end{uwaga}

\section{Zbiór Cantora i przykłady wymiarów ułamkowych}
\begin{twierdzenie}[Wymiar klasycznego zbioru Cantora]
\label{thm:cantor-dimension-26}
Niech $C\subseteq[0,1]$ będzie klasycznym zbiorem
Cantora uzyskanym przez kolejne usuwanie otwartych
środkowych trzecich części przedziałów. Wtedy
\[
 \dim_HC=\frac{\log2}{\log3}.
\]
\end{twierdzenie}
\begin{proof}
Niech $s=\log2/\log3$, zatem $2\cdot3^{-s}=1$.
Po $n$ krokach zbiór $C$ zawiera się w sumie
$2^n$ domkniętych przedziałów długości $3^{-n}$.
Dla $t>s$ otrzymujemy
$\mathcal H^t_{3^{-n}}(C)\le
c_t2^n3^{-nt}=c_t(2\cdot3^{-t})^n\to0$,
stąd $\dim_HC\le s$.

Dla nierówności odwrotnej skonstruujmy miarę
Cantora. Dla każdego $m$ wybieramy po jednym punkcie
$C$ z każdego przedziału poziomu $m$ (na przykład
jego lewy koniec) i nadajemy mu masę $2^{-m}$.
Otrzymujemy miarę probabilistyczną $\eta_m$ na zwartym
$C$. Z twierdzenia o zwartości miar z rozdziału
\ref{ch:weak-measures-22} wybieramy podciąg zbieżny
wąsko do miary probabilistycznej $\mu$ na $C$.
Jeżeli $I$ jest przedziałem poziomu $n$, to
$I\cap C$ jest otwarty i domknięty w $C$,
gdyż pozostałe przedziały tego poziomu są od niego
dodatnio oddalone. Funkcja $\chi_{I\cap C}$ jest
więc ciągła na $C$, a dla $m\ge n$ zachodzi
$\eta_m(I\cap C)=2^{m-n}2^{-m}=2^{-n}$.
Przechodząc do granicy w całce tej funkcji, otrzymujemy
$\mu(I\cap C)=2^{-n}$. Tak skonstruowana miara
jest skoncentrowana na $C$.
Pokażemy, że $\mu(B(x,r))\le Cr^s$.
Dla $0<r<1/3$ wybierzmy $n$ tak, aby
$3^{-(n+1)}<r\le3^{-n}$. Przedziały poziomu $n$
mają długość $\ell=3^{-n}$ i odległości między
różnymi przedziałami są co najmniej $\ell$.
Przedział $B(x,r)$ ma długość $2r\le2\ell$,
więc przecina co najwyżej cztery przedziały poziomu
$n$. W konsekwencji
$\mu(B(x,r))\le4\cdot2^{-n}
=4\cdot3^{-ns}\le4\cdot3^sr^s$.
Dla $r\ge1/3$ powiększamy stałą.
Zasada rozkładu masy~\ref{thm:mass-distribution-26}
daje $\mathcal H^s(C)>0$, a w szczególności
$\dim_HC\ge s$. Ponadto pokrycia przedziałami
poziomu $n$ dają $\mathcal H^s(C)\le c_s$, ponieważ
$2^n3^{-ns}=1$. Otrzymujemy więc również
$0<\mathcal H^s(C)<\infty$.
\end{proof}

\begin{wniosek}[Miara Cantora jako unormowana miara Hausdorffa]\label{it:cor:cantor-hausdorff}
Niech $K$ będzie klasycznym zbiorem Cantora i $s=\log2/\log3$. Wówczas
\[
 \kappa=\frac{\mathcal H^s|_K}{\mathcal H^s(K)}
\]
jako miary borelowskie na $K$, gdzie $\kappa$ jest miarą Cantora z przykładu \ref{ex:cantor-measure-17}, ograniczoną do $K$. Dla każdej borelowskiej funkcji $f$ całkowalnej względem $\kappa$ zachodzi
\[
 \int_K f\,\mathrm d\kappa
 =\frac{1}{\mathcal H^s(K)}\int_K f\,\mathrm d\mathcal H^s.
\]
\end{wniosek}
\begin{proof}
Z twierdzenia \ref{thm:cantor-dimension-26} wynika $0<\mathcal H^s(K)<\infty$. Dla podobieństwa $T:\mathbb{R}\to\mathbb{R}$ o skali $q>0$ obraz dowolnego pokrycia ma średnice pomnożone przez $q$. Stosując to zarówno do $T$, jak i do $T^{-1}$, otrzymujemy $\mathcal H^s_{q\delta}(T(E))=q^s\mathcal H^s_\delta(E)$, a po przejściu do granicy $\mathcal H^s(T(E))=q^s\mathcal H^s(E)$. Każdy przedział podstawowy $I$ poziomu $n$ przecina $K$ w podobnej kopii $K$ o skali $3^{-n}$. Zatem $\mathcal H^s(K\cap I)=3^{-ns}\mathcal H^s(K)=2^{-n}\mathcal H^s(K)$. Unormowana miara Hausdorffa i $\kappa$ mają tę samą masę $2^{-n}$ na każdym takim zbiorze. Zbiory podstawowe wraz ze zbiorem pustym tworzą rodzinę zamkniętą na skończone przecięcia, generującą $\mathcal B(K)$: są względnie otwarte i domknięte oraz tworzą bazę, gdyż ich średnice dążą do zera. Jednoznaczność skończonych miar na takiej rodzinie daje równość miar. Wzór całkowy wynika z tej równości, najpierw dla funkcji charakterystycznych, następnie prostych i nieujemnych przez zbieżność monotoniczną, a dla całkowalnych przez rozkład na części dodatnią i ujemną (w przypadku zespolonym także rzeczywistą i urojoną).
\end{proof}

\begin{center}
\begin{tikzpicture}[x=10cm,y=0.8cm]
\foreach \i/\y in {0/0,1/-1,2/-2,3/-3}{
 \ifnum\i=0
  \draw[line width=2pt,blue!65!black] (0,\y)--(1,\y);
 \fi
 \ifnum\i=1
  \draw[line width=2pt,blue!65!black] (0,\y)--(1/3,\y);
  \draw[line width=2pt,blue!65!black] (2/3,\y)--(1,\y);
 \fi
 \ifnum\i=2
  \foreach \a in {0,2/9,2/3,8/9}{\draw[line width=2pt,blue!65!black] (\a,\y)--({\a+1/9},\y);}
 \fi
 \ifnum\i=3
  \foreach \a in {0,2/27,2/9,8/27,2/3,20/27,8/9,26/27}{\draw[line width=2pt,blue!65!black] (\a,\y)--({\a+1/27},\y);}
 \fi
 \node[anchor=east,font=\scriptsize] at (-0.025,\y) {etap \i};
}
\end{tikzpicture}
\end{center}
\noindent\textit{Pierwsze etapy konstrukcji zbioru Cantora.}
Każde przejście zastępuje przedział dwoma
przedziałami o jednej trzeciej jego długości.

\begin{przyklad}[Gruby zbiór Cantora]
Zaczynamy od $F_0=[0,1]$. W kroku $n\ge1$
usuwamy ze środka każdego z $2^{n-1}$ pozostałych
przedziałów otwarty przedział długości $4^{-n}$.
Łączna długość usunięta do kroku $n$ wynosi
$\sum_{j=1}^n2^{j-1}4^{-j}=(1-2^{-n})/2$.
Wszystkie $2^n$ przedziały tworzące $F_n$ mają więc
długość $\ell_n=2^{-n-1}(1+2^{-n})$.
W szczególności $\ell_{n-1}>4^{-n}$, toteż każdy
krok jest wykonalny, a $\ell_n\to0$.
Zbiór $F=\bigcap_{n\ge0}F_n$ jest zwarty; z ciągłości
miary od góry otrzymujemy $\lambda_1(F)=1/2$.
Nie zawiera niepustego otwartego przedziału, ponieważ
taki przedział musiałby mieścić się w jednym przedziale
każdego poziomu $n$, a ich długości dążą do zera.
Zatem $F$ jest nigdzie gęsty. Na mocy twierdzenia~\ref{thm:euclidean-dimension-26}
ma wymiar Hausdorffa równy $1$, choć ma puste wnętrze.
Wymiar Hausdorffa nie charakteryzuje topologicznego
wymiaru ani posiadania wnętrza.
\end{przyklad}

\begin{wniosek}[Wykres funkcji Lipschitzowskiej]
\label{cor:graph-dimension-26}
Jeżeli $A\subseteq\mathbb R^m$ i
$f:A\to\mathbb R^k$ jest Lipschitzowska, to
$\dim_H\{(x,f(x)):x\in A\}=\dim_HA$.
\end{wniosek}
\begin{proof}
Odwzorowanie $x\mapsto(x,f(x))$ jest Lipschitzowskie,
więc nie zwiększa wymiaru Hausdorffa na mocy
twierdzenia~\ref{thm:hausdorff-lipschitz-26}.
Rzut $(x,y)\mapsto x$ jest również Lipschitzowski
i odwzorowuje wykres na $A$. Stosując tę samą
nierówność w drugą stronę, otrzymujemy równość.
\end{proof}

\section{Energie potencjałowe i wymiar}
\begin{definicja}[$t$-energia miary]
\label{def:energy-26}
Niech $t>0$, a $\mu$ będzie skończoną dodatnią miarą
borelowską na $\mathbb R^d$. Jej $t$-energią nazywamy
\[
 I_t(\mu)=\int_{\mathbb R^d}\!\int_{\mathbb R^d}
 |x-y|^{-t}\,\mathrm d\mu(y)\,\mathrm d\mu(x),
\]
przyjmując $|0|^{-t}=+\infty$. Całka należy do
$[0,+\infty]$; nie zakładamy jej skończoności.
\end{definicja}

\begin{twierdzenie}[Oszacowanie energii przez warunek Frostmana]
\label{thm:frostman-energy-26}
Jeżeli $\mu$ jest skończoną dodatnią miarą o
$\mu(B(x,r))\le Cr^s$ dla wszystkich $x,r$,
gdzie $s>0$, to $I_t(\mu)<\infty$ dla każdego
$0<t<s$.
\end{twierdzenie}
\begin{proof}
Z tożsamości warstwowej dla nieujemnych funkcji
$|x-y|^{-t}=t\int_{|x-y|}^{\infty}r^{-t-1}\,\mathrm dr$
i z twierdzenia Tonellego otrzymujemy
\[
 \int_{\mathbb R^d}|x-y|^{-t}\,\mathrm d\mu(y)
 =t\int_0^\infty r^{-t-1}\mu(B(x,r))\,\mathrm dr.
\]
W przedziale $(0,1)$ całka jest ograniczona przez
$Ct\int_0^1r^{s-t-1}\,\mathrm dr=Ct/(s-t)$.
W przedziale $(1,\infty)$ nie przekracza
$t\mu(\mathbb R^d)\int_1^\infty r^{-t-1}\,\mathrm dr
=\mu(\mathbb R^d)$. Oszacowanie jest jednostajne
względem $x$. Całkowanie po $x$ dowodzi tezy.
\end{proof}

\begin{twierdzenie}[Skończona energia a dolne oszacowanie wymiaru]
\label{thm:energy-dimension-26}
Niech $E\subseteq\mathbb R^d$ będzie zbiorem borelowskim,
$t>0$, a $\mu$ niezerową skończoną dodatnią miarą
borelowską skoncentrowaną na $E$. Jeśli $I_t(\mu)<\infty$,
to $\mathcal H^t(E)>0$, a zatem $\dim_HE\ge t$.
\end{twierdzenie}
\begin{proof}
Niech $U(x)=\int|x-y|^{-t}\,\mathrm d\mu(y)$.
Funkcja $U$ jest borelowska (jest półciągła z dołu
na mocy lematu Fatou) i $\int U\,\mathrm d\mu<\infty$,
więc $U$ jest skończona $\mu$-prawie wszędzie.
Istnieje liczba całkowita $M\ge1$ taka, że
$F=E\cap\{U\le M\}$ ma dodatnią $\mu$-miarę.
Kładziemy $\nu=\mu|_F$. Dla $z\in F$ i $r>0$
nierówność $|z-y|^{-t}\ge r^{-t}$ na $B(z,r)$
daje $\mu(B(z,r))\le Mr^t$.
Jeśli $B(x,r)\cap F=\varnothing$, to $\nu(B(x,r))=0$.
W przeciwnym przypadku wybieramy $z\in B(x,r)\cap F$,
a z inkluzji $B(x,r)\subseteq B(z,2r)$ otrzymujemy
$\nu(B(x,r))\le M2^tr^t$.
Zasada rozkładu masy~\ref{thm:mass-distribution-26}
zastosowana do $\nu$ na $F$ daje
$\mathcal H^t(E)\ge\mathcal H^t(F)>0$.
\end{proof}

\begin{wniosek}[Energetyczna charakteryzacja wymiaru zwartego zbioru]
\label{cor:energy-dimension-26}
Dla niepustego zwartego $K\subseteq\mathbb R^d$ zachodzi
\[
 \dim_H K=\sup\{t>0:\text{ istnieje miara probabilistyczna }
 \mu\text{ na }K,\ I_t(\mu)<\infty\},
\]
gdzie supremum pustego zbioru wynosi zero.
\end{wniosek}
\begin{proof}
Nierówność $\ge$ wynika z poprzedniego twierdzenia.
Jeżeli $0<t<\dim_H K$, wybieramy $s$ takie, że
$t<s<\dim_H K$. Wówczas $\mathcal H^s(K)=+\infty$,
więc twierdzenie Frostmana daje miarę probabilistyczną
na $K$ spełniającą warunek wzrostu masy kul z wykładnikiem
$s$. Twierdzenie~\ref{thm:frostman-energy-26}
zapewnia $I_t(\mu)<\infty$. Przechodzimy z $t$ do
$\dim_H K$ z dołu. Gdy $\dim_H K=0$, pierwsza
nierówność i konwencja o supremum kończą dowód.
\end{proof}

\begin{uwaga}[Dalsza perspektywa]
W rozdziałach dotyczących geometrycznej teorii miary
pojawią się miary prostowalne, wzór na pole powierzchni
i wzór koobszaru. Miara Hausdorffa nadaje im naturalną
skalę. Związek energii z pojemnością stanowi natomiast
punkt wyjścia do teorii potencjału, której pełne
rozwinięcie pozostawiamy do osobnego opracowania.
\end{uwaga}

% ROZDZIAŁ 27 -- ujęcie miarowe; definicje geometryczne w osobnym skrypcie.
\chapter[Miary na rozmaitościach]{Miary na rozmaitościach: ujęcie miarowo-teoretyczne}
\label{ch:miary-rozmaitosci-27}

\section{Zakres rozdziału i ustalone obiekty}
W opracowaniu \emph{Geometria i topologia różniczkowa}, w rozdziale
\emph{Tensor metryczny, długość i miara}, skonstruowano już
miarę riemannowską, wyprowadzono jej niezależność od mapy,
wprowadzono metrykę indukowaną i obliczono całki po gładkich
podrozmaitościach. W rozdziale o formach udowodniono twierdzenie
Stokesa. Nie powtarzamy tych dowodów ani przykładów
sfery i paraboloidy. Zamiast tego badamy miarę riemannowską
jako obiekt teorii miary: regularność, zbiory zerowe,
gęstości Radona--Nikodyma, transport miar i produkty.

Przez $(M^n,g)$ rozumiemy w tym rozdziale niepustą, gładką,
Hausdorffowską rozmaitość bez brzegu, spełniającą drugi aksjomat przeliczalności,
wymiaru $n\ge1$, wyposażoną w gładką metrykę riemannowską.
Nie wymagamy orientowalności. Symbol $\mathcal B(M)$ oznacza
$\sigma$-algebrę borelowską, a $\mu_g$ oznacza miarę
borelowską określoną lokalnie przez
\begin{equation}\label{eq:g-density-27}
  \mu_g(A\cap U)=\int_{x(A\cap U)}
  \sqrt{\det G_x(u)}\,\mathrm d\lambda_n(u),
\end{equation}
gdzie $(U,x)$ jest mapą, $A\in\mathcal B(M)$,
a $G_x(u)=(g_{ij}(x^{-1}(u)))$ jest macierzą metryki
w bazie współrzędnościowej w punkcie $x^{-1}(u)$.
Jest to przypomnienie definicji, nie nowa konstrukcja.
Jej uzupełnienie będziemy oznaczać przez
$(\overline{\mathcal B(M)}^{\,\mu_g},\overline\mu_g)$.

\section{Regularność i $\sigma$-skończoność}
\begin{twierdzenie}[Miara objętościowa jest miarą Radona]
\label{thm:riemann-radon-27}
Niech $(M^n,g)$ będzie gładką rozmaitością riemannowską
bez brzegu, Hausdorffa, spełniającą drugi aksjomat przeliczalności,
wymiaru $n\ge1$. Wówczas
$\mu_g$ jest lokalnie skończoną, $\sigma$-skończoną
miarą Radona na $\mathcal B(M)$. W szczególności
$\mu_g(K)<\infty$ dla każdego zwartego $K\subseteq M$.
\end{twierdzenie}
\begin{proof}
Weźmy punkt $p\in M$ i mapę $(U,x)$ wokół $p$.
Wybieramy kulę euklidesową $B$ tak, aby
$\overline B\subset x(U)$, a jej przeciwobraz $K=x^{-1}(\overline B)$
był zwartym otoczeniem punktu $p$. Gęstość
$\rho(u)=\sqrt{\det G_x(u)}$ jest dodatnia
i ciągła, zatem na $\overline B$ jest ograniczona
z góry i z dołu przez dodatnie stałe. Z
\eqref{eq:g-density-27} wynika $\mu_g(K)<\infty$.
Dzięki zwartości dowolny zwarty $L\subset M$ pokrywamy
skończenie wieloma takimi otoczeniami, skąd
$\mu_g(L)<\infty$.

Rozmaitość $M$ jest lokalnie zwarta i ma przeliczalną
bazę, więc jest $\sigma$-zwarta: wybieramy
przeliczalnie wiele otoczeń o zwartych domknięciach,
które ją pokrywają. W konsekwencji $\mu_g$ jest
$\sigma$-skończona. Ponadto $M$ jest przestrzenią
metryzowalną i lokalnie zwartą. Stosujemy
udowodnione wcześniej twierdzenie o regularności
lokalnie skończonych miar borelowskich na takich
przestrzeniach (Rozdział~\ref{ch:radon-measures-20}),
aby otrzymać regularność wewnętrzną przez zbiory
zwarte i zewnętrzną przez zbiory otwarte.
\end{proof}

\begin{uwaga}[Miara borelowska a jej uzupełnienie]
Nie należy utożsamiać $\mathcal B(M)$ z jej
uzupełnieniem. Podzbiory borelowskiego zbioru
miary zero mogą nie być borelowskie, ale należą do
$\overline{\mathcal B(M)}^{\,\mu_g}$.
Definicja miary Radona dotyczy w tym rozdziale
miary borelowskiej; do całkowania funkcji
mierzalnych względem jej uzupełnienia używamy
standardowego rozszerzenia całki.
\end{uwaga}

\section{Zbiory o zerowej objętości}
\begin{twierdzenie}[Lokalne kryterium zbiorów zerowych]
\label{thm:null-local-27}
Niech $(M^n,g)$ będzie gładką rozmaitością riemannowską
bez brzegu, Hausdorffa, spełniającą drugi aksjomat przeliczalności,
wymiaru $n\ge1$. Niech $A\in\mathcal B(M)$, a
$\{(U_j,x_j)\}_{j\ge1}$ będzie przeliczalnym atlasem
pokrywającym $M$. Wówczas
\[
 \mu_g(A)=0\quad\Longleftrightarrow\quad
 \lambda_n\bigl(x_j(A\cap U_j)\bigr)=0
 \quad\text{dla każdego }j\ge1.
\]
To samo kryterium działa dla zbiorów mierzalnych
w uzupełnionej $\sigma$-algebrze, gdy ich obrazy
rozumiemy jako zbiory mierzalne w sensie Lebesgue'a.
\end{twierdzenie}
\begin{proof}
Na każdej mapie gęstość $\rho_j$ jest dodatnia
i ciągła. Dla zbioru borelowskiego $E\subseteq x_j(U_j)$
mamy $\int_E\rho_j\,\mathrm d\lambda_n=0$
wtedy i tylko wtedy, gdy $\lambda_n(E)=0$:
wystarcza wyczerpać $x_j(U_j)$ zbiorami zwartymi,
na których $\rho_j$ ma dodatnie minimum.
Z \eqref{eq:g-density-27} otrzymujemy kryterium
lokalne, a przeliczalność atlasu pozwala
zsumować je na całym $M$. Jeżeli $A$ należy do uzupełnienia, istnieją zbiory borelowskie
$B,Z\subseteq M$ takie, że $A\triangle B\subseteq Z$ i $\mu_g(Z)=0$.
W każdej mapie obraz różnicy symetrycznej jest zawarty w
$x_j(Z\cap U_j)$, który na mocy części borelowskiej ma miarę
Lebesgue'a zero. Zupełność miary Lebesgue'a zapewnia zatem
mierzalność $x_j(A\cap U_j)$ oraz równość jego miary z miarą
$x_j(B\cap U_j)$. Ponadto $\overline\mu_g(A)=\mu_g(B)$.
Stosując część borelowską do $B$, otrzymujemy obie implikacje
dla uzupełnienia.
\end{proof}

\begin{wniosek}[Obrazy immersji niższego wymiaru]
\label{cor:submanifold-null-27}
Niech $(M^n,g)$ będzie gładką rozmaitością riemannowską
bez brzegu, Hausdorffa, spełniającą drugi aksjomat przeliczalności,
wymiaru $n\ge1$. Niech $N^k$ będzie gładką rozmaitością bez brzegu,
Hausdorffa, spełniającą drugi aksjomat przeliczalności, $0\le k<n$,
a $F:N\to M$ immersją klasy $C^1$. Wówczas $S=F(N)$
jest zbiorem borelowskim w $M$ i $\mu_g(S)=0$.
W szczególności teza zachodzi dla każdej osadzonej
podrozmaitości gładkiej wymiaru $k<n$.
\end{wniosek}
\begin{proof}
Dla $k=0$ przestrzeń $N$ jest dyskretna i przeliczalna,
więc jej obraz jest przeliczalną sumą punktów o zerowej
$\mu_g$-mierze. Niech zatem $k\ge1$.
Ustalmy $p\in N$ i gładkie mapy źródła i celu wokół $p$
oraz $F(p)$. Po ograniczeniu dziedziny mapy źródła tak, aby jej obraz
przez $F$ mieścił się w dziedzinie mapy celu, otrzymujemy
współrzędnościowe przedstawienie $F$, którego macierz
pochodnej w punkcie odpowiadającym $p$ ma rząd $k$. Po ewentualnej permutacji współrzędnych celu
pierwsze $k$ składowych tworzy odwzorowanie $P$ o odwracalnej
pochodnej w punkcie odpowiadającym $p$.
Z twierdzenia o odwzorowaniu odwrotnym, po ograniczeniu
dziedziny do otoczenia $V_p$, odwzorowanie $P$ jest
dyfeomorfizmem klasy $C^1$ na otwarty zbiór $W\subseteq\mathbb R^k$.
Pozostałe składowe, złożone z $P^{-1}$, tworzą funkcję
ciągłą $\psi:W\to\mathbb R^{n-k}$. Obraz $F(V_p)$
w wybranej gładkiej mapie celu jest więc wykresem $\psi$.
Wykres jest domknięty względem $W\times\mathbb R^{n-k}$,
zatem borelowski w $\mathbb R^n$. Każdy jego przekrój
pionowy nad punktem $W$ jest singletonem, a poza $W$ pusty.
Ponieważ $n-k\ge1$, twierdzenie Tonellego daje zerową
$n$-wymiarową miarę Lebesgue'a wykresu. Wzór
\eqref{eq:g-density-27} daje $\mu_g(F(V_p))=0$.
Takie uzasadnienie nie wymaga użycia map klasy $C^1$
jako map gładkiego atlasu celu.
Ponieważ $N$ ma przeliczalną bazę, z pokrycia
$(V_p)_{p\in N}$ wybieramy podpokrycie
przeliczalne $(V_j)$. Zatem
$S=\bigcup_jF(V_j)$. Każdy $F(V_j)$ jest borelowski
w dziedzinie gładkiej mapy celu, a więc także w $M$,
ponieważ dziedzina mapy jest otwarta. Stąd $S$ jest
borelowski i z przeliczalnej podaddytywności
$\mu_g(S)\le\sum_j\mu_g(F(V_j))=0$.
Argument nie zakłada, że \emph{cały} obraz immersji
jest lokalnie osadzoną podrozmaitością w otoczeniu
każdego punktu obrazu: przy samoprzecięciach lub
nagromadzeniach może to być nieprawdą.
\end{proof}

\begin{uwaga}[Dwie różne miary na tym samym zbiorze]
Wniosek nie oznacza, że długość krzywej lub pole
powierzchni znika. Dla osadzonej $k$-wymiarowej
podrozmaitości $S$, miara $\mu_{g|_S}$ może być
dodatnia, choć $\mu_g(S)=0$, gdy $k<n$.
To porównanie prowadzi do miar Hausdorffa;
identyfikację $\mu_{g|_S}$ z odpowiednio
unormowaną $\mathcal H^k$ uzasadnimy
w Rozdziale~\ref{ch:area-coarea-28}.
\end{uwaga}

\section{Miary z gęstością i rozkład Lebesgue'a}
\begin{twierdzenie}[Miary bezwzględnie ciągłe względem objętości]
\label{thm:rn-manifold-27}
Niech $(M^n,g)$ będzie gładką rozmaitością riemannowską
bez brzegu, Hausdorffa, spełniającą drugi aksjomat przeliczalności,
wymiaru $n\ge1$. Niech
$\nu$ będzie dodatnią, $\sigma$-skończoną
miarą na $\mathcal B(M)$, przy czym
$\nu\ll\mu_g$. Wtedy istnieje nieujemna funkcja
borelowska $h$, jednoznaczna $\mu_g$-prawie
wszędzie, taka że dla każdego $A\in\mathcal B(M)$
\[
 \nu(A)=\int_A h(p)\,\mathrm d\mu_g(p).
\]
Jeżeli dodatkowo $\nu$ jest lokalnie skończona,
to $h\in L^1_{\mathrm{loc}}(M,\mu_g)$,
tj. $\int_K h\,\mathrm d\mu_g<\infty$
dla każdego zwartego $K\subseteq M$.
\end{twierdzenie}
\begin{proof}
Miara $\mu_g$ jest $\sigma$-skończona na mocy
Twierdzenia~\ref{thm:riemann-radon-27}.
Stosujemy twierdzenie Radona--Nikodyma
(Rozdział~\ref{ch:radon-nikodym-18})
do pary $(\nu,\mu_g)$. Otrzymujemy nieujemną funkcję $h$ mierzalną względem $\mathcal B(M)$,
ponieważ twierdzenie stosujemy na tej właśnie $\sigma$-algebrze. Jeżeli $\nu$ jest lokalnie
skończona, każdy zwarty $K$ można przykryć
skończenie wieloma otoczeniami o skończonej
$\nu$-mierze, więc $\nu(K)<\infty$.
Wstawienie $A=K$ daje lokalną całkowalność.
\end{proof}

\begin{przyklad}[Gęstość w mapie]
Niech $(U,x)$ będzie mapą rozmaitości riemannowskiej
$(M^n,g)$, $V=x(U)\subseteq\mathbb R^n$, a $\nu$ dodatnią
$\sigma$-skończoną miarą borelowską na $M$.
Załóżmy, że miara obrazowa $x_*(\nu|_U)$ ma gęstość
$q:V\to[0,+\infty]$ względem $\lambda_n|_V$.
Wtedy gęstość $\nu|_U$ względem $\mu_g|_U$ wynosi
$h(p)=q(x(p))/\sqrt{\det G_x(x(p))}$ dla
$\mu_g$-prawie wszystkich $p\in U$. Istotnie, dla każdego
borelowskiego $A\subseteq U$ z lokalnego wzoru na miarę
riemannowską i definicji miary obrazowej otrzymujemy
\[
 \int_A h\,\mathrm d\mu_g
 =\int_{x(A)}q(a)\,\mathrm d\lambda_n(a)=\nu(A).
\]
Dodatniość gęstości riemannowskiej uzasadnia dzielenie.
Wyrażenie $q$ zmienia się po zmianie mapy;
$h$ jest funkcją skalarną na rozmaitości.
\end{przyklad}

\begin{twierdzenie}[Rozkład na część objętościową i osobliwą]
\label{thm:leb-manifold-27}
Niech $(M^n,g)$ będzie gładką rozmaitością riemannowską
bez brzegu, Hausdorffa, spełniającą drugi aksjomat przeliczalności,
wymiaru $n\ge1$. Jeżeli $\nu$ jest dodatnią, $\sigma$-skończoną
miarą na $\mathcal B(M)$, to ma jednoznaczny rozkład
$\nu=h\mu_g+\nu_{\mathrm s}$, gdzie
$h\ge0$ jest borelowska, a
$\nu_{\mathrm s}\perp\mu_g$. Jednoznaczność $h$ rozumiemy
z dokładnością do równości $\mu_g$-prawie wszędzie.
\end{twierdzenie}
\begin{proof}
Jest to zastosowanie twierdzenia o rozkładzie
Lebesgue'a (Rozdział~\ref{ch:ac-singular-17})
i twierdzenia Radona--Nikodyma do części
bezwzględnie ciągłej. Jednoznaczność obu składników
wynika z tych twierdzeń.
\end{proof}

\section{Przenoszenie objętości przez dyfeomorfizm}
\begin{definicja}[Jakobian objętościowy]
\label{def:volume-jacobian-27}
Niech $(M^n,g)$ i $(N^n,h)$ będą rozmaitościami
riemannowskimi wymiaru $n\ge1$, gładkimi, bez brzegu,
Hausdorffa i spełniającymi drugi aksjomat przeliczalności,
a $F:M\to N$ dyfeomorfizmem klasy $C^1$.
Dla $p\in M$ wybieramy ortonormalne bazy
w $T_pM$ i $T_{F(p)}N$, a przez $A_p$ oznaczamy
macierz $\mathrm dF_p$ w tych bazach. Definiujemy
$J_{g,h}F(p)=|\det A_p|$. Zmiana każdej z baz mnoży macierz z odpowiedniej strony
przez macierz ortogonalną, której wyznacznik ma moduł jeden.
Wartość nie zależy więc od wyboru baz.
Niech $u,v$ będą mapami wokół $p,F(p)$, a $G=G_u$ i $H=H_v$
macierzami metryk w tych współrzędnych. Wówczas
\[
 J_{g,h}F(p)=|\det D(v\circ F\circ u^{-1})(u(p))|
 \frac{\sqrt{\det H(v(F(p)))}}{\sqrt{\det G(u(p))}}.
\]
\end{definicja}

\begin{twierdzenie}[Zamiana zmiennych na rozmaitościach]
\label{thm:manifold-change-27}
Niech $(M^n,g)$ i $(N^n,h)$ będą gładkimi rozmaitościami
riemannowskimi bez brzegu, Hausdorffa, spełniającymi drugi
aksjomat przeliczalności, wymiaru $n\ge1$.
Niech $F:M\to N$ będzie dyfeomorfizmem klasy $C^1$,
a $J_{g,h}F$ jego jakobianem objętościowym.
Dla każdej funkcji borelowskiej $f:N\to[0,+\infty]$ zachodzi
\begin{equation}\label{eq:manifold-change-27}
 \int_N f(q)\,\mathrm d\mu_h(q)=
 \int_M f(F(p))J_{g,h}F(p)\,\mathrm d\mu_g(p).
\end{equation}
Ta sama równość zachodzi dla funkcji ze znakiem,
gdy $f\in L^1(N,\mu_h)$.
W szczególności $(F_*\mu_g)\ll\mu_h$ i
\[
 \frac{\mathrm d(F_*\mu_g)}{\mathrm d\mu_h}(q)
 =\frac{1}{J_{g,h}F(F^{-1}(q))}
 \quad\text{dla }\mu_h\text{-prawie wszystkich }q.
\]
\end{twierdzenie}
\begin{proof}
Uzasadnijmy najpierw wzór współrzędnościowy z definicji
jakobianu. Ustalmy $p\in M$ i mapy $u,v$ wokół $p,F(p)$,
ograniczając w razie potrzeby dziedzinę $u$ tak, aby jej obraz
przez $F$ mieścił się w dziedzinie $v$.
Niech $T=v\circ F\circ u^{-1}$, $D=DT(u(p))$,
$G=G_u(u(p))$ oraz $H=H_v(v(F(p)))$.
Niech kolumny macierzy $P$ i $Q$ będą współrzędnymi
wektorów wybranych baz ortonormalnych względem baz
współrzędnościowych źródła i celu. Wtedy
$P^{\mathsf T}GP=I$ oraz $Q^{\mathsf T}HQ=I$, a macierz
różniczki w bazach ortonormalnych to $A_p=Q^{-1}DP$.
Z równości wyznaczników dostajemy
$|\det P|=(\det G)^{-1/2}$ i
$|\det Q|=(\det H)^{-1/2}$, skąd
\[
 |\det A_p|=|\det D|\frac{\sqrt{\det H}}{\sqrt{\det G}}.
\]
W szczególności $J_{g,h}F$ jest ciągły i ściśle dodatni.
Dla różniczki odwrotnej macierz w tych samych bazach,
zamienionych rolami, to $A_p^{-1}$, zatem
$J_{h,g}(F^{-1})(F(p))=1/J_{g,h}F(p)$.

Najpierw załóżmy, że $f$ znika poza otwartym zbiorem $V$
zawartym w dziedzinie mapy celu, a $F^{-1}(V)$
mieści się w dziedzinie mapy źródła. Z definicji
$\mu_g,\mu_h$ oraz twierdzenia o zamianie
zmiennych w $\mathbb R^n$ otrzymujemy
\eqref{eq:manifold-change-27}, po skróceniu
wyznaczników macierzy metryk: jeśli $T=v\circ F\circ u^{-1}$,
to iloczyn $J_{g,h}F(u^{-1}(a))\sqrt{\det G(a)}$
jest równy $|\det DT(a)|\sqrt{\det H(T(a))}$.
Jest to dokładnie gęstość występująca po euklidesowej zamianie zmiennych.

W ogólności dla każdego punktu $q\in N$ wybieramy otwarte
$V$ zawarte w dziedzinie mapy celu tak, aby $F^{-1}(V)$
mieścił się w dziedzinie mapy źródła. Istnienie takiego $V$
wynika z ciągłości $F^{-1}$. Z drugiego aksjomatu przeliczalności
wybieramy przeliczalne podpokrycie $(V_j)$. Tworzymy rozłączny
podział borelowski $B_1=V_1$, $B_j=V_j\setminus\bigcup_{i<j}V_i$.
Dla $f\chi_{B_j}$ stosujemy argument lokalny na $V_j$,
a następnie sumujemy równości, korzystając z twierdzenia
o zbieżności monotonicznej.

Aby obliczyć miarę obrazową, dla $A\in\mathcal B(N)$
stosujemy równanie do $F^{-1}$ i funkcji $\chi_{F^{-1}(A)}$ na $M$.
Ponieważ
$J_{h,g}(F^{-1})(q)=[J_{g,h}F(F^{-1}q)]^{-1}$,
dostajemy podaną pochodną Radona--Nikodyma.
Dla funkcji całkowalnych stosujemy równość
oddzielnie do części dodatniej i ujemnej.
\end{proof}

\begin{przyklad}[Konforemna zmiana metryki]
Niech $\varphi\in C^\infty(M;\mathbb R)$ i
$g'=e^{2\varphi}g$. Macierze spełniają
$G'=e^{2\varphi}G$, więc
$\det G'=e^{2n\varphi}\det G$. W konsekwencji
\[
 \mathrm d\mu_{g'}(p)=e^{n\varphi(p)}\,\mathrm d\mu_g(p).
\]
Nie oznacza to, że krzywizna obu metryk jest
taka sama: wzór dotyczy tylko miary objętościowej.
\end{przyklad}

\section{Miara na produkcie rozmaitości}
\begin{twierdzenie}[Miara produktu riemannowskiego]
\label{thm:product-riemann-27}
Niech $(M^n,g)$ i $(N^m,h)$ będą gładkimi rozmaitościami
riemannowskimi bez brzegu, Hausdorffa, spełniającymi drugi
aksjomat przeliczalności, wymiarów odpowiednio $n,m\ge1$. Na $M\times N$ rozważmy metrykę
$g\oplus h$. Wtedy na $\mathcal B(M\times N)$
zachodzi $\mu_{g\oplus h}=\mu_g\otimes\mu_h$.
Dla każdej funkcji borelowskiej $f:M\times N\to[0,+\infty]$ mamy
\[
 \int_{M\times N} f(p,q)\,\mathrm d\mu_{g\oplus h}(p,q)
 =\int_M\!\left(\int_N f(p,q)\,\mathrm d\mu_h(q)\right)
   \mathrm d\mu_g(p).
\]
\end{twierdzenie}
\begin{proof}
Ponieważ obie rozmaitości spełniają drugi aksjomat przeliczalności,
$\mathcal B(M\times N)=\mathcal B(M)\otimes\mathcal B(N)$:
każdy zbiór otwarty produktu jest przeliczalną sumą otwartych
prostokątów z baz obu czynników, a prostokąty borelowskie są
borelowskie dzięki ciągłości rzutów.
W mapach produktowych macierz metryki jest
blokowo diagonalna $\operatorname{diag}(G,H)$.
Stąd $\sqrt{\det(G\oplus H)}=
\sqrt{\det G}\sqrt{\det H}$. Dla zbiorów borelowskich $A\subseteq U$ i $B\subseteq V$,
gdzie $(U,u)$ i $(V,v)$ są mapami czynników, wzór na
gęstość i twierdzenie Tonellego dają
\[
 \mu_{g\oplus h}(A\times B)=
 \int_{u(A)}\int_{v(B)}
 \sqrt{\det G_u(a)}\sqrt{\det H_v(b)}
 \,\mathrm d\lambda_m(b)\,\mathrm d\lambda_n(a)
 =\mu_g(A)\mu_h(B).
\]
W iloczynie wartości miar stosujemy konwencję $0\cdot\infty=0$.
Obie miary ograniczone do $U\times V$ są więc równe
na prostokątach, a przez $\sigma$-skończoność czynników
i jednoznaczność miary produktowej są równe na całej
$\mathcal B(U\times V)$. Wybieramy przeliczalne atlasy
obu czynników i numerujemy wszystkie ich prostokąty
jako $(R_j)_{j\ge1}$. Dla dowolnego zbioru borelowskiego
$E\subseteq M\times N$ zbiory
$E_j=E\cap(R_j\setminus\bigcup_{i<j}R_i)$ są rozłączne,
borelowskie i pokrywają $E$. Równość miar na każdym
$E_j\subseteq R_j$ oraz przeliczalna addytywność
dają równość miar na $E$.
Ostatnia równość jest twierdzeniem Tonellego.
\end{proof}

\section{Od gładkich podrozmaitości do zbiorów prostowalnych}
W geometrii rozważaliśmy miarę indukowaną
na \emph{gładkich} podrozmaitościach. Obraz
odwzorowania Lipschitzowskiego nie musi być
gładką podrozmaitością, lecz nadal może
posiadać sensowną miarę $k$-wymiarową.
W następnym rozdziale pokażemy, jak
$\mathcal H^k$ mierzy takie obrazy i dlaczego
przy odwzorowaniach nieiniektywnych konieczny
jest czynnik krotności. Wprowadzone tam wzory
są ściśle miarowe: nie potrzebują globalnej
orientacji ani globalnych parametryzacji.

% ROZDZIAŁ 28 -- wzory miarowo-geometryczne, nie powtórzenie całek po podrozmaitościach.
\chapter{Wzór na pole powierzchni, krotności i koobszar}
\label{ch:area-coarea-28}

\section{Od wzoru parametrycznego do wzoru miarowego}
W skrypcie \emph{Geometria i topologia różniczkowa}
udowodniono wzór na całkę po kawałku gładkiej
podrozmaitości opisanej parametryzacją, która
pokrywa ten kawałek dokładnie raz. Tutaj rozważamy
\emph{odwzorowania nieiniektywne} oraz ich włókna.
Nowością są krotności obrazów, całkowanie po
przeciwobrazach punktów i formuły dla
odwzorowań Lipschitzowskich.

W tym rozdziale $\lambda_k$ oznacza miarę Lebesgue'a
na $\mathbb R^k$, a $\mathcal H^k$ -- miarę
Hausdorffa z Rozdziału~\ref{ch:hausdorff-26}.
Stosujemy konwencję z
Uwagi~\ref{rem:hausdorff-convention-26}: dla wymiaru
całkowitego $k\ge1$ wybieramy $c_k=\kappa_k^{-1}$,
tak aby $\mathcal H^k|_{\mathbb R^k}=\lambda_k$;
$\mathcal H^0$ jest miarą liczącą.
Dzięki temu poniższe wzory nie zakładają
nieudowodnionej nierówności izodiametrycznej.
Przy innej stałej normalizującej ich prawe strony
należy odpowiednio przeskalować.

\begin{definicja}[Jakobian wymiaru $k$]
\label{def:k-jacobian-28}
Niech $1\le k\le N$. Dla dowolnego odwzorowania
liniowego $A:\mathbb R^k\to\mathbb R^N$ definiujemy
\[
 J_k(A)=\sqrt{\det(A^{\mathsf T}A)}.
\]
Jeżeli $U\subseteq\mathbb R^k$ jest otwarty,
a $F:U\to\mathbb R^N$ jest różniczkowalne
w punkcie $x$, to kładziemy $J_kF(x)=J_k(DF(x))$.
W szczególności dotyczy to wszystkich punktów odwzorowania
$C^1$ oraz prawie wszystkich punktów odwzorowania
lokalnie Lipschitzowskiego. Macierz $A^{\mathsf T}A$ jest symetryczna i nieujemnie
określona, ponieważ $v^{\mathsf T}A^{\mathsf T}Av=|Av|^2$.
Ponadto $\ker(A^{\mathsf T}A)=\ker A$: jedna inkluzja
wynika z mnożenia przez $A^{\mathsf T}$, a druga z
poprzedniej równości. Zatem $J_k(A)=0$ wtedy i tylko
wtedy, gdy $\operatorname{rank}A<k$. Dla $k=N$
mamy $\det(A^{\mathsf T}A)=(\det A)^2$, więc
$J_k(A)=|\det A|$.
\end{definicja}

\begin{definicja}[Krotność odwzorowania]
\label{def:multiplicity-28}
Dla dowolnego odwzorowania $F:X\to Y$
i zbioru $E\subseteq X$ oraz dla każdego $y\in Y$ definiujemy
\[
 N(F,E,y)=\#\{x\in E:F(x)=y\}\in\mathbb N_0\cup\{+\infty\}.
\]
Symbol $\#$ oznacza liczbę elementów zbioru
(z wartością $+\infty$ dla zbiorów nieskończonych).
Pojawia się tu również wartość $+\infty$,
gdy włókno ma nieskończenie wiele punktów.
\end{definicja}

\section{Miara Hausdorffa gładkiego wykresu}
\begin{lemat}[Małe zniekształcenie wykresu]
\label{lem:graph-hausdorff-28}
Niech $1\le k\le N$, $V\subseteq\mathbb R^k$ będzie otwarty,
a $h:V\to\mathbb R^{N-k}$ funkcją klasy $C^1$.
Niech $\Gamma(z)=(z,h(z))$ oraz
$\nu(A)=\mathcal H^k(\Gamma(A))$
dla borelowskich $A\subseteq V$. Wtedy $\nu$
jest miarą Radona na $V$ i dla każdego
$z_0\in V$ zachodzi
\[
 \lim_{r\downarrow0}
 \frac{\nu(B(z_0,r))}{\lambda_k(B(z_0,r))}
 =\sqrt{\det\bigl(I_k+(Dh(z_0))^{\mathsf T}Dh(z_0)\bigr)}.
\]
\end{lemat}
\begin{proof}
Na każdej kuli domkniętej $K\subset V$ wykres $\Gamma$
jest Lipschitzowski, ponieważ $Dh$ jest ograniczona
na $K$. Zatem, na mocy
Twierdzenia~\ref{thm:hausdorff-lipschitz-26},
$\mathcal H^k(\Gamma(K))<\infty$, gdyż
$\mathcal H^k(K)=\lambda_k(K)$ jest skończona.
Z drugiej strony rzut ortogonalny $\pi$ na
pierwszych $k$ współrzędnych jest odwrotnością
$\Gamma$ na obrazie i jest $1$-Lipschitzowski.
Wykres $\Gamma(V)$ jest domknięty względem
$V\times\mathbb R^{N-k}$, więc jest borelowski w $\mathbb R^N$.
Dla borelowskiego $A\subseteq V$ mamy
$\Gamma(A)=\Gamma(V)\cap\pi^{-1}(A)$, zatem
$\Gamma(A)$ jest borelowski w $\mathbb R^N$.
Iniektywność $\Gamma$ zapewnia, że obrazy parami
rozłącznych zbiorów są parami rozłączne. Przeliczalna
addytywność $\mathcal H^k$ na zbiorach borelowskich
daje więc przeliczalną addytywność $\nu$.
Poprzednie oszacowanie zapewnia jej lokalną skończoność. Stosujemy twierdzenie o
regularności miar lokalnie skończonych na lokalnie
zwartych ośrodkowych przestrzeniach metrycznych
z Rozdziału~\ref{ch:radon-measures-20}.
Otrzymujemy, że $\nu$ jest miarą Radona,
bez zakładania z góry regularności miary
Hausdorffa na całym wykresie.

Ustalmy $z_0\in V$ i afiniczne odwzorowanie
$L(z)=(z,h(z_0)+Dh(z_0)(z-z_0))$.
Wybieramy $r>0$ tak, aby $\overline B(z_0,r)\subset V$,
i kładziemy
$\varepsilon_r=\sup_{w\in\overline B(z_0,r)}
\|Dh(w)-Dh(z_0)\|$. Z ciągłości $Dh$ wynika
$\varepsilon_r\to0$ dla $r\downarrow0$.
Dla $u,v\in B(z_0,r)$ odcinek między nimi
jest zawarty w tej kuli, więc podstawowe
twierdzenie rachunku całkowego daje
\[
 h(u)-h(v)-Dh(z_0)(u-v)
 =\int_{[0,1]}(Dh(v+t(u-v))-Dh(z_0))(u-v)
 \,\mathrm d\lambda_1(t).
\]
Norma prawej strony nie przekracza
$\varepsilon_r|u-v|$.
Ponieważ $|L(u)-L(v)|\ge|u-v|$,
odwzorowanie $\Gamma\circ L^{-1}$ między
odpowiednimi fragmentami wykresów jest
dwulipschitzowskie: dla $u,v$ z rozpatrywanej kuli
\[
 (1-\varepsilon_r)|L(u)-L(v)|
 \le|\Gamma(u)-\Gamma(v)|
 \le(1+\varepsilon_r)|L(u)-L(v)|.
\]
Dla dostatecznie małego $r$ mamy $\varepsilon_r<1$.
Z oszacowania miary Hausdorffa przez stałą
Lipschitza otrzymujemy
\[
 (1-\varepsilon_r)^k\mathcal H^k(L(B(z_0,r)))
 \le\nu(B(z_0,r))\le
 (1+\varepsilon_r)^k\mathcal H^k(L(B(z_0,r))).
\]
Obraz afinicznego odwzorowania $L$ jest
$k$-wymiarową płaszczyzną afiniczną. Wybierzmy izometrię
afiniczną tej płaszczyzny na $\mathbb R^k$ i oznaczmy przez
$B$ macierz części liniowej złożenia tej izometrii z $L$.
Wówczas $B^{\mathsf T}B=(DL)^{\mathsf T}DL=
I_k+Dh(z_0)^{\mathsf T}Dh(z_0)$.
Z niezmienniczości $\mathcal H^k$ względem izometrii,
kalibracji $\mathcal H^k=\lambda_k$ w $\mathbb R^k$
i liniowej zamiany zmiennych wynika
$\mathcal H^k(L(A))=|\det B|\lambda_k(A)$
dla borelowskiego $A\subseteq\mathbb R^k$.
Ponieważ $|\det B|=\sqrt{\det(B^{\mathsf T}B)}$,
jest to dokładnie $J_k(DL)\lambda_k(A)$.
Dzielimy przez $\lambda_k(B(z_0,r))$ i
przechodzimy do granicy.
\end{proof}

\begin{twierdzenie}[Miara na gładkim wykresie]
\label{thm:graph-area-28}
Niech $1\le k\le N$, $V\subseteq\mathbb R^k$ będzie otwarty,
$h:V\to\mathbb R^{N-k}$ będzie klasy $C^1$ i
$\Gamma(z)=(z,h(z))$. Dla każdej funkcji borelowskiej
$\psi:\Gamma(V)\to[0,+\infty]$ mamy
\[
 \int_{\Gamma(V)}\psi(y)\,\mathrm d\mathcal H^k(y)
 =\int_V\psi(\Gamma(z))
 \sqrt{\det(I_k+Dh(z)^{\mathsf T}Dh(z))}
 \,\mathrm d\lambda_k(z).
\]
To samo zachodzi dla funkcji całkowalnych.
\end{twierdzenie}
\begin{proof}
Z Lematu~\ref{lem:graph-hausdorff-28} wynika, że miara
$\nu(A)=\mathcal H^k(\Gamma(A))$ na zbiorach borelowskich
$A\subseteq V$ jest lokalnie skończona.
Pokrywamy $V$ przeliczalnie wieloma kulami o domknięciach
zawartych w $V$. Na każdej takiej kuli $\Gamma$ jest
Lipschitzowska. Jeśli $\lambda_k(A)=0$, to kalibracja
miary Hausdorffa w $\mathbb R^k$ i
Twierdzenie~\ref{thm:hausdorff-lipschitz-26} dają
zerową miarę obrazu przecięcia $A$ z każdą kulą.
Podaddytywność daje $\mathcal H^k(\Gamma(A))=0$,
czyli $\nu\ll\lambda_k$. Obie miary są
$\sigma$-skończone: każda z kul ma skończoną
miarę zarówno względem $\lambda_k$, jak i $\nu$.
Twierdzenie~\ref{thm:RN-sigfin-18} daje więc
gęstość $a=\mathrm d\nu/\mathrm d\lambda_k$.
Jest ona lokalnie całkowalna na $V$, gdyż
$\int_K a\,\mathrm d\lambda_k=\nu(K)<\infty$
dla każdego zwartego $K\subseteq V$.
Stosujemy Twierdzenie~\ref{thm:lebesgue-diff-25}
o różniczkowaniu całek lokalnie: na każdej kuli
o domknięciu zawartym w $V$ ograniczamy $a$ do
tej kuli i przedłużamy przez zero do funkcji
całkowalnej na $\mathbb R^k$. W każdym punkcie
wewnętrznym dostatecznie małe kule nie wychodzą
poza wybraną kulę. Przeliczalne pokrycie $V$
takimi kulami daje dla prawie każdego $z\in V$
\[
 a(z)=\lim_{r\downarrow0}
 \frac{\int_{B(z,r)}a\,\mathrm d\lambda_k}
 {\lambda_k(B(z,r))}
 =\lim_{r\downarrow0}\frac{\nu(B(z,r))}
 {\lambda_k(B(z,r))}=J_k\Gamma(z),
\]
gdzie ostatnia równość wynika z
Lematu~\ref{lem:graph-hausdorff-28}.
Tak otrzymujemy $\nu=J_k\Gamma\,\lambda_k$.
Dla borelowskiego $D\subseteq\Gamma(V)$,
z definicji $\nu$ mamy
$\nu(\Gamma^{-1}(D))=\mathcal H^k(D)$.
Stąd wzór na całki dla $\psi=\chi_D$, następnie
dla nieujemnych funkcji prostych przez liniowość,
a dla dowolnej nieujemnej borelowskiej $\psi$
przez rosnące przybliżenie funkcjami prostymi
i twierdzenie o zbieżności monotonicznej.
Dla funkcji całkowalnej rzeczywistej stosujemy
wzór do jej części dodatniej i ujemnej;
dla zespolonej do części rzeczywistej i urojonej.
\end{proof}

\section{Wzór na pole powierzchni dla immersji}
\begin{twierdzenie}[Wzór z krotnością dla immersji $C^1$]
\label{thm:immersion-area-28}
Niech $1\le k\le N$, $U\subseteq\mathbb R^k$
będzie otwarty, a $F:U\to\mathbb R^N$
immersją klasy $C^1$, tj.
$\operatorname{rank}DF(x)=k$ dla każdego $x\in U$.
Niech $E\subseteq U$ będzie borelowski,
a $f:E\to[0,+\infty]$ borelowska.
Wówczas funkcja
$y\mapsto\sum_{x\in E\cap F^{-1}(\{y\})}f(x)$
jest $\mathcal H^k$-mierzalna i
\begin{equation}\label{eq:area-immersion-28}
 \int_E f(x)J_kF(x)\,\mathrm d\lambda_k(x)
 =\int_{\mathbb R^N}\left(
 \sum_{x\in E\cap F^{-1}(\{y\})}f(x)\right)
 \,\mathrm d\mathcal H^k(y).
\end{equation}
Przy $f\equiv1$ na $E$ (czyli $f=\chi_E$ na $U$)
prawa strona redukuje się do
\[
 \int_{\mathbb R^N} N(F,E,y)\,\mathrm d\mathcal H^k(y).
\]
\end{twierdzenie}
\begin{proof}
Dla każdego $x\in U$ wybieramy $k$ składowych $F$,
których macierz pochodnej ma niezerowy wyznacznik w $x$.
Po permutacji współrzędnych celu oznaczmy rzut na te
składowe przez $\pi$. Twierdzenie o odwzorowaniu odwrotnym
daje otoczenie $U_x\subseteq U$, na którym
$\pi\circ F$ jest dyfeomorfizmem klasy $C^1$ na otwarty
zbiór w $\mathbb R^k$. W tych współrzędnych obraz $F(U_x)$
jest wykresem funkcji klasy $C^1$, a $F|_{U_x}$ jest osadzeniem.
Przeliczalność bazy pozwala wybrać przeliczalne
pokrycie takimi otoczeniami $U_j$. Tworzymy
rozłączny podział borelowski
$B_1=U_1$, $B_j=U_j\setminus\bigcup_{i<j}U_i$.

Na każdym $U_j$ korzystamy z opisanego wyżej wykresu
i dyfeomorfizmu $\phi_j=(\pi\circ F|_{U_j})^{-1}$
(z rzutem $\pi$ wybranym dla $U_j$).
Permutacja współrzędnych celu nie zmienia miary Hausdorffa
ani jakobianu $J_kF$.
Twierdzenie~\ref{thm:graph-area-28}, zastosowane
w tych współrzędnych, oraz klasyczny wzór
podstawienia na dziedzinie $\mathbb R^k$
dają dla borelowskich $A\subseteq B_j\cap E$
\[
 \int_A f(x)J_kF(x)\,\mathrm d\lambda_k(x)
 =\int_{F(A)}f((F|_{U_j})^{-1}(y))
 \,\mathrm d\mathcal H^k(y).
\]
Różniczka zmiany parametrów wprowadza dokładnie
ten czynnik, który przekształca macierz Grama
$F$ w macierz Grama wykresu; jest to tożsamość
$J_k(F\circ\phi)=J_kF(\phi)|\det D\phi|$
dla lokalnego dyfeomorfizmu $\phi$, wynikająca
z reguły łańcuchowej i równości
$\det((D\phi)^{\mathsf T}(DF)^{\mathsf T}DF D\phi)
=(\det D\phi)^2\det((DF)^{\mathsf T}DF)$,
gdzie $DF$ jest obliczane w punkcie $\phi(z)$.

Każdy $F(B_j\cap E)$ jest zbiorem borelowskim
w lokalnie osadzonym wykresie, a zatem również
w $\mathbb R^N$. Funkcja
$f\circ(F|_{U_j})^{-1}$ na $F(B_j\cap E)$
jest borelowska, gdyż odwrotność osadzenia jest ciągła.
Przedłużamy ją przez zero poza $F(B_j\cap E)$;
tak otrzymana funkcja jest borelowska na $\mathbb R^N$.
Sumujemy równości po $j$; po lewej używamy
rozłączności $B_j$, a po prawej twierdzenia
o zbieżności monotonicznej dla nieujemnych sum
częściowych. Nie wymaga ono $\sigma$-skończoności
miary Hausdorffa w przestrzeni otaczającej. Dla każdego $y$ suma
wkładów z kolejnych $B_j$ jest dokładnie sumą
po całym włóknie $E\cap F^{-1}(\{y\})$.
\end{proof}

\begin{przyklad}[Dlaczego trzeba liczyć krotności?]
Niech $F:(0,4\pi)\to\mathbb R^2$,
$F(t)=(\cos t,\sin t)$. Wówczas $J_1F=1$,
więc lewa strona \eqref{eq:area-immersion-28}
dla $f\equiv1$ wynosi $4\pi$.
Dla $y\in S^1\setminus\{(1,0)\}$ istnieje dokładnie
jeden $s\in(0,2\pi)$ z $F(s)=y$, a wszystkimi
przeciwobrazami w $(0,4\pi)$ są $s$ i $s+2\pi$.
Punkt $(1,0)$ ma jeden przeciwobraz $2\pi$.
Poza $S^1$ krotność wynosi zero.
Aby niezależnie obliczyć miarę okręgu, stosujemy
wzór immersji do różnowartościowego ograniczenia
$F|_{(0,2\pi)}$ z funkcją $1$. Otrzymujemy
$\mathcal H^1(S^1\setminus\{(1,0)\})=2\pi$.
Singleton ma zerową $\mathcal H^1$-miarę, więc
$\mathcal H^1(S^1)=2\pi$. Prawa strona dla
$(0,4\pi)$ wynosi zatem $2\cdot2\pi=4\pi$.
Bez krotności otrzymalibyśmy wynik dwa razy za mały.
\end{przyklad}

\section{Jakobian normalny i całkowanie po włóknach}
\begin{definicja}[Jakobian normalny]
\label{def:normal-jacobian-28}
Niech $n\ge m\ge1$. Dla dowolnego odwzorowania
liniowego $A:\mathbb R^n\to\mathbb R^m$ definiujemy
\[
 \operatorname{Jac}_m(A)=\sqrt{\det(AA^{\mathsf T})}.
\]
Dla otwartego $U\subseteq\mathbb R^n$ i odwzorowania
$F:U\to\mathbb R^m$ różniczkowalnego
w punkcie $x$ przyjmujemy
$\operatorname{Jac}_mF(x)=\operatorname{Jac}_m(DF(x))$.
Definicja obejmuje zatem odwzorowania $C^1$ i punkty
różniczkowalności odwzorowań Lipschitzowskich.
Ponieważ $w^{\mathsf T}AA^{\mathsf T}w=|A^{\mathsf T}w|^2$,
macierz $AA^{\mathsf T}$ jest nieujemnie określona,
a jej jądro jest równe $\ker A^{\mathsf T}$.
Jej wyznacznik jest dodatni dokładnie wtedy, gdy
$A$ ma rangę $m$. Stąd $\operatorname{Jac}_mF(x)=0$
dla rangi mniejszej niż $m$, a dla pełnej rangi
jest dodatni. Dla $m=1$ macierz $DF(x)$ jest
wierszem pochodnych cząstkowych, więc
$DF(x)DF(x)^{\mathsf T}=|\nabla F(x)|^2$.
Dla $n=m$ wyznacznik $AA^{\mathsf T}$ wynosi
$(\det A)^2$, zatem $\operatorname{Jac}_n(A)=|\det A|$.
\end{definicja}

\begin{lemat}[Tożsamość liniowo-algebraiczna koobszaru]
\label{lem:linear-coarea-28}
Niech $n\ge m\ge1$, $A:\mathbb R^n\to\mathbb R^m$
będzie surjektywnym odwzorowaniem liniowym, a
$G:\mathbb R^m\times\mathbb R^{n-m}\to\mathbb R^n$
odwracalnym odwzorowaniem liniowym, dla którego
$AG(y,z)=y$. Wówczas
\[
 \sqrt{\det(AA^{\mathsf T})}\,|\det DG|
 =\sqrt{\det((D_zG)^{\mathsf T}D_zG)}.
\]
\end{lemat}
\begin{proof}
Wybieramy ortonormalne rozkłady
$\mathbb R^n=(\ker A)^\perp\oplus\ker A$
i zapisujemy macierze względem tych baz.
Ograniczenie $A_0=A|_{(\ker A)^\perp}$
jest iniekcją, ponieważ jego jądro jest przecięciem
$\ker A$ z $(\ker A)^\perp$, czyli $\{0\}$.
Obie przestrzenie mają wymiar $m$, więc $A_0$
jest izomorfizmem. W tych bazach $A=(A_0\ \ 0)$.
Z warunku $AG(y,z)=y$ wynika, że pierwszy
wiersz blokowy macierzy $G$ to $(A_0^{-1}\ \ 0)$.
Macierz $G$ ma zatem postać
$\begin{psmallmatrix} A_0^{-1}&0\\B&C\end{psmallmatrix}$,
gdzie $C$ jest odwracalną macierzą kwadratową
rozmiaru $n-m$. Stąd
$|\det DG|=|\det A_0|^{-1}|\det C|$.
Ponieważ $AA^{\mathsf T}=A_0A_0^{\mathsf T}$,
lewy pierwiastek jest równy $|\det A_0|$.
Kolumny $D_zG$ leżą w $\ker A$, a ich macierz
w wybranej ortonormalnej bazie jest $C$,
więc prawy pierwiastek wynosi $|\det C|$.
Dla $n=m$ przyjmujemy wyznacznik pustej
macierzy równy $1$ i dowód pozostaje poprawny.
\end{proof}

\begin{twierdzenie}[Wzór koobszaru dla submersji]
\label{thm:coarea-submersion-28}
Niech $n\ge m\ge1$, $U\subseteq\mathbb R^n$
będzie otwarty, a $F:U\to\mathbb R^m$
submersją klasy $C^1$, tj.
$\operatorname{rank}DF(x)=m$ dla każdego $x\in U$.
Dla każdej borelowskiej funkcji $g:U\to[0,+\infty]$
funkcja całki wewnętrznej w poniższym wzorze jest
borelowska i zachodzi
\begin{equation}\label{eq:coarea-smooth-28}
 \int_U g(x)\operatorname{Jac}_mF(x)\,\mathrm d\lambda_n(x)
 =\int_{\mathbb R^m}\left(
 \int_{F^{-1}(\{y\})}g(x)\,\mathrm d\mathcal H^{n-m}(x)
 \right)\mathrm d\lambda_m(y).
\end{equation}
Wewnętrzna całka dotyczy miary Hausdorffa
na włóknie, które, o ile nie jest puste, jest
$(n-m)$-wymiarową podrozmaitością klasy $C^1$ (a dla
$n=m$ jest zliczaniem punktów). Jeżeli jedna ze stron
jest skończona dla $|g|$, równość zachodzi także
dla rzeczywistych i zespolonych funkcji całkowalnych.
\end{twierdzenie}
\begin{proof}
Gdy $n=m$, różniczka $DF(x)$ jest odwracalna
w każdym punkcie $U$. Twierdzenie o odwzorowaniu
odwrotnym daje pokrycie $U$ otwartymi zbiorami,
na których $F$ jest dyfeomorfizmem klasy $C^1$
na otwarty obraz. Wybieramy przeliczalne podpokrycie
$(U_j)$ i rozłączny borelowski podział
$B_j=U_j\setminus\bigcup_{i<j}U_i$.
Ponieważ $\operatorname{Jac}_nF=|\det DF|$,
zamiana zmiennych na $U_j$ daje
\[
 \int_{B_j}g(x)|\det DF(x)|\,\mathrm d\lambda_n(x)
 =\int_{F(B_j)}g((F|_{U_j})^{-1}(y))
 \,\mathrm d\lambda_n(y).
\]
Zbiory $F(B_j)$ są borelowskie, ponieważ
$F|_{U_j}$ jest homeomorfizmem na otwarty obraz.
Funkcje po prawej, przedłużone przez zero poza
$F(B_j)$, są borelowskie. Ich suma jest dla każdego
$y$ sumą $g$ po $F^{-1}(\{y\})$, gdyż z każdego
$B_j$ pochodzi co najwyżej jeden punkt włókna.
Ponieważ $\mathcal H^0$ jest miarą liczącą,
zbieżność monotoniczna sum częściowych daje wzór
oraz borelowskość całki po włóknie.
W dalszym dowodzie przyjmujemy $n>m$.

\emph{Krok 1: lokalne współrzędne włóknowe.}
Ustalmy $x_0\in U$. Ponieważ $DF(x_0)$ ma rangę $m$,
po permutacji współrzędnych źródła możemy zapisać
$x=(a,z)\in\mathbb R^m\times\mathbb R^{n-m}$
tak, aby minor $D_aF(x_0)$ był odwracalny.
Odwzorowanie $x\mapsto(F(x),z)$ jest klasy $C^1$,
a jego różniczka w $x_0$ ma postać
$\begin{psmallmatrix}D_aF&D_zF\\0&I_{n-m}\end{psmallmatrix}$.
Jej wyznacznik jest niezerowy. Twierdzenie o odwzorowaniu
odwrotnym daje otwarte otoczenie $V$ punktu $x_0$
i dyfeomorfizm odwrotny $G:W\to V$ klasy $C^1$,
gdzie $W\subseteq\mathbb R^m\times\mathbb R^{n-m}$
jest otwarty. Z definicji tej odwrotności
$F(G(y,z))=y$ dla $(y,z)\in W$.
Najpierw rozważamy funkcję $g$ znikającą poza $V$.

\emph{Krok 2: zamiana zmiennych i Fubini.}
Przez klasyczną zamianę zmiennych dostajemy
\[
 \int_V g(x)\operatorname{Jac}_mF(x)\,\mathrm d\lambda_n(x)
 =\int_W g(G(y,z))\operatorname{Jac}_mF(G(y,z))
 |\det DG(y,z)|\,\mathrm d\lambda_n(y,z).
\]
Ponieważ $DF(G)D_{(y,z)}G=(I_m\ \ 0)$,
Lemat~\ref{lem:linear-coarea-28} zastosowany
punktowo do różniczek daje
$\operatorname{Jac}_mF(G)|\det DG|=
J_{n-m}(z\mapsto G(y,z))$.
Przedłużamy nieujemną borelowską funkcję
$(y,z)\mapsto g(G(y,z))J_{n-m}(z\mapsto G(y,z))$
przez zero poza $W$. Ponieważ obie euklidesowe miary
borelowskie są $\sigma$-skończone, twierdzenie Tonellego
zapewnia borelowskość całki po $z$ i daje
\[
 \int_{\mathbb R^m}\left(\int_{W_y}
 g(G(y,z))J_{n-m}(z\mapsto G(y,z))
 \,\mathrm d\lambda_{n-m}(z)\right)\mathrm d\lambda_m(y),
\]
gdzie $W_y=\{z:(y,z)\in W\}$.
Dla ustalonego $y$ zbiór $W_y$ jest otwarty, a
$z\mapsto G(y,z)$ jest osadzeniem klasy $C^1$
o obrazie $V\cap F^{-1}(\{y\})$. Jest iniekcją
jako ograniczenie dyfeomorfizmu $G$, a jego
różniczka ma rangę $n-m$ jako ograniczenie
odwracalnego $DG$ do ostatnich $n-m$ współrzędnych.
Twierdzenie~\ref{thm:immersion-area-28} zastosowane
na $W_y$ do funkcji $z\mapsto g(G(y,z))$ daje
\[
 \int_{W_y}g(G(y,z))J_{n-m}(z\mapsto G(y,z))
 \,\mathrm d\lambda_{n-m}(z)
 =\int_{V\cap F^{-1}(\{y\})}g\,\mathrm d\mathcal H^{n-m}.
\]
Każda krotność na obrazie wynosi $1$; jeśli $W_y$
jest pusty, obie strony są zerowe. Ponieważ $g$
znika poza $V$, otrzymujemy lokalnie prawą stronę
\eqref{eq:coarea-smooth-28}.

\emph{Krok 3: sklejenie.}
Wybieramy przeliczalne pokrycie $U$ otoczeniami
$V_j$ jak w kroku 1 i rozłączny podział borelowski
$B_1=V_1$, $B_j=V_j\setminus\bigcup_{i<j}V_i$.
Zastępujemy $g$ przez $g\chi_{B_j}$ i stosujemy
lokalną równość. Dla każdego $y$ zbiory
$B_j\cap F^{-1}(\{y\})$ tworzą rozłączny
borelowski podział włókna, więc całka po włóknie
jest sumą lokalnych całek. Każda z nich jest
borelowską funkcją $y$ na mocy kroku 2;
ich suma również jest borelowska. Zbieżność
monotoniczna po obu stronach daje wzór globalny.

Jeżeli całka dla $|g|$ jest skończona, wzór dla
$|g|$ pokazuje, że całka po włóknie z $|g|$
jest skończona dla prawie wszystkich $y$.
Dla rzeczywistej funkcji $g$ stosujemy wzór
do $g^+$ i $g^-$ i odejmujemy skończone całki.
Dla zespolonej $g$ stosujemy ten argument
do części rzeczywistej i urojonej; ich moduły
są ograniczone przez $|g|$.
\end{proof}

\begin{wniosek}[Różniczkowanie poziomic]
\label{cor:coarea-level-28}
Dla $U\subseteq\mathbb R^n$ otwartego i
funkcji $F\in C^1(U;\mathbb R)$
bez punktów krytycznych oraz dowolnej
nieujemnej funkcji borelowskiej $g$ zachodzi
\[
 \int_U g(x)|\nabla F(x)|\,\mathrm d\lambda_n(x)
 =\int_{\mathbb R}\left(\int_{F^{-1}(\{t\})}
 g(x)\,\mathrm d\mathcal H^{n-1}(x)\right)
 \,\mathrm d\lambda_1(t).
\]
\end{wniosek}

\section{Zastosowanie: rozkład miary obrazowej}
\begin{twierdzenie}[Gęstość obrazu przez submersję]
\label{thm:pushforward-coarea-28}
Niech $n\ge m\ge1$, $U\subseteq\mathbb R^n$
będzie otwarty, $F:U\to\mathbb R^m$ submersją
klasy $C^1$, a $w:U\to[0,+\infty]$ funkcją
borelowską.
Niech $\nu=w\lambda_n$ i $\eta=F_*\nu$.
Wówczas $\eta\ll\lambda_m$, a jej gęstość jest
równa
\[
 q(y)=\int_{F^{-1}(\{y\})}
 \frac{w(x)}{\operatorname{Jac}_mF(x)}
 \,\mathrm d\mathcal H^{n-m}(x)
\]
dla $\lambda_m$-prawie wszystkich $y$,
z możliwością przyjmowania wartości $+\infty$.
\end{twierdzenie}
\begin{proof}
Dla borelowskiego $A\subseteq\mathbb R^m$
stosujemy wzór koobszaru do funkcji
\[
 g(x)=\frac{\chi_{F^{-1}(A)}(x)w(x)}
 {\operatorname{Jac}_mF(x)}.
\]
Mianownik jest dodatni na całym $U$, więc
jest to funkcja borelowska nieujemna.
Wynik ma postać
$\eta(A)=\int_A q(y)\,\mathrm d\lambda_m(y)$.
To dowodzi bezwzględnej ciągłości i istnienia
nieujemnej borelowskiej gęstości $q$, również gdy
$\eta$ nie jest $\sigma$-skończona. Jeżeli $\eta$
jest $\sigma$-skończona, $q$ jest jej pochodną
Radona--Nikodyma względem $\lambda_m$ w sensie
Twierdzenia~\ref{thm:RN-sigfin-18}.
\end{proof}

\begin{przyklad}[Poziomice promienia w pierścieniu]
Niech $U=\{x\in\mathbb R^2:a<|x|<b\}$,
gdzie $0<a<b$, i $F(x)=|x|^2$.
Wtedy $|\nabla F(x)|=2|x|$. Coarea daje
\[
 \int_U 2|x|\,\mathrm d\lambda_2(x)
 =\int_{a^2}^{b^2}\mathcal H^1(F^{-1}(\{t\}))
 \,\mathrm d\lambda_1(t)
 =\int_{a^2}^{b^2}2\pi\sqrt t\,\mathrm d\lambda_1(t)
 =\frac{4\pi}{3}(b^3-a^3).
\]
Jest to ilustracja całkowania po \emph{włóknach},
a nie kolejny sposób obliczania powierzchni sfery.
\end{przyklad}


\section{Twierdzenie Rademachera: skąd bierze się różniczka prawie wszędzie?}
\label{sec:rademacher-28}

Wzory geometryczne wymagają jakobianu, czyli różniczki. W przypadku
odwzorowań klasy $C^1$ istnieje ona wszędzie, natomiast odwzorowanie
Lipschitzowskie może nie być różniczkowalne w pewnych punktach.
Twierdzenie Rademachera dowodzi, że taki zbiór wyjątkowy ma zerową
miarę Lebesgue'a. Samo istnienie różniczki prawie wszędzie nie dowodzi
jednak ogólnych wzorów area i coarea: do ich uzasadnienia potrzebne są
jeszcze twierdzenia o pokrywaniu i mierzeniu obrazów zbiorów krytycznych.

\begin{definicja}[Różniczkowalność Frécheta]
\label{def:frechet-rademacher-28}
Niech $n,m\ge1$, $U\subseteq\mathbb R^n$ będzie otwarty,
$f:U\to\mathbb R^m$ i $x\in U$. Mówimy, że $f$ jest
\emph{różniczkowalna w sensie Frécheta} w $x$, jeżeli istnieje
liniowe $A:\mathbb R^n\to\mathbb R^m$ takie, że
\[
 \lim_{\substack{h\to0,\ h\ne0\\x+h\in U}}
 \frac{|f(x+h)-f(x)-Ah|}{|h|}=0.
\]
Takie $A$ jest jednoznaczne: jeżeli $A,B$ spełniają
warunek graniczny, to dla każdego wektora jednostkowego $v$
podstawienie $h=tv$ i przejście do granicy daje
$|(A-B)v|=0$. Piszemy $A=Df(x)$.
Odwzorowanie $f$ nazywamy \emph{lokalnie Lipschitzowskim},
jeżeli dla każdego $x\in U$ istnieje otwarte otoczenie
$V_x\subseteq U$ i stała $L_x\ge0$ takie, że
$|f(y)-f(z)|\le L_x|y-z|$ dla wszystkich $y,z\in V_x$.
\end{definicja}

\begin{lemat}[Jednowymiarowa różniczkowalność]
\label{lem:rademacher-line-28}
Niech $a<b$, a $f:[a,b]\to\mathbb R$ będzie funkcją
Lipschitzowską ze stałą $L$. Wtedy istnieje
$g\in L^\infty((a,b),\lambda_1)$ takie, że
$|g|\le L$ prawie wszędzie oraz
\[
 f(t)-f(a)=\int_{(a,t)}g(s)\,\mathrm d\lambda_1(s)
 \quad (t\in[a,b]).
\]
W konsekwencji $f'=g$ prawie wszędzie.
\end{lemat}
\begin{proof}
Ponieważ $f$ jest Lipschitzowska, jest absolutnie ciągła.
Dla przejrzystości odzyskamy gęstość bez powoływania się na
odwrotne podstawowe twierdzenie rachunku całkowego.
Połóżmy $p(t)=f(t)+Lt$. Dla $s<t$ mamy
$p(t)-p(s)\ge -L(t-s)+L(t-s)=0$, więc $p$ jest
niemalejąca i ciągła. Jej przyrosty na półotwartych
przedziałach określają dodatnią miarę borelowską
$\mu_p((s,t])=p(t)-p(s)$; istnienie tej miary wynika
z konstrukcji Lebesgue'a--Stieltjesa przez rozszerzenie
Carathéodory'ego (ciągłość $p$ zapewnia ciągłość
premiary z góry na przedziałach). Dla każdego przedziału
$[s,t]$ zachodzi
$0\le\mu_p((s,t])\le2L(t-s)$.
Nierówność przenosi się z przedziałów półotwartych na
ich skończone rozłączne sumy, czyli na algebrę generującą
zbiory borelowskie $(a,b]$. Obie miary są skończone na
$(a,b]$, więc rodzina zbiorów borelowskich $E$, dla których
$\mu_p(E)\le2L\lambda_1(E)$, jest zamknięta na granice
rosnących i malejących ciągów zbiorów.
Twierdzenie o klasie monotonicznej daje
$\mu_p\le2L\lambda_1$ na całej $\sigma$-algebrze borelowskiej
$(a,b]$. Dominacja przez $2L\lambda_1$ zapewnia brak atomów,
w tym w $b$. Obie miary na $(a,b]$ są skończone,
a $\mu_p\ll\lambda_1$, więc Twierdzenie~\ref{thm:RN-finite-18}
daje nieujemną całkowalną gęstość $q$.
Gdyby $\lambda_1(\{q>2L\})>0$, to dla pewnego
$j\ge1$ zbiór $E_j=\{q\ge2L+1/j\}$ miałby dodatnią
miarę. Wówczas
$\mu_p(E_j)=\int_{E_j}q\,\mathrm d\lambda_1
\ge(2L+1/j)\lambda_1(E_j)>2L\lambda_1(E_j)$,
wbrew dominacji. Stąd $0\le q\le2L$ prawie wszędzie.
Dla $a<t\le b$ mamy
\[
 f(t)-f(a)=p(t)-p(a)-L(t-a)
 =\int_{(a,t)}(q(s)-L)\,\mathrm d\lambda_1(s),
\]
ponieważ punkt $t$ ma miarę zero. Dla $t=a$
obie strony wynoszą zero. Funkcja $g=q-L$ spełnia
$|g|\le L$ prawie wszędzie, więc jest całkowalna
na $(a,b)$. Twierdzenie~\ref{thm:FTC-lebesgue-25}
zastosowane do tej funkcji daje $f'=g$ prawie wszędzie.
\end{proof}

\begin{lemat}[Słabe i kierunkowe pochodne funkcji Lipschitzowskiej]
\label{lem:rademacher-direction-28}
Niech $n\ge1$, $U\subseteq\mathbb R^n$ będzie otwarty,
a $f:U\to\mathbb R$ będzie lokalnie Lipschitzowska.
Istnieją mierzalne funkcje $g_i\in L^\infty_{\mathrm{loc}}(U)$,
$i=1,\ldots,n$, takie, że dla $\varphi\in C_c^1(U)$
\[
 \int_U f(x)\partial_i\varphi(x)\,\mathrm d\lambda_n(x)
 =-\int_U g_i(x)\varphi(x)\,\mathrm d\lambda_n(x).
\]
Dla każdego ustalonego $v\in\mathbb R^n$ pochodna
kierunkowa $\partial_v f(x)=\lim_{t\to0}
(f(x+tv)-f(x))/t$ istnieje dla prawie wszystkich
$x\in U$ i jest tam równa $\sum_i v_i g_i(x)$.
\end{lemat}
\begin{proof}
Przypuśćmy najpierw, że $f$ jest Lipschitzowska
ze stałą $L$ na domknięciu otwartego prostopadłościanu
$Q\Subset U$. Dla ustalonego $i$ oraz $t\in\mathbb Q\setminus\{0\}$
niech $q_t(x)=(f(x+te_i)-f(x))/t$ tam, gdzie $x+te_i\in Q$.
Zbiór punktów istnienia pochodnej $\partial_i f$ w $Q$
jest borelowski. Istotnie, warunek Cauchy'ego dla
$q_t(x)$ przy wymiernych $t\to0$ można zapisać jako
\[
 \forall l\ge1\ \exists j\ge1\
 \forall t,s\in\mathbb Q,\quad
 0<|t|,|s|<1/j\ \Longrightarrow\ |q_t(x)-q_s(x)|\le1/l,
\]
gdzie porównujemy tylko przyrosty z $x+te_i,x+se_i\in Q$.
Każdy warunek dotyczący ustalonych $t,s$ jest borelowski,
a wszystkie kwantyfikatory przebiegają zbiory przeliczalne.
Ponieważ $Q$ jest otwarty, a iloraz różnicowy jest ciągły
w rzeczywistym parametrze $t\ne0$, warunek ten jest
równoważny istnieniu skończonej granicy dla rzeczywistych
$t\to0$. Z Lematu~\ref{lem:rademacher-line-28} każdy
przekrój zbioru braku pochodnej w kierunku osi $i$ ma
jednowymiarową miarę zero. Stosując Tonellego do funkcji
charakterystycznej tego zbioru borelowskiego, otrzymujemy
jego $\lambda_n$-miarę zero (dla $n=1$ wystarcza sam lemat).
Na zbiorze istnienia pochodnej jest ona granicą
ilorazów różnicowych dla $t=1/j$, zdefiniowanych przez
zero tam, gdzie $x+e_i/j\notin Q$; dla każdego $x\in Q$
ostatecznie te przyrosty należą do $Q$.
Przypisując pochodnej wartość zero w pozostałych punktach,
otrzymujemy borelowską funkcję o module nie większym
niż $L$. Niech $\varphi\in C_c^1(Q)$.
Na każdym przekroju $Q$ równoległym do osi $i$
piszemy zmienną tej osi jako $t\in(a,b)$.
Niech $u(t)$ i $\psi(t)$ będą odpowiednimi
ograniczeniami $f$ i $\varphi$. Lemat
\ref{lem:rademacher-line-28} daje
$u(t)=u(a)+\int_{(a,t)}u'(s)\,\mathrm d\lambda_1(s)$.
Funkcja $\psi$ znika w otoczeniach obu końców,
gdyż $\varphi$ ma zwarty nośnik w $Q$.
Stąd $\int_{(a,b)}u(a)\psi'(t)\,\mathrm d\lambda_1(t)=0$.
Dla pozostałego składnika stosujemy Fubiniego:
całkowana funkcja jest bezwzględnie całkowalna,
bo $|u'|\le L$ prawie wszędzie, a $\psi'$ jest
ograniczona i przedział jest skończony. Otrzymujemy
\begin{align*}
 \int_{(a,b)}u(t)\psi'(t)\,\mathrm d\lambda_1(t)
 &=\int_{(a,b)}u'(s)
    \left(\int_{(s,b)}\psi'(t)\,\mathrm d\lambda_1(t)\right)
    \mathrm d\lambda_1(s)\\
 &=-\int_{(a,b)}u'(s)\psi(s)\,\mathrm d\lambda_1(s).
\end{align*}
Całkując tę równość po pozostałych współrzędnych,
stosujemy Fubiniego jeszcze raz; $f$ i $\partial_i f$
są ograniczone na $Q$, a $\varphi$ i jej pochodna
mają zwarty nośnik. Dostajemy
$\int_Q f\partial_i\varphi\,\mathrm d\lambda_n
=-\int_Q(\partial_i f)\varphi\,\mathrm d\lambda_n$.
Dla $n=1$ wystarcza już równość jednowymiarowa.
Wybieramy przeliczalne pokrycie $U$ wnętrzami takich
prostopadłościanów. Na przecięciach lokalne pochodne
cząstkowe są równe prawie wszędzie, bo obie są zwykłą
pochodną cząstkową tam, gdzie istnieje. Definiujemy $g_i$
na rozłącznych borelowskich częściach pokrycia, wybierając
odpowiedni lokalny reprezentant. Przeliczalność pokrycia
zapewnia zgodność prawie wszędzie na każdym prostopadłościanie.
Funkcja $g_i$ jest lokalnie istotnie ograniczona, ponieważ
każdy zwarty podzbiór $U$ pokrywa skończenie wiele wnętrz
tych prostopadłościanów. Dla $\varphi\in C_c^1(U)$ wybieramy
gładki podział jedności podporządkowany skończonemu pokryciu
jej zwartego nośnika tymi wnętrzami. Zapisujemy $\varphi$
jako skończoną sumę funkcji testowych o nośnikach w tych
wnętrzach i sumujemy lokalne tożsamości całkowania przez
części. Daje to żądaną tożsamość na całym $U$.

Dla $v=0$ obie strony żądanej równości są zerowe.
Dla ustalonego $v\ne0$ wybieramy ortonormalną bazę
z ostatnim wektorem $e=v/|v|$. Stosujemy dotychczas
udowodnioną część lematu w tych współrzędnych.
Otrzymujemy istnienie $\partial_e f$ prawie wszędzie,
jej lokalne istotne ograniczenie i tożsamość całkowania
przez części. Ortogonalna zmiana współrzędnych
zachowuje miarę Lebesgue'a. Ponieważ
$\partial_v f=|v|\partial_e f$, dla każdej
$\varphi\in C_c^1(U)$ mamy
\begin{align*}
 \int_U(\partial_v f)\varphi\,\mathrm d\lambda_n
 &=-\int_U f\sum_{i=1}^n v_i\partial_i\varphi\,\mathrm d\lambda_n\\
 &=\int_U\Bigl(\sum_{i=1}^n v_i g_i\Bigr)\varphi\,\mathrm d\lambda_n.
\end{align*}
Niech $w=\partial_v f-\sum_i v_i g_i$, gdzie
pochodną kierunkową przedłużamy przez zero na zbiór
braku pochodnej. Jest to funkcja lokalnie całkowalna
i $\int_U w\varphi\,\mathrm d\lambda_n=0$
dla każdej $\varphi\in C_c^1(U)$.
Aby uzasadnić jednoznaczność, wybieramy nieujemną
$\zeta\in C_c^\infty(B(0,1))$ o całce $1$
i kładziemy $\zeta_r(z)=r^{-n}\zeta(z/r)$.
Dla $x\in U$ oraz $\overline B(x,r)\subset U$
funkcja $y\mapsto\zeta_r(x-y)$ należy do $C_c^1(U)$,
więc jej całka z $w$ wynosi zero.
Dla każdego punktu Lebesgue'a $x$ funkcji $w$ mamy
\[
 |w(x)|\le \|\zeta\|_\infty r^{-n}
 \int_{B(x,r)}|w(y)-w(x)|\,\mathrm d\lambda_n(y).
\]
Prawa strona dąży do zera na mocy
Twierdzenia~\ref{thm:lebesgue-diff-25},
zastosowanego lokalnie do $w$. Prawie każdy punkt
jest punktem Lebesgue'a, a zatem $w=0$
prawie wszędzie. To daje
$\partial_v f=\sum_i v_i g_i$ prawie wszędzie.
\end{proof}

\begin{twierdzenie}[Rademacher]
\label{thm:rademacher-28}
Niech $n,m\ge1$, $U\subseteq\mathbb R^n$ będzie
otwarty, a $f:U\to\mathbb R^m$ lokalnie Lipschitzowska.
Wtedy istnieje zbiór borelowski $N\subseteq U$ o
$\lambda_n(N)=0$ taki, że $f$ jest różniczkowalna
w sensie Frécheta w każdym punkcie $U\setminus N$.
Ponadto w każdym punkcie różniczkowalności
$\|Df(x)\|\le\operatorname{Lip}(f;V)$ dla
każdego otoczenia $V$ punktu $x$, na którym $f$
jest Lipschitzowska.
\end{twierdzenie}
\begin{proof}
Najpierw dowodzimy różniczkowalności prawie wszędzie
dla $m=1$. Po jej udowodnieniu stosujemy ten wynik
do każdej z $m$ składowych $f_\ell$ odwzorowania
wektorowego. Poza skończoną sumą zbiorów zerowych
wszystkie $Df_\ell(x)$ istnieją. Macierz $A$ o wierszach
$Df_\ell(x)$ spełnia wtedy
\[
 \frac{|f(x+h)-f(x)-Ah|^2}{|h|^2}
 =\sum_{\ell=1}^m
 \frac{|f_\ell(x+h)-f_\ell(x)-Df_\ell(x)h|^2}{|h|^2}
 \longrightarrow0,
\]
co dowodzi różniczkowalności wektorowej.
Niech $g_1,\ldots,g_n$ będą słabymi pochodnymi z
Lematu~\ref{lem:rademacher-direction-28}. Rozważmy
przeliczalny zbiór kierunków $D=\mathbb Q^n\setminus\{0\}$.
Dla każdego $v\in D$ na zbiorze pełnej miary
istnieje $\partial_v f(x)=\sum_i v_i g_i(x)$.
Przecięcie tych przeliczalnie wielu zbiorów nadal
ma pełną miarę. Ustalmy punkt $x$ tego przecięcia,
w którym wszystkie $g_i(x)$ są skończone, i liniowe
odwzorowanie $A_xh=\sum_i g_i(x)h_i$.

Pokażemy, że $f(x+h)=f(x)+A_xh+o(|h|)$.
Ustalmy $\varepsilon>0$. Z gęstości racjonalnych
kierunków na sferze można wybrać \emph{skończenie}
wiele wektorów jednostkowych
$q_1,\ldots,q_M\in\{v/|v|:v\in D\}$, tak aby
każdy $u\in S^{n-1}$ był w odległości mniejszej
od $\delta$ od któregoś $q_j$, gdzie
$\delta>0$ wybierzemy później. Dla $h=tu$, $t=|h|$,
wybieramy taki $q_j$. Jeżeli $L$ jest stałą
Lipschitza $f$ w otoczeniu $x$, to
\begin{align*}
 |f(x+h)-f(x)-A_xh|
 &\le |f(x+tu)-f(x+tq_j)|\\
 &\quad+|f(x+tq_j)-f(x)-tA_xq_j|
       +|tA_x(q_j-u)|\\
 &\le t(L+\|A_x\|)\delta+o_j(t).
\end{align*}
Ponieważ lista kierunków $q_j$ jest skończona,
po podzieleniu przez $t$ ilorazy $o_j(t)/t$
zbiegają do zera \emph{jednostajnie względem $j$}.
Wybieramy $\delta$ tak, aby
$(L+\|A_x\|)\delta<\varepsilon$, a potem
dostatecznie małe $t$. Wynika stąd
$\limsup_{h\to0}|f(x+h)-f(x)-A_xh|/|h|\le\varepsilon$.
Dowolność $\varepsilon$ daje różniczkowalność
Frécheta z $Df(x)=A_x$.

Pozostaje uzasadnić borelowskość zbioru wyjątkowego.
Dla $A\in\mathbb Q^{m\times n}$, $r>0$ i $x\in U$
połóżmy
\[
 R_{A,r}(x)=\sup_{\substack{h\in\mathbb Q^n,\;0<|h|<r\\
                                    x+h\in U}}
 \frac{|f(x+h)-f(x)-Ah|}{|h|},
\]
przy czym supremum pustego zbioru przyjmujemy równe
zeru. Dla ustalonych $A,r$ funkcja $R_{A,r}$ jest
borelowska: supremum liczymy po przeliczalnie wielu
borelowskich ilorazach na otwartych dziedzinach
$\{x\in U:x+h\in U\}$.
Punkt $x$ jest punktem różniczkowalności wtedy i tylko wtedy,
gdy dla każdego $q\ge1$ istnieją
$A\in\mathbb Q^{m\times n}$ i $j\ge1$ takie, że
$R_{A,1/j}(x)<1/q$. Istotnie, jeśli $Df(x)$ istnieje,
wybieramy macierz wymierną $A$ z
$\|A-Df(x)\|<1/(3q)$, a następnie $j$ tak duże,
aby reszta dla $Df(x)$, podzielona przez $|h|$,
była nie większa niż $1/(3q)$ dla $0<|h|<1/j$.
Wtedy $R_{A,1/j}(x)\le2/(3q)<1/q$.

W drugą stronę dla każdego $q$ wybieramy odpowiednią
macierz $A_q$ i promień $1/j_q$. Dla $p,q\ge1$
i dowolnego niezerowego $v\in\mathbb Q^n$ wybieramy
dodatnie wymierne $t$ tak małe, aby $h=tv$ należało
do obu dopuszczalnych zakresów przyrostów. Odejmując
reszty dla $A_q$ i $A_p$, otrzymujemy
\[
 |(A_q-A_p)h|\le (1/q+1/p)|h|.
\]
Po podzieleniu przez $t$ i użyciu gęstości $\mathbb Q^n$
nierówność obowiązuje dla wszystkich $v\in\mathbb R^n$,
więc $\|A_q-A_p\|\le1/q+1/p$. Z zupełności
skończenie wymiarowej przestrzeni macierzy otrzymujemy
$A_q\to A$. Dla ustalonego $q$ i dostatecznie małych
wymiernych $h\ne0$ mamy
\[
 \frac{|f(x+h)-f(x)-Ah|}{|h|}
 \le1/q+\|A_q-A\|.
\]
Przez ciągłość względem $h\ne0$ ta nierówność obowiązuje
również dla rzeczywistych $h$ w tym samym otwartym
zakresie. Dla każdej dodatniej tolerancji wybieramy
$q$ tak duże, aby prawa strona była od niej mniejsza,
co dowodzi różniczkowalności z różniczką $A$.
Zatem zbiór punktów różniczkowalności ma postać
\[
 \bigcap_{q\ge1}\ \bigcup_{A\in\mathbb Q^{m\times n}}
 \ \bigcup_{j\ge1}\{x\in U:R_{A,1/j}(x)<1/q\}.
\]
Jest borelowski jako przeliczalne przecięcie
przeliczalnych sum zbiorów borelowskich. Jego dopełnienie ma zerową
miarę na mocy poprzedniego argumentu.
Jeśli $f$ jest Lipschitzowska na otoczeniu $V$ punktu
różniczkowalności $x$ ze stałą $L$, to dla wektora
jednostkowego $u$ i dostatecznie małych $t>0$
mamy $|f(x+tu)-f(x)|/t\le L$.
Przechodząc $t\downarrow0$, dostajemy
$|Df(x)u|\le L$. Biorąc supremum po $|u|=1$,
a następnie infimum po dopuszczalnych stałych $L$,
otrzymujemy $\|Df(x)\|\le\operatorname{Lip}(f;V)$.
\end{proof}

\begin{uwaga}[Dlaczego pochodne cząstkowe nie wystarczają?]
Na przykład $f(x,y)=|x|$ ma pochodne cząstkowe
w punktach z $x\ne0$, ale nie ma różniczki na prostej
$x=0$; zbiór ten jest zerowy dla $\lambda_2$.
Ponadto istnienie wszystkich dwustronnych pochodnych
kierunkowych w \emph{jednym} punkcie nie zapewnia
różniczkowalności Frécheta. Dla funkcji określonej przez
$f(0,0)=0$ i $f(x,y)=x^3/(x^2+y^2)$ poza początkiem
mamy $\partial_v f(0,0)=v_1^3/(v_1^2+v_2^2)$ dla
$v=(v_1,v_2)\ne0$ (iloraz różnicowy ma tę wartość
dla każdego $t\ne0$), a $\partial_0 f(0,0)=0$.
Nie jest to funkcja liniowa kierunku: wartości dla
$(1,0)$ i $(0,1)$ wynoszą $1$ i $0$, lecz dla $(1,1)$
otrzymujemy $1/2$. Istotna w dowodzie
Rademachera jest zgodność pochodnych kierunkowych
z jedną \emph{słabą} różniczką, a następnie
argument skończonej siatki kierunków.
\end{uwaga}

\begin{przyklad}[Wszystkie pochodne cząstkowe w punkcie, ale brak różniczki]
Niech $f:\mathbb R^2\to\mathbb R$ będzie zadana przez
$f(0,0)=0$ oraz $f(x,y)=xy/\sqrt{x^2+y^2}$
dla $(x,y)\ne(0,0)$. Poza początkiem ma ona
ograniczony gradient: jest jednorodna stopnia $1$,
a jej pochodne pierwszego rzędu są ciągłe na
okręgu jednostkowym i jednorodne stopnia $0$.
W związku z tym $f$ jest globalnie Lipschitzowska
(dla odcinka przechodzącego przez $0$ dzielimy go
na dwa odcinki). Obie pochodne cząstkowe w $(0,0)$
istnieją i są równe $0$, ponieważ $f(x,0)=f(0,y)=0$.
Jednak dla $h=(t,t)$, $t>0$, mamy
$f(t,t)=t/\sqrt2$, a więc
$|f(h)-f(0)-0\cdot h|/|h|=1/2$.
Nie istnieje zatem różniczka Frécheta w $(0,0)$.
Pokazuje to, dlaczego sam argument z istnienia
pochodnych cząstkowych prawie wszędzie jest
niewystarczający dla twierdzenia Rademachera.
\end{przyklad}

\begin{przyklad}[Różniczkowalność prawie wszędzie a brak gładkości]
Niech $F:\mathbb R^2\to\mathbb R^2$,
$F(x,y)=(x,y+|x|)$. Jest globalnie Lipschitzowskie,
choć nie jest klasy $C^1$. Dla $x>0$ i $x<0$
ma odpowiednio macierze $\begin{psmallmatrix}1&0\\1&1\end{psmallmatrix}$
i $\begin{psmallmatrix}1&0\\-1&1\end{psmallmatrix}$,
a na prostej $x=0$ nie ma różniczki. Mimo to
$|\det DF|=1$ prawie wszędzie. Wzór zamiany
zmiennych jest poprawny, ale jego uzasadnienie
w klasie Lipschitzowskiej nie może ograniczać się
do zwykłego twierdzenia o dyfeomorfizmach $C^1$.
\end{przyklad}

\begin{uwaga}[Związek z miarami wektorowymi]
Twierdzenie Rademachera wykorzystuje skończony wymiar
\emph{przestrzeni wartości}. Dla dowolnej przestrzeni
Banacha $X$ nie musi być prawdą, że każda funkcja
Lipschitzowska $[0,1]\to X$ ma pochodną prawie wszędzie.
Dokładną granicę wyznacza \emph{własność Radona--Nikodyma},
która zostanie zdefiniowana i scharakteryzowana w
Rozdziale~\ref{ch:vector-measures-30}.
\end{uwaga}

\section{Aproksymacja typu Łuzina--Whitneya w klasie $C^1$}
Aby przejść od wersji gładkich do odwzorowań Lipschitzowskich,
nie wystarczy różniczkowalność prawie wszędzie. Potrzebujemy
możliwości zastąpienia odwzorowania na dużych zbiorach
\emph{dokładnie} funkcją klasy $C^1$, wraz z jej różniczką.
Poniżej dowodzimy odpowiedniego faktu, nie przyjmując
jako gotowego twierdzenia o przedłużaniu Whitneya.

\begin{lemat}[Przedłużenie zgodnych danych pierwszego rzędu]
\label{lem:whitney-c1-28}
Niech $K\subseteq\mathbb R^n$ będzie niepustym zbiorem zwartym,
$a:K\to\mathbb R^N$ i
$A:K\to\mathcal L(\mathbb R^n,\mathbb R^N)$ funkcjami ciągłymi.
Załóżmy, że
\begin{equation}\label{eq:whitney-condition-28}
 \lim_{r\downarrow0}\ \sup_{\substack{x,y\in K,\;0<|x-y|\le r}}
 \frac{|a(y)-a(x)-A(x)(y-x)|}{|y-x|}=0,
\end{equation}
przy czym supremum pustego zbioru jest równe zeru.
Istnieje funkcja $\widetilde a\in C^1(\mathbb R^n;\mathbb R^N)$
taka, że $\widetilde a|_K=a$ i
$D\widetilde a(x)=A(x)$ dla $x\in K$.
\end{lemat}
\begin{proof}
Wprowadzamy dwa moduły ciągłości: $\omega(t)$ jest
supremum z \eqref{eq:whitney-condition-28} z $r=t$,
a $\alpha(t)=\sup\{\|A(x)-A(y)\|:x,y\in K,
|x-y|\le t\}$. Obydwa dążą do zera dla $t\downarrow0$.

\emph{Krok 1: kostki otwartego dopełnienia.}
Zbiór $\mathbb R^n\setminus K$ rozkładamy na kostki
diadyczne $Q$ o rozłącznych wnętrzach, dobierając
kostki maksymalne (względem włączenia w rodzinie
kostek diadycznych) spośród tych, dla których
$20\sqrt n\,\ell(Q)\le\operatorname{dist}(Q,K)$.
Tutaj $\ell(Q)$ oznacza długość boku kostki.
Używamy domkniętych kostek standardowej siatki
diadycznej, tj. zbiorów $2^{-j}(b+[0,1]^n)$ dla
$j\in\mathbb Z$ i $b\in\mathbb Z^n$.
Dwie kostki tej siatki o przecinających się wnętrzach
są zagnieżdżone: w każdej współrzędnej końce przedziału
o większej długości należą do siatki wyznaczonej przez
mniejszą długość, więc wnętrze mniejszego przedziału
nie może przecinać końca większego. Stąd różne
maksymalne kostki mają rozłączne wnętrza.
Dla $x\notin K$ kładziemy $\delta_x=\operatorname{dist}(x,K)>0$.
Jeśli kostka $Q$ zawiera $x$ i
$21\sqrt n\,\ell(Q)\le\delta_x$, to
$\operatorname{dist}(Q,K)\ge\delta_x-\sqrt n\,\ell(Q)
\ge20\sqrt n\,\ell(Q)$.
Każdy punkt dopełnienia należy więc do dostatecznie
małej kostki spełniającej nierówność. Kolejne kostki
macierzyste zawierają ten sam punkt $x$, więc ich odległość
od $K$ nie przekracza $\operatorname{dist}(x,K)$, natomiast
ich boki rosną bez ograniczenia. Każda dopuszczalna kostka
ma więc maksymalnego przodka: jeżeli kostka nie spełnia
warunku, to jej więksi przodkowie też go nie spełniają,
ponieważ odległość od $K$ nie rośnie, a bok rośnie.
Maksymalne kostki pokrywają zatem całe dopełnienie $K$,
a każda z nich leży w tym dopełnieniu dzięki dodatniej
odległości od $K$. Kostka macierzysta $Q'$ kostki
maksymalnej $Q$ nie spełnia nierówności, skąd
\[
 \operatorname{dist}(Q,K)
 \le\operatorname{dist}(Q',K)+\operatorname{diam}Q'
 <42\sqrt n\,\ell(Q).
\]
Oznaczmy przez $Q^*$ współśrodkową kostkę o boku
$(9/8)\ell(Q)$. Dla $x\in Q^*$ mamy
\[
 19\sqrt n\,\ell(Q)\le\operatorname{dist}(x,K)
 \le44\sqrt n\,\ell(Q).
\]
Dolne oszacowanie wynika z odległości $x$ od $Q$,
nie większej niż $\sqrt n\,\ell(Q)/16$; górne
wynika z poprzedniej nierówności i średnicy $Q^*$.
Zatem gdy $Q^*\cap R^*\ne\varnothing$, stosunek boków
$Q,R$ zawiera się w $[19/44,44/19]$.
Dla ustalonego $x\notin K$ wszystkie kostki, dla których
$x\in Q^*$, mają boki porównywalne z
$\delta=\operatorname{dist}(x,K)>0$ i mieszczą się
w kuli $B(x,C_n\delta)$. Ich wnętrza są rozłączne,
a objętość każdej jest nie mniejsza niż $c_n\delta^n$.
Porównanie z objętością tej kuli ogranicza ich liczbę
przez stałą zależną tylko od $n$. Ten sam argument na
kuli $B(x,\delta/4)$, na której odległość od $K$ pozostaje
między $3\delta/4$ a $5\delta/4$, dowodzi lokalnej
skończoności rodziny powiększonych kostek w dopełnieniu $K$.

Dla każdej kostki wybieramy $p_Q\in K$ minimalizujący
odległość od jej środka; minimum istnieje dzięki
niepustości i zwartości $K$. Wybieramy $C^\infty$-funkcje
$\psi_Q\ge0$, równe $1$ na $Q$, o nośnikach zawartych
w kostkach współśrodkowych o boku $(9/8)\ell(Q)$,
i spełniające $\|D\psi_Q\|\le C_n/\ell(Q)$.
Dokładniej, wybieramy nieujemną funkcję
$\theta\in C_c^\infty((-9/16,9/16))$, równą $1$
na $[-1/2,1/2]$ i nie większą niż $1$.
Dla środka $c_Q$ kostki $Q$ kładziemy
$\psi_Q(x)=\prod_{i=1}^n\theta((x_i-(c_Q)_i)/\ell(Q))$.
Reguła łańcuchowa daje
$\|D\psi_Q\|\le\sqrt n\|\theta'\|_\infty/\ell(Q)$.
Suma $\sum_Q\psi_Q$ jest lokalnie skończona i nie
mniejsza niż $1$ na $\mathbb R^n\setminus K$.
Oznaczmy $S(x)=\sum_R\psi_R(x)$. Funkcja $S$
jest gładka na dopełnieniu $K$ i $S\ge1$.
Funkcje $\varphi_Q=\psi_Q/S$ są nieujemne, gładkie,
mają nośniki w $Q^*$ i sumują się do $1$.
Dla $x\in Q^*$ wszystkie kostki $R$, których składnik
może mieć niezerową pochodną w $x$, mają boki
porównywalne z $\ell(Q)$, a ich liczba jest ograniczona
przez stałą zależną tylko od $n$. Stąd
$\|DS(x)\|\le C_n/\ell(Q)$. Z reguły ilorazu,
$S\ge1$ oraz $0\le\psi_Q\le1$ wynika
\[
 \|D\varphi_Q(x)\|
 \le\frac{\|D\psi_Q(x)\|}{S(x)}
 +\frac{\psi_Q(x)\|DS(x)\|}{S(x)^2}
 \le C_n/\ell(Q).
\]
Poza $Q^*$ pochodna jest zerowa, więc oszacowanie
obowiązuje w całym dopełnieniu $K$.

\emph{Krok 2: sklejanie przybliżeń afinicznych.}
Dla $x\notin K$ definiujemy
\[
 P_Q(x)=a(p_Q)+A(p_Q)(x-p_Q),\qquad
 \widetilde a(x)=\sum_Q\varphi_Q(x)P_Q(x),
\]
a dla $x\in K$ przyjmujemy $\widetilde a(x)=a(x)$.
Niech $x\notin K$, $\delta=\operatorname{dist}(x,K)$,
a $p\in K$ będzie punktem najbliższym $x$.
Dla każdej kostki, dla której $x\in Q^*$,
mamy $\ell(Q)\asymp\delta$. Ponieważ $p_Q$ jest
punktem najbliższym środkowi $c_Q$, mamy
$|c_Q-p_Q|\le|c_Q-p|\le |c_Q-x|+\delta$.
Zatem $|x-p_Q|\le 2|x-c_Q|+\delta\le C_n\delta$
i $|p_Q-p|\le |p_Q-x|+\delta\le C_n\delta$.
Rozpisujemy różnicę przybliżeń afinicznych:
\begin{align*}
 P_Q(x)-a(p)-A(p)(x-p)
 &=a(p_Q)-a(p)-A(p)(p_Q-p)\\
 &\quad+(A(p_Q)-A(p))(x-p_Q).
\end{align*}
Pierwszy składnik ma normę nie większą niż
$|p_Q-p|\omega(|p_Q-p|)$ (jest zerowy dla $p_Q=p$),
a drugi nie większą niż
$\alpha(|p_Q-p|)|x-p_Q|$.
Z monotoniczności obu modułów otrzymujemy
\[
 |P_Q(x)-[a(p)+A(p)(x-p)]|
 \le C_n\delta\,[\omega(C_n\delta)+\alpha(C_n\delta)],
\]
zaś $\|A(p_Q)-A(p)\|\le\alpha(C_n\delta)$.
Ponieważ $\sum_Q\varphi_Q=1$, pierwsze z poniższych
oszacowań otrzymujemy przez sumowanie poprzedniego
oszacowania z wagami $\varphi_Q(x)$.
Ponadto $\sum_QD\varphi_Q=0$, więc w ustalonym punkcie $x$
(wybór $p$ jest tu ustalony, nie różniczkujemy go) mamy
\begin{align*}
 D\widetilde a(x)-A(p)
 &=\sum_Q\varphi_Q(x)(A(p_Q)-A(p))\\
 &\quad+\sum_Q [P_Q(x)-a(p)-A(p)(x-p)]D\varphi_Q(x).
\end{align*}
Druga suma oznacza sumę odwzorowań liniowych
$h\mapsto [P_Q(x)-a(p)-A(p)(x-p)]D\varphi_Q(x)h$.
Liczba niezerowych składników jest jednostajnie ograniczona,
a $\|D\varphi_Q(x)\|\le C_n/\delta$ dla uczestniczących
kostek, zatem otrzymujemy
\begin{align*}
 |\widetilde a(x)-a(p)-A(p)(x-p)|
 &\le C_n\delta[\omega(C_n\delta)+\alpha(C_n\delta)],\\
 \|D\widetilde a(x)-A(p)\|
 &\le C_n[\omega(C_n\delta)+\alpha(C_n\delta)].
\end{align*}
Stałe mogą również zależeć od wyboru funkcji
odcinających, ale nie od $x$.

Gdy $x\to p_0\in K$ od zewnątrz, mamy
$\delta\le|x-p_0|$ i $|p-p_0|\le2|x-p_0|$.
Stąd $p\to p_0$ i $\delta\to0$. Drugie oszacowanie
oraz ciągłość $A$ dają $D\widetilde a(x)\to A(p_0)$.
Dla pierwszego rozwinięcia piszemy
\begin{align*}
 \widetilde a(x)-a(p_0)-A(p_0)(x-p_0)
 &=\widetilde a(x)-a(p)-A(p)(x-p)\\
 &\quad+a(p)-a(p_0)-A(p_0)(p-p_0)\\
 &\quad+(A(p)-A(p_0))(x-p).
\end{align*}
Po podzieleniu przez $|x-p_0|$ normy kolejnych
składników są ograniczone odpowiednio przez
$C_n[\omega(C_n\delta)+\alpha(C_n\delta)]$,
$2\omega(2|x-p_0|)$ i $\alpha(2|x-p_0|)$.
Wszystkie trzy ograniczenia dążą do zera,
co daje żądany rozwój z resztą $o(|x-p_0|)$.
Ten sam warunek bezpośrednio daje taki rozwój,
gdy $x\in K$. Zatem $\widetilde a$ jest
różniczkowalna także na $K$, jej różniczka jest
tam równa $A$ i jest ciągła na całej przestrzeni.
\end{proof}

\begin{lemat}[Zwarte fragmenty klasy $C^1$]
\label{lem:lusin-whitney-28}
Niech $U\subseteq\mathbb R^n$ będzie otwarty,
$F:U\to\mathbb R^N$ lokalnie Lipschitzowski,
$Q\subseteq U$ zwarty i $\varepsilon>0$.
Istnieją zbiór zwarty $K\subseteq Q$ oraz
$F_K\in C^1(\mathbb R^n;\mathbb R^N)$ takie, że
\[
 \lambda_n(Q\setminus K)<\varepsilon,\qquad
 F_K(x)=F(x),\quad DF_K(x)=DF(x)\quad(x\in K).
\]
Różniczkę $DF$ rozumiemy tylko w punktach,
w których istnieje; wszystkie punkty $K$ są
takimi punktami. Ponadto istnieje przeliczalna rodzina zbiorów zwartych
$K_j\subseteq U$, z których każdy ma opisane wyżej
przedłużenie i dla której
$\lambda_n(U\setminus\bigcup_jK_j)=0$.
\end{lemat}
\begin{proof}
Niech $D\subseteq U$ będzie borelowskim zbiorem
punktów różniczkowalności $F$. Na mocy
Twierdzenia~\ref{thm:rademacher-28} dopełnienie
$D$ jest zerowe. Funkcja $x\mapsto DF(x)$ jest
borelowska na $D$, ponieważ jej współczynniki
są granicami ilorazów różnicowych w kierunkach
wektorów bazowych.

Dla $j,l\ge1$ oznaczmy przez $H_{j,l}$ zbiór
punktów $x\in D\cap Q$, dla których
\[
 |F(x+h)-F(x)-DF(x)h|\le j^{-1}|h|
\]
dla każdego $h\in\mathbb Q^n$ spełniającego
$0<|h|<l^{-1}$ i $x+h\in U$. Jest to zbiór
borelowski, bo warunek sprawdzamy dla
przeliczalnej liczby przyrostów. Z różniczkowalności wynika
$\bigcup_lH_{j,l}=D\cap Q$, przy czym dla ustalonego
$j$ zbiory $H_{j,l}$ rosną wraz z $l$.
Ponieważ $Q$ jest zwarty, $\lambda_n(Q)<\infty$,
więc ciągłość miary z dołu pozwala dla każdego $j$
wybrać $l_j$ tak duże,
aby $\lambda_n((D\cap Q)\setminus H_{j,l_j})
<\varepsilon2^{-j-2}$. Zbiór
$G=\bigcap_{j\ge1}H_{j,l_j}$ spełnia
\[
 \lambda_n(Q\setminus G)
 \le\lambda_n(Q\setminus D)
 +\sum_{j\ge1}\lambda_n((D\cap Q)\setminus H_{j,l_j})
 \le\varepsilon/4<\varepsilon/2.
\]
Z wewnętrznej regularności miary Lebesgue'a
wybieramy zwarty $K_0\subseteq G$ z
$\lambda_n(G\setminus K_0)<\varepsilon/4$.
Ograniczenie $DF$ do $K_0$ jest borelowskim
odwzorowaniem o wartościach w skończenie wymiarowej
przestrzeni macierzy, a miara Lebesgue'a ograniczona
do $K_0$ jest skończona i regularna. Stosujemy Twierdzenie~\ref{thm:lusin-10}
do tego odwzorowania (lub do wszystkich jego
$Nn$ współczynników, wybierając dla każdego zbiór zwarty
ze stratą mniejszą niż $\varepsilon/(4Nn)$ i biorąc
ich skończone przecięcie). Otrzymujemy zwarty
$K\subseteq K_0$, na którym $DF|_K$ jest ciągła
i $\lambda_n(K_0\setminus K)<\varepsilon/4$.
W konsekwencji $\lambda_n(G\setminus K)<\varepsilon/2$. Jeśli $K=\varnothing$,
przyjmujemy $F_K\equiv0$ i pierwsza teza jest
trywialna; w dalszym ciągu zakładamy $K\ne\varnothing$.

Warunek w definicji $H_{j,l_j}$, przez ciągłość
$F$ względem przyrostu $h$, pozostaje prawdziwy
również dla rzeczywistych $h$ o tej samej
ograniczonej długości. Stąd dla każdego $j$ i
$x,y\in K$ o $0<|x-y|<l_j^{-1}$ zachodzi
\[
 |F(y)-F(x)-DF(x)(y-x)|\le j^{-1}|x-y|.
\]
Dla każdej zadanej dodatniej tolerancji wybieramy $j$
tak, aby $1/j$ było od niej mniejsze; dla $r<1/l_j$
supremum ilorazów w warunku Whitneya nie przekracza
$1/j$. Ponadto $a=F|_K$ jest ciągłe dzięki lokalnej
Lipschitzowskości $F$, a $A=DF|_K$ jest ciągłe
z wyboru $K$. Zatem dane $a,A$ spełniają wszystkie
założenia Lematu~\ref{lem:whitney-c1-28}.
Jego zastosowanie daje żądane przedłużenie.

Aby uzyskać ostatnią tezę, wybieramy przeliczalne
pokrycie $U$ zwartymi kostkami $Q_s\Subset U$.
Dla każdego $s$ i $j$ wykonujemy konstrukcję
z $\varepsilon=1/j$ na $Q_s$, uzyskując
$K_{s,j}$. Ponieważ
$\lambda_n(Q_s\setminus\bigcup_jK_{s,j})\le1/j$
dla każdego $j$, miara ta jest zerowa.
Przeliczalna suma pozostałych zbiorów zerowych
ma miarę zero, co kończy dowód.
\end{proof}

\section{Dwa oszacowania dla zbiorów krytycznych}
\begin{lemat}[Obraz zbioru o niedostatecznej randze]
\label{lem:critical-area-null-28}
Niech $1\le k\le N$, $V\subseteq\mathbb R^k$ będzie
otwarty, $G\in C^1(V;\mathbb R^N)$, a
$K\subseteq V$ zwarty. Jeżeli
$\operatorname{rank}DG(x)<k$ dla wszystkich $x\in K$,
to $\mathcal H^k(G(K))=0$.
\end{lemat}
\begin{proof}
Jeśli $K=\varnothing$, teza wynika z $\mathcal H^k(\varnothing)=0$.
W przeciwnym razie wybieramy $\rho>0$ tak, aby
$K_\rho=\{z:\operatorname{dist}(z,K)\le\rho\}$ był
zwartym podzbiorem $V$. Istnieje takie $\rho$, ponieważ
$K$ jest zwarty i $V$ otwarty. Niech $M$ będzie
ograniczeniem normy $DG$ na $K_\rho$. Ustalmy
$0<\eta\le1$. Z jednostajnej ciągłości $DG$ na $K_\rho$
wybieramy $r>0$ dostatecznie małe, aby $\sqrt k\,r<\rho$
i aby dla punktów tego otoczenia odległych o co najwyżej
$\sqrt k\,r$ różnica ich różniczek miała normę co najwyżej
$\eta$. Warunki te zachodzą również dla wszystkich
mniejszych dodatnich $r$.

Rozpatrujemy domknięte kostki siatki o boku $r$,
przecinające $K$. Dla każdej takiej kostki $Q$
wybieramy $x_Q\in K\cap Q$. Dla $x\in K\cap Q$
cały odcinek $[x_Q,x]$ leży w $K_\rho$, więc
\begin{align*}
 G(x)-G(x_Q)-DG(x_Q)(x-x_Q)
 &=\int_0^1[DG(x_Q+t(x-x_Q))-DG(x_Q)]\\
 &\qquad\cdot(x-x_Q)\,\mathrm d\lambda_1(t).
\end{align*}
Norma tej różnicy nie przekracza $\eta\sqrt k\,r$.
Niech $P_Q=\operatorname{im}DG(x_Q)$ oraz
$j=\dim P_Q\le k-1$. Po odjęciu $G(x_Q)$
obraz $G(K\cap Q)$ leży w $\eta\sqrt k\,r$-otoczeniu
kuli w $P_Q$ o promieniu $M\sqrt k\,r$.

Dla $j\ge1$ wybieramy ortonormalne współrzędne w $P_Q$.
Kula ta mieści się w $j$-wymiarowej kostce o boku
$2M\sqrt k\,r$. Dzielimy ją siatką o boku $\eta r$.
Liczba przecinających ją kostek nie przekracza
$(2M\sqrt k/\eta+2)^j\le C_{k,M}\eta^{-j}$.
Każda z nich ma średnicę $\sqrt j\,\eta r$.
Jeśli $j=0$, część liniowa jest zerowa i wystarcza
jeden punkt $0\in P_Q$, czyli ten sam szacunek
liczby zbiorów przy wykładniku $j=0$.
Po powiększeniu każdej kostki do kuli w $\mathbb R^N$
o promieniu $(\sqrt k+\sqrt j)\eta r$ uzyskujemy
pokrycie przesuniętego obrazu. Ponieważ $j\le k-1$
i $\eta\le1$, koszt tego pokrycia nie przekracza
$C_{k,N,M}\eta^{-j}(\eta r)^k
\le C_{k,N,M}\eta r^k$.

Zwartość $K$ zapewnia, że leży ono w ustalonej
kostce $[-R,R]^k$ dla pewnego $R>0$. Dla $r\le1$
liczba kostek siatki przecinających $K$ nie przekracza
$(2R/r+2)^k\le C_{k,R}r^{-k}$. Łącząc powyższe
pokrycia, otrzymujemy pokrycie $G(K)$ zbiorami
o średnicach co najwyżej $C_k\eta r$ i koszcie
(co uwzględnia stałą normalizującą $c_k$)
nie większym niż $C_{k,N,M,R}\eta$.
Dla każdej zadanej skali $\delta>0$ możemy wybrać
$r$ tak małe, że $C_k\eta r\le\delta$ i wszystkie
wcześniejsze warunki są spełnione. Zatem
$\mathcal H^k_\delta(G(K))\le C_{k,N,M,R}\eta$.
Biorąc supremum po $\delta>0$, a następnie
przechodząc $\eta\downarrow0$, otrzymujemy
$\mathcal H^k(G(K))=0$.
\end{proof}

\begin{lemat}[Zbiory zerowe i ich włókna]
\label{lem:null-fibres-lipschitz-28}
Niech $n\ge m\ge1$, $F:\mathbb R^n\to\mathbb R^m$
będzie lokalnie Lipschitzowska, a
$Z\subseteq\mathbb R^n$ zbiorem o zerowej zewnętrznej
mierze Lebesgue'a. Wówczas
\[
 \mathcal H^{n-m}(Z\cap F^{-1}(\{y\}))=0
 \quad\text{dla }\lambda_m\text{-prawie wszystkich }y.
\]
\end{lemat}
\begin{proof}
Przez przeliczalne pokrycie dziedziny otwartymi
kostkami o domknięciach zawartych w większych otoczeniach,
na których $F$ jest Lipschitzowska, wystarczy rozpatrzyć
część $Z$ zawartą w jednej takiej kostce: dla każdej
części otrzymamy zbiór wyjątkowy o zewnętrznej mierze
zero, a ich przeliczalna suma nadal będzie zerowa.
Poza nią podaddytywność $\mathcal H^{n-m}$ da tezę
dla całego $Z$. Wybieramy dodatnią
stałą Lipschitza $L>0$ na większym otoczeniu (również
gdy $F$ jest na nim stałe).
Piszemy $d=n-m$. Ustalmy $\delta,\varepsilon>0$.
Z warunku $\lambda_n^*(Z)=0$ i kalibracji
$\mathcal H^n=\lambda_n$ (na zbiorach mierzalnych,
w szczególności zerowych) wynika $\mathcal H^n(Z)=0$.
Stąd dla dowolnej dodatniej skali $\tau$ istnieje
przeliczalne pokrycie $Z$ niepustymi zbiorami
$D_i$ o średnicach co najwyżej $\tau$ i sumie
$n$-tych potęg średnic dowolnie małej. Odrzucamy
zbiory nieprzecinające $Z$ i wybieramy $x_i\in D_i\cap Z$.
Jeśli $a_i=\operatorname{diam}D_i>0$, domknięta
kostka o środku $x_i$ i boku $2a_i$ pokrywa $D_i$,
a jej średnica wynosi $2\sqrt n\,a_i$.
Jeśli $a_i=0$, zbiór jest singletonem; pokrywamy
go kostką o dodatnim boku, wybranym tak małym,
aby jej średnica była mniejsza niż $2\sqrt n\,\tau$
i jej koszt nie przekraczał $\varepsilon 2^{-i-1}$.
Wybierając pierwotną sumę $n$-tych potęg średnic
mniejszą niż $\varepsilon/(2(2\sqrt n)^n)$,
otrzymujemy kostki $Q_i$ z
$\sum_i(\operatorname{diam}Q_i)^n<\varepsilon$.
Możemy wybrać $\tau$ tak małe, aby średnice były
mniejsze niż $\delta$ oraz niż odległość domknięcia
rozpatrywanej kostki od dopełnienia większego otoczenia,
na którym działa $L$ (jeżeli dopełnienie jest puste,
wystarczy pierwszy warunek). Wszystkie $Q_i$ leżą więc
w tym otoczeniu.
Niech $z_i$ będzie środkiem $Q_i$ i niech
$V_i=B(F(z_i),L\operatorname{diam}Q_i)$.
Jeśli $x\in Z\cap Q_i$ i $F(x)=y$, to
$|y-F(z_i)|\le L|x-z_i|\le L\operatorname{diam}Q_i/2
<L\operatorname{diam}Q_i$, a więc $y\in V_i$.
Dla ustalonego $y$ kostki o indeksach spełniających
$y\in V_i$ pokrywają całe włókno $Z\cap F^{-1}(\{y\})$.
Ich średnice są mniejsze niż $\delta$, dlatego
z definicji kosztu pokrycia otrzymujemy punktowo
\[
 \mathcal H^d_\delta(Z\cap F^{-1}(\{y\}))
 \le C_{n,m}\sum_i(\operatorname{diam}Q_i)^d\chi_{V_i}(y).
\]
Dla $d=0$ koszt niepustego zbioru w pokryciu
wynosi $1$, zgodnie z prawą stroną dla $C_{n,m}\ge1$.
Dla $d>0$ wystarcza $C_{n,m}\ge c_d$.
Prawa strona jest borelowska. Z twierdzenia
o zbieżności monotonicznej, ze wzoru
$\lambda_m(V_i)=\lambda_m(B(0,1))L^m
(\operatorname{diam}Q_i)^m$ i z $d+m=n$ otrzymujemy
\[
 \int_{\mathbb R^m}C_{n,m}\sum_i
 (\operatorname{diam}Q_i)^d\chi_{V_i}(y)
 \,\mathrm d\lambda_m(y)
 \le C_{n,m}L^m\varepsilon.
\]
Zatem dla dowolnego $t>0$ zbiór, na którym
$\mathcal H^d_\delta$ rozpatrywanego włókna
przekracza $t$, ma zewnętrzną miarę nie większą
niż $C_{n,m}L^m\varepsilon/t$. Ponieważ
$\varepsilon$ jest dowolne, miara ta jest zerowa.
Stosujemy to do $\delta=1/j$ i $t=1/l$ dla
$j,l\ge1$. Przeliczalna suma otrzymanych zbiorów
wyjątkowych ma zewnętrzną miarę zero. Poza nią
dla każdego $j$ zawartość $\mathcal H^d_{1/j}$
rozpatrywanego włókna nie przekracza $1/l$ dla
każdego $l$, więc jest zerowa. Przechodząc
$j\to\infty$, otrzymujemy $\mathcal H^d=0$
na prawie wszystkich włóknach.
\end{proof}

\begin{lemat}[Włókna zwartego zbioru krytycznego]
\label{lem:critical-coarea-null-28}
Niech $n\ge m\ge1$, $V\subseteq\mathbb R^n$
będzie otwarty, $G\in C^1(V;\mathbb R^m)$,
a $K\subseteq V$ zwarty. Jeżeli
$\operatorname{rank}DG(x)<m$ dla wszystkich $x\in K$,
to
\[
 \mathcal H^{n-m}(K\cap G^{-1}(\{y\}))=0
 \quad\text{dla }\lambda_m\text{-prawie wszystkich }y.
\]
\end{lemat}
\begin{proof}
Piszemy $d=n-m$. Dla ustalonej skali $\delta$ nie
zakładamy mierzalności funkcji
$y\mapsto\mathcal H^d_\delta(K\cap G^{-1}(\{y\}))$.
W definicji z Rozdziału~\ref{ch:hausdorff-26} dopuszcza
się dowolne zbiory o średnicy \emph{co najwyżej}
$\delta$; zastąpienie ich otwartymi zbiorami o średnicy
ściśle mniejszej niż $\delta$ może zmienić infimum.
Wystarczy oszacowanie zewnętrznej miary zbiorów
parametrów, jak w Lemacie~\ref{lem:null-fibres-lipschitz-28}.

Jeśli $K=\varnothing$, wszystkie włókna są puste.
Niech $K\ne\varnothing$ i $0<\eta\le1$.
Pokryjmy $K$ domkniętymi kostkami siatki
o boku $r$, dla których wybieramy $x_Q\in K\cap Q$.
Wybieramy zwarte otoczenie $K$ zawarte w $V$,
zawierające wszystkie punkty odległe od $K$
o co najwyżej pewną dodatnią liczbę.
Dla dostatecznie małego $r$ odcinki $[x_Q,x]$
dla $x\in K\cap Q$ leżą w tym otoczeniu,
a jednostajna ciągłość $DG$ zapewnia na nich
$\|DG(z)-DG(x_Q)\|\le\eta$.
Całkując różniczkę wzdłuż odcinka, dostajemy
\[
 |G(x)-G(x_Q)-DG(x_Q)(x-x_Q)|\le C_n\eta r
 \quad(x\in K\cap Q).
\]
Ponieważ ranga $DG(x_Q)$ jest mniejsza niż $m$,
zbiór $G(K\cap Q)$ zawiera się w
$C_n\eta r$-otoczeniu afinicznej podprzestrzeni
wymiaru $j=\operatorname{rank}DG(x_Q)\le m-1$,
w obrębie kuli o promieniu $C_n(M+1)r$, gdzie
$M$ ogranicza normę $DG$ na wybranym zwartym otoczeniu.
Po przesunięciu o $G(x_Q)$ i ortogonalnej zmianie
współrzędnych obraz jest zawarty w produkcie
$j$-wymiarowej kuli o promieniu $C_n(M+1)r$
i $(m-j)$-wymiarowej kuli o promieniu $C_n\eta r$.
Przy $j=0$ pierwszy czynnik jest jednopunktową
przestrzenią o zerowym wymiarze i mierze liczącej $1$.
Objętość produktu nie przekracza
$C_{n,m,M}r^j(\eta r)^{m-j}$, a
$\eta^{m-j}\le\eta$ dla $0<\eta\le1$.
Ponieważ $K\cap Q$ jest zwarty i $G$ ciągłe,
jego obraz jest zwarty, a zatem borelowski. Otrzymujemy
\[
 \lambda_m(G(K\cap Q))\le C_{n,m,M}\eta r^m.
\]
Dla każdego $y$ zbiór $K\cap G^{-1}(\{y\})$
można pokryć fragmentami leżącymi w tych
kostkach $Q$, dla których $y\in G(K\cap Q)$.
Każdy fragment ma średnicę nie większą niż
$\sqrt n\,r$, więc dla $\delta_r=2\sqrt n\,r$
\[
 \mathcal H^d_{\delta_r}(K\cap G^{-1}(\{y\}))
 \le C_{n,m}\sum_Q r^d\chi_{G(K\cap Q)}(y).
\]
W sumie występuje najwyżej $C_K r^{-n}$ kostek
dla $0<r\le1$, ponieważ $K$ leży w ustalonej
ograniczonej kostce. Oznaczmy prawą stronę
nierówności przez $b_r(y)$. Jest to nieujemna
funkcja borelowska. Z oszacowania objętości
obrazów i z równości $d+m=n$ wynika
\[
 \int_{\mathbb R^m}b_r\,\mathrm d\lambda_m
 \le C_{n,m,M,K}r^{-n}r^d\eta r^m
 =C_{n,m,M,K}\eta.
\]
Ustalmy teraz skalę $\delta>0$ i próg $t>0$.
Dla każdego $0<\eta\le1$ możemy dobrać $r$
tak małe, że powyższe oszacowania obowiązują
oraz $\delta_r\le\delta$.
Z monotoniczności zawartości względem skali mamy
$\mathcal H^d_\delta(K\cap G^{-1}(\{y\}))
\le\mathcal H^d_{\delta_r}(K\cap G^{-1}(\{y\}))
\le b_r(y)$.
W konsekwencji
\[
 \lambda_m^*\{y:\mathcal H^d_\delta
 (K\cap G^{-1}(\{y\}))>t\}
 \le\lambda_m(\{b_r>t\})
 \le C_{n,m,M,K}\eta/t.
\]
Stała po prawej nie zależy od $\eta,r,\delta,t$.
Przechodząc $\eta\downarrow0$, otrzymujemy
zewnętrzną miarę zero dla ustalonych $\delta,t$.
Bierzemy przeliczalną sumę tych zbiorów dla
$\delta=1/j$ oraz $t=1/l$, $j,l\ge1$.
Poza nią wszystkie zawartości
$\mathcal H^d_{1/j}(K\cap G^{-1}(\{y\}))$ są zerowe.
Ich granica dla $j\to\infty$ jest z definicji
miarą $\mathcal H^d$ włókna i również wynosi zero.
To daje tezę bez całkowania funkcji o nieustalonej
mierzalności.
\end{proof}


\section{Wzory area i coarea dla odwzorowań Lipschitzowskich}
W poniższych wzorach $DF$ oznacza różniczkę z twierdzenia
Rademachera, istniejącą prawie wszędzie. Jakobian $J_kF$
oraz jakobian normalny $\operatorname{Jac}_mF$ obliczamy
w punktach różniczkowalności, podstawiając $DF(x)$
do funkcji jakobianowych z
Definicji~\ref{def:k-jacobian-28} i
Definicji~\ref{def:normal-jacobian-28},
a na pozostałym zbiorze zerowym przypisujemy im wartość $0$.
Zbiór punktów różniczkowalności jest borelowski
na mocy dowodu Twierdzenia~\ref{thm:rademacher-28},
a współczynniki $DF$ są na nim granicami borelowskich
ilorazów różnicowych w kierunkach wektorów bazowych.
Funkcje jakobianowe macierzy są ciągłe, więc opisane
jakobiany są borelowskie. Wybór wartości na zbiorze
wyjątkowym nie zmienia całek. W całkach przyjmujemy
$0\cdot(+\infty)=0$. Sumę nieujemnych liczb po
dowolnym, także nieprzeliczalnym zbiorze indeksów,
rozumiemy jako supremum sum po jego skończonych
podzbiorach (suma pusta wynosi zero).
\begin{twierdzenie}[Wzory area i coarea dla odwzorowań Lipschitzowskich]
\label{thm:lipschitz-area-coarea-28}
Niech $1\le k\le N$, $E\subseteq\mathbb R^k$ będzie
mierzalny w sensie Lebesgue'a, a $F:\mathbb R^k\to\mathbb R^N$
będzie lokalnie Lipschitzowski.
Wówczas różniczka $DF(x)$ istnieje dla prawie
wszystkich $x$ (twierdzenie Rademachera)
i dla każdej funkcji $f:E\to[0,+\infty]$
mierzalnej w sensie Lebesgue'a zachodzi
\[
 \int_E f(x)J_kF(x)\,\mathrm d\lambda_k(x)
 =\int_{\mathbb R^N}\left(\sum_{x\in E\cap F^{-1}(\{y\})}
 f(x)\right)\mathrm d\mathcal H^k(y).
\]
Funkcja pod całką po $\mathbb R^N$ jest mierzalna
względem zupełnej miary Hausdorffa.
Ponadto jeżeli $n\ge m\ge1$ i
$F:\mathbb R^n\to\mathbb R^m$ jest lokalnie
Lipschitzowski, to dla każdej funkcji
$g:\mathbb R^n\to[0,+\infty]$ mierzalnej
w sensie Lebesgue'a jej ograniczenie do włókna
jest $\mathcal H^{n-m}$-mierzalne dla
$\lambda_m$-prawie wszystkich $y$ i zachodzi
\[
 \int_{\mathbb R^n}g(x)\operatorname{Jac}_mF(x)
 \,\mathrm d\lambda_n(x)
 =\int_{\mathbb R^m}\left(
 \int_{F^{-1}(\{y\})}g(x)
 \,\mathrm d\mathcal H^{n-m}(x)\right)
 \,\mathrm d\lambda_m(y).
\]
Wewnętrzna całka określa funkcję mierzalną względem
miary Lebesgue'a po przypisaniu dowolnej wartości
nieujemnej na ewentualnym zbiorze wyjątkowym miary zero.
W przypadku $n=m$ wewnętrzna całka po włóknie
jest sumą względem miary liczącej.
\end{twierdzenie}

\begin{lemat}[Wzór area dla dowolnego odwzorowania $C^1$]
\label{lem:c1-general-area-28}
Niech $1\le k\le N$, $V\subseteq\mathbb R^k$ będzie
otwarty, $G\in C^1(V;\mathbb R^N)$, a $A\subseteq V$
borelowski. Dla każdej nieujemnej funkcji borelowskiej
$f:A\to[0,+\infty]$ zachodzi
\[
 \int_Af(x)J_kG(x)\,\mathrm d\lambda_k(x)
 =\int_{\mathbb R^N}\left(
 \sum_{x\in A\cap G^{-1}(\{y\})}f(x)\right)
 \,\mathrm d\mathcal H^k(y).
\]
\end{lemat}
\begin{proof}
Zbiór $V_{\mathrm{reg}}=\{x\in V:\operatorname{rank}DG(x)=k\}$
jest otwarty, a $G|_{V_{\mathrm{reg}}}$ jest immersją.
Stosujemy Twierdzenie~\ref{thm:immersion-area-28}
do $G|_{V_{\mathrm{reg}}}$, zbioru borelowskiego
$A\cap V_{\mathrm{reg}}$ i ograniczenia $f$ do tego zbioru.
Jego dopełnienie w $V$ jest względnie domkniętym
zbiorem punktów krytycznych, który można wyczerpać
przeliczalną rodziną zwartych podzbiorów.
Istotnie, wystarczy przeciąć zbiór krytyczny ze
zwartymi zbiorami wyczerpującymi $V$, np. kulami
domkniętymi o promieniach $l$ ograniczonymi dodatkowo
warunkiem $\operatorname{dist}(x,\mathbb R^k\setminus V)
\ge1/l$, z pominięciem tego warunku gdy $V=\mathbb R^k$.
Każdy taki zbiór jest zwarty i spełnia hipotezy
Lematu~\ref{lem:critical-area-null-28}. Podaddytywność
pokazuje, że obraz całego zbioru krytycznego ma
$\mathcal H^k$-miarę zero.
Na części krytycznej $J_kG=0$, a suma po włóknach
może być różna od zera tylko na wspomnianym
zbiorze zerowym. Jest ona mierzalna względem zupełnej
miary Hausdorffa, ponieważ każdy podzbiór zbioru
zerowego jest mierzalny. Obie całki części krytycznej
wynoszą więc zero. Dodając równość dla części
regularnej, otrzymujemy żądany wzór.
\end{proof}

\begin{proof}[Dowód ogólnych wzorów area i coarea z
Twierdzenia~\ref{thm:lipschitz-area-coarea-28}]
\emph{Krok 1: wzór area.}
Z Lematu~\ref{lem:lusin-whitney-28} otrzymujemy
zwarte zbiory $K_j\subseteq\mathbb R^k$ oraz
odwzorowania $F_j\in C^1(\mathbb R^k;\mathbb R^N)$
takie, że $F_j=F$ i $DF_j=DF$ na $K_j$, a
$Z=\mathbb R^k\setminus\bigcup_jK_j$ jest
$\lambda_k$-zerowy. Połóżmy
$A_1=K_1$ i
$A_j=K_j\setminus\bigcup_{i<j}K_i$ dla $j\ge2$.
Są to parami rozłączne zbiory borelowskie,
których suma jest dopełnieniem $Z$.

Dla $E$ i $f$ borelowskich stosujemy
Lemat~\ref{lem:c1-general-area-28} do
$F_j$ na borelowskim zbiorze $A_j\cap E$
z ograniczeniem $f$ do tego zbioru. Ponieważ
$A_j\subseteq K_j$, dla każdego $x\in A_j$ mamy
$J_kF_j(x)=J_kF(x)$ oraz $F_j(x)=F(x)$.
Zatem dla każdego $y$ zachodzi
$A_j\cap E\cap F_j^{-1}(\{y\})
=A_j\cap E\cap F^{-1}(\{y\})$.
Podstawiając obie równości do wzoru z lematu, dostajemy
\[
 \int_{A_j\cap E}f(x)J_kF(x)\,\mathrm d\lambda_k(x)
 =\int_{\mathbb R^N}
 \sum_{x\in A_j\cap E\cap F^{-1}(\{y\})}f(x)
 \,\mathrm d\mathcal H^k(y).
\]
Po lewej stronie sumujemy całki po rozłącznych
zbiorach $A_j\cap E$. Po prawej każda funkcja
sumy po włóknie jest nieujemna i mierzalna
względem zupełnej miary Hausdorffa na mocy dowodu
lematu. Dla ustalonego $y$ połóżmy
$T_j=A_j\cap E\cap F^{-1}(\{y\})$.
Zbiory $T_j$ są rozłączne. Każdy skończony podzbiór
ich sumy przecina tylko skończenie wiele $T_j$,
więc jego suma wartości $f$ nie przekracza
$\sum_j\sum_{x\in T_j}f(x)$.
Biorąc supremum po skończonych podzbiorach, otrzymujemy
jedną nierówność. Odwrotnie, dla ustalonego $l$
i skończonych podzbiorów $D_j\subseteq T_j$,
$j=1,\ldots,l$, suma wartości $f$ na ich rozłącznej
sumie nie przekracza sumy po całym włóknie
z pominięciem $Z$. Biorąc kolejno supremum po
$D_1,\ldots,D_l$, a następnie po $l$, otrzymujemy
nierówność odwrotną. To dowodzi równości sum
także wtedy, gdy któraś ma wartość $+\infty$. Stosujemy twierdzenie
o zbieżności monotonicznej do sum częściowych,
które nie wymaga $\sigma$-skończoności $\mathcal H^k$
na $\mathbb R^N$.
Zbiór $Z$ nie wnosi nic do lewej strony.
Ponieważ $F$ jest lokalnie Lipschitzowskie,
wybieramy przeliczalne pokrycie $\mathbb R^k$
otwartymi kulami $U_i$, na których $F$ jest
Lipschitzowskie ze stałą $L_i>0$. Dla każdego $i$
kalibracja daje $\mathcal H^k(Z\cap U_i)=0$,
a Twierdzenie~\ref{thm:hausdorff-lipschitz-26} daje
$\mathcal H^k(F(Z\cap U_i))\le
L_i^k\mathcal H^k(Z\cap U_i)=0$.
Podaddytywność daje $\mathcal H^k(F(Z))=0$.
Zatem również prawa strona nie zmienia się
po dołączeniu punktów z $Z$.

Jeżeli $E$ i $f$ są jedynie mierzalne w sensie
Lebesgue'a, przedłużamy $f$ przez zero poza $E$
do funkcji $u:\mathbb R^k\to[0,+\infty]$.
Wybieramy borelowski reprezentant $u_0$ funkcji $u$
równy jej poza zbiorem $\lambda_k$-zerowym $D$.
Stosujemy dotychczasowy wzór do $E=\mathbb R^k$
i funkcji $u_0$. Całki po dziedzinie dla $u,u_0$
są równe. Jak wykazano wyżej, $\mathcal H^k(F(D))=0$.
Dla $y\notin F(D)$ wszystkie wartości $u$ i $u_0$
na włóknie są identyczne, a zerowe wartości poza $E$
nie wpływają na sumę. Zupełność miary Hausdorffa
zapewnia mierzalność rzeczywistej funkcji sumy
i równość całek po $\mathbb R^N$.

\emph{Krok 2: wzór coarea na częściach pełnej rangi.}
Przyjmijmy na razie, że $g$ jest borelowska.
Dla odwzorowania $F:\mathbb R^n\to\mathbb R^m$
stosujemy Lemat~\ref{lem:lusin-whitney-28},
otrzymując zwarte zbiory $K_j\subseteq\mathbb R^n$
i odwzorowania $F_j\in C^1(\mathbb R^n;\mathbb R^m)$
takie, że $F_j=F$ i $DF_j=DF$ na $K_j$, a
$\lambda_n(\mathbb R^n\setminus\bigcup_jK_j)=0$.
Kładziemy $A_j=K_j\setminus\bigcup_{i<j}K_i$,
więc są to parami rozłączne zbiory borelowskie.
Niech
$V_j=\{x:\operatorname{rank}DF_j(x)=m\}$.
Jest to zbiór otwarty, a $F_j|_{V_j}$ jest
submersją. Twierdzenie~\ref{thm:coarea-submersion-28}
stosujemy do $F_j|_{V_j}$ i nieujemnej borelowskiej
funkcji $(g\chi_{A_j})|_{V_j}$. Po lewej stronie
na $A_j\cap V_j$ jakobiany $F_j$ i $F$ są równe,
gdyż $DF_j=DF$ na $K_j$. Po prawej stronie
funkcja znika poza $A_j$, a na tym zbiorze
$F_j=F$, więc
$A_j\cap V_j\cap F_j^{-1}(\{y\})
=A_j\cap V_j\cap F^{-1}(\{y\})$.
Wzór dla submersji daje zatem
\begin{align*}
 &\int_{A_j\cap V_j}g(x)\operatorname{Jac}_mF(x)
 \,\mathrm d\lambda_n(x)\\
 &\quad=\int_{\mathbb R^m}\left(
 \int_{A_j\cap V_j\cap F^{-1}(\{y\})}
 g(x)\,\mathrm d\mathcal H^{n-m}(x)\right)
 \,\mathrm d\lambda_m(y).
\end{align*}

\emph{Krok 3: części krytyczne i zbiór wyjątkowy.}
Zbiór
$C_j=K_j\cap\{x:\operatorname{rank}DF_j(x)<m\}$
jest zwarty, ponieważ warunek rangi mniejszej niż $m$
jest domknięty (wszystkie minory rzędu $m$ znikają).
Na $C_j\subseteq K_j$ mamy $DF_j=DF$, więc
$\operatorname{Jac}_mF=0$. Lemat~\ref{lem:critical-coarea-null-28}
stosujemy do odwzorowania $F_j$ i zbioru $C_j$:
wszystkie punkty $C_j$ są krytyczne dla $F_j$.
Ponieważ $F_j=F$ na $C_j$, otrzymujemy
$\mathcal H^{n-m}(C_j\cap F^{-1}(\{y\}))=0$
dla prawie wszystkich $y$. Na pozostałym zbiorze
$Z=\mathbb R^n\setminus\bigcup_jK_j$ jakobian
również nie wnosi nic do całki, bo
$\lambda_n(Z)=0$. Z kolei
Lemat~\ref{lem:null-fibres-lipschitz-28} dowodzi,
że $Z$ jest $\mathcal H^{n-m}$-zerowy na
prawie wszystkich włóknach. Bierzemy przeliczalną sumę zbiorów wyjątkowych
dla wszystkich $C_j$ i dla $Z$. Poza nią brakujące
części włókna mają zerową $\mathcal H^{n-m}$-miarę.
Dla borelowskiego $g$ ograniczenia do włókien są
mierzalne, ponieważ $F$ jest ciągłe, włókna domknięte,
a miara Hausdorffa mierzy zbiory borelowskie.
Dla każdego $y$ spoza zbioru wyjątkowego
całka po włóknie jest sumą całek po rozłącznych
$A_j\cap V_j\cap F^{-1}(\{y\})$.
Każda funkcja tej lokalnej całki jest mierzalna
względem $\lambda_m$ na mocy dowodu wzoru
dla submersji (Tonelli i rozłączny podział).
Suma jest mierzalna, a rzeczywista całka po włóknie
różni się od niej co najwyżej na zbiorze zerowym;
stosujemy zupełność miary Lebesgue'a.
Po obu stronach sumujemy równości z kroku 2
przez twierdzenie o zbieżności monotonicznej.
Pominięte części dają całki zerowe, co dowodzi wzoru.

W przypadku funkcji $g$ jedynie mierzalnej
w sensie Lebesgue'a zastępujemy ją wersją
borelowską poza zbiorem zerowym. Niech $D$ będzie zbiorem zerowym, poza którym
$g$ i jego borelowski reprezentant $g_0$ są równe.
Lemat~\ref{lem:null-fibres-lipschitz-28} daje
$\mathcal H^{n-m}(D\cap F^{-1}(\{y\}))=0$
dla prawie wszystkich $y$. Zupełność miary Hausdorffa
zapewnia dla tych $y$ mierzalność ograniczenia $g$
do włókna i równość jego całki z całką $g_0$.
Na pozostałych włóknach całce możemy przypisać
dowolną nieujemną wartość; zupełność miary Lebesgue'a
zapewnia mierzalność tak otrzymanej funkcji i równość
jej całki z całką odpowiadającą $g_0$.
Nie trzeba w żadnym kroku zakładać, że odwzorowanie
$F$ jest globalnie różnowartościowe lub że ma
stałą rangę.
\end{proof}

\begin{uwaga}[Zakres uzyskanego dowodu]
W porównaniu z wersjami gładkimi nowym elementem
jest jednoczesna kontrola zbiorów krytycznych i ich
włókien. Lemat~\ref{lem:whitney-c1-28} wyprowadza
potrzebne przedłużenie $C^1$, a lematy
\ref{lem:critical-area-null-28}--\ref{lem:critical-coarea-null-28}
wyprowadzają oszacowania bez powoływania się na
same wzory area/coarea. Dzięki temu dowód nie jest
błędnym kołem. Nadal poza zakresem pozostają
analogiczne wzory na ogólnych przestrzeniach
metrycznych oraz dla odwzorowań, które nie są
lokalnie Lipschitzowskie.
\end{uwaga}

\section{Postać na rozmaitościach}
\begin{twierdzenie}[Wzór koobszaru na rozmaitościach riemannowskich]
Niech $(M^n,g)$ i $(N^m,h)$ będą gładkimi rozmaitościami
riemannowskimi bez brzegu, Hausdorffa, spełniającymi drugi
aksjomat przeliczalności, wymiarów $n\ge m\ge1$.
Niech $F:M\to N$ będzie submersją klasy $C^1$, tj.
$\mathrm dF_p:T_pM\to T_{F(p)}N$ jest surjekcją dla każdego
$p\in M$. Przez $\mathrm dF_p^*$ oznaczamy sprzężenie
różniczki względem iloczynów skalarnych $g_p,h_{F(p)}$.
Jakobian normalny definiujemy jako
$\operatorname{Jac}_{g,h}F(p)=
\sqrt{\det(\mathrm dF_p\mathrm dF_p^*)}$,
gdzie wyznacznik liczymy na $T_{F(p)}N$.
Dla każdej funkcji borelowskiej $\varphi:M\to[0,+\infty]$
funkcja wewnętrznej całki w poniższym wzorze jest
borelowska i zachodzi
\begin{equation}\label{eq:coarea-manifolds-28}
 \int_M\varphi(p)\operatorname{Jac}_{g,h}F(p)
 \,\mathrm d\mu_g(p)
 =\int_N\left(\int_{F^{-1}(\{q\})}\varphi(p)
 \,\mathrm d\mu_{g|_{F^{-1}(\{q\})}}(p)\right)
 \,\mathrm d\mu_h(q).
\end{equation}
Na niepustym włóknie, będącym osadzoną podrozmaitością
klasy $C^1$, miarę indukowaną definiujemy za pomocą
ciągłej dodatnio określonej macierzy metryki w lokalnych
parametryzacjach klasy $C^1$. Dla $n=m$ jest to miara
licząca, a całka po pustym włóknie wynosi zero.
\end{twierdzenie}
\begin{proof}
\emph{Krok 1: współrzędne dostosowane do submersji.}
Ustalmy $p_0\in M$ i mapy $u$ wokół $p_0$ oraz $v$ wokół
$F(p_0)$. Ograniczamy dziedzinę $u$, aby jej obraz przez $F$
mieścił się w dziedzinie $v$. Odwzorowanie
$T=v\circ F\circ u^{-1}$ jest klasy $C^1$, a jego różniczka
w $u(p_0)$ ma rangę $m$. Po permutacji współrzędnych
źródła piszemy $a=(a',a'')\in\mathbb R^m\times\mathbb R^{n-m}$
tak, aby minor $D_{a'}T$ był odwracalny w $u(p_0)$.
Różniczka odwzorowania $a\mapsto(T(a),a'')$
ma wtedy postać blokowo trójkątną z odwracalnymi
blokami diagonalnymi. Z twierdzenia o odwzorowaniu
odwrotnym otrzymujemy jego lokalną odwrotność
$\Psi:W\to u(V)$ klasy $C^1$, gdzie $W$ jest otwarty
w $\mathbb R^m\times\mathbb R^{n-m}$, a $V$ jest
otwartym otoczeniem $p_0$ w $M$.
Kładziemy $P=u^{-1}\circ\Psi:W\to V$.
Dla $(y,z)\in W$ mamy $v(F(P(y,z)))=y$.
W szczególności dla $q$ z dziedziny mapy $v$,
przy $y=v(q)$ i $W_y=\{z:(y,z)\in W\}$,
odwzorowanie $z\mapsto P(y,z)$ parametryzuje
$V\cap F^{-1}(\{q\})$. Gdy $W_y$ jest puste,
ten fragment włókna jest pusty.

\emph{Krok 2: tożsamość macierzy metryk.}
W punkcie $(y,z)\in W$ niech $B$ oznacza macierz
metryki $g$ w parametryzacji $P$, czyli
$B=D\Psi^{\mathsf T}G_u(\Psi)D\Psi$, z notacją
macierzy $G_u$ z \eqref{eq:g-density-27}.
Niech $H=H_v(y)$ będzie macierzą metryki $h$
w mapie $v$. Piszemy
$B=\begin{psmallmatrix}B_{11}&B_{12}\\B_{21}&C\end{psmallmatrix}$,
gdzie $C$ jest blokiem rozmiaru $d\times d$, $d=n-m$.
Dla $d>0$ jest on dodatnio określony i jest właśnie
macierzą metryki indukowanej w parametryzacji włókna.
W tych współrzędnych różniczka $F$ ma macierz
$L=(I_m\ \ 0)$. Macierz jej sprzężenia wynosi
$B^{-1}L^{\mathsf T}H$: dla dowolnych wektorów
$a\in\mathbb R^n,b\in\mathbb R^m$ zachodzi
$(La)^{\mathsf T}Hb=a^{\mathsf T}B(B^{-1}L^{\mathsf T}Hb)$.
Zatem kwadrat jakobianu normalnego wynosi
$\det(LB^{-1}L^{\mathsf T})\det H$.

Dla $d>0$ kładziemy $S=B_{11}-B_{12}C^{-1}B_{21}$.
Mnożąc blokowo, otrzymujemy
\[
 \begin{pmatrix}I_m&-B_{12}C^{-1}\\0&I_d\end{pmatrix}
 B
 \begin{pmatrix}I_m&0\\-C^{-1}B_{21}&I_d\end{pmatrix}
 =\begin{pmatrix}S&0\\0&C\end{pmatrix}.
\]
Macierze mnożące mają wyznacznik $1$, skąd
$\det B=\det S\det C$. Dodatnia określoność $B$
i powyższa kongruencja dają dodatnią określoność $S$.
Jeśli $R$ oznacza prawą macierz mnożącą,
to lewa macierz mnożąca jest równa $R^{\mathsf T}$.
Zatem $R^{\mathsf T}BR=\operatorname{diag}(S,C)$
i po odwróceniu oraz pomnożeniu przez $R$ i
$R^{\mathsf T}$ dostajemy
\[
 B^{-1}=R\operatorname{diag}(S^{-1},C^{-1})R^{\mathsf T}.
\]
Pierwszy wiersz blokowy $R$ jest równy $(I_m\ \ 0)$,
więc lewy górny blok $B^{-1}$ jest równy $S^{-1}$. Ponieważ $L$ wybiera
ten blok, mamy $\det(LB^{-1}L^{\mathsf T})=
1/\det S=\det C/\det B$.
Dla $d=0$ mamy $L=I_m$, a przy konwencji $\det C=1$
ta ostatnia równość również obowiązuje. Stąd w obu przypadkach
\[
 \operatorname{Jac}_{g,h}F(P(y,z))\sqrt{\det B}
 =\sqrt{\det H}\sqrt{\det C}.
\]
Wszystkie macierze metryk są ciągłe i dodatnio określone;
w szczególności jakobian normalny jest dodatni i ciągły.

\emph{Krok 3: całkowanie lokalne.}
Niech $\varphi$ znika poza $V$. Współrzędnościowy wzór
na $\mu_g$ i zamiana zmiennych przez dyfeomorfizm
$\Psi$ klasy $C^1$ dają gęstość
$\sqrt{\det G_u(\Psi)}|\det D\Psi|=\sqrt{\det B}$
w parametryzacji $P$. Analogicznie wzór na miarę
indukowaną we włóknie daje gęstość $\sqrt{\det C}$.
Ta definicja jest niezależna od parametryzacji:
przy zmianie parametrów przez dyfeomorfizm klasy $C^1$
macierz metryki zmienia się z $C$ na $D^{\mathsf T}CD$,
gdzie $D$ jest macierzą różniczki zmiany parametrów
(z $C$ obliczanym w odpowiednim punkcie).
Pierwiastek z wyznacznika mnoży się przez $|\det D|$,
a euklidesowa zamiana zmiennych usuwa tę zmianę.
Dla $d=0$ nie występuje zmiana gęstości: każdy punkt
ma miarę $1$.

Dla $d>0$ definiujemy na $W$ nieujemną funkcję borelowską
\[
 b(y,z)=\varphi(P(y,z))\sqrt{\det C(y,z)}.
\]
Przedłużamy $b$ przez zero na całe $\mathbb R^m\times
\mathbb R^d$. Korzystając z tożsamości z kroku 2
i twierdzenia Tonellego dla $b$ (z czynnikiem
$\sqrt{\det H_v(y)}$ w całce po $y$), otrzymujemy
\begin{align*}
 &\int_V\varphi(p)\operatorname{Jac}_{g,h}F(p)\,\mathrm d\mu_g(p)\\
 &\quad=\int_W\varphi(P(y,z))\sqrt{\det H_v(y)}
      \sqrt{\det C(y,z)}\,\mathrm d\lambda_n(y,z)\\
 &\quad=\int_{v(\operatorname{dom} v)}
       \sqrt{\det H_v(y)}
       \left(\int_{W_y}\varphi(P(y,z))
              \sqrt{\det C(y,z)}\,\mathrm d\lambda_d(z)\right)
       \mathrm d\lambda_m(y).
\end{align*}
Wewnętrzna całka jest borelowską funkcją $y$, gdyż
jest równa $\int_{\mathbb R^d}b(y,z)\,\mathrm d\lambda_d(z)$,
a $b$ jest borelowska i nieujemna oraz $\lambda_d$
jest $\sigma$-skończoną miarą borelowską.
Współrzędnościowy wzór na $\mu_h$ i definicja miary
indukowanej identyfikują ostatnią całkę z prawą stroną
\eqref{eq:coarea-manifolds-28} dla funkcji znikającej
poza $V$. Gdy $d=0$, $P$ jest lokalną odwrotnością
$F$ zapisaną w mapie $v$, a każdy fragment włókna w $V$
jest pusty albo jednopunktowy. Ta sama tożsamość
wyznaczników i zamiana zmiennych dają równość,
w której całka wewnętrzna jest wartością funkcji
w jedynym punkcie fragmentu albo zerem.

\emph{Krok 4: przejście do całego $M$.}
Z otoczeń $V$ skonstruowanych w kroku 1 wybieramy
przeliczalne podpokrycie $(V_j)_{j\ge1}$, korzystając
z drugiego aksjomatu przeliczalności. Zbiory
$B_j=V_j\setminus\bigcup_{i<j}V_i$ są borelowskie,
parami rozłączne i pokrywają $M$.
Stosujemy lokalny wzór do $\varphi\chi_{B_j}$.
Dla każdego $q\in N$ zbiory $B_j\cap F^{-1}(\{q\})$
tworzą rozłączny borelowski podział włókna, więc
całka po włóknie jest sumą lokalnych całek.
Każda funkcja lokalnej całki jest borelowska
(równa zero poza odpowiednią dziedziną mapy celu),
a ich suma również jest borelowska.
Twierdzenie o zbieżności monotonicznej pozwala
zsumować lokalne równości po obu stronach.
Otrzymujemy \eqref{eq:coarea-manifolds-28}.
Nie wykorzystaliśmy orientowalności.
\end{proof}

\begin{wniosek}[Miara riemannowska i Hausdorffa]
\label{cor:riemann-hausdorff-28}
Niech $(M^n,g)$ będzie gładką rozmaitością riemannowską
bez brzegu, Hausdorffa, spełniającą drugi aksjomat
przeliczalności, wymiaru $n\ge1$. Na każdej składowej
spójności $M$ niech $d_g$ oznacza odległość riemannowską,
a $\mathcal H^n_{d_g}$ miarę Hausdorffa skonstruowaną
z tej odległości przy kalibracji z
Uwagi~\ref{rem:hausdorff-convention-26}.
Na zbiorach borelowskich zachodzi
$\mu_g=\mathcal H^n_{d_g}$ (dla rozmaitości niespójnej
miary sumujemy po składowych spójności).
Jeżeli $S^k\subseteq M$ jest osadzoną gładką
podrozmaitością bez brzegu, $0\le k\le n$, to
$\mu_{g|_S}=\mathcal H^k_{d_g|_{S\times S}}$
na zbiorach borelowskich $S$, również z sumowaniem
po składowych spójności $M$. Dla $k=0$ obie
miary na $S$ rozumiemy jako miary liczące.
\end{wniosek}
\begin{proof}
Wystarczy pracować na spójnej składowej $M$.
Składowe są otwarte, ponieważ rozmaitość jest
lokalnie spójna drogowo; jest ich co najwyżej
przeliczalnie wiele, gdyż każda zawiera inny element
przeliczalnej bazy. Odległość $d_g$ na każdej
składowej jest infimum długości odcinkami gładkich
krzywych łączących dane punkty.
Najpierw rozpatrujemy $1\le k\le n$, przy czym
przypadek $S=M$ odpowiada $k=n$.

\emph{Krok 1: lokalne porównanie odległości.}
Ustalmy $p\in S$ i $0<\varepsilon<1$. Wybieramy gładką
mapę $u$ otoczenia $p$ w $M$ taką, że $u(p)=0$
oraz $S$ w tej mapie odpowiada przecięciu jej obrazu
z podprzestrzenią $E=\mathbb R^k\times\{0\}$.
Gdy $S=M$, bierzemy $E=\mathbb R^n$.
Niech $G(a)=G_u(a)$ będzie macierzą metryki,
$G_0=G(0)$ oraz $|v|_0=\sqrt{v^{\mathsf T}G_0v}$.
Z ciągłości macierzy i dodatniej określoności $G_0$
wybieramy $R>0$, tak aby domknięta kula
$\{a:|a|_0\le R\}$ była zawarta w obrazie mapy
i dla $|a|_0\le R$, $v\in\mathbb R^n$ zachodziło
\[
 (1-\varepsilon)^2|v|_0^2
 \le v^{\mathsf T}G(a)v
 \le(1+\varepsilon)^2|v|_0^2.
\]
Istotnie, na zwartej sferze $|v|_0=1$ różnica
$v^{\mathsf T}(G(a)-G_0)v$ dąży jednostajnie do zera
dla $a\to0$; dla dowolnych $v$ stosujemy jednorodność.
Niech $0<r<R/3$ i $U_r=u^{-1}(\{a:|a|_0<r\})$.

Dla $x,y\in U_r$ odcinek między $u(x)$ i $u(y)$
pozostaje w kuli o promieniu $r$, ponieważ jest ona
wypukła. Jego przeciwobraz przez $u$ jest krzywą
łączącą $x,y$ o długości nie większej niż
$(1+\varepsilon)|u(x)-u(y)|_0$.
Z definicji infimum wynika takie samo górne
oszacowanie $d_g(x,y)$.

Aby otrzymać dolne oszacowanie, rozważamy dowolną
odcinkami gładką krzywą $\gamma$ łączącą $x,y$.
Jeśli pozostaje w przeciwobrazie kuli o promieniu
$R$, to jej długość jest co najmniej
$(1-\varepsilon)$ razy długość krzywej $u\circ\gamma$
w normie $|\cdot|_0$, a ta długość nie jest mniejsza
niż $|u(x)-u(y)|_0$ (nierówność trójkąta dla
całki z prędkości). Jeśli krzywa opuszcza tę kulę,
jej początkowy odcinek do pierwszego punktu na
sferze o promieniu $R$ leży w dziedzinie mapy:
domknięta kula jest w niej zawarta i jej przeciwobraz
jest zwarty. Długość tego odcinka jest co najmniej
$(1-\varepsilon)(R-r)$, gdyż w normie $|\cdot|_0$
odległość punktu początkowego od sfery wynosi
co najmniej $R-r$. Ponieważ $R-r>2r\ge|u(x)-u(y)|_0$,
również taka krzywa spełnia potrzebne dolne
oszacowanie. Biorąc infimum po wszystkich krzywych,
otrzymujemy
\[
 (1-\varepsilon)|u(x)-u(y)|_0
 \le d_g(x,y)
 \le(1+\varepsilon)|u(x)-u(y)|_0
 \quad(x,y\in U_r).
\]
Nie zakładamy istnienia krzywej minimalizującej.
Dla $x,y\in S\cap U_r$ górne oszacowanie pochodzi
nawet od krzywej w $S$, bo odcinek współrzędnych
leży w $E$. Dolne oszacowanie obowiązuje dla wszystkich
krzywych w $M$, a zatem dotyczy dokładnie odległości
zewnętrznej $d_g$ ograniczonej do $S$.

\emph{Krok 2: porównanie obu miar na małym otoczeniu.}
Przez $C(a)$ oznaczamy macierz ograniczenia $G(a)$
do $E$ w standardowej bazie $E$ i kładziemy $C_0=C(0)$.
Identyfikujemy $E$ z $\mathbb R^k$. Ograniczając
nierówność dla $G(a)$ do wektorów w $E$, otrzymujemy
oszacowanie wszystkich wartości własnych symetrycznej
macierzy $C_0^{-1/2}C(a)C_0^{-1/2}$ pomiędzy
$(1-\varepsilon)^2$ i $(1+\varepsilon)^2$.
Iloczyn tych wartości jest jej wyznacznikiem,
równym $\det C(a)/\det C_0$. Zatem
\[
 (1-\varepsilon)^k\sqrt{\det C_0}
 \le\sqrt{\det C(a)}
 \le(1+\varepsilon)^k\sqrt{\det C_0}.
\]
Wzór na gęstość miary indukowanej daje dla
borelowskiego $A\subseteq S\cap U_r$
\[
 (1-\varepsilon)^k D_A\le\mu_{g|_S}(A)
 \le(1+\varepsilon)^k D_A,\qquad
 D_A=\sqrt{\det C_0}\,\lambda_k(u(A)),
\]
gdzie w ostatniej mierze identyfikujemy $E$ z $\mathbb R^k$.
Dla $S=M$ jest to ten sam wzór z $C=G$.

Rozważmy na $E$ odległość $d_0(a,b)=|a-b|_0$.
Odwzorowanie $a\mapsto C_0^{1/2}a$ w bazie $E$
jest izometrią na zwykłą przestrzeń euklidesową.
Z definicji miary Hausdorffa wynika jej niezmienniczość
względem izometrii (pokrycia i średnice odpowiadają
sobie w obie strony). Kalibracja i liniowa zamiana
zmiennych dają więc $\mathcal H^k_{d_0}(u(A))=D_A$.
Miara Hausdorffa zbioru zawartego w podprzestrzeni
nie zależy od tego, czy pokrycia wybieramy w tej
podprzestrzeni, czy w przestrzeni otaczającej:
w drugim przypadku przecinamy każdy zbiór pokrycia
z podprzestrzenią, nie zwiększając średnic ani kosztów.
Zatem stosujemy Twierdzenie~\ref{thm:hausdorff-lipschitz-26}
do ograniczenia $u$ i jego odwrotności pomiędzy
$S\cap U_r$ z odległością $d_g$ oraz jego obrazem
z odległością $d_0$. Krok 1 daje odpowiednie stałe
Lipschitza $(1-\varepsilon)^{-1}$ i $1+\varepsilon$.
Otrzymujemy
\[
 (1-\varepsilon)^kD_A
 \le\mathcal H^k_{d_g|_{S\times S}}(A)
 \le(1+\varepsilon)^kD_A.
\]
Łącząc oba porównania (także dla $D_A=0$ lub
$D_A=+\infty$), dostajemy
\[
 \left(\frac{1-\varepsilon}{1+\varepsilon}\right)^k
 \mu_{g|_S}(A)
 \le\mathcal H^k_{d_g|_{S\times S}}(A)
 \le
 \left(\frac{1+\varepsilon}{1-\varepsilon}\right)^k
 \mu_{g|_S}(A).
\]

\emph{Krok 3: przejście do całej podrozmaitości.}
Dla ustalonego $\varepsilon$ otoczenia $S\cap U_r$
z kroku 1 pokrywają $S$. Podrozmaitość $S$ z topologią
podprzestrzeni spełnia drugi aksjomat przeliczalności,
więc wybieramy przeliczalne podpokrycie $(O_j)_{j\ge1}$.
Zbiory $O_j\setminus\bigcup_{i<j}O_i$ tworzą rozłączny
borelowski podział $S$. Obie miary mierzą zbiory
borelowskie w tej topologii: dla Hausdorffa wynika
to z konstrukcji miary na przestrzeni metrycznej
i faktu, że odległość $d_g$ indukuje topologię $M$
(co potwierdzają także lokalne porównania z kroku 1).
Dla każdego borelowskiego $A\subseteq S$ rozkładamy
$A$ według powyższego podziału i sumujemy nierówności
z kroku 2. Otrzymujemy te same globalne nierówności,
ze stałymi zależnymi jedynie od $k,\varepsilon$.
Jeśli $\mu_{g|_S}(A)<\infty$, przechodząc
$\varepsilon\downarrow0$, otrzymujemy równość obu miar.
Jeśli $\mu_{g|_S}(A)=+\infty$, już dolna nierówność
dla dowolnego $\varepsilon\in(0,1)$ daje nieskończoność
miary Hausdorffa. To kończy dowód dla $k\ge1$.
Dla $k=0$ osadzona podrozmaitość jest dyskretna
i przeliczalna. Miara indukowana nadaje każdemu
punktowi masę $1$, a $\mathcal H^0$ jest miarą
liczącą z definicji i wyniku z Rozdziału~\ref{ch:hausdorff-26}.
Obie miary są więc równe również w tym przypadku.
\end{proof}

\begin{uwaga}[Co przechodzi do geometrycznej teorii miary?]
Następny rozdział zastosuje miary Hausdorffa,
formuły area i coarea oraz twierdzenia o
różniczkowaniu miar do zbiorów prostowalnych
i zbiorów o skończonym obwodzie. Te obiekty
nie muszą być gładkimi podrozmaitościami,
a zatem stanowią właściwe rozszerzenie
geometrii różniczkowej.
\end{uwaga}

\begin{uwaga}[Dżety przybliżone a wzory miarowo-geometryczne]
Pochodna $DF(x)$ z twierdzenia Rademachera jest częścią
pierwszego dżetu odwzorowania $F$ w punktach różniczkowalności.
Przy osłabieniu pojęcia granicy do granicy przybliżonej otrzymujemy
\emph{dżety przybliżone}, określone miarowo prawie wszędzie.
W rozdziale~\ref{ch:dzety} rozwijamy geometryczną teorię wiązek
dżetów, omawiamy twierdzenie Whitneya o przedłużaniu danych
dżetowych oraz jego miarową wersję Whitneya--Łuzina.
Jest to uzupełnienie obecnego rozdziału, a nie ponowny dowód
wzorów area i coarea.
\end{uwaga}

% ROZDZIAŁ 29 -- podstawy geometrycznej teorii miary
\chapter[Prostowalność i zbiory o skończonym obwodzie]{Zbiory prostowalne, funkcje BV i zbiory o skończonym obwodzie}
\label{ch:geometric-measure-29}

\section{Cel i zakres teorii}
W geometrii różniczkowej badaliśmy gładkie podrozmaitości, ich
metryki indukowane oraz normalne. W tym rozdziale będziemy mierzyć
również brzegi zbiorów, które mają narożniki, pęknięcia lub
nieregularności i nie są podrozmaitościami. Zamiast zakładać
istnienie klasycznej normalnej w każdym punkcie, zdefiniujemy ją
\emph{miarowo}. Przygotujemy w ten sposób język geometrycznej
teorii miary, teorii wariacyjnej i równań różniczkowych cząstkowych.

Pracujemy na $\mathbb R^n$, gdzie $n\ge1$, z miarą Lebesgue'a
$\lambda_n$ i miarami Hausdorffa $\mathcal H^s$
z Rozdziału~\ref{ch:hausdorff-26}; dla wymiarów
całkowitych stosujemy dokładnie kalibrację
z Uwagi~\ref{rem:hausdorff-convention-26}.
W szczególności $\mathcal H^{n-1}$ ma tę samą
normalizację co we wzorze na pole z
Rozdziału~\ref{ch:area-coarea-28}, a gdy $n=1$,
$\mathcal H^0$ jest miarą liczącą. Oznaczenie $\chi_E$ oznacza
funkcję charakterystyczną $E$. Wszystkie przestrzenie $L^1$
są przestrzeniami klas funkcji równych prawie wszędzie.
Stosujemy przy tym konwencję dystrybucyjną
$\langle D_i u,\varphi\rangle=-\int_{\mathbb R^n}u(x)
\partial_i\varphi(x)\,\mathrm d\lambda_n(x)$ dla
$\varphi\in C_c^1(\mathbb R^n)$.

\section{Prostowalność: co pozostaje z gładkiej powierzchni?}
\begin{definicja}[Zbiory przeliczalnie prostowalne]
\label{def:rectifiable-29}
Niech $1\le k\le n$. Zbiór $E\subseteq\mathbb R^n$ nazywamy
\emph{przeliczalnie $k$-prostowalnym}, jeśli istnieją
zbiory $A_j\subseteq\mathbb R^k$ oraz odwzorowania
Lipschitzowskie $f_j:A_j\to\mathbb R^n$, $j\in\mathbb N$,
takie że
\[
 \mathcal H^k\!\left(E\setminus\bigcup_{j=1}^\infty f_j(A_j)\right)=0.
\]
Mówimy, że $E$ jest \emph{$k$-prostowalny o skończonej mierze},
jeżeli dodatkowo $\mathcal H^k(E)<\infty$.
Zbiór $P\subseteq\mathbb R^n$ jest \emph{czysto
$k$-nieprostowalny}, gdy
$\mathcal H^k(P\cap f(A))=0$ dla każdego
$A\subseteq\mathbb R^k$ i każdego odwzorowania
Lipschitzowskiego $f:A\to\mathbb R^n$.
\end{definicja}

Na dowolnych zbiorach symbol $\mathcal H^k$ oznacza
miarę zewnętrzną Hausdorffa; gdy korzystamy z addytywności,
jawnie zakładamy mierzalność w sensie Carathéodory'ego.
Dla $k=0$ zbiór nazywamy przeliczalnie $0$-prostowalnym,
jeśli jest co najwyżej przeliczalny; wtedy $\mathcal H^0$
jest miarą liczącą. Jest to zgodne z definicją dla
$\mathbb R^0=\{0\}$, ponieważ każdy obraz podzbioru
$\mathbb R^0$ zawiera co najwyżej jeden punkt,
a jedynym zbiorem o zerowej mierze liczącej jest zbiór pusty.

W tej definicji nie utożsamiamy prostowalności z topologiczną
gładkością. Funkcja Lipschitzowska może mieć punkty załamania,
a zbiór prostowalny może być przeliczalną sumą rozłącznych
fragmentów powierzchni. Warunek z miarą Hausdorffa pozwala
pominąć tylko zbiór $k$-wymiarowo zerowy, nie dowolny zbiór
mały w sensie miary $n$-wymiarowej.

\begin{przyklad}[Wykres z narożnikiem]
\label{ex:lip-graph-29}
Zbiór $E=\{(t,|t|):-1\le t\le1\}\subseteq\mathbb R^2$
jest $1$-prostowalny: to obraz przedziału przez odwzorowanie
$F(t)=(t,|t|)$. Dla $s,t\in[-1,1]$ mamy
$|F(t)-F(s)|^2=|t-s|^2+(|t|-|s|)^2\le2|t-s|^2$,
więc $F$ jest $\sqrt2$-Lipschitzowskie.
Gdyby $E$ było podrozmaitością klasy $C^1$ w otoczeniu
początku, lokalna regularna parametryzacja $\gamma$
z $\gamma(0)=0$ spełniałaby
$\gamma(h)=h\gamma'(0)+o(|h|)$, gdzie $\gamma'(0)\ne0$.
Wszystkie graniczne kierunki znormalizowanych siecznych
należałyby do prostej rozpiętej przez $\gamma'(0)$.
Tymczasem gałęzie $E$ mają niewspółliniowe kierunki
$(1,1)$ i $(-1,1)$, co daje sprzeczność. Mimo to $\mathcal H^1(E)$
jest skończona i dodatnia: rzut $E\to[-1,1]$ jest
$1$-Lipschitzowski, a odwzorowanie $F$ jest
$\sqrt2$-Lipschitzowskie, więc
$2\le\mathcal H^1(E)\le2\sqrt2$.
Każda z parametryzacji $t\mapsto(t,t)$ oraz
$t\mapsto(-t,t)$ dla $t\in[0,1]$ mnoży odległości
przez $\sqrt2$. Stosując Twierdzenie~\ref{thm:hausdorff-lipschitz-26}
do każdej parametryzacji i jej odwrotności, dostajemy
dla każdego odcinka miarę $\sqrt2\,\mathcal H^1([0,1])=\sqrt2$.
Oba odcinki są zwarte, a ich przecięcie jest singletonem
o zerowej $\mathcal H^1$-mierze. Zatem z addytywności
$\mathcal H^1(E)=2\sqrt2$.
\end{przyklad}

\begin{lemat}[Podstawowe własności prostowalności]
\label{lem:rectifiable-union-29}
Niech $1\le k\le n$. Wtedy:
\begin{enumerate}[label=(\roman*)]
\item przeliczalna suma zbiorów przeliczalnie
$k$-prostowalnych jest przeliczalnie $k$-prostowalna;
\item każdy podzbiór zbioru przeliczalnie
$k$-prostowalnego jest przeliczalnie $k$-prostowalny;
\item dla $N\ge k$, otwartego $U\subseteq\mathbb R^n$,
przeliczalnie $k$-prostowalnego $E\subseteq U$ i lokalnie
Lipschitzowskiego $G:U\to\mathbb R^N$ obraz $G(E)$
jest przeliczalnie $k$-prostowalny.
\end{enumerate}
\end{lemat}
\begin{proof}
W (i) niech $S_i=\bigcup_j f_{ij}(A_{ij})$ będzie
sumą obrazów odpowiadającą zbiorowi $E_i$. Wtedy
$(\bigcup_i E_i)\setminus(\bigcup_i S_i)
\subseteq\bigcup_i(E_i\setminus S_i)$.
Prawa strona ma zerową miarę zewnętrzną na mocy
podaddytywności, a rodzina obrazów jest przeliczalna.
W (ii) dla $E'\subseteq E$ i sumy obrazów $S$ mamy
$E'\setminus S\subseteq E\setminus S$, co daje tezę.

W (iii) niech $Z=E\setminus\bigcup_j f_j(A_j)$.
Lokalna własność Lipschitza i drugi aksjomat
przeliczalności dają przeliczalne pokrycie $U$
otwartymi zbiorami $U_i$, na których $G$ jest
Lipschitzowskie ze stałymi $L_i>0$.
Kładziemy $A_{ij}=\{x\in A_j:f_j(x)\in E\cap U_i\}$.
Każde złożenie $G\circ f_j|_{A_{ij}}$ jest określone
i Lipschitzowskie. Jeżeli $y=G(x)$ dla $x\in E\setminus Z$,
to $x$ należy do pewnego $U_i$ i jest wartością
pewnego $f_j$ na $A_{ij}$. Zatem
\[
 G(E)\setminus\bigcup_{i,j}(G\circ f_j)(A_{ij})
 \subseteq G(Z)\subseteq\bigcup_i G(Z\cap U_i).
\]
Twierdzenie~\ref{thm:hausdorff-lipschitz-26} daje
$\mathcal H^k(G(Z\cap U_i))\le L_i^k\mathcal H^k(Z\cap U_i)=0$.
Podaddytywność kończy dowód.
\end{proof}

\begin{twierdzenie}[Rozkład na część prostowalną i czysto nieprostowalną]
\label{thm:rectifiable-decomposition-29}
Niech $1\le k\le n$, a $E\subseteq\mathbb R^n$ będzie
$\mathcal H^k$-mierzalny i $\mathcal H^k(E)<\infty$.
Istnieją mierzalne, rozłączne zbiory $R,P$ o sumie $E$,
gdzie $R$ jest przeliczalnie $k$-prostowalny, a $P$
czysto $k$-nieprostowalny. Rozkład jest jednoznaczny
z dokładnością do zbiorów o zerowej $\mathcal H^k$-mierze.
\end{twierdzenie}
\begin{proof}
Niech $f:A\to\mathbb R^n$, $A\subseteq\mathbb R^k$,
będzie Lipschitzowskie ze stałą $L>0$.
Dla $x\in\overline A$ wybieramy ciąg $(x_i)$ w $A$
zbieżny do $x$. Nierówność
$|f(x_i)-f(x_j)|\le L|x_i-x_j|$ pokazuje, że
$(f(x_i))$ jest ciągiem Cauchy'ego, a zupełność
$\mathbb R^n$ daje granicę. Ta sama nierówność
dla dwóch ciągów zbieżnych do $x$ zapewnia niezależność
granicy od wyboru ciągu. Tak definiujemy przedłużenie
$\overline f$ na $\overline A$, nadal $L$-Lipschitzowskie
(przechodzimy do granicy dla ciągów zbieżnych do dwóch
dowolnych punktów domknięcia).
Wobec tego $f(A)$ zawiera się w przeliczalnej sumie
zwartych obrazów $\overline f(\overline A\cap[-l,l]^k)$,
$l\ge1$. Dla pustego $A$ wystarczy pusty obraz.

Niech $\mathcal S$ będzie rodziną przeliczalnych sum
obrazów zwartych podzbiorów $\mathbb R^k$ przez
odwzorowania Lipschitzowskie do $\mathbb R^n$.
Dopuszczamy zbiór pusty. Każdy element rodziny jest
borelowski i prostowalny, a rodzina jest zamknięta
na przeliczalne sumy. Kładziemy
$a=\sup_{S\in\mathcal S}\mathcal H^k(E\cap S)$.
Mamy $0\le a\le\mathcal H^k(E)<\infty$.
Wybieramy $S_j\in\mathcal S$ z
$\mathcal H^k(E\cap S_j)>a-1/j$ i kładziemy
$S_*=\bigcup_j S_j\in\mathcal S$.
Z monotoniczności oraz definicji supremum wynika
$\mathcal H^k(E\cap S_*)=a$.
Niech $R=E\cap S_*$, $P=E\setminus S_*$.
Są to rozłączne zbiory mierzalne o sumie $E$,
a $R$ jest prostowalny jako podzbiór $S_*$.

Gdyby $\mathcal H^k(P\cap f(A))>0$ dla pewnego
odwzorowania Lipschitzowskiego $f:A\to\mathbb R^n$,
to na mocy pierwszej części dowodu istniałby
$S\in\mathcal S$ zawierający $f(A)$.
Zatem $\mathcal H^k(P\cap S)>0$. Zbiory $R$ i $P\cap S$
są mierzalne, rozłączne, a ich suma to $E\cap(S_*\cup S)$.
Addytywność daje
\[
 \mathcal H^k(E\cap(S_*\cup S))
 =a+\mathcal H^k(P\cap S)>a,
\]
wbrew $S_*\cup S\in\mathcal S$ i definicji $a$.
Stąd $P$ jest czysto $k$-nieprostowalny.

Dla innego rozkładu $E=R'\,\dot\cup\,P'$ zbiór
$R\cap P'$ ma miarę zero: poza zbiorem zerowym
$R$ pokrywa przeliczalna rodzina obrazów Lipschitzowskich,
a przecięcie każdego z nich z $P'$ jest zerowe.
Podaddytywność daje tezę. Analogicznie $R'\cap P$
ma miarę zero. Ponieważ $R\setminus R'=R\cap P'$
i $R'\setminus R=R'\cap P$, różnica symetryczna
$R\triangle R'$ ma miarę zero, podobnie $P\triangle P'$.
\end{proof}

Nie potrzeba tu teorii zbiorów analitycznych:
nie twierdzimy, że dowolny obraz Lipschitzowski
zbioru borelowskiego jest borelowski, lecz korzystamy
z borelowskich nadzbiorów otrzymanych przez
przedłużenie odwzorowania na domknięcie.

\section{Dystrybucyjna pochodna funkcji całkowalnej lokalnie}
\begin{definicja}[Wariacja funkcji]
\label{def:BV-variation-29}
Niech $n\ge1$, $\Omega\subseteq\mathbb R^n$ będzie
zbiorem otwartym, a $u\in L^1_{\mathrm{loc}}(\Omega;\mathbb R)$.
Definiujemy
\[
 |Du|(\Omega)=\sup\left\{
 \int_\Omega u(x)\operatorname{div}\varphi(x)
   \,\mathrm d\lambda_n(x):
 \varphi\in C_c^1(\Omega;\mathbb R^n),\quad
 |\varphi(x)|\le1\ \text{dla wszystkich }x\in\Omega
 \right\}.
\]
Przez $BV(\Omega)$ oznaczamy przestrzeń
klas funkcji $u\in L^1(\Omega;\mathbb R)$, dla których
$|Du|(\Omega)<\infty$. Normę definiujemy wzorem
$\|u\|_{BV}=\|u\|_{L^1}+|Du|(\Omega)$.
W zastosowaniach do obszarów nieograniczonych
wygodna jest również lokalna klasa $BV_{\mathrm{loc}}(\Omega)$:
wymagamy, aby $u\in L^1_{\mathrm{loc}}(\Omega;\mathbb R)$
oraz $|Du|(V)<\infty$ dla każdego otwartego
$V$ o zwartym domknięciu zawartym w $\Omega$.
Przez $|Du|(V)$ rozumiemy supremum z tej samej
definicji z testami o zwartym nośniku w $V$.
Dzięki symetrii klasy testów (razem z $\varphi$
dopuszczalne jest $-\varphi$) w supremum można
równoważnie wziąć moduł całki.
\end{definicja}

W definicji nie jest wymagane istnienie klasycznej
pochodnej $u$. Symbol $Du$ na początku oznacza
wektorową dystrybucję. Zapis $|Du|(V)$ dla
otwartego $V\subseteq\Omega$ oznacza na tym etapie
wartość supremum, nie wariację już istniejącej
miary. Dopiero twierdzenie o reprezentacji
uzasadni, że w przypadku lokalnie skończonej
wariacji $Du$ jest lokalnie skończoną wektorową
miarą Radona, a wartości tych supremów są
masami jej dodatniej miary wariacji.

\begin{lemat}[Wariacja gładkiej funkcji]
\label{lem:BV-smooth-29}
Niech $\Omega\subseteq\mathbb R^n$ będzie zbiorem otwartym.
Jeżeli $u\in C^1(\Omega;\mathbb R)$,
to
\[
 |Du|(\Omega)=\int_\Omega|\nabla u(x)|
 \,\mathrm d\lambda_n(x)\in[0,+\infty].
\]
\end{lemat}
\begin{proof}
Całkowanie przez części dla testów o zwartym nośniku
daje $\int_\Omega u\operatorname{div}\varphi\,\mathrm d\lambda_n=-\int_\Omega\nabla u\cdot\varphi\,\mathrm d\lambda_n$.
Z nierówności Cauchy'ego--Schwarza wynika
$|Du|(\Omega)\le\int_\Omega|\nabla u|\,\mathrm d\lambda_n$.
Dla nierówności odwrotnej ustalmy $K\Subset\Omega$
i $\varepsilon>0$. Wybieramy funkcję $\theta\in C_c^\infty(\Omega)$,
$0\le\theta\le1$, równą $1$ na $K$. Pole
$w=-\theta\nabla u/(|\nabla u|+\varepsilon)$
jest ciągłe, ma zwarty nośnik i normę nie większą niż $1$.
Przedłużamy je przez zero na $\mathbb R^n$ i wygładzamy
splotem z nieujemnym jądrem wygładzającym o całce $1$.
Dla dostatecznie małego promienia splotu pola $w_\delta$
mają nośniki w jednym zwartym podzbiorze $\Omega$,
normę nie większą niż $1$ i zbiegają jednostajnie do $w$.
Ponieważ gradient jest ograniczony na tym zwartym zbiorze,
przechodząc do granicy w całkach testowych, otrzymujemy
\[
 |Du|(\Omega)\ge -\int_\Omega\nabla u\cdot w\,\mathrm d\lambda_n
 =\int_\Omega\theta\frac{|\nabla u|^2}{|\nabla u|+\varepsilon}
 \,\mathrm d\lambda_n
 \ge\int_K\frac{|\nabla u|^2}{|\nabla u|+\varepsilon}\,\mathrm d\lambda_n.
\]
Przez $\varepsilon\downarrow0$ i wyczerpanie $\Omega$
zbiorami zwartymi otrzymujemy przeciwną nierówność.
\end{proof}

Przez skończoną miarę wektorową o wartościach
w $\mathbb R^n$ rozumiemy tutaj przeliczalnie
addytywną w normie funkcję określoną na
$\sigma$-algebrze borelowskiej. W przypadku lokalnym
ograniczamy ją do względnie zwartych zbiorów
borelowskich; przeliczalna addytywność w normie
obowiązuje wówczas dla rozłącznych rodzin
o względnie zwartej sumie. Ogólna definicja
i teoria wariacji miar wektorowych znajduje się w
Rozdziale~\ref{ch:vector-measures-30}.

\begin{twierdzenie}[Pochodna $Du$ jako miara wektorowa]
\label{thm:BV-measure-29}
Niech $\Omega\subseteq\mathbb R^n$ będzie zbiorem otwartym,
$u\in L^1_{\mathrm{loc}}(\Omega;\mathbb R)$ oraz
$|Du|(\Omega)<\infty$. Wtedy istnieje jedyna
skończona miara wektorowa $Du=(D_1u,\ldots,D_nu)$,
której współrzędne są skończonymi miarami ze znakiem,
i
\[
 \int_\Omega u(x)\operatorname{div}\varphi(x)
 \,\mathrm d\lambda_n(x)
 =-\int_\Omega\varphi(x)\cdot\mathrm dDu(x)
 \qquad(\varphi\in C_c^1(\Omega;\mathbb R^n)).
\]
Jej wariacja całkowita, zdefiniowana dla borelowskiego
$B\subseteq\Omega$ jako supremum
$\sum_j|Du(B_j)|$ po wszystkich skończonych
borelowskich podziałach $B=\dot\bigcup_j B_j$,
spełnia dla każdego otwartego $V\subseteq\Omega$
\[
 |Du|(V)=\sup_{\substack{\varphi\in C_c^1(V;\mathbb R^n)\\
                         |\varphi|\le1}}
 \left|\int_V\varphi\cdot\mathrm dDu\right|.
\]
W szczególności jest to wartość supremum z definicji
wariacji funkcji ograniczonej do $V$.
Jeżeli $u\in BV_{\mathrm{loc}}(\Omega)$,
otrzymujemy lokalnie skończoną wektorową miarę
Radona, rozumianą na zbiorach borelowskich o
zwartym domknięciu. Jej dodatnia miara wariacji
jest określona na wszystkich zbiorach borelowskich
$\Omega$ i może mieć nieskończoną masę całkowitą.
\end{twierdzenie}
\begin{proof}
Oznaczmy przez $V=|Du|(\Omega)$ wartość supremum.
Dla każdego $i$ funkcjonał
$T_i(\psi)=-\int_\Omega u\partial_i\psi\,\mathrm d\lambda_n$
spełnia $|T_i(\psi)|\le V\|\psi\|_\infty$,
gdyż testujemy pole $\varphi=\psi e_i/\|\psi\|_\infty$
(dla $\psi=0$ nierówność zachodzi bezpośrednio). Funkcjonał ten
przedłuża się jednoznacznie i ciągle z $C_c^1(\Omega)$
do $C_0(\Omega)$. Istotnie, dla $\psi\in C_0(\Omega)$
i $\varepsilon>0$ zbiór $\{|\psi|\ge\varepsilon\}$
jest zwarty. Mnożymy $\psi$ przez gładką funkcję
odcinającą o wartościach w $[0,1]$, równą $1$
na tym zbiorze i o zwartym nośniku w $\Omega$.
Błąd w normie supremum nie przekracza $\varepsilon$.
Tak otrzymaną funkcję ciągłą o zwartym nośniku
przedłużamy przez zero i wygładzamy splotem
o dostatecznie małym promieniu. Splot zbiega
jednostajnie do tej funkcji i ma nośnik w $\Omega$.
To dowodzi wymaganej gęstości.
Twierdzenie Riesza--Markowa--Kakutaniego
(Rozdział~\ref{ch:riesz-markov-21}) daje jedyną
skończoną miarę Radona ze znakiem $D_i u$.
Złożenie współrzędnych daje miarę wektorową:
przeliczalna addytywność każdej współrzędnej
pociąga addytywność w normie euklidesowej.
Sumując tożsamości dla składowych pola testowego,
otrzymujemy podaną równość. Jeżeli inna miara
spełniałaby tę równość, to testy $\psi e_i$
wyznaczałyby te same funkcjonały na gęstej
podprzestrzeni $C_0(\Omega)$, zatem jednoznaczność
reprezentacji daje równość wszystkich współrzędnych.

Aby określić dokładną wariację, przyjmijmy
$\mu=\sum_{i=1}^n|D_i u|$, skończoną dodatnią miarę
Radona. Każde $D_i u$ jest bezwzględnie ciągłe
względem $\mu$, więc twierdzenie Radona--Nikodyma
dla miar ze znakiem daje $D_i u=a_i\mu$,
gdzie $a_i\in L^1(\mu;\mathbb R)$. Wybieramy borelowskie
reprezentanty i piszemy $a=(a_1,\ldots,a_n)$.
Ponieważ $|a|\le\sum_i|a_i|$, funkcja $|a|$
jest całkowalna. Niech $\nu=|a|\mu$.
Jest to skończona miara borelowska, a z regularności
skończonych miar borelowskich na $\Omega$
wynika, że jest miarą Radona.

Najpierw utożsamiamy $\nu$ z wariacją podziałową.
Dla każdego skończonego borelowskiego podziału
$B=\dot\bigcup_j B_j$ zachodzi
\[
 \sum_j|Du(B_j)|=\sum_j\left|\int_{B_j}a\,\mathrm d\mu\right|
 \le\int_B|a|\,\mathrm d\mu=\nu(B).
\]
Dla nierówności odwrotnej niech $p=a/|a|$ na
$\{|a|>0\}$ i $p=0$ w pozostałych punktach.
Dla $\varepsilon>0$ wybieramy skończoną
$\varepsilon$-siatkę $v_1,\ldots,v_M$ w domkniętej
kuli jednostkowej $\mathbb R^n$, korzystając z jej
zwartości. Rozdzielamy $B$ na zbiory borelowskie $B_j$,
przypisując każdemu punktowi pierwszy indeks $j$,
dla którego $|p-v_j|\le\varepsilon$.
Na $B_j$ zachodzi
$v_j\cdot a=p\cdot a+(v_j-p)\cdot a
\ge(1-\varepsilon)|a|$.
W konsekwencji
\[
 \sum_j|Du(B_j)|\ge\sum_j v_j\cdot Du(B_j)
 \ge(1-\varepsilon)\nu(B).
\]
Przechodząc $\varepsilon\downarrow0$, otrzymujemy
$|Du|(B)=\nu(B)$ dla każdego borelowskiego $B$.
W szczególności jest to miara i zachodzi
$Du=p|Du|$, gdzie $|p|=1$ prawie wszędzie względem
$|Du|$.

Ustalmy otwarty $V\subseteq\Omega$.
Dla każdego testu $|\varphi|\le1$ mamy
$|\int_V\varphi\cdot\mathrm dDu|\le\nu(V)$.
Aby uzyskać odwrotną nierówność, stosujemy opisaną
skończoną siatkę do $p$ na $V$ i otrzymujemy
borelowski podział $V=\dot\bigcup_{j=1}^M B_j$.
Z regularności $\nu$ wybieramy zwarte
$K_j\subseteq B_j$ tak, aby
$\nu(B_j\setminus K_j)<\varepsilon/M$.
Niepuste $K_j$ są parami rozłączne, więc mają
dodatnie wzajemne odległości, a ponadto każde
ma dodatnią odległość od dopełnienia $V$,
jeśli dopełnienie nie jest puste.
Wybieramy parami rozłączne otwarte otoczenia
$K_j$ o zwartych domknięciach w $V$ oraz
funkcje $\theta_j\in C_c^\infty(V)$ o wartościach
w $[0,1]$, równe $1$ na $K_j$ i o nośnikach
w tych otoczeniach (dla pustego $K_j$ bierzemy zero).
Pole $\varphi=\sum_j\theta_j v_j$ ma normę nie
większą niż $1$ w każdym punkcie, gdyż co najwyżej
jeden składnik jest niezerowy. Na $\bigcup_j K_j$
mamy $|\varphi-p|\le\varepsilon$, a na dopełnieniu
w $V$ mamy $|\varphi-p|\le2$. Zatem
\[
 \int_V|\varphi-p|\,\mathrm d\nu
 \le\varepsilon\nu(V)+2\varepsilon.
\]
Ponieważ $Du=p\nu$ oraz $\int_V p\cdot p\,\mathrm d\nu
=\nu(V)$, całki $\int_V\varphi\cdot\mathrm dDu$
mogą być dowolnie bliskie $\nu(V)$.
Daje to żądaną równość dla każdego otwartego $V$.

W przypadku lokalnym wybieramy rosnące otwarte
wyczerpanie $(V_j)$ zbioru $\Omega$ takie, że
$\overline V_j$ jest zwarty i zawarty w $V_{j+1}$.
Na każdym $V_j$ otrzymujemy skończoną miarę
wektorową $\eta_j$. Dla $i<j$ ograniczenie
$\eta_j$ do $V_i$ i $\eta_i$ reprezentują ten
sam funkcjonał na testach w $V_i$, więc są równe
z jednoznaczności części skończonej. W szczególności
zgadzają się ich wariacje na $V_i$.
Dla borelowskiego $B$ o zwartym domknięciu
definiujemy $Du(B)=\eta_j(B)$, wybierając
$V_j$ zawierający $\overline B$. Zgodność
ograniczeń zapewnia niezależność wyboru $j$.
Na wszystkich zbiorach borelowskich definiujemy
dodatnią miarę
$|Du|(B)=\lim_j|\eta_j|(B\cap V_j)$.
Granica istnieje, ponieważ składniki rosną.
Przeliczalna addytywność wynika z przeliczalnej
addytywności miar $B\mapsto|\eta_j|(B\cap V_j)$
i z zamiany granicy rosnącej z sumą
nieujemnych wyrazów. Miara jest lokalnie skończona,
zatem jest miarą Radona. Dla otwartego $V$
każdy test ma zwarty nośnik zawarty w pewnym $V_j$,
a supremum po testach z $V\cap V_j$ wynosi
$|\eta_j|(V\cap V_j)$. Biorąc supremum po $j$,
otrzymujemy tożsamość również dla lokalnej wariacji,
bez przypisywania wartości typu $+\infty-\infty$
składowym miary na nieograniczonych zbiorach.
\end{proof}

\begin{twierdzenie}[Zupełność przestrzeni $BV$]
\label{thm:BV-banach-29}
Niech $\Omega\subseteq\mathbb R^n$ będzie zbiorem otwartym.
Przestrzeń $BV(\Omega)$ z normą
$\|u\|_{BV}=\|u\|_{L^1(\Omega)}+|Du|(\Omega)$
jest przestrzenią Banacha.
\end{twierdzenie}
\begin{proof}
Najpierw sprawdzamy, że $BV(\Omega)$ jest przestrzenią
liniową, a podany wzór określa normę. Wariacja zależy
wyłącznie od klasy funkcji równych prawie wszędzie.
Dla $u,v\in BV(\Omega)$ i $c\in\mathbb R$ z liniowości
całki oraz symetrii klasy testów wynika
$|D(cu)|(\Omega)=|c|\,|Du|(\Omega)$ oraz
$|D(u+v)|(\Omega)\le |Du|(\Omega)+|Dv|(\Omega)$.
Istotnie, dla każdego dopuszczalnego testu $\varphi$
całka z $(u+v)\operatorname{div}\varphi$ nie przekracza
sumy obu wariacji. Bierzemy supremum po testach.
Jednorodność dla $c<0$ otrzymujemy zastępując
$\varphi$ przez $-\varphi$; dla $c=0$ jest oczywista.
Wraz z odpowiednimi własnościami normy $L^1$ daje to
domknięcie na kombinacje liniowe, jednorodność
i nierówność trójkąta dla normy $BV$.
Jeśli $\|u\|_{BV}=0$, to $\|u\|_{L^1}=0$,
czyli $u$ jest zerowym elementem przestrzeni klas
funkcji równych prawie wszędzie. Zatem jest to norma.

Niech $(u_j)$ będzie ciągiem Cauchy'ego w normie
$BV$. W szczególności jest to ciąg Cauchy'ego w
$L^1(\Omega)$, więc istnieje $u\in L^1(\Omega)$,
$u_j\to u$ w $L^1$. Z twierdzenia o reprezentacji
pochodnych mamy skończone miary wektorowe $Du_j$.
Z liniowości pochodnej dystrybucyjnej i jednoznaczności
Twierdzenia~\ref{thm:BV-measure-29} mamy
$D(u_j-u_\ell)=Du_j-Du_\ell$. Wobec tego
$|Du_j-Du_\ell|(\Omega)\le\|u_j-u_\ell\|_{BV}$,
czyli ciąg miar jest Cauchy'ego w normie wariacji.

Uzasadnimy zupełność przestrzeni takich miar.
Dla skończonej miary wektorowej $\xi=(\xi_1,\ldots,\xi_n)$
kładziemy $\rho=\sum_i|\xi_i|$, a przez $b_i$
oznaczamy gęstość $\xi_i$ względem $\rho$.
Argument podziałowy z dowodu
Twierdzenia~\ref{thm:BV-measure-29}, zastosowany
do dowolnej skończonej miary wektorowej, daje
$|\xi|=|b|\rho$. Analogicznie w wymiarze jeden
$|\xi_i|=|b_i|\rho$. Z nierówności dla norm w $\mathbb R^n$
$|b|\le\sum_i|b_i|\le\sqrt n\,|b|$ i całkowania
otrzymujemy
\[
 \frac1{\sqrt n}\sum_i|\xi_i|(\Omega)
 \le|\xi|(\Omega)\le\sum_i|\xi_i|(\Omega).
\]
W szczególności każda współrzędna ciągu $(Du_j)$
jest ciągiem Cauchy'ego w normie wariacji miar ze znakiem.
Z zupełności tej przestrzeni
(Rozdział~\ref{ch:miary-ze-znakiem}) istnieją skończone
miary ze znakiem $\eta_i$ będące granicami odpowiednich
współrzędnych. Dla $\eta=(\eta_1,\ldots,\eta_n)$
górna nierówność daje
$|Du_j-\eta|(\Omega)\le\sum_i|(Du_j)_i-\eta_i|(\Omega)\to0$.

Dla pola $\varphi\in C_c^1(\Omega;\mathbb R^n)$
przechodzimy do granicy w równości
$\int_\Omega u_j\operatorname{div}\varphi\,\mathrm d\lambda_n
=-\int_\Omega\varphi\cdot\mathrm dDu_j$.
Różnica całek po lewej stronie ma moduł nie większy
niż $\|\operatorname{div}\varphi\|_\infty\|u_j-u\|_{L^1}$,
a po prawej nie większy niż
$\|\varphi\|_\infty|Du_j-\eta|(\Omega)$.
Oba oszacowania dążą do zera. Otrzymujemy
$\int_\Omega u\operatorname{div}\varphi\,\mathrm d\lambda_n
=-\int_\Omega\varphi\cdot\mathrm d\eta$.
Dla testów o normie co najwyżej $1$ moduł tej całki
nie przekracza $|\eta|(\Omega)<\infty$.
Zatem $|Du|(\Omega)<\infty$, a ponieważ
$u\in L^1(\Omega)$, mamy $u\in BV(\Omega)$.
Dopiero teraz stosujemy Twierdzenie~\ref{thm:BV-measure-29}
i z jednoznaczności otrzymujemy $Du=\eta$. W konsekwencji
$\|u_j-u\|_{BV}=\|u_j-u\|_{L^1}+|Du_j-\eta|(\Omega)\to0$.
\end{proof}

\section{Zbiory o skończonym obwodzie}
\begin{definicja}[Obwód w sensie De Giorgiego]
\label{def:finite-perimeter-29}
Niech $E\subseteq\mathbb R^n$ będzie zbiorem mierzalnym
w sensie Lebesgue'a, a $\Omega\subseteq\mathbb R^n$ otwartym.
Definiujemy \emph{obwód $E$ w $\Omega$} przez
\[
 P(E;\Omega):=|D\chi_E|(\Omega)
 =\sup\left\{\int_{E\cap\Omega}\operatorname{div}\varphi(x)
 \,\mathrm d\lambda_n(x):
 \varphi\in C_c^1(\Omega;\mathbb R^n),\quad|\varphi|\le1
 \right\}.
\]
Zbiór $E$ ma \emph{lokalnie skończony obwód},
gdy $P(E;V)<\infty$ dla wszystkich otwartych $V$
o zwartym domknięciu w $\mathbb R^n$.
Równoważnie $\chi_E\in BV_{\mathrm{loc}}(\mathbb R^n)$,
gdyż $\chi_E$ jest lokalnie całkowalna niezależnie
od $\lambda_n(E)$.
Jeśli $P(E;\mathbb R^n)<\infty$, mówimy o skończonym
obwodzie globalnym. Nie wymagamy $\lambda_n(E)<\infty$:
półprzestrzeń ma lokalnie skończony obwód, ale
nieskończony obwód globalny przy $n\ge2$.
Istotnie, dla $E=\{x\in\mathbb R^n:x_n>0\}$
i testu $\varphi\in C_c^1(\mathbb R^n;\mathbb R^n)$
całkowanie po współrzędnych stycznych daje zero,
a całkowanie po ostatniej współrzędnej daje
$-\int_{\mathbb R^{n-1}}\varphi_n(x',0)
\,\mathrm d\lambda_{n-1}(x')$.
Zatem $D\chi_E=e_n\mathcal H^{n-1}\llcorner\{x_n=0\}$
lokalnie, a jego wariacja jest miarą Hausdorffa
hiperpłaszczyzny. Jej masa na zbiorach ograniczonych
jest skończona, natomiast masa całkowita jest
nieskończona dla $n\ge2$. Dla $n=1$ półprosta
ma pochodną $\delta_0$ i obwód globalny $1$.
Tu i dalej dla $n=1$ stosujemy konwencję
$\mathbb R^0=\{0\}$ z miarą jednostkową.
\end{definicja}

\begin{twierdzenie}[Granica klasycznej geometrii]
\label{thm:smooth-perimeter-29}
Niech $E\subseteq\mathbb R^n$ będzie ograniczonym
obszarem z brzegiem klasy $C^1$ (lokalnie $E$ leży
po jednej stronie wykresu klasy $C^1$) i niech
$\nu_E$ będzie jednostkowym polem normalnych
zewnętrznych na $\partial E$. Wtedy jako skończone
miary wektorowe na $\mathbb R^n$ mamy
\[
 D\chi_E=-\nu_E\,\mathcal H^{n-1}\llcorner\partial E.
\]
Dla każdego otwartego $\Omega\subseteq\mathbb R^n$
zachodzi $P(E;\Omega)=\mathcal H^{n-1}(\partial E\cap\Omega)$;
pochodną na $\Omega$ otrzymujemy przez ograniczenie
powyższej miary do $\Omega$.
Zapis $\mu\llcorner A$ oznacza ograniczenie miary
$\mu$ do zbioru $A$: $(\mu\llcorner A)(B)=\mu(B\cap A)$.
\end{twierdzenie}
\begin{proof}
Dla $n=1$ ograniczony obszar jest przedziałem
$(a,b)$. Z podstawowego twierdzenia rachunku całkowego
otrzymujemy $\int_a^b\varphi'\,\mathrm d\lambda_1
=\varphi(b)-\varphi(a)$, czyli $D\chi_E=\delta_a-\delta_b$.
Normalne zewnętrzne w $a,b$ to odpowiednio $-1,1$,
więc oba wzory wynikają z konwencji miary liczącej.
Od tej chwili zakładamy $n\ge2$.

\emph{Krok 1: rachunek lokalny.}
Wokół każdego punktu brzegu wybieramy, po izometrii
afinicznej, otoczenie cylindryczne $U=W\times(c,d)$,
gdzie $W\subseteq\mathbb R^{n-1}$ jest otwarty,
oraz funkcję $g\in C^1(W;\mathbb R)$ o wartościach
w $(c,d)$, tak że w $U$ brzeg jest wykresem $g$,
a $E\cap U=\{(x',t)\in U:t<g(x')\}$.
Jeżeli obszar leży nad wykresem, odbicie ostatniej
współrzędnej sprowadza nas do powyższej sytuacji.
Izometria $x=b+Ry$ z macierzą ortogonalną $R$
nie zmienia całek w rozpatrywanej tożsamości:
pole w nowych współrzędnych to
$R^{\mathsf T}\varphi(b+Ry)$, jego macierz różniczki
to $R^{\mathsf T}D\varphi(b+Ry)R$, a ślad jest
niezmieniony. Normalne zmieniają się przez
$R^{\mathsf T}$ i zachowane są iloczyny skalarne
oraz miary Lebesgue'a i Hausdorffa.

Niech $\psi\in C_c^1(U;\mathbb R^n)$.
Dla $i<n$ definiujemy funkcję
$H_i(x')=\int_c^{g(x')}\psi_i(x',t)\,\mathrm d\lambda_1(t)$.
Jest ona klasy $C^1$ i ma zwarty nośnik w $W$,
ponieważ rzut nośnika $\psi$ na $W$ jest zwarty.
Różniczkowanie całki ze zmiennym końcem daje
\[
 \partial_i H_i(x')=
 \int_c^{g(x')}\partial_i\psi_i(x',t)\,\mathrm d\lambda_1(t)
 +\psi_i(x',g(x'))\partial_i g(x').
\]
Całka z $\partial_i H_i$ po $W$ jest zerowa:
przedłużamy $H_i$ przez zero i stosujemy Fubiniego
oraz podstawowe twierdzenie rachunku całkowego
na każdej prostej równoległej do $i$-tej osi.
Zatem całka z $\partial_i\psi_i$ po $E\cap U$
jest równa
$-\int_W\psi_i(x',g(x'))\partial_i g(x')
\,\mathrm d\lambda_{n-1}(x')$.
Dla ostatniej składowej otrzymujemy
$\int_c^{g(x')}\partial_n\psi_n(x',t)\,\mathrm d\lambda_1(t)
=\psi_n(x',g(x'))$, gdyż pole znika przy dolnym końcu.
Sumując i stosując Fubiniego do funkcji o zwartym
nośniku, dostajemy
\[
 \int_{E\cap U}\operatorname{div}\psi\,\mathrm d\lambda_n
 =\int_W\psi(x',g(x'))\cdot(-\nabla g(x'),1)
 \,\mathrm d\lambda_{n-1}(x').
\]
Wektor $(-\nabla g,1)$ jest gradientem funkcji
$t-g(x')$, ujemnej w $E\cap U$, więc po normalizacji
jest normalną zewnętrzną. Parametryzacja wykresu
$x'\mapsto(x',g(x'))$ ma macierz Grama
$I_{n-1}+\nabla g\,\nabla g^{\mathsf T}$,
o wyznaczniku $1+|\nabla g|^2$ (w bazie zawierającej
kierunek gradientu, gdy jest niezerowy, macierz jest
diagonalna; dla zerowego gradientu jest jednostkowa).
Ze wzoru na pole wykresu
(Twierdzenie~\ref{thm:graph-area-28}) otrzymujemy
$\int_{E\cap U}\operatorname{div}\psi\,\mathrm d\lambda_n
=\int_{\partial E\cap U}\psi\cdot\nu_E
\,\mathrm d\mathcal H^{n-1}$.

\emph{Krok 2: sklejanie.}
Brzeg $\partial E$ jest zwarty, gdyż $E$ jest ograniczony.
Wybieramy skończone pokrycie brzegu mniejszymi
otoczeniami wykresowymi o zwartych domknięciach
w otoczeniach z kroku 1. Na każdym takim fragmencie
miara $\mathcal H^{n-1}$ jest skończona:
rzut jego domknięcia leży w zwartym podzbiorze $W$,
gradient jest tam ograniczony, a wzór wykresowy
ogranicza pole przez stałą razy miara ograniczonego rzutu.
Stąd $\mathcal H^{n-1}(\partial E)<\infty$.

Do rodziny otoczeń wykresowych dołączamy $E$
oraz $\mathbb R^n\setminus\overline E$.
Otrzymujemy otwarte pokrycie $\mathbb R^n$.
Dla $\varphi\in C_c^1(\mathbb R^n;\mathbb R^n)$
stosujemy gładki podział jedności podporządkowany
temu pokryciu (można ograniczyć konstrukcję do
otoczenia zwartego nośnika $\varphi$).
Pole jest skończoną sumą pól $\psi_j=\theta_j\varphi$
o zwartych nośnikach w odpowiednich otoczeniach.
Różniczkując tę równość, dostajemy
$\operatorname{div}\varphi=\sum_j\operatorname{div}\psi_j$.
Dla składnika o nośniku w $E$ całka z dywergencji
po $E$ jest zerowa, gdyż jest całką po całej przestrzeni;
całka brzegowa również jest zerowa.
Dla składnika o nośniku poza $\overline E$
obie całki są zerowe. Dla pozostałych składników
stosujemy krok 1. Sumowanie daje
\[
 \int_E\operatorname{div}\varphi\,\mathrm d\lambda_n
 =\int_{\partial E}\varphi\cdot\nu_E\,\mathrm d\mathcal H^{n-1}.
\]

\emph{Krok 3: identyfikacja miary.}
Prawa strona dla $|\varphi|\le1$ jest co do modułu
ograniczona przez skończoną miarę brzegu. Zatem
$|D\chi_E|(\mathbb R^n)<\infty$ i
Twierdzenie~\ref{thm:BV-measure-29} ma zastosowanie.
Jednoznaczność reprezentacji oraz znak w definicji
pochodnej dystrybucyjnej dają
$D\chi_E=-\nu_E\mathcal H^{n-1}\llcorner\partial E$.
Wariacja tej miary jest miarą powierzchniową,
ponieważ $|\nu_E|=1$. Dla testów w $\Omega$
przedłużenie przez zero jest klasy $C^1$ i pozwala
ograniczyć miarę do $\Omega$; tożsamość wariacji
na otwartych zbiorach daje wzór na $P(E;\Omega)$.
\end{proof}

\begin{przyklad}[Kostka ma skończony obwód]
\label{ex:cube-perimeter-29}
Niech $Q=(0,a)^n$, $a>0$, i niech
$\varphi\in C_c^1(\mathbb R^n;\mathbb R^n)$.
Najpierw dla $n\ge2$, całkując każdą składową
$\partial_i\varphi_i$ po $i$-tej współrzędnej,
stosujemy Fubiniego (całkowane funkcje są ograniczone
na zbiorze o skończonej mierze) oraz podstawowe
twierdzenie rachunku całkowego. Otrzymujemy
\[
 \int_Q\operatorname{div}\varphi(x)\,\mathrm d\lambda_n(x)
 =\sum_{i=1}^n\int_{(0,a)^{n-1}}
 [\varphi_i(x_1,\ldots,a,\ldots,x_n)-
  \varphi_i(x_1,\ldots,0,\ldots,x_n)]
 \,\mathrm d\lambda_{n-1}(x_{\widehat i}).
\]
Wyraz $x_{\widehat i}$ oznacza wszystkie współrzędne
poza $i$-tą. Niech $S_i^+$ i $S_i^-$ oznaczają względne wnętrza
ścian kostki odpowiednio o $x_i=a$ i $x_i=0$.
Parametryzacja każdej ściany pozostałymi współrzędnymi
jest izometrią, a więc jej $\mathcal H^{n-1}$-miara
wynosi $a^{n-1}$. Z powyższej tożsamości i znaku
w definicji pochodnej dystrybucyjnej otrzymujemy
\[
 D\chi_Q=\sum_{i=1}^n e_i
 \bigl(\mathcal H^{n-1}\llcorner S_i^-
       -\mathcal H^{n-1}\llcorner S_i^+\bigr).
\]
Ściany są parami rozłączne, a gęstość wektorowa
na każdej z nich jest wektorem jednostkowym
o znaku przeciwnym do normalnej zewnętrznej.
Miara po prawej jest skończona, więc tożsamość
testowa daje skończoność wariacji i jednoznacznie
identyfikuje $D\chi_Q$. Na mocy dowodu
Twierdzenia~\ref{thm:BV-measure-29} jej wariacja
jest sumą miar poszczególnych ścian.
Pozostałe fragmenty brzegu leżą w skończonej sumie
płaszczyzn afinicznych wymiaru co najwyżej $n-2$,
więc mają zerową $\mathcal H^{n-1}$-miarę.
Stąd $P(Q;\mathbb R^n)=2na^{n-1}$.
Dla $n=1$ bezpośrednio
$\int_{(0,a)}\varphi'\,\mathrm d\lambda_1
=\varphi(a)-\varphi(0)$, zatem
$D\chi_{(0,a)}=\delta_0-\delta_a$ i obwód wynosi $2$,
zgodnie z tym samym wzorem. Nie wymagaliśmy
gładkości brzegu w narożnikach. Dla prostokąta
$(0,a)\times(0,b)$ z $a,b>0$ ten sam rachunek
daje dwie ściany długości $a$ i dwie długości $b$,
stąd jego obwód wynosi $2(a+b)$.
\end{przyklad}

\begin{center}
\begin{tikzpicture}[scale=1.02,>=stealth]
  \fill[blue!8] (0,0) rectangle (3.5,2.4);
  \draw[blue!65!black,thick] (0,0) rectangle (3.5,2.4);
  \node at (1.75,1.2) {$E$};
  \foreach \x in {0.65,1.75,2.85} {
    \draw[->,very thick,red!75!black] (\x,2.4)--(\x,2.99);
    \draw[->,very thick,red!75!black] (\x,0)--(\x,-0.59);
  }
  \foreach \y in {0.6,1.2,1.8} {
    \draw[->,very thick,red!75!black] (0,\y)--(-0.6,\y);
    \draw[->,very thick,red!75!black] (3.5,\y)--(4.1,\y);
  }
  \node[font=\small] at (1.75,3.25) {$\nu_E$};
  \node[font=\small] at (1.75,-.95) {$a$};
  \node[font=\small] at (4.4,1.2) {$b$};
\end{tikzpicture}
\end{center}
\noindent\emph{Rysunek:} miarowo-teoretyczna normalna i obwód
prostokąta w $\mathbb R^2$. Dla boków długości $a,b$
mamy $P(E;\mathbb R^2)=2(a+b)$ oraz
$D\chi_E=-\nu_E\mathcal H^1\llcorner\partial E$,
gdzie normalną na wierzchołkach można określić dowolnie:
cztery wierzchołki mają zerową $\mathcal H^1$-miarę
i nie wnoszą osobnej masy do obwodu.

\begin{lemat}[Pochodna i dopełnienie]
\label{lem:perimeter-complement-29}
Niech $E,F\subseteq\mathbb R^n$ będą mierzalne
w sensie Lebesgue'a, a $\Omega\subseteq\mathbb R^n$
będzie zbiorem otwartym. Wtedy
$D\chi_{E^c}=-D\chi_E$ jako dystrybucje na $\Omega$
oraz $P(E^c;\Omega)=P(E;\Omega)$.
Jeżeli $\lambda_n(E\triangle F)=0$,
to $P(E;\Omega)=P(F;\Omega)$.
\end{lemat}
\begin{proof}
Dla $\varphi\in C_c^1(\Omega;\mathbb R^n)$
przedłużenie przez zero jest klasy $C^1$ na całej
przestrzeni. Dla każdego $i$ całka z $\partial_i\varphi_i$
po $\mathbb R^n$ wynosi zero na mocy Fubiniego
i podstawowego twierdzenia rachunku całkowego,
bo całkowane funkcje mają zwarty nośnik.
Z tożsamości $\chi_{E^c}=1-\chi_E$ wynika więc
$\int_\Omega\chi_{E^c}\operatorname{div}\varphi
\,\mathrm d\lambda_n=-\int_\Omega\chi_E
\operatorname{div}\varphi\,\mathrm d\lambda_n$.
Daje to równość dystrybucji ze znakiem przeciwnym.
Ponieważ dopuszczalny test pozostaje dopuszczalny
po zmianie znaku, oba suprema są równe, także gdy
mają wartość $+\infty$.
Dla zbiorów równych prawie wszędzie ich funkcje
charakterystyczne są równe w $L^1_{\mathrm{loc}}$,
więc dają tę samą dystrybucję pochodnej.
\end{proof}

\section{Zwartość i półciągłość: dwa pełne dowody}
\begin{twierdzenie}[Dolna półciągłość obwodu]
\label{thm:perimeter-lsc-29}
Niech $\Omega\subseteq\mathbb R^n$ będzie zbiorem otwartym,
a $E,E_j\subseteq\mathbb R^n$, $j\ge1$, będą
mierzalne w sensie Lebesgue'a.
Jeśli $\chi_{E_j}\to\chi_E$ w
$L^1_{\mathrm{loc}}(\Omega)$, to
\[
 P(E;\Omega)\le\liminf_{j\to\infty}P(E_j;\Omega).
\]
Ogólniej, dla $u,u_j\in L^1_{\mathrm{loc}}(\Omega;\mathbb R)$,
jeśli $u_j\to u$ w $L^1_{\mathrm{loc}}(\Omega)$, to
$|Du|(\Omega)\le\liminf_j|Du_j|(\Omega)$.
\end{twierdzenie}
\begin{proof}
Ustalmy pole $\varphi\in C_c^1(\Omega;\mathbb R^n)$,
$|\varphi|\le1$. Niech $K=\operatorname{supp}\varphi\subseteq\Omega$.
Różnica całek testowych ma moduł nie większy
niż $\|\operatorname{div}\varphi\|_\infty
\int_K|u_j-u|\,\mathrm d\lambda_n$, a to dąży do zera
z lokalnej zbieżności w $L^1$.
Każda z całek po lewej jest nie większa niż
$|Du_j|(\Omega)$. Zatem całka dla $u$
jest nie większa niż $\liminf_j|Du_j|(\Omega)$.
Bierzemy supremum po wszystkich testach,
uzyskując tezę. Przypadek zbiorów otrzymujemy
po podstawieniu $u_j=\chi_{E_j}$.
\end{proof}

\begin{lemat}[Oszacowanie przesunięcia]
\label{lem:BV-translation-29}
Niech $n\ge1$ i $u\in BV(\mathbb R^n)$.
Wówczas dla każdego $h\in\mathbb R^n$
\[
 \int_{\mathbb R^n}|u(x+h)-u(x)|
 \,\mathrm d\lambda_n(x)\le|h|\,|Du|(\mathbb R^n).
\]
W szczególności dla mierzalnego zbioru $E$ o
skończonej mierze i skończonym obwodzie zachodzi
$\lambda_n((E+h)\triangle E)\le|h|P(E;\mathbb R^n)$.
\end{lemat}
\begin{proof}
Wybierzmy nieujemny gładki jądro wygładzające
$\rho\in C_c^\infty(B(0,1))$ z
$\int_{\mathbb R^n}\rho(x)\,\mathrm d\lambda_n(x)=1$ i
$\rho_\varepsilon(x)=\varepsilon^{-n}\rho(x/\varepsilon)$.
Niech $u_\varepsilon=u*\rho_\varepsilon$, czyli
$u_\varepsilon(x)=\int_{\mathbb R^n}u(y)
\rho_\varepsilon(x-y)\,\mathrm d\lambda_n(y)$.
Ponieważ $u\in L^1(\mathbb R^n)$, a wszystkie
pochodne $\rho_\varepsilon$ są ograniczone,
różniczkowanie pod całką jest uzasadnione
twierdzeniem o zbieżności zdominowanej
(stosowanym do ilorazów różnicowych).
W szczególności $u_\varepsilon$ jest gładka.
Dla ustalonego $x$ funkcja
$y\mapsto\rho_\varepsilon(x-y)$ jest testem gładkim
o zwartym nośniku. Z definicji pochodnej
dystrybucyjnej i Twierdzenia~\ref{thm:BV-measure-29}
otrzymujemy, dla $i=1,\ldots,n$,
\begin{align*}
 \partial_i u_\varepsilon(x)
 &=\int_{\mathbb R^n}u(y)
 \partial_{x_i}\rho_\varepsilon(x-y)\,\mathrm d\lambda_n(y)\\
 &=-\int_{\mathbb R^n}u(y)
 \partial_{y_i}\rho_\varepsilon(x-y)\,\mathrm d\lambda_n(y)\\
 &=\int_{\mathbb R^n}\rho_\varepsilon(x-y)\,\mathrm dD_i u(y).
\end{align*}
Stąd $\nabla u_\varepsilon=(Du)*\rho_\varepsilon$.
Jeżeli $Du=p|Du|$ jest rozkładem polarnym, to
$|p|=1$ prawie wszędzie względem $|Du|$.
Nieujemność jądra daje punktowo
\[
 |\nabla u_\varepsilon(x)|
 \le\int_{\mathbb R^n}\rho_\varepsilon(x-y)\,\mathrm d|Du|(y).
\]
Całkując względem $x$ i stosując Tonellego do
nieujemnej funkcji podcałkowej, otrzymujemy
\[
 \int_{\mathbb R^n}|\nabla u_\varepsilon|\,\mathrm d\lambda_n
 \le\int_{\mathbb R^n}
 \left(\int_{\mathbb R^n}\rho_\varepsilon(x-y)
 \,\mathrm d\lambda_n(x)\right)\mathrm d|Du|(y)
 =|Du|(\mathbb R^n).
\]
Podstawowe twierdzenie rachunku całkowego zastosowane
do funkcji $t\mapsto u_\varepsilon(x+th)$ daje
\[
 u_\varepsilon(x+h)-u_\varepsilon(x)
 =\int_0^1\nabla u_\varepsilon(x+th)\cdot h
 \,\mathrm d\lambda_1(t).
\]
Biorąc moduł, stosując nierówność Cauchy'ego--Schwarza,
Tonellego oraz niezmienniczość miary Lebesgue'a
względem translacji, otrzymujemy
\begin{align*}
 \int_{\mathbb R^n}|u_\varepsilon(x+h)-u_\varepsilon(x)|
 \,\mathrm d\lambda_n(x)
 &\le |h|\int_0^1\int_{\mathbb R^n}
 |\nabla u_\varepsilon(x+th)|\,\mathrm d\lambda_n(x)
 \,\mathrm d\lambda_1(t)\\
 &\le |h|\,|Du|(\mathbb R^n).
\end{align*}
Uzasadnijmy przejście do granicy.
Translacje są ciągłe w $L^1(\mathbb R^n)$: dla
funkcji ciągłych o zwartym nośniku wynika to
z jednostajnej ciągłości i wspólnego ograniczonego
nośnika przy małych przesunięciach; dla dowolnego
$u\in L^1$ korzystamy z gęstości tych funkcji
oraz z faktu, że translacja zachowuje normę $L^1$.
Z nierówności trójkąta i Tonellego wynika
\[
 \|u_\varepsilon-u\|_{L^1}
 \le\int_{\mathbb R^n}\rho(z)
 \|u(\,\cdot-\varepsilon z)-u\|_{L^1}\,\mathrm d\lambda_n(z)
 \longrightarrow0,
\]
gdzie stosujemy zbieżność zdominowaną z dominantą
$2\|u\|_{L^1}\rho(z)$.
Różnica między $u_\varepsilon(\cdot+h)-u_\varepsilon$
a $u(\cdot+h)-u$ ma normę nie większą niż
$2\|u_\varepsilon-u\|_{L^1}$. Przechodzimy zatem
do granicy po lewej stronie oszacowania, którego
prawa strona nie zależy od $\varepsilon$.

Dla $u=\chi_E$ założenia zapewniają $u\in BV$,
a $\chi_E(x+h)=\chi_{E-h}(x)$.
Norma różnicy jest więc równa
$\lambda_n((E-h)\triangle E)=
\lambda_n(E\triangle(E+h))$,
gdzie ostatnia równość wynika z translacji o $h$.
\end{proof}

\begin{twierdzenie}[Zwartość rodzin zbiorów o ograniczonym obwodzie]
\label{thm:perimeter-compact-29}
Niech $n\ge1$, $Q\subseteq\mathbb R^n$ będzie
ograniczonym prostopadłościanem, a $(E_j)_{j\ge1}$
ciągiem zbiorów mierzalnych w sensie Lebesgue'a
zawartych w $Q$ i spełniających
$\sup_j P(E_j;\mathbb R^n)<\infty$.
Wówczas istnieje podciąg $E_{j_\ell}$ oraz
mierzalny $E\subseteq\overline Q$, dla których
$\chi_{E_{j_\ell}}\to\chi_E$ w $L^1(\mathbb R^n)$,
a $P(E;\mathbb R^n)\le\liminf_\ell P(E_{j_\ell};\mathbb R^n)$.
\end{twierdzenie}
\begin{proof}
Niech $P_0=\sup_jP(E_j;\mathbb R^n)$.
Wybierzmy duży domknięty sześcian $S$, którego wnętrze
zawiera $\overline Q$. Dla $m\ge1$ podzielmy $S$
na $N_m=2^{mn}$ jednakowych kostek o boku
$\delta_m=2^{-m}\ell(S)$, gdzie $\ell(S)$ jest
długością boku $S$. Przyjmujemy kostki półotwarte,
pomijając zbiór miary zero na brzegu $S$.
Są one parami rozłączne i pokrywają $S$ prawie wszędzie.
Każdy $u_j=\chi_{E_j}$ należy do $BV(\mathbb R^n)$,
gdyż $\|u_j\|_{L^1}\le\lambda_n(Q)<\infty$ i
$|Du_j|(\mathbb R^n)\le P_0$. Zastępujemy go funkcją
$A_m u_j$, stałą na każdej kostce i równą
średniej $u_j$ na tej kostce, oraz równą $0$ poza $S$.
Współczynniki $A_m u_j$ należą do $[0,1]$.

Oszacujmy błąd przybliżenia. Dla kostki $C$
i jej średniej $u_C$ nierówność trójkąta i Fubini dają
\[
 \int_C|u_j(x)-u_C|\,\mathrm d\lambda_n(x)
 \le\frac1{\lambda_n(C)}\int_C\int_C
 |u_j(x)-u_j(y)|\,\mathrm d\lambda_n(y)
 \,\mathrm d\lambda_n(x).
\]
Dla każdej kostki podstawiamy $y=x+h$.
Dla ustalonego $h$ warunki $x,y\in C$ są równoważne
$x\in C\cap(C-h)$. Różnica dwóch punktów kostki
leży w $[-\delta_m,\delta_m]^n$. Wobec tego
\begin{align*}
 &\sum_C\int_C\int_C |u_j(x)-u_j(y)|
 \,\mathrm d\lambda_n(y)\,\mathrm d\lambda_n(x)\\
 &=\int_{[-\delta_m,\delta_m]^n}\sum_C
 \int_{C\cap(C-h)}|u_j(x+h)-u_j(x)|
 \,\mathrm d\lambda_n(x)\,\mathrm d\lambda_n(h).
\end{align*}
Zbiory $C\cap(C-h)$ są parami rozłączne dla
ustalonego $h$, ponieważ każdy z nich zawiera się
w odpowiedniej kostce $C$. Możemy więc zastąpić
sumę całek po nich jedną całką po $\mathbb R^n$
jako górnym oszacowaniem, bez dodatkowego czynnika
$N_m$. Ponieważ $\lambda_n(C)=\delta_m^n$,
otrzymujemy
\[
 \sum_C\int_C|u_j-u_C|\,\mathrm d\lambda_n
 \le\delta_m^{-n}\int_{[-\delta_m,\delta_m]^n}
 \int_{\mathbb R^n}|u_j(x+h)-u_j(x)|
 \,\mathrm d\lambda_n(x)\,\mathrm d\lambda_n(h)
 \le 2^n\sqrt n\,\delta_m P_0.
\]
Ostatnia nierówność wynika z Lematu
\ref{lem:BV-translation-29}, ponieważ $|h|\le\sqrt n\,\delta_m$
na całym $[-\delta_m,\delta_m]^n$, którego objętość
wynosi $(2\delta_m)^n$. Poza $S$ obie funkcje
są zerowe prawie wszędzie. Zatem
$\sup_j\|u_j-A_mu_j\|_{L^1}\le
2^n\sqrt n\,\delta_m P_0\to0$.

Dla ustalonego $m$ wektor średnich na kostkach
należy do zwartego sześcianu $[0,1]^{N_m}$.
Jeśli dwa wektory średnich mają współrzędne
$a_C,b_C$, to odpowiadające im funkcje schodkowe
różnią się w normie $L^1$ o
$\delta_m^n\sum_C|a_C-b_C|$. Zbieżność współrzędnych
daje zatem zbieżność w $L^1$.
Wybieramy dla $m=1$ podciąg indeksów, dla którego
$A_1u_j$ jest zbieżny, następnie jego podciąg dla
$m=2$ i tak dalej. Oznaczmy przez $j^{(m)}_\ell$
rosnący ciąg indeksów wybrany w $m$-tym kroku.
Kładziemy $j_\ell=j^{(\ell)}_\ell$. Indeksy
$j_\ell$ są ściśle rosnące, a dla każdego
ustalonego $m$ wszystkie $j_\ell$ przy $\ell\ge m$
należą do podciągu wybranego w kroku $m$.
Stąd $A_m u_{j_\ell}$ jest zbieżny w $L^1$.

Dla $\varepsilon>0$ wybieramy $m$ tak duże, aby
$\sup_j\|u_j-A_mu_j\|_{L^1}<\varepsilon/3$.
Dla dostatecznie dużych $\ell,r$ mamy także
$\|A_mu_{j_\ell}-A_mu_{j_r}\|_{L^1}<\varepsilon/3$.
Nierówność trójkąta daje
\[
 \|u_{j_\ell}-u_{j_r}\|_{L^1}
 \le\|u_{j_\ell}-A_mu_{j_\ell}\|_{L^1}
 +\|A_mu_{j_\ell}-A_mu_{j_r}\|_{L^1}
 +\|A_mu_{j_r}-u_{j_r}\|_{L^1}<\varepsilon.
\]
Z zupełności $L^1$ istnieje $u\in L^1(\mathbb R^n)$,
do którego ten podciąg zbiega w normie.

Wybieramy mierzalnego reprezentanta $u$.
Z faktu $u_{j_\ell}(x)\in\{0,1\}$ wynika
$\operatorname{dist}(u(x),\{0,1\})
\le|u(x)-u_{j_\ell}(x)|$.
Po scałkowaniu i przejściu do granicy otrzymujemy
$\int_{\mathbb R^n}\operatorname{dist}(u(x),\{0,1\})
\,\mathrm d\lambda_n(x)=0$.
Nieujemność funkcji podcałkowej daje
$u(x)\in\{0,1\}$ prawie wszędzie.
Ponadto
$\int_{\mathbb R^n\setminus\overline Q}|u|
\,\mathrm d\lambda_n\le\|u-u_{j_\ell}\|_{L^1}\to0$.
Definiujemy mierzalny zbiór
$E=\{x\in\overline Q:u(x)=1\}$, wówczas
$u=\chi_E$ prawie wszędzie. Nie musimy wybierać
dalszego podciągu. Zbieżność w $L^1(\mathbb R^n)$
implikuje zbieżność lokalną, więc
Twierdzenie~\ref{thm:perimeter-lsc-29} daje
$P(E;\mathbb R^n)\le\liminf_\ell P(E_{j_\ell};\mathbb R^n)
\le P_0$, co kończy dowód.
\end{proof}

\section{Normalna miarowo-teoretyczna i brzeg zredukowany}
\begin{definicja}[Normalna z pochodnej miary]
\label{def:normal-BV-29}
Niech $n\ge1$ i niech $E\subseteq\mathbb R^n$
będzie mierzalnym w sensie Lebesgue'a zbiorem
lokalnie skończonego obwodu.
Miara wektorowa $D\chi_E$ ma rozkład polarny
$D\chi_E=-\nu_E|D\chi_E|$, gdzie
$\nu_E:\mathbb R^n\to\mathbb R^n$ jest
$|D\chi_E|$-mierzalna i
$|\nu_E(x)|=1$ dla $|D\chi_E|$-prawie wszystkich $x$.
Funkcję $\nu_E$ nazywamy \emph{miarowo-teoretyczną
normalną zewnętrzną}. Jest ona wyznaczona jednoznacznie
z dokładnością do zbioru $|D\chi_E|$-miary zero.
\end{definicja}

\begin{proof}[Uzasadnienie istnienia normalnej]
Na każdej kuli $B(0,j)$, $j\ge1$, stosujemy
Twierdzenie~\ref{thm:BV-measure-29} do ograniczenia
$\chi_E$. Otrzymujemy rozkład
$D\chi_E=p_j|D\chi_E|$ na tej kuli, gdzie $p_j$
jest borelowskie i ma normę $1$ prawie wszędzie
względem wariacji. Na przecięciu dwóch kul
pola te są równe prawie wszędzie względem wariacji
z jednoznaczności pochodnych Radona--Nikodyma
każdej składowej. Definiujemy $p$ na rozłącznych
borelowskich warstwach $B(0,j)\setminus B(0,j-1)$
(przy $B(0,0)=\varnothing$) przez $p_j$.
Zgodność ograniczeń miar i przeliczalność warstw
dają $D\chi_E=p|D\chi_E|$ na każdym względnie
zwartym zbiorze borelowskim i $|p|=1$ prawie
wszędzie względem wariacji. Kładziemy $\nu_E=-p$.
Na ewentualnym zbiorze zerowym, gdzie norma pola
nie wynosi $1$, możemy przypisać wartość zero.
Jednoznaczność wynika z jednoznaczności gęstości
na każdej z kul, a przeliczalna suma zbiorów
wyjątkowych nadal ma miarę zero.
\end{proof}

\begin{definicja}[Punkty gęstości i brzeg miarowy]
Dla\label{def:measure-boundary-29} $n\ge1$ i zbioru mierzalnego w sensie
Lebesgue'a $E\subseteq\mathbb R^n$
niech $E^{(1)}$ i $E^{(0)}$ oznaczają zbiory punktów,
w których granica
\[
 \theta(E,x)=\lim_{r\downarrow0}
 \frac{\lambda_n(E\cap B(x,r))}{\lambda_n(B(x,r))}
\]
istnieje i wynosi odpowiednio $1$ lub $0$.
\emph{Brzegiem miarowym} nazywamy
$\partial^{\mathrm m}E=\mathbb R^n\setminus(E^{(0)}\cup E^{(1)})$.
Punkty, w których powyższa granica nie istnieje,
również należą do $\partial^{\mathrm m}E$.
\end{definicja}

\begin{lemat}[Brzeg miarowy jest objętościowo mały]
\label{lem:essential-boundary-null-29}
Dla każdego mierzalnego $E\subseteq\mathbb R^n$
zachodzi $\lambda_n(\partial^{\mathrm m}E)=0$.
\end{lemat}
\begin{proof}
Najpierw sprawdzamy borelowskość rozważanych zbiorów,
nawet jeśli sam $E$ nie jest borelowski.
Dla każdego $r>0$ funkcja
$f_r(x)=\lambda_n(E\cap B(x,r))/(\omega_n r^n)$,
gdzie $\omega_n=\lambda_n(B(0,1))$, jest ciągła:
moduł różnicy liczników w punktach $x,y$ nie
przekracza $\lambda_n(B(x,r)\triangle B(y,r))$,
która dąży do zera dla $y\to x$ (różnica kul
mieści się w pierścieniu o grubości dążącej do zera).
Dla $i\in\{0,1\}$ mamy
\[
 E^{(i)}=\bigcap_{l=1}^\infty\bigcup_{m=1}^\infty
 \bigcap_{\substack{q\in\mathbb Q\\0<q<1/m}}
 \{x:|f_q(x)-i|\le1/l\}.
\]
Równość wynika z ciągłości $r\mapsto f_r(x)$
dla $r>0$: przy zmianie promienia różnica objętości
kul dąży do zera. Zatem warunki dla wymiernych
promieni przenoszą się na wszystkie promienie
w tym samym przedziale. Prawa strona jest borelowska.
W konsekwencji także $\partial^{\mathrm m}E$
jest borelowski.
Twierdzenie~\ref{thm:lebesgue-diff-25} o różniczkowaniu
Lebesgue'a stosujemy do
$\chi_E\in L^1_{\mathrm{loc}}(\mathbb R^n)$.
Dla prawie każdego $x$ średnie funkcji $\chi_E$
na kulach $B(x,r)$ zbiegają do $\chi_E(x)$,
a więc do $0$ lub $1$. Takie punkty należą
do $E^{(0)}\cup E^{(1)}$. Ponadto
$\lambda_n(E\triangle E^{(1)})=0$.
Jeżeli $\lambda_n(E\triangle F)=0$, to dla każdej
kuli liczniki dla $E$ i $F$ są identyczne, więc
oba zbiory punktów gęstości i brzegi miarowe
pokrywają się punktowo.
\end{proof}

\begin{definicja}[Brzeg zredukowany]
\label{def:reduced-boundary-29}
Niech $n\ge1$ i niech $E\subseteq\mathbb R^n$
będzie mierzalnym zbiorem lokalnie skończonego obwodu.
\emph{Brzeg zredukowany} $\partial^*E$ składa się
z tych $x\in\mathbb R^n$, dla których
$|D\chi_E|(B(x,r))>0$ dla każdego $r>0$
oraz istnieje granica
\[
 \nu_E^*(x)=-\lim_{r\downarrow0}
 \frac{D\chi_E(B(x,r))}{|D\chi_E|(B(x,r))}
 \quad\text{o normie }1.
\]
To definicja punktowa, nie tylko klasy funkcji
wyznaczonej prawie wszędzie.
\end{definicja}

\begin{lemat}[Mierzalność brzegu zredukowanego i skupienie wariacji]
\label{lem:reduced-boundary-measurable-29}
Niech $E\subseteq\mathbb R^n$, $n\ge1$, będzie
mierzalnym zbiorem lokalnie skończonego obwodu
i niech $\mu=|D\chi_E|$.
Wówczas $\partial^*E$ jest borelowski,
$\mu(\mathbb R^n\setminus\partial^*E)=0$,
a $\nu_E^*$ jest borelowska na $\partial^*E$
i jest równa $\nu_E$ dla $\mu$-prawie wszystkich
punktów. Wyniki te nie wymagają twierdzenia
strukturalnego De Giorgiego.
\end{lemat}
\begin{proof}
Niech $S=\{x:\mu(B(x,r))>0\text{ dla każdego }r>0\}$.
Dopełnienie $S$ jest otwarte: jeśli $\mu(B(x,r))=0$,
to każdy punkt tej kuli ma mniejszą kulę
o mierze zero. Z drugiego aksjomatu przeliczalności
wynika, że dopełnienie $S$ można pokryć przeliczalnie
wieloma kulami miary zero, stąd $\mu(S^c)=0$.

Dla ustalonego $q>0$ funkcje
$x\mapsto\mu(B(x,q))$ oraz $x\mapsto D\chi_E(B(x,q))$
są borelowskie. Dla pierwszej stosujemy mierzalność
całki parametrycznej do borelowskiej funkcji
$(x,y)\mapsto\chi_{\{|x-y|<q\}}(x,y)$
względem $\sigma$-skończonej miary $\mu$.
Dla drugiej stosujemy tę samą własność do składowych
rozkładu polarnego, rozdzielając części dodatnie
i ujemne. Wszystkie całki po ograniczonej kuli są
skończone. Na $S$ definiujemy więc borelowskie pole
$w_q(x)=-D\chi_E(B(x,q))/\mu(B(x,q))$ o normie
nie większej niż $1$.
Dla każdego $x\in S$ i $r>0$, gdy $q_j\uparrow r$,
kule $B(x,q_j)$ rosną do $B(x,r)$.
Ciągłość miar od dołu (dla składowych ze znakiem
stosowana do obu części na ograniczonej kuli)
daje $w_{q_j}(x)\to w_r(x)$.
Zatem kryterium Cauchy'ego dla promieni wymiernych
wystarcza do istnienia granicy przy wszystkich
promieniach $r\downarrow0$.
Dokładniej, zbiór punktów $S$, dla których
granica istnieje, jest równy
\[
 S\cap\bigcap_{l=1}^\infty\bigcup_{m=1}^\infty
 \bigcap_{\substack{p,q\in\mathbb Q\\0<p,q<1/m}}
 \{x\in S:|w_p(x)-w_q(x)|\le1/l\}.
\]
Jest on borelowski. Na nim granica jest borelowska
jako granica ciągu pól $w_{1/j}$.
Warunek normy $1$ wybiera borelowski podzbiór
tego zbioru, dokładnie $\partial^*E$,
a granica jest tam równa $\nu_E^*$.

Na mocy Twierdzenia~\ref{thm:relative-lebesgue-25}
o różniczkowaniu całek względem miary Radona,
zastosowanego lokalnie do każdej składowej
ograniczonego pola $\nu_E$, dla $\mu$-prawie
każdego $x\in S$ zachodzi
\[
 \lim_{r\downarrow0}\frac{1}{\mu(B(x,r))}
 \int_{B(x,r)}\nu_E(y)\,\mathrm d\mu(y)=\nu_E(x).
\]
Lewa strona to właśnie granica $w_r(x)$.
Ponieważ $|\nu_E|=1$ prawie wszędzie względem
$\mu$, punkty te należą do $\partial^*E$
i zachodzi $\nu_E^*(x)=\nu_E(x)$.
Wraz z $\mu(S^c)=0$ daje to tezę.
\end{proof}

\begin{twierdzenie}[Twierdzenie Gaussa--Greena w postaci miarowej]
\label{thm:gauss-green-BV-29}
Niech $n\ge1$, $E\subseteq\mathbb R^n$ będzie
mierzalnym zbiorem lokalnie skończonego obwodu,
a $\varphi\in C_c^1(\mathbb R^n;\mathbb R^n)$.
Wtedy
\begin{equation}\label{eq:gauss-green-BV-29}
 \int_E\operatorname{div}\varphi(x)
 \,\mathrm d\lambda_n(x)
 =\int_{\mathbb R^n}\varphi(x)\cdot\nu_E(x)
 \,\mathrm d|D\chi_E|(x).
\end{equation}
\end{twierdzenie}
\begin{proof}
Definicja dystrybucyjnej pochodnej daje
$\int_E\operatorname{div}\varphi\,\mathrm d\lambda_n=-\int_{\mathbb R^n}\varphi\cdot\mathrm dD\chi_E$.
Podstawiamy rozkład polarny z
Definicji~\ref{def:normal-BV-29} i otrzymujemy
prawą stronę. Obie całki są skończone:
pierwsza ma ograniczoną funkcję podcałkową
o zwartym nośniku, druga jest ograniczona przez
$\|\varphi\|_\infty|D\chi_E|(\operatorname{supp}\varphi)$.
Z poprzedniego stwierdzenia można także zastąpić
prawą stronę całką po $\partial^*E$ względem
$|D\chi_E|$, z polem $\nu_E^*$ zamiast $\nu_E$.
Nie zastępujemy jednak jeszcze miary wariacji
miarą Hausdorffa: wymaga to dalszego twierdzenia
strukturalnego. Nie trzeba zakładać gładkości brzegu.
\end{proof}

\subsection*{Narzędzia do dowodu twierdzenia De Giorgiego}
Główny trud polega na wyprowadzeniu z samej miary $D\chi_E$
zarówno lokalnej płaskości brzegu, jak i jego miary powierzchniowej.
Przedstawimy ten argument w kilku krokach. W tej części zakładamy
$n\ge2$; przypadek $n=1$ omówimy na końcu dowodu.
Przez $\omega_j$ oznaczamy objętość kuli jednostkowej
w $\mathbb R^j$.

\begin{lemat}[Nierówności Sobolewa i względna nierówność izoperymetryczna]
\label{lem:dg-relative-isoperimetric}
Dla każdego $n\ge2$ istnieje $C_n>0$ taka, że
\begin{enumerate}[label=(\roman*)]
\item jeśli $v\in C_c^1(\mathbb R^n)$, to
 $\|v\|_{L^{n/(n-1)}(\mathbb R^n)}\le
 C_n\int_{\mathbb R^n}|\nabla v(x)|\,\mathrm d\lambda_n(x)$;
\item dla każdej kuli $B=B(x,r)$, $x\in\mathbb R^n$,
 $r>0$, i $v\in BV(B)$ zachodzi
 $\|v-v_B\|_{L^{n/(n-1)}(B)}\le C_n|Dv|(B)$,
 gdzie $v_B=\lambda_n(B)^{-1}\int_Bv\,\mathrm d\lambda_n$;
\item dla mierzalnego $E$ o lokalnie skończonym obwodzie
 i dowolnej kuli $B=B(x,r)$ zachodzi
 \[
 \min\{\lambda_n(E\cap B),\lambda_n(B\setminus E)\}^{(n-1)/n}
 \le C_n P(E;B).
 \]
 Ponadto, jeżeli $A\subset\mathbb R^n$ jest mierzalny,
 ma skończoną miarę i skończony obwód, to
 $\lambda_n(A)^{(n-1)/n}\le C_nP(A;\mathbb R^n)$.
\end{enumerate}
\end{lemat}
\begin{proof}
Dla nieujemnych mierzalnych funkcji
$g_i\in L^1(\mathbb R^{n-1})$, $i=1,\ldots,n$,
oznaczamy przez $x_{\widehat i}$ wektor otrzymany
przez usunięcie $i$-tej współrzędnej z $x$.
Zachodzi nierówność
\[
 \int_{\mathbb R^n}\prod_{i=1}^n
 g_i(x_{\widehat i})^{1/(n-1)}\,\mathrm d\lambda_n(x)
 \le\prod_{i=1}^n
 \left(\int_{\mathbb R^{n-1}}g_i(y)\,
 \mathrm d\lambda_{n-1}(y)\right)^{1/(n-1)}.
\]
Oznaczmy lewą stronę przez $I$.
Dla $n=2$ nierówność jest równością Tonellego.
W kroku indukcyjnym, przy $n\ge3$, oznaczamy
$h_i(x'_{\widehat i})=\int_{\mathbb R}
 g_i(x_{\widehat i})\,\mathrm d\lambda_1(x_n)$
przy $i<n$, gdzie $x'=(x_1,\ldots,x_{n-1})$.
Z twierdzenia Tonellego funkcje $h_i$ są mierzalne
i całkowalne, z całkami równymi całkom $g_i$.
Nierówność Höldera dla $n-1$ czynników z jednakowymi
wykładnikami $n-1$ przy całkowaniu po $x_n$ daje
\[
 I\le\int_{\mathbb R^{n-1}}g_n(x')^{1/(n-1)}
 \prod_{i<n}h_i(x'_{\widehat i})^{1/(n-1)}
 \,\mathrm d\lambda_{n-1}(x').
\]
Stosując ponownie Höldera, z wykładnikami
$n-1$ oraz $(n-1)/(n-2)$, a następnie
założenie indukcyjne dla $(n-1)$ funkcji $h_i$,
otrzymujemy następujący rachunek. Oznaczmy
$G_i=\int_{\mathbb R^{n-1}}g_i(y)\,
\mathrm d\lambda_{n-1}(y)$ oraz
\[
 K=\int_{\mathbb R^{n-1}}\prod_{i<n}
 h_i(x'_{\widehat i})^{1/(n-2)}
 \,\mathrm d\lambda_{n-1}(x').
\]
Wtedy
\[
 I\le G_n^{1/(n-1)}K^{(n-2)/(n-1)}
 \le\prod_{i=1}^n G_i^{1/(n-1)}.
\]
To zamyka indukcję.
Dla $v\in C_c^1$ kładziemy
$g_i(x_{\widehat i})=\int_{\mathbb R}
|\partial_i v(x)|\,\mathrm d\lambda_1(x_i)$.
Z zasadniczego twierdzenia rachunku całkowego wynika
$|v(x)|\le g_i(x_{\widehat i})$ dla każdego $i$.
Mnożymy te nierówności i podnosimy do potęgi
$1/(n-1)$, otrzymując
$|v|^{n/(n-1)}\le\prod_i g_i^{1/(n-1)}$.
Całkujemy i stosujemy poprzednie oszacowanie,
a następnie podnosimy obie strony do potęgi
$(n-1)/n$. Wobec $\int g_i=\|\partial_i v\|_1$
oraz $|\partial_i v|\le|\nabla v|$ mamy
$\|v\|_{n/(n-1)}\le\prod_i\|\partial_i v\|_1^{1/n}
\le\|\nabla v\|_1$. To dowodzi (i).

Dla (ii) najpierw rozważamy funkcje klasy $C^1$
na otoczeniu domknięcia kuli jednostkowej $B=B(0,1)$.
Z definicji średniej i nierówności trójkąta
\[
 \|v-v_B\|_{L^1(B)}
 \le\lambda_n(B)^{-1}\int_B\int_B|v(x)-v(y)|
 \,\mathrm d\lambda_n(x)\,\mathrm d\lambda_n(y).
\]
Wypukłość kuli i podstawowe twierdzenie rachunku
całkowego na odcinku dają
$|v(x)-v(y)|\le2\int_0^1
|\nabla v((1-t)x+ty)|\,\mathrm d\lambda_1(t)$.
Dla ustalonego $y$ i $0\le t\le1/2$ odwzorowanie
$x\mapsto(1-t)x+ty$ ma wyznacznik $(1-t)^n$,
a obraz $B$ zawiera się w $B$. Stąd
\[
 \int_B|\nabla v((1-t)x+ty)|\,\mathrm d\lambda_n(x)
 \le2^n\int_B|\nabla v(z)|\,\mathrm d\lambda_n(z).
\]
Dla $1/2\le t\le1$ otrzymujemy to samo oszacowanie
całki po $y$ przy ustalonym $x$ przez zmianę
$y\mapsto(1-t)x+ty$ o wyznaczniku $t^n$.
Tonelli, całkowanie po pozostałej zmiennej oraz
po $t$ dają
$\|v-v_B\|_{L^1(B)}\le2^{n+1}\|\nabla v\|_{L^1(B)}$.

Aby podnieść wykładnik, kładziemy $w=v-v_B$
oraz definiujemy funkcję $\widetilde w$ równą $w$
na domkniętej kuli jednostkowej, a na pierścieniu
$1<|x|<3/2$ równą $w(R(x))$, gdzie
$R(x)=(2-|x|)x/|x|$. Odwzorowanie $R$ jest
dyfeomorfizmem pierścienia na $\{1/2<|y|<1\}$.
Różniczka w kierunku radialnym ma wartość własną
$-1$, a w kierunkach stycznych $(2-|x|)/|x|$,
zawartą w przedziale $(1/3,1)$. W szczególności
$\|DR\|\le1$ oraz $|\det DR|\ge3^{-(n-1)}$.
Wybieramy $\theta\in C_c^\infty(B(0,3/2))$
o wartościach w $[0,1]$, równą $1$ na $\overline B$,
i definiujemy $V=\theta\widetilde w$, przedłużając
przez zero poza kulę o promieniu $3/2$.
Funkcja $V$ jest ciągła, ma zwarty nośnik,
a w obu częściach oddzielonych sferą jednostkową
jest klasy $C^1$ z ograniczonymi pochodnymi.
Jej pochodne dystrybucyjne są równe klasycznym
pochodnym na obu częściach: całkując przez części
w kuli i w pierścieniu (twierdzenie o dywergencji
dla brzegu klasy $C^1$), widzimy, że wyrazy brzegowe
na wspólnej sferze znoszą się, ponieważ wartości
$V$ są tam zgodne, a normalne przeciwne.
Przy zewnętrznym brzegu $V$ znika.
Można też zastosować ten argument do pierścieni
odległych od sfery o $\delta>0$ i przejść
$\delta\downarrow0$, korzystając z ciągłości
i ograniczoności funkcji oraz pochodnych.
Wobec tego $DV=\nabla V\,\lambda_n$, gdzie
gradient jest określony kawałkami, a sfera
ma zerową miarę Lebesgue'a.
Reguła iloczynu i zamiana zmiennych przez $R$
na pierścieniu dają
\[
 \|\nabla V\|_{L^1(\mathbb R^n)}
 \le C_n\bigl(\|\nabla v\|_{L^1(B)}+\|w\|_{L^1(B)}\bigr)
 \le C_n\|\nabla v\|_{L^1(B)}.
\]
Stała zależy tylko od wymiaru i ustalonej funkcji
odcinającej, nie od $v$.
Wygładzenia $V_\varepsilon=V*\rho_\varepsilon$
są gładkie i mają zwarty nośnik; według rachunku
splotów z dowodu Lematu~\ref{lem:BV-translation-29}
$\|\nabla V_\varepsilon\|_1\le\|\nabla V\|_1$,
a $V_\varepsilon\to V$ w $L^1$.
Wybieramy podciąg zbieżny prawie wszędzie.
Z (i) i lematu Fatou dla potęgi $p=n/(n-1)$
wynika $\|V\|_{L^p}\le C_n\|\nabla V\|_{L^1}$.
Ponieważ $V=w$ na $B$, otrzymujemy (ii) dla
funkcji klasy $C^1$ na otoczeniu $\overline B$.

Dla dowolnej kuli $B(x,r)$ stosujemy ten wynik
do funkcji $z\mapsto v(x+rz)$ na kuli jednostkowej.
Jej średnia jest równa $v_{B(x,r)}$. Norma $L^p$
skaluje się przez czynnik $r^{-n/p}=r^{1-n}$,
a norma $L^1$ gradientu przez $r^{1-n}$, co
dowodzi nierówności z tą samą stałą na każdej kuli. Dla $v\in BV(B(x,r))$ ustalamy $0<s<r$.
Wybieramy nieujemne jądro $\rho_\varepsilon$
o całce $1$ i nośniku w $B(0,\varepsilon)$,
gdzie $\varepsilon<r-s$. Funkcję $v$ możemy
przedłużyć przez zero jako funkcję $L^1$,
ale nie zakładamy, że przedłużenie należy do $BV$.
Dla $z\in B(x,s)$ test
$y\mapsto\rho_\varepsilon(z-y)$ ma zwarty
nośnik w $B(x,r)$, dlatego tożsamość splotu
z pochodną miarową zachodzi w tej kuli.
Rozkład polarny i Tonelli dają
\[
 \int_{B(x,s)}|\nabla v_\varepsilon(z)|\,\mathrm d\lambda_n(z)
 \le |Dv|(B(x,s+\varepsilon))
 \le |Dv|(B(x,r)).
\]
W pierwszej nierówności wykorzystujemy, że jądro
jest zerowe dla $|z-y|\ge\varepsilon$, a całka
jądra po $z\in B(x,s)$ nie przekracza $1$.
Funkcja $v_\varepsilon$ jest gładka na otoczeniu
domknięcia $B(x,s)$, więc stosujemy udowodnioną
nierówność na tej kuli. Zbieżność splotu w $L^1$
daje $v_\varepsilon\to v$ w $L^1(B(x,s))$,
a w szczególności średnie
$(v_\varepsilon)_{B(x,s)}\to v_{B(x,s)}$.
Wybierając podciąg zbieżny prawie wszędzie i
stosując lemat Fatou do potęgi $p$, uzyskujemy
\[
 \|v-v_{B(x,s)}\|_{L^p(B(x,s))}
 \le C_n|Dv|(B(x,r)).
\]
Nie potrzebujemy zakładać, że sfera o promieniu
$s$ ma zerową miarę wariacji.
Teraz wybieramy $s_j\uparrow r$.
Ponieważ $v\in L^1(B(x,r))$ i
$\lambda_n(B(x,s_j))\to\lambda_n(B(x,r))$,
średnie $v_{B(x,s_j)}$ zbiegają do $v_{B(x,r)}$.
Funkcje
$\chi_{B(x,s_j)}|v-v_{B(x,s_j)}|^p$ zbiegają
punktowo prawie wszędzie w $B(x,r)$ do
$|v-v_{B(x,r)}|^p$. Ponowne zastosowanie lematu
Fatou daje (ii) z prawą stroną
$C_n|Dv|(B(x,r))$, bez żadnego składnika brzegowego.

W (iii) bierzemy $v=\chi_E$, a w (ii)
$v_B=\theta:=\lambda_n(E\cap B)/\lambda_n(B)$.
Jeśli $\theta\le1/2$, to $|\chi_E-\theta|\ge1/2$
na $E\cap B$, a jeżeli $\theta\ge1/2$, to
$|\chi_E-\theta|\ge1/2$ na $B\setminus E$.
W obu przypadkach
\[
 \|\chi_E-\theta\|_{L^p(B)}
 \ge\frac12\min\{\lambda_n(E\cap B),
                   \lambda_n(B\setminus E)\}^{1/p}.
\]
Po zastosowaniu (ii) i powiększeniu stałej
otrzymujemy względną nierówność izoperymetryczną.

Dla ostatniej nierówności nie możemy bezpośrednio
stosować (i) do wygładzenia $\chi_A$, ponieważ
$A$ nie musi być ograniczony. Niech
$w_\varepsilon=\chi_A*\rho_\varepsilon$.
Z Lematu~\ref{lem:BV-translation-29} i jego dowodu
mamy $\|w_\varepsilon\|_1\le\lambda_n(A)$,
$\|\nabla w_\varepsilon\|_1\le P(A;\mathbb R^n)$
oraz zbieżność $w_\varepsilon\to\chi_A$ w $L^1$.
Wybieramy $\zeta\in C_c^\infty(B(0,2))$
równą $1$ na $B(0,1)$, o wartościach w $[0,1]$,
i kładziemy $\zeta_R(x)=\zeta(x/R)$.
Wówczas $\zeta_R w_\varepsilon\in C_c^1$,
a z (i) i reguły iloczynu wynika
\[
 \|\zeta_R w_\varepsilon\|_{L^p}
 \le C_n\left(P(A;\mathbb R^n)
 +R^{-1}\|\nabla\zeta\|_\infty\lambda_n(A)\right).
\]
Najpierw $R\to\infty$; lemat Fatou daje
$\|w_\varepsilon\|_{L^p}\le C_nP(A;\mathbb R^n)$.
Następnie wybieramy $\varepsilon_j\downarrow0$
z $w_{\varepsilon_j}\to\chi_A$ prawie wszędzie,
co jest możliwe dzięki zbieżności w $L^1$.
Jeszcze raz stosując lemat Fatou, otrzymujemy
$\lambda_n(A)^{1/p}=\|\chi_A\|_{L^p}
\le C_nP(A;\mathbb R^n)$, zgodnie z tezą.
\end{proof}

\begin{lemat}[Przecięcie zbioru skończonego obwodu z kulą]
\label{lem:dg-ball-cut}
Niech $n\ge2$, $E\subseteq\mathbb R^n$ będzie
mierzalnym w sensie Lebesgue'a zbiorem lokalnie
skończonego obwodu, $x\in\mathbb R^n$ oraz
$\mu=|D\chi_E|$. Dla $\lambda_1$-prawie wszystkich
$r>0$ zbiór $E\cap B(x,r)$ ma skończony obwód
globalny i zachodzi równość skończonych miar wektorowych
\[
 D\chi_{E\cap B(x,r)}
 =\chi_{B(x,r)}D\chi_E
 -\frac{y-x}{r}\,\chi_E(y)
 \mathcal H^{n-1}\llcorner\partial B(x,r).
\]
W ostatnim członie możemy użyć dowolnego
borelowskiego reprezentanta $\chi_E$ (dwa takie
reprezentanty zgadzają się $\mathcal H^{n-1}$-prawie
wszędzie na sferach dla prawie wszystkich promieni).
Ponadto funkcje $v(r)=\lambda_n(E\cap B(x,r))$
oraz $w(r)=\lambda_n(B(x,r)\setminus E)$ są
lokalnie bezwzględnie ciągłe na $[0,\infty)$
i spełniają
\[
 v'(r)=\mathcal H^{n-1}(E\cap\partial B(x,r)),\qquad
 w'(r)=\mathcal H^{n-1}(E^c\cap\partial B(x,r))
\]
dla prawie wszystkich $r>0$.
\end{lemat}
\begin{proof}
Wybieramy borelowski zbiór różniący się od $E$
o zbiór miary Lebesgue'a zero i zastępujemy nim
$E$ w rachunku. Nie zmienia to pochodnej
dystrybucyjnej ani funkcji $v,w$.
Dla nieujemnej funkcji borelowskiej $f$ obowiązuje
polarny wzór całkowy
\[
 \int_{\mathbb R^n} f(y)\,\mathrm d\lambda_n(y)
 =\int_0^\infty\int_{\partial B(x,t)} f(y)
 \,\mathrm d\mathcal H^{n-1}(y)\,\mathrm d\lambda_1(t).
\]
Jest to wzór koobszaru z Rozdziału~\ref{ch:area-coarea-28}
dla funkcji $y\mapsto|y-x|$, której gradient ma normę
$1$ poza singletonem $\{x\}$.
Zastosowanie go do funkcji charakterystycznej
różnicy symetrycznej dwóch borelowskich reprezentantów
pokazuje ich zgodność prawie wszędzie na sferach
dla prawie wszystkich promieni. W dalszym rachunku
ustalamy jeden reprezentant.

Wybieramy nieujemną $\kappa\in C_c^\infty((0,1))$
o całce $1$ i definiujemy
$\eta(s)=\int_s^\infty\kappa(t)\,\mathrm d\lambda_1(t)$.
Wtedy $\eta=1$ na $(-\infty,0]$, $\eta=0$
na $[1,\infty)$ oraz $\eta'=-\kappa$.
Dla $r>0$ i $\varepsilon>0$ kładziemy
$\zeta_\varepsilon(y)=\eta((|y-x|-r)/\varepsilon)$.
Funkcja jest gładka także w otoczeniu $x$,
gdyż tam jest stała, i ma zwarty nośnik.
Dla każdego pola $\varphi\in C_c^1(\mathbb R^n;\mathbb R^n)$
reguła iloczynu dla dywergencji daje bezpośrednio
\begin{align*}
 -\int \chi_E\zeta_\varepsilon\operatorname{div}\varphi
 \,\mathrm d\lambda_n
 &=-\int\chi_E\operatorname{div}(\zeta_\varepsilon\varphi)
 \,\mathrm d\lambda_n
 +\int\chi_E\varphi\cdot\nabla\zeta_\varepsilon
 \,\mathrm d\lambda_n\\
 &=\int\zeta_\varepsilon\varphi\cdot\mathrm dD\chi_E
 +\int\chi_E\varphi\cdot\nabla\zeta_\varepsilon
 \,\mathrm d\lambda_n.
\end{align*}
Wszystkie całki bez jawnej dziedziny w tym rachunku
są po $\mathbb R^n$.
Nie trzeba wygładzać $\chi_E$ ani zakładać
zbieżności jej wygładzonych pochodnych.
Po lewej stronie zbieżność zdominowana daje
granicę $-\int_{E\cap B(x,r)}\operatorname{div}\varphi
\,\mathrm d\lambda_n$, ponieważ sfera ma zerową
miarę Lebesgue'a.

Zbiór promieni $r>0$, dla których
$\mu(\partial B(x,r))>0$, jest co najwyżej
przeliczalny. Istotnie, na każdej kuli domkniętej
$\overline B(x,M)$ miara $\mu$ jest skończona,
a sfery są rozłączne. Dla $r\le M$ tylko skończenie
wiele sfer może mieć masę co najmniej $1/k$,
dla ustalonego $k\ge1$. Sumujemy po $M,k\in\mathbb N$.
Dla pozostałych promieni pierwszy składnik po
prawej stronie zbiega do
$\int_{B(x,r)}\varphi\cdot\mathrm dD\chi_E$,
na mocy zbieżności zdominowanej względem
lokalnie skończonej wariacji $\mu$.

Dla drugiego składnika definiujemy funkcję
\[
 F_\varphi(t)=\int_{\partial B(x,t)}
 \chi_E(y)\varphi(y)\cdot\frac{y-x}{t}
 \,\mathrm d\mathcal H^{n-1}(y),\qquad t>0.
\]
Jest ona mierzalna, co widać po parametryzacji
sfer przez $y=x+t\theta$, $\theta\in\partial B(0,1)$,
ze skalowaniem miary powierzchniowej przez $t^{n-1}$,
i zastosowaniu Fubiniego. Ponadto
$|F_\varphi(t)|\le\|\varphi\|_\infty n\omega_n t^{n-1}$.
Zatem $F_\varphi$ jest lokalnie całkowalna.
Ze wzoru na gradient $\zeta_\varepsilon$ oraz
polarnego wzoru całkowego otrzymujemy
\[
 \int\chi_E\varphi\cdot\nabla\zeta_\varepsilon
 \,\mathrm d\lambda_n
 =-\int_r^{r+\varepsilon}
 \frac1\varepsilon\kappa((t-r)/\varepsilon)
 F_\varphi(t)\,\mathrm d\lambda_1(t).
\]
Dla każdego punktu Lebesgue'a $r>0$ funkcji
$F_\varphi$ różnica między prawą stroną a
$-F_\varphi(r)$ ma moduł nie większy niż
\[
 \frac{\|\kappa\|_\infty}{\varepsilon}
 \int_r^{r+\varepsilon}|F_\varphi(t)-F_\varphi(r)|
 \,\mathrm d\lambda_1(t)\longrightarrow0.
\]
Zbieżność wynika z twierdzenia o różniczkowaniu
Lebesgue'a w wymiarze jeden (średnią jednostronną
ograniczamy przez dwukrotność średniej symetrycznej).

Aby zbiór dobrych promieni nie zależał od wyboru
pola testowego, wybieramy przeliczalną rodzinę
$\mathcal T\subset C_c^1(\mathbb R^n;\mathbb R^n)$
taką, że każde pole testowe można przybliżyć
w normie $C^1$ przez pola z $\mathcal T$,
o nośnikach we wspólnej kuli.
Istnienie takiej rodziny wynika następująco.
Dla każdego całkowitego $M\ge1$ przestrzeń pól
$C_c^1$ o nośniku w $\overline B(0,M)$
z normą $C^1$ zanurza się izometrycznie
(przy normie sumy norm supremum) w skończony
produkt przestrzeni $C(\overline B(0,M))$,
przypisując polu jego składowe oraz ich pierwsze
pochodne. Przestrzeń $C$ na zwartym podzbiorze
euklidesowym jest ośrodkowa, np. dzięki jednostajnej
aproksymacji wielomianami o współczynnikach wymiernych.
Podprzestrzeń przestrzeni metrycznej ośrodkowej
jest ośrodkowa. Wybieramy więc przeliczalny
gęsty podzbiór dla każdego $M$ i bierzemy ich sumę.
Usuwamy z $(0,\infty)$ przeliczalny zbiór promieni
o dodatniej masie sfery oraz zbiory punktów,
które nie są punktami Lebesgue'a $F_\varphi$,
dla wszystkich $\varphi\in\mathcal T$.
Pozostaje zbiór pełnej miary.

Dla każdego promienia z tego zbioru otrzymana
tożsamość zachodzi jednocześnie dla wszystkich
$\varphi\in\mathcal T$. Przenosimy ją na dowolne
pole testowe przez aproksymację w normie $C^1$.
Lewa strona jest ciągła w tej normie, z oszacowaniem
przez $\lambda_n(B(x,r))$ razy normę supremum
dywergencji; obie całki miarowe są ciągłe nawet
w normie supremum, z oszacowaniami odpowiednio
przez $\mu(B(x,r))$ i $n\omega_n r^{n-1}$.
W konsekwencji pochodną dystrybucyjną
$\chi_{E\cap B(x,r)}$ reprezentuje skończona
miara wektorowa podana w tezie. Tożsamość testowa
daje skończoność wariacji, więc twierdzenie o
reprezentacji i jego jednoznaczność identyfikują
tę miarę z $D\chi_{E\cap B(x,r)}$.

Na koniec polarny wzór całkowy daje dla każdego $r>0$
\[
 v(r)=\int_0^r\mathcal H^{n-1}(E\cap\partial B(x,t))
 \,\mathrm d\lambda_1(t)
\]
i analogiczny wzór dla $w$, z $E^c$ zamiast $E$.
Funkcje podcałkowe są mierzalne i ograniczone przez
$n\omega_n t^{n-1}$, zatem są całkowalne na każdym
ograniczonym przedziale. Podstawowe twierdzenie
rachunku całkowego dla całki Lebesgue'a daje
lokalną bezwzględną ciągłość i podane wzory
na pochodne dla prawie wszystkich $r>0$.
Przecinając z wcześniejszym zbiorem pełnej miary,
otrzymujemy wszystkie te tożsamości jednocześnie.
\end{proof}

\begin{lemat}[Oszacowania gęstości na brzegu zredukowanym]
\label{lem:dg-density}
Niech $n\ge2$, $E\subseteq\mathbb R^n$ będzie
mierzalnym zbiorem lokalnie skończonego obwodu,
$\mu=|D\chi_E|$, a $x\in\partial^*E$.
Istnieją liczby $r_x>0$ i $0<c_n<C_n<\infty$,
zależne wyłącznie od wymiaru poza $r_x$, takie że
\[
 c_nr^{n-1}\le\mu(B(x,r))\le C_nr^{n-1},\qquad
 \lambda_n(E\cap B(x,r))\ge c_nr^n,\quad
 \lambda_n(E^c\cap B(x,r))\ge c_nr^n
\]
dla każdego $0<r<r_x$.
\end{lemat}
\begin{proof}
Niech $\nu_0=\nu_E^*(x)$ oraz
$D\chi_E=-\nu_E\mu$. Z definicji brzegu
zredukowanego otrzymujemy
\[
 \frac{1}{\mu(B(x,r))}\int_{B(x,r)}\nu_E(y)
 \,\mathrm d\mu(y)\longrightarrow\nu_0.
\]
Ponieważ $|\nu_E|=|\nu_0|=1$,
\[
 \frac1{\mu(B(x,r))}\int_{B(x,r)}
 |\nu_E(y)-\nu_0|^2\,\mathrm d\mu(y)
 =2-2\nu_0\cdot\frac{\int_{B(x,r)}\nu_E(y)\,\mathrm d\mu(y)}
 {\mu(B(x,r))}\longrightarrow0.
\]
W szczególności, dla małych $r$ mamy
$\int_{B(x,r)}\nu_0\cdot\nu_E(y)\,\mathrm d\mu(y)
\ge\frac12\mu(B(x,r))$.

Ustalmy $r_x>0$ tak, aby oszacowanie z czynnikiem
$1/2$ zachodziło dla wszystkich $0<r<r_x$.
Przez $v,w$ oznaczamy funkcje objętości z
Lematu~\ref{lem:dg-ball-cut}.
Dla promieni dobrych w tym lemacie zarówno dla
$E$, jak i dla $E^c$ mamy skończone miary pochodnych
obu przecięć, o nośniku w $\overline B(x,r)$.
Ich całkowite masy wektorowe są zerowe:
dla dowolnego stałego wektora $e$ wybieramy
$\theta\in C_c^\infty(\mathbb R^n)$ równą $1$
na otoczeniu $\overline B(x,r)$.
W tożsamości testowej dla pola $\theta e$
dywergencja znika na obu przecięciach,
a całka względem miary pochodnej jest równa
iloczynowi skalarnemu $e$ z jej masą wektorową.
Z dowolności $e$ wynika zerowanie masy.

Wobec $D\chi_E=-\nu_E\mu$ wzór przecięcia daje
\[
 \int_{B(x,r)}\nu_E\,\mathrm d\mu
 =-\int_{\partial B(x,r)}\chi_E(y)\frac{y-x}{r}
 \,\mathrm d\mathcal H^{n-1}(y).
\]
Iloczyn skalarny z $\nu_0$, a następnie
$|\nu_0\cdot(y-x)/r|\le1$, dają
$\mu(B(x,r))/2\le v'(r)$ dla prawie wszystkich
$r\in(0,r_x)$. Analogicznie, ponieważ
$D\chi_{E^c}=-D\chi_E$, normalna dla dopełnienia
jest przeciwna, a w analogicznym oszacowaniu
używamy $-\nu_0$. Otrzymujemy
$\mu(B(x,r))/2\le w'(r)$ na wspólnym zbiorze
promieni pełnej miary.
Ponieważ $v'(r),w'(r)\le n\omega_n r^{n-1}$,
mamy $\mu(B(x,r))\le2n\omega_n r^{n-1}$
dla prawie wszystkich małych $r$.
Dla dowolnego $r\in(0,r_x)$ wybieramy ciąg
dobrych promieni $r_j\downarrow r$ z $r_j<r_x$.
Monotoniczność $\mu(B(x,r))\le\mu(B(x,r_j))$
i przejście do granicy w oszacowaniu liczbowym
dają tę samą nierówność dla każdego takiego $r$.

Dla dobrego promienia miara pochodnej przecięcia
jest sumą miary ograniczonej do otwartej kuli
i miary powierzchniowej na sferze. Z nierówności
trójkąta dla wariacji wynika
$P(E\cap B(x,r);\mathbb R^n)\le
\mu(B(x,r))+v'(r)$. Objętość przecięcia jest
skończona, więc globalna nierówność izoperymetryczna
z Lematu~\ref{lem:dg-relative-isoperimetric} daje
\[
 v(r)^{(n-1)/n}\le C_n
 \bigl(\mu(B(x,r))+v'(r)\bigr)
 \le C'_n v'(r).
\]
Analogicznie $w(r)^{(n-1)/n}\le C'_nw'(r)$
dla prawie wszystkich $r\in(0,r_x)$.
Z warunku $\mu(B(x,r))>0$ wynika, że
$v(r),w(r)>0$ dla każdego $r>0$.
W przeciwnym razie $\chi_E$ byłaby równa prawie
wszędzie odpowiednio $0$ lub $1$ w tej kuli.
Całka z dywergencji każdego pola testowego
o nośniku w kuli byłaby wtedy zerowa,
a z definicji wariacji i tożsamości na otwartych
zbiorach wynikałoby $\mu(B(x,r))=0$.

Ustalmy $0<\varepsilon<r<r_x$.
Funkcja $v$ jest bezwzględnie ciągła i niemalejąca,
a $v(\varepsilon)>0$. Na przedziale
$[\varepsilon,r]$ funkcja $v^{1/n}$ jest
bezwzględnie ciągła, gdyż składamy $v$ z funkcją
klasy $C^1$ o ograniczonej pochodnej na
$[v(\varepsilon),v(r)]$. Reguła łańcuchowa dla
funkcji bezwzględnie ciągłych daje
\[
 (v^{1/n})'(t)=\frac1n
 v(t)^{-(n-1)/n}v'(t)\ge\frac1{nC'_n}
\]
dla prawie wszystkich $t\in(\varepsilon,r)$.
Całkowanie daje
$v(r)^{1/n}\ge v(\varepsilon)^{1/n}+
(r-\varepsilon)/(nC'_n)$. Ponieważ
$0\le v(\varepsilon)\le\omega_n\varepsilon^n$,
przejście $\varepsilon\downarrow0$ daje
$v(r)\ge(r/(nC'_n))^n$. Identyczny rachunek
dla $w$ daje $w(r)\ge(r/(nC'_n))^n$.
Obie nierówności zachodzą dla każdego
$r\in(0,r_x)$, mimo że nierówności różniczkowe
zachodziły jedynie prawie wszędzie.
Względna nierówność izoperymetryczna z
Lematu~\ref{lem:dg-relative-isoperimetric}
daje teraz $\mu(B(x,r))\ge c_nr^{n-1}$.
\end{proof}

\begin{lemat}[Półprzestrzeń jako jedyna granica po przeskalowaniu]
\label{lem:dg-blowup}
Niech $n\ge2$, $E\subseteq\mathbb R^n$ będzie
mierzalnym zbiorem lokalnie skończonego obwodu,
$\mu=|D\chi_E|$, $x\in\partial^*E$,
$\nu=\nu_E^*(x)$ oraz, dla $r>0$,
$E_{x,r}=\{z\in\mathbb R^n:x+rz\in E\}$.
Wtedy przy $r\downarrow0$
\[
 \chi_{E_{x,r}}\longrightarrow\chi_{H_\nu}
 \quad\text{w }L^1_{\rm loc}(\mathbb R^n),\qquad
 H_\nu=\{z\in\mathbb R^n:z\cdot\nu<0\}.
\]
Ponadto dodatnie miary Radona określone dla
borelowskich $A\subseteq\mathbb R^n$ wzorem
$\mu_{x,r}(A)=r^{1-n}\mu(x+rA)$
zbiegają do $\mathcal H^{n-1}\llcorner\nu^\perp$
względem testów z $C_c(\mathbb R^n)$, a więc
\[
 \lim_{r\downarrow0}
 \frac{\mu(B(x,r))}{\omega_{n-1}r^{n-1}}=1.
\]
\end{lemat}
\begin{proof}
\emph{Krok 1: skalowanie.}
Niech $T_r(z)=x+rz$ i $u_r=\chi_{E_{x,r}}$.
Dla $\varphi\in C_c^1(\mathbb R^n;\mathbb R^n)$
kładziemy $\psi(y)=\varphi((y-x)/r)$.
Mamy $\operatorname{div}\psi(y)=
r^{-1}\operatorname{div}\varphi((y-x)/r)$.
Zmiana zmiennych $y=T_r(z)$ daje
\begin{align*}
 -\int_{\mathbb R^n}u_r(z)\operatorname{div}\varphi(z)
 \,\mathrm d\lambda_n(z)
 &=-r^{1-n}\int_E\operatorname{div}\psi(y)
 \,\mathrm d\lambda_n(y)\\
 &=r^{1-n}\int_{\mathbb R^n}\psi\cdot\mathrm dD\chi_E.
\end{align*}
Prawa strona reprezentuje lokalnie skończoną
miarę wektorową, więc z jednoznaczności
reprezentacji i definicji wariacji wynika
$u_r\in BV_{\rm loc}(\mathbb R^n)$ oraz
\[
 Du_r=-\nu_E(T_r(\cdot))\mu_{x,r},
 \qquad |Du_r|=\mu_{x,r}.
\]
W ostatniej równości wykorzystujemy
$|\nu_E\circ T_r|=1$ prawie wszędzie względem
$\mu_{x,r}$ i wzór na wariację z
Twierdzenia~\ref{thm:BV-measure-29}.

\emph{Krok 2: lokalna zwartość.}
Weźmy dowolny ciąg dodatnich promieni $r_j\to0$.
Dla każdego ustalonego $R>0$ Lemat~\ref{lem:dg-density}
daje
$\mu_{x,r_j}(B(0,R))\le C_nR^{n-1}$
dla wszystkich dostatecznie dużych $j$,
gdyż $Rr_j<r_x$. Skończenie wiele pozostałych
wyrazów również ma skończoną masę na tej kuli,
więc supremum po wszystkich $j$ jest skończone.

Dla całkowitego $R\ge1$ wybieramy
$\theta_R\in C_c^\infty(B(0,R+1))$,
$0\le\theta_R\le1$, równą $1$ na $\overline B(0,R)$.
Niech $v_j=\theta_R u_{r_j}$.
Z dystrybucyjnej reguły iloczynu (testujemy
$Du_{r_j}$ polem $\theta_R\varphi$) wynika
$Dv_j=\theta_R Du_{r_j}+u_{r_j}\nabla\theta_R\,\lambda_n$.
Stąd $v_j\in BV(\mathbb R^n)$ i
\[
 \|v_j\|_1\le\lambda_n(B(0,R+1)),\qquad
 |Dv_j|(\mathbb R^n)
 \le\mu_{x,r_j}(B(0,R+1))+\|\nabla\theta_R\|_1.
\]
W szczególności obie normy są jednostajnie
ograniczone. Funkcje $v_j$ mają wartości w $[0,1]$
i znikają poza wspólną ograniczoną kulą.
Argument z dowodu Twierdzenia~\ref{thm:perimeter-compact-29}
stosuje się do nich aż do uzyskania zbieżności
w $L^1$: średnie na kostkach należą do $[0,1]$,
oszacowanie błędu zależy tylko od wariacji,
a argument przekątniowy i kryterium Cauchy'ego
nie wymagają dwuwartościowości.
Otrzymujemy podciąg zbieżny na $B(0,R)$,
gdzie $v_j=u_{r_j}$. Wybierając zagnieżdżone
podciągi dla $R=1,2,\ldots$ i ich przekątną,
dostajemy zbieżność $u_{r_j}$ w $L^1$ na każdej
z tych kul. Granice na mniejszych kulach są
ograniczeniami granic na większych z jednoznaczności
granicy $L^1$. Określają zatem mierzalną funkcję $u$
na $\mathbb R^n$ z dokładnością do zbioru zerowego.
Na każdej kuli całka z
$\operatorname{dist}(u,\{0,1\})$ jest ograniczona
przez $\|u-u_{r_j}\|_{L^1}$ na tej kuli,
więc jest zerowa. Ponieważ kul jest przeliczalnie
wiele, $u=\chi_F$ prawie wszędzie dla pewnego
mierzalnego $F$.

Z tych samych lokalnych ograniczeń masy możemy
wybrać dalszy podciąg miar $\mu_{x,r_j}$ zbieżny
względem $C_c(\mathbb R^n)$ do dodatniej miary Radona
$\sigma$. Dla pełności, na każdej kuli domkniętej
$\overline B(0,M)$ przestrzeń funkcji ciągłych
o nośniku w tej kuli jest ośrodkowa w normie
supremum. Wybieramy przeliczalną gęstą rodzinę,
a następnie podciąg, na którym całki wszystkich
funkcji z sumy tych rodzin są zbieżne.
Jest to możliwe przez argument przekątniowy,
gdyż każda z tych całek jest ciągiem ograniczonym.
Jednostajne lokalne ograniczenia mas pozwalają
przenieść zbieżność na dowolny test ciągły
o zwartym nośniku. Granica całek określa dodatni
funkcjonał liniowy na $C_c$, ograniczony na każdej
przestrzeni testów o nośniku w ustalonej kuli
przez stałą razy normę supremum.
Twierdzenie Riesza--Markowa--Kakutaniego
(Rozdział~\ref{ch:riesz-markov-21}) reprezentuje
ten funkcjonał przez lokalnie skończoną dodatnią
miarę Radona $\sigma$. Dalszy podciąg zachowuje
wcześniejszą zbieżność w $L^1_{\rm loc}$.

\emph{Krok 3: kierunek pochodnej granicy.}
W dowodzie Lematu~\ref{lem:dg-density} wykazano
\[
 \frac{\int_{B(x,s)}|\nu_E-\nu|^2\,\mathrm d\mu}
 {\mu(B(x,s))}\longrightarrow0
 \qquad(s\downarrow0).
\]
Dla każdego $R>0$ nierówność Cauchy'ego--Schwarza
i definicja przeskalowanej miary dają
\[
 \int_{B(0,R)}|\nu_E\circ T_{r_j}-\nu|
 \,\mathrm d\mu_{x,r_j}
 \le \mu_{x,r_j}(B(0,R))
 \left(
 \frac{\int_{B(x,Rr_j)}|\nu_E-\nu|^2\,\mathrm d\mu}
 {\mu(B(x,Rr_j))}
 \right)^{1/2}\longrightarrow0.
\]
Pierwszy czynnik jest jednostajnie ograniczony,
a mianownik jest dodatni z definicji $\partial^*E$.
Dla dowolnego pola testowego $\varphi$ wybieramy
$R$ z nośnikiem $\varphi$ w $B(0,R)$.
W tożsamości testowej dla $Du_{r_j}$ lewa
strona zbiega do odpowiedniej całki dla $\chi_F$
z lokalnej zbieżności w $L^1$.
Prawa różni się od
$-\int\varphi\cdot\nu\,\mathrm d\mu_{x,r_j}$
o składnik, którego moduł jest ograniczony przez
$\|\varphi\|_\infty$ razy powyższa całka.
Zbieżność miar względem $C_c$ daje
$D\chi_F=-\nu\sigma$ jako dystrybucje.
Miara po prawej jest lokalnie skończona, więc
$\chi_F\in BV_{\rm loc}$ i jest to równość
lokalnych miar wektorowych.

\emph{Krok 4: identyfikacja zbioru granicznego.}
Po ortogonalnej zmianie współrzędnych możemy
przyjąć $\nu=e_n$. Wtedy $D_i\chi_F=0$ dla $i<n$
oraz $D_n\chi_F=-\sigma$. Niech $\rho_\varepsilon$
będzie nieujemnym gładkim jądrem o masie $1$
i nośniku w $B(0,\varepsilon)$.
Funkcja $f_\varepsilon=\chi_F*\rho_\varepsilon$
jest gładka i ma wartości w $[0,1]$.
Choć $\chi_F$ nie musi być globalnie całkowalna,
rachunek splotu z miarą działa lokalnie:
dla ustalonego punktu $z$ test
$\rho_\varepsilon(z-\cdot)$ ma zwarty nośnik,
a różniczkowanie pod całką wymaga jedynie
całkowalności na wspólnym ograniczonym otoczeniu.
Dostajemy $\partial_i f_\varepsilon=0$ dla $i<n$
oraz $\partial_n f_\varepsilon=-\sigma*\rho_\varepsilon\le0$.
Pierwsze równości, przez całkowanie wzdłuż
odcinków w płaszczyznach o stałym $z_n$, dają
$f_\varepsilon(z',t)=h_\varepsilon(t)$.
Druga nierówność pokazuje, że $h_\varepsilon$
jest nierosnąca.

Sploty zbiegają do $\chi_F$ w $L^1_{\rm loc}$:
na dowolnym zwartym zbiorze zastępujemy $\chi_F$
funkcją równą jej na większej kuli i znikającą poza
tą kulą, po czym stosujemy zbieżność splotów w $L^1$.
Wybieramy jednostkową kostkę $C\subset\mathbb R^{n-1}$
i definiujemy $h(t)=\int_C\chi_F(z',t)\,
\mathrm d\lambda_{n-1}(z')$ dla prawie wszystkich $t$.
Dla każdego ograniczonego przedziału $I$ mamy
\[
 \|h_\varepsilon-h\|_{L^1(I)}
 \le\|f_\varepsilon-\chi_F\|_{L^1(C\times I)}
 \longrightarrow0.
\]
Dla dowolnej ograniczonej kostki $C'\subset\mathbb R^{n-1}$
nierówność trójkąta daje
\[
 \|\chi_F-h\|_{L^1(C'\times I)}
 \le\|\chi_F-f_\varepsilon\|_{L^1(C'\times I)}
 +\lambda_{n-1}(C')\|h_\varepsilon-h\|_{L^1(I)}
 \longrightarrow0.
\]
Przeliczalne wyczerpanie kostkami i przedziałami
pokazuje, że $\chi_F(z',t)=h(t)$ prawie wszędzie.
Z Fubiniego wynika $h(t)\in\{0,1\}$ prawie wszędzie.
Wybieramy ciąg $\varepsilon_l\downarrow0$,
dla którego $h_{\varepsilon_l}\to h$ prawie wszędzie
na $\mathbb R$ (wybierając przekątniowo
podciągi zbieżne prawie wszędzie na $[-m,m]$).
Na zbiorze $T$ pełnej miary, na którym zachodzi
ta zbieżność i $h$ jest dwuwartościowa, dla $s<t$
mamy $h(s)\ge h(t)$ jako granicę nierówności
dla $h_{\varepsilon_l}$. Jeżeli obie wartości
występują na $T$, to wszystkie punkty $T$ o wartości
$1$ leżą przed wszystkimi punktami o wartości $0$.
Kładąc $a=\sup\{t\in T:h(t)=1\}$, dostajemy
$a\in\mathbb R$ i $h=\chi_{(-\infty,a)}$
prawie wszędzie. Jeśli jedna wartość nie występuje,
$h$ jest stała prawie wszędzie.
Po powrocie do pierwotnych współrzędnych $F$
jest zatem albo półprzestrzenią $\{z:z\cdot\nu<a\}$,
albo zbiorem pustym, albo całą przestrzenią,
z dokładnością do zbiorów zerowych.

Dla dowolnego $R>0$ dolne oszacowania objętości
z Lematu~\ref{lem:dg-density}, zastosowane
dla promieni $Rr_j$, dają
\[
 \lambda_n(E_{x,r_j}\cap B(0,R))\ge c_nR^n,
 \qquad
 \lambda_n(E_{x,r_j}^c\cap B(0,R))\ge c_nR^n
\]
dla wszystkich dostatecznie dużych $j$.
Lokalna zbieżność w $L^1$ przenosi obie nierówności
na $F$. Wyklucza to oba przypadki stałe.
Gdyby $a\ne0$, kula o promieniu $|a|/2$
leżałaby całkowicie po jednej stronie hiperpłaszczyzny
$\{z:z\cdot\nu=a\}$, co przeczy jednej z nierówności.
Zatem $a=0$ i $F=H_\nu$ prawie wszędzie.

\emph{Krok 5: identyfikacja miary i pełnej granicy.}
Dla $\nu=e_n$ całki z pochodnych stycznych
pola testowego po półprzestrzeni znikają przez
Fubiniego i całkowanie po całej prostej.
Całkowanie ostatniej pochodnej od $-\infty$
do $0$ daje całkę z ostatniej składowej pola
na $\{z_n=0\}$. Z kalibracji miary Hausdorffa
i znaku pochodnej dystrybucyjnej otrzymujemy
$D\chi_{H_{e_n}}=-e_n\mathcal H^{n-1}
\llcorner\{z_n=0\}$. Ortogonalna zmiana współrzędnych
daje dla dowolnego $\nu$
$D\chi_{H_\nu}=-\nu\mathcal H^{n-1}\llcorner\nu^\perp$.
Porównując z $D\chi_F=-\nu\sigma$ i biorąc
iloczyn skalarny z $-\nu$, dostajemy
$\sigma=\mathcal H^{n-1}\llcorner\nu^\perp$.

Powyższy argument zaczynał się od dowolnego ciągu
$r_j\to0$ i dawał dalszy podciąg z obiema
wskazanymi granicami. Gdyby zbieżność w $L^1$
na pewnej ustalonej kuli nie zachodziła, można
byłoby wybrać ciąg promieni dążących do zera,
na którym odległość od $\chi_{H_\nu}$ jest
ograniczona od dołu dodatnią stałą. Otrzymany
dalszy podciąg przeczyłby temu. Analogicznie,
dla każdego testu z $C_c$, brak zbieżności
całek pozwalałby wybrać ciąg sprzeczny ze
zbieżnością na dalszym podciągu.
Daje to obie granice przy $r\downarrow0$.

Niech $\sigma=\mathcal H^{n-1}\llcorner\nu^\perp$.
Dla $0<\delta<1$ wybieramy ciągłe funkcje
odcinające o wartościach w $[0,1]$,
z których pierwsza jest równa $1$ na
$\overline B(0,1-\delta)$ i znika poza $B(0,1)$,
a druga jest równa $1$ na $\overline B(0,1)$
i znika poza $B(0,1+\delta)$.
Ograniczają one $\chi_{B(0,1)}$ odpowiednio
od dołu i od góry. Zbieżność całek testowych daje
\[
 \omega_{n-1}(1-\delta)^{n-1}
 \le\liminf_{r\downarrow0}\mu_{x,r}(B(0,1))
 \le\limsup_{r\downarrow0}\mu_{x,r}(B(0,1))
 \le\omega_{n-1}(1+\delta)^{n-1}.
\]
Tutaj wykorzystujemy kalibrację $\sigma$ na
hiperpłaszczyźnie oraz jej zerową masę na sferach
o dodatnim promieniu, co wynika z zerowej
$(n-1)$-wymiarowej objętości sfery w tej
hiperpłaszczyźnie. Przez $\delta\downarrow0$
otrzymujemy $\mu_{x,r}(B(0,1))\to\omega_{n-1}$.
Ponieważ $\mu_{x,r}(B(0,1))=r^{1-n}\mu(B(x,r))$,
dowodzi to ostatniej tezy.
\end{proof}

\begin{lemat}[Prostowalność brzegu zredukowanego]
\label{lem:dg-rectifiability}
Niech $n\ge2$, a $E\subseteq\mathbb R^n$ będzie
mierzalnym zbiorem lokalnie skończonego obwodu.
Wówczas cały zbiór $\partial^*E$ zawiera się
w przeliczalnej sumie pełnych wykresów funkcji
Lipschitzowskich nad $(n-1)$-wymiarowymi
hiperpłaszczyznami. W szczególności jest
przeliczalnie $(n-1)$-prostowalny.
\end{lemat}
\begin{proof}
Niech $S=\partial^*E$. Ustalmy $0<\alpha<1/4$
i skończoną siatkę $e_1,\ldots,e_N$ na sferze
jednostkowej, taką, że każda jednostkowa normalna
ma odległość mniejszą niż $\alpha/8$ od któregoś
$e_i$. Istnienie siatki wynika ze zwartości sfery.
Ustalamy $\delta=\omega_n(\alpha/16)^n/8>0$.
Z Lematu~\ref{lem:dg-blowup} każdy $x\in S$
spełnia, dla dostatecznie małych $r$,
\[
 \lambda_n\bigl((E\triangle(x+H_{\nu_E^*(x)}))
 \cap B(x,2r)\bigr)\le\delta r^n,
\]
co wynika z zastosowania zbieżności
$L^1_{\rm loc}$ na $B(0,2)$ i zmiany zmiennych
$y=x+rz$. Pokryjmy zatem $S$ przeliczalną rodziną
zbiorów $S_{i,k,m}$, na których $|x|\le k$,
$|\nu_E^*(x)-e_i|<\alpha/8$ oraz powyższe
oszacowanie zachodzi dla każdego $0<r<1/m$.
Nie pomijamy żadnego punktu $S$.

Pokażemy, że jeśli $x,y\in S_{i,k,m}$ i
$0<|x-y|<1/(2m)$, to
$|(y-x)\cdot e_i|<\alpha|y-x|$.
Gdyby nie, połóżmy $r=|x-y|$ i
$\rho=\alpha r/16$. Oznaczając $\nu_x=\nu_E^*(x)$,
otrzymujemy
\[
 |(y-x)\cdot\nu_x|
 \ge |(y-x)\cdot e_i|-r|\nu_x-e_i|
 >7\alpha r/8.
\]
Dla $z\in B(y,\rho)$ zmiana iloczynu skalarnego
z $\nu_x$ jest co do modułu mniejsza niż $\rho$,
więc cała kula leży ściśle po jednej stronie
hiperpłaszczyzny $x+\nu_x^\perp$.
Ponadto $B(y,\rho)\subset B(x,2r)$.
Niech $V=\lambda_n(B(y,\rho))=\omega_n\rho^n$.
Jeśli kula leży w półprzestrzeni $x+H_{\nu_x}$,
to $\lambda_n(B(y,\rho)\setminus E)\le\delta r^n$;
jeśli leży po przeciwnej stronie, to
$\lambda_n(E\cap B(y,\rho))\le\delta r^n$.
Zatem minimum tych dwóch objętości jest
nie większe niż $\delta r^n=V/8$.

Z drugiej strony w punkcie $y$ stosujemy
oszacowanie przy skali $\rho/2<1/m$.
Półprzestrzeń $y+H_{\nu_E^*(y)}$ dzieli
$B(y,\rho)$ na dwie części o tej samej objętości.
Objętość różnicy symetrycznej z $E$ w tej kuli
jest nie większa niż $\delta(\rho/2)^n$.
Stąd każda z obu objętości jest nie mniejsza
niż $V/2-\delta(\rho/2)^n\ge3V/8$,
gdzie wykorzystano $\delta(\rho/2)^n\le V/8$.
Jest to sprzeczność z poprzednim oszacowaniem.
Dowodzi to żądanej nierówności stożkowej.

W konsekwencji rzut ortogonalny
$\pi_i:\mathbb R^n\to e_i^\perp$ jest
różnowartościowy na każdym fragmencie
$S_{i,k,m}$ o średnicy mniejszej niż $1/(2m)$.
Istotnie, dla różnicy $d=y-x$ z nierówności
stożkowej i twierdzenia Pitagorasa wynika
$|\pi_i(d)|\ge\sqrt{1-\alpha^2}|d|$.
Jeśli rzuty się pokrywają, musi być $d=0$.
Pokrywamy $S_{i,k,m}$ przeliczalnie wieloma
fragmentami $A$ otrzymanymi przez przecięcie
z kostkami siatki o średnicy mniejszej niż
$1/(2m)$. Dla każdego niepustego fragmentu
odwrotność rzutu na $\pi_i(A)$ ma postać
$z\mapsto z+h(z)e_i$, gdzie $h(\pi_i(x))=x\cdot e_i$.
Dla $x,y\in A$ otrzymujemy
\[
 |h(\pi_i(y))-h(\pi_i(x))|
 \le\frac{\alpha}{\sqrt{1-\alpha^2}}|\pi_i(y)-\pi_i(x)|.
\]
Niech $L=\alpha/\sqrt{1-\alpha^2}$.
Przedłużamy $h$ na całe $e_i^\perp$ wzorem
\[
 \widetilde h(z)=\inf_{w\in\pi_i(A)}(h(w)+L|z-w|).
\]
Infimum jest skończone: dla ustalonego
$w_0\in\pi_i(A)$ jest ograniczone od góry
przez $h(w_0)+L|z-w_0|$, a od dołu przez
$h(w_0)-L|z-w_0|$ na mocy nierówności
Lipschitza dla $h$ i nierówności trójkąta.
Na $\pi_i(A)$ wzór daje dokładnie $h$,
gdyż każdy składnik infimum jest nie mniejszy
niż $h(z)$ i można wziąć $w=z$.
Nierówność trójkąta dla $|z-w|$ pokazuje
$\widetilde h(z)\le\widetilde h(z')+L|z-z'|$,
a zamiana $z,z'$ daje własność Lipschitza.
Każdy fragment zawiera się więc w pełnym wykresie
$\widetilde h$. Po identyfikacji $e_i^\perp$
z $\mathbb R^{n-1}$ wykres jest obrazem
Lipschitzowskim całej tej przestrzeni.
Przeliczalność indeksów i fragmentów dowodzi tezy.
\end{proof}

\begin{lemat}[Gęstość miary na wykresach prostowalnych]
\label{lem:dg-density-identification}
Niech $n\ge2$, $S\subseteq\mathbb R^n$ będzie
borelowski i zawarty w przeliczalnej sumie pełnych
wykresów funkcji Lipschitzowskich nad hiperpłaszczyznami
$(n-1)$-wymiarowymi, a $\mu$ niech będzie dodatnią
miarą Radona z $\mu(\mathbb R^n\setminus S)=0$.
Załóżmy, że dla każdego $x\in S$ zachodzi
$\lim_{r\downarrow0}\mu(B(x,r))/
(\omega_{n-1}r^{n-1})=1$.
Wtedy $\mu=\mathcal H^{n-1}\llcorner S$.
\end{lemat}
\begin{proof}
Niech $(\Gamma_j)$ będzie pokryciem $S$ wykresami
z założenia. Każdy pełny wykres jest domknięty,
ponieważ funkcja Lipschitzowska jest ciągła.
Kładziemy $S_j=(S\cap\Gamma_j)\setminus
\bigcup_{i<j}\Gamma_i$. Są to parami rozłączne
zbiory borelowskie, a ich suma jest dokładnie $S$.
Borelowskość $S$ jest założeniem tego lematu,
a nie wnioskiem z samej prostowalności.

Dla $M\in\mathbb N$ niech $S_M$ oznacza
zbiór punktów $x\in S$, dla których
\[
 \frac12\omega_{n-1}r^{n-1}
 \le\mu(B(x,r))\le2\omega_{n-1}r^{n-1}
 \qquad(0<r<1/M).
\]
Z założenia $S=\bigcup_MS_M$.
Zbiory $S_M$ są borelowskie: funkcje mas kul
dla ustalonych promieni są borelowskie,
a nierówności wystarczy sprawdzać dla promieni
wymiernych, gdyż przy $q\uparrow r$ zachodzi
$\mu(B(x,q))\uparrow\mu(B(x,r))$.
Oznaczmy $d=n-1>0$. Dla $x\in S_M$ mamy
$\mu(\{x\})=0$ z górnego oszacowania przy $r\downarrow0$.
Jeżeli $A\subseteq S_M$ jest borelowski, pokrywamy
go zbiorami $U_i$ o średnicach $d_i<1/(2M)$.
Pomijamy zbiory nieprzecinające $A$ i wybieramy
$x_i\in A\cap U_i$. Gdy $d_i>0$, zachodzi
$A\cap U_i\subset B(x_i,2d_i)$ i
$\mu(B(x_i,2d_i))\le2\omega_d(2d_i)^d$.
Zbiory o średnicy zero wnoszą masę zero.
Podaddytywność, infimum po pokryciach i przejście
ze średnicą maksymalną do zera dają
$\mu(A)\le C_n\mathcal H^d(A)$, przy czym stała
uwzględnia kalibrację miary Hausdorffa.

Pokażemy również odwrotną bezwzględną ciągłość
na każdym $S_j$. Niech
$K\subseteq S_M\cap S_j$ będzie zwarty.
Dla $0<\delta<1/M$ bierzemy rodzinę wszystkich
kul o środkach w $K$ i promieniach mniejszych
od $\delta$. Rodzina ma ograniczone promienie,
a jej suma zawiera $K$. Lemat selekcji kul z
Rozdziału~\ref{ch:covering-24} pozwala wybrać
co najwyżej przeliczalną rodzinę parami rozłącznych
kul $B(x_i,r_i)$, których pięciokrotne powiększenia
pokrywają $K$. Każda wybrana kula leży w $K_\delta$
i spełnia $r_i<1/M$ oraz $x_i\in S_M$.
Definicja miary Hausdorffa i dolne oszacowanie
masy dają
\[
 \mathcal H^{n-1}_{10\delta}(K)
 \le C_n\sum_i r_i^{n-1}
 \le C'_n\sum_i\mu(B(x_i,r_i))
 \le C'_n\mu(K_\delta),
\]
gdzie $K_\delta=\{z:\operatorname{dist}(z,K)<\delta\}$.
Przez $\delta\downarrow0$ i skończoność
$\mu$ w pewnym otoczeniu zwartego $K$
otrzymujemy $\mathcal H^{n-1}(K)\le C'_n\mu(K)$.
Miara $\mathcal H^d\llcorner\Gamma_j$ jest
borelowska i lokalnie skończona: parametryzacja
wykresu przez $G(z)=(z,h(z))$ jest Lipschitzowska,
a przeciwobraz części wykresu zawartej w zwartej
kuli jest ograniczony dzięki rzutowi na przestrzeń
parametrów. Nierówność Lipschitzowska dla miary
Hausdorffa ogranicza więc jej masę na każdej
zwartej kuli przez skończoną stałą razy objętość
ograniczonego zbioru parametrów. Z regularności
lokalnie skończonych miar borelowskich w przestrzeni
euklidesowej wynika, że jest to miara Radona.
Dla borelowskiego $A\subseteq S_M\cap S_j$
bierzemy supremum oszacowania po zwartych
$K\subseteq A$, uzyskując
$\mathcal H^d(A)\le C'_n\mu(A)$.
Zbiory $S_M$ są rosnące względem $M$, a ich suma
jest równa $S$. Dla dowolnego borelowskiego
$A\subseteq S_j$ stosujemy obie nierówności
do $A\cap S_M$ i przechodzimy $M\to\infty$
na mocy ciągłości miar od dołu. W szczególności
$\mu\llcorner S_j\ll\mathcal H^d\llcorner S_j$
oraz $\mathcal H^d\llcorner S_j\ll\mu\llcorner S_j$.

Uzasadnimy dokładnie gęstość miary Hausdorffa
na $S_j$. Po izometrii możemy zapisać wykres
jako obraz $G:\mathbb R^d\to\mathbb R^n$,
$G(z)=(z,h(z))$, gdzie $h$ jest Lipschitzowska.
Niech $A=G^{-1}(S_j)$; jest to zbiór borelowski.
Odwzorowanie $G$ jest różnowartościowe,
a jego odwrotność na wykresie jest rzutem.
Z twierdzenia Rademachera
(Twierdzenie~\ref{thm:rademacher-28})
$G$ jest różniczkowalne prawie wszędzie.
Niech $J(z)=\sqrt{\det(DG(z)^{\mathsf T}DG(z))}$
w punktach różniczkowalności, z wartością $0$
na pozostałym zbiorze zerowym.
Funkcja $J$ jest mierzalna i ograniczona, ponieważ
różniczki mają normy ograniczone przez stałą
Lipschitza $G$. Wzór na pole dla odwzorowań
Lipschitzowskich
(Twierdzenie~\ref{thm:lipschitz-area-coarea-28})
daje dla każdego $r>0$ i $z_0\in A$
\[
 \mathcal H^d(S_j\cap B(G(z_0),r))
 =\int_{\{z:|G(z)-G(z_0)|<r\}}\chi_A(z)J(z)
 \,\mathrm d\lambda_d(z).
\]
Rozważamy te $z_0\in A$, w których istnieje
$L=DG(z_0)$ i które są punktami Lebesgue'a
funkcji $\chi_A J$ o wartości $J(z_0)$.
Oba warunki zachodzą dla prawie wszystkich
$z_0\in A$ względem $\lambda_d$.
Macierz $L$ spełnia $|Lv|\ge|v|$ dla wszystkich
$v\in\mathbb R^d$, ponieważ pierwsze $d$
składowych $Lv$ to $v$.
Dla każdego $0<\varepsilon<1$ i dostatecznie
małego $|v|$ różniczkowalność daje
\[
 (1-\varepsilon)|Lv|
 \le |G(z_0+v)-G(z_0)|
 \le (1+\varepsilon)|Lv|.
\]
Ponadto $|G(z)-G(z_0)|\ge|z-z_0|$ globalnie.
Dlatego przy dostatecznie małym $r$
zbiór całkowania, po przesunięciu o $-z_0$,
zawiera elipsoidę $\{|Lv|<r/(1+\varepsilon)\}$
i jest zawarty w elipsoidzie
$\{|Lv|<r/(1-\varepsilon)\}$, a także w $B(0,r)$.
Ponieważ $L^{\mathsf T}L$ jest dodatnio określona,
liniowa zmiana zmiennych przez
$(L^{\mathsf T}L)^{1/2}$ daje
\[
 \lambda_d(\{v:|Lv|<t\})
 =\frac{\omega_d t^d}{J(z_0)}.
\]
Błąd zastąpienia $\chi_A(z)J(z)$ przez $J(z_0)$
w powyższej całce jest co do modułu nie większy
niż
$\int_{B(z_0,r)}|\chi_A(z)J(z)-J(z_0)|
\,\mathrm d\lambda_d(z)=o(r^d)$
z własności punktu Lebesgue'a.
Po podzieleniu przez $\omega_d r^d$ otrzymujemy
dolną granicę nie mniejszą niż $(1+\varepsilon)^{-d}$
i górną nie większą niż $(1-\varepsilon)^{-d}$.
Przez $\varepsilon\downarrow0$ dostajemy
\[
 \lim_{r\downarrow0}
 \frac{\mathcal H^d(S_j\cap B(G(z_0),r))}
 {\omega_d r^d}=1.
\]
Zbiór wyjątkowych parametrów ma zerową
$\lambda_d$-miarę; jego obraz ma zerową
$\mathcal H^d$-miarę na mocy oszacowania
Lipschitzowskiego dla miary Hausdorffa.
Zatem wzór zachodzi dla $\mathcal H^d$-prawie
każdego punktu $S_j$.
Stosujemy Twierdzenie~\ref{thm:relative-lebesgue-25}
do miary $\mu$ i funkcji $\chi_{S_j}$, która jest
lokalnie całkowalna względem $\mu$.
Dla $\mu$-prawie każdego $x\in S_j$ otrzymujemy
$\mu(B(x,r)\cap S_j)/\mu(B(x,r))\to1$.
Z wzajemnej bezwzględnej ciągłości ograniczeń
miar do $S_j$ wynika, że ta granica zachodzi
również dla $\mathcal H^d$-prawie każdego $x\in S_j$.
Łącząc te dwa rezultaty z założeniem o gęstości
$\mu$, uzyskujemy
\[
 \frac{\mu(B(x,r)\cap S_j)}
 {\mathcal H^{n-1}(S_j\cap B(x,r))}
 \longrightarrow1
 \quad\text{dla }\mathcal H^{n-1}\text{-prawie wszystkich }x\in S_j.
\]
Niech $\tau_j=\mathcal H^d\llcorner S_j$
oraz $\eta_j=\mu\llcorner S_j$. Obie miary są
lokalnie skończonymi miarami borelowskimi
na $\mathbb R^n$, a więc miarami Radona;
nie potrzebujemy domkniętości samego $S_j$.
Są $\sigma$-skończone i $\eta_j\ll\tau_j$,
stąd istnieje nieujemna lokalnie całkowalna
gęstość $f_j$ z $\eta_j=f_j\tau_j$.
Twierdzenie~\ref{thm:relative-radon-25} dla pary
$(\tau_j,\eta_j)$ identyfikuje powyższy iloraz
z $f_j(x)$ dla $\tau_j$-prawie każdego $x$.
Zatem $f_j=1$ $\tau_j$-prawie wszędzie i
$\eta_j=\tau_j$, czyli
$\mu\llcorner S_j=\mathcal H^d\llcorner S_j$.
Sumując po rozłącznych zbiorach $S_j$,
których suma jest $S$, oraz korzystając
ze skupienia $\mu$ na $S$, dostajemy tezę.
\end{proof}

\begin{twierdzenie}[Twierdzenie strukturalne De Giorgiego]
\label{thm:de-giorgi-structure-29}
Niech $n\ge1$, a $E\subseteq\mathbb R^n$ będzie
mierzalnym w sensie Lebesgue'a zbiorem lokalnie
skończonego obwodu. Wtedy
$\partial^*E$ jest przeliczalnie $(n-1)$-prostowalny,
$|D\chi_E|=\mathcal H^{n-1}\llcorner\partial^*E$
oraz dla każdego $x\in\partial^*E$ i każdego $R>0$
zachodzi
\[
 \lim_{r\downarrow0}\int_{B(0,R)}
 \bigl|\chi_E(x+rz)-\chi_{H_{\nu_E^*(x)}}(z)\bigr|
 \,\mathrm d\lambda_n(z)=0\qquad(R>0),
\]
gdzie $H_\nu=\{z:z\cdot\nu<0\}$.
W szczególności, dla każdego pola
$\varphi\in C_c^1(\mathbb R^n;\mathbb R^n)$
\[
 \int_E\operatorname{div}\varphi(x)\,\mathrm d\lambda_n(x)
 =\int_{\partial^*E}\varphi(x)\cdot\nu_E^*(x)
 \,\mathrm d\mathcal H^{n-1}(x).
\]
\end{twierdzenie}
\begin{proof}
Dla $n\ge2$ przyjmijmy $\mu=|D\chi_E|$.
Z wcześniejszego lematu o mierzalności
brzegu zredukowanego i skupieniu wariacji wynika,
że $S=\partial^*E$ jest borelowski,
$\mu(\mathbb R^n\setminus S)=0$ oraz
$\nu_E^*=\nu_E$ $\mu$-prawie wszędzie.
Lemat~\ref{lem:dg-blowup} dowodzi dla każdego
$x\in S$ granicy po przeskalowaniu i dokładnej
gęstości $(n-1)$-wymiarowej miary $\mu$ równej $1$.
Lemat~\ref{lem:dg-rectifiability} pokrywa cały $S$
przeliczalnie wieloma pełnymi wykresami
Lipschitzowskimi, a więc daje prostowalność
oraz brak zbioru wyjątkowego w tym pokryciu.
Spełnione są zatem wszystkie założenia
Lematu~\ref{lem:dg-density-identification},
który daje $\mu=\mathcal H^{n-1}\llcorner S$.
W Twierdzeniu~\ref{thm:gauss-green-BV-29}
zastępujemy miarę $\mu$ tym ograniczeniem
oraz pole $\nu_E$ równym mu prawie wszędzie
polem $\nu_E^*$. Otrzymujemy wzór całkowy.
Ponieważ miara wariacji jest lokalnie skończona,
prawa strona jest dobrze określona i skończona
dla każdego pola testowego.

\emph{Przypadek jednowymiarowy.}
Niech teraz $n=1$ i niech $\eta=D\chi_E$
będzie lokalnie skończoną miarą ze znakiem.
Na dowolnym ograniczonym otwartym przedziale
$I=(a,b)$ jej wariacja jest skończona.
Definiujemy funkcję $F:I\to\mathbb R$ wzorem
$F(t)=\eta((a,t])$. Jest ona prawostronnie
ciągła (z ciągłości skończonych miar od góry),
ma granicę lewostronną
$F(t-)=\eta((a,t))$ i dla dowolnego podziału
$a<t_0<\cdots<t_m<b$ spełnia
\[
 \sum_{i=1}^m|F(t_i)-F(t_{i-1})|
 =\sum_{i=1}^m|\eta((t_{i-1},t_i])|
 \le|\eta|(I).
\]
Dla testu $\psi\in C_c^1(I)$ Fubini daje
\begin{align*}
 \int_I F(t)\psi'(t)\,\mathrm d\lambda_1(t)
 &=\int_I\int_{[s,b)}\psi'(t)
 \,\mathrm d\lambda_1(t)\,\mathrm d\eta(s)\\
 &=-\int_I\psi(s)\,\mathrm d\eta(s).
\end{align*}
Zastosowanie Fubiniego do miary ze znakiem
jest uzasadnione, gdyż całka z modułu
podcałkowego względem $|\eta|$ i $\lambda_1$
jest ograniczona przez
$(b-a)\|\psi'\|_\infty|\eta|(I)$.
Zatem $DF=\eta$ na $I$ i
$D(\chi_E-F)=0$ w sensie dystrybucji.

Funkcja lokalnie całkowalna na przedziale
o zerowej pochodnej dystrybucyjnej jest stała
prawie wszędzie. Istotnie, na każdym przedziale
o zwartym domknięciu w $I$ jej wygładzenia
jądrem o dostatecznie małym nośniku mają
pochodną zero, więc są stałe. Zbieżność
w $L^1$ na tym przedziale wymusza zbieżność
tych stałych do jednej liczby, a więc granica
jest stała prawie wszędzie. Przez rosnące
wyczerpanie $I$ przedziałami wnioskujemy,
że stała jest wspólna. Stąd istnieje $c\in\mathbb R$
takie, że $g=c+F$ jest reprezentantem $\chi_E$
na $I$.

Z prawostronnej ciągłości $g$ i faktu, że
$g\in\{0,1\}$ prawie wszędzie, wynika
$g(t)\in\{0,1\}$ dla każdego $t\in I$:
wybieramy ciąg punktów pełnej miary z prawej
strony $t$ zbieżny do $t$, a zbiór $\{0,1\}$
jest domknięty. Tak samo granica lewostronna
$g(t-)$ należy do $\{0,1\}$.
Punkt $t$ jest punktem nieciągłości wtedy
i tylko wtedy, gdy $g(t)\ne g(t-)$.
Każdy taki skok ma moduł $1$, a
$g(t)-g(t-)=\eta(\{t\})$.
Liczba punktów skoku w $I$ jest skończona:
dla każdego skończonego zbioru $J$ tych punktów
addytywność wariacji daje
$\#J=\sum_{t\in J}|\eta(\{t\})|\le|\eta|(I)$.
Gdyby punktów skoku było nieskończenie wiele,
moglibyśmy wybrać dowolnie duże skończone $J$,
co przeczyłoby temu oszacowaniu.

Pomiędzy kolejnymi punktami skoku $g$ jest
ciągła i przyjmuje wartości w $\{0,1\}$, więc
jest stała na każdym z tych przedziałów.
Całkowanie przez części dla funkcji schodkowej
daje na $I$
\[
 \eta=\sum_{t\in J_I}(g(t)-g(t-))\,\delta_t,
 \qquad |\eta|=\sum_{t\in J_I}\delta_t,
\]
gdzie $J_I$ jest skończonym zbiorem skoków.
Dokładniej, w całce z $g\psi'$ sumujemy
różnice wartości testu na końcach kolejnych
przedziałów stałości; współczynnik przy
$\psi(t)$ wynosi $g(t-)-g(t)$, a przy końcach
$I$ test znika. Jednoznaczność reprezentacji
miarowej daje powyższą tożsamość.
Przez wyczerpanie $\mathbb R$ przedziałami
$(-m,m)$ otrzymujemy lokalnie skończony,
co najwyżej przeliczalny zbiór $J$ skoków
i lokalne tożsamości tych miar.
Wybór przedziału nie ma znaczenia: miary
na przecięciach zgadzają się, a punkty skoku
są dokładnie atomami wariacji.

Jeśli $x\notin J$, to istnieje otoczenie
$x$ niezawierające żadnego punktu $J$,
gdyż $J$ jest lokalnie skończony.
W tym otoczeniu $|\eta|$ jest zerowa,
więc $x\notin\partial^*E$.
Jeśli $x\in J$, to każda kula o środku $x$
ma dodatnią masę wariacji, a dla wszystkich
dostatecznie małych promieni zawiera tylko
atom w $x$. Wówczas
$\nu_E^*(x)=-\eta(\{x\})\in\{-1,1\}$.
Zatem $\partial^*E=J$ i
$|\eta|=\mathcal H^0\llcorner\partial^*E$.
Dla skoku z $1$ do $0$ normalna wynosi $1$
i lokalnie zbiór $E$ zgadza się prawie wszędzie
z lewą półprostą o końcu $x$.
Dla skoku z $0$ do $1$ normalna wynosi $-1$
i odpowiada prawym półprostom.
Dla każdego ustalonego $R>0$ i dostatecznie
małego $r>0$ funkcje
$\chi_E(x+rz)$ i $\chi_{H_{\nu_E^*(x)}}(z)$
są zatem równe prawie wszędzie na $(-R,R)$.
Daje to granicę po przeskalowaniu, a
przeliczalność $J$ daje $0$-prostowalność.
Wzór Gaussa--Greena wynika z reprezentacji
miary atomowej lub z
Twierdzenia~\ref{thm:gauss-green-BV-29}
i powyższej identyfikacji miary.
\end{proof}

\begin{center}
\begin{tikzpicture}[>=stealth,scale=0.84]
  \begin{scope}[shift={(0,0)}]
    \clip (0,0) circle (1.35);
    \fill[blue!13] (-1.5,-1.5)--(-0.10,-1.5)--(-0.12,-1.3)--
      (0.08,-0.9)--(-0.10,-0.5)--(0.13,0)--(-0.08,0.45)--
      (0.09,0.9)--(-0.02,1.5)--(-1.5,1.5)--cycle;
    \draw[blue!75!black,very thick] plot[smooth] coordinates
      {(-0.12,-1.35) (0.08,-0.9) (-0.10,-0.5) (0.13,0)
       (-0.08,0.45) (0.09,0.9) (-0.02,1.35)};
  \end{scope}
  \draw (0,0) circle (1.35);
  \node at (0,-1.68) {$E_{x,r}\cap B(0,1)$};
  \draw[->,thick] (1.75,0)--(3.3,0)
    node[pos=.5,above] {$r\downarrow0$};
  \begin{scope}[shift={(5.15,0)}]
    \clip (0,0) circle (1.35);
    \fill[blue!13] (-1.5,-1.5) rectangle (0,1.5);
    \draw[blue!75!black,very thick] (0,-1.4)--(0,1.4);
    \draw[->,red!70!black,thick] (0,0)--(0.9,0)
      node[above,font=\scriptsize] {$\nu_E^*(x)$};
  \end{scope}
  \draw (5.15,0) circle (1.35);
  \node at (5.15,-1.68) {$H_{\nu_E^*(x)}\cap B(0,1)$};
\end{tikzpicture}

\small\emph{Schemat granicy De Giorgiego:}
po przeskalowaniu przy punkcie brzegu zredukowanego
zbiór staje się półprzestrzenią w sensie zbieżności
$L^1_{\mathrm{loc}}$; rysunek nie oznacza zbieżności
brzegów w odległości Hausdorffa.
\end{center}

\begin{uwaga}[Brzeg topologiczny a zredukowany]
Twierdzenie nie orzeka, że $\partial E=\partial^*E$.
Zmiana $E$ na zbiorze $\lambda_n$-zerowym może
zupełnie zmienić jego brzeg topologiczny,
lecz nie zmienia $D\chi_E$, obwodu ani wzoru
Gaussa--Greena. W szczególności właściwą
miarą powierzchniową dla ogólnych zbiorów
skończonego obwodu jest
$\mathcal H^{n-1}\llcorner\partial^*E$.
\end{uwaga}

\begin{lemat}[Mierzalność obwodów zbiorów nad- i podpoziomicowych]
\label{lem:BV-level-meas-auditVI}
Niech $\Omega\subseteq\mathbb R^n$ będzie zbiorem otwartym,
a $u:\Omega\to\mathbb R$ będzie mierzalna
w sensie Lebesgue'a. Dla $t\in\mathbb R$ oznaczmy
$E_t=\{x\in\Omega:u(x)>t\}$ oraz
$F_t=\{x\in\Omega:u(x)<-t\}$.
Obie funkcje $t\mapsto P(E_t;\Omega)$ i
$t\mapsto P(F_t;\Omega)$ o wartościach
w $[0,+\infty]$ są borelowskie.
Nie zakładamy całkowalności $u$ ani skończoności
obwodów. Wartości tych funkcji nie zmieniają się
po zmianie $u$ na zbiorze miary Lebesgue'a zero.
\end{lemat}
\begin{proof}
Jeśli $\Omega=\varnothing$, oba obwody są zerowe.
W przeciwnym razie wybieramy wyczerpanie
$\Omega$ zbiorami zwartymi $(K_j)_{j\ge1}$
takie, że $K_j\subseteq\operatorname{int}K_{j+1}$.
Niech $\mathcal A_j$ oznacza zbiór pól
$\varphi\in C_c^1(\Omega;\mathbb R^n)$
o nośniku zawartym w $K_j$ i
$|\varphi(x)|\le1$ dla każdego $x$.
Wyposażamy go w metrykę normy $C^1$:
sumy norm supremum składowych i ich pierwszych
pochodnych. Odwzorowanie przypisujące polu
ograniczenia tych funkcji do $K_j$ zanurza
$\mathcal A_j$ izometrycznie w skończonym
produkcie przestrzeni $C(K_j;\mathbb R)$.
Istotnie, składowe i ich pochodne znikają
poza $K_j$, więc wystarczy mierzyć ich normy
na $K_j$. Przestrzeń $C(K_j;\mathbb R)$
jest ośrodkowa: wielomiany o współczynnikach
wymiernych tworzą przeliczalną rodzinę gęstą
na mocy twierdzenia Stone'a--Weierstrassa.
Skończony produkt przestrzeni metrycznych
ośrodkowych jest ośrodkowy, a każda jego
podprzestrzeń również jest ośrodkowa
(dziedziczy drugi aksjomat przeliczalności,
a wybór punktu z każdego niepustego elementu
przeliczalnej bazy daje zbiór gęsty).
Wybieramy więc przeliczalny gęsty podzbiór
$\mathcal D_j$ samego $\mathcal A_j$.
W szczególności każdy wybrany test już spełnia
ograniczenie punktowe normy; nie trzeba go
ponownie skalować.

Niech $\mathcal D=\{0\}\cup\bigcup_j
(\mathcal D_j\cup(-\mathcal D_j))$.
Jest to przeliczalna symetryczna rodzina
dopuszczalnych pól, którą wyliczamy jako
$(\varphi_l)_{l\ge1}$, z powtórzeniami w razie
potrzeby. Dla każdego mierzalnego $E\subseteq\Omega$
zachodzi
\[
 P(E;\Omega)=\sup_{l\ge1}
 \int_E\operatorname{div}\varphi_l(x)
 \,\mathrm d\lambda_n(x).
\]
Nierówność z prawą stroną nie większą od lewej
wynika z dopuszczalności testów.
Dla nierówności odwrotnej ustalamy dowolny test
$\varphi$ z definicji obwodu. Jego zwarty nośnik
jest zawarty w pewnym $K_j$, więc istnieją
$\psi_m\in\mathcal D_j$ z $\psi_m\to\varphi$
w normie $C^1$. Wszystkie dywergencje znikają
poza $K_j$, a zatem
\[
 \left|\int_E\operatorname{div}(\psi_m-\varphi)
 \,\mathrm d\lambda_n\right|
 \le\lambda_n(K_j)
 \|\operatorname{div}(\psi_m-\varphi)\|_\infty
 \longrightarrow0.
\]
Całka dla $\varphi$ nie przekracza zatem
supremum po $\mathcal D$. Biorąc supremum
po wszystkich $\varphi$, otrzymujemy równość,
także gdy obwód jest nieskończony.
Rodzina nie zależy od $E$ ani od $t$.

Dla ustalonego $l$ rozkładamy
$\operatorname{div}\varphi_l=g_l^+-g_l^-$,
gdzie $g_l^\pm$ są nieujemnymi funkcjami
całkowalnymi. Dla każdego znaku funkcja
\[
 t\longmapsto\int_{E_t}g_l^\pm(x)\,
 \mathrm d\lambda_n(x)
\]
jest skończona i nierosnąca, ponieważ
$E_{t_2}\subseteq E_{t_1}$ przy $t_1<t_2$.
Jest zatem borelowska. Jej różnica z analogiczną
funkcją dla drugiego znaku jest borelowska,
a supremum przeliczalnie wielu takich różnic
daje borelowskość $t\mapsto P(E_t;\Omega)$.
Rodzina $F_t$ również maleje przy wzroście $t$,
więc ten sam argument stosujemy do $F_t$.
Nie korzystamy tutaj z borelowskości $u$
ani z twierdzenia o mierzalności całek
parametrycznych: wystarcza monotoniczność
całek z nieujemnych funkcji.

Jeśli $u=\widetilde u$ prawie wszędzie, to dla
każdego $t$ odpowiednie zbiory nad- i podpoziomicowe różnią
się co najwyżej na jednym ustalonym zbiorze
miary Lebesgue'a zero. Z Lematu~\ref{lem:perimeter-complement-29}
ich obwody są równe. Dowodzi to ostatniej tezy.
\end{proof}

\section{Koobszar, funkcje BV i warstwowanie}
\begin{twierdzenie}[Jedna nierówność formuły koobszaru dla $BV$]
\label{thm:BV-layercake-29}
Niech $\Omega\subseteq\mathbb R^n$ będzie zbiorem otwartym,
$u\in L^1(\Omega;\mathbb R)$ oraz dla $t>0$
niech $E_t^+=\{x\in\Omega:u(x)>t\}$ i
$E_t^-=\{x\in\Omega:u(x)<-t\}$.
Jeśli
\[
 \int_0^\infty\!P(E_t^+;\Omega)\,\mathrm d\lambda_1(t)
 +\int_0^\infty\!P(E_t^-;\Omega)\,\mathrm d\lambda_1(t)<\infty,
\]
to $u\in BV(\Omega)$ i
\[
 |Du|(\Omega)\le\int_0^\infty
 [P(E_t^+;\Omega)+P(E_t^-;\Omega)]\,
 \mathrm d\lambda_1(t)
 =\int_{\mathbb R}P(\{u>s\};\Omega)
 \,\mathrm d\lambda_1(s).
\]
W szczególności założenie całkowe samo pociąga
skończoność obu obwodów dla prawie wszystkich $t>0$.
\end{twierdzenie}
\begin{proof}
Wybieramy wszędzie skończony borelowski
reprezentant $u$; poprzedni lemat zapewnia,
że nie zmienia to żadnego z obwodów.
Dla każdej liczby $a\in\mathbb R$ zachodzi
$\int_0^\infty\chi_{\{a>t\}}\,\mathrm d\lambda_1(t)
=\max\{a,0\}$ oraz
$\int_0^\infty\chi_{\{a<-t\}}\,\mathrm d\lambda_1(t)
=\max\{-a,0\}$. Odejmując te równości dla
$a=u(x)$, otrzymujemy punktowo wzór warstwowania
\[
 u(x)=\int_0^\infty
 \bigl(\chi_{E_t^+}(x)-\chi_{E_t^-}(x)\bigr)
 \,\mathrm d\lambda_1(t).
\]
Funkcje charakterystyczne na przestrzeni
$\Omega\times(0,\infty)$ są borelowskie.
Ponieważ dla $t>0$ zbiory $E_t^+,E_t^-$
są rozłączne, z Tonellego otrzymujemy
\begin{align*}
 \int_0^\infty\int_\Omega
 (\chi_{E_t^+}(x)+\chi_{E_t^-}(x))
 \,\mathrm d\lambda_n(x)\,\mathrm d\lambda_1(t)
 &=\int_\Omega|u(x)|\,\mathrm d\lambda_n(x)<\infty.
\end{align*}
Dla dowolnego $\varphi\in C_c^1(\Omega;\mathbb R^n)$
z $|\varphi|\le1$ całka z modułu funkcji
$(\chi_{E_t^+}-\chi_{E_t^-})\operatorname{div}\varphi$
na przestrzeni produktowej nie przekracza
$\|\operatorname{div}\varphi\|_\infty\|u\|_{L^1}$.
Spełnione są więc hipotezy Fubiniego.
Otrzymujemy
\[
 \int_\Omega u(x)\operatorname{div}\varphi(x)
 \,\mathrm d\lambda_n(x)
 =\int_0^\infty\!\left(
 \int_{E_t^+}\operatorname{div}\varphi(x)
 \,\mathrm d\lambda_n(x)
 -\int_{E_t^-}\operatorname{div}\varphi(x)
 \,\mathrm d\lambda_n(x)\right)\mathrm d\lambda_1(t).
\]
Dla każdego $t>0$ moduł pierwszej całki
w nawiasie nie przekracza $P(E_t^+;\Omega)$,
a moduł drugiej nie przekracza $P(E_t^-;\Omega)$:
w definicji obwodu dopuszczalne są oba pola
$\varphi$ i $-\varphi$.
Całka po $t$ z sumy obwodów jest dobrze
określona dzięki poprzedniemu lematowi.
Otrzymujemy oszacowanie niezależne od pola
testowego. Biorąc supremum po wszystkich
dopuszczalnych polach, dostajemy pierwszą
nierówność. Jej prawa strona jest skończona
z założenia, więc $u\in BV(\Omega)$.

Pozostaje uzasadnić przejście od dodatnich
i ujemnych warstw do całki po całej prostej.
Zbiory $\{u=-t\}$ są parami rozłączne dla
różnych $t>0$. Ponieważ $\lambda_n$ na $\Omega$
jest $\sigma$-skończona, co najwyżej przeliczalnie
wiele z tych zbiorów ma dodatnią miarę:
na każdym zbiorze wyczerpania o skończonej mierze
tylko skończenie wiele rozłącznych poziomic
może mieć miarę co najmniej $1/k$, dla
ustalonego $k\ge1$. Bierzemy sumę po wyczerpaniu
i po $k$. Dla pozostałych $t>0$ zbiory
$\{u<-t\}$ i $\{u\le-t\}$ są równe prawie wszędzie.
Drugi z nich jest dopełnieniem $\{u>-t\}$
w $\Omega$. Lemat~\ref{lem:perimeter-complement-29}
(dla dopełnień w $\mathbb R^n$ ograniczonych
do $\Omega$) daje
$P(E_t^-;\Omega)=P(\{u>-t\};\Omega)$
dla prawie wszystkich $t>0$.
Zamiana zmiennej $s=-t$ w całce z obwodów
ujemnych warstw daje równość obu wyrażeń
w tezie. Nie jest przy tym wymagana skończona
miara samego $\Omega$.
\end{proof}

\section{Pełna formuła koobszaru dla funkcji \texorpdfstring{$BV$}{BV}}
Związek między obwodem zbiorów nadpoziomicowych a wariacją
funkcji nie wymaga twierdzenia strukturalnego De Giorgiego.
Wystarczą ścisła aproksymacja w $BV$, gładki wzór
koobszaru z Rozdziału~\ref{ch:area-coarea-28} i dolna
półciągłość wariacji.

\begin{lemat}[Ścisła aproksymacja gładka w $BV$]
\label{lem:BV-strict-approx}
Niech $n\ge1$, $\Omega\subseteq\mathbb R^n$ będzie
zbiorem otwartym, a $u\in BV(\Omega)$.
Istnieje ciąg funkcji o wartościach rzeczywistych
$u_j\in C^\infty(\Omega)\cap W^{1,1}(\Omega)$ taki, że
\[
 \|u_j-u\|_{L^1(\Omega)}\longrightarrow0,
 \qquad
 \int_\Omega|\nabla u_j(x)|\,\mathrm d\lambda_n(x)
 \longrightarrow |Du|(\Omega).
\]
Nie wymaga się regularności brzegu $\Omega$.
\end{lemat}
\begin{proof}
Jeśli $\Omega=\varnothing$, teza jest oczywista.
Zakładamy dalej, że $\Omega\ne\varnothing$.

\emph{Krok 1: lokalnie skończone otoczenia nośników.}
Wybieramy rosnące wyczerpanie $(U_i)_{i\ge1}$
zbioru $\Omega$ zbiorami otwartymi takimi, że
$\overline U_i$ jest zwarty i
$\overline U_i\subseteq U_{i+1}\subseteq\Omega$.
Możemy na przykład przyjąć
$U_i=\{x\in\Omega:|x|<i,\,
\operatorname{dist}(x,\mathbb R^n\setminus\Omega)>1/i\}$,
z konwencją odległości $+\infty$ od zbioru pustego.
Kładziemy $U_i=\varnothing$ dla $i\le0$.
Rodzina otwartych zbiorów
$V_i=U_{i+1}\setminus\overline U_{i-1}$,
$i\ge1$, pokrywa $\Omega$: dla punktu $x$
bierzemy najmniejszy indeks $m$ taki, że
$x\in U_m$. Jeżeli $x\in\overline U_{m-1}$, to $m\ge2$,
$x\notin U_{m-1}$ i
$\overline U_{m-2}\subset U_{m-1}$,
więc $x\in V_{m-1}$.
W przeciwnym przypadku $x\in V_m$.
Ponadto dla każdego zwartego $K\subseteq\Omega$
istnieje $M$ z $K\subset U_M$, a wtedy
$K\cap V_i=\varnothing$ dla $i\ge M+1$.
Rodzina $(V_i)$ jest zatem lokalnie skończona.
Każdy $V_i$ ma zwarte domknięcie w $\Omega$.

Wybieramy gładki podział jedności $(\psi_i)$
podporządkowany temu pokryciu, czyli
$\psi_i\ge0$, $\operatorname{supp}\psi_i\subset V_i$
oraz $\sum_i\psi_i=1$ na $\Omega$.
Nośniki są zwarte, ponieważ są domknięte
w $\Omega$ i zawierają się w $\overline U_{i+1}$.
Do skonstruowania takiego podziału można użyć
gładkiego podziału jedności podporządkowanego
lokalnie skończonemu pokryciu otwartemu.
Istotna jest tutaj nie tylko lokalna skończoność
samych nośników, ale również otwartych zbiorów
$V_i$, w których będziemy je powiększać.
Indeksy z $\psi_i=0$ nie wnoszą wkładu do sum.
Dla każdego pozostałego indeksu istnieje $d_i>0$
takie, że zamknięte $d_i$-otoczenie
$\operatorname{supp}\psi_i$ zawiera się w $V_i$.
Wynika to ze zwartości nośnika i otwartości $V_i$.

\emph{Krok 2: dobór lokalnych wygładzeń.}
Niech $\rho\in C_c^\infty(B(0,1))$ będzie
nieujemna i niech $\int\rho\,\mathrm d\lambda_n=1$.
Dla $\varepsilon>0$ kładziemy
$\rho_\varepsilon(x)=\varepsilon^{-n}\rho(x/\varepsilon)$.
Funkcje $f_i=\psi_i u$ i $g_i=u\nabla\psi_i$
przedłużamy przez zero poza $\Omega$.
Obie należą do $L^1(\mathbb R^n)$ (druga
ma wartości wektorowe), ponieważ odpowiednie
mnożniki są ograniczone i mają zwarte nośniki.
Dla każdego $j\ge1$ oraz każdego $i$
z niezerowym $\psi_i$ wybieramy
$0<\varepsilon_{i,j}<d_i$ tak, aby
\begin{align*}
 \|\rho_{\varepsilon_{i,j}}*f_i-f_i\|_{L^1(\mathbb R^n)}
 &<2^{-i-j},\\
 \|\rho_{\varepsilon_{i,j}}*g_i-g_i\|_{L^1(\mathbb R^n)}
 &<2^{-i-j}.
\end{align*}
Jest to możliwe dzięki zbieżności splotów w $L^1$
z dowodu Lematu~\ref{lem:BV-translation-29},
stosowanej również do każdej składowej $g_i$.
Dla $\psi_i=0$ wybieramy dowolne dodatnie
$\varepsilon_{i,j}$; oba błędy są wtedy zerowe.
Nośniki wszystkich niezerowych wygładzeń
mieszczą się w $V_i$, ponieważ jądro ma nośnik
w kuli o promieniu $\varepsilon_{i,j}$.
Definiujemy
\[
 u_j=\sum_{i=1}^\infty\rho_{\varepsilon_{i,j}}*f_i
 \quad\text{na }\Omega.
\]
Lokalna skończoność $(V_i)$ pozwala różniczkować
sumę dowolną liczbę razy, więc $u_j\in C^\infty(\Omega)$.
Ponadto z nieujemności jądra i funkcji $\psi_i$
oraz z Tonellego mamy
\[
 \sum_i\|\rho_{\varepsilon_{i,j}}*f_i\|_1
 \le\sum_i\|f_i\|_1
 =\int_\Omega |u|\sum_i\psi_i\,\mathrm d\lambda_n
 =\|u\|_{L^1(\Omega)}.
\]
Suma zatem zbiega także w $L^1$.
Ponieważ $u=\sum_i f_i$ prawie wszędzie
na $\Omega$, nierówność trójkąta i dobór
$\varepsilon_{i,j}$ dają
$\|u_j-u\|_{L^1(\Omega)}\le\sum_i2^{-i-j}=2^{-j}$.

\emph{Krok 3: kontrola gradientu.}
Niech $Du=p\,|Du|$ będzie rozkładem polarnym.
Dla każdego $i$ funkcja $f_i$, przedłużona
przez zero, należy do $BV(\mathbb R^n)$.
Istotnie, dla pola $\varphi\in
C_c^1(\mathbb R^n;\mathbb R^n)$ test
$\psi_i\varphi$ ma zwarty nośnik w $\Omega$.
Stąd
\begin{align*}
 -\int_{\mathbb R^n} f_i\operatorname{div}\varphi
 \,\mathrm d\lambda_n
 &=-\int_\Omega u\operatorname{div}(\psi_i\varphi)
 \,\mathrm d\lambda_n
 +\int_\Omega u\nabla\psi_i\cdot\varphi\,\mathrm d\lambda_n\\
 &=\int_\Omega\psi_i\varphi\cdot\mathrm dDu
 +\int_\Omega u\nabla\psi_i\cdot\varphi\,\mathrm d\lambda_n.
\end{align*}
Prawa strona reprezentuje skończoną miarę
wektorową na $\mathbb R^n$ (przedłużoną przez zero
poza $\Omega$), co dowodzi skończoności wariacji
$f_i$ i tożsamości
$Df_i=\psi_i Du+g_i\lambda_n$.
Nie przedłużamy przez zero całej funkcji $u$
jako funkcji BV, a jedynie jej lokalne iloczyny.

Z rachunku pochodnych splotów wynika
$\nabla(\rho_{\varepsilon_{i,j}}*f_i)
=\rho_{\varepsilon_{i,j}}*(\psi_i Du)
+\rho_{\varepsilon_{i,j}}*g_i$.
W każdym punkcie $\Omega$ zachodzi
$\sum_i\nabla\psi_i=0$ jako pochodna lokalnie
skończonej sumy $\sum_i\psi_i=1$.
Zatem $\sum_i g_i=0$ prawie wszędzie, w sensie
lokalnie skończonej sumy (nie zakładamy
$\sum_i\|g_i\|_1<\infty$).
Otrzymujemy
\[
 \nabla u_j
 =\sum_i\rho_{\varepsilon_{i,j}}*(\psi_i Du)
 +\sum_i(\rho_{\varepsilon_{i,j}}*g_i-g_i)
 \quad\text{prawie wszędzie na }\Omega.
\]
Pierwszą sumę szacujemy przez rozkład polarny,
nierówność trójkąta i Tonellego:
\begin{align*}
 \int_\Omega\left|\sum_i
 \rho_{\varepsilon_{i,j}}*(\psi_i Du)(x)\right|
 \,\mathrm d\lambda_n(x)
 &\le\sum_i\int_{\mathbb R^n}\int_\Omega
 \rho_{\varepsilon_{i,j}}(x-y)\psi_i(y)
 \,\mathrm d|Du|(y)\,\mathrm d\lambda_n(x)\\
 &=\sum_i\int_\Omega\psi_i(y)\,\mathrm d|Du|(y)
 =|Du|(\Omega).
\end{align*}
W szczególności suma norm $L^1$ jej składników
jest skończona. Suma norm $L^1$ składników
drugiej sumy nie przekracza $2^{-j}$ z doboru
promieni wygładzeń. Stąd
\[
 \int_\Omega|\nabla u_j|\,\mathrm d\lambda_n
 \le |Du|(\Omega)+2^{-j}.
\]
Funkcja $u_j$ jest całkowalna, a jej klasyczne
pierwsze pochodne są całkowalne. Z całkowania
przez części względem testów o zwartym nośniku
wynika, że są to również jej pochodne słabe,
więc $u_j\in W^{1,1}(\Omega)$.

\emph{Krok 4: dokładna granica wariacji.}
Dla funkcji gładkiej z całkowalnym gradientem
wariacja jest równa całce z normy gradientu
(Lemat~\ref{lem:BV-smooth-29}).
Dolna półciągłość wariacji przy zbieżności
w $L^1$ z Twierdzenia~\ref{thm:perimeter-lsc-29}
daje
\[
 |Du|(\Omega)
 \le\liminf_{j\to\infty}|Du_j|(\Omega)
 =\liminf_{j\to\infty}\int_\Omega|\nabla u_j|
 \,\mathrm d\lambda_n
 \le\limsup_{j\to\infty}\int_\Omega|\nabla u_j|
 \,\mathrm d\lambda_n
 \le|Du|(\Omega).
\]
Wszystkie nierówności są zatem równościami.
Dowód nie wykorzystuje żadnej regularności
brzegu $\Omega$. Zbieżność w $L^1$ wraz ze
zbieżnością całkowitych wariacji nazywamy
\emph{zbieżnością ścisłą w $BV$}; nie oznacza
ona na ogół zbieżności w normie przestrzeni $BV$.
Na przykład dla $u=\chi_{(0,1)}$ na $\mathbb R$
miara $Du=\delta_0-\delta_1$ jest skupiona
na dwóch punktach, podczas gdy $Du_j$ dla
gładkich $u_j\in BV$ jest bezwzględnie ciągła
względem $\lambda_1$. Miary ich wariacji
są więc wzajemnie osobliwe i
$|D(u_j-u)|(\mathbb R)=|Du_j|(\mathbb R)+2\ge2$,
choć można dobrać $u_j$ zbieżne ściśle do $u$.
\end{proof}

\begin{lemat}[Obwody zbiorów nadpoziomicowych funkcji gładkiej]
\label{lem:BV-smooth-levels}
Niech $n\ge1$, $\Omega\subseteq\mathbb R^n$ będzie
zbiorem otwartym, a $v\in C^1(\Omega;\mathbb R)$
spełnia $\int_\Omega|\nabla v(x)|\,\mathrm d\lambda_n(x)<\infty$.
Dla $t\in\mathbb R$ oznaczmy $E_t=\{x\in\Omega:v(x)>t\}$.
Wówczas dla $\lambda_1$-prawie każdego $t$
zbiór $E_t$ ma skończony obwód w $\Omega$ oraz
jako miary wektorowe na $\Omega$ zachodzi
\[
 D\chi_{E_t}=
 \frac{\nabla v}{|\nabla v|}\,
 \mathcal H^{n-1}\llcorner
 \{x\in\Omega:v(x)=t,\ \nabla v(x)\ne0\}.
\]
Ponadto
\[
 \int_{-\infty}^{\infty}P(\{v>t\};\Omega)
 \,\mathrm d\lambda_1(t)
 \le\int_\Omega|\nabla v(x)|\,\mathrm d\lambda_n(x).
\]
\end{lemat}
\begin{proof}
Wybierzmy gładką, niemalejącą funkcję
$\theta:\mathbb R\to[0,1]$, równą $0$ na
$(-\infty,0]$ i $1$ na $[1,\infty)$, i niech
$\theta_\varepsilon(s)=\theta(s/\varepsilon)$.
Dla pola $\varphi\in C_c^1(\Omega;\mathbb R^n)$
całkowanie przez części daje
\[
 \int_\Omega\theta_\varepsilon(v(x)-t)
 \operatorname{div}\varphi(x)\,\mathrm d\lambda_n(x)
 =-\int_\Omega\theta'_\varepsilon(v(x)-t)
 \nabla v(x)\cdot\varphi(x)\,\mathrm d\lambda_n(x).
\]
Dla każdego $x$ i każdego $t$ funkcja
$\theta_\varepsilon(v(x)-t)$ zbiega do
$\chi_{E_t}(x)$, również gdy $v(x)=t$,
ponieważ $\theta(0)=0$. Całkowanie odbywa się
na zwartym nośniku dywergencji pola testowego,
więc twierdzenie o zbieżności zdominowanej
daje granicę $\int_{E_t}\operatorname{div}\varphi$
dla każdego $t$. Nie trzeba wykluczać poziomic
o dodatniej mierze Lebesgue'a.
Z kolei wzór koobszaru dla odwzorowania
$v:\Omega\to\mathbb R$ pozwala zapisać
prawą stronę jako
\[
 -\int_{\mathbb R}\theta'_\varepsilon(s-t)
 \left(\int_{\{v=s,\,\nabla v\ne0\}}
 \frac{\nabla v(x)}{|\nabla v(x)|}\cdot\varphi(x)
 \,\mathrm d\mathcal H^{n-1}(x)\right)
 \,\mathrm d\lambda_1(s).
\]
Wzór koobszaru stosujemy do borelowskich
funkcji wagowych, rozdzielając części dodatnią
i ujemną iloczynu skalarnego. Na dowolnym
wyczerpaniu zbiorami względnie zwartymi
odwzorowanie $v$ jest lokalnie Lipschitzowskie;
wzór z Rozdziału~\ref{ch:area-coarea-28}
można stosować na przeliczalnym pokryciu
kulami o zwartych domknięciach w $\Omega$,
a następnie sumować po rozłącznych borelowskich
fragmentach tych kul. W szczególności funkcja
$M(t)=\mathcal H^{n-1}(\{v=t,\nabla v\ne0\})$
jest mierzalna względem $\lambda_1$ i
$\int_{\mathbb R}M(t)\,\mathrm d\lambda_1(t)
=\int_\Omega|\nabla v|\,\mathrm d\lambda_n<\infty$.
Oznaczmy funkcję w nawiasie przez $F_\varphi(s)$.
Jest ona mierzalna względem $\lambda_1$,
a jej moduł jest ograniczony przez
$\|\varphi\|_\infty M(s)$, zatem należy do $L^1$.
Wartości na zbiorze wyjątkowym, na którym
całka powierzchniowa nie jest skończona,
ustalamy jako zero.
Ponieważ $\theta'\ge0$, jej nośnik leży
w $[0,1]$, a całka wynosi $1$,
$\theta'_\varepsilon(s-t)$ jest jądrem
jednostronnej średniej. W każdym punkcie
Lebesgue'a $t$ funkcji $F_\varphi$ mamy
\[
 \left|\int_{\mathbb R}\theta'_\varepsilon(s-t)
 F_\varphi(s)\,\mathrm d\lambda_1(s)-F_\varphi(t)\right|
 \le\frac{\|\theta'\|_\infty}{\varepsilon}
 \int_t^{t+\varepsilon}|F_\varphi(s)-F_\varphi(t)|
 \,\mathrm d\lambda_1(s)\longrightarrow0.
\]
Twierdzenie Lebesgue'a o różniczkowaniu daje
zatem dla prawie wszystkich $t$
\[
 \int_{\{v>t\}}\operatorname{div}\varphi(x)
 \,\mathrm d\lambda_n(x)
 =-\int_{\{v=t,\,\nabla v\ne0\}}
 \frac{\nabla v(x)}{|\nabla v(x)|}\cdot\varphi(x)
 \,\mathrm d\mathcal H^{n-1}(x).
\]
Aby uzyskać jeden zbiór poziomów dla wszystkich
testów, stosujemy konstrukcję ośrodkowości
z dowodu Lematu~\ref{lem:BV-level-meas-auditVI},
tym razem bez ograniczenia normy pól.
Na każdym zwartym zbiorze wyczerpania
wybieramy przeliczalną rodzinę gęstą w normie
$C^1$ w przestrzeni pól o nośniku w tym zbiorze.
Usuwamy z prostej zbiór zerowy, na którym
$M(t)=+\infty$, oraz zbiory wyjątkowe
dla wszystkich pól z sumy tych rodzin.
Dla każdego pozostałego $t$ równość zachodzi
dla całej przeliczalnej rodziny. Lewa strona
jest ciągła w normie $C^1$ na ustalonym
zwartym nośniku, a prawa w normie supremum,
z oszacowaniem przez $M(t)$ razy normę pola.
Równość przenosi się zatem na dowolny test.
Prawa strona reprezentuje, z uwzględnieniem
znaku w definicji pochodnej dystrybucyjnej,
miarę wektorową z tezy. Jest ona skończona,
a jej gęstość względem miary powierzchniowej
ma normę $1$. Z twierdzenia o reprezentacji
pochodnej i wariacji otrzymujemy
$P(E_t;\Omega)=M(t)$ dla tych poziomów.
Zatem nawet równość całek zachodzi w podanym
oszacowaniu. Dowód nie wymaga twierdzenia
Sarda ani tego, aby $v$ należała do $L^1(\Omega)$.
\end{proof}

\begin{twierdzenie}[Pełna formuła koobszaru dla $BV$]
\label{thm:BV-coarea-full-29}
Niech $\Omega\subseteq\mathbb R^n$ będzie zbiorem otwartym.
Jeżeli $u\in BV(\Omega)$, to dla prawie wszystkich
$t\in\mathbb R$ zbiór $\{x\in\Omega:u(x)>t\}$
ma skończony obwód w $\Omega$ oraz
\begin{equation}\label{eq:BV-coarea-full-29}
 |Du|(\Omega)=\int_{-\infty}^{\infty}
 P(\{u>t\};\Omega)\,\mathrm d\lambda_1(t).
\end{equation}
W szczególności dla każdego $u\in BV_{\mathrm{loc}}(\Omega)$
i otwartego $V\subseteq\Omega$ zachodzi równość,
być może z dwiema stronami równymi $+\infty$,
\[
 |Du|(V)=\int_{-\infty}^{\infty}
 P(\{u>t\};V)\,\mathrm d\lambda_1(t).
\]
\end{twierdzenie}
\begin{proof}
Najpierw niech $u\in BV(\Omega)$.
Wybierzmy z Lematu~\ref{lem:BV-strict-approx}
ciąg $u_j\in C^\infty(\Omega)\cap W^{1,1}(\Omega)$
zbiegający do $u$ w $L^1$, ze zbieżnością
całek z gradientu do $|Du|(\Omega)$.
Z konstrukcji w Lemacie~\ref{lem:BV-strict-approx}
możemy przyjąć $\|u_j-u\|_{L^1}\le2^{-j}$,
więc $\sum_j\|u_j-u\|_{L^1}<\infty$.
Wybieramy skończony borelowski reprezentant $u$,
bez zmiany żadnego obwodu.
Dla $a,b\in\mathbb R$ obie funkcje
$t\mapsto\chi_{\{a>t\}}$ i
$t\mapsto\chi_{\{b>t\}}$ różnią się na
przedziale między $a$ i $b$, z ewentualną
różnicą na końcach. Dlatego tożsamość liczbowa
$\int_{\mathbb R}|\chi_{\{a>t\}}-
\chi_{\{b>t\}}|\,\mathrm d\lambda_1(t)=|a-b|$
i twierdzenie Tonellego dają
\[
 \int_{\mathbb R}\|\chi_{\{u_j>t\}}-
 \chi_{\{u>t\}}\|_{L^1(\Omega)}
 \,\mathrm d\lambda_1(t)=\|u_j-u\|_{L^1(\Omega)}.
\]
Niech $a_j(t)$ oznacza całkę z modułu różnicy
funkcji charakterystycznych po $\Omega$.
Z Tonellego i powyższej tożsamości
$\int_{\mathbb R}\sum_j a_j(t)\,\mathrm d\lambda_1(t)
=\sum_j\|u_j-u\|_1<\infty$.
Zatem dla prawie wszystkich $t$ wszystkie
$a_j(t)$ są skończone i $a_j(t)\to0$.
W szczególności funkcje charakterystyczne
zbiegają w $L^1_{\rm loc}(\Omega)$.
Nie twierdzimy, że każda z nich należy
do globalnego $L^1(\Omega)$: dla ujemnych
poziomów i zbioru $\Omega$ o nieskończonej
mierze może to być fałszywe.
Półciągłość obwodu,
Lemat~\ref{lem:BV-level-meas-auditVI},
nierówność Fatou i Lemat~\ref{lem:BV-smooth-levels}
dają
\begin{align*}
 \int_{\mathbb R}P(\{u>t\};\Omega)\,
 \mathrm d\lambda_1(t)
 &\le\liminf_j\int_{\mathbb R}
 P(\{u_j>t\};\Omega)\,\mathrm d\lambda_1(t)\\
 &\le\liminf_j\int_\Omega|\nabla u_j(x)|
 \,\mathrm d\lambda_n(x)=|Du|(\Omega).
\end{align*}
W szczególności dla prawie wszystkich $t$
zbiór nadpoziomicowy ma skończony obwód.

Otrzymana skończoność całki jest dokładnie
założeniem Twierdzenia~\ref{thm:BV-layercake-29},
w którym wykazano również równość całki
po całej prostej z sumą całek po dodatnich
i ujemnych warstwach. Stosując to twierdzenie
do $u\in L^1(\Omega)$, dostajemy
$|Du|(\Omega)\le\int_{\mathbb R}
P(\{u>t\};\Omega)\,\mathrm d\lambda_1(t)$.
To kończy dowód równości.

Dla $u\in BV_{\mathrm{loc}}(\Omega)$ i otwartego
$V\subseteq\Omega$ wybieramy wyczerpanie $V$
przez rosnące zbiory otwarte $V_j$ o zwartych
domknięciach w $\Omega$. Na każdym $V_j$
funkcja $u$ jest całkowalna i ma skończoną
wariację; stosujemy właśnie udowodnioną równość.
Mamy $|Du|(V_j)\uparrow|Du|(V)$ z ciągłości
miary od dołu. Dla każdego $t$, również
wyjątkowego, zachodzi
$P(\{u>t\};V_j)\uparrow P(\{u>t\};V)$:
obwody są niemalejące z rozszerzaniem klasy
testów, a każde pole o zwartym nośniku w $V$
ma nośnik zawarty w pewnym $V_j$.
Supremum po wszystkich testach w $V$
jest zatem supremum po rosnącej rodzinie
klas testowych w $V_j$.
Lemat~\ref{lem:BV-level-meas-auditVI}
zapewnia mierzalność funkcji obwodu
dla każdego $V_j$ i dla $V$.
Twierdzenie o zbieżności monotonicznej
daje tezę, także gdy wspólna wartość
obu stron wynosi $+\infty$.
\end{proof}

\begin{wniosek}[Wzór koobszaru z wagą]
\label{cor:BV-coarea-weighted-29}
Niech $n\ge1$, $\Omega\subseteq\mathbb R^n$
będzie zbiorem otwartym i $u\in BV_{\mathrm{loc}}(\Omega)$.
Dla prawie wszystkich $t\in\mathbb R$ funkcja
$\chi_{\{u>t\}}$ należy do $BV_{\mathrm{loc}}(\Omega)$.
Dla pozostałych $t$ przyjmujemy w całce po $t$
miarę zerową zamiast niezdefiniowanej wariacji.
Dla każdej borelowskiej $g:\Omega\to[0,+\infty]$
zachodzi
\[
 \int_\Omega g(x)\,\mathrm d|Du|(x)
 =\int_{\mathbb R}\left(\int_\Omega g(x)
 \,\mathrm d|D\chi_{\{u>t\}}|(x)\right)
 \,\mathrm d\lambda_1(t).
\]
\end{wniosek}
\begin{proof}
Wybieramy rosnące wyczerpanie $\Omega$
otwartymi zbiorami $V_j$ o zwartych domknięciach
w $\Omega$. Na mocy Twierdzenia~\ref{thm:BV-coarea-full-29}
oraz Lematu~\ref{lem:BV-level-meas-auditVI}
zbiór
\[
 T=\bigcap_{j=1}^\infty
 \{t\in\mathbb R:P(\{u>t\};V_j)<\infty\}
\]
jest borelowski i ma dopełnienie miary zero.
Dla $t\in T$ funkcja charakterystyczna
zbioru nadpoziomicowego jest lokalnie całkowalna
i ma skończoną wariację na każdym $V_j$.
Każdy względnie zwarty otwarty podzbiór $\Omega$
zawiera się w pewnym $V_j$, więc funkcja ta
należy do $BV_{\rm loc}(\Omega)$.
Definiujemy dodatnią lokalnie skończoną
miarę Radona $\mu_t$ na $\Omega$ przez
$\mu_t=|D\chi_{\{u>t\}}|$ dla $t\in T$,
a $\mu_t=0$ dla $t\notin T$.

Dla każdego otwartego $V\subseteq\Omega$
funkcja $t\mapsto\mu_t(V)$ jest borelowska:
na $T$ jest równa $P(\{u>t\};V)$,
a poza $T$ jest zerowa.
Uzasadnimy rozszerzenie tej mierzalności
na wszystkie zbiory borelowskie.
Ustalmy $j$ i rozważmy klasę $\mathcal L_j$
borelowskich podzbiorów $A\subseteq V_j$,
dla których $t\mapsto\mu_t(A)$ jest borelowska.
Mamy $V_j\in\mathcal L_j$, a $\mu_t(V_j)$
jest skończona dla każdego $t$ z definicji $T$
i zerowego przedłużenia.
Jeżeli $A\in\mathcal L_j$, to
$\mu_t(V_j\setminus A)=\mu_t(V_j)-\mu_t(A)$,
więc również $V_j\setminus A\in\mathcal L_j$.
Dla ciągu parami rozłącznych zbiorów z
$\mathcal L_j$ ich suma należy do $\mathcal L_j$,
ponieważ jej masa jest sumą nieujemnych
funkcji borelowskich.
Jest to zatem układ Dynkina.
Zawiera rodzinę wszystkich względnie otwartych
podzbiorów $V_j$, zamkniętą na skończone
przecięcia, która generuje borelowskie zbiory
$V_j$. Twierdzenie o układach $\pi$ i $\lambda$
daje $\mathcal L_j=\mathcal B(V_j)$.
Dla dowolnego borelowskiego $A\subseteq\Omega$
mamy $\mu_t(A)=\lim_j\mu_t(A\cap V_j)$
z ciągłości od dołu, więc $t\mapsto\mu_t(A)$
jest borelowska.

Definiujemy na borelowskich zbiorach $\Omega$
$\eta(A)=\int_{\mathbb R}\mu_t(A)\,\mathrm d\lambda_1(t)$.
Mamy $\eta(\varnothing)=0$, a dla rozłącznych
borelowskich $A_i$ addytywność $\mu_t$
i twierdzenie o zbieżności monotonicznej dają
\[
 \eta\Bigl(\bigcup_iA_i\Bigr)
 =\int_{\mathbb R}\sum_i\mu_t(A_i)\,\mathrm d\lambda_1(t)
 =\sum_i\eta(A_i).
\]
Zatem $\eta$ jest miarą. Dla otwartego $V$
pominięcie $\mathbb R\setminus T$ nie zmienia
całki z nieujemnej funkcji obwodu, więc
Twierdzenie~\ref{thm:BV-coarea-full-29} daje
$\eta(V)=|Du|(V)$.
W szczególności $\eta(V_j)<\infty$.
Na każdym $V_j$ ograniczenia obu miar są
skończone i zgadzają się na względnie otwartych
podzbiorach, które tworzą układ $\pi$
generujący zbiory borelowskie. Z jednoznaczności
miar skończonych wynika ich zgodność na
$\mathcal B(V_j)$. Przechodząc do granicy
po wyczerpaniu dostajemy $\eta=|Du|$ na
$\mathcal B(\Omega)$.

Dla nieujemnej borelowskiej funkcji prostej
$g=\sum_{i=1}^m a_i\chi_{A_i}$ z nieujemnymi
skończonymi współczynnikami $a_i$ równość
z tezy wynika z liniowości dla skończonych
sum nieujemnych całek i tożsamości miar
$\eta(A_i)=|Du|(A_i)$. Funkcja
$t\mapsto\int g\,\mathrm d\mu_t$ jest wtedy
borelowska. Dla dowolnej nieujemnej borelowskiej
$g$ wybieramy rosnący ciąg nieujemnych funkcji
prostych $g_m\uparrow g$. Zbieżność monotoniczna
najpierw względem $\mu_t$ daje
$\int g_m\,\mathrm d\mu_t\uparrow\int g\,\mathrm d\mu_t$
dla każdego $t$, stąd mierzalność całki wewnętrznej.
Następnie zbieżność monotoniczna względem
$\lambda_1$ i $|Du|$ daje żądaną równość
także dla całek nieskończonych.
\end{proof}

\begin{przyklad}[Funkcja ze skokiem i formuła koobszaru]
Niech $\Omega=(-2,2)$ i $u=\chi_{(0,1)}$.
Wtedy $Du=\delta_0-\delta_1$ jako miara
ze znakiem, zatem $|Du|(\Omega)=2$.
Dla $0\le t<1$ mamy $\{u>t\}=(0,1)$ i
$P(\{u>t\};\Omega)=2$.
Dla $t<0$ zbiór nadpoziomicowy jest równy
$\Omega$, a dla $t\ge1$ jest pusty; w obu
przypadkach obwód względny w $\Omega$ wynosi zero. Stąd
$\int_{\mathbb R}P(\{u>t\};\Omega)
\,\mathrm d\lambda_1(t)=\int_0^1 2\,
\mathrm d\lambda_1(t)=2$, zgodnie z
Twierdzeniem~\ref{thm:BV-coarea-full-29}.
\end{przyklad}

\begin{przyklad}[Gęsta zmiana reprezentanta nie zmienia brzegu zredukowanego]
Dla $E=(0,1)^n\setminus\mathbb Q^n$
zachodzi $\chi_E=\chi_{(0,1)^n}$ prawie wszędzie.
Stąd $|D\chi_E|$ jest miarą powierzchniową
na ścianach sześcianu, a jego brzeg zredukowany
zgadza się z sumą wnętrz ścian, z dokładnością
do zbioru $\mathcal H^{n-1}$-zerowego.
Natomiast brzeg topologiczny $E$ zawiera
cały domknięty sześcian, ponieważ zarówno
$E$, jak i jego dopełnienie są gęste
wewnątrz $[0,1]^n$. To pokazuje, że
brzeg topologiczny nie może zastępować
$\partial^*E$ w twierdzeniu De Giorgiego.
\end{przyklad}

\section{Elementy teorii prądów}
\begin{definicja}[Prąd i jego brzeg]
\label{def:current-29}
Niech $n\ge1$, $\Omega\subseteq\mathbb R^n$
będzie zbiorem otwartym i $0\le k\le n$.
Przez $\mathcal D^k(\Omega)$ oznaczamy przestrzeń
gładkich rzeczywistych $k$-form o zwartym nośniku.
W standardowych współrzędnych zapisujemy
$\omega=\sum_I\omega_I\,\mathrm dx^I$, gdzie
$I=(i_1,\ldots,i_k)$ przebiega rosnące układy
indeksów, a $\mathrm dx^I=\mathrm dx^{i_1}\wedge
\cdots\wedge\mathrm dx^{i_k}$. Dla $k=0$
jest to po prostu przestrzeń funkcji testowych.
Nośnik formy jest sumą nośników jej współczynników.

Dla zwartego $K\subseteq\Omega$ i liczby
całkowitej $m\ge0$ kładziemy
\[
 p_{K,m}(\omega)=
 \max_{I,\ |\alpha|\le m}\sup_{x\in K}
 |\partial^\alpha\omega_I(x)|,
\]
gdzie $\alpha\in\mathbb N_0^n$ jest wielowskaźnikiem.
Na formach o nośniku w ustalonym $K$ używamy
topologii wyznaczonej przez te półnormy, a na
całym $\mathcal D^k(\Omega)$ topologii
lokalnie wypukłej granicy induktywnej tych
przestrzeni przy zwartym wyczerpaniu $\Omega$.

\emph{Prądem wymiaru $k$} nazywamy ciągły
rzeczywisty funkcjonał liniowy
$T:\mathcal D^k(\Omega)\to\mathbb R$.
Równoważnie, dla każdego zwartego $K\subseteq\Omega$
istnieją $C_K>0$ oraz całkowite $m_K\ge0$ takie, że
\[
 |T(\omega)|\le C_Kp_{K,m_K}(\omega)
 \qquad(\operatorname{supp}\omega\subseteq K).
\]
Jeśli $T$ jest prądem wymiaru $k\ge1$, jego
\emph{brzeg} definiujemy przez
$(\partial T)(\omega)=T(\mathrm d\omega)$
dla $\omega\in\mathcal D^{k-1}(\Omega)$.
Dla $k=0$ przyjmujemy brzeg równy zeru
w zerowej przestrzeni prądów wymiaru $-1$.
\end{definicja}

\begin{uwaga}[Ciągłość i wymiar brzegu]
Kryterium półnormowe z definicji wynika z tego,
że ciągłość funkcjonału na przestrzeni form
o nośniku w $K$ oznacza ograniczoność przez
skończenie wiele półnorm, które można zastąpić
jedną, ponieważ $p_{K,m}$ rośnie wraz z $m$.
Dokładniej, z ciągłości w zerze istnieją
$m$ i $\delta>0$, dla których
$p_{K,m}(\omega)<\delta$ pociąga
$|T(\omega)|<1$. Skalowanie daje
$|T(\omega)|\le 2p_{K,m}(\omega)/\delta$,
a gdy półnorma jest zerowa, skalowanie
dowolnym czynnikiem daje $T(\omega)=0$.
Odwrotnie, takie oszacowania zapewniają
ciągłość ograniczeń do przestrzeni o ustalonym
zwartym nośniku, a własność definiująca
lokalnie wypukłą granicę induktywną
zapewnia ciągłość funkcjonału na całej
przestrzeni testowej.
\end{uwaga}

\begin{lemat}[Brzeg brzegu znika]
\label{lem:current-boundary-square-29}
Niech $\Omega\subseteq\mathbb R^n$ będzie zbiorem otwartym,
$n\ge1$, $1\le k\le n$, a $T$ będzie prądem
wymiaru $k$ na $\Omega$. Wtedy $\partial T$
jest prądem wymiaru $k-1$ oraz
$\partial(\partial T)=0$ (z przyjętą konwencją
dla $k=1$).
\end{lemat}
\begin{proof}
Operator różniczki zewnętrznej jest liniowy
i nie powiększa nośnika formy. Dla testu
$\omega\in\mathcal D^{k-1}(\Omega)$ o nośniku
w $K$ współczynniki $\mathrm d\omega$
są sumami z odpowiednimi znakami pierwszych
pochodnych współczynników $\omega$.
Stąd dla każdego $m\ge0$ zachodzi
$p_{K,m}(\mathrm d\omega)\le n p_{K,m+1}(\omega)$.
Wybierając dla $T$ stałą $C_K$ i rząd $m_K$
z definicji, otrzymujemy
\[
 |(\partial T)(\omega)|
 \le C_K p_{K,m_K}(\mathrm d\omega)
 \le n C_K p_{K,m_K+1}(\omega).
\]
Funkcjonał $\partial T$ jest liniowy i spełnia
kryterium ciągłości na każdym zwartym nośniku.
Jest więc prądem wymiaru $k-1$.

Dla $k=1$ druga teza jest konsekwencją przyjętej
konwencji. Jeśli $k\ge2$, to dla
$\omega\in\mathcal D^{k-2}(\Omega)$ mamy
\[
 (\partial(\partial T))(\omega)
 =(\partial T)(\mathrm d\omega)
 =T(\mathrm d^2\omega)=0.
\]
Ostatnia równość wynika z przemienności
mieszanych pochodnych drugiego rzędu i
antysymetrii iloczynu zewnętrznego:
w wyrażeniu $\mathrm d^2\omega$ składniki
z powtórzonym indeksem znikają, a składniki
z różnymi indeksami występują parami
z przeciwnymi znakami. Z dowolności testu
wynika $\partial(\partial T)=0$.
\end{proof}

\begin{przyklad}[Prąd zadany przez powierzchnię]
Niech $\Omega\subseteq\mathbb R^n$ będzie
otwarty, $1\le k\le n$ i niech $S\subseteq\Omega$
będzie zorientowaną zwartą gładką
$k$-wymiarową podrozmaitością z brzegiem,
z orientacją brzegu indukowaną przez normalną
zewnętrzną w $TS$ (normalna zewnętrzna
jest pierwszym wektorem dodatniej bazy $TS$
na brzegu). Jeżeli $\iota:S\to\Omega$ jest
inkluzją, definiujemy
$T_S(\omega)=\int_S\iota^*\omega$.
Sprawdzimy, że jest to prąd wymiaru $k$.
Niech $\tau_S(x)$ oznacza jednostkowy prosty
$k$-wektor wyznaczający orientację $T_xS$.
Wówczas
\[
 T_S(\omega)=\int_S\omega_x(\tau_S(x))
 \,\mathrm d\mathcal H^k(x).
\]
Tożsamość wynika ze wzoru na pole dla lokalnych
parametryzacji i definicji całkowania form:
w dodatnio zorientowanej parametryzacji
iloczyn zewnętrzny wektorów stycznych jest
$\tau_S$ pomnożonym przez pierwiastek
z wyznacznika macierzy Grama.
W standardowej bazie ortonormalnej przestrzeni
$k$-wektorów współrzędne jednostkowego
$\tau_S$ mają sumę kwadratów równą $1$.
Zatem dla formy o nośniku w $K$
nierówność Cauchy'ego--Schwarza daje
\[
 |T_S(\omega)|
 \le\binom nk^{1/2}p_{K,0}(\omega)\,
 \mathcal H^k(S\cap K).
\]
Miara $\mathcal H^k(S)$ jest skończona:
zwartość $S$ pozwala wybrać skończenie wiele
lokalnych parametryzacji, również przy brzegu,
ograniczonych do zwartych fragmentów dziedzin;
na każdym z nich wzór na pole i ograniczoność
różniczki dają skończoną miarę obrazu.
Otrzymane oszacowanie dowodzi ciągłości $T_S$,
a liniowość wynika z liniowości całki.

Dla $\omega\in\mathcal D^{k-1}(\Omega)$
różniczka zewnętrzna komutuje z cofnięciem
formy przez inkluzję. Twierdzenie Stokesa,
z opisaną orientacją brzegu, daje
\[
 (\partial T_S)(\omega)
 =\int_S\mathrm d(\iota^*\omega)
 =\int_{\partial S}(\iota|_{\partial S})^*\omega
 =T_{\partial S}(\omega).
\]
Jeśli brzeg jest pusty, ostatni prąd jest zerowy.
Właśnie w tej postaci klasyczny rachunek na
rozmaitościach staje się szczególnym przypadkiem
teorii prądów. Dla przykładu na prostej,
gdy $S=[a,b]$ jest zorientowany zgodnie
z rosnącą współrzędną, dla testu $f$ mamy
$(\partial T_S)(f)=\int_a^b f'(t)\,
\mathrm d\lambda_1(t)=f(b)-f(a)$.
Stąd $\partial T_S=\delta_b-\delta_a$,
a nie $\delta_a-\delta_b$.
\end{przyklad}

\begin{przyklad}[Prąd zbioru o lokalnie skończonym obwodzie]
Niech $E\subseteq\mathbb R^n$ będzie mierzalnym
zbiorem lokalnie skończonego obwodu, a
$\Omega\subseteq\mathbb R^n$ zbiorem otwartym.
Oznaczmy standardową formę objętości przez
$\mathrm{vol}=\mathrm dx^1\wedge\cdots\wedge\mathrm dx^n$.
Definiujemy prąd wymiaru $n$ przez
$T_E(f\,\mathrm{vol})=\int_{E\cap\Omega}
f\,\mathrm d\lambda_n$.
Dla testu o nośniku w zwartym $K\subseteq\Omega$
jego wartość ma moduł nie większy niż
$\lambda_n(K)\|f\|_\infty$, więc jest to prąd
bez założenia skończonej objętości $E$.

Każdą $(n-1)$-formę testową można jednoznacznie
zapisać jako $\omega=\iota_\varphi\mathrm{vol}$,
gdzie $\varphi\in C_c^\infty(\Omega;\mathbb R^n)$,
a $\iota_\varphi$ oznacza kontrakcję przez wektor.
Współrzędnie
\[
 \omega=\sum_{i=1}^n(-1)^{i-1}\varphi_i\,
 \mathrm dx^1\wedge\cdots\wedge
 \widehat{\mathrm dx^i}\wedge\cdots\wedge\mathrm dx^n.
\]
Daszek oznacza pominięcie czynnika. Różniczkując
ten wzór, widzimy, że przeżywają tylko pochodne
$\partial_i\varphi_i$; przestawienie $\mathrm dx^i$
na właściwe miejsce daje drugi czynnik
$(-1)^{i-1}$. Dlatego
$\mathrm d\omega=(\operatorname{div}\varphi)\mathrm{vol}$.
Po przedłużeniu pola przez zero poza $\Omega$
twierdzenie strukturalne De Giorgiego daje
\[
 (\partial T_E)(\omega)
 =\int_{E\cap\Omega}\operatorname{div}\varphi\,
 \mathrm d\lambda_n
 =\int_{\partial^*E\cap\Omega}\varphi\cdot\nu_E^*
 \,\mathrm d\mathcal H^{n-1}.
\]
Znak jest zgodny z orientacją brzegu:
$(\iota_\varphi\mathrm{vol})(\tau)
=\mathrm{vol}(\varphi,\tau)$, gdzie
$\tau$ jest jednostkowym $(n-1)$-wektorem
orientującym $\nu_E^*(x)^\perp$ tak,
aby $\mathrm{vol}(\nu_E^*,\tau)=1$.
Rozkładając $\varphi$ na składową normalną
i styczną, dostajemy
$\mathrm{vol}(\varphi,\tau)=\varphi\cdot\nu_E^*$,
bo składnik styczny daje wyznacznik zerowy.
Formuła jest więc miarowym odpowiednikiem
klasycznej całki brzegowej.
Nie pojawia się tu dodatkowy wkład od
$\partial\Omega$, ponieważ testy mają nośnik
zwarty wewnątrz $\Omega$.
Na prostej dla $\Omega=\mathbb R$ i $E=(a,b)$
otrzymujemy $\partial T_E=\delta_b-\delta_a$,
podczas gdy $D\chi_E=\delta_a-\delta_b$:
znak wynika z różnych definicji obu operatorów.
\end{przyklad}

\section{Podsumowanie zależności i dalsze kierunki}
W rozdziale otrzymaliśmy reprezentację pochodnych
$BV$ przez miary, półciągłość i zwartość obwodu,
miarową formułę Gaussa--Greena, twierdzenie strukturalne
De Giorgiego wraz z geometrycznym opisem brzegu
zredukowanego oraz pełną formułę koobszaru dla $BV$.
Dowód twierdzenia strukturalnego opierał się na
nierównościach izoperymetrycznych, zwartości w $BV$,
granicach zbiorów po przeskalowaniu, prostowalności
wyprowadzonej z lokalnej płaskości oraz
identyfikacji gęstości miar. Skupienie wariacji
na brzegu zredukowanym uzyskaliśmy niezależnie
przez różniczkowanie względem miary Radona.
Dowód formuły koobszaru był niezależny od
tej struktury: wykorzystał ścisłą aproksymację
gładką i wzór koobszaru dla odwzorowań Lipschitzowskich.
Dalsze zagadnienia to m.in. dokładni reprezentanci
funkcji $BV$, twierdzenia o śladach, reguły
łańcuchowe oraz teoria powierzchni minimalnych.
Teoria prądów wiąże wzór Gaussa--Greena z operatorem
brzegu, ale nie była potrzebna do dowodu twierdzenia
strukturalnego ani wzoru koobszaru.

\begin{uwaga}[Miary wektorowe w dalszym ciągu skryptu]
Dla $u\in BV(\Omega)$ miara $Du$ ma skończoną wariację
i jest szczególnym przypadkiem miary o wartościach
w przestrzeni Banacha, omawianej w
Rozdziale~\ref{ch:vector-measures-30}.
W przypadku $u\in BV_{\mathrm{loc}}(\Omega)$ jej wariacja
jest tylko lokalnie skończona; nie musimy wtedy otrzymać
skończonego wektora $Du(\Omega)$.
Rozdział 30 interpretuje rozkład polarny poprzez RNP
przestrzeni $\mathbb R^n$. Nie należy jednak utożsamiać
tego rozkładu z istnieniem gęstości $Du$ względem
miary Lebesgue'a.
\end{uwaga}

% ROZDZIAL 30: miary wektorowe i wlasnosc Radona-Nikodyma
\chapter[Miary wektorowe i własność Radona--Nikodyma]{Miary wektorowe, całka Bochnera i własność Radona--Nikodyma}
\label{ch:vector-measures-30}

\providecommand{\atlasref}[2]{#2}


\section{Od miar skalarnych do miar w przestrzeni Banacha}
\label{sec:vector-intro-30}
W Rozdziale~\ref{ch:geometric-measure-29} pojawiła się
pochodna $Du$ o wartościach w $\mathbb R^n$,
a w Rozdziale~\ref{ch:area-coarea-28} różniczka
odwzorowania między przestrzeniami skończenie wymiarowymi.
W pierwszym przypadku dla $u\in BV(\Omega)$ mamy
skończoną wariację miary $Du$, natomiast dla
$u\in BV_{\mathrm{loc}}(\Omega)$ tylko wariację
lokalnie skończoną. Definicja miary wektorowej
poniżej wymaga, aby wartość na każdym zbiorze
mierzalnym należała do $X$; nie należy bez dodatkowych
założeń podstawiać do niej lokalnie skończonej
miary $Du$ na całym $\Omega$. Teraz pozwalamy, aby miara przyjmowała
wartości w \emph{dowolnej} rzeczywistej lub zespolonej przestrzeni
Banacha $X$. Szczególnie ważne okaże się pytanie, kiedy istnieje
funkcja $g:\Omega\to X$ taka, że $\nu(A)=\int_Ag\,\mathrm d\mu$.
W przypadku skalarnym odpowiedź daje twierdzenie Radona--Nikodyma;
w przypadku przestrzeni Banacha zależy ona od geometrii $X$.

\begin{definicja}[Miara wektorowa]
\label{def:vector-measure-30}
Niech $(\Omega,\mathcal A)$ będzie przestrzenią mierzalną,
a $X$ przestrzenią Banacha. \emph{Miarą wektorową o wartościach w $X$}
nazywamy odwzorowanie $\nu:\mathcal A\to X$, dla którego
$\nu(\varnothing)=0$ i dla każdej rodziny parami rozłącznych
$A_j\in\mathcal A$ zachodzi
\[
 \nu\left(\bigcup_{j=1}^{\infty}A_j\right)
 =\lim_{N\to\infty}\sum_{j=1}^{N}\nu(A_j)
 \quad\text{w normie }X.
\]
Przez $\|\nu(A)\|$ oznaczamy normę \emph{wektora} $\nu(A)$.
Nie należy jej mylić z wariacją miary na zbiorze $A$.
\end{definicja}

\begin{przyklad}[Dwie podstawowe konstrukcje]
\label{ex:basic-vector-measures-30}
Jeżeli $\mu$ jest dodatnią miarą skończoną, a $x_0\in X$,
to $\nu(A)=\mu(A)x_0$ jest miarą wektorową.
Jeśli $(\Omega,\mathcal A,\mu)$ jest przestrzenią o mierze
skończonej, a $g:\Omega\to X$ jest całkowalna w sensie
Bochnera (definicja poniżej), to
$\nu_g(A)=\int_Ag\,\mathrm d\mu$ jest miarą wektorową.
Zbieżność w normie sum po rozłącznych zbiorach
wynika z oszacowania
$\|\int_Bg\,\mathrm d\mu\|\le\int_B\|g\|\,\mathrm d\mu$
i bezwzględnej ciągłości całki skalarnej.
\end{przyklad}

\section{Wariacja całkowita i półwariacja}
\label{sec:vector-variation-30}
\begin{definicja}[Wariacja i półwariacja]
\label{def:vector-variation-30}
Niech $\nu:\mathcal A\to X$ będzie miarą wektorową.
Piszemy $\nu\ll\mu$ dla dodatniej miary $\mu$
na $(\Omega,\mathcal A)$, gdy $\mu(A)=0$ implikuje
$\nu(A)=0$ dla każdego $A\in\mathcal A$.
Dla $A\in\mathcal A$ definiujemy \emph{wariację całkowitą}
\[
 |\nu|(A)=\sup\left\{\sum_{j=1}^{N}\|\nu(A_j)\|:
 A=\dot\bigcup_{j=1}^{N}A_j,\quad A_j\in\mathcal A,\ N\ge1\right\}.
\]
Jeżeli $|\nu|(\Omega)<\infty$, mówimy, że $\nu$ ma
\emph{skończoną wariację}. Dla $x^*\in X^*$ odwzorowanie
$A\mapsto x^*(\nu(A))$ jest miarą skalarną.
\emph{Półwariacją} nazywamy
$\|\nu\|_{\mathrm{sv}}(A)
=\sup_{\|x^*\|\le1}|x^*\nu|(A)$.
\end{definicja}

\begin{twierdzenie}[Wariacja jest miarą dodatnią]
\label{thm:vector-variation-measure-30}
Dla każdej miary wektorowej $\nu:\mathcal A\to X$
funkcja $|\nu|:\mathcal A\to[0,+\infty]$ jest miarą.
Ponadto dla każdego $A\in\mathcal A$ zachodzi
$\|\nu(A)\|\le\|\nu\|_{\mathrm{sv}}(A)\le|\nu|(A)$.
\end{twierdzenie}
\begin{proof}
Mamy $|\nu|(\varnothing)=0$, ponieważ
$\nu(\varnothing)=0$. Monotoniczność wariacji wynika z dopisania zbioru
$B\setminus A$ do każdego podziału $A\subseteq B$.
Niech $A,B$ będą rozłączne. Łącząc podziały $A$ i $B$,
otrzymujemy $|\nu|(A\cup B)\ge|\nu|(A)+|\nu|(B)$.
W drugą stronę dowolny skończony podział
$A\cup B=\dot\bigcup_i C_i$ rozcinamy przez $A$ i $B$.
Trójkąt daje
$\sum_i\|\nu(C_i)\|
\le\sum_i\|\nu(C_i\cap A)\|+\sum_i\|\nu(C_i\cap B)\|
\le|\nu|(A)+|\nu|(B)$.
Zatem wariacja jest skończenie addytywna, także wtedy,
gdy przyjmuje wartość nieskończoną.

Jeśli $A=\dot\bigcup_{j=1}^{\infty}A_j$,
to z monotoniczności i skończonej addytywności
$|\nu|(A)\ge\sum_j|\nu|(A_j)$.
Dla nierówności odwrotnej wybierzmy dowolny
skończony podział $A=\dot\bigcup_i C_i$.
Z przeliczalnej addytywności $\nu$ w normie wynika
\[
 \|\nu(C_i)\|\le\sum_j\|\nu(C_i\cap A_j)\|.
\]
Sumujemy po skończenie wielu $i$ i zamieniamy kolejność
sum nieujemnych, uzyskując
$\sum_i\|\nu(C_i)\|\le\sum_j|\nu|(A_j)$.
Supremum po podziałach daje nierówność odwrotną.
Dla dowolnego $x^*$ z $\|x^*\|\le1$ mamy
$|x^*\nu|(A)\le|\nu|(A)$. Supremum po $x^*$
daje drugą nierówność z tezy, zaś pierwsza wynika
z twierdzenia Hahna--Banacha:
$\|\nu(A)\|=\sup_{\|x^*\|\le1}|x^*(\nu(A))|$.
\end{proof}

\begin{przyklad}[Skończona półwariacja nie implikuje skończonej wariacji]
\label{ex:semivariation-not-variation-30}
Na $\Omega=\mathbb N$ z $\mathcal A=2^{\mathbb N}$
niech $X=\ell^2$ i $\nu(A)=\sum_{n\in A}e_n/n$,
gdzie $(e_n)$ jest bazą kanoniczną. Szereg po dowolnym
zbiorze zbiega bezwarunkowo w $\ell^2$, więc $\nu$ jest dobrze określona. Dla rozłącznych
$A_j\subseteq\mathbb N$ kwadrat normy reszty spełnia
\[
 \left\|\nu\left(\bigcup_j A_j\right)
 -\sum_{j=1}^N\nu(A_j)\right\|^2
 =\sum_{n\in\bigcup_{j>N}A_j}n^{-2}\longrightarrow0,
\]
na mocy zbieżności szeregu $\sum_n n^{-2}$ i tego,
że zbiory po prawej stronie maleją do zbioru pustego.
Zatem $\nu$ jest miarą wektorową. Dla podziału na
pierwszych $N$ singletonów i ich dopełnienie
$|\nu|(\mathbb N)\ge\sum_{n=1}^{N}1/n$; zatem
$|\nu|(\mathbb N)=+\infty$.
Natomiast dla $x^*=(a_n)\in\ell^2$ z $\|x^*\|\le1$
\[
 |x^*\nu|(\mathbb N)=\sum_{n=1}^{\infty}\frac{|a_n|}{n}
 \le\left(\sum_n|a_n|^2\right)^{1/2}
 \left(\sum_n\frac1{n^2}\right)^{1/2}.
\]
Wybierając $a_n=c/n$, gdzie
$c=(\sum_n n^{-2})^{-1/2}$, otrzymujemy funkcjonał
o normie $1$ i równość w oszacowaniu.
Stąd $\|\nu\|_{\mathrm{sv}}(\mathbb N)
=(\sum_n n^{-2})^{1/2}$.
Półwariacja jest skończona, mimo że wariacja nie jest.
\end{przyklad}

\begin{lemat}[Bezwzględna ciągłość wariacji]
\label{lem:vector-variation-ac-30}
Niech $(\Omega,\mathcal A)$ będzie przestrzenią
mierzalną, $X$ przestrzenią Banacha, $\nu:\mathcal A\to X$
miarą wektorową, a $\mu$ dodatnią miarą na $\mathcal A$.
Jeśli $\nu\ll\mu$, to $|\nu|\ll\mu$.
\end{lemat}
\begin{proof}
Jeśli $\mu(A)=0$, to $\nu(B)=0$ dla każdego
mierzalnego $B\subseteq A$. Wszystkie sumy definiujące
$|\nu|(A)$ są więc zerowe.
\end{proof}

\section{Całka Bochnera i miary z gęstością}
\label{sec:bochner-30}

\noindent\emph{Powiązania tematyczne w Bibliotece matematycznej:} \atlasref{FA-15B::thm:bochner-criterion15b}{FA-15B: kryterium całkowalności Bochnera}.

\begin{definicja}[Silna mierzalność i całka Bochnera]
\label{def:bochner-30}
Niech $(\Omega,\mathcal A,\mu)$ będzie przestrzenią
z dodatnią miarą, a $X$ przestrzenią Banacha.
Funkcja $g:\Omega\to X$ jest \emph{silnie mierzalna},
jeżeli jest granicą punktową prawie wszędzie ciągu
funkcji prostych o wartościach w $X$.
Dla miary niezupełnej wybieramy mierzalnego przedstawiciela:
na mierzalnym zbiorze zbieżności ciągu funkcji prostych
przyjmujemy jego granicę, a poza nim wartość $0$.
Zbiór zbieżności jest mierzalny z kryterium Cauchy'ego
i zupełności $X$, a jego dopełnienie ma miarę zero.
Granica jest borelowsko mierzalna, ponieważ dla każdego
niepustego domkniętego $C\subseteq X$ odległość od $C$
jest ciągła, a zatem odległość granicy od $C$ jest
mierzalną funkcją skalarną. W dalszym ciągu całki
rozumiemy dla takiego przedstawiciela.
Jest \emph{całkowalna w sensie Bochnera}, gdy dodatkowo
$\int_\Omega\|g(\omega)\|\,\mathrm d\mu(\omega)<\infty$.
Dla funkcji prostej $s=\sum_{j=1}^{N}x_j\chi_{A_j}$
o rozłącznych zbiorach $A_j$ o skończonej mierze
kładziemy $\int_\Omega s\,\mathrm d\mu=\sum_j\mu(A_j)x_j$.
Dla ogólnej $g$ całkę określa granica tych sum
w normie $X$ dla aproksymacji zbieżnej w
$L^1(\mu;X)$. Przestrzeń tych klas funkcji
oznaczamy $L^1(\mu;X)$.
\end{definicja}

\begin{twierdzenie}[Konstrukcja i oszacowanie całki Bochnera]
\label{thm:bochner-basic-30}
Niech $(\Omega,\mathcal A,\mu)$ będzie przestrzenią
z dodatnią miarą, a $X$ przestrzenią Banacha.
Przestrzeń $L^1(\mu;X)$ jest zupełna, a wzór
$g\mapsto\int_\Omega g\,\mathrm d\mu$ wyznacza
ciągły operator liniowy $L^1(\mu;X)\to X$
o normie co najwyżej $1$. Dla każdego $g\in L^1(\mu;X)$
zachodzi
\[
 \left\|\int_\Omega g\,\mathrm d\mu\right\|
 \le\int_\Omega\|g\|\,\mathrm d\mu.
\]
\end{twierdzenie}
\begin{proof}
Dla funkcji prostych nierówność jest zwykłą nierównością
trójkąta. Niech $g$ będzie silnie mierzalna i
$\int_\Omega\|g\|\,\mathrm d\mu<\infty$.
Wybieramy mierzalnego przedstawiciela $g$.
Z definicji silnej mierzalności wybierzmy proste
$s_n\to g$ prawie wszędzie. Możemy dobrać je tak, by
$\|s_n(\omega)\|\le2\|g(\omega)\|$ prawie wszędzie:
w punkcie, w którym pierwotna prosta jest zbyt duża,
zastępujemy jej wartość przez $0$. Gdy $g(\omega)=0$,
nowa prosta jest zerowa, a gdy $\|g(\omega)\|>0$
i pierwotne proste zbiegają do $g(\omega)$,
od pewnego miejsca nie są zmieniane. Zbieżność
prawie wszędzie zostaje zachowana. Otrzymany zbiór
wycięcia jest mierzalny, a po podziale na skończenie
wiele wartości prostej nadal otrzymujemy funkcję prostą.
Twierdzenie Lebesgue'a o zbieżności zmajoryzowanej
zastosowane do norm daje $\|s_n-g\|_{L^1}\to0$.
Wybór przybliżeń z nośnikami o skończonej mierze
wymaga jeszcze uzasadnienia, gdy $\mu(\Omega)=\infty$.
Dla $k\ge1$ zbiór $\{\|g\|>1/k\}$ ma skończoną
miarę z nierówności Markowa. Ich suma pokrywa
$\{\|g\|>0\}$, a z twierdzenia Lebesgue'a
o zbieżności zmajoryzowanej wynika,
że $g\chi_{\{\|g\|>1/k\}}\to g$ w $L^1$.
Na każdym z tych zbiorów o skończonej mierze
powyższe funkcje proste można obcinać bez zmiany
argumentu dominacji. Zatem funkcje proste
o całkowalnych normach są gęste także bez
założenia $\sigma$-skończoności całej miary $\mu$.

Dla ciągu prostych Cauchy'ego w $L^1$ wybieramy
podciąg $s_{n_k}$ taki, że
$\sum_k\|s_{n_{k+1}}-s_{n_k}\|_{L^1}<\infty$.
Twierdzenie Beppo Leviego daje
$\sum_k\|s_{n_{k+1}}(\omega)-s_{n_k}(\omega)\|<\infty$
prawie wszędzie. Zupełność $X$ zapewnia punktową
zbieżność do silnie mierzalnej $g$; nierówność Fatou
daje
\[
 \|g-s_{n_k}\|_{L^1}
 \le \sum_{j\ge k}\|s_{n_{j+1}}-s_{n_j}\|_{L^1}.
\]
W szczególności $g$ jest całkowalna, bo $s_{n_1}$
jest całkowalna. Dla dowolnego ciągu Cauchy'ego
$(g_n)$ w $L^1(\mu;X)$ wybieramy proste $s_n$
takie, że $\|s_n-g_n\|_{L^1}<1/n$.
Ciąg $(s_n)$ jest Cauchy'ego, więc zbiega do $g$
w $L^1$ na mocy poprzedniego argumentu.
Wówczas $\|g_n-g\|_{L^1}\le1/n+\|s_n-g\|_{L^1}\to0$.
To dowodzi zupełności. Nierówność dla całki
prostej i gęstość jednoznacznie przedłużają całkę
i oszacowanie na całe $L^1(\mu;X)$.
\end{proof}

\begin{twierdzenie}[Wariacja miary o gęstości Bochnera]
Niech \label{thm:variation-density-30}$(\Omega,\mathcal A,\mu)$ będzie przestrzenią
z dodatnią miarą, $X$ przestrzenią Banacha,
$g\in L^1(\mu;X)$, a
$\nu_g(A)=\int_Ag\,\mathrm d\mu$ dla $A\in\mathcal A$.
Wtedy $\nu_g$ jest miarą wektorową o skończonej wariacji
oraz
\[
 |\nu_g|(A)=\int_A\|g(\omega)\|\,\mathrm d\mu(\omega)
 \quad(A\in\mathcal A).
\]
\end{twierdzenie}
\begin{proof}
W przypadku zupełnej miary skończonej por. \atlasref{FA-15B::thm:bochner-measure15b}{FA-15B: miara generowana przez całkę Bochnera}. Poniższy dowód nie wymaga tego ograniczenia.

Liniowość całki i tożsamość funkcji charakterystycznych
dla rozłącznych zbiorów dają skończoną addytywność
$\nu_g$ oraz $\nu_g(\varnothing)=0$. Dla parami
rozłącznych $A_j\in\mathcal A$ kładziemy
$B_N=\bigcup_{j>N}A_j$. Wówczas $B_N\downarrow\varnothing$
i
\[
 \left\|\nu_g\left(\bigcup_jA_j\right)
 -\sum_{j=1}^N\nu_g(A_j)\right\|
 =\|\nu_g(B_N)\|
 \le\int_{B_N}\|g\|\,\mathrm d\mu\longrightarrow0
\]
z twierdzenia Lebesgue'a o zbieżności zmajoryzowanej.
To dowodzi przeliczalnej addytywności.
Dla każdego podziału $A=\dot\bigcup_jA_j$ mamy
$\sum_j\|\nu_g(A_j)\|\le\int_A\|g\|\,\mathrm d\mu$.
W drugą stronę niech $s=\sum_jx_j\chi_{B_j}$
będzie całkowalną funkcją prostą z rozłącznymi $B_j$
o skończonej mierze. Podział $A$ na $A\cap B_j$
i $A\setminus\bigcup_j B_j$ daje
\[
 |\nu_s|(A)\ge\sum_j\mu(A\cap B_j)\|x_j\|
 =\int_A\|s\|\,\mathrm d\mu.
\]
Zatem dla $s$ mamy równość. Dla ogólnej $g$
i takiej funkcji $s$ z nierówności trójkąta na
każdym skończonym podziale $A$ otrzymujemy
$|\nu_g|(A)\le|\nu_s|(A)+|\nu_{g-s}|(A)$.
Zamieniając $g$ i $s$ miejscami, dostajemy
\[
 \big||\nu_g|(A)-|\nu_s|(A)\big|
 \le|\nu_{g-s}|(A)\le\int_A\|g-s\|\,\mathrm d\mu.
\]
Wszystkie odejmowane wariacje są skończone na mocy
już udowodnionej nierówności górnej.
Ponadto $|\|g\|-\|s\||\le\|g-s\|$ punktowo.
Wybierając proste $s_k\to g$ w $L^1$, uzyskujemy
zarówno $|\nu_{s_k}|(A)\to|\nu_g|(A)$, jak i
$\int_A\|s_k\|\,\mathrm d\mu\to\int_A\|g\|\,\mathrm d\mu$.
Równość dla funkcji prostych przechodzi więc na $g$.
\end{proof}

\begin{lemat}[Różniczkowanie całki Bochnera]
\label{lem:bochner-differentiation-30}
Niech $a,b\in\mathbb R$, $a<b$, a $X$ będzie
przestrzenią Banacha. Jeśli $g\in L^1((a,b),\lambda_1;X)$, to
\[
 \lim_{h\to0}\left\|
 \frac1h\int_{(t,t+h)}g(s)\,\mathrm d\lambda_1(s)-g(t)
 \right\|=0
\]
dla prawie każdego $t\in(a,b)$; dla $h<0$ całka
po przedziale jest rozumiana ze znakiem, tzn.
$\int_{(t,t+h)}=-\int_{(t+h,t)}$.
\end{lemat}
\begin{proof}
Por. \atlasref{FA-15B::thm:bochner-lebesgue-differentiation15b}{FA-15B: różniczkowanie całki Bochnera i bezwzględna ciągłość jej funkcji pierwotnej}.

Wystarczy wykazać, że prawie każdy punkt jest
punktem Lebesgue'a funkcji $X$-wartościowej, czyli
średnia z $\|g(s)-g(t)\|$ na małych symetrycznych
przedziałach dąży do zera. Funkcje ciągłe
$\phi:[a,b]\to X$ są gęste w $L^1((a,b);X)$:
najpierw aproksymujemy $g$ funkcjami prostymi,
a funkcje charakterystyczne zbiorów mierzalnych
aproksymujemy w $L^1$ funkcjami ciągłymi,
korzystając z regularności miary Lebesgue'a.
Niech $\phi_j$ zbiega do $g$ w $L^1$ tak szybko, aby
$\sum_j2^j\|g-\phi_j\|_{L^1}<\infty$.
Dla $a_j(t)=\|g(t)-\phi_j(t)\|$ mamy
$\sum_j2^j a_j(t)<\infty$ dla prawie wszystkich $t$,
na mocy twierdzenia Beppo Leviego. Przedłużamy
$a_j$ przez $0$ poza $(a,b)$; dla ustalonego
$t\in(a,b)$ dostatecznie małe przedziały
symetryczne są zawarte w $(a,b)$.
Słaba nierówność maksymalna z Rozdziału
\ref{ch:differentiation-25} daje dla ustalonego
$\varepsilon>0$ oszacowanie
\[
 \lambda_1\!\left(\left\{t:
 \sup_{r>0}\frac1{2r}\int_{(t-r,t+r)}a_j(s)
 \,\mathrm d\lambda_1(s)>\varepsilon\right\}\right)
 \le C\varepsilon^{-1}\|g-\phi_j\|_{L^1}.
\]
Oznaczmy zbiory z lewej strony oszacowania przez
$E_j(\varepsilon)$. Ponieważ
$\sum_j\lambda_1(E_j(\varepsilon))<\infty$,
\[
 \lambda_1\left(\bigcap_N\bigcup_{j\ge N}E_j(\varepsilon)\right)=0.
\]
Istotnie, miara tego zbioru nie przekracza żadnego
ogona zbieżnego szeregu miar. Poza nim średnie
$a_j$ są nie większe niż $\varepsilon$ dla wszystkich
dostatecznie dużych $j$. Ponadto $a_j(t)\to0$
prawie wszędzie. Dla każdego takiego $t$ i każdego
dostatecznie dużego $j$ nierówność trójkąta daje
\[
 \limsup_{r\downarrow0}\frac1{2r}
 \int_{(t-r,t+r)}\|g(s)-g(t)\|\,\mathrm d\lambda_1(s)
 \le\varepsilon+a_j(t),
\]
ponieważ średnia z $\|\phi_j(s)-\phi_j(t)\|$
dąży do zera na mocy ciągłości $\phi_j$.
Przechodzimy z $j$ do nieskończoności. Stosujemy
ten argument dla przeliczalnie wielu $\varepsilon=1/m$,
$m\ge1$, i usuwamy sumę wszystkich otrzymanych
zbiorów zerowych. Na pozostałym zbiorze granica
górna jest nie większa niż $1/m$ dla każdego $m$,
a więc jest równa zero.
Średnią po przedziale jednostronnym ograniczamy
przez dwukrotność średniej symetrycznej.
\end{proof}

\section{Własność Radona--Nikodyma przestrzeni Banacha}
\label{sec:rnp-30}

\noindent\emph{Powiązania tematyczne w Bibliotece matematycznej:} \atlasref{FA-15B::def:rnp15b}{FA-15B: własność Radona--Nikodyma}.

\begin{definicja}[Własność Radona--Nikodyma]
\label{def:RNP-30}
Przestrzeń Banacha $X$ ma \emph{własność Radona--Nikodyma}
(w skrócie RNP), jeśli dla każdej przestrzeni z miarą
skończoną $(\Omega,\mathcal A,\mu)$ oraz każdej
$X$-wartościowej miary wektorowej $\nu$ o skończonej
wariacji, spełniającej $\nu\ll\mu$, istnieje
$g\in L^1(\mu;X)$ takie, że
\begin{equation}\label{eq:RNP-30}
 \nu(A)=\int_Ag(\omega)\,\mathrm d\mu(\omega)
 \qquad(A\in\mathcal A).
\end{equation}
Funkcja $g$ jest wyznaczona jednoznacznie
$\mu$-prawie wszędzie.
\end{definicja}
\begin{proof}[Jednoznaczność]
Niech $g,h\in L^1(\mu;X)$ mają takie same całki
po wszystkich $A\in\mathcal A$. Dla $k=g-h$
miara $\nu_k$ jest identycznie zerowa. Na mocy
Twierdzenia~\ref{thm:variation-density-30}
$0=|\nu_k|(\Omega)=\int_\Omega\|k\|\,\mathrm d\mu$.
Nieujemność normy daje $\|k\|=0$ prawie wszędzie,
a zatem $g=h$ prawie wszędzie.

Można też wykazać jednoznaczność przez funkcjonały
skalarne. Wybieramy ciąg funkcji prostych zbieżny
do $k$ prawie wszędzie. Domknięta podprzestrzeń $Y$
rozpięta przez wszystkie wartości tych prostych
jest separowalna i zawiera $k(\omega)$ poza jednym
zbiorem zerowym. Jeśli $Y=\{0\}$, teza jest oczywista.
W przeciwnym razie wybieramy ciąg $(y_j)$ gęsty
w sferze jednostkowej $Y$. Z twierdzenia Hahna--Banacha
istnieją $x_j^*\in X^*$ o normie $1$ i
$x_j^*(y_j)=1$. Dla $y\in Y$ o normie $1$ mamy
$|x_j^*(y)|\ge1-\|y-y_j\|$, a więc
$\|y\|=\sup_j|x_j^*(y)|$ dla każdego $y\in Y$
(po przeskalowaniu, również dla $y=0$).
Ciągły funkcjonał komutuje z całką Bochnera:
równość wynika dla prostych z liniowości, a dla
funkcji całkowalnych przez aproksymację w $L^1$.
Zatem $\int_A x_j^*(k)\,\mathrm d\mu=0$ dla każdego
$A$ i każdego $j$. Jednoznaczność skalarnej gęstości
Radona--Nikodyma daje $x_j^*(k)=0$ prawie wszędzie
(w przypadku zespolonym stosujemy ją do obu części).
Po usunięciu przeliczalnej sumy zbiorów zerowych
i skorzystaniu z normującej rodziny otrzymujemy
$\|k\|=0$. Nie potrzeba separowalności całego $X$.
\end{proof}

\begin{przyklad}[Przestrzenie skończenie wymiarowe]
Każda przestrzeń Banacha $X$ o skończonym wymiarze
ma RNP. Wybierając bazę $e_1,\ldots,e_d$,
zapisujemy $\nu=\sum_j\nu_je_j$.
Twierdzenie Radona--Nikodyma dla miar skalarnych
daje $\nu_j=g_j\mu$. Funkcja
$g=\sum_jg_je_j$ jest silnie mierzalna i całkowalna,
gdyż $\|g\|\le C\sum_j|g_j|$ dla pewnej stałej $C$.
Wtedy $\nu=g\mu$. Ta argumentacja nie działa
bezpośrednio przy nieskończenie wielu współrzędnych.
\end{przyklad}

\begin{lemat}[Przybliżanie miary przez skończone podziały]
\label{lem:rnp-partition-30}
Niech $(\Omega,\mathcal A,\mu)$ będzie przestrzenią
o mierze skończonej, a $\nu:\mathcal A\to X$ miarą
wektorową spełniającą $\|\nu(A)\|\le\mu(A)$
dla każdego $A\in\mathcal A$.
Dla każdego skończonego podziału mierzalnego
$\pi=\{A_1,\ldots,A_N\}$ definiujemy
\[
 g_\pi=\sum_{j:\mu(A_j)>0}
 \frac{\nu(A_j)}{\mu(A_j)}\chi_{A_j}
 \quad\in L^1(\mu;X).
\]
Wówczas $\|g_\pi\|\le1$ prawie wszędzie,
a następujące warunki są równoważne:
\begin{enumerate}[label=(\roman*)]
 \item istnieje $g\in L^1(\mu;X)$ z
       $\nu(A)=\int_Ag\,\mathrm d\mu$ dla każdego $A$;
 \item sieć $(g_\pi)$, uporządkowana przez
       rozdrobnienie skończonych podziałów,
       jest siecią Cauchy'ego w $L^1(\mu;X)$.
\end{enumerate}
Jeśli $\rho$ rozdrabnia $\pi$, to
$g_\pi=\mathbb E_\pi g_\rho$, gdzie
$\mathbb E_\pi$ oznacza uśrednienie na komórkach
$\pi$. Operator $\mathbb E_\pi$ jest kontrakcją
na $L^1(\mu;X)$.
\end{lemat}
\begin{proof}
Na komórce $A_j$ uśrednienie funkcji $h$ jest równe
$\mu(A_j)^{-1}\int_{A_j}h\,\mathrm d\mu$
(na komórkach zerowych przyjmujemy $0$).
Z nierówności trójkąta dla całki Bochnera mamy
$\|\mathbb E_\pi h\|_{L^1}\le\|h\|_{L^1}$.
Jeżeli $\rho$ rozdrabnia $\pi$, to
$\int_{A_j}g_\rho\,\mathrm d\mu=\nu(A_j)$,
co daje równanie z tezy.

Załóżmy (i). Wtedy $g_\pi=\mathbb E_\pi g$.
Dla $\varepsilon>0$ wybierzmy funkcję prostą
$s$ z $\|g-s\|_{L^1}<\varepsilon$ i podział
$\pi_0$, na którego komórkach $s$ jest stała.
Dla każdego $\pi$ rozdrabniającego $\pi_0$
zachodzi $\mathbb E_\pi s=s$ w $L^1(\mu;X)$
(wartości na komórkach zerowych są nieistotne), skąd
\[
 \|g_\pi-g\|_{L^1}
 \le\|\mathbb E_\pi(g-s)\|_{L^1}+\|s-g\|_{L^1}
 <2\varepsilon.
\]
Zatem sieć jest zbieżna, a więc Cauchy'ego.

Jeśli zachodzi (ii), wybieramy podziały $\pi_n$
tak, aby $\pi_{n+1}$ rozdrabniał $\pi_n$ i aby
dla każdych $\alpha,\beta$ rozdrabniających $\pi_n$
zachodziło $\|g_\alpha-g_\beta\|_{L^1}<2^{-n}$.
Można to zrobić, biorąc wspólne rozdrobnienie
poprzedniego podziału i podziału z warunku Cauchy'ego.
Ciąg $(g_{\pi_n})$ jest Cauchy'ego, więc zbiega
do $g$ w zupełnej przestrzeni $L^1(\mu;X)$.
Dla każdego $\pi$ rozdrabniającego $\pi_n$ mamy
$\|g_\pi-g\|_{L^1}\le2^{-n}+\|g_{\pi_n}-g\|_{L^1}$,
co dowodzi zbieżności całej sieci. Ustalmy $A\in\mathcal A$.
Dla każdego podziału $\pi$ rozdrabniającego
$\{A,\Omega\setminus A\}$ mamy
$\int_Ag_\pi\,\mathrm d\mu=\nu(A)$.
Przechodząc do granicy w $L^1$ otrzymujemy
$\int_Ag\,\mathrm d\mu=\nu(A)$.
\end{proof}

\begin{lemat}[Całka po krzywej z pochodną Banachowską]
\label{lem:FTC-Banach-30}
Niech $a<b$ będą liczbami rzeczywistymi, $X$
przestrzenią Banacha, a $F:[a,b]\to X$ będzie Lipschitzowska i
różniczkowalna w normie $X$ dla prawie wszystkich
$t\in(a,b)$. Wtedy jej pochodna $F'$ ma silnie
mierzalnego przedstawiciela, jest istotnie ograniczona,
a dla wszystkich $a\le s<t\le b$ zachodzi
\[
 F(t)-F(s)=\int_{(s,t)}F'(r)\,\mathrm d\lambda_1(r).
\]
\end{lemat}
\begin{proof}
Dla $n\ge1$ kładziemy $h_n=(b-a)/2^n$ i
\[
 q_n(t)=
 \begin{cases}
 (F(t+h_n)-F(t))/h_n,&a\le t\le b-h_n,\\
 0,&b-h_n<t\le b.
 \end{cases}
\]
Na pierwszym przedziale $q_n$ jest ciągła,
więc jednostajna ciągłość pozwala przybliżać ją
dowolnie dokładnie w normie supremum funkcjami
prostymi stałymi na podprzedziałach.
Na drugim przedziale przybliżenia przyjmują zero.
W szczególności istnieje prosta $s_n$ z
$\sup_{t\in[a,b]}\|s_n(t)-q_n(t)\|\le1/n$.
W każdym punkcie różniczkowalności $t\in(a,b)$
mamy $q_n(t)\to F'(t)$, a więc $s_n(t)\to F'(t)$.
Na mierzalnym zbiorze zbieżności $(s_n)$ bierzemy
jego granicę, a poza nim przyjmujemy zero, zgodnie
z Definicją~\ref{def:bochner-30}. Uzyskujemy silnie
mierzalnego przedstawiciela pochodnej.
Ponieważ $\|q_n(t)\|\le\operatorname{Lip}(F)$,
mamy $\|F'\|\le\operatorname{Lip}(F)$ prawie wszędzie.
Niech $x^*\in X^*$. Funkcja skalarna
$x^*\circ F$ jest Lipschitzowska, a prawie wszędzie
$(x^*\circ F)'=x^*(F')$. Lemat~\ref{lem:rademacher-line-28}, zastosowany
do funkcji rzeczywistej $x^*\circ F$ lub do jej
części rzeczywistej i urojonej w przypadku zespolonym, daje
\[
 x^*(F(t)-F(s))
 =\int_{(s,t)}x^*(F'(r))\,\mathrm d\lambda_1(r)
 =x^*\!\left(\int_{(s,t)}F'(r)\,\mathrm d\lambda_1(r)\right).
\]
Ponieważ zachodzi to dla każdego $x^*$,
Hahn--Banach dowodzi równości wektorów.
\end{proof}

\begin{lemat}[Dzielenie miary bezatomowej]
\label{lem:nonatomic-divisibility-30}
Niech $\mu$ będzie skończoną miarą bezatomową.
Dla każdego mierzalnego $A$ i każdej liczby
$0\le t\le\mu(A)$ istnieje mierzalny $B\subseteq A$
taki, że $\mu(B)=t$. W konsekwencji dla $\eta>0$
każdy mierzalny $A$ można podzielić na skończenie
wiele zbiorów miary co najwyżej $\eta$.
\end{lemat}
\begin{proof}
Najpierw zauważmy, że każdy zbiór dodatniej miary
zawiera podzbiór o dowolnie małej dodatniej mierze.
Istotnie, jeśli $E$ ma dodatnią miarę, z braku atomów
istnieje $D\subset E$ z $0<\mu(D)<\mu(E)$.
Mniejszy z $D$ i $E\setminus D$ ma dodatnią miarę
nie większą niż $\mu(E)/2$. Iterując otrzymujemy
zbiory dodatniej miary mniejszej niż zadane $\eta$.

Dla $0<t<\mu(A)$ rozważmy rodzinę klas
\[
 \mathcal C=\{[B]: B\subseteq A,\ B\in\mathcal A,\ \mu(B)\le t\},
\]
gdzie $[B]=[C]$ wtedy i tylko wtedy, gdy
$\mu(B\triangle C)=0$. Relacja $[B]\le[C]$
oznacza $\mu(B\setminus C)=0$; jest to porządek
częściowy na klasach. Niech $\mathcal L$ będzie
niepustym łańcuchem w $\mathcal C$ i niech
$s=\sup\{\mu(B):[B]\in\mathcal L\}$.
Wybieramy klasy o miarach dążących do $s$,
a następnie wśród pierwszych $j$ klas wybieramy
największą w porządku łańcucha. Otrzymujemy
rosnący ciąg klas z przedstawicielami $B_j$
takimi, że $\mu(B_j)\uparrow s$.
Zastępując $B_j$ przez $\bigcup_{i\le j}B_i$,
otrzymujemy rzeczywiście rosnące zbiory bez zmiany
ich klas. Dla $B=\bigcup_jB_j$ mamy $\mu(B)=s\le t$.
Każda klasa $[C]$ z łańcucha o $\mu(C)<s$
jest mniejsza od pewnego $[B_j]$, zatem $[C]\le[B]$.
Jeśli $\mu(C)=s$, to $[B_j]\le[C]$ dla każdego $j$,
gdyż przy odwrotnej nierówności klasy są równe.
Stąd $\mu(B\setminus C)=0$, a z równości miar
także $\mu(C\setminus B)=0$. Zatem $[C]=[B]$.
Każdy niepusty łańcuch ma więc ograniczenie górne
$[B]$; dla pustego łańcucha jest nim $[\varnothing]$.
Z lematu Zorna istnieje element maksymalny $[B]$.
Wybieramy jego przedstawiciela $B\subseteq A$.
Jeśli $\mu(B)<t$, to $\mu(A\setminus B)>0$;
wybierając w $A\setminus B$ zbiór $C$ o mierze
$0<\mu(C)<t-\mu(B)$, otrzymalibyśmy większy
element $[B\cup C]\in\mathcal C$, sprzeczność.
Stąd $\mu(B)=t$. Przypadki $t=0$ i $t=\mu(A)$
są bezpośrednie. Dzieląc kolejno $A$ na części
miary $\eta$ aż do uzyskania reszty o mierze
co najwyżej $\eta$, dostajemy skończony podział.
\end{proof}

\begin{twierdzenie}[RNP a różniczkowalność krzywych Lipschitzowskich]
\label{thm:rnp-lipschitz-equiv-30}
Dla rzeczywistej lub zespolonej przestrzeni Banacha $X$
następujące warunki są równoważne:
\begin{enumerate}[label=(\roman*)]
 \item $X$ ma własność Radona--Nikodyma z
       Definicji~\ref{def:RNP-30};
 \item każda funkcja Lipschitzowska
       $F:[0,1]\to X$ jest różniczkowalna w normie $X$
       dla $\lambda_1$-prawie każdego $t\in(0,1)$.
\end{enumerate}
Wtedy każda taka $F$ spełnia wzór
$F(t)=F(0)+\int_{(0,t)}F'(s)\,\mathrm d\lambda_1(s)$.
\end{twierdzenie}
\begin{proof}
Por. \atlasref{FA-15B::thm:rnp-lipschitz15b}{FA-15B: charakteryzacja RNP}, gdzie uwzględniono również krzywe bezwzględnie ciągłe.

\emph{Kierunek (i)$\Rightarrow$(ii).}
Niech $F:[0,1]\to X$ ma stałą Lipschitza $L$.
Na algebrze podzbiorów $[0,1]$ złożonych ze
skończonych sum rozłącznych przedziałów $(s,t]$
oraz ewentualnego singletonu $\{0\}$ definiujemy
$\nu((s,t])=F(t)-F(s)$ i $\nu(\{0\})=0$. Definicja jest skończenie
addytywna, ponieważ przy dzieleniu przedziału przyrosty
się teleskopują. Z nierówności Lipschitza wynika
$\|\nu(A)\|\le L\lambda_1(A)$ dla takich zbiorów.
Każdy zbiór mierzalny $A\subseteq[0,1]$ można
przybliżyć w mierze Lebesgue'a skończonymi sumami
przedziałów półotwartych. Przybliżające wartości
$\nu(A_j)$ są ciągiem Cauchy'ego w $X$:
$\|\nu(A_j)-\nu(A_k)\|\le L\lambda_1(A_j\triangle A_k)$.
Ich granica nie zależy od aproksymacji i określa
rozszerzenie $\nu$ na $\sigma$-algebrę Lebesgue'a z
$\|\nu(A)\|\le L\lambda_1(A)$. Sprawdzamy najpierw skończoną addytywność.
Dla rozłącznych zbiorów mierzalnych $A,B$ wybieramy
aproksymacje $P_j,Q_j$ z algebry z błędami w mierze
dążącymi do zera. Zastępujemy $Q_j$ przez $Q_j\setminus P_j$.
Ponieważ
\[
 \lambda_1((Q_j\setminus P_j)\triangle B)
 \le\lambda_1(Q_j\triangle B)+\lambda_1(P_j\triangle A),
\]
rozłączne aproksymacje nadal zbiegają w mierze
do $A,B$, a ich suma przybliża $A\cup B$.
Addytywność na algebrze i przejście do granicy w normie
dają $\nu(A\cup B)=\nu(A)+\nu(B)$.
Dla rozłącznych $A_j$ mamy więc
$\nu(\bigcup_jA_j)-\sum_{j=1}^N\nu(A_j)
=\nu(\bigcup_{j>N}A_j)$. Norma prawej strony
jest ograniczona przez $L\lambda_1(\bigcup_{j>N}A_j)$,
co dąży do zera. Rozszerzenie jest przeliczalnie addytywne. Ma skończoną wariację:
$|\nu|(A)\le L\lambda_1(A)$.

Na mocy RNP istnieje $g\in L^1(\lambda_1;X)$ takie,
że $\nu(A)=\int_Ag\,\mathrm d\lambda_1$.
Dla $t\in[0,1]$ otrzymujemy
$F(t)=F(0)+\int_{(0,t)}g\,\mathrm d\lambda_1$.
Lemat~\ref{lem:bochner-differentiation-30} daje
$F'=g$ prawie wszędzie.

\emph{Kierunek (ii)$\Rightarrow$(i): redukcja do części bezatomowej.}
Niech $\nu:\mathcal A\to X$ ma skończoną wariację
oraz $\nu\ll\mu_0$ dla skończonej miary $\mu_0$.
Wybierzmy miarę $\mu=|\nu|$. Jest skończona i
spełnia $\|\nu(A)\|\le\mu(A)$.
Jeśli $\nu=g\mu$ dla pewnego $g\in L^1(\mu;X)$,
to na mocy Lematu~\ref{lem:vector-variation-ac-30}
mamy $\mu\ll\mu_0$. Skalarne twierdzenie
Radona--Nikodyma daje $w=\mathrm d\mu/\mathrm d\mu_0$;
wybieramy mierzalnych przedstawicieli $g,w$,
przy czym $w\ge0$ i $w$ jest skończona.
Na $\{w>0\}$ zbiory $\mu$-zerowe są $\mu_0$-zerowe,
bo $\mu(B)=\int_Bw\,\mathrm d\mu_0$. Na $\{w=0\}$
przyjmujemy $gw=0$. Mnożąc proste aproksymacje
$g$ i $w$ (te drugie zerowe na $\{w=0\}$), otrzymujemy
silną mierzalność $gw$ względem $\mu_0$.
Ponadto skalarny wzór zmiany miary daje
$\int\|gw\|\,\mathrm d\mu_0=\int\|g\|\,\mathrm d\mu<\infty$.
Dla prostych $s$ całkowalnych względem $\mu$ mamy
$\int_A sw\,\mathrm d\mu_0=\int_A s\,\mathrm d\mu$
z definicji $w$. Jeśli $s_j\to g$ w $L^1(\mu;X)$, to
$\|s_jw-gw\|_{L^1(\mu_0;X)}=\|s_j-g\|_{L^1(\mu;X)}\to0$.
Z ciągłości całki wynika, że $gw$ reprezentuje $\nu$
względem $\mu_0$. Wystarczy zatem zbudować
gęstość względem $\mu=|\nu|$.

Atomy dodatniej miary skończonej $\mu$ rozpatrujemy
modulo zbiory zerowe. Dwa atomy o przecięciu
dodatniej miary mają tę samą klasę: ich przecięcie
ma pełną miarę w każdym z nich. Dla każdego $k\ge1$
klas atomów miary większej niż $1/k$ jest skończenie
wiele, gdyż dowolna skończona rodzina różnych
takich klas ma przedstawicieli o przecięciach
miary zero, a suma ich miar nie przekracza
$\mu(\Omega)$. Każdy atom dodatniej miary należy
do jednej z tych rodzin. Wybieramy zatem co najwyżej
przeliczalnie wielu przedstawicieli $B_j$ wszystkich
klas atomów. Zbiory $A_j=B_j\setminus\bigcup_{i<j}B_i$
są parami rozłączne, mają te same klasy co $B_j$
i pozostają atomami.
Niech $\Omega_{\mathrm{at}}=\bigcup_j A_j$ oraz
$\Omega_{\mathrm{c}}=\Omega\setminus\Omega_{\mathrm{at}}$.
Ograniczenie $\mu$ do $\Omega_{\mathrm{c}}$ jest
bezatomowe: atom na tym zbiorze byłby również
atomem dla $\mu$, a więc należałby do jednej
z wybranych klas, co przeczy rozłączności.

Dla mierzalnego $C\subseteq A_j$ mamy $\mu(C)=0$
lub $\mu(A_j\setminus C)=0$. Z dominacji przez $\mu$
i addytywności $\nu$ wynika odpowiednio $\nu(C)=0$
lub $\nu(C)=\nu(A_j)$. Definiujemy
$g_{\mathrm{at}}=\sum_j\frac{\nu(A_j)}{\mu(A_j)}\chi_{A_j}$,
z wartością zero poza $\Omega_{\mathrm{at}}$.
Sumy częściowe zbiegają punktowo i w $L^1$, gdyż
$\sum_j\|\nu(A_j)\|\le\sum_j\mu(A_j)\le\mu(\Omega)$.
Dla każdego mierzalnego $C$ całka z $g_{\mathrm{at}}$
po $C$ jest równa $\sum_j\nu(C\cap A_j)
=\nu(C\cap\Omega_{\mathrm{at}})$ na mocy przeliczalnej
addytywności. Wystarczy więc znaleźć gęstość
$g_{\mathrm{c}}$ dla ograniczenia $\nu$ do
$\Omega_{\mathrm{c}}$ i przedłużyć ją przez zero;
suma $g_{\mathrm{at}}+g_{\mathrm{c}}$ reprezentuje $\nu$.
Jeśli $\mu(\Omega_{\mathrm{c}})=0$, wystarczy
$g_{\mathrm{c}}=0$. W pozostałym przypadku możemy
zatem zakładać, że $\mu$ jest bezatomowa
i $\mu(\Omega)=M>0$.

\emph{Przypuszczenie o braku gęstości i wybór podziałów.}
Załóżmy przeciwnie, że $\nu$ nie ma gęstości
względem $\mu$. Z Lematu~\ref{lem:rnp-partition-30}
wynika, że sieć funkcji $g_\pi$ nie jest Cauchy'ego.
Istnieje $\varepsilon>0$ takie, że dla każdego
podziału $\pi_0$ znajdują się jego rozdrobnienia
$\alpha,\beta$ z
$\|g_\alpha-g_\beta\|_{L^1}\ge\varepsilon$.
Niech $\rho$ będzie wspólnym rozdrobnieniem
$\alpha$ i $\beta$. Ponieważ uśrednienia są
kontrakcjami i $\mathbb E_\alpha g_\rho=g_\alpha$,
$\mathbb E_\beta g_\rho=g_\beta$,
$\mathbb E_\alpha g_{\pi_0}=g_{\pi_0}$ oraz
$\mathbb E_\beta g_{\pi_0}=g_{\pi_0}$,
otrzymujemy
$\varepsilon\le2\|g_\rho-g_{\pi_0}\|_{L^1}$.
Na mocy Lematu~\ref{lem:nonatomic-divisibility-30}
każdy mierzalny zbiór można rozciąć na skończenie
wiele części o arbitralnie małej mierze. Jeśli dalej rozdrobnimy
$\rho$ do $\pi$, to ponownie przez kontrakcję
$\|g_\pi-g_{\pi_0}\|_{L^1}
\ge\|g_\rho-g_{\pi_0}\|_{L^1}$.
Możemy więc zbudować zagnieżdżone podziały
$\pi_1\prec\pi_2\prec\cdots$ takie, że
\begin{equation}\label{eq:rnp-partition-jumps-30}
 \max_{A\in\pi_n}\mu(A)\longrightarrow0,
 \qquad
 \|g_{\pi_{n+1}}-g_{\pi_n}\|_{L^1}
 \ge\varepsilon/2\quad(n\ge1).
\end{equation}
Komórki o mierze zero zachowujemy w podziałach
przestrzeni $\Omega$, lecz pomijamy przy kodowaniu:
mają wartość $\nu(A)=0$ i nie wpływają na normy $L^1$.

\emph{Kodowanie podziałów jako przedziałów.}
Każdej komórce $A\in\pi_1$ o dodatniej mierze
przypisujemy przedział $I_A=(c,d]\subseteq(0,M]$
o długości $\mu(A)$, tak aby przedziały te
były rozłączne i w kolejności wypełniały $(0,M]$.
Punkt $0$ traktujemy osobno i przypisujemy mu
wartość zero w funkcjach prostych.
Jeśli komórka $A$ jest podzielona na
$A_1,\ldots,A_r$ przez $\pi_2$, dzielimy
$I_A$ na kolejne podprzedziały długości
$\mu(A_1),\ldots,\mu(A_r)$. Powtarzamy to
rekurencyjnie. W ten sposób dostajemy zagnieżdżone
podziały przedziału, których maksymalna długość
komórki zmierza do zera. Niech $h_n:[0,M]\to X$
będzie funkcją stałą na $I_A$ dla $A\in\pi_n$
o wartości $\nu(A)/\mu(A)$. Z konstrukcji
\[
 \|h_n\|\le1\ \text{prawie wszędzie},\qquad
 \int_{I_A}h_n(t)\,\mathrm d\lambda_1(t)=\nu(A).
\]
Ponadto
\[
 \|h_{n+1}-h_n\|_{L^1([0,M];X)}
 =\|g_{\pi_{n+1}}-g_{\pi_n}\|_{L^1(\mu;X)}.
\]
Istotnie, dla komórki $B\in\pi_{n+1}$ dodatniej
miary oznaczmy przez $A(B)$ zawierającą ją
komórkę $\pi_n$. Obie normy są równe sumie
\[
 \sum_{\substack{B\in\pi_{n+1}\\\mu(B)>0}}
 \mu(B)\left\|\frac{\nu(B)}{\mu(B)}
 -\frac{\nu(A(B))}{\mu(A(B))}\right\|.
\]
Po stronie przedziałów wynika to z długości $I_B$
i inkluzji $I_B\subseteq I_{A(B)}$, a po stronie
$\Omega$ z definicji funkcji $g_\pi$.

Niech $D$ będzie zbiorem końców wszystkich
przedziałów z wszystkich poziomów. Jest on
gęsty w $[0,M]$. Na $D$ definiujemy $F(0)=0$
i określamy $F(t)$ jako sumę $\nu(A)$ po
komórkach leżących na lewo od $t$, na dowolnym
poziomie, na którym $t$ jest końcem przedziału.
Definicja jest spójna: na późniejszym poziomie
każda komórka po lewej stronie danego końca
zastąpiona jest rodziną swoich podkomórek, których
wartości $\nu$ sumują się do wartości komórki
pierwotnej. Dla $s<t$ z $D$ wybieramy wspólny
poziom, na którym oba punkty są końcami przedziałów.
Jeśli $\mathcal P$ jest rodziną komórek tego poziomu
o $I_B\subseteq(s,t]$, to
\[
 \|F(t)-F(s)\|=\left\|\sum_{B\in\mathcal P}\nu(B)\right\|
 \le\sum_{B\in\mathcal P}\mu(B)=t-s.
\]
Dla dowolnego $t\in[0,M]$ wybieramy ciąg $t_j\in D$
zbieżny do $t$. Z powyższej nierówności $(F(t_j))$
jest ciągiem Cauchy'ego, więc ma granicę w $X$,
niezależną od wyboru ciągu. Przechodząc do granicy
w oszacowaniu Lipschitza, otrzymujemy jednoznaczne
$1$-Lipschitzowskie przedłużenie na $[0,M]$.
Dla każdej komórki $I_A=(c,d]$ zachodzi
\begin{equation}\label{eq:rnp-interval-average-30}
 F(d)-F(c)=\nu(A)
 =\int_{I_A}h_n(t)\,\mathrm d\lambda_1(t).
\end{equation}

\emph{Różniczkowanie wymusza sprzeczność.}
Warunek (ii), po liniowym przeskalowaniu $[0,M]$
do $[0,1]$, zapewnia istnienie $F'$ prawie wszędzie.
Lemat~\ref{lem:FTC-Banach-30} daje
$F(d)-F(c)=\int_{(c,d)}F'(t)\,\mathrm d\lambda_1(t)$.
Porównanie z \eqref{eq:rnp-interval-average-30}
pokazuje, że na każdej komórce $I_A$ funkcja $h_n$
jest średnią $F'$; inaczej mówiąc,
$h_n=\mathbb E_{\mathcal I_n}F'$,
gdzie $\mathcal I_n$ oznacza skończony podział
$[0,M]$ na opisane przedziały.
Ponieważ ich maksymalna długość dąży do zera,
$h_n\to F'$ w $L^1([0,M];X)$. Istotnie, dla ciągłej
$\phi:[0,M]\to X$ średnie po przedziałach różnią się
od $\phi$ o nie więcej niż jej moduł ciągłości
przy maksymalnej długości komórki. Dla ogólnej $F'\in L^1$ wybieramy ciągłą
$\phi$ z $\|F'-\phi\|_{L^1}<\delta$. Oznaczając
maksymalną długość komórki przez $m_n$ i moduł
ciągłości $\phi$ przez $\omega_\phi$, otrzymujemy
z kontrakcji operatorów uśredniających
\[
 \|h_n-F'\|_{L^1}
 \le 2\delta+M\omega_\phi(m_n).
\]
Najpierw $n\to\infty$, następnie $\delta\downarrow0$,
co daje żądaną zbieżność w $L^1$. W szczególności
$\|h_{n+1}-h_n\|_{L^1}\to0$.
Otrzymana zbieżność przeczy nierównościom
\eqref{eq:rnp-partition-jumps-30}. Zatem $\nu$
ma gęstość względem $|\nu|$, a stąd, jak wykazano
na początku, względem $\mu_0$.
\end{proof}

\section{Przykłady: przestrzenie z RNP i bez RNP}
\label{sec:rnp-examples-30}
\begin{twierdzenie}[Przestrzenie Hilberta mają RNP]
\label{thm:hilbert-rnp-30}
Każda rzeczywista lub zespolona przestrzeń Hilberta $H$
ma własność Radona--Nikodyma.
\end{twierdzenie}
\begin{proof}
Niech $(\Omega,\mathcal A,\mu_0)$ będzie przestrzenią
z miarą skończoną, a $\nu:\mathcal A\to H$ miarą
wektorową o skończonej wariacji, taką że $\nu\ll\mu_0$.
Kładziemy $\mu=|\nu|$. Dla podziału $\pi$ funkcja
$g_\pi$ z Lematu~\ref{lem:rnp-partition-30} spełnia
$\|g_\pi\|_{L^\infty}\le1$, zatem jest ograniczona
w przestrzeni Hilberta $L^2(\mu;H)$:
$\|g_\pi\|_{L^2}^2\le\mu(\Omega)$.
Istotnie, iloczyn skalarny w $L^2(\mu;H)$ definiujemy
przez całkę z iloczynu skalarnego w $H$ (przyjmujemy
liniowość w pierwszej zmiennej). Nierówność
Cauchy'ego--Schwarza zapewnia jego istnienie.
Zupełność wykazujemy tak jak dla $L^1$: z ciągu
Cauchy'ego wybieramy podciąg o sumowalnych normach
$L^2$ kolejnych różnic. Nierówność Minkowskiego
dla skończonych sum norm punktowych i lemat Fatou
dają, że suma tych norm należy do skalarnego $L^2$.
W szczególności podciąg jest punktowo Cauchy'ego
prawie wszędzie. Jego granica jest silnie mierzalna
jako granica prawie wszędzie silnie mierzalnych funkcji
(można wybrać proste przybliżenia każdej z nich
z błędem większym niż $1/n$ na zbiorze miary
co najwyżej $2^{-n}$, gdyż $\mu$ jest skończona).
Ponowne zastosowanie lematu Fatou do ogonów sumy
norm różnic daje zbieżność w $L^2$, a więc zbieżność
całego ciągu Cauchy'ego. Otrzymujemy przestrzeń Hilberta
bez założenia separowalności $H$.

Domknięta kula przestrzeni Hilberta jest słabo zwarta
(z refleksywności i twierdzenia Banacha--Alaoglu).
Zatem sieć $(g_\pi)$ w tej kuli ma słabo zbieżną
podsieć $(g_{\pi_\alpha})$. Nie zastępujemy jej
dowolnym podciągiem: potrzebujemy własności podsieci,
że dla każdego ustalonego podziału $\pi_0$
od pewnego miejsca $\pi_\alpha$ rozdrabnia $\pi_0$.
Niech $g\in L^2(\mu;H)$ będzie jej granicą.
Ponieważ $\mu$ jest skończona, nierówność
Cauchy'ego--Schwarza daje
$\|g\|_{L^1}\le\mu(\Omega)^{1/2}\|g\|_{L^2}$,
więc całki Bochnera poniżej istnieją.
Dla każdego $A\in\mathcal A$ i $h\in H$ funkcja
$\chi_A h$ należy do $L^2(\mu;H)$ i
$\|\chi_Ah\|_{L^2}=\mu(A)^{1/2}\|h\|$.
Z własności podsieci od pewnego miejsca
$\pi_\alpha$ rozdrabnia $\{A,\Omega\setminus A\}$
(z pominięciem pustych komórek), a wtedy
$\int_Ag_{\pi_\alpha}\,\mathrm d\mu=\nu(A)$.
Testując słabą zbieżność przez $\chi_Ah$
i przenosząc ciągły funkcjonał $x\mapsto\langle x,h\rangle_H$
przez całkę Bochnera, otrzymujemy
\[
 \left\langle\int_Ag\,\mathrm d\mu,h\right\rangle_H
 =\lim_\alpha\left\langle
 \int_Ag_{\pi_\alpha}\,\mathrm d\mu,h\right\rangle_H
 =\langle\nu(A),h\rangle_H.
\]
Dowolność $h$ dowodzi $\nu(A)=\int_Ag\,\mathrm d\mu$.
Dla pierwotnej miary dominującej $\mu_0$
przechodzimy z $\mu=|\nu|$ do $\mu_0$ dokładnie jak
w dowodzie Twierdzenia~\ref{thm:rnp-lipschitz-equiv-30}.
\end{proof}

\begin{przyklad}[Dlaczego przestrzeń $L^1$ nie ma RNP]
\label{ex:L1-no-RNP-30}
Niech $X=L^1([0,1],\lambda_1)$ oraz
$F:[0,1]\to X$ będzie określona przez
$F(t)=\chi_{[0,t]}$ (jako klasa funkcji w zmiennej $s$).
Dla $t>u$ mamy $\|F(t)-F(u)\|_{L^1}=t-u$,
a więc $F$ jest $1$-Lipschitzowska.
W żadnym punkcie $t\in(0,1)$ nie istnieje jednak
pochodna w normie $L^1$. Istotnie, dla $h>0$
iloraz różnicowy to
$h^{-1}\chi_{(t,t+h]}$, którego norma wynosi $1$.
Dla $0<h<1-t$ oznaczmy ten iloraz przez $q_h$.
Wówczas $q_{h/2}-q_h$ przyjmuje wartość $1/h$
na $(t,t+h/2]$, wartość $-1/h$ na $(t+h/2,t+h]$
i zero poza $(t,t+h]$. Zatem
\[
 \|q_{h/2}-q_h\|_{L^1}
 =\frac{h/2}{h}+\frac{h/2}{h}=1.
\]
Gdyby $q_h$ zbiegał w normie $L^1$ przy $h\downarrow0$,
to również $q_{h/2}$ zbiegałby do tej samej granicy,
a norma ich różnicy dążyłaby do zera. Sprzeczność.
Z Twierdzenia~\ref{thm:rnp-lipschitz-equiv-30}
wynika, że $L^1([0,1])$ nie ma RNP.

Równoważnie, miara wektorowa
$\nu(A)=\chi_A\in L^1([0,1])$ spełnia
$\|\nu(A)\|_{L^1}=\lambda_1(A)$ i
$|\nu|=\lambda_1$. Istotnie, dla rozłącznych
mierzalnych $A_j$ różnica między $\chi_{\bigcup_j A_j}$
a sumą pierwszych $N$ funkcji charakterystycznych
jest funkcją charakterystyczną $\bigcup_{j>N}A_j$,
której norma $L^1$ dąży do zera. Stąd $\nu$ jest
miarą wektorową. Dla każdego skończonego podziału
$A=\dot\bigcup_j A_j$ mamy
$\sum_j\|\nu(A_j)\|_{L^1}=\sum_j\lambda_1(A_j)=\lambda_1(A)$,
co daje równość wariacji z $\lambda_1$.
W szczególności $\nu$ ma skończoną wariację
i $\nu\ll\lambda_1$, ale \emph{nie} ma gęstości
Bochnera $g$ względem $\lambda_1$.
Gdyby miała, jej funkcja pierwotna $F(t)=\nu([0,t])$
musiałaby być różniczkowalna prawie wszędzie.
\end{przyklad}

\begin{przyklad}[Przestrzeń funkcji ciągłych bez RNP]
Niech $X=C([0,1])$ z normą supremum i
$F(t)(s)=|s-t|$ dla $s,t\in[0,1]$.
Wtedy $\|F(t)-F(u)\|_\infty\le|t-u|$.
Dla ustalonego $t\in(0,1)$ skalarne ilorazy różnicowe
przy $h\downarrow0$ dążą punktowo (dla $s\ne t$)
do funkcji równej $+1$ dla $s<t$ oraz $-1$ dla $s>t$.
Gdyby ilorazy różnicowe zbiegały w normie
supremum do $v\in C([0,1])$, musielibyśmy mieć
$v(s)=1$ dla $s<t$ i $v(s)=-1$ dla $s>t$.
Ciągłość w $t$ wymagałaby jednocześnie $v(t)=1$
i $v(t)=-1$, co jest niemożliwe.
Nie istnieje zatem pochodna $F'(t)$ w normie supremum,
a $C([0,1])$ nie ma RNP.
\end{przyklad}

\begin{uwaga}[Twierdzenie Rademachera a własność Radona--Nikodyma]
Twierdzenie~\ref{thm:rademacher-28} dowodzi
różniczkowalności odwzorowań Lipschitzowskich
$U\subseteq\mathbb R^n\to\mathbb R^m$ prawie
wszędzie, bez dodatkowych założeń. Dla ogólnej
przestrzeni Banacha $X$ nawet odwzorowanie
$[0,1]\to X$ może być nigdzie nieróżniczkowalne.
Twierdzenie~\ref{thm:rnp-lipschitz-equiv-30}
wykazuje, że dokładną przeszkodą jest brak RNP.
W szczególności nie wolno wnioskować o RNP
z samego istnienia skalarnej pochodnej
Radona--Nikodyma dla $x^*\nu$ przy każdym
$x^*\in X^*$: te skalarne gęstości nie muszą
składać się w silnie mierzalną gęstość $X$-wartościową.
\end{uwaga}

\section{Miary wektorowe w analizie i geometrii}
\label{sec:vector-applications-30}
\begin{przyklad}[Miara gradientowa funkcji $BV$]
Jeśli $u\in BV(\Omega)$, to z
Twierdzenia~\ref{thm:BV-measure-29} dystrybucyjna
pochodna $Du=(D_1u,\ldots,D_nu)$ jest skończoną
miarą wektorową o wartościach w $\mathbb R^n$.
Dla $u\in BV_{\mathrm{loc}}(\Omega)$ stosujemy
poniższe wnioski o gęstościach lokalnie, po ograniczeniu
do części o skończonej wariacji.
Z RNP tej przestrzeni i twierdzenia
Radona--Nikodyma otrzymujemy gęstość
$\sigma_u=\mathrm dDu/\mathrm d|Du|$.
Z Twierdzenia~\ref{thm:variation-density-30}
$|Du|(A)=\int_A\|\sigma_u\|\,\mathrm d|Du|$
dla każdego mierzalnego $A$. Jednoznaczność
skalarnej gęstości względem $|Du|$ daje
$\|\sigma_u\|=1$ $|Du|$-prawie wszędzie. Dla funkcji gładkiej
$Du=\nabla u\,\lambda_n$, natomiast dla funkcji
charakterystycznej zbioru o skończonym obwodzie
miara $D\chi_E$ może być osobliwa względem
$\lambda_n$ i skoncentrowana na brzegu miarowym.
To wyjaśnia, dlaczego naturalną miarą dominującą
jest $|Du|$, a nie zawsze $\lambda_n$.
\end{przyklad}

\begin{wniosek}[Przestrzenie Banacha i całki z krzywych]\label{idx:30:wniosek:15}

Niech $a,b\in\mathbb R$, $a<b$, a $X$ będzie
przestrzenią Banacha z własnością Radona--Nikodyma.
Jeśli $F:[a,b]\to X$ jest Lipschitzowska, to istnieje
silnie mierzalna $g\in L^\infty((a,b),\lambda_1;X)$
taka, że dla każdego $t\in[a,b]$ zachodzi
$F(t)-F(a)=\int_{(a,t)}g\,\mathrm d\lambda_1$
oraz $F'=g$ prawie wszędzie na $(a,b)$.
Ponadto $\|g\|_{L^\infty}\le\operatorname{Lip}(F)$.
\end{wniosek}
\begin{proof}
Funkcja $G:[0,1]\to X$ dana wzorem
$G(r)=F(a+(b-a)r)$ jest Lipschitzowska.
Twierdzenie~\ref{thm:rnp-lipschitz-equiv-30} zapewnia
istnienie $G'$ prawie wszędzie. Przy $t=a+(b-a)r$
ilorazy różnicowe dają $F'(t)=G'(r)/(b-a)$
w każdym punkcie istnienia $G'(r)$; afiniczna
zmiana zmiennej przenosi zbiory zerowe na zerowe.
Lemat~\ref{lem:FTC-Banach-30} zastosowany do $F$
daje przedstawiciela $g$ oraz wzór całkowy.
Z oszacowania ilorazów różnicowych przez
$\operatorname{Lip}(F)$ wynika żądane ograniczenie normy.
\end{proof}

\section{Granice zakresu i dalsze kierunki}
Teoria miar wektorowych ma również zagadnienia,
których nie potrzebowaliśmy do dowodu RNP:
całkę Pettisa, miary o skończonej półwariacji,
operatory słabo zwarte oraz twierdzenie
Bartle'a--Dunforda--Schwartza.
Są one naturalną kontynuacją w analizie funkcjonalnej.
W tym rozdziale udowodniliśmy od podstaw
charakteryzację RNP przez różniczkowalność krzywych
Lipschitzowskich. Nie zakładaliśmy ani twierdzenia
o dentowalności zbiorów wypukłych, ani ogólnego
twierdzenia o zbieżności martyngałów Banachowskich.

% Rozdział współdzielony tematycznie z geometrią różniczkową; wersja audytowana 2026-11-19.
% Samodzielny rozdział do włączenia przez \input{rozdzial_dzety.tex}.
% Wymaga: amsmath, amssymb, amsthm, mathtools, tikz, hyperref.
% Środowiska: definicja, twierdzenie, lemat, wniosek, uwaga, przyklad, proof.
\chapter{Dżety: od geometrii różniczkowej do teorii miary}
\label{ch:dzety}

\section{Motywacja: ile informacji niesie punkt?}
W rozdziale~\ref{ch:area-coarea-28} wyprowadziliśmy twierdzenie
Rademachera i wzory miarowo-geometryczne. Teraz interpretujemy
różniczkowalność prawie wszędzie językiem dżetów, zachowując też
pełne ujęcie geometryczne, formy kontaktowe i twierdzenia Whitneya.
Wartość $f(p)$ opisuje jedynie punkt obrazu, pochodna $Df(p)$ ---
przybliżenie liniowe, a kolejne pochodne --- coraz dokładniejsze
przybliżenia lokalne. \emph{Dżet rzędu $k$} jest sposobem zebrania
tej informacji w jeden obiekt. Rozróżniamy trzy sytuacje:
\begin{enumerate}
\item w geometrii gładkiej dżet jest klasą równoważności
kiełków odwzorowań lub przekrojów;
\item w równaniach różniczkowych współrzędne dżetowe
traktuje się jako formalnie niezależne, lecz rozwiązania
muszą spełniać warunki holonomiczności;
\item w teorii miary dżet może istnieć tylko prawie wszędzie,
a zwykłą granicę zastępuje się \emph{granicą przybliżoną}.
\end{enumerate}
Te trzy punkty widzenia spotykają się w twierdzeniu Whitneya:
\emph{kiedy formalnie zadane pochodne dają się zrealizować przez
prawdziwą funkcję gładką?} Wersja Whitneya--Luzina mówi zaś,
że wiele funkcji mierzalnych staje się gładkimi po pominięciu
zbioru o dowolnie małej mierze.

\section{Kiełki, klasy równoważności i wielomiany Taylora}
\begin{definicja}[Kiełek odwzorowania]
\label{def:dzety-kielek}
Niech $M,N$ będą gładkimi rozmaitościami i $p\in M$.
Dwa gładkie odwzorowania $f:U\to N$ i $g:V\to N$,
gdzie $p\in U\cap V$, mają ten sam \emph{kiełek w $p$},
jeżeli pokrywają się na pewnym otoczeniu $p$ zawartym w $U\cap V$.
Zbiór kiełków oznaczamy przez $\mathcal C^\infty_p(M,N)$.
\end{definicja}

\begin{definicja}[$k$-dżet odwzorowania]
\label{def:dzety-klasa}
Niech $M,N$ będą gładkimi rozmaitościami, $p\in M$,
a $k\in\mathbb N_0$. Na zbiorze kiełków
$\mathcal C^\infty_p(M,N)$ definiujemy relację:
kiełki o różnych wartościach w $p$ nie są równoważne.
Jeżeli $f(p)=g(p)=q\in N$, piszemy $f\sim_{p,k}g$, gdy
po wybraniu mapy $x$ wokół $p$ i mapy $y$ wokół $q$ zachodzi
\[
 \partial^\alpha (y\circ f\circ x^{-1})(x(p))
 =\partial^\alpha (y\circ g\circ x^{-1})(x(p)),
 \qquad |\alpha|\le k.
\]
Dziedziny przedstawicieli ograniczamy tak, aby obrazy
leżały w mapie celu. Zgodność pochodnych nie zależy
od przedstawicieli kiełków, gdyż są oni równi w pewnym
otoczeniu $p$. Niezależność od map wykażemy poniżej.
Klasę równoważności oznaczamy przez $j_p^k f$ i nazywamy
\emph{$k$-dżetem $f$ w $p$}. Dla $k=0$ jest to po prostu para
$(p,f(p))$. Wartości pochodnych liczymy składowo.
\end{definicja}

\begin{lemat}[Niezależność od map]
\label{lem:dzety-mapy}
Relacja w Definicji~\ref{def:dzety-klasa} nie zależy od wyboru
map źródła i celu.
\end{lemat}
\begin{proof}
Niech $\varphi,\psi$ będą zmianami współrzędnych.
Dla odwzorowania współrzędnego $F$ nowy zapis ma postać
$\psi\circ F\circ\varphi^{-1}$. Reguła łańcuchowa mówi,
że jego pochodna rzędu $r\le k$ jest skończoną sumą wyrażeń
zawierających pochodne $F$ rzędów od $1$ do $r$, pochodne
$\psi,\varphi^{-1}$ i wartości funkcji w rozpatrywanym punkcie.
Jeżeli odpowiednie dane dla $F$ i $G$ są równe do rzędu $k$,
to równe są także wszystkie nowe pochodne. Zmieniając mapy
w przeciwnym kierunku, otrzymujemy równoważność.
\end{proof}

\begin{przyklad}[Różne funkcje, ten sam dżet]
\label{ex:dzety-taylor}
Dla $f(x)=\sin x$ i $g(x)=x$ mamy $j_0^2f=j_0^2g$,
ale $j_0^3f\ne j_0^3g$. Natomiast $f(x)=x+x^2$
i $g(x)=x-x^2$ mają identyczne $1$-dżety w zerze,
ale różne $2$-dżety.
\end{przyklad}

\begin{uwaga}[Nie należy utożsamiać drugiej pochodnej z tensorem]
\label{rem:dzety-hess}
Pierwsza pochodna funkcji skalarnych definiuje kanoniczny
kowektor $\mathrm df_p\in T_p^*M$. Dla drugiej pochodnej jest inaczej.
Współrzędna $x$ na $\mathbb R$ i współrzędna $y=x+x^2$ wokół zera
dałyby dla funkcji $f(x)=x$ odpowiednio $f''(0)=0$
oraz $\frac{\mathrm d^2 f}{\mathrm dy^2}(0)=-2$.
Oznacza to, że $2$-dżet jest obiektem niezmienniczym,
ale same współrzędne drugich pochodnych nie są tensorami.
Tensor Hessego powstaje kanonicznie dopiero po określeniu
koneksji (lub w punkcie krytycznym funkcji).
\end{uwaga}

\begin{definicja}[Algebra dżetów funkcji skalarnych]
\label{def:dzety-algebra}
Niech $\mathcal E_p$ oznacza pierścień kiełków gładkich
funkcji rzeczywistych na $M$ w punkcie $p$ i
$\mathfrak m_p=\{[f]:f(p)=0\}$ jego ideał maksymalny.
Definiujemy $\mathcal J_p^k=\mathcal E_p/\mathfrak m_p^{k+1}$.
\end{definicja}

\begin{twierdzenie}[Algebra i wielomiany Taylora]
\label{thm:dzety-algebra}
Dla $M=\mathbb R^n$ i $p=0$ odwzorowanie
\[
 [f]\longmapsto \sum_{|\alpha|\le k}
 \frac{\partial^\alpha f(0)}{\alpha!}X^\alpha
\]
indukuje izomorfizm algebr
$\mathcal J_0^k\cong\mathbb R[X_1,\ldots,X_n]/(X_1,\ldots,X_n)^{k+1}$.
W szczególności $\dim\mathcal J_0^k=\binom{n+k}{k}$.
\end{twierdzenie}
\begin{proof}
Wielomian Taylora rzędu $k$ iloczynu jest iloczynem
wielomianów Taylora po odrzuceniu wyrazów stopnia większego niż $k$,
co wynika z reguły Leibniza. Jądro odwzorowania składa się
z kiełków o pochodnych rzędów $0,\ldots,k$ równych zeru.
Wybieramy przedstawiciela $f$ w kuli wokół zera.
Dla ustalonego $x$ stosujemy jednowymiarowy wzór
Taylora z resztą całkową do $t\mapsto f(tx)$.
Ponieważ pochodne do rzędu $k$ w zerze znikają,
\[
 f(x)=\sum_{|\alpha|=k+1}x^\alpha
 \frac{k+1}{\alpha!}\int_0^1(1-t)^k
 \partial^\alpha f(tx)\,\mathrm d\lambda_1(t).
\]
Współczynniki przy $x^\alpha$ są gładkie:
na każdej mniejszej domkniętej kuli wszystkie
pochodne całkowanej funkcji względem $x$ są
ciągłe i ograniczone, więc wolno różniczkować
pod całką. Każdy $x^\alpha$ jest iloczynem $k+1$
kiełków z $\mathfrak m_0$, co daje inkluzję
jądra w $\mathfrak m_0^{k+1}$.
Odwrotnie, w każdej pochodnej rzędu co najwyżej $k$
iloczynu $k+1$ kiełków z $\mathfrak m_0$ reguła
Leibniza pozostawia co najmniej jeden czynnik
niezróżniczkowany, który znika w zerze.
To samo dotyczy skończonych sum takich iloczynów.
Zatem jądrem jest dokładnie $\mathfrak m_0^{k+1}$.
Każdy wielomian stopnia co najwyżej $k$ sam jest gładki,
więc homomorfizm jest surjektywny. Liczba jednomianów
o $|\alpha|\le k$ wynosi $\binom{n+k}{k}$.
\end{proof}

\begin{przyklad}[Liczby dualne]
\label{ex:dzety-dualne}
Dla $n=k=1$ dostajemy
$\mathcal J_0^1\cong\mathbb R[\varepsilon]/(\varepsilon^2)$.
Dżet funkcji $f$ ma zapis $f(0)+f'(0)\varepsilon$;
reguła mnożenia automatycznie daje wzór na pochodną iloczynu.
To algebraiczne źródło arytmetyki liczb dualnych używanej
w automatycznym różniczkowaniu.
\end{przyklad}

\section{Wiązki dżetów i ich naturalne odwzorowania}
\begin{definicja}[Wiązka dżetów odwzorowań]
\label{def:dzety-bundle}
Dla gładkich rozmaitości $M,N$ oznaczamy przez $J^k(M,N)$
zbiór wszystkich $k$-dżetów lokalnych gładkich odwzorowań
$M\to N$. Naturalne odwzorowania to
$\pi_0^k(j_p^k f)=(p,f(p))$ oraz dla $0\le r\le k$
$\pi_r^k(j_p^k f)=j_p^r f$.
Dla ustalonego $f$ odwzorowanie $j^kf:p\mapsto j_p^kf$
nazywamy \emph{przedłużeniem dżetowym}.
\end{definicja}

W lokalnych współrzędnych na $M$ i $N$, o wymiarach $n,m$,
przestrzeń $J^k(M,N)$ ma współrzędne $(x^i,u^a,u^a_\alpha)$,
$1\le|\alpha|\le k$, gdzie $u^a_\alpha$ odpowiada pochodnej
$\partial^\alpha u^a$. Jej wymiar wynosi
$n+m\binom{n+k}{k}$. Każdy zestaw współrzędnych nad ustalonymi punktami
map jest realizowany przez wielomian Taylora
o zadanych współczynnikach, po ograniczeniu
otoczenia punktu źródłowego tak, aby obraz leżał
w mapie celu. Współrzędne są więc bijekcjami
na otwarte podzbiory przestrzeni euklidesowej.
Zmiany współrzędnych są gładkie na mocy reguły
łańcuchowej i mają gładkie odwrotności.
Topologia określona przez te mapy jest Hausdorffa:
dżety o różnych punktach bazowych w $M\times N$
rozdzielamy przez przeciwobrazy rozłącznych otoczeń,
a różne dżety nad tym samym punktem rozdzielamy
w jednej wspólnej mapie dżetowej.
Przeliczalne atlasy $M,N$ dają przeliczalną rodzinę
map dżetowych, każdą z przeliczalną bazą; stąd drugi
aksjomat przeliczalności. Otrzymujemy rozmaitość gładką.

\begin{twierdzenie}[Pierwsze dżety i struktura afiniczna]
\label{thm:dzety-affine}
Dla dowolnych $(p,q)\in M\times N$ włókno
$(\pi_0^1)^{-1}(p,q)$ utożsamia się kanonicznie z
$\operatorname{Hom}(T_pM,T_qN)$. Dla $k\ge2$ włókna
$\pi_{k-1}^k:J^k(M,N)\to J^{k-1}(M,N)$ są przestrzeniami
afinicznymi modelowanymi na
$S^kT_p^*M\otimes T_qN$, gdzie $(p,q)$ jest
wartością rozpatrywanego $(k-1)$-dżetu.
\end{twierdzenie}
\begin{proof}
Wartość i różniczka $f$ w $p$ jednoznacznie wyznaczają
$1$-dżet, a dowolne odwzorowanie liniowe
$T_pM\to T_qN$ jest różniczką odpowiedniej lokalnej funkcji
w mapach. Dla $k\ge2$ dwa przedłużenia tego samego
$(k-1)$-dżetu różnią się tylko pochodną rzędu $k$.
W mapach oznaczmy różnicę najwyższych pochodnych
przez $B=D^kF-D^kG$. Po zmianie map $\varphi,\psi$
różnica przyjmuje postać
\[
 \widetilde B(v_1,\ldots,v_k)
 =D\psi\,B(D\varphi^{-1}v_1,\ldots,D\varphi^{-1}v_k),
\]
gdzie różniczki liczymy w odpowiednich punktach
bazowych. Jedyny wyraz z pochodną rzędu $k$
powstaje, gdy wszystkie $k$ różniczkowań padają
na $F$ z pierwszymi pochodnymi zmian map.
Pozostałe składniki zależą tylko od niższych
pochodnych, więc skracają się przy odejmowaniu.
Otrzymujemy regułę transformacji elementu
$S^kT_p^*M\otimes T_qN$. Dodanie dowolnego takiego
tensora do najwyższej pochodnej daje nowy dżet,
realizowany lokalnie przez wielomian Taylora.
Działanie jest wolne i przechodnie, co dowodzi
struktury afinicznej bez wyboru wyróżnionego zera.
\end{proof}

\begin{uwaga}
Nie należy nazywać całej $J^k(M,N)$ wiązką wektorową
nad $M\times N$ dla $k\ge2$: zmiany współrzędnych
zawierają wyrazy niższych rzędów. Różnica między
\emph{wiązka wektorowa} a \emph{wiązka afiniczna}
jest w tym miejscu istotna.
\end{uwaga}

\section{Dżety przekrojów, holonomiczność i formy kontaktowe}
Niech $\pi:E\to M$ będzie gładką wiązką lokalnie trywialną.
Dżety lokalnych przekrojów $s$ tworzą $J^kE$.
W lokalnej trywializacji współrzędne mają postać
$(x^i,u^a,u^a_\alpha)$ z $1\le|\alpha|\le k$.
Oznaczamy też $u^a_0=u^a$. Przekroje klasy $C^k$
mają dżety w tej samej przestrzeni, ponieważ ich
skończone dane pochodnych w ustalonym punkcie
można zrealizować przez gładki lokalny przekrój
zadany wielomianem w mapie włókna.
Każde $j^ks$ jest przekrojem $J^kE\to M$.

\begin{definicja}[Przekrój holonomiczny]
\label{def:dzety-holonomic}
Przekrój ciągły $\sigma:U\to J^kE$ nazywamy
\emph{holonomicznym}, jeżeli istnieje przekrój
$s:U\to E$ klasy $C^k$ taki, że $\sigma=j^ks$.
Gdy ponadto $\sigma$ jest klasy $C^1$, to $s$ jest
klasy $C^{k+1}$; wynika to kolejno ze współrzędnych
$u^a_\alpha\circ\sigma$, które są pochodnymi $s$.
\end{definicja}

Na $J^kE$ występują \emph{formy kontaktowe}
\[
 \theta^a_\alpha=\mathrm d u^a_\alpha
 -\sum_{i=1}^n u^a_{\alpha+e_i}\,\mathrm d x^i,
 \qquad |\alpha|\le k-1.
\]
Dla $\alpha=0$ przyjmujemy $u_0^a=u^a$.

\begin{twierdzenie}[Kryterium holonomiczności]
\label{thm:dzety-contact}
Niech $U\subseteq M$ będzie otwarty, $k\ge1$,
a $\sigma:U\to J^kE$ będzie przekrojem klasy $C^1$.
Wtedy $\sigma=j^ks$ dla pewnego przekroju
$s:U\to E$ klasy $C^{k+1}$ wtedy i tylko wtedy, gdy
$\sigma^*\theta^a_\alpha=0$ dla wszystkich form
kontaktowych, w dowolnych lokalnych trywializacjach.
\end{twierdzenie}
\begin{proof}
Dla $\sigma=j^ks$ zachodzi
$u^a_\alpha\circ\sigma=\partial^\alpha s^a$,
więc $\sigma^*\theta^a_\alpha=0$ wynika z różniczkowania.
Odwrotnie, niech $s=\pi_0^k\circ\sigma$, tj. składowa
zerowego rzędu. Równości $\sigma^*\theta^a_0=0$ dają
$\partial_i s^a=u^a_{e_i}\circ\sigma$. Indukcyjnie,
$\sigma^*\theta^a_\alpha=0$ implikuje
$\partial_i(u^a_\alpha\circ\sigma)
=u^a_{\alpha+e_i}\circ\sigma$. Otrzymujemy
$u^a_\alpha\circ\sigma=\partial^\alpha s^a$
dla $|\alpha|\le k$. Ponieważ wszystkie współrzędne $u^a_\alpha\circ\sigma$
są klasy $C^1$, otrzymujemy ponadto $s\in C^{k+1}$.
W szczególności regularność $C^1$ pola dżetów nie jest
tożsama z regularnością $C^1$ przekroju wyjściowego.
Konstrukcja jest zgodna na nakładaniu się map.
\end{proof}

\begin{przyklad}[Formalne dane nie muszą być pochodnymi]
W trywialnej wiązce $\mathbb R\times\mathbb R\to\mathbb R$
przekrój $\sigma(x)=(x,u(x),p(x))=(x,0,1)$
przestrzeni $J^1E$ ma zerową składową wartości i
stałą składową pochodnej równą $1$. Nie pochodzi od
żadnej funkcji $u$, ponieważ
$\sigma^*(\mathrm du-p\,\mathrm dx)=-\mathrm dx\ne0$.
\end{przyklad}

\begin{center}
\begin{tikzpicture}[>=stealth, every node/.style={font=\small}]
 \node (jk) at (0,1) {$J^kE$};
 \node (jm) at (2.2,1) {$J^{k-1}E$};
 \node (je) at (3.8,1) {$\cdots$};
 \node (e) at (5.2,1) {$E$};
 \node (m) at (6.6,1) {$M$};
 \node (u) at (0,-.25) {$U\subset M$};
 \draw[->] (jk) -- node[above] {$\pi_{k-1}^k$} (jm);
 \draw[->] (jm) -- (je);
 \draw[->] (je) -- (e);
 \draw[->] (e) -- node[above] {$\pi$} (m);
 \draw[->] (u) -- node[left] {$j^k s$} (jk);
\end{tikzpicture}

{\small Rzutowania w wieży dżetów odrzucają pochodne najwyższego rzędu.}
\end{center}

\section{Równania różniczkowe jako podzbiory przestrzeni dżetów}
\begin{definicja}[Równanie różniczkowe rzędu $k$]
\label{def:dzety-pde}
Niech $\pi:E\to M$ będzie gładką wiązką.
Układ równań różniczkowych rzędu nie większego niż $k$
możemy opisać \emph{lokalnie} zbiorem zer
$\mathcal R\cap V=\{z\in V:F(z)=0\}$, gdzie
$V\subseteq J^kE$ jest otwarty, a
$F:V\to\mathbb R^r$ jest gładkie. Nie każda globalna
podrozmaitość równań musi być zbiorem zer jednego
globalnie określonego odwzorowania do $\mathbb R^r$.
Przekrój $s$ klasy co najmniej $C^k$ jest rozwiązaniem,
gdy $j^ks(x)\in\mathcal R$ dla każdego $x$ z jego
dziedziny określenia.
\end{definicja}

\begin{przyklad}[Hamilton--Jacobi i równanie Laplace'a]
Dla rzeczywistej funkcji $u$ na $\mathbb R^n$
równanie Hamiltona--Jacobiego
$H(x,\nabla u(x))=0$ jest podzbiorem $J^1E$
opisanym przez $H(x,p)=0$. Równanie Laplace'a
$\sum_i\partial_{ii}u=0$ jest hiperpłaszczyzną
w każdym włóknie $J^2E\to J^1E$ o równaniu
$\sum_i u_{2e_i}=0$. W obu przypadkach sam wybór
punktu z $\mathcal R$ nie wystarcza: rozwiązanie
musi ponadto wyznaczać przekrój holonomiczny.
\end{przyklad}

\begin{definicja}[Symbol główny operatora liniowego]
\label{def:dzety-symbol}
Dla skalarnego liniowego operatora
$Lu=\sum_{|\alpha|\le k}a_\alpha(x)\partial^\alpha u$
o gładkich współczynnikach, na otwartym $U\subset\mathbb R^n$,
jego symbol główny jest wielomianem
$\sigma_k(L)(x,\xi)=\sum_{|\alpha|=k}a_\alpha(x)\xi^\alpha$.
\end{definicja}

\begin{uwaga}
Symbol główny jest zależny wyłącznie od najwyższej części
dżetu; pod zmianą map jego współczynniki przekształcają się
zgodnie z tensorową częścią struktury afinicznej
Twierdzenia~\ref{thm:dzety-affine}. Dla Laplasjanu
$\sigma_2(\Delta)(x,\xi)=|\xi|^2$.
\end{uwaga}

\begin{przyklad}[Równanie Eulera--Lagrange'a]
Niech $\Omega\subset\mathbb R^n$ będzie ograniczonym obszarem,
$u\in C^2(\Omega)$, a $L(x,z,p)$ będzie funkcją klasy $C^2$.
Zakładamy, że $L(x,u(x),\nabla u(x))$ jest całkowalna na $\Omega$.
Funkcjonał zależący od $1$-dżetu ma postać
\[
 \mathcal A(u)=\int_\Omega L(x,u(x),\nabla u(x))
 \,\mathrm d\lambda_n(x).
\]
Dla $\varphi\in C_c^\infty(\Omega)$ i zróżniczkowanego
pod całką wyrażenia (na zwartym nośniku $\varphi$) mamy
\[
 \frac{\mathrm d}{\mathrm dt}\mathcal A(u+t\varphi)\big|_{t=0}
 =\int_\Omega\left(L_z\varphi+
 \sum_i L_{p_i}\partial_i\varphi\right)\mathrm d\lambda_n(x).
\]
Jeżeli $u$ jest punktem stacjonarnym, czyli powyższa
pierwsza wariacja znika dla wszystkich
$\varphi\in C_c^\infty(\Omega)$, to po całkowaniu
przez części i zastosowaniu lematu podstawowego rachunku
wariacyjnego otrzymujemy równanie
$L_z-\sum_i\partial_i L_{p_i}=0$ w $\Omega$.
Jest ono zazwyczaj równaniem drugiego rzędu, a zatem
jego geometria znajduje się w $J^2E$, mimo że
sam funkcjonał zależy tylko od $J^1E$.
\end{przyklad}

\section{Od gładkości do miary: gęstość i granice przybliżone}
Odtąd pracujemy na $\mathbb R^n$. Przez $\lambda_n$ oznaczamy
miarę Lebesgue'a, a przez $B(x,r)$ kulę otwartą.

\begin{definicja}[Punkt gęstości i granica przybliżona]
\label{def:dzety-ap-limit}
Dla zbioru mierzalnego $A\subset\mathbb R^n$ punkt $x$
ma gęstość $1$ względem $A$, jeżeli
\[
 \lim_{r\downarrow0}
 \frac{\lambda_n(A\cap B(x,r))}{\lambda_n(B(x,r))}=1.
\]
Jeżeli $v$ jest funkcją mierzalną w otoczeniu $x$,
to piszemy $\operatorname{ap}\!\lim_{y\to x}v(y)=L$,
gdy dla każdego $\varepsilon>0$ zbiór
$\{y:|v(y)-L|>\varepsilon\}$ ma w $x$ gęstość $0$.
W przypadku funkcji określonej tylko na zbiorze mierzalnym
$A$ zakładamy, że $x$ jest punktem gęstości $1$ zbioru $A$
i liczymy gęstość zbiorów wyjątkowych w $A\cap B(x,r)$.
\end{definicja}

\begin{definicja}[Przybliżony $k$-dżet]
\label{def:dzety-approx-jet}
Niech $A\subset\mathbb R^n$ będzie mierzalny,
$x\in A$ punktem jego gęstości $1$, a
$f:A\to\mathbb R^m$ funkcją mierzalną.
Mówimy, że $f$ ma \emph{przybliżony dżet Peana rzędu $k\ge1$}
w $x$, jeżeli istnieje wielomian wektorowy
$P_x(h)=\sum_{|\alpha|\le k}a_\alpha(x)h^\alpha/\alpha!$,
przy $a_0(x)=f(x)$, taki że
\[
 \operatorname{ap}\!\lim_{\substack{y\to x\\y\in A}}
 \frac{f(y)-P_x(y-x)}{|y-x|^k}=0.
\]
Dla $k=1$ współczynniki liniowe wyznaczają
\emph{różniczkę przybliżoną} $D_{\mathrm{ap}}f(x)$.
W tym sformułowaniu wartość $f(x)$ jest częścią danych:
zmiana wartości funkcji w samym punkcie $x$ może więc
zniszczyć istnienie dżetu w tym punkcie. Natomiast
\emph{prawie wszędzie} pojęcie to zależy wyłącznie od
klasy funkcji modulo zbiory miary zero: dwa równe prawie
wszędzie przedstawiciele mają te same dżety w punktach
gęstości zbioru, na którym są równe.
\end{definicja}

\begin{lemat}[Jednoznaczność przybliżonego dżetu]
\label{lem:dzety-approx-unique}
Jeżeli $x$ jest punktem gęstości $1$ zbioru $A$,
to wielomian $P_x$ spełniający
Definicję~\ref{def:dzety-approx-jet} jest jednoznaczny.
\end{lemat}
\begin{proof}
Niech $Q=P_x-\widetilde P_x$ będzie różnicą dwóch takich
wielomianów. Wtedy $Q(0)=0$ i
$\operatorname{ap}\!\lim_{h\to0}Q(h)/|h|^k=0$.
Gdyby $Q\ne0$, jego pierwszy niezerowy jednorodny składnik
$Q_j$, $j\le k$, byłby niezerowy na pewnym otwartym
podzbiorze sfery jednostkowej. Na stożku o dodatnim kącie
bryłowym zachodziłoby
$|Q(h)|/|h|^k\ge c|h|^{j-k}$ dla dostatecznie małych $h$.
Na mniejszym otwartym fragmencie sfery $|Q_j|\ge2c>0$,
a wyższe składniki po podzieleniu przez $|h|^j$
są jednostajnie $O(|h|)$. Stąd dla dostatecznie
małych $|h|\le1$ iloraz jest nie mniejszy niż $c$.
Stożek $C$ ma względną miarę $\gamma>0$ w każdej
kuli wokół zera. Ponieważ $A-x$ ma gęstość $1$,
\[
 \liminf_{r\downarrow0}
 \frac{\lambda_n(C\cap(A-x)\cap B(0,r))}
 {\lambda_n(B(0,r))}\ge\gamma.
\]
Istotnie, od względnej miary stożka odejmujemy
co najwyżej względną miarę dopełnienia $A-x$,
która dąży do zera. Z nierówności trójkąta
zbiór, na którym $|Q(h)|/|h|^k>c/2$ w $A-x$,
zawiera się w sumie dwóch zbiorów wyjątkowych
dla obu wielomianów z progiem $c/4$.
Ich gęstości są zerowe, co daje sprzeczność.
\end{proof}

\begin{przyklad}[Funkcja nieciągła z przybliżoną pochodną]
\label{ex:dzety-approx}
Niech $A=\bigcup_{j\ge1}(2^{-j},2^{-j}+2^{-3j})$
i $f(x)=\chi_A(x)$, gdzie $\chi_A$ oznacza funkcję
charakterystyczną zbioru $A$ i $f(0)=0$. Wtedy $A$ ma
w zerze gęstość $0$, zatem $f$ ma w zerze przybliżoną
granicę równą $0$. Ponadto
$\lambda_1(A\cap(-r,r))=O(r^3)$ dla $r\downarrow0$,
a więc również $D_{\mathrm{ap}}f(0)=0$.
Zwykła ciągłość (a tym bardziej różniczkowalność) nie zachodzi,
ponieważ w każdym otoczeniu zera są punkty $f=1$.
\end{przyklad}

\begin{przyklad}[Dlaczego trzeba ustalić przedstawiciela?]
Niech $f=\chi_{\mathbb Q}$ na $\mathbb R$. W punktach
niewymiernych funkcja ma przybliżony dżet każdego rzędu
równy zeru, ponieważ zbiór $\mathbb Q$ ma miarę zero.
W punktach wymiernych $f(x)=1$, więc nie ma nawet
przybliżonej granicy równej własnej wartości $f(x)$.
Przedstawiciel zerowy tej samej klasy w $L^0$ ma natomiast
dżet zerowy w każdym punkcie. Rozróżnienie to nie ma
znaczenia dla twierdzeń prawie wszędzie.
\end{przyklad}

\begin{uwaga}[Gdzie tkwi istota mierzalności?]
Równość $D_{\mathrm{ap}}f(x)=A$ nie oznacza
zbieżności po \emph{wszystkich} punktach $y\to x$.
Wystarcza odrzucanie zbioru o gęstości zerowej,
który może być gęsty topologicznie. Dlatego dżety
przybliżone nie są dżetami kiełków w sensie
Definicji~\ref{def:dzety-klasa}.
\end{uwaga}

\section{Rademacher, Sobolew i $BV$}
\begin{twierdzenie}[Rademacher i mierzalny pierwszy dżet]
\label{thm:dzety-rademacher}
Niech $U\subset\mathbb R^n$ będzie otwarty, a
$f:U\to\mathbb R^m$ lokalnie Lipschitzowskie.
Wtedy $f$ jest różniczkowalne w $\lambda_n$-prawie każdym
punkcie $U$. W tych punktach istnieje przybliżony
$1$-dżet i $D_{\mathrm{ap}}f(x)=Df(x)$.
Ponadto $x\mapsto Df(x)$ ma mierzalny przedstawiciel.
\end{twierdzenie}
\begin{proof}
Istnienie zwykłej różniczki prawie wszędzie jest treścią
twierdzenia Rademachera. Zwykła granica implikuje
przybliżoną, więc równość dżetów jest natychmiastowa.
Zbiór różniczkowalności jest borelowski. Istotnie,
na zbiorze istnienia wszystkich pochodnych cząstkowych
ich wartości są granicami ilorazów różnicowych z
przyrostami $1/j$, a więc są funkcjami borelowskimi.
Warunek
$|f(x+h)-f(x)-A(x)h|\le |h|/m$
dla wszystkich dostatecznie małych niezerowych $h$
można, dzięki ciągłości $f$, sprawdzać na przeliczalnym
zbiorze przyrostów wymiernych. Kwantyfikatory
$\forall m\in\mathbb N\,\exists j\in\mathbb N\,
\forall h\in\mathbb Q^n\setminus\{0\}$ z $|h|<1/j$
i $x+h\in U$ opisują zatem borelowski podzbiór $S$,
który jest dokładnie zbiorem różniczkowalności.
Wynika to z ciągłości wyrażenia względem przyrostu $h$
poza $h=0$. Przedłużenie macierzy
$Df$ przez zero na dopełnienie tego zbioru jest mierzalne.
\end{proof}

\begin{przyklad}[Punkt o dwóch geometrycznych dżetach]
Funkcja $f(x)=|x|$ na $\mathbb R$ ma zwykłą różniczkę
$Df(x)=\operatorname{sgn}(x)$ dla $x\ne0$.
W zerze nie ma ani zwykłego, ani przybliżonego
$1$-dżetu: dodatnia i ujemna półprosta mają dodatnią
gęstość i dawałyby sprzeczne pochodne. Brak dżetu w jednym
punkcie nie narusza twierdzenia Rademachera.
\end{przyklad}

\begin{twierdzenie}[Dżet przybliżony funkcji Sobolewa]
\label{thm:dzety-sobolev}
Niech $U\subset\mathbb R^n$ będzie otwarty,
$1\le p\le\infty$, a $u\in W_{\mathrm{loc}}^{1,p}(U)$.
Wówczas istnieje przedstawiciel $u$, który ma
przybliżony $1$-dżet w $\lambda_n$-prawie każdym
punkcie $U$, przy czym
$D_{\mathrm{ap}}u(x)=\nabla u(x)$ prawie wszędzie.
Po prawej stronie występuje gradient słaby.
\end{twierdzenie}
\begin{proof}
Korzystamy z następującego twierdzenia o precyzyjnych
przedstawicielach Sobolewa: funkcja klasy $W^{1,1}_{\rm loc}$
jest przybliżenie różniczkowalna prawie wszędzie, a jej
różniczka przybliżona pokrywa się z gradientem słabym.
Standardowy dowód tego wyniku opiera się na różniczkowaniu
całek i przybliżeniu funkcji Sobolewa przez funkcje gładkie.
Dla $p>1$ można go uzyskać także z oszacowań funkcji
maksymalnej Hardy'ego--Littlewooda; przypadek $p=1$
wymaga wersji słabego typu $(1,1)$. Ponieważ
$W^{1,p}_{\rm loc}\subset W^{1,1}_{\rm loc}$,
otrzymujemy tezę. Podany rezultat o przedstawicielach
Sobolewa jest w tym dowodzie \emph{jawnie użytym
wcześniejszym twierdzeniem}, a nie wynikiem wyprowadzonym
z samej definicji dżetu.
\end{proof}

\begin{uwaga}[Funkcje $BV$ i punkty skoku]
\label{rem:dzety-bv}
Dla $u\in BV_{\rm loc}(U)$ istnieje przybliżona
różniczka w prawie każdym punkcie względem $\lambda_n$,
lecz pochodna dystrybucyjna jest miarą wektorową, która
może mieć część osobliwą. W pełnym twierdzeniu
strukturalnym dla $BV$ rozkłada się ją na
\[
 Du=\nabla u\,\lambda_n
 +(u^+-u^-)\nu_u\,\mathcal H^{n-1}\llcorner J_u+D^cu,
\]
gdzie $J_u$ jest przybliżonym zbiorem skoków,
$u^\pm$ jednostronnymi śladami, $\nu_u$ normalną,
a $D^cu$ częścią Cantora. Jest to odrębny, głębszy wynik:
nie należy wyprowadzać go wyłącznie z istnienia
przybliżonego $1$-dżetu prawie wszędzie.
\end{uwaga}

\begin{przyklad}[Skok a przybliżony dżet]
Dla $E=\{x\in\mathbb R^n:x_1>0\}$ funkcja
$u=\chi_E$ jest lokalnie stała poza hiperpłaszczyzną
$x_1=0$, więc $D_{\mathrm{ap}}u=0$ prawie wszędzie.
Jednak $Du=e_1\,\mathcal H^{n-1}\llcorner\{x_1=0\}$
(z orientacją odpowiadającą przejściu z $0$ do $1$).
Zatem przybliżony dżet pierwszego rzędu nie opisuje
całej miary pochodnej funkcji $BV$.
\end{przyklad}

\section{Jakobiany i wzory area/coarea odczytywane z dżetów}
Niech $f:U\subseteq\mathbb R^k\to\mathbb R^N$,
$k\le N$, będzie lokalnie Lipschitzowskie.
W punkcie różniczkowalności określamy
$J_kf(x)=\sqrt{\det(Df(x)^{\mathsf T}Df(x))}$.
Jest to funkcja samego \emph{liniowego składnika $1$-dżetu}.
Dla mierzalnego $A\subset U$ i nieujemnej
mierzalnej funkcji $h:U\to[0,\infty]$ wzór area ma postać
\[
 \int_A h(x)J_kf(x)\,\mathrm d\lambda_k(x)
 =\int_{\mathbb R^N}\sum_{x\in A\cap f^{-1}(\{y\})}
 h(x)\,\mathrm d\mathcal H^k(y).
\]
Dla lokalnie Lipschitzowskiego $F:V\subset\mathbb R^n\to
\mathbb R^m$, $n\ge m$, gdzie $V$ jest otwarty,
i nieujemnej mierzalnej $g:V\to[0,\infty]$, wzór
koobszaru brzmi
\[
 \int_V g(x)\sqrt{\det(DF(x)DF(x)^{\mathsf T})}
 \,\mathrm d\lambda_n(x)
 =\int_{\mathbb R^m}\left(
 \int_{F^{-1}(\{y\})}g(x)\,\mathrm d\mathcal H^{n-m}(x)
 \right)\mathrm d\lambda_m(y).
\]
W obu wzorach miary Hausdorffa są znormalizowane tak,
że $\mathcal H^j|_{\mathbb R^j}=\lambda_j$ dla całkowitych
$j\ge1$ oraz $\mathcal H^0$ jest miarą liczącą.
Dla zbiorów mierzalnych w sensie Lebesgue'a konieczne
są właściwe sformułowania twierdzeń, uwzględniające
ukończenie miar i mierzalność funkcji krotnościowej.
W szczególności sama teoria dżetów nie rozwiązuje
kwestii mierzalności włókien i całek po nich.

\begin{uwaga}
Teoria dżetów wyjaśnia, \emph{jakie dane} wyznaczają
lokalny współczynnik geometryczny. Sam wzór area/coarea
wymaga jednak niezależnie twierdzeń o mierzalności,
różniczkowaniu i krotnościach włókien. Nie jest więc
czysto formalną konsekwencją teorii dżetów.
\end{uwaga}

\section{Pola dżetów i twierdzenie Whitneya}
\begin{definicja}[Pole dżetów Whitneya klasy $C^k$]
\label{def:dzety-whitney-field}
Niech $K\subset\mathbb R^n$ będzie domknięty,
$k\in\mathbb N_0$, a dla każdego $|\alpha|\le k$
niech $a_\alpha:K\to\mathbb R$ będzie funkcją.
Kładziemy $P_x(z)=\sum_{|\alpha|\le k}
 a_\alpha(x)(z-x)^\alpha/\alpha!$.
Rodzinę $(P_x)_{x\in K}$ nazywamy
\emph{polem dżetów Whitneya klasy $C^k$}, jeżeli
wszystkie $a_\alpha$ są ciągłe i dla każdego zwartego
$K_0\subset K$ oraz każdego $|\alpha|\le k$ zachodzi
\[
 \lim_{\delta\downarrow0}
 \sup_{\substack{x,y\in K_0,\;0<|x-y|<\delta}}
 \frac{|D^\alpha P_x(y)-a_\alpha(y)|}
 {|x-y|^{k-|\alpha|}}=0,
\]
przy umowie, że supremum po zbiorze pustym wynosi $0$.
\end{definicja}

\begin{uwaga}
Dla $|\alpha|=k$ mianownik jest równy $1$:
warunek mówi o ciągłości najwyższych pochodnych.
Dla $k=1$ warunki na zwartym $K$ sprowadzają się do
ciągłości $a_0,v=(a_{e_1},\ldots,a_{e_n})$ i
$|a_0(y)-a_0(x)-v(x)\cdot(y-x)|=o(|x-y|)$
jednostajnie dla $x,y\in K$.
\end{uwaga}

\begin{lemat}[Kostki Whitneya i rozkład jedności]
\label{lem:dzety-whitney-cubes}
Niech $K\subset\mathbb R^n$ będzie niepustym zbiorem
domkniętym, różnym od całej przestrzeni, a
$\Omega=\mathbb R^n\setminus K$. Istnieje przeliczalna
rodzina domkniętych kostek diadycznych $Q\subset\Omega$,
o parami rozłącznych wnętrzach, których suma jest
równa $\Omega$, taka że
$\ell(Q)\asymp\operatorname{dist}(Q,K)$. Istnieje
również gładki rozkład jedności $(\varphi_Q)$ na $\Omega$
o nośnikach w nieznacznie powiększonych kostkach $Q^*$,
takich że
$|D^\alpha\varphi_Q|\le C_\alpha\ell(Q)^{-|\alpha|}$.
W każdym punkcie co najwyżej $N(n)$ funkcji
$\varphi_Q$ jest niezerowych.
\end{lemat}
\begin{proof}
\emph{Wybór kostek.}
Przez kostkę diadyczną rozumiemy domknięcie zbioru
$2^{-j}(m+[0,1)^n)$, gdzie $j\in\mathbb Z$,
$m\in\mathbb Z^n$. Wybieramy maksymalne względem
inkluzji kostki spełniające
\[
 8\sqrt n\,\ell(Q)\le\operatorname{dist}(Q,K).
\]
Dla $z\in\Omega$ niech $d=\operatorname{dist}(z,K)>0$.
Każda dostatecznie mała kostka zawierająca $z$
spełnia warunek wyboru, ponieważ jej odległość od $K$
jest nie mniejsza niż $d-\sqrt n\,\ell(Q)$.
Wśród kolejnych kostek rodzicielskich warunek
w końcu przestaje zachodzić: każda zawiera $z$,
więc jej odległość od $K$ jest nie większa niż $d$,
a długości boków rosną nieograniczenie.
Ostatnia dopuszczalna kostka jest maksymalna.
Dwie kostki diadyczne o przecinających się wnętrzach
są zagnieżdżone, więc różne kostki maksymalne
mają rozłączne wnętrza. Każda leży w $\Omega$,
a wszystkie razem pokrywają $\Omega$.

Jeśli $Q'$ jest kostką rodzicielską dla wybranej $Q$,
to $\ell(Q')=2\ell(Q)$, a z maksymalności
$\operatorname{dist}(Q',K)<16\sqrt n\,\ell(Q)$.
Każdy punkt $Q'$ jest odległy od $Q$ co najwyżej
o średnicę $Q'$. Stąd
\[
 8\sqrt n\,\ell(Q)\le\operatorname{dist}(Q,K)
 <18\sqrt n\,\ell(Q).
\]

\emph{Powiększenia i krotność pokrycia.}
Niech $Q^*$ będzie otwartą kostką o tym samym środku,
o boku $(5/4)\ell(Q)$. Dla $z\in\overline{Q^*}$
odległość od $Q$ jest co najwyżej
$\sqrt n\,\ell(Q)/8$. Z dolnego oszacowania
odległości $Q$ od $K$ oraz nierówności trójkąta mamy
$\operatorname{dist}(z,K)\ge(8-1/8)\sqrt n\,\ell(Q)$.
Z drugiej strony, dla dowolnego $q\in Q$
mamy $|z-q|\le(9/8)\sqrt n\,\ell(Q)$.
Biorąc kres dolny odległości $|q-b|$, $b\in K$,
uzyskujemy wygodne oszacowania
\[
 7\sqrt n\,\ell(Q)\le\operatorname{dist}(z,K)
 \le20\sqrt n\,\ell(Q)\qquad(z\in\overline{Q^*}).
\]
W szczególności $\overline{Q^*}\subset\Omega$.
Jeżeli powiększenia dwóch kostek zawierają ten
sam punkt, stosunek ich boków mieści się między
$7/20$ a $20/7$.

Ustalmy $z\in\Omega$ i $d=\operatorname{dist}(z,K)$.
Wszystkie kostki, których powiększenia zawierają $z$,
mają boki co najmniej $d/(20\sqrt n)$ i co najwyżej
$d/(7\sqrt n)$. Każda z nich leży w $B(z,d)$,
gdyż jej punkty są odległe od $z$ najwyżej
o $(9/8)\sqrt n\,\ell(Q)<d$.
Z rozłączności wnętrz, dla dowolnej skończonej
podrodziny liczącej $N$ kostek otrzymujemy
\[
 N\left(\frac{d}{20\sqrt n}\right)^n
 \le\lambda_n(B(z,d))=\lambda_n(B(0,1))d^n.
\]
Daje to ograniczenie liczności całej rodziny
przez stałą zależną tylko od $n$.
Ponadto pokrycie powiększeniami jest lokalnie
skończone: jeżeli zwarty $L\subset\Omega$ jest
niepusty, to odległość jego punktów od $K$
ma dodatnie minimum i skończone maksimum.
Każde $\overline{Q^*}$ przecinające $L$ ma więc
bok ograniczony z dołu i z góry, a środek leżący
w pewnym ograniczonym zbiorze. Rozłączne wnętrza
i dodatnie dolne ograniczenie objętości kostek
dają skończenie wiele takich kostek.

\emph{Rozkład jedności.}
Wybieramy ustaloną gładką funkcję obcięcia
$0\le\eta\le1$, równą $1$ na kostce jednostkowej
o środku zero, o zwartym nośniku w jej
powiększeniu o czynnik $5/4$. Dla środka $c_Q$
definiujemy $\eta_Q(z)=\eta((z-c_Q)/\ell(Q))$.
Wtedy
$|D^\alpha\eta_Q|\le C_\alpha\ell(Q)^{-|\alpha|}$.
Suma $S=\sum_Q\eta_Q$ jest lokalnie skończona
i gładka, a pokrycie przez $Q$ oraz ograniczona
krotność dają $1\le S\le N(n)$.
Kładziemy $\varphi_Q=\eta_Q/S$.
Są to nieujemne funkcje gładkie o sumie $1$,
z nośnikami zawartymi w $Q^*$.

Na $Q^*$ pochodne $S$ są ograniczone przez
$C_\alpha'\ell(Q)^{-|\alpha|}$, ponieważ w każdym
punkcie jedynie ograniczona liczba składników
może mieć tam niezerowe pochodne, a ich boki
są porównywalne z $\ell(Q)$ (stosujemy oszacowania
również na domknięciach powiększeń).
Z tożsamości $S\cdot S^{-1}=1$ i reguły Leibniza,
dla $|\alpha|>0$,
\[
 D^\alpha(S^{-1})
 =-S^{-1}\sum_{0<\beta\le\alpha}
 \binom{\alpha}{\beta}D^\beta S\,
 D^{\alpha-\beta}(S^{-1}).
\]
Indukcja względem $|\alpha|$, zaczynając od
$|S^{-1}|\le1$, daje na $Q^*$
$|D^\alpha(S^{-1})|\le C_\alpha''\ell(Q)^{-|\alpha|}$.
Reguła Leibniza dla $\eta_Q S^{-1}$ daje żądane
oszacowanie pochodnych $\varphi_Q$; poza $Q^*$
funkcja i wszystkie jej pochodne znikają.
\end{proof}

\begin{lemat}[Sklejanie ciągłych pochodnych przez zbiór domknięty]
\label{lem:dzety-glue}
Niech $K\subset\mathbb R^n$ będzie domknięty. Załóżmy, że
$H:\mathbb R^n\to\mathbb R$ jest ciągła, klasy $C^1$ poza $K$,
a $V:\mathbb R^n\to\mathbb R^n$ jest ciągła i spełnia
$V=\nabla H$ poza $K$. Jeżeli dla każdego $a\in K$
zachodzi $H(a+h)=H(a)+V(a)\cdot h+o(|h|)$, to
$H\in C^1(\mathbb R^n)$ i $\nabla H=V$ wszędzie.
\end{lemat}
\begin{proof}
Poza $K$ teza jest znana. W punkcie $a\in K$ podany
warunek jest definicją różniczkowalności i daje
$\nabla H(a)=V(a)$. Ponieważ $V$ jest ciągła na całej
przestrzeni, otrzymany gradient jest ciągły, czyli
$H\in C^1$. Nie trzeba różniczkować po krzywych
zawartych w $K$.
\end{proof}

\begin{twierdzenie}[Twierdzenie Whitneya o przedłużaniu]
\label{thm:dzety-whitney-extension}
Niech $K\subset\mathbb R^n$ będzie domknięty,
$k\in\mathbb N_0$, a $(a_\alpha)_{|\alpha|\le k}$
polem dżetów Whitneya klasy $C^k$ na $K$.
Wtedy istnieje $F\in C^k(\mathbb R^n)$ spełniająca
$D^\alpha F(x)=a_\alpha(x)$ dla wszystkich
$x\in K$ i $|\alpha|\le k$.
Warunki Definicji~\ref{def:dzety-whitney-field} są
zarazem konieczne dla takich przedłużeń.
\end{twierdzenie}

\begin{proof}
Konieczność: dla $F\in C^k(\mathbb R^n)$ kładziemy
$a_\alpha=D^\alpha F|_K$. Ustalmy zwarty $K_0\subset K$
i $|\alpha|\le k$, niech $m=k-|\alpha|$.
Dla $m=0$ żądane oszacowanie wynika z jednostajnej
ciągłości $D^\alpha F$ na $K_0$.
Dla $m\ge1$ stosujemy wzór Taylora z resztą całkową
do funkcji $t\mapsto D^\alpha F(x+t(y-x))$,
odejmując od reszty wyraz stopnia $m$ w punkcie $x$.
Dostajemy
\[
 a_\alpha(y)-D^\alpha P_x(y)
 =\sum_{|\gamma|=m}\frac{m(y-x)^\gamma}{\gamma!}
 \int_0^1(1-t)^{m-1}
 [D^{\alpha+\gamma}F(x+t(y-x))-D^{\alpha+\gamma}F(x)]
 \,\mathrm d\lambda_1(t).
\]
Dla $|x-y|\le1$ wszystkie punkty odcinka należą
do zwartego zbioru $\{z:\operatorname{dist}(z,K_0)\le1\}$.
Moduł ciągłości pochodnych rzędu $k$ na tym zbiorze
ogranicza więc iloraz reszty przez $|y-x|^m$
jednostajnie przez wielokrotność funkcji dążącej do zera.
Przedstawimy konstrukcję wystarczalności; dla $K=\varnothing$
można przyjąć $F=0$.

\emph{Krok 1: kostki Whitneya.}
Jeżeli $K=\mathbb R^n$, przyjmujemy $F=a_0$.
Dla $|\beta|<k$ warunki Whitneya w otoczeniu
dowolnego $x$ dają
$a_\beta(x+h)=a_\beta(x)+\sum_i a_{\beta+e_i}(x)h_i+o(|h|)$:
wyrazy stopni co najmniej $2$ są $O(|h|^2)$,
a reszta jest $o(|h|^{k-|\beta|})$.
Zatem $a_\beta$ jest różniczkowalna i jej pochodne
pierwszego rzędu są równe odpowiednim ciągłym
$a_{\beta+e_i}$. Indukcyjnie $a_0\in C^k$
i $D^\beta a_0=a_\beta$. Dla $k=0$ wystarcza
ciągłość $a_0$.
W przeciwnym przypadku Lemat~\ref{lem:dzety-whitney-cubes}
dostarcza rozkładu zbioru
$\Omega=\mathbb R^n\setminus K$
na kostki $Q$, wraz z rozkładem jedności
$(\varphi_Q)$ i oszacowaniami ich pochodnych.

\emph{Krok 2: sklejanie wielomianów.}
Dla każdej kostki wybierzmy $b_Q\in K$ takie, że
$\operatorname{dist}(b_Q,Q)=\operatorname{dist}(Q,K)$, i połóżmy
$F(z)=\sum_Q\varphi_Q(z)P_{b_Q}(z)$ dla $z\in\Omega$,
a na $K$ niech $F(x)=a_0(x)$. Suma jest lokalnie skończona
i funkcja jest gładka poza $K$.

\emph{Krok 3: oszacowanie w pobliżu brzegu.}
Ustalmy $x\in K$. Wszystkie punkty $a,b\in K$
potrzebne w oszacowaniach dla $z$ dostatecznie
bliskiego $x$ leżą w zwartym $K\cap\overline{B(x,1)}$,
jak wynika z poniższych oszacowań odległości.
Na tym zbiorze niech $\omega(t)$ będzie maksimum
po $|\beta|\le k$ supremów ilorazów z definicji
Whitneya dla $0<|a-b|\le t$, z wartością zero
dla pustej rodziny. Dla dostatecznie małych $t>0$
jest skończona, niemalejąca i dąży do zera
przy $t\downarrow0$. Dla $a,b$ z tego zbioru,
$0<|a-b|\le t$, mamy
\[
 |D^\beta P_a(b)-a_\beta(b)|
 \le\omega(t)|a-b|^{k-|\beta|},\qquad |\beta|\le k.
\]
Dla $a=b$ różnica wielomianów jest zerowa.
Dla $a\ne b$ rozwijamy wielomian $D^\beta(P_a-P_b)$
dokładnie w punkcie $b$:
\[
 D^\beta(P_a-P_b)(z)
 =\sum_{|\gamma|\le k-|\beta|}
 \frac{D^{\beta+\gamma}P_a(b)-a_{\beta+\gamma}(b)}
 {\gamma!}(z-b)^\gamma.
\]
Jeżeli $|z-a|+|z-b|\le Ct$, to $|a-b|\le Ct$.
Każdy składnik sumy jest ograniczony w module
przez stałą razy
$\omega(Ct)t^{k-|\beta|-|\gamma|}t^{|\gamma|}$.
Skończenie wiele składników daje
$|D^\beta(P_a-P_b)(z)|\le C'\omega(Ct)t^{k-|\beta|}$.
Niech teraz $z\in\Omega$, $r=\operatorname{dist}(z,K)$,
a $a\in K$ jest punktem najbliższym $z$, więc
$|z-a|=r$ (odległość od niepustego zbioru domkniętego
w $\mathbb R^n$ jest osiągana). Dla każdej kostki
$Q$ z $z\in\overline{Q^*}$ mamy $\ell(Q)\asymp r$,
a z wyboru $b_Q$ i oszacowań kostek
$|b_Q-z|\le C\ell(Q)\le C'r$. W konsekwencji
$|b_Q-a|\le C''r$. Ponieważ $r\le|z-x|$,
punkty $a,b_Q$ należą do ustalonego zwartego
fragmentu $K$ dla dostatecznie małego $|z-x|$.

Traktując $a$ jako ustalony punkt podczas
różniczkowania, z $\sum_Q\varphi_Q=1$ mamy
$F-P_a=\sum_Q\varphi_Q(P_{b_Q}-P_a)$ na $\Omega$.
Reguła Leibniza daje w punkcie $z$
\[
 D^\beta(F-P_a)(z)
 =\sum_{Q:\ z\in\overline{Q^*}}\sum_{\gamma\le\beta}
 \binom{\beta}{\gamma}
 D^\gamma\varphi_Q(z)\,
 D^{\beta-\gamma}(P_{b_Q}-P_a)(z).
\]
Każdy iloczyn jest ograniczony w module przez
$C r^{-|\gamma|}\omega(Cr)r^{k-|\beta|+|\gamma|}$.
Liczba kostek jest ograniczona przez stałą
wymiarową, a liczba wielowskaźników przez stałą
zależną od $n,k$. Ostatecznie dla $|\beta|\le k$
\[
 |D^\beta F(z)-D^\beta P_a(z)|
 \le C\omega(Cr)r^{k-|\beta|}.
\]
W szczególności pochodne do rzędu $k$ mają przy $K$
granice zadane przez współczynniki $a_\beta$.

\emph{Krok 4: regularność na $K$.}
Dla każdego $|\beta|\le k$ definiujemy na całej
przestrzeni funkcję $H_\beta$ równą $D^\beta F$
na $\Omega$ i $a_\beta$ na $K$. W szczególności
$H_0=F$. Ustalmy $x\in K$ i oznaczmy $d=|z-x|>0$.
Dla $z\in\Omega$ wybieramy $a$ jak wyżej;
$r\le d$ i $|a-x|\le2d$. Wobec tego
\[
 \begin{split}
 |H_\beta(z)-D^\beta P_x(z)|
 &\le |D^\beta F(z)-D^\beta P_a(z)|\\
 &\quad+|D^\beta(P_a-P_x)(z)|\\
 &\le C\omega(Cr)r^{k-|\beta|}
       +C'\omega(C'd)d^{k-|\beta|}\\
 &\le C''\omega(C''d)d^{k-|\beta|}.
 \end{split}
\]
Ostatnia nierówność wykorzystuje niemalejący
charakter $\omega$ i nieujemność $k-|\beta|$.
Dla $z\in K$ to samo oszacowanie wynika
bezpośrednio z warunku Whitneya dla $x,z$.
Zatem dla dowolnego kierunku zbliżania $z\to x$
mamy $H_\beta(z)-D^\beta P_x(z)=o(d^{k-|\beta|})$.
Dla $|\beta|=k$ interpretujemy $o(d^0)$ jako $o(1)$.
Ponieważ $D^\beta P_x(x)=a_\beta(x)$,
wszystkie $H_\beta$ są ciągłe w $x$; poza $K$
są gładkie. Dla $k=0$ kończy to dowód.

Jeżeli $|\beta|<k$, rozwijamy wielomian w punkcie $x$:
\[
 D^\beta P_x(x+h)
 =a_\beta(x)+\sum_i a_{\beta+e_i}(x)h_i+O(|h|^2).
\]
W połączeniu z poprzednim oszacowaniem, ponieważ
$k-|\beta|\ge1$, otrzymujemy
$H_\beta(x+h)=H_\beta(x)+\sum_i H_{\beta+e_i}(x)h_i+o(|h|)$.
Poza $K$ mamy już $\partial_i H_\beta=H_{\beta+e_i}$.
Wektor $(H_{\beta+e_i})_i$ jest ciągły na całej przestrzeni. Stosując
Lemat~\ref{lem:dzety-glue} kolejno do pochodnych
rzędów $0,1,\ldots,k-1$, wykazujemy
$F\in C^k(\mathbb R^n)$ oraz $D^\alpha F|_K=a_\alpha$.
W ten sposób przejście przez $K$ nie wymaga
nieuzasadnionego różniczkowania funkcji na samym zbiorze
$K$, który może mieć puste wnętrze.
\end{proof}

\begin{przyklad}[Dlaczego sama ciągłość dżetów nie wystarcza?]
Niech $K=\{0\}\cup\{1/j:j\ge1\}\subset\mathbb R$,
$a_0(0)=0$, $a_0(1/j)=1/j$ i $a_1\equiv0$.
Obie funkcje $a_0,a_1$ są ciągłe na $K$, ale
$|a_0(1/j)-a_0(0)-a_1(0)/j|/(1/j)=1$.
Nie istnieje więc funkcja $C^1$ o tych wartościach i
pochodnych. Jest to konkretne naruszenie warunku Whitneya.
\end{przyklad}

\section{Twierdzenie Whitneya--Luzina: od dżetów przybliżonych do $C^1$}
\begin{lemat}[Mierzalność różniczki przybliżonej]
\label{lem:dzety-ap-meas}
Niech $U\subset\mathbb R^n$ będzie otwarty,
a $f:U\to\mathbb R^m$ mierzalna. Na mierzalnym
zbiorze punktów, w których $f$ jest przybliżenie
różniczkowalna, odwzorowanie $D_{\mathrm{ap}}f$
ma mierzalny przedstawiciel.
\end{lemat}
\begin{proof}
Zmienimy wartości na zbiorze miary zero tylko wtedy,
gdy jest to potrzebne do zdefiniowania funkcji na całym $U$.
Dla $M\in\mathbb N$ i $r>0$ rozważmy macierz
\[
 T_{M,r}(x)=\frac{n}{\lambda_n(B(0,r))}
 \int_{B(0,r)}\frac{f(x+h)-f(x)}{|h|}
 \frac{h^{\mathsf T}}{|h|}
 \chi_{\{|f(x+h)-f(x)|\le M|h|\}}
 \,\mathrm d\lambda_n(h),
\]
gdzie dla $x+h\notin U$ całkowaną funkcję przyjmujemy za zero.
Jest to odwzorowanie mierzalne, ponieważ całkowana funkcja
jest mierzalna i ograniczona przez $M$.
Jeżeli $f$ ma różniczkę przybliżoną $A$ w $x$ i
$M>\|A\|+1$, to na kulach malejących do zera
iloraz różnicowy jest w mierze bliski
$A(h/|h|)$; obcięcie nie usuwa prawie wszystkich
kierunków. Ponieważ całkowana funkcja jest ograniczona,
własność zbieżności w mierze daje
$\lim_{r\downarrow0}T_{M,r}(x)=A$.
Wykorzystaliśmy tożsamość
$\lambda_n(B(0,r))^{-1}\int_{B(0,r)}
 hh^{\mathsf T}/|h|^2\,\mathrm d\lambda_n(h)=I_n/n$.
Dla ciągu $r=1/j$ granice macierzy $T_{M,1/j}$ są
mierzalne na zbiorach swojego istnienia. W każdym punkcie
przybliżonej różniczkowalności wszystkie dostatecznie
duże wartości $M$ dają tę samą granicę $A$.
Wybór pierwszego stabilnego ogona tych granic jest
operacją na przeliczalnej rodzinie funkcji mierzalnych,
a więc daje mierzalny przedstawiciel $A$.
\end{proof}

\begin{twierdzenie}[Whitney--Luzin dla pierwszego dżetu]
\label{thm:dzety-whitney-lusin}
Niech $U\subset\mathbb R^n$ będzie otwarty,
a $f:U\to\mathbb R^m$ funkcją mierzalną,
przybliżenie różniczkowalną w $\lambda_n$-prawie każdym
punkcie $U$. Dla każdego zbioru mierzalnego $A\subset U$
o skończonej mierze i każdego $\varepsilon>0$ istnieją
zbiór zwarty $K\subset A$ i odwzorowanie
$F\in C^1(\mathbb R^n,\mathbb R^m)$ takie, że
\[
 \lambda_n(A\setminus K)<\varepsilon,
 \qquad F(x)=f(x),\quad DF(x)=D_{\mathrm{ap}}f(x)
 \quad (x\in K).
\]
Odwrotnie, jeżeli dla każdego mierzalnego $A\subset U$
o skończonej mierze i każdego $\varepsilon>0$ istnieją
$K,F$ spełniające powyższą równość wartości (bez wymogu
dotyczącego pochodnych), to $f$ jest przybliżenie
różniczkowalna prawie wszędzie na $U$.
\end{twierdzenie}

\begin{proof}
Wystarczy dowodzić dla funkcji rzeczywistej: dla wektorowych
stosujemy ten sam wybór zwartego zbioru i wspólne
oszacowania normy w $\mathbb R^m$.

\emph{Krok 1: duży zbiór z ciągłymi danymi.}
Niech $D$ będzie pełnomiarowym podzbiorem $U$, gdzie
istnieje różniczka przybliżona $G(x)=D_{\mathrm{ap}}f(x)$.
Dla sformułowania „prawie wszędzie” wybieramy mierzalny
zbiór $D$ o pełnej mierze, na którym dżet istnieje;
na nim działa Lemat~\ref{lem:dzety-ap-meas}.
Funkcje $f,G$ są więc mierzalne na $D$. Z twierdzenia Luzina i regularności
miary Lebesgue'a wybieramy zwarty $K_0\subset A\cap D$,
$\lambda_n(A\setminus K_0)<\varepsilon/3$, tak że
$f|_{K_0}$ i $G|_{K_0}$ są ciągłe.
W szczególności $G|_{K_0}$ jest jednostajnie ciągła.
Ponieważ $K_0\Subset U$, ma dodatnią odległość od
$\mathbb R^n\setminus U$.

\emph{Krok 2: jednostajna kontrola zbiorów złych.}
Dla $j\in\mathbb N$ i $r>0$ oznaczmy przez
$\vartheta_{j,r}(x)$ iloraz
\[
 \frac{\lambda_n\bigl(\{z\in B(x,r)\cap U:
 |f(z)-f(x)-G(x)(z-x)|>j^{-1}|z-x|\}\bigr)}
 {\lambda_n(B(x,r))}.
\]
Z przybliżonej różniczkowalności wynika
$\vartheta_{j,r}(x)\to0$ przy $r\downarrow0$ dla
każdego $x\in D$. Dla ustalonego $j$ funkcje
$H_{j,l}(x)=\sup_{0<r<1/l}\vartheta_{j,r}(x)$ są
mierzalne: dla ustalonego $x$ licznik zależy ciągle
od $r$ (sfery mają miarę zero), toteż supremum można
brać po promieniach wymiernych; mierzalność całki
zależącej od $x$ daje twierdzenie Tonellego.
Mamy $H_{j,l}(x)\downarrow0$.

Wybierzmy stałą $\delta=8^{-n}$. Z twierdzenia Jegorowa
stosowanego kolejno dla $j=1,2,\ldots$ wybieramy
mierzalne podzbiory $E_j\subset K_0$ tak, że
$\lambda_n(K_0\setminus E_j)<\varepsilon/(3\cdot 2^j)$
oraz $H_{j,l}\to0$ jednostajnie na $E_j$.
Niech $E=\bigcap_j E_j$. Następnie wybierzmy zwarty
$K\subset E$ z $\lambda_n(E\setminus K)<\varepsilon/3$.
Wtedy $\lambda_n(A\setminus K)<\varepsilon$.
Dla każdego $j$ istnieje $\rho_j>0$ taka, że
$\vartheta_{j,r}(x)<\delta$ jednocześnie dla
$x\in K$ i $0<r<\rho_j$.

\emph{Krok 3: z kontroli w mierze do warunku Whitneya.}
Niech $x,y\in K$, $d=|x-y|>0$ i
$2d<\rho_j$; wymagamy także, aby $B(x,2d),B(y,2d)$
były zawarte w $U$. Przecięcie tych kul zawiera
$B((x+y)/2,3d/2)$, więc jego miara jest większa
niż suma dwóch górnych oszacowań
$\delta\lambda_n(B(0,2d))$ na zbiory złe
(dla $\delta=8^{-n}$). Istnieje zatem punkt $z$
w przecięciu, dla którego jednocześnie
\[
 |f(z)-f(x)-G(x)(z-x)|\le 2d/j,
 \quad |f(z)-f(y)-G(y)(z-y)|\le 2d/j.
\]
Odejmując równości otrzymujemy
\[
 |f(y)-f(x)-G(x)(y-x)|
 \le4d/j+2d\,\|G(x)-G(y)\|.
\]
Jednostajna ciągłość $G$ na $K$ i dowolność $j$
dowodzą, że iloraz lewej strony przez $|x-y|$
dąży jednostajnie do zera. Tak więc $(f|_K,G|_K)$
jest polem dżetów Whitneya klasy $C^1$.

\emph{Krok 4: przedłużenie.}
Twierdzenie~\ref{thm:dzety-whitney-extension}
daje $F\in C^1(\mathbb R^n)$ o danych wartościach
i różniczkach na $K$.

\emph{Odwrotność.}
Dla każdego wyczerpującego ciągu zbiorów
$A_l=U\cap B(0,l)$ i każdego $l$ możemy wybrać
zbiory zwarte $K_{l,j}\subset A_l$, na których
$f=F_{l,j}\in C^1$, z
$\lambda_n(A_l\setminus K_{l,j})<1/j$.
Prawie każdy punkt $A_l$ należy do przynajmniej
jednego z tych zbiorów i jest jego punktem gęstości.
W takim punkcie $f$ ma przybliżoną różniczkę
$DF_{l,j}$, ponieważ obie funkcje zgadzają się na
zbiorze gęstości $1$. Biorąc sumę po $l$,
otrzymujemy przybliżoną różniczkowalność prawie wszędzie.
\end{proof}

\begin{wniosek}[Luzinowskie wygładzenie odwzorowania Lipschitzowskiego]
\label{cor:dzety-lusin-lipschitz}
Niech $U\subset\mathbb R^n$ będzie otwarty,
$f:U\to\mathbb R^m$ lokalnie Lipschitzowskie,
a $A\subset U$ mierzalny o skończonej mierze.
Dla każdego $\varepsilon>0$ istnieje
$F\in C^1(\mathbb R^n,\mathbb R^m)$ takie, że
\[
 \lambda_n(\{x\in A:F(x)\ne f(x)\})<\varepsilon,
\]
a na pewnym zwartym $K\subset A$ o
$\lambda_n(A\setminus K)<\varepsilon$ zachodzi
$DF=Df$.
\end{wniosek}
\begin{proof}
Twierdzenie Rademachera daje przybliżoną
różniczkowalność prawie wszędzie.
Stosujemy Twierdzenie~\ref{thm:dzety-whitney-lusin}.
\end{proof}

\begin{uwaga}[Granice twierdzenia]
Twierdzenie~\ref{thm:dzety-whitney-lusin} nie mówi,
że $f$ jest funkcją $C^1$ po zmianie na zbiorze miary zero.
Wolno odrzucić zbiór o \emph{dowolnie małej}, lecz niekoniecznie
zerowej mierze, zależny od $\varepsilon$.
Na przykład $f(x)=|x|$ na $(-1,1)$ nie ma globalnego
przedstawiciela $C^1$ równego $f$ prawie wszędzie.
\end{uwaga}

\section{Dżety wyższych rzędów i zgodność danych}
Dla $k\ge2$ sam formalny zapis danych
$(a_\alpha(x))_{|\alpha|\le k}$ nie gwarantuje,
że są one pochodnymi jednej funkcji. W teorii miary
istotne jest, czy można znaleźć duży zbiór, na którym
dane spełniają \emph{wszystkie} warunki Whitneya.

\begin{twierdzenie}[Luzinowskie przedłużenie zgodnych dżetów rzędu $k$]
\label{thm:dzety-lusin-higher}
Niech $U\subset\mathbb R^n$ będzie otwarty, $k\ge1$,
a $a_\alpha:U\to\mathbb R$, $|\alpha|\le k$,
będą funkcjami mierzalnymi. Załóżmy, że dla
$\lambda_n$-prawie każdego $x\in U$ i każdego
$|\alpha|\le k$ zachodzi zwykła, względna do
pełnomiarowego zbioru $D\subset U$, ekspansja
\[
 a_\alpha(y)=\sum_{|\beta|\le k-|\alpha|}
 \frac{a_{\alpha+\beta}(x)}{\beta!}(y-x)^\beta
 +o(|y-x|^{k-|\alpha|})\quad(D\ni y\to x).
\]
Wówczas dla każdego mierzalnego $A\subset U$
o skończonej mierze i każdego $\varepsilon>0$
istnieją $K\subset A$ zwarty i
$F\in C^k(\mathbb R^n)$ takie, że
$\lambda_n(A\setminus K)<\varepsilon$ oraz
$D^\alpha F|_K=a_\alpha|_K$ dla wszystkich
$|\alpha|\le k$.
\end{twierdzenie}
\begin{proof}
Twierdzenie Luzina zastosowane do skończonej rodziny
$a_\alpha$ daje zwarty $K_0\subset A\cap D$ o dowolnie
małym dopełnieniu i ciągłości wszystkich współczynników.
Dla każdej wielopochodnej $\alpha$ i promienia $r$
weźmy supremum po $y\in K_0$ z $0<|x-y|<r$
ilorazu reszty Taylora przez $|x-y|^{k-|\alpha|}$.
Jest to funkcja mierzalna względem $x\in K_0$:
możemy zastąpić wszystkie $y$ przeliczalnym zbiorem
gęstym w $K_0$, ponieważ na $K_0$ współczynniki
są ciągłe. Z hipotezy supremum dąży do zera punktowo przy
$r\downarrow0$, ponieważ $K_0\subset D$ i każda
reszta jest oceniana wzdłuż całego zbioru $D$.
To właśnie tu korzystamy z silniejszej, względnej
wersji założenia, a nie jedynie ze zbieżności przybliżonej. Twierdzenie Jegorowa
pozwala przejść do zwartego $K\subset K_0$ z
arbitralnie małą utratą miary, na którym zbieżność
jest jednostajna dla skończenie wielu $\alpha$.
Właśnie to są warunki pola Whitneya z
Definicji~\ref{def:dzety-whitney-field}. Zastosowanie
Twierdzenia~\ref{thm:dzety-whitney-extension}
kończy dowód.
\end{proof}

\begin{uwaga}[Ważne rozróżnienie]
Warunek z Twierdzenia~\ref{thm:dzety-lusin-higher}
jest mocniejszy od samego istnienia przybliżonego
wielomianu Peana rzędu $k$ dla jednej funkcji.
Dla ogólnego twierdzenia o przybliżonych dżetach
wyższych rzędów potrzebny jest dodatkowy krok
przechodzący od informacji w mierze do zgodności
wszystkich pochodnych na wybranym zbiorze.
Nie wolno pomijać go, stosując twierdzenie Whitneya.
\end{uwaga}

\section{Dodatkowe przykłady i powiązania}
\begin{przyklad}[Dżet holomorficzny]
Niech $U\subset\mathbb C^n$ będzie obszarem,
$p\in U$ i $f\in\mathcal O(U)$. Jego holomorficzny
$k$-dżet zapisujemy w lokalnych współrzędnych jako
\[
 j_p^kf\equiv\sum_{|\alpha|\le k}
 \frac{\partial_z^\alpha f(p)}{\alpha!}(z-p)^\alpha
 \pmod{\mathfrak m_p^{k+1}}.
\]
Warunki interpolacji wartości i pochodnych
(np. w zadaniach typu Carathéodory'ego--Fejéra)
są warunkami na holomorficzne dżety.
W przeciwieństwie do dowolnych gładkich funkcji,
holomorficzne dżety podlegają globalnym ograniczeniom
wynikającym z istnienia funkcji holomorficznej na
ustalonym obszarze. Twierdzenie Whitneya w klasie $C^k$
nie zapewnia holomorficznego przedłużenia.
\end{przyklad}

\begin{przyklad}[Dżety a metoda elementów skończonych]
Przybliżenia typu Hermite'a w jednym wymiarze
łączą zadane wartości $u(x_i)$ i pochodne $u'(x_i)$
na końcach przedziałów. Są więc interpolacją
$1$-dżetów na skończonym zbiorze węzłów.
Warunki zgodności na węzłach decydują,
czy powstaje aproksymacja klasy $C^1$.
\end{przyklad}

\begin{przyklad}[Dżety i geometryczne oszacowania rzędu]
Dwie gładkie funkcje o identycznym $k$-dżecie w $p$
mogą różnić się globalnie, ale ich różnica
spełnia $f(x)-g(x)=o(|x-p|^k)$ przy $x\to p$.
Jeżeli mają również równe pochodne rzędu $k+1$,
to ich $(k+1)$-dżety są równe i otrzymujemy
odpowiednią poprawę lokalnego rzędu reszty.
\end{przyklad}

\section{Mapa zależności i zakres udowodnionych rezultatów}
\begin{center}
\begin{tikzpicture}[>=stealth, every node/.style={font=\small,align=center},
 box/.style={draw,rounded corners=2pt,inner sep=5pt,text width=3.8cm}]
 \node[box] (j) at (-2.4,2.8) {Dżety gładkie\\wiązki i formy kontaktowe};
 \node[box] (a) at (2.4,2.8) {Gęstość Lebesgue'a\\dżety przybliżone};
 \node[box] (w) at (-2.4,0.7) {Warunki Whitneya\\przedłużenie $C^k$};
 \node[box] (l) at (2.4,0.7) {Whitney--Luzin\\wybór dużego $K$};
 \node[box] (c) at (0,-1.45) {Zastosowania: Sobolew, $BV$,\\area/coarea, PDE};
 \draw[->] (j) -- (w);
 \draw[->] (a) -- (l);
 \draw[->] (w) -- (l);
 \draw[->] (l) -- (c);
\end{tikzpicture}
\end{center}

\begin{uwaga}[Co zostało udowodnione, a co przyjęte?]
W tym rozdziale szczegółowo wyprowadzono algebraiczny opis
dżetów, strukturę afiniczną, kryterium holonomiczności,
jednoznaczność dżetów przybliżonych, twierdzenie Whitneya
metodą kostek i twierdzenie Whitneya--Luzina dla $C^1$.
Osobno przyjęto znane twierdzenia Rademachera, Luzina,
Jegorowa i Sobolewa o przybliżonej różniczkowalności
oraz ogólne wzory area/coarea i rozkład $BV$.
Ich przytoczenie pokazuje zastosowania, nie stanowi
nowego dowodu tych wyników.
\end{uwaga}


\section*{Wskazówki bibliograficzne}
Konstrukcja wiązek dżetów i form kontaktowych jest standardowo
przedstawiana w monografii D.~J.~Saundersa,
\emph{The Geometry of Jet Bundles}. Dla teorii przedłużeń
podstawową pracą jest H.~Whitney,
\emph{Analytic extensions of differentiable functions defined
in closed sets} (1934). Przybliżona różniczkowalność,
twierdzenia typu Luzina i ich związek z teorią Sobolewa i $BV$
są omawiane w książce L.~C.~Evansa i R.~F.~Gariepy'ego,
\emph{Measure Theory and Fine Properties of Functions}.
W tekście rozdzielono twierdzenia dowodzone tutaj od
rezultatów zewnętrznych (Rademacher, własności Sobolewa,
wzory area i koobszaru). W konstrukcji Whitneya wykorzystuje
się standardowy rozkład na kostki i rozkład jedności;
podano ich własności oraz argument sklejania na zbiorze
domkniętym. Dowody powinny być czytane z zachowaniem
wskazanych założeń regularności.

% Materiał przeniesiony z CAL-12; zachowano oryginalne etykiety.
\providecommand{\atlasref}[2]{#2}
\chapter{Całka Daniella--Stone'a}
\label{ch:daniell-stone-32}

Korzystamy z \atlasref{ME-03::thm:caratheodory}{ME-03: twierdzenie Carathéodory'ego}
oraz \atlasref{ME-02::thm:uniqueness-finite}{ME-02: jednoznaczność miar skończonych}.
Nie powtarzamy ich dowodów. Szczegółowo wykażemy natomiast,
że funkcjonał Daniella pozwala skonstruować miarę zewnętrzną,
że jej ograniczenie reprezentuje funkcjonał i kiedy
reprezentująca miara jest jednoznaczna.
Dla porównania reprezentację dodatnich funkcjonałów
na przestrzeniach funkcji ciągłych omawia
\atlasref{ME-21::thm:riesz-positive-21}{ME-21: twierdzenie Riesza--Markowa--Kakutaniego}.

W tej konstrukcji zaczynamy od funkcjonału całkowego na kracie funkcji, a dopiero potem odzyskujemy miarę. Nie jest to kolejne rozszerzenie całki niebezwzględnej, lecz abstrakcyjny sposób konstruowania całki Lebesgue'a.
\section{Krata funkcji i ciągłość porządkowa}
\begin{definicja}[Krata i krata wektorowa]\label{it:def:daniell-lattice}
Kratą jest zbiór częściowo uporządkowany $E$, w którym każda para $u,v$ ma w $E$ supremum i infimum, oznaczane $u\vee v$ i $u\wedge v$. Jeżeli ponadto $E$ jest rzeczywistą przestrzenią wektorową, a porządek jest zgodny z dodawaniem i mnożeniem przez skalary nieujemne, to $E$ jest \emph{kratą wektorową} (przestrzenią Riesza). Dla zbioru $\Omega$ kratą funkcji nazywamy podprzestrzeń liniową $H\subset\mathbb{R}^\Omega$ zamkniętą na punktowe maksimum i minimum. Równoważnie: $f\in H\Rightarrow|f|\in H$.
\end{definicja}
Istotnie, $f\vee g=\tfrac12(f+g+|f-g|)$, $f\wedge g=\tfrac12(f+g-|f-g|)$, $f^+=f\vee0$, $f^-=(-f)\vee0$. W przeciwnym kierunku $|f|=f\vee(-f)$, więc zamkniętość na operacje kratowe implikuje zamkniętość na moduł. Nie wymagamy, aby każde dwa elementy kraty były porównywalne.
\begin{przyklad}[Funkcje ciągłe i wielomiany]
Przestrzeń $C[a,b]$ jest kratą wektorową: maksimum, minimum i moduł funkcji ciągłej są ciągłe. Natomiast wielomiany na $[-1,1]$ tworzą przestrzeń liniową, ale nie kratę: dla $f(x)=x$ funkcja $|f(x)|=|x|$ nie jest wielomianem, ponieważ nie jest różniczkowalna w zerze.
\end{przyklad}
\begin{uwaga}[Rola warunku kraty]
Zamkniętość na operacje kratowe pozwala obcinać funkcje i porównywać je bez opuszczania dziedziny funkcjonału. Obecność funkcji stałej $1$ jest osobnym założeniem. Warunek kraty nie oznacza zamkniętości na granice ciągów monotonicznych.
\end{uwaga}
\begin{definicja}[Funkcjonał Daniella]
\label{it:def:daniell-functional}
Niech $\Omega\ne\varnothing$ i niech $H$ będzie kratą wektorową ograniczonych funkcji rzeczywistych na $\Omega$, zawierającą $\chi_\Omega=1$. Liniowy funkcjonał $I_0:H\to\mathbb{R}$ nazywamy funkcjonałem Daniella, jeśli jest dodatni ($h\ge0\Rightarrow I_0(h)\ge0$) i ciągły porządkowo z góry ($h_n\in H$, $h_n\downarrow0$ punktowo $\Rightarrow I_0(h_n)\downarrow0$). Symbolem $\sigma(H)$ oznaczamy najmniejszą $\sigma$-algebrę, względem której wszystkie funkcje z $H$ są mierzalne. Dodatniość i liniowość implikują monotoniczność: jeśli $f,g\in H$ i $f\le g$, to $I_0(g)-I_0(f)=I_0(g-f)\ge0$. Warunek ciągłości porządkowej dotyczy ciągów malejących do zera punktowo na całej dziedzinie, nie tylko do zera w porządku wewnętrznym kraty. Na początku zakładamy ograniczoność i obecność jedynki; dalej rozpatrzymy także przypadki bez tych założeń.
\end{definicja}
\begin{lemat}[Pierwszy krok przedłużenia]\label{it:lem:daniell-monotone-extension}
Jeśli $0\le h_n\uparrow f$ punktowo, $h_n\in H$, to wartość $\lim_n I_0(h_n)\in[0,\infty]$ zależy tylko od $f$, a nie od wybranego ciągu przybliżającego.
\end{lemat}
\begin{proof}
Niech również $0\le g_n\uparrow f$. Dla ustalonego $k$ mamy $(h_k-g_n)^+\downarrow0$, przy czym te funkcje należą do $H$ dzięki własności kraty. Z ciągłości porządkowej $I_0((h_k-g_n)^+)\downarrow0$. Ponieważ $h_k\le g_n+(h_k-g_n)^+$, dostajemy $I_0(h_k)\le I_0(g_n)+I_0((h_k-g_n)^+)$. Przechodząc z $n\to\infty$, a następnie biorąc supremum po $k$, uzyskujemy $\lim_kI_0(h_k)\le\lim_nI_0(g_n)$. Zamiana ról daje nierówność przeciwną. Argument działa również dla wspólnej wartości $+\infty$.
\end{proof}
\section{Konstrukcja miary zewnętrznej i reprezentacja}
\begin{lemat}[Oszacowanie przez przeliczalne pokrycie]\label{it:lem:daniell-cover-domination}
Niech $u\in H$, $u\ge0$, oraz $h_n\in H$, $h_n\ge0$. Jeżeli $u\le\sum_{n\ge1}h_n$ punktowo (dopuszczamy nieskończoną sumę), to
\[
 I_0(u)\le\sum_{n\ge1}I_0(h_n).
\]
\end{lemat}
\begin{proof}
Funkcje $r_N=(u-\sum_{n=1}^N h_n)^+$ należą do $H$ i maleją punktowo do zera. Z ciągłości porządkowej $I_0(r_N)\to0$. Ponieważ $u\le\sum_{n=1}^N h_n+r_N$, dodatniość i liniowość dają
$I_0(u)\le\sum_{n=1}^N I_0(h_n)+I_0(r_N)$. Granica po $N$ daje tezę również wtedy, gdy szereg kosztów jest rozbieżny do $+\infty$.
\end{proof}

\begin{definicja}[Miara zewnętrzna Daniella]\label{it:def:daniell-outer}
Dla $A\subset\Omega$ określamy
\begin{equation}\label{it:eq:daniell-outer}
 \mu^*(A)=\inf\left\{\sum_{n=1}^\infty I_0(h_n):
 h_n\in H,\ h_n\ge0,\ \chi_A\le\sum_{n=1}^\infty h_n\right\}.
\end{equation}
Pokrycia skończone traktujemy jako ciągi uzupełnione funkcjami zerowymi. Przyjmujemy $\inf\varnothing=+\infty$, co będzie istotne w wersji bez jedynki.
\end{definicja}
\begin{lemat}[Własności miary zewnętrznej]\label{it:lem:daniell-outer}
Funkcja $\mu^*$ jest miarą zewnętrzną, a $\mu^*(\Omega)=I_0(1)<\infty$. Dla $u\in H$ z $0\le u\le\chi_A$ mamy $I_0(u)\le\mu^*(A)$.
\end{lemat}
\begin{proof}
Pokrycie funkcją $1$ daje $0\le\mu^*(A)\le I_0(1)$; pokrycie zerowe daje $\mu^*(\varnothing)=0$. Monotoniczność wynika z inkluzji rodzin dopuszczalnych pokryć. Jeśli $A=\bigcup_{j\ge1}A_j$, dla $\varepsilon>0$ wybieramy pokrycie $(h_{j,n})_{n\ge1}$ każdego $A_j$ o koszcie nie większym niż $\mu^*(A_j)+\varepsilon 2^{-j}$. Przenumerowany ciąg wszystkich $h_{j,n}$ pokrywa $A$. Ponieważ wszystkie koszty są nieujemne, jego koszt jest równy sumie kosztów pokryć, więc
$\mu^*(A)\le\sum_j\mu^*(A_j)+\varepsilon$. Przejście z $\varepsilon\downarrow0$ daje podaddytywność przeliczalną. Jeśli $0\le u\le\chi_A$, każde pokrycie $A$ dominuje $u$, więc lemat \ref{it:lem:daniell-cover-domination} daje $I_0(u)\le\mu^*(A)$ po wzięciu infimum. Dla $A=\Omega$ i $u=1$ otrzymujemy nierówność odwrotną do oszacowania z pokrycia funkcją $1$.
\end{proof}

\begin{lemat}[Addytywność przy rozdzieleniu wartościami funkcji]\label{it:lem:daniell-separated}
Jeżeli $h\in H$, $s<t$, $A\subset\{h\le s\}$ i $B\subset\{h\ge t\}$, to $\mu^*(A\cup B)=\mu^*(A)+\mu^*(B)$.
\end{lemat}
\begin{proof}
Funkcja $v=1\wedge((h-s)^+/(t-s))$ należy do $H$
i spełnia $0\le v\le1$, $v|_A=0$, $v|_B=1$.
Weźmy dowolne pokrycie $(h_n)$ zbioru $A\cup B$
z definicji \ref{it:def:daniell-outer}.
Kładziemy $k_n=h_n\wedge1$. Jest to nadal pokrycie:
w każdym punkcie $A\cup B$ albo pewien składnik $h_n$
jest co najmniej równy $1$, albo żaden składnik
nie zmienia się przy obcięciu. Ponadto
$I_0(k_n)\le I_0(h_n)$.
Określamy
\[
 b_n=k_n\wedge v,\qquad a_n=k_n-b_n.
\]
Obie funkcje należą do $H$ i są nieujemne.
Na $A$ mamy $a_n=k_n$, na $B$ mamy $b_n=k_n$,
więc $(a_n)$ pokrywa $A$, a $(b_n)$ pokrywa $B$.
Z liniowości wynika
\[
 \mu^*(A)+\mu^*(B)\le
 \sum_n\bigl(I_0(a_n)+I_0(b_n)\bigr)
 =\sum_n I_0(k_n)\le\sum_n I_0(h_n).
\]
Bierzemy infimum po pokryciach. Nierówność odwrotna
wynika z podaddytywności.
Użyliśmy tylko działań liniowych i kratowych,
bez mnożenia funkcji. Ten argument wykorzystamy
także w wersji z warunkiem Stone'a, bez jedynki.
\end{proof}

\begin{lemat}[Mierzalność zbiorów wyznaczonych przez kratę]\label{it:lem:daniell-measurability}
Każdy zbiór $\{h>t\}$, $h\in H$, $t\in\mathbb{R}$, jest $\mu^*$-mierzalny w sensie Carathéodory'ego. W konsekwencji $\sigma(H)$ zawiera się w $\sigma$-algebrze zbiorów $\mu^*$-mierzalnych.
\end{lemat}
\begin{proof}
Ustalmy $E\subset\Omega$ i połóżmy $B=E\cap\{h\le t\}$, $C=E\cap\{h>t\}$ oraz $C_n=E\cap\{h>t+2^{-n}\}$ dla $n\ge0$. Z lematu \ref{it:lem:daniell-separated} i monotoniczności mamy
\[
 \mu^*(E)\ge\mu^*(B)+\mu^*(C_n).
\]
Nie wolno jeszcze stosować ciągłości miary z dołu do $C_n$, ponieważ ich mierzalności dopiero dowodzimy. Rozważmy warstwy
\[
 S_k=E\cap\{t+2^{-(k+1)}<h\le t+2^{-k}\},\quad k\ge0.
\]
Dwie warstwy o indeksach różniących się co najmniej o dwa są rozdzielone dodatnią przerwą między zakresami wartości $h$. Dla skończonej rodziny warstw parzystych numerujemy indeksy rosnąco. Gdy dołączamy warstwę $S_k$, $k\ge2$, mamy $h\le t+2^{-k}$ na $S_k$, natomiast na sumie wcześniej dołączonych warstw zachodzi $h>t+2^{-(k-1)}$. Progi są różne, ponieważ $2^{-k}<2^{-(k-1)}$. Lemat \ref{it:lem:daniell-separated} pozwala więc dodać miarę zewnętrzną nowej warstwy do miary zewnętrznej sumy wcześniejszych. Indukcja daje addytywność dla każdej skończonej rodziny warstw parzystych; dokładnie to samo rozumowanie działa dla nieparzystych. Monotoniczność daje
$\sum_{k\ge0}\mu^*(S_{2k})\le\mu^*(E)$ oraz
$\sum_{k\ge0}\mu^*(S_{2k+1})\le\mu^*(E)$.
Obie sumy są skończone na mocy lematu \ref{it:lem:daniell-outer}. Ponieważ $C\setminus C_n=\bigcup_{k\ge n}S_k$, mamy
\[
 \mu^*(C_n)\le\mu^*(C)\le
 \mu^*(C_n)+\sum_{k\ge n}\mu^*(S_k).
\]
Ogon dąży do zera, zatem $\mu^*(C_n)\to\mu^*(C)$. W oszacowaniu $\mu^*(E)\ge\mu^*(B)+\mu^*(C_n)$ przechodzimy do granicy i otrzymujemy $\mu^*(E)\ge\mu^*(B)+\mu^*(C)$. Odwrotna nierówność jest podaddytywnością. Jest to kryterium Carathéodory'ego dla $\{h>t\}$. Rodzina zbiorów spełniających to kryterium jest $\sigma$-algebrą na mocy lematu \ref{lem:hausdorff-caratheodory-selfcontained-26}; zawiera wszystkie generatory $\sigma(H)$, ponieważ półproste $(t,\infty)$ generują $\mathcal B(\mathbb{R})$.
\end{proof}

\begin{lemat}[Odzyskanie masy zbioru poziomicowego]\label{it:lem:daniell-level}
Niech $h\in H$, $t\in\mathbb{R}$ i $u_n=1\wedge n(h-t)^+$. Wówczas
\[
 \mu^*(\{h>t\})=\lim_{n\to\infty}I_0(u_n).
\]
\end{lemat}
\begin{proof}
Mamy $u_n\in H$ oraz $0\le u_n\uparrow\chi_{\{h>t\}}$ punktowo. Lemat \ref{it:lem:daniell-outer} daje $I_0(u_n)\le\mu^*(\{h>t\})$. Odwrotnie, funkcje $u_1,u_2-u_1,u_3-u_2,\ldots$ należą do $H$, są nieujemne, a ich punktowa suma jest równa $\chi_{\{h>t\}}$. Tworzą więc pokrycie o koszcie $\lim_n I_0(u_n)$ dzięki teleskopowaniu skończonych sum. Definicja \ref{it:def:daniell-outer} daje nierówność przeciwną.
\end{proof}

\begin{twierdzenie}[Daniella--Stone'a: istnienie i jednoznaczność miary]
\label{it:thm:daniell-stone}
Dla funkcjonału $I_0$ z definicji \ref{it:def:daniell-functional} istnieje dokładnie jedna skończona miara $\mu$ na $\sigma(H)$ taka, że $I_0(h)=\int_\Omega h\,\mathrm d\mu$ dla wszystkich $h\in H$. Ponadto $\mu(\Omega)=I_0(1)$. Po uzupełnieniu miary całkę rozszerzamy na nieujemne funkcje mierzalne i następnie na funkcje z $\mathcal L^1(\mu)$.
\end{twierdzenie}
\begin{proof}[Dowód istnienia i reprezentacji]
Niech $\mathcal M_{\mu^*}$ będzie $\sigma$-algebrą zbiorów $\mu^*$-mierzalnych w sensie Carathéodory'ego. Z lematu \ref{lem:hausdorff-caratheodory-selfcontained-26} i lematu \ref{it:lem:daniell-measurability} wynika, że $\mu=\mu^*|_{\sigma(H)}$ jest miarą skończoną, o masie $\mu(\Omega)=I_0(1)$.

Ustalmy $h\in H$, $h\ge0$. Dla dowolnego $\delta>0$ wybieramy liczbę naturalną $m\ge1$ tak, aby $h\le m\delta$. Niech
\[
 L_\delta=\delta\sum_{k=1}^{m}\chi_{\{h>k\delta\}},\qquad
 U_\delta=\delta\sum_{k=0}^{m-1}\chi_{\{h>k\delta\}}.
\]
Zliczając progi mniejsze od wartości $h(\omega)$, także gdy $h(\omega)$ jest równa progowi, otrzymujemy
$L_\delta\le h\le U_\delta$ oraz
$0\le U_\delta-L_\delta\le\delta$.
Dla każdego $n\ge1$ zachodzi
$\delta\sum_{k=1}^{m}(1\wedge n(h-k\delta)^+)\le L_\delta\le h$.
Stosujemy $I_0$ do pierwszej i ostatniej funkcji (obie należą do $H$). Liniowość, dodatniość i lemat \ref{it:lem:daniell-level} po przejściu z $n\to\infty$ dają
\[
 \int_\Omega L_\delta\,\mathrm d\mu\le I_0(h).
\]
W drugą stronę korzystamy z dokładnego rozkładu na obcięte przyrosty
\[
 h=\delta\sum_{k=0}^{m-1}
 \left(1\wedge\frac{(h-k\delta)^+}{\delta}\right).
\]
Istotnie, dla wartości $h(\omega)\in[j\delta,(j+1)\delta]$ z $0\le j<m$ pierwszych $j$ składników w nawiasach jest równych $1$, następny wynosi $(h(\omega)-j\delta)/\delta$, a pozostałe są zerowe; wzór obejmuje także $h(\omega)=m\delta$. Każdy składnik w nawiasie należy do $H$ i nie przekracza $\chi_{\{h>k\delta\}}$. Z lematu \ref{it:lem:daniell-outer} wynika zatem
$I_0(h)\le\delta\sum_{k=0}^{m-1}\mu(\{h>k\delta\})=\int_\Omega U_\delta\,\mathrm d\mu$.
Całka $h$ również leży między tymi dwiema całkami, które różnią się najwyżej o $\delta\mu(\Omega)$. Dowolność $\delta>0$ daje $I_0(h)=\int_\Omega h\,\mathrm d\mu$. Dla dowolnego $h\in H$ stosujemy wynik do $h^+$ i $h^-$, które należą do $H$, i odejmujemy skończone wartości. Otrzymaliśmy reprezentację.

Miara $\mu$ może nie być zupełna na $\sigma(H)$.
Stosujemy \atlasref{ME-02::thm:completion}{ME-02: uzupełnienie miary}.
Uzupełnienie nie zmienia masy całkowitej ani całek
funkcji mierzalnych względem pierwotnej $\sigma$-algebry:
najpierw dla funkcji prostych, następnie przez
twierdzenie \ref{thm:mct} i rozkład na części
dodatnią i ujemną. Całkę na uzupełnieniu definiujemy
zgodnie z \atlasref{ME-07::ch:calkowanie-lebesgue}{Teorią miary, rozdział 7: konstrukcja całki Lebesgue'a}.
\end{proof}
\begin{proof}[Dowód jednoznaczności]
Niech $\mu,\nu$ reprezentują $I_0$. Dla $h\in H$, $t\in\mathbb{R}$, funkcje $u_n=1\wedge n(h-t)^+$ należą do $H$ i rosną do $\chi_{\{h>t\}}$. Stąd $\mu(\{h>t\})=\lim_nI_0(u_n)=\nu(\{h>t\})$. Aby odzyskać skończone przecięcia zbiorów tego rodzaju, stosujemy funkcje $v_n=1\wedge n\bigwedge_{j=1}^m(h_j-t_j)^+$, które również należą do $H$ i rosną do $\chi_{\bigcap_{j=1}^m\{h_j>t_j\}}$. Zbieżność całek w obu przypadkach wynika z twierdzenia \ref{thm:mct}. Zatem miary zgadzają się na układzie $\pi$ skończonych przecięć tych zbiorów (wraz z $\Omega$, odpowiadającym pustemu przecięciu), generującym $\sigma(H)$, a obie są skończone i mają masę $I_0(1)$. Z \atlasref{ME-02::thm:uniqueness-finite}{ME-02: jednoznaczność miar skończonych}
otrzymujemy $\mu=\nu$ na $\sigma(H)$.
\end{proof}
\begin{wniosek}[Otoczka mierzalna i uzupełnienie miary Daniella]
\label{it:cor:daniell-completion}
Niech $\mu^*$ i $\mu$ będą określone odpowiednio
w definicji \ref{it:def:daniell-outer} i twierdzeniu
\ref{it:thm:daniell-stone}. Dla każdego $A\subseteq\Omega$
istnieje $B\in\sigma(H)$ takie, że
\[
 A\subseteq B,\qquad \mu(B)=\mu^*(A).
\]
W szczególności $\mu^*$ jest miarą zewnętrzną indukowaną
przez $\mu$, a $\sigma$-algebra zbiorów
$\mu^*$-mierzalnych w sensie Carathéodory'ego
jest dokładnie uzupełnieniem $\sigma(H)$ względem $\mu$.
\end{wniosek}
\begin{proof}
Dla każdego $k\ge1$ wybieramy pokrycie
$(h_{k,n})_{n\ge1}$ z definicji miary zewnętrznej
o koszcie mniejszym od $\mu^*(A)+2^{-k}$.
Wybór jest możliwy, ponieważ $\mu^*(A)$ jest
skończona i rodzina pokryć jest niepusta.
Funkcja
$g_k=\sum_{n\ge1}h_{k,n}:\Omega\to[0,\infty]$
jest $\sigma(H)$-mierzalna jako granica rosnąca
mierzalnych sum częściowych. Nie musi należeć do $H$.
Zbiór $B_k=\{g_k\ge1\}$ jest mierzalny i zawiera $A$.
Z reprezentacji $I_0$ oraz twierdzenia
\ref{thm:sum-nonnegative-integrals} wynika
\[
 \mu(B_k)\le \int_\Omega g_k\,\mathrm d\mu
 =\sum_{n\ge1}I_0(h_{k,n})<\mu^*(A)+2^{-k}.
\]
Kładziemy $B=\bigcap_{k\ge1}B_k$.
Z monotoniczności oraz zgodności $\mu$ z $\mu^*$
na $\sigma(H)$ mamy
$\mu^*(A)\le\mu(B)\le\mu(B_k)<\mu^*(A)+2^{-k}$
dla każdego $k$, więc $\mu(B)=\mu^*(A)$.
Dowodzi to także wzoru
$\mu^*(A)=\inf\{\mu(D):A\subseteq D,\ D\in\sigma(H)\}$,
czyli równości z miarą zewnętrzną indukowaną przez
$\mu$ (każde przeliczalne pokrycie mierzalne można
zastąpić jego sumą, nie zwiększając kosztu).

Jeżeli $A$ jest mierzalny w sensie Carathéodory'ego,
to stosując kryterium do otoczki $B$ otrzymujemy
$\mu^*(B)=\mu^*(A)+\mu^*(B\setminus A)$.
Wszystkie wartości są skończone i
$\mu^*(B)=\mu^*(A)$, zatem $\mu^*(B\setminus A)=0$.
Pierwsza część dowodu daje $N\in\sigma(H)$
z $B\setminus A\subseteq N$ i $\mu(N)=0$.
Wówczas $B\setminus N\subseteq A\subseteq B$,
a $\mu(B\setminus(B\setminus N))=0$.
Z definicji uzupełnienia miary wynika, że $A$
należy do uzupełnionej $\sigma$-algebry.
Odwrotnie, na mocy
\atlasref{ME-03::thm:caratheodory}{ME-03: twierdzenie Carathéodory'ego}
miara na $\mathcal M_{\mu^*}$ jest zupełna.
Ponieważ $\sigma(H)\subseteq\mathcal M_{\mu^*}$
i obie miary zgadzają się na $\sigma(H)$,
każdy podzbiór $\mu$-zerowego zbioru z $\sigma(H)$
należy do $\mathcal M_{\mu^*}$.
Z charakterystyki uzupełnienia
(\atlasref{ME-02::thm:completion}{ME-02: uzupełnienie miary})
otrzymujemy przeciwną inkluzję.
\end{proof}

\begin{wniosek}[Konieczność ciągłości porządkowej]\label{it:cor:daniell-necessity}
Jeżeli dodatni liniowy funkcjonał $I_0:H\to\mathbb{R}$ na kracie
rzeczywistych funkcji jest reprezentowany przez miarę
na $\sigma(H)$, względem której wszystkie funkcje z $H$
są bezwzględnie całkowalne, to $I_0$ jest ciągły
porządkowo z góry. Nie wymagamy tu $1\in H$ ani
skończoności miary.
\end{wniosek}
\begin{proof}
Jeśli $h_n\downarrow0$ punktowo, to $0\le h_n\le h_1$,
a $h_1$ jest całkowalna na mocy założenia o reprezentacji.
Twierdzenie \ref{thm:dct} daje $\int_\Omega h_n\,\mathrm d\mu\to0$,
czyli $I_0(h_n)\to0$. Ponadto z dodatniości
$I_0(h_n-h_{n+1})\ge0$, więc wartości maleją.
W szczególności dla kraty ograniczonych funkcji
zawierającej $1$ całkowalność wynika ze skończoności
$\mu(\Omega)=I_0(1)$.
\end{proof}

\section{Gęstość kraty i jednoznaczność rozszerzenia}

\begin{twierdzenie}[Gęstość początkowej kraty w przestrzeni $L^1$]
\label{it:thm:daniell-l1-density}
Niech $H,I_0,\mu$ spełniają założenia i tezę
twierdzenia \ref{it:thm:daniell-stone}.
Obraz $H$ w rzeczywistej przestrzeni $L^1(\mu)$
jest gęsty, zarówno dla $\sigma(H)$, jak i dla jej
uzupełnienia. Ponadto dla każdego zbioru mierzalnego
$A$ w jednej z tych $\sigma$-algebr istnieje ciąg
$(h_n)\subset H$ taki, że
\[
 0\le h_n\le1,\qquad
 \int_\Omega|h_n-\chi_A|\,\mathrm d\mu\longrightarrow0.
\]
\end{twierdzenie}
\begin{proof}
Najpierw pracujemy na $\sigma(H)$.
Niech $\mathcal C$ będzie rodziną zbiorów $A\in\sigma(H)$
spełniających drugą tezę. Rodzina ta zawiera
$\varnothing$ i $\Omega$, gdyż $0,1\in H$.
Jeżeli $h_n$ przybliża $\chi_A$, to $1-h_n\in H$
przybliża $\chi_{\Omega\setminus A}$ z tym samym błędem.
Jeżeli $h_n,k_n$ przybliżają odpowiednio
$\chi_A,\chi_B$, to $h_n\vee k_n\in H$, przyjmuje
wartości w $[0,1]$, a z nierówności
\[
 |\max(u,v)-\max(u',v')|\le |u-u'|+|v-v'|
\]
wynika
\[
 \|h_n\vee k_n-\chi_{A\cup B}\|_1
 \le\|h_n-\chi_A\|_1+\|k_n-\chi_B\|_1\to0.
\]
Nierówność dla maksimum wynika stąd, że jeśli
$e=\max(|u-u'|,|v-v'|)$, to
$\max(u',v')-e\le\max(u,v)\le\max(u',v')+e$.
Zatem $\mathcal C$ jest zamknięta na skończone sumy.

Dla $A_j\in\mathcal C$ połóżmy
$B_n=\bigcup_{j=1}^n A_j$ oraz $B=\bigcup_{j\ge1}A_j$.
Z ciągłości miary z dołu i skończoności $\mu$ mamy
$\|\chi_B-\chi_{B_n}\|_1=\mu(B\setminus B_n)\to0$.
Dla każdego $n$ wybieramy $h_n\in H$ z $0\le h_n\le1$
i $\|h_n-\chi_{B_n}\|_1<1/n$.
Nierówność trójkąta dowodzi, że $B\in\mathcal C$.
Rodzina $\mathcal C$ jest więc $\sigma$-algebrą.

Dla $h\in H$ i $t\in\mathbb{R}$ funkcje
$u_n=1\wedge n(h-t)^+$ należą do $H$, przyjmują
wartości w $[0,1]$ i rosną do $\chi_{\{h>t\}}$.
Z lematu \ref{it:lem:daniell-level} i reprezentacji
$I_0$ otrzymujemy ponadto dokładny błąd:
\[
 \|u_n-\chi_{\{h>t\}}\|_1
 =\mu(\{h>t\})-I_0(u_n)\longrightarrow0.
\]
Wszystkie składniki są skończone, gdyż $\mu(\Omega)<\infty$.
Stąd $\{h>t\}\in\mathcal C$, a te zbiory generują
$\sigma(H)$. Otrzymujemy $\mathcal C=\sigma(H)$.

Jeżeli $s=\sum_{j=1}^m a_j\chi_{A_j}$ jest rzeczywistą
mierzalną funkcją prostą, wybieramy przybliżenia
$h_{j,n}$ dla $\chi_{A_j}$ jak wyżej. Z liniowości
kraty $\sum_j a_j h_{j,n}\in H$, a błąd w $L^1$
jest nie większy niż $\sum_j |a_j|\|h_{j,n}-\chi_{A_j}\|_1$,
co dąży do zera. Funkcje proste są gęste w $L^1$ na mocy
twierdzenia \atlasref{ME-14::thm:simple-dense-lp-14}{ME-14: gęstość funkcji prostych}
(por. \atlasref{ME-14::thm:simple-dense-lp-14}{ME-14: gęstość funkcji prostych}).
Dla $f\in L^1$ wybieramy funkcje proste $s_n$
z $\|f-s_n\|_1<1/n$, a następnie $h_n\in H$
z $\|s_n-h_n\|_1<1/n$.
Stąd $\|f-h_n\|_1<2/n$.

Na uzupełnieniu miary każdy mierzalny zbiór $A$
różni się od pewnego $B\in\sigma(H)$ tylko na podzbiorze
mierzalnego zbioru zerowego. Przybliżenia funkcji
$\chi_B$ przybliżają też $\chi_A$ z tym samym błędem,
ponieważ $\chi_A=\chi_B$ prawie wszędzie. Argument dla funkcji prostych
i całkowalnych powtarzamy już na uzupełnieniu.
\end{proof}

\begin{wniosek}[Rozszerzenie ciągłe funkcjonału Daniella]
\label{it:cor:daniell-continuous-extension}
Przy założeniach twierdzenia \ref{it:thm:daniell-stone}
na $H$ określmy półnormę $p(h)=I_0(|h|)$.
Niech $N=\{h\in H:p(h)=0\}$. Uzupełnienie przestrzeni
normowanej $H/N$ względem normy indukowanej przez $p$
jest izometrycznie izomorficzne z rzeczywistą
przestrzenią $L^1(\mu)$.
Funkcjonał $I_0$ ma dokładnie jedno ciągłe liniowe
rozszerzenie na to uzupełnienie, dane wzorem
$I(f)=\int_\Omega f\,\mathrm d\mu$.
Jeżeli $\mu(\Omega)>0$, norma tego funkcjonału wynosi $1$;
jeśli $\mu(\Omega)=0$, przestrzeń i funkcjonał są zerowe.
\end{wniosek}
\begin{proof}
Z reprezentacji \ref{it:thm:daniell-stone} mamy
$p(h)=\int_\Omega |h|\,\mathrm d\mu=\|[h]\|_1$.
Własności półnormy i podprzestrzeni $N$ wynikają więc
z definicji \atlasref{ME-14::def:lp-14}{ME-14: definicja przestrzeni $L^1$}. Przyporządkowanie
$h+N\mapsto[h]$ jest dobrze określoną liniową izometrią:
$h-k\in N$ wtedy i tylko wtedy, gdy $h=k$ prawie
wszędzie. Jej obraz jest gęsty na mocy twierdzenia
\ref{it:thm:daniell-l1-density}, a $L^1$ jest zupełna na mocy twierdzenia
\atlasref{ME-14::thm:riesz-fischer-14}{ME-14: zupełność przestrzeni $L^1$}.
Stosujemy
\atlasref{FA-03::prop:completion-criterion}{FA-03: kryterium uzupełnienia}.
Ponieważ zanurzenie jest izometrią, utożsamienie
uzupełnienia również zachowuje normę: dla dowolnego
elementu uzupełnienia wybieramy zbieżny do niego
ciąg elementów $H/N$ i przechodzimy do granicy
w równości norm ich obrazów.

Z $- |h|\le h\le|h|$ i dodatniości wynika
$|I_0(h)|\le p(h)$, zatem $I_0$ znika na $N$.
Całka Lebesgue'a jest liniowa i ciągła na mocy
twierdzenia \atlasref{ME-08::thm:integral-linear-real}{ME-08: liniowość całki}, z normą
nie większą od $1$ i na obrazie $H$ jest równa $I_0$.
Jednoznaczność wynika z gęstości obrazu $H/N$
w $L^1$: dla dowolnego $f\in L^1$ i dowolnego ciągu
$h_n\in H$ spełniającego $\|[h_n]-f\|_1\to0$
każde ciągłe rozszerzenie musi przyjąć wartość
\[
 I(f)=\lim_{n\to\infty}I_0(h_n),
 \qquad
 |I(f)-I_0(h_n)|\le\|f-[h_n]\|_1.
\]
Oszacowanie wynika z ograniczenia modułu całki
przez całkę modułu. W szczególności dla $f\ge0$
można zastąpić $h_n$ przez $h_n^+$, gdyż
$|h_n^+-f|\le |h_n-f|$ prawie wszędzie.
Wówczas $I_0(h_n^+)\ge0$, co daje dodatniość
rozszerzenia również bez odwoływania się do jego
jawnego wzoru.
Gdy $\mu(\Omega)>0$, element
$[1]/\mu(\Omega)$ ma normę $1$ i całkę $1$.
Gdy $\mu(\Omega)=0$, każda funkcja całkowalna
jest zerowa prawie wszędzie, co daje ostatnią tezę.
\end{proof}

\section{Krata funkcji nieograniczonych zawierająca jedynkę}

\begin{twierdzenie}[Reprezentacja Daniella dla kraty z jedynką]
\label{it:thm:daniell-unbounded-unit}
Niech $H\subseteq\mathbb{R}^\Omega$ będzie kratą wektorową
funkcji o skończonych wartościach (niekoniecznie
ograniczonych) i niech $1\in H$. Niech $I_0:H\to\mathbb{R}$
będzie dodatnim liniowym funkcjonałem, ciągłym
porządkowo z góry w sensie definicji
\ref{it:def:daniell-functional}.
Wtedy istnieje dokładnie jedna skończona miara $\mu$
na $\sigma(H)$ taka, że wszystkie $h\in H$ są
bezwzględnie całkowalne i
\[
 I_0(h)=\int_\Omega h\,\mathrm d\mu,\qquad
 \mu(\Omega)=I_0(1).
\]
Miara zewnętrzna z definicji \ref{it:def:daniell-outer},
obliczona z użyciem całej kraty $H$, jest taka sama
jak miara zewnętrzna obliczona z użyciem jej
podkraty funkcji ograniczonych $H_b$.
W szczególności wniosek \ref{it:cor:daniell-completion}
pozostaje prawdziwy. Obraz $H$ jest gęsty w $L^1(\mu)$,
także po uzupełnieniu miary, a wniosek
\ref{it:cor:daniell-continuous-extension}
pozostaje prawdziwy.
\end{twierdzenie}
\begin{proof}
Niech $H_b=\{h\in H:\sup_{\omega\in\Omega}|h(\omega)|<\infty\}$.
Jest to krata wektorowa ograniczonych funkcji, zawierająca
$1$. Ograniczenie $I_b=I_0|_{H_b}$ jest dodatnie,
liniowe i ciągłe porządkowo z góry: każdy ciąg
występujący w warunku dla $H_b$ jest też ciągiem w $H$.
Twierdzenie \ref{it:thm:daniell-stone} daje więc
skończoną miarę $\mu$ na $\sigma(H_b)$.

Wykażemy $\sigma(H_b)=\sigma(H)$.
Jedna inkluzja wynika z $H_b\subseteq H$.
Dla każdego $h\in H$ i $n\ge1$ funkcja
$h^{(n)}=(-n)\vee(h\wedge n)$ należy do $H_b$.
Ponieważ $h$ ma skończone wartości, $h^{(n)}\to h$
punktowo. Stąd $h$ jest $\sigma(H_b)$-mierzalna,
co daje drugą inkluzję.

Dla $h\in H$, $h\ge0$, mamy $h\wedge n\in H_b$,
$h\wedge n\uparrow h$ oraz
$h-h\wedge n=(h-n)^+\downarrow0$.
Z ciągłości porządkowej i liniowości otrzymujemy
\[
 I_0(h\wedge n)=I_0(h)-I_0((h-n)^+)\uparrow I_0(h).
\]
Z twierdzenia \ref{thm:mct} oraz reprezentacji na $H_b$
wynika
\[
 \int_\Omega h\,\mathrm d\mu
 =\lim_n\int_\Omega(h\wedge n)\,\mathrm d\mu
 =\lim_n I_b(h\wedge n)=I_0(h)<\infty.
\]
Dla ogólnego $h\in H$ stosujemy wynik do $|h|,h^+,h^-$,
które należą do $H$, i odejmujemy skończone całki.
Jednoznaczność wynika z jednoznaczności reprezentacji
na $H_b$, ponieważ obie $\sigma$-algebry są równe.

Niech $\mu_H^*$ oraz $\mu_{H_b}^*$ oznaczają
miary zewnętrzne określone wzorem
\eqref{it:eq:daniell-outer} odpowiednio dla $H$ i $H_b$.
Z $H_b\subseteq H$ mamy $\mu_H^*(A)\le\mu_{H_b}^*(A)$
dla każdego $A\subseteq\Omega$.
Odwrotnie, każde pokrycie $\chi_A\le\sum_n h_n$
z $h_n\in H$, $h_n\ge0$, można zastąpić pokryciem
$\chi_A\le\sum_n(h_n\wedge1)$, jak w dowodzie lematu
\ref{it:lem:daniell-separated}. Wszystkie obcięte
funkcje należą do $H_b$, a monotoniczność $I_0$
zapewnia, że koszt nie wzrasta. Biorąc infimum
otrzymujemy $\mu_{H_b}^*(A)\le\mu_H^*(A)$.
Wobec równości obu miar zewnętrznych i
$\sigma(H_b)=\sigma(H)$ stosuje się wniosek
\ref{it:cor:daniell-completion}.

Gęstość obrazu $H_b$ wynika z twierdzenia
\ref{it:thm:daniell-l1-density}, zarówno dla
$\sigma(H)$, jak i jej uzupełnienia; obraz większej
kraty $H$ jest zatem również gęsty.
Ponadto każde $h\in H$ można przybliżyć jego
obcięciami $h^{(n)}\in H_b$ w półnormie Daniella:
z tożsamości $|h-h^{(n)}|=(|h|-n)^+$ wynika
\[
 p(h-h^{(n)})=I_0((|h|-n)^+)\downarrow0
\]
na mocy ciągłości porządkowej. W szczególności
$I_0(h^{(n)})\to I_0(h)$, gdyż
$|I_0(h)-I_0(h^{(n)})|\le p(h-h^{(n)})$.
Dowód wniosku \ref{it:cor:daniell-continuous-extension}
korzysta wyłącznie z reprezentacji, gęstości,
obecności jedynki i zupełności $L^1$, więc
działa bez zmian.
\end{proof}

\begin{uwaga}[Zakres wersji z jedynką]
Obecność $1$ i skończoność $I_0(1)$ wymuszają
skończoność miary reprezentującej. Nie obejmuje to
na przykład całki Lebesgue'a na całej $\mathbb{R}$ dla kraty
$C_c(\mathbb{R})$, która nie zawiera $1$. Wersja bez jedynki
wymaga osobnych założeń i osobnego dowodu.
\end{uwaga}

\section{Warunek Stone'a i reprezentacja bez jedynki}

\begin{twierdzenie}[Twierdzenie Daniella--Stone'a bez założenia skończoności miary]
\label{it:thm:daniell-stone-general}
Niech $H\subseteq\mathbb{R}^\Omega$ będzie kratą wektorową
funkcji o skończonych wartościach, spełniającą
\emph{warunek Stone'a}: $h\wedge1\in H$ dla $h\in H$,
$h\ge0$. Niech $I_0:H\to\mathbb{R}$ będzie dodatnim
liniowym funkcjonałem ciągłym porządkowo z góry.
Wówczas istnieje miara $\mu$ na $\sigma(H)$ taka, że
każda funkcja z $H$ jest bezwzględnie całkowalna i
\[
 I_0(h)=\int_\Omega h\,\mathrm d\mu,\qquad h\in H.
\]
Można wybrać $\mu$ jako ograniczenie miary zewnętrznej
\eqref{it:eq:daniell-outer}, przyjmując $\inf\varnothing=\infty$.
Tak wybrana miara jest największą miarą reprezentującą
na $\sigma(H)$ w porządku punktowym na zbiorach.
Jeśli ponadto istnieje ciąg $e_j\in H$, $e_j\ge0$, taki, że
\[
 \Omega=\bigcup_{j\ge1}\{e_j>0\},
\]
to miara jest $\sigma$-skończona i jednoznaczna.
Odwrotnie, $\sigma$-skończoność największej miary
reprezentującej wymusza istnienie takiego pokrycia.
\end{twierdzenie}
\begin{proof}
\emph{Krok 1. Obcięcia i miara zewnętrzna.}
Dla $h\in H$, $h\ge0$ i $c>0$ mamy
$h\wedge c=c((h/c)\wedge1)\in H$, a więc również
$(h-c)^+=h-h\wedge c\in H$; dla $c=0$ ostatnia
funkcja to $h$. Dowód lematu
\ref{it:lem:daniell-cover-domination} nie wymaga
jedynki ani ograniczoności funkcji: reszty
$(u-\sum_{n=1}^N h_n)^+$ nadal maleją do zera.
Dlatego dla $\mu^*$ określonej przez
\eqref{it:eq:daniell-outer} mamy
$I_0(u)\le\mu^*(A)$, jeśli $u\in H$ i
$0\le u\le\chi_A$.
Dowód własności miary zewnętrznej z lematu
\ref{it:lem:daniell-outer} działa z następującą
zmianą: przy podaddytywności wystarczy rozważać
rodziny $(A_j)$ z $\sum_j\mu^*(A_j)<\infty$,
gdyż w przeciwnym razie nierówność jest automatyczna.
W przypadku skończonej sumy wybieramy pokrycia
o kosztach mniejszych od $\mu^*(A_j)+\varepsilon2^{-j}$,
co jest możliwe nawet bez jedynki. Pokrycie zerowe
i monotoniczność pozostają bez zmian.

\emph{Krok 2. Rozdzielenie poziomów.}
Dla $h\in H$, $h\ge0$, $0\le s<t$ oraz
$A\subseteq\{h\le s\}$, $B\subseteq\{h\ge t\}$
pokażemy $\mu^*(A\cup B)=\mu^*(A)+\mu^*(B)$.
Wystarczy przypadek $\mu^*(A\cup B)<\infty$;
w przeciwnym przypadku podaddytywność daje,
że również suma po prawej jest nieskończona.
Stosujemy dowód lematu \ref{it:lem:daniell-separated}
z $v=1\wedge((h-s)^+/(t-s))$, obciętymi funkcjami
$k_n=h_n\wedge1$ i rozkładem
$k_n=(k_n-k_n\wedge v)+(k_n\wedge v)$.
Wszystkie te funkcje należą do $H$ dzięki krokowi 1,
warunkowi Stone'a i własności kraty. Dowód nie wymaga
ani obecności jedynki w $H$, ani ograniczoności
pierwotnych funkcji pokrywających.

\emph{Krok 3. Mierzalność.}
Dla $h\in H$, $h\ge0$ i $t\ge0$ stosujemy argument
warstw z lematu \ref{it:lem:daniell-measurability}
do dowolnego $E\subseteq\Omega$ o $\mu^*(E)<\infty$.
Wszystkie progi w tym argumencie są nieujemne,
więc krok 2 zapewnia wymagane rozdzielenie.
Sumy miar zewnętrznych warstw parzystych
i nieparzystych są ograniczone przez $\mu^*(E)$,
a ich ogony dążą do zera; otrzymujemy kryterium
Carathéodory'ego dla $\{h>t\}$.
Jeśli $\mu^*(E)=\infty$, nierówność
$\mu^*(E)\ge\mu^*(E\cap\{h>t\})+
\mu^*(E\setminus\{h>t\})$ jest automatyczna,
a przeciwna wynika z podaddytywności.
Zatem wszystkie zbiory $\{h>t\}$ dla $h\ge0$,
$t\ge0$, są mierzalne w sensie Carathéodory'ego.
Generują one $\sigma(H)$: każda nieujemna funkcja
z $H$ jest względem nich mierzalna, a dowolne
$h=h^+-h^-$ jest różnicą dwóch takich funkcji.
Z lematu \ref{lem:hausdorff-caratheodory-selfcontained-26}
wynika, że $\mu=\mu^*|_{\sigma(H)}$ jest miarą.

\emph{Krok 4. Reprezentacja.}
Dla $h\ge0$ w $H$ i $t\ge0$ funkcje
$u_n=1\wedge n(h-t)^+$ należą do $H$.
Dowód lematu \ref{it:lem:daniell-level} działa bez
zmian, bo opiera się jedynie na obcięciach,
oszacowaniu przez pokrycie i teleskopowaniu.
Otrzymujemy $\mu(\{h>t\})=\lim_n I_0(u_n)$,
także gdy granica jest nieskończona.
Dla $t>0$ miara ta jest skończona, ponieważ
$t^{-1}h\ge\chi_{\{h>t\}}$ stanowi dopuszczalne
pokrycie, a więc $\mu(\{h>t\})\le I_0(h)/t$.

Jeśli $h$ jest ograniczona, wybieramy $\delta>0$
oraz $m\in\mathbb{N}$, $h\le m\delta$ i definiujemy
$L_\delta=\delta\sum_{k=1}^m\chi_{\{h>k\delta\}}$.
Jak w dowodzie twierdzenia
\ref{it:thm:daniell-stone}, przybliżając każdy
składnik od dołu przez obcięcia z $H$, otrzymujemy
$\int_\Omega L_\delta\,\mathrm d\mu\le I_0(h)$.
Ponadto $0\le h-L_\delta\le\delta$.
Dla $\delta=1/n$ lemat Fatou
\ref{thm:fatou} daje $\int_\Omega h\,\mathrm d\mu
\le I_0(h)$, bez założenia skończoności $\mu(\Omega)$.

W drugą stronę rozkładamy $h$ na obcięte przyrosty,
wydzielając pierwszy składnik:
\[
 h=(h\wedge\delta)+
 \delta\sum_{k=1}^{m-1}
 \left(1\wedge\frac{(h-k\delta)^+}{\delta}\right).
\]
Każda funkcja w nawiasie nie przekracza
$\chi_{\{h>k\delta\}}$, więc oszacowanie z kroku 1
i liniowość dają
\[
 I_0(h)\le I_0(h\wedge\delta)+
 \int_\Omega L_\delta\,\mathrm d\mu
 \le I_0(h\wedge\delta)+\int_\Omega h\,\mathrm d\mu.
\]
Przy $\delta=1/n\downarrow0$ funkcje $h\wedge\delta$
maleją do zera w $H$, więc ich całki Daniella
dążą do zera. Otrzymujemy równość dla ograniczonych
$h\ge0$. Dla dowolnego $h\ge0$ stosujemy ją do
$h\wedge n$; reszty $h-h\wedge n$ maleją do zera,
a twierdzenie \ref{thm:mct} daje
$\int_\Omega h\,\mathrm d\mu=I_0(h)<\infty$.
Rozkład na $h^+,h^-$ daje reprezentację i bezwzględną
całkowalność każdej funkcji z $H$.

\emph{Krok 5. Maksymalność i jednoznaczność.}
Jeśli $\nu$ jest inną miarą reprezentującą na
$\sigma(H)$, dla każdego mierzalnego $A$ i każdego
pokrycia $\chi_A\le\sum_n h_n$ mamy
$\nu(A)\le\sum_n\int_\Omega h_n\,\mathrm d\nu
=\sum_n I_0(h_n)$ na mocy twierdzenia
\ref{thm:sum-nonnegative-integrals}.
Infimum daje $\nu(A)\le\mu(A)$.

Przy założeniu o $(e_j)$ zbiory
$D_{j,k}=\{e_j>1/k\}$, $j,k\ge1$, pokrywają $\Omega$,
a $\mu(D_{j,k})\le k I_0(e_j)<\infty$.
Niech $\mathcal P$ będzie rodziną skończonych przecięć
zbiorów $\{h>t\}$, $h\in H$, $h\ge0$, $t\ge0$,
z dołączonym $\Omega$. Jest to układ typu $\pi$
generujący $\sigma(H)$.
Dla przecięcia $P=\bigcap_{i=1}^m\{h_i>t_i\}$
co najmniej jednego generatora ($m\ge1$), także gdy
$P=\varnothing$, funkcje
\[
 v_n=1\wedge n\bigwedge_{i=1}^m(h_i-t_i)^+
\]
rosną do $\chi_P$, należą do $H$ i są nieujemne.
Dla $t_i>0$ mamy $(h_i-t_i)^+=h_i-h_i\wedge t_i$,
a dla $t_i=0$ otrzymujemy $h_i$. Skończone minimum
tych funkcji jest dodatnie dokładnie na $P$ i należy
do kraty; ostatnie obcięcie zapewnia warunek Stone'a.
Twierdzenie \ref{thm:mct} zastosowane do obu miar
daje $\mu(P)=\lim_n I_0(v_n)=\nu(P)$.
Nie potrzebujemy równości masy $\Omega$ na tym etapie.
Dla ustalonego $D=D_{j,k}$ miary
$A\mapsto\mu(A\cap D)$ i $A\mapsto\nu(A\cap D)$
są skończone i zgodne na $\mathcal P$, włącznie
z $A=\Omega$, ponieważ $A\cap D$ jest niepustym
skończonym przecięciem generatorów w sensie liczby
czynników (sam zbiór może być pusty).
Twierdzenie
\atlasref{ME-02::thm:uniqueness-finite}{ME-02: jednoznaczność miar skończonych}
daje zgodność tych ograniczeń na całej $\sigma(H)$.
Numerując zbiory $D_{j,k}$ i rozłączając je przez
odejmowanie sum wcześniejszych zbiorów, otrzymujemy
mierzalny podział $\Omega$ na zbiory zawarte
w odpowiednich $D_{j,k}$. Sumowanie równości
ograniczeń obu miar po tym podziale daje
$\mu=\nu$.

Odwrotnie, jeżeli największa miara $\mu$ jest
$\sigma$-skończona, wybieramy pokrycie
$\Omega=\bigcup_j A_j$, $A_j\in\sigma(H)$,
$\mu(A_j)<\infty$. Z definicji $\mu^*$ dla każdego
$j$ istnieje pokrycie $\chi_{A_j}\le\sum_n h_{j,n}$
funkcjami nieujemnymi z $H$ o skończonym koszcie.
Każdy punkt $A_j$ należy do pewnego
$\{h_{j,n}>0\}$, więc przeliczalna rodzina
$(h_{j,n})_{j,n\ge1}$ daje żądane pokrycie.
\end{proof}

\begin{wniosek}[Gęstość kraty bez jedynki i uzupełnienie miary]
\label{it:cor:daniell-general-density}
Niech $\mu$ będzie największą miarą reprezentującą
z twierdzenia \ref{it:thm:daniell-stone-general}.
Wówczas obraz $H$ jest gęsty w rzeczywistej
przestrzeni $L^1(\mu)$, także po uzupełnieniu miary,
bez założenia $\sigma$-skończoności.
Jeżeli $\mu$ jest $\sigma$-skończona, to
$\mathcal M_{\mu^*}$ jest dokładnie uzupełnieniem
$\sigma(H)$ względem $\mu$.
\end{wniosek}
\begin{proof}
Dla $A\in\sigma(H)$ z $\mu(A)<\infty$ i $\varepsilon>0$
wybieramy pokrycie $\chi_A\le\sum_n h_n$, gdzie
$0\le h_n\in H$ i $\sum_n I_0(h_n)<\mu(A)+\varepsilon$.
Funkcje $v_N=1\wedge\sum_{n=1}^N h_n$ należą do $H$
dzięki warunkowi Stone'a i rosną do mierzalnej funkcji $v$.
Mamy $\chi_A\le v\le\sum_n h_n$, zatem z reprezentacji
i twierdzenia \ref{thm:sum-nonnegative-integrals}
\[
 \|v-\chi_A\|_1=\int_\Omega v\,\mathrm d\mu-\mu(A)
 <\varepsilon.
\]
Ponieważ $\int v\,\mathrm d\mu<\infty$, z twierdzenia
\ref{thm:mct} i liniowości całki otrzymujemy
$\|v-v_N\|_1=\int(v-v_N)\,\mathrm d\mu\to0$.
Dla dostatecznie dużego $N$ mamy więc
$\|v_N-\chi_A\|_1<2\varepsilon$.
Liniowość $H$ oraz twierdzenie
\atlasref{ME-14::thm:simple-dense-lp-14}{ME-14: gęstość funkcji prostych} o gęstości funkcji prostych
o zbiorze niezerowości skończonej miary dają gęstość $H$.
Na uzupełnieniu korzystamy z tego, że każdy zbiór
mierzalny skończonej miary jest równy prawie wszędzie
zbiorowi z $\sigma(H)$ o tej samej mierze.

Dla dowolnego $A\subseteq\Omega$ o $\mu^*(A)<\infty$
argument z wniosku \ref{it:cor:daniell-completion}
daje otoczkę $B\in\sigma(H)$ z
$A\subseteq B$ i $\mu(B)=\mu^*(A)$:
wybieramy pokrycia z błędami $2^{-k}$,
a następnie przecinamy zbiory, na których ich sumy
są co najmniej równe $1$. Nie wymaga to skończoności
$\mu(\Omega)$, lecz jedynie skończoności $\mu^*(A)$.
Jeżeli dodatkowo $A$ jest mierzalny w sensie
Carathéodory'ego, to $\mu^*(B\setminus A)=0$,
a więc $A$ należy do uzupełnienia, jak w cytowanym
wniosku. Przy $\sigma$-skończoności wybieramy
mierzalne pokrycie $\Omega=\bigcup_j E_j$ o
$\mu(E_j)<\infty$. Dla $A\in\mathcal M_{\mu^*}$
stosujemy właśnie wykazany wynik do $A\cap E_j$
i bierzemy ich przeliczalną sumę.
Przeciwna inkluzja wynika z zupełności miary
Carathéodory'ego.
\end{proof}

\begin{uwaga}[Rozszerzenie funkcjonału bez jedynki]
Dla największej miary reprezentującej $\mu$ wniosek
\ref{it:cor:daniell-continuous-extension} o
utożsamieniu uzupełnienia $H/\ker p$ z $L^1(\mu)$
oraz jednoznacznym ciągłym rozszerzeniu $I_0$
obowiązuje również bez jedynki, z $p(h)=I_0(|h|)$.
Dowód korzysta z gęstości wykazanej powyżej.
Norma rozszerzenia jest równa $1$, jeśli
$L^1(\mu)\ne\{0\}$: wybieramy niezerowy element
$f\in L^1$ i stosujemy funkcjonał do
$|f|/\|f\|_1$. W przeciwnym przypadku norma
wynosi zero. Bez $\sigma$-skończoności nie wolno
zastępować warunku $L^1(\mu)\ne\{0\}$ samym
$\mu\ne0$, co ilustruje poniższy przykład.
\end{uwaga}

\begin{przyklad}[Bez przeliczalnego pokrycia jednoznaczność może zawodzić]
\label{it:ex:daniell-nonunique}
Niech $\Omega$ będzie nieprzeliczalny, a $H$ niech
składa się ze wszystkich rzeczywistych funkcji
o skończonym nośniku. Warunek Stone'a jest spełniony.
Funkcjonał $I_0(h)=0$ jest dodatni, liniowy i ciągły
porządkowo z góry. $\sigma(H)$ jest $\sigma$-algebrą
zbiorów przeliczalnych i dopełnień zbiorów przeliczalnych:
funkcje z $H$ są względem niej mierzalne, a funkcje
charakterystyczne singletonów generują tę $\sigma$-algebrę.
Zarówno miara zerowa, jak i miara $\nu$ przypisująca
zbiorom przeliczalnym masę zero, a ich dopełnieniom
masę jeden, reprezentują $I_0$.
Istotnie, w rozłącznej przeliczalnej rodzinie
zbiorów mierzalnych najwyżej jeden może mieć
przeliczalne dopełnienie; jeśli żaden go nie ma,
suma zbiorów jest przeliczalna, w przeciwnym razie
jej dopełnienie jest przeliczalne.
To dowodzi przeliczalnej addytywności $\nu$.
Każda funkcja z $H$ znika $\nu$-prawie wszędzie.
W konstrukcji przez pokrycia $\mu^*(A)=0$ dla
$A$ przeliczalnego: wystarczy pokrycie funkcjami
charakterystycznymi jego singletonów, każda o koszcie
zero (dla zbioru pustego używamy pokrycia zerowego).
Dla $A$ nieprzeliczalnego nie istnieje dopuszczalne
pokrycie, ponieważ przeliczalna suma skończonych
nośników jest przeliczalna. Stąd $\mu^*(A)=\infty$.
Jest to inna, największa miara reprezentująca.
Jej przestrzeń $L^1(\mu)$ jest zerowa mimo $\mu\ne0$:
dla całkowalnej $f$ każdy zbiór $\{|f|>1/k\}$
ma skończoną miarę, a zatem jest przeliczalny
i ma miarę zero. Ich przeliczalna suma to zbiór
niezerowości $f$, więc $f=0$ prawie wszędzie.
Nie jest ona $\sigma$-skończona: każdy zbiór jej
skończonej miary jest przeliczalny, więc przeliczalna
suma takich zbiorów nie pokrywa $\Omega$.
Obie pozostałe miary reprezentujące są jednak skończone.
Sama $\sigma$-skończoność pewnej miary reprezentującej
nie zastępuje zatem warunku przeliczalnego pokrycia
z twierdzenia \ref{it:thm:daniell-stone-general}.
Gęstość z wniosku \ref{it:cor:daniell-general-density}
również zależy od wyboru miary: dla miary $\nu$
obraz $H$ w $L^1(\nu)$ jest zerowy, podczas gdy
$\|[1]\|_1=1$. Nie można więc zastąpić największej
miary dowolną miarą reprezentującą w tym wniosku.
\end{przyklad}

\begin{przyklad}[Znaczenie warunku Stone'a dla jednoznaczności]
\label{it:ex:daniell-without-stone}
Niech $\Omega=(0,1]$, $e(x)=x$ i
$H=\{ae:a\in\mathbb{R}\}$. Jest to krata wektorowa,
a $I_0(ae)=a$ jest dodatnim liniowym funkcjonałem
ciągłym porządkowo z góry: jeśli $a_ne\downarrow0$
punktowo, to $a_n\downarrow0$ przez ewaluację w $1$.
Mamy $\{e>0\}=\Omega$ i $\sigma(H)=\mathcal B(\Omega)$,
ale $H$ nie spełnia warunku Stone'a, ponieważ
$(2e)\wedge1$ nie jest wielokrotnością $e$.
Dla każdego $t\in(0,1]$ miara $t^{-1}\delta_t$
reprezentuje $I_0$, gdyż
$\int_\Omega ax\,\mathrm d(t^{-1}\delta_t)(x)=a$.
Miary te są skończone i parami różne. Samo pokrycie
przez zbiór niezerowości nie wystarcza zatem do
jednoznaczności bez warunku Stone'a.
\end{przyklad}

\begin{przyklad}[Zastosowanie do $C_c(X)$]
\label{it:ex:daniell-cc}
Niech $X$ będzie lokalnie zwartą, $\sigma$-zwartą
przestrzenią metryczną, a $H=C_c(X;\mathbb{R})$.
Krata spełnia warunek Stone'a, gdyż dla $h\ge0$
obcięcie $h\wedge1$ jest ciągłe i znika poza nośnikiem $h$.
Każdy dodatni liniowy funkcjonał $I_0$ na $H$
jest ciągły porządkowo z góry. Jeśli $h_n\downarrow0$,
wszystkie $h_n$ znikają poza zwartym nośnikiem $K$
funkcji $h_1$. Z twierdzenia Diniego na $K$ wynika
$\|h_n\|_\infty\to0$. Istotnie, dla $\varepsilon>0$
względnie domknięte zbiory
$K_n=\{x\in K:h_n(x)\ge\varepsilon\}$ maleją
i mają puste przecięcie. Zwartość $K$ daje
$K_n=\varnothing$ dla pewnego $n$. Od tego indeksu
$h_n<\varepsilon$ na $K$, a poza $K$ funkcje są zerowe. Wybieramy funkcję odcinającą
$e\in C_c(X)$ równą $1$ na $K$ i nieujemną
(por. lemat \ref{lem:cutoff-radon-20}); wtedy
$0\le h_n\le\|h_n\|_\infty e$, więc
$0\le I_0(h_n)\le\|h_n\|_\infty I_0(e)\to0$.
Jeśli $K$ jest pusty, wszystkie $h_n$ są zerowe.

Dla wyczerpania $X$ przeliczalną rodziną zbiorów
zwartych wybieramy analogiczne funkcje odcinające
$e_j$; zbiory $\{e_j>0\}$ pokrywają $X$.
Ponadto $\sigma(C_c(X))=\mathcal B(X)$:
jedna inkluzja wynika z ciągłości. Przestrzeń
$\sigma$-zwarta metryczna jest ośrodkowa, bo każdy
zwarty zbiór metryczny ma przeliczalny zbiór gęsty.
Jest więc drugiej przeliczalności i każdy jej otwarty
podzbiór jest przestrzenią Lindelöfa.
Dla każdego otwartego $U$ i każdego $x\in U$
wybieramy $e_x\in C_c(X)$ o wartościach w $[0,1]$,
równą $1$ w $x$ i z nośnikiem zawartym w $U$
(z tego samego lematu o funkcji odcinającej).
Zbiory $\{e_x>0\}$ tworzą otwarte pokrycie $U$,
z którego wybieramy przeliczalne podpokrycie.
Zatem $U\in\sigma(C_c(X))$.
Twierdzenie \ref{it:thm:daniell-stone-general}
daje jednoznaczną $\sigma$-skończoną miarę borelowską
reprezentującą $I_0$. Dla każdego zwartego $K$
jej masa jest nie większa od $I_0(e)$ dla funkcji
odcinającej równej $1$ na $K$, czyli jest skończona.
Miara jest lokalnie skończona, ponieważ każdy punkt
ma otoczenie o zwartym domknięciu. Z
\atlasref{ME-20::thm:lch-secondcountable-radon-20}{ME-20: lokalnie skończone miary borelowskie}
wynika, że jest miarą Radona. Otrzymujemy w ten sposób
reprezentację z \atlasref{ME-21::thm:riesz-positive-21}{ME-21: twierdzenie Riesza--Markowa--Kakutaniego}.
Założenia lokalnej zwartości i ośrodkowości w wersji
metrycznej są równoważne lokalnej zwartości i
$\sigma$-zwartości: ośrodkowa przestrzeń metryczna
jest przestrzenią Lindelöfa, więc z pokrycia otwartymi
otoczeniami o zwartych domknięciach wybieramy
przeliczalne podpokrycie. Odwrotną implikację wykazano
wyżej. Nie powtarzamy dowodu regularności miary
z ME-20; dalsze własności całki Radona omawiamy
w rozdziale \ref{ch:radon-measures-20}.
\end{przyklad}

\section{Przykłady i konieczność założeń}
\begin{przyklad}[Od całki Riemanna do miary Lebesgue'a]
Niech $\Omega=[0,1]$, $H=C[0,1]$ i $I_0(h)=\int_0^1h(x)\,\mathrm dx$ w sensie Riemanna. Jeśli $h_n\downarrow0$ i $h_n\in H$, z twierdzenia Diniego wynika zbieżność jednostajna, więc $I_0(h_n)\downarrow0$. Mamy $\sigma(H)=\mathcal B([0,1])$: wszystkie funkcje ciągłe są borelowskie, a funkcja identycznościowa należy do $H$ i generuje $\mathcal B([0,1])$. Ograniczenie $\lambda_1$ do zbiorów borelowskich reprezentuje $I_0$ dzięki twierdzeniu \atlasref{ME-08::thm:riemann-to-lebesgue}{ME-08: zgodność całek Riemanna i Lebesgue'a}. Jednoznaczność z twierdzenia \ref{it:thm:daniell-stone} daje równość otrzymanej miary z tym ograniczeniem. Po uzupełnieniu jest zwykłą miarą Lebesgue'a.
\end{przyklad}
\begin{przyklad}[Dodatniość nie zastępuje ciągłości porządkowej]
Niech $\Omega=\mathbb{N}$, $H=c$ będzie kratą ciągów zbieżnych, a $I_0(h)=\lim_{k\to\infty}h(k)$. Jest to dodatni funkcjonał liniowy. Dla $h_n(k)=\chi_{\{j\in\mathbb{N}:j>n\}}(k)$ mamy $h_n\downarrow0$ punktowo, ale $I_0(h_n)=1$ dla wszystkich $n$. Ponieważ funkcje $\chi_{\{k\}}$ należą do $c$, mamy $\sigma(H)=2^\mathbb{N}$. Nie ma reprezentującej miary przeliczalnie addytywnej na $\mathbb{N}$: musiałaby przypisywać zerową masę każdemu punktowi, a jednocześnie całej $\mathbb{N}$ masę $1$.
\end{przyklad}

% ============================================================
% UZUPELNIENIE MATEMATYCZNE: ROZDZIAL 12
% ============================================================
\subsection{Dlaczego ciągłość porządkowa jest konieczna}
\begin{przyklad}[Funkcjonał dodatni niepochodzący od miary na dziedzinie]\label{prz:daniel-koniecznosc}
Niech $X=(0,1]$, a $H$ będzie kratą funkcji na $X$ będących ograniczeniami funkcji z $C[0,1]$. Każda $h\in H$ ma jednoznaczne ciągłe przedłużenie $\widetilde h$ do $[0,1]$. Definiujemy dodatni funkcjonał liniowy $I(h)=\widetilde h(0)$. Dla $h_n(t)=\max\{1-nt,0\}$, $t\in X$, zachodzi $h_n\downarrow0$ punktowo na $X$, lecz $I(h_n)=1$ dla każdego $n$. Funkcjonał nie jest zatem ciągły porządkowo.

Gdyby istniała miara $\mu$ na $X$ reprezentująca $I$ dla wszystkich $h\in H$, mielibyśmy $\mu(X)=I(1)=1$. Z twierdzenia o zbieżności monotonicznej zastosowanego do $1-h_n\uparrow1$ wynikałoby $\int_Xh_n\,\mathrm d\mu\downarrow0$, sprzecznie z $I(h_n)=1$. Zatem dodatniość i liniowość bez ciągłości porządkowej nie wystarczają do reprezentacji przez miarę na danej przestrzeni.

Po dołączeniu punktu $0$ funkcjonał na $C[0,1]$
jest reprezentowany przez $\delta_0$. Te same
funkcje $h_n(t)=\max(1-nt,0)$ nie maleją wtedy
punktowo do zera, lecz do $\chi_{\{0\}}$.
Ich infimum w porządku wewnętrznym kraty $C[0,1]$
jest jednak zerem: każda ciągła funkcja będąca
ograniczeniem dolnym jest niedodatnia na $(0,1]$,
a przez ciągłość również w $0$. To pokazuje, dlaczego
w definicji \ref{it:def:daniell-functional} wymagamy
zbieżności punktowej, a nie istnienia zerowego
infimum w samej kracie.
\end{przyklad}
\begin{przyklad}[Jawna konstrukcja miary z funkcjonału Daniella]\label{prz:daniell-atoms}
Na $\Omega=[0,1]$ weźmy kratę $H=C[0,1]$ zawierającą funkcję stałą $1$ oraz funkcjonał
$I(h)=\sum_{n=1}^\infty 2^{-n}h(1/n)$. Szereg jest bezwzględnie zbieżny dla każdego $h\in H$, ponieważ $|I(h)|\le\|h\|_\infty\sum_{n\ge1}2^{-n}=\|h\|_\infty$. Funkcjonał $I$ jest dodatni i liniowy. Jeśli $h_k\downarrow0$ punktowo w $H$, to twierdzenie Diniego daje $\|h_k\|_\infty\to0$, a więc $0\le I(h_k)\le\|h_k\|_\infty\to0$. Funkcjonał jest zatem ciągły porządkowo.

Miara reprezentująca $I$ ma postać $\mu=\sum_{n=1}^\infty 2^{-n}\delta_{1/n}$. Przeliczalna addytywność $\mu$ wynika z możliwości zamiany dwóch szeregów nieujemnych. Dla funkcji nieujemnej wzór na całkę względem $\mu$ otrzymujemy najpierw dla funkcji prostych, a następnie z twierdzeń \atlasref{ME-06::thm:simple-up}{ME-06: aproksymacja funkcjami prostymi} i \ref{thm:mct}, zamieniając rosnącą granicę z sumą nieujemną. Dla ograniczonej $h$ stosujemy wzór do $h^+$ i $h^-$; oba szeregi mają skończoną sumę. Stąd $I(h)=\int_{[0,1]}h\,\mathrm d\mu$, zaś $\mu([0,1])=\sum_n2^{-n}=1$. Jednoznaczność wynika z twierdzenia
\ref{it:thm:daniell-stone}, ponieważ
$\sigma(C[0,1])=\mathcal B([0,1])$.
Chociaż punkty $1/n$ gromadzą się w $0$, mamy
$\mu(\{0\})=0$. Punkt $0$ należy jednak do nośnika
topologicznego miary: każde jego względne otoczenie
zawiera pewien punkt $1/n$ o dodatniej masie.
Dokładnie $\operatorname{supp}\mu=\{0\}\cup\{1/n:n\ge1\}$,
gdyż dopełnienie tego domkniętego zbioru jest otwarte
i ma miarę zero (por.
\atlasref{ME-20::def:support-radon-20}{ME-20: nośnik miary}). Przykład pokazuje, że konstrukcja Daniella--Stone'a obejmuje miary czysto atomowe bez rozpoczynania od miary zewnętrznej na przedziałach.
\end{przyklad}
\begin{uwaga}[Trzy różne poziomy konstrukcji]
Trzeba oddzielić: (i) sprawdzenie dodatniości i liniowości $I$ na kracie, (ii) ciągłość porządkową $h_n\downarrow0\Rightarrow I(h_n)\downarrow0$ oraz (iii) rozszerzenie funkcjonału i konstrukcję miary na odpowiedniej $\sigma$-algebrze. Punkt (ii) jest warunkiem umożliwiającym przeliczalną addytywność. Samo zdefiniowanie $\mu(A)=I(\chi_A)$ nie wystarcza, jeżeli funkcje charakterystyczne zbiorów $A$ nie należą do początkowej kraty $H$.
\end{uwaga}


\backmatter
% Indeks generowany na podstawie \begin{definicja|twierdzenie|lemat|wniosek}.
\chapter*{Indeks definicji i twierdzeń}
\markboth{INDEKS DEFINICJI I TWIERDZEŃ}{INDEKS DEFINICJI I TWIERDZEŃ}
\phantomsection\addcontentsline{toc}{chapter}{Indeks definicji i twierdzeń}
Indeks podaje nazwę rezultatu, jego numer i stronę. Nazwy i numery są aktywnymi odsyłaczami.
Uwzględniono także lematy i wnioski.

\section*{Definicje}
\begingroup\small
\setlength\LTleft{0pt}\setlength\LTright{0pt}
\begin{longtable}{@{}p{0.72\linewidth}p{0.135\linewidth}p{0.08\linewidth}@{}}
\textbf{Nazwa} & \textbf{Nr} & \textbf{Str.}\\\hline
\endfirsthead
\multicolumn{3}{@{}l}{\emph{Definicje -- ciąg dalszy}}\\[2pt]
\textbf{Nazwa} & \textbf{Nr} & \textbf{Str.}\\\hline
\endhead
\hyperref[def:atom]{Atom i miara bezatomowa} & \ref*{def:atom} & \pageref*{def:atom} \\
\hyperref[def:ac-measure-17]{Bezwzględna ciągłość miary} & \ref*{def:ac-measure-17} & \pageref*{def:ac-measure-17} \\
\hyperref[def:reduced-boundary-29]{Brzeg zredukowany} & \ref*{def:reduced-boundary-29} & \pageref*{def:reduced-boundary-29} \\
\hyperref[def:signed-integral-over-set-7]{Całka funkcji całkowalnej po zbiorze} & \ref*{def:signed-integral-over-set-7} & \pageref*{def:signed-integral-over-set-7} \\
\hyperref[def:integral-simple]{Całka funkcji prostej} & \ref*{def:integral-simple} & \pageref*{def:integral-simple} \\
\hyperref[def:complex-integral]{Całka funkcji zespolonej} & \ref*{def:complex-integral} & \pageref*{def:complex-integral} \\
\hyperref[def:integral-positive]{Całka Lebesgue'a funkcji nieujemnej} & \ref*{def:integral-positive} & \pageref*{def:integral-positive} \\
\hyperref[def:integral-over-set]{Całka po podzbiorze} & \ref*{def:integral-over-set} & \pageref*{def:integral-over-set} \\
\hyperref[def:tight-single-20]{Ciasność pojedynczej miary} & \ref*{def:tight-single-20} & \pageref*{def:tight-single-20} \\
\hyperref[def:cauchy-measure-9]{Ciąg Cauchy'ego według miary} & \ref*{def:cauchy-measure-9} & \pageref*{def:cauchy-measure-9} \\
\hyperref[def:positive-negative-parts]{Części dodatnia i ujemna} & \ref*{def:positive-negative-parts} & \pageref*{def:positive-negative-parts} \\
\hyperref[def:vitali-fine-24]{Drobne pokrycie Vitaliego} & \ref*{def:vitali-fine-24} & \pageref*{def:vitali-fine-24} \\
\hyperref[def:dual-lp-15]{Dualność i norma funkcjonału} & \ref*{def:dual-lp-15} & \pageref*{def:dual-lp-15} \\
\hyperref[def:HL-max-24]{Funkcja maksymalna Hardy'ego--Littlewooda} & \ref*{def:HL-max-24} & \pageref*{def:HL-max-24} \\
\hyperref[def:simple-function]{Funkcja prosta} & \ref*{def:simple-function} & \pageref*{def:simple-function} \\
\hyperref[idx:06:definicja:11]{Funkcja zespolona mierzalna} & \ref*{idx:06:definicja:11} & \pageref*{idx:06:definicja:11} \\
\hyperref[def:limsup-sets]{Granica górna i dolna zbiorów} & \ref*{def:limsup-sets} & \pageref*{def:limsup-sets} \\
\hyperref[def:topological-group-23]{Grupa topologiczna} & \ref*{def:topological-group-23} & \pageref*{def:topological-group-23} \\
\hyperref[def:normal-jacobian-28]{Jakobian normalny} & \ref*{def:normal-jacobian-28} & \pageref*{def:normal-jacobian-28} \\
\hyperref[def:volume-jacobian-27]{Jakobian objętościowy} & \ref*{def:volume-jacobian-27} & \pageref*{def:volume-jacobian-27} \\
\hyperref[def:k-jacobian-28]{Jakobian wymiaru $k$} & \ref*{def:k-jacobian-28} & \pageref*{def:k-jacobian-28} \\
\hyperref[def:uniform-integrability]{Jednostajna całkowalność} & \ref*{def:uniform-integrability} & \pageref*{def:uniform-integrability} \\
\hyperref[def:uniform-tight-22]{Jednostajna ciasność} & \ref*{def:uniform-tight-22} & \pageref*{def:uniform-tight-22} \\
\hyperref[def:monotone-class]{Klasa monotoniczna} & \ref*{def:monotone-class} & \pageref*{def:monotone-class} \\
\hyperref[def:measure-finiteness]{Klasy miar} & \ref*{def:measure-finiteness} & \pageref*{def:measure-finiteness} \\
\hyperref[def:hausdorff-measure-26]{Koszt pokrycia i miara Hausdorffa} & \ref*{def:hausdorff-measure-26} & \pageref*{def:hausdorff-measure-26} \\
\hyperref[def:linf-14]{Kres górny istotny i przestrzeń $L^\infty$} & \ref*{def:linf-14} & \pageref*{def:linf-14} \\
\hyperref[def:multiplicity-28]{Krotność odwzorowania} & \ref*{def:multiplicity-28} & \pageref*{def:multiplicity-28} \\
\hyperref[def:balls-24]{Kule i krotność pokrycia} & \ref*{def:balls-24} & \pageref*{def:balls-24} \\
\hyperref[def:haar-23]{Lewostronna i prawostronna miara Haara} & \ref*{def:haar-23} & \pageref*{def:haar-23} \\
\hyperref[def:local-majorant-11]{Lokalna dominacja całkowalna} & \ref*{def:local-majorant-11} & \pageref*{def:local-majorant-11} \\
\hyperref[def:max-measure-25]{Maksymalna funkcja miary} & \ref*{def:max-measure-25} & \pageref*{def:max-measure-25} \\
\hyperref[def:measure]{Miara} & \ref*{def:measure} & \pageref*{def:measure} \\
\hyperref[def:borel-locfinite-20]{Miara borelowska i lokalna skończoność} & \ref*{def:borel-locfinite-20} & \pageref*{def:borel-locfinite-20} \\
\hyperref[def:pushforward]{Miara obrazowa} & \ref*{def:pushforward} & \pageref*{def:pushforward} \\
\hyperref[def:radon-measure-20]{Miara Radona --- przyjęta konwencja} & \ref*{def:radon-measure-20} & \pageref*{def:radon-measure-20} \\
\hyperref[def:restriction]{Miara śladowa i ograniczenie} & \ref*{def:restriction} & \pageref*{def:restriction} \\
\hyperref[def:vector-measure-30]{Miara wektorowa} & \ref*{def:vector-measure-30} & \pageref*{def:vector-measure-30} \\
\hyperref[def:complex-measure-16]{Miara zespolona} & \ref*{def:complex-measure-16} & \pageref*{def:complex-measure-16} \\
\hyperref[def:outer-measure]{Miara zewnętrzna} & \ref*{def:outer-measure} & \pageref*{def:outer-measure} \\
\hyperref[def:lebesgue-outer]{Miara zewnętrzna Lebesgue'a} & \ref*{def:lebesgue-outer} & \pageref*{def:lebesgue-outer} \\
\hyperref[def:metric-outer-21]{Miara zewnętrzna metryczna} & \ref*{def:metric-outer-21} & \pageref*{def:metric-outer-21} \\
\hyperref[def:complete-measure]{Miara zupełna} & \ref*{def:complete-measure} & \pageref*{def:complete-measure} \\
\hyperref[def:extended-measurable]{Mierzalność funkcji rozszerzonej} & \ref*{def:extended-measurable} & \pageref*{def:extended-measurable} \\
\hyperref[def:normal-BV-29]{Normalna z pochodnej miary} & \ref*{def:normal-BV-29} & \pageref*{def:normal-BV-29} \\
\hyperref[def:support-radon-20]{Nośnik miary} & \ref*{def:support-radon-20} & \pageref*{def:support-radon-20} \\
\hyperref[def:finite-perimeter-29]{Obwód w sensie De Giorgiego} & \ref*{def:finite-perimeter-29} & \pageref*{def:finite-perimeter-29} \\
\hyperref[def:measurable-map]{Odwzorowanie mierzalne} & \ref*{def:measurable-map} & \pageref*{def:measurable-map} \\
\hyperref[def:algebry]{Pierścień, algebra, $\sigma$-algebra} & \ref*{def:algebry} & \pageref*{def:algebry} \\
\hyperref[def:five-convergences]{Pięć rodzajów zbieżności} & \ref*{def:five-convergences} & \pageref*{def:five-convergences} \\
\hyperref[def:RN-derivative-18]{Pochodna Radona--Nikodyma} & \ref*{def:RN-derivative-18} & \pageref*{def:RN-derivative-18} \\
\hyperref[def:current-29]{Prąd i jego brzeg} & \ref*{def:current-29} & \pageref*{def:current-29} \\
\hyperref[def:premeasure]{Premiara} & \ref*{def:premeasure} & \pageref*{def:premeasure} \\
\hyperref[def:halfopen-box]{Prostopadłościan półotwarty} & \ref*{def:halfopen-box} & \pageref*{def:halfopen-box} \\
\hyperref[def:complex-measures-space-16]{Przestrzeń $\mathcal M(X,\mathcal F)$} & \ref*{def:complex-measures-space-16} & \pageref*{def:complex-measures-space-16} \\
\hyperref[def:standard-borel]{Przestrzeń polska i standardowa borelowska} & \ref*{def:standard-borel} & \pageref*{def:standard-borel} \\
\hyperref[def:measure-boundary-29]{Punkty gęstości i brzeg miarowy} & \ref*{def:measure-boundary-29} & \pageref*{def:measure-boundary-29} \\
\hyperref[def:weak-lp-15]{Refleksywność i zbieżność słaba} & \ref*{def:weak-lp-15} & \pageref*{def:weak-lp-15} \\
\hyperref[def:regularity-20]{Regularność względem topologii} & \ref*{def:regularity-20} & \pageref*{def:regularity-20} \\
\hyperref[def:ae-functions]{Równość i zbieżność prawie wszędzie} & \ref*{def:ae-functions} & \pageref*{def:ae-functions} \\
\hyperref[def:lp-14]{Równość prawie wszędzie i przestrzeń $L^p$} & \ref*{def:lp-14} & \pageref*{def:lp-14} \\
\hyperref[def:frechet-rademacher-28]{Różniczkowalność Frécheta} & \ref*{def:frechet-rademacher-28} & \pageref*{def:frechet-rademacher-28} \\
\hyperref[def:borel]{$\sigma$-algebra borelowska} & \ref*{def:borel} & \pageref*{def:borel} \\
\hyperref[def:sigma-generated]{$\sigma$-algebra generowana} & \ref*{def:sigma-generated} & \pageref*{def:sigma-generated} \\
\hyperref[def:sigma-of-function]{$\sigma$-algebra generowana przez odwzorowanie} & \ref*{def:sigma-of-function} & \pageref*{def:sigma-of-function} \\
\hyperref[def:product-sigma]{$\sigma$-algebra produktowa} & \ref*{def:product-sigma} & \pageref*{def:product-sigma} \\
\hyperref[idx:01:definicja:04]{$\sigma$-algebra śladowa} & \ref*{idx:01:definicja:04} & \pageref*{idx:01:definicja:04} \\
\hyperref[idx:03:definicja:09]{\texorpdfstring{$\sigma$}{sigma}-skończona premiara} & \ref*{idx:03:definicja:09} & \pageref*{idx:03:definicja:09} \\
\hyperref[def:bochner-30]{Silna mierzalność i całka Bochnera} & \ref*{def:bochner-30} & \pageref*{def:bochner-30} \\
\hyperref[def:signed-measure-16]{Skończona miara ze znakiem} & \ref*{def:signed-measure-16} & \pageref*{def:signed-measure-16} \\
\hyperref[idx:21:definicja:08]{Skończona zespolona miara Radona} & \ref*{idx:21:definicja:08} & \pageref*{idx:21:definicja:08} \\
\hyperref[def:cantor-set]{Środkowotrzeci zbiór Cantora} & \ref*{def:cantor-set} & \pageref*{def:cantor-set} \\
\hyperref[def:test-19]{Test losowy i jego moc} & \ref*{def:test-19} & \pageref*{def:test-19} \\
\hyperref[def:energy-26]{$t$-energia miary} & \ref*{def:energy-26} & \pageref*{def:energy-26} \\
\hyperref[def:pi-lambda]{Układ $\pi$ i układ Dynkina} & \ref*{def:pi-lambda} & \pageref*{def:pi-lambda} \\
\hyperref[def:variation-signed-16]{Wariacja całkowita} & \ref*{def:variation-signed-16} & \pageref*{def:variation-signed-16} \\
\hyperref[def:BV-variation-29]{Wariacja funkcji} & \ref*{def:BV-variation-29} & \pageref*{def:BV-variation-29} \\
\hyperref[def:vector-variation-30]{Wariacja i półwariacja} & \ref*{def:vector-variation-30} & \pageref*{def:vector-variation-30} \\
\hyperref[def:lebesgue-point-25]{Wartość średnia i punkt Lebesgue'a} & \ref*{def:lebesgue-point-25} & \pageref*{def:lebesgue-point-25} \\
\hyperref[def:conditional-exp-18]{Warunkowa wartość oczekiwana} & \ref*{def:conditional-exp-18} & \pageref*{def:conditional-exp-18} \\
\hyperref[def:almost-everywhere]{Własność prawie wszędzie} & \ref*{def:almost-everywhere} & \pageref*{def:almost-everywhere} \\
\hyperref[def:RNP-30]{Własność Radona--Nikodyma} & \ref*{def:RNP-30} & \pageref*{def:RNP-30} \\
\hyperref[def:hausdorff-dimension-26]{Wymiar Hausdorffa} & \ref*{def:hausdorff-dimension-26} & \pageref*{def:hausdorff-dimension-26} \\
\hyperref[def:singular-16]{Wzajemna osobliwość} & \ref*{def:singular-16} & \pageref*{def:singular-16} \\
\hyperref[def:equivalent-measures-17]{Wzajemna równoważność miar} & \ref*{def:equivalent-measures-17} & \pageref*{def:equivalent-measures-17} \\
\hyperref[def:vague-weakstar-22]{Zbieżność słaba-* i mglista} & \ref*{def:vague-weakstar-22} & \pageref*{def:vague-weakstar-22} \\
\hyperref[def:narrow-measures-22]{Zbieżność wąska} & \ref*{def:narrow-measures-22} & \pageref*{def:narrow-measures-22} \\
\hyperref[def:convergence-measure]{Zbieżność według miary} & \ref*{def:convergence-measure} & \pageref*{def:convergence-measure} \\
\hyperref[def:lp-convergence-9]{Zbieżność w średniej rzędu $p$} & \ref*{def:lp-convergence-9} & \pageref*{def:lp-convergence-9} \\
\hyperref[def:rectifiable-29]{Zbiory przeliczalnie prostowalne} & \ref*{def:rectifiable-29} & \pageref*{def:rectifiable-29} \\
\hyperref[def:continuity-set-22]{Zbiór ciągłości miary} & \ref*{def:continuity-set-22} & \pageref*{def:continuity-set-22} \\
\hyperref[def:positive-negative-16]{Zbiór dodatni i ujemny} & \ref*{def:positive-negative-16} & \pageref*{def:positive-negative-16} \\
\hyperref[def:caratheodory-measurable]{Zbiór mierzalny w sensie Carathéodory'ego} & \ref*{def:caratheodory-measurable} & \pageref*{def:caratheodory-measurable} \\
\end{longtable}
\endgroup

\section*{Twierdzenia}
\begingroup\small
\setlength\LTleft{0pt}\setlength\LTright{0pt}
\begin{longtable}{@{}p{0.72\linewidth}p{0.135\linewidth}p{0.08\linewidth}@{}}
\textbf{Nazwa} & \textbf{Nr} & \textbf{Str.}\\\hline
\endfirsthead
\multicolumn{3}{@{}l}{\emph{Twierdzenia -- ciąg dalszy}}\\[2pt]
\textbf{Nazwa} & \textbf{Nr} & \textbf{Str.}\\\hline
\endhead
\hyperref[thm:positive-additivity]{Addytywność całki nieujemnej} & \ref*{thm:positive-additivity} & \pageref*{thm:positive-additivity} \\
\hyperref[thm:cc-approx-in-measure-10]{Aproksymacja funkcji o nośniku mierzalnym skończonej miary} & \ref*{thm:cc-approx-in-measure-10} & \pageref*{thm:cc-approx-in-measure-10} \\
\hyperref[thm:simple-real-approx]{Aproksymacja funkcji o wartościach rzeczywistych} & \ref*{thm:simple-real-approx} & \pageref*{thm:simple-real-approx} \\
\hyperref[thm:lebesgue-outer-regular]{Aproksymacja otwartymi zbiorami} & \ref*{thm:lebesgue-outer-regular} & \pageref*{thm:lebesgue-outer-regular} \\
\hyperref[thm:lebesgue-inner-regular]{Aproksymacja zbiorami zwartymi} & \ref*{thm:lebesgue-inner-regular} & \pageref*{thm:lebesgue-inner-regular} \\
\hyperref[thm:finite-metric-regular-20]{Aproksymacja zbiorów borelowskich} & \ref*{thm:finite-metric-regular-20} & \pageref*{thm:finite-metric-regular-20} \\
\hyperref[thm:atom-sigma-finite]{Atomy miar $\sigma$-skończonych mają skończoną masę} & \ref*{thm:atom-sigma-finite} & \pageref*{thm:atom-sigma-finite} \\
\hyperref[thm:bayes-finite-19]{Bayes dla skończonej przestrzeni parametrów} & \ref*{thm:bayes-finite-19} & \pageref*{thm:bayes-finite-19} \\
\hyperref[thm:image-ac-19]{Bezwzględna ciągłość zachowana przez obraz} & \ref*{thm:image-ac-19} & \pageref*{thm:image-ac-19} \\
\hyperref[thm:lebesgue-completion]{Borelowskie aproksymacje i uzupełnienie} & \ref*{thm:lebesgue-completion} & \pageref*{thm:lebesgue-completion} \\
\hyperref[thm:borel-Rd]{Borelowskie prostokąty} & \ref*{thm:borel-Rd} & \pageref*{thm:borel-Rd} \\
\hyperref[thm:moving-domains-11]{Całkowanie po zależnej od parametru dziedzinie} & \ref*{thm:moving-domains-11} & \pageref*{thm:moving-domains-11} \\
\hyperref[thm:change-measure-19]{Całkowanie po zmianie miary} & \ref*{thm:change-measure-19} & \pageref*{thm:change-measure-19} \\
\hyperref[thm:sum-absolutely-integrable]{Całkowanie szeregów absolutnie sumowalnych} & \ref*{thm:sum-absolutely-integrable} & \pageref*{thm:sum-absolutely-integrable} \\
\hyperref[thm:pushforward-integral-13]{Całkowanie względem miary obrazowej} & \ref*{thm:pushforward-integral-13} & \pageref*{thm:pushforward-integral-13} \\
\hyperref[thm:product-projections]{Charakteryzacja przez współrzędne} & \ref*{thm:product-projections} & \pageref*{thm:product-projections} \\
\hyperref[thm:polish-tight-20]{Ciasność skończonych miar na przestrzeniach polskich} & \ref*{thm:polish-tight-20} & \pageref*{thm:polish-tight-20} \\
\hyperref[thm:parameter-continuity]{Ciągłość pod znakiem całki} & \ref*{thm:parameter-continuity} & \pageref*{thm:parameter-continuity} \\
\hyperref[thm:translation-lp-14]{Ciągłość translacji w $L^p(\mathbb R^d)$} & \ref*{thm:translation-lp-14} & \pageref*{thm:translation-lp-14} \\
\hyperref[thm:measure-continuity]{Ciągłość z dołu i z góry} & \ref*{thm:measure-continuity} & \pageref*{thm:measure-continuity} \\
\hyperref[thm:atomic-cont-singular-17]{Dalszy rozkład miary osobliwej na prostej} & \ref*{thm:atomic-cont-singular-17} & \pageref*{thm:atomic-cont-singular-17} \\
\hyperref[thm:weak-lsc-15]{Dolna półciągłość normy} & \ref*{thm:weak-lsc-15} & \pageref*{thm:weak-lsc-15} \\
\hyperref[thm:perimeter-lsc-29]{Dolna półciągłość obwodu} & \ref*{thm:perimeter-lsc-29} & \pageref*{thm:perimeter-lsc-29} \\
\hyperref[thm:dual-lp-15]{Dualność $L^p$ dla $1<p<\infty$} & \ref*{thm:dual-lp-15} & \pageref*{thm:dual-lp-15} \\
\hyperref[thm:dual-l1-15]{Dualność $L^1$ przy $\sigma$-skończoności} & \ref*{thm:dual-l1-15} & \pageref*{thm:dual-l1-15} \\
\hyperref[thm:tv-dual-tests-16]{Dualny opis normy w wariacji całkowitej} & \ref*{thm:tv-dual-tests-16} & \pageref*{thm:tv-dual-tests-16} \\
\hyperref[thm:ui-two-criteria-8]{Dwa równoważne opisy jednostajnej całkowalności} & \ref*{thm:ui-two-criteria-8} & \pageref*{thm:ui-two-criteria-8} \\
\hyperref[thm:measurable-algebra]{Działania algebraiczne} & \ref*{thm:measurable-algebra} & \pageref*{thm:measurable-algebra} \\
\hyperref[thm:total-variation-formula-19]{Formuła wariacyjna dla dwóch prawdopodobieństw} & \ref*{thm:total-variation-formula-19} & \pageref*{thm:total-variation-formula-19} \\
\hyperref[thm:frostman-26]{Frostmana, wersja dla zbiorów zwartych} & \ref*{thm:frostman-26} & \pageref*{thm:frostman-26} \\
\hyperref[thm:fubini-12]{Fubini} & \ref*{thm:fubini-12} & \pageref*{thm:fubini-12} \\
\hyperref[thm:cdf-derivative-25]{Funkcja rozkładu a absolutna ciągłość miary} & \ref*{thm:cdf-derivative-25} & \pageref*{thm:cdf-derivative-25} \\
\hyperref[thm:measurable-real]{Funkcje rzeczywiste i działania} & \ref*{thm:measurable-real} & \pageref*{thm:measurable-real} \\
\hyperref[thm:pushforward-composition]{Funktorialność miar obrazowych} & \ref*{thm:pushforward-composition} & \pageref*{thm:pushforward-composition} \\
\hyperref[thm:cc-dense-lp-radon-20]{Gęstość $C_c(X)$ w przestrzeniach $L^p$} & \ref*{thm:cc-dense-lp-radon-20} & \pageref*{thm:cc-dense-lp-radon-20} \\
\hyperref[thm:cc-in-lp-10]{Gęstość funkcji gładkich} & \ref*{thm:cc-in-lp-10} & \pageref*{thm:cc-in-lp-10} \\
\hyperref[thm:cc-dense-lp-14]{Gęstość funkcji gładkich w przestrzeni euklidesowej} & \ref*{thm:cc-dense-lp-14} & \pageref*{thm:cc-dense-lp-14} \\
\hyperref[thm:simple-dense-lp-14]{Gęstość funkcji prostych} & \ref*{thm:simple-dense-lp-14} & \pageref*{thm:simple-dense-lp-14} \\
\hyperref[thm:pushforward-coarea-28]{Gęstość obrazu przez submersję} & \ref*{thm:pushforward-coarea-28} & \pageref*{thm:pushforward-coarea-28} \\
\hyperref[thm:smooth-perimeter-29]{Granica klasycznej geometrii} & \ref*{thm:smooth-perimeter-29} & \pageref*{thm:smooth-perimeter-29} \\
\hyperref[thm:haar-main-23]{Haara: istnienie i jednoznaczność} & \ref*{thm:haar-main-23} & \pageref*{thm:haar-main-23} \\
\hyperref[thm:laplace-holo-11]{Holomorficzność transformaty Laplace'a} & \ref*{thm:laplace-holo-11} & \pageref*{thm:laplace-holo-11} \\
\hyperref[thm:holder-14]{Hölder} & \ref*{thm:holder-14} & \pageref*{thm:holder-14} \\
\hyperref[thm:finiteness-implications]{Implikacje} & \ref*{thm:finiteness-implications} & \pageref*{thm:finiteness-implications} \\
\hyperref[thm:lebesgue-exists]{Istnienie i podstawowe własności} & \ref*{thm:lebesgue-exists} & \pageref*{thm:lebesgue-exists} \\
\hyperref[thm:vitali]{Istnienie zbioru niemierzalnego} & \ref*{thm:vitali} & \pageref*{thm:vitali} \\
\hyperref[thm:CE-existence-18]{Istnienie, jednoznaczność i własności podstawowe} & \ref*{thm:CE-existence-18} & \pageref*{thm:CE-existence-18} \\
\hyperref[thm:BV-layercake-29]{Jedna nierówność formuły koobszaru dla $BV$} & \ref*{thm:BV-layercake-29} & \pageref*{thm:BV-layercake-29} \\
\hyperref[thm:uniform-convexity-15]{Jednostajna wypukłość} & \ref*{thm:uniform-convexity-15} & \pageref*{thm:uniform-convexity-15} \\
\hyperref[thm:haar-unique-23]{Jednoznaczność miary Haara} & \ref*{thm:haar-unique-23} & \pageref*{thm:haar-unique-23} \\
\hyperref[thm:uniqueness-cc-20]{Jednoznaczność miary Radona przez całki} & \ref*{thm:uniqueness-cc-20} & \pageref*{thm:uniqueness-cc-20} \\
\hyperref[thm:uniqueness-finite]{Jednoznaczność miar skończonych} & \ref*{thm:uniqueness-finite} & \pageref*{thm:uniqueness-finite} \\
\hyperref[thm:uniqueness-sigma-finite]{Jednoznaczność przy wspólnym wyczerpaniu} & \ref*{thm:uniqueness-sigma-finite} & \pageref*{thm:uniqueness-sigma-finite} \\
\hyperref[thm:extension-uniqueness]{Jednoznaczność rozszerzenia} & \ref*{thm:extension-uniqueness} & \pageref*{thm:extension-uniqueness} \\
\hyperref[thm:egorov-9]{Jegorow} & \ref*{thm:egorov-9} & \pageref*{thm:egorov-9} \\
\hyperref[thm:jensen-14]{Jensen dla miary probabilistycznej} & \ref*{thm:jensen-14} & \pageref*{thm:jensen-14} \\
\hyperref[thm:hausdorff-calibration-26]{Kalibracja wymiaru całkowitego bez twierdzeń geometrycznych} & \ref*{thm:hausdorff-calibration-26} & \pageref*{thm:hausdorff-calibration-26} \\
\hyperref[thm:bochner-basic-30]{Konstrukcja i oszacowanie całki Bochnera} & \ref*{thm:bochner-basic-30} & \pageref*{thm:bochner-basic-30} \\
\hyperref[thm:product-measure-12]{Konstrukcja miary produktowej} & \ref*{thm:product-measure-12} & \pageref*{thm:product-measure-12} \\
\hyperref[thm:outer-from-cover]{Konstrukcja przez pokrycia} & \ref*{thm:outer-from-cover} & \pageref*{thm:outer-from-cover} \\
\hyperref[thm:local-to-global-measure-9]{Kontrola ucieczki masy przy zbieżności lokalnej} & \ref*{thm:local-to-global-measure-9} & \pageref*{thm:local-to-global-measure-9} \\
\hyperref[thm:pointwise-limits]{Kresy i granice} & \ref*{thm:pointwise-limits} & \pageref*{thm:pointwise-limits} \\
\hyperref[thm:integrability-criterion]{Kryterium absolutnej całkowalności} & \ref*{thm:integrability-criterion} & \pageref*{thm:integrability-criterion} \\
\hyperref[thm:premeasure-continuity]{Kryterium ciągłości premiary skończonej} & \ref*{thm:premeasure-continuity} & \pageref*{thm:premeasure-continuity} \\
\hyperref[thm:eps-delta-ac-17]{Kryterium epsilon--delta} & \ref*{thm:eps-delta-ac-17} & \pageref*{thm:eps-delta-ac-17} \\
\hyperref[thm:tails-11]{Kryterium jednostajnej kontroli ogonów} & \ref*{thm:tails-11} & \pageref*{thm:tails-11} \\
\hyperref[thm:riemann-criterion]{Kryterium Lebesgue'a całkowalności Riemanna} & \ref*{thm:riemann-criterion} & \pageref*{thm:riemann-criterion} \\
\hyperref[thm:measurable-generator]{Kryterium przez generatory} & \ref*{thm:measurable-generator} & \pageref*{thm:measurable-generator} \\
\hyperref[thm:weak-lp-15]{Kryterium słabej zbieżności w $L^p$} & \ref*{thm:weak-lp-15} & \pageref*{thm:weak-lp-15} \\
\hyperref[thm:param-vitali-11]{Kryterium Vitaliego dla parametru} & \ref*{thm:param-vitali-11} & \pageref*{thm:param-vitali-11} \\
\hyperref[thm:series-lp-14]{Kryterium zbieżności szeregów w $L^p$} & \ref*{thm:series-lp-14} & \pageref*{thm:series-lp-14} \\
\hyperref[thm:lebesgue-diff-25]{Lebesgue'a o różniczkowaniu} & \ref*{thm:lebesgue-diff-25} & \pageref*{thm:lebesgue-diff-25} \\
\hyperref[thm:borel-cantelli-1]{Lemat Borela--Cantellego, część pierwsza} & \ref*{thm:borel-cantelli-1} & \pageref*{thm:borel-cantelli-1} \\
\hyperref[thm:fatou]{Lemat Fatou} & \ref*{thm:fatou} & \pageref*{thm:fatou} \\
\hyperref[thm:fatou-integrable-lower-8]{Lemat Fatou z całkowalnym ograniczeniem od dołu} & \ref*{thm:fatou-integrable-lower-8} & \pageref*{thm:fatou-integrable-lower-8} \\
\hyperref[thm:neyman-pearson-19]{Lemat Neymana--Pearsona} & \ref*{thm:neyman-pearson-19} & \pageref*{thm:neyman-pearson-19} \\
\hyperref[thm:5r-24]{Lemat o pięciokrotnym powiększeniu} & \ref*{thm:5r-24} & \pageref*{thm:5r-24} \\
\hyperref[thm:scheffe-9]{Lemat Scheffégo} & \ref*{thm:scheffe-9} & \pageref*{thm:scheffe-9} \\
\hyperref[thm:integral-linear-real]{Liniowość całki funkcji całkowalnych} & \ref*{thm:integral-linear-real} & \pageref*{thm:integral-linear-real} \\
\hyperref[thm:integral-complex-linear]{Liniowość zespolona i oszacowanie modułu} & \ref*{thm:integral-complex-linear} & \pageref*{thm:integral-complex-linear} \\
\hyperref[thm:linear-change-13]{Liniowy wzór na objętość} & \ref*{thm:linear-change-13} & \pageref*{thm:linear-change-13} \\
\hyperref[thm:local-convergence-measure]{Lokalna zbieżność według miary} & \ref*{thm:local-convergence-measure} & \pageref*{thm:local-convergence-measure} \\
\hyperref[thm:local-continuity-11]{Lokalne kryterium ciągłości} & \ref*{thm:local-continuity-11} & \pageref*{thm:local-continuity-11} \\
\hyperref[thm:null-local-27]{Lokalne kryterium zbiorów zerowych} & \ref*{thm:null-local-27} & \pageref*{thm:null-local-27} \\
\hyperref[thm:lch-secondcountable-radon-20]{Lokalnie skończone miary borelowskie} & \ref*{thm:lch-secondcountable-radon-20} & \pageref*{thm:lch-secondcountable-radon-20} \\
\hyperref[thm:lusin-10]{Łuzin} & \ref*{thm:lusin-10} & \pageref*{thm:lusin-10} \\
\hyperref[thm:measure-metric-9]{Metryka zbieżności według miary} & \ref*{thm:measure-metric-9} & \pageref*{thm:measure-metric-9} \\
\hyperref[thm:graph-area-28]{Miara na gładkim wykresie} & \ref*{thm:graph-area-28} & \pageref*{thm:graph-area-28} \\
\hyperref[thm:riemann-radon-27]{Miara objętościowa jest miarą Radona} & \ref*{thm:riemann-radon-27} & \pageref*{thm:riemann-radon-27} \\
\hyperref[thm:product-riemann-27]{Miara produktu riemannowskiego} & \ref*{thm:product-riemann-27} & \pageref*{thm:product-riemann-27} \\
\hyperref[thm:rn-manifold-27]{Miary absolutnie ciągłe względem objętości} & \ref*{thm:rn-manifold-27} & \pageref*{thm:rn-manifold-27} \\
\hyperref[thm:restriction-meas-5]{Mierzalność ograniczenia} & \ref*{thm:restriction-meas-5} & \pageref*{thm:restriction-meas-5} \\
\hyperref[thm:minimize-lp-15]{Minimalizacja wypukłej energii} & \ref*{thm:minimize-lp-15} & \pageref*{thm:minimize-lp-15} \\
\hyperref[thm:initial-sigma]{Minimalna struktura mierzalna} & \ref*{thm:initial-sigma} & \pageref*{thm:initial-sigma} \\
\hyperref[thm:minkowski-14]{Minkowski} & \ref*{thm:minkowski-14} & \pageref*{thm:minkowski-14} \\
\hyperref[thm:null-modifications]{Modyfikacje na zbiorach zerowych} & \ref*{thm:null-modifications} & \pageref*{thm:null-modifications} \\
\hyperref[thm:simple-up]{Monotoniczna aproksymacja funkcji nieujemnej} & \ref*{thm:simple-up} & \pageref*{thm:simple-up} \\
\hyperref[thm:fatou-for-sets]{Nierówności dla granic zbiorów} & \ref*{thm:fatou-for-sets} & \pageref*{thm:fatou-for-sets} \\
\hyperref[thm:lp-implies-measure-9]{Nierówność Czebyszewa i zbieżność według miary} & \ref*{thm:lp-implies-measure-9} & \pageref*{thm:lp-implies-measure-9} \\
\hyperref[thm:lp-interpolation-14]{Nierówność interpolacyjna} & \ref*{thm:lp-interpolation-14} & \pageref*{thm:lp-interpolation-14} \\
\hyperref[thm:markov-nonnegative]{Nierówność Markowa dla całki} & \ref*{thm:markov-nonnegative} & \pageref*{thm:markov-nonnegative} \\
\hyperref[thm:HL-weak-24]{Nierówność słabego typu $(1,1)$} & \ref*{thm:HL-weak-24} & \pageref*{thm:HL-weak-24} \\
\hyperref[thm:gibbs-19]{Nieujemność entropii względnej} & \ref*{thm:gibbs-19} & \pageref*{thm:gibbs-19} \\
\hyperref[thm:ae-integral-invariance]{Niezmienniczość całki przy zmianie na zbiorze zerowym} & \ref*{thm:ae-integral-invariance} & \pageref*{thm:ae-integral-invariance} \\
\hyperref[thm:support-fullmass-20]{Nośnik na przestrzeni drugiej przeliczalności} & \ref*{thm:support-fullmass-20} & \pageref*{thm:support-fullmass-20} \\
\hyperref[thm:riesz-c0-positive-21]{Ograniczone dodatnie funkcjonały na $C_0$} & \ref*{thm:riesz-c0-positive-21} & \pageref*{thm:riesz-c0-positive-21} \\
\hyperref[thm:extended-operations]{Operacje na funkcjach rozszerzonych} & \ref*{thm:extended-operations} & \pageref*{thm:extended-operations} \\
\hyperref[thm:partition-variation-16]{Opis wariacji przez podziały} & \ref*{thm:partition-variation-16} & \pageref*{thm:partition-variation-16} \\
\hyperref[thm:bounded-integral-complex-16]{Oszacowanie całki względem miary zespolonej} & \ref*{thm:bounded-integral-complex-16} & \pageref*{thm:bounded-integral-complex-16} \\
\hyperref[thm:frostman-energy-26]{Oszacowanie energii przez warunek Frostmana} & \ref*{thm:frostman-energy-26} & \pageref*{thm:frostman-energy-26} \\
\hyperref[thm:weak-measure-25]{Oszacowanie słabego typu dla miar} & \ref*{thm:weak-measure-25} & \pageref*{thm:weak-measure-25} \\
\hyperref[thm:mct]{O zbieżności monotonicznej; Beppo Levi} & \ref*{thm:mct} & \pageref*{thm:mct} \\
\hyperref[thm:dct]{O zbieżności zmajoryzowanej; Lebesgue} & \ref*{thm:dct} & \pageref*{thm:dct} \\
\hyperref[thm:BV-coarea-full-29]{Pełna formuła koobszaru dla $BV$} & \ref*{thm:BV-coarea-full-29} & \pageref*{thm:BV-coarea-full-29} \\
\hyperref[thm:local-differentiation-11]{Pierwsza pochodna — wersja lokalna} & \ref*{thm:local-differentiation-11} & \pageref*{thm:local-differentiation-11} \\
\hyperref[thm:nonnegative-first-properties]{Pierwsze własności całki nieujemnej} & \ref*{thm:nonnegative-first-properties} & \pageref*{thm:nonnegative-first-properties} \\
\hyperref[thm:BV-measure-29]{Pochodna $Du$ jako miara wektorowa} & \ref*{thm:BV-measure-29} & \pageref*{thm:BV-measure-29} \\
\hyperref[thm:RN-inverse-18]{Pochodna odwrotna} & \ref*{thm:RN-inverse-18} & \pageref*{thm:RN-inverse-18} \\
\hyperref[thm:higher-differentiation-11]{Pochodne wyższego rzędu} & \ref*{thm:higher-differentiation-11} & \pageref*{thm:higher-differentiation-11} \\
\hyperref[thm:weak-subseq-15]{Podciąg słabo zbieżny: wariant separowalny} & \ref*{thm:weak-subseq-15} & \pageref*{thm:weak-subseq-15} \\
\hyperref[thm:measure-subsequence]{Podciąg zbieżny prawie wszędzie} & \ref*{thm:measure-subsequence} & \pageref*{thm:measure-subsequence} \\
\hyperref[thm:conv-implications-9]{Podstawowe implikacje} & \ref*{thm:conv-implications-9} & \pageref*{thm:conv-implications-9} \\
\hyperref[thm:FTC-lebesgue-25]{Podstawowe twierdzenie rachunku całkowego dla całki Lebesgue'a} & \ref*{thm:FTC-lebesgue-25} & \pageref*{thm:FTC-lebesgue-25} \\
\hyperref[thm:basic-measure]{Podstawowe własności miary} & \ref*{thm:basic-measure} & \pageref*{thm:basic-measure} \\
\hyperref[thm:cantor-properties]{Podstawowe własności zbioru Cantora} & \ref*{thm:cantor-properties} & \pageref*{thm:cantor-properties} \\
\hyperref[thm:vitali-cover-24]{Pokryciowe twierdzenie Vitaliego} & \ref*{thm:vitali-cover-24} & \pageref*{thm:vitali-cover-24} \\
\hyperref[thm:portmanteau-22]{Portmanteau} & \ref*{thm:portmanteau-22} & \pageref*{thm:portmanteau-22} \\
\hyperref[thm:haar-modular-23]{Prawa miara Haara i modularność} & \ref*{thm:haar-modular-23} & \pageref*{thm:haar-modular-23} \\
\hyperref[thm:ae-smooth-approx-10]{Prawie wszędzie gładka aproksymacja funkcji mierzalnej} & \ref*{thm:ae-smooth-approx-10} & \pageref*{thm:ae-smooth-approx-10} \\
\hyperref[thm:ae-to-measure-finite]{Prawie wszędzie implikuje zbieżność według miary przy mierze skończonej} & \ref*{thm:ae-to-measure-finite} & \pageref*{thm:ae-to-measure-finite} \\
\hyperref[thm:volume-premeasure]{Premiara na ograniczonym prostopadłościanie} & \ref*{thm:volume-premeasure} & \pageref*{thm:volume-premeasure} \\
\hyperref[thm:prokhorov-22]{Prochorow na przestrzeni polskiej} & \ref*{thm:prokhorov-22} & \pageref*{thm:prokhorov-22} \\
\hyperref[thm:prokhorov-forward-22]{Prochorow: kierunek ciasność $\Rightarrow$ zwartość} & \ref*{thm:prokhorov-forward-22} & \pageref*{thm:prokhorov-forward-22} \\
\hyperref[thm:borel-product]{Produkt borelowski przestrzeni drugoprzeliczalnych} & \ref*{thm:borel-product} & \pageref*{thm:borel-product} \\
\hyperref[thm:lebesgue-products-12]{Produkt miar Lebesgue'a} & \ref*{thm:lebesgue-products-12} & \pageref*{thm:lebesgue-products-12} \\
\hyperref[thm:compact-extension-10]{Przedłużenie z kompaktu w $\mathbb R^d$} & \ref*{thm:compact-extension-10} & \pageref*{thm:compact-extension-10} \\
\hyperref[thm:intersection-sigma]{Przekroje $\sigma$-algebr} & \ref*{thm:intersection-sigma} & \pageref*{thm:intersection-sigma} \\
\hyperref[thm:borel-generators-R]{Przeliczalne generatory na $\mathbb R$} & \ref*{thm:borel-generators-R} & \pageref*{thm:borel-generators-R} \\
\hyperref[thm:standard-countably-generated]{Przeliczalne generowanie standardowych przestrzeni} & \ref*{thm:standard-countably-generated} & \pageref*{thm:standard-countably-generated} \\
\hyperref[thm:hilbert-rnp-30]{Przestrzenie Hilberta mają RNP} & \ref*{thm:hilbert-rnp-30} & \pageref*{thm:hilbert-rnp-30} \\
\hyperref[thm:CE-orthogonal-18]{Przypadek $L^2$: rzut ortogonalny} & \ref*{thm:CE-orthogonal-18} & \pageref*{thm:CE-orthogonal-18} \\
\hyperref[thm:simple-properties]{Rachunek dla funkcji prostych} & \ref*{thm:simple-properties} & \pageref*{thm:simple-properties} \\
\hyperref[thm:rademacher-28]{Rademacher} & \ref*{thm:rademacher-28} & \pageref*{thm:rademacher-28} \\
\hyperref[thm:RN-finite-18]{Radon--Nikodym dla miar skończonych} & \ref*{thm:RN-finite-18} & \pageref*{thm:RN-finite-18} \\
\hyperref[thm:RN-complex-18]{Radon--Nikodym dla miar zespolonych} & \ref*{thm:RN-complex-18} & \pageref*{thm:RN-complex-18} \\
\hyperref[thm:reflexive-lp-15]{Refleksywność $L^p$} & \ref*{thm:reflexive-lp-15} & \pageref*{thm:reflexive-lp-15} \\
\hyperref[thm:hausdorff-gdelta-26]{Regularizacja przez zbiory $G_\delta$} & \ref*{thm:hausdorff-gdelta-26} & \pageref*{thm:hausdorff-gdelta-26} \\
\hyperref[thm:RN-chain-18]{Reguła łańcuchowa} & \ref*{thm:RN-chain-18} & \pageref*{thm:RN-chain-18} \\
\hyperref[thm:riesz-positive-21]{Reprezentacja dodatnich funkcjonałów} & \ref*{thm:riesz-positive-21} & \pageref*{thm:riesz-positive-21} \\
\hyperref[thm:compact-riesz-21]{Reprezentacja na zwartej przestrzeni metrycznej} & \ref*{thm:compact-riesz-21} & \pageref*{thm:compact-riesz-21} \\
\hyperref[thm:riesz-complex-21]{Reprezentacja zespolona i norma wariacji} & \ref*{thm:riesz-complex-21} & \pageref*{thm:riesz-complex-21} \\
\hyperref[thm:riesz-fischer-14]{Riesz--Fischer: zupełność $L^p$} & \ref*{thm:riesz-fischer-14} & \pageref*{thm:riesz-fischer-14} \\
\hyperref[thm:rnp-lipschitz-equiv-30]{RNP a różniczkowalność krzywych Lipschitzowskich} & \ref*{thm:rnp-lipschitz-equiv-30} & \pageref*{thm:rnp-lipschitz-equiv-30} \\
\hyperref[thm:polar-complex-measure-18]{Rozkład biegunowy miary zespolonej} & \ref*{thm:polar-complex-measure-18} & \pageref*{thm:polar-complex-measure-18} \\
\hyperref[thm:hahn-16]{Rozkład Hahna} & \ref*{thm:hahn-16} & \pageref*{thm:hahn-16} \\
\hyperref[thm:jordan-16]{Rozkład Jordana} & \ref*{thm:jordan-16} & \pageref*{thm:jordan-16} \\
\hyperref[thm:lebesgue-decomp-finite-17]{Rozkład Lebesgue'a dla miar skończonych} & \ref*{thm:lebesgue-decomp-finite-17} & \pageref*{thm:lebesgue-decomp-finite-17} \\
\hyperref[thm:lebesgue-decomp-sigfin-17]{Rozkład Lebesgue'a w przypadku $\sigma$-skończonym} & \ref*{thm:lebesgue-decomp-sigfin-17} & \pageref*{thm:lebesgue-decomp-sigfin-17} \\
\hyperref[thm:leb-manifold-27]{Rozkład na część objętościową i osobliwą} & \ref*{thm:leb-manifold-27} & \pageref*{thm:leb-manifold-27} \\
\hyperref[thm:extension-caratheodory]{Rozszerzenie Carathéodory'ego} & \ref*{thm:extension-caratheodory} & \pageref*{thm:extension-caratheodory} \\
\hyperref[thm:measure-differentiation-25]{Różniczkowanie miary względem Lebesgue'a} & \ref*{thm:measure-differentiation-25} & \pageref*{thm:measure-differentiation-25} \\
\hyperref[thm:differentiating-integrals]{Różniczkowanie pod znakiem całki} & \ref*{thm:differentiating-integrals} & \pageref*{thm:differentiating-integrals} \\
\hyperref[thm:hausdorff-borel-26]{Samodzielna konstrukcja miary Hausdorffa} & \ref*{thm:hausdorff-borel-26} & \pageref*{thm:hausdorff-borel-26} \\
\hyperref[thm:separable-lp-14]{Separowalność $L^p$} & \ref*{thm:separable-lp-14} & \pageref*{thm:separable-lp-14} \\
\hyperref[thm:hausdorff-lipschitz-26]{Skalowanie i odwzorowania Lipschitzowskie} & \ref*{thm:hausdorff-lipschitz-26} & \pageref*{thm:hausdorff-lipschitz-26} \\
\hyperref[thm:besicovitch-finite-24]{Skończona postać lematu Besicovitcha} & \ref*{thm:besicovitch-finite-24} & \pageref*{thm:besicovitch-finite-24} \\
\hyperref[thm:measure-operations]{Sumy i mnożenie przez stałą} & \ref*{thm:measure-operations} & \pageref*{thm:measure-operations} \\
\hyperref[thm:sum-nonnegative-integrals]{Szeregi nieujemnych funkcji mierzalnych} & \ref*{thm:sum-nonnegative-integrals} & \pageref*{thm:sum-nonnegative-integrals} \\
\hyperref[thm:extended-tests]{Testy mierzalności na półprostych} & \ref*{thm:extended-tests} & \pageref*{thm:extended-tests} \\
\hyperref[thm:complex-test]{Test na częściach rzeczywistej i urojonej} & \ref*{thm:complex-test} & \pageref*{thm:complex-test} \\
\hyperref[thm:tonelli-12]{Tonelli} & \ref*{thm:tonelli-12} & \pageref*{thm:tonelli-12} \\
\hyperref[thm:completed-product-fubini-12]{Tonelli i Fubini dla uzupełnionej miary produktowej} & \ref*{thm:completed-product-fubini-12} & \pageref*{thm:completed-product-fubini-12} \\
\hyperref[thm:section-identity-12]{Tożsamość przekrojowa przy $\sigma$-skończoności} & \ref*{thm:section-identity-12} & \pageref*{thm:section-identity-12} \\
\hyperref[thm:density-transform-13]{Transformacja gęstości} & \ref*{thm:density-transform-13} & \pageref*{thm:density-transform-13} \\
\hyperref[thm:invariance-lebesgue]{Translacje i jednolite skalowanie} & \ref*{thm:invariance-lebesgue} & \pageref*{thm:invariance-lebesgue} \\
\hyperref[thm:caratheodory]{Twierdzenie Carathéodory'ego} & \ref*{thm:caratheodory} & \pageref*{thm:caratheodory} \\
\hyperref[thm:gauss-green-BV-29]{Twierdzenie Gaussa--Greena w postaci miarowej} & \ref*{thm:gauss-green-BV-29} & \pageref*{thm:gauss-green-BV-29} \\
\hyperref[thm:monotone-class]{Twierdzenie o klasie monotonicznej} & \ref*{thm:monotone-class} & \pageref*{thm:monotone-class} \\
\hyperref[thm:pi-lambda]{Twierdzenie $\pi$--$\lambda$ Dynkina} & \ref*{thm:pi-lambda} & \pageref*{thm:pi-lambda} \\
\hyperref[thm:besicovitch-24]{Twierdzenie pokryciowe Besicovitcha -- pełny dowód} & \ref*{thm:besicovitch-24} & \pageref*{thm:besicovitch-24} \\
\hyperref[thm:RN-sigfin-18]{Twierdzenie Radona--Nikodyma} & \ref*{thm:RN-sigfin-18} & \pageref*{thm:RN-sigfin-18} \\
\hyperref[thm:de-giorgi-structure-29]{Twierdzenie strukturalne De Giorgiego} & \ref*{thm:de-giorgi-structure-29} & \pageref*{thm:de-giorgi-structure-29} \\
\hyperref[thm:completion]{Uzupełnienie przestrzeni z miarą} & \ref*{thm:completion} & \pageref*{thm:completion} \\
\hyperref[thm:caratheodory-completion]{Uzupełnienie rozszerzenia} & \ref*{thm:caratheodory-completion} & \pageref*{thm:caratheodory-completion} \\
\hyperref[thm:vector-variation-measure-30]{Wariacja jest miarą dodatnią} & \ref*{thm:vector-variation-measure-30} & \pageref*{thm:vector-variation-measure-30} \\
\hyperref[thm:density-variation-16]{Wariacja miary określonej całką} & \ref*{thm:density-variation-16} & \pageref*{thm:density-variation-16} \\
\hyperref[thm:variation-density-30]{Wariacja miary o gęstości Bochnera} & \ref*{thm:variation-density-30} & \pageref*{thm:variation-density-30} \\
\hyperref[thm:complex-variation-16]{Wariacja miary zespolonej} & \ref*{thm:complex-variation-16} & \pageref*{thm:complex-variation-16} \\
\hyperref[thm:conditional-jensen-19]{Warunkowa nierówność Jensena} & \ref*{thm:conditional-jensen-19} & \pageref*{thm:conditional-jensen-19} \\
\hyperref[thm:lp-embedding-finite-14]{Włożenia przy skończonej mierze} & \ref*{thm:lp-embedding-finite-14} & \pageref*{thm:lp-embedding-finite-14} \\
\hyperref[thm:cantor-dimension-26]{Wymiar klasycznego zbioru Cantora} & \ref*{thm:cantor-dimension-26} & \pageref*{thm:cantor-dimension-26} \\
\hyperref[thm:euclidean-dimension-26]{Wymiar przestrzeni euklidesowej} & \ref*{thm:euclidean-dimension-26} & \pageref*{thm:euclidean-dimension-26} \\
\hyperref[thm:lipschitz-area-coarea-28]{Wzory area i coarea dla odwzorowań Lipschitzowskich} & \ref*{thm:lipschitz-area-coarea-28} & \pageref*{thm:lipschitz-area-coarea-28} \\
\hyperref[thm:coarea-submersion-28]{Wzór koobszaru dla submersji} & \ref*{thm:coarea-submersion-28} & \pageref*{thm:coarea-submersion-28} \\
\hyperref[thm:layer-cake-8]{Wzór warstwowy na całkę nieujemną} & \ref*{thm:layer-cake-8} & \pageref*{thm:layer-cake-8} \\
\hyperref[thm:immersion-area-28]{Wzór z krotnością dla immersji $C^1$} & \ref*{thm:immersion-area-28} & \pageref*{thm:immersion-area-28} \\
\hyperref[thm:hausdorff-line-26]{W wymiarze jeden otrzymujemy miarę Lebesgue'a} & \ref*{thm:hausdorff-line-26} & \pageref*{thm:hausdorff-line-26} \\
\hyperref[thm:change-variables-13]{Zamiana zmiennych dla dyfeomorfizmu $C^1$} & \ref*{thm:change-variables-13} & \pageref*{thm:change-variables-13} \\
\hyperref[thm:manifold-change-27]{Zamiana zmiennych na rozmaitościach} & \ref*{thm:manifold-change-27} & \pageref*{thm:manifold-change-27} \\
\hyperref[thm:kl-data-processing-19]{Zasada przetwarzania danych dla entropii} & \ref*{thm:kl-data-processing-19} & \pageref*{thm:kl-data-processing-19} \\
\hyperref[thm:mass-distribution-26]{Zasada rozkładu masy} & \ref*{thm:mass-distribution-26} & \pageref*{thm:mass-distribution-26} \\
\hyperref[thm:uniform-parameter-11]{Zbieżność jednostajna po scałkowaniu} & \ref*{thm:uniform-parameter-11} & \pageref*{thm:uniform-parameter-11} \\
\hyperref[thm:mct-ae]{Zbieżność monotoniczna prawie wszędzie} & \ref*{thm:mct-ae} & \pageref*{thm:mct-ae} \\
\hyperref[thm:vitali-l1]{Zbieżność Vitaliego na przestrzeni o mierze skończonej} & \ref*{thm:vitali-l1} & \pageref*{thm:vitali-l1} \\
\hyperref[thm:lp-to-measure-14]{Zbieżność w $L^p$ a zbieżność według miary} & \ref*{thm:lp-to-measure-14} & \pageref*{thm:lp-to-measure-14} \\
\hyperref[thm:tv-implies-narrow-22]{Zbieżność w wariacji implikuje wąską} & \ref*{thm:tv-implies-narrow-22} & \pageref*{thm:tv-implies-narrow-22} \\
\hyperref[thm:borel-trace]{Zbiory borelowskie podprzestrzeni} & \ref*{thm:borel-trace} & \pageref*{thm:borel-trace} \\
\hyperref[thm:riemann-to-lebesgue]{Zgodność całki Riemanna i Lebesgue'a} & \ref*{thm:riemann-to-lebesgue} & \pageref*{thm:riemann-to-lebesgue} \\
\hyperref[thm:composition-measurable]{Złożenia i funkcje ciągłe} & \ref*{thm:composition-measurable} & \pageref*{thm:composition-measurable} \\
\hyperref[thm:BV-banach-29]{Zupełność przestrzeni $BV$} & \ref*{thm:BV-banach-29} & \pageref*{thm:BV-banach-29} \\
\hyperref[thm:complete-measures-16]{Zupełność przestrzeni miar} & \ref*{thm:complete-measures-16} & \pageref*{thm:complete-measures-16} \\
\hyperref[thm:cauchy-measure-complete-9]{Zupełność zbieżności według miary} & \ref*{thm:cauchy-measure-complete-9} & \pageref*{thm:cauchy-measure-complete-9} \\
\hyperref[thm:perimeter-compact-29]{Zwartość rodzin zbiorów o ograniczonym obwodzie} & \ref*{thm:perimeter-compact-29} & \pageref*{thm:perimeter-compact-29} \\
\end{longtable}
\endgroup

\section*{Lematy}
\begingroup\small
\setlength\LTleft{0pt}\setlength\LTright{0pt}
\begin{longtable}{@{}p{0.72\linewidth}p{0.135\linewidth}p{0.08\linewidth}@{}}
\textbf{Nazwa} & \textbf{Nr} & \textbf{Str.}\\\hline
\endfirsthead
\multicolumn{3}{@{}l}{\emph{Lematy -- ciąg dalszy}}\\[2pt]
\textbf{Nazwa} & \textbf{Nr} & \textbf{Str.}\\\hline
\endhead
\hyperref[lem:hausdorff-metricity-26]{Addytywność na zbiorach oddalonych} & \ref*{lem:hausdorff-metricity-26} & \pageref*{lem:hausdorff-metricity-26} \\
\hyperref[lem:outer-approx-algebra]{Aproksymacja od góry} & \ref*{lem:outer-approx-algebra} & \pageref*{lem:outer-approx-algebra} \\
\hyperref[lem:vector-variation-ac-30]{Bezwzględna ciągłość wariacji} & \ref*{lem:vector-variation-ac-30} & \pageref*{lem:vector-variation-ac-30} \\
\hyperref[lem:metric-outer-borel-21]{Borelowska mierzalność} & \ref*{lem:metric-outer-borel-21} & \pageref*{lem:metric-outer-borel-21} \\
\hyperref[lem:current-boundary-square-29]{Brzeg brzegu znika} & \ref*{lem:current-boundary-square-29} & \pageref*{lem:current-boundary-square-29} \\
\hyperref[lem:essential-boundary-null-29]{Brzeg miarowy jest objętościowo mały} & \ref*{lem:essential-boundary-null-29} & \pageref*{lem:essential-boundary-null-29} \\
\hyperref[lem:FTC-Banach-30]{Całka po krzywej z pochodną Banachowską} & \ref*{lem:FTC-Banach-30} & \pageref*{lem:FTC-Banach-30} \\
\hyperref[lem:integration-density-18]{Całkowanie względem miary o gęstości} & \ref*{lem:integration-density-18} & \pageref*{lem:integration-density-18} \\
\hyperref[lem:positive-content-26]{Dodatnia zawartość zbioru zwartego} & \ref*{lem:positive-content-26} & \pageref*{lem:positive-content-26} \\
\hyperref[lem:midpoint-lp-15]{Dwie nierówności dla punktu środkowego} & \ref*{lem:midpoint-lp-15} & \pageref*{lem:midpoint-lp-15} \\
\hyperref[lem:nonatomic-divisibility-30]{Dzielenie miary bezatomowej} & \ref*{lem:nonatomic-divisibility-30} & \pageref*{lem:nonatomic-divisibility-30} \\
\hyperref[lem:singularity-17]{Elementarne własności osobliwości} & \ref*{lem:singularity-17} & \pageref*{lem:singularity-17} \\
\hyperref[lem:haar-fubini-compact-23]{Fubini dla ciągłych jąder o zwartym nośniku} & \ref*{lem:haar-fubini-compact-23} & \pageref*{lem:haar-fubini-compact-23} \\
\hyperref[lem:cutoff-radon-20]{Funkcja odcinająca} & \ref*{lem:cutoff-radon-20} & \pageref*{lem:cutoff-radon-20} \\
\hyperref[lem:density-cc-c0-21]{Gęstość $C_c(X)$ w $C_0(X)$} & \ref*{lem:density-cc-c0-21} & \pageref*{lem:density-cc-c0-21} \\
\hyperref[lem:dg-density-identification]{Gęstość miary na wykresach prostowalnych} & \ref*{lem:dg-density-identification} & \pageref*{lem:dg-density-identification} \\
\hyperref[lem:smooth-cutoff-10]{Gładka funkcja odcinająca} & \ref*{lem:smooth-cutoff-10} & \pageref*{lem:smooth-cutoff-10} \\
\hyperref[lem:canonical-lpq-15]{Izometryczne zanurzenie do dualnej} & \ref*{lem:canonical-lpq-15} & \pageref*{lem:canonical-lpq-15} \\
\hyperref[lem:haar-ratio-23]{Jednolite oszacowania ilorazów} & \ref*{lem:haar-ratio-23} & \pageref*{lem:haar-ratio-23} \\
\hyperref[lem:weak-sequence-tight-22]{Jednostajna ciasność ciągu wąsko zbieżnego} & \ref*{lem:weak-sequence-tight-22} & \pageref*{lem:weak-sequence-tight-22} \\
\hyperref[lem:haar-uniform-23]{Jednostajna ciągłość translacji funkcji} & \ref*{lem:haar-uniform-23} & \pageref*{lem:haar-uniform-23} \\
\hyperref[lem:rademacher-line-28]{Jednowymiarowa różniczkowalność} & \ref*{lem:rademacher-line-28} & \pageref*{lem:rademacher-line-28} \\
\hyperref[lem:uniform-absolute-continuity]{Kontrola całek na małych zbiorach} & \ref*{lem:uniform-absolute-continuity} & \pageref*{lem:uniform-absolute-continuity} \\
\hyperref[lem:hausdorff-caratheodory-selfcontained-26]{Kryterium Carathéodory'ego, z dowodem} & \ref*{lem:hausdorff-caratheodory-selfcontained-26} & \pageref*{lem:hausdorff-caratheodory-selfcontained-26} \\
\hyperref[lem:dynkin-intersection]{Lemat z rozdziału 1} & \ref*{lem:dynkin-intersection} & \pageref*{lem:dynkin-intersection} \\
\hyperref[lem:lipschitz-approx-compact-10]{Lipschitzowska aproksymacja na zbiorze zwartym} & \ref*{lem:lipschitz-approx-compact-10} & \pageref*{lem:lipschitz-approx-compact-10} \\
\hyperref[lem:lusin-simple-10]{Łuzin dla funkcji prostych} & \ref*{lem:lusin-simple-10} & \pageref*{lem:lusin-simple-10} \\
\hyperref[lem:haar-neigh-23]{Małe otoczenia, zwarte rozdrobnienia i pokrycia} & \ref*{lem:haar-neigh-23} & \pageref*{lem:haar-neigh-23} \\
\hyperref[lem:graph-hausdorff-28]{Małe zniekształcenie wykresu} & \ref*{lem:graph-hausdorff-28} & \pageref*{lem:graph-hausdorff-28} \\
\hyperref[lem:hausdorff-metric-open-26]{Metryczna miara zewnętrzna a zbiory domknięte} & \ref*{lem:hausdorff-metric-open-26} & \pageref*{lem:hausdorff-metric-open-26} \\
\hyperref[lem:BV-level-meas-auditVI]{Mierzalność obwodów poziomic} & \ref*{lem:BV-level-meas-auditVI} & \pageref*{lem:BV-level-meas-auditVI} \\
\hyperref[lem:sections-measurable-12]{Mierzalność wszystkich przekrojów} & \ref*{lem:sections-measurable-12} & \pageref*{lem:sections-measurable-12} \\
\hyperref[lem:variational-lp-15]{Minimum energii} & \ref*{lem:variational-lp-15} & \pageref*{lem:variational-lp-15} \\
\hyperref[lem:dg-relative-isoperimetric]{Nierówności Sobolewa i względna nierówność izoperymetryczna} & \ref*{lem:dg-relative-isoperimetric} & \pageref*{lem:dg-relative-isoperimetric} \\
\hyperref[lem:rect-cover-ineq]{Nierówność objętościowa dla pokryć} & \ref*{lem:rect-cover-ineq} & \pageref*{lem:rect-cover-ineq} \\
\hyperref[lem:simple-refinement]{Niezależność od reprezentacji} & \ref*{lem:simple-refinement} & \pageref*{lem:simple-refinement} \\
\hyperref[lem:critical-area-null-28]{Obraz zbioru o niedostatecznej randze} & \ref*{lem:critical-area-null-28} & \pageref*{lem:critical-area-null-28} \\
\hyperref[lem:BV-smooth-levels]{Obwody poziomic funkcji gładkiej} & \ref*{lem:BV-smooth-levels} & \pageref*{lem:BV-smooth-levels} \\
\hyperref[lem:lipschitz-null-13]{Odwzorowania Lipschitzowskie zachowują zbiory zerowe} & \ref*{lem:lipschitz-null-13} & \pageref*{lem:lipschitz-null-13} \\
\hyperref[lem:haar-content-to-measure-23]{Od zawartości na kompaktach do miary borelowskiej} & \ref*{lem:haar-content-to-measure-23} & \pageref*{lem:haar-content-to-measure-23} \\
\hyperref[lem:bounded-signed-16]{Ograniczoność zakresu} & \ref*{lem:bounded-signed-16} & \pageref*{lem:bounded-signed-16} \\
\hyperref[lem:dg-density]{Oszacowania gęstości na brzegu zredukowanym} & \ref*{lem:dg-density} & \pageref*{lem:dg-density} \\
\hyperref[lem:small-cube-13]{Oszacowanie obrazu małego sześcianu} & \ref*{lem:small-cube-13} & \pageref*{lem:small-cube-13} \\
\hyperref[lem:BV-translation-29]{Oszacowanie przesunięcia} & \ref*{lem:BV-translation-29} & \pageref*{lem:BV-translation-29} \\
\hyperref[lem:perimeter-complement-29]{Pochodna i dopełnienie} & \ref*{lem:perimeter-complement-29} & \pageref*{lem:perimeter-complement-29} \\
\hyperref[lem:premeasure-cover]{Podaddytywność premiary na pokryciach} & \ref*{lem:premeasure-cover} & \pageref*{lem:premeasure-cover} \\
\hyperref[lem:rectifiable-union-29]{Podstawowe własności prostowalności} & \ref*{lem:rectifiable-union-29} & \pageref*{lem:rectifiable-union-29} \\
\hyperref[lem:dg-blowup]{Półprzestrzeń jako jedyna granica po przeskalowaniu} & \ref*{lem:dg-blowup} & \pageref*{lem:dg-blowup} \\
\hyperref[lem:dg-rectifiability]{Prostowalność brzegu zredukowanego} & \ref*{lem:dg-rectifiability} & \pageref*{lem:dg-rectifiability} \\
\hyperref[lem:hausdorff-threshold-26]{Próg zerowania} & \ref*{lem:hausdorff-threshold-26} & \pageref*{lem:hausdorff-threshold-26} \\
\hyperref[lem:dg-ball-cut]{Przecięcie zbioru skończonego obwodu z kulą} & \ref*{lem:dg-ball-cut} & \pageref*{lem:dg-ball-cut} \\
\hyperref[lem:whitney-c1-28]{Przedłużenie zgodnych danych pierwszego rzędu} & \ref*{lem:whitney-c1-28} & \pageref*{lem:whitney-c1-28} \\
\hyperref[lem:finite-rectangle-sections-12]{Przekroje na prostokącie skończonej miary} & \ref*{lem:finite-rectangle-sections-12} & \pageref*{lem:finite-rectangle-sections-12} \\
\hyperref[lem:lipschitz-extend-22]{Przybliżanie lipschitzowskie i przedłużanie} & \ref*{lem:lipschitz-extend-22} & \pageref*{lem:lipschitz-extend-22} \\
\hyperref[lem:continuous-dense-25]{Przybliżanie lokalne funkcją ciągłą} & \ref*{lem:continuous-dense-25} & \pageref*{lem:continuous-dense-25} \\
\hyperref[lem:rnp-partition-30]{Przybliżanie miary przez skończone podziały} & \ref*{lem:rnp-partition-30} & \pageref*{lem:rnp-partition-30} \\
\hyperref[lem:compact-inside-10]{Przybliżenie mierzalnego zbioru przez zbiór zwarty} & \ref*{lem:compact-inside-10} & \pageref*{lem:compact-inside-10} \\
\hyperref[lem:radon-sigmafinite-inner-20]{Regularność na zbiorach $\sigma$-skończonych} & \ref*{lem:radon-sigmafinite-inner-20} & \pageref*{lem:radon-sigmafinite-inner-20} \\
\hyperref[lem:bochner-differentiation-30]{Różniczkowanie całki Bochnera} & \ref*{lem:bochner-differentiation-30} & \pageref*{lem:bochner-differentiation-30} \\
\hyperref[lem:besicovitch-angle-24]{Skończoność konfiguracji Besicovitcha} & \ref*{lem:besicovitch-angle-24} & \pageref*{lem:besicovitch-angle-24} \\
\hyperref[lem:rect-grid]{Skończony podział na siatkę} & \ref*{lem:rect-grid} & \pageref*{lem:rect-grid} \\
\hyperref[lem:rademacher-direction-28]{Słabe i kierunkowe pochodne funkcji Lipschitzowskiej} & \ref*{lem:rademacher-direction-28} & \pageref*{lem:rademacher-direction-28} \\
\hyperref[lem:BV-strict-approx]{Ścisła aproksymacja gładka w $BV$} & \ref*{lem:BV-strict-approx} & \pageref*{lem:BV-strict-approx} \\
\hyperref[lem:trace-generation]{Ślad i generowanie} & \ref*{lem:trace-generation} & \pageref*{lem:trace-generation} \\
\hyperref[lem:linear-coarea-28]{Tożsamość liniowo-algebraiczna koobszaru} & \ref*{lem:linear-coarea-28} & \pageref*{lem:linear-coarea-28} \\
\hyperref[lem:shear-13]{Translacje, permutacje i ścinanie} & \ref*{lem:shear-13} & \pageref*{lem:shear-13} \\
\hyperref[lem:BV-smooth-29]{Wariacja gładkiej funkcji} & \ref*{lem:BV-smooth-29} & \pageref*{lem:BV-smooth-29} \\
\hyperref[lem:hilbert-cube-22]{Włożenie separowalnej przestrzeni metrycznej w zwartą przestrzeń metryczną} & \ref*{lem:hilbert-cube-22} & \pageref*{lem:hilbert-cube-22} \\
\hyperref[lem:critical-coarea-null-28]{Włókna zwartego zbioru krytycznego} & \ref*{lem:critical-coarea-null-28} & \pageref*{lem:critical-coarea-null-28} \\
\hyperref[lem:conditional-pullout-19]{Wyciąganie czynnika mierzalnego} & \ref*{lem:conditional-pullout-19} & \pageref*{lem:conditional-pullout-19} \\
\hyperref[lem:exhaustion-21]{Wyczerpanie przez zbiory relatywnie zwarte} & \ref*{lem:exhaustion-21} & \pageref*{lem:exhaustion-21} \\
\hyperref[lem:mollify-cc-10]{Wygładzanie funkcji ciągłej o zwartym nośniku} & \ref*{lem:mollify-cc-10} & \pageref*{lem:mollify-cc-10} \\
\hyperref[lem:c1-general-area-28]{Wzór area dla dowolnego odwzorowania $C^1$} & \ref*{lem:c1-general-area-28} & \pageref*{lem:c1-general-area-28} \\
\hyperref[lem:young-numbers-14]{Young dla iloczynu dwóch liczb} & \ref*{lem:young-numbers-14} & \pageref*{lem:young-numbers-14} \\
\hyperref[lem:singular-density-25]{Zanik gęstości miary osobliwej} & \ref*{lem:singular-density-25} & \pageref*{lem:singular-density-25} \\
\hyperref[lem:haar-content-23]{Zawartość niezmiennicza na kompaktach} & \ref*{lem:haar-content-23} & \pageref*{lem:haar-content-23} \\
\hyperref[lem:simple-exhaustion]{Zbieżność całek prostej funkcji na zbiorach rosnących} & \ref*{lem:simple-exhaustion} & \pageref*{lem:simple-exhaustion} \\
\hyperref[lem:hyperplane-null]{Zbiory zawarte w hiperpłaszczyźnie} & \ref*{lem:hyperplane-null} & \pageref*{lem:hyperplane-null} \\
\hyperref[lem:null-fibres-lipschitz-28]{Zbiory zerowe i ich włókna} & \ref*{lem:null-fibres-lipschitz-28} & \pageref*{lem:null-fibres-lipschitz-28} \\
\hyperref[lem:hausdorff-outer-26]{Zewnętrzne własności kosztu pokryć} & \ref*{lem:hausdorff-outer-26} & \pageref*{lem:hausdorff-outer-26} \\
\hyperref[lem:restriction-integral]{Zgodność z ograniczeniem miary} & \ref*{lem:restriction-integral} & \pageref*{lem:restriction-integral} \\
\hyperref[lem:esssup-14]{Znaczenie kresu istotnego} & \ref*{lem:esssup-14} & \pageref*{lem:esssup-14} \\
\hyperref[lem:lusin-whitney-28]{Zwarte fragmenty klasy $C^1$} & \ref*{lem:lusin-whitney-28} & \pageref*{lem:lusin-whitney-28} \\
\hyperref[lem:haar-product-compact-23]{Zwartość dowolnego produktu przedziałów} & \ref*{lem:haar-product-compact-23} & \pageref*{lem:haar-product-compact-23} \\
\hyperref[lem:compact-prob-22]{Zwartość prawdopodobieństw na przestrzeni zwartej} & \ref*{lem:compact-prob-22} & \pageref*{lem:compact-prob-22} \\
\end{longtable}
\endgroup

\section*{Wnioski}
\begingroup\small
\setlength\LTleft{0pt}\setlength\LTright{0pt}
\begin{longtable}{@{}p{0.72\linewidth}p{0.135\linewidth}p{0.08\linewidth}@{}}
\textbf{Nazwa} & \textbf{Nr} & \textbf{Str.}\\\hline
\endfirsthead
\multicolumn{3}{@{}l}{\emph{Wnioski -- ciąg dalszy}}\\[2pt]
\textbf{Nazwa} & \textbf{Nr} & \textbf{Str.}\\\hline
\endhead
\hyperref[cor:smooth-approx-in-measure-10]{Aproksymacja w mierze funkcjami gładkimi} & \ref*{cor:smooth-approx-in-measure-10} & \pageref*{cor:smooth-approx-in-measure-10} \\
\hyperref[cor:besicovitch-measure-infinite-24]{Besicovitch dla dodatnich miar borelowskich} & \ref*{cor:besicovitch-measure-infinite-24} & \pageref*{cor:besicovitch-measure-infinite-24} \\
\hyperref[cor:borel-representative-10]{Borelowska reprezentacja funkcji Lebesgue'a} & \ref*{cor:borel-representative-10} & \pageref*{cor:borel-representative-10} \\
\hyperref[cor:coordinate-measurable]{Borelowskie odwzorowania wektorowe} & \ref*{cor:coordinate-measurable} & \pageref*{cor:coordinate-measurable} \\
\hyperref[cor:meas-simple-characterization]{Charakteryzacja mierzalności przez funkcje proste} & \ref*{cor:meas-simple-characterization} & \pageref*{cor:meas-simple-characterization} \\
\hyperref[cor:coercive-tight-22]{Ciasność z oszacowania momentów} & \ref*{cor:coercive-tight-22} & \pageref*{cor:coercive-tight-22} \\
\hyperref[cor:uniform-abs-cont-int-17]{Ciągłość całki względem zbioru całkowania} & \ref*{cor:uniform-abs-cont-int-17} & \pageref*{cor:uniform-abs-cont-int-17} \\
\hyperref[cor:l1-param-11]{Ciągłość w normie $L^1$} & \ref*{cor:l1-param-11} & \pageref*{cor:l1-param-11} \\
\hyperref[cor:haar-open-positive-23]{Dodatniość miary zbiorów otwartych} & \ref*{cor:haar-open-positive-23} & \pageref*{cor:haar-open-positive-23} \\
\hyperref[cor:holder-sharp-15]{Dokładne maksimum w nierówności Höldera} & \ref*{cor:holder-sharp-15} & \pageref*{cor:holder-sharp-15} \\
\hyperref[cor:lusin-finite-ae-10]{Funkcje skończone prawie wszędzie} & \ref*{cor:lusin-finite-ae-10} & \pageref*{cor:lusin-finite-ae-10} \\
\hyperref[cor:lebesgue-regularity-characterization]{Geometryczna charakterystyka mierzalności} & \ref*{cor:lebesgue-regularity-characterization} & \pageref*{cor:lebesgue-regularity-characterization} \\
\hyperref[cor:countable-lusin-10]{Jednoczesne twierdzenie Łuzina dla przeliczalnej rodziny} & \ref*{cor:countable-lusin-10} & \pageref*{cor:countable-lusin-10} \\
\hyperref[cor:measure-limit-unique]{Jednoznaczność granicy według miary} & \ref*{cor:measure-limit-unique} & \pageref*{cor:measure-limit-unique} \\
\hyperref[cor:vague-c0-bounded-22]{Kiedy zbieżność mglista jest słabą-*?} & \ref*{cor:vague-c0-bounded-22} & \pageref*{cor:vague-c0-bounded-22} \\
\hyperref[cor:NP-exact-alpha-19]{Konstrukcja testu optymalnego o zadanym poziomie} & \ref*{cor:NP-exact-alpha-19} & \pageref*{cor:NP-exact-alpha-19} \\
\hyperref[cor:simple-approx-integral]{Konstruktywne obliczanie całki} & \ref*{cor:simple-approx-integral} & \pageref*{cor:simple-approx-integral} \\
\hyperref[cor:conditional-lp-19]{Kontrakcja w $L^p$} & \ref*{cor:conditional-lp-19} & \pageref*{cor:conditional-lp-19} \\
\hyperref[cor:HL-level-24]{Kontrola zbioru, na którym średnie są duże} & \ref*{cor:HL-level-24} & \pageref*{cor:HL-level-24} \\
\hyperref[cor:tonelli-integrability-12]{Kryterium nieujemnej całkowalności} & \ref*{cor:tonelli-integrability-12} & \pageref*{cor:tonelli-integrability-12} \\
\hyperref[cor:practical-fubini-12]{Kryterium praktyczne} & \ref*{cor:practical-fubini-12} & \pageref*{cor:practical-fubini-12} \\
\hyperref[cor:every-subseq-9]{Kryterium za pomocą wszystkich podciągów} & \ref*{cor:every-subseq-9} & \pageref*{cor:every-subseq-9} \\
\hyperref[cor:nonborel-lebesgue]{Lebesgue'owskie zbiory nieborelowskie} & \ref*{cor:nonborel-lebesgue} & \pageref*{cor:nonborel-lebesgue} \\
\hyperref[cor:associative-product-12]{Łączność skończonych produktów} & \ref*{cor:associative-product-12} & \pageref*{cor:associative-product-12} \\
\hyperref[cor:riemann-hausdorff-28]{Miara riemannowska i Hausdorffa} & \ref*{cor:riemann-hausdorff-28} & \pageref*{cor:riemann-hausdorff-28} \\
\hyperref[cor:measures-separate-21]{Miary odróżniane przez funkcje ciągłe} & \ref*{cor:measures-separate-21} & \pageref*{cor:measures-separate-21} \\
\hyperref[cor:haar-adjoint-23]{Modularność grupy Liego} & \ref*{cor:haar-adjoint-23} & \pageref*{cor:haar-adjoint-23} \\
\hyperref[idx:06:wniosek:06]{Monotoniczne ciągi funkcji} & \ref*{idx:06:wniosek:06} & \pageref*{idx:06:wniosek:06} \\
\hyperref[cor:minimal-variation-16]{Najmniejsza miara dominująca} & \ref*{cor:minimal-variation-16} & \pageref*{cor:minimal-variation-16} \\
\hyperref[cor:integral-triangle]{Nierówność trójkąta i struktura $L^1$} & \ref*{cor:integral-triangle} & \pageref*{cor:integral-triangle} \\
\hyperref[cor:no-all-subsets-invariant]{Nie istnieje miara długości wszystkich podzbiorów} & \ref*{cor:no-all-subsets-invariant} & \pageref*{cor:no-all-subsets-invariant} \\
\hyperref[cor:submanifold-null-27]{Obrazy immersji niższego wymiaru} & \ref*{cor:submanifold-null-27} & \pageref*{cor:submanifold-null-27} \\
\hyperref[cor:reverse-fatou]{Odwrotny lemat Fatou przy dominacji} & \ref*{cor:reverse-fatou} & \pageref*{cor:reverse-fatou} \\
\hyperref[cor:5r-volume-24]{Oszacowanie miary sumy kul} & \ref*{cor:5r-volume-24} & \pageref*{cor:5r-volume-24} \\
\hyperref[cor:finite-besicovitch-measure-24]{Oszacowanie sumy miar w rodzinie skończonej} & \ref*{cor:finite-besicovitch-measure-24} & \pageref*{cor:finite-besicovitch-measure-24} \\
\hyperref[cor:RN-lebesgue-sum-18]{Pełny rozkład miary na część z gęstością i osobliwą} & \ref*{cor:RN-lebesgue-sum-18} & \pageref*{cor:RN-lebesgue-sum-18} \\
\hyperref[cor:riesz-order-21]{Porządek dodatnich funkcjonałów i miar} & \ref*{cor:riesz-order-21} & \pageref*{cor:riesz-order-21} \\
\hyperref[cor:null-sections-12]{Prawie wszystkie przekroje zbioru zerowego} & \ref*{cor:null-sections-12} & \pageref*{cor:null-sections-12} \\
\hyperref[cor:real-borel-tests]{Prawostronny test mierzalności} & \ref*{cor:real-borel-tests} & \pageref*{cor:real-borel-tests} \\
\hyperref[cor:null-union]{Przeliczalne sumy zbiorów miary zero} & \ref*{cor:null-union} & \pageref*{cor:null-union} \\
\hyperref[cor:countably-many-disjoint-atoms]{Przeliczalność rozłącznych atomów} & \ref*{cor:countably-many-disjoint-atoms} & \pageref*{cor:countably-many-disjoint-atoms} \\
\hyperref[idx:30:wniosek:15]{Przestrzenie Banacha i całki z krzywych} & \ref*{idx:30:wniosek:15} & \pageref*{idx:30:wniosek:15} \\
\hyperref[cor:lp-normed-14]{Przestrzeń normowana} & \ref*{cor:lp-normed-14} & \pageref*{cor:lp-normed-14} \\
\hyperref[cor:tv-data-processing-19]{Przetwarzanie danych nie zwiększa wariacji całkowitej} & \ref*{cor:tv-data-processing-19} & \pageref*{cor:tv-data-processing-19} \\
\hyperref[cor:indicator-smooth-10]{Przybliżanie funkcji charakterystycznej} & \ref*{cor:indicator-smooth-10} & \pageref*{cor:indicator-smooth-10} \\
\hyperref[cor:vitali-open-24]{Przybliżanie zbioru rozłącznymi kulami} & \ref*{cor:vitali-open-24} & \pageref*{cor:vitali-open-24} \\
\hyperref[cor:hausdorff-singletons-auditVI]{Punkty i zbiory przeliczalne} & \ref*{cor:hausdorff-singletons-auditVI} & \pageref*{cor:hausdorff-singletons-auditVI} \\
\hyperref[cor:integral-countable-partition]{Rozkład całki na rozłączne zbiory} & \ref*{cor:integral-countable-partition} & \pageref*{cor:integral-countable-partition} \\
\hyperref[cor:complex-lebesgue-decomp-17]{Rozkład miary zespolonej} & \ref*{cor:complex-lebesgue-decomp-17} & \pageref*{cor:complex-lebesgue-decomp-17} \\
\hyperref[cor:RN-linearity-18]{Różnica i suma miar z gęstościami} & \ref*{cor:RN-linearity-18} & \pageref*{cor:RN-linearity-18} \\
\hyperref[cor:coarea-level-28]{Różniczkowanie poziomic} & \ref*{cor:coarea-level-28} & \pageref*{cor:coarea-level-28} \\
\hyperref[cor:cubes-differentiation-25]{Średnie po kostkach} & \ref*{cor:cubes-differentiation-25} & \pageref*{cor:cubes-differentiation-25} \\
\hyperref[cor:bct]{Twierdzenie o zbieżności ograniczonej} & \ref*{cor:bct} & \pageref*{cor:bct} \\
\hyperref[idx:08:wniosek:20]{Uogólniona formuła rozkładu na poziomice} & \ref*{idx:08:wniosek:20} & \pageref*{idx:08:wniosek:20} \\
\hyperref[cor:generalized-holder-14]{Uogólniony Hölder} & \ref*{cor:generalized-holder-14} & \pageref*{cor:generalized-holder-14} \\
\hyperref[cor:diff-basis-25]{Wartość funkcji z otaczających ją średnich} & \ref*{cor:diff-basis-25} & \pageref*{cor:diff-basis-25} \\
\hyperref[idx:09:wniosek:14]{Wersja dla gęstości prawdopodobieństwa} & \ref*{idx:09:wniosek:14} & \pageref*{idx:09:wniosek:14} \\
\hyperref[cor:local-egorov-9]{Wersja lokalna dla przestrzeni $\sigma$-skończonej} & \ref*{cor:local-egorov-9} & \pageref*{cor:local-egorov-9} \\
\hyperref[cor:lusin-extension-10]{Wersja Łuzina przez funkcję ciągłą na całej przestrzeni} & \ref*{cor:lusin-extension-10} & \pageref*{cor:lusin-extension-10} \\
\hyperref[idx:06:wniosek:15]{Wersje mierzalne funkcji określonych p.w.} & \ref*{idx:06:wniosek:15} & \pageref*{idx:06:wniosek:15} \\
\hyperref[cor:graph-dimension-26]{Wykres funkcji Lipschitzowskiej} & \ref*{cor:graph-dimension-26} & \pageref*{cor:graph-dimension-26} \\
\hyperref[cor:hausdorff-union-26]{Wymiar sumy przeliczalnej} & \ref*{cor:hausdorff-union-26} & \pageref*{cor:hausdorff-union-26} \\
\hyperref[cor:BV-coarea-weighted-29]{Wzór koobszaru z wagą} & \ref*{cor:BV-coarea-weighted-29} & \pageref*{cor:BV-coarea-weighted-29} \\
\hyperref[cor:change-subdomain-13]{Zamiana na mierzalnym podzbiorze} & \ref*{cor:change-subdomain-13} & \pageref*{cor:change-subdomain-13} \\
\hyperref[cor:vague-narrow-22]{Zbieżność mglista bez ucieczki masy} & \ref*{cor:vague-narrow-22} & \pageref*{cor:vague-narrow-22} \\
\hyperref[cor:dct-l1]{Zbieżność zmajoryzowana w $L^1$} & \ref*{cor:dct-l1} & \pageref*{cor:dct-l1} \\
\hyperref[cor:measure-metric-complete-9]{Zupełność metryki zbieżności według miary} & \ref*{cor:measure-metric-complete-9} & \pageref*{cor:measure-metric-complete-9} \\
\end{longtable}
\endgroup


% Dżety dodane jako rozdział 31 — bez przenumerowania indeksu historycznego.

\end{document}
